% Complete English EGA IV3 editable text.
%
\PassOptionsToPackage{hyperfootnotes=false}{hyperref}
% BEGIN NATIVE INPUT preamble
\documentclass[oneside]{amsart}

\usepackage[all]{xy}
\usepackage{tikz-cd}
\usepackage[T1]{fontenc}
\usepackage{xstring}
\usepackage{xparse}
\usepackage{xcolor}
\definecolor{brightmaroon}{rgb}{0.76, 0.13, 0.28}
\usepackage[linktocpage=true,colorlinks=true,hyperindex,hypertexnames=false,citecolor=blue,linkcolor=brightmaroon]{hyperref}
\usepackage[left=1.25in,right=1.25in,top=0.75in,bottom=0.75in]{geometry}
%\usepackage[charter,greekfamily=didot]{mathdesign}
%\usepackage{Baskervaldx}
\usepackage{amssymb}
\usepackage{mathrsfs}
\usepackage{mathpazo}
\linespread{1.05}

\usepackage[nobottomtitles]{titlesec}
\usepackage{marginnote}
\usepackage[shortlabels]{enumitem}
\usepackage{longtable}
\usepackage{multicol}
\usepackage{aurical}
\usepackage{microtype}

\newtheoremstyle{ega-env-style}%
  {}{}{\rmfamily}{}{\bfseries}{.}{ }{\thmnote{(#3)}}%

\newtheoremstyle{ega-thm-env-style}%
  {}{}{\itshape}{}{\bfseries}{. --- }{ }{\thmname{#1}\thmnote{ (#3)}}%

\newtheoremstyle{ega-defn-env-style}%
  {}{}{\rmfamily}{}{\bfseries}{. --- }{ }{\thmname{#1}\thmnote{ (#3)}}%

\theoremstyle{ega-env-style}
\newtheorem*{env}{---}

\theoremstyle{ega-thm-env-style}
\newtheorem*{theorem}{Theorem}
\newtheorem*{proposition}{Proposition}
\newtheorem*{lemma}{Lemma}
\newtheorem*{corollary}{Corollary}

\theoremstyle{ega-defn-env-style}
\newtheorem*{definition}{Definition}
\newtheorem*{example}{Example}
\newtheorem*{examples}{Examples}
\newtheorem*{remark}{Remark}
\newtheorem*{remarks}{Remarks}
\newtheorem*{scholium}{Scholium}
\newtheorem*{notation}{Notation}

% indent subsections, see https://tex.stackexchange.com/questions/177290/.
% also make section titles bigger.
% also add § to \thesection, https://tex.stackexchange.com/questions/119667/ and https://tex.stackexchange.com/questions/308737/.
\makeatletter
\def\l@subsection{\@tocline{2}{0pt}{2.5pc}{2.2pc}{}}
\def\section{\@startsection{section}{1}%
  \z@{.7\linespacing\@plus\linespacing}{.5\linespacing}%
  {\normalfont\bfseries\Large\scshape\centering}}
\renewcommand{\@seccntformat}[1]{%
  \ifnum\pdfstrcmp{#1}{section}=0\textsection\fi%
  \csname the#1\endcsname.~}
\makeatother

%\allowdisplaybreaks[1]
%\binoppenalty=9999
%\relpenalty=9999

% for Chapter 0, Chapter I, etc.
% credit for ZeroRoman https://tex.stackexchange.com/questions/211414/
% added into scripts/make_book.py
%\newcommand{\ZeroRoman}[1]{\ifcase\value{#1}\relax 0\else\Roman{#1}\fi}
%\renewcommand{\thechapter}{\ZeroRoman{chapter}}

\def\mathcal{\mathscr}
% BEGIN NATIVE INPUT preamble-base
\def\sh{\mathcal}                   % sheaf font
\def\bb{\mathbf}                    % bold font
\def\cat{\mathtt}                   % category font
\def\leq{\leqslant}                 % <=
\def\geq{\geqslant}                 % >=
\def\setmin{-}                      % set minus
\def\rad{\mathfrak{r}}              % radical
\def\nilrad{\mathfrak{N}}           % nilradical
\def\emp{\varnothing}               % empty set
\def\vphi{\phi}                     % for switching \phi and \varphi, change if needed
\def\HH{\mathrm{H}}                 % cohomology H
\def\CHH{\check{\HH}}               % Čech cohomology H
\def\RR{\mathrm{R}}                 % right derived R
\def\LL{\mathrm{L}}                 % left derived L
\def\dual#1{{#1}^\vee}              % dual
\def\kres{k}                        % residue field k
\def\C{\cat{C}}                     % category C
\def\op{^\cat{op}}                  % opposite category
\def\Set{\cat{Set}}                 % category of sets
\def\CHom{\cat{Hom}}                % functor category
\def\supertilde{{\,\widetilde{\,}\,}}   % use \supertilde instead of ^\sim
\def\GL{\bb{GL}}
\def\bbGamma{\boldsymbol\Gamma}
\def\red{\mathrm{red}}
\def\rg{\operatorname{rg}}
\def\gr{\operatorname{gr}}
\def\Gr{\operatorname{Gr}}
\def\Hom{\operatorname{Hom}}
\def\Proj{\operatorname{Proj}}
\def\Tor{\operatorname{Tor}}
\def\Ext{\operatorname{Ext}}
\def\Supp{\operatorname{Supp}}
\def\Ker{\operatorname{Ker}}
\def\Im{\operatorname{Im}}
\def\Coker{\operatorname{Coker}}
\def\Spec{\operatorname{Spec}}
\def\Spm{\operatorname{Spm}}
\def\Spf{\operatorname{Spf}}
\def\grad{\operatorname{grad}}
\def\dimc{\operatorname{dimc}}
\def\codim{\operatorname{codim}}
\def\id{\operatorname{id}}
\def\Der{\operatorname{Der}}
\def\Diff{\operatorname{Diff}}
\def\Hyp{\operatorname{Hyp}}
\def\Tr{\operatorname{Tr}}

\renewcommand{\to}{\mathchoice{\longrightarrow}{\rightarrow}{\rightarrow}{\rightarrow}}
\newcommand{\from}{\mathchoice{\longleftarrow}{\leftarrow}{\leftarrow}{\leftarrow}}
\let\mapstoo\mapsto
\renewcommand{\mapsto}{\mathchoice{\longmapsto}{\mapstoo}{\mapstoo}{\mapstoo}}
\def\isoto{\simeq} % isomorphism
\def\simto{\xrightarrow{\sim}}

\renewcommand{\hat}[1]{\widehat{#1}}
\renewcommand{\tilde}[1]{\widetilde{#1}}
% END NATIVE INPUT preamble-base

\def\shHom{\sh{H}\textup{\kern-2.2pt{\Fontauri\slshape om}}}   % sheaf Hom
\def\shProj{\sh{P}\textup{\kern-2.2pt{\Fontauri\slshape roj}}} % sheaf Proj
\def\shExt{\sh{E}\textup{\kern-2.2pt{\Fontauri\slshape xt}}}   % sheaf Ext
\def\shTor{\sh{T}\textup{\kern-2.2pt{\Fontauri\slshape or}}}   % sheaf Tor
\def\shGr{\sh{G}\textup{\kern-2.2pt{\Fontauri\slshape r}}}     % sheaf Gr
\def\shDer{\sh{D}\,\textup{\kern-2.2pt{\Fontauri\slshape er}}} % sheaf Der
\def\shDiff{\sh{D}\,\textup{\kern-2.2pt{\Fontauri\slshape if{}f}}\,} % sheaf Diff
\def\shHomcont{\sh{H}\textup{\kern-2.2pt{\Fontauri\slshape om.\,cont}}}   % sheaf Hom.cont
\def\shAut{\sh{A}\textup{\kern-2.2pt{\Fontauri\slshape ut}}}   % sheaf Aut

% if unsure of a translation
%\newcommand{\unsure}[2][]{\hl{#2}\marginpar{#1}}
%\newcommand{\completelyunsure}{\unsure{[\ldots]}}
\def\unsure#1{#1 {\color{red}(?)}}
\def\completelyunsure{{\color{red}(???)}}

% use to mark where original page starts
\newcommand{\oldpage}[2][]{{\marginnote{\normalfont{\textbf{#1}~|~#2}}}\ignorespaces}
\def\sectionbreak{\begin{center}***\end{center}}

% Give every source label its own zero-size PDF destination.  The inherited
% unnumbered theorem environments otherwise leave many labels pointing to the
% preceding subsection or display anchor rather than to their exact statement.
\let\EGAOriginalLabel\label
\renewcommand{\label}[1]{\phantomsection\EGAOriginalLabel{#1}}

% for referencing environments.
% use as \sref{chapter-number.x.y.z}, with optional args
% for volume and indices, e.g. \sref[volume]{chapter-number.x.y.z}[i].
\makeatletter
\newcommand{\EGAUnavailableHyperref}[2]{%
  \begingroup
  \let\@latex@warning\@gobble
  \let\@latex@warning@no@line\@gobble
  \let\G@refundefinedtrue\relax
  \hyperref[#1]{#2}%
  \endgroup
}
\NewDocumentCommand{\sref}{o m o}{%
  \@ifundefined{r@#2}%
    {\IfNoValueTF{#1}%
      {\IfNoValueTF{#3}%
        {\EGAUnavailableHyperref{#2}{\normalfont{(\StrBehind{#2}{.})}}}%
        {\EGAUnavailableHyperref{#2}{\normalfont{(\StrBehind{#2}{.},~{#3})}}}}%
      {\IfNoValueTF{#3}%
        {\EGAUnavailableHyperref{#2}{\normalfont{(\textbf{#1},~\StrBehind{#2}{.})}}}%
        {\EGAUnavailableHyperref{#2}{\normalfont{(\textbf{#1},~\StrBehind{#2}{.},~{#3})}}}}}%
    {\IfNoValueTF{#1}%
      {\IfNoValueTF{#3}%
        {\hyperref[#2]{\normalfont{(\StrBehind{#2}{.})}}}%
        {\hyperref[#2]{\normalfont{(\StrBehind{#2}{.},~{#3})}}}}%
      {\IfNoValueTF{#3}%
        {\hyperref[#2]{\normalfont{(\textbf{#1},~\StrBehind{#2}{.})}}}%
        {\hyperref[#2]{\normalfont{(\textbf{#1},~\StrBehind{#2}{.},~{#3})}}}}}%
}
\makeatother

% for marking changes following errata
% use as \erratum[volume]{correction} to say where the erratum is given
\newcommand{\erratum}[2][]{{#2}\marginpar{\footnotesize\textbf{Err}\textsubscript{#1}}}

% for referencing equations, use as \eref{eq:chapter-number.x.y.z}.
\newcommand{\eref}[1]{\hyperref[#1]{\normalfont{(\StrBehind{#1}{.})}}}
% END NATIVE INPUT preamble
% BEGIN NATIVE INPUT EGA4_reference_v2_registry
% Generated stable semantic destinations for the EGA IV reader.
% Visible text and numbering are unchanged.
\makeatletter
\newcommand{\sgaStableTarget}[1]{%
  \Hy@raisedlink{\hyper@anchorstart{#1}\hyper@anchorend}}
\expandafter\def\csname sga@target@IV.1.1.1\endcsname{ega4.c04.def.1-1-1}
\expandafter\def\csname sga@target@IV.1.1.2\endcsname{ega4.c04.prop.1-1-2}
\expandafter\def\csname sga@target@IV.1.1.3\endcsname{ega4.c04.prop.1-1-3}
\expandafter\def\csname sga@target@IV.1.1.4\endcsname{ega4.c04.cor.1-1-4}
\expandafter\def\csname sga@target@IV.1.1.5\endcsname{ega4.c04.prop.1-1-5}
\expandafter\def\csname sga@target@IV.1.1.6\endcsname{ega4.c04.prop.1-1-6}
\expandafter\def\csname sga@target@IV.1.10.1\endcsname{ega4.c04.thm.1-10-1}
\expandafter\def\csname sga@target@IV.1.10.2\endcsname{ega4.c04.cor.1-10-2}
\expandafter\def\csname sga@target@IV.1.10.3\endcsname{ega4.c04.prop.1-10-3}
\expandafter\def\csname sga@target@IV.1.10.4\endcsname{ega4.c04.cor.1-10-4}
\expandafter\def\csname sga@target@IV.1.2.1\endcsname{ega4.c04.def.1-2-1}
\expandafter\def\csname sga@target@IV.1.2.2\endcsname{ega4.c04.prop.1-2-2}
\expandafter\def\csname sga@target@IV.1.2.3\endcsname{ega4.c04.cor.1-2-3}
\expandafter\def\csname sga@target@IV.1.2.4\endcsname{ega4.c04.prop.1-2-4}
\expandafter\def\csname sga@target@IV.1.2.4.1\endcsname{ega4.c04.eq.1-2-4-1}
\expandafter\def\csname sga@target@IV.1.2.5\endcsname{ega4.c04.prop.1-2-5}
\expandafter\def\csname sga@target@IV.1.2.6\endcsname{ega4.c04.prop.1-2-6}
\expandafter\def\csname sga@target@IV.1.2.7\endcsname{ega4.c04.prop.1-2-7}
\expandafter\def\csname sga@target@IV.1.2.8\endcsname{ega4.c04.cor.1-2-8}
\expandafter\def\csname sga@target@IV.1.2.9\endcsname{ega4.c04.prop.1-2-9}
\expandafter\def\csname sga@target@IV.1.3.1\endcsname{ega4.c04.statement.1-3-1}
\expandafter\def\csname sga@target@IV.1.3.10\endcsname{ega4.c04.statement.1-3-10}
\expandafter\def\csname sga@target@IV.1.3.11\endcsname{ega4.c04.prop.1-3-11}
\expandafter\def\csname sga@target@IV.1.3.2\endcsname{ega4.c04.def.1-3-2}
\expandafter\def\csname sga@target@IV.1.3.3\endcsname{ega4.c04.prop.1-3-3}
\expandafter\def\csname sga@target@IV.1.3.4\endcsname{ega4.c04.prop.1-3-4}
\expandafter\def\csname sga@target@IV.1.3.5\endcsname{ega4.c04.cor.1-3-5}
\expandafter\def\csname sga@target@IV.1.3.6\endcsname{ega4.c04.prop.1-3-6}
\expandafter\def\csname sga@target@IV.1.3.7\endcsname{ega4.c04.prop.1-3-7}
\expandafter\def\csname sga@target@IV.1.3.8\endcsname{ega4.c04.statement.1-3-8}
\expandafter\def\csname sga@target@IV.1.3.9\endcsname{ega4.c04.prop.1-3-9}
\expandafter\def\csname sga@target@IV.1.4.1\endcsname{ega4.c04.statement.1-4-1}
\expandafter\def\csname sga@target@IV.1.4.2\endcsname{ega4.c04.def.1-4-2}
\expandafter\def\csname sga@target@IV.1.4.3\endcsname{ega4.c04.prop.1-4-3}
\expandafter\def\csname sga@target@IV.1.4.3.1\endcsname{ega4.c04.cor.1-4-3-1}
\expandafter\def\csname sga@target@IV.1.4.4\endcsname{ega4.c04.prop.1-4-4}
\expandafter\def\csname sga@target@IV.1.4.4.1\endcsname{ega4.c04.lem.1-4-4-1}
\expandafter\def\csname sga@target@IV.1.4.5\endcsname{ega4.c04.cor.1-4-5}
\expandafter\def\csname sga@target@IV.1.4.6\endcsname{ega4.c04.prop.1-4-6}
\expandafter\def\csname sga@target@IV.1.4.7\endcsname{ega4.c04.prop.1-4-7}
\expandafter\def\csname sga@target@IV.1.4.7.1\endcsname{ega4.c04.lem.1-4-7-1}
\expandafter\def\csname sga@target@IV.1.5.1\endcsname{ega4.c04.prop.1-5-1}
\expandafter\def\csname sga@target@IV.1.5.2\endcsname{ega4.c04.def.1-5-2}
\expandafter\def\csname sga@target@IV.1.5.3\endcsname{ega4.c04.prop.1-5-3}
\expandafter\def\csname sga@target@IV.1.5.4\endcsname{ega4.c04.prop.1-5-4}
\expandafter\def\csname sga@target@IV.1.5.5\endcsname{ega4.c04.cor.1-5-5}
\expandafter\def\csname sga@target@IV.1.5.6\endcsname{ega4.c04.cor.1-5-6}
\expandafter\def\csname sga@target@IV.1.5.7\endcsname{ega4.c04.prop.1-5-7}
\expandafter\def\csname sga@target@IV.1.6.1\endcsname{ega4.c04.def.1-6-1}
\expandafter\def\csname sga@target@IV.1.6.2\endcsname{ega4.c04.prop.1-6-2}
\expandafter\def\csname sga@target@IV.1.6.3\endcsname{ega4.c04.cor.1-6-3}
\expandafter\def\csname sga@target@IV.1.6.4\endcsname{ega4.c04.rem.1-6-4}
\expandafter\def\csname sga@target@IV.1.6.5\endcsname{ega4.c04.prop.1-6-5}
\expandafter\def\csname sga@target@IV.1.7.1\endcsname{ega4.c04.statement.1-7-1}
\expandafter\def\csname sga@target@IV.1.7.10\endcsname{ega4.c04.statement.1-7-10}
\expandafter\def\csname sga@target@IV.1.7.11\endcsname{ega4.c04.statement.1-7-11}
\expandafter\def\csname sga@target@IV.1.7.12\endcsname{ega4.c04.statement.1-7-12}
\expandafter\def\csname sga@target@IV.1.7.13\endcsname{ega4.c04.statement.1-7-13}
\expandafter\def\csname sga@target@IV.1.7.14\endcsname{ega4.c04.statement.1-7-14}
\expandafter\def\csname sga@target@IV.1.7.15\endcsname{ega4.c04.statement.1-7-15}
\expandafter\def\csname sga@target@IV.1.7.16\endcsname{ega4.c04.statement.1-7-16}
\expandafter\def\csname sga@target@IV.1.7.17\endcsname{ega4.c04.statement.1-7-17}
\expandafter\def\csname sga@target@IV.1.7.18\endcsname{ega4.c04.statement.1-7-18}
\expandafter\def\csname sga@target@IV.1.7.19\endcsname{ega4.c04.statement.1-7-19}
\expandafter\def\csname sga@target@IV.1.7.2\endcsname{ega4.c04.statement.1-7-2}
\expandafter\def\csname sga@target@IV.1.7.20\endcsname{ega4.c04.statement.1-7-20}
\expandafter\def\csname sga@target@IV.1.7.21\endcsname{ega4.c04.statement.1-7-21}
\expandafter\def\csname sga@target@IV.1.7.22\endcsname{ega4.c04.statement.1-7-22}
\expandafter\def\csname sga@target@IV.1.7.3\endcsname{ega4.c04.statement.1-7-3}
\expandafter\def\csname sga@target@IV.1.7.4\endcsname{ega4.c04.statement.1-7-4}
\expandafter\def\csname sga@target@IV.1.7.5\endcsname{ega4.c04.statement.1-7-5}
\expandafter\def\csname sga@target@IV.1.7.6\endcsname{ega4.c04.statement.1-7-6}
\expandafter\def\csname sga@target@IV.1.7.7\endcsname{ega4.c04.statement.1-7-7}
\expandafter\def\csname sga@target@IV.1.7.8\endcsname{ega4.c04.statement.1-7-8}
\expandafter\def\csname sga@target@IV.1.7.9\endcsname{ega4.c04.statement.1-7-9}
\expandafter\def\csname sga@target@IV.1.8.1\endcsname{ega4.c04.statement.1-8-1}
\expandafter\def\csname sga@target@IV.1.8.2\endcsname{ega4.c04.prop.1-8-2}
\expandafter\def\csname sga@target@IV.1.8.3\endcsname{ega4.c04.lem.1-8-3}
\expandafter\def\csname sga@target@IV.1.8.4\endcsname{ega4.c04.thm.1-8-4}
\expandafter\def\csname sga@target@IV.1.8.4.1\endcsname{ega4.c04.lem.1-8-4-1}
\expandafter\def\csname sga@target@IV.1.8.4.2\endcsname{ega4.c04.lem.1-8-4-2}
\expandafter\def\csname sga@target@IV.1.8.5\endcsname{ega4.c04.cor.1-8-5}
\expandafter\def\csname sga@target@IV.1.8.5.1\endcsname{ega4.c04.lem.1-8-5-1}
\expandafter\def\csname sga@target@IV.1.8.6\endcsname{ega4.c04.cor.1-8-6}
\expandafter\def\csname sga@target@IV.1.8.7\endcsname{ega4.c04.prop.1-8-7}
\expandafter\def\csname sga@target@IV.1.8.7.1\endcsname{ega4.c04.lem.1-8-7-1}
\expandafter\def\csname sga@target@IV.1.9.1\endcsname{ega4.c04.lem.1-9-1}
\expandafter\def\csname sga@target@IV.1.9.10\endcsname{ega4.c04.prop.1-9-10}
\expandafter\def\csname sga@target@IV.1.9.11\endcsname{ega4.c04.prop.1-9-11}
\expandafter\def\csname sga@target@IV.1.9.12\endcsname{ega4.c04.cor.1-9-12}
\expandafter\def\csname sga@target@IV.1.9.13\endcsname{ega4.c04.statement.1-9-13}
\expandafter\def\csname sga@target@IV.1.9.14\endcsname{ega4.c04.prop.1-9-14}
\expandafter\def\csname sga@target@IV.1.9.15\endcsname{ega4.c04.prop.1-9-15}
\expandafter\def\csname sga@target@IV.1.9.16\endcsname{ega4.c04.statement.1-9-16}
\expandafter\def\csname sga@target@IV.1.9.2\endcsname{ega4.c04.prop.1-9-2}
\expandafter\def\csname sga@target@IV.1.9.2.1\endcsname{ega4.c04.statement.1-9-2-1}
\expandafter\def\csname sga@target@IV.1.9.2.i\endcsname{ega4.c04.statement.1-9-2-i}
\expandafter\def\csname sga@target@IV.1.9.2.ii\endcsname{ega4.c04.statement.1-9-2-ii}
\expandafter\def\csname sga@target@IV.1.9.3\endcsname{ega4.c04.prop.1-9-3}
\expandafter\def\csname sga@target@IV.1.9.3.1\endcsname{ega4.c04.lem.1-9-3-1}
\expandafter\def\csname sga@target@IV.1.9.3.2\endcsname{ega4.c04.lem.1-9-3-2}
\expandafter\def\csname sga@target@IV.1.9.4\endcsname{ega4.c04.def.1-9-4}
\expandafter\def\csname sga@target@IV.1.9.5\endcsname{ega4.c04.prop.1-9-5}
\expandafter\def\csname sga@target@IV.1.9.6\endcsname{ega4.c04.statement.1-9-6}
\expandafter\def\csname sga@target@IV.1.9.7\endcsname{ega4.c04.thm.1-9-7}
\expandafter\def\csname sga@target@IV.1.9.8\endcsname{ega4.c04.cor.1-9-8}
\expandafter\def\csname sga@target@IV.1.9.9\endcsname{ega4.c04.cor.1-9-9}
\expandafter\def\csname sga@target@IV.10.1.1\endcsname{ega4.c04.statement.10-1-1}
\expandafter\def\csname sga@target@IV.10.1.2\endcsname{ega4.c04.prop.10-1-2}
\expandafter\def\csname sga@target@IV.10.1.3\endcsname{ega4.c04.def.10-1-3}
\expandafter\def\csname sga@target@IV.10.1.4\endcsname{ega4.c04.cor.10-1-4}
\expandafter\def\csname sga@target@IV.10.10.1\endcsname{ega4.c04.statement.10-10-1}
\expandafter\def\csname sga@target@IV.10.10.2\endcsname{ega4.c04.statement.10-10-2}
\expandafter\def\csname sga@target@IV.10.10.3\endcsname{ega4.c04.statement.10-10-3}
\expandafter\def\csname sga@target@IV.10.10.4\endcsname{ega4.c04.statement.10-10-4}
\expandafter\def\csname sga@target@IV.10.10.5\endcsname{ega4.c04.statement.10-10-5}
\expandafter\def\csname sga@target@IV.10.2.1\endcsname{ega4.c04.prop.10-2-1}
\expandafter\def\csname sga@target@IV.10.2.2\endcsname{ega4.c04.def.10-2-2}
\expandafter\def\csname sga@target@IV.10.2.3\endcsname{ega4.c04.cor.10-2-3}
\expandafter\def\csname sga@target@IV.10.2.4\endcsname{ega4.c04.cor.10-2-4}
\expandafter\def\csname sga@target@IV.10.2.5\endcsname{ega4.c04.cor.10-2-5}
\expandafter\def\csname sga@target@IV.10.2.6\endcsname{ega4.c04.cor.10-2-6}
\expandafter\def\csname sga@target@IV.10.2.7\endcsname{ega4.c04.prop.10-2-7}
\expandafter\def\csname sga@target@IV.10.2.8\endcsname{ega4.c04.item.10-2-8}
\expandafter\def\csname sga@target@IV.10.2.8.1\endcsname{ega4.c04.eq.10-2-8-1}
\expandafter\def\csname sga@target@IV.10.3.1\endcsname{ega4.c04.def.10-3-1}
\expandafter\def\csname sga@target@IV.10.3.2\endcsname{ega4.c04.prop.10-3-2}
\expandafter\def\csname sga@target@IV.10.3.3\endcsname{ega4.c04.prop.10-3-3}
\expandafter\def\csname sga@target@IV.10.4.1\endcsname{ega4.c04.def.10-4-1}
\expandafter\def\csname sga@target@IV.10.4.10\endcsname{ega4.c04.statement.10-4-10}
\expandafter\def\csname sga@target@IV.10.4.11\endcsname{ega4.c04.statement.10-4-11}
\expandafter\def\csname sga@target@IV.10.4.11.1\endcsname{ega4.c04.lem.10-4-11-1}
\expandafter\def\csname sga@target@IV.10.4.12\endcsname{ega4.c04.cor.10-4-12}
\expandafter\def\csname sga@target@IV.10.4.12.1\endcsname{ega4.c04.sec.10-4-12-1}
\expandafter\def\csname sga@target@IV.10.4.12.1.errata\endcsname{ega4.c04.eq.10-4-12-1-errata}
\expandafter\def\csname sga@target@IV.10.4.13\endcsname{ega4.c04.cor.10-4-13}
\expandafter\def\csname sga@target@IV.10.4.14\endcsname{ega4.c04.rem.10-4-14}
\expandafter\def\csname sga@target@IV.10.4.2\endcsname{ega4.c04.prop.10-4-2}
\expandafter\def\csname sga@target@IV.10.4.3\endcsname{ega4.c04.statement.10-4-3}
\expandafter\def\csname sga@target@IV.10.4.4\endcsname{ega4.c04.prop.10-4-4}
\expandafter\def\csname sga@target@IV.10.4.5\endcsname{ega4.c04.prop.10-4-5}
\expandafter\def\csname sga@target@IV.10.4.5.1\endcsname{ega4.c04.lem.10-4-5-1}
\expandafter\def\csname sga@target@IV.10.4.6\endcsname{ega4.c04.cor.10-4-6}
\expandafter\def\csname sga@target@IV.10.4.7\endcsname{ega4.c04.cor.10-4-7}
\expandafter\def\csname sga@target@IV.10.4.8\endcsname{ega4.c04.cor.10-4-8}
\expandafter\def\csname sga@target@IV.10.4.9\endcsname{ega4.c04.statement.10-4-9}
\expandafter\def\csname sga@target@IV.10.5.1\endcsname{ega4.c04.prop.10-5-1}
\expandafter\def\csname sga@target@IV.10.5.10\endcsname{ega4.c04.prop.10-5-10}
\expandafter\def\csname sga@target@IV.10.5.2\endcsname{ega4.c04.cor.10-5-2}
\expandafter\def\csname sga@target@IV.10.5.3\endcsname{ega4.c04.cor.10-5-3}
\expandafter\def\csname sga@target@IV.10.5.4\endcsname{ega4.c04.rem.10-5-4}
\expandafter\def\csname sga@target@IV.10.5.5\endcsname{ega4.c04.prop.10-5-5}
\expandafter\def\csname sga@target@IV.10.5.6\endcsname{ega4.c04.rem.10-5-6}
\expandafter\def\csname sga@target@IV.10.5.7\endcsname{ega4.c04.prop.10-5-7}
\expandafter\def\csname sga@target@IV.10.5.8\endcsname{ega4.c04.cor.10-5-8}
\expandafter\def\csname sga@target@IV.10.5.9\endcsname{ega4.c04.cor.10-5-9}
\expandafter\def\csname sga@target@IV.10.6.1\endcsname{ega4.c04.prop.10-6-1}
\expandafter\def\csname sga@target@IV.10.6.1.1\endcsname{ega4.c04.eq.10-6-1-1}
\expandafter\def\csname sga@target@IV.10.6.1.2\endcsname{ega4.c04.eq.10-6-1-2}
\expandafter\def\csname sga@target@IV.10.6.1.3\endcsname{ega4.c04.eq.10-6-1-3}
\expandafter\def\csname sga@target@IV.10.6.1.4\endcsname{ega4.c04.eq.10-6-1-4}
\expandafter\def\csname sga@target@IV.10.6.2\endcsname{ega4.c04.cor.10-6-2}
\expandafter\def\csname sga@target@IV.10.6.3\endcsname{ega4.c04.prop.10-6-3}
\expandafter\def\csname sga@target@IV.10.6.3.1\endcsname{ega4.c04.eq.10-6-3-1}
\expandafter\def\csname sga@target@IV.10.6.3.2\endcsname{ega4.c04.eq.10-6-3-2}
\expandafter\def\csname sga@target@IV.10.6.3.3\endcsname{ega4.c04.eq.10-6-3-3}
\expandafter\def\csname sga@target@IV.10.6.4\endcsname{ega4.c04.cor.10-6-4}
\expandafter\def\csname sga@target@IV.10.6.4.1\endcsname{ega4.c04.eq.10-6-4-1}
\expandafter\def\csname sga@target@IV.10.7.1\endcsname{ega4.c04.statement.10-7-1}
\expandafter\def\csname sga@target@IV.10.7.2\endcsname{ega4.c04.statement.10-7-2}
\expandafter\def\csname sga@target@IV.10.7.3\endcsname{ega4.c04.statement.10-7-3}
\expandafter\def\csname sga@target@IV.10.8.1\endcsname{ega4.c04.def.10-8-1}
\expandafter\def\csname sga@target@IV.10.8.1.1\endcsname{ega4.c04.eq.10-8-1-1}
\expandafter\def\csname sga@target@IV.10.8.1.2\endcsname{ega4.c04.eq.10-8-1-2}
\expandafter\def\csname sga@target@IV.10.8.2\endcsname{ega4.c04.statement.10-8-2}
\expandafter\def\csname sga@target@IV.10.8.3\endcsname{ega4.c04.lem.10-8-3}
\expandafter\def\csname sga@target@IV.10.8.3.1\endcsname{ega4.c04.eq.10-8-3-1}
\expandafter\def\csname sga@target@IV.10.8.4\endcsname{ega4.c04.cor.10-8-4}
\expandafter\def\csname sga@target@IV.10.8.5\endcsname{ega4.c04.prop.10-8-5}
\expandafter\def\csname sga@target@IV.10.8.6\endcsname{ega4.c04.cor.10-8-6}
\expandafter\def\csname sga@target@IV.10.8.7\endcsname{ega4.c04.prop.10-8-7}
\expandafter\def\csname sga@target@IV.10.8.7.1\endcsname{ega4.c04.eq.10-8-7-1}
\expandafter\def\csname sga@target@IV.10.8.8\endcsname{ega4.c04.cor.10-8-8}
\expandafter\def\csname sga@target@IV.10.9.1\endcsname{ega4.c04.statement.10-9-1}
\expandafter\def\csname sga@target@IV.10.9.2\endcsname{ega4.c04.statement.10-9-2}
\expandafter\def\csname sga@target@IV.10.9.3\endcsname{ega4.c04.statement.10-9-3}
\expandafter\def\csname sga@target@IV.10.9.4\endcsname{ega4.c04.prop.10-9-4}
\expandafter\def\csname sga@target@IV.10.9.5\endcsname{ega4.c04.statement.10-9-5}
\expandafter\def\csname sga@target@IV.10.9.6\endcsname{ega4.c04.prop.10-9-6}
\expandafter\def\csname sga@target@IV.11.1.1\endcsname{ega4.c04.thm.11-1-1}
\expandafter\def\csname sga@target@IV.11.1.1.1\endcsname{ega4.c04.sec.11-1-1-1}
\expandafter\def\csname sga@target@IV.11.1.1.2\endcsname{ega4.c04.lem.11-1-1-2}
\expandafter\def\csname sga@target@IV.11.1.2\endcsname{ega4.c04.cor.11-1-2}
\expandafter\def\csname sga@target@IV.11.1.3\endcsname{ega4.c04.cor.11-1-3}
\expandafter\def\csname sga@target@IV.11.1.4\endcsname{ega4.c04.cor.11-1-4}
\expandafter\def\csname sga@target@IV.11.1.5\endcsname{ega4.c04.cor.11-1-5}
\expandafter\def\csname sga@target@IV.11.1.6\endcsname{ega4.c04.cor.11-1-6}
\expandafter\def\csname sga@target@IV.11.10.1\endcsname{ega4.c04.prop.11-10-1}
\expandafter\def\csname sga@target@IV.11.10.1.1\endcsname{ega4.c04.eq.11-10-1-1}
\expandafter\def\csname sga@target@IV.11.10.1.2\endcsname{ega4.c04.eq.11-10-1-2}
\expandafter\def\csname sga@target@IV.11.10.10\endcsname{ega4.c04.prop.11-10-10}
\expandafter\def\csname sga@target@IV.11.10.2\endcsname{ega4.c04.def.11-10-2}
\expandafter\def\csname sga@target@IV.11.10.3\endcsname{ega4.c04.item.11-10-3}
\expandafter\def\csname sga@target@IV.11.10.4\endcsname{ega4.c04.prop.11-10-4}
\expandafter\def\csname sga@target@IV.11.10.5\endcsname{ega4.c04.thm.11-10-5}
\expandafter\def\csname sga@target@IV.11.10.6\endcsname{ega4.c04.prop.11-10-6}
\expandafter\def\csname sga@target@IV.11.10.7\endcsname{ega4.c04.cor.11-10-7}
\expandafter\def\csname sga@target@IV.11.10.8\endcsname{ega4.c04.def.11-10-8}
\expandafter\def\csname sga@target@IV.11.10.9\endcsname{ega4.c04.thm.11-10-9}
\expandafter\def\csname sga@target@IV.11.2.1\endcsname{ega4.c04.statement.11-2-1}
\expandafter\def\csname sga@target@IV.11.2.1.1\endcsname{ega4.c04.eq.11-2-1-1}
\expandafter\def\csname sga@target@IV.11.2.1.2\endcsname{ega4.c04.eq.11-2-1-2}
\expandafter\def\csname sga@target@IV.11.2.1.3\endcsname{ega4.c04.eq.11-2-1-3}
\expandafter\def\csname sga@target@IV.11.2.1.4\endcsname{ega4.c04.eq.11-2-1-4}
\expandafter\def\csname sga@target@IV.11.2.1.5\endcsname{ega4.c04.eq.11-2-1-5}
\expandafter\def\csname sga@target@IV.11.2.10\endcsname{ega4.c04.rem.11-2-10}
\expandafter\def\csname sga@target@IV.11.2.2\endcsname{ega4.c04.statement.11-2-2}
\expandafter\def\csname sga@target@IV.11.2.2.1\endcsname{ega4.c04.eq.11-2-2-1}
\expandafter\def\csname sga@target@IV.11.2.2.2\endcsname{ega4.c04.lem.11-2-2-2}
\expandafter\def\csname sga@target@IV.11.2.3\endcsname{ega4.c04.statement.11-2-3}
\expandafter\def\csname sga@target@IV.11.2.3.1\endcsname{ega4.c04.eq.11-2-3-1}
\expandafter\def\csname sga@target@IV.11.2.3.2\endcsname{ega4.c04.eq.11-2-3-2}
\expandafter\def\csname sga@target@IV.11.2.3.3\endcsname{ega4.c04.eq.11-2-3-3}
\expandafter\def\csname sga@target@IV.11.2.4\endcsname{ega4.c04.lem.11-2-4}
\expandafter\def\csname sga@target@IV.11.2.5\endcsname{ega4.c04.lem.11-2-5}
\expandafter\def\csname sga@target@IV.11.2.6\endcsname{ega4.c04.thm.11-2-6}
\expandafter\def\csname sga@target@IV.11.2.6.1\endcsname{ega4.c04.cor.11-2-6-1}
\expandafter\def\csname sga@target@IV.11.2.7\endcsname{ega4.c04.cor.11-2-7}
\expandafter\def\csname sga@target@IV.11.2.8\endcsname{ega4.c04.prop.11-2-8}
\expandafter\def\csname sga@target@IV.11.2.9\endcsname{ega4.c04.thm.11-2-9}
\expandafter\def\csname sga@target@IV.11.2.9.1\endcsname{ega4.c04.statement.11-2-9-1}
\expandafter\def\csname sga@target@IV.11.2.9.2\endcsname{ega4.c04.lem.11-2-9-2}
\expandafter\def\csname sga@target@IV.11.2.9.3\endcsname{ega4.c04.lem.11-2-9-3}
\expandafter\def\csname sga@target@IV.11.2.9.4\endcsname{ega4.c04.cor.11-2-9-4}
\expandafter\def\csname sga@target@IV.11.2.9.5\endcsname{ega4.c04.lem.11-2-9-5}
\expandafter\def\csname sga@target@IV.11.2.9.6\endcsname{ega4.c04.lem.11-2-9-6}
\expandafter\def\csname sga@target@IV.11.2.9.7\endcsname{ega4.c04.eq.11-2-9-7}
\expandafter\def\csname sga@target@IV.11.2.9.8\endcsname{ega4.c04.eq.11-2-9-8}
\expandafter\def\csname sga@target@IV.11.2.9.9\endcsname{ega4.c04.eq.11-2-9-9}
\expandafter\def\csname sga@target@IV.11.3\endcsname{ega4.c04.subsec.11-3.key-iv-11-3}
\expandafter\def\csname sga@target@IV.11.3.1\endcsname{ega4.c04.thm.11-3-1}
\expandafter\def\csname sga@target@IV.11.3.10\endcsname{ega4.c04.thm.11-3-10}
\expandafter\def\csname sga@target@IV.11.3.10.1\endcsname{ega4.c04.lem.11-3-10-1}
\expandafter\def\csname sga@target@IV.11.3.10.2\endcsname{ega4.c04.lem.11-3-10-2}
\expandafter\def\csname sga@target@IV.11.3.10.3\endcsname{ega4.c04.eq.11-3-10-3}
\expandafter\def\csname sga@target@IV.11.3.10.4\endcsname{ega4.c04.sec.11-3-10-4}
\expandafter\def\csname sga@target@IV.11.3.11\endcsname{ega4.c04.cor.11-3-11}
\expandafter\def\csname sga@target@IV.11.3.12\endcsname{ega4.c04.rem.11-3-12}
\expandafter\def\csname sga@target@IV.11.3.13\endcsname{ega4.c04.prop.11-3-13}
\expandafter\def\csname sga@target@IV.11.3.14\endcsname{ega4.c04.cor.11-3-14}
\expandafter\def\csname sga@target@IV.11.3.15\endcsname{ega4.c04.prop.11-3-15}
\expandafter\def\csname sga@target@IV.11.3.16\endcsname{ega4.c04.prop.11-3-16}
\expandafter\def\csname sga@target@IV.11.3.17\endcsname{ega4.c04.cor.11-3-17}
\expandafter\def\csname sga@target@IV.11.3.2\endcsname{ega4.c04.cor.11-3-2}
\expandafter\def\csname sga@target@IV.11.3.3\endcsname{ega4.c04.cor.11-3-3}
\expandafter\def\csname sga@target@IV.11.3.4\endcsname{ega4.c04.prop.11-3-4}
\expandafter\def\csname sga@target@IV.11.3.5\endcsname{ega4.c04.cor.11-3-5}
\expandafter\def\csname sga@target@IV.11.3.6\endcsname{ega4.c04.prop.11-3-6}
\expandafter\def\csname sga@target@IV.11.3.7\endcsname{ega4.c04.prop.11-3-7}
\expandafter\def\csname sga@target@IV.11.3.8\endcsname{ega4.c04.thm.11-3-8}
\expandafter\def\csname sga@target@IV.11.3.9\endcsname{ega4.c04.prop.11-3-9}
\expandafter\def\csname sga@target@IV.11.3.9.1\endcsname{ega4.c04.lem.11-3-9-1}
\expandafter\def\csname sga@target@IV.11.4.1\endcsname{ega4.c04.thm.11-4-1}
\expandafter\def\csname sga@target@IV.11.4.1.1\endcsname{ega4.c04.lem.11-4-1-1}
\expandafter\def\csname sga@target@IV.11.4.1.2\endcsname{ega4.c04.lem.11-4-1-2}
\expandafter\def\csname sga@target@IV.11.4.10\endcsname{ega4.c04.cor.11-4-10}
\expandafter\def\csname sga@target@IV.11.4.11\endcsname{ega4.c04.prop.11-4-11}
\expandafter\def\csname sga@target@IV.11.4.12\endcsname{ega4.c04.prop.11-4-12}
\expandafter\def\csname sga@target@IV.11.4.13\endcsname{ega4.c04.cor.11-4-13}
\expandafter\def\csname sga@target@IV.11.4.2\endcsname{ega4.c04.rem.11-4-2}
\expandafter\def\csname sga@target@IV.11.4.3\endcsname{ega4.c04.cor.11-4-3}
\expandafter\def\csname sga@target@IV.11.4.4\endcsname{ega4.c04.cor.11-4-4}
\expandafter\def\csname sga@target@IV.11.4.5\endcsname{ega4.c04.prop.11-4-5}
\expandafter\def\csname sga@target@IV.11.4.6\endcsname{ega4.c04.cor.11-4-6}
\expandafter\def\csname sga@target@IV.11.4.7\endcsname{ega4.c04.prop.11-4-7}
\expandafter\def\csname sga@target@IV.11.4.8\endcsname{ega4.c04.rem.11-4-8}
\expandafter\def\csname sga@target@IV.11.4.9\endcsname{ega4.c04.cor.11-4-9}
\expandafter\def\csname sga@target@IV.11.5.1\endcsname{ega4.c04.thm.11-5-1}
\expandafter\def\csname sga@target@IV.11.5.2\endcsname{ega4.c04.cor.11-5-2}
\expandafter\def\csname sga@target@IV.11.5.3\endcsname{ega4.c04.prop.11-5-3}
\expandafter\def\csname sga@target@IV.11.5.4\endcsname{ega4.c04.cor.11-5-4}
\expandafter\def\csname sga@target@IV.11.5.5\endcsname{ega4.c04.prop.11-5-5}
\expandafter\def\csname sga@target@IV.11.5.5.1\endcsname{ega4.c04.lem.11-5-5-1}
\expandafter\def\csname sga@target@IV.11.5.5.2\endcsname{ega4.c04.lem.11-5-5-2}
\expandafter\def\csname sga@target@IV.11.6.1\endcsname{ega4.c04.thm.11-6-1}
\expandafter\def\csname sga@target@IV.11.6.2\endcsname{ega4.c04.cor.11-6-2}
\expandafter\def\csname sga@target@IV.11.7.1\endcsname{ega4.c04.statement.11-7-1}
\expandafter\def\csname sga@target@IV.11.7.2\endcsname{ega4.c04.statement.11-7-2}
\expandafter\def\csname sga@target@IV.11.7.3\endcsname{ega4.c04.statement.11-7-3}
\expandafter\def\csname sga@target@IV.11.7.4\endcsname{ega4.c04.statement.11-7-4}
\expandafter\def\csname sga@target@IV.11.7.5\endcsname{ega4.c04.statement.11-7-5}
\expandafter\def\csname sga@target@IV.11.8.1\endcsname{ega4.c04.thm.11-8-1}
\expandafter\def\csname sga@target@IV.11.8.2\endcsname{ega4.c04.subsec.11-8-2}
\expandafter\def\csname sga@target@IV.11.9.1\endcsname{ega4.c04.statement.11-9-1}
\expandafter\def\csname sga@target@IV.11.9.1.1\endcsname{ega4.c04.statement.11-9-1-1}
\expandafter\def\csname sga@target@IV.11.9.10\endcsname{ega4.c04.thm.11-9-10}
\expandafter\def\csname sga@target@IV.11.9.11\endcsname{ega4.c04.rem.11-9-11}
\expandafter\def\csname sga@target@IV.11.9.12\endcsname{ega4.c04.lem.11-9-12}
\expandafter\def\csname sga@target@IV.11.9.13\endcsname{ega4.c04.prop.11-9-13}
\expandafter\def\csname sga@target@IV.11.9.14\endcsname{ega4.c04.statement.11-9-14}
\expandafter\def\csname sga@target@IV.11.9.15\endcsname{ega4.c04.prop.11-9-15}
\expandafter\def\csname sga@target@IV.11.9.16\endcsname{ega4.c04.thm.11-9-16}
\expandafter\def\csname sga@target@IV.11.9.17\endcsname{ega4.c04.prop.11-9-17}
\expandafter\def\csname sga@target@IV.11.9.17.1\endcsname{ega4.c04.lem.11-9-17-1}
\expandafter\def\csname sga@target@IV.11.9.18\endcsname{ega4.c04.item.11-9-18}
\expandafter\def\csname sga@target@IV.11.9.2\endcsname{ega4.c04.statement.11-9-2}
\expandafter\def\csname sga@target@IV.11.9.3\endcsname{ega4.c04.statement.11-9-3}
\expandafter\def\csname sga@target@IV.11.9.4\endcsname{ega4.c04.statement.11-9-4}
\expandafter\def\csname sga@target@IV.11.9.5\endcsname{ega4.c04.statement.11-9-5}
\expandafter\def\csname sga@target@IV.11.9.6\endcsname{ega4.c04.statement.11-9-6}
\expandafter\def\csname sga@target@IV.11.9.7\endcsname{ega4.c04.prop.11-9-7}
\expandafter\def\csname sga@target@IV.11.9.7.1\endcsname{ega4.c04.lem.11-9-7-1}
\expandafter\def\csname sga@target@IV.11.9.8\endcsname{ega4.c04.subsec.11-9-8}
\expandafter\def\csname sga@target@IV.11.9.9\endcsname{ega4.c04.cor.11-9-9}
\expandafter\def\csname sga@target@IV.12.0.1\endcsname{ega4.c04.statement.12-0-1}
\expandafter\def\csname sga@target@IV.12.0.2\endcsname{ega4.c04.statement.12-0-2}
\expandafter\def\csname sga@target@IV.12.1.1\endcsname{ega4.c04.thm.12-1-1}
\expandafter\def\csname sga@target@IV.12.1.1.1\endcsname{ega4.c04.statement.12-1-1-1}
\expandafter\def\csname sga@target@IV.12.1.1.2\endcsname{ega4.c04.statement.12-1-1-2}
\expandafter\def\csname sga@target@IV.12.1.1.3\endcsname{ega4.c04.statement.12-1-1-3}
\expandafter\def\csname sga@target@IV.12.1.1.4\endcsname{ega4.c04.statement.12-1-1-4}
\expandafter\def\csname sga@target@IV.12.1.1.5\endcsname{ega4.c04.prop.12-1-1-5}
\expandafter\def\csname sga@target@IV.12.1.2\endcsname{ega4.c04.item.12-1-2}
\expandafter\def\csname sga@target@IV.12.1.3\endcsname{ega4.c04.cor.12-1-3}
\expandafter\def\csname sga@target@IV.12.1.4\endcsname{ega4.c04.rem.12-1-4}
\expandafter\def\csname sga@target@IV.12.1.5\endcsname{ega4.c04.cor.12-1-5}
\expandafter\def\csname sga@target@IV.12.1.6\endcsname{ega4.c04.thm.12-1-6}
\expandafter\def\csname sga@target@IV.12.1.7\endcsname{ega4.c04.cor.12-1-7}
\expandafter\def\csname sga@target@IV.12.1.8\endcsname{ega4.c04.item.12-1-8}
\expandafter\def\csname sga@target@IV.12.2.1\endcsname{ega4.c04.thm.12-2-1}
\expandafter\def\csname sga@target@IV.12.2.1.1\endcsname{ega4.c04.statement.12-2-1-1}
\expandafter\def\csname sga@target@IV.12.2.1.2\endcsname{ega4.c04.statement.12-2-1-2}
\expandafter\def\csname sga@target@IV.12.2.1.3\endcsname{ega4.c04.statement.12-2-1-3}
\expandafter\def\csname sga@target@IV.12.2.1.4\endcsname{ega4.c04.lem.12-2-1-4}
\expandafter\def\csname sga@target@IV.12.2.2\endcsname{ega4.c04.cor.12-2-2}
\expandafter\def\csname sga@target@IV.12.2.3\endcsname{ega4.c04.item.12-2-3}
\expandafter\def\csname sga@target@IV.12.2.4\endcsname{ega4.c04.thm.12-2-4}
\expandafter\def\csname sga@target@IV.12.3.1\endcsname{ega4.c04.lem.12-3-1}
\expandafter\def\csname sga@target@IV.12.3.2\endcsname{ega4.c04.prop.12-3-2}
\expandafter\def\csname sga@target@IV.12.3.2.1\endcsname{ega4.c04.eq.12-3-2-1}
\expandafter\def\csname sga@target@IV.12.3.3\endcsname{ega4.c04.prop.12-3-3}
\expandafter\def\csname sga@target@IV.12.3.3.1\endcsname{ega4.c04.lem.12-3-3-1}
\expandafter\def\csname sga@target@IV.12.3.3.2\endcsname{ega4.c04.eq.12-3-3-2}
\expandafter\def\csname sga@target@IV.12.3.3.3\endcsname{ega4.c04.lem.12-3-3-3}
\expandafter\def\csname sga@target@IV.12.3.3.4\endcsname{ega4.c04.eq.12-3-3-4}
\expandafter\def\csname sga@target@IV.12.3.3.5\endcsname{ega4.c04.lem.12-3-3-5}
\expandafter\def\csname sga@target@IV.12.3.4\endcsname{ega4.c04.cor.12-3-4}
\expandafter\def\csname sga@target@IV.13.1.1\endcsname{ega4.c04.lem.13-1-1}
\expandafter\def\csname sga@target@IV.13.1.2\endcsname{ega4.c04.lem.13-1-2}
\expandafter\def\csname sga@target@IV.13.1.3\endcsname{ega4.c04.thm.13-1-3}
\expandafter\def\csname sga@target@IV.13.1.4\endcsname{ega4.c04.cor.13-1-4}
\expandafter\def\csname sga@target@IV.13.1.5\endcsname{ega4.c04.cor.13-1-5}
\expandafter\def\csname sga@target@IV.13.1.6\endcsname{ega4.c04.cor.13-1-6}
\expandafter\def\csname sga@target@IV.13.1.7\endcsname{ega4.c04.prop.13-1-7}
\expandafter\def\csname sga@target@IV.13.2.1\endcsname{ega4.c04.statement.13-2-1}
\expandafter\def\csname sga@target@IV.13.2.10\endcsname{ega4.c04.prop.13-2-10}
\expandafter\def\csname sga@target@IV.13.2.10.1\endcsname{ega4.c04.eq.13-2-10-1}
\expandafter\def\csname sga@target@IV.13.2.10.2\endcsname{ega4.c04.eq.13-2-10-2}
\expandafter\def\csname sga@target@IV.13.2.11\endcsname{ega4.c04.cor.13-2-11}
\expandafter\def\csname sga@target@IV.13.2.12\endcsname{ega4.c04.item.13-2-12}
\expandafter\def\csname sga@target@IV.13.2.2\endcsname{ega4.c04.def.13-2-2}
\expandafter\def\csname sga@target@IV.13.2.3\endcsname{ega4.c04.prop.13-2-3}
\expandafter\def\csname sga@target@IV.13.2.3.1\endcsname{ega4.c04.prop.13-2-3-1}
\expandafter\def\csname sga@target@IV.13.2.4\endcsname{ega4.c04.statement.13-2-4}
\expandafter\def\csname sga@target@IV.13.2.4.1\endcsname{ega4.c04.statement.13-2-4-1}
\expandafter\def\csname sga@target@IV.13.2.4.2\endcsname{ega4.c04.eq.13-2-4-2}
\expandafter\def\csname sga@target@IV.13.2.4.3\endcsname{ega4.c04.eq.13-2-4-3}
\expandafter\def\csname sga@target@IV.13.2.4.4\endcsname{ega4.c04.eq.13-2-4-4}
\expandafter\def\csname sga@target@IV.13.2.4.5\endcsname{ega4.c04.eq.13-2-4-5}
\expandafter\def\csname sga@target@IV.13.2.5\endcsname{ega4.c04.lem.13-2-5}
\expandafter\def\csname sga@target@IV.13.2.6\endcsname{ega4.c04.statement.13-2-6}
\expandafter\def\csname sga@target@IV.13.2.6.1\endcsname{ega4.c04.statement.13-2-6-1}
\expandafter\def\csname sga@target@IV.13.2.7\endcsname{ega4.c04.prop.13-2-7}
\expandafter\def\csname sga@target@IV.13.2.7.1\endcsname{ega4.c04.prop.13-2-7-1}
\expandafter\def\csname sga@target@IV.13.2.8\endcsname{ega4.c04.cor.13-2-8}
\expandafter\def\csname sga@target@IV.13.2.9\endcsname{ega4.c04.cor.13-2-9}
\expandafter\def\csname sga@target@IV.13.2.9.1\endcsname{ega4.c04.eq.13-2-9-1}
\expandafter\def\csname sga@target@IV.13.2.9.2\endcsname{ega4.c04.eq.13-2-9-2}
\expandafter\def\csname sga@target@IV.13.3.1\endcsname{ega4.c04.prop.13-3-1}
\expandafter\def\csname sga@target@IV.13.3.1.1\endcsname{ega4.c04.lem.13-3-1-1}
\expandafter\def\csname sga@target@IV.13.3.2\endcsname{ega4.c04.def.13-3-2}
\expandafter\def\csname sga@target@IV.13.3.3\endcsname{ega4.c04.prop.13-3-3}
\expandafter\def\csname sga@target@IV.13.3.4\endcsname{ega4.c04.cor.13-3-4}
\expandafter\def\csname sga@target@IV.13.3.4.1\endcsname{ega4.c04.cor.13-3-4-1}
\expandafter\def\csname sga@target@IV.13.3.5\endcsname{ega4.c04.cor.13-3-5}
\expandafter\def\csname sga@target@IV.13.3.6\endcsname{ega4.c04.prop.13-3-6}
\expandafter\def\csname sga@target@IV.13.3.6.1\endcsname{ega4.c04.prop.13-3-6-1}
\expandafter\def\csname sga@target@IV.13.3.6.2\endcsname{ega4.c04.eq.13-3-6-2}
\expandafter\def\csname sga@target@IV.13.3.7\endcsname{ega4.c04.prop.13-3-7}
\expandafter\def\csname sga@target@IV.13.3.7.1\endcsname{ega4.c04.prop.13-3-7-1}
\expandafter\def\csname sga@target@IV.13.3.8\endcsname{ega4.c04.prop.13-3-8}
\expandafter\def\csname sga@target@IV.13.3.9\endcsname{ega4.c04.statement.13-3-9}
\expandafter\def\csname sga@target@IV.14.1.1\endcsname{ega4.c04.statement.14-1-1}
\expandafter\def\csname sga@target@IV.14.1.2\endcsname{ega4.c04.prop.14-1-2}
\expandafter\def\csname sga@target@IV.14.1.3\endcsname{ega4.c04.item.14-1-3}
\expandafter\def\csname sga@target@IV.14.1.4\endcsname{ega4.c04.prop.14-1-4}
\expandafter\def\csname sga@target@IV.14.1.5\endcsname{ega4.c04.rem.14-1-5}
\expandafter\def\csname sga@target@IV.14.1.6\endcsname{ega4.c04.statement.14-1-6}
\expandafter\def\csname sga@target@IV.14.2.1\endcsname{ega4.c04.thm.14-2-1}
\expandafter\def\csname sga@target@IV.14.2.1.1\endcsname{ega4.c04.eq.14-2-1-1}
\expandafter\def\csname sga@target@IV.14.2.2\endcsname{ega4.c04.cor.14-2-2}
\expandafter\def\csname sga@target@IV.14.2.3\endcsname{ega4.c04.prop.14-2-3}
\expandafter\def\csname sga@target@IV.14.2.4\endcsname{ega4.c04.cor.14-2-4}
\expandafter\def\csname sga@target@IV.14.2.5\endcsname{ega4.c04.cor.14-2-5}
\expandafter\def\csname sga@target@IV.14.2.6\endcsname{ega4.c04.cor.14-2-6}
\expandafter\def\csname sga@target@IV.14.2.7\endcsname{ega4.c04.rem.14-2-7}
\expandafter\def\csname sga@target@IV.14.3\endcsname{ega4.c04.subsec.14-3.key-iv-14-3}
\expandafter\def\csname sga@target@IV.14.3.1\endcsname{ega4.c04.statement.14-3-1}
\expandafter\def\csname sga@target@IV.14.3.10\endcsname{ega4.c04.lem.14-3-10}
\expandafter\def\csname sga@target@IV.14.3.11\endcsname{ega4.c04.prop.14-3-11}
\expandafter\def\csname sga@target@IV.14.3.12\endcsname{ega4.c04.rem.14-3-12}
\expandafter\def\csname sga@target@IV.14.3.13\endcsname{ega4.c04.prop.14-3-13}
\expandafter\def\csname sga@target@IV.14.3.14\endcsname{ega4.c04.rem.14-3-14}
\expandafter\def\csname sga@target@IV.14.3.2\endcsname{ega4.c04.prop.14-3-2}
\expandafter\def\csname sga@target@IV.14.3.3\endcsname{ega4.c04.def.14-3-3}
\expandafter\def\csname sga@target@IV.14.3.3.1\endcsname{ega4.c04.subsec.14-3-3-1}
\expandafter\def\csname sga@target@IV.14.3.4\endcsname{ega4.c04.prop.14-3-4}
\expandafter\def\csname sga@target@IV.14.3.5\endcsname{ega4.c04.prop.14-3-5}
\expandafter\def\csname sga@target@IV.14.3.6\endcsname{ega4.c04.prop.14-3-6}
\expandafter\def\csname sga@target@IV.14.3.7\endcsname{ega4.c04.cor.14-3-7}
\expandafter\def\csname sga@target@IV.14.3.8\endcsname{ega4.c04.prop.14-3-8}
\expandafter\def\csname sga@target@IV.14.3.9\endcsname{ega4.c04.item.14-3-9}
\expandafter\def\csname sga@target@IV.14.4.1\endcsname{ega4.c04.thm.14-4-1}
\expandafter\def\csname sga@target@IV.14.4.1.1\endcsname{ega4.c04.lem.14-4-1-1}
\expandafter\def\csname sga@target@IV.14.4.1.2\endcsname{ega4.c04.lem.14-4-1-2}
\expandafter\def\csname sga@target@IV.14.4.1.3\endcsname{ega4.c04.lem.14-4-1-3}
\expandafter\def\csname sga@target@IV.14.4.10\endcsname{ega4.c04.item.14-4-10}
\expandafter\def\csname sga@target@IV.14.4.2\endcsname{ega4.c04.cor.14-4-2}
\expandafter\def\csname sga@target@IV.14.4.3\endcsname{ega4.c04.cor.14-4-3}
\expandafter\def\csname sga@target@IV.14.4.4\endcsname{ega4.c04.cor.14-4-4}
\expandafter\def\csname sga@target@IV.14.4.5\endcsname{ega4.c04.rem.14-4-5}
\expandafter\def\csname sga@target@IV.14.4.6\endcsname{ega4.c04.cor.14-4-6}
\expandafter\def\csname sga@target@IV.14.4.7\endcsname{ega4.c04.prop.14-4-7}
\expandafter\def\csname sga@target@IV.14.4.7.1\endcsname{ega4.c04.eq.14-4-7-1}
\expandafter\def\csname sga@target@IV.14.4.7.2\endcsname{ega4.c04.eq.14-4-7-2}
\expandafter\def\csname sga@target@IV.14.4.8\endcsname{ega4.c04.prop.14-4-8}
\expandafter\def\csname sga@target@IV.14.4.8.1\endcsname{ega4.c04.lem.14-4-8-1}
\expandafter\def\csname sga@target@IV.14.4.9\endcsname{ega4.c04.cor.14-4-9}
\expandafter\def\csname sga@target@IV.14.5\endcsname{ega4.c04.subsec.14-5.key-iv-14-5}
\expandafter\def\csname sga@target@IV.14.5.1\endcsname{ega4.c04.lem.14-5-1}
\expandafter\def\csname sga@target@IV.14.5.1.1\endcsname{ega4.c04.eq.14-5-1-1}
\expandafter\def\csname sga@target@IV.14.5.1.2\endcsname{ega4.c04.eq.14-5-1-2}
\expandafter\def\csname sga@target@IV.14.5.1.3\endcsname{ega4.c04.eq.14-5-1-3}
\expandafter\def\csname sga@target@IV.14.5.1.4\endcsname{ega4.c04.eq.14-5-1-4}
\expandafter\def\csname sga@target@IV.14.5.10\endcsname{ega4.c04.prop.14-5-10}
\expandafter\def\csname sga@target@IV.14.5.11\endcsname{ega4.c04.cor.14-5-11}
\expandafter\def\csname sga@target@IV.14.5.11.1\endcsname{ega4.c04.lem.14-5-11-1}
\expandafter\def\csname sga@target@IV.14.5.12\endcsname{ega4.c04.cor.14-5-12}
\expandafter\def\csname sga@target@IV.14.5.2\endcsname{ega4.c04.cor.14-5-2}
\expandafter\def\csname sga@target@IV.14.5.2.1\endcsname{ega4.c04.eq.14-5-2-1}
\expandafter\def\csname sga@target@IV.14.5.3\endcsname{ega4.c04.prop.14-5-3}
\expandafter\def\csname sga@target@IV.14.5.4\endcsname{ega4.c04.cor.14-5-4}
\expandafter\def\csname sga@target@IV.14.5.5\endcsname{ega4.c04.rem.14-5-5}
\expandafter\def\csname sga@target@IV.14.5.6\endcsname{ega4.c04.prop.14-5-6}
\expandafter\def\csname sga@target@IV.14.5.7\endcsname{ega4.c04.rem.14-5-7}
\expandafter\def\csname sga@target@IV.14.5.8\endcsname{ega4.c04.prop.14-5-8}
\expandafter\def\csname sga@target@IV.14.5.9\endcsname{ega4.c04.prop.14-5-9}
\expandafter\def\csname sga@target@IV.15.1.1\endcsname{ega4.c04.prop.15-1-1}
\expandafter\def\csname sga@target@IV.15.1.1-proof\endcsname{ega4.c04.statement.15-1-1-proof}
\expandafter\def\csname sga@target@IV.15.1.1.1\endcsname{ega4.c04.eq.15-1-1-1}
\expandafter\def\csname sga@target@IV.15.1.1.2\endcsname{ega4.c04.eq.15-1-1-2}
\expandafter\def\csname sga@target@IV.15.1.1.3\endcsname{ega4.c04.eq.15-1-1-3}
\expandafter\def\csname sga@target@IV.15.1.1.4\endcsname{ega4.c04.eq.15-1-1-4}
\expandafter\def\csname sga@target@IV.15.1.1.5\endcsname{ega4.c04.eq.15-1-1-5}
\expandafter\def\csname sga@target@IV.15.1.1.6\endcsname{ega4.c04.lem.15-1-1-6}
\expandafter\def\csname sga@target@IV.15.1.2\endcsname{ega4.c04.cor.15-1-2}
\expandafter\def\csname sga@target@IV.15.1.3\endcsname{ega4.c04.cor.15-1-3}
\expandafter\def\csname sga@target@IV.15.1.4\endcsname{ega4.c04.statement.15-1-4}
\expandafter\def\csname sga@target@IV.15.2.1\endcsname{ega4.c04.statement.15-2-1}
\expandafter\def\csname sga@target@IV.15.2.2\endcsname{ega4.c04.thm.15-2-2}
\expandafter\def\csname sga@target@IV.15.2.2-proof\endcsname{ega4.c04.item.15-2-2-proof}
\expandafter\def\csname sga@target@IV.15.2.2.1\endcsname{ega4.c04.lem.15-2-2-1}
\expandafter\def\csname sga@target@IV.15.2.3\endcsname{ega4.c04.cor.15-2-3}
\expandafter\def\csname sga@target@IV.15.2.4\endcsname{ega4.c04.item.15-2-4}
\expandafter\def\csname sga@target@IV.15.3.1\endcsname{ega4.c04.prop.15-3-1}
\expandafter\def\csname sga@target@IV.15.3.1-proof\endcsname{ega4.c04.item.15-3-1-proof}
\expandafter\def\csname sga@target@IV.15.3.2\endcsname{ega4.c04.prop.15-3-2}
\expandafter\def\csname sga@target@IV.15.3.2-proof\endcsname{ega4.c04.item.15-3-2-proof}
\expandafter\def\csname sga@target@IV.15.3.3\endcsname{ega4.c04.cor.15-3-3}
\expandafter\def\csname sga@target@IV.15.3.3-proof\endcsname{ega4.c04.item.15-3-3-proof}
\expandafter\def\csname sga@target@IV.15.3.4\endcsname{ega4.c04.item.15-3-4}
\expandafter\def\csname sga@target@IV.15.4.1\endcsname{ega4.c04.prop.15-4-1}
\expandafter\def\csname sga@target@IV.15.4.2\endcsname{ega4.c04.prop.15-4-2}
\expandafter\def\csname sga@target@IV.15.4.3\endcsname{ega4.c04.prop.15-4-3}
\expandafter\def\csname sga@target@IV.15.4.3-proof\endcsname{ega4.c04.item.15-4-3-proof}
\expandafter\def\csname sga@target@IV.15.5.1\endcsname{ega4.c04.prop.15-5-1}
\expandafter\def\csname sga@target@IV.15.5.10\endcsname{ega4.c04.item.15-5-10}
\expandafter\def\csname sga@target@IV.15.5.2\endcsname{ega4.c04.cor.15-5-2}
\expandafter\def\csname sga@target@IV.15.5.2.1\endcsname{ega4.c04.eq.15-5-2-1}
\expandafter\def\csname sga@target@IV.15.5.2.2\endcsname{ega4.c04.lem.15-5-2-2}
\expandafter\def\csname sga@target@IV.15.5.2.3\endcsname{ega4.c04.eq.15-5-2-3}
\expandafter\def\csname sga@target@IV.15.5.2.4\endcsname{ega4.c04.eq.15-5-2-4}
\expandafter\def\csname sga@target@IV.15.5.3\endcsname{ega4.c04.prop.15-5-3}
\expandafter\def\csname sga@target@IV.15.5.4\endcsname{ega4.c04.cor.15-5-4}
\expandafter\def\csname sga@target@IV.15.5.5\endcsname{ega4.c04.statement.15-5-5}
\expandafter\def\csname sga@target@IV.15.5.6\endcsname{ega4.c04.lem.15-5-6}
\expandafter\def\csname sga@target@IV.15.5.6.1\endcsname{ega4.c04.lem.15-5-6-1}
\expandafter\def\csname sga@target@IV.15.5.7\endcsname{ega4.c04.prop.15-5-7}
\expandafter\def\csname sga@target@IV.15.5.7.1\endcsname{ega4.c04.eq.15-5-7-1}
\expandafter\def\csname sga@target@IV.15.5.8\endcsname{ega4.c04.statement.15-5-8}
\expandafter\def\csname sga@target@IV.15.5.9\endcsname{ega4.c04.prop.15-5-9}
\expandafter\def\csname sga@target@IV.15.6.1\endcsname{ega4.c04.prop.15-6-1}
\expandafter\def\csname sga@target@IV.15.6.1.1\endcsname{ega4.c04.eq.15-6-1-1}
\expandafter\def\csname sga@target@IV.15.6.10\endcsname{ega4.c04.item.15-6-10}
\expandafter\def\csname sga@target@IV.15.6.2\endcsname{ega4.c04.cor.15-6-2}
\expandafter\def\csname sga@target@IV.15.6.3\endcsname{ega4.c04.prop.15-6-3}
\expandafter\def\csname sga@target@IV.15.6.4\endcsname{ega4.c04.prop.15-6-4}
\expandafter\def\csname sga@target@IV.15.6.5\endcsname{ega4.c04.cor.15-6-5}
\expandafter\def\csname sga@target@IV.15.6.6\endcsname{ega4.c04.prop.15-6-6}
\expandafter\def\csname sga@target@IV.15.6.7\endcsname{ega4.c04.cor.15-6-7}
\expandafter\def\csname sga@target@IV.15.6.8\endcsname{ega4.c04.prop.15-6-8}
\expandafter\def\csname sga@target@IV.15.6.9\endcsname{ega4.c04.cor.15-6-9}
\expandafter\def\csname sga@target@IV.15.7.1\endcsname{ega4.c04.def.15-7-1}
\expandafter\def\csname sga@target@IV.15.7.10\endcsname{ega4.c04.cor.15-7-10}
\expandafter\def\csname sga@target@IV.15.7.11\endcsname{ega4.c04.cor.15-7-11}
\expandafter\def\csname sga@target@IV.15.7.11.1\endcsname{ega4.c04.lem.15-7-11-1}
\expandafter\def\csname sga@target@IV.15.7.2\endcsname{ega4.c04.prop.15-7-2}
\expandafter\def\csname sga@target@IV.15.7.3\endcsname{ega4.c04.rem.15-7-3}
\expandafter\def\csname sga@target@IV.15.7.4\endcsname{ega4.c04.cor.15-7-4}
\expandafter\def\csname sga@target@IV.15.7.5\endcsname{ega4.c04.cor.15-7-5}
\expandafter\def\csname sga@target@IV.15.7.6\endcsname{ega4.c04.cor.15-7-6}
\expandafter\def\csname sga@target@IV.15.7.7\endcsname{ega4.c04.item.15-7-7}
\expandafter\def\csname sga@target@IV.15.7.8\endcsname{ega4.c04.statement.15-7-8}
\expandafter\def\csname sga@target@IV.15.7.9\endcsname{ega4.c04.prop.15-7-9}
\expandafter\def\csname sga@target@IV.16.1\endcsname{ega4.c04.subsec.16-1}
\expandafter\def\csname sga@target@IV.16.1.1\endcsname{ega4.c04.statement.16-1-1}
\expandafter\def\csname sga@target@IV.16.1.10\endcsname{ega4.c04.cor.16-1-10}
\expandafter\def\csname sga@target@IV.16.1.11\endcsname{ega4.c04.rem.16-1-11}
\expandafter\def\csname sga@target@IV.16.1.2\endcsname{ega4.c04.def.16-1-2}
\expandafter\def\csname sga@target@IV.16.1.2.1\endcsname{ega4.c04.def.16-1-2-1}
\expandafter\def\csname sga@target@IV.16.1.2.2\endcsname{ega4.c04.eq.16-1-2-2}
\expandafter\def\csname sga@target@IV.16.1.3\endcsname{ega4.c04.statement.16-1-3}
\expandafter\def\csname sga@target@IV.16.1.4\endcsname{ega4.c04.statement.16-1-4}
\expandafter\def\csname sga@target@IV.16.1.5\endcsname{ega4.c04.prop.16-1-5}
\expandafter\def\csname sga@target@IV.16.1.6\endcsname{ega4.c04.cor.16-1-6}
\expandafter\def\csname sga@target@IV.16.1.7\endcsname{ega4.c04.cor.16-1-7}
\expandafter\def\csname sga@target@IV.16.1.8\endcsname{ega4.c04.prop.16-1-8}
\expandafter\def\csname sga@target@IV.16.1.9\endcsname{ega4.c04.prop.16-1-9}
\expandafter\def\csname sga@target@IV.16.10\endcsname{ega4.c04.subsec.16-10}
\expandafter\def\csname sga@target@IV.16.10.1\endcsname{ega4.c04.def.16-10-1}
\expandafter\def\csname sga@target@IV.16.10.2\endcsname{ega4.c04.prop.16-10-2}
\expandafter\def\csname sga@target@IV.16.10.3\endcsname{ega4.c04.prop.16-10-3}
\expandafter\def\csname sga@target@IV.16.10.4\endcsname{ega4.c04.prop.16-10-4}
\expandafter\def\csname sga@target@IV.16.10.5\endcsname{ega4.c04.prop.16-10-5}
\expandafter\def\csname sga@target@IV.16.10.6\endcsname{ega4.c04.rem.16-10-6}
\expandafter\def\csname sga@target@IV.16.11\endcsname{ega4.c04.subsec.16-11}
\expandafter\def\csname sga@target@IV.16.11.1\endcsname{ega4.c04.statement.16-11-1}
\expandafter\def\csname sga@target@IV.16.11.1.1\endcsname{ega4.c04.statement.16-11-1-1}
\expandafter\def\csname sga@target@IV.16.11.1.2\endcsname{ega4.c04.statement.16-11-1-2}
\expandafter\def\csname sga@target@IV.16.11.1.3\endcsname{ega4.c04.eq.16-11-1-3}
\expandafter\def\csname sga@target@IV.16.11.1.4\endcsname{ega4.c04.eq.16-11-1-4}
\expandafter\def\csname sga@target@IV.16.11.1.5\endcsname{ega4.c04.eq.16-11-1-5}
\expandafter\def\csname sga@target@IV.16.11.1.6\endcsname{ega4.c04.eq.16-11-1-6}
\expandafter\def\csname sga@target@IV.16.11.1.7\endcsname{ega4.c04.eq.16-11-1-7}
\expandafter\def\csname sga@target@IV.16.11.2\endcsname{ega4.c04.thm.16-11-2}
\expandafter\def\csname sga@target@IV.16.11.2.1\endcsname{ega4.c04.eq.16-11-2-1}
\expandafter\def\csname sga@target@IV.16.11.2.2\endcsname{ega4.c04.eq.16-11-2-2}
\expandafter\def\csname sga@target@IV.16.11.2.3\endcsname{ega4.c04.eq.16-11-2-3}
\expandafter\def\csname sga@target@IV.16.11.3\endcsname{ega4.c04.statement.16-11-3}
\expandafter\def\csname sga@target@IV.16.12\endcsname{ega4.c04.subsec.16-12}
\expandafter\def\csname sga@target@IV.16.12.1\endcsname{ega4.c04.statement.16-12-1}
\expandafter\def\csname sga@target@IV.16.12.2\endcsname{ega4.c04.thm.16-12-2}
\expandafter\def\csname sga@target@IV.16.12.3\endcsname{ega4.c04.statement.16-12-3}
\expandafter\def\csname sga@target@IV.16.2\endcsname{ega4.c04.subsec.16-2}
\expandafter\def\csname sga@target@IV.16.2.1\endcsname{ega4.c04.statement.16-2-1}
\expandafter\def\csname sga@target@IV.16.2.1.1\endcsname{ega4.c04.statement.16-2-1-1}
\expandafter\def\csname sga@target@IV.16.2.1.2\endcsname{ega4.c04.eq.16-2-1-2}
\expandafter\def\csname sga@target@IV.16.2.1.3\endcsname{ega4.c04.eq.16-2-1-3}
\expandafter\def\csname sga@target@IV.16.2.1.4\endcsname{ega4.c04.eq.16-2-1-4}
\expandafter\def\csname sga@target@IV.16.2.2\endcsname{ega4.c04.prop.16-2-2}
\expandafter\def\csname sga@target@IV.16.2.3\endcsname{ega4.c04.cor.16-2-3}
\expandafter\def\csname sga@target@IV.16.2.3.1\endcsname{ega4.c04.eq.16-2-3-1}
\expandafter\def\csname sga@target@IV.16.2.3.2\endcsname{ega4.c04.eq.16-2-3-2}
\expandafter\def\csname sga@target@IV.16.2.4\endcsname{ega4.c04.cor.16-2-4}
\expandafter\def\csname sga@target@IV.16.2.5\endcsname{ega4.c04.statement.16-2-5}
\expandafter\def\csname sga@target@IV.16.2.6\endcsname{ega4.c04.statement.16-2-6}
\expandafter\def\csname sga@target@IV.16.2.6.1\endcsname{ega4.c04.statement.16-2-6-1}
\expandafter\def\csname sga@target@IV.16.2.6.2\endcsname{ega4.c04.eq.16-2-6-2}
\expandafter\def\csname sga@target@IV.16.2.7\endcsname{ega4.c04.prop.16-2-7}
\expandafter\def\csname sga@target@IV.16.2.7.1\endcsname{ega4.c04.prop.16-2-7-1}
\expandafter\def\csname sga@target@IV.16.3\endcsname{ega4.c04.subsec.16-3}
\expandafter\def\csname sga@target@IV.16.3.1\endcsname{ega4.c04.def.16-3-1}
\expandafter\def\csname sga@target@IV.16.3.1.1\endcsname{ega4.c04.eq.16-3-1-1}
\expandafter\def\csname sga@target@IV.16.3.2\endcsname{ega4.c04.statement.16-3-2}
\expandafter\def\csname sga@target@IV.16.3.3\endcsname{ega4.c04.statement.16-3-3}
\expandafter\def\csname sga@target@IV.16.3.3.1\endcsname{ega4.c04.statement.16-3-3-1}
\expandafter\def\csname sga@target@IV.16.3.4\endcsname{ega4.c04.prop.16-3-4}
\expandafter\def\csname sga@target@IV.16.3.5\endcsname{ega4.c04.statement.16-3-5}
\expandafter\def\csname sga@target@IV.16.3.6\endcsname{ega4.c04.def.16-3-6}
\expandafter\def\csname sga@target@IV.16.3.6.1\endcsname{ega4.c04.eq.16-3-6-1}
\expandafter\def\csname sga@target@IV.16.3.7\endcsname{ega4.c04.statement.16-3-7}
\expandafter\def\csname sga@target@IV.16.3.7.1\endcsname{ega4.c04.eq.16-3-7-1}
\expandafter\def\csname sga@target@IV.16.3.8\endcsname{ega4.c04.prop.16-3-8}
\expandafter\def\csname sga@target@IV.16.3.9\endcsname{ega4.c04.prop.16-3-9}
\expandafter\def\csname sga@target@IV.16.4\endcsname{ega4.c04.subsec.16-4}
\expandafter\def\csname sga@target@IV.16.4.1\endcsname{ega4.c04.statement.16-4-1}
\expandafter\def\csname sga@target@IV.16.4.1.1\endcsname{ega4.c04.statement.16-4-1-1}
\expandafter\def\csname sga@target@IV.16.4.1.2\endcsname{ega4.c04.statement.16-4-1-2}
\expandafter\def\csname sga@target@IV.16.4.1.3\endcsname{ega4.c04.eq.16-4-1-3}
\expandafter\def\csname sga@target@IV.16.4.1.4\endcsname{ega4.c04.eq.16-4-1-4}
\expandafter\def\csname sga@target@IV.16.4.1.5\endcsname{ega4.c04.eq.16-4-1-5}
\expandafter\def\csname sga@target@IV.16.4.10\endcsname{ega4.c04.cor.16-4-10}
\expandafter\def\csname sga@target@IV.16.4.11\endcsname{ega4.c04.prop.16-4-11}
\expandafter\def\csname sga@target@IV.16.4.11.1\endcsname{ega4.c04.prop.16-4-11-1}
\expandafter\def\csname sga@target@IV.16.4.11.2\endcsname{ega4.c04.prop.16-4-11-2}
\expandafter\def\csname sga@target@IV.16.4.11.3\endcsname{ega4.c04.eq.16-4-11-3}
\expandafter\def\csname sga@target@IV.16.4.11.4\endcsname{ega4.c04.eq.16-4-11-4}
\expandafter\def\csname sga@target@IV.16.4.12\endcsname{ega4.c04.cor.16-4-12}
\expandafter\def\csname sga@target@IV.16.4.13\endcsname{ega4.c04.cor.16-4-13}
\expandafter\def\csname sga@target@IV.16.4.14\endcsname{ega4.c04.prop.16-4-14}
\expandafter\def\csname sga@target@IV.16.4.14.1\endcsname{ega4.c04.prop.16-4-14-1}
\expandafter\def\csname sga@target@IV.16.4.15\endcsname{ega4.c04.cor.16-4-15}
\expandafter\def\csname sga@target@IV.16.4.15.1\endcsname{ega4.c04.cor.16-4-15-1}
\expandafter\def\csname sga@target@IV.16.4.15.2\endcsname{ega4.c04.eq.16-4-15-2}
\expandafter\def\csname sga@target@IV.16.4.16\endcsname{ega4.c04.cor.16-4-16}
\expandafter\def\csname sga@target@IV.16.4.16.1\endcsname{ega4.c04.cor.16-4-16-1}
\expandafter\def\csname sga@target@IV.16.4.16.2\endcsname{ega4.c04.eq.16-4-16-2}
\expandafter\def\csname sga@target@IV.16.4.17\endcsname{ega4.c04.cor.16-4-17}
\expandafter\def\csname sga@target@IV.16.4.18\endcsname{ega4.c04.statement.16-4-18}
\expandafter\def\csname sga@target@IV.16.4.18.1\endcsname{ega4.c04.statement.16-4-18-1}
\expandafter\def\csname sga@target@IV.16.4.18.2\endcsname{ega4.c04.statement.16-4-18-2}
\expandafter\def\csname sga@target@IV.16.4.18.3\endcsname{ega4.c04.eq.16-4-18-3}
\expandafter\def\csname sga@target@IV.16.4.19\endcsname{ega4.c04.cor.16-4-19}
\expandafter\def\csname sga@target@IV.16.4.19.1\endcsname{ega4.c04.cor.16-4-19-1}
\expandafter\def\csname sga@target@IV.16.4.2\endcsname{ega4.c04.statement.16-4-2}
\expandafter\def\csname sga@target@IV.16.4.20\endcsname{ega4.c04.prop.16-4-20}
\expandafter\def\csname sga@target@IV.16.4.21\endcsname{ega4.c04.cor.16-4-21}
\expandafter\def\csname sga@target@IV.16.4.21.1\endcsname{ega4.c04.cor.16-4-21-1}
\expandafter\def\csname sga@target@IV.16.4.22\endcsname{ega4.c04.cor.16-4-22}
\expandafter\def\csname sga@target@IV.16.4.23\endcsname{ega4.c04.prop.16-4-23}
\expandafter\def\csname sga@target@IV.16.4.23.1\endcsname{ega4.c04.prop.16-4-23-1}
\expandafter\def\csname sga@target@IV.16.4.24\endcsname{ega4.c04.statement.16-4-24}
\expandafter\def\csname sga@target@IV.16.4.3\endcsname{ega4.c04.statement.16-4-3}
\expandafter\def\csname sga@target@IV.16.4.3.1\endcsname{ega4.c04.statement.16-4-3-1}
\expandafter\def\csname sga@target@IV.16.4.3.2\endcsname{ega4.c04.statement.16-4-3-2}
\expandafter\def\csname sga@target@IV.16.4.3.3\endcsname{ega4.c04.eq.16-4-3-3}
\expandafter\def\csname sga@target@IV.16.4.3.4\endcsname{ega4.c04.eq.16-4-3-4}
\expandafter\def\csname sga@target@IV.16.4.3.5\endcsname{ega4.c04.eq.16-4-3-5}
\expandafter\def\csname sga@target@IV.16.4.3.6\endcsname{ega4.c04.eq.16-4-3-6}
\expandafter\def\csname sga@target@IV.16.4.3.7\endcsname{ega4.c04.eq.16-4-3-7}
\expandafter\def\csname sga@target@IV.16.4.4\endcsname{ega4.c04.statement.16-4-4}
\expandafter\def\csname sga@target@IV.16.4.5\endcsname{ega4.c04.prop.16-4-5}
\expandafter\def\csname sga@target@IV.16.4.6\endcsname{ega4.c04.cor.16-4-6}
\expandafter\def\csname sga@target@IV.16.4.7\endcsname{ega4.c04.statement.16-4-7}
\expandafter\def\csname sga@target@IV.16.4.7.1\endcsname{ega4.c04.statement.16-4-7-1}
\expandafter\def\csname sga@target@IV.16.4.7.2\endcsname{ega4.c04.eq.16-4-7-2}
\expandafter\def\csname sga@target@IV.16.4.7.3\endcsname{ega4.c04.eq.16-4-7-3}
\expandafter\def\csname sga@target@IV.16.4.8\endcsname{ega4.c04.prop.16-4-8}
\expandafter\def\csname sga@target@IV.16.4.8.1\endcsname{ega4.c04.prop.16-4-8-1}
\expandafter\def\csname sga@target@IV.16.4.8.2\endcsname{ega4.c04.eq.16-4-8-2}
\expandafter\def\csname sga@target@IV.16.4.8.3\endcsname{ega4.c04.lem.16-4-8-3}
\expandafter\def\csname sga@target@IV.16.4.8.4\endcsname{ega4.c04.eq.16-4-8-4}
\expandafter\def\csname sga@target@IV.16.4.9\endcsname{ega4.c04.cor.16-4-9}
\expandafter\def\csname sga@target@IV.16.4.9.1\endcsname{ega4.c04.cor.16-4-9-1}
\expandafter\def\csname sga@target@IV.16.5\endcsname{ega4.c04.subsec.16-5}
\expandafter\def\csname sga@target@IV.16.5.1\endcsname{ega4.c04.statement.16-5-1}
\expandafter\def\csname sga@target@IV.16.5.1.1\endcsname{ega4.c04.eq.16-5-1-1}
\expandafter\def\csname sga@target@IV.16.5.1.2\endcsname{ega4.c04.eq.16-5-1-2}
\expandafter\def\csname sga@target@IV.16.5.10\endcsname{ega4.c04.statement.16-5-10}
\expandafter\def\csname sga@target@IV.16.5.10.1\endcsname{ega4.c04.statement.16-5-10-1}
\expandafter\def\csname sga@target@IV.16.5.10.2\endcsname{ega4.c04.statement.16-5-10-2}
\expandafter\def\csname sga@target@IV.16.5.11\endcsname{ega4.c04.prop.16-5-11}
\expandafter\def\csname sga@target@IV.16.5.12\endcsname{ega4.c04.statement.16-5-12}
\expandafter\def\csname sga@target@IV.16.5.12.1\endcsname{ega4.c04.statement.16-5-12-1}
\expandafter\def\csname sga@target@IV.16.5.12.2\endcsname{ega4.c04.eq.16-5-12-2}
\expandafter\def\csname sga@target@IV.16.5.13\endcsname{ega4.c04.statement.16-5-13}
\expandafter\def\csname sga@target@IV.16.5.13.1\endcsname{ega4.c04.statement.16-5-13-1}
\expandafter\def\csname sga@target@IV.16.5.13.2\endcsname{ega4.c04.eq.16-5-13-2}
\expandafter\def\csname sga@target@IV.16.5.13.3\endcsname{ega4.c04.eq.16-5-13-3}
\expandafter\def\csname sga@target@IV.16.5.13.4\endcsname{ega4.c04.eq.16-5-13-4}
\expandafter\def\csname sga@target@IV.16.5.13.5\endcsname{ega4.c04.eq.16-5-13-5}
\expandafter\def\csname sga@target@IV.16.5.14\endcsname{ega4.c04.statement.16-5-14}
\expandafter\def\csname sga@target@IV.16.5.14.1\endcsname{ega4.c04.statement.16-5-14-1}
\expandafter\def\csname sga@target@IV.16.5.14.2\endcsname{ega4.c04.eq.16-5-14-2}
\expandafter\def\csname sga@target@IV.16.5.14.3\endcsname{ega4.c04.eq.16-5-14-3}
\expandafter\def\csname sga@target@IV.16.5.15\endcsname{ega4.c04.statement.16-5-15}
\expandafter\def\csname sga@target@IV.16.5.16\endcsname{ega4.c04.prop.16-5-16}
\expandafter\def\csname sga@target@IV.16.5.17\endcsname{ega4.c04.prop.16-5-17}
\expandafter\def\csname sga@target@IV.16.5.18\endcsname{ega4.c04.cor.16-5-18}
\expandafter\def\csname sga@target@IV.16.5.2\endcsname{ega4.c04.prop.16-5-2}
\expandafter\def\csname sga@target@IV.16.5.3\endcsname{ega4.c04.prop.16-5-3}
\expandafter\def\csname sga@target@IV.16.5.3.1\endcsname{ega4.c04.eq.16-5-3-1}
\expandafter\def\csname sga@target@IV.16.5.4\endcsname{ega4.c04.statement.16-5-4}
\expandafter\def\csname sga@target@IV.16.5.5\endcsname{ega4.c04.cor.16-5-5}
\expandafter\def\csname sga@target@IV.16.5.5.1\endcsname{ega4.c04.cor.16-5-5-1}
\expandafter\def\csname sga@target@IV.16.5.6\endcsname{ega4.c04.cor.16-5-6}
\expandafter\def\csname sga@target@IV.16.5.7\endcsname{ega4.c04.statement.16-5-7}
\expandafter\def\csname sga@target@IV.16.5.7.1\endcsname{ega4.c04.statement.16-5-7-1}
\expandafter\def\csname sga@target@IV.16.5.8\endcsname{ega4.c04.statement.16-5-8}
\expandafter\def\csname sga@target@IV.16.5.8.1\endcsname{ega4.c04.eq.16-5-8-1}
\expandafter\def\csname sga@target@IV.16.5.8.2\endcsname{ega4.c04.eq.16-5-8-2}
\expandafter\def\csname sga@target@IV.16.5.9\endcsname{ega4.c04.statement.16-5-9}
\expandafter\def\csname sga@target@IV.16.5.9.1\endcsname{ega4.c04.statement.16-5-9-1}
\expandafter\def\csname sga@target@IV.16.6\endcsname{ega4.c04.subsec.16-6}
\expandafter\def\csname sga@target@IV.16.6.1\endcsname{ega4.c04.statement.16-6-1}
\expandafter\def\csname sga@target@IV.16.6.1.1\endcsname{ega4.c04.statement.16-6-1-1}
\expandafter\def\csname sga@target@IV.16.6.1.2\endcsname{ega4.c04.statement.16-6-1-2}
\expandafter\def\csname sga@target@IV.16.6.2\endcsname{ega4.c04.thm.16-6-2}
\expandafter\def\csname sga@target@IV.16.6.2.1\endcsname{ega4.c04.eq.16-6-2-1}
\expandafter\def\csname sga@target@IV.16.6.2.2\endcsname{ega4.c04.eq.16-6-2-2}
\expandafter\def\csname sga@target@IV.16.6.2.3\endcsname{ega4.c04.eq.16-6-2-3}
\expandafter\def\csname sga@target@IV.16.6.2.4\endcsname{ega4.c04.eq.16-6-2-4}
\expandafter\def\csname sga@target@IV.16.6.2.5\endcsname{ega4.c04.eq.16-6-2-5}
\expandafter\def\csname sga@target@IV.16.6.3\endcsname{ega4.c04.def.16-6-3}
\expandafter\def\csname sga@target@IV.16.6.4\endcsname{ega4.c04.prop.16-6-4}
\expandafter\def\csname sga@target@IV.16.6.4.1\endcsname{ega4.c04.prop.16-6-4-1}
\expandafter\def\csname sga@target@IV.16.6.5\endcsname{ega4.c04.statement.16-6-5}
\expandafter\def\csname sga@target@IV.16.6.5.1\endcsname{ega4.c04.statement.16-6-5-1}
\expandafter\def\csname sga@target@IV.16.7\endcsname{ega4.c04.subsec.16-7}
\expandafter\def\csname sga@target@IV.16.7.1\endcsname{ega4.c04.statement.16-7-1}
\expandafter\def\csname sga@target@IV.16.7.1.1\endcsname{ega4.c04.eq.16-7-1-1}
\expandafter\def\csname sga@target@IV.16.7.1.2\endcsname{ega4.c04.eq.16-7-1-2}
\expandafter\def\csname sga@target@IV.16.7.10\endcsname{ega4.c04.rem.16-7-10}
\expandafter\def\csname sga@target@IV.16.7.2\endcsname{ega4.c04.statement.16-7-2}
\expandafter\def\csname sga@target@IV.16.7.2.1\endcsname{ega4.c04.statement.16-7-2-1}
\expandafter\def\csname sga@target@IV.16.7.3\endcsname{ega4.c04.prop.16-7-3}
\expandafter\def\csname sga@target@IV.16.7.4\endcsname{ega4.c04.statement.16-7-4}
\expandafter\def\csname sga@target@IV.16.7.4.1\endcsname{ega4.c04.statement.16-7-4-1}
\expandafter\def\csname sga@target@IV.16.7.5\endcsname{ega4.c04.statement.16-7-5}
\expandafter\def\csname sga@target@IV.16.7.5.1\endcsname{ega4.c04.statement.16-7-5-1}
\expandafter\def\csname sga@target@IV.16.7.5.2\endcsname{ega4.c04.statement.16-7-5-2}
\expandafter\def\csname sga@target@IV.16.7.5.3\endcsname{ega4.c04.statement.16-7-5-3}
\expandafter\def\csname sga@target@IV.16.7.5.4\endcsname{ega4.c04.statement.16-7-5-4}
\expandafter\def\csname sga@target@IV.16.7.6\endcsname{ega4.c04.prop.16-7-6}
\expandafter\def\csname sga@target@IV.16.7.7\endcsname{ega4.c04.statement.16-7-7}
\expandafter\def\csname sga@target@IV.16.7.7.1\endcsname{ega4.c04.eq.16-7-7-1}
\expandafter\def\csname sga@target@IV.16.7.7.2\endcsname{ega4.c04.eq.16-7-7-2}
\expandafter\def\csname sga@target@IV.16.7.8\endcsname{ega4.c04.statement.16-7-8}
\expandafter\def\csname sga@target@IV.16.7.8.1\endcsname{ega4.c04.statement.16-7-8-1}
\expandafter\def\csname sga@target@IV.16.7.8.2\endcsname{ega4.c04.eq.16-7-8-2}
\expandafter\def\csname sga@target@IV.16.7.9\endcsname{ega4.c04.statement.16-7-9}
\expandafter\def\csname sga@target@IV.16.7.9.1\endcsname{ega4.c04.statement.16-7-9-1}
\expandafter\def\csname sga@target@IV.16.8\endcsname{ega4.c04.subsec.16-8}
\expandafter\def\csname sga@target@IV.16.8.1\endcsname{ega4.c04.def.16-8-1}
\expandafter\def\csname sga@target@IV.16.8.10\endcsname{ega4.c04.cor.16-8-10}
\expandafter\def\csname sga@target@IV.16.8.11\endcsname{ega4.c04.statement.16-8-11}
\expandafter\def\csname sga@target@IV.16.8.2\endcsname{ega4.c04.statement.16-8-2}
\expandafter\def\csname sga@target@IV.16.8.3\endcsname{ega4.c04.statement.16-8-3}
\expandafter\def\csname sga@target@IV.16.8.3.1\endcsname{ega4.c04.eq.16-8-3-1}
\expandafter\def\csname sga@target@IV.16.8.4\endcsname{ega4.c04.prop.16-8-4}
\expandafter\def\csname sga@target@IV.16.8.4.1\endcsname{ega4.c04.prop.16-8-4-1}
\expandafter\def\csname sga@target@IV.16.8.5\endcsname{ega4.c04.statement.16-8-5}
\expandafter\def\csname sga@target@IV.16.8.5.1\endcsname{ega4.c04.eq.16-8-5-1}
\expandafter\def\csname sga@target@IV.16.8.5.2\endcsname{ega4.c04.eq.16-8-5-2}
\expandafter\def\csname sga@target@IV.16.8.6\endcsname{ega4.c04.prop.16-8-6}
\expandafter\def\csname sga@target@IV.16.8.7\endcsname{ega4.c04.statement.16-8-7}
\expandafter\def\csname sga@target@IV.16.8.7.1\endcsname{ega4.c04.eq.16-8-7-1}
\expandafter\def\csname sga@target@IV.16.8.8\endcsname{ega4.c04.prop.16-8-8}
\expandafter\def\csname sga@target@IV.16.8.8.1\endcsname{ega4.c04.eq.16-8-8-1}
\expandafter\def\csname sga@target@IV.16.8.8.2\endcsname{ega4.c04.eq.16-8-8-2}
\expandafter\def\csname sga@target@IV.16.8.9\endcsname{ega4.c04.prop.16-8-9}
\expandafter\def\csname sga@target@IV.16.8.9.1\endcsname{ega4.c04.lem.16-8-9-1}
\expandafter\def\csname sga@target@IV.16.8.9.2\endcsname{ega4.c04.lem.16-8-9-2}
\expandafter\def\csname sga@target@IV.16.8.9.3\endcsname{ega4.c04.lem.16-8-9-3}
\expandafter\def\csname sga@target@IV.16.8.9.4\endcsname{ega4.c04.eq.16-8-9-4}
\expandafter\def\csname sga@target@IV.16.9\endcsname{ega4.c04.subsec.16-9}
\expandafter\def\csname sga@target@IV.16.9.1\endcsname{ega4.c04.def.16-9-1}
\expandafter\def\csname sga@target@IV.16.9.10\endcsname{ega4.c04.prop.16-9-10}
\expandafter\def\csname sga@target@IV.16.9.11\endcsname{ega4.c04.cor.16-9-11}
\expandafter\def\csname sga@target@IV.16.9.12\endcsname{ega4.c04.statement.16-9-12}
\expandafter\def\csname sga@target@IV.16.9.13\endcsname{ega4.c04.prop.16-9-13}
\expandafter\def\csname sga@target@IV.16.9.13.1\endcsname{ega4.c04.prop.16-9-13-1}
\expandafter\def\csname sga@target@IV.16.9.13.2\endcsname{ega4.c04.lem.16-9-13-2}
\expandafter\def\csname sga@target@IV.16.9.13.3\endcsname{ega4.c04.lem.16-9-13-3}
\expandafter\def\csname sga@target@IV.16.9.13.4\endcsname{ega4.c04.eq.16-9-13-4}
\expandafter\def\csname sga@target@IV.16.9.13.5\endcsname{ega4.c04.cor.16-9-13-5}
\expandafter\def\csname sga@target@IV.16.9.13.6\endcsname{ega4.c04.cor.16-9-13-6}
\expandafter\def\csname sga@target@IV.16.9.13.7\endcsname{ega4.c04.eq.16-9-13-7}
\expandafter\def\csname sga@target@IV.16.9.2\endcsname{ega4.c04.def.16-9-2}
\expandafter\def\csname sga@target@IV.16.9.3\endcsname{ega4.c04.prop.16-9-3}
\expandafter\def\csname sga@target@IV.16.9.4\endcsname{ega4.c04.cor.16-9-4}
\expandafter\def\csname sga@target@IV.16.9.4.1\endcsname{ega4.c04.eq.16-9-4-1}
\expandafter\def\csname sga@target@IV.16.9.5\endcsname{ega4.c04.cor.16-9-5}
\expandafter\def\csname sga@target@IV.16.9.6\endcsname{ega4.c04.statement.16-9-6}
\expandafter\def\csname sga@target@IV.16.9.7\endcsname{ega4.c04.statement.16-9-7}
\expandafter\def\csname sga@target@IV.16.9.8\endcsname{ega4.c04.prop.16-9-8}
\expandafter\def\csname sga@target@IV.16.9.9\endcsname{ega4.c04.cor.16-9-9}
\expandafter\def\csname sga@target@IV.17.1.1\endcsname{ega4.c04.def.17-1-1}
\expandafter\def\csname sga@target@IV.17.1.1.1\endcsname{ega4.c04.def.17-1-1-1}
\expandafter\def\csname sga@target@IV.17.1.2\endcsname{ega4.c04.rem.17-1-2}
\expandafter\def\csname sga@target@IV.17.1.2.1\endcsname{ega4.c04.eq.17-1-2-1}
\expandafter\def\csname sga@target@IV.17.1.3\endcsname{ega4.c04.prop.17-1-3}
\expandafter\def\csname sga@target@IV.17.1.4\endcsname{ega4.c04.prop.17-1-4}
\expandafter\def\csname sga@target@IV.17.1.5\endcsname{ega4.c04.cor.17-1-5}
\expandafter\def\csname sga@target@IV.17.1.6\endcsname{ega4.c04.prop.17-1-6}
\expandafter\def\csname sga@target@IV.17.10.1\endcsname{ega4.c04.def.17-10-1}
\expandafter\def\csname sga@target@IV.17.10.2\endcsname{ega4.c04.prop.17-10-2}
\expandafter\def\csname sga@target@IV.17.10.3\endcsname{ega4.c04.cor.17-10-3}
\expandafter\def\csname sga@target@IV.17.10.3.1\endcsname{ega4.c04.cor.17-10-3-1}
\expandafter\def\csname sga@target@IV.17.10.4\endcsname{ega4.c04.cor.17-10-4}
\expandafter\def\csname sga@target@IV.17.10.4.1\endcsname{ega4.c04.eq.17-10-4-1}
\expandafter\def\csname sga@target@IV.17.11.1\endcsname{ega4.c04.thm.17-11-1}
\expandafter\def\csname sga@target@IV.17.11.1.1\endcsname{ega4.c04.eq.17-11-1-1}
\expandafter\def\csname sga@target@IV.17.11.1.2\endcsname{ega4.c04.eq.17-11-1-2}
\expandafter\def\csname sga@target@IV.17.11.2\endcsname{ega4.c04.cor.17-11-2}
\expandafter\def\csname sga@target@IV.17.11.3\endcsname{ega4.c04.cor.17-11-3}
\expandafter\def\csname sga@target@IV.17.11.4\endcsname{ega4.c04.cor.17-11-4}
\expandafter\def\csname sga@target@IV.17.11.5\endcsname{ega4.c04.prop.17-11-5}
\expandafter\def\csname sga@target@IV.17.12.1\endcsname{ega4.c04.thm.17-12-1}
\expandafter\def\csname sga@target@IV.17.12.2\endcsname{ega4.c04.cor.17-12-2}
\expandafter\def\csname sga@target@IV.17.12.3\endcsname{ega4.c04.cor.17-12-3}
\expandafter\def\csname sga@target@IV.17.12.4\endcsname{ega4.c04.prop.17-12-4}
\expandafter\def\csname sga@target@IV.17.12.5\endcsname{ega4.c04.prop.17-12-5}
\expandafter\def\csname sga@target@IV.17.12.6\endcsname{ega4.c04.cor.17-12-6}
\expandafter\def\csname sga@target@IV.17.12.7\endcsname{ega4.c04.ex.17-12-7}
\expandafter\def\csname sga@target@IV.17.12.8\endcsname{ega4.c04.rem.17-12-8}
\expandafter\def\csname sga@target@IV.17.13\endcsname{ega4.c04.subsec.17-13.key-iv-17-13}
\expandafter\def\csname sga@target@IV.17.13.1\endcsname{ega4.c04.statement.17-13-1}
\expandafter\def\csname sga@target@IV.17.13.1.1\endcsname{ega4.c04.statement.17-13-1-1}
\expandafter\def\csname sga@target@IV.17.13.1.2\endcsname{ega4.c04.statement.17-13-1-2}
\expandafter\def\csname sga@target@IV.17.13.1.3\endcsname{ega4.c04.eq.17-13-1-3}
\expandafter\def\csname sga@target@IV.17.13.10\endcsname{ega4.c04.cor.17-13-10}
\expandafter\def\csname sga@target@IV.17.13.11\endcsname{ega4.c04.statement.17-13-11}
\expandafter\def\csname sga@target@IV.17.13.11.1\endcsname{ega4.c04.statement.17-13-11-1}
\expandafter\def\csname sga@target@IV.17.13.12\endcsname{ega4.c04.cor.17-13-12}
\expandafter\def\csname sga@target@IV.17.13.12.1\endcsname{ega4.c04.cor.17-13-12-1}
\expandafter\def\csname sga@target@IV.17.13.12.2\endcsname{ega4.c04.eq.17-13-12-2}
\expandafter\def\csname sga@target@IV.17.13.2\endcsname{ega4.c04.prop.17-13-2}
\expandafter\def\csname sga@target@IV.17.13.2.1\endcsname{ega4.c04.eq.17-13-2-1}
\expandafter\def\csname sga@target@IV.17.13.2.2\endcsname{ega4.c04.eq.17-13-2-2}
\expandafter\def\csname sga@target@IV.17.13.3\endcsname{ega4.c04.def.17-13-3}
\expandafter\def\csname sga@target@IV.17.13.4\endcsname{ega4.c04.statement.17-13-4}
\expandafter\def\csname sga@target@IV.17.13.5\endcsname{ega4.c04.statement.17-13-5}
\expandafter\def\csname sga@target@IV.17.13.5.1\endcsname{ega4.c04.statement.17-13-5-1}
\expandafter\def\csname sga@target@IV.17.13.5.2\endcsname{ega4.c04.eq.17-13-5-2}
\expandafter\def\csname sga@target@IV.17.13.5.3\endcsname{ega4.c04.eq.17-13-5-3}
\expandafter\def\csname sga@target@IV.17.13.5.4\endcsname{ega4.c04.eq.17-13-5-4}
\expandafter\def\csname sga@target@IV.17.13.6\endcsname{ega4.c04.cor.17-13-6}
\expandafter\def\csname sga@target@IV.17.13.6.1\endcsname{ega4.c04.eq.17-13-6-1}
\expandafter\def\csname sga@target@IV.17.13.6.2\endcsname{ega4.c04.eq.17-13-6-2}
\expandafter\def\csname sga@target@IV.17.13.7\endcsname{ega4.c04.statement.17-13-7}
\expandafter\def\csname sga@target@IV.17.13.7.1\endcsname{ega4.c04.eq.17-13-7-1}
\expandafter\def\csname sga@target@IV.17.13.8\endcsname{ega4.c04.prop.17-13-8}
\expandafter\def\csname sga@target@IV.17.13.8.1\endcsname{ega4.c04.eq.17-13-8-1}
\expandafter\def\csname sga@target@IV.17.13.9\endcsname{ega4.c04.statement.17-13-9}
\expandafter\def\csname sga@target@IV.17.13.9.1\endcsname{ega4.c04.eq.17-13-9-1}
\expandafter\def\csname sga@target@IV.17.13.9.2\endcsname{ega4.c04.eq.17-13-9-2}
\expandafter\def\csname sga@target@IV.17.14.1\endcsname{ega4.c04.prop.17-14-1}
\expandafter\def\csname sga@target@IV.17.14.2\endcsname{ega4.c04.prop.17-14-2}
\expandafter\def\csname sga@target@IV.17.14.3\endcsname{ega4.c04.rem.17-14-3}
\expandafter\def\csname sga@target@IV.17.15.1\endcsname{ega4.c04.prop.17-15-1}
\expandafter\def\csname sga@target@IV.17.15.10\endcsname{ega4.c04.statement.17-15-10}
\expandafter\def\csname sga@target@IV.17.15.11\endcsname{ega4.c04.prop.17-15-11}
\expandafter\def\csname sga@target@IV.17.15.12\endcsname{ega4.c04.prop.17-15-12}
\expandafter\def\csname sga@target@IV.17.15.13\endcsname{ega4.c04.cor.17-15-13}
\expandafter\def\csname sga@target@IV.17.15.14\endcsname{ega4.c04.prop.17-15-14}
\expandafter\def\csname sga@target@IV.17.15.14.1\endcsname{ega4.c04.lem.17-15-14-1}
\expandafter\def\csname sga@target@IV.17.15.14.2\endcsname{ega4.c04.lem.17-15-14-2}
\expandafter\def\csname sga@target@IV.17.15.15\endcsname{ega4.c04.prop.17-15-15}
\expandafter\def\csname sga@target@IV.17.15.2\endcsname{ega4.c04.cor.17-15-2}
\expandafter\def\csname sga@target@IV.17.15.3\endcsname{ega4.c04.prop.17-15-3}
\expandafter\def\csname sga@target@IV.17.15.4\endcsname{ega4.c04.cor.17-15-4}
\expandafter\def\csname sga@target@IV.17.15.5\endcsname{ega4.c04.prop.17-15-5}
\expandafter\def\csname sga@target@IV.17.15.6\endcsname{ega4.c04.cor.17-15-6}
\expandafter\def\csname sga@target@IV.17.15.7\endcsname{ega4.c04.cor.17-15-7}
\expandafter\def\csname sga@target@IV.17.15.8\endcsname{ega4.c04.prop.17-15-8}
\expandafter\def\csname sga@target@IV.17.15.8.1\endcsname{ega4.c04.eq.17-15-8-1}
\expandafter\def\csname sga@target@IV.17.15.8.2\endcsname{ega4.c04.eq.17-15-8-2}
\expandafter\def\csname sga@target@IV.17.15.9\endcsname{ega4.c04.cor.17-15-9}
\expandafter\def\csname sga@target@IV.17.16\endcsname{ega4.c04.subsec.17-16.key-iv-17-16}
\expandafter\def\csname sga@target@IV.17.16.1\endcsname{ega4.c04.prop.17-16-1}
\expandafter\def\csname sga@target@IV.17.16.2\endcsname{ega4.c04.cor.17-16-2}
\expandafter\def\csname sga@target@IV.17.16.3\endcsname{ega4.c04.cor.17-16-3}
\expandafter\def\csname sga@target@IV.17.16.4\endcsname{ega4.c04.prop.17-16-4}
\expandafter\def\csname sga@target@IV.17.16.5\endcsname{ega4.c04.cor.17-16-5}
\expandafter\def\csname sga@target@IV.17.16.6\endcsname{ega4.c04.prop.17-16-6}
\expandafter\def\csname sga@target@IV.17.2.1\endcsname{ega4.c04.prop.17-2-1}
\expandafter\def\csname sga@target@IV.17.2.2\endcsname{ega4.c04.cor.17-2-2}
\expandafter\def\csname sga@target@IV.17.2.3\endcsname{ega4.c04.prop.17-2-3}
\expandafter\def\csname sga@target@IV.17.2.3.1\endcsname{ega4.c04.eq.17-2-3-1}
\expandafter\def\csname sga@target@IV.17.2.4\endcsname{ega4.c04.cor.17-2-4}
\expandafter\def\csname sga@target@IV.17.2.5\endcsname{ega4.c04.prop.17-2-5}
\expandafter\def\csname sga@target@IV.17.2.5.1\endcsname{ega4.c04.prop.17-2-5-1}
\expandafter\def\csname sga@target@IV.17.2.6\endcsname{ega4.c04.prop.17-2-6}
\expandafter\def\csname sga@target@IV.17.3.1\endcsname{ega4.c04.def.17-3-1}
\expandafter\def\csname sga@target@IV.17.3.2\endcsname{ega4.c04.subsec.17-3-2}
\expandafter\def\csname sga@target@IV.17.3.3\endcsname{ega4.c04.prop.17-3-3}
\expandafter\def\csname sga@target@IV.17.3.4\endcsname{ega4.c04.prop.17-3-4}
\expandafter\def\csname sga@target@IV.17.3.5\endcsname{ega4.c04.cor.17-3-5}
\expandafter\def\csname sga@target@IV.17.3.6\endcsname{ega4.c04.prop.17-3-6}
\expandafter\def\csname sga@target@IV.17.3.7\endcsname{ega4.c04.def.17-3-7}
\expandafter\def\csname sga@target@IV.17.3.8\endcsname{ega4.c04.prop.17-3-8}
\expandafter\def\csname sga@target@IV.17.3.9\endcsname{ega4.c04.cor.17-3-9}
\expandafter\def\csname sga@target@IV.17.4.1\endcsname{ega4.c04.thm.17-4-1}
\expandafter\def\csname sga@target@IV.17.4.1.1\endcsname{ega4.c04.eq.17-4-1-1}
\expandafter\def\csname sga@target@IV.17.4.1.2\endcsname{ega4.c04.rem.17-4-1-2}
\expandafter\def\csname sga@target@IV.17.4.10\endcsname{ega4.c04.rem.17-4-10}
\expandafter\def\csname sga@target@IV.17.4.2\endcsname{ega4.c04.cor.17-4-2}
\expandafter\def\csname sga@target@IV.17.4.3\endcsname{ega4.c04.cor.17-4-3}
\expandafter\def\csname sga@target@IV.17.4.4\endcsname{ega4.c04.prop.17-4-4}
\expandafter\def\csname sga@target@IV.17.4.5\endcsname{ega4.c04.statement.17-4-5}
\expandafter\def\csname sga@target@IV.17.4.6\endcsname{ega4.c04.prop.17-4-6}
\expandafter\def\csname sga@target@IV.17.4.7\endcsname{ega4.c04.cor.17-4-7}
\expandafter\def\csname sga@target@IV.17.4.8\endcsname{ega4.c04.cor.17-4-8}
\expandafter\def\csname sga@target@IV.17.4.9\endcsname{ega4.c04.prop.17-4-9}
\expandafter\def\csname sga@target@IV.17.5.1\endcsname{ega4.c04.thm.17-5-1}
\expandafter\def\csname sga@target@IV.17.5.1.1\endcsname{ega4.c04.eq.17-5-1-1}
\expandafter\def\csname sga@target@IV.17.5.1.2\endcsname{ega4.c04.eq.17-5-1-2}
\expandafter\def\csname sga@target@IV.17.5.2\endcsname{ega4.c04.cor.17-5-2}
\expandafter\def\csname sga@target@IV.17.5.3\endcsname{ega4.c04.prop.17-5-3}
\expandafter\def\csname sga@target@IV.17.5.4\endcsname{ega4.c04.rem.17-5-4}
\expandafter\def\csname sga@target@IV.17.5.5\endcsname{ega4.c04.prop.17-5-5}
\expandafter\def\csname sga@target@IV.17.5.6\endcsname{ega4.c04.cor.17-5-6}
\expandafter\def\csname sga@target@IV.17.5.7\endcsname{ega4.c04.prop.17-5-7}
\expandafter\def\csname sga@target@IV.17.5.8\endcsname{ega4.c04.prop.17-5-8}
\expandafter\def\csname sga@target@IV.17.6.1\endcsname{ega4.c04.thm.17-6-1}
\expandafter\def\csname sga@target@IV.17.6.2\endcsname{ega4.c04.cor.17-6-2}
\expandafter\def\csname sga@target@IV.17.6.3\endcsname{ega4.c04.prop.17-6-3}
\expandafter\def\csname sga@target@IV.17.6.4\endcsname{ega4.c04.prop.17-6-4}
\expandafter\def\csname sga@target@IV.17.7\endcsname{ega4.c04.subsec.17-7.key-iv-17-7}
\expandafter\def\csname sga@target@IV.17.7.1\endcsname{ega4.c04.prop.17-7-1}
\expandafter\def\csname sga@target@IV.17.7.10\endcsname{ega4.c04.prop.17-7-10}
\expandafter\def\csname sga@target@IV.17.7.11\endcsname{ega4.c04.prop.17-7-11}
\expandafter\def\csname sga@target@IV.17.7.11.1\endcsname{ega4.c04.lem.17-7-11-1}
\expandafter\def\csname sga@target@IV.17.7.2\endcsname{ega4.c04.cor.17-7-2}
\expandafter\def\csname sga@target@IV.17.7.3\endcsname{ega4.c04.cor.17-7-3}
\expandafter\def\csname sga@target@IV.17.7.4\endcsname{ega4.c04.prop.17-7-4}
\expandafter\def\csname sga@target@IV.17.7.5\endcsname{ega4.c04.lem.17-7-5}
\expandafter\def\csname sga@target@IV.17.7.6\endcsname{ega4.c04.rem.17-7-6}
\expandafter\def\csname sga@target@IV.17.7.7\endcsname{ega4.c04.prop.17-7-7}
\expandafter\def\csname sga@target@IV.17.7.8\endcsname{ega4.c04.prop.17-7-8}
\expandafter\def\csname sga@target@IV.17.7.9\endcsname{ega4.c04.cor.17-7-9}
\expandafter\def\csname sga@target@IV.17.8\endcsname{ega4.c04.subsec.17-8.key-iv-17-8}
\expandafter\def\csname sga@target@IV.17.8.1\endcsname{ega4.c04.prop.17-8-1}
\expandafter\def\csname sga@target@IV.17.8.2\endcsname{ega4.c04.prop.17-8-2}
\expandafter\def\csname sga@target@IV.17.8.3\endcsname{ega4.c04.rem.17-8-3}
\expandafter\def\csname sga@target@IV.17.9\endcsname{ega4.c04.subsec.17-9.key-iv-17-9}
\expandafter\def\csname sga@target@IV.17.9.1\endcsname{ega4.c04.thm.17-9-1}
\expandafter\def\csname sga@target@IV.17.9.2\endcsname{ega4.c04.cor.17-9-2}
\expandafter\def\csname sga@target@IV.17.9.3\endcsname{ega4.c04.cor.17-9-3}
\expandafter\def\csname sga@target@IV.17.9.4\endcsname{ega4.c04.cor.17-9-4}
\expandafter\def\csname sga@target@IV.17.9.5\endcsname{ega4.c04.cor.17-9-5}
\expandafter\def\csname sga@target@IV.17.9.6\endcsname{ega4.c04.prop.17-9-6}
\expandafter\def\csname sga@target@IV.18\endcsname{ega4.c04.sec.18.key-iv-18}
\expandafter\def\csname sga@target@IV.18.1.1\endcsname{ega4.c04.prop.18-1-1}
\expandafter\def\csname sga@target@IV.18.1.1.1\endcsname{ega4.c04.eq.18-1-1-1}
\expandafter\def\csname sga@target@IV.18.1.2\endcsname{ega4.c04.thm.18-1-2}
\expandafter\def\csname sga@target@IV.18.1.3\endcsname{ega4.c04.cor.18-1-3}
\expandafter\def\csname sga@target@IV.18.10.1\endcsname{ega4.c04.thm.18-10-1}
\expandafter\def\csname sga@target@IV.18.10.10\endcsname{ega4.c04.statement.18-10-10}
\expandafter\def\csname sga@target@IV.18.10.11\endcsname{ega4.c04.rem.18-10-11}
\expandafter\def\csname sga@target@IV.18.10.12\endcsname{ega4.c04.cor.18-10-12}
\expandafter\def\csname sga@target@IV.18.10.13\endcsname{ega4.c04.prop.18-10-13}
\expandafter\def\csname sga@target@IV.18.10.14\endcsname{ega4.c04.cor.18-10-14}
\expandafter\def\csname sga@target@IV.18.10.15\endcsname{ega4.c04.statement.18-10-15}
\expandafter\def\csname sga@target@IV.18.10.16\endcsname{ega4.c04.thm.18-10-16}
\expandafter\def\csname sga@target@IV.18.10.17\endcsname{ega4.c04.rem.18-10-17}
\expandafter\def\csname sga@target@IV.18.10.17.1\endcsname{ega4.c04.statement.18-10-17-1}
\expandafter\def\csname sga@target@IV.18.10.17.2\endcsname{ega4.c04.statement.18-10-17-2}
\expandafter\def\csname sga@target@IV.18.10.18\endcsname{ega4.c04.lem.18-10-18}
\expandafter\def\csname sga@target@IV.18.10.19\endcsname{ega4.c04.prop.18-10-19}
\expandafter\def\csname sga@target@IV.18.10.2\endcsname{ega4.c04.rem.18-10-2}
\expandafter\def\csname sga@target@IV.18.10.20\endcsname{ega4.c04.rem.18-10-20}
\expandafter\def\csname sga@target@IV.18.10.21\endcsname{ega4.c04.prop.18-10-21}
\expandafter\def\csname sga@target@IV.18.10.3\endcsname{ega4.c04.cor.18-10-3}
\expandafter\def\csname sga@target@IV.18.10.3.1\endcsname{ega4.c04.rem.18-10-3-1}
\expandafter\def\csname sga@target@IV.18.10.4\endcsname{ega4.c04.cor.18-10-4}
\expandafter\def\csname sga@target@IV.18.10.5\endcsname{ega4.c04.cor.18-10-5}
\expandafter\def\csname sga@target@IV.18.10.6\endcsname{ega4.c04.lem.18-10-6}
\expandafter\def\csname sga@target@IV.18.10.7\endcsname{ega4.c04.prop.18-10-7}
\expandafter\def\csname sga@target@IV.18.10.8\endcsname{ega4.c04.prop.18-10-8}
\expandafter\def\csname sga@target@IV.18.10.9\endcsname{ega4.c04.cor.18-10-9}
\expandafter\def\csname sga@target@IV.18.11.1\endcsname{ega4.c04.lem.18-11-1}
\expandafter\def\csname sga@target@IV.18.11.10\endcsname{ega4.c04.thm.18-11-10}
\expandafter\def\csname sga@target@IV.18.11.11\endcsname{ega4.c04.rem.18-11-11}
\expandafter\def\csname sga@target@IV.18.11.12\endcsname{ega4.c04.cor.18-11-12}
\expandafter\def\csname sga@target@IV.18.11.12.1\endcsname{ega4.c04.lem.18-11-12-1}
\expandafter\def\csname sga@target@IV.18.11.2\endcsname{ega4.c04.lem.18-11-2}
\expandafter\def\csname sga@target@IV.18.11.2.1\endcsname{ega4.c04.eq.18-11-2-1}
\expandafter\def\csname sga@target@IV.18.11.2.2\endcsname{ega4.c04.eq.18-11-2-2}
\expandafter\def\csname sga@target@IV.18.11.3\endcsname{ega4.c04.lem.18-11-3}
\expandafter\def\csname sga@target@IV.18.11.3.1\endcsname{ega4.c04.eq.18-11-3-1}
\expandafter\def\csname sga@target@IV.18.11.3.2\endcsname{ega4.c04.eq.18-11-3-2}
\expandafter\def\csname sga@target@IV.18.11.3.3\endcsname{ega4.c04.lem.18-11-3-3}
\expandafter\def\csname sga@target@IV.18.11.3.4\endcsname{ega4.c04.lem.18-11-3-4}
\expandafter\def\csname sga@target@IV.18.11.3.5\endcsname{ega4.c04.lem.18-11-3-5}
\expandafter\def\csname sga@target@IV.18.11.3.6\endcsname{ega4.c04.eq.18-11-3-6}
\expandafter\def\csname sga@target@IV.18.11.4\endcsname{ega4.c04.lem.18-11-4}
\expandafter\def\csname sga@target@IV.18.11.5\endcsname{ega4.c04.prop.18-11-5}
\expandafter\def\csname sga@target@IV.18.11.5.1\endcsname{ega4.c04.lem.18-11-5-1}
\expandafter\def\csname sga@target@IV.18.11.5.2\endcsname{ega4.c04.lem.18-11-5-2}
\expandafter\def\csname sga@target@IV.18.11.6\endcsname{ega4.c04.rem.18-11-6}
\expandafter\def\csname sga@target@IV.18.11.7\endcsname{ega4.c04.cor.18-11-7}
\expandafter\def\csname sga@target@IV.18.11.8\endcsname{ega4.c04.cor.18-11-8}
\expandafter\def\csname sga@target@IV.18.11.9\endcsname{ega4.c04.cor.18-11-9}
\expandafter\def\csname sga@target@IV.18.12.1\endcsname{ega4.c04.thm.18-12-1}
\expandafter\def\csname sga@target@IV.18.12.10\endcsname{ega4.c04.cor.18-12-10}
\expandafter\def\csname sga@target@IV.18.12.11\endcsname{ega4.c04.cor.18-12-11}
\expandafter\def\csname sga@target@IV.18.12.12\endcsname{ega4.c04.prop.18-12-12}
\expandafter\def\csname sga@target@IV.18.12.13\endcsname{ega4.c04.cor.18-12-13}
\expandafter\def\csname sga@target@IV.18.12.14\endcsname{ega4.c04.rem.18-12-14}
\expandafter\def\csname sga@target@IV.18.12.15\endcsname{ega4.c04.prop.18-12-15}
\expandafter\def\csname sga@target@IV.18.12.16\endcsname{ega4.c04.prop.18-12-16}
\expandafter\def\csname sga@target@IV.18.12.17\endcsname{ega4.c04.cor.18-12-17}
\expandafter\def\csname sga@target@IV.18.12.2\endcsname{ega4.c04.rem.18-12-2}
\expandafter\def\csname sga@target@IV.18.12.3\endcsname{ega4.c04.cor.18-12-3}
\expandafter\def\csname sga@target@IV.18.12.4\endcsname{ega4.c04.cor.18-12-4}
\expandafter\def\csname sga@target@IV.18.12.5\endcsname{ega4.c04.rem.18-12-5}
\expandafter\def\csname sga@target@IV.18.12.6\endcsname{ega4.c04.cor.18-12-6}
\expandafter\def\csname sga@target@IV.18.12.7\endcsname{ega4.c04.prop.18-12-7}
\expandafter\def\csname sga@target@IV.18.12.7.1\endcsname{ega4.c04.lem.18-12-7-1}
\expandafter\def\csname sga@target@IV.18.12.8\endcsname{ega4.c04.prop.18-12-8}
\expandafter\def\csname sga@target@IV.18.12.9\endcsname{ega4.c04.rem.18-12-9}
\expandafter\def\csname sga@target@IV.18.2.1\endcsname{ega4.c04.statement.18-2-1}
\expandafter\def\csname sga@target@IV.18.2.1.1\endcsname{ega4.c04.statement.18-2-1-1}
\expandafter\def\csname sga@target@IV.18.2.1.2\endcsname{ega4.c04.eq.18-2-1-2}
\expandafter\def\csname sga@target@IV.18.2.10\endcsname{ega4.c04.statement.18-2-10}
\expandafter\def\csname sga@target@IV.18.2.2\endcsname{ega4.c04.statement.18-2-2}
\expandafter\def\csname sga@target@IV.18.2.2.1\endcsname{ega4.c04.eq.18-2-2-1}
\expandafter\def\csname sga@target@IV.18.2.2.2\endcsname{ega4.c04.eq.18-2-2-2}
\expandafter\def\csname sga@target@IV.18.2.2.3\endcsname{ega4.c04.eq.18-2-2-3}
\expandafter\def\csname sga@target@IV.18.2.2.4\endcsname{ega4.c04.eq.18-2-2-4}
\expandafter\def\csname sga@target@IV.18.2.3\endcsname{ega4.c04.prop.18-2-3}
\expandafter\def\csname sga@target@IV.18.2.4\endcsname{ega4.c04.cor.18-2-4}
\expandafter\def\csname sga@target@IV.18.2.5\endcsname{ega4.c04.cor.18-2-5}
\expandafter\def\csname sga@target@IV.18.2.6\endcsname{ega4.c04.cor.18-2-6}
\expandafter\def\csname sga@target@IV.18.2.7\endcsname{ega4.c04.item.18-2-7}
\expandafter\def\csname sga@target@IV.18.2.7.1\endcsname{ega4.c04.eq.18-2-7-1}
\expandafter\def\csname sga@target@IV.18.2.7.2\endcsname{ega4.c04.eq.18-2-7-2}
\expandafter\def\csname sga@target@IV.18.2.8\endcsname{ega4.c04.prop.18-2-8}
\expandafter\def\csname sga@target@IV.18.2.9\endcsname{ega4.c04.cor.18-2-9}
\expandafter\def\csname sga@target@IV.18.3.1\endcsname{ega4.c04.prop.18-3-1}
\expandafter\def\csname sga@target@IV.18.3.2\endcsname{ega4.c04.prop.18-3-2}
\expandafter\def\csname sga@target@IV.18.3.2.1\endcsname{ega4.c04.lem.18-3-2-1}
\expandafter\def\csname sga@target@IV.18.3.2.2\endcsname{ega4.c04.eq.18-3-2-2}
\expandafter\def\csname sga@target@IV.18.3.3\endcsname{ega4.c04.statement.18-3-3}
\expandafter\def\csname sga@target@IV.18.3.4\endcsname{ega4.c04.thm.18-3-4}
\expandafter\def\csname sga@target@IV.18.4.1\endcsname{ega4.c04.lem.18-4-1}
\expandafter\def\csname sga@target@IV.18.4.10\endcsname{ega4.c04.prop.18-4-10}
\expandafter\def\csname sga@target@IV.18.4.11\endcsname{ega4.c04.statement.18-4-11}
\expandafter\def\csname sga@target@IV.18.4.12\endcsname{ega4.c04.cor.18-4-12}
\expandafter\def\csname sga@target@IV.18.4.13\endcsname{ega4.c04.cor.18-4-13}
\expandafter\def\csname sga@target@IV.18.4.14\endcsname{ega4.c04.cor.18-4-14}
\expandafter\def\csname sga@target@IV.18.4.2\endcsname{ega4.c04.prop.18-4-2}
\expandafter\def\csname sga@target@IV.18.4.3\endcsname{ega4.c04.cor.18-4-3}
\expandafter\def\csname sga@target@IV.18.4.4\endcsname{ega4.c04.lem.18-4-4}
\expandafter\def\csname sga@target@IV.18.4.5\endcsname{ega4.c04.prop.18-4-5}
\expandafter\def\csname sga@target@IV.18.4.6\endcsname{ega4.c04.thm.18-4-6}
\expandafter\def\csname sga@target@IV.18.4.6.1\endcsname{ega4.c04.lem.18-4-6-1}
\expandafter\def\csname sga@target@IV.18.4.7\endcsname{ega4.c04.cor.18-4-7}
\expandafter\def\csname sga@target@IV.18.4.8\endcsname{ega4.c04.cor.18-4-8}
\expandafter\def\csname sga@target@IV.18.4.9\endcsname{ega4.c04.prop.18-4-9}
\expandafter\def\csname sga@target@IV.18.5.1\endcsname{ega4.c04.statement.18-5-1}
\expandafter\def\csname sga@target@IV.18.5.1.1\endcsname{ega4.c04.eq.18-5-1-1}
\expandafter\def\csname sga@target@IV.18.5.10\endcsname{ega4.c04.prop.18-5-10}
\expandafter\def\csname sga@target@IV.18.5.11\endcsname{ega4.c04.thm.18-5-11}
\expandafter\def\csname sga@target@IV.18.5.12\endcsname{ega4.c04.cor.18-5-12}
\expandafter\def\csname sga@target@IV.18.5.13\endcsname{ega4.c04.rem.18-5-13}
\expandafter\def\csname sga@target@IV.18.5.13.1\endcsname{ega4.c04.lem.18-5-13-1}
\expandafter\def\csname sga@target@IV.18.5.14\endcsname{ega4.c04.prop.18-5-14}
\expandafter\def\csname sga@target@IV.18.5.15\endcsname{ega4.c04.prop.18-5-15}
\expandafter\def\csname sga@target@IV.18.5.16\endcsname{ega4.c04.item.18-5-16}
\expandafter\def\csname sga@target@IV.18.5.17\endcsname{ega4.c04.thm.18-5-17}
\expandafter\def\csname sga@target@IV.18.5.18\endcsname{ega4.c04.rem.18-5-18}
\expandafter\def\csname sga@target@IV.18.5.19\endcsname{ega4.c04.prop.18-5-19}
\expandafter\def\csname sga@target@IV.18.5.2\endcsname{ega4.c04.statement.18-5-2}
\expandafter\def\csname sga@target@IV.18.5.3\endcsname{ega4.c04.lem.18-5-3}
\expandafter\def\csname sga@target@IV.18.5.3.1\endcsname{ega4.c04.lem.18-5-3-1}
\expandafter\def\csname sga@target@IV.18.5.4\endcsname{ega4.c04.prop.18-5-4}
\expandafter\def\csname sga@target@IV.18.5.4.1\endcsname{ega4.c04.eq.18-5-4-1}
\expandafter\def\csname sga@target@IV.18.5.4.2\endcsname{ega4.c04.eq.18-5-4-2}
\expandafter\def\csname sga@target@IV.18.5.4.3\endcsname{ega4.c04.lem.18-5-4-3}
\expandafter\def\csname sga@target@IV.18.5.4.4\endcsname{ega4.c04.rem.18-5-4-4}
\expandafter\def\csname sga@target@IV.18.5.5\endcsname{ega4.c04.def.18-5-5}
\expandafter\def\csname sga@target@IV.18.5.6\endcsname{ega4.c04.prop.18-5-6}
\expandafter\def\csname sga@target@IV.18.5.7\endcsname{ega4.c04.rem.18-5-7}
\expandafter\def\csname sga@target@IV.18.5.8\endcsname{ega4.c04.def.18-5-8}
\expandafter\def\csname sga@target@IV.18.5.9\endcsname{ega4.c04.prop.18-5-9}
\expandafter\def\csname sga@target@IV.18.6.1\endcsname{ega4.c04.statement.18-6-1}
\expandafter\def\csname sga@target@IV.18.6.10\endcsname{ega4.c04.cor.18-6-10}
\expandafter\def\csname sga@target@IV.18.6.11\endcsname{ega4.c04.cor.18-6-11}
\expandafter\def\csname sga@target@IV.18.6.12\endcsname{ega4.c04.prop.18-6-12}
\expandafter\def\csname sga@target@IV.18.6.13\endcsname{ega4.c04.cor.18-6-13}
\expandafter\def\csname sga@target@IV.18.6.14\endcsname{ega4.c04.prop.18-6-14}
\expandafter\def\csname sga@target@IV.18.6.15\endcsname{ega4.c04.cor.18-6-15}
\expandafter\def\csname sga@target@IV.18.6.2\endcsname{ega4.c04.lem.18-6-2}
\expandafter\def\csname sga@target@IV.18.6.3\endcsname{ega4.c04.lem.18-6-3}
\expandafter\def\csname sga@target@IV.18.6.4\endcsname{ega4.c04.statement.18-6-4}
\expandafter\def\csname sga@target@IV.18.6.5\endcsname{ega4.c04.def.18-6-5}
\expandafter\def\csname sga@target@IV.18.6.6\endcsname{ega4.c04.thm.18-6-6}
\expandafter\def\csname sga@target@IV.18.6.7\endcsname{ega4.c04.statement.18-6-7}
\expandafter\def\csname sga@target@IV.18.6.8\endcsname{ega4.c04.prop.18-6-8}
\expandafter\def\csname sga@target@IV.18.6.9\endcsname{ega4.c04.thm.18-6-9}
\expandafter\def\csname sga@target@IV.18.6.9.1\endcsname{ega4.c04.lem.18-6-9-1}
\expandafter\def\csname sga@target@IV.18.6.9.2\endcsname{ega4.c04.lem.18-6-9-2}
\expandafter\def\csname sga@target@IV.18.7.1\endcsname{ega4.c04.statement.18-7-1}
\expandafter\def\csname sga@target@IV.18.7.2\endcsname{ega4.c04.prop.18-7-2}
\expandafter\def\csname sga@target@IV.18.7.3\endcsname{ega4.c04.cor.18-7-3}
\expandafter\def\csname sga@target@IV.18.7.4\endcsname{ega4.c04.cor.18-7-4}
\expandafter\def\csname sga@target@IV.18.7.5\endcsname{ega4.c04.prop.18-7-5}
\expandafter\def\csname sga@target@IV.18.7.5.1\endcsname{ega4.c04.lem.18-7-5-1}
\expandafter\def\csname sga@target@IV.18.7.6\endcsname{ega4.c04.cor.18-7-6}
\expandafter\def\csname sga@target@IV.18.7.7\endcsname{ega4.c04.rem.18-7-7}
\expandafter\def\csname sga@target@IV.18.8.1\endcsname{ega4.c04.prop.18-8-1}
\expandafter\def\csname sga@target@IV.18.8.10\endcsname{ega4.c04.prop.18-8-10}
\expandafter\def\csname sga@target@IV.18.8.11\endcsname{ega4.c04.rem.18-8-11}
\expandafter\def\csname sga@target@IV.18.8.12\endcsname{ega4.c04.prop.18-8-12}
\expandafter\def\csname sga@target@IV.18.8.13\endcsname{ega4.c04.cor.18-8-13}
\expandafter\def\csname sga@target@IV.18.8.14\endcsname{ega4.c04.cor.18-8-14}
\expandafter\def\csname sga@target@IV.18.8.15\endcsname{ega4.c04.prop.18-8-15}
\expandafter\def\csname sga@target@IV.18.8.16\endcsname{ega4.c04.cor.18-8-16}
\expandafter\def\csname sga@target@IV.18.8.17\endcsname{ega4.c04.prop.18-8-17}
\expandafter\def\csname sga@target@IV.18.8.18\endcsname{ega4.c04.prop.18-8-18}
\expandafter\def\csname sga@target@IV.18.8.2\endcsname{ega4.c04.def.18-8-2}
\expandafter\def\csname sga@target@IV.18.8.3\endcsname{ega4.c04.rem.18-8-3}
\expandafter\def\csname sga@target@IV.18.8.4\endcsname{ega4.c04.lem.18-8-4}
\expandafter\def\csname sga@target@IV.18.8.5\endcsname{ega4.c04.lem.18-8-5}
\expandafter\def\csname sga@target@IV.18.8.6\endcsname{ega4.c04.statement.18-8-6}
\expandafter\def\csname sga@target@IV.18.8.6.1\endcsname{ega4.c04.eq.18-8-6-1}
\expandafter\def\csname sga@target@IV.18.8.6.2\endcsname{ega4.c04.eq.18-8-6-2}
\expandafter\def\csname sga@target@IV.18.8.7\endcsname{ega4.c04.def.18-8-7}
\expandafter\def\csname sga@target@IV.18.8.8\endcsname{ega4.c04.prop.18-8-8}
\expandafter\def\csname sga@target@IV.18.8.8.1\endcsname{ega4.c04.statement.18-8-8-1}
\expandafter\def\csname sga@target@IV.18.8.9\endcsname{ega4.c04.statement.18-8-9}
\expandafter\def\csname sga@target@IV.18.9.1\endcsname{ega4.c04.thm.18-9-1}
\expandafter\def\csname sga@target@IV.18.9.10\endcsname{ega4.c04.rem.18-9-10}
\expandafter\def\csname sga@target@IV.18.9.11\endcsname{ega4.c04.cor.18-9-11}
\expandafter\def\csname sga@target@IV.18.9.2\endcsname{ega4.c04.cor.18-9-2}
\expandafter\def\csname sga@target@IV.18.9.3\endcsname{ega4.c04.cor.18-9-3}
\expandafter\def\csname sga@target@IV.18.9.4\endcsname{ega4.c04.cor.18-9-4}
\expandafter\def\csname sga@target@IV.18.9.5\endcsname{ega4.c04.cor.18-9-5}
\expandafter\def\csname sga@target@IV.18.9.5.1\endcsname{ega4.c04.eq.18-9-5-1}
\expandafter\def\csname sga@target@IV.18.9.6\endcsname{ega4.c04.statement.18-9-6}
\expandafter\def\csname sga@target@IV.18.9.7\endcsname{ega4.c04.thm.18-9-7}
\expandafter\def\csname sga@target@IV.18.9.7.1\endcsname{ega4.c04.rem.18-9-7-1}
\expandafter\def\csname sga@target@IV.18.9.7.2\endcsname{ega4.c04.statement.18-9-7-2}
\expandafter\def\csname sga@target@IV.18.9.7.3\endcsname{ega4.c04.lem.18-9-7-3}
\expandafter\def\csname sga@target@IV.18.9.7.4\endcsname{ega4.c04.eq.18-9-7-4}
\expandafter\def\csname sga@target@IV.18.9.7.5\endcsname{ega4.c04.lem.18-9-7-5}
\expandafter\def\csname sga@target@IV.18.9.7.6\endcsname{ega4.c04.rem.18-9-7-6}
\expandafter\def\csname sga@target@IV.18.9.7.7\endcsname{ega4.c04.lem.18-9-7-7}
\expandafter\def\csname sga@target@IV.18.9.7.8\endcsname{ega4.c04.rem.18-9-7-8}
\expandafter\def\csname sga@target@IV.18.9.8\endcsname{ega4.c04.cor.18-9-8}
\expandafter\def\csname sga@target@IV.18.9.9\endcsname{ega4.c04.cor.18-9-9}
\expandafter\def\csname sga@target@IV.19.1.1\endcsname{ega4.c04.prop.19-1-1}
\expandafter\def\csname sga@target@IV.19.1.2\endcsname{ega4.c04.cor.19-1-2}
\expandafter\def\csname sga@target@IV.19.1.3\endcsname{ega4.c04.def.19-1-3}
\expandafter\def\csname sga@target@IV.19.1.4\endcsname{ega4.c04.prop.19-1-4}
\expandafter\def\csname sga@target@IV.19.1.4.1\endcsname{ega4.c04.eq.19-1-4-1}
\expandafter\def\csname sga@target@IV.19.1.5\endcsname{ega4.c04.prop.19-1-5}
\expandafter\def\csname sga@target@IV.19.1.5.1\endcsname{ega4.c04.item.19-1-5-1}
\expandafter\def\csname sga@target@IV.19.1.6\endcsname{ega4.c04.cor.19-1-6}
\expandafter\def\csname sga@target@IV.19.1.6.1\endcsname{ega4.c04.lem.19-1-6-1}
\expandafter\def\csname sga@target@IV.19.1.7\endcsname{ega4.c04.item.19-1-7}
\expandafter\def\csname sga@target@IV.19.1.8\endcsname{ega4.c04.statement.19-1-8}
\expandafter\def\csname sga@target@IV.19.2.1\endcsname{ega4.c04.def.19-2-1}
\expandafter\def\csname sga@target@IV.19.2.2\endcsname{ega4.c04.def.19-2-2}
\expandafter\def\csname sga@target@IV.19.2.3\endcsname{ega4.c04.prop.19-2-3}
\expandafter\def\csname sga@target@IV.19.2.4\endcsname{ega4.c04.prop.19-2-4}
\expandafter\def\csname sga@target@IV.19.2.5\endcsname{ega4.c04.item.19-2-5}
\expandafter\def\csname sga@target@IV.19.2.6\endcsname{ega4.c04.cor.19-2-6}
\expandafter\def\csname sga@target@IV.19.2.7\endcsname{ega4.c04.prop.19-2-7}
\expandafter\def\csname sga@target@IV.19.2.8\endcsname{ega4.c04.prop.19-2-8}
\expandafter\def\csname sga@target@IV.19.2.9\endcsname{ega4.c04.prop.19-2-9}
\expandafter\def\csname sga@target@IV.19.3.1\endcsname{ega4.c04.def.19-3-1}
\expandafter\def\csname sga@target@IV.19.3.10\endcsname{ega4.c04.cor.19-3-10}
\expandafter\def\csname sga@target@IV.19.3.2\endcsname{ega4.c04.prop.19-3-2}
\expandafter\def\csname sga@target@IV.19.3.2.1\endcsname{ega4.c04.lem.19-3-2-1}
\expandafter\def\csname sga@target@IV.19.3.3\endcsname{ega4.c04.cor.19-3-3}
\expandafter\def\csname sga@target@IV.19.3.4\endcsname{ega4.c04.cor.19-3-4}
\expandafter\def\csname sga@target@IV.19.3.5\endcsname{ega4.c04.cor.19-3-5}
\expandafter\def\csname sga@target@IV.19.3.6\endcsname{ega4.c04.def.19-3-6}
\expandafter\def\csname sga@target@IV.19.3.7\endcsname{ega4.c04.prop.19-3-7}
\expandafter\def\csname sga@target@IV.19.3.8\endcsname{ega4.c04.cor.19-3-8}
\expandafter\def\csname sga@target@IV.19.3.9\endcsname{ega4.c04.prop.19-3-9}
\expandafter\def\csname sga@target@IV.19.4.1\endcsname{ega4.c04.statement.19-4-1}
\expandafter\def\csname sga@target@IV.19.4.1.1\endcsname{ega4.c04.subsec.19-4-1-1}
\expandafter\def\csname sga@target@IV.19.4.10\endcsname{ega4.c04.rem.19-4-10}
\expandafter\def\csname sga@target@IV.19.4.11\endcsname{ega4.c04.prop.19-4-11}
\expandafter\def\csname sga@target@IV.19.4.11.1\endcsname{ega4.c04.eq.19-4-11-1}
\expandafter\def\csname sga@target@IV.19.4.2\endcsname{ega4.c04.lem.19-4-2}
\expandafter\def\csname sga@target@IV.19.4.3\endcsname{ega4.c04.cor.19-4-3}
\expandafter\def\csname sga@target@IV.19.4.4\endcsname{ega4.c04.prop.19-4-4}
\expandafter\def\csname sga@target@IV.19.4.5\endcsname{ega4.c04.cor.19-4-5}
\expandafter\def\csname sga@target@IV.19.4.6\endcsname{ega4.c04.prop.19-4-6}
\expandafter\def\csname sga@target@IV.19.4.7\endcsname{ega4.c04.cor.19-4-7}
\expandafter\def\csname sga@target@IV.19.4.8\endcsname{ega4.c04.prop.19-4-8}
\expandafter\def\csname sga@target@IV.19.4.9\endcsname{ega4.c04.cor.19-4-9}
\expandafter\def\csname sga@target@IV.19.5.1\endcsname{ega4.c04.thm.19-5-1}
\expandafter\def\csname sga@target@IV.19.5.1.1\endcsname{ega4.c04.eq.19-5-1-1}
\expandafter\def\csname sga@target@IV.19.5.2\endcsname{ega4.c04.cor.19-5-2}
\expandafter\def\csname sga@target@IV.19.5.3\endcsname{ega4.c04.prop.19-5-3}
\expandafter\def\csname sga@target@IV.19.5.4\endcsname{ega4.c04.statement.19-5-4}
\expandafter\def\csname sga@target@IV.19.5.5\endcsname{ega4.c04.thm.19-5-5}
\expandafter\def\csname sga@target@IV.19.5.5.1\endcsname{ega4.c04.statement.19-5-5-1}
\expandafter\def\csname sga@target@IV.19.5.5.2\endcsname{ega4.c04.statement.19-5-5-2}
\expandafter\def\csname sga@target@IV.19.5.5.3\endcsname{ega4.c04.statement.19-5-5-3}
\expandafter\def\csname sga@target@IV.19.5.5.4\endcsname{ega4.c04.statement.19-5-5-4}
\expandafter\def\csname sga@target@IV.19.5.5.5\endcsname{ega4.c04.statement.19-5-5-5}
\expandafter\def\csname sga@target@IV.19.5.6\endcsname{ega4.c04.statement.19-5-6}
\expandafter\def\csname sga@target@IV.19.6.1\endcsname{ega4.c04.statement.19-6-1}
\expandafter\def\csname sga@target@IV.19.6.1.1\endcsname{ega4.c04.eq.19-6-1-1}
\expandafter\def\csname sga@target@IV.19.6.2\endcsname{ega4.c04.statement.19-6-2}
\expandafter\def\csname sga@target@IV.19.6.2.1\endcsname{ega4.c04.statement.19-6-2-1}
\expandafter\def\csname sga@target@IV.19.6.3\endcsname{ega4.c04.prop.19-6-3}
\expandafter\def\csname sga@target@IV.19.6.4\endcsname{ega4.c04.cor.19-6-4}
\expandafter\def\csname sga@target@IV.19.7\endcsname{ega4.c04.subsec.19-7.key-iv-19-7}
\expandafter\def\csname sga@target@IV.19.7.1\endcsname{ega4.c04.thm.19-7-1}
\expandafter\def\csname sga@target@IV.19.7.1.1\endcsname{ega4.c04.statement.19-7-1-1}
\expandafter\def\csname sga@target@IV.19.7.1.10\endcsname{ega4.c04.statement.19-7-1-10}
\expandafter\def\csname sga@target@IV.19.7.1.11\endcsname{ega4.c04.statement.19-7-1-11}
\expandafter\def\csname sga@target@IV.19.7.1.12\endcsname{ega4.c04.statement.19-7-1-12}
\expandafter\def\csname sga@target@IV.19.7.1.13\endcsname{ega4.c04.statement.19-7-1-13}
\expandafter\def\csname sga@target@IV.19.7.1.14\endcsname{ega4.c04.statement.19-7-1-14}
\expandafter\def\csname sga@target@IV.19.7.1.15\endcsname{ega4.c04.statement.19-7-1-15}
\expandafter\def\csname sga@target@IV.19.7.1.16\endcsname{ega4.c04.statement.19-7-1-16}
\expandafter\def\csname sga@target@IV.19.7.1.17\endcsname{ega4.c04.statement.19-7-1-17}
\expandafter\def\csname sga@target@IV.19.7.1.18\endcsname{ega4.c04.statement.19-7-1-18}
\expandafter\def\csname sga@target@IV.19.7.1.19\endcsname{ega4.c04.statement.19-7-1-19}
\expandafter\def\csname sga@target@IV.19.7.1.2\endcsname{ega4.c04.statement.19-7-1-2}
\expandafter\def\csname sga@target@IV.19.7.1.20\endcsname{ega4.c04.statement.19-7-1-20}
\expandafter\def\csname sga@target@IV.19.7.1.21\endcsname{ega4.c04.statement.19-7-1-21}
\expandafter\def\csname sga@target@IV.19.7.1.22\endcsname{ega4.c04.statement.19-7-1-22}
\expandafter\def\csname sga@target@IV.19.7.1.23\endcsname{ega4.c04.statement.19-7-1-23}
\expandafter\def\csname sga@target@IV.19.7.1.3\endcsname{ega4.c04.statement.19-7-1-3}
\expandafter\def\csname sga@target@IV.19.7.1.5\endcsname{ega4.c04.statement.19-7-1-5}
\expandafter\def\csname sga@target@IV.19.7.1.6\endcsname{ega4.c04.statement.19-7-1-6}
\expandafter\def\csname sga@target@IV.19.7.1.7\endcsname{ega4.c04.lem.19-7-1-7}
\expandafter\def\csname sga@target@IV.19.7.1.8\endcsname{ega4.c04.statement.19-7-1-8}
\expandafter\def\csname sga@target@IV.19.7.1.9\endcsname{ega4.c04.statement.19-7-1-9}
\expandafter\def\csname sga@target@IV.19.7.2\endcsname{ega4.c04.statement.19-7-2}
\expandafter\def\csname sga@target@IV.19.7.3\endcsname{ega4.c04.cor.19-7-3}
\expandafter\def\csname sga@target@IV.19.7.3.1\endcsname{ega4.c04.eq.19-7-3-1}
\expandafter\def\csname sga@target@IV.19.7.4\endcsname{ega4.c04.cor.19-7-4}
\expandafter\def\csname sga@target@IV.19.7.4.1\endcsname{ega4.c04.eq.19-7-4-1}
\expandafter\def\csname sga@target@IV.19.7.5\endcsname{ega4.c04.cor.19-7-5}
\expandafter\def\csname sga@target@IV.19.8.1\endcsname{ega4.c04.prop.19-8-1}
\expandafter\def\csname sga@target@IV.19.8.2\endcsname{ega4.c04.prop.19-8-2}
\expandafter\def\csname sga@target@IV.19.8.3\endcsname{ega4.c04.rem.19-8-3}
\expandafter\def\csname sga@target@IV.19.9.1\endcsname{ega4.c04.statement.19-9-1}
\expandafter\def\csname sga@target@IV.19.9.1.1\endcsname{ega4.c04.eq.19-9-1-1}
\expandafter\def\csname sga@target@IV.19.9.1.2\endcsname{ega4.c04.eq.19-9-1-2}
\expandafter\def\csname sga@target@IV.19.9.1.3\endcsname{ega4.c04.eq.19-9-1-3}
\expandafter\def\csname sga@target@IV.19.9.2\endcsname{ega4.c04.lem.19-9-2}
\expandafter\def\csname sga@target@IV.19.9.3\endcsname{ega4.c04.prop.19-9-3}
\expandafter\def\csname sga@target@IV.19.9.4\endcsname{ega4.c04.cor.19-9-4}
\expandafter\def\csname sga@target@IV.19.9.5\endcsname{ega4.c04.prop.19-9-5}
\expandafter\def\csname sga@target@IV.19.9.6\endcsname{ega4.c04.prop.19-9-6}
\expandafter\def\csname sga@target@IV.19.9.7\endcsname{ega4.c04.cor.19-9-7}
\expandafter\def\csname sga@target@IV.19.9.8\endcsname{ega4.c04.prop.19-9-8}
\expandafter\def\csname sga@target@IV.19.9.9\endcsname{ega4.c04.rem.19-9-9}
\expandafter\def\csname sga@target@IV.2.1.1\endcsname{ega4.c04.statement.2-1-1}
\expandafter\def\csname sga@target@IV.2.1.10\endcsname{ega4.c04.cor.2-1-10}
\expandafter\def\csname sga@target@IV.2.1.11\endcsname{ega4.c04.prop.2-1-11}
\expandafter\def\csname sga@target@IV.2.1.12\endcsname{ega4.c04.prop.2-1-12}
\expandafter\def\csname sga@target@IV.2.1.13\endcsname{ega4.c04.prop.2-1-13}
\expandafter\def\csname sga@target@IV.2.1.14\endcsname{ega4.c04.prop.2-1-14}
\expandafter\def\csname sga@target@IV.2.1.2\endcsname{ega4.c04.prop.2-1-2}
\expandafter\def\csname sga@target@IV.2.1.3\endcsname{ega4.c04.prop.2-1-3}
\expandafter\def\csname sga@target@IV.2.1.4\endcsname{ega4.c04.prop.2-1-4}
\expandafter\def\csname sga@target@IV.2.1.5\endcsname{ega4.c04.prop.2-1-5}
\expandafter\def\csname sga@target@IV.2.1.6\endcsname{ega4.c04.cor.2-1-6}
\expandafter\def\csname sga@target@IV.2.1.7\endcsname{ega4.c04.cor.2-1-7}
\expandafter\def\csname sga@target@IV.2.1.8\endcsname{ega4.c04.prop.2-1-8}
\expandafter\def\csname sga@target@IV.2.1.9\endcsname{ega4.c04.cor.2-1-9}
\expandafter\def\csname sga@target@IV.2.2.1\endcsname{ega4.c04.prop.2-2-1}
\expandafter\def\csname sga@target@IV.2.2.10\endcsname{ega4.c04.prop.2-2-10}
\expandafter\def\csname sga@target@IV.2.2.10.1\endcsname{ega4.c04.eq.2-2-10-1}
\expandafter\def\csname sga@target@IV.2.2.11\endcsname{ega4.c04.cor.2-2-11}
\expandafter\def\csname sga@target@IV.2.2.12\endcsname{ega4.c04.cor.2-2-12}
\expandafter\def\csname sga@target@IV.2.2.13\endcsname{ega4.c04.cor.2-2-13}
\expandafter\def\csname sga@target@IV.2.2.14\endcsname{ega4.c04.prop.2-2-14}
\expandafter\def\csname sga@target@IV.2.2.15\endcsname{ega4.c04.prop.2-2-15}
\expandafter\def\csname sga@target@IV.2.2.16\endcsname{ega4.c04.cor.2-2-16}
\expandafter\def\csname sga@target@IV.2.2.17\endcsname{ega4.c04.prop.2-2-17}
\expandafter\def\csname sga@target@IV.2.2.2\endcsname{ega4.c04.prop.2-2-2}
\expandafter\def\csname sga@target@IV.2.2.2.1\endcsname{ega4.c04.eq.2-2-2-1}
\expandafter\def\csname sga@target@IV.2.2.3\endcsname{ega4.c04.cor.2-2-3}
\expandafter\def\csname sga@target@IV.2.2.4\endcsname{ega4.c04.def.2-2-4}
\expandafter\def\csname sga@target@IV.2.2.5\endcsname{ega4.c04.cor.2-2-5}
\expandafter\def\csname sga@target@IV.2.2.6\endcsname{ega4.c04.cor.2-2-6}
\expandafter\def\csname sga@target@IV.2.2.7\endcsname{ega4.c04.prop.2-2-7}
\expandafter\def\csname sga@target@IV.2.2.8\endcsname{ega4.c04.cor.2-2-8}
\expandafter\def\csname sga@target@IV.2.2.8.1\endcsname{ega4.c04.eq.2-2-8-1}
\expandafter\def\csname sga@target@IV.2.2.9\endcsname{ega4.c04.rem.2-2-9}
\expandafter\def\csname sga@target@IV.2.3.1\endcsname{ega4.c04.lem.2-3-1}
\expandafter\def\csname sga@target@IV.2.3.1.1\endcsname{ega4.c04.eq.2-3-1-1}
\expandafter\def\csname sga@target@IV.2.3.10\endcsname{ega4.c04.thm.2-3-10}
\expandafter\def\csname sga@target@IV.2.3.11\endcsname{ega4.c04.cor.2-3-11}
\expandafter\def\csname sga@target@IV.2.3.12\endcsname{ega4.c04.cor.2-3-12}
\expandafter\def\csname sga@target@IV.2.3.13\endcsname{ega4.c04.cor.2-3-13}
\expandafter\def\csname sga@target@IV.2.3.14\endcsname{ega4.c04.cor.2-3-14}
\expandafter\def\csname sga@target@IV.2.3.15\endcsname{ega4.c04.rem.2-3-15}
\expandafter\def\csname sga@target@IV.2.3.2\endcsname{ega4.c04.prop.2-3-2}
\expandafter\def\csname sga@target@IV.2.3.3\endcsname{ega4.c04.statement.2-3-3}
\expandafter\def\csname sga@target@IV.2.3.4\endcsname{ega4.c04.prop.2-3-4}
\expandafter\def\csname sga@target@IV.2.3.5\endcsname{ega4.c04.cor.2-3-5}
\expandafter\def\csname sga@target@IV.2.3.6\endcsname{ega4.c04.prop.2-3-6}
\expandafter\def\csname sga@target@IV.2.3.7\endcsname{ega4.c04.prop.2-3-7}
\expandafter\def\csname sga@target@IV.2.3.8\endcsname{ega4.c04.cor.2-3-8}
\expandafter\def\csname sga@target@IV.2.3.9\endcsname{ega4.c04.rem.2-3-9}
\expandafter\def\csname sga@target@IV.2.4.1\endcsname{ega4.c04.statement.2-4-1}
\expandafter\def\csname sga@target@IV.2.4.10\endcsname{ega4.c04.cor.2-4-10}
\expandafter\def\csname sga@target@IV.2.4.11\endcsname{ega4.c04.rem.2-4-11}
\expandafter\def\csname sga@target@IV.2.4.2\endcsname{ega4.c04.def.2-4-2}
\expandafter\def\csname sga@target@IV.2.4.3\endcsname{ega4.c04.prop.2-4-3}
\expandafter\def\csname sga@target@IV.2.4.4\endcsname{ega4.c04.prop.2-4-4}
\expandafter\def\csname sga@target@IV.2.4.5\endcsname{ega4.c04.prop.2-4-5}
\expandafter\def\csname sga@target@IV.2.4.6\endcsname{ega4.c04.thm.2-4-6}
\expandafter\def\csname sga@target@IV.2.4.7\endcsname{ega4.c04.cor.2-4-7}
\expandafter\def\csname sga@target@IV.2.4.8\endcsname{ega4.c04.rem.2-4-8}
\expandafter\def\csname sga@target@IV.2.4.9\endcsname{ega4.c04.prop.2-4-9}
\expandafter\def\csname sga@target@IV.2.5.1\endcsname{ega4.c04.prop.2-5-1}
\expandafter\def\csname sga@target@IV.2.5.2\endcsname{ega4.c04.prop.2-5-2}
\expandafter\def\csname sga@target@IV.2.5.3\endcsname{ega4.c04.lem.2-5-3}
\expandafter\def\csname sga@target@IV.2.5.4\endcsname{ega4.c04.rem.2-5-4}
\expandafter\def\csname sga@target@IV.2.5.4.1\endcsname{ega4.c04.lem.2-5-4-1}
\expandafter\def\csname sga@target@IV.2.5.5\endcsname{ega4.c04.prop.2-5-5}
\expandafter\def\csname sga@target@IV.2.5.5.1\endcsname{ega4.c04.statement.2-5-5-1}
\expandafter\def\csname sga@target@IV.2.5.5.2\endcsname{ega4.c04.lem.2-5-5-2}
\expandafter\def\csname sga@target@IV.2.5.5.3\endcsname{ega4.c04.statement.2-5-5-3}
\expandafter\def\csname sga@target@IV.2.5.5.4\endcsname{ega4.c04.rem.2-5-5-4}
\expandafter\def\csname sga@target@IV.2.5.6\endcsname{ega4.c04.lem.2-5-6}
\expandafter\def\csname sga@target@IV.2.5.7\endcsname{ega4.c04.cor.2-5-7}
\expandafter\def\csname sga@target@IV.2.5.8\endcsname{ega4.c04.prop.2-5-8}
\expandafter\def\csname sga@target@IV.2.5.8.1\endcsname{ega4.c04.statement.2-5-8-1}
\expandafter\def\csname sga@target@IV.2.5.8.2\endcsname{ega4.c04.statement.2-5-8-2}
\expandafter\def\csname sga@target@IV.2.5.8.3\endcsname{ega4.c04.statement.2-5-8-3}
\expandafter\def\csname sga@target@IV.2.5.8.4\endcsname{ega4.c04.statement.2-5-8-4}
\expandafter\def\csname sga@target@IV.2.5.8.5\endcsname{ega4.c04.statement.2-5-8-5}
\expandafter\def\csname sga@target@IV.2.5.8.6\endcsname{ega4.c04.statement.2-5-8-6}
\expandafter\def\csname sga@target@IV.2.5.8.7\endcsname{ega4.c04.statement.2-5-8-7}
\expandafter\def\csname sga@target@IV.2.5.8.8\endcsname{ega4.c04.statement.2-5-8-8}
\expandafter\def\csname sga@target@IV.2.5.8.9\endcsname{ega4.c04.statement.2-5-8-9}
\expandafter\def\csname sga@target@IV.2.6.1\endcsname{ega4.c04.prop.2-6-1}
\expandafter\def\csname sga@target@IV.2.6.2\endcsname{ega4.c04.prop.2-6-2}
\expandafter\def\csname sga@target@IV.2.6.3\endcsname{ega4.c04.rem.2-6-3}
\expandafter\def\csname sga@target@IV.2.6.4\endcsname{ega4.c04.cor.2-6-4}
\expandafter\def\csname sga@target@IV.2.7.1\endcsname{ega4.c04.prop.2-7-1}
\expandafter\def\csname sga@target@IV.2.7.1.1\endcsname{ega4.c04.lem.2-7-1-1}
\expandafter\def\csname sga@target@IV.2.7.1.viii\endcsname{ega4.c04.item.2-7-1-viii}
\expandafter\def\csname sga@target@IV.2.7.2\endcsname{ega4.c04.cor.2-7-2}
\expandafter\def\csname sga@target@IV.2.7.3\endcsname{ega4.c04.sec.2-7-3}
\expandafter\def\csname sga@target@IV.2.8.1\endcsname{ega4.c04.prop.2-8-1}
\expandafter\def\csname sga@target@IV.2.8.1.1\endcsname{ega4.c04.lem.2-8-1-1}
\expandafter\def\csname sga@target@IV.2.8.2\endcsname{ega4.c04.cor.2-8-2}
\expandafter\def\csname sga@target@IV.2.8.3\endcsname{ega4.c04.statement.2-8-3}
\expandafter\def\csname sga@target@IV.2.8.4\endcsname{ega4.c04.prop.2-8-4}
\expandafter\def\csname sga@target@IV.2.8.4.1\endcsname{ega4.c04.statement.2-8-4-1}
\expandafter\def\csname sga@target@IV.2.8.5\endcsname{ega4.c04.prop.2-8-5}
\expandafter\def\csname sga@target@IV.2.8.6\endcsname{ega4.c04.cor.2-8-6}
\expandafter\def\csname sga@target@IV.2.8.6.1\endcsname{ega4.c04.statement.2-8-6-1}
\expandafter\def\csname sga@target@IV.20\endcsname{ega4.c04.sec.20}
\expandafter\def\csname sga@target@IV.20.0\endcsname{ega4.c04.subsec.20-0}
\expandafter\def\csname sga@target@IV.20.1\endcsname{ega4.c04.subsec.20-1}
\expandafter\def\csname sga@target@IV.20.1.1\endcsname{ega4.c04.statement.20-1-1}
\expandafter\def\csname sga@target@IV.20.1.1.1\endcsname{ega4.c04.eq.20-1-1-1}
\expandafter\def\csname sga@target@IV.20.1.10\endcsname{ega4.c04.statement.20-1-10}
\expandafter\def\csname sga@target@IV.20.1.11\endcsname{ega4.c04.statement.20-1-11}
\expandafter\def\csname sga@target@IV.20.1.11.1\endcsname{ega4.c04.eq.20-1-11-1}
\expandafter\def\csname sga@target@IV.20.1.11.2\endcsname{ega4.c04.eq.20-1-11-2}
\expandafter\def\csname sga@target@IV.20.1.12\endcsname{ega4.c04.prop.20-1-12}
\expandafter\def\csname sga@target@IV.20.1.13\endcsname{ega4.c04.rem.20-1-13}
\expandafter\def\csname sga@target@IV.20.1.2\endcsname{ega4.c04.statement.20-1-2}
\expandafter\def\csname sga@target@IV.20.1.2.1\endcsname{ega4.c04.statement.20-1-2-1}
\expandafter\def\csname sga@target@IV.20.1.2.2\endcsname{ega4.c04.statement.20-1-2-2}
\expandafter\def\csname sga@target@IV.20.1.3\endcsname{ega4.c04.statement.20-1-3}
\expandafter\def\csname sga@target@IV.20.1.3.1\endcsname{ega4.c04.statement.20-1-3-1}
\expandafter\def\csname sga@target@IV.20.1.4\endcsname{ega4.c04.statement.20-1-4}
\expandafter\def\csname sga@target@IV.20.1.4.1\endcsname{ega4.c04.statement.20-1-4-1}
\expandafter\def\csname sga@target@IV.20.1.5\endcsname{ega4.c04.statement.20-1-5}
\expandafter\def\csname sga@target@IV.20.1.5.1\endcsname{ega4.c04.statement.20-1-5-1}
\expandafter\def\csname sga@target@IV.20.1.6\endcsname{ega4.c04.prop.20-1-6}
\expandafter\def\csname sga@target@IV.20.1.7\endcsname{ega4.c04.statement.20-1-7}
\expandafter\def\csname sga@target@IV.20.1.8\endcsname{ega4.c04.statement.20-1-8}
\expandafter\def\csname sga@target@IV.20.1.9\endcsname{ega4.c04.statement.20-1-9}
\expandafter\def\csname sga@target@IV.20.2\endcsname{ega4.c04.subsec.20-2}
\expandafter\def\csname sga@target@IV.20.2.1\endcsname{ega4.c04.statement.20-2-1}
\expandafter\def\csname sga@target@IV.20.2.10\endcsname{ega4.c04.statement.20-2-10}
\expandafter\def\csname sga@target@IV.20.2.10.1\endcsname{ega4.c04.eq.20-2-10-1}
\expandafter\def\csname sga@target@IV.20.2.10.2\endcsname{ega4.c04.eq.20-2-10-2}
\expandafter\def\csname sga@target@IV.20.2.11\endcsname{ega4.c04.prop.20-2-11}
\expandafter\def\csname sga@target@IV.20.2.11.1\endcsname{ega4.c04.eq.20-2-11-1}
\expandafter\def\csname sga@target@IV.20.2.12\endcsname{ega4.c04.cor.20-2-12}
\expandafter\def\csname sga@target@IV.20.2.13\endcsname{ega4.c04.statement.20-2-13}
\expandafter\def\csname sga@target@IV.20.2.14\endcsname{ega4.c04.statement.20-2-14}
\expandafter\def\csname sga@target@IV.20.2.15\endcsname{ega4.c04.prop.20-2-15}
\expandafter\def\csname sga@target@IV.20.2.16\endcsname{ega4.c04.rem.20-2-16}
\expandafter\def\csname sga@target@IV.20.2.2\endcsname{ega4.c04.statement.20-2-2}
\expandafter\def\csname sga@target@IV.20.2.3\endcsname{ega4.c04.statement.20-2-3}
\expandafter\def\csname sga@target@IV.20.2.3.1\endcsname{ega4.c04.statement.20-2-3-1}
\expandafter\def\csname sga@target@IV.20.2.4\endcsname{ega4.c04.prop.20-2-4}
\expandafter\def\csname sga@target@IV.20.2.5\endcsname{ega4.c04.cor.20-2-5}
\expandafter\def\csname sga@target@IV.20.2.6\endcsname{ega4.c04.cor.20-2-6}
\expandafter\def\csname sga@target@IV.20.2.7\endcsname{ega4.c04.rem.20-2-7}
\expandafter\def\csname sga@target@IV.20.2.8\endcsname{ega4.c04.statement.20-2-8}
\expandafter\def\csname sga@target@IV.20.2.9\endcsname{ega4.c04.lem.20-2-9}
\expandafter\def\csname sga@target@IV.20.3\endcsname{ega4.c04.subsec.20-3}
\expandafter\def\csname sga@target@IV.20.3.1\endcsname{ega4.c04.statement.20-3-1}
\expandafter\def\csname sga@target@IV.20.3.10\endcsname{ega4.c04.lem.20-3-10}
\expandafter\def\csname sga@target@IV.20.3.11\endcsname{ega4.c04.prop.20-3-11}
\expandafter\def\csname sga@target@IV.20.3.11.1\endcsname{ega4.c04.eq.20-3-11-1}
\expandafter\def\csname sga@target@IV.20.3.12\endcsname{ega4.c04.cor.20-3-12}
\expandafter\def\csname sga@target@IV.20.3.13\endcsname{ega4.c04.rem.20-3-13}
\expandafter\def\csname sga@target@IV.20.3.2\endcsname{ega4.c04.statement.20-3-2}
\expandafter\def\csname sga@target@IV.20.3.2.1\endcsname{ega4.c04.eq.20-3-2-1}
\expandafter\def\csname sga@target@IV.20.3.3\endcsname{ega4.c04.statement.20-3-3}
\expandafter\def\csname sga@target@IV.20.3.4\endcsname{ega4.c04.statement.20-3-4}
\expandafter\def\csname sga@target@IV.20.3.5\endcsname{ega4.c04.prop.20-3-5}
\expandafter\def\csname sga@target@IV.20.3.6\endcsname{ega4.c04.statement.20-3-6}
\expandafter\def\csname sga@target@IV.20.3.6.1\endcsname{ega4.c04.eq.20-3-6-1}
\expandafter\def\csname sga@target@IV.20.3.7\endcsname{ega4.c04.prop.20-3-7}
\expandafter\def\csname sga@target@IV.20.3.8\endcsname{ega4.c04.prop.20-3-8}
\expandafter\def\csname sga@target@IV.20.3.8.1\endcsname{ega4.c04.eq.20-3-8-1}
\expandafter\def\csname sga@target@IV.20.3.8.2\endcsname{ega4.c04.lem.20-3-8-2}
\expandafter\def\csname sga@target@IV.20.3.8.3\endcsname{ega4.c04.lem.20-3-8-3}
\expandafter\def\csname sga@target@IV.20.3.8.4\endcsname{ega4.c04.lem.20-3-8-4}
\expandafter\def\csname sga@target@IV.20.3.8.5\endcsname{ega4.c04.sec.20-3-8-5}
\expandafter\def\csname sga@target@IV.20.3.8.6\endcsname{ega4.c04.sec.20-3-8-6}
\expandafter\def\csname sga@target@IV.20.3.8.7\endcsname{ega4.c04.sec.20-3-8-7}
\expandafter\def\csname sga@target@IV.20.3.8.8\endcsname{ega4.c04.lem.20-3-8-8}
\expandafter\def\csname sga@target@IV.20.3.9\endcsname{ega4.c04.cor.20-3-9}
\expandafter\def\csname sga@target@IV.20.4\endcsname{ega4.c04.subsec.20-4}
\expandafter\def\csname sga@target@IV.20.4.1\endcsname{ega4.c04.statement.20-4-1}
\expandafter\def\csname sga@target@IV.20.4.10\endcsname{ega4.c04.cor.20-4-10}
\expandafter\def\csname sga@target@IV.20.4.11\endcsname{ega4.c04.prop.20-4-11}
\expandafter\def\csname sga@target@IV.20.4.12\endcsname{ega4.c04.cor.20-4-12}
\expandafter\def\csname sga@target@IV.20.4.2\endcsname{ega4.c04.statement.20-4-2}
\expandafter\def\csname sga@target@IV.20.4.3\endcsname{ega4.c04.statement.20-4-3}
\expandafter\def\csname sga@target@IV.20.4.4\endcsname{ega4.c04.prop.20-4-4}
\expandafter\def\csname sga@target@IV.20.4.5\endcsname{ega4.c04.prop.20-4-5}
\expandafter\def\csname sga@target@IV.20.4.6\endcsname{ega4.c04.prop.20-4-6}
\expandafter\def\csname sga@target@IV.20.4.7\endcsname{ega4.c04.cor.20-4-7}
\expandafter\def\csname sga@target@IV.20.4.8\endcsname{ega4.c04.rem.20-4-8}
\expandafter\def\csname sga@target@IV.20.4.9\endcsname{ega4.c04.cor.20-4-9}
\expandafter\def\csname sga@target@IV.20.5\endcsname{ega4.c04.subsec.20-5}
\expandafter\def\csname sga@target@IV.20.5.1\endcsname{ega4.c04.statement.20-5-1}
\expandafter\def\csname sga@target@IV.20.5.2\endcsname{ega4.c04.statement.20-5-2}
\expandafter\def\csname sga@target@IV.20.5.2.1\endcsname{ega4.c04.statement.20-5-2-1}
\expandafter\def\csname sga@target@IV.20.5.3\endcsname{ega4.c04.prop.20-5-3}
\expandafter\def\csname sga@target@IV.20.5.4\endcsname{ega4.c04.statement.20-5-4}
\expandafter\def\csname sga@target@IV.20.5.5\endcsname{ega4.c04.statement.20-5-5}
\expandafter\def\csname sga@target@IV.20.5.6\endcsname{ega4.c04.statement.20-5-6}
\expandafter\def\csname sga@target@IV.20.5.7\endcsname{ega4.c04.statement.20-5-7}
\expandafter\def\csname sga@target@IV.20.5.8\endcsname{ega4.c04.statement.20-5-8}
\expandafter\def\csname sga@target@IV.20.6\endcsname{ega4.c04.subsec.20-6}
\expandafter\def\csname sga@target@IV.20.6.1\endcsname{ega4.c04.statement.20-6-1}
\expandafter\def\csname sga@target@IV.20.6.1.1\endcsname{ega4.c04.eq.20-6-1-1}
\expandafter\def\csname sga@target@IV.20.6.2\endcsname{ega4.c04.statement.20-6-2}
\expandafter\def\csname sga@target@IV.20.6.2.1\endcsname{ega4.c04.statement.20-6-2-1}
\expandafter\def\csname sga@target@IV.20.6.2.2\endcsname{ega4.c04.eq.20-6-2-2}
\expandafter\def\csname sga@target@IV.20.6.3\endcsname{ega4.c04.prop.20-6-3}
\expandafter\def\csname sga@target@IV.20.6.4\endcsname{ega4.c04.prop.20-6-4}
\expandafter\def\csname sga@target@IV.20.6.5\endcsname{ega4.c04.statement.20-6-5}
\expandafter\def\csname sga@target@IV.20.6.5.1\endcsname{ega4.c04.eq.20-6-5-1}
\expandafter\def\csname sga@target@IV.20.6.6\endcsname{ega4.c04.prop.20-6-6}
\expandafter\def\csname sga@target@IV.20.6.7\endcsname{ega4.c04.statement.20-6-7}
\expandafter\def\csname sga@target@IV.20.6.8\endcsname{ega4.c04.prop.20-6-8}
\expandafter\def\csname sga@target@IV.20.6.8.1\endcsname{ega4.c04.eq.20-6-8-1}
\expandafter\def\csname sga@target@IV.20.6.9\endcsname{ega4.c04.statement.20-6-9}
\expandafter\def\csname sga@target@IV.20.6.9.1\endcsname{ega4.c04.statement.20-6-9-1}
\expandafter\def\csname sga@target@IV.21.1\endcsname{ega4.c04.subsec.21-1.key-iv-21-1}
\expandafter\def\csname sga@target@IV.21.1.1\endcsname{ega4.c04.statement.21-1-1}
\expandafter\def\csname sga@target@IV.21.1.10\endcsname{ega4.c04.prop.21-1-10}
\expandafter\def\csname sga@target@IV.21.1.2\endcsname{ega4.c04.statement.21-1-2}
\expandafter\def\csname sga@target@IV.21.1.3\endcsname{ega4.c04.statement.21-1-3}
\expandafter\def\csname sga@target@IV.21.1.3.1\endcsname{ega4.c04.statement.21-1-3-1}
\expandafter\def\csname sga@target@IV.21.1.3.2\endcsname{ega4.c04.statement.21-1-3-2}
\expandafter\def\csname sga@target@IV.21.1.3.3\endcsname{ega4.c04.statement.21-1-3-3}
\expandafter\def\csname sga@target@IV.21.1.3.4\endcsname{ega4.c04.statement.21-1-3-4}
\expandafter\def\csname sga@target@IV.21.1.4\endcsname{ega4.c04.statement.21-1-4}
\expandafter\def\csname sga@target@IV.21.1.4.1\endcsname{ega4.c04.eq.21-1-4-1}
\expandafter\def\csname sga@target@IV.21.1.4.2\endcsname{ega4.c04.eq.21-1-4-2}
\expandafter\def\csname sga@target@IV.21.1.4.3\endcsname{ega4.c04.eq.21-1-4-3}
\expandafter\def\csname sga@target@IV.21.1.5\endcsname{ega4.c04.statement.21-1-5}
\expandafter\def\csname sga@target@IV.21.1.5.1\endcsname{ega4.c04.statement.21-1-5-1}
\expandafter\def\csname sga@target@IV.21.1.5.2\endcsname{ega4.c04.eq.21-1-5-2}
\expandafter\def\csname sga@target@IV.21.1.6\endcsname{ega4.c04.def.21-1-6}
\expandafter\def\csname sga@target@IV.21.1.6.1\endcsname{ega4.c04.eq.21-1-6-1}
\expandafter\def\csname sga@target@IV.21.1.6.2\endcsname{ega4.c04.eq.21-1-6-2}
\expandafter\def\csname sga@target@IV.21.1.6.3\endcsname{ega4.c04.eq.21-1-6-3}
\expandafter\def\csname sga@target@IV.21.1.6.4\endcsname{ega4.c04.eq.21-1-6-4}
\expandafter\def\csname sga@target@IV.21.1.7\endcsname{ega4.c04.statement.21-1-7}
\expandafter\def\csname sga@target@IV.21.1.7.1\endcsname{ega4.c04.statement.21-1-7-1}
\expandafter\def\csname sga@target@IV.21.1.8\endcsname{ega4.c04.prop.21-1-8}
\expandafter\def\csname sga@target@IV.21.1.9\endcsname{ega4.c04.cor.21-1-9}
\expandafter\def\csname sga@target@IV.21.10.1\endcsname{ega4.c04.statement.21-10-1}
\expandafter\def\csname sga@target@IV.21.10.1.1\endcsname{ega4.c04.eq.21-10-1-1}
\expandafter\def\csname sga@target@IV.21.10.1.2\endcsname{ega4.c04.eq.21-10-1-2}
\expandafter\def\csname sga@target@IV.21.10.10\endcsname{ega4.c04.statement.21-10-10}
\expandafter\def\csname sga@target@IV.21.10.11\endcsname{ega4.c04.statement.21-10-11}
\expandafter\def\csname sga@target@IV.21.10.11.1\endcsname{ega4.c04.eq.21-10-11-1}
\expandafter\def\csname sga@target@IV.21.10.12\endcsname{ega4.c04.rem.21-10-12}
\expandafter\def\csname sga@target@IV.21.10.13\endcsname{ega4.c04.lem.21-10-13}
\expandafter\def\csname sga@target@IV.21.10.13.1\endcsname{ega4.c04.prop.21-10-13-1}
\expandafter\def\csname sga@target@IV.21.10.13bis\endcsname{ega4.c04.prop.21-10-13bis}
\expandafter\def\csname sga@target@IV.21.10.14\endcsname{ega4.c04.statement.21-10-14}
\expandafter\def\csname sga@target@IV.21.10.14.1\endcsname{ega4.c04.eq.21-10-14-1}
\expandafter\def\csname sga@target@IV.21.10.14.2\endcsname{ega4.c04.eq.21-10-14-2}
\expandafter\def\csname sga@target@IV.21.10.15\endcsname{ega4.c04.prop.21-10-15}
\expandafter\def\csname sga@target@IV.21.10.16\endcsname{ega4.c04.prop.21-10-16}
\expandafter\def\csname sga@target@IV.21.10.16.1\endcsname{ega4.c04.prop.21-10-16-1}
\expandafter\def\csname sga@target@IV.21.10.16.2\endcsname{ega4.c04.eq.21-10-16-2}
\expandafter\def\csname sga@target@IV.21.10.17\endcsname{ega4.c04.prop.21-10-17}
\expandafter\def\csname sga@target@IV.21.10.17.1\endcsname{ega4.c04.prop.21-10-17-1}
\expandafter\def\csname sga@target@IV.21.10.17.10\endcsname{ega4.c04.eq.21-10-17-10}
\expandafter\def\csname sga@target@IV.21.10.17.2\endcsname{ega4.c04.eq.21-10-17-2}
\expandafter\def\csname sga@target@IV.21.10.17.3\endcsname{ega4.c04.lem.21-10-17-3}
\expandafter\def\csname sga@target@IV.21.10.17.4\endcsname{ega4.c04.lem.21-10-17-4}
\expandafter\def\csname sga@target@IV.21.10.17.5\endcsname{ega4.c04.eq.21-10-17-5}
\expandafter\def\csname sga@target@IV.21.10.17.6\endcsname{ega4.c04.eq.21-10-17-6}
\expandafter\def\csname sga@target@IV.21.10.17.7\endcsname{ega4.c04.lem.21-10-17-7}
\expandafter\def\csname sga@target@IV.21.10.17.8\endcsname{ega4.c04.lem.21-10-17-8}
\expandafter\def\csname sga@target@IV.21.10.17.9\endcsname{ega4.c04.eq.21-10-17-9}
\expandafter\def\csname sga@target@IV.21.10.18\endcsname{ega4.c04.prop.21-10-18}
\expandafter\def\csname sga@target@IV.21.10.18.1\endcsname{ega4.c04.prop.21-10-18-1}
\expandafter\def\csname sga@target@IV.21.10.19\endcsname{ega4.c04.cor.21-10-19}
\expandafter\def\csname sga@target@IV.21.10.19.1\endcsname{ega4.c04.cor.21-10-19-1}
\expandafter\def\csname sga@target@IV.21.10.2\endcsname{ega4.c04.statement.21-10-2}
\expandafter\def\csname sga@target@IV.21.10.3\endcsname{ega4.c04.statement.21-10-3}
\expandafter\def\csname sga@target@IV.21.10.3.1\endcsname{ega4.c04.statement.21-10-3-1}
\expandafter\def\csname sga@target@IV.21.10.4\endcsname{ega4.c04.prop.21-10-4}
\expandafter\def\csname sga@target@IV.21.10.4.1\endcsname{ega4.c04.prop.21-10-4-1}
\expandafter\def\csname sga@target@IV.21.10.5\endcsname{ega4.c04.statement.21-10-5}
\expandafter\def\csname sga@target@IV.21.10.6\endcsname{ega4.c04.prop.21-10-6}
\expandafter\def\csname sga@target@IV.21.10.7\endcsname{ega4.c04.rem.21-10-7}
\expandafter\def\csname sga@target@IV.21.10.8\endcsname{ega4.c04.prop.21-10-8}
\expandafter\def\csname sga@target@IV.21.10.9\endcsname{ega4.c04.statement.21-10-9}
\expandafter\def\csname sga@target@IV.21.11\endcsname{ega4.c04.subsec.21-11.key-iv-21-11}
\expandafter\def\csname sga@target@IV.21.11.1\endcsname{ega4.c04.thm.21-11-1}
\expandafter\def\csname sga@target@IV.21.11.1.1\endcsname{ega4.c04.eq.21-11-1-1}
\expandafter\def\csname sga@target@IV.21.11.2\endcsname{ega4.c04.lem.21-11-2}
\expandafter\def\csname sga@target@IV.21.12\endcsname{ega4.c04.subsec.21-12.key-iv-21-12}
\expandafter\def\csname sga@target@IV.21.12.1\endcsname{ega4.c04.statement.21-12-1}
\expandafter\def\csname sga@target@IV.21.12.1.1\endcsname{ega4.c04.eq.21-12-1-1}
\expandafter\def\csname sga@target@IV.21.12.10\endcsname{ega4.c04.prop.21-12-10}
\expandafter\def\csname sga@target@IV.21.12.12\endcsname{ega4.c04.thm.21-12-12}
\expandafter\def\csname sga@target@IV.21.12.13\endcsname{ega4.c04.cor.21-12-13}
\expandafter\def\csname sga@target@IV.21.12.14\endcsname{ega4.c04.item.21-12-14}
\expandafter\def\csname sga@target@IV.21.12.15\endcsname{ega4.c04.prop.21-12-15}
\expandafter\def\csname sga@target@IV.21.12.16\endcsname{ega4.c04.cor.21-12-16}
\expandafter\def\csname sga@target@IV.21.12.17\endcsname{ega4.c04.item.21-12-17}
\expandafter\def\csname sga@target@IV.21.12.2\endcsname{ega4.c04.prop.21-12-2}
\expandafter\def\csname sga@target@IV.21.12.2.1\endcsname{ega4.c04.prop.21-12-2-1}
\expandafter\def\csname sga@target@IV.21.12.3\endcsname{ega4.c04.cor.21-12-3}
\expandafter\def\csname sga@target@IV.21.12.4\endcsname{ega4.c04.eq.21-12-4}
\expandafter\def\csname sga@target@IV.21.12.5\endcsname{ega4.c04.statement.21-12-5}
\expandafter\def\csname sga@target@IV.21.12.5.1\endcsname{ega4.c04.eq.21-12-5-1}
\expandafter\def\csname sga@target@IV.21.12.5.2\endcsname{ega4.c04.eq.21-12-5-2}
\expandafter\def\csname sga@target@IV.21.12.5.3\endcsname{ega4.c04.eq.21-12-5-3}
\expandafter\def\csname sga@target@IV.21.12.6\endcsname{ega4.c04.thm.21-12-6}
\expandafter\def\csname sga@target@IV.21.12.7\endcsname{ega4.c04.cor.21-12-7}
\expandafter\def\csname sga@target@IV.21.12.8\endcsname{ega4.c04.statement.21-12-8}
\expandafter\def\csname sga@target@IV.21.12.9\endcsname{ega4.c04.statement.21-12-9}
\expandafter\def\csname sga@target@IV.21.13\endcsname{ega4.c04.subsec.21-13.key-iv-21-13}
\expandafter\def\csname sga@target@IV.21.13.1\endcsname{ega4.c04.def.21-13-1}
\expandafter\def\csname sga@target@IV.21.13.10\endcsname{ega4.c04.prop.21-13-10}
\expandafter\def\csname sga@target@IV.21.13.11\endcsname{ega4.c04.cor.21-13-11}
\expandafter\def\csname sga@target@IV.21.13.12\endcsname{ega4.c04.prop.21-13-12}
\expandafter\def\csname sga@target@IV.21.13.12.1\endcsname{ega4.c04.statement.21-13-12-1}
\expandafter\def\csname sga@target@IV.21.13.13\endcsname{ega4.c04.rem.21-13-13}
\expandafter\def\csname sga@target@IV.21.13.14\endcsname{ega4.c04.prop.21-13-14}
\expandafter\def\csname sga@target@IV.21.13.15\endcsname{ega4.c04.cor.21-13-15}
\expandafter\def\csname sga@target@IV.21.13.16\endcsname{ega4.c04.rem.21-13-16}
\expandafter\def\csname sga@target@IV.21.13.2\endcsname{ega4.c04.lem.21-13-2}
\expandafter\def\csname sga@target@IV.21.13.2.1\endcsname{ega4.c04.eq.21-13-2-1}
\expandafter\def\csname sga@target@IV.21.13.3\endcsname{ega4.c04.cor.21-13-3}
\expandafter\def\csname sga@target@IV.21.13.4\endcsname{ega4.c04.lem.21-13-4}
\expandafter\def\csname sga@target@IV.21.13.5\endcsname{ega4.c04.prop.21-13-5}
\expandafter\def\csname sga@target@IV.21.13.6\endcsname{ega4.c04.cor.21-13-6}
\expandafter\def\csname sga@target@IV.21.13.7\endcsname{ega4.c04.def.21-13-7}
\expandafter\def\csname sga@target@IV.21.13.8\endcsname{ega4.c04.prop.21-13-8}
\expandafter\def\csname sga@target@IV.21.13.9\endcsname{ega4.c04.statement.21-13-9}
\expandafter\def\csname sga@target@IV.21.13.9.1\endcsname{ega4.c04.eq.21-13-9-1}
\expandafter\def\csname sga@target@IV.21.14\endcsname{ega4.c04.subsec.21-14.key-iv-21-14}
\expandafter\def\csname sga@target@IV.21.14.1\endcsname{ega4.c04.thm.21-14-1}
\expandafter\def\csname sga@target@IV.21.14.2\endcsname{ega4.c04.cor.21-14-2}
\expandafter\def\csname sga@target@IV.21.14.3\endcsname{ega4.c04.prop.21-14-3}
\expandafter\def\csname sga@target@IV.21.14.4\endcsname{ega4.c04.statement.21-14-4}
\expandafter\def\csname sga@target@IV.21.14.4bis\endcsname{ega4.c04.statement.21-14-4bis}
\expandafter\def\csname sga@target@IV.21.15\endcsname{ega4.c04.subsec.21-15.key-iv-21-15}
\expandafter\def\csname sga@target@IV.21.15.1\endcsname{ega4.c04.statement.21-15-1}
\expandafter\def\csname sga@target@IV.21.15.2\endcsname{ega4.c04.def.21-15-2}
\expandafter\def\csname sga@target@IV.21.15.3\endcsname{ega4.c04.statement.21-15-3}
\expandafter\def\csname sga@target@IV.21.15.3.1\endcsname{ega4.c04.eq.21-15-3-1}
\expandafter\def\csname sga@target@IV.21.15.3.2\endcsname{ega4.c04.eq.21-15-3-2}
\expandafter\def\csname sga@target@IV.21.15.3bis\endcsname{ega4.c04.statement.21-15-3bis}
\expandafter\def\csname sga@target@IV.21.15.4\endcsname{ega4.c04.prop.21-15-4}
\expandafter\def\csname sga@target@IV.21.15.5\endcsname{ega4.c04.prop.21-15-5}
\expandafter\def\csname sga@target@IV.21.15.6\endcsname{ega4.c04.prop.21-15-6}
\expandafter\def\csname sga@target@IV.21.15.7\endcsname{ega4.c04.statement.21-15-7}
\expandafter\def\csname sga@target@IV.21.15.8\endcsname{ega4.c04.statement.21-15-8}
\expandafter\def\csname sga@target@IV.21.15.9\endcsname{ega4.c04.statement.21-15-9}
\expandafter\def\csname sga@target@IV.21.15.9.1\endcsname{ega4.c04.eq.21-15-9-1}
\expandafter\def\csname sga@target@IV.21.2.1\endcsname{ega4.c04.statement.21-2-1}
\expandafter\def\csname sga@target@IV.21.2.10\endcsname{ega4.c04.statement.21-2-10}
\expandafter\def\csname sga@target@IV.21.2.10.1\endcsname{ega4.c04.eq.21-2-10-1}
\expandafter\def\csname sga@target@IV.21.2.11\endcsname{ega4.c04.prop.21-2-11}
\expandafter\def\csname sga@target@IV.21.2.11.1\endcsname{ega4.c04.prop.21-2-11-1}
\expandafter\def\csname sga@target@IV.21.2.12\endcsname{ega4.c04.statement.21-2-12}
\expandafter\def\csname sga@target@IV.21.2.2\endcsname{ega4.c04.prop.21-2-2}
\expandafter\def\csname sga@target@IV.21.2.3\endcsname{ega4.c04.cor.21-2-3}
\expandafter\def\csname sga@target@IV.21.2.4\endcsname{ega4.c04.statement.21-2-4}
\expandafter\def\csname sga@target@IV.21.2.5\endcsname{ega4.c04.statement.21-2-5}
\expandafter\def\csname sga@target@IV.21.2.5.1\endcsname{ega4.c04.statement.21-2-5-1}
\expandafter\def\csname sga@target@IV.21.2.5.2\endcsname{ega4.c04.eq.21-2-5-2}
\expandafter\def\csname sga@target@IV.21.2.5.3\endcsname{ega4.c04.eq.21-2-5-3}
\expandafter\def\csname sga@target@IV.21.2.5.4\endcsname{ega4.c04.eq.21-2-5-4}
\expandafter\def\csname sga@target@IV.21.2.6\endcsname{ega4.c04.prop.21-2-6}
\expandafter\def\csname sga@target@IV.21.2.6.1\endcsname{ega4.c04.eq.21-2-6-1}
\expandafter\def\csname sga@target@IV.21.2.7\endcsname{ega4.c04.statement.21-2-7}
\expandafter\def\csname sga@target@IV.21.2.7.1\endcsname{ega4.c04.statement.21-2-7-1}
\expandafter\def\csname sga@target@IV.21.2.7.2\endcsname{ega4.c04.eq.21-2-7-2}
\expandafter\def\csname sga@target@IV.21.2.8\endcsname{ega4.c04.statement.21-2-8}
\expandafter\def\csname sga@target@IV.21.2.8.1\endcsname{ega4.c04.statement.21-2-8-1}
\expandafter\def\csname sga@target@IV.21.2.8.2\endcsname{ega4.c04.statement.21-2-8-2}
\expandafter\def\csname sga@target@IV.21.2.8.3\endcsname{ega4.c04.statement.21-2-8-3}
\expandafter\def\csname sga@target@IV.21.2.9\endcsname{ega4.c04.statement.21-2-9}
\expandafter\def\csname sga@target@IV.21.2.9.1\endcsname{ega4.c04.statement.21-2-9-1}
\expandafter\def\csname sga@target@IV.21.2.9.2\endcsname{ega4.c04.eq.21-2-9-2}
\expandafter\def\csname sga@target@IV.21.2.9.3\endcsname{ega4.c04.eq.21-2-9-3}
\expandafter\def\csname sga@target@IV.21.2.9.4\endcsname{ega4.c04.eq.21-2-9-4}
\expandafter\def\csname sga@target@IV.21.3.1\endcsname{ega4.c04.statement.21-3-1}
\expandafter\def\csname sga@target@IV.21.3.2\endcsname{ega4.c04.statement.21-3-2}
\expandafter\def\csname sga@target@IV.21.3.2.1\endcsname{ega4.c04.statement.21-3-2-1}
\expandafter\def\csname sga@target@IV.21.3.2.2\endcsname{ega4.c04.eq.21-3-2-2}
\expandafter\def\csname sga@target@IV.21.3.2.3\endcsname{ega4.c04.eq.21-3-2-3}
\expandafter\def\csname sga@target@IV.21.3.2.4\endcsname{ega4.c04.eq.21-3-2-4}
\expandafter\def\csname sga@target@IV.21.3.3\endcsname{ega4.c04.prop.21-3-3}
\expandafter\def\csname sga@target@IV.21.3.3.1\endcsname{ega4.c04.eq.21-3-3-1}
\expandafter\def\csname sga@target@IV.21.3.4\endcsname{ega4.c04.prop.21-3-4}
\expandafter\def\csname sga@target@IV.21.3.5\endcsname{ega4.c04.cor.21-3-5}
\expandafter\def\csname sga@target@IV.21.3.6\endcsname{ega4.c04.rem.21-3-6}
\expandafter\def\csname sga@target@IV.21.3.6.1\endcsname{ega4.c04.eq.21-3-6-1}
\expandafter\def\csname sga@target@IV.21.4\endcsname{ega4.c04.subsec.21-4.key-iv-21-4}
\expandafter\def\csname sga@target@IV.21.4.1\endcsname{ega4.c04.statement.21-4-1}
\expandafter\def\csname sga@target@IV.21.4.10\endcsname{ega4.c04.cor.21-4-10}
\expandafter\def\csname sga@target@IV.21.4.11\endcsname{ega4.c04.cor.21-4-11}
\expandafter\def\csname sga@target@IV.21.4.12\endcsname{ega4.c04.rem.21-4-12}
\expandafter\def\csname sga@target@IV.21.4.13\endcsname{ega4.c04.cor.21-4-13}
\expandafter\def\csname sga@target@IV.21.4.2\endcsname{ega4.c04.def.21-4-2}
\expandafter\def\csname sga@target@IV.21.4.2.1\endcsname{ega4.c04.eq.21-4-2-1}
\expandafter\def\csname sga@target@IV.21.4.3\endcsname{ega4.c04.statement.21-4-3}
\expandafter\def\csname sga@target@IV.21.4.4\endcsname{ega4.c04.statement.21-4-4}
\expandafter\def\csname sga@target@IV.21.4.5\endcsname{ega4.c04.prop.21-4-5}
\expandafter\def\csname sga@target@IV.21.4.6\endcsname{ega4.c04.cor.21-4-6}
\expandafter\def\csname sga@target@IV.21.4.6.1\endcsname{ega4.c04.eq.21-4-6-1}
\expandafter\def\csname sga@target@IV.21.4.7\endcsname{ega4.c04.statement.21-4-7}
\expandafter\def\csname sga@target@IV.21.4.8\endcsname{ega4.c04.prop.21-4-8}
\expandafter\def\csname sga@target@IV.21.4.9\endcsname{ega4.c04.prop.21-4-9}
\expandafter\def\csname sga@target@IV.21.5\endcsname{ega4.c04.subsec.21-5.key-iv-21-5}
\expandafter\def\csname sga@target@IV.21.5.1\endcsname{ega4.c04.statement.21-5-1}
\expandafter\def\csname sga@target@IV.21.5.2\endcsname{ega4.c04.lem.21-5-2}
\expandafter\def\csname sga@target@IV.21.5.3\endcsname{ega4.c04.statement.21-5-3}
\expandafter\def\csname sga@target@IV.21.5.3.1\endcsname{ega4.c04.eq.21-5-3-1}
\expandafter\def\csname sga@target@IV.21.5.3.2\endcsname{ega4.c04.eq.21-5-3-2}
\expandafter\def\csname sga@target@IV.21.5.4\endcsname{ega4.c04.statement.21-5-4}
\expandafter\def\csname sga@target@IV.21.5.5\endcsname{ega4.c04.def.21-5-5}
\expandafter\def\csname sga@target@IV.21.5.5.1\endcsname{ega4.c04.eq.21-5-5-1}
\expandafter\def\csname sga@target@IV.21.5.5.2\endcsname{ega4.c04.eq.21-5-5-2}
\expandafter\def\csname sga@target@IV.21.5.5.3\endcsname{ega4.c04.eq.21-5-5-3}
\expandafter\def\csname sga@target@IV.21.5.6\endcsname{ega4.c04.prop.21-5-6}
\expandafter\def\csname sga@target@IV.21.5.6.1\endcsname{ega4.c04.prop.21-5-6-1}
\expandafter\def\csname sga@target@IV.21.5.7\endcsname{ega4.c04.prop.21-5-7}
\expandafter\def\csname sga@target@IV.21.5.7.1\endcsname{ega4.c04.prop.21-5-7-1}
\expandafter\def\csname sga@target@IV.21.5.7.2\endcsname{ega4.c04.lem.21-5-7-2}
\expandafter\def\csname sga@target@IV.21.5.7.3\endcsname{ega4.c04.lem.21-5-7-3}
\expandafter\def\csname sga@target@IV.21.5.7.4\endcsname{ega4.c04.eq.21-5-7-4}
\expandafter\def\csname sga@target@IV.21.5.8\endcsname{ega4.c04.prop.21-5-8}
\expandafter\def\csname sga@target@IV.21.5.8.1\endcsname{ega4.c04.eq.21-5-8-1}
\expandafter\def\csname sga@target@IV.21.6\endcsname{ega4.c04.subsec.21-6.key-iv-21-6}
\expandafter\def\csname sga@target@IV.21.6.1\endcsname{ega4.c04.statement.21-6-1}
\expandafter\def\csname sga@target@IV.21.6.10\endcsname{ega4.c04.cor.21-6-10}
\expandafter\def\csname sga@target@IV.21.6.10.1\endcsname{ega4.c04.eq.21-6-10-1}
\expandafter\def\csname sga@target@IV.21.6.11\endcsname{ega4.c04.cor.21-6-11}
\expandafter\def\csname sga@target@IV.21.6.12\endcsname{ega4.c04.prop.21-6-12}
\expandafter\def\csname sga@target@IV.21.6.12.1\endcsname{ega4.c04.eq.21-6-12-1}
\expandafter\def\csname sga@target@IV.21.6.12.2\endcsname{ega4.c04.eq.21-6-12-2}
\expandafter\def\csname sga@target@IV.21.6.13\endcsname{ega4.c04.cor.21-6-13}
\expandafter\def\csname sga@target@IV.21.6.14\endcsname{ega4.c04.cor.21-6-14}
\expandafter\def\csname sga@target@IV.21.6.2\endcsname{ega4.c04.statement.21-6-2}
\expandafter\def\csname sga@target@IV.21.6.3\endcsname{ega4.c04.statement.21-6-3}
\expandafter\def\csname sga@target@IV.21.6.4\endcsname{ega4.c04.statement.21-6-4}
\expandafter\def\csname sga@target@IV.21.6.4.1\endcsname{ega4.c04.statement.21-6-4-1}
\expandafter\def\csname sga@target@IV.21.6.4.2\endcsname{ega4.c04.eq.21-6-4-2}
\expandafter\def\csname sga@target@IV.21.6.5\endcsname{ega4.c04.statement.21-6-5}
\expandafter\def\csname sga@target@IV.21.6.5.1\endcsname{ega4.c04.eq.21-6-5-1}
\expandafter\def\csname sga@target@IV.21.6.6\endcsname{ega4.c04.prop.21-6-6}
\expandafter\def\csname sga@target@IV.21.6.6.2\endcsname{ega4.c04.statement.21-6-6-2}
\expandafter\def\csname sga@target@IV.21.6.7\endcsname{ega4.c04.statement.21-6-7}
\expandafter\def\csname sga@target@IV.21.6.7.1\endcsname{ega4.c04.eq.21-6-7-1}
\expandafter\def\csname sga@target@IV.21.6.7.2\endcsname{ega4.c04.eq.21-6-7-2}
\expandafter\def\csname sga@target@IV.21.6.7.3\endcsname{ega4.c04.eq.21-6-7-3}
\expandafter\def\csname sga@target@IV.21.6.7.4\endcsname{ega4.c04.eq.21-6-7-4}
\expandafter\def\csname sga@target@IV.21.6.7.5\endcsname{ega4.c04.eq.21-6-7-5}
\expandafter\def\csname sga@target@IV.21.6.8\endcsname{ega4.c04.statement.21-6-8}
\expandafter\def\csname sga@target@IV.21.6.9\endcsname{ega4.c04.thm.21-6-9}
\expandafter\def\csname sga@target@IV.21.6.9.1\endcsname{ega4.c04.eq.21-6-9-1}
\expandafter\def\csname sga@target@IV.21.6.9.2\endcsname{ega4.c04.statement.21-6-9-2}
\expandafter\def\csname sga@target@IV.21.7\endcsname{ega4.c04.subsec.21-7.key-iv-21-7}
\expandafter\def\csname sga@target@IV.21.7.1\endcsname{ega4.c04.statement.21-7-1}
\expandafter\def\csname sga@target@IV.21.7.2\endcsname{ega4.c04.prop.21-7-2}
\expandafter\def\csname sga@target@IV.21.7.2.1\endcsname{ega4.c04.eq.21-7-2-1}
\expandafter\def\csname sga@target@IV.21.7.3\endcsname{ega4.c04.cor.21-7-3}
\expandafter\def\csname sga@target@IV.21.7.3.1\endcsname{ega4.c04.statement.21-7-3-1}
\expandafter\def\csname sga@target@IV.21.7.4\endcsname{ega4.c04.prop.21-7-4}
\expandafter\def\csname sga@target@IV.21.7.5\endcsname{ega4.c04.rem.21-7-5}
\expandafter\def\csname sga@target@IV.21.7.6\endcsname{ega4.c04.cor.21-7-6}
\expandafter\def\csname sga@target@IV.21.7.7\endcsname{ega4.c04.rem.21-7-7}
\expandafter\def\csname sga@target@IV.21.8.1\endcsname{ega4.c04.lem.21-8-1}
\expandafter\def\csname sga@target@IV.21.8.2\endcsname{ega4.c04.cor.21-8-2}
\expandafter\def\csname sga@target@IV.21.8.3\endcsname{ega4.c04.prop.21-8-3}
\expandafter\def\csname sga@target@IV.21.8.3.1\endcsname{ega4.c04.eq.21-8-3-1}
\expandafter\def\csname sga@target@IV.21.8.4\endcsname{ega4.c04.cor.21-8-4}
\expandafter\def\csname sga@target@IV.21.8.5\endcsname{ega4.c04.prop.21-8-5}
\expandafter\def\csname sga@target@IV.21.8.5.1\endcsname{ega4.c04.eq.21-8-5-1}
\expandafter\def\csname sga@target@IV.21.8.5.2\endcsname{ega4.c04.eq.21-8-5-2}
\expandafter\def\csname sga@target@IV.21.8.5.3\endcsname{ega4.c04.eq.21-8-5-3}
\expandafter\def\csname sga@target@IV.21.8.6\endcsname{ega4.c04.rem.21-8-6}
\expandafter\def\csname sga@target@IV.21.9\endcsname{ega4.c04.subsec.21-9.key-iv-21-9}
\expandafter\def\csname sga@target@IV.21.9.1\endcsname{ega4.c04.statement.21-9-1}
\expandafter\def\csname sga@target@IV.21.9.1.1\endcsname{ega4.c04.eq.21-9-1-1}
\expandafter\def\csname sga@target@IV.21.9.1.2\endcsname{ega4.c04.eq.21-9-1-2}
\expandafter\def\csname sga@target@IV.21.9.10\endcsname{ega4.c04.statement.21-9-10}
\expandafter\def\csname sga@target@IV.21.9.11\endcsname{ega4.c04.prop.21-9-11}
\expandafter\def\csname sga@target@IV.21.9.12\endcsname{ega4.c04.cor.21-9-12}
\expandafter\def\csname sga@target@IV.21.9.13\endcsname{ega4.c04.rem.21-9-13}
\expandafter\def\csname sga@target@IV.21.9.2\endcsname{ega4.c04.lem.21-9-2}
\expandafter\def\csname sga@target@IV.21.9.3\endcsname{ega4.c04.rem.21-9-3}
\expandafter\def\csname sga@target@IV.21.9.4\endcsname{ega4.c04.prop.21-9-4}
\expandafter\def\csname sga@target@IV.21.9.4.1\endcsname{ega4.c04.prop.21-9-4-1}
\expandafter\def\csname sga@target@IV.21.9.5\endcsname{ega4.c04.cor.21-9-5}
\expandafter\def\csname sga@target@IV.21.9.6\endcsname{ega4.c04.cor.21-9-6}
\expandafter\def\csname sga@target@IV.21.9.7\endcsname{ega4.c04.cor.21-9-7}
\expandafter\def\csname sga@target@IV.21.9.8\endcsname{ega4.c04.statement.21-9-8}
\expandafter\def\csname sga@target@IV.21.9.9\endcsname{ega4.c04.prop.21-9-9}
\expandafter\def\csname sga@target@IV.3.1.1\endcsname{ega4.c04.def.3-1-1}
\expandafter\def\csname sga@target@IV.3.1.1.1\endcsname{ega4.c04.eq.3-1-1-1}
\expandafter\def\csname sga@target@IV.3.1.10\endcsname{ega4.c04.prop.3-1-10}
\expandafter\def\csname sga@target@IV.3.1.10.1\endcsname{ega4.c04.lem.3-1-10-1}
\expandafter\def\csname sga@target@IV.3.1.10.2\endcsname{ega4.c04.rem.3-1-10-2}
\expandafter\def\csname sga@target@IV.3.1.10.3\endcsname{ega4.c04.eq.3-1-10-3}
\expandafter\def\csname sga@target@IV.3.1.10.4\endcsname{ega4.c04.eq.3-1-10-4}
\expandafter\def\csname sga@target@IV.3.1.10.5\endcsname{ega4.c04.rem.3-1-10-5}
\expandafter\def\csname sga@target@IV.3.1.11\endcsname{ega4.c04.cor.3-1-11}
\expandafter\def\csname sga@target@IV.3.1.12\endcsname{ega4.c04.rem.3-1-12}
\expandafter\def\csname sga@target@IV.3.1.13\endcsname{ega4.c04.prop.3-1-13}
\expandafter\def\csname sga@target@IV.3.1.2\endcsname{ega4.c04.prop.3-1-2}
\expandafter\def\csname sga@target@IV.3.1.3\endcsname{ega4.c04.prop.3-1-3}
\expandafter\def\csname sga@target@IV.3.1.4\endcsname{ega4.c04.cor.3-1-4}
\expandafter\def\csname sga@target@IV.3.1.5\endcsname{ega4.c04.cor.3-1-5}
\expandafter\def\csname sga@target@IV.3.1.6\endcsname{ega4.c04.prop.3-1-6}
\expandafter\def\csname sga@target@IV.3.1.7\endcsname{ega4.c04.prop.3-1-7}
\expandafter\def\csname sga@target@IV.3.1.8\endcsname{ega4.c04.prop.3-1-8}
\expandafter\def\csname sga@target@IV.3.1.9\endcsname{ega4.c04.cor.3-1-9}
\expandafter\def\csname sga@target@IV.3.2.1\endcsname{ega4.c04.prop.3-2-1}
\expandafter\def\csname sga@target@IV.3.2.2\endcsname{ega4.c04.def.3-2-2}
\expandafter\def\csname sga@target@IV.3.2.3\endcsname{ega4.c04.prop.3-2-3}
\expandafter\def\csname sga@target@IV.3.2.4\endcsname{ega4.c04.def.3-2-4}
\expandafter\def\csname sga@target@IV.3.2.5\endcsname{ega4.c04.def.3-2-5}
\expandafter\def\csname sga@target@IV.3.2.6\endcsname{ega4.c04.prop.3-2-6}
\expandafter\def\csname sga@target@IV.3.2.7\endcsname{ega4.c04.cor.3-2-7}
\expandafter\def\csname sga@target@IV.3.2.8\endcsname{ega4.c04.cor.3-2-8}
\expandafter\def\csname sga@target@IV.3.2.9\endcsname{ega4.c04.rem.3-2-9}
\expandafter\def\csname sga@target@IV.3.3.1\endcsname{ega4.c04.prop.3-3-1}
\expandafter\def\csname sga@target@IV.3.3.1.1\endcsname{ega4.c04.prop.3-3-1-1}
\expandafter\def\csname sga@target@IV.3.3.2\endcsname{ega4.c04.cor.3-3-2}
\expandafter\def\csname sga@target@IV.3.3.3\endcsname{ega4.c04.cor.3-3-3}
\expandafter\def\csname sga@target@IV.3.3.4\endcsname{ega4.c04.prop.3-3-4}
\expandafter\def\csname sga@target@IV.3.3.5\endcsname{ega4.c04.cor.3-3-5}
\expandafter\def\csname sga@target@IV.3.3.6\endcsname{ega4.c04.prop.3-3-6}
\expandafter\def\csname sga@target@IV.3.3.6.1\endcsname{ega4.c04.eq.3-3-6-1}
\expandafter\def\csname sga@target@IV.3.3.6.2\endcsname{ega4.c04.eq.3-3-6-2}
\expandafter\def\csname sga@target@IV.3.3.6.3\endcsname{ega4.c04.eq.3-3-6-3}
\expandafter\def\csname sga@target@IV.3.3.6.4\endcsname{ega4.c04.eq.3-3-6-4}
\expandafter\def\csname sga@target@IV.3.3.7\endcsname{ega4.c04.cor.3-3-7}
\expandafter\def\csname sga@target@IV.3.3.7.1\endcsname{ega4.c04.cor.3-3-7-1}
\expandafter\def\csname sga@target@IV.3.3.8\endcsname{ega4.c04.sec.3-3-8}
\expandafter\def\csname sga@target@IV.3.3.9\endcsname{ega4.c04.cor.3-3-9}
\expandafter\def\csname sga@target@IV.3.3.9.1\endcsname{ega4.c04.cor.3-3-9-1}
\expandafter\def\csname sga@target@IV.3.4\endcsname{ega4.c04.subsec.3-4.key-iv-3-4}
\expandafter\def\csname sga@target@IV.3.4.1\endcsname{ega4.c04.prop.3-4-1}
\expandafter\def\csname sga@target@IV.3.4.1.1\endcsname{ega4.c04.eq.3-4-1-1}
\expandafter\def\csname sga@target@IV.3.4.1.2\endcsname{ega4.c04.lem.3-4-1-2}
\expandafter\def\csname sga@target@IV.3.4.1.3\endcsname{ega4.c04.lem.3-4-1-3}
\expandafter\def\csname sga@target@IV.3.4.1.4\endcsname{ega4.c04.lem.3-4-1-4}
\expandafter\def\csname sga@target@IV.3.4.1.5\endcsname{ega4.c04.eq.3-4-1-5}
\expandafter\def\csname sga@target@IV.3.4.2\endcsname{ega4.c04.cor.3-4-2}
\expandafter\def\csname sga@target@IV.3.4.3\endcsname{ega4.c04.prop.3-4-3}
\expandafter\def\csname sga@target@IV.3.4.4\endcsname{ega4.c04.prop.3-4-4}
\expandafter\def\csname sga@target@IV.3.4.5\endcsname{ega4.c04.prop.3-4-5}
\expandafter\def\csname sga@target@IV.3.4.6\endcsname{ega4.c04.prop.3-4-6}
\expandafter\def\csname sga@target@IV.3.4.6.1\endcsname{ega4.c04.lem.3-4-6-1}
\expandafter\def\csname sga@target@IV.3.4.7\endcsname{ega4.c04.cor.3-4-7}
\expandafter\def\csname sga@target@IV.3.4.8\endcsname{ega4.c04.rem.3-4-8}
\expandafter\def\csname sga@target@IV.3.4.9\endcsname{ega4.c04.prop.3-4-9}
\expandafter\def\csname sga@target@IV.4.1.1\endcsname{ega4.c04.def.4-1-1}
\expandafter\def\csname sga@target@IV.4.1.1.1\endcsname{ega4.c04.eq.4-1-1-1}
\expandafter\def\csname sga@target@IV.4.1.1.2\endcsname{ega4.c04.eq.4-1-1-2}
\expandafter\def\csname sga@target@IV.4.1.1.3\endcsname{ega4.c04.eq.4-1-1-3}
\expandafter\def\csname sga@target@IV.4.1.2\endcsname{ega4.c04.thm.4-1-2}
\expandafter\def\csname sga@target@IV.4.1.2.1\endcsname{ega4.c04.cor.4-1-2-1}
\expandafter\def\csname sga@target@IV.4.1.3\endcsname{ega4.c04.rem.4-1-3}
\expandafter\def\csname sga@target@IV.4.1.4\endcsname{ega4.c04.cor.4-1-4}
\expandafter\def\csname sga@target@IV.4.1.5\endcsname{ega4.c04.cor.4-1-5}
\expandafter\def\csname sga@target@IV.4.2.1\endcsname{ega4.c04.prop.4-2-1}
\expandafter\def\csname sga@target@IV.4.2.1.1\endcsname{ega4.c04.eq.4-2-1-1}
\expandafter\def\csname sga@target@IV.4.2.1.2\endcsname{ega4.c04.eq.4-2-1-2}
\expandafter\def\csname sga@target@IV.4.2.1.3\endcsname{ega4.c04.eq.4-2-1-3}
\expandafter\def\csname sga@target@IV.4.2.1.4\endcsname{ega4.c04.rem.4-2-1-4}
\expandafter\def\csname sga@target@IV.4.2.1.5\endcsname{ega4.c04.eq.4-2-1-5}
\expandafter\def\csname sga@target@IV.4.2.10\endcsname{ega4.c04.cor.4-2-10}
\expandafter\def\csname sga@target@IV.4.2.2\endcsname{ega4.c04.cor.4-2-2}
\expandafter\def\csname sga@target@IV.4.2.3\endcsname{ega4.c04.cor.4-2-3}
\expandafter\def\csname sga@target@IV.4.2.4\endcsname{ega4.c04.prop.4-2-4}
\expandafter\def\csname sga@target@IV.4.2.4.i\endcsname{ega4.c04.sec.4-2-4-i}
\expandafter\def\csname sga@target@IV.4.2.4.ii\endcsname{ega4.c04.sec.4-2-4-ii}
\expandafter\def\csname sga@target@IV.4.2.4.iii\endcsname{ega4.c04.sec.4-2-4-iii}
\expandafter\def\csname sga@target@IV.4.2.5\endcsname{ega4.c04.prop.4-2-5}
\expandafter\def\csname sga@target@IV.4.2.6\endcsname{ega4.c04.cor.4-2-6}
\expandafter\def\csname sga@target@IV.4.2.7\endcsname{ega4.c04.prop.4-2-7}
\expandafter\def\csname sga@target@IV.4.2.7.i\endcsname{ega4.c04.sec.4-2-7-i}
\expandafter\def\csname sga@target@IV.4.2.7.ii\endcsname{ega4.c04.sec.4-2-7-ii}
\expandafter\def\csname sga@target@IV.4.2.7.iii\endcsname{ega4.c04.sec.4-2-7-iii}
\expandafter\def\csname sga@target@IV.4.2.7.iv\endcsname{ega4.c04.sec.4-2-7-iv}
\expandafter\def\csname sga@target@IV.4.2.8\endcsname{ega4.c04.cor.4-2-8}
\expandafter\def\csname sga@target@IV.4.2.9\endcsname{ega4.c04.prop.4-2-9}
\expandafter\def\csname sga@target@IV.4.3.1\endcsname{ega4.c04.statement.4-3-1}
\expandafter\def\csname sga@target@IV.4.3.2\endcsname{ega4.c04.prop.4-3-2}
\expandafter\def\csname sga@target@IV.4.3.3\endcsname{ega4.c04.cor.4-3-3}
\expandafter\def\csname sga@target@IV.4.3.4\endcsname{ega4.c04.cor.4-3-4}
\expandafter\def\csname sga@target@IV.4.3.5\endcsname{ega4.c04.prop.4-3-5}
\expandafter\def\csname sga@target@IV.4.3.6\endcsname{ega4.c04.cor.4-3-6}
\expandafter\def\csname sga@target@IV.4.3.7\endcsname{ega4.c04.cor.4-3-7}
\expandafter\def\csname sga@target@IV.4.3.8\endcsname{ega4.c04.cor.4-3-8}
\expandafter\def\csname sga@target@IV.4.3.9\endcsname{ega4.c04.prop.4-3-9}
\expandafter\def\csname sga@target@IV.4.4.1\endcsname{ega4.c04.statement.4-4-1}
\expandafter\def\csname sga@target@IV.4.4.2\endcsname{ega4.c04.lem.4-4-2}
\expandafter\def\csname sga@target@IV.4.4.2.i\endcsname{ega4.c04.sec.4-4-2-i}
\expandafter\def\csname sga@target@IV.4.4.2.ii\endcsname{ega4.c04.sec.4-4-2-ii}
\expandafter\def\csname sga@target@IV.4.4.3\endcsname{ega4.c04.rem.4-4-3}
\expandafter\def\csname sga@target@IV.4.4.4\endcsname{ega4.c04.thm.4-4-4}
\expandafter\def\csname sga@target@IV.4.4.5\endcsname{ega4.c04.cor.4-4-5}
\expandafter\def\csname sga@target@IV.4.4.6\endcsname{ega4.c04.cor.4-4-6}
\expandafter\def\csname sga@target@IV.4.5\endcsname{ega4.c04.subsec.4-5.key-iv-4-5}
\expandafter\def\csname sga@target@IV.4.5.1\endcsname{ega4.c04.prop.4-5-1}
\expandafter\def\csname sga@target@IV.4.5.10\endcsname{ega4.c04.cor.4-5-10}
\expandafter\def\csname sga@target@IV.4.5.11\endcsname{ega4.c04.cor.4-5-11}
\expandafter\def\csname sga@target@IV.4.5.12\endcsname{ega4.c04.sec.4-5-12}
\expandafter\def\csname sga@target@IV.4.5.12.i\endcsname{ega4.c04.sec.4-5-12-i}
\expandafter\def\csname sga@target@IV.4.5.12.ii\endcsname{ega4.c04.sec.4-5-12-ii}
\expandafter\def\csname sga@target@IV.4.5.13\endcsname{ega4.c04.prop.4-5-13}
\expandafter\def\csname sga@target@IV.4.5.13.1\endcsname{ega4.c04.cor.4-5-13-1}
\expandafter\def\csname sga@target@IV.4.5.13.1.i\endcsname{ega4.c04.sec.4-5-13-1-i}
\expandafter\def\csname sga@target@IV.4.5.13.1.ii\endcsname{ega4.c04.sec.4-5-13-1-ii}
\expandafter\def\csname sga@target@IV.4.5.14\endcsname{ega4.c04.cor.4-5-14}
\expandafter\def\csname sga@target@IV.4.5.15\endcsname{ega4.c04.prop.4-5-15}
\expandafter\def\csname sga@target@IV.4.5.15.i\endcsname{ega4.c04.sec.4-5-15-i}
\expandafter\def\csname sga@target@IV.4.5.15.ii\endcsname{ega4.c04.sec.4-5-15-ii}
\expandafter\def\csname sga@target@IV.4.5.16\endcsname{ega4.c04.cor.4-5-16}
\expandafter\def\csname sga@target@IV.4.5.16.i\endcsname{ega4.c04.sec.4-5-16-i}
\expandafter\def\csname sga@target@IV.4.5.16.ii\endcsname{ega4.c04.sec.4-5-16-ii}
\expandafter\def\csname sga@target@IV.4.5.17\endcsname{ega4.c04.cor.4-5-17}
\expandafter\def\csname sga@target@IV.4.5.18\endcsname{ega4.c04.rem.4-5-18}
\expandafter\def\csname sga@target@IV.4.5.19\endcsname{ega4.c04.prop.4-5-19}
\expandafter\def\csname sga@target@IV.4.5.19.1\endcsname{ega4.c04.lem.4-5-19-1}
\expandafter\def\csname sga@target@IV.4.5.19.2\endcsname{ega4.c04.eq.4-5-19-2}
\expandafter\def\csname sga@target@IV.4.5.19.3\endcsname{ega4.c04.cor.4-5-19-3}
\expandafter\def\csname sga@target@IV.4.5.2\endcsname{ega4.c04.def.4-5-2}
\expandafter\def\csname sga@target@IV.4.5.2.a\endcsname{ega4.c04.sec.4-5-2-a}
\expandafter\def\csname sga@target@IV.4.5.2.b\endcsname{ega4.c04.sec.4-5-2-b}
\expandafter\def\csname sga@target@IV.4.5.2.c\endcsname{ega4.c04.sec.4-5-2-c}
\expandafter\def\csname sga@target@IV.4.5.20\endcsname{ega4.c04.rem.4-5-20}
\expandafter\def\csname sga@target@IV.4.5.21\endcsname{ega4.c04.prop.4-5-21}
\expandafter\def\csname sga@target@IV.4.5.3\endcsname{ega4.c04.statement.4-5-3}
\expandafter\def\csname sga@target@IV.4.5.4\endcsname{ega4.c04.prop.4-5-4}
\expandafter\def\csname sga@target@IV.4.5.5\endcsname{ega4.c04.def.4-5-5}
\expandafter\def\csname sga@target@IV.4.5.6\endcsname{ega4.c04.prop.4-5-6}
\expandafter\def\csname sga@target@IV.4.5.6.i\endcsname{ega4.c04.sec.4-5-6-i}
\expandafter\def\csname sga@target@IV.4.5.6.ii\endcsname{ega4.c04.sec.4-5-6-ii}
\expandafter\def\csname sga@target@IV.4.5.7\endcsname{ega4.c04.prop.4-5-7}
\expandafter\def\csname sga@target@IV.4.5.8\endcsname{ega4.c04.cor.4-5-8}
\expandafter\def\csname sga@target@IV.4.5.8.i\endcsname{ega4.c04.sec.4-5-8-i}
\expandafter\def\csname sga@target@IV.4.5.8.ii\endcsname{ega4.c04.sec.4-5-8-ii}
\expandafter\def\csname sga@target@IV.4.5.9\endcsname{ega4.c04.prop.4-5-9}
\expandafter\def\csname sga@target@IV.4.5.9.a\endcsname{ega4.c04.sec.4-5-9-a}
\expandafter\def\csname sga@target@IV.4.5.9.b\endcsname{ega4.c04.sec.4-5-9-b}
\expandafter\def\csname sga@target@IV.4.5.9.c\endcsname{ega4.c04.sec.4-5-9-c}
\expandafter\def\csname sga@target@IV.4.6\endcsname{ega4.c04.subsec.4-6.key-iv-4-6}
\expandafter\def\csname sga@target@IV.4.6.1\endcsname{ega4.c04.prop.4-6-1}
\expandafter\def\csname sga@target@IV.4.6.1.a\endcsname{ega4.c04.sec.4-6-1-a}
\expandafter\def\csname sga@target@IV.4.6.1.b\endcsname{ega4.c04.sec.4-6-1-b}
\expandafter\def\csname sga@target@IV.4.6.1.c\endcsname{ega4.c04.sec.4-6-1-c}
\expandafter\def\csname sga@target@IV.4.6.1.d\endcsname{ega4.c04.sec.4-6-1-d}
\expandafter\def\csname sga@target@IV.4.6.1.e\endcsname{ega4.c04.sec.4-6-1-e}
\expandafter\def\csname sga@target@IV.4.6.10\endcsname{ega4.c04.prop.4-6-10}
\expandafter\def\csname sga@target@IV.4.6.11\endcsname{ega4.c04.cor.4-6-11}
\expandafter\def\csname sga@target@IV.4.6.12\endcsname{ega4.c04.prop.4-6-12}
\expandafter\def\csname sga@target@IV.4.6.12.a\endcsname{ega4.c04.sec.4-6-12-a}
\expandafter\def\csname sga@target@IV.4.6.12.b\endcsname{ega4.c04.sec.4-6-12-b}
\expandafter\def\csname sga@target@IV.4.6.12.c\endcsname{ega4.c04.sec.4-6-12-c}
\expandafter\def\csname sga@target@IV.4.6.12.d\endcsname{ega4.c04.sec.4-6-12-d}
\expandafter\def\csname sga@target@IV.4.6.12.e\endcsname{ega4.c04.sec.4-6-12-e}
\expandafter\def\csname sga@target@IV.4.6.13\endcsname{ega4.c04.cor.4-6-13}
\expandafter\def\csname sga@target@IV.4.6.13.f\endcsname{ega4.c04.sec.4-6-13-f}
\expandafter\def\csname sga@target@IV.4.6.14\endcsname{ega4.c04.statement.4-6-14}
\expandafter\def\csname sga@target@IV.4.6.15\endcsname{ega4.c04.prop.4-6-15}
\expandafter\def\csname sga@target@IV.4.6.15.i\endcsname{ega4.c04.sec.4-6-15-i}
\expandafter\def\csname sga@target@IV.4.6.15.ii\endcsname{ega4.c04.sec.4-6-15-ii}
\expandafter\def\csname sga@target@IV.4.6.16\endcsname{ega4.c04.prop.4-6-16}
\expandafter\def\csname sga@target@IV.4.6.16.a\endcsname{ega4.c04.sec.4-6-16-a}
\expandafter\def\csname sga@target@IV.4.6.16.b\endcsname{ega4.c04.sec.4-6-16-b}
\expandafter\def\csname sga@target@IV.4.6.16.c\endcsname{ega4.c04.sec.4-6-16-c}
\expandafter\def\csname sga@target@IV.4.6.16.d\endcsname{ega4.c04.sec.4-6-16-d}
\expandafter\def\csname sga@target@IV.4.6.17\endcsname{ega4.c04.def.4-6-17}
\expandafter\def\csname sga@target@IV.4.6.18\endcsname{ega4.c04.prop.4-6-18}
\expandafter\def\csname sga@target@IV.4.6.18.a\endcsname{ega4.c04.sec.4-6-18-a}
\expandafter\def\csname sga@target@IV.4.6.18.b\endcsname{ega4.c04.sec.4-6-18-b}
\expandafter\def\csname sga@target@IV.4.6.18.c\endcsname{ega4.c04.sec.4-6-18-c}
\expandafter\def\csname sga@target@IV.4.6.18.d\endcsname{ega4.c04.sec.4-6-18-d}
\expandafter\def\csname sga@target@IV.4.6.19\endcsname{ega4.c04.def.4-6-19}
\expandafter\def\csname sga@target@IV.4.6.2\endcsname{ega4.c04.def.4-6-2}
\expandafter\def\csname sga@target@IV.4.6.20\endcsname{ega4.c04.prop.4-6-20}
\expandafter\def\csname sga@target@IV.4.6.21\endcsname{ega4.c04.prop.4-6-21}
\expandafter\def\csname sga@target@IV.4.6.21.i\endcsname{ega4.c04.sec.4-6-21-i}
\expandafter\def\csname sga@target@IV.4.6.21.ii\endcsname{ega4.c04.sec.4-6-21-ii}
\expandafter\def\csname sga@target@IV.4.6.22\endcsname{ega4.c04.def.4-6-22}
\expandafter\def\csname sga@target@IV.4.6.3\endcsname{ega4.c04.cor.4-6-3}
\expandafter\def\csname sga@target@IV.4.6.4\endcsname{ega4.c04.cor.4-6-4}
\expandafter\def\csname sga@target@IV.4.6.5\endcsname{ega4.c04.prop.4-6-5}
\expandafter\def\csname sga@target@IV.4.6.5.i\endcsname{ega4.c04.sec.4-6-5-i}
\expandafter\def\csname sga@target@IV.4.6.5.ii\endcsname{ega4.c04.sec.4-6-5-ii}
\expandafter\def\csname sga@target@IV.4.6.6\endcsname{ega4.c04.prop.4-6-6}
\expandafter\def\csname sga@target@IV.4.6.7\endcsname{ega4.c04.cor.4-6-7}
\expandafter\def\csname sga@target@IV.4.6.8\endcsname{ega4.c04.cor.4-6-8}
\expandafter\def\csname sga@target@IV.4.6.9\endcsname{ega4.c04.def.4-6-9}
\expandafter\def\csname sga@target@IV.4.7.1\endcsname{ega4.c04.lem.4-7-1}
\expandafter\def\csname sga@target@IV.4.7.1.1\endcsname{ega4.c04.eq.4-7-1-1}
\expandafter\def\csname sga@target@IV.4.7.10\endcsname{ega4.c04.prop.4-7-10}
\expandafter\def\csname sga@target@IV.4.7.11\endcsname{ega4.c04.prop.4-7-11}
\expandafter\def\csname sga@target@IV.4.7.12\endcsname{ega4.c04.def.4-7-12}
\expandafter\def\csname sga@target@IV.4.7.13\endcsname{ega4.c04.prop.4-7-13}
\expandafter\def\csname sga@target@IV.4.7.2\endcsname{ega4.c04.cor.4-7-2}
\expandafter\def\csname sga@target@IV.4.7.2.1\endcsname{ega4.c04.eq.4-7-2-1}
\expandafter\def\csname sga@target@IV.4.7.3\endcsname{ega4.c04.prop.4-7-3}
\expandafter\def\csname sga@target@IV.4.7.3.1\endcsname{ega4.c04.eq.4-7-3-1}
\expandafter\def\csname sga@target@IV.4.7.3.i\endcsname{ega4.c04.sec.4-7-3-i}
\expandafter\def\csname sga@target@IV.4.7.3.ii\endcsname{ega4.c04.sec.4-7-3-ii}
\expandafter\def\csname sga@target@IV.4.7.3.iii\endcsname{ega4.c04.sec.4-7-3-iii}
\expandafter\def\csname sga@target@IV.4.7.4\endcsname{ega4.c04.def.4-7-4}
\expandafter\def\csname sga@target@IV.4.7.5\endcsname{ega4.c04.def.4-7-5}
\expandafter\def\csname sga@target@IV.4.7.6\endcsname{ega4.c04.prop.4-7-6}
\expandafter\def\csname sga@target@IV.4.7.6.1\endcsname{ega4.c04.eq.4-7-6-1}
\expandafter\def\csname sga@target@IV.4.7.7\endcsname{ega4.c04.statement.4-7-7}
\expandafter\def\csname sga@target@IV.4.7.8\endcsname{ega4.c04.prop.4-7-8}
\expandafter\def\csname sga@target@IV.4.7.9\endcsname{ega4.c04.cor.4-7-9}
\expandafter\def\csname sga@target@IV.4.8.1\endcsname{ega4.c04.statement.4-8-1}
\expandafter\def\csname sga@target@IV.4.8.10\endcsname{ega4.c04.cor.4-8-10}
\expandafter\def\csname sga@target@IV.4.8.11\endcsname{ega4.c04.cor.4-8-11}
\expandafter\def\csname sga@target@IV.4.8.12\endcsname{ega4.c04.cor.4-8-12}
\expandafter\def\csname sga@target@IV.4.8.13\endcsname{ega4.c04.prop.4-8-13}
\expandafter\def\csname sga@target@IV.4.8.14\endcsname{ega4.c04.prop.4-8-14}
\expandafter\def\csname sga@target@IV.4.8.15\endcsname{ega4.c04.sec.4-8-15}
\expandafter\def\csname sga@target@IV.4.8.15.i\endcsname{ega4.c04.sec.4-8-15-i}
\expandafter\def\csname sga@target@IV.4.8.15.ii\endcsname{ega4.c04.sec.4-8-15-ii}
\expandafter\def\csname sga@target@IV.4.8.2\endcsname{ega4.c04.statement.4-8-2}
\expandafter\def\csname sga@target@IV.4.8.2.1\endcsname{ega4.c04.eq.4-8-2-1}
\expandafter\def\csname sga@target@IV.4.8.2.2\endcsname{ega4.c04.eq.4-8-2-2}
\expandafter\def\csname sga@target@IV.4.8.2.3\endcsname{ega4.c04.eq.4-8-2-3}
\expandafter\def\csname sga@target@IV.4.8.2.4\endcsname{ega4.c04.eq.4-8-2-4}
\expandafter\def\csname sga@target@IV.4.8.2.5\endcsname{ega4.c04.eq.4-8-2-5}
\expandafter\def\csname sga@target@IV.4.8.2.6\endcsname{ega4.c04.eq.4-8-2-6}
\expandafter\def\csname sga@target@IV.4.8.3\endcsname{ega4.c04.prop.4-8-3}
\expandafter\def\csname sga@target@IV.4.8.4\endcsname{ega4.c04.def.4-8-4}
\expandafter\def\csname sga@target@IV.4.8.5\endcsname{ega4.c04.prop.4-8-5}
\expandafter\def\csname sga@target@IV.4.8.6\endcsname{ega4.c04.lem.4-8-6}
\expandafter\def\csname sga@target@IV.4.8.6.a\endcsname{ega4.c04.sec.4-8-6-a}
\expandafter\def\csname sga@target@IV.4.8.6.b\endcsname{ega4.c04.sec.4-8-6-b}
\expandafter\def\csname sga@target@IV.4.8.6.c\endcsname{ega4.c04.sec.4-8-6-c}
\expandafter\def\csname sga@target@IV.4.8.7\endcsname{ega4.c04.cor.4-8-7}
\expandafter\def\csname sga@target@IV.4.8.8\endcsname{ega4.c04.lem.4-8-8}
\expandafter\def\csname sga@target@IV.4.8.8.i\endcsname{ega4.c04.sec.4-8-8-i}
\expandafter\def\csname sga@target@IV.4.8.8.ii\endcsname{ega4.c04.sec.4-8-8-ii}
\expandafter\def\csname sga@target@IV.4.8.8.iii\endcsname{ega4.c04.sec.4-8-8-iii}
\expandafter\def\csname sga@target@IV.4.8.9\endcsname{ega4.c04.prop.4-8-9}
\expandafter\def\csname sga@target@IV.4.9.1\endcsname{ega4.c04.prop.4-9-1}
\expandafter\def\csname sga@target@IV.4.9.1.a\endcsname{ega4.c04.sec.4-9-1-a}
\expandafter\def\csname sga@target@IV.4.9.1.b\endcsname{ega4.c04.sec.4-9-1-b}
\expandafter\def\csname sga@target@IV.4.9.1.bprime\endcsname{ega4.c04.sec.4-9-1-bprime}
\expandafter\def\csname sga@target@IV.4.9.1.c\endcsname{ega4.c04.sec.4-9-1-c}
\expandafter\def\csname sga@target@IV.4.9.2\endcsname{ega4.c04.cor.4-9-2}
\expandafter\def\csname sga@target@IV.4.9.3\endcsname{ega4.c04.rem.4-9-3}
\expandafter\def\csname sga@target@IV.4.9.4\endcsname{ega4.c04.prop.4-9-4}
\expandafter\def\csname sga@target@IV.4.9.5\endcsname{ega4.c04.cor.4-9-5}
\expandafter\def\csname sga@target@IV.4.9.5.a\endcsname{ega4.c04.sec.4-9-5-a}
\expandafter\def\csname sga@target@IV.4.9.5.b\endcsname{ega4.c04.sec.4-9-5-b}
\expandafter\def\csname sga@target@IV.4.9.6\endcsname{ega4.c04.prop.4-9-6}
\expandafter\def\csname sga@target@IV.4.9.6.a\endcsname{ega4.c04.sec.4-9-6-a}
\expandafter\def\csname sga@target@IV.4.9.6.aprime\endcsname{ega4.c04.sec.4-9-6-aprime}
\expandafter\def\csname sga@target@IV.4.9.6.b\endcsname{ega4.c04.sec.4-9-6-b}
\expandafter\def\csname sga@target@IV.4.9.6.bprime\endcsname{ega4.c04.sec.4-9-6-bprime}
\expandafter\def\csname sga@target@IV.4.9.7\endcsname{ega4.c04.cor.4-9-7}
\expandafter\def\csname sga@target@IV.5.1.1\endcsname{ega4.c04.statement.5-1-1}
\expandafter\def\csname sga@target@IV.5.1.1.1\endcsname{ega4.c04.eq.5-1-1-1}
\expandafter\def\csname sga@target@IV.5.1.1.2\endcsname{ega4.c04.eq.5-1-1-2}
\expandafter\def\csname sga@target@IV.5.1.1.3\endcsname{ega4.c04.eq.5-1-1-3}
\expandafter\def\csname sga@target@IV.5.1.1.4\endcsname{ega4.c04.sec.5-1-1-4}
\expandafter\def\csname sga@target@IV.5.1.1.5\endcsname{ega4.c04.sec.5-1-1-5}
\expandafter\def\csname sga@target@IV.5.1.1.6\endcsname{ega4.c04.sec.5-1-1-6}
\expandafter\def\csname sga@target@IV.5.1.10\endcsname{ega4.c04.prop.5-1-10}
\expandafter\def\csname sga@target@IV.5.1.10.1\endcsname{ega4.c04.lem.5-1-10-1}
\expandafter\def\csname sga@target@IV.5.1.10.i\endcsname{ega4.c04.sec.5-1-10-i}
\expandafter\def\csname sga@target@IV.5.1.10.ii\endcsname{ega4.c04.sec.5-1-10-ii}
\expandafter\def\csname sga@target@IV.5.1.11\endcsname{ega4.c04.cor.5-1-11}
\expandafter\def\csname sga@target@IV.5.1.12\endcsname{ega4.c04.statement.5-1-12}
\expandafter\def\csname sga@target@IV.5.1.12.1\endcsname{ega4.c04.eq.5-1-12-1}
\expandafter\def\csname sga@target@IV.5.1.12.2\endcsname{ega4.c04.eq.5-1-12-2}
\expandafter\def\csname sga@target@IV.5.1.12.3\endcsname{ega4.c04.eq.5-1-12-3}
\expandafter\def\csname sga@target@IV.5.1.13\endcsname{ega4.c04.prop.5-1-13}
\expandafter\def\csname sga@target@IV.5.1.13.1\endcsname{ega4.c04.eq.5-1-13-1}
\expandafter\def\csname sga@target@IV.5.1.2\endcsname{ega4.c04.prop.5-1-2}
\expandafter\def\csname sga@target@IV.5.1.2.1\endcsname{ega4.c04.eq.5-1-2-1}
\expandafter\def\csname sga@target@IV.5.1.3\endcsname{ega4.c04.cor.5-1-3}
\expandafter\def\csname sga@target@IV.5.1.3.1\endcsname{ega4.c04.eq.5-1-3-1}
\expandafter\def\csname sga@target@IV.5.1.3.2\endcsname{ega4.c04.eq.5-1-3-2}
\expandafter\def\csname sga@target@IV.5.1.4\endcsname{ega4.c04.prop.5-1-4}
\expandafter\def\csname sga@target@IV.5.1.4.1\endcsname{ega4.c04.eq.5-1-4-1}
\expandafter\def\csname sga@target@IV.5.1.4.2\endcsname{ega4.c04.eq.5-1-4-2}
\expandafter\def\csname sga@target@IV.5.1.5\endcsname{ega4.c04.cor.5-1-5}
\expandafter\def\csname sga@target@IV.5.1.6\endcsname{ega4.c04.statement.5-1-6}
\expandafter\def\csname sga@target@IV.5.1.7\endcsname{ega4.c04.prop.5-1-7}
\expandafter\def\csname sga@target@IV.5.1.8\endcsname{ega4.c04.cor.5-1-8}
\expandafter\def\csname sga@target@IV.5.1.9\endcsname{ega4.c04.prop.5-1-9}
\expandafter\def\csname sga@target@IV.5.1.9.1\endcsname{ega4.c04.statement.5-1-9-1}
\expandafter\def\csname sga@target@IV.5.10.1\endcsname{ega4.c04.statement.5-10-1}
\expandafter\def\csname sga@target@IV.5.10.1.1\endcsname{ega4.c04.eq.5-10-1-1}
\expandafter\def\csname sga@target@IV.5.10.10\endcsname{ega4.c04.prop.5-10-10}
\expandafter\def\csname sga@target@IV.5.10.11\endcsname{ega4.c04.statement.5-10-11}
\expandafter\def\csname sga@target@IV.5.10.11.1\endcsname{ega4.c04.eq.5-10-11-1}
\expandafter\def\csname sga@target@IV.5.10.12\endcsname{ega4.c04.prop.5-10-12}
\expandafter\def\csname sga@target@IV.5.10.13\endcsname{ega4.c04.statement.5-10-13}
\expandafter\def\csname sga@target@IV.5.10.14\endcsname{ega4.c04.prop.5-10-14}
\expandafter\def\csname sga@target@IV.5.10.15\endcsname{ega4.c04.cor.5-10-15}
\expandafter\def\csname sga@target@IV.5.10.16\endcsname{ega4.c04.prop.5-10-16}
\expandafter\def\csname sga@target@IV.5.10.17\endcsname{ega4.c04.prop.5-10-17}
\expandafter\def\csname sga@target@IV.5.10.2\endcsname{ega4.c04.prop.5-10-2}
\expandafter\def\csname sga@target@IV.5.10.3\endcsname{ega4.c04.cor.5-10-3}
\expandafter\def\csname sga@target@IV.5.10.4\endcsname{ega4.c04.cor.5-10-4}
\expandafter\def\csname sga@target@IV.5.10.5\endcsname{ega4.c04.thm.5-10-5}
\expandafter\def\csname sga@target@IV.5.10.6\endcsname{ega4.c04.cor.5-10-6}
\expandafter\def\csname sga@target@IV.5.10.7\endcsname{ega4.c04.thm.5-10-7}
\expandafter\def\csname sga@target@IV.5.10.8\endcsname{ega4.c04.cor.5-10-8}
\expandafter\def\csname sga@target@IV.5.10.8.1\endcsname{ega4.c04.lem.5-10-8-1}
\expandafter\def\csname sga@target@IV.5.10.9\endcsname{ega4.c04.cor.5-10-9}
\expandafter\def\csname sga@target@IV.5.11.1\endcsname{ega4.c04.prop.5-11-1}
\expandafter\def\csname sga@target@IV.5.11.1.1\endcsname{ega4.c04.lem.5-11-1-1}
\expandafter\def\csname sga@target@IV.5.11.2\endcsname{ega4.c04.cor.5-11-2}
\expandafter\def\csname sga@target@IV.5.11.3\endcsname{ega4.c04.rem.5-11-3}
\expandafter\def\csname sga@target@IV.5.11.4\endcsname{ega4.c04.cor.5-11-4}
\expandafter\def\csname sga@target@IV.5.11.5\endcsname{ega4.c04.cor.5-11-5}
\expandafter\def\csname sga@target@IV.5.11.6\endcsname{ega4.c04.cor.5-11-6}
\expandafter\def\csname sga@target@IV.5.11.7\endcsname{ega4.c04.item.5-11-7}
\expandafter\def\csname sga@target@IV.5.12\endcsname{ega4.c04.subsec.5-12.key-iv-5-12}
\expandafter\def\csname sga@target@IV.5.12.1\endcsname{ega4.c04.prop.5-12-1}
\expandafter\def\csname sga@target@IV.5.12.1.1\endcsname{ega4.c04.eq.5-12-1-1}
\expandafter\def\csname sga@target@IV.5.12.10\endcsname{ega4.c04.cor.5-12-10}
\expandafter\def\csname sga@target@IV.5.12.2\endcsname{ega4.c04.prop.5-12-2}
\expandafter\def\csname sga@target@IV.5.12.3\endcsname{ega4.c04.rem.5-12-3}
\expandafter\def\csname sga@target@IV.5.12.4\endcsname{ega4.c04.cor.5-12-4}
\expandafter\def\csname sga@target@IV.5.12.5\endcsname{ega4.c04.prop.5-12-5}
\expandafter\def\csname sga@target@IV.5.12.6\endcsname{ega4.c04.rem.5-12-6}
\expandafter\def\csname sga@target@IV.5.12.7\endcsname{ega4.c04.cor.5-12-7}
\expandafter\def\csname sga@target@IV.5.12.8\endcsname{ega4.c04.prop.5-12-8}
\expandafter\def\csname sga@target@IV.5.12.9\endcsname{ega4.c04.rem.5-12-9}
\expandafter\def\csname sga@target@IV.5.13.1\endcsname{ega4.c04.prop.5-13-1}
\expandafter\def\csname sga@target@IV.5.13.2\endcsname{ega4.c04.cor.5-13-2}
\expandafter\def\csname sga@target@IV.5.13.3\endcsname{ega4.c04.prop.5-13-3}
\expandafter\def\csname sga@target@IV.5.13.4\endcsname{ega4.c04.cor.5-13-4}
\expandafter\def\csname sga@target@IV.5.13.5\endcsname{ega4.c04.statement.5-13-5}
\expandafter\def\csname sga@target@IV.5.13.6\endcsname{ega4.c04.cor.5-13-6}
\expandafter\def\csname sga@target@IV.5.13.7\endcsname{ega4.c04.prop.5-13-7}
\expandafter\def\csname sga@target@IV.5.13.7.1\endcsname{ega4.c04.lem.5-13-7-1}
\expandafter\def\csname sga@target@IV.5.13.8\endcsname{ega4.c04.cor.5-13-8}
\expandafter\def\csname sga@target@IV.5.2.1\endcsname{ega4.c04.prop.5-2-1}
\expandafter\def\csname sga@target@IV.5.2.1.1\endcsname{ega4.c04.eq.5-2-1-1}
\expandafter\def\csname sga@target@IV.5.2.2\endcsname{ega4.c04.cor.5-2-2}
\expandafter\def\csname sga@target@IV.5.2.3\endcsname{ega4.c04.cor.5-2-3}
\expandafter\def\csname sga@target@IV.5.2.3.1\endcsname{ega4.c04.eq.5-2-3-1}
\expandafter\def\csname sga@target@IV.5.2.4\endcsname{ega4.c04.cor.5-2-4}
\expandafter\def\csname sga@target@IV.5.2.5\endcsname{ega4.c04.sec.5-2-5}
\expandafter\def\csname sga@target@IV.5.2.5.i\endcsname{ega4.c04.sec.5-2-5-i}
\expandafter\def\csname sga@target@IV.5.2.5.ii\endcsname{ega4.c04.sec.5-2-5-ii}
\expandafter\def\csname sga@target@IV.5.3.1\endcsname{ega4.c04.prop.5-3-1}
\expandafter\def\csname sga@target@IV.5.3.1.1\endcsname{ega4.c04.eq.5-3-1-1}
\expandafter\def\csname sga@target@IV.5.3.1.2\endcsname{ega4.c04.eq.5-3-1-2}
\expandafter\def\csname sga@target@IV.5.3.1.3\endcsname{ega4.c04.lem.5-3-1-3}
\expandafter\def\csname sga@target@IV.5.4.1\endcsname{ega4.c04.prop.5-4-1}
\expandafter\def\csname sga@target@IV.5.4.1.i\endcsname{ega4.c04.sec.5-4-1-i}
\expandafter\def\csname sga@target@IV.5.4.1.ii\endcsname{ega4.c04.sec.5-4-1-ii}
\expandafter\def\csname sga@target@IV.5.4.2\endcsname{ega4.c04.cor.5-4-2}
\expandafter\def\csname sga@target@IV.5.4.3\endcsname{ega4.c04.sec.5-4-3}
\expandafter\def\csname sga@target@IV.5.4.3.1\endcsname{ega4.c04.eq.5-4-3-1}
\expandafter\def\csname sga@target@IV.5.4.3.i\endcsname{ega4.c04.sec.5-4-3-i}
\expandafter\def\csname sga@target@IV.5.4.3.ii\endcsname{ega4.c04.sec.5-4-3-ii}
\expandafter\def\csname sga@target@IV.5.5\endcsname{ega4.c04.sec.5-5}
\expandafter\def\csname sga@target@IV.5.5.1\endcsname{ega4.c04.statement.5-5-1}
\expandafter\def\csname sga@target@IV.5.5.1.1\endcsname{ega4.c04.eq.5-5-1-1}
\expandafter\def\csname sga@target@IV.5.5.2\endcsname{ega4.c04.prop.5-5-2}
\expandafter\def\csname sga@target@IV.5.5.2.1\endcsname{ega4.c04.eq.5-5-2-1}
\expandafter\def\csname sga@target@IV.5.5.2.2\endcsname{ega4.c04.eq.5-5-2-2}
\expandafter\def\csname sga@target@IV.5.5.3\endcsname{ega4.c04.prop.5-5-3}
\expandafter\def\csname sga@target@IV.5.5.3.1\endcsname{ega4.c04.eq.5-5-3-1}
\expandafter\def\csname sga@target@IV.5.5.4\endcsname{ega4.c04.cor.5-5-4}
\expandafter\def\csname sga@target@IV.5.5.4.1\endcsname{ega4.c04.eq.5-5-4-1}
\expandafter\def\csname sga@target@IV.5.5.5\endcsname{ega4.c04.cor.5-5-5}
\expandafter\def\csname sga@target@IV.5.5.6\endcsname{ega4.c04.cor.5-5-6}
\expandafter\def\csname sga@target@IV.5.5.6.1\endcsname{ega4.c04.eq.5-5-6-1}
\expandafter\def\csname sga@target@IV.5.5.7\endcsname{ega4.c04.lem.5-5-7}
\expandafter\def\csname sga@target@IV.5.5.7.1\endcsname{ega4.c04.eq.5-5-7-1}
\expandafter\def\csname sga@target@IV.5.5.7.2\endcsname{ega4.c04.statement.5-5-7-2}
\expandafter\def\csname sga@target@IV.5.5.7.3\endcsname{ega4.c04.eq.5-5-7-3}
\expandafter\def\csname sga@target@IV.5.5.7.i\endcsname{ega4.c04.sec.5-5-7-i}
\expandafter\def\csname sga@target@IV.5.5.7.ii\endcsname{ega4.c04.sec.5-5-7-ii}
\expandafter\def\csname sga@target@IV.5.5.8\endcsname{ega4.c04.thm.5-5-8}
\expandafter\def\csname sga@target@IV.5.5.8.1\endcsname{ega4.c04.eq.5-5-8-1}
\expandafter\def\csname sga@target@IV.5.5.8.2\endcsname{ega4.c04.eq.5-5-8-2}
\expandafter\def\csname sga@target@IV.5.5.8.3\endcsname{ega4.c04.eq.5-5-8-3}
\expandafter\def\csname sga@target@IV.5.6\endcsname{ega4.c04.sec.5-6}
\expandafter\def\csname sga@target@IV.5.6.1\endcsname{ega4.c04.prop.5-6-1}
\expandafter\def\csname sga@target@IV.5.6.1.1\endcsname{ega4.c04.eq.5-6-1-1}
\expandafter\def\csname sga@target@IV.5.6.1.2\endcsname{ega4.c04.eq.5-6-1-2}
\expandafter\def\csname sga@target@IV.5.6.1.3\endcsname{ega4.c04.statement.5-6-1-3}
\expandafter\def\csname sga@target@IV.5.6.1.4\endcsname{ega4.c04.statement.5-6-1-4}
\expandafter\def\csname sga@target@IV.5.6.1.5\endcsname{ega4.c04.statement.5-6-1-5}
\expandafter\def\csname sga@target@IV.5.6.1.a\endcsname{ega4.c04.sec.5-6-1-a}
\expandafter\def\csname sga@target@IV.5.6.1.b\endcsname{ega4.c04.sec.5-6-1-b}
\expandafter\def\csname sga@target@IV.5.6.1.c\endcsname{ega4.c04.sec.5-6-1-c}
\expandafter\def\csname sga@target@IV.5.6.10\endcsname{ega4.c04.prop.5-6-10}
\expandafter\def\csname sga@target@IV.5.6.11\endcsname{ega4.c04.ex.5-6-11}
\expandafter\def\csname sga@target@IV.5.6.11.1\endcsname{ega4.c04.lem.5-6-11-1}
\expandafter\def\csname sga@target@IV.5.6.11.1.i\endcsname{ega4.c04.sec.5-6-11-1-i}
\expandafter\def\csname sga@target@IV.5.6.11.1.ii\endcsname{ega4.c04.sec.5-6-11-1-ii}
\expandafter\def\csname sga@target@IV.5.6.2\endcsname{ega4.c04.def.5-6-2}
\expandafter\def\csname sga@target@IV.5.6.3\endcsname{ega4.c04.sec.5-6-3}
\expandafter\def\csname sga@target@IV.5.6.3.i\endcsname{ega4.c04.sec.5-6-3-i}
\expandafter\def\csname sga@target@IV.5.6.3.ii\endcsname{ega4.c04.sec.5-6-3-ii}
\expandafter\def\csname sga@target@IV.5.6.3.iii\endcsname{ega4.c04.sec.5-6-3-iii}
\expandafter\def\csname sga@target@IV.5.6.3.iv\endcsname{ega4.c04.sec.5-6-3-iv}
\expandafter\def\csname sga@target@IV.5.6.4\endcsname{ega4.c04.prop.5-6-4}
\expandafter\def\csname sga@target@IV.5.6.5\endcsname{ega4.c04.prop.5-6-5}
\expandafter\def\csname sga@target@IV.5.6.5.1\endcsname{ega4.c04.eq.5-6-5-1}
\expandafter\def\csname sga@target@IV.5.6.5.2\endcsname{ega4.c04.eq.5-6-5-2}
\expandafter\def\csname sga@target@IV.5.6.5.3\endcsname{ega4.c04.eq.5-6-5-3}
\expandafter\def\csname sga@target@IV.5.6.6\endcsname{ega4.c04.cor.5-6-6}
\expandafter\def\csname sga@target@IV.5.6.6.1\endcsname{ega4.c04.eq.5-6-6-1}
\expandafter\def\csname sga@target@IV.5.6.6.2\endcsname{ega4.c04.eq.5-6-6-2}
\expandafter\def\csname sga@target@IV.5.6.7\endcsname{ega4.c04.cor.5-6-7}
\expandafter\def\csname sga@target@IV.5.6.7.1\endcsname{ega4.c04.eq.5-6-7-1}
\expandafter\def\csname sga@target@IV.5.6.8\endcsname{ega4.c04.cor.5-6-8}
\expandafter\def\csname sga@target@IV.5.6.8.1\endcsname{ega4.c04.eq.5-6-8-1}
\expandafter\def\csname sga@target@IV.5.6.8.2\endcsname{ega4.c04.eq.5-6-8-2}
\expandafter\def\csname sga@target@IV.5.6.9\endcsname{ega4.c04.sec.5-6-9}
\expandafter\def\csname sga@target@IV.5.6.9.i\endcsname{ega4.c04.sec.5-6-9-i}
\expandafter\def\csname sga@target@IV.5.6.9.ii\endcsname{ega4.c04.sec.5-6-9-ii}
\expandafter\def\csname sga@target@IV.5.7.1\endcsname{ega4.c04.def.5-7-1}
\expandafter\def\csname sga@target@IV.5.7.1.1\endcsname{ega4.c04.eq.5-7-1-1}
\expandafter\def\csname sga@target@IV.5.7.10\endcsname{ega4.c04.cor.5-7-10}
\expandafter\def\csname sga@target@IV.5.7.11\endcsname{ega4.c04.cor.5-7-11}
\expandafter\def\csname sga@target@IV.5.7.11.i\endcsname{ega4.c04.sec.5-7-11-i}
\expandafter\def\csname sga@target@IV.5.7.11.ii\endcsname{ega4.c04.sec.5-7-11-ii}
\expandafter\def\csname sga@target@IV.5.7.12\endcsname{ega4.c04.statement.5-7-12}
\expandafter\def\csname sga@target@IV.5.7.13\endcsname{ega4.c04.cor.5-7-13}
\expandafter\def\csname sga@target@IV.5.7.13.1\endcsname{ega4.c04.eq.5-7-13-1}
\expandafter\def\csname sga@target@IV.5.7.2\endcsname{ega4.c04.def.5-7-2}
\expandafter\def\csname sga@target@IV.5.7.2.1\endcsname{ega4.c04.eq.5-7-2-1}
\expandafter\def\csname sga@target@IV.5.7.2.2\endcsname{ega4.c04.eq.5-7-2-2}
\expandafter\def\csname sga@target@IV.5.7.3\endcsname{ega4.c04.sec.5-7-3}
\expandafter\def\csname sga@target@IV.5.7.3.i\endcsname{ega4.c04.sec.5-7-3-i}
\expandafter\def\csname sga@target@IV.5.7.3.ii\endcsname{ega4.c04.sec.5-7-3-ii}
\expandafter\def\csname sga@target@IV.5.7.3.iii\endcsname{ega4.c04.sec.5-7-3-iii}
\expandafter\def\csname sga@target@IV.5.7.3.iv\endcsname{ega4.c04.sec.5-7-3-iv}
\expandafter\def\csname sga@target@IV.5.7.3.v\endcsname{ega4.c04.sec.5-7-3-v}
\expandafter\def\csname sga@target@IV.5.7.3.vi\endcsname{ega4.c04.sec.5-7-3-vi}
\expandafter\def\csname sga@target@IV.5.7.4\endcsname{ega4.c04.prop.5-7-4}
\expandafter\def\csname sga@target@IV.5.7.4.1\endcsname{ega4.c04.eq.5-7-4-1}
\expandafter\def\csname sga@target@IV.5.7.4.2\endcsname{ega4.c04.eq.5-7-4-2}
\expandafter\def\csname sga@target@IV.5.7.4.i\endcsname{ega4.c04.sec.5-7-4-i}
\expandafter\def\csname sga@target@IV.5.7.4.ii\endcsname{ega4.c04.sec.5-7-4-ii}
\expandafter\def\csname sga@target@IV.5.7.5\endcsname{ega4.c04.prop.5-7-5}
\expandafter\def\csname sga@target@IV.5.7.6\endcsname{ega4.c04.cor.5-7-6}
\expandafter\def\csname sga@target@IV.5.7.7\endcsname{ega4.c04.cor.5-7-7}
\expandafter\def\csname sga@target@IV.5.7.8\endcsname{ega4.c04.sec.5-7-8}
\expandafter\def\csname sga@target@IV.5.7.9\endcsname{ega4.c04.prop.5-7-9}
\expandafter\def\csname sga@target@IV.5.7.9.1\endcsname{ega4.c04.eq.5-7-9-1}
\expandafter\def\csname sga@target@IV.5.7.9.2\endcsname{ega4.c04.eq.5-7-9-2}
\expandafter\def\csname sga@target@IV.5.8.1\endcsname{ega4.c04.statement.5-8-1}
\expandafter\def\csname sga@target@IV.5.8.2\endcsname{ega4.c04.def.5-8-2}
\expandafter\def\csname sga@target@IV.5.8.3\endcsname{ega4.c04.prop.5-8-3}
\expandafter\def\csname sga@target@IV.5.8.4\endcsname{ega4.c04.statement.5-8-4}
\expandafter\def\csname sga@target@IV.5.8.5\endcsname{ega4.c04.prop.5-8-5}
\expandafter\def\csname sga@target@IV.5.8.6\endcsname{ega4.c04.thm.5-8-6}
\expandafter\def\csname sga@target@IV.5.8.6.1\endcsname{ega4.c04.eq.5-8-6-1}
\expandafter\def\csname sga@target@IV.5.8.6.i\endcsname{ega4.c04.item.5-8-6-i}
\expandafter\def\csname sga@target@IV.5.8.6.ii\endcsname{ega4.c04.item.5-8-6-ii}
\expandafter\def\csname sga@target@IV.5.9.1\endcsname{ega4.c04.statement.5-9-1}
\expandafter\def\csname sga@target@IV.5.9.1.1\endcsname{ega4.c04.eq.5-9-1-1}
\expandafter\def\csname sga@target@IV.5.9.1.2\endcsname{ega4.c04.eq.5-9-1-2}
\expandafter\def\csname sga@target@IV.5.9.1.3\endcsname{ega4.c04.eq.5-9-1-3}
\expandafter\def\csname sga@target@IV.5.9.10\endcsname{ega4.c04.cor.5-9-10}
\expandafter\def\csname sga@target@IV.5.9.10.i\endcsname{ega4.c04.item.5-9-10-i}
\expandafter\def\csname sga@target@IV.5.9.10.ii\endcsname{ega4.c04.item.5-9-10-ii}
\expandafter\def\csname sga@target@IV.5.9.11\endcsname{ega4.c04.prop.5-9-11}
\expandafter\def\csname sga@target@IV.5.9.12\endcsname{ega4.c04.item.5-9-12}
\expandafter\def\csname sga@target@IV.5.9.12.i\endcsname{ega4.c04.item.5-9-12-i}
\expandafter\def\csname sga@target@IV.5.9.12.ii\endcsname{ega4.c04.item.5-9-12-ii}
\expandafter\def\csname sga@target@IV.5.9.12.iii\endcsname{ega4.c04.item.5-9-12-iii}
\expandafter\def\csname sga@target@IV.5.9.2\endcsname{ega4.c04.prop.5-9-2}
\expandafter\def\csname sga@target@IV.5.9.2.i\endcsname{ega4.c04.item.5-9-2-i}
\expandafter\def\csname sga@target@IV.5.9.2.ii\endcsname{ega4.c04.item.5-9-2-ii}
\expandafter\def\csname sga@target@IV.5.9.3\endcsname{ega4.c04.rem.5-9-3}
\expandafter\def\csname sga@target@IV.5.9.3.1\endcsname{ega4.c04.eq.5-9-3-1}
\expandafter\def\csname sga@target@IV.5.9.4\endcsname{ega4.c04.prop.5-9-4}
\expandafter\def\csname sga@target@IV.5.9.4.1\endcsname{ega4.c04.eq.5-9-4-1}
\expandafter\def\csname sga@target@IV.5.9.5\endcsname{ega4.c04.cor.5-9-5}
\expandafter\def\csname sga@target@IV.5.9.6\endcsname{ega4.c04.cor.5-9-6}
\expandafter\def\csname sga@target@IV.5.9.6.1\endcsname{ega4.c04.eq.5-9-6-1}
\expandafter\def\csname sga@target@IV.5.9.7\endcsname{ega4.c04.statement.5-9-7}
\expandafter\def\csname sga@target@IV.5.9.7.1\endcsname{ega4.c04.eq.5-9-7-1}
\expandafter\def\csname sga@target@IV.5.9.8\endcsname{ega4.c04.prop.5-9-8}
\expandafter\def\csname sga@target@IV.5.9.8.1\endcsname{ega4.c04.eq.5-9-8-1}
\expandafter\def\csname sga@target@IV.5.9.8.a\endcsname{ega4.c04.item.5-9-8-a}
\expandafter\def\csname sga@target@IV.5.9.8.a-prime\endcsname{ega4.c04.item.5-9-8-a-prime}
\expandafter\def\csname sga@target@IV.5.9.8.b\endcsname{ega4.c04.item.5-9-8-b}
\expandafter\def\csname sga@target@IV.5.9.8.b-prime\endcsname{ega4.c04.item.5-9-8-b-prime}
\expandafter\def\csname sga@target@IV.5.9.9\endcsname{ega4.c04.def.5-9-9}
\expandafter\def\csname sga@target@IV.6.1.1\endcsname{ega4.c04.prop.6-1-1}
\expandafter\def\csname sga@target@IV.6.1.1.1\endcsname{ega4.c04.eq.6-1-1-1}
\expandafter\def\csname sga@target@IV.6.1.2\endcsname{ega4.c04.cor.6-1-2}
\expandafter\def\csname sga@target@IV.6.1.2.1\endcsname{ega4.c04.eq.6-1-2-1}
\expandafter\def\csname sga@target@IV.6.1.3\endcsname{ega4.c04.cor.6-1-3}
\expandafter\def\csname sga@target@IV.6.1.4\endcsname{ega4.c04.cor.6-1-4}
\expandafter\def\csname sga@target@IV.6.1.4.1\endcsname{ega4.c04.eq.6-1-4-1}
\expandafter\def\csname sga@target@IV.6.1.5\endcsname{ega4.c04.prop.6-1-5}
\expandafter\def\csname sga@target@IV.6.10.1\endcsname{ega4.c04.statement.6-10-1}
\expandafter\def\csname sga@target@IV.6.10.2\endcsname{ega4.c04.prop.6-10-2}
\expandafter\def\csname sga@target@IV.6.10.3\endcsname{ega4.c04.prop.6-10-3}
\expandafter\def\csname sga@target@IV.6.10.3.1\endcsname{ega4.c04.eq.6-10-3-1}
\expandafter\def\csname sga@target@IV.6.10.3.2\endcsname{ega4.c04.eq.6-10-3-2}
\expandafter\def\csname sga@target@IV.6.10.3.3\endcsname{ega4.c04.eq.6-10-3-3}
\expandafter\def\csname sga@target@IV.6.10.4\endcsname{ega4.c04.cor.6-10-4}
\expandafter\def\csname sga@target@IV.6.10.5\endcsname{ega4.c04.prop.6-10-5}
\expandafter\def\csname sga@target@IV.6.10.5.1\endcsname{ega4.c04.eq.6-10-5-1}
\expandafter\def\csname sga@target@IV.6.10.5.2\endcsname{ega4.c04.eq.6-10-5-2}
\expandafter\def\csname sga@target@IV.6.10.5.3\endcsname{ega4.c04.eq.6-10-5-3}
\expandafter\def\csname sga@target@IV.6.10.5.4\endcsname{ega4.c04.eq.6-10-5-4}
\expandafter\def\csname sga@target@IV.6.10.6\endcsname{ega4.c04.prop.6-10-6}
\expandafter\def\csname sga@target@IV.6.10.6.1\endcsname{ega4.c04.eq.6-10-6-1}
\expandafter\def\csname sga@target@IV.6.10.6.2\endcsname{ega4.c04.eq.6-10-6-2}
\expandafter\def\csname sga@target@IV.6.10.6.3\endcsname{ega4.c04.eq.6-10-6-3}
\expandafter\def\csname sga@target@IV.6.11\endcsname{ega4.c04.subsec.6-11.key-iv-6-11}
\expandafter\def\csname sga@target@IV.6.11.1\endcsname{ega4.c04.lem.6-11-1}
\expandafter\def\csname sga@target@IV.6.11.2\endcsname{ega4.c04.prop.6-11-2}
\expandafter\def\csname sga@target@IV.6.11.2.1\endcsname{ega4.c04.eq.6-11-2-1}
\expandafter\def\csname sga@target@IV.6.11.2.2\endcsname{ega4.c04.eq.6-11-2-2}
\expandafter\def\csname sga@target@IV.6.11.3\endcsname{ega4.c04.cor.6-11-3}
\expandafter\def\csname sga@target@IV.6.11.4\endcsname{ega4.c04.statement.6-11-4}
\expandafter\def\csname sga@target@IV.6.11.5\endcsname{ega4.c04.prop.6-11-5}
\expandafter\def\csname sga@target@IV.6.11.5.1\endcsname{ega4.c04.eq.6-11-5-1}
\expandafter\def\csname sga@target@IV.6.11.6\endcsname{ega4.c04.prop.6-11-6}
\expandafter\def\csname sga@target@IV.6.11.6.1\endcsname{ega4.c04.lem.6-11-6-1}
\expandafter\def\csname sga@target@IV.6.11.7\endcsname{ega4.c04.cor.6-11-7}
\expandafter\def\csname sga@target@IV.6.11.8\endcsname{ega4.c04.prop.6-11-8}
\expandafter\def\csname sga@target@IV.6.11.8.1\endcsname{ega4.c04.eq.6-11-8-1}
\expandafter\def\csname sga@target@IV.6.11.9\endcsname{ega4.c04.statement.6-11-9}
\expandafter\def\csname sga@target@IV.6.12\endcsname{ega4.c04.subsec.6-12.key-iv-6-12}
\expandafter\def\csname sga@target@IV.6.12.1\endcsname{ega4.c04.statement.6-12-1}
\expandafter\def\csname sga@target@IV.6.12.10\endcsname{ega4.c04.prop.6-12-10}
\expandafter\def\csname sga@target@IV.6.12.2\endcsname{ega4.c04.prop.6-12-2}
\expandafter\def\csname sga@target@IV.6.12.3\endcsname{ega4.c04.cor.6-12-3}
\expandafter\def\csname sga@target@IV.6.12.4\endcsname{ega4.c04.thm.6-12-4}
\expandafter\def\csname sga@target@IV.6.12.4.1\endcsname{ega4.c04.lem.6-12-4-1}
\expandafter\def\csname sga@target@IV.6.12.5\endcsname{ega4.c04.cor.6-12-5}
\expandafter\def\csname sga@target@IV.6.12.6\endcsname{ega4.c04.cor.6-12-6}
\expandafter\def\csname sga@target@IV.6.12.7\endcsname{ega4.c04.thm.6-12-7}
\expandafter\def\csname sga@target@IV.6.12.8\endcsname{ega4.c04.cor.6-12-8}
\expandafter\def\csname sga@target@IV.6.12.9\endcsname{ega4.c04.prop.6-12-9}
\expandafter\def\csname sga@target@IV.6.13\endcsname{ega4.c04.subsec.6-13.key-iv-6-13}
\expandafter\def\csname sga@target@IV.6.13.1\endcsname{ega4.c04.statement.6-13-1}
\expandafter\def\csname sga@target@IV.6.13.2\endcsname{ega4.c04.prop.6-13-2}
\expandafter\def\csname sga@target@IV.6.13.3\endcsname{ega4.c04.cor.6-13-3}
\expandafter\def\csname sga@target@IV.6.13.4\endcsname{ega4.c04.prop.6-13-4}
\expandafter\def\csname sga@target@IV.6.13.5\endcsname{ega4.c04.cor.6-13-5}
\expandafter\def\csname sga@target@IV.6.13.6\endcsname{ega4.c04.prop.6-13-6}
\expandafter\def\csname sga@target@IV.6.13.7\endcsname{ega4.c04.prop.6-13-7}
\expandafter\def\csname sga@target@IV.6.13.7.1\endcsname{ega4.c04.lem.6-13-7-1}
\expandafter\def\csname sga@target@IV.6.14.1\endcsname{ega4.c04.prop.6-14-1}
\expandafter\def\csname sga@target@IV.6.14.1.1\endcsname{ega4.c04.lem.6-14-1-1}
\expandafter\def\csname sga@target@IV.6.14.1.2\endcsname{ega4.c04.lem.6-14-1-2}
\expandafter\def\csname sga@target@IV.6.14.1.2.env\endcsname{ega4.c04.statement.6-14-1-2-env}
\expandafter\def\csname sga@target@IV.6.14.1.3\endcsname{ega4.c04.eq.6-14-1-3}
\expandafter\def\csname sga@target@IV.6.14.1.4\endcsname{ega4.c04.lem.6-14-1-4}
\expandafter\def\csname sga@target@IV.6.14.2\endcsname{ega4.c04.cor.6-14-2}
\expandafter\def\csname sga@target@IV.6.14.3\endcsname{ega4.c04.cor.6-14-3}
\expandafter\def\csname sga@target@IV.6.14.4\endcsname{ega4.c04.prop.6-14-4}
\expandafter\def\csname sga@target@IV.6.14.4.1\endcsname{ega4.c04.lem.6-14-4-1}
\expandafter\def\csname sga@target@IV.6.14.5\endcsname{ega4.c04.cor.6-14-5}
\expandafter\def\csname sga@target@IV.6.15.1\endcsname{ega4.c04.statement.6-15-1}
\expandafter\def\csname sga@target@IV.6.15.10\endcsname{ega4.c04.prop.6-15-10}
\expandafter\def\csname sga@target@IV.6.15.11\endcsname{ega4.c04.item.6-15-11}
\expandafter\def\csname sga@target@IV.6.15.2\endcsname{ega4.c04.statement.6-15-2}
\expandafter\def\csname sga@target@IV.6.15.3\endcsname{ega4.c04.statement.6-15-3}
\expandafter\def\csname sga@target@IV.6.15.3.1\endcsname{ega4.c04.lem.6-15-3-1}
\expandafter\def\csname sga@target@IV.6.15.4\endcsname{ega4.c04.statement.6-15-4}
\expandafter\def\csname sga@target@IV.6.15.4.1\endcsname{ega4.c04.lem.6-15-4-1}
\expandafter\def\csname sga@target@IV.6.15.5\endcsname{ega4.c04.prop.6-15-5}
\expandafter\def\csname sga@target@IV.6.15.6\endcsname{ega4.c04.prop.6-15-6}
\expandafter\def\csname sga@target@IV.6.15.7\endcsname{ega4.c04.prop.6-15-7}
\expandafter\def\csname sga@target@IV.6.15.8\endcsname{ega4.c04.lem.6-15-8}
\expandafter\def\csname sga@target@IV.6.15.9\endcsname{ega4.c04.cor.6-15-9}
\expandafter\def\csname sga@target@IV.6.2.1\endcsname{ega4.c04.prop.6-2-1}
\expandafter\def\csname sga@target@IV.6.2.1.1\endcsname{ega4.c04.eq.6-2-1-1}
\expandafter\def\csname sga@target@IV.6.2.1.2\endcsname{ega4.c04.eq.6-2-1-2}
\expandafter\def\csname sga@target@IV.6.2.2\endcsname{ega4.c04.cor.6-2-2}
\expandafter\def\csname sga@target@IV.6.2.3\endcsname{ega4.c04.prop.6-2-3}
\expandafter\def\csname sga@target@IV.6.2.3.1\endcsname{ega4.c04.eq.6-2-3-1}
\expandafter\def\csname sga@target@IV.6.2.4\endcsname{ega4.c04.prop.6-2-4}
\expandafter\def\csname sga@target@IV.6.2.4.1\endcsname{ega4.c04.eq.6-2-4-1}
\expandafter\def\csname sga@target@IV.6.3.1\endcsname{ega4.c04.prop.6-3-1}
\expandafter\def\csname sga@target@IV.6.3.1.1\endcsname{ega4.c04.eq.6-3-1-1}
\expandafter\def\csname sga@target@IV.6.3.2\endcsname{ega4.c04.cor.6-3-2}
\expandafter\def\csname sga@target@IV.6.3.2.1\endcsname{ega4.c04.eq.6-3-2-1}
\expandafter\def\csname sga@target@IV.6.3.3\endcsname{ega4.c04.cor.6-3-3}
\expandafter\def\csname sga@target@IV.6.3.4\endcsname{ega4.c04.cor.6-3-4}
\expandafter\def\csname sga@target@IV.6.3.4.1\endcsname{ega4.c04.eq.6-3-4-1}
\expandafter\def\csname sga@target@IV.6.3.4.2\endcsname{ega4.c04.eq.6-3-4-2}
\expandafter\def\csname sga@target@IV.6.3.5\endcsname{ega4.c04.cor.6-3-5}
\expandafter\def\csname sga@target@IV.6.3.6\endcsname{ega4.c04.cor.6-3-6}
\expandafter\def\csname sga@target@IV.6.3.7\endcsname{ega4.c04.cor.6-3-7}
\expandafter\def\csname sga@target@IV.6.3.8\endcsname{ega4.c04.prop.6-3-8}
\expandafter\def\csname sga@target@IV.6.4.1\endcsname{ega4.c04.prop.6-4-1}
\expandafter\def\csname sga@target@IV.6.4.2\endcsname{ega4.c04.cor.6-4-2}
\expandafter\def\csname sga@target@IV.6.4.3\endcsname{ega4.c04.rem.6-4-3}
\expandafter\def\csname sga@target@IV.6.5.1\endcsname{ega4.c04.prop.6-5-1}
\expandafter\def\csname sga@target@IV.6.5.2\endcsname{ega4.c04.cor.6-5-2}
\expandafter\def\csname sga@target@IV.6.5.3\endcsname{ega4.c04.prop.6-5-3}
\expandafter\def\csname sga@target@IV.6.5.4\endcsname{ega4.c04.cor.6-5-4}
\expandafter\def\csname sga@target@IV.6.5.5\endcsname{ega4.c04.item.6-5-5}
\expandafter\def\csname sga@target@IV.6.6.1\endcsname{ega4.c04.prop.6-6-1}
\expandafter\def\csname sga@target@IV.6.6.2\endcsname{ega4.c04.rem.6-6-2}
\expandafter\def\csname sga@target@IV.6.7.1\endcsname{ega4.c04.prop.6-7-1}
\expandafter\def\csname sga@target@IV.6.7.1.1\endcsname{ega4.c04.lem.6-7-1-1}
\expandafter\def\csname sga@target@IV.6.7.2\endcsname{ega4.c04.cor.6-7-2}
\expandafter\def\csname sga@target@IV.6.7.3\endcsname{ega4.c04.cor.6-7-3}
\expandafter\def\csname sga@target@IV.6.7.4\endcsname{ega4.c04.prop.6-7-4}
\expandafter\def\csname sga@target@IV.6.7.4.1\endcsname{ega4.c04.lem.6-7-4-1}
\expandafter\def\csname sga@target@IV.6.7.5\endcsname{ega4.c04.item.6-7-5}
\expandafter\def\csname sga@target@IV.6.7.6\endcsname{ega4.c04.def.6-7-6}
\expandafter\def\csname sga@target@IV.6.7.7\endcsname{ega4.c04.prop.6-7-7}
\expandafter\def\csname sga@target@IV.6.7.8\endcsname{ega4.c04.cor.6-7-8}
\expandafter\def\csname sga@target@IV.6.8.1\endcsname{ega4.c04.def.6-8-1}
\expandafter\def\csname sga@target@IV.6.8.2\endcsname{ega4.c04.prop.6-8-2}
\expandafter\def\csname sga@target@IV.6.8.3\endcsname{ega4.c04.prop.6-8-3}
\expandafter\def\csname sga@target@IV.6.8.4\endcsname{ega4.c04.item.6-8-4}
\expandafter\def\csname sga@target@IV.6.8.5\endcsname{ega4.c04.prop.6-8-5}
\expandafter\def\csname sga@target@IV.6.8.6\endcsname{ega4.c04.thm.6-8-6}
\expandafter\def\csname sga@target@IV.6.8.7\endcsname{ega4.c04.cor.6-8-7}
\expandafter\def\csname sga@target@IV.6.8.8\endcsname{ega4.c04.rem.6-8-8}
\expandafter\def\csname sga@target@IV.6.9.1\endcsname{ega4.c04.thm.6-9-1}
\expandafter\def\csname sga@target@IV.6.9.2\endcsname{ega4.c04.lem.6-9-2}
\expandafter\def\csname sga@target@IV.6.9.3\endcsname{ega4.c04.cor.6-9-3}
\expandafter\def\csname sga@target@IV.7.1.1\endcsname{ega4.c04.def.7-1-1}
\expandafter\def\csname sga@target@IV.7.1.10\endcsname{ega4.c04.prop.7-1-10}
\expandafter\def\csname sga@target@IV.7.1.11\endcsname{ega4.c04.prop.7-1-11}
\expandafter\def\csname sga@target@IV.7.1.12\endcsname{ega4.c04.item.7-1-12}
\expandafter\def\csname sga@target@IV.7.1.13\endcsname{ega4.c04.cor.7-1-13}
\expandafter\def\csname sga@target@IV.7.1.2\endcsname{ega4.c04.prop.7-1-2}
\expandafter\def\csname sga@target@IV.7.1.3\endcsname{ega4.c04.prop.7-1-3}
\expandafter\def\csname sga@target@IV.7.1.3.1\endcsname{ega4.c04.eq.7-1-3-1}
\expandafter\def\csname sga@target@IV.7.1.3.2\endcsname{ega4.c04.eq.7-1-3-2}
\expandafter\def\csname sga@target@IV.7.1.3.3\endcsname{ega4.c04.eq.7-1-3-3}
\expandafter\def\csname sga@target@IV.7.1.3.4\endcsname{ega4.c04.eq.7-1-3-4}
\expandafter\def\csname sga@target@IV.7.1.3.5\endcsname{ega4.c04.eq.7-1-3-5}
\expandafter\def\csname sga@target@IV.7.1.3.6\endcsname{ega4.c04.eq.7-1-3-6}
\expandafter\def\csname sga@target@IV.7.1.3.7\endcsname{ega4.c04.eq.7-1-3-7}
\expandafter\def\csname sga@target@IV.7.1.4\endcsname{ega4.c04.cor.7-1-4}
\expandafter\def\csname sga@target@IV.7.1.5\endcsname{ega4.c04.cor.7-1-5}
\expandafter\def\csname sga@target@IV.7.1.6\endcsname{ega4.c04.rem.7-1-6}
\expandafter\def\csname sga@target@IV.7.1.7\endcsname{ega4.c04.lem.7-1-7}
\expandafter\def\csname sga@target@IV.7.1.8\endcsname{ega4.c04.thm.7-1-8}
\expandafter\def\csname sga@target@IV.7.1.9\endcsname{ega4.c04.def.7-1-9}
\expandafter\def\csname sga@target@IV.7.2.1\endcsname{ega4.c04.statement.7-2-1}
\expandafter\def\csname sga@target@IV.7.2.1.1\endcsname{ega4.c04.statement.7-2-1-1}
\expandafter\def\csname sga@target@IV.7.2.10\endcsname{ega4.c04.statement.7-2-10}
\expandafter\def\csname sga@target@IV.7.2.2\endcsname{ega4.c04.prop.7-2-2}
\expandafter\def\csname sga@target@IV.7.2.3\endcsname{ega4.c04.prop.7-2-3}
\expandafter\def\csname sga@target@IV.7.2.4\endcsname{ega4.c04.prop.7-2-4}
\expandafter\def\csname sga@target@IV.7.2.5\endcsname{ega4.c04.thm.7-2-5}
\expandafter\def\csname sga@target@IV.7.2.5.1\endcsname{ega4.c04.eq.7-2-5-1}
\expandafter\def\csname sga@target@IV.7.2.6\endcsname{ega4.c04.def.7-2-6}
\expandafter\def\csname sga@target@IV.7.2.7\endcsname{ega4.c04.cor.7-2-7}
\expandafter\def\csname sga@target@IV.7.2.8\endcsname{ega4.c04.cor.7-2-8}
\expandafter\def\csname sga@target@IV.7.2.9\endcsname{ega4.c04.cor.7-2-9}
\expandafter\def\csname sga@target@IV.7.3.1\endcsname{ega4.c04.statement.7-3-1}
\expandafter\def\csname sga@target@IV.7.3.10\endcsname{ega4.c04.statement.7-3-10}
\expandafter\def\csname sga@target@IV.7.3.11\endcsname{ega4.c04.rem.7-3-11}
\expandafter\def\csname sga@target@IV.7.3.12\endcsname{ega4.c04.prop.7-3-12}
\expandafter\def\csname sga@target@IV.7.3.12.1\endcsname{ega4.c04.eq.7-3-12-1}
\expandafter\def\csname sga@target@IV.7.3.13\endcsname{ega4.c04.statement.7-3-13}
\expandafter\def\csname sga@target@IV.7.3.14\endcsname{ega4.c04.prop.7-3-14}
\expandafter\def\csname sga@target@IV.7.3.15\endcsname{ega4.c04.prop.7-3-15}
\expandafter\def\csname sga@target@IV.7.3.16\endcsname{ega4.c04.prop.7-3-16}
\expandafter\def\csname sga@target@IV.7.3.16.1\endcsname{ega4.c04.eq.7-3-16-1}
\expandafter\def\csname sga@target@IV.7.3.16.2\endcsname{ega4.c04.lem.7-3-16-2}
\expandafter\def\csname sga@target@IV.7.3.17\endcsname{ega4.c04.prop.7-3-17}
\expandafter\def\csname sga@target@IV.7.3.18\endcsname{ega4.c04.prop.7-3-18}
\expandafter\def\csname sga@target@IV.7.3.19\endcsname{ega4.c04.statement.7-3-19}
\expandafter\def\csname sga@target@IV.7.3.2\endcsname{ega4.c04.lem.7-3-2}
\expandafter\def\csname sga@target@IV.7.3.3\endcsname{ega4.c04.cor.7-3-3}
\expandafter\def\csname sga@target@IV.7.3.4\endcsname{ega4.c04.statement.7-3-4}
\expandafter\def\csname sga@target@IV.7.3.5\endcsname{ega4.c04.item.7-3-5}
\expandafter\def\csname sga@target@IV.7.3.6\endcsname{ega4.c04.statement.7-3-6}
\expandafter\def\csname sga@target@IV.7.3.7\endcsname{ega4.c04.lem.7-3-7}
\expandafter\def\csname sga@target@IV.7.3.8\endcsname{ega4.c04.statement.7-3-8}
\expandafter\def\csname sga@target@IV.7.3.9\endcsname{ega4.c04.item.7-3-9}
\expandafter\def\csname sga@target@IV.7.4.1\endcsname{ega4.c04.statement.7-4-1}
\expandafter\def\csname sga@target@IV.7.4.2\endcsname{ega4.c04.lem.7-4-2}
\expandafter\def\csname sga@target@IV.7.4.3\endcsname{ega4.c04.cor.7-4-3}
\expandafter\def\csname sga@target@IV.7.4.4\endcsname{ega4.c04.thm.7-4-4}
\expandafter\def\csname sga@target@IV.7.4.4.1\endcsname{ega4.c04.lem.7-4-4-1}
\expandafter\def\csname sga@target@IV.7.4.4.2\endcsname{ega4.c04.lem.7-4-4-2}
\expandafter\def\csname sga@target@IV.7.4.5\endcsname{ega4.c04.cor.7-4-5}
\expandafter\def\csname sga@target@IV.7.4.6\endcsname{ega4.c04.prop.7-4-6}
\expandafter\def\csname sga@target@IV.7.4.7\endcsname{ega4.c04.cor.7-4-7}
\expandafter\def\csname sga@target@IV.7.4.8\endcsname{ega4.c04.rem.7-4-8}
\expandafter\def\csname sga@target@IV.7.5.0\endcsname{ega4.c04.statement.7-5-0}
\expandafter\def\csname sga@target@IV.7.5.1\endcsname{ega4.c04.thm.7-5-1}
\expandafter\def\csname sga@target@IV.7.5.1.1\endcsname{ega4.c04.lem.7-5-1-1}
\expandafter\def\csname sga@target@IV.7.5.2\endcsname{ega4.c04.cor.7-5-2}
\expandafter\def\csname sga@target@IV.7.5.3\endcsname{ega4.c04.item.7-5-3}
\expandafter\def\csname sga@target@IV.7.5.4\endcsname{ega4.c04.item.7-5-4}
\expandafter\def\csname sga@target@IV.7.5.5\endcsname{ega4.c04.lem.7-5-5}
\expandafter\def\csname sga@target@IV.7.5.6\endcsname{ega4.c04.prop.7-5-6}
\expandafter\def\csname sga@target@IV.7.5.7\endcsname{ega4.c04.cor.7-5-7}
\expandafter\def\csname sga@target@IV.7.5.7.1\endcsname{ega4.c04.lem.7-5-7-1}
\expandafter\def\csname sga@target@IV.7.5.8\endcsname{ega4.c04.rem.7-5-8}
\expandafter\def\csname sga@target@IV.7.6.1\endcsname{ega4.c04.prop.7-6-1}
\expandafter\def\csname sga@target@IV.7.6.2\endcsname{ega4.c04.cor.7-6-2}
\expandafter\def\csname sga@target@IV.7.6.3\endcsname{ega4.c04.cor.7-6-3}
\expandafter\def\csname sga@target@IV.7.6.4\endcsname{ega4.c04.thm.7-6-4}
\expandafter\def\csname sga@target@IV.7.6.4.1\endcsname{ega4.c04.lem.7-6-4-1}
\expandafter\def\csname sga@target@IV.7.6.5\endcsname{ega4.c04.cor.7-6-5}
\expandafter\def\csname sga@target@IV.7.6.6\endcsname{ega4.c04.cor.7-6-6}
\expandafter\def\csname sga@target@IV.7.6.7\endcsname{ega4.c04.item.7-6-7}
\expandafter\def\csname sga@target@IV.7.7.1\endcsname{ega4.c04.statement.7-7-1}
\expandafter\def\csname sga@target@IV.7.7.2\endcsname{ega4.c04.thm.7-7-2}
\expandafter\def\csname sga@target@IV.7.7.3\endcsname{ega4.c04.cor.7-7-3}
\expandafter\def\csname sga@target@IV.7.7.4\endcsname{ega4.c04.cor.7-7-4}
\expandafter\def\csname sga@target@IV.7.8.1\endcsname{ega4.c04.statement.7-8-1}
\expandafter\def\csname sga@target@IV.7.8.2\endcsname{ega4.c04.def.7-8-2}
\expandafter\def\csname sga@target@IV.7.8.3\endcsname{ega4.c04.statement.7-8-3}
\expandafter\def\csname sga@target@IV.7.8.3.1\endcsname{ega4.c04.statement.7-8-3-1}
\expandafter\def\csname sga@target@IV.7.8.4\endcsname{ega4.c04.item.7-8-4}
\expandafter\def\csname sga@target@IV.7.8.4.1\endcsname{ega4.c04.statement.7-8-4-1}
\expandafter\def\csname sga@target@IV.7.8.5\endcsname{ega4.c04.statement.7-8-5}
\expandafter\def\csname sga@target@IV.7.8.6\endcsname{ega4.c04.prop.7-8-6}
\expandafter\def\csname sga@target@IV.7.9.1\endcsname{ega4.c04.statement.7-9-1}
\expandafter\def\csname sga@target@IV.7.9.10\endcsname{ega4.c04.item.7-9-10}
\expandafter\def\csname sga@target@IV.7.9.2\endcsname{ega4.c04.statement.7-9-2}
\expandafter\def\csname sga@target@IV.7.9.3\endcsname{ega4.c04.prop.7-9-3}
\expandafter\def\csname sga@target@IV.7.9.3.1\endcsname{ega4.c04.lem.7-9-3-1}
\expandafter\def\csname sga@target@IV.7.9.4\endcsname{ega4.c04.cor.7-9-4}
\expandafter\def\csname sga@target@IV.7.9.5\endcsname{ega4.c04.prop.7-9-5}
\expandafter\def\csname sga@target@IV.7.9.6\endcsname{ega4.c04.rem.7-9-6}
\expandafter\def\csname sga@target@IV.7.9.7\endcsname{ega4.c04.statement.7-9-7}
\expandafter\def\csname sga@target@IV.7.9.8\endcsname{ega4.c04.prop.7-9-8}
\expandafter\def\csname sga@target@IV.7.9.8.1\endcsname{ega4.c04.statement.7-9-8-1}
\expandafter\def\csname sga@target@IV.7.9.9\endcsname{ega4.c04.cor.7-9-9}
\expandafter\def\csname sga@target@IV.8.1.1\endcsname{ega4.c04.statement.8-1-1}
\expandafter\def\csname sga@target@IV.8.1.2\endcsname{ega4.c04.statement.8-1-2}
\expandafter\def\csname sga@target@IV.8.1.3\endcsname{ega4.c04.statement.8-1-3}
\expandafter\def\csname sga@target@IV.8.10.1\endcsname{ega4.c04.prop.8-10-1}
\expandafter\def\csname sga@target@IV.8.10.2\endcsname{ega4.c04.cor.8-10-2}
\expandafter\def\csname sga@target@IV.8.10.3\endcsname{ega4.c04.prop.8-10-3}
\expandafter\def\csname sga@target@IV.8.10.3.1\endcsname{ega4.c04.rem.8-10-3-1}
\expandafter\def\csname sga@target@IV.8.10.4\endcsname{ega4.c04.prop.8-10-4}
\expandafter\def\csname sga@target@IV.8.10.5\endcsname{ega4.c04.thm.8-10-5}
\expandafter\def\csname sga@target@IV.8.10.5.1\endcsname{ega4.c04.prop.8-10-5-1}
\expandafter\def\csname sga@target@IV.8.10.5.2\endcsname{ega4.c04.lem.8-10-5-2}
\expandafter\def\csname sga@target@IV.8.11\endcsname{ega4.c04.subsec.8-11.key-iv-8-11}
\expandafter\def\csname sga@target@IV.8.11.1\endcsname{ega4.c04.thm.8-11-1}
\expandafter\def\csname sga@target@IV.8.11.2\endcsname{ega4.c04.thm.8-11-2}
\expandafter\def\csname sga@target@IV.8.11.3\endcsname{ega4.c04.rem.8-11-3}
\expandafter\def\csname sga@target@IV.8.11.4\endcsname{ega4.c04.statement.8-11-4}
\expandafter\def\csname sga@target@IV.8.11.5\endcsname{ega4.c04.prop.8-11-5}
\expandafter\def\csname sga@target@IV.8.11.5.1\endcsname{ega4.c04.rem.8-11-5-1}
\expandafter\def\csname sga@target@IV.8.11.6\endcsname{ega4.c04.prop.8-11-6}
\expandafter\def\csname sga@target@IV.8.12\endcsname{ega4.c04.subsec.8-12.key-iv-8-12}
\expandafter\def\csname sga@target@IV.8.12.1\endcsname{ega4.c04.lem.8-12-1}
\expandafter\def\csname sga@target@IV.8.12.10\endcsname{ega4.c04.cor.8-12-10}
\expandafter\def\csname sga@target@IV.8.12.10.1\endcsname{ega4.c04.lem.8-12-10-1}
\expandafter\def\csname sga@target@IV.8.12.11\endcsname{ega4.c04.cor.8-12-11}
\expandafter\def\csname sga@target@IV.8.12.12\endcsname{ega4.c04.cor.8-12-12}
\expandafter\def\csname sga@target@IV.8.12.2\endcsname{ega4.c04.lem.8-12-2}
\expandafter\def\csname sga@target@IV.8.12.3\endcsname{ega4.c04.prop.8-12-3}
\expandafter\def\csname sga@target@IV.8.12.4\endcsname{ega4.c04.cor.8-12-4}
\expandafter\def\csname sga@target@IV.8.12.5\endcsname{ega4.c04.statement.8-12-5}
\expandafter\def\csname sga@target@IV.8.12.6\endcsname{ega4.c04.thm.8-12-6}
\expandafter\def\csname sga@target@IV.8.12.6.1\endcsname{ega4.c04.eq.8-12-6-1}
\expandafter\def\csname sga@target@IV.8.12.7\endcsname{ega4.c04.rem.8-12-7}
\expandafter\def\csname sga@target@IV.8.12.8\endcsname{ega4.c04.cor.8-12-8}
\expandafter\def\csname sga@target@IV.8.12.9\endcsname{ega4.c04.cor.8-12-9}
\expandafter\def\csname sga@target@IV.8.13.1\endcsname{ega4.c04.prop.8-13-1}
\expandafter\def\csname sga@target@IV.8.13.1.1\endcsname{ega4.c04.eq.8-13-1-1}
\expandafter\def\csname sga@target@IV.8.13.2\endcsname{ega4.c04.cor.8-13-2}
\expandafter\def\csname sga@target@IV.8.13.2.1\endcsname{ega4.c04.eq.8-13-2-1}
\expandafter\def\csname sga@target@IV.8.13.3\endcsname{ega4.c04.statement.8-13-3}
\expandafter\def\csname sga@target@IV.8.13.4\endcsname{ega4.c04.statement.8-13-4}
\expandafter\def\csname sga@target@IV.8.13.4.1\endcsname{ega4.c04.eq.8-13-4-1}
\expandafter\def\csname sga@target@IV.8.13.4.2\endcsname{ega4.c04.eq.8-13-4-2}
\expandafter\def\csname sga@target@IV.8.13.5\endcsname{ega4.c04.prop.8-13-5}
\expandafter\def\csname sga@target@IV.8.13.6\endcsname{ega4.c04.rem.8-13-6}
\expandafter\def\csname sga@target@IV.8.14.1\endcsname{ega4.c04.statement.8-14-1}
\expandafter\def\csname sga@target@IV.8.14.2\endcsname{ega4.c04.prop.8-14-2}
\expandafter\def\csname sga@target@IV.8.14.2.1\endcsname{ega4.c04.eq.8-14-2-1}
\expandafter\def\csname sga@target@IV.8.14.2.2\endcsname{ega4.c04.cor.8-14-2-2}
\expandafter\def\csname sga@target@IV.8.14.2.3\endcsname{ega4.c04.eq.8-14-2-3}
\expandafter\def\csname sga@target@IV.8.14.2.4\endcsname{ega4.c04.eq.8-14-2-4}
\expandafter\def\csname sga@target@IV.8.14.3\endcsname{ega4.c04.rem.8-14-3}
\expandafter\def\csname sga@target@IV.8.2.1\endcsname{ega4.c04.statement.8-2-1}
\expandafter\def\csname sga@target@IV.8.2.1.1\endcsname{ega4.c04.eq.8-2-1-1}
\expandafter\def\csname sga@target@IV.8.2.10\endcsname{ega4.c04.cor.8-2-10}
\expandafter\def\csname sga@target@IV.8.2.11\endcsname{ega4.c04.cor.8-2-11}
\expandafter\def\csname sga@target@IV.8.2.12\endcsname{ega4.c04.cor.8-2-12}
\expandafter\def\csname sga@target@IV.8.2.12.1\endcsname{ega4.c04.eq.8-2-12-1}
\expandafter\def\csname sga@target@IV.8.2.13\endcsname{ega4.c04.prop.8-2-13}
\expandafter\def\csname sga@target@IV.8.2.14\endcsname{ega4.c04.rem.8-2-14}
\expandafter\def\csname sga@target@IV.8.2.2\endcsname{ega4.c04.statement.8-2-2}
\expandafter\def\csname sga@target@IV.8.2.3\endcsname{ega4.c04.prop.8-2-3}
\expandafter\def\csname sga@target@IV.8.2.3.1\endcsname{ega4.c04.eq.8-2-3-1}
\expandafter\def\csname sga@target@IV.8.2.4\endcsname{ega4.c04.lem.8-2-4}
\expandafter\def\csname sga@target@IV.8.2.5\endcsname{ega4.c04.prop.8-2-5}
\expandafter\def\csname sga@target@IV.8.2.6\endcsname{ega4.c04.lem.8-2-6}
\expandafter\def\csname sga@target@IV.8.2.7\endcsname{ega4.c04.rem.8-2-7}
\expandafter\def\csname sga@target@IV.8.2.8\endcsname{ega4.c04.statement.8-2-8}
\expandafter\def\csname sga@target@IV.8.2.9\endcsname{ega4.c04.prop.8-2-9}
\expandafter\def\csname sga@target@IV.8.3.1\endcsname{ega4.c04.statement.8-3-1}
\expandafter\def\csname sga@target@IV.8.3.10\endcsname{ega4.c04.statement.8-3-10}
\expandafter\def\csname sga@target@IV.8.3.10.1\endcsname{ega4.c04.eq.8-3-10-1}
\expandafter\def\csname sga@target@IV.8.3.11\endcsname{ega4.c04.thm.8-3-11}
\expandafter\def\csname sga@target@IV.8.3.12\endcsname{ega4.c04.cor.8-3-12}
\expandafter\def\csname sga@target@IV.8.3.13\endcsname{ega4.c04.prop.8-3-13}
\expandafter\def\csname sga@target@IV.8.3.2\endcsname{ega4.c04.thm.8-3-2}
\expandafter\def\csname sga@target@IV.8.3.2.1\endcsname{ega4.c04.eq.8-3-2-1}
\expandafter\def\csname sga@target@IV.8.3.3\endcsname{ega4.c04.cor.8-3-3}
\expandafter\def\csname sga@target@IV.8.3.4\endcsname{ega4.c04.cor.8-3-4}
\expandafter\def\csname sga@target@IV.8.3.5\endcsname{ega4.c04.cor.8-3-5}
\expandafter\def\csname sga@target@IV.8.3.6\endcsname{ega4.c04.cor.8-3-6}
\expandafter\def\csname sga@target@IV.8.3.7\endcsname{ega4.c04.cor.8-3-7}
\expandafter\def\csname sga@target@IV.8.3.7.1\endcsname{ega4.c04.eq.8-3-7-1}
\expandafter\def\csname sga@target@IV.8.3.8\endcsname{ega4.c04.prop.8-3-8}
\expandafter\def\csname sga@target@IV.8.3.9\endcsname{ega4.c04.statement.8-3-9}
\expandafter\def\csname sga@target@IV.8.3.9.1\endcsname{ega4.c04.eq.8-3-9-1}
\expandafter\def\csname sga@target@IV.8.3.9.2\endcsname{ega4.c04.eq.8-3-9-2}
\expandafter\def\csname sga@target@IV.8.3.9.3\endcsname{ega4.c04.eq.8-3-9-3}
\expandafter\def\csname sga@target@IV.8.3.9.4\endcsname{ega4.c04.eq.8-3-9-4}
\expandafter\def\csname sga@target@IV.8.3.9.5\endcsname{ega4.c04.eq.8-3-9-5}
\expandafter\def\csname sga@target@IV.8.4.1\endcsname{ega4.c04.prop.8-4-1}
\expandafter\def\csname sga@target@IV.8.4.2\endcsname{ega4.c04.prop.8-4-2}
\expandafter\def\csname sga@target@IV.8.4.2.1\endcsname{ega4.c04.lem.8-4-2-1}
\expandafter\def\csname sga@target@IV.8.4.3\endcsname{ega4.c04.cor.8-4-3}
\expandafter\def\csname sga@target@IV.8.4.4\endcsname{ega4.c04.prop.8-4-4}
\expandafter\def\csname sga@target@IV.8.4.5\endcsname{ega4.c04.cor.8-4-5}
\expandafter\def\csname sga@target@IV.8.4.6\endcsname{ega4.c04.rem.8-4-6}
\expandafter\def\csname sga@target@IV.8.5.1\endcsname{ega4.c04.statement.8-5-1}
\expandafter\def\csname sga@target@IV.8.5.1.1\endcsname{ega4.c04.eq.8-5-1-1}
\expandafter\def\csname sga@target@IV.8.5.1.2\endcsname{ega4.c04.eq.8-5-1-2}
\expandafter\def\csname sga@target@IV.8.5.1.3\endcsname{ega4.c04.eq.8-5-1-3}
\expandafter\def\csname sga@target@IV.8.5.1.4\endcsname{ega4.c04.eq.8-5-1-4}
\expandafter\def\csname sga@target@IV.8.5.10\endcsname{ega4.c04.statement.8-5-10}
\expandafter\def\csname sga@target@IV.8.5.10.1\endcsname{ega4.c04.eq.8-5-10-1}
\expandafter\def\csname sga@target@IV.8.5.10.2\endcsname{ega4.c04.eq.8-5-10-2}
\expandafter\def\csname sga@target@IV.8.5.11\endcsname{ega4.c04.prop.8-5-11}
\expandafter\def\csname sga@target@IV.8.5.2\endcsname{ega4.c04.thm.8-5-2}
\expandafter\def\csname sga@target@IV.8.5.2.1\endcsname{ega4.c04.lem.8-5-2-1}
\expandafter\def\csname sga@target@IV.8.5.2.2\endcsname{ega4.c04.eq.8-5-2-2}
\expandafter\def\csname sga@target@IV.8.5.2.3\endcsname{ega4.c04.eq.8-5-2-3}
\expandafter\def\csname sga@target@IV.8.5.2.4\endcsname{ega4.c04.cor.8-5-2-4}
\expandafter\def\csname sga@target@IV.8.5.2.5\endcsname{ega4.c04.cor.8-5-2-5}
\expandafter\def\csname sga@target@IV.8.5.3\endcsname{ega4.c04.statement.8-5-3}
\expandafter\def\csname sga@target@IV.8.5.4\endcsname{ega4.c04.cor.8-5-4}
\expandafter\def\csname sga@target@IV.8.5.5\endcsname{ega4.c04.prop.8-5-5}
\expandafter\def\csname sga@target@IV.8.5.6\endcsname{ega4.c04.prop.8-5-6}
\expandafter\def\csname sga@target@IV.8.5.7\endcsname{ega4.c04.cor.8-5-7}
\expandafter\def\csname sga@target@IV.8.5.8\endcsname{ega4.c04.cor.8-5-8}
\expandafter\def\csname sga@target@IV.8.5.9\endcsname{ega4.c04.lem.8-5-9}
\expandafter\def\csname sga@target@IV.8.6.1\endcsname{ega4.c04.statement.8-6-1}
\expandafter\def\csname sga@target@IV.8.6.2\endcsname{ega4.c04.statement.8-6-2}
\expandafter\def\csname sga@target@IV.8.6.2.1\endcsname{ega4.c04.eq.8-6-2-1}
\expandafter\def\csname sga@target@IV.8.6.2.2\endcsname{ega4.c04.eq.8-6-2-2}
\expandafter\def\csname sga@target@IV.8.6.2.3\endcsname{ega4.c04.eq.8-6-2-3}
\expandafter\def\csname sga@target@IV.8.6.3\endcsname{ega4.c04.prop.8-6-3}
\expandafter\def\csname sga@target@IV.8.6.3.1\endcsname{ega4.c04.cor.8-6-3-1}
\expandafter\def\csname sga@target@IV.8.6.4\endcsname{ega4.c04.cor.8-6-4}
\expandafter\def\csname sga@target@IV.8.7.1\endcsname{ega4.c04.prop.8-7-1}
\expandafter\def\csname sga@target@IV.8.7.2\endcsname{ega4.c04.prop.8-7-2}
\expandafter\def\csname sga@target@IV.8.7.3\endcsname{ega4.c04.cor.8-7-3}
\expandafter\def\csname sga@target@IV.8.7.4\endcsname{ega4.c04.rem.8-7-4}
\expandafter\def\csname sga@target@IV.8.8.1\endcsname{ega4.c04.statement.8-8-1}
\expandafter\def\csname sga@target@IV.8.8.1.1\endcsname{ega4.c04.eq.8-8-1-1}
\expandafter\def\csname sga@target@IV.8.8.2\endcsname{ega4.c04.thm.8-8-2}
\expandafter\def\csname sga@target@IV.8.8.2.1\endcsname{ega4.c04.lem.8-8-2-1}
\expandafter\def\csname sga@target@IV.8.8.2.2\endcsname{ega4.c04.lem.8-8-2-2}
\expandafter\def\csname sga@target@IV.8.8.2.3\endcsname{ega4.c04.lem.8-8-2-3}
\expandafter\def\csname sga@target@IV.8.8.2.4\endcsname{ega4.c04.cor.8-8-2-4}
\expandafter\def\csname sga@target@IV.8.8.2.5\endcsname{ega4.c04.cor.8-8-2-5}
\expandafter\def\csname sga@target@IV.8.8.3\endcsname{ega4.c04.statement.8-8-3}
\expandafter\def\csname sga@target@IV.8.9.1\endcsname{ega4.c04.prop.8-9-1}
\expandafter\def\csname sga@target@IV.8.9.2\endcsname{ega4.c04.statement.8-9-2}
\expandafter\def\csname sga@target@IV.8.9.3\endcsname{ega4.c04.prop.8-9-3}
\expandafter\def\csname sga@target@IV.8.9.4\endcsname{ega4.c04.prop.8-9-4}
\expandafter\def\csname sga@target@IV.8.9.4.1\endcsname{ega4.c04.lem.8-9-4-1}
\expandafter\def\csname sga@target@IV.8.9.5\endcsname{ega4.c04.cor.8-9-5}
\expandafter\def\csname sga@target@IV.9.1.1\endcsname{ega4.c04.prop.9-1-1}
\expandafter\def\csname sga@target@IV.9.1.2\endcsname{ega4.c04.cor.9-1-2}
\expandafter\def\csname sga@target@IV.9.1.3\endcsname{ega4.c04.rem.9-1-3}
\expandafter\def\csname sga@target@IV.9.1.4\endcsname{ega4.c04.prop.9-1-4}
\expandafter\def\csname sga@target@IV.9.1.5\endcsname{ega4.c04.cor.9-1-5}
\expandafter\def\csname sga@target@IV.9.1.6\endcsname{ega4.c04.prop.9-1-6}
\expandafter\def\csname sga@target@IV.9.2\endcsname{ega4.c04.subsec.9-2.key-iv-9-2}
\expandafter\def\csname sga@target@IV.9.2.1\endcsname{ega4.c04.def.9-2-1}
\expandafter\def\csname sga@target@IV.9.2.2\endcsname{ega4.c04.item.9-2-2}
\expandafter\def\csname sga@target@IV.9.2.3\endcsname{ega4.c04.prop.9-2-3}
\expandafter\def\csname sga@target@IV.9.2.4\endcsname{ega4.c04.cor.9-2-4}
\expandafter\def\csname sga@target@IV.9.2.5\endcsname{ega4.c04.rem.9-2-5}
\expandafter\def\csname sga@target@IV.9.2.6\endcsname{ega4.c04.prop.9-2-6}
\expandafter\def\csname sga@target@IV.9.2.6.1\endcsname{ega4.c04.prop.9-2-6-1}
\expandafter\def\csname sga@target@IV.9.2.6.2\endcsname{ega4.c04.cor.9-2-6-2}
\expandafter\def\csname sga@target@IV.9.3.1\endcsname{ega4.c04.prop.9-3-1}
\expandafter\def\csname sga@target@IV.9.3.2\endcsname{ega4.c04.prop.9-3-2}
\expandafter\def\csname sga@target@IV.9.3.3\endcsname{ega4.c04.prop.9-3-3}
\expandafter\def\csname sga@target@IV.9.3.4\endcsname{ega4.c04.rem.9-3-4}
\expandafter\def\csname sga@target@IV.9.3.5\endcsname{ega4.c04.prop.9-3-5}
\expandafter\def\csname sga@target@IV.9.3.6\endcsname{ega4.c04.statement.9-3-6}
\expandafter\def\csname sga@target@IV.9.4.1\endcsname{ega4.c04.statement.9-4-1}
\expandafter\def\csname sga@target@IV.9.4.2\endcsname{ega4.c04.prop.9-4-2}
\expandafter\def\csname sga@target@IV.9.4.3\endcsname{ega4.c04.cor.9-4-3}
\expandafter\def\csname sga@target@IV.9.4.3.1\endcsname{ega4.c04.eq.9-4-3-1}
\expandafter\def\csname sga@target@IV.9.4.4\endcsname{ega4.c04.prop.9-4-4}
\expandafter\def\csname sga@target@IV.9.4.5\endcsname{ega4.c04.cor.9-4-5}
\expandafter\def\csname sga@target@IV.9.4.6\endcsname{ega4.c04.cor.9-4-6}
\expandafter\def\csname sga@target@IV.9.4.7\endcsname{ega4.c04.prop.9-4-7}
\expandafter\def\csname sga@target@IV.9.4.7.1\endcsname{ega4.c04.lem.9-4-7-1}
\expandafter\def\csname sga@target@IV.9.4.7.2\endcsname{ega4.c04.rem.9-4-7-2}
\expandafter\def\csname sga@target@IV.9.4.8\endcsname{ega4.c04.prop.9-4-8}
\expandafter\def\csname sga@target@IV.9.4.9\endcsname{ega4.c04.rem.9-4-9}
\expandafter\def\csname sga@target@IV.9.5.1\endcsname{ega4.c04.prop.9-5-1}
\expandafter\def\csname sga@target@IV.9.5.2\endcsname{ega4.c04.cor.9-5-2}
\expandafter\def\csname sga@target@IV.9.5.3\endcsname{ega4.c04.prop.9-5-3}
\expandafter\def\csname sga@target@IV.9.5.3.case2\endcsname{ega4.c04.item.9-5-3-case2}
\expandafter\def\csname sga@target@IV.9.5.4\endcsname{ega4.c04.cor.9-5-4}
\expandafter\def\csname sga@target@IV.9.5.5\endcsname{ega4.c04.prop.9-5-5}
\expandafter\def\csname sga@target@IV.9.5.6\endcsname{ega4.c04.cor.9-5-6}
\expandafter\def\csname sga@target@IV.9.5.7\endcsname{ega4.c04.rem.9-5-7}
\expandafter\def\csname sga@target@IV.9.6.1\endcsname{ega4.c04.prop.9-6-1}
\expandafter\def\csname sga@target@IV.9.6.2\endcsname{ega4.c04.prop.9-6-2}
\expandafter\def\csname sga@target@IV.9.6.3\endcsname{ega4.c04.prop.9-6-3}
\expandafter\def\csname sga@target@IV.9.6.3.1\endcsname{ega4.c04.eq.9-6-3-1}
\expandafter\def\csname sga@target@IV.9.6.4\endcsname{ega4.c04.cor.9-6-4}
\expandafter\def\csname sga@target@IV.9.6.5\endcsname{ega4.c04.cor.9-6-5}
\expandafter\def\csname sga@target@IV.9.6.6\endcsname{ega4.c04.sec.9-6-6}
\expandafter\def\csname sga@target@IV.9.6.7\endcsname{ega4.c04.sec.9-6-7}
\expandafter\def\csname sga@target@IV.9.7\endcsname{ega4.c04.subsec.9-7.key-iv-9-7}
\expandafter\def\csname sga@target@IV.9.7.1\endcsname{ega4.c04.lem.9-7-1}
\expandafter\def\csname sga@target@IV.9.7.10\endcsname{ega4.c04.cor.9-7-10}
\expandafter\def\csname sga@target@IV.9.7.11\endcsname{ega4.c04.rem.9-7-11}
\expandafter\def\csname sga@target@IV.9.7.12\endcsname{ega4.c04.prop.9-7-12}
\expandafter\def\csname sga@target@IV.9.7.2\endcsname{ega4.c04.lem.9-7-2}
\expandafter\def\csname sga@target@IV.9.7.3\endcsname{ega4.c04.subsec.9-7-3}
\expandafter\def\csname sga@target@IV.9.7.4\endcsname{ega4.c04.lem.9-7-4}
\expandafter\def\csname sga@target@IV.9.7.5\endcsname{ega4.c04.lem.9-7-5}
\expandafter\def\csname sga@target@IV.9.7.5.1\endcsname{ega4.c04.lem.9-7-5-1}
\expandafter\def\csname sga@target@IV.9.7.6\endcsname{ega4.c04.prop.9-7-6}
\expandafter\def\csname sga@target@IV.9.7.7\endcsname{ega4.c04.thm.9-7-7}
\expandafter\def\csname sga@target@IV.9.7.8\endcsname{ega4.c04.prop.9-7-8}
\expandafter\def\csname sga@target@IV.9.7.9\endcsname{ega4.c04.cor.9-7-9}
\expandafter\def\csname sga@target@IV.9.8\endcsname{ega4.c04.subsec.9-8.key-iv-9-8}
\expandafter\def\csname sga@target@IV.9.8.1\endcsname{ega4.c04.prop.9-8-1}
\expandafter\def\csname sga@target@IV.9.8.1-proof\endcsname{ega4.c04.subsec.9-8-1-proof}
\expandafter\def\csname sga@target@IV.9.8.2\endcsname{ega4.c04.subsec.9-8-2}
\expandafter\def\csname sga@target@IV.9.8.3\endcsname{ega4.c04.prop.9-8-3}
\expandafter\def\csname sga@target@IV.9.8.3-proof\endcsname{ega4.c04.subsec.9-8-3-proof}
\expandafter\def\csname sga@target@IV.9.8.4\endcsname{ega4.c04.cor.9-8-4}
\expandafter\def\csname sga@target@IV.9.8.4-proof\endcsname{ega4.c04.item.9-8-4-proof}
\expandafter\def\csname sga@target@IV.9.8.5\endcsname{ega4.c04.prop.9-8-5}
\expandafter\def\csname sga@target@IV.9.8.5-proof\endcsname{ega4.c04.item.9-8-5-proof}
\expandafter\def\csname sga@target@IV.9.8.6\endcsname{ega4.c04.prop.9-8-6}
\expandafter\def\csname sga@target@IV.9.8.6-proof\endcsname{ega4.c04.item.9-8-6-proof}
\expandafter\def\csname sga@target@IV.9.8.7\endcsname{ega4.c04.cor.9-8-7}
\expandafter\def\csname sga@target@IV.9.8.7-proof\endcsname{ega4.c04.item.9-8-7-proof}
\expandafter\def\csname sga@target@IV.9.8.8\endcsname{ega4.c04.prop.9-8-8}
\expandafter\def\csname sga@target@IV.9.8.8-proof\endcsname{ega4.c04.item.9-8-8-proof}
\expandafter\def\csname sga@target@IV.9.8.9\endcsname{ega4.c04.rem.9-8-9}
\expandafter\def\csname sga@target@IV.9.8.9.1\endcsname{ega4.c04.item.9-8-9-1}
\expandafter\def\csname sga@target@IV.9.9\endcsname{ega4.c04.subsec.9-9.key-iv-9-9}
\expandafter\def\csname sga@target@IV.9.9.1\endcsname{ega4.c04.prop.9-9-1}
\expandafter\def\csname sga@target@IV.9.9.1-proof\endcsname{ega4.c04.item.9-9-1-proof}
\expandafter\def\csname sga@target@IV.9.9.1.1\endcsname{ega4.c04.eq.9-9-1-1}
\expandafter\def\csname sga@target@IV.9.9.2\endcsname{ega4.c04.prop.9-9-2}
\expandafter\def\csname sga@target@IV.9.9.2-proof\endcsname{ega4.c04.item.9-9-2-proof}
\expandafter\def\csname sga@target@IV.9.9.2.1\endcsname{ega4.c04.eq.9-9-2-1}
\expandafter\def\csname sga@target@IV.9.9.2.2\endcsname{ega4.c04.eq.9-9-2-2}
\expandafter\def\csname sga@target@IV.9.9.3\endcsname{ega4.c04.cor.9-9-3}
\expandafter\def\csname sga@target@IV.9.9.3-proof\endcsname{ega4.c04.statement.9-9-3-proof}
\expandafter\def\csname sga@target@IV.9.9.4\endcsname{ega4.c04.prop.9-9-4}
\expandafter\def\csname sga@target@IV.9.9.4-proof\endcsname{ega4.c04.item.9-9-4-proof}
\expandafter\def\csname sga@target@IV.9.9.5\endcsname{ega4.c04.cor.9-9-5}
\expandafter\def\csname sga@target@IV.9.9.5-proof\endcsname{ega4.c04.item.9-9-5-proof}
\expandafter\def\csname sga@target@IV.9.9.6\endcsname{ega4.c04.prop.9-9-6}
\expandafter\def\csname sga@target@IV.9.9.6-proof\endcsname{ega4.c04.item.9-9-6-proof}
\expandafter\def\csname sga@target@section:IV.1\endcsname{ega4.c04.sec.1}
\expandafter\def\csname sga@target@section:IV.10\endcsname{ega4.c04.sec.10}
\expandafter\def\csname sga@target@section:IV.11\endcsname{ega4.c04.sec.11}
\expandafter\def\csname sga@target@section:IV.12\endcsname{ega4.c04.sec.12}
\expandafter\def\csname sga@target@section:IV.13\endcsname{ega4.c04.sec.13}
\expandafter\def\csname sga@target@section:IV.14\endcsname{ega4.c04.sec.14}
\expandafter\def\csname sga@target@section:IV.15\endcsname{ega4.c04.sec.15}
\expandafter\def\csname sga@target@section:IV.16\endcsname{ega4.c04.sec.16}
\expandafter\def\csname sga@target@section:IV.17\endcsname{ega4.c04.sec.17}
\expandafter\def\csname sga@target@section:IV.18\endcsname{ega4.c04.sec.18}
\expandafter\def\csname sga@target@section:IV.19\endcsname{ega4.c04.sec.19}
\expandafter\def\csname sga@target@section:IV.2\endcsname{ega4.c04.sec.2}
\expandafter\def\csname sga@target@section:IV.21\endcsname{ega4.c04.sec.21}
\expandafter\def\csname sga@target@section:IV.3\endcsname{ega4.c04.sec.3}
\expandafter\def\csname sga@target@section:IV.4\endcsname{ega4.c04.sec.4}
\expandafter\def\csname sga@target@section:IV.5\endcsname{ega4.c04.sec.5}
\expandafter\def\csname sga@target@section:IV.6\endcsname{ega4.c04.sec.6}
\expandafter\def\csname sga@target@section:IV.7\endcsname{ega4.c04.sec.7}
\expandafter\def\csname sga@target@section:IV.8\endcsname{ega4.c04.sec.8}
\expandafter\def\csname sga@target@section:IV.9\endcsname{ega4.c04.sec.9}
\expandafter\def\csname sga@target@subsection:IV.1.1\endcsname{ega4.c04.subsec.1-1}
\expandafter\def\csname sga@target@subsection:IV.1.10\endcsname{ega4.c04.subsec.1-10}
\expandafter\def\csname sga@target@subsection:IV.1.2\endcsname{ega4.c04.subsec.1-2}
\expandafter\def\csname sga@target@subsection:IV.1.3\endcsname{ega4.c04.subsec.1-3}
\expandafter\def\csname sga@target@subsection:IV.1.4\endcsname{ega4.c04.subsec.1-4}
\expandafter\def\csname sga@target@subsection:IV.1.5\endcsname{ega4.c04.subsec.1-5}
\expandafter\def\csname sga@target@subsection:IV.1.6\endcsname{ega4.c04.subsec.1-6}
\expandafter\def\csname sga@target@subsection:IV.1.7\endcsname{ega4.c04.subsec.1-7}
\expandafter\def\csname sga@target@subsection:IV.1.8\endcsname{ega4.c04.subsec.1-8}
\expandafter\def\csname sga@target@subsection:IV.1.9\endcsname{ega4.c04.subsec.1-9}
\expandafter\def\csname sga@target@subsection:IV.10.1\endcsname{ega4.c04.subsec.10-1}
\expandafter\def\csname sga@target@subsection:IV.10.10\endcsname{ega4.c04.subsec.10-10}
\expandafter\def\csname sga@target@subsection:IV.10.2\endcsname{ega4.c04.subsec.10-2}
\expandafter\def\csname sga@target@subsection:IV.10.3\endcsname{ega4.c04.subsec.10-3}
\expandafter\def\csname sga@target@subsection:IV.10.4\endcsname{ega4.c04.subsec.10-4}
\expandafter\def\csname sga@target@subsection:IV.10.5\endcsname{ega4.c04.subsec.10-5}
\expandafter\def\csname sga@target@subsection:IV.10.6\endcsname{ega4.c04.subsec.10-6}
\expandafter\def\csname sga@target@subsection:IV.10.7\endcsname{ega4.c04.subsec.10-7}
\expandafter\def\csname sga@target@subsection:IV.10.8\endcsname{ega4.c04.subsec.10-8}
\expandafter\def\csname sga@target@subsection:IV.10.9\endcsname{ega4.c04.subsec.10-9}
\expandafter\def\csname sga@target@subsection:IV.11.1\endcsname{ega4.c04.subsec.11-1}
\expandafter\def\csname sga@target@subsection:IV.11.10\endcsname{ega4.c04.subsec.11-10}
\expandafter\def\csname sga@target@subsection:IV.11.2\endcsname{ega4.c04.subsec.11-2}
\expandafter\def\csname sga@target@subsection:IV.11.3\endcsname{ega4.c04.subsec.11-3}
\expandafter\def\csname sga@target@subsection:IV.11.4\endcsname{ega4.c04.subsec.11-4}
\expandafter\def\csname sga@target@subsection:IV.11.5\endcsname{ega4.c04.subsec.11-5}
\expandafter\def\csname sga@target@subsection:IV.11.6\endcsname{ega4.c04.subsec.11-6}
\expandafter\def\csname sga@target@subsection:IV.11.7\endcsname{ega4.c04.subsec.11-7}
\expandafter\def\csname sga@target@subsection:IV.11.8\endcsname{ega4.c04.subsec.11-8}
\expandafter\def\csname sga@target@subsection:IV.11.9\endcsname{ega4.c04.subsec.11-9}
\expandafter\def\csname sga@target@subsection:IV.12.0\endcsname{ega4.c04.subsec.12-0}
\expandafter\def\csname sga@target@subsection:IV.12.1\endcsname{ega4.c04.subsec.12-1}
\expandafter\def\csname sga@target@subsection:IV.12.2\endcsname{ega4.c04.subsec.12-2}
\expandafter\def\csname sga@target@subsection:IV.12.3\endcsname{ega4.c04.subsec.12-3}
\expandafter\def\csname sga@target@subsection:IV.13.1\endcsname{ega4.c04.subsec.13-1}
\expandafter\def\csname sga@target@subsection:IV.13.2\endcsname{ega4.c04.subsec.13-2}
\expandafter\def\csname sga@target@subsection:IV.13.3\endcsname{ega4.c04.subsec.13-3}
\expandafter\def\csname sga@target@subsection:IV.14.1\endcsname{ega4.c04.subsec.14-1}
\expandafter\def\csname sga@target@subsection:IV.14.2\endcsname{ega4.c04.subsec.14-2}
\expandafter\def\csname sga@target@subsection:IV.14.3\endcsname{ega4.c04.subsec.14-3}
\expandafter\def\csname sga@target@subsection:IV.14.4\endcsname{ega4.c04.subsec.14-4}
\expandafter\def\csname sga@target@subsection:IV.14.5\endcsname{ega4.c04.subsec.14-5}
\expandafter\def\csname sga@target@subsection:IV.15.1\endcsname{ega4.c04.subsec.15-1}
\expandafter\def\csname sga@target@subsection:IV.15.2\endcsname{ega4.c04.subsec.15-2}
\expandafter\def\csname sga@target@subsection:IV.15.3\endcsname{ega4.c04.subsec.15-3}
\expandafter\def\csname sga@target@subsection:IV.15.4\endcsname{ega4.c04.subsec.15-4}
\expandafter\def\csname sga@target@subsection:IV.15.5\endcsname{ega4.c04.subsec.15-5}
\expandafter\def\csname sga@target@subsection:IV.15.6\endcsname{ega4.c04.subsec.15-6}
\expandafter\def\csname sga@target@subsection:IV.15.7\endcsname{ega4.c04.subsec.15-7}
\expandafter\def\csname sga@target@subsection:IV.17.1\endcsname{ega4.c04.subsec.17-1}
\expandafter\def\csname sga@target@subsection:IV.17.10\endcsname{ega4.c04.subsec.17-10}
\expandafter\def\csname sga@target@subsection:IV.17.11\endcsname{ega4.c04.subsec.17-11}
\expandafter\def\csname sga@target@subsection:IV.17.12\endcsname{ega4.c04.subsec.17-12}
\expandafter\def\csname sga@target@subsection:IV.17.13\endcsname{ega4.c04.subsec.17-13}
\expandafter\def\csname sga@target@subsection:IV.17.14\endcsname{ega4.c04.subsec.17-14}
\expandafter\def\csname sga@target@subsection:IV.17.15\endcsname{ega4.c04.subsec.17-15}
\expandafter\def\csname sga@target@subsection:IV.17.16\endcsname{ega4.c04.subsec.17-16}
\expandafter\def\csname sga@target@subsection:IV.17.2\endcsname{ega4.c04.subsec.17-2}
\expandafter\def\csname sga@target@subsection:IV.17.3\endcsname{ega4.c04.subsec.17-3}
\expandafter\def\csname sga@target@subsection:IV.17.4\endcsname{ega4.c04.subsec.17-4}
\expandafter\def\csname sga@target@subsection:IV.17.5\endcsname{ega4.c04.subsec.17-5}
\expandafter\def\csname sga@target@subsection:IV.17.6\endcsname{ega4.c04.subsec.17-6}
\expandafter\def\csname sga@target@subsection:IV.17.7\endcsname{ega4.c04.subsec.17-7}
\expandafter\def\csname sga@target@subsection:IV.17.8\endcsname{ega4.c04.subsec.17-8}
\expandafter\def\csname sga@target@subsection:IV.17.9\endcsname{ega4.c04.subsec.17-9}
\expandafter\def\csname sga@target@subsection:IV.18.1\endcsname{ega4.c04.subsec.18-1}
\expandafter\def\csname sga@target@subsection:IV.18.10\endcsname{ega4.c04.subsec.18-10}
\expandafter\def\csname sga@target@subsection:IV.18.11\endcsname{ega4.c04.subsec.18-11}
\expandafter\def\csname sga@target@subsection:IV.18.12\endcsname{ega4.c04.subsec.18-12}
\expandafter\def\csname sga@target@subsection:IV.18.2\endcsname{ega4.c04.subsec.18-2}
\expandafter\def\csname sga@target@subsection:IV.18.3\endcsname{ega4.c04.subsec.18-3}
\expandafter\def\csname sga@target@subsection:IV.18.4\endcsname{ega4.c04.subsec.18-4}
\expandafter\def\csname sga@target@subsection:IV.18.5\endcsname{ega4.c04.subsec.18-5}
\expandafter\def\csname sga@target@subsection:IV.18.6\endcsname{ega4.c04.subsec.18-6}
\expandafter\def\csname sga@target@subsection:IV.18.7\endcsname{ega4.c04.subsec.18-7}
\expandafter\def\csname sga@target@subsection:IV.18.8\endcsname{ega4.c04.subsec.18-8}
\expandafter\def\csname sga@target@subsection:IV.18.9\endcsname{ega4.c04.subsec.18-9}
\expandafter\def\csname sga@target@subsection:IV.19.1\endcsname{ega4.c04.subsec.19-1}
\expandafter\def\csname sga@target@subsection:IV.19.2\endcsname{ega4.c04.subsec.19-2}
\expandafter\def\csname sga@target@subsection:IV.19.3\endcsname{ega4.c04.subsec.19-3}
\expandafter\def\csname sga@target@subsection:IV.19.4\endcsname{ega4.c04.subsec.19-4}
\expandafter\def\csname sga@target@subsection:IV.19.5\endcsname{ega4.c04.subsec.19-5}
\expandafter\def\csname sga@target@subsection:IV.19.6\endcsname{ega4.c04.subsec.19-6}
\expandafter\def\csname sga@target@subsection:IV.19.7\endcsname{ega4.c04.subsec.19-7}
\expandafter\def\csname sga@target@subsection:IV.19.8\endcsname{ega4.c04.subsec.19-8}
\expandafter\def\csname sga@target@subsection:IV.19.9\endcsname{ega4.c04.subsec.19-9}
\expandafter\def\csname sga@target@subsection:IV.2.1\endcsname{ega4.c04.subsec.2-1}
\expandafter\def\csname sga@target@subsection:IV.2.2\endcsname{ega4.c04.subsec.2-2}
\expandafter\def\csname sga@target@subsection:IV.2.3\endcsname{ega4.c04.subsec.2-3}
\expandafter\def\csname sga@target@subsection:IV.2.4\endcsname{ega4.c04.subsec.2-4}
\expandafter\def\csname sga@target@subsection:IV.2.5\endcsname{ega4.c04.subsec.2-5}
\expandafter\def\csname sga@target@subsection:IV.2.6\endcsname{ega4.c04.subsec.2-6}
\expandafter\def\csname sga@target@subsection:IV.2.7\endcsname{ega4.c04.subsec.2-7}
\expandafter\def\csname sga@target@subsection:IV.2.8\endcsname{ega4.c04.subsec.2-8}
\expandafter\def\csname sga@target@subsection:IV.21.1\endcsname{ega4.c04.subsec.21-1}
\expandafter\def\csname sga@target@subsection:IV.21.10\endcsname{ega4.c04.subsec.21-10}
\expandafter\def\csname sga@target@subsection:IV.21.11\endcsname{ega4.c04.subsec.21-11}
\expandafter\def\csname sga@target@subsection:IV.21.12\endcsname{ega4.c04.subsec.21-12}
\expandafter\def\csname sga@target@subsection:IV.21.13\endcsname{ega4.c04.subsec.21-13}
\expandafter\def\csname sga@target@subsection:IV.21.14\endcsname{ega4.c04.subsec.21-14}
\expandafter\def\csname sga@target@subsection:IV.21.15\endcsname{ega4.c04.subsec.21-15}
\expandafter\def\csname sga@target@subsection:IV.21.2\endcsname{ega4.c04.subsec.21-2}
\expandafter\def\csname sga@target@subsection:IV.21.3\endcsname{ega4.c04.subsec.21-3}
\expandafter\def\csname sga@target@subsection:IV.21.4\endcsname{ega4.c04.subsec.21-4}
\expandafter\def\csname sga@target@subsection:IV.21.5\endcsname{ega4.c04.subsec.21-5}
\expandafter\def\csname sga@target@subsection:IV.21.6\endcsname{ega4.c04.subsec.21-6}
\expandafter\def\csname sga@target@subsection:IV.21.7\endcsname{ega4.c04.subsec.21-7}
\expandafter\def\csname sga@target@subsection:IV.21.8\endcsname{ega4.c04.subsec.21-8}
\expandafter\def\csname sga@target@subsection:IV.21.9\endcsname{ega4.c04.subsec.21-9}
\expandafter\def\csname sga@target@subsection:IV.3.1\endcsname{ega4.c04.subsec.3-1}
\expandafter\def\csname sga@target@subsection:IV.3.2\endcsname{ega4.c04.subsec.3-2}
\expandafter\def\csname sga@target@subsection:IV.3.3\endcsname{ega4.c04.subsec.3-3}
\expandafter\def\csname sga@target@subsection:IV.3.4\endcsname{ega4.c04.subsec.3-4}
\expandafter\def\csname sga@target@subsection:IV.4.1\endcsname{ega4.c04.subsec.4-1}
\expandafter\def\csname sga@target@subsection:IV.4.2\endcsname{ega4.c04.subsec.4-2}
\expandafter\def\csname sga@target@subsection:IV.4.3\endcsname{ega4.c04.subsec.4-3}
\expandafter\def\csname sga@target@subsection:IV.4.4\endcsname{ega4.c04.subsec.4-4}
\expandafter\def\csname sga@target@subsection:IV.4.5\endcsname{ega4.c04.subsec.4-5}
\expandafter\def\csname sga@target@subsection:IV.4.6\endcsname{ega4.c04.subsec.4-6}
\expandafter\def\csname sga@target@subsection:IV.4.7\endcsname{ega4.c04.subsec.4-7}
\expandafter\def\csname sga@target@subsection:IV.4.8\endcsname{ega4.c04.subsec.4-8}
\expandafter\def\csname sga@target@subsection:IV.4.9\endcsname{ega4.c04.subsec.4-9}
\expandafter\def\csname sga@target@subsection:IV.5.1\endcsname{ega4.c04.subsec.5-1}
\expandafter\def\csname sga@target@subsection:IV.5.10\endcsname{ega4.c04.subsec.5-10}
\expandafter\def\csname sga@target@subsection:IV.5.11\endcsname{ega4.c04.subsec.5-11}
\expandafter\def\csname sga@target@subsection:IV.5.12\endcsname{ega4.c04.subsec.5-12}
\expandafter\def\csname sga@target@subsection:IV.5.13\endcsname{ega4.c04.subsec.5-13}
\expandafter\def\csname sga@target@subsection:IV.5.2\endcsname{ega4.c04.subsec.5-2}
\expandafter\def\csname sga@target@subsection:IV.5.3\endcsname{ega4.c04.subsec.5-3}
\expandafter\def\csname sga@target@subsection:IV.5.4\endcsname{ega4.c04.subsec.5-4}
\expandafter\def\csname sga@target@subsection:IV.5.5\endcsname{ega4.c04.subsec.5-5}
\expandafter\def\csname sga@target@subsection:IV.5.6\endcsname{ega4.c04.subsec.5-6}
\expandafter\def\csname sga@target@subsection:IV.5.7\endcsname{ega4.c04.subsec.5-7}
\expandafter\def\csname sga@target@subsection:IV.5.8\endcsname{ega4.c04.subsec.5-8}
\expandafter\def\csname sga@target@subsection:IV.5.9\endcsname{ega4.c04.subsec.5-9}
\expandafter\def\csname sga@target@subsection:IV.6.1\endcsname{ega4.c04.subsec.6-1}
\expandafter\def\csname sga@target@subsection:IV.6.10\endcsname{ega4.c04.subsec.6-10}
\expandafter\def\csname sga@target@subsection:IV.6.11\endcsname{ega4.c04.subsec.6-11}
\expandafter\def\csname sga@target@subsection:IV.6.12\endcsname{ega4.c04.subsec.6-12}
\expandafter\def\csname sga@target@subsection:IV.6.13\endcsname{ega4.c04.subsec.6-13}
\expandafter\def\csname sga@target@subsection:IV.6.14\endcsname{ega4.c04.subsec.6-14}
\expandafter\def\csname sga@target@subsection:IV.6.15\endcsname{ega4.c04.subsec.6-15}
\expandafter\def\csname sga@target@subsection:IV.6.2\endcsname{ega4.c04.subsec.6-2}
\expandafter\def\csname sga@target@subsection:IV.6.3\endcsname{ega4.c04.subsec.6-3}
\expandafter\def\csname sga@target@subsection:IV.6.4\endcsname{ega4.c04.subsec.6-4}
\expandafter\def\csname sga@target@subsection:IV.6.5\endcsname{ega4.c04.subsec.6-5}
\expandafter\def\csname sga@target@subsection:IV.6.6\endcsname{ega4.c04.subsec.6-6}
\expandafter\def\csname sga@target@subsection:IV.6.7\endcsname{ega4.c04.subsec.6-7}
\expandafter\def\csname sga@target@subsection:IV.6.8\endcsname{ega4.c04.subsec.6-8}
\expandafter\def\csname sga@target@subsection:IV.6.9\endcsname{ega4.c04.subsec.6-9}
\expandafter\def\csname sga@target@subsection:IV.7.1\endcsname{ega4.c04.subsec.7-1}
\expandafter\def\csname sga@target@subsection:IV.7.2\endcsname{ega4.c04.subsec.7-2}
\expandafter\def\csname sga@target@subsection:IV.7.3\endcsname{ega4.c04.subsec.7-3}
\expandafter\def\csname sga@target@subsection:IV.7.4\endcsname{ega4.c04.subsec.7-4}
\expandafter\def\csname sga@target@subsection:IV.7.5\endcsname{ega4.c04.subsec.7-5}
\expandafter\def\csname sga@target@subsection:IV.7.6\endcsname{ega4.c04.subsec.7-6}
\expandafter\def\csname sga@target@subsection:IV.7.7\endcsname{ega4.c04.subsec.7-7}
\expandafter\def\csname sga@target@subsection:IV.7.8\endcsname{ega4.c04.subsec.7-8}
\expandafter\def\csname sga@target@subsection:IV.7.9\endcsname{ega4.c04.subsec.7-9}
\expandafter\def\csname sga@target@subsection:IV.8.1\endcsname{ega4.c04.subsec.8-1}
\expandafter\def\csname sga@target@subsection:IV.8.10\endcsname{ega4.c04.subsec.8-10}
\expandafter\def\csname sga@target@subsection:IV.8.11\endcsname{ega4.c04.subsec.8-11}
\expandafter\def\csname sga@target@subsection:IV.8.12\endcsname{ega4.c04.subsec.8-12}
\expandafter\def\csname sga@target@subsection:IV.8.13\endcsname{ega4.c04.subsec.8-13}
\expandafter\def\csname sga@target@subsection:IV.8.14\endcsname{ega4.c04.subsec.8-14}
\expandafter\def\csname sga@target@subsection:IV.8.2\endcsname{ega4.c04.subsec.8-2}
\expandafter\def\csname sga@target@subsection:IV.8.3\endcsname{ega4.c04.subsec.8-3}
\expandafter\def\csname sga@target@subsection:IV.8.4\endcsname{ega4.c04.subsec.8-4}
\expandafter\def\csname sga@target@subsection:IV.8.5\endcsname{ega4.c04.subsec.8-5}
\expandafter\def\csname sga@target@subsection:IV.8.6\endcsname{ega4.c04.subsec.8-6}
\expandafter\def\csname sga@target@subsection:IV.8.7\endcsname{ega4.c04.subsec.8-7}
\expandafter\def\csname sga@target@subsection:IV.8.8\endcsname{ega4.c04.subsec.8-8}
\expandafter\def\csname sga@target@subsection:IV.8.9\endcsname{ega4.c04.subsec.8-9}
\expandafter\def\csname sga@target@subsection:IV.9.1\endcsname{ega4.c04.subsec.9-1}
\expandafter\def\csname sga@target@subsection:IV.9.2\endcsname{ega4.c04.subsec.9-2}
\expandafter\def\csname sga@target@subsection:IV.9.3\endcsname{ega4.c04.subsec.9-3}
\expandafter\def\csname sga@target@subsection:IV.9.4\endcsname{ega4.c04.subsec.9-4}
\expandafter\def\csname sga@target@subsection:IV.9.5\endcsname{ega4.c04.subsec.9-5}
\expandafter\def\csname sga@target@subsection:IV.9.6\endcsname{ega4.c04.subsec.9-6}
\expandafter\def\csname sga@target@subsection:IV.9.7\endcsname{ega4.c04.subsec.9-7}
\expandafter\def\csname sga@target@subsection:IV.9.8\endcsname{ega4.c04.subsec.9-8}
\expandafter\def\csname sga@target@subsection:IV.9.9\endcsname{ega4.c04.subsec.9-9}
\expandafter\def\csname sga@bib@I-4\endcsname{ega4.reader.bib.i-4}
\expandafter\def\csname sga@bib@IV-32\endcsname{ega4.reader.bib.iv-32}
\expandafter\def\csname sga@bib@IV-34\endcsname{ega4.reader.bib.iv-34}
\expandafter\def\csname sga@bib@IV-37\endcsname{ega4.reader.bib.iv-37}
\expandafter\def\csname sga@bib@IV-41\endcsname{ega4.reader.bib.iv-41}
\expandafter\def\csname sga@bib@IV-42\endcsname{ega4.reader.bib.iv-42}
\expandafter\def\csname sga@bib@IV-43\endcsname{ega4.reader.bib.iv-43}
\expandafter\def\csname sga@bib@IV-44\endcsname{ega4.reader.bib.iv-44}
\expandafter\def\csname sga@bib@SGA1\endcsname{ega4.reader.bib.sga1}
\let\sgaOriginalLabel\label
\renewcommand{\label}[1]{%
  \sgaOriginalLabel{#1}%
  \ifcsname sga@target@#1\endcsname
    \expandafter\sgaStableTarget\expandafter{\csname sga@target@#1\endcsname}%
  \fi
}
\let\sgaOriginalBibitem\bibitem
\newcommand{\sgaBibTarget}[1]{%
  \ifcsname sga@bib@#1\endcsname
    \expandafter\sgaStableTarget\expandafter{\csname sga@bib@#1\endcsname}%
  \fi
}
\newcommand{\sgaBibitemPlain}[1]{%
  \sgaOriginalBibitem{#1}\sgaBibTarget{#1}}
\newcommand{\sgaBibitemOptional}[2][]{%
  \sgaOriginalBibitem[#1]{#2}\sgaBibTarget{#2}}
\renewcommand{\bibitem}{%
  \@ifnextchar[\sgaBibitemOptional\sgaBibitemPlain}
\makeatother
% END NATIVE INPUT EGA4_reference_v2_registry
% BEGIN NATIVE INPUT EGA4_reference_v2_links
% Generated stable-link dispatch for the EGA IV reader.
% Stable destinations are emitted by EGA4_reference_v2_registry.tex.
\makeatletter
\expandafter\def\csname sga@target@IV.1.1.1\endcsname{ega4.c04.def.1-1-1}
\expandafter\def\csname sga@target@IV.1.1.2\endcsname{ega4.c04.prop.1-1-2}
\expandafter\def\csname sga@target@IV.1.1.3\endcsname{ega4.c04.prop.1-1-3}
\expandafter\def\csname sga@target@IV.1.1.4\endcsname{ega4.c04.cor.1-1-4}
\expandafter\def\csname sga@target@IV.1.1.5\endcsname{ega4.c04.prop.1-1-5}
\expandafter\def\csname sga@target@IV.1.1.6\endcsname{ega4.c04.prop.1-1-6}
\expandafter\def\csname sga@target@IV.1.10.1\endcsname{ega4.c04.thm.1-10-1}
\expandafter\def\csname sga@target@IV.1.10.2\endcsname{ega4.c04.cor.1-10-2}
\expandafter\def\csname sga@target@IV.1.10.3\endcsname{ega4.c04.prop.1-10-3}
\expandafter\def\csname sga@target@IV.1.10.4\endcsname{ega4.c04.cor.1-10-4}
\expandafter\def\csname sga@target@IV.1.2.1\endcsname{ega4.c04.def.1-2-1}
\expandafter\def\csname sga@target@IV.1.2.2\endcsname{ega4.c04.prop.1-2-2}
\expandafter\def\csname sga@target@IV.1.2.3\endcsname{ega4.c04.cor.1-2-3}
\expandafter\def\csname sga@target@IV.1.2.4\endcsname{ega4.c04.prop.1-2-4}
\expandafter\def\csname sga@target@IV.1.2.4.1\endcsname{ega4.c04.eq.1-2-4-1}
\expandafter\def\csname sga@target@IV.1.2.5\endcsname{ega4.c04.prop.1-2-5}
\expandafter\def\csname sga@target@IV.1.2.6\endcsname{ega4.c04.prop.1-2-6}
\expandafter\def\csname sga@target@IV.1.2.7\endcsname{ega4.c04.prop.1-2-7}
\expandafter\def\csname sga@target@IV.1.2.8\endcsname{ega4.c04.cor.1-2-8}
\expandafter\def\csname sga@target@IV.1.2.9\endcsname{ega4.c04.prop.1-2-9}
\expandafter\def\csname sga@target@IV.1.3.1\endcsname{ega4.c04.statement.1-3-1}
\expandafter\def\csname sga@target@IV.1.3.10\endcsname{ega4.c04.statement.1-3-10}
\expandafter\def\csname sga@target@IV.1.3.11\endcsname{ega4.c04.prop.1-3-11}
\expandafter\def\csname sga@target@IV.1.3.2\endcsname{ega4.c04.def.1-3-2}
\expandafter\def\csname sga@target@IV.1.3.3\endcsname{ega4.c04.prop.1-3-3}
\expandafter\def\csname sga@target@IV.1.3.4\endcsname{ega4.c04.prop.1-3-4}
\expandafter\def\csname sga@target@IV.1.3.5\endcsname{ega4.c04.cor.1-3-5}
\expandafter\def\csname sga@target@IV.1.3.6\endcsname{ega4.c04.prop.1-3-6}
\expandafter\def\csname sga@target@IV.1.3.7\endcsname{ega4.c04.prop.1-3-7}
\expandafter\def\csname sga@target@IV.1.3.8\endcsname{ega4.c04.statement.1-3-8}
\expandafter\def\csname sga@target@IV.1.3.9\endcsname{ega4.c04.prop.1-3-9}
\expandafter\def\csname sga@target@IV.1.4.1\endcsname{ega4.c04.statement.1-4-1}
\expandafter\def\csname sga@target@IV.1.4.2\endcsname{ega4.c04.def.1-4-2}
\expandafter\def\csname sga@target@IV.1.4.3\endcsname{ega4.c04.prop.1-4-3}
\expandafter\def\csname sga@target@IV.1.4.3.1\endcsname{ega4.c04.cor.1-4-3-1}
\expandafter\def\csname sga@target@IV.1.4.4\endcsname{ega4.c04.prop.1-4-4}
\expandafter\def\csname sga@target@IV.1.4.4.1\endcsname{ega4.c04.lem.1-4-4-1}
\expandafter\def\csname sga@target@IV.1.4.5\endcsname{ega4.c04.cor.1-4-5}
\expandafter\def\csname sga@target@IV.1.4.6\endcsname{ega4.c04.prop.1-4-6}
\expandafter\def\csname sga@target@IV.1.4.7\endcsname{ega4.c04.prop.1-4-7}
\expandafter\def\csname sga@target@IV.1.4.7.1\endcsname{ega4.c04.lem.1-4-7-1}
\expandafter\def\csname sga@target@IV.1.5.1\endcsname{ega4.c04.prop.1-5-1}
\expandafter\def\csname sga@target@IV.1.5.2\endcsname{ega4.c04.def.1-5-2}
\expandafter\def\csname sga@target@IV.1.5.3\endcsname{ega4.c04.prop.1-5-3}
\expandafter\def\csname sga@target@IV.1.5.4\endcsname{ega4.c04.prop.1-5-4}
\expandafter\def\csname sga@target@IV.1.5.5\endcsname{ega4.c04.cor.1-5-5}
\expandafter\def\csname sga@target@IV.1.5.6\endcsname{ega4.c04.cor.1-5-6}
\expandafter\def\csname sga@target@IV.1.5.7\endcsname{ega4.c04.prop.1-5-7}
\expandafter\def\csname sga@target@IV.1.6.1\endcsname{ega4.c04.def.1-6-1}
\expandafter\def\csname sga@target@IV.1.6.2\endcsname{ega4.c04.prop.1-6-2}
\expandafter\def\csname sga@target@IV.1.6.3\endcsname{ega4.c04.cor.1-6-3}
\expandafter\def\csname sga@target@IV.1.6.4\endcsname{ega4.c04.rem.1-6-4}
\expandafter\def\csname sga@target@IV.1.6.5\endcsname{ega4.c04.prop.1-6-5}
\expandafter\def\csname sga@target@IV.1.7.1\endcsname{ega4.c04.statement.1-7-1}
\expandafter\def\csname sga@target@IV.1.7.10\endcsname{ega4.c04.statement.1-7-10}
\expandafter\def\csname sga@target@IV.1.7.11\endcsname{ega4.c04.statement.1-7-11}
\expandafter\def\csname sga@target@IV.1.7.12\endcsname{ega4.c04.statement.1-7-12}
\expandafter\def\csname sga@target@IV.1.7.13\endcsname{ega4.c04.statement.1-7-13}
\expandafter\def\csname sga@target@IV.1.7.14\endcsname{ega4.c04.statement.1-7-14}
\expandafter\def\csname sga@target@IV.1.7.15\endcsname{ega4.c04.statement.1-7-15}
\expandafter\def\csname sga@target@IV.1.7.16\endcsname{ega4.c04.statement.1-7-16}
\expandafter\def\csname sga@target@IV.1.7.17\endcsname{ega4.c04.statement.1-7-17}
\expandafter\def\csname sga@target@IV.1.7.18\endcsname{ega4.c04.statement.1-7-18}
\expandafter\def\csname sga@target@IV.1.7.19\endcsname{ega4.c04.statement.1-7-19}
\expandafter\def\csname sga@target@IV.1.7.2\endcsname{ega4.c04.statement.1-7-2}
\expandafter\def\csname sga@target@IV.1.7.20\endcsname{ega4.c04.statement.1-7-20}
\expandafter\def\csname sga@target@IV.1.7.21\endcsname{ega4.c04.statement.1-7-21}
\expandafter\def\csname sga@target@IV.1.7.22\endcsname{ega4.c04.statement.1-7-22}
\expandafter\def\csname sga@target@IV.1.7.3\endcsname{ega4.c04.statement.1-7-3}
\expandafter\def\csname sga@target@IV.1.7.4\endcsname{ega4.c04.statement.1-7-4}
\expandafter\def\csname sga@target@IV.1.7.5\endcsname{ega4.c04.statement.1-7-5}
\expandafter\def\csname sga@target@IV.1.7.6\endcsname{ega4.c04.statement.1-7-6}
\expandafter\def\csname sga@target@IV.1.7.7\endcsname{ega4.c04.statement.1-7-7}
\expandafter\def\csname sga@target@IV.1.7.8\endcsname{ega4.c04.statement.1-7-8}
\expandafter\def\csname sga@target@IV.1.7.9\endcsname{ega4.c04.statement.1-7-9}
\expandafter\def\csname sga@target@IV.1.8.1\endcsname{ega4.c04.statement.1-8-1}
\expandafter\def\csname sga@target@IV.1.8.2\endcsname{ega4.c04.prop.1-8-2}
\expandafter\def\csname sga@target@IV.1.8.3\endcsname{ega4.c04.lem.1-8-3}
\expandafter\def\csname sga@target@IV.1.8.4\endcsname{ega4.c04.thm.1-8-4}
\expandafter\def\csname sga@target@IV.1.8.4.1\endcsname{ega4.c04.lem.1-8-4-1}
\expandafter\def\csname sga@target@IV.1.8.4.2\endcsname{ega4.c04.lem.1-8-4-2}
\expandafter\def\csname sga@target@IV.1.8.5\endcsname{ega4.c04.cor.1-8-5}
\expandafter\def\csname sga@target@IV.1.8.5.1\endcsname{ega4.c04.lem.1-8-5-1}
\expandafter\def\csname sga@target@IV.1.8.6\endcsname{ega4.c04.cor.1-8-6}
\expandafter\def\csname sga@target@IV.1.8.7\endcsname{ega4.c04.prop.1-8-7}
\expandafter\def\csname sga@target@IV.1.8.7.1\endcsname{ega4.c04.lem.1-8-7-1}
\expandafter\def\csname sga@target@IV.1.9.1\endcsname{ega4.c04.lem.1-9-1}
\expandafter\def\csname sga@target@IV.1.9.10\endcsname{ega4.c04.prop.1-9-10}
\expandafter\def\csname sga@target@IV.1.9.11\endcsname{ega4.c04.prop.1-9-11}
\expandafter\def\csname sga@target@IV.1.9.12\endcsname{ega4.c04.cor.1-9-12}
\expandafter\def\csname sga@target@IV.1.9.13\endcsname{ega4.c04.statement.1-9-13}
\expandafter\def\csname sga@target@IV.1.9.14\endcsname{ega4.c04.prop.1-9-14}
\expandafter\def\csname sga@target@IV.1.9.15\endcsname{ega4.c04.prop.1-9-15}
\expandafter\def\csname sga@target@IV.1.9.16\endcsname{ega4.c04.statement.1-9-16}
\expandafter\def\csname sga@target@IV.1.9.2\endcsname{ega4.c04.prop.1-9-2}
\expandafter\def\csname sga@target@IV.1.9.2.1\endcsname{ega4.c04.statement.1-9-2-1}
\expandafter\def\csname sga@target@IV.1.9.2.i\endcsname{ega4.c04.statement.1-9-2-i}
\expandafter\def\csname sga@target@IV.1.9.2.ii\endcsname{ega4.c04.statement.1-9-2-ii}
\expandafter\def\csname sga@target@IV.1.9.3\endcsname{ega4.c04.prop.1-9-3}
\expandafter\def\csname sga@target@IV.1.9.3.1\endcsname{ega4.c04.lem.1-9-3-1}
\expandafter\def\csname sga@target@IV.1.9.3.2\endcsname{ega4.c04.lem.1-9-3-2}
\expandafter\def\csname sga@target@IV.1.9.4\endcsname{ega4.c04.def.1-9-4}
\expandafter\def\csname sga@target@IV.1.9.5\endcsname{ega4.c04.prop.1-9-5}
\expandafter\def\csname sga@target@IV.1.9.6\endcsname{ega4.c04.statement.1-9-6}
\expandafter\def\csname sga@target@IV.1.9.7\endcsname{ega4.c04.thm.1-9-7}
\expandafter\def\csname sga@target@IV.1.9.8\endcsname{ega4.c04.cor.1-9-8}
\expandafter\def\csname sga@target@IV.1.9.9\endcsname{ega4.c04.cor.1-9-9}
\expandafter\def\csname sga@target@IV.10.1.1\endcsname{ega4.c04.statement.10-1-1}
\expandafter\def\csname sga@target@IV.10.1.2\endcsname{ega4.c04.prop.10-1-2}
\expandafter\def\csname sga@target@IV.10.1.3\endcsname{ega4.c04.def.10-1-3}
\expandafter\def\csname sga@target@IV.10.1.4\endcsname{ega4.c04.cor.10-1-4}
\expandafter\def\csname sga@target@IV.10.10.1\endcsname{ega4.c04.statement.10-10-1}
\expandafter\def\csname sga@target@IV.10.10.2\endcsname{ega4.c04.statement.10-10-2}
\expandafter\def\csname sga@target@IV.10.10.3\endcsname{ega4.c04.statement.10-10-3}
\expandafter\def\csname sga@target@IV.10.10.4\endcsname{ega4.c04.statement.10-10-4}
\expandafter\def\csname sga@target@IV.10.10.5\endcsname{ega4.c04.statement.10-10-5}
\expandafter\def\csname sga@target@IV.10.2.1\endcsname{ega4.c04.prop.10-2-1}
\expandafter\def\csname sga@target@IV.10.2.2\endcsname{ega4.c04.def.10-2-2}
\expandafter\def\csname sga@target@IV.10.2.3\endcsname{ega4.c04.cor.10-2-3}
\expandafter\def\csname sga@target@IV.10.2.4\endcsname{ega4.c04.cor.10-2-4}
\expandafter\def\csname sga@target@IV.10.2.5\endcsname{ega4.c04.cor.10-2-5}
\expandafter\def\csname sga@target@IV.10.2.6\endcsname{ega4.c04.cor.10-2-6}
\expandafter\def\csname sga@target@IV.10.2.7\endcsname{ega4.c04.prop.10-2-7}
\expandafter\def\csname sga@target@IV.10.2.8\endcsname{ega4.c04.item.10-2-8}
\expandafter\def\csname sga@target@IV.10.2.8.1\endcsname{ega4.c04.eq.10-2-8-1}
\expandafter\def\csname sga@target@IV.10.3.1\endcsname{ega4.c04.def.10-3-1}
\expandafter\def\csname sga@target@IV.10.3.2\endcsname{ega4.c04.prop.10-3-2}
\expandafter\def\csname sga@target@IV.10.3.3\endcsname{ega4.c04.prop.10-3-3}
\expandafter\def\csname sga@target@IV.10.4.1\endcsname{ega4.c04.def.10-4-1}
\expandafter\def\csname sga@target@IV.10.4.10\endcsname{ega4.c04.statement.10-4-10}
\expandafter\def\csname sga@target@IV.10.4.11\endcsname{ega4.c04.statement.10-4-11}
\expandafter\def\csname sga@target@IV.10.4.11.1\endcsname{ega4.c04.lem.10-4-11-1}
\expandafter\def\csname sga@target@IV.10.4.12\endcsname{ega4.c04.cor.10-4-12}
\expandafter\def\csname sga@target@IV.10.4.12.1\endcsname{ega4.c04.sec.10-4-12-1}
\expandafter\def\csname sga@target@IV.10.4.12.1.errata\endcsname{ega4.c04.eq.10-4-12-1-errata}
\expandafter\def\csname sga@target@IV.10.4.13\endcsname{ega4.c04.cor.10-4-13}
\expandafter\def\csname sga@target@IV.10.4.14\endcsname{ega4.c04.rem.10-4-14}
\expandafter\def\csname sga@target@IV.10.4.2\endcsname{ega4.c04.prop.10-4-2}
\expandafter\def\csname sga@target@IV.10.4.3\endcsname{ega4.c04.statement.10-4-3}
\expandafter\def\csname sga@target@IV.10.4.4\endcsname{ega4.c04.prop.10-4-4}
\expandafter\def\csname sga@target@IV.10.4.5\endcsname{ega4.c04.prop.10-4-5}
\expandafter\def\csname sga@target@IV.10.4.5.1\endcsname{ega4.c04.lem.10-4-5-1}
\expandafter\def\csname sga@target@IV.10.4.6\endcsname{ega4.c04.cor.10-4-6}
\expandafter\def\csname sga@target@IV.10.4.7\endcsname{ega4.c04.cor.10-4-7}
\expandafter\def\csname sga@target@IV.10.4.8\endcsname{ega4.c04.cor.10-4-8}
\expandafter\def\csname sga@target@IV.10.4.9\endcsname{ega4.c04.statement.10-4-9}
\expandafter\def\csname sga@target@IV.10.5.1\endcsname{ega4.c04.prop.10-5-1}
\expandafter\def\csname sga@target@IV.10.5.10\endcsname{ega4.c04.prop.10-5-10}
\expandafter\def\csname sga@target@IV.10.5.2\endcsname{ega4.c04.cor.10-5-2}
\expandafter\def\csname sga@target@IV.10.5.3\endcsname{ega4.c04.cor.10-5-3}
\expandafter\def\csname sga@target@IV.10.5.4\endcsname{ega4.c04.rem.10-5-4}
\expandafter\def\csname sga@target@IV.10.5.5\endcsname{ega4.c04.prop.10-5-5}
\expandafter\def\csname sga@target@IV.10.5.6\endcsname{ega4.c04.rem.10-5-6}
\expandafter\def\csname sga@target@IV.10.5.7\endcsname{ega4.c04.prop.10-5-7}
\expandafter\def\csname sga@target@IV.10.5.8\endcsname{ega4.c04.cor.10-5-8}
\expandafter\def\csname sga@target@IV.10.5.9\endcsname{ega4.c04.cor.10-5-9}
\expandafter\def\csname sga@target@IV.10.6.1\endcsname{ega4.c04.prop.10-6-1}
\expandafter\def\csname sga@target@IV.10.6.1.1\endcsname{ega4.c04.eq.10-6-1-1}
\expandafter\def\csname sga@target@IV.10.6.1.2\endcsname{ega4.c04.eq.10-6-1-2}
\expandafter\def\csname sga@target@IV.10.6.1.3\endcsname{ega4.c04.eq.10-6-1-3}
\expandafter\def\csname sga@target@IV.10.6.1.4\endcsname{ega4.c04.eq.10-6-1-4}
\expandafter\def\csname sga@target@IV.10.6.2\endcsname{ega4.c04.cor.10-6-2}
\expandafter\def\csname sga@target@IV.10.6.3\endcsname{ega4.c04.prop.10-6-3}
\expandafter\def\csname sga@target@IV.10.6.3.1\endcsname{ega4.c04.eq.10-6-3-1}
\expandafter\def\csname sga@target@IV.10.6.3.2\endcsname{ega4.c04.eq.10-6-3-2}
\expandafter\def\csname sga@target@IV.10.6.3.3\endcsname{ega4.c04.eq.10-6-3-3}
\expandafter\def\csname sga@target@IV.10.6.4\endcsname{ega4.c04.cor.10-6-4}
\expandafter\def\csname sga@target@IV.10.6.4.1\endcsname{ega4.c04.eq.10-6-4-1}
\expandafter\def\csname sga@target@IV.10.7.1\endcsname{ega4.c04.statement.10-7-1}
\expandafter\def\csname sga@target@IV.10.7.2\endcsname{ega4.c04.statement.10-7-2}
\expandafter\def\csname sga@target@IV.10.7.3\endcsname{ega4.c04.statement.10-7-3}
\expandafter\def\csname sga@target@IV.10.8.1\endcsname{ega4.c04.def.10-8-1}
\expandafter\def\csname sga@target@IV.10.8.1.1\endcsname{ega4.c04.eq.10-8-1-1}
\expandafter\def\csname sga@target@IV.10.8.1.2\endcsname{ega4.c04.eq.10-8-1-2}
\expandafter\def\csname sga@target@IV.10.8.2\endcsname{ega4.c04.statement.10-8-2}
\expandafter\def\csname sga@target@IV.10.8.3\endcsname{ega4.c04.lem.10-8-3}
\expandafter\def\csname sga@target@IV.10.8.3.1\endcsname{ega4.c04.eq.10-8-3-1}
\expandafter\def\csname sga@target@IV.10.8.4\endcsname{ega4.c04.cor.10-8-4}
\expandafter\def\csname sga@target@IV.10.8.5\endcsname{ega4.c04.prop.10-8-5}
\expandafter\def\csname sga@target@IV.10.8.6\endcsname{ega4.c04.cor.10-8-6}
\expandafter\def\csname sga@target@IV.10.8.7\endcsname{ega4.c04.prop.10-8-7}
\expandafter\def\csname sga@target@IV.10.8.7.1\endcsname{ega4.c04.eq.10-8-7-1}
\expandafter\def\csname sga@target@IV.10.8.8\endcsname{ega4.c04.cor.10-8-8}
\expandafter\def\csname sga@target@IV.10.9.1\endcsname{ega4.c04.statement.10-9-1}
\expandafter\def\csname sga@target@IV.10.9.2\endcsname{ega4.c04.statement.10-9-2}
\expandafter\def\csname sga@target@IV.10.9.3\endcsname{ega4.c04.statement.10-9-3}
\expandafter\def\csname sga@target@IV.10.9.4\endcsname{ega4.c04.prop.10-9-4}
\expandafter\def\csname sga@target@IV.10.9.5\endcsname{ega4.c04.statement.10-9-5}
\expandafter\def\csname sga@target@IV.10.9.6\endcsname{ega4.c04.prop.10-9-6}
\expandafter\def\csname sga@target@IV.11.1.1\endcsname{ega4.c04.thm.11-1-1}
\expandafter\def\csname sga@target@IV.11.1.1.1\endcsname{ega4.c04.sec.11-1-1-1}
\expandafter\def\csname sga@target@IV.11.1.1.2\endcsname{ega4.c04.lem.11-1-1-2}
\expandafter\def\csname sga@target@IV.11.1.2\endcsname{ega4.c04.cor.11-1-2}
\expandafter\def\csname sga@target@IV.11.1.3\endcsname{ega4.c04.cor.11-1-3}
\expandafter\def\csname sga@target@IV.11.1.4\endcsname{ega4.c04.cor.11-1-4}
\expandafter\def\csname sga@target@IV.11.1.5\endcsname{ega4.c04.cor.11-1-5}
\expandafter\def\csname sga@target@IV.11.1.6\endcsname{ega4.c04.cor.11-1-6}
\expandafter\def\csname sga@target@IV.11.10.1\endcsname{ega4.c04.prop.11-10-1}
\expandafter\def\csname sga@target@IV.11.10.1.1\endcsname{ega4.c04.eq.11-10-1-1}
\expandafter\def\csname sga@target@IV.11.10.1.2\endcsname{ega4.c04.eq.11-10-1-2}
\expandafter\def\csname sga@target@IV.11.10.10\endcsname{ega4.c04.prop.11-10-10}
\expandafter\def\csname sga@target@IV.11.10.2\endcsname{ega4.c04.def.11-10-2}
\expandafter\def\csname sga@target@IV.11.10.3\endcsname{ega4.c04.item.11-10-3}
\expandafter\def\csname sga@target@IV.11.10.4\endcsname{ega4.c04.prop.11-10-4}
\expandafter\def\csname sga@target@IV.11.10.5\endcsname{ega4.c04.thm.11-10-5}
\expandafter\def\csname sga@target@IV.11.10.6\endcsname{ega4.c04.prop.11-10-6}
\expandafter\def\csname sga@target@IV.11.10.7\endcsname{ega4.c04.cor.11-10-7}
\expandafter\def\csname sga@target@IV.11.10.8\endcsname{ega4.c04.def.11-10-8}
\expandafter\def\csname sga@target@IV.11.10.9\endcsname{ega4.c04.thm.11-10-9}
\expandafter\def\csname sga@target@IV.11.2.1\endcsname{ega4.c04.statement.11-2-1}
\expandafter\def\csname sga@target@IV.11.2.1.1\endcsname{ega4.c04.eq.11-2-1-1}
\expandafter\def\csname sga@target@IV.11.2.1.2\endcsname{ega4.c04.eq.11-2-1-2}
\expandafter\def\csname sga@target@IV.11.2.1.3\endcsname{ega4.c04.eq.11-2-1-3}
\expandafter\def\csname sga@target@IV.11.2.1.4\endcsname{ega4.c04.eq.11-2-1-4}
\expandafter\def\csname sga@target@IV.11.2.1.5\endcsname{ega4.c04.eq.11-2-1-5}
\expandafter\def\csname sga@target@IV.11.2.10\endcsname{ega4.c04.rem.11-2-10}
\expandafter\def\csname sga@target@IV.11.2.2\endcsname{ega4.c04.statement.11-2-2}
\expandafter\def\csname sga@target@IV.11.2.2.1\endcsname{ega4.c04.eq.11-2-2-1}
\expandafter\def\csname sga@target@IV.11.2.2.2\endcsname{ega4.c04.lem.11-2-2-2}
\expandafter\def\csname sga@target@IV.11.2.3\endcsname{ega4.c04.statement.11-2-3}
\expandafter\def\csname sga@target@IV.11.2.3.1\endcsname{ega4.c04.eq.11-2-3-1}
\expandafter\def\csname sga@target@IV.11.2.3.2\endcsname{ega4.c04.eq.11-2-3-2}
\expandafter\def\csname sga@target@IV.11.2.3.3\endcsname{ega4.c04.eq.11-2-3-3}
\expandafter\def\csname sga@target@IV.11.2.4\endcsname{ega4.c04.lem.11-2-4}
\expandafter\def\csname sga@target@IV.11.2.5\endcsname{ega4.c04.lem.11-2-5}
\expandafter\def\csname sga@target@IV.11.2.6\endcsname{ega4.c04.thm.11-2-6}
\expandafter\def\csname sga@target@IV.11.2.6.1\endcsname{ega4.c04.cor.11-2-6-1}
\expandafter\def\csname sga@target@IV.11.2.7\endcsname{ega4.c04.cor.11-2-7}
\expandafter\def\csname sga@target@IV.11.2.8\endcsname{ega4.c04.prop.11-2-8}
\expandafter\def\csname sga@target@IV.11.2.9\endcsname{ega4.c04.thm.11-2-9}
\expandafter\def\csname sga@target@IV.11.2.9.1\endcsname{ega4.c04.statement.11-2-9-1}
\expandafter\def\csname sga@target@IV.11.2.9.2\endcsname{ega4.c04.lem.11-2-9-2}
\expandafter\def\csname sga@target@IV.11.2.9.3\endcsname{ega4.c04.lem.11-2-9-3}
\expandafter\def\csname sga@target@IV.11.2.9.4\endcsname{ega4.c04.cor.11-2-9-4}
\expandafter\def\csname sga@target@IV.11.2.9.5\endcsname{ega4.c04.lem.11-2-9-5}
\expandafter\def\csname sga@target@IV.11.2.9.6\endcsname{ega4.c04.lem.11-2-9-6}
\expandafter\def\csname sga@target@IV.11.2.9.7\endcsname{ega4.c04.eq.11-2-9-7}
\expandafter\def\csname sga@target@IV.11.2.9.8\endcsname{ega4.c04.eq.11-2-9-8}
\expandafter\def\csname sga@target@IV.11.2.9.9\endcsname{ega4.c04.eq.11-2-9-9}
\expandafter\def\csname sga@target@IV.11.3\endcsname{ega4.c04.subsec.11-3.key-iv-11-3}
\expandafter\def\csname sga@target@IV.11.3.1\endcsname{ega4.c04.thm.11-3-1}
\expandafter\def\csname sga@target@IV.11.3.10\endcsname{ega4.c04.thm.11-3-10}
\expandafter\def\csname sga@target@IV.11.3.10.1\endcsname{ega4.c04.lem.11-3-10-1}
\expandafter\def\csname sga@target@IV.11.3.10.2\endcsname{ega4.c04.lem.11-3-10-2}
\expandafter\def\csname sga@target@IV.11.3.10.3\endcsname{ega4.c04.eq.11-3-10-3}
\expandafter\def\csname sga@target@IV.11.3.10.4\endcsname{ega4.c04.sec.11-3-10-4}
\expandafter\def\csname sga@target@IV.11.3.11\endcsname{ega4.c04.cor.11-3-11}
\expandafter\def\csname sga@target@IV.11.3.12\endcsname{ega4.c04.rem.11-3-12}
\expandafter\def\csname sga@target@IV.11.3.13\endcsname{ega4.c04.prop.11-3-13}
\expandafter\def\csname sga@target@IV.11.3.14\endcsname{ega4.c04.cor.11-3-14}
\expandafter\def\csname sga@target@IV.11.3.15\endcsname{ega4.c04.prop.11-3-15}
\expandafter\def\csname sga@target@IV.11.3.16\endcsname{ega4.c04.prop.11-3-16}
\expandafter\def\csname sga@target@IV.11.3.17\endcsname{ega4.c04.cor.11-3-17}
\expandafter\def\csname sga@target@IV.11.3.2\endcsname{ega4.c04.cor.11-3-2}
\expandafter\def\csname sga@target@IV.11.3.3\endcsname{ega4.c04.cor.11-3-3}
\expandafter\def\csname sga@target@IV.11.3.4\endcsname{ega4.c04.prop.11-3-4}
\expandafter\def\csname sga@target@IV.11.3.5\endcsname{ega4.c04.cor.11-3-5}
\expandafter\def\csname sga@target@IV.11.3.6\endcsname{ega4.c04.prop.11-3-6}
\expandafter\def\csname sga@target@IV.11.3.7\endcsname{ega4.c04.prop.11-3-7}
\expandafter\def\csname sga@target@IV.11.3.8\endcsname{ega4.c04.thm.11-3-8}
\expandafter\def\csname sga@target@IV.11.3.9\endcsname{ega4.c04.prop.11-3-9}
\expandafter\def\csname sga@target@IV.11.3.9.1\endcsname{ega4.c04.lem.11-3-9-1}
\expandafter\def\csname sga@target@IV.11.4.1\endcsname{ega4.c04.thm.11-4-1}
\expandafter\def\csname sga@target@IV.11.4.1.1\endcsname{ega4.c04.lem.11-4-1-1}
\expandafter\def\csname sga@target@IV.11.4.1.2\endcsname{ega4.c04.lem.11-4-1-2}
\expandafter\def\csname sga@target@IV.11.4.10\endcsname{ega4.c04.cor.11-4-10}
\expandafter\def\csname sga@target@IV.11.4.11\endcsname{ega4.c04.prop.11-4-11}
\expandafter\def\csname sga@target@IV.11.4.12\endcsname{ega4.c04.prop.11-4-12}
\expandafter\def\csname sga@target@IV.11.4.13\endcsname{ega4.c04.cor.11-4-13}
\expandafter\def\csname sga@target@IV.11.4.2\endcsname{ega4.c04.rem.11-4-2}
\expandafter\def\csname sga@target@IV.11.4.3\endcsname{ega4.c04.cor.11-4-3}
\expandafter\def\csname sga@target@IV.11.4.4\endcsname{ega4.c04.cor.11-4-4}
\expandafter\def\csname sga@target@IV.11.4.5\endcsname{ega4.c04.prop.11-4-5}
\expandafter\def\csname sga@target@IV.11.4.6\endcsname{ega4.c04.cor.11-4-6}
\expandafter\def\csname sga@target@IV.11.4.7\endcsname{ega4.c04.prop.11-4-7}
\expandafter\def\csname sga@target@IV.11.4.8\endcsname{ega4.c04.rem.11-4-8}
\expandafter\def\csname sga@target@IV.11.4.9\endcsname{ega4.c04.cor.11-4-9}
\expandafter\def\csname sga@target@IV.11.5.1\endcsname{ega4.c04.thm.11-5-1}
\expandafter\def\csname sga@target@IV.11.5.2\endcsname{ega4.c04.cor.11-5-2}
\expandafter\def\csname sga@target@IV.11.5.3\endcsname{ega4.c04.prop.11-5-3}
\expandafter\def\csname sga@target@IV.11.5.4\endcsname{ega4.c04.cor.11-5-4}
\expandafter\def\csname sga@target@IV.11.5.5\endcsname{ega4.c04.prop.11-5-5}
\expandafter\def\csname sga@target@IV.11.5.5.1\endcsname{ega4.c04.lem.11-5-5-1}
\expandafter\def\csname sga@target@IV.11.5.5.2\endcsname{ega4.c04.lem.11-5-5-2}
\expandafter\def\csname sga@target@IV.11.6.1\endcsname{ega4.c04.thm.11-6-1}
\expandafter\def\csname sga@target@IV.11.6.2\endcsname{ega4.c04.cor.11-6-2}
\expandafter\def\csname sga@target@IV.11.7.1\endcsname{ega4.c04.statement.11-7-1}
\expandafter\def\csname sga@target@IV.11.7.2\endcsname{ega4.c04.statement.11-7-2}
\expandafter\def\csname sga@target@IV.11.7.3\endcsname{ega4.c04.statement.11-7-3}
\expandafter\def\csname sga@target@IV.11.7.4\endcsname{ega4.c04.statement.11-7-4}
\expandafter\def\csname sga@target@IV.11.7.5\endcsname{ega4.c04.statement.11-7-5}
\expandafter\def\csname sga@target@IV.11.8.1\endcsname{ega4.c04.thm.11-8-1}
\expandafter\def\csname sga@target@IV.11.8.2\endcsname{ega4.c04.subsec.11-8-2}
\expandafter\def\csname sga@target@IV.11.9.1\endcsname{ega4.c04.statement.11-9-1}
\expandafter\def\csname sga@target@IV.11.9.1.1\endcsname{ega4.c04.statement.11-9-1-1}
\expandafter\def\csname sga@target@IV.11.9.10\endcsname{ega4.c04.thm.11-9-10}
\expandafter\def\csname sga@target@IV.11.9.11\endcsname{ega4.c04.rem.11-9-11}
\expandafter\def\csname sga@target@IV.11.9.12\endcsname{ega4.c04.lem.11-9-12}
\expandafter\def\csname sga@target@IV.11.9.13\endcsname{ega4.c04.prop.11-9-13}
\expandafter\def\csname sga@target@IV.11.9.14\endcsname{ega4.c04.statement.11-9-14}
\expandafter\def\csname sga@target@IV.11.9.15\endcsname{ega4.c04.prop.11-9-15}
\expandafter\def\csname sga@target@IV.11.9.16\endcsname{ega4.c04.thm.11-9-16}
\expandafter\def\csname sga@target@IV.11.9.17\endcsname{ega4.c04.prop.11-9-17}
\expandafter\def\csname sga@target@IV.11.9.17.1\endcsname{ega4.c04.lem.11-9-17-1}
\expandafter\def\csname sga@target@IV.11.9.18\endcsname{ega4.c04.item.11-9-18}
\expandafter\def\csname sga@target@IV.11.9.2\endcsname{ega4.c04.statement.11-9-2}
\expandafter\def\csname sga@target@IV.11.9.3\endcsname{ega4.c04.statement.11-9-3}
\expandafter\def\csname sga@target@IV.11.9.4\endcsname{ega4.c04.statement.11-9-4}
\expandafter\def\csname sga@target@IV.11.9.5\endcsname{ega4.c04.statement.11-9-5}
\expandafter\def\csname sga@target@IV.11.9.6\endcsname{ega4.c04.statement.11-9-6}
\expandafter\def\csname sga@target@IV.11.9.7\endcsname{ega4.c04.prop.11-9-7}
\expandafter\def\csname sga@target@IV.11.9.7.1\endcsname{ega4.c04.lem.11-9-7-1}
\expandafter\def\csname sga@target@IV.11.9.8\endcsname{ega4.c04.subsec.11-9-8}
\expandafter\def\csname sga@target@IV.11.9.9\endcsname{ega4.c04.cor.11-9-9}
\expandafter\def\csname sga@target@IV.12.0.1\endcsname{ega4.c04.statement.12-0-1}
\expandafter\def\csname sga@target@IV.12.0.2\endcsname{ega4.c04.statement.12-0-2}
\expandafter\def\csname sga@target@IV.12.1.1\endcsname{ega4.c04.thm.12-1-1}
\expandafter\def\csname sga@target@IV.12.1.1.1\endcsname{ega4.c04.statement.12-1-1-1}
\expandafter\def\csname sga@target@IV.12.1.1.2\endcsname{ega4.c04.statement.12-1-1-2}
\expandafter\def\csname sga@target@IV.12.1.1.3\endcsname{ega4.c04.statement.12-1-1-3}
\expandafter\def\csname sga@target@IV.12.1.1.4\endcsname{ega4.c04.statement.12-1-1-4}
\expandafter\def\csname sga@target@IV.12.1.1.5\endcsname{ega4.c04.prop.12-1-1-5}
\expandafter\def\csname sga@target@IV.12.1.2\endcsname{ega4.c04.item.12-1-2}
\expandafter\def\csname sga@target@IV.12.1.3\endcsname{ega4.c04.cor.12-1-3}
\expandafter\def\csname sga@target@IV.12.1.4\endcsname{ega4.c04.rem.12-1-4}
\expandafter\def\csname sga@target@IV.12.1.5\endcsname{ega4.c04.cor.12-1-5}
\expandafter\def\csname sga@target@IV.12.1.6\endcsname{ega4.c04.thm.12-1-6}
\expandafter\def\csname sga@target@IV.12.1.7\endcsname{ega4.c04.cor.12-1-7}
\expandafter\def\csname sga@target@IV.12.1.8\endcsname{ega4.c04.item.12-1-8}
\expandafter\def\csname sga@target@IV.12.2.1\endcsname{ega4.c04.thm.12-2-1}
\expandafter\def\csname sga@target@IV.12.2.1.1\endcsname{ega4.c04.statement.12-2-1-1}
\expandafter\def\csname sga@target@IV.12.2.1.2\endcsname{ega4.c04.statement.12-2-1-2}
\expandafter\def\csname sga@target@IV.12.2.1.3\endcsname{ega4.c04.statement.12-2-1-3}
\expandafter\def\csname sga@target@IV.12.2.1.4\endcsname{ega4.c04.lem.12-2-1-4}
\expandafter\def\csname sga@target@IV.12.2.2\endcsname{ega4.c04.cor.12-2-2}
\expandafter\def\csname sga@target@IV.12.2.3\endcsname{ega4.c04.item.12-2-3}
\expandafter\def\csname sga@target@IV.12.2.4\endcsname{ega4.c04.thm.12-2-4}
\expandafter\def\csname sga@target@IV.12.3.1\endcsname{ega4.c04.lem.12-3-1}
\expandafter\def\csname sga@target@IV.12.3.2\endcsname{ega4.c04.prop.12-3-2}
\expandafter\def\csname sga@target@IV.12.3.2.1\endcsname{ega4.c04.eq.12-3-2-1}
\expandafter\def\csname sga@target@IV.12.3.3\endcsname{ega4.c04.prop.12-3-3}
\expandafter\def\csname sga@target@IV.12.3.3.1\endcsname{ega4.c04.lem.12-3-3-1}
\expandafter\def\csname sga@target@IV.12.3.3.2\endcsname{ega4.c04.eq.12-3-3-2}
\expandafter\def\csname sga@target@IV.12.3.3.3\endcsname{ega4.c04.lem.12-3-3-3}
\expandafter\def\csname sga@target@IV.12.3.3.4\endcsname{ega4.c04.eq.12-3-3-4}
\expandafter\def\csname sga@target@IV.12.3.3.5\endcsname{ega4.c04.lem.12-3-3-5}
\expandafter\def\csname sga@target@IV.12.3.4\endcsname{ega4.c04.cor.12-3-4}
\expandafter\def\csname sga@target@IV.13.1.1\endcsname{ega4.c04.lem.13-1-1}
\expandafter\def\csname sga@target@IV.13.1.2\endcsname{ega4.c04.lem.13-1-2}
\expandafter\def\csname sga@target@IV.13.1.3\endcsname{ega4.c04.thm.13-1-3}
\expandafter\def\csname sga@target@IV.13.1.4\endcsname{ega4.c04.cor.13-1-4}
\expandafter\def\csname sga@target@IV.13.1.5\endcsname{ega4.c04.cor.13-1-5}
\expandafter\def\csname sga@target@IV.13.1.6\endcsname{ega4.c04.cor.13-1-6}
\expandafter\def\csname sga@target@IV.13.1.7\endcsname{ega4.c04.prop.13-1-7}
\expandafter\def\csname sga@target@IV.13.2.1\endcsname{ega4.c04.statement.13-2-1}
\expandafter\def\csname sga@target@IV.13.2.10\endcsname{ega4.c04.prop.13-2-10}
\expandafter\def\csname sga@target@IV.13.2.10.1\endcsname{ega4.c04.eq.13-2-10-1}
\expandafter\def\csname sga@target@IV.13.2.10.2\endcsname{ega4.c04.eq.13-2-10-2}
\expandafter\def\csname sga@target@IV.13.2.11\endcsname{ega4.c04.cor.13-2-11}
\expandafter\def\csname sga@target@IV.13.2.12\endcsname{ega4.c04.item.13-2-12}
\expandafter\def\csname sga@target@IV.13.2.2\endcsname{ega4.c04.def.13-2-2}
\expandafter\def\csname sga@target@IV.13.2.3\endcsname{ega4.c04.prop.13-2-3}
\expandafter\def\csname sga@target@IV.13.2.3.1\endcsname{ega4.c04.prop.13-2-3-1}
\expandafter\def\csname sga@target@IV.13.2.4\endcsname{ega4.c04.statement.13-2-4}
\expandafter\def\csname sga@target@IV.13.2.4.1\endcsname{ega4.c04.statement.13-2-4-1}
\expandafter\def\csname sga@target@IV.13.2.4.2\endcsname{ega4.c04.eq.13-2-4-2}
\expandafter\def\csname sga@target@IV.13.2.4.3\endcsname{ega4.c04.eq.13-2-4-3}
\expandafter\def\csname sga@target@IV.13.2.4.4\endcsname{ega4.c04.eq.13-2-4-4}
\expandafter\def\csname sga@target@IV.13.2.4.5\endcsname{ega4.c04.eq.13-2-4-5}
\expandafter\def\csname sga@target@IV.13.2.5\endcsname{ega4.c04.lem.13-2-5}
\expandafter\def\csname sga@target@IV.13.2.6\endcsname{ega4.c04.statement.13-2-6}
\expandafter\def\csname sga@target@IV.13.2.6.1\endcsname{ega4.c04.statement.13-2-6-1}
\expandafter\def\csname sga@target@IV.13.2.7\endcsname{ega4.c04.prop.13-2-7}
\expandafter\def\csname sga@target@IV.13.2.7.1\endcsname{ega4.c04.prop.13-2-7-1}
\expandafter\def\csname sga@target@IV.13.2.8\endcsname{ega4.c04.cor.13-2-8}
\expandafter\def\csname sga@target@IV.13.2.9\endcsname{ega4.c04.cor.13-2-9}
\expandafter\def\csname sga@target@IV.13.2.9.1\endcsname{ega4.c04.eq.13-2-9-1}
\expandafter\def\csname sga@target@IV.13.2.9.2\endcsname{ega4.c04.eq.13-2-9-2}
\expandafter\def\csname sga@target@IV.13.3.1\endcsname{ega4.c04.prop.13-3-1}
\expandafter\def\csname sga@target@IV.13.3.1.1\endcsname{ega4.c04.lem.13-3-1-1}
\expandafter\def\csname sga@target@IV.13.3.2\endcsname{ega4.c04.def.13-3-2}
\expandafter\def\csname sga@target@IV.13.3.3\endcsname{ega4.c04.prop.13-3-3}
\expandafter\def\csname sga@target@IV.13.3.4\endcsname{ega4.c04.cor.13-3-4}
\expandafter\def\csname sga@target@IV.13.3.4.1\endcsname{ega4.c04.cor.13-3-4-1}
\expandafter\def\csname sga@target@IV.13.3.5\endcsname{ega4.c04.cor.13-3-5}
\expandafter\def\csname sga@target@IV.13.3.6\endcsname{ega4.c04.prop.13-3-6}
\expandafter\def\csname sga@target@IV.13.3.6.1\endcsname{ega4.c04.prop.13-3-6-1}
\expandafter\def\csname sga@target@IV.13.3.6.2\endcsname{ega4.c04.eq.13-3-6-2}
\expandafter\def\csname sga@target@IV.13.3.7\endcsname{ega4.c04.prop.13-3-7}
\expandafter\def\csname sga@target@IV.13.3.7.1\endcsname{ega4.c04.prop.13-3-7-1}
\expandafter\def\csname sga@target@IV.13.3.8\endcsname{ega4.c04.prop.13-3-8}
\expandafter\def\csname sga@target@IV.13.3.9\endcsname{ega4.c04.statement.13-3-9}
\expandafter\def\csname sga@target@IV.14.1.1\endcsname{ega4.c04.statement.14-1-1}
\expandafter\def\csname sga@target@IV.14.1.2\endcsname{ega4.c04.prop.14-1-2}
\expandafter\def\csname sga@target@IV.14.1.3\endcsname{ega4.c04.item.14-1-3}
\expandafter\def\csname sga@target@IV.14.1.4\endcsname{ega4.c04.prop.14-1-4}
\expandafter\def\csname sga@target@IV.14.1.5\endcsname{ega4.c04.rem.14-1-5}
\expandafter\def\csname sga@target@IV.14.1.6\endcsname{ega4.c04.statement.14-1-6}
\expandafter\def\csname sga@target@IV.14.2.1\endcsname{ega4.c04.thm.14-2-1}
\expandafter\def\csname sga@target@IV.14.2.1.1\endcsname{ega4.c04.eq.14-2-1-1}
\expandafter\def\csname sga@target@IV.14.2.2\endcsname{ega4.c04.cor.14-2-2}
\expandafter\def\csname sga@target@IV.14.2.3\endcsname{ega4.c04.prop.14-2-3}
\expandafter\def\csname sga@target@IV.14.2.4\endcsname{ega4.c04.cor.14-2-4}
\expandafter\def\csname sga@target@IV.14.2.5\endcsname{ega4.c04.cor.14-2-5}
\expandafter\def\csname sga@target@IV.14.2.6\endcsname{ega4.c04.cor.14-2-6}
\expandafter\def\csname sga@target@IV.14.2.7\endcsname{ega4.c04.rem.14-2-7}
\expandafter\def\csname sga@target@IV.14.3\endcsname{ega4.c04.subsec.14-3.key-iv-14-3}
\expandafter\def\csname sga@target@IV.14.3.1\endcsname{ega4.c04.statement.14-3-1}
\expandafter\def\csname sga@target@IV.14.3.10\endcsname{ega4.c04.lem.14-3-10}
\expandafter\def\csname sga@target@IV.14.3.11\endcsname{ega4.c04.prop.14-3-11}
\expandafter\def\csname sga@target@IV.14.3.12\endcsname{ega4.c04.rem.14-3-12}
\expandafter\def\csname sga@target@IV.14.3.13\endcsname{ega4.c04.prop.14-3-13}
\expandafter\def\csname sga@target@IV.14.3.14\endcsname{ega4.c04.rem.14-3-14}
\expandafter\def\csname sga@target@IV.14.3.2\endcsname{ega4.c04.prop.14-3-2}
\expandafter\def\csname sga@target@IV.14.3.3\endcsname{ega4.c04.def.14-3-3}
\expandafter\def\csname sga@target@IV.14.3.3.1\endcsname{ega4.c04.subsec.14-3-3-1}
\expandafter\def\csname sga@target@IV.14.3.4\endcsname{ega4.c04.prop.14-3-4}
\expandafter\def\csname sga@target@IV.14.3.5\endcsname{ega4.c04.prop.14-3-5}
\expandafter\def\csname sga@target@IV.14.3.6\endcsname{ega4.c04.prop.14-3-6}
\expandafter\def\csname sga@target@IV.14.3.7\endcsname{ega4.c04.cor.14-3-7}
\expandafter\def\csname sga@target@IV.14.3.8\endcsname{ega4.c04.prop.14-3-8}
\expandafter\def\csname sga@target@IV.14.3.9\endcsname{ega4.c04.item.14-3-9}
\expandafter\def\csname sga@target@IV.14.4.1\endcsname{ega4.c04.thm.14-4-1}
\expandafter\def\csname sga@target@IV.14.4.1.1\endcsname{ega4.c04.lem.14-4-1-1}
\expandafter\def\csname sga@target@IV.14.4.1.2\endcsname{ega4.c04.lem.14-4-1-2}
\expandafter\def\csname sga@target@IV.14.4.1.3\endcsname{ega4.c04.lem.14-4-1-3}
\expandafter\def\csname sga@target@IV.14.4.10\endcsname{ega4.c04.item.14-4-10}
\expandafter\def\csname sga@target@IV.14.4.2\endcsname{ega4.c04.cor.14-4-2}
\expandafter\def\csname sga@target@IV.14.4.3\endcsname{ega4.c04.cor.14-4-3}
\expandafter\def\csname sga@target@IV.14.4.4\endcsname{ega4.c04.cor.14-4-4}
\expandafter\def\csname sga@target@IV.14.4.5\endcsname{ega4.c04.rem.14-4-5}
\expandafter\def\csname sga@target@IV.14.4.6\endcsname{ega4.c04.cor.14-4-6}
\expandafter\def\csname sga@target@IV.14.4.7\endcsname{ega4.c04.prop.14-4-7}
\expandafter\def\csname sga@target@IV.14.4.7.1\endcsname{ega4.c04.eq.14-4-7-1}
\expandafter\def\csname sga@target@IV.14.4.7.2\endcsname{ega4.c04.eq.14-4-7-2}
\expandafter\def\csname sga@target@IV.14.4.8\endcsname{ega4.c04.prop.14-4-8}
\expandafter\def\csname sga@target@IV.14.4.8.1\endcsname{ega4.c04.lem.14-4-8-1}
\expandafter\def\csname sga@target@IV.14.4.9\endcsname{ega4.c04.cor.14-4-9}
\expandafter\def\csname sga@target@IV.14.5\endcsname{ega4.c04.subsec.14-5.key-iv-14-5}
\expandafter\def\csname sga@target@IV.14.5.1\endcsname{ega4.c04.lem.14-5-1}
\expandafter\def\csname sga@target@IV.14.5.1.1\endcsname{ega4.c04.eq.14-5-1-1}
\expandafter\def\csname sga@target@IV.14.5.1.2\endcsname{ega4.c04.eq.14-5-1-2}
\expandafter\def\csname sga@target@IV.14.5.1.3\endcsname{ega4.c04.eq.14-5-1-3}
\expandafter\def\csname sga@target@IV.14.5.1.4\endcsname{ega4.c04.eq.14-5-1-4}
\expandafter\def\csname sga@target@IV.14.5.10\endcsname{ega4.c04.prop.14-5-10}
\expandafter\def\csname sga@target@IV.14.5.11\endcsname{ega4.c04.cor.14-5-11}
\expandafter\def\csname sga@target@IV.14.5.11.1\endcsname{ega4.c04.lem.14-5-11-1}
\expandafter\def\csname sga@target@IV.14.5.12\endcsname{ega4.c04.cor.14-5-12}
\expandafter\def\csname sga@target@IV.14.5.2\endcsname{ega4.c04.cor.14-5-2}
\expandafter\def\csname sga@target@IV.14.5.2.1\endcsname{ega4.c04.eq.14-5-2-1}
\expandafter\def\csname sga@target@IV.14.5.3\endcsname{ega4.c04.prop.14-5-3}
\expandafter\def\csname sga@target@IV.14.5.4\endcsname{ega4.c04.cor.14-5-4}
\expandafter\def\csname sga@target@IV.14.5.5\endcsname{ega4.c04.rem.14-5-5}
\expandafter\def\csname sga@target@IV.14.5.6\endcsname{ega4.c04.prop.14-5-6}
\expandafter\def\csname sga@target@IV.14.5.7\endcsname{ega4.c04.rem.14-5-7}
\expandafter\def\csname sga@target@IV.14.5.8\endcsname{ega4.c04.prop.14-5-8}
\expandafter\def\csname sga@target@IV.14.5.9\endcsname{ega4.c04.prop.14-5-9}
\expandafter\def\csname sga@target@IV.15.1.1\endcsname{ega4.c04.prop.15-1-1}
\expandafter\def\csname sga@target@IV.15.1.1-proof\endcsname{ega4.c04.statement.15-1-1-proof}
\expandafter\def\csname sga@target@IV.15.1.1.1\endcsname{ega4.c04.eq.15-1-1-1}
\expandafter\def\csname sga@target@IV.15.1.1.2\endcsname{ega4.c04.eq.15-1-1-2}
\expandafter\def\csname sga@target@IV.15.1.1.3\endcsname{ega4.c04.eq.15-1-1-3}
\expandafter\def\csname sga@target@IV.15.1.1.4\endcsname{ega4.c04.eq.15-1-1-4}
\expandafter\def\csname sga@target@IV.15.1.1.5\endcsname{ega4.c04.eq.15-1-1-5}
\expandafter\def\csname sga@target@IV.15.1.1.6\endcsname{ega4.c04.lem.15-1-1-6}
\expandafter\def\csname sga@target@IV.15.1.2\endcsname{ega4.c04.cor.15-1-2}
\expandafter\def\csname sga@target@IV.15.1.3\endcsname{ega4.c04.cor.15-1-3}
\expandafter\def\csname sga@target@IV.15.1.4\endcsname{ega4.c04.statement.15-1-4}
\expandafter\def\csname sga@target@IV.15.2.1\endcsname{ega4.c04.statement.15-2-1}
\expandafter\def\csname sga@target@IV.15.2.2\endcsname{ega4.c04.thm.15-2-2}
\expandafter\def\csname sga@target@IV.15.2.2-proof\endcsname{ega4.c04.item.15-2-2-proof}
\expandafter\def\csname sga@target@IV.15.2.2.1\endcsname{ega4.c04.lem.15-2-2-1}
\expandafter\def\csname sga@target@IV.15.2.3\endcsname{ega4.c04.cor.15-2-3}
\expandafter\def\csname sga@target@IV.15.2.4\endcsname{ega4.c04.item.15-2-4}
\expandafter\def\csname sga@target@IV.15.3.1\endcsname{ega4.c04.prop.15-3-1}
\expandafter\def\csname sga@target@IV.15.3.1-proof\endcsname{ega4.c04.item.15-3-1-proof}
\expandafter\def\csname sga@target@IV.15.3.2\endcsname{ega4.c04.prop.15-3-2}
\expandafter\def\csname sga@target@IV.15.3.2-proof\endcsname{ega4.c04.item.15-3-2-proof}
\expandafter\def\csname sga@target@IV.15.3.3\endcsname{ega4.c04.cor.15-3-3}
\expandafter\def\csname sga@target@IV.15.3.3-proof\endcsname{ega4.c04.item.15-3-3-proof}
\expandafter\def\csname sga@target@IV.15.3.4\endcsname{ega4.c04.item.15-3-4}
\expandafter\def\csname sga@target@IV.15.4.1\endcsname{ega4.c04.prop.15-4-1}
\expandafter\def\csname sga@target@IV.15.4.2\endcsname{ega4.c04.prop.15-4-2}
\expandafter\def\csname sga@target@IV.15.4.3\endcsname{ega4.c04.prop.15-4-3}
\expandafter\def\csname sga@target@IV.15.4.3-proof\endcsname{ega4.c04.item.15-4-3-proof}
\expandafter\def\csname sga@target@IV.15.5.1\endcsname{ega4.c04.prop.15-5-1}
\expandafter\def\csname sga@target@IV.15.5.10\endcsname{ega4.c04.item.15-5-10}
\expandafter\def\csname sga@target@IV.15.5.2\endcsname{ega4.c04.cor.15-5-2}
\expandafter\def\csname sga@target@IV.15.5.2.1\endcsname{ega4.c04.eq.15-5-2-1}
\expandafter\def\csname sga@target@IV.15.5.2.2\endcsname{ega4.c04.lem.15-5-2-2}
\expandafter\def\csname sga@target@IV.15.5.2.3\endcsname{ega4.c04.eq.15-5-2-3}
\expandafter\def\csname sga@target@IV.15.5.2.4\endcsname{ega4.c04.eq.15-5-2-4}
\expandafter\def\csname sga@target@IV.15.5.3\endcsname{ega4.c04.prop.15-5-3}
\expandafter\def\csname sga@target@IV.15.5.4\endcsname{ega4.c04.cor.15-5-4}
\expandafter\def\csname sga@target@IV.15.5.5\endcsname{ega4.c04.statement.15-5-5}
\expandafter\def\csname sga@target@IV.15.5.6\endcsname{ega4.c04.lem.15-5-6}
\expandafter\def\csname sga@target@IV.15.5.6.1\endcsname{ega4.c04.lem.15-5-6-1}
\expandafter\def\csname sga@target@IV.15.5.7\endcsname{ega4.c04.prop.15-5-7}
\expandafter\def\csname sga@target@IV.15.5.7.1\endcsname{ega4.c04.eq.15-5-7-1}
\expandafter\def\csname sga@target@IV.15.5.8\endcsname{ega4.c04.statement.15-5-8}
\expandafter\def\csname sga@target@IV.15.5.9\endcsname{ega4.c04.prop.15-5-9}
\expandafter\def\csname sga@target@IV.15.6.1\endcsname{ega4.c04.prop.15-6-1}
\expandafter\def\csname sga@target@IV.15.6.1.1\endcsname{ega4.c04.eq.15-6-1-1}
\expandafter\def\csname sga@target@IV.15.6.10\endcsname{ega4.c04.item.15-6-10}
\expandafter\def\csname sga@target@IV.15.6.2\endcsname{ega4.c04.cor.15-6-2}
\expandafter\def\csname sga@target@IV.15.6.3\endcsname{ega4.c04.prop.15-6-3}
\expandafter\def\csname sga@target@IV.15.6.4\endcsname{ega4.c04.prop.15-6-4}
\expandafter\def\csname sga@target@IV.15.6.5\endcsname{ega4.c04.cor.15-6-5}
\expandafter\def\csname sga@target@IV.15.6.6\endcsname{ega4.c04.prop.15-6-6}
\expandafter\def\csname sga@target@IV.15.6.7\endcsname{ega4.c04.cor.15-6-7}
\expandafter\def\csname sga@target@IV.15.6.8\endcsname{ega4.c04.prop.15-6-8}
\expandafter\def\csname sga@target@IV.15.6.9\endcsname{ega4.c04.cor.15-6-9}
\expandafter\def\csname sga@target@IV.15.7.1\endcsname{ega4.c04.def.15-7-1}
\expandafter\def\csname sga@target@IV.15.7.10\endcsname{ega4.c04.cor.15-7-10}
\expandafter\def\csname sga@target@IV.15.7.11\endcsname{ega4.c04.cor.15-7-11}
\expandafter\def\csname sga@target@IV.15.7.11.1\endcsname{ega4.c04.lem.15-7-11-1}
\expandafter\def\csname sga@target@IV.15.7.2\endcsname{ega4.c04.prop.15-7-2}
\expandafter\def\csname sga@target@IV.15.7.3\endcsname{ega4.c04.rem.15-7-3}
\expandafter\def\csname sga@target@IV.15.7.4\endcsname{ega4.c04.cor.15-7-4}
\expandafter\def\csname sga@target@IV.15.7.5\endcsname{ega4.c04.cor.15-7-5}
\expandafter\def\csname sga@target@IV.15.7.6\endcsname{ega4.c04.cor.15-7-6}
\expandafter\def\csname sga@target@IV.15.7.7\endcsname{ega4.c04.item.15-7-7}
\expandafter\def\csname sga@target@IV.15.7.8\endcsname{ega4.c04.statement.15-7-8}
\expandafter\def\csname sga@target@IV.15.7.9\endcsname{ega4.c04.prop.15-7-9}
\expandafter\def\csname sga@target@IV.16.1\endcsname{ega4.c04.subsec.16-1}
\expandafter\def\csname sga@target@IV.16.1.1\endcsname{ega4.c04.statement.16-1-1}
\expandafter\def\csname sga@target@IV.16.1.10\endcsname{ega4.c04.cor.16-1-10}
\expandafter\def\csname sga@target@IV.16.1.11\endcsname{ega4.c04.rem.16-1-11}
\expandafter\def\csname sga@target@IV.16.1.2\endcsname{ega4.c04.def.16-1-2}
\expandafter\def\csname sga@target@IV.16.1.2.1\endcsname{ega4.c04.def.16-1-2-1}
\expandafter\def\csname sga@target@IV.16.1.2.2\endcsname{ega4.c04.eq.16-1-2-2}
\expandafter\def\csname sga@target@IV.16.1.3\endcsname{ega4.c04.statement.16-1-3}
\expandafter\def\csname sga@target@IV.16.1.4\endcsname{ega4.c04.statement.16-1-4}
\expandafter\def\csname sga@target@IV.16.1.5\endcsname{ega4.c04.prop.16-1-5}
\expandafter\def\csname sga@target@IV.16.1.6\endcsname{ega4.c04.cor.16-1-6}
\expandafter\def\csname sga@target@IV.16.1.7\endcsname{ega4.c04.cor.16-1-7}
\expandafter\def\csname sga@target@IV.16.1.8\endcsname{ega4.c04.prop.16-1-8}
\expandafter\def\csname sga@target@IV.16.1.9\endcsname{ega4.c04.prop.16-1-9}
\expandafter\def\csname sga@target@IV.16.10\endcsname{ega4.c04.subsec.16-10}
\expandafter\def\csname sga@target@IV.16.10.1\endcsname{ega4.c04.def.16-10-1}
\expandafter\def\csname sga@target@IV.16.10.2\endcsname{ega4.c04.prop.16-10-2}
\expandafter\def\csname sga@target@IV.16.10.3\endcsname{ega4.c04.prop.16-10-3}
\expandafter\def\csname sga@target@IV.16.10.4\endcsname{ega4.c04.prop.16-10-4}
\expandafter\def\csname sga@target@IV.16.10.5\endcsname{ega4.c04.prop.16-10-5}
\expandafter\def\csname sga@target@IV.16.10.6\endcsname{ega4.c04.rem.16-10-6}
\expandafter\def\csname sga@target@IV.16.11\endcsname{ega4.c04.subsec.16-11}
\expandafter\def\csname sga@target@IV.16.11.1\endcsname{ega4.c04.statement.16-11-1}
\expandafter\def\csname sga@target@IV.16.11.1.1\endcsname{ega4.c04.statement.16-11-1-1}
\expandafter\def\csname sga@target@IV.16.11.1.2\endcsname{ega4.c04.statement.16-11-1-2}
\expandafter\def\csname sga@target@IV.16.11.1.3\endcsname{ega4.c04.eq.16-11-1-3}
\expandafter\def\csname sga@target@IV.16.11.1.4\endcsname{ega4.c04.eq.16-11-1-4}
\expandafter\def\csname sga@target@IV.16.11.1.5\endcsname{ega4.c04.eq.16-11-1-5}
\expandafter\def\csname sga@target@IV.16.11.1.6\endcsname{ega4.c04.eq.16-11-1-6}
\expandafter\def\csname sga@target@IV.16.11.1.7\endcsname{ega4.c04.eq.16-11-1-7}
\expandafter\def\csname sga@target@IV.16.11.2\endcsname{ega4.c04.thm.16-11-2}
\expandafter\def\csname sga@target@IV.16.11.2.1\endcsname{ega4.c04.eq.16-11-2-1}
\expandafter\def\csname sga@target@IV.16.11.2.2\endcsname{ega4.c04.eq.16-11-2-2}
\expandafter\def\csname sga@target@IV.16.11.2.3\endcsname{ega4.c04.eq.16-11-2-3}
\expandafter\def\csname sga@target@IV.16.11.3\endcsname{ega4.c04.statement.16-11-3}
\expandafter\def\csname sga@target@IV.16.12\endcsname{ega4.c04.subsec.16-12}
\expandafter\def\csname sga@target@IV.16.12.1\endcsname{ega4.c04.statement.16-12-1}
\expandafter\def\csname sga@target@IV.16.12.2\endcsname{ega4.c04.thm.16-12-2}
\expandafter\def\csname sga@target@IV.16.12.3\endcsname{ega4.c04.statement.16-12-3}
\expandafter\def\csname sga@target@IV.16.2\endcsname{ega4.c04.subsec.16-2}
\expandafter\def\csname sga@target@IV.16.2.1\endcsname{ega4.c04.statement.16-2-1}
\expandafter\def\csname sga@target@IV.16.2.1.1\endcsname{ega4.c04.statement.16-2-1-1}
\expandafter\def\csname sga@target@IV.16.2.1.2\endcsname{ega4.c04.eq.16-2-1-2}
\expandafter\def\csname sga@target@IV.16.2.1.3\endcsname{ega4.c04.eq.16-2-1-3}
\expandafter\def\csname sga@target@IV.16.2.1.4\endcsname{ega4.c04.eq.16-2-1-4}
\expandafter\def\csname sga@target@IV.16.2.2\endcsname{ega4.c04.prop.16-2-2}
\expandafter\def\csname sga@target@IV.16.2.3\endcsname{ega4.c04.cor.16-2-3}
\expandafter\def\csname sga@target@IV.16.2.3.1\endcsname{ega4.c04.eq.16-2-3-1}
\expandafter\def\csname sga@target@IV.16.2.3.2\endcsname{ega4.c04.eq.16-2-3-2}
\expandafter\def\csname sga@target@IV.16.2.4\endcsname{ega4.c04.cor.16-2-4}
\expandafter\def\csname sga@target@IV.16.2.5\endcsname{ega4.c04.statement.16-2-5}
\expandafter\def\csname sga@target@IV.16.2.6\endcsname{ega4.c04.statement.16-2-6}
\expandafter\def\csname sga@target@IV.16.2.6.1\endcsname{ega4.c04.statement.16-2-6-1}
\expandafter\def\csname sga@target@IV.16.2.6.2\endcsname{ega4.c04.eq.16-2-6-2}
\expandafter\def\csname sga@target@IV.16.2.7\endcsname{ega4.c04.prop.16-2-7}
\expandafter\def\csname sga@target@IV.16.2.7.1\endcsname{ega4.c04.prop.16-2-7-1}
\expandafter\def\csname sga@target@IV.16.3\endcsname{ega4.c04.subsec.16-3}
\expandafter\def\csname sga@target@IV.16.3.1\endcsname{ega4.c04.def.16-3-1}
\expandafter\def\csname sga@target@IV.16.3.1.1\endcsname{ega4.c04.eq.16-3-1-1}
\expandafter\def\csname sga@target@IV.16.3.2\endcsname{ega4.c04.statement.16-3-2}
\expandafter\def\csname sga@target@IV.16.3.3\endcsname{ega4.c04.statement.16-3-3}
\expandafter\def\csname sga@target@IV.16.3.3.1\endcsname{ega4.c04.statement.16-3-3-1}
\expandafter\def\csname sga@target@IV.16.3.4\endcsname{ega4.c04.prop.16-3-4}
\expandafter\def\csname sga@target@IV.16.3.5\endcsname{ega4.c04.statement.16-3-5}
\expandafter\def\csname sga@target@IV.16.3.6\endcsname{ega4.c04.def.16-3-6}
\expandafter\def\csname sga@target@IV.16.3.6.1\endcsname{ega4.c04.eq.16-3-6-1}
\expandafter\def\csname sga@target@IV.16.3.7\endcsname{ega4.c04.statement.16-3-7}
\expandafter\def\csname sga@target@IV.16.3.7.1\endcsname{ega4.c04.eq.16-3-7-1}
\expandafter\def\csname sga@target@IV.16.3.8\endcsname{ega4.c04.prop.16-3-8}
\expandafter\def\csname sga@target@IV.16.3.9\endcsname{ega4.c04.prop.16-3-9}
\expandafter\def\csname sga@target@IV.16.4\endcsname{ega4.c04.subsec.16-4}
\expandafter\def\csname sga@target@IV.16.4.1\endcsname{ega4.c04.statement.16-4-1}
\expandafter\def\csname sga@target@IV.16.4.1.1\endcsname{ega4.c04.statement.16-4-1-1}
\expandafter\def\csname sga@target@IV.16.4.1.2\endcsname{ega4.c04.statement.16-4-1-2}
\expandafter\def\csname sga@target@IV.16.4.1.3\endcsname{ega4.c04.eq.16-4-1-3}
\expandafter\def\csname sga@target@IV.16.4.1.4\endcsname{ega4.c04.eq.16-4-1-4}
\expandafter\def\csname sga@target@IV.16.4.1.5\endcsname{ega4.c04.eq.16-4-1-5}
\expandafter\def\csname sga@target@IV.16.4.10\endcsname{ega4.c04.cor.16-4-10}
\expandafter\def\csname sga@target@IV.16.4.11\endcsname{ega4.c04.prop.16-4-11}
\expandafter\def\csname sga@target@IV.16.4.11.1\endcsname{ega4.c04.prop.16-4-11-1}
\expandafter\def\csname sga@target@IV.16.4.11.2\endcsname{ega4.c04.prop.16-4-11-2}
\expandafter\def\csname sga@target@IV.16.4.11.3\endcsname{ega4.c04.eq.16-4-11-3}
\expandafter\def\csname sga@target@IV.16.4.11.4\endcsname{ega4.c04.eq.16-4-11-4}
\expandafter\def\csname sga@target@IV.16.4.12\endcsname{ega4.c04.cor.16-4-12}
\expandafter\def\csname sga@target@IV.16.4.13\endcsname{ega4.c04.cor.16-4-13}
\expandafter\def\csname sga@target@IV.16.4.14\endcsname{ega4.c04.prop.16-4-14}
\expandafter\def\csname sga@target@IV.16.4.14.1\endcsname{ega4.c04.prop.16-4-14-1}
\expandafter\def\csname sga@target@IV.16.4.15\endcsname{ega4.c04.cor.16-4-15}
\expandafter\def\csname sga@target@IV.16.4.15.1\endcsname{ega4.c04.cor.16-4-15-1}
\expandafter\def\csname sga@target@IV.16.4.15.2\endcsname{ega4.c04.eq.16-4-15-2}
\expandafter\def\csname sga@target@IV.16.4.16\endcsname{ega4.c04.cor.16-4-16}
\expandafter\def\csname sga@target@IV.16.4.16.1\endcsname{ega4.c04.cor.16-4-16-1}
\expandafter\def\csname sga@target@IV.16.4.16.2\endcsname{ega4.c04.eq.16-4-16-2}
\expandafter\def\csname sga@target@IV.16.4.17\endcsname{ega4.c04.cor.16-4-17}
\expandafter\def\csname sga@target@IV.16.4.18\endcsname{ega4.c04.statement.16-4-18}
\expandafter\def\csname sga@target@IV.16.4.18.1\endcsname{ega4.c04.statement.16-4-18-1}
\expandafter\def\csname sga@target@IV.16.4.18.2\endcsname{ega4.c04.statement.16-4-18-2}
\expandafter\def\csname sga@target@IV.16.4.18.3\endcsname{ega4.c04.eq.16-4-18-3}
\expandafter\def\csname sga@target@IV.16.4.19\endcsname{ega4.c04.cor.16-4-19}
\expandafter\def\csname sga@target@IV.16.4.19.1\endcsname{ega4.c04.cor.16-4-19-1}
\expandafter\def\csname sga@target@IV.16.4.2\endcsname{ega4.c04.statement.16-4-2}
\expandafter\def\csname sga@target@IV.16.4.20\endcsname{ega4.c04.prop.16-4-20}
\expandafter\def\csname sga@target@IV.16.4.21\endcsname{ega4.c04.cor.16-4-21}
\expandafter\def\csname sga@target@IV.16.4.21.1\endcsname{ega4.c04.cor.16-4-21-1}
\expandafter\def\csname sga@target@IV.16.4.22\endcsname{ega4.c04.cor.16-4-22}
\expandafter\def\csname sga@target@IV.16.4.23\endcsname{ega4.c04.prop.16-4-23}
\expandafter\def\csname sga@target@IV.16.4.23.1\endcsname{ega4.c04.prop.16-4-23-1}
\expandafter\def\csname sga@target@IV.16.4.24\endcsname{ega4.c04.statement.16-4-24}
\expandafter\def\csname sga@target@IV.16.4.3\endcsname{ega4.c04.statement.16-4-3}
\expandafter\def\csname sga@target@IV.16.4.3.1\endcsname{ega4.c04.statement.16-4-3-1}
\expandafter\def\csname sga@target@IV.16.4.3.2\endcsname{ega4.c04.statement.16-4-3-2}
\expandafter\def\csname sga@target@IV.16.4.3.3\endcsname{ega4.c04.eq.16-4-3-3}
\expandafter\def\csname sga@target@IV.16.4.3.4\endcsname{ega4.c04.eq.16-4-3-4}
\expandafter\def\csname sga@target@IV.16.4.3.5\endcsname{ega4.c04.eq.16-4-3-5}
\expandafter\def\csname sga@target@IV.16.4.3.6\endcsname{ega4.c04.eq.16-4-3-6}
\expandafter\def\csname sga@target@IV.16.4.3.7\endcsname{ega4.c04.eq.16-4-3-7}
\expandafter\def\csname sga@target@IV.16.4.4\endcsname{ega4.c04.statement.16-4-4}
\expandafter\def\csname sga@target@IV.16.4.5\endcsname{ega4.c04.prop.16-4-5}
\expandafter\def\csname sga@target@IV.16.4.6\endcsname{ega4.c04.cor.16-4-6}
\expandafter\def\csname sga@target@IV.16.4.7\endcsname{ega4.c04.statement.16-4-7}
\expandafter\def\csname sga@target@IV.16.4.7.1\endcsname{ega4.c04.statement.16-4-7-1}
\expandafter\def\csname sga@target@IV.16.4.7.2\endcsname{ega4.c04.eq.16-4-7-2}
\expandafter\def\csname sga@target@IV.16.4.7.3\endcsname{ega4.c04.eq.16-4-7-3}
\expandafter\def\csname sga@target@IV.16.4.8\endcsname{ega4.c04.prop.16-4-8}
\expandafter\def\csname sga@target@IV.16.4.8.1\endcsname{ega4.c04.prop.16-4-8-1}
\expandafter\def\csname sga@target@IV.16.4.8.2\endcsname{ega4.c04.eq.16-4-8-2}
\expandafter\def\csname sga@target@IV.16.4.8.3\endcsname{ega4.c04.lem.16-4-8-3}
\expandafter\def\csname sga@target@IV.16.4.8.4\endcsname{ega4.c04.eq.16-4-8-4}
\expandafter\def\csname sga@target@IV.16.4.9\endcsname{ega4.c04.cor.16-4-9}
\expandafter\def\csname sga@target@IV.16.4.9.1\endcsname{ega4.c04.cor.16-4-9-1}
\expandafter\def\csname sga@target@IV.16.5\endcsname{ega4.c04.subsec.16-5}
\expandafter\def\csname sga@target@IV.16.5.1\endcsname{ega4.c04.statement.16-5-1}
\expandafter\def\csname sga@target@IV.16.5.1.1\endcsname{ega4.c04.eq.16-5-1-1}
\expandafter\def\csname sga@target@IV.16.5.1.2\endcsname{ega4.c04.eq.16-5-1-2}
\expandafter\def\csname sga@target@IV.16.5.10\endcsname{ega4.c04.statement.16-5-10}
\expandafter\def\csname sga@target@IV.16.5.10.1\endcsname{ega4.c04.statement.16-5-10-1}
\expandafter\def\csname sga@target@IV.16.5.10.2\endcsname{ega4.c04.statement.16-5-10-2}
\expandafter\def\csname sga@target@IV.16.5.11\endcsname{ega4.c04.prop.16-5-11}
\expandafter\def\csname sga@target@IV.16.5.12\endcsname{ega4.c04.statement.16-5-12}
\expandafter\def\csname sga@target@IV.16.5.12.1\endcsname{ega4.c04.statement.16-5-12-1}
\expandafter\def\csname sga@target@IV.16.5.12.2\endcsname{ega4.c04.eq.16-5-12-2}
\expandafter\def\csname sga@target@IV.16.5.13\endcsname{ega4.c04.statement.16-5-13}
\expandafter\def\csname sga@target@IV.16.5.13.1\endcsname{ega4.c04.statement.16-5-13-1}
\expandafter\def\csname sga@target@IV.16.5.13.2\endcsname{ega4.c04.eq.16-5-13-2}
\expandafter\def\csname sga@target@IV.16.5.13.3\endcsname{ega4.c04.eq.16-5-13-3}
\expandafter\def\csname sga@target@IV.16.5.13.4\endcsname{ega4.c04.eq.16-5-13-4}
\expandafter\def\csname sga@target@IV.16.5.13.5\endcsname{ega4.c04.eq.16-5-13-5}
\expandafter\def\csname sga@target@IV.16.5.14\endcsname{ega4.c04.statement.16-5-14}
\expandafter\def\csname sga@target@IV.16.5.14.1\endcsname{ega4.c04.statement.16-5-14-1}
\expandafter\def\csname sga@target@IV.16.5.14.2\endcsname{ega4.c04.eq.16-5-14-2}
\expandafter\def\csname sga@target@IV.16.5.14.3\endcsname{ega4.c04.eq.16-5-14-3}
\expandafter\def\csname sga@target@IV.16.5.15\endcsname{ega4.c04.statement.16-5-15}
\expandafter\def\csname sga@target@IV.16.5.16\endcsname{ega4.c04.prop.16-5-16}
\expandafter\def\csname sga@target@IV.16.5.17\endcsname{ega4.c04.prop.16-5-17}
\expandafter\def\csname sga@target@IV.16.5.18\endcsname{ega4.c04.cor.16-5-18}
\expandafter\def\csname sga@target@IV.16.5.2\endcsname{ega4.c04.prop.16-5-2}
\expandafter\def\csname sga@target@IV.16.5.3\endcsname{ega4.c04.prop.16-5-3}
\expandafter\def\csname sga@target@IV.16.5.3.1\endcsname{ega4.c04.eq.16-5-3-1}
\expandafter\def\csname sga@target@IV.16.5.4\endcsname{ega4.c04.statement.16-5-4}
\expandafter\def\csname sga@target@IV.16.5.5\endcsname{ega4.c04.cor.16-5-5}
\expandafter\def\csname sga@target@IV.16.5.5.1\endcsname{ega4.c04.cor.16-5-5-1}
\expandafter\def\csname sga@target@IV.16.5.6\endcsname{ega4.c04.cor.16-5-6}
\expandafter\def\csname sga@target@IV.16.5.7\endcsname{ega4.c04.statement.16-5-7}
\expandafter\def\csname sga@target@IV.16.5.7.1\endcsname{ega4.c04.statement.16-5-7-1}
\expandafter\def\csname sga@target@IV.16.5.8\endcsname{ega4.c04.statement.16-5-8}
\expandafter\def\csname sga@target@IV.16.5.8.1\endcsname{ega4.c04.eq.16-5-8-1}
\expandafter\def\csname sga@target@IV.16.5.8.2\endcsname{ega4.c04.eq.16-5-8-2}
\expandafter\def\csname sga@target@IV.16.5.9\endcsname{ega4.c04.statement.16-5-9}
\expandafter\def\csname sga@target@IV.16.5.9.1\endcsname{ega4.c04.statement.16-5-9-1}
\expandafter\def\csname sga@target@IV.16.6\endcsname{ega4.c04.subsec.16-6}
\expandafter\def\csname sga@target@IV.16.6.1\endcsname{ega4.c04.statement.16-6-1}
\expandafter\def\csname sga@target@IV.16.6.1.1\endcsname{ega4.c04.statement.16-6-1-1}
\expandafter\def\csname sga@target@IV.16.6.1.2\endcsname{ega4.c04.statement.16-6-1-2}
\expandafter\def\csname sga@target@IV.16.6.2\endcsname{ega4.c04.thm.16-6-2}
\expandafter\def\csname sga@target@IV.16.6.2.1\endcsname{ega4.c04.eq.16-6-2-1}
\expandafter\def\csname sga@target@IV.16.6.2.2\endcsname{ega4.c04.eq.16-6-2-2}
\expandafter\def\csname sga@target@IV.16.6.2.3\endcsname{ega4.c04.eq.16-6-2-3}
\expandafter\def\csname sga@target@IV.16.6.2.4\endcsname{ega4.c04.eq.16-6-2-4}
\expandafter\def\csname sga@target@IV.16.6.2.5\endcsname{ega4.c04.eq.16-6-2-5}
\expandafter\def\csname sga@target@IV.16.6.3\endcsname{ega4.c04.def.16-6-3}
\expandafter\def\csname sga@target@IV.16.6.4\endcsname{ega4.c04.prop.16-6-4}
\expandafter\def\csname sga@target@IV.16.6.4.1\endcsname{ega4.c04.prop.16-6-4-1}
\expandafter\def\csname sga@target@IV.16.6.5\endcsname{ega4.c04.statement.16-6-5}
\expandafter\def\csname sga@target@IV.16.6.5.1\endcsname{ega4.c04.statement.16-6-5-1}
\expandafter\def\csname sga@target@IV.16.7\endcsname{ega4.c04.subsec.16-7}
\expandafter\def\csname sga@target@IV.16.7.1\endcsname{ega4.c04.statement.16-7-1}
\expandafter\def\csname sga@target@IV.16.7.1.1\endcsname{ega4.c04.eq.16-7-1-1}
\expandafter\def\csname sga@target@IV.16.7.1.2\endcsname{ega4.c04.eq.16-7-1-2}
\expandafter\def\csname sga@target@IV.16.7.10\endcsname{ega4.c04.rem.16-7-10}
\expandafter\def\csname sga@target@IV.16.7.2\endcsname{ega4.c04.statement.16-7-2}
\expandafter\def\csname sga@target@IV.16.7.2.1\endcsname{ega4.c04.statement.16-7-2-1}
\expandafter\def\csname sga@target@IV.16.7.3\endcsname{ega4.c04.prop.16-7-3}
\expandafter\def\csname sga@target@IV.16.7.4\endcsname{ega4.c04.statement.16-7-4}
\expandafter\def\csname sga@target@IV.16.7.4.1\endcsname{ega4.c04.statement.16-7-4-1}
\expandafter\def\csname sga@target@IV.16.7.5\endcsname{ega4.c04.statement.16-7-5}
\expandafter\def\csname sga@target@IV.16.7.5.1\endcsname{ega4.c04.statement.16-7-5-1}
\expandafter\def\csname sga@target@IV.16.7.5.2\endcsname{ega4.c04.statement.16-7-5-2}
\expandafter\def\csname sga@target@IV.16.7.5.3\endcsname{ega4.c04.statement.16-7-5-3}
\expandafter\def\csname sga@target@IV.16.7.5.4\endcsname{ega4.c04.statement.16-7-5-4}
\expandafter\def\csname sga@target@IV.16.7.6\endcsname{ega4.c04.prop.16-7-6}
\expandafter\def\csname sga@target@IV.16.7.7\endcsname{ega4.c04.statement.16-7-7}
\expandafter\def\csname sga@target@IV.16.7.7.1\endcsname{ega4.c04.eq.16-7-7-1}
\expandafter\def\csname sga@target@IV.16.7.7.2\endcsname{ega4.c04.eq.16-7-7-2}
\expandafter\def\csname sga@target@IV.16.7.8\endcsname{ega4.c04.statement.16-7-8}
\expandafter\def\csname sga@target@IV.16.7.8.1\endcsname{ega4.c04.statement.16-7-8-1}
\expandafter\def\csname sga@target@IV.16.7.8.2\endcsname{ega4.c04.eq.16-7-8-2}
\expandafter\def\csname sga@target@IV.16.7.9\endcsname{ega4.c04.statement.16-7-9}
\expandafter\def\csname sga@target@IV.16.7.9.1\endcsname{ega4.c04.statement.16-7-9-1}
\expandafter\def\csname sga@target@IV.16.8\endcsname{ega4.c04.subsec.16-8}
\expandafter\def\csname sga@target@IV.16.8.1\endcsname{ega4.c04.def.16-8-1}
\expandafter\def\csname sga@target@IV.16.8.10\endcsname{ega4.c04.cor.16-8-10}
\expandafter\def\csname sga@target@IV.16.8.11\endcsname{ega4.c04.statement.16-8-11}
\expandafter\def\csname sga@target@IV.16.8.2\endcsname{ega4.c04.statement.16-8-2}
\expandafter\def\csname sga@target@IV.16.8.3\endcsname{ega4.c04.statement.16-8-3}
\expandafter\def\csname sga@target@IV.16.8.3.1\endcsname{ega4.c04.eq.16-8-3-1}
\expandafter\def\csname sga@target@IV.16.8.4\endcsname{ega4.c04.prop.16-8-4}
\expandafter\def\csname sga@target@IV.16.8.4.1\endcsname{ega4.c04.prop.16-8-4-1}
\expandafter\def\csname sga@target@IV.16.8.5\endcsname{ega4.c04.statement.16-8-5}
\expandafter\def\csname sga@target@IV.16.8.5.1\endcsname{ega4.c04.eq.16-8-5-1}
\expandafter\def\csname sga@target@IV.16.8.5.2\endcsname{ega4.c04.eq.16-8-5-2}
\expandafter\def\csname sga@target@IV.16.8.6\endcsname{ega4.c04.prop.16-8-6}
\expandafter\def\csname sga@target@IV.16.8.7\endcsname{ega4.c04.statement.16-8-7}
\expandafter\def\csname sga@target@IV.16.8.7.1\endcsname{ega4.c04.eq.16-8-7-1}
\expandafter\def\csname sga@target@IV.16.8.8\endcsname{ega4.c04.prop.16-8-8}
\expandafter\def\csname sga@target@IV.16.8.8.1\endcsname{ega4.c04.eq.16-8-8-1}
\expandafter\def\csname sga@target@IV.16.8.8.2\endcsname{ega4.c04.eq.16-8-8-2}
\expandafter\def\csname sga@target@IV.16.8.9\endcsname{ega4.c04.prop.16-8-9}
\expandafter\def\csname sga@target@IV.16.8.9.1\endcsname{ega4.c04.lem.16-8-9-1}
\expandafter\def\csname sga@target@IV.16.8.9.2\endcsname{ega4.c04.lem.16-8-9-2}
\expandafter\def\csname sga@target@IV.16.8.9.3\endcsname{ega4.c04.lem.16-8-9-3}
\expandafter\def\csname sga@target@IV.16.8.9.4\endcsname{ega4.c04.eq.16-8-9-4}
\expandafter\def\csname sga@target@IV.16.9\endcsname{ega4.c04.subsec.16-9}
\expandafter\def\csname sga@target@IV.16.9.1\endcsname{ega4.c04.def.16-9-1}
\expandafter\def\csname sga@target@IV.16.9.10\endcsname{ega4.c04.prop.16-9-10}
\expandafter\def\csname sga@target@IV.16.9.11\endcsname{ega4.c04.cor.16-9-11}
\expandafter\def\csname sga@target@IV.16.9.12\endcsname{ega4.c04.statement.16-9-12}
\expandafter\def\csname sga@target@IV.16.9.13\endcsname{ega4.c04.prop.16-9-13}
\expandafter\def\csname sga@target@IV.16.9.13.1\endcsname{ega4.c04.prop.16-9-13-1}
\expandafter\def\csname sga@target@IV.16.9.13.2\endcsname{ega4.c04.lem.16-9-13-2}
\expandafter\def\csname sga@target@IV.16.9.13.3\endcsname{ega4.c04.lem.16-9-13-3}
\expandafter\def\csname sga@target@IV.16.9.13.4\endcsname{ega4.c04.eq.16-9-13-4}
\expandafter\def\csname sga@target@IV.16.9.13.5\endcsname{ega4.c04.cor.16-9-13-5}
\expandafter\def\csname sga@target@IV.16.9.13.6\endcsname{ega4.c04.cor.16-9-13-6}
\expandafter\def\csname sga@target@IV.16.9.13.7\endcsname{ega4.c04.eq.16-9-13-7}
\expandafter\def\csname sga@target@IV.16.9.2\endcsname{ega4.c04.def.16-9-2}
\expandafter\def\csname sga@target@IV.16.9.3\endcsname{ega4.c04.prop.16-9-3}
\expandafter\def\csname sga@target@IV.16.9.4\endcsname{ega4.c04.cor.16-9-4}
\expandafter\def\csname sga@target@IV.16.9.4.1\endcsname{ega4.c04.eq.16-9-4-1}
\expandafter\def\csname sga@target@IV.16.9.5\endcsname{ega4.c04.cor.16-9-5}
\expandafter\def\csname sga@target@IV.16.9.6\endcsname{ega4.c04.statement.16-9-6}
\expandafter\def\csname sga@target@IV.16.9.7\endcsname{ega4.c04.statement.16-9-7}
\expandafter\def\csname sga@target@IV.16.9.8\endcsname{ega4.c04.prop.16-9-8}
\expandafter\def\csname sga@target@IV.16.9.9\endcsname{ega4.c04.cor.16-9-9}
\expandafter\def\csname sga@target@IV.17.1.1\endcsname{ega4.c04.def.17-1-1}
\expandafter\def\csname sga@target@IV.17.1.1.1\endcsname{ega4.c04.def.17-1-1-1}
\expandafter\def\csname sga@target@IV.17.1.2\endcsname{ega4.c04.rem.17-1-2}
\expandafter\def\csname sga@target@IV.17.1.2.1\endcsname{ega4.c04.eq.17-1-2-1}
\expandafter\def\csname sga@target@IV.17.1.3\endcsname{ega4.c04.prop.17-1-3}
\expandafter\def\csname sga@target@IV.17.1.4\endcsname{ega4.c04.prop.17-1-4}
\expandafter\def\csname sga@target@IV.17.1.5\endcsname{ega4.c04.cor.17-1-5}
\expandafter\def\csname sga@target@IV.17.1.6\endcsname{ega4.c04.prop.17-1-6}
\expandafter\def\csname sga@target@IV.17.10.1\endcsname{ega4.c04.def.17-10-1}
\expandafter\def\csname sga@target@IV.17.10.2\endcsname{ega4.c04.prop.17-10-2}
\expandafter\def\csname sga@target@IV.17.10.3\endcsname{ega4.c04.cor.17-10-3}
\expandafter\def\csname sga@target@IV.17.10.3.1\endcsname{ega4.c04.cor.17-10-3-1}
\expandafter\def\csname sga@target@IV.17.10.4\endcsname{ega4.c04.cor.17-10-4}
\expandafter\def\csname sga@target@IV.17.10.4.1\endcsname{ega4.c04.eq.17-10-4-1}
\expandafter\def\csname sga@target@IV.17.11.1\endcsname{ega4.c04.thm.17-11-1}
\expandafter\def\csname sga@target@IV.17.11.1.1\endcsname{ega4.c04.eq.17-11-1-1}
\expandafter\def\csname sga@target@IV.17.11.1.2\endcsname{ega4.c04.eq.17-11-1-2}
\expandafter\def\csname sga@target@IV.17.11.2\endcsname{ega4.c04.cor.17-11-2}
\expandafter\def\csname sga@target@IV.17.11.3\endcsname{ega4.c04.cor.17-11-3}
\expandafter\def\csname sga@target@IV.17.11.4\endcsname{ega4.c04.cor.17-11-4}
\expandafter\def\csname sga@target@IV.17.11.5\endcsname{ega4.c04.prop.17-11-5}
\expandafter\def\csname sga@target@IV.17.12.1\endcsname{ega4.c04.thm.17-12-1}
\expandafter\def\csname sga@target@IV.17.12.2\endcsname{ega4.c04.cor.17-12-2}
\expandafter\def\csname sga@target@IV.17.12.3\endcsname{ega4.c04.cor.17-12-3}
\expandafter\def\csname sga@target@IV.17.12.4\endcsname{ega4.c04.prop.17-12-4}
\expandafter\def\csname sga@target@IV.17.12.5\endcsname{ega4.c04.prop.17-12-5}
\expandafter\def\csname sga@target@IV.17.12.6\endcsname{ega4.c04.cor.17-12-6}
\expandafter\def\csname sga@target@IV.17.12.7\endcsname{ega4.c04.ex.17-12-7}
\expandafter\def\csname sga@target@IV.17.12.8\endcsname{ega4.c04.rem.17-12-8}
\expandafter\def\csname sga@target@IV.17.13\endcsname{ega4.c04.subsec.17-13.key-iv-17-13}
\expandafter\def\csname sga@target@IV.17.13.1\endcsname{ega4.c04.statement.17-13-1}
\expandafter\def\csname sga@target@IV.17.13.1.1\endcsname{ega4.c04.statement.17-13-1-1}
\expandafter\def\csname sga@target@IV.17.13.1.2\endcsname{ega4.c04.statement.17-13-1-2}
\expandafter\def\csname sga@target@IV.17.13.1.3\endcsname{ega4.c04.eq.17-13-1-3}
\expandafter\def\csname sga@target@IV.17.13.10\endcsname{ega4.c04.cor.17-13-10}
\expandafter\def\csname sga@target@IV.17.13.11\endcsname{ega4.c04.statement.17-13-11}
\expandafter\def\csname sga@target@IV.17.13.11.1\endcsname{ega4.c04.statement.17-13-11-1}
\expandafter\def\csname sga@target@IV.17.13.12\endcsname{ega4.c04.cor.17-13-12}
\expandafter\def\csname sga@target@IV.17.13.12.1\endcsname{ega4.c04.cor.17-13-12-1}
\expandafter\def\csname sga@target@IV.17.13.12.2\endcsname{ega4.c04.eq.17-13-12-2}
\expandafter\def\csname sga@target@IV.17.13.2\endcsname{ega4.c04.prop.17-13-2}
\expandafter\def\csname sga@target@IV.17.13.2.1\endcsname{ega4.c04.eq.17-13-2-1}
\expandafter\def\csname sga@target@IV.17.13.2.2\endcsname{ega4.c04.eq.17-13-2-2}
\expandafter\def\csname sga@target@IV.17.13.3\endcsname{ega4.c04.def.17-13-3}
\expandafter\def\csname sga@target@IV.17.13.4\endcsname{ega4.c04.statement.17-13-4}
\expandafter\def\csname sga@target@IV.17.13.5\endcsname{ega4.c04.statement.17-13-5}
\expandafter\def\csname sga@target@IV.17.13.5.1\endcsname{ega4.c04.statement.17-13-5-1}
\expandafter\def\csname sga@target@IV.17.13.5.2\endcsname{ega4.c04.eq.17-13-5-2}
\expandafter\def\csname sga@target@IV.17.13.5.3\endcsname{ega4.c04.eq.17-13-5-3}
\expandafter\def\csname sga@target@IV.17.13.5.4\endcsname{ega4.c04.eq.17-13-5-4}
\expandafter\def\csname sga@target@IV.17.13.6\endcsname{ega4.c04.cor.17-13-6}
\expandafter\def\csname sga@target@IV.17.13.6.1\endcsname{ega4.c04.eq.17-13-6-1}
\expandafter\def\csname sga@target@IV.17.13.6.2\endcsname{ega4.c04.eq.17-13-6-2}
\expandafter\def\csname sga@target@IV.17.13.7\endcsname{ega4.c04.statement.17-13-7}
\expandafter\def\csname sga@target@IV.17.13.7.1\endcsname{ega4.c04.eq.17-13-7-1}
\expandafter\def\csname sga@target@IV.17.13.8\endcsname{ega4.c04.prop.17-13-8}
\expandafter\def\csname sga@target@IV.17.13.8.1\endcsname{ega4.c04.eq.17-13-8-1}
\expandafter\def\csname sga@target@IV.17.13.9\endcsname{ega4.c04.statement.17-13-9}
\expandafter\def\csname sga@target@IV.17.13.9.1\endcsname{ega4.c04.eq.17-13-9-1}
\expandafter\def\csname sga@target@IV.17.13.9.2\endcsname{ega4.c04.eq.17-13-9-2}
\expandafter\def\csname sga@target@IV.17.14.1\endcsname{ega4.c04.prop.17-14-1}
\expandafter\def\csname sga@target@IV.17.14.2\endcsname{ega4.c04.prop.17-14-2}
\expandafter\def\csname sga@target@IV.17.14.3\endcsname{ega4.c04.rem.17-14-3}
\expandafter\def\csname sga@target@IV.17.15.1\endcsname{ega4.c04.prop.17-15-1}
\expandafter\def\csname sga@target@IV.17.15.10\endcsname{ega4.c04.statement.17-15-10}
\expandafter\def\csname sga@target@IV.17.15.11\endcsname{ega4.c04.prop.17-15-11}
\expandafter\def\csname sga@target@IV.17.15.12\endcsname{ega4.c04.prop.17-15-12}
\expandafter\def\csname sga@target@IV.17.15.13\endcsname{ega4.c04.cor.17-15-13}
\expandafter\def\csname sga@target@IV.17.15.14\endcsname{ega4.c04.prop.17-15-14}
\expandafter\def\csname sga@target@IV.17.15.14.1\endcsname{ega4.c04.lem.17-15-14-1}
\expandafter\def\csname sga@target@IV.17.15.14.2\endcsname{ega4.c04.lem.17-15-14-2}
\expandafter\def\csname sga@target@IV.17.15.15\endcsname{ega4.c04.prop.17-15-15}
\expandafter\def\csname sga@target@IV.17.15.2\endcsname{ega4.c04.cor.17-15-2}
\expandafter\def\csname sga@target@IV.17.15.3\endcsname{ega4.c04.prop.17-15-3}
\expandafter\def\csname sga@target@IV.17.15.4\endcsname{ega4.c04.cor.17-15-4}
\expandafter\def\csname sga@target@IV.17.15.5\endcsname{ega4.c04.prop.17-15-5}
\expandafter\def\csname sga@target@IV.17.15.6\endcsname{ega4.c04.cor.17-15-6}
\expandafter\def\csname sga@target@IV.17.15.7\endcsname{ega4.c04.cor.17-15-7}
\expandafter\def\csname sga@target@IV.17.15.8\endcsname{ega4.c04.prop.17-15-8}
\expandafter\def\csname sga@target@IV.17.15.8.1\endcsname{ega4.c04.eq.17-15-8-1}
\expandafter\def\csname sga@target@IV.17.15.8.2\endcsname{ega4.c04.eq.17-15-8-2}
\expandafter\def\csname sga@target@IV.17.15.9\endcsname{ega4.c04.cor.17-15-9}
\expandafter\def\csname sga@target@IV.17.16\endcsname{ega4.c04.subsec.17-16.key-iv-17-16}
\expandafter\def\csname sga@target@IV.17.16.1\endcsname{ega4.c04.prop.17-16-1}
\expandafter\def\csname sga@target@IV.17.16.2\endcsname{ega4.c04.cor.17-16-2}
\expandafter\def\csname sga@target@IV.17.16.3\endcsname{ega4.c04.cor.17-16-3}
\expandafter\def\csname sga@target@IV.17.16.4\endcsname{ega4.c04.prop.17-16-4}
\expandafter\def\csname sga@target@IV.17.16.5\endcsname{ega4.c04.cor.17-16-5}
\expandafter\def\csname sga@target@IV.17.16.6\endcsname{ega4.c04.prop.17-16-6}
\expandafter\def\csname sga@target@IV.17.2.1\endcsname{ega4.c04.prop.17-2-1}
\expandafter\def\csname sga@target@IV.17.2.2\endcsname{ega4.c04.cor.17-2-2}
\expandafter\def\csname sga@target@IV.17.2.3\endcsname{ega4.c04.prop.17-2-3}
\expandafter\def\csname sga@target@IV.17.2.3.1\endcsname{ega4.c04.eq.17-2-3-1}
\expandafter\def\csname sga@target@IV.17.2.4\endcsname{ega4.c04.cor.17-2-4}
\expandafter\def\csname sga@target@IV.17.2.5\endcsname{ega4.c04.prop.17-2-5}
\expandafter\def\csname sga@target@IV.17.2.5.1\endcsname{ega4.c04.prop.17-2-5-1}
\expandafter\def\csname sga@target@IV.17.2.6\endcsname{ega4.c04.prop.17-2-6}
\expandafter\def\csname sga@target@IV.17.3.1\endcsname{ega4.c04.def.17-3-1}
\expandafter\def\csname sga@target@IV.17.3.2\endcsname{ega4.c04.subsec.17-3-2}
\expandafter\def\csname sga@target@IV.17.3.3\endcsname{ega4.c04.prop.17-3-3}
\expandafter\def\csname sga@target@IV.17.3.4\endcsname{ega4.c04.prop.17-3-4}
\expandafter\def\csname sga@target@IV.17.3.5\endcsname{ega4.c04.cor.17-3-5}
\expandafter\def\csname sga@target@IV.17.3.6\endcsname{ega4.c04.prop.17-3-6}
\expandafter\def\csname sga@target@IV.17.3.7\endcsname{ega4.c04.def.17-3-7}
\expandafter\def\csname sga@target@IV.17.3.8\endcsname{ega4.c04.prop.17-3-8}
\expandafter\def\csname sga@target@IV.17.3.9\endcsname{ega4.c04.cor.17-3-9}
\expandafter\def\csname sga@target@IV.17.4.1\endcsname{ega4.c04.thm.17-4-1}
\expandafter\def\csname sga@target@IV.17.4.1.1\endcsname{ega4.c04.eq.17-4-1-1}
\expandafter\def\csname sga@target@IV.17.4.1.2\endcsname{ega4.c04.rem.17-4-1-2}
\expandafter\def\csname sga@target@IV.17.4.10\endcsname{ega4.c04.rem.17-4-10}
\expandafter\def\csname sga@target@IV.17.4.2\endcsname{ega4.c04.cor.17-4-2}
\expandafter\def\csname sga@target@IV.17.4.3\endcsname{ega4.c04.cor.17-4-3}
\expandafter\def\csname sga@target@IV.17.4.4\endcsname{ega4.c04.prop.17-4-4}
\expandafter\def\csname sga@target@IV.17.4.5\endcsname{ega4.c04.statement.17-4-5}
\expandafter\def\csname sga@target@IV.17.4.6\endcsname{ega4.c04.prop.17-4-6}
\expandafter\def\csname sga@target@IV.17.4.7\endcsname{ega4.c04.cor.17-4-7}
\expandafter\def\csname sga@target@IV.17.4.8\endcsname{ega4.c04.cor.17-4-8}
\expandafter\def\csname sga@target@IV.17.4.9\endcsname{ega4.c04.prop.17-4-9}
\expandafter\def\csname sga@target@IV.17.5.1\endcsname{ega4.c04.thm.17-5-1}
\expandafter\def\csname sga@target@IV.17.5.1.1\endcsname{ega4.c04.eq.17-5-1-1}
\expandafter\def\csname sga@target@IV.17.5.1.2\endcsname{ega4.c04.eq.17-5-1-2}
\expandafter\def\csname sga@target@IV.17.5.2\endcsname{ega4.c04.cor.17-5-2}
\expandafter\def\csname sga@target@IV.17.5.3\endcsname{ega4.c04.prop.17-5-3}
\expandafter\def\csname sga@target@IV.17.5.4\endcsname{ega4.c04.rem.17-5-4}
\expandafter\def\csname sga@target@IV.17.5.5\endcsname{ega4.c04.prop.17-5-5}
\expandafter\def\csname sga@target@IV.17.5.6\endcsname{ega4.c04.cor.17-5-6}
\expandafter\def\csname sga@target@IV.17.5.7\endcsname{ega4.c04.prop.17-5-7}
\expandafter\def\csname sga@target@IV.17.5.8\endcsname{ega4.c04.prop.17-5-8}
\expandafter\def\csname sga@target@IV.17.6.1\endcsname{ega4.c04.thm.17-6-1}
\expandafter\def\csname sga@target@IV.17.6.2\endcsname{ega4.c04.cor.17-6-2}
\expandafter\def\csname sga@target@IV.17.6.3\endcsname{ega4.c04.prop.17-6-3}
\expandafter\def\csname sga@target@IV.17.6.4\endcsname{ega4.c04.prop.17-6-4}
\expandafter\def\csname sga@target@IV.17.7\endcsname{ega4.c04.subsec.17-7.key-iv-17-7}
\expandafter\def\csname sga@target@IV.17.7.1\endcsname{ega4.c04.prop.17-7-1}
\expandafter\def\csname sga@target@IV.17.7.10\endcsname{ega4.c04.prop.17-7-10}
\expandafter\def\csname sga@target@IV.17.7.11\endcsname{ega4.c04.prop.17-7-11}
\expandafter\def\csname sga@target@IV.17.7.11.1\endcsname{ega4.c04.lem.17-7-11-1}
\expandafter\def\csname sga@target@IV.17.7.2\endcsname{ega4.c04.cor.17-7-2}
\expandafter\def\csname sga@target@IV.17.7.3\endcsname{ega4.c04.cor.17-7-3}
\expandafter\def\csname sga@target@IV.17.7.4\endcsname{ega4.c04.prop.17-7-4}
\expandafter\def\csname sga@target@IV.17.7.5\endcsname{ega4.c04.lem.17-7-5}
\expandafter\def\csname sga@target@IV.17.7.6\endcsname{ega4.c04.rem.17-7-6}
\expandafter\def\csname sga@target@IV.17.7.7\endcsname{ega4.c04.prop.17-7-7}
\expandafter\def\csname sga@target@IV.17.7.8\endcsname{ega4.c04.prop.17-7-8}
\expandafter\def\csname sga@target@IV.17.7.9\endcsname{ega4.c04.cor.17-7-9}
\expandafter\def\csname sga@target@IV.17.8\endcsname{ega4.c04.subsec.17-8.key-iv-17-8}
\expandafter\def\csname sga@target@IV.17.8.1\endcsname{ega4.c04.prop.17-8-1}
\expandafter\def\csname sga@target@IV.17.8.2\endcsname{ega4.c04.prop.17-8-2}
\expandafter\def\csname sga@target@IV.17.8.3\endcsname{ega4.c04.rem.17-8-3}
\expandafter\def\csname sga@target@IV.17.9\endcsname{ega4.c04.subsec.17-9.key-iv-17-9}
\expandafter\def\csname sga@target@IV.17.9.1\endcsname{ega4.c04.thm.17-9-1}
\expandafter\def\csname sga@target@IV.17.9.2\endcsname{ega4.c04.cor.17-9-2}
\expandafter\def\csname sga@target@IV.17.9.3\endcsname{ega4.c04.cor.17-9-3}
\expandafter\def\csname sga@target@IV.17.9.4\endcsname{ega4.c04.cor.17-9-4}
\expandafter\def\csname sga@target@IV.17.9.5\endcsname{ega4.c04.cor.17-9-5}
\expandafter\def\csname sga@target@IV.17.9.6\endcsname{ega4.c04.prop.17-9-6}
\expandafter\def\csname sga@target@IV.18\endcsname{ega4.c04.sec.18.key-iv-18}
\expandafter\def\csname sga@target@IV.18.1.1\endcsname{ega4.c04.prop.18-1-1}
\expandafter\def\csname sga@target@IV.18.1.1.1\endcsname{ega4.c04.eq.18-1-1-1}
\expandafter\def\csname sga@target@IV.18.1.2\endcsname{ega4.c04.thm.18-1-2}
\expandafter\def\csname sga@target@IV.18.1.3\endcsname{ega4.c04.cor.18-1-3}
\expandafter\def\csname sga@target@IV.18.10.1\endcsname{ega4.c04.thm.18-10-1}
\expandafter\def\csname sga@target@IV.18.10.10\endcsname{ega4.c04.statement.18-10-10}
\expandafter\def\csname sga@target@IV.18.10.11\endcsname{ega4.c04.rem.18-10-11}
\expandafter\def\csname sga@target@IV.18.10.12\endcsname{ega4.c04.cor.18-10-12}
\expandafter\def\csname sga@target@IV.18.10.13\endcsname{ega4.c04.prop.18-10-13}
\expandafter\def\csname sga@target@IV.18.10.14\endcsname{ega4.c04.cor.18-10-14}
\expandafter\def\csname sga@target@IV.18.10.15\endcsname{ega4.c04.statement.18-10-15}
\expandafter\def\csname sga@target@IV.18.10.16\endcsname{ega4.c04.thm.18-10-16}
\expandafter\def\csname sga@target@IV.18.10.17\endcsname{ega4.c04.rem.18-10-17}
\expandafter\def\csname sga@target@IV.18.10.17.1\endcsname{ega4.c04.statement.18-10-17-1}
\expandafter\def\csname sga@target@IV.18.10.17.2\endcsname{ega4.c04.statement.18-10-17-2}
\expandafter\def\csname sga@target@IV.18.10.18\endcsname{ega4.c04.lem.18-10-18}
\expandafter\def\csname sga@target@IV.18.10.19\endcsname{ega4.c04.prop.18-10-19}
\expandafter\def\csname sga@target@IV.18.10.2\endcsname{ega4.c04.rem.18-10-2}
\expandafter\def\csname sga@target@IV.18.10.20\endcsname{ega4.c04.rem.18-10-20}
\expandafter\def\csname sga@target@IV.18.10.21\endcsname{ega4.c04.prop.18-10-21}
\expandafter\def\csname sga@target@IV.18.10.3\endcsname{ega4.c04.cor.18-10-3}
\expandafter\def\csname sga@target@IV.18.10.3.1\endcsname{ega4.c04.rem.18-10-3-1}
\expandafter\def\csname sga@target@IV.18.10.4\endcsname{ega4.c04.cor.18-10-4}
\expandafter\def\csname sga@target@IV.18.10.5\endcsname{ega4.c04.cor.18-10-5}
\expandafter\def\csname sga@target@IV.18.10.6\endcsname{ega4.c04.lem.18-10-6}
\expandafter\def\csname sga@target@IV.18.10.7\endcsname{ega4.c04.prop.18-10-7}
\expandafter\def\csname sga@target@IV.18.10.8\endcsname{ega4.c04.prop.18-10-8}
\expandafter\def\csname sga@target@IV.18.10.9\endcsname{ega4.c04.cor.18-10-9}
\expandafter\def\csname sga@target@IV.18.11.1\endcsname{ega4.c04.lem.18-11-1}
\expandafter\def\csname sga@target@IV.18.11.10\endcsname{ega4.c04.thm.18-11-10}
\expandafter\def\csname sga@target@IV.18.11.11\endcsname{ega4.c04.rem.18-11-11}
\expandafter\def\csname sga@target@IV.18.11.12\endcsname{ega4.c04.cor.18-11-12}
\expandafter\def\csname sga@target@IV.18.11.12.1\endcsname{ega4.c04.lem.18-11-12-1}
\expandafter\def\csname sga@target@IV.18.11.2\endcsname{ega4.c04.lem.18-11-2}
\expandafter\def\csname sga@target@IV.18.11.2.1\endcsname{ega4.c04.eq.18-11-2-1}
\expandafter\def\csname sga@target@IV.18.11.2.2\endcsname{ega4.c04.eq.18-11-2-2}
\expandafter\def\csname sga@target@IV.18.11.3\endcsname{ega4.c04.lem.18-11-3}
\expandafter\def\csname sga@target@IV.18.11.3.1\endcsname{ega4.c04.eq.18-11-3-1}
\expandafter\def\csname sga@target@IV.18.11.3.2\endcsname{ega4.c04.eq.18-11-3-2}
\expandafter\def\csname sga@target@IV.18.11.3.3\endcsname{ega4.c04.lem.18-11-3-3}
\expandafter\def\csname sga@target@IV.18.11.3.4\endcsname{ega4.c04.lem.18-11-3-4}
\expandafter\def\csname sga@target@IV.18.11.3.5\endcsname{ega4.c04.lem.18-11-3-5}
\expandafter\def\csname sga@target@IV.18.11.3.6\endcsname{ega4.c04.eq.18-11-3-6}
\expandafter\def\csname sga@target@IV.18.11.4\endcsname{ega4.c04.lem.18-11-4}
\expandafter\def\csname sga@target@IV.18.11.5\endcsname{ega4.c04.prop.18-11-5}
\expandafter\def\csname sga@target@IV.18.11.5.1\endcsname{ega4.c04.lem.18-11-5-1}
\expandafter\def\csname sga@target@IV.18.11.5.2\endcsname{ega4.c04.lem.18-11-5-2}
\expandafter\def\csname sga@target@IV.18.11.6\endcsname{ega4.c04.rem.18-11-6}
\expandafter\def\csname sga@target@IV.18.11.7\endcsname{ega4.c04.cor.18-11-7}
\expandafter\def\csname sga@target@IV.18.11.8\endcsname{ega4.c04.cor.18-11-8}
\expandafter\def\csname sga@target@IV.18.11.9\endcsname{ega4.c04.cor.18-11-9}
\expandafter\def\csname sga@target@IV.18.12.1\endcsname{ega4.c04.thm.18-12-1}
\expandafter\def\csname sga@target@IV.18.12.10\endcsname{ega4.c04.cor.18-12-10}
\expandafter\def\csname sga@target@IV.18.12.11\endcsname{ega4.c04.cor.18-12-11}
\expandafter\def\csname sga@target@IV.18.12.12\endcsname{ega4.c04.prop.18-12-12}
\expandafter\def\csname sga@target@IV.18.12.13\endcsname{ega4.c04.cor.18-12-13}
\expandafter\def\csname sga@target@IV.18.12.14\endcsname{ega4.c04.rem.18-12-14}
\expandafter\def\csname sga@target@IV.18.12.15\endcsname{ega4.c04.prop.18-12-15}
\expandafter\def\csname sga@target@IV.18.12.16\endcsname{ega4.c04.prop.18-12-16}
\expandafter\def\csname sga@target@IV.18.12.17\endcsname{ega4.c04.cor.18-12-17}
\expandafter\def\csname sga@target@IV.18.12.2\endcsname{ega4.c04.rem.18-12-2}
\expandafter\def\csname sga@target@IV.18.12.3\endcsname{ega4.c04.cor.18-12-3}
\expandafter\def\csname sga@target@IV.18.12.4\endcsname{ega4.c04.cor.18-12-4}
\expandafter\def\csname sga@target@IV.18.12.5\endcsname{ega4.c04.rem.18-12-5}
\expandafter\def\csname sga@target@IV.18.12.6\endcsname{ega4.c04.cor.18-12-6}
\expandafter\def\csname sga@target@IV.18.12.7\endcsname{ega4.c04.prop.18-12-7}
\expandafter\def\csname sga@target@IV.18.12.7.1\endcsname{ega4.c04.lem.18-12-7-1}
\expandafter\def\csname sga@target@IV.18.12.8\endcsname{ega4.c04.prop.18-12-8}
\expandafter\def\csname sga@target@IV.18.12.9\endcsname{ega4.c04.rem.18-12-9}
\expandafter\def\csname sga@target@IV.18.2.1\endcsname{ega4.c04.statement.18-2-1}
\expandafter\def\csname sga@target@IV.18.2.1.1\endcsname{ega4.c04.statement.18-2-1-1}
\expandafter\def\csname sga@target@IV.18.2.1.2\endcsname{ega4.c04.eq.18-2-1-2}
\expandafter\def\csname sga@target@IV.18.2.10\endcsname{ega4.c04.statement.18-2-10}
\expandafter\def\csname sga@target@IV.18.2.2\endcsname{ega4.c04.statement.18-2-2}
\expandafter\def\csname sga@target@IV.18.2.2.1\endcsname{ega4.c04.eq.18-2-2-1}
\expandafter\def\csname sga@target@IV.18.2.2.2\endcsname{ega4.c04.eq.18-2-2-2}
\expandafter\def\csname sga@target@IV.18.2.2.3\endcsname{ega4.c04.eq.18-2-2-3}
\expandafter\def\csname sga@target@IV.18.2.2.4\endcsname{ega4.c04.eq.18-2-2-4}
\expandafter\def\csname sga@target@IV.18.2.3\endcsname{ega4.c04.prop.18-2-3}
\expandafter\def\csname sga@target@IV.18.2.4\endcsname{ega4.c04.cor.18-2-4}
\expandafter\def\csname sga@target@IV.18.2.5\endcsname{ega4.c04.cor.18-2-5}
\expandafter\def\csname sga@target@IV.18.2.6\endcsname{ega4.c04.cor.18-2-6}
\expandafter\def\csname sga@target@IV.18.2.7\endcsname{ega4.c04.item.18-2-7}
\expandafter\def\csname sga@target@IV.18.2.7.1\endcsname{ega4.c04.eq.18-2-7-1}
\expandafter\def\csname sga@target@IV.18.2.7.2\endcsname{ega4.c04.eq.18-2-7-2}
\expandafter\def\csname sga@target@IV.18.2.8\endcsname{ega4.c04.prop.18-2-8}
\expandafter\def\csname sga@target@IV.18.2.9\endcsname{ega4.c04.cor.18-2-9}
\expandafter\def\csname sga@target@IV.18.3.1\endcsname{ega4.c04.prop.18-3-1}
\expandafter\def\csname sga@target@IV.18.3.2\endcsname{ega4.c04.prop.18-3-2}
\expandafter\def\csname sga@target@IV.18.3.2.1\endcsname{ega4.c04.lem.18-3-2-1}
\expandafter\def\csname sga@target@IV.18.3.2.2\endcsname{ega4.c04.eq.18-3-2-2}
\expandafter\def\csname sga@target@IV.18.3.3\endcsname{ega4.c04.statement.18-3-3}
\expandafter\def\csname sga@target@IV.18.3.4\endcsname{ega4.c04.thm.18-3-4}
\expandafter\def\csname sga@target@IV.18.4.1\endcsname{ega4.c04.lem.18-4-1}
\expandafter\def\csname sga@target@IV.18.4.10\endcsname{ega4.c04.prop.18-4-10}
\expandafter\def\csname sga@target@IV.18.4.11\endcsname{ega4.c04.statement.18-4-11}
\expandafter\def\csname sga@target@IV.18.4.12\endcsname{ega4.c04.cor.18-4-12}
\expandafter\def\csname sga@target@IV.18.4.13\endcsname{ega4.c04.cor.18-4-13}
\expandafter\def\csname sga@target@IV.18.4.14\endcsname{ega4.c04.cor.18-4-14}
\expandafter\def\csname sga@target@IV.18.4.2\endcsname{ega4.c04.prop.18-4-2}
\expandafter\def\csname sga@target@IV.18.4.3\endcsname{ega4.c04.cor.18-4-3}
\expandafter\def\csname sga@target@IV.18.4.4\endcsname{ega4.c04.lem.18-4-4}
\expandafter\def\csname sga@target@IV.18.4.5\endcsname{ega4.c04.prop.18-4-5}
\expandafter\def\csname sga@target@IV.18.4.6\endcsname{ega4.c04.thm.18-4-6}
\expandafter\def\csname sga@target@IV.18.4.6.1\endcsname{ega4.c04.lem.18-4-6-1}
\expandafter\def\csname sga@target@IV.18.4.7\endcsname{ega4.c04.cor.18-4-7}
\expandafter\def\csname sga@target@IV.18.4.8\endcsname{ega4.c04.cor.18-4-8}
\expandafter\def\csname sga@target@IV.18.4.9\endcsname{ega4.c04.prop.18-4-9}
\expandafter\def\csname sga@target@IV.18.5.1\endcsname{ega4.c04.statement.18-5-1}
\expandafter\def\csname sga@target@IV.18.5.1.1\endcsname{ega4.c04.eq.18-5-1-1}
\expandafter\def\csname sga@target@IV.18.5.10\endcsname{ega4.c04.prop.18-5-10}
\expandafter\def\csname sga@target@IV.18.5.11\endcsname{ega4.c04.thm.18-5-11}
\expandafter\def\csname sga@target@IV.18.5.12\endcsname{ega4.c04.cor.18-5-12}
\expandafter\def\csname sga@target@IV.18.5.13\endcsname{ega4.c04.rem.18-5-13}
\expandafter\def\csname sga@target@IV.18.5.13.1\endcsname{ega4.c04.lem.18-5-13-1}
\expandafter\def\csname sga@target@IV.18.5.14\endcsname{ega4.c04.prop.18-5-14}
\expandafter\def\csname sga@target@IV.18.5.15\endcsname{ega4.c04.prop.18-5-15}
\expandafter\def\csname sga@target@IV.18.5.16\endcsname{ega4.c04.item.18-5-16}
\expandafter\def\csname sga@target@IV.18.5.17\endcsname{ega4.c04.thm.18-5-17}
\expandafter\def\csname sga@target@IV.18.5.18\endcsname{ega4.c04.rem.18-5-18}
\expandafter\def\csname sga@target@IV.18.5.19\endcsname{ega4.c04.prop.18-5-19}
\expandafter\def\csname sga@target@IV.18.5.2\endcsname{ega4.c04.statement.18-5-2}
\expandafter\def\csname sga@target@IV.18.5.3\endcsname{ega4.c04.lem.18-5-3}
\expandafter\def\csname sga@target@IV.18.5.3.1\endcsname{ega4.c04.lem.18-5-3-1}
\expandafter\def\csname sga@target@IV.18.5.4\endcsname{ega4.c04.prop.18-5-4}
\expandafter\def\csname sga@target@IV.18.5.4.1\endcsname{ega4.c04.eq.18-5-4-1}
\expandafter\def\csname sga@target@IV.18.5.4.2\endcsname{ega4.c04.eq.18-5-4-2}
\expandafter\def\csname sga@target@IV.18.5.4.3\endcsname{ega4.c04.lem.18-5-4-3}
\expandafter\def\csname sga@target@IV.18.5.4.4\endcsname{ega4.c04.rem.18-5-4-4}
\expandafter\def\csname sga@target@IV.18.5.5\endcsname{ega4.c04.def.18-5-5}
\expandafter\def\csname sga@target@IV.18.5.6\endcsname{ega4.c04.prop.18-5-6}
\expandafter\def\csname sga@target@IV.18.5.7\endcsname{ega4.c04.rem.18-5-7}
\expandafter\def\csname sga@target@IV.18.5.8\endcsname{ega4.c04.def.18-5-8}
\expandafter\def\csname sga@target@IV.18.5.9\endcsname{ega4.c04.prop.18-5-9}
\expandafter\def\csname sga@target@IV.18.6.1\endcsname{ega4.c04.statement.18-6-1}
\expandafter\def\csname sga@target@IV.18.6.10\endcsname{ega4.c04.cor.18-6-10}
\expandafter\def\csname sga@target@IV.18.6.11\endcsname{ega4.c04.cor.18-6-11}
\expandafter\def\csname sga@target@IV.18.6.12\endcsname{ega4.c04.prop.18-6-12}
\expandafter\def\csname sga@target@IV.18.6.13\endcsname{ega4.c04.cor.18-6-13}
\expandafter\def\csname sga@target@IV.18.6.14\endcsname{ega4.c04.prop.18-6-14}
\expandafter\def\csname sga@target@IV.18.6.15\endcsname{ega4.c04.cor.18-6-15}
\expandafter\def\csname sga@target@IV.18.6.2\endcsname{ega4.c04.lem.18-6-2}
\expandafter\def\csname sga@target@IV.18.6.3\endcsname{ega4.c04.lem.18-6-3}
\expandafter\def\csname sga@target@IV.18.6.4\endcsname{ega4.c04.statement.18-6-4}
\expandafter\def\csname sga@target@IV.18.6.5\endcsname{ega4.c04.def.18-6-5}
\expandafter\def\csname sga@target@IV.18.6.6\endcsname{ega4.c04.thm.18-6-6}
\expandafter\def\csname sga@target@IV.18.6.7\endcsname{ega4.c04.statement.18-6-7}
\expandafter\def\csname sga@target@IV.18.6.8\endcsname{ega4.c04.prop.18-6-8}
\expandafter\def\csname sga@target@IV.18.6.9\endcsname{ega4.c04.thm.18-6-9}
\expandafter\def\csname sga@target@IV.18.6.9.1\endcsname{ega4.c04.lem.18-6-9-1}
\expandafter\def\csname sga@target@IV.18.6.9.2\endcsname{ega4.c04.lem.18-6-9-2}
\expandafter\def\csname sga@target@IV.18.7.1\endcsname{ega4.c04.statement.18-7-1}
\expandafter\def\csname sga@target@IV.18.7.2\endcsname{ega4.c04.prop.18-7-2}
\expandafter\def\csname sga@target@IV.18.7.3\endcsname{ega4.c04.cor.18-7-3}
\expandafter\def\csname sga@target@IV.18.7.4\endcsname{ega4.c04.cor.18-7-4}
\expandafter\def\csname sga@target@IV.18.7.5\endcsname{ega4.c04.prop.18-7-5}
\expandafter\def\csname sga@target@IV.18.7.5.1\endcsname{ega4.c04.lem.18-7-5-1}
\expandafter\def\csname sga@target@IV.18.7.6\endcsname{ega4.c04.cor.18-7-6}
\expandafter\def\csname sga@target@IV.18.7.7\endcsname{ega4.c04.rem.18-7-7}
\expandafter\def\csname sga@target@IV.18.8.1\endcsname{ega4.c04.prop.18-8-1}
\expandafter\def\csname sga@target@IV.18.8.10\endcsname{ega4.c04.prop.18-8-10}
\expandafter\def\csname sga@target@IV.18.8.11\endcsname{ega4.c04.rem.18-8-11}
\expandafter\def\csname sga@target@IV.18.8.12\endcsname{ega4.c04.prop.18-8-12}
\expandafter\def\csname sga@target@IV.18.8.13\endcsname{ega4.c04.cor.18-8-13}
\expandafter\def\csname sga@target@IV.18.8.14\endcsname{ega4.c04.cor.18-8-14}
\expandafter\def\csname sga@target@IV.18.8.15\endcsname{ega4.c04.prop.18-8-15}
\expandafter\def\csname sga@target@IV.18.8.16\endcsname{ega4.c04.cor.18-8-16}
\expandafter\def\csname sga@target@IV.18.8.17\endcsname{ega4.c04.prop.18-8-17}
\expandafter\def\csname sga@target@IV.18.8.18\endcsname{ega4.c04.prop.18-8-18}
\expandafter\def\csname sga@target@IV.18.8.2\endcsname{ega4.c04.def.18-8-2}
\expandafter\def\csname sga@target@IV.18.8.3\endcsname{ega4.c04.rem.18-8-3}
\expandafter\def\csname sga@target@IV.18.8.4\endcsname{ega4.c04.lem.18-8-4}
\expandafter\def\csname sga@target@IV.18.8.5\endcsname{ega4.c04.lem.18-8-5}
\expandafter\def\csname sga@target@IV.18.8.6\endcsname{ega4.c04.statement.18-8-6}
\expandafter\def\csname sga@target@IV.18.8.6.1\endcsname{ega4.c04.eq.18-8-6-1}
\expandafter\def\csname sga@target@IV.18.8.6.2\endcsname{ega4.c04.eq.18-8-6-2}
\expandafter\def\csname sga@target@IV.18.8.7\endcsname{ega4.c04.def.18-8-7}
\expandafter\def\csname sga@target@IV.18.8.8\endcsname{ega4.c04.prop.18-8-8}
\expandafter\def\csname sga@target@IV.18.8.8.1\endcsname{ega4.c04.statement.18-8-8-1}
\expandafter\def\csname sga@target@IV.18.8.9\endcsname{ega4.c04.statement.18-8-9}
\expandafter\def\csname sga@target@IV.18.9.1\endcsname{ega4.c04.thm.18-9-1}
\expandafter\def\csname sga@target@IV.18.9.10\endcsname{ega4.c04.rem.18-9-10}
\expandafter\def\csname sga@target@IV.18.9.11\endcsname{ega4.c04.cor.18-9-11}
\expandafter\def\csname sga@target@IV.18.9.2\endcsname{ega4.c04.cor.18-9-2}
\expandafter\def\csname sga@target@IV.18.9.3\endcsname{ega4.c04.cor.18-9-3}
\expandafter\def\csname sga@target@IV.18.9.4\endcsname{ega4.c04.cor.18-9-4}
\expandafter\def\csname sga@target@IV.18.9.5\endcsname{ega4.c04.cor.18-9-5}
\expandafter\def\csname sga@target@IV.18.9.5.1\endcsname{ega4.c04.eq.18-9-5-1}
\expandafter\def\csname sga@target@IV.18.9.6\endcsname{ega4.c04.statement.18-9-6}
\expandafter\def\csname sga@target@IV.18.9.7\endcsname{ega4.c04.thm.18-9-7}
\expandafter\def\csname sga@target@IV.18.9.7.1\endcsname{ega4.c04.rem.18-9-7-1}
\expandafter\def\csname sga@target@IV.18.9.7.2\endcsname{ega4.c04.statement.18-9-7-2}
\expandafter\def\csname sga@target@IV.18.9.7.3\endcsname{ega4.c04.lem.18-9-7-3}
\expandafter\def\csname sga@target@IV.18.9.7.4\endcsname{ega4.c04.eq.18-9-7-4}
\expandafter\def\csname sga@target@IV.18.9.7.5\endcsname{ega4.c04.lem.18-9-7-5}
\expandafter\def\csname sga@target@IV.18.9.7.6\endcsname{ega4.c04.rem.18-9-7-6}
\expandafter\def\csname sga@target@IV.18.9.7.7\endcsname{ega4.c04.lem.18-9-7-7}
\expandafter\def\csname sga@target@IV.18.9.7.8\endcsname{ega4.c04.rem.18-9-7-8}
\expandafter\def\csname sga@target@IV.18.9.8\endcsname{ega4.c04.cor.18-9-8}
\expandafter\def\csname sga@target@IV.18.9.9\endcsname{ega4.c04.cor.18-9-9}
\expandafter\def\csname sga@target@IV.19.1.1\endcsname{ega4.c04.prop.19-1-1}
\expandafter\def\csname sga@target@IV.19.1.2\endcsname{ega4.c04.cor.19-1-2}
\expandafter\def\csname sga@target@IV.19.1.3\endcsname{ega4.c04.def.19-1-3}
\expandafter\def\csname sga@target@IV.19.1.4\endcsname{ega4.c04.prop.19-1-4}
\expandafter\def\csname sga@target@IV.19.1.4.1\endcsname{ega4.c04.eq.19-1-4-1}
\expandafter\def\csname sga@target@IV.19.1.5\endcsname{ega4.c04.prop.19-1-5}
\expandafter\def\csname sga@target@IV.19.1.5.1\endcsname{ega4.c04.item.19-1-5-1}
\expandafter\def\csname sga@target@IV.19.1.6\endcsname{ega4.c04.cor.19-1-6}
\expandafter\def\csname sga@target@IV.19.1.6.1\endcsname{ega4.c04.lem.19-1-6-1}
\expandafter\def\csname sga@target@IV.19.1.7\endcsname{ega4.c04.item.19-1-7}
\expandafter\def\csname sga@target@IV.19.1.8\endcsname{ega4.c04.statement.19-1-8}
\expandafter\def\csname sga@target@IV.19.2.1\endcsname{ega4.c04.def.19-2-1}
\expandafter\def\csname sga@target@IV.19.2.2\endcsname{ega4.c04.def.19-2-2}
\expandafter\def\csname sga@target@IV.19.2.3\endcsname{ega4.c04.prop.19-2-3}
\expandafter\def\csname sga@target@IV.19.2.4\endcsname{ega4.c04.prop.19-2-4}
\expandafter\def\csname sga@target@IV.19.2.5\endcsname{ega4.c04.item.19-2-5}
\expandafter\def\csname sga@target@IV.19.2.6\endcsname{ega4.c04.cor.19-2-6}
\expandafter\def\csname sga@target@IV.19.2.7\endcsname{ega4.c04.prop.19-2-7}
\expandafter\def\csname sga@target@IV.19.2.8\endcsname{ega4.c04.prop.19-2-8}
\expandafter\def\csname sga@target@IV.19.2.9\endcsname{ega4.c04.prop.19-2-9}
\expandafter\def\csname sga@target@IV.19.3.1\endcsname{ega4.c04.def.19-3-1}
\expandafter\def\csname sga@target@IV.19.3.10\endcsname{ega4.c04.cor.19-3-10}
\expandafter\def\csname sga@target@IV.19.3.2\endcsname{ega4.c04.prop.19-3-2}
\expandafter\def\csname sga@target@IV.19.3.2.1\endcsname{ega4.c04.lem.19-3-2-1}
\expandafter\def\csname sga@target@IV.19.3.3\endcsname{ega4.c04.cor.19-3-3}
\expandafter\def\csname sga@target@IV.19.3.4\endcsname{ega4.c04.cor.19-3-4}
\expandafter\def\csname sga@target@IV.19.3.5\endcsname{ega4.c04.cor.19-3-5}
\expandafter\def\csname sga@target@IV.19.3.6\endcsname{ega4.c04.def.19-3-6}
\expandafter\def\csname sga@target@IV.19.3.7\endcsname{ega4.c04.prop.19-3-7}
\expandafter\def\csname sga@target@IV.19.3.8\endcsname{ega4.c04.cor.19-3-8}
\expandafter\def\csname sga@target@IV.19.3.9\endcsname{ega4.c04.prop.19-3-9}
\expandafter\def\csname sga@target@IV.19.4.1\endcsname{ega4.c04.statement.19-4-1}
\expandafter\def\csname sga@target@IV.19.4.1.1\endcsname{ega4.c04.subsec.19-4-1-1}
\expandafter\def\csname sga@target@IV.19.4.10\endcsname{ega4.c04.rem.19-4-10}
\expandafter\def\csname sga@target@IV.19.4.11\endcsname{ega4.c04.prop.19-4-11}
\expandafter\def\csname sga@target@IV.19.4.11.1\endcsname{ega4.c04.eq.19-4-11-1}
\expandafter\def\csname sga@target@IV.19.4.2\endcsname{ega4.c04.lem.19-4-2}
\expandafter\def\csname sga@target@IV.19.4.3\endcsname{ega4.c04.cor.19-4-3}
\expandafter\def\csname sga@target@IV.19.4.4\endcsname{ega4.c04.prop.19-4-4}
\expandafter\def\csname sga@target@IV.19.4.5\endcsname{ega4.c04.cor.19-4-5}
\expandafter\def\csname sga@target@IV.19.4.6\endcsname{ega4.c04.prop.19-4-6}
\expandafter\def\csname sga@target@IV.19.4.7\endcsname{ega4.c04.cor.19-4-7}
\expandafter\def\csname sga@target@IV.19.4.8\endcsname{ega4.c04.prop.19-4-8}
\expandafter\def\csname sga@target@IV.19.4.9\endcsname{ega4.c04.cor.19-4-9}
\expandafter\def\csname sga@target@IV.19.5.1\endcsname{ega4.c04.thm.19-5-1}
\expandafter\def\csname sga@target@IV.19.5.1.1\endcsname{ega4.c04.eq.19-5-1-1}
\expandafter\def\csname sga@target@IV.19.5.2\endcsname{ega4.c04.cor.19-5-2}
\expandafter\def\csname sga@target@IV.19.5.3\endcsname{ega4.c04.prop.19-5-3}
\expandafter\def\csname sga@target@IV.19.5.4\endcsname{ega4.c04.statement.19-5-4}
\expandafter\def\csname sga@target@IV.19.5.5\endcsname{ega4.c04.thm.19-5-5}
\expandafter\def\csname sga@target@IV.19.5.5.1\endcsname{ega4.c04.statement.19-5-5-1}
\expandafter\def\csname sga@target@IV.19.5.5.2\endcsname{ega4.c04.statement.19-5-5-2}
\expandafter\def\csname sga@target@IV.19.5.5.3\endcsname{ega4.c04.statement.19-5-5-3}
\expandafter\def\csname sga@target@IV.19.5.5.4\endcsname{ega4.c04.statement.19-5-5-4}
\expandafter\def\csname sga@target@IV.19.5.5.5\endcsname{ega4.c04.statement.19-5-5-5}
\expandafter\def\csname sga@target@IV.19.5.6\endcsname{ega4.c04.statement.19-5-6}
\expandafter\def\csname sga@target@IV.19.6.1\endcsname{ega4.c04.statement.19-6-1}
\expandafter\def\csname sga@target@IV.19.6.1.1\endcsname{ega4.c04.eq.19-6-1-1}
\expandafter\def\csname sga@target@IV.19.6.2\endcsname{ega4.c04.statement.19-6-2}
\expandafter\def\csname sga@target@IV.19.6.2.1\endcsname{ega4.c04.statement.19-6-2-1}
\expandafter\def\csname sga@target@IV.19.6.3\endcsname{ega4.c04.prop.19-6-3}
\expandafter\def\csname sga@target@IV.19.6.4\endcsname{ega4.c04.cor.19-6-4}
\expandafter\def\csname sga@target@IV.19.7\endcsname{ega4.c04.subsec.19-7.key-iv-19-7}
\expandafter\def\csname sga@target@IV.19.7.1\endcsname{ega4.c04.thm.19-7-1}
\expandafter\def\csname sga@target@IV.19.7.1.1\endcsname{ega4.c04.statement.19-7-1-1}
\expandafter\def\csname sga@target@IV.19.7.1.10\endcsname{ega4.c04.statement.19-7-1-10}
\expandafter\def\csname sga@target@IV.19.7.1.11\endcsname{ega4.c04.statement.19-7-1-11}
\expandafter\def\csname sga@target@IV.19.7.1.12\endcsname{ega4.c04.statement.19-7-1-12}
\expandafter\def\csname sga@target@IV.19.7.1.13\endcsname{ega4.c04.statement.19-7-1-13}
\expandafter\def\csname sga@target@IV.19.7.1.14\endcsname{ega4.c04.statement.19-7-1-14}
\expandafter\def\csname sga@target@IV.19.7.1.15\endcsname{ega4.c04.statement.19-7-1-15}
\expandafter\def\csname sga@target@IV.19.7.1.16\endcsname{ega4.c04.statement.19-7-1-16}
\expandafter\def\csname sga@target@IV.19.7.1.17\endcsname{ega4.c04.statement.19-7-1-17}
\expandafter\def\csname sga@target@IV.19.7.1.18\endcsname{ega4.c04.statement.19-7-1-18}
\expandafter\def\csname sga@target@IV.19.7.1.19\endcsname{ega4.c04.statement.19-7-1-19}
\expandafter\def\csname sga@target@IV.19.7.1.2\endcsname{ega4.c04.statement.19-7-1-2}
\expandafter\def\csname sga@target@IV.19.7.1.20\endcsname{ega4.c04.statement.19-7-1-20}
\expandafter\def\csname sga@target@IV.19.7.1.21\endcsname{ega4.c04.statement.19-7-1-21}
\expandafter\def\csname sga@target@IV.19.7.1.22\endcsname{ega4.c04.statement.19-7-1-22}
\expandafter\def\csname sga@target@IV.19.7.1.23\endcsname{ega4.c04.statement.19-7-1-23}
\expandafter\def\csname sga@target@IV.19.7.1.3\endcsname{ega4.c04.statement.19-7-1-3}
\expandafter\def\csname sga@target@IV.19.7.1.5\endcsname{ega4.c04.statement.19-7-1-5}
\expandafter\def\csname sga@target@IV.19.7.1.6\endcsname{ega4.c04.statement.19-7-1-6}
\expandafter\def\csname sga@target@IV.19.7.1.7\endcsname{ega4.c04.lem.19-7-1-7}
\expandafter\def\csname sga@target@IV.19.7.1.8\endcsname{ega4.c04.statement.19-7-1-8}
\expandafter\def\csname sga@target@IV.19.7.1.9\endcsname{ega4.c04.statement.19-7-1-9}
\expandafter\def\csname sga@target@IV.19.7.2\endcsname{ega4.c04.statement.19-7-2}
\expandafter\def\csname sga@target@IV.19.7.3\endcsname{ega4.c04.cor.19-7-3}
\expandafter\def\csname sga@target@IV.19.7.3.1\endcsname{ega4.c04.eq.19-7-3-1}
\expandafter\def\csname sga@target@IV.19.7.4\endcsname{ega4.c04.cor.19-7-4}
\expandafter\def\csname sga@target@IV.19.7.4.1\endcsname{ega4.c04.eq.19-7-4-1}
\expandafter\def\csname sga@target@IV.19.7.5\endcsname{ega4.c04.cor.19-7-5}
\expandafter\def\csname sga@target@IV.19.8.1\endcsname{ega4.c04.prop.19-8-1}
\expandafter\def\csname sga@target@IV.19.8.2\endcsname{ega4.c04.prop.19-8-2}
\expandafter\def\csname sga@target@IV.19.8.3\endcsname{ega4.c04.rem.19-8-3}
\expandafter\def\csname sga@target@IV.19.9.1\endcsname{ega4.c04.statement.19-9-1}
\expandafter\def\csname sga@target@IV.19.9.1.1\endcsname{ega4.c04.eq.19-9-1-1}
\expandafter\def\csname sga@target@IV.19.9.1.2\endcsname{ega4.c04.eq.19-9-1-2}
\expandafter\def\csname sga@target@IV.19.9.1.3\endcsname{ega4.c04.eq.19-9-1-3}
\expandafter\def\csname sga@target@IV.19.9.2\endcsname{ega4.c04.lem.19-9-2}
\expandafter\def\csname sga@target@IV.19.9.3\endcsname{ega4.c04.prop.19-9-3}
\expandafter\def\csname sga@target@IV.19.9.4\endcsname{ega4.c04.cor.19-9-4}
\expandafter\def\csname sga@target@IV.19.9.5\endcsname{ega4.c04.prop.19-9-5}
\expandafter\def\csname sga@target@IV.19.9.6\endcsname{ega4.c04.prop.19-9-6}
\expandafter\def\csname sga@target@IV.19.9.7\endcsname{ega4.c04.cor.19-9-7}
\expandafter\def\csname sga@target@IV.19.9.8\endcsname{ega4.c04.prop.19-9-8}
\expandafter\def\csname sga@target@IV.19.9.9\endcsname{ega4.c04.rem.19-9-9}
\expandafter\def\csname sga@target@IV.2.1.1\endcsname{ega4.c04.statement.2-1-1}
\expandafter\def\csname sga@target@IV.2.1.10\endcsname{ega4.c04.cor.2-1-10}
\expandafter\def\csname sga@target@IV.2.1.11\endcsname{ega4.c04.prop.2-1-11}
\expandafter\def\csname sga@target@IV.2.1.12\endcsname{ega4.c04.prop.2-1-12}
\expandafter\def\csname sga@target@IV.2.1.13\endcsname{ega4.c04.prop.2-1-13}
\expandafter\def\csname sga@target@IV.2.1.14\endcsname{ega4.c04.prop.2-1-14}
\expandafter\def\csname sga@target@IV.2.1.2\endcsname{ega4.c04.prop.2-1-2}
\expandafter\def\csname sga@target@IV.2.1.3\endcsname{ega4.c04.prop.2-1-3}
\expandafter\def\csname sga@target@IV.2.1.4\endcsname{ega4.c04.prop.2-1-4}
\expandafter\def\csname sga@target@IV.2.1.5\endcsname{ega4.c04.prop.2-1-5}
\expandafter\def\csname sga@target@IV.2.1.6\endcsname{ega4.c04.cor.2-1-6}
\expandafter\def\csname sga@target@IV.2.1.7\endcsname{ega4.c04.cor.2-1-7}
\expandafter\def\csname sga@target@IV.2.1.8\endcsname{ega4.c04.prop.2-1-8}
\expandafter\def\csname sga@target@IV.2.1.9\endcsname{ega4.c04.cor.2-1-9}
\expandafter\def\csname sga@target@IV.2.2.1\endcsname{ega4.c04.prop.2-2-1}
\expandafter\def\csname sga@target@IV.2.2.10\endcsname{ega4.c04.prop.2-2-10}
\expandafter\def\csname sga@target@IV.2.2.10.1\endcsname{ega4.c04.eq.2-2-10-1}
\expandafter\def\csname sga@target@IV.2.2.11\endcsname{ega4.c04.cor.2-2-11}
\expandafter\def\csname sga@target@IV.2.2.12\endcsname{ega4.c04.cor.2-2-12}
\expandafter\def\csname sga@target@IV.2.2.13\endcsname{ega4.c04.cor.2-2-13}
\expandafter\def\csname sga@target@IV.2.2.14\endcsname{ega4.c04.prop.2-2-14}
\expandafter\def\csname sga@target@IV.2.2.15\endcsname{ega4.c04.prop.2-2-15}
\expandafter\def\csname sga@target@IV.2.2.16\endcsname{ega4.c04.cor.2-2-16}
\expandafter\def\csname sga@target@IV.2.2.17\endcsname{ega4.c04.prop.2-2-17}
\expandafter\def\csname sga@target@IV.2.2.2\endcsname{ega4.c04.prop.2-2-2}
\expandafter\def\csname sga@target@IV.2.2.2.1\endcsname{ega4.c04.eq.2-2-2-1}
\expandafter\def\csname sga@target@IV.2.2.3\endcsname{ega4.c04.cor.2-2-3}
\expandafter\def\csname sga@target@IV.2.2.4\endcsname{ega4.c04.def.2-2-4}
\expandafter\def\csname sga@target@IV.2.2.5\endcsname{ega4.c04.cor.2-2-5}
\expandafter\def\csname sga@target@IV.2.2.6\endcsname{ega4.c04.cor.2-2-6}
\expandafter\def\csname sga@target@IV.2.2.7\endcsname{ega4.c04.prop.2-2-7}
\expandafter\def\csname sga@target@IV.2.2.8\endcsname{ega4.c04.cor.2-2-8}
\expandafter\def\csname sga@target@IV.2.2.8.1\endcsname{ega4.c04.eq.2-2-8-1}
\expandafter\def\csname sga@target@IV.2.2.9\endcsname{ega4.c04.rem.2-2-9}
\expandafter\def\csname sga@target@IV.2.3.1\endcsname{ega4.c04.lem.2-3-1}
\expandafter\def\csname sga@target@IV.2.3.1.1\endcsname{ega4.c04.eq.2-3-1-1}
\expandafter\def\csname sga@target@IV.2.3.10\endcsname{ega4.c04.thm.2-3-10}
\expandafter\def\csname sga@target@IV.2.3.11\endcsname{ega4.c04.cor.2-3-11}
\expandafter\def\csname sga@target@IV.2.3.12\endcsname{ega4.c04.cor.2-3-12}
\expandafter\def\csname sga@target@IV.2.3.13\endcsname{ega4.c04.cor.2-3-13}
\expandafter\def\csname sga@target@IV.2.3.14\endcsname{ega4.c04.cor.2-3-14}
\expandafter\def\csname sga@target@IV.2.3.15\endcsname{ega4.c04.rem.2-3-15}
\expandafter\def\csname sga@target@IV.2.3.2\endcsname{ega4.c04.prop.2-3-2}
\expandafter\def\csname sga@target@IV.2.3.3\endcsname{ega4.c04.statement.2-3-3}
\expandafter\def\csname sga@target@IV.2.3.4\endcsname{ega4.c04.prop.2-3-4}
\expandafter\def\csname sga@target@IV.2.3.5\endcsname{ega4.c04.cor.2-3-5}
\expandafter\def\csname sga@target@IV.2.3.6\endcsname{ega4.c04.prop.2-3-6}
\expandafter\def\csname sga@target@IV.2.3.7\endcsname{ega4.c04.prop.2-3-7}
\expandafter\def\csname sga@target@IV.2.3.8\endcsname{ega4.c04.cor.2-3-8}
\expandafter\def\csname sga@target@IV.2.3.9\endcsname{ega4.c04.rem.2-3-9}
\expandafter\def\csname sga@target@IV.2.4.1\endcsname{ega4.c04.statement.2-4-1}
\expandafter\def\csname sga@target@IV.2.4.10\endcsname{ega4.c04.cor.2-4-10}
\expandafter\def\csname sga@target@IV.2.4.11\endcsname{ega4.c04.rem.2-4-11}
\expandafter\def\csname sga@target@IV.2.4.2\endcsname{ega4.c04.def.2-4-2}
\expandafter\def\csname sga@target@IV.2.4.3\endcsname{ega4.c04.prop.2-4-3}
\expandafter\def\csname sga@target@IV.2.4.4\endcsname{ega4.c04.prop.2-4-4}
\expandafter\def\csname sga@target@IV.2.4.5\endcsname{ega4.c04.prop.2-4-5}
\expandafter\def\csname sga@target@IV.2.4.6\endcsname{ega4.c04.thm.2-4-6}
\expandafter\def\csname sga@target@IV.2.4.7\endcsname{ega4.c04.cor.2-4-7}
\expandafter\def\csname sga@target@IV.2.4.8\endcsname{ega4.c04.rem.2-4-8}
\expandafter\def\csname sga@target@IV.2.4.9\endcsname{ega4.c04.prop.2-4-9}
\expandafter\def\csname sga@target@IV.2.5.1\endcsname{ega4.c04.prop.2-5-1}
\expandafter\def\csname sga@target@IV.2.5.2\endcsname{ega4.c04.prop.2-5-2}
\expandafter\def\csname sga@target@IV.2.5.3\endcsname{ega4.c04.lem.2-5-3}
\expandafter\def\csname sga@target@IV.2.5.4\endcsname{ega4.c04.rem.2-5-4}
\expandafter\def\csname sga@target@IV.2.5.4.1\endcsname{ega4.c04.lem.2-5-4-1}
\expandafter\def\csname sga@target@IV.2.5.5\endcsname{ega4.c04.prop.2-5-5}
\expandafter\def\csname sga@target@IV.2.5.5.1\endcsname{ega4.c04.statement.2-5-5-1}
\expandafter\def\csname sga@target@IV.2.5.5.2\endcsname{ega4.c04.lem.2-5-5-2}
\expandafter\def\csname sga@target@IV.2.5.5.3\endcsname{ega4.c04.statement.2-5-5-3}
\expandafter\def\csname sga@target@IV.2.5.5.4\endcsname{ega4.c04.rem.2-5-5-4}
\expandafter\def\csname sga@target@IV.2.5.6\endcsname{ega4.c04.lem.2-5-6}
\expandafter\def\csname sga@target@IV.2.5.7\endcsname{ega4.c04.cor.2-5-7}
\expandafter\def\csname sga@target@IV.2.5.8\endcsname{ega4.c04.prop.2-5-8}
\expandafter\def\csname sga@target@IV.2.5.8.1\endcsname{ega4.c04.statement.2-5-8-1}
\expandafter\def\csname sga@target@IV.2.5.8.2\endcsname{ega4.c04.statement.2-5-8-2}
\expandafter\def\csname sga@target@IV.2.5.8.3\endcsname{ega4.c04.statement.2-5-8-3}
\expandafter\def\csname sga@target@IV.2.5.8.4\endcsname{ega4.c04.statement.2-5-8-4}
\expandafter\def\csname sga@target@IV.2.5.8.5\endcsname{ega4.c04.statement.2-5-8-5}
\expandafter\def\csname sga@target@IV.2.5.8.6\endcsname{ega4.c04.statement.2-5-8-6}
\expandafter\def\csname sga@target@IV.2.5.8.7\endcsname{ega4.c04.statement.2-5-8-7}
\expandafter\def\csname sga@target@IV.2.5.8.8\endcsname{ega4.c04.statement.2-5-8-8}
\expandafter\def\csname sga@target@IV.2.5.8.9\endcsname{ega4.c04.statement.2-5-8-9}
\expandafter\def\csname sga@target@IV.2.6.1\endcsname{ega4.c04.prop.2-6-1}
\expandafter\def\csname sga@target@IV.2.6.2\endcsname{ega4.c04.prop.2-6-2}
\expandafter\def\csname sga@target@IV.2.6.3\endcsname{ega4.c04.rem.2-6-3}
\expandafter\def\csname sga@target@IV.2.6.4\endcsname{ega4.c04.cor.2-6-4}
\expandafter\def\csname sga@target@IV.2.7.1\endcsname{ega4.c04.prop.2-7-1}
\expandafter\def\csname sga@target@IV.2.7.1.1\endcsname{ega4.c04.lem.2-7-1-1}
\expandafter\def\csname sga@target@IV.2.7.1.viii\endcsname{ega4.c04.item.2-7-1-viii}
\expandafter\def\csname sga@target@IV.2.7.2\endcsname{ega4.c04.cor.2-7-2}
\expandafter\def\csname sga@target@IV.2.7.3\endcsname{ega4.c04.sec.2-7-3}
\expandafter\def\csname sga@target@IV.2.8.1\endcsname{ega4.c04.prop.2-8-1}
\expandafter\def\csname sga@target@IV.2.8.1.1\endcsname{ega4.c04.lem.2-8-1-1}
\expandafter\def\csname sga@target@IV.2.8.2\endcsname{ega4.c04.cor.2-8-2}
\expandafter\def\csname sga@target@IV.2.8.3\endcsname{ega4.c04.statement.2-8-3}
\expandafter\def\csname sga@target@IV.2.8.4\endcsname{ega4.c04.prop.2-8-4}
\expandafter\def\csname sga@target@IV.2.8.4.1\endcsname{ega4.c04.statement.2-8-4-1}
\expandafter\def\csname sga@target@IV.2.8.5\endcsname{ega4.c04.prop.2-8-5}
\expandafter\def\csname sga@target@IV.2.8.6\endcsname{ega4.c04.cor.2-8-6}
\expandafter\def\csname sga@target@IV.2.8.6.1\endcsname{ega4.c04.statement.2-8-6-1}
\expandafter\def\csname sga@target@IV.20\endcsname{ega4.c04.sec.20}
\expandafter\def\csname sga@target@IV.20.0\endcsname{ega4.c04.subsec.20-0}
\expandafter\def\csname sga@target@IV.20.1\endcsname{ega4.c04.subsec.20-1}
\expandafter\def\csname sga@target@IV.20.1.1\endcsname{ega4.c04.statement.20-1-1}
\expandafter\def\csname sga@target@IV.20.1.1.1\endcsname{ega4.c04.eq.20-1-1-1}
\expandafter\def\csname sga@target@IV.20.1.10\endcsname{ega4.c04.statement.20-1-10}
\expandafter\def\csname sga@target@IV.20.1.11\endcsname{ega4.c04.statement.20-1-11}
\expandafter\def\csname sga@target@IV.20.1.11.1\endcsname{ega4.c04.eq.20-1-11-1}
\expandafter\def\csname sga@target@IV.20.1.11.2\endcsname{ega4.c04.eq.20-1-11-2}
\expandafter\def\csname sga@target@IV.20.1.12\endcsname{ega4.c04.prop.20-1-12}
\expandafter\def\csname sga@target@IV.20.1.13\endcsname{ega4.c04.rem.20-1-13}
\expandafter\def\csname sga@target@IV.20.1.2\endcsname{ega4.c04.statement.20-1-2}
\expandafter\def\csname sga@target@IV.20.1.2.1\endcsname{ega4.c04.statement.20-1-2-1}
\expandafter\def\csname sga@target@IV.20.1.2.2\endcsname{ega4.c04.statement.20-1-2-2}
\expandafter\def\csname sga@target@IV.20.1.3\endcsname{ega4.c04.statement.20-1-3}
\expandafter\def\csname sga@target@IV.20.1.3.1\endcsname{ega4.c04.statement.20-1-3-1}
\expandafter\def\csname sga@target@IV.20.1.4\endcsname{ega4.c04.statement.20-1-4}
\expandafter\def\csname sga@target@IV.20.1.4.1\endcsname{ega4.c04.statement.20-1-4-1}
\expandafter\def\csname sga@target@IV.20.1.5\endcsname{ega4.c04.statement.20-1-5}
\expandafter\def\csname sga@target@IV.20.1.5.1\endcsname{ega4.c04.statement.20-1-5-1}
\expandafter\def\csname sga@target@IV.20.1.6\endcsname{ega4.c04.prop.20-1-6}
\expandafter\def\csname sga@target@IV.20.1.7\endcsname{ega4.c04.statement.20-1-7}
\expandafter\def\csname sga@target@IV.20.1.8\endcsname{ega4.c04.statement.20-1-8}
\expandafter\def\csname sga@target@IV.20.1.9\endcsname{ega4.c04.statement.20-1-9}
\expandafter\def\csname sga@target@IV.20.2\endcsname{ega4.c04.subsec.20-2}
\expandafter\def\csname sga@target@IV.20.2.1\endcsname{ega4.c04.statement.20-2-1}
\expandafter\def\csname sga@target@IV.20.2.10\endcsname{ega4.c04.statement.20-2-10}
\expandafter\def\csname sga@target@IV.20.2.10.1\endcsname{ega4.c04.eq.20-2-10-1}
\expandafter\def\csname sga@target@IV.20.2.10.2\endcsname{ega4.c04.eq.20-2-10-2}
\expandafter\def\csname sga@target@IV.20.2.11\endcsname{ega4.c04.prop.20-2-11}
\expandafter\def\csname sga@target@IV.20.2.11.1\endcsname{ega4.c04.eq.20-2-11-1}
\expandafter\def\csname sga@target@IV.20.2.12\endcsname{ega4.c04.cor.20-2-12}
\expandafter\def\csname sga@target@IV.20.2.13\endcsname{ega4.c04.statement.20-2-13}
\expandafter\def\csname sga@target@IV.20.2.14\endcsname{ega4.c04.statement.20-2-14}
\expandafter\def\csname sga@target@IV.20.2.15\endcsname{ega4.c04.prop.20-2-15}
\expandafter\def\csname sga@target@IV.20.2.16\endcsname{ega4.c04.rem.20-2-16}
\expandafter\def\csname sga@target@IV.20.2.2\endcsname{ega4.c04.statement.20-2-2}
\expandafter\def\csname sga@target@IV.20.2.3\endcsname{ega4.c04.statement.20-2-3}
\expandafter\def\csname sga@target@IV.20.2.3.1\endcsname{ega4.c04.statement.20-2-3-1}
\expandafter\def\csname sga@target@IV.20.2.4\endcsname{ega4.c04.prop.20-2-4}
\expandafter\def\csname sga@target@IV.20.2.5\endcsname{ega4.c04.cor.20-2-5}
\expandafter\def\csname sga@target@IV.20.2.6\endcsname{ega4.c04.cor.20-2-6}
\expandafter\def\csname sga@target@IV.20.2.7\endcsname{ega4.c04.rem.20-2-7}
\expandafter\def\csname sga@target@IV.20.2.8\endcsname{ega4.c04.statement.20-2-8}
\expandafter\def\csname sga@target@IV.20.2.9\endcsname{ega4.c04.lem.20-2-9}
\expandafter\def\csname sga@target@IV.20.3\endcsname{ega4.c04.subsec.20-3}
\expandafter\def\csname sga@target@IV.20.3.1\endcsname{ega4.c04.statement.20-3-1}
\expandafter\def\csname sga@target@IV.20.3.10\endcsname{ega4.c04.lem.20-3-10}
\expandafter\def\csname sga@target@IV.20.3.11\endcsname{ega4.c04.prop.20-3-11}
\expandafter\def\csname sga@target@IV.20.3.11.1\endcsname{ega4.c04.eq.20-3-11-1}
\expandafter\def\csname sga@target@IV.20.3.12\endcsname{ega4.c04.cor.20-3-12}
\expandafter\def\csname sga@target@IV.20.3.13\endcsname{ega4.c04.rem.20-3-13}
\expandafter\def\csname sga@target@IV.20.3.2\endcsname{ega4.c04.statement.20-3-2}
\expandafter\def\csname sga@target@IV.20.3.2.1\endcsname{ega4.c04.eq.20-3-2-1}
\expandafter\def\csname sga@target@IV.20.3.3\endcsname{ega4.c04.statement.20-3-3}
\expandafter\def\csname sga@target@IV.20.3.4\endcsname{ega4.c04.statement.20-3-4}
\expandafter\def\csname sga@target@IV.20.3.5\endcsname{ega4.c04.prop.20-3-5}
\expandafter\def\csname sga@target@IV.20.3.6\endcsname{ega4.c04.statement.20-3-6}
\expandafter\def\csname sga@target@IV.20.3.6.1\endcsname{ega4.c04.eq.20-3-6-1}
\expandafter\def\csname sga@target@IV.20.3.7\endcsname{ega4.c04.prop.20-3-7}
\expandafter\def\csname sga@target@IV.20.3.8\endcsname{ega4.c04.prop.20-3-8}
\expandafter\def\csname sga@target@IV.20.3.8.1\endcsname{ega4.c04.eq.20-3-8-1}
\expandafter\def\csname sga@target@IV.20.3.8.2\endcsname{ega4.c04.lem.20-3-8-2}
\expandafter\def\csname sga@target@IV.20.3.8.3\endcsname{ega4.c04.lem.20-3-8-3}
\expandafter\def\csname sga@target@IV.20.3.8.4\endcsname{ega4.c04.lem.20-3-8-4}
\expandafter\def\csname sga@target@IV.20.3.8.5\endcsname{ega4.c04.sec.20-3-8-5}
\expandafter\def\csname sga@target@IV.20.3.8.6\endcsname{ega4.c04.sec.20-3-8-6}
\expandafter\def\csname sga@target@IV.20.3.8.7\endcsname{ega4.c04.sec.20-3-8-7}
\expandafter\def\csname sga@target@IV.20.3.8.8\endcsname{ega4.c04.lem.20-3-8-8}
\expandafter\def\csname sga@target@IV.20.3.9\endcsname{ega4.c04.cor.20-3-9}
\expandafter\def\csname sga@target@IV.20.4\endcsname{ega4.c04.subsec.20-4}
\expandafter\def\csname sga@target@IV.20.4.1\endcsname{ega4.c04.statement.20-4-1}
\expandafter\def\csname sga@target@IV.20.4.10\endcsname{ega4.c04.cor.20-4-10}
\expandafter\def\csname sga@target@IV.20.4.11\endcsname{ega4.c04.prop.20-4-11}
\expandafter\def\csname sga@target@IV.20.4.12\endcsname{ega4.c04.cor.20-4-12}
\expandafter\def\csname sga@target@IV.20.4.2\endcsname{ega4.c04.statement.20-4-2}
\expandafter\def\csname sga@target@IV.20.4.3\endcsname{ega4.c04.statement.20-4-3}
\expandafter\def\csname sga@target@IV.20.4.4\endcsname{ega4.c04.prop.20-4-4}
\expandafter\def\csname sga@target@IV.20.4.5\endcsname{ega4.c04.prop.20-4-5}
\expandafter\def\csname sga@target@IV.20.4.6\endcsname{ega4.c04.prop.20-4-6}
\expandafter\def\csname sga@target@IV.20.4.7\endcsname{ega4.c04.cor.20-4-7}
\expandafter\def\csname sga@target@IV.20.4.8\endcsname{ega4.c04.rem.20-4-8}
\expandafter\def\csname sga@target@IV.20.4.9\endcsname{ega4.c04.cor.20-4-9}
\expandafter\def\csname sga@target@IV.20.5\endcsname{ega4.c04.subsec.20-5}
\expandafter\def\csname sga@target@IV.20.5.1\endcsname{ega4.c04.statement.20-5-1}
\expandafter\def\csname sga@target@IV.20.5.2\endcsname{ega4.c04.statement.20-5-2}
\expandafter\def\csname sga@target@IV.20.5.2.1\endcsname{ega4.c04.statement.20-5-2-1}
\expandafter\def\csname sga@target@IV.20.5.3\endcsname{ega4.c04.prop.20-5-3}
\expandafter\def\csname sga@target@IV.20.5.4\endcsname{ega4.c04.statement.20-5-4}
\expandafter\def\csname sga@target@IV.20.5.5\endcsname{ega4.c04.statement.20-5-5}
\expandafter\def\csname sga@target@IV.20.5.6\endcsname{ega4.c04.statement.20-5-6}
\expandafter\def\csname sga@target@IV.20.5.7\endcsname{ega4.c04.statement.20-5-7}
\expandafter\def\csname sga@target@IV.20.5.8\endcsname{ega4.c04.statement.20-5-8}
\expandafter\def\csname sga@target@IV.20.6\endcsname{ega4.c04.subsec.20-6}
\expandafter\def\csname sga@target@IV.20.6.1\endcsname{ega4.c04.statement.20-6-1}
\expandafter\def\csname sga@target@IV.20.6.1.1\endcsname{ega4.c04.eq.20-6-1-1}
\expandafter\def\csname sga@target@IV.20.6.2\endcsname{ega4.c04.statement.20-6-2}
\expandafter\def\csname sga@target@IV.20.6.2.1\endcsname{ega4.c04.statement.20-6-2-1}
\expandafter\def\csname sga@target@IV.20.6.2.2\endcsname{ega4.c04.eq.20-6-2-2}
\expandafter\def\csname sga@target@IV.20.6.3\endcsname{ega4.c04.prop.20-6-3}
\expandafter\def\csname sga@target@IV.20.6.4\endcsname{ega4.c04.prop.20-6-4}
\expandafter\def\csname sga@target@IV.20.6.5\endcsname{ega4.c04.statement.20-6-5}
\expandafter\def\csname sga@target@IV.20.6.5.1\endcsname{ega4.c04.eq.20-6-5-1}
\expandafter\def\csname sga@target@IV.20.6.6\endcsname{ega4.c04.prop.20-6-6}
\expandafter\def\csname sga@target@IV.20.6.7\endcsname{ega4.c04.statement.20-6-7}
\expandafter\def\csname sga@target@IV.20.6.8\endcsname{ega4.c04.prop.20-6-8}
\expandafter\def\csname sga@target@IV.20.6.8.1\endcsname{ega4.c04.eq.20-6-8-1}
\expandafter\def\csname sga@target@IV.20.6.9\endcsname{ega4.c04.statement.20-6-9}
\expandafter\def\csname sga@target@IV.20.6.9.1\endcsname{ega4.c04.statement.20-6-9-1}
\expandafter\def\csname sga@target@IV.21.1\endcsname{ega4.c04.subsec.21-1.key-iv-21-1}
\expandafter\def\csname sga@target@IV.21.1.1\endcsname{ega4.c04.statement.21-1-1}
\expandafter\def\csname sga@target@IV.21.1.10\endcsname{ega4.c04.prop.21-1-10}
\expandafter\def\csname sga@target@IV.21.1.2\endcsname{ega4.c04.statement.21-1-2}
\expandafter\def\csname sga@target@IV.21.1.3\endcsname{ega4.c04.statement.21-1-3}
\expandafter\def\csname sga@target@IV.21.1.3.1\endcsname{ega4.c04.statement.21-1-3-1}
\expandafter\def\csname sga@target@IV.21.1.3.2\endcsname{ega4.c04.statement.21-1-3-2}
\expandafter\def\csname sga@target@IV.21.1.3.3\endcsname{ega4.c04.statement.21-1-3-3}
\expandafter\def\csname sga@target@IV.21.1.3.4\endcsname{ega4.c04.statement.21-1-3-4}
\expandafter\def\csname sga@target@IV.21.1.4\endcsname{ega4.c04.statement.21-1-4}
\expandafter\def\csname sga@target@IV.21.1.4.1\endcsname{ega4.c04.eq.21-1-4-1}
\expandafter\def\csname sga@target@IV.21.1.4.2\endcsname{ega4.c04.eq.21-1-4-2}
\expandafter\def\csname sga@target@IV.21.1.4.3\endcsname{ega4.c04.eq.21-1-4-3}
\expandafter\def\csname sga@target@IV.21.1.5\endcsname{ega4.c04.statement.21-1-5}
\expandafter\def\csname sga@target@IV.21.1.5.1\endcsname{ega4.c04.statement.21-1-5-1}
\expandafter\def\csname sga@target@IV.21.1.5.2\endcsname{ega4.c04.eq.21-1-5-2}
\expandafter\def\csname sga@target@IV.21.1.6\endcsname{ega4.c04.def.21-1-6}
\expandafter\def\csname sga@target@IV.21.1.6.1\endcsname{ega4.c04.eq.21-1-6-1}
\expandafter\def\csname sga@target@IV.21.1.6.2\endcsname{ega4.c04.eq.21-1-6-2}
\expandafter\def\csname sga@target@IV.21.1.6.3\endcsname{ega4.c04.eq.21-1-6-3}
\expandafter\def\csname sga@target@IV.21.1.6.4\endcsname{ega4.c04.eq.21-1-6-4}
\expandafter\def\csname sga@target@IV.21.1.7\endcsname{ega4.c04.statement.21-1-7}
\expandafter\def\csname sga@target@IV.21.1.7.1\endcsname{ega4.c04.statement.21-1-7-1}
\expandafter\def\csname sga@target@IV.21.1.8\endcsname{ega4.c04.prop.21-1-8}
\expandafter\def\csname sga@target@IV.21.1.9\endcsname{ega4.c04.cor.21-1-9}
\expandafter\def\csname sga@target@IV.21.10.1\endcsname{ega4.c04.statement.21-10-1}
\expandafter\def\csname sga@target@IV.21.10.1.1\endcsname{ega4.c04.eq.21-10-1-1}
\expandafter\def\csname sga@target@IV.21.10.1.2\endcsname{ega4.c04.eq.21-10-1-2}
\expandafter\def\csname sga@target@IV.21.10.10\endcsname{ega4.c04.statement.21-10-10}
\expandafter\def\csname sga@target@IV.21.10.11\endcsname{ega4.c04.statement.21-10-11}
\expandafter\def\csname sga@target@IV.21.10.11.1\endcsname{ega4.c04.eq.21-10-11-1}
\expandafter\def\csname sga@target@IV.21.10.12\endcsname{ega4.c04.rem.21-10-12}
\expandafter\def\csname sga@target@IV.21.10.13\endcsname{ega4.c04.lem.21-10-13}
\expandafter\def\csname sga@target@IV.21.10.13.1\endcsname{ega4.c04.prop.21-10-13-1}
\expandafter\def\csname sga@target@IV.21.10.13bis\endcsname{ega4.c04.prop.21-10-13bis}
\expandafter\def\csname sga@target@IV.21.10.14\endcsname{ega4.c04.statement.21-10-14}
\expandafter\def\csname sga@target@IV.21.10.14.1\endcsname{ega4.c04.eq.21-10-14-1}
\expandafter\def\csname sga@target@IV.21.10.14.2\endcsname{ega4.c04.eq.21-10-14-2}
\expandafter\def\csname sga@target@IV.21.10.15\endcsname{ega4.c04.prop.21-10-15}
\expandafter\def\csname sga@target@IV.21.10.16\endcsname{ega4.c04.prop.21-10-16}
\expandafter\def\csname sga@target@IV.21.10.16.1\endcsname{ega4.c04.prop.21-10-16-1}
\expandafter\def\csname sga@target@IV.21.10.16.2\endcsname{ega4.c04.eq.21-10-16-2}
\expandafter\def\csname sga@target@IV.21.10.17\endcsname{ega4.c04.prop.21-10-17}
\expandafter\def\csname sga@target@IV.21.10.17.1\endcsname{ega4.c04.prop.21-10-17-1}
\expandafter\def\csname sga@target@IV.21.10.17.10\endcsname{ega4.c04.eq.21-10-17-10}
\expandafter\def\csname sga@target@IV.21.10.17.2\endcsname{ega4.c04.eq.21-10-17-2}
\expandafter\def\csname sga@target@IV.21.10.17.3\endcsname{ega4.c04.lem.21-10-17-3}
\expandafter\def\csname sga@target@IV.21.10.17.4\endcsname{ega4.c04.lem.21-10-17-4}
\expandafter\def\csname sga@target@IV.21.10.17.5\endcsname{ega4.c04.eq.21-10-17-5}
\expandafter\def\csname sga@target@IV.21.10.17.6\endcsname{ega4.c04.eq.21-10-17-6}
\expandafter\def\csname sga@target@IV.21.10.17.7\endcsname{ega4.c04.lem.21-10-17-7}
\expandafter\def\csname sga@target@IV.21.10.17.8\endcsname{ega4.c04.lem.21-10-17-8}
\expandafter\def\csname sga@target@IV.21.10.17.9\endcsname{ega4.c04.eq.21-10-17-9}
\expandafter\def\csname sga@target@IV.21.10.18\endcsname{ega4.c04.prop.21-10-18}
\expandafter\def\csname sga@target@IV.21.10.18.1\endcsname{ega4.c04.prop.21-10-18-1}
\expandafter\def\csname sga@target@IV.21.10.19\endcsname{ega4.c04.cor.21-10-19}
\expandafter\def\csname sga@target@IV.21.10.19.1\endcsname{ega4.c04.cor.21-10-19-1}
\expandafter\def\csname sga@target@IV.21.10.2\endcsname{ega4.c04.statement.21-10-2}
\expandafter\def\csname sga@target@IV.21.10.3\endcsname{ega4.c04.statement.21-10-3}
\expandafter\def\csname sga@target@IV.21.10.3.1\endcsname{ega4.c04.statement.21-10-3-1}
\expandafter\def\csname sga@target@IV.21.10.4\endcsname{ega4.c04.prop.21-10-4}
\expandafter\def\csname sga@target@IV.21.10.4.1\endcsname{ega4.c04.prop.21-10-4-1}
\expandafter\def\csname sga@target@IV.21.10.5\endcsname{ega4.c04.statement.21-10-5}
\expandafter\def\csname sga@target@IV.21.10.6\endcsname{ega4.c04.prop.21-10-6}
\expandafter\def\csname sga@target@IV.21.10.7\endcsname{ega4.c04.rem.21-10-7}
\expandafter\def\csname sga@target@IV.21.10.8\endcsname{ega4.c04.prop.21-10-8}
\expandafter\def\csname sga@target@IV.21.10.9\endcsname{ega4.c04.statement.21-10-9}
\expandafter\def\csname sga@target@IV.21.11\endcsname{ega4.c04.subsec.21-11.key-iv-21-11}
\expandafter\def\csname sga@target@IV.21.11.1\endcsname{ega4.c04.thm.21-11-1}
\expandafter\def\csname sga@target@IV.21.11.1.1\endcsname{ega4.c04.eq.21-11-1-1}
\expandafter\def\csname sga@target@IV.21.11.2\endcsname{ega4.c04.lem.21-11-2}
\expandafter\def\csname sga@target@IV.21.12\endcsname{ega4.c04.subsec.21-12.key-iv-21-12}
\expandafter\def\csname sga@target@IV.21.12.1\endcsname{ega4.c04.statement.21-12-1}
\expandafter\def\csname sga@target@IV.21.12.1.1\endcsname{ega4.c04.eq.21-12-1-1}
\expandafter\def\csname sga@target@IV.21.12.10\endcsname{ega4.c04.prop.21-12-10}
\expandafter\def\csname sga@target@IV.21.12.12\endcsname{ega4.c04.thm.21-12-12}
\expandafter\def\csname sga@target@IV.21.12.13\endcsname{ega4.c04.cor.21-12-13}
\expandafter\def\csname sga@target@IV.21.12.14\endcsname{ega4.c04.item.21-12-14}
\expandafter\def\csname sga@target@IV.21.12.15\endcsname{ega4.c04.prop.21-12-15}
\expandafter\def\csname sga@target@IV.21.12.16\endcsname{ega4.c04.cor.21-12-16}
\expandafter\def\csname sga@target@IV.21.12.17\endcsname{ega4.c04.item.21-12-17}
\expandafter\def\csname sga@target@IV.21.12.2\endcsname{ega4.c04.prop.21-12-2}
\expandafter\def\csname sga@target@IV.21.12.2.1\endcsname{ega4.c04.prop.21-12-2-1}
\expandafter\def\csname sga@target@IV.21.12.3\endcsname{ega4.c04.cor.21-12-3}
\expandafter\def\csname sga@target@IV.21.12.4\endcsname{ega4.c04.eq.21-12-4}
\expandafter\def\csname sga@target@IV.21.12.5\endcsname{ega4.c04.statement.21-12-5}
\expandafter\def\csname sga@target@IV.21.12.5.1\endcsname{ega4.c04.eq.21-12-5-1}
\expandafter\def\csname sga@target@IV.21.12.5.2\endcsname{ega4.c04.eq.21-12-5-2}
\expandafter\def\csname sga@target@IV.21.12.5.3\endcsname{ega4.c04.eq.21-12-5-3}
\expandafter\def\csname sga@target@IV.21.12.6\endcsname{ega4.c04.thm.21-12-6}
\expandafter\def\csname sga@target@IV.21.12.7\endcsname{ega4.c04.cor.21-12-7}
\expandafter\def\csname sga@target@IV.21.12.8\endcsname{ega4.c04.statement.21-12-8}
\expandafter\def\csname sga@target@IV.21.12.9\endcsname{ega4.c04.statement.21-12-9}
\expandafter\def\csname sga@target@IV.21.13\endcsname{ega4.c04.subsec.21-13.key-iv-21-13}
\expandafter\def\csname sga@target@IV.21.13.1\endcsname{ega4.c04.def.21-13-1}
\expandafter\def\csname sga@target@IV.21.13.10\endcsname{ega4.c04.prop.21-13-10}
\expandafter\def\csname sga@target@IV.21.13.11\endcsname{ega4.c04.cor.21-13-11}
\expandafter\def\csname sga@target@IV.21.13.12\endcsname{ega4.c04.prop.21-13-12}
\expandafter\def\csname sga@target@IV.21.13.12.1\endcsname{ega4.c04.statement.21-13-12-1}
\expandafter\def\csname sga@target@IV.21.13.13\endcsname{ega4.c04.rem.21-13-13}
\expandafter\def\csname sga@target@IV.21.13.14\endcsname{ega4.c04.prop.21-13-14}
\expandafter\def\csname sga@target@IV.21.13.15\endcsname{ega4.c04.cor.21-13-15}
\expandafter\def\csname sga@target@IV.21.13.16\endcsname{ega4.c04.rem.21-13-16}
\expandafter\def\csname sga@target@IV.21.13.2\endcsname{ega4.c04.lem.21-13-2}
\expandafter\def\csname sga@target@IV.21.13.2.1\endcsname{ega4.c04.eq.21-13-2-1}
\expandafter\def\csname sga@target@IV.21.13.3\endcsname{ega4.c04.cor.21-13-3}
\expandafter\def\csname sga@target@IV.21.13.4\endcsname{ega4.c04.lem.21-13-4}
\expandafter\def\csname sga@target@IV.21.13.5\endcsname{ega4.c04.prop.21-13-5}
\expandafter\def\csname sga@target@IV.21.13.6\endcsname{ega4.c04.cor.21-13-6}
\expandafter\def\csname sga@target@IV.21.13.7\endcsname{ega4.c04.def.21-13-7}
\expandafter\def\csname sga@target@IV.21.13.8\endcsname{ega4.c04.prop.21-13-8}
\expandafter\def\csname sga@target@IV.21.13.9\endcsname{ega4.c04.statement.21-13-9}
\expandafter\def\csname sga@target@IV.21.13.9.1\endcsname{ega4.c04.eq.21-13-9-1}
\expandafter\def\csname sga@target@IV.21.14\endcsname{ega4.c04.subsec.21-14.key-iv-21-14}
\expandafter\def\csname sga@target@IV.21.14.1\endcsname{ega4.c04.thm.21-14-1}
\expandafter\def\csname sga@target@IV.21.14.2\endcsname{ega4.c04.cor.21-14-2}
\expandafter\def\csname sga@target@IV.21.14.3\endcsname{ega4.c04.prop.21-14-3}
\expandafter\def\csname sga@target@IV.21.14.4\endcsname{ega4.c04.statement.21-14-4}
\expandafter\def\csname sga@target@IV.21.14.4bis\endcsname{ega4.c04.statement.21-14-4bis}
\expandafter\def\csname sga@target@IV.21.15\endcsname{ega4.c04.subsec.21-15.key-iv-21-15}
\expandafter\def\csname sga@target@IV.21.15.1\endcsname{ega4.c04.statement.21-15-1}
\expandafter\def\csname sga@target@IV.21.15.2\endcsname{ega4.c04.def.21-15-2}
\expandafter\def\csname sga@target@IV.21.15.3\endcsname{ega4.c04.statement.21-15-3}
\expandafter\def\csname sga@target@IV.21.15.3.1\endcsname{ega4.c04.eq.21-15-3-1}
\expandafter\def\csname sga@target@IV.21.15.3.2\endcsname{ega4.c04.eq.21-15-3-2}
\expandafter\def\csname sga@target@IV.21.15.3bis\endcsname{ega4.c04.statement.21-15-3bis}
\expandafter\def\csname sga@target@IV.21.15.4\endcsname{ega4.c04.prop.21-15-4}
\expandafter\def\csname sga@target@IV.21.15.5\endcsname{ega4.c04.prop.21-15-5}
\expandafter\def\csname sga@target@IV.21.15.6\endcsname{ega4.c04.prop.21-15-6}
\expandafter\def\csname sga@target@IV.21.15.7\endcsname{ega4.c04.statement.21-15-7}
\expandafter\def\csname sga@target@IV.21.15.8\endcsname{ega4.c04.statement.21-15-8}
\expandafter\def\csname sga@target@IV.21.15.9\endcsname{ega4.c04.statement.21-15-9}
\expandafter\def\csname sga@target@IV.21.15.9.1\endcsname{ega4.c04.eq.21-15-9-1}
\expandafter\def\csname sga@target@IV.21.2.1\endcsname{ega4.c04.statement.21-2-1}
\expandafter\def\csname sga@target@IV.21.2.10\endcsname{ega4.c04.statement.21-2-10}
\expandafter\def\csname sga@target@IV.21.2.10.1\endcsname{ega4.c04.eq.21-2-10-1}
\expandafter\def\csname sga@target@IV.21.2.11\endcsname{ega4.c04.prop.21-2-11}
\expandafter\def\csname sga@target@IV.21.2.11.1\endcsname{ega4.c04.prop.21-2-11-1}
\expandafter\def\csname sga@target@IV.21.2.12\endcsname{ega4.c04.statement.21-2-12}
\expandafter\def\csname sga@target@IV.21.2.2\endcsname{ega4.c04.prop.21-2-2}
\expandafter\def\csname sga@target@IV.21.2.3\endcsname{ega4.c04.cor.21-2-3}
\expandafter\def\csname sga@target@IV.21.2.4\endcsname{ega4.c04.statement.21-2-4}
\expandafter\def\csname sga@target@IV.21.2.5\endcsname{ega4.c04.statement.21-2-5}
\expandafter\def\csname sga@target@IV.21.2.5.1\endcsname{ega4.c04.statement.21-2-5-1}
\expandafter\def\csname sga@target@IV.21.2.5.2\endcsname{ega4.c04.eq.21-2-5-2}
\expandafter\def\csname sga@target@IV.21.2.5.3\endcsname{ega4.c04.eq.21-2-5-3}
\expandafter\def\csname sga@target@IV.21.2.5.4\endcsname{ega4.c04.eq.21-2-5-4}
\expandafter\def\csname sga@target@IV.21.2.6\endcsname{ega4.c04.prop.21-2-6}
\expandafter\def\csname sga@target@IV.21.2.6.1\endcsname{ega4.c04.eq.21-2-6-1}
\expandafter\def\csname sga@target@IV.21.2.7\endcsname{ega4.c04.statement.21-2-7}
\expandafter\def\csname sga@target@IV.21.2.7.1\endcsname{ega4.c04.statement.21-2-7-1}
\expandafter\def\csname sga@target@IV.21.2.7.2\endcsname{ega4.c04.eq.21-2-7-2}
\expandafter\def\csname sga@target@IV.21.2.8\endcsname{ega4.c04.statement.21-2-8}
\expandafter\def\csname sga@target@IV.21.2.8.1\endcsname{ega4.c04.statement.21-2-8-1}
\expandafter\def\csname sga@target@IV.21.2.8.2\endcsname{ega4.c04.statement.21-2-8-2}
\expandafter\def\csname sga@target@IV.21.2.8.3\endcsname{ega4.c04.statement.21-2-8-3}
\expandafter\def\csname sga@target@IV.21.2.9\endcsname{ega4.c04.statement.21-2-9}
\expandafter\def\csname sga@target@IV.21.2.9.1\endcsname{ega4.c04.statement.21-2-9-1}
\expandafter\def\csname sga@target@IV.21.2.9.2\endcsname{ega4.c04.eq.21-2-9-2}
\expandafter\def\csname sga@target@IV.21.2.9.3\endcsname{ega4.c04.eq.21-2-9-3}
\expandafter\def\csname sga@target@IV.21.2.9.4\endcsname{ega4.c04.eq.21-2-9-4}
\expandafter\def\csname sga@target@IV.21.3.1\endcsname{ega4.c04.statement.21-3-1}
\expandafter\def\csname sga@target@IV.21.3.2\endcsname{ega4.c04.statement.21-3-2}
\expandafter\def\csname sga@target@IV.21.3.2.1\endcsname{ega4.c04.statement.21-3-2-1}
\expandafter\def\csname sga@target@IV.21.3.2.2\endcsname{ega4.c04.eq.21-3-2-2}
\expandafter\def\csname sga@target@IV.21.3.2.3\endcsname{ega4.c04.eq.21-3-2-3}
\expandafter\def\csname sga@target@IV.21.3.2.4\endcsname{ega4.c04.eq.21-3-2-4}
\expandafter\def\csname sga@target@IV.21.3.3\endcsname{ega4.c04.prop.21-3-3}
\expandafter\def\csname sga@target@IV.21.3.3.1\endcsname{ega4.c04.eq.21-3-3-1}
\expandafter\def\csname sga@target@IV.21.3.4\endcsname{ega4.c04.prop.21-3-4}
\expandafter\def\csname sga@target@IV.21.3.5\endcsname{ega4.c04.cor.21-3-5}
\expandafter\def\csname sga@target@IV.21.3.6\endcsname{ega4.c04.rem.21-3-6}
\expandafter\def\csname sga@target@IV.21.3.6.1\endcsname{ega4.c04.eq.21-3-6-1}
\expandafter\def\csname sga@target@IV.21.4\endcsname{ega4.c04.subsec.21-4.key-iv-21-4}
\expandafter\def\csname sga@target@IV.21.4.1\endcsname{ega4.c04.statement.21-4-1}
\expandafter\def\csname sga@target@IV.21.4.10\endcsname{ega4.c04.cor.21-4-10}
\expandafter\def\csname sga@target@IV.21.4.11\endcsname{ega4.c04.cor.21-4-11}
\expandafter\def\csname sga@target@IV.21.4.12\endcsname{ega4.c04.rem.21-4-12}
\expandafter\def\csname sga@target@IV.21.4.13\endcsname{ega4.c04.cor.21-4-13}
\expandafter\def\csname sga@target@IV.21.4.2\endcsname{ega4.c04.def.21-4-2}
\expandafter\def\csname sga@target@IV.21.4.2.1\endcsname{ega4.c04.eq.21-4-2-1}
\expandafter\def\csname sga@target@IV.21.4.3\endcsname{ega4.c04.statement.21-4-3}
\expandafter\def\csname sga@target@IV.21.4.4\endcsname{ega4.c04.statement.21-4-4}
\expandafter\def\csname sga@target@IV.21.4.5\endcsname{ega4.c04.prop.21-4-5}
\expandafter\def\csname sga@target@IV.21.4.6\endcsname{ega4.c04.cor.21-4-6}
\expandafter\def\csname sga@target@IV.21.4.6.1\endcsname{ega4.c04.eq.21-4-6-1}
\expandafter\def\csname sga@target@IV.21.4.7\endcsname{ega4.c04.statement.21-4-7}
\expandafter\def\csname sga@target@IV.21.4.8\endcsname{ega4.c04.prop.21-4-8}
\expandafter\def\csname sga@target@IV.21.4.9\endcsname{ega4.c04.prop.21-4-9}
\expandafter\def\csname sga@target@IV.21.5\endcsname{ega4.c04.subsec.21-5.key-iv-21-5}
\expandafter\def\csname sga@target@IV.21.5.1\endcsname{ega4.c04.statement.21-5-1}
\expandafter\def\csname sga@target@IV.21.5.2\endcsname{ega4.c04.lem.21-5-2}
\expandafter\def\csname sga@target@IV.21.5.3\endcsname{ega4.c04.statement.21-5-3}
\expandafter\def\csname sga@target@IV.21.5.3.1\endcsname{ega4.c04.eq.21-5-3-1}
\expandafter\def\csname sga@target@IV.21.5.3.2\endcsname{ega4.c04.eq.21-5-3-2}
\expandafter\def\csname sga@target@IV.21.5.4\endcsname{ega4.c04.statement.21-5-4}
\expandafter\def\csname sga@target@IV.21.5.5\endcsname{ega4.c04.def.21-5-5}
\expandafter\def\csname sga@target@IV.21.5.5.1\endcsname{ega4.c04.eq.21-5-5-1}
\expandafter\def\csname sga@target@IV.21.5.5.2\endcsname{ega4.c04.eq.21-5-5-2}
\expandafter\def\csname sga@target@IV.21.5.5.3\endcsname{ega4.c04.eq.21-5-5-3}
\expandafter\def\csname sga@target@IV.21.5.6\endcsname{ega4.c04.prop.21-5-6}
\expandafter\def\csname sga@target@IV.21.5.6.1\endcsname{ega4.c04.prop.21-5-6-1}
\expandafter\def\csname sga@target@IV.21.5.7\endcsname{ega4.c04.prop.21-5-7}
\expandafter\def\csname sga@target@IV.21.5.7.1\endcsname{ega4.c04.prop.21-5-7-1}
\expandafter\def\csname sga@target@IV.21.5.7.2\endcsname{ega4.c04.lem.21-5-7-2}
\expandafter\def\csname sga@target@IV.21.5.7.3\endcsname{ega4.c04.lem.21-5-7-3}
\expandafter\def\csname sga@target@IV.21.5.7.4\endcsname{ega4.c04.eq.21-5-7-4}
\expandafter\def\csname sga@target@IV.21.5.8\endcsname{ega4.c04.prop.21-5-8}
\expandafter\def\csname sga@target@IV.21.5.8.1\endcsname{ega4.c04.eq.21-5-8-1}
\expandafter\def\csname sga@target@IV.21.6\endcsname{ega4.c04.subsec.21-6.key-iv-21-6}
\expandafter\def\csname sga@target@IV.21.6.1\endcsname{ega4.c04.statement.21-6-1}
\expandafter\def\csname sga@target@IV.21.6.10\endcsname{ega4.c04.cor.21-6-10}
\expandafter\def\csname sga@target@IV.21.6.10.1\endcsname{ega4.c04.eq.21-6-10-1}
\expandafter\def\csname sga@target@IV.21.6.11\endcsname{ega4.c04.cor.21-6-11}
\expandafter\def\csname sga@target@IV.21.6.12\endcsname{ega4.c04.prop.21-6-12}
\expandafter\def\csname sga@target@IV.21.6.12.1\endcsname{ega4.c04.eq.21-6-12-1}
\expandafter\def\csname sga@target@IV.21.6.12.2\endcsname{ega4.c04.eq.21-6-12-2}
\expandafter\def\csname sga@target@IV.21.6.13\endcsname{ega4.c04.cor.21-6-13}
\expandafter\def\csname sga@target@IV.21.6.14\endcsname{ega4.c04.cor.21-6-14}
\expandafter\def\csname sga@target@IV.21.6.2\endcsname{ega4.c04.statement.21-6-2}
\expandafter\def\csname sga@target@IV.21.6.3\endcsname{ega4.c04.statement.21-6-3}
\expandafter\def\csname sga@target@IV.21.6.4\endcsname{ega4.c04.statement.21-6-4}
\expandafter\def\csname sga@target@IV.21.6.4.1\endcsname{ega4.c04.statement.21-6-4-1}
\expandafter\def\csname sga@target@IV.21.6.4.2\endcsname{ega4.c04.eq.21-6-4-2}
\expandafter\def\csname sga@target@IV.21.6.5\endcsname{ega4.c04.statement.21-6-5}
\expandafter\def\csname sga@target@IV.21.6.5.1\endcsname{ega4.c04.eq.21-6-5-1}
\expandafter\def\csname sga@target@IV.21.6.6\endcsname{ega4.c04.prop.21-6-6}
\expandafter\def\csname sga@target@IV.21.6.6.2\endcsname{ega4.c04.statement.21-6-6-2}
\expandafter\def\csname sga@target@IV.21.6.7\endcsname{ega4.c04.statement.21-6-7}
\expandafter\def\csname sga@target@IV.21.6.7.1\endcsname{ega4.c04.eq.21-6-7-1}
\expandafter\def\csname sga@target@IV.21.6.7.2\endcsname{ega4.c04.eq.21-6-7-2}
\expandafter\def\csname sga@target@IV.21.6.7.3\endcsname{ega4.c04.eq.21-6-7-3}
\expandafter\def\csname sga@target@IV.21.6.7.4\endcsname{ega4.c04.eq.21-6-7-4}
\expandafter\def\csname sga@target@IV.21.6.7.5\endcsname{ega4.c04.eq.21-6-7-5}
\expandafter\def\csname sga@target@IV.21.6.8\endcsname{ega4.c04.statement.21-6-8}
\expandafter\def\csname sga@target@IV.21.6.9\endcsname{ega4.c04.thm.21-6-9}
\expandafter\def\csname sga@target@IV.21.6.9.1\endcsname{ega4.c04.eq.21-6-9-1}
\expandafter\def\csname sga@target@IV.21.6.9.2\endcsname{ega4.c04.statement.21-6-9-2}
\expandafter\def\csname sga@target@IV.21.7\endcsname{ega4.c04.subsec.21-7.key-iv-21-7}
\expandafter\def\csname sga@target@IV.21.7.1\endcsname{ega4.c04.statement.21-7-1}
\expandafter\def\csname sga@target@IV.21.7.2\endcsname{ega4.c04.prop.21-7-2}
\expandafter\def\csname sga@target@IV.21.7.2.1\endcsname{ega4.c04.eq.21-7-2-1}
\expandafter\def\csname sga@target@IV.21.7.3\endcsname{ega4.c04.cor.21-7-3}
\expandafter\def\csname sga@target@IV.21.7.3.1\endcsname{ega4.c04.statement.21-7-3-1}
\expandafter\def\csname sga@target@IV.21.7.4\endcsname{ega4.c04.prop.21-7-4}
\expandafter\def\csname sga@target@IV.21.7.5\endcsname{ega4.c04.rem.21-7-5}
\expandafter\def\csname sga@target@IV.21.7.6\endcsname{ega4.c04.cor.21-7-6}
\expandafter\def\csname sga@target@IV.21.7.7\endcsname{ega4.c04.rem.21-7-7}
\expandafter\def\csname sga@target@IV.21.8.1\endcsname{ega4.c04.lem.21-8-1}
\expandafter\def\csname sga@target@IV.21.8.2\endcsname{ega4.c04.cor.21-8-2}
\expandafter\def\csname sga@target@IV.21.8.3\endcsname{ega4.c04.prop.21-8-3}
\expandafter\def\csname sga@target@IV.21.8.3.1\endcsname{ega4.c04.eq.21-8-3-1}
\expandafter\def\csname sga@target@IV.21.8.4\endcsname{ega4.c04.cor.21-8-4}
\expandafter\def\csname sga@target@IV.21.8.5\endcsname{ega4.c04.prop.21-8-5}
\expandafter\def\csname sga@target@IV.21.8.5.1\endcsname{ega4.c04.eq.21-8-5-1}
\expandafter\def\csname sga@target@IV.21.8.5.2\endcsname{ega4.c04.eq.21-8-5-2}
\expandafter\def\csname sga@target@IV.21.8.5.3\endcsname{ega4.c04.eq.21-8-5-3}
\expandafter\def\csname sga@target@IV.21.8.6\endcsname{ega4.c04.rem.21-8-6}
\expandafter\def\csname sga@target@IV.21.9\endcsname{ega4.c04.subsec.21-9.key-iv-21-9}
\expandafter\def\csname sga@target@IV.21.9.1\endcsname{ega4.c04.statement.21-9-1}
\expandafter\def\csname sga@target@IV.21.9.1.1\endcsname{ega4.c04.eq.21-9-1-1}
\expandafter\def\csname sga@target@IV.21.9.1.2\endcsname{ega4.c04.eq.21-9-1-2}
\expandafter\def\csname sga@target@IV.21.9.10\endcsname{ega4.c04.statement.21-9-10}
\expandafter\def\csname sga@target@IV.21.9.11\endcsname{ega4.c04.prop.21-9-11}
\expandafter\def\csname sga@target@IV.21.9.12\endcsname{ega4.c04.cor.21-9-12}
\expandafter\def\csname sga@target@IV.21.9.13\endcsname{ega4.c04.rem.21-9-13}
\expandafter\def\csname sga@target@IV.21.9.2\endcsname{ega4.c04.lem.21-9-2}
\expandafter\def\csname sga@target@IV.21.9.3\endcsname{ega4.c04.rem.21-9-3}
\expandafter\def\csname sga@target@IV.21.9.4\endcsname{ega4.c04.prop.21-9-4}
\expandafter\def\csname sga@target@IV.21.9.4.1\endcsname{ega4.c04.prop.21-9-4-1}
\expandafter\def\csname sga@target@IV.21.9.5\endcsname{ega4.c04.cor.21-9-5}
\expandafter\def\csname sga@target@IV.21.9.6\endcsname{ega4.c04.cor.21-9-6}
\expandafter\def\csname sga@target@IV.21.9.7\endcsname{ega4.c04.cor.21-9-7}
\expandafter\def\csname sga@target@IV.21.9.8\endcsname{ega4.c04.statement.21-9-8}
\expandafter\def\csname sga@target@IV.21.9.9\endcsname{ega4.c04.prop.21-9-9}
\expandafter\def\csname sga@target@IV.3.1.1\endcsname{ega4.c04.def.3-1-1}
\expandafter\def\csname sga@target@IV.3.1.1.1\endcsname{ega4.c04.eq.3-1-1-1}
\expandafter\def\csname sga@target@IV.3.1.10\endcsname{ega4.c04.prop.3-1-10}
\expandafter\def\csname sga@target@IV.3.1.10.1\endcsname{ega4.c04.lem.3-1-10-1}
\expandafter\def\csname sga@target@IV.3.1.10.2\endcsname{ega4.c04.rem.3-1-10-2}
\expandafter\def\csname sga@target@IV.3.1.10.3\endcsname{ega4.c04.eq.3-1-10-3}
\expandafter\def\csname sga@target@IV.3.1.10.4\endcsname{ega4.c04.eq.3-1-10-4}
\expandafter\def\csname sga@target@IV.3.1.10.5\endcsname{ega4.c04.rem.3-1-10-5}
\expandafter\def\csname sga@target@IV.3.1.11\endcsname{ega4.c04.cor.3-1-11}
\expandafter\def\csname sga@target@IV.3.1.12\endcsname{ega4.c04.rem.3-1-12}
\expandafter\def\csname sga@target@IV.3.1.13\endcsname{ega4.c04.prop.3-1-13}
\expandafter\def\csname sga@target@IV.3.1.2\endcsname{ega4.c04.prop.3-1-2}
\expandafter\def\csname sga@target@IV.3.1.3\endcsname{ega4.c04.prop.3-1-3}
\expandafter\def\csname sga@target@IV.3.1.4\endcsname{ega4.c04.cor.3-1-4}
\expandafter\def\csname sga@target@IV.3.1.5\endcsname{ega4.c04.cor.3-1-5}
\expandafter\def\csname sga@target@IV.3.1.6\endcsname{ega4.c04.prop.3-1-6}
\expandafter\def\csname sga@target@IV.3.1.7\endcsname{ega4.c04.prop.3-1-7}
\expandafter\def\csname sga@target@IV.3.1.8\endcsname{ega4.c04.prop.3-1-8}
\expandafter\def\csname sga@target@IV.3.1.9\endcsname{ega4.c04.cor.3-1-9}
\expandafter\def\csname sga@target@IV.3.2.1\endcsname{ega4.c04.prop.3-2-1}
\expandafter\def\csname sga@target@IV.3.2.2\endcsname{ega4.c04.def.3-2-2}
\expandafter\def\csname sga@target@IV.3.2.3\endcsname{ega4.c04.prop.3-2-3}
\expandafter\def\csname sga@target@IV.3.2.4\endcsname{ega4.c04.def.3-2-4}
\expandafter\def\csname sga@target@IV.3.2.5\endcsname{ega4.c04.def.3-2-5}
\expandafter\def\csname sga@target@IV.3.2.6\endcsname{ega4.c04.prop.3-2-6}
\expandafter\def\csname sga@target@IV.3.2.7\endcsname{ega4.c04.cor.3-2-7}
\expandafter\def\csname sga@target@IV.3.2.8\endcsname{ega4.c04.cor.3-2-8}
\expandafter\def\csname sga@target@IV.3.2.9\endcsname{ega4.c04.rem.3-2-9}
\expandafter\def\csname sga@target@IV.3.3.1\endcsname{ega4.c04.prop.3-3-1}
\expandafter\def\csname sga@target@IV.3.3.1.1\endcsname{ega4.c04.prop.3-3-1-1}
\expandafter\def\csname sga@target@IV.3.3.2\endcsname{ega4.c04.cor.3-3-2}
\expandafter\def\csname sga@target@IV.3.3.3\endcsname{ega4.c04.cor.3-3-3}
\expandafter\def\csname sga@target@IV.3.3.4\endcsname{ega4.c04.prop.3-3-4}
\expandafter\def\csname sga@target@IV.3.3.5\endcsname{ega4.c04.cor.3-3-5}
\expandafter\def\csname sga@target@IV.3.3.6\endcsname{ega4.c04.prop.3-3-6}
\expandafter\def\csname sga@target@IV.3.3.6.1\endcsname{ega4.c04.eq.3-3-6-1}
\expandafter\def\csname sga@target@IV.3.3.6.2\endcsname{ega4.c04.eq.3-3-6-2}
\expandafter\def\csname sga@target@IV.3.3.6.3\endcsname{ega4.c04.eq.3-3-6-3}
\expandafter\def\csname sga@target@IV.3.3.6.4\endcsname{ega4.c04.eq.3-3-6-4}
\expandafter\def\csname sga@target@IV.3.3.7\endcsname{ega4.c04.cor.3-3-7}
\expandafter\def\csname sga@target@IV.3.3.7.1\endcsname{ega4.c04.cor.3-3-7-1}
\expandafter\def\csname sga@target@IV.3.3.8\endcsname{ega4.c04.sec.3-3-8}
\expandafter\def\csname sga@target@IV.3.3.9\endcsname{ega4.c04.cor.3-3-9}
\expandafter\def\csname sga@target@IV.3.3.9.1\endcsname{ega4.c04.cor.3-3-9-1}
\expandafter\def\csname sga@target@IV.3.4\endcsname{ega4.c04.subsec.3-4.key-iv-3-4}
\expandafter\def\csname sga@target@IV.3.4.1\endcsname{ega4.c04.prop.3-4-1}
\expandafter\def\csname sga@target@IV.3.4.1.1\endcsname{ega4.c04.eq.3-4-1-1}
\expandafter\def\csname sga@target@IV.3.4.1.2\endcsname{ega4.c04.lem.3-4-1-2}
\expandafter\def\csname sga@target@IV.3.4.1.3\endcsname{ega4.c04.lem.3-4-1-3}
\expandafter\def\csname sga@target@IV.3.4.1.4\endcsname{ega4.c04.lem.3-4-1-4}
\expandafter\def\csname sga@target@IV.3.4.1.5\endcsname{ega4.c04.eq.3-4-1-5}
\expandafter\def\csname sga@target@IV.3.4.2\endcsname{ega4.c04.cor.3-4-2}
\expandafter\def\csname sga@target@IV.3.4.3\endcsname{ega4.c04.prop.3-4-3}
\expandafter\def\csname sga@target@IV.3.4.4\endcsname{ega4.c04.prop.3-4-4}
\expandafter\def\csname sga@target@IV.3.4.5\endcsname{ega4.c04.prop.3-4-5}
\expandafter\def\csname sga@target@IV.3.4.6\endcsname{ega4.c04.prop.3-4-6}
\expandafter\def\csname sga@target@IV.3.4.6.1\endcsname{ega4.c04.lem.3-4-6-1}
\expandafter\def\csname sga@target@IV.3.4.7\endcsname{ega4.c04.cor.3-4-7}
\expandafter\def\csname sga@target@IV.3.4.8\endcsname{ega4.c04.rem.3-4-8}
\expandafter\def\csname sga@target@IV.3.4.9\endcsname{ega4.c04.prop.3-4-9}
\expandafter\def\csname sga@target@IV.4.1.1\endcsname{ega4.c04.def.4-1-1}
\expandafter\def\csname sga@target@IV.4.1.1.1\endcsname{ega4.c04.eq.4-1-1-1}
\expandafter\def\csname sga@target@IV.4.1.1.2\endcsname{ega4.c04.eq.4-1-1-2}
\expandafter\def\csname sga@target@IV.4.1.1.3\endcsname{ega4.c04.eq.4-1-1-3}
\expandafter\def\csname sga@target@IV.4.1.2\endcsname{ega4.c04.thm.4-1-2}
\expandafter\def\csname sga@target@IV.4.1.2.1\endcsname{ega4.c04.cor.4-1-2-1}
\expandafter\def\csname sga@target@IV.4.1.3\endcsname{ega4.c04.rem.4-1-3}
\expandafter\def\csname sga@target@IV.4.1.4\endcsname{ega4.c04.cor.4-1-4}
\expandafter\def\csname sga@target@IV.4.1.5\endcsname{ega4.c04.cor.4-1-5}
\expandafter\def\csname sga@target@IV.4.2.1\endcsname{ega4.c04.prop.4-2-1}
\expandafter\def\csname sga@target@IV.4.2.1.1\endcsname{ega4.c04.eq.4-2-1-1}
\expandafter\def\csname sga@target@IV.4.2.1.2\endcsname{ega4.c04.eq.4-2-1-2}
\expandafter\def\csname sga@target@IV.4.2.1.3\endcsname{ega4.c04.eq.4-2-1-3}
\expandafter\def\csname sga@target@IV.4.2.1.4\endcsname{ega4.c04.rem.4-2-1-4}
\expandafter\def\csname sga@target@IV.4.2.1.5\endcsname{ega4.c04.eq.4-2-1-5}
\expandafter\def\csname sga@target@IV.4.2.10\endcsname{ega4.c04.cor.4-2-10}
\expandafter\def\csname sga@target@IV.4.2.2\endcsname{ega4.c04.cor.4-2-2}
\expandafter\def\csname sga@target@IV.4.2.3\endcsname{ega4.c04.cor.4-2-3}
\expandafter\def\csname sga@target@IV.4.2.4\endcsname{ega4.c04.prop.4-2-4}
\expandafter\def\csname sga@target@IV.4.2.4.i\endcsname{ega4.c04.sec.4-2-4-i}
\expandafter\def\csname sga@target@IV.4.2.4.ii\endcsname{ega4.c04.sec.4-2-4-ii}
\expandafter\def\csname sga@target@IV.4.2.4.iii\endcsname{ega4.c04.sec.4-2-4-iii}
\expandafter\def\csname sga@target@IV.4.2.5\endcsname{ega4.c04.prop.4-2-5}
\expandafter\def\csname sga@target@IV.4.2.6\endcsname{ega4.c04.cor.4-2-6}
\expandafter\def\csname sga@target@IV.4.2.7\endcsname{ega4.c04.prop.4-2-7}
\expandafter\def\csname sga@target@IV.4.2.7.i\endcsname{ega4.c04.sec.4-2-7-i}
\expandafter\def\csname sga@target@IV.4.2.7.ii\endcsname{ega4.c04.sec.4-2-7-ii}
\expandafter\def\csname sga@target@IV.4.2.7.iii\endcsname{ega4.c04.sec.4-2-7-iii}
\expandafter\def\csname sga@target@IV.4.2.7.iv\endcsname{ega4.c04.sec.4-2-7-iv}
\expandafter\def\csname sga@target@IV.4.2.8\endcsname{ega4.c04.cor.4-2-8}
\expandafter\def\csname sga@target@IV.4.2.9\endcsname{ega4.c04.prop.4-2-9}
\expandafter\def\csname sga@target@IV.4.3.1\endcsname{ega4.c04.statement.4-3-1}
\expandafter\def\csname sga@target@IV.4.3.2\endcsname{ega4.c04.prop.4-3-2}
\expandafter\def\csname sga@target@IV.4.3.3\endcsname{ega4.c04.cor.4-3-3}
\expandafter\def\csname sga@target@IV.4.3.4\endcsname{ega4.c04.cor.4-3-4}
\expandafter\def\csname sga@target@IV.4.3.5\endcsname{ega4.c04.prop.4-3-5}
\expandafter\def\csname sga@target@IV.4.3.6\endcsname{ega4.c04.cor.4-3-6}
\expandafter\def\csname sga@target@IV.4.3.7\endcsname{ega4.c04.cor.4-3-7}
\expandafter\def\csname sga@target@IV.4.3.8\endcsname{ega4.c04.cor.4-3-8}
\expandafter\def\csname sga@target@IV.4.3.9\endcsname{ega4.c04.prop.4-3-9}
\expandafter\def\csname sga@target@IV.4.4.1\endcsname{ega4.c04.statement.4-4-1}
\expandafter\def\csname sga@target@IV.4.4.2\endcsname{ega4.c04.lem.4-4-2}
\expandafter\def\csname sga@target@IV.4.4.2.i\endcsname{ega4.c04.sec.4-4-2-i}
\expandafter\def\csname sga@target@IV.4.4.2.ii\endcsname{ega4.c04.sec.4-4-2-ii}
\expandafter\def\csname sga@target@IV.4.4.3\endcsname{ega4.c04.rem.4-4-3}
\expandafter\def\csname sga@target@IV.4.4.4\endcsname{ega4.c04.thm.4-4-4}
\expandafter\def\csname sga@target@IV.4.4.5\endcsname{ega4.c04.cor.4-4-5}
\expandafter\def\csname sga@target@IV.4.4.6\endcsname{ega4.c04.cor.4-4-6}
\expandafter\def\csname sga@target@IV.4.5\endcsname{ega4.c04.subsec.4-5.key-iv-4-5}
\expandafter\def\csname sga@target@IV.4.5.1\endcsname{ega4.c04.prop.4-5-1}
\expandafter\def\csname sga@target@IV.4.5.10\endcsname{ega4.c04.cor.4-5-10}
\expandafter\def\csname sga@target@IV.4.5.11\endcsname{ega4.c04.cor.4-5-11}
\expandafter\def\csname sga@target@IV.4.5.12\endcsname{ega4.c04.sec.4-5-12}
\expandafter\def\csname sga@target@IV.4.5.12.i\endcsname{ega4.c04.sec.4-5-12-i}
\expandafter\def\csname sga@target@IV.4.5.12.ii\endcsname{ega4.c04.sec.4-5-12-ii}
\expandafter\def\csname sga@target@IV.4.5.13\endcsname{ega4.c04.prop.4-5-13}
\expandafter\def\csname sga@target@IV.4.5.13.1\endcsname{ega4.c04.cor.4-5-13-1}
\expandafter\def\csname sga@target@IV.4.5.13.1.i\endcsname{ega4.c04.sec.4-5-13-1-i}
\expandafter\def\csname sga@target@IV.4.5.13.1.ii\endcsname{ega4.c04.sec.4-5-13-1-ii}
\expandafter\def\csname sga@target@IV.4.5.14\endcsname{ega4.c04.cor.4-5-14}
\expandafter\def\csname sga@target@IV.4.5.15\endcsname{ega4.c04.prop.4-5-15}
\expandafter\def\csname sga@target@IV.4.5.15.i\endcsname{ega4.c04.sec.4-5-15-i}
\expandafter\def\csname sga@target@IV.4.5.15.ii\endcsname{ega4.c04.sec.4-5-15-ii}
\expandafter\def\csname sga@target@IV.4.5.16\endcsname{ega4.c04.cor.4-5-16}
\expandafter\def\csname sga@target@IV.4.5.16.i\endcsname{ega4.c04.sec.4-5-16-i}
\expandafter\def\csname sga@target@IV.4.5.16.ii\endcsname{ega4.c04.sec.4-5-16-ii}
\expandafter\def\csname sga@target@IV.4.5.17\endcsname{ega4.c04.cor.4-5-17}
\expandafter\def\csname sga@target@IV.4.5.18\endcsname{ega4.c04.rem.4-5-18}
\expandafter\def\csname sga@target@IV.4.5.19\endcsname{ega4.c04.prop.4-5-19}
\expandafter\def\csname sga@target@IV.4.5.19.1\endcsname{ega4.c04.lem.4-5-19-1}
\expandafter\def\csname sga@target@IV.4.5.19.2\endcsname{ega4.c04.eq.4-5-19-2}
\expandafter\def\csname sga@target@IV.4.5.19.3\endcsname{ega4.c04.cor.4-5-19-3}
\expandafter\def\csname sga@target@IV.4.5.2\endcsname{ega4.c04.def.4-5-2}
\expandafter\def\csname sga@target@IV.4.5.2.a\endcsname{ega4.c04.sec.4-5-2-a}
\expandafter\def\csname sga@target@IV.4.5.2.b\endcsname{ega4.c04.sec.4-5-2-b}
\expandafter\def\csname sga@target@IV.4.5.2.c\endcsname{ega4.c04.sec.4-5-2-c}
\expandafter\def\csname sga@target@IV.4.5.20\endcsname{ega4.c04.rem.4-5-20}
\expandafter\def\csname sga@target@IV.4.5.21\endcsname{ega4.c04.prop.4-5-21}
\expandafter\def\csname sga@target@IV.4.5.3\endcsname{ega4.c04.statement.4-5-3}
\expandafter\def\csname sga@target@IV.4.5.4\endcsname{ega4.c04.prop.4-5-4}
\expandafter\def\csname sga@target@IV.4.5.5\endcsname{ega4.c04.def.4-5-5}
\expandafter\def\csname sga@target@IV.4.5.6\endcsname{ega4.c04.prop.4-5-6}
\expandafter\def\csname sga@target@IV.4.5.6.i\endcsname{ega4.c04.sec.4-5-6-i}
\expandafter\def\csname sga@target@IV.4.5.6.ii\endcsname{ega4.c04.sec.4-5-6-ii}
\expandafter\def\csname sga@target@IV.4.5.7\endcsname{ega4.c04.prop.4-5-7}
\expandafter\def\csname sga@target@IV.4.5.8\endcsname{ega4.c04.cor.4-5-8}
\expandafter\def\csname sga@target@IV.4.5.8.i\endcsname{ega4.c04.sec.4-5-8-i}
\expandafter\def\csname sga@target@IV.4.5.8.ii\endcsname{ega4.c04.sec.4-5-8-ii}
\expandafter\def\csname sga@target@IV.4.5.9\endcsname{ega4.c04.prop.4-5-9}
\expandafter\def\csname sga@target@IV.4.5.9.a\endcsname{ega4.c04.sec.4-5-9-a}
\expandafter\def\csname sga@target@IV.4.5.9.b\endcsname{ega4.c04.sec.4-5-9-b}
\expandafter\def\csname sga@target@IV.4.5.9.c\endcsname{ega4.c04.sec.4-5-9-c}
\expandafter\def\csname sga@target@IV.4.6\endcsname{ega4.c04.subsec.4-6.key-iv-4-6}
\expandafter\def\csname sga@target@IV.4.6.1\endcsname{ega4.c04.prop.4-6-1}
\expandafter\def\csname sga@target@IV.4.6.1.a\endcsname{ega4.c04.sec.4-6-1-a}
\expandafter\def\csname sga@target@IV.4.6.1.b\endcsname{ega4.c04.sec.4-6-1-b}
\expandafter\def\csname sga@target@IV.4.6.1.c\endcsname{ega4.c04.sec.4-6-1-c}
\expandafter\def\csname sga@target@IV.4.6.1.d\endcsname{ega4.c04.sec.4-6-1-d}
\expandafter\def\csname sga@target@IV.4.6.1.e\endcsname{ega4.c04.sec.4-6-1-e}
\expandafter\def\csname sga@target@IV.4.6.10\endcsname{ega4.c04.prop.4-6-10}
\expandafter\def\csname sga@target@IV.4.6.11\endcsname{ega4.c04.cor.4-6-11}
\expandafter\def\csname sga@target@IV.4.6.12\endcsname{ega4.c04.prop.4-6-12}
\expandafter\def\csname sga@target@IV.4.6.12.a\endcsname{ega4.c04.sec.4-6-12-a}
\expandafter\def\csname sga@target@IV.4.6.12.b\endcsname{ega4.c04.sec.4-6-12-b}
\expandafter\def\csname sga@target@IV.4.6.12.c\endcsname{ega4.c04.sec.4-6-12-c}
\expandafter\def\csname sga@target@IV.4.6.12.d\endcsname{ega4.c04.sec.4-6-12-d}
\expandafter\def\csname sga@target@IV.4.6.12.e\endcsname{ega4.c04.sec.4-6-12-e}
\expandafter\def\csname sga@target@IV.4.6.13\endcsname{ega4.c04.cor.4-6-13}
\expandafter\def\csname sga@target@IV.4.6.13.f\endcsname{ega4.c04.sec.4-6-13-f}
\expandafter\def\csname sga@target@IV.4.6.14\endcsname{ega4.c04.statement.4-6-14}
\expandafter\def\csname sga@target@IV.4.6.15\endcsname{ega4.c04.prop.4-6-15}
\expandafter\def\csname sga@target@IV.4.6.15.i\endcsname{ega4.c04.sec.4-6-15-i}
\expandafter\def\csname sga@target@IV.4.6.15.ii\endcsname{ega4.c04.sec.4-6-15-ii}
\expandafter\def\csname sga@target@IV.4.6.16\endcsname{ega4.c04.prop.4-6-16}
\expandafter\def\csname sga@target@IV.4.6.16.a\endcsname{ega4.c04.sec.4-6-16-a}
\expandafter\def\csname sga@target@IV.4.6.16.b\endcsname{ega4.c04.sec.4-6-16-b}
\expandafter\def\csname sga@target@IV.4.6.16.c\endcsname{ega4.c04.sec.4-6-16-c}
\expandafter\def\csname sga@target@IV.4.6.16.d\endcsname{ega4.c04.sec.4-6-16-d}
\expandafter\def\csname sga@target@IV.4.6.17\endcsname{ega4.c04.def.4-6-17}
\expandafter\def\csname sga@target@IV.4.6.18\endcsname{ega4.c04.prop.4-6-18}
\expandafter\def\csname sga@target@IV.4.6.18.a\endcsname{ega4.c04.sec.4-6-18-a}
\expandafter\def\csname sga@target@IV.4.6.18.b\endcsname{ega4.c04.sec.4-6-18-b}
\expandafter\def\csname sga@target@IV.4.6.18.c\endcsname{ega4.c04.sec.4-6-18-c}
\expandafter\def\csname sga@target@IV.4.6.18.d\endcsname{ega4.c04.sec.4-6-18-d}
\expandafter\def\csname sga@target@IV.4.6.19\endcsname{ega4.c04.def.4-6-19}
\expandafter\def\csname sga@target@IV.4.6.2\endcsname{ega4.c04.def.4-6-2}
\expandafter\def\csname sga@target@IV.4.6.20\endcsname{ega4.c04.prop.4-6-20}
\expandafter\def\csname sga@target@IV.4.6.21\endcsname{ega4.c04.prop.4-6-21}
\expandafter\def\csname sga@target@IV.4.6.21.i\endcsname{ega4.c04.sec.4-6-21-i}
\expandafter\def\csname sga@target@IV.4.6.21.ii\endcsname{ega4.c04.sec.4-6-21-ii}
\expandafter\def\csname sga@target@IV.4.6.22\endcsname{ega4.c04.def.4-6-22}
\expandafter\def\csname sga@target@IV.4.6.3\endcsname{ega4.c04.cor.4-6-3}
\expandafter\def\csname sga@target@IV.4.6.4\endcsname{ega4.c04.cor.4-6-4}
\expandafter\def\csname sga@target@IV.4.6.5\endcsname{ega4.c04.prop.4-6-5}
\expandafter\def\csname sga@target@IV.4.6.5.i\endcsname{ega4.c04.sec.4-6-5-i}
\expandafter\def\csname sga@target@IV.4.6.5.ii\endcsname{ega4.c04.sec.4-6-5-ii}
\expandafter\def\csname sga@target@IV.4.6.6\endcsname{ega4.c04.prop.4-6-6}
\expandafter\def\csname sga@target@IV.4.6.7\endcsname{ega4.c04.cor.4-6-7}
\expandafter\def\csname sga@target@IV.4.6.8\endcsname{ega4.c04.cor.4-6-8}
\expandafter\def\csname sga@target@IV.4.6.9\endcsname{ega4.c04.def.4-6-9}
\expandafter\def\csname sga@target@IV.4.7.1\endcsname{ega4.c04.lem.4-7-1}
\expandafter\def\csname sga@target@IV.4.7.1.1\endcsname{ega4.c04.eq.4-7-1-1}
\expandafter\def\csname sga@target@IV.4.7.10\endcsname{ega4.c04.prop.4-7-10}
\expandafter\def\csname sga@target@IV.4.7.11\endcsname{ega4.c04.prop.4-7-11}
\expandafter\def\csname sga@target@IV.4.7.12\endcsname{ega4.c04.def.4-7-12}
\expandafter\def\csname sga@target@IV.4.7.13\endcsname{ega4.c04.prop.4-7-13}
\expandafter\def\csname sga@target@IV.4.7.2\endcsname{ega4.c04.cor.4-7-2}
\expandafter\def\csname sga@target@IV.4.7.2.1\endcsname{ega4.c04.eq.4-7-2-1}
\expandafter\def\csname sga@target@IV.4.7.3\endcsname{ega4.c04.prop.4-7-3}
\expandafter\def\csname sga@target@IV.4.7.3.1\endcsname{ega4.c04.eq.4-7-3-1}
\expandafter\def\csname sga@target@IV.4.7.3.i\endcsname{ega4.c04.sec.4-7-3-i}
\expandafter\def\csname sga@target@IV.4.7.3.ii\endcsname{ega4.c04.sec.4-7-3-ii}
\expandafter\def\csname sga@target@IV.4.7.3.iii\endcsname{ega4.c04.sec.4-7-3-iii}
\expandafter\def\csname sga@target@IV.4.7.4\endcsname{ega4.c04.def.4-7-4}
\expandafter\def\csname sga@target@IV.4.7.5\endcsname{ega4.c04.def.4-7-5}
\expandafter\def\csname sga@target@IV.4.7.6\endcsname{ega4.c04.prop.4-7-6}
\expandafter\def\csname sga@target@IV.4.7.6.1\endcsname{ega4.c04.eq.4-7-6-1}
\expandafter\def\csname sga@target@IV.4.7.7\endcsname{ega4.c04.statement.4-7-7}
\expandafter\def\csname sga@target@IV.4.7.8\endcsname{ega4.c04.prop.4-7-8}
\expandafter\def\csname sga@target@IV.4.7.9\endcsname{ega4.c04.cor.4-7-9}
\expandafter\def\csname sga@target@IV.4.8.1\endcsname{ega4.c04.statement.4-8-1}
\expandafter\def\csname sga@target@IV.4.8.10\endcsname{ega4.c04.cor.4-8-10}
\expandafter\def\csname sga@target@IV.4.8.11\endcsname{ega4.c04.cor.4-8-11}
\expandafter\def\csname sga@target@IV.4.8.12\endcsname{ega4.c04.cor.4-8-12}
\expandafter\def\csname sga@target@IV.4.8.13\endcsname{ega4.c04.prop.4-8-13}
\expandafter\def\csname sga@target@IV.4.8.14\endcsname{ega4.c04.prop.4-8-14}
\expandafter\def\csname sga@target@IV.4.8.15\endcsname{ega4.c04.sec.4-8-15}
\expandafter\def\csname sga@target@IV.4.8.15.i\endcsname{ega4.c04.sec.4-8-15-i}
\expandafter\def\csname sga@target@IV.4.8.15.ii\endcsname{ega4.c04.sec.4-8-15-ii}
\expandafter\def\csname sga@target@IV.4.8.2\endcsname{ega4.c04.statement.4-8-2}
\expandafter\def\csname sga@target@IV.4.8.2.1\endcsname{ega4.c04.eq.4-8-2-1}
\expandafter\def\csname sga@target@IV.4.8.2.2\endcsname{ega4.c04.eq.4-8-2-2}
\expandafter\def\csname sga@target@IV.4.8.2.3\endcsname{ega4.c04.eq.4-8-2-3}
\expandafter\def\csname sga@target@IV.4.8.2.4\endcsname{ega4.c04.eq.4-8-2-4}
\expandafter\def\csname sga@target@IV.4.8.2.5\endcsname{ega4.c04.eq.4-8-2-5}
\expandafter\def\csname sga@target@IV.4.8.2.6\endcsname{ega4.c04.eq.4-8-2-6}
\expandafter\def\csname sga@target@IV.4.8.3\endcsname{ega4.c04.prop.4-8-3}
\expandafter\def\csname sga@target@IV.4.8.4\endcsname{ega4.c04.def.4-8-4}
\expandafter\def\csname sga@target@IV.4.8.5\endcsname{ega4.c04.prop.4-8-5}
\expandafter\def\csname sga@target@IV.4.8.6\endcsname{ega4.c04.lem.4-8-6}
\expandafter\def\csname sga@target@IV.4.8.6.a\endcsname{ega4.c04.sec.4-8-6-a}
\expandafter\def\csname sga@target@IV.4.8.6.b\endcsname{ega4.c04.sec.4-8-6-b}
\expandafter\def\csname sga@target@IV.4.8.6.c\endcsname{ega4.c04.sec.4-8-6-c}
\expandafter\def\csname sga@target@IV.4.8.7\endcsname{ega4.c04.cor.4-8-7}
\expandafter\def\csname sga@target@IV.4.8.8\endcsname{ega4.c04.lem.4-8-8}
\expandafter\def\csname sga@target@IV.4.8.8.i\endcsname{ega4.c04.sec.4-8-8-i}
\expandafter\def\csname sga@target@IV.4.8.8.ii\endcsname{ega4.c04.sec.4-8-8-ii}
\expandafter\def\csname sga@target@IV.4.8.8.iii\endcsname{ega4.c04.sec.4-8-8-iii}
\expandafter\def\csname sga@target@IV.4.8.9\endcsname{ega4.c04.prop.4-8-9}
\expandafter\def\csname sga@target@IV.4.9.1\endcsname{ega4.c04.prop.4-9-1}
\expandafter\def\csname sga@target@IV.4.9.1.a\endcsname{ega4.c04.sec.4-9-1-a}
\expandafter\def\csname sga@target@IV.4.9.1.b\endcsname{ega4.c04.sec.4-9-1-b}
\expandafter\def\csname sga@target@IV.4.9.1.bprime\endcsname{ega4.c04.sec.4-9-1-bprime}
\expandafter\def\csname sga@target@IV.4.9.1.c\endcsname{ega4.c04.sec.4-9-1-c}
\expandafter\def\csname sga@target@IV.4.9.2\endcsname{ega4.c04.cor.4-9-2}
\expandafter\def\csname sga@target@IV.4.9.3\endcsname{ega4.c04.rem.4-9-3}
\expandafter\def\csname sga@target@IV.4.9.4\endcsname{ega4.c04.prop.4-9-4}
\expandafter\def\csname sga@target@IV.4.9.5\endcsname{ega4.c04.cor.4-9-5}
\expandafter\def\csname sga@target@IV.4.9.5.a\endcsname{ega4.c04.sec.4-9-5-a}
\expandafter\def\csname sga@target@IV.4.9.5.b\endcsname{ega4.c04.sec.4-9-5-b}
\expandafter\def\csname sga@target@IV.4.9.6\endcsname{ega4.c04.prop.4-9-6}
\expandafter\def\csname sga@target@IV.4.9.6.a\endcsname{ega4.c04.sec.4-9-6-a}
\expandafter\def\csname sga@target@IV.4.9.6.aprime\endcsname{ega4.c04.sec.4-9-6-aprime}
\expandafter\def\csname sga@target@IV.4.9.6.b\endcsname{ega4.c04.sec.4-9-6-b}
\expandafter\def\csname sga@target@IV.4.9.6.bprime\endcsname{ega4.c04.sec.4-9-6-bprime}
\expandafter\def\csname sga@target@IV.4.9.7\endcsname{ega4.c04.cor.4-9-7}
\expandafter\def\csname sga@target@IV.5.1.1\endcsname{ega4.c04.statement.5-1-1}
\expandafter\def\csname sga@target@IV.5.1.1.1\endcsname{ega4.c04.eq.5-1-1-1}
\expandafter\def\csname sga@target@IV.5.1.1.2\endcsname{ega4.c04.eq.5-1-1-2}
\expandafter\def\csname sga@target@IV.5.1.1.3\endcsname{ega4.c04.eq.5-1-1-3}
\expandafter\def\csname sga@target@IV.5.1.1.4\endcsname{ega4.c04.sec.5-1-1-4}
\expandafter\def\csname sga@target@IV.5.1.1.5\endcsname{ega4.c04.sec.5-1-1-5}
\expandafter\def\csname sga@target@IV.5.1.1.6\endcsname{ega4.c04.sec.5-1-1-6}
\expandafter\def\csname sga@target@IV.5.1.10\endcsname{ega4.c04.prop.5-1-10}
\expandafter\def\csname sga@target@IV.5.1.10.1\endcsname{ega4.c04.lem.5-1-10-1}
\expandafter\def\csname sga@target@IV.5.1.10.i\endcsname{ega4.c04.sec.5-1-10-i}
\expandafter\def\csname sga@target@IV.5.1.10.ii\endcsname{ega4.c04.sec.5-1-10-ii}
\expandafter\def\csname sga@target@IV.5.1.11\endcsname{ega4.c04.cor.5-1-11}
\expandafter\def\csname sga@target@IV.5.1.12\endcsname{ega4.c04.statement.5-1-12}
\expandafter\def\csname sga@target@IV.5.1.12.1\endcsname{ega4.c04.eq.5-1-12-1}
\expandafter\def\csname sga@target@IV.5.1.12.2\endcsname{ega4.c04.eq.5-1-12-2}
\expandafter\def\csname sga@target@IV.5.1.12.3\endcsname{ega4.c04.eq.5-1-12-3}
\expandafter\def\csname sga@target@IV.5.1.13\endcsname{ega4.c04.prop.5-1-13}
\expandafter\def\csname sga@target@IV.5.1.13.1\endcsname{ega4.c04.eq.5-1-13-1}
\expandafter\def\csname sga@target@IV.5.1.2\endcsname{ega4.c04.prop.5-1-2}
\expandafter\def\csname sga@target@IV.5.1.2.1\endcsname{ega4.c04.eq.5-1-2-1}
\expandafter\def\csname sga@target@IV.5.1.3\endcsname{ega4.c04.cor.5-1-3}
\expandafter\def\csname sga@target@IV.5.1.3.1\endcsname{ega4.c04.eq.5-1-3-1}
\expandafter\def\csname sga@target@IV.5.1.3.2\endcsname{ega4.c04.eq.5-1-3-2}
\expandafter\def\csname sga@target@IV.5.1.4\endcsname{ega4.c04.prop.5-1-4}
\expandafter\def\csname sga@target@IV.5.1.4.1\endcsname{ega4.c04.eq.5-1-4-1}
\expandafter\def\csname sga@target@IV.5.1.4.2\endcsname{ega4.c04.eq.5-1-4-2}
\expandafter\def\csname sga@target@IV.5.1.5\endcsname{ega4.c04.cor.5-1-5}
\expandafter\def\csname sga@target@IV.5.1.6\endcsname{ega4.c04.statement.5-1-6}
\expandafter\def\csname sga@target@IV.5.1.7\endcsname{ega4.c04.prop.5-1-7}
\expandafter\def\csname sga@target@IV.5.1.8\endcsname{ega4.c04.cor.5-1-8}
\expandafter\def\csname sga@target@IV.5.1.9\endcsname{ega4.c04.prop.5-1-9}
\expandafter\def\csname sga@target@IV.5.1.9.1\endcsname{ega4.c04.statement.5-1-9-1}
\expandafter\def\csname sga@target@IV.5.10.1\endcsname{ega4.c04.statement.5-10-1}
\expandafter\def\csname sga@target@IV.5.10.1.1\endcsname{ega4.c04.eq.5-10-1-1}
\expandafter\def\csname sga@target@IV.5.10.10\endcsname{ega4.c04.prop.5-10-10}
\expandafter\def\csname sga@target@IV.5.10.11\endcsname{ega4.c04.statement.5-10-11}
\expandafter\def\csname sga@target@IV.5.10.11.1\endcsname{ega4.c04.eq.5-10-11-1}
\expandafter\def\csname sga@target@IV.5.10.12\endcsname{ega4.c04.prop.5-10-12}
\expandafter\def\csname sga@target@IV.5.10.13\endcsname{ega4.c04.statement.5-10-13}
\expandafter\def\csname sga@target@IV.5.10.14\endcsname{ega4.c04.prop.5-10-14}
\expandafter\def\csname sga@target@IV.5.10.15\endcsname{ega4.c04.cor.5-10-15}
\expandafter\def\csname sga@target@IV.5.10.16\endcsname{ega4.c04.prop.5-10-16}
\expandafter\def\csname sga@target@IV.5.10.17\endcsname{ega4.c04.prop.5-10-17}
\expandafter\def\csname sga@target@IV.5.10.2\endcsname{ega4.c04.prop.5-10-2}
\expandafter\def\csname sga@target@IV.5.10.3\endcsname{ega4.c04.cor.5-10-3}
\expandafter\def\csname sga@target@IV.5.10.4\endcsname{ega4.c04.cor.5-10-4}
\expandafter\def\csname sga@target@IV.5.10.5\endcsname{ega4.c04.thm.5-10-5}
\expandafter\def\csname sga@target@IV.5.10.6\endcsname{ega4.c04.cor.5-10-6}
\expandafter\def\csname sga@target@IV.5.10.7\endcsname{ega4.c04.thm.5-10-7}
\expandafter\def\csname sga@target@IV.5.10.8\endcsname{ega4.c04.cor.5-10-8}
\expandafter\def\csname sga@target@IV.5.10.8.1\endcsname{ega4.c04.lem.5-10-8-1}
\expandafter\def\csname sga@target@IV.5.10.9\endcsname{ega4.c04.cor.5-10-9}
\expandafter\def\csname sga@target@IV.5.11.1\endcsname{ega4.c04.prop.5-11-1}
\expandafter\def\csname sga@target@IV.5.11.1.1\endcsname{ega4.c04.lem.5-11-1-1}
\expandafter\def\csname sga@target@IV.5.11.2\endcsname{ega4.c04.cor.5-11-2}
\expandafter\def\csname sga@target@IV.5.11.3\endcsname{ega4.c04.rem.5-11-3}
\expandafter\def\csname sga@target@IV.5.11.4\endcsname{ega4.c04.cor.5-11-4}
\expandafter\def\csname sga@target@IV.5.11.5\endcsname{ega4.c04.cor.5-11-5}
\expandafter\def\csname sga@target@IV.5.11.6\endcsname{ega4.c04.cor.5-11-6}
\expandafter\def\csname sga@target@IV.5.11.7\endcsname{ega4.c04.item.5-11-7}
\expandafter\def\csname sga@target@IV.5.12\endcsname{ega4.c04.subsec.5-12.key-iv-5-12}
\expandafter\def\csname sga@target@IV.5.12.1\endcsname{ega4.c04.prop.5-12-1}
\expandafter\def\csname sga@target@IV.5.12.1.1\endcsname{ega4.c04.eq.5-12-1-1}
\expandafter\def\csname sga@target@IV.5.12.10\endcsname{ega4.c04.cor.5-12-10}
\expandafter\def\csname sga@target@IV.5.12.2\endcsname{ega4.c04.prop.5-12-2}
\expandafter\def\csname sga@target@IV.5.12.3\endcsname{ega4.c04.rem.5-12-3}
\expandafter\def\csname sga@target@IV.5.12.4\endcsname{ega4.c04.cor.5-12-4}
\expandafter\def\csname sga@target@IV.5.12.5\endcsname{ega4.c04.prop.5-12-5}
\expandafter\def\csname sga@target@IV.5.12.6\endcsname{ega4.c04.rem.5-12-6}
\expandafter\def\csname sga@target@IV.5.12.7\endcsname{ega4.c04.cor.5-12-7}
\expandafter\def\csname sga@target@IV.5.12.8\endcsname{ega4.c04.prop.5-12-8}
\expandafter\def\csname sga@target@IV.5.12.9\endcsname{ega4.c04.rem.5-12-9}
\expandafter\def\csname sga@target@IV.5.13.1\endcsname{ega4.c04.prop.5-13-1}
\expandafter\def\csname sga@target@IV.5.13.2\endcsname{ega4.c04.cor.5-13-2}
\expandafter\def\csname sga@target@IV.5.13.3\endcsname{ega4.c04.prop.5-13-3}
\expandafter\def\csname sga@target@IV.5.13.4\endcsname{ega4.c04.cor.5-13-4}
\expandafter\def\csname sga@target@IV.5.13.5\endcsname{ega4.c04.statement.5-13-5}
\expandafter\def\csname sga@target@IV.5.13.6\endcsname{ega4.c04.cor.5-13-6}
\expandafter\def\csname sga@target@IV.5.13.7\endcsname{ega4.c04.prop.5-13-7}
\expandafter\def\csname sga@target@IV.5.13.7.1\endcsname{ega4.c04.lem.5-13-7-1}
\expandafter\def\csname sga@target@IV.5.13.8\endcsname{ega4.c04.cor.5-13-8}
\expandafter\def\csname sga@target@IV.5.2.1\endcsname{ega4.c04.prop.5-2-1}
\expandafter\def\csname sga@target@IV.5.2.1.1\endcsname{ega4.c04.eq.5-2-1-1}
\expandafter\def\csname sga@target@IV.5.2.2\endcsname{ega4.c04.cor.5-2-2}
\expandafter\def\csname sga@target@IV.5.2.3\endcsname{ega4.c04.cor.5-2-3}
\expandafter\def\csname sga@target@IV.5.2.3.1\endcsname{ega4.c04.eq.5-2-3-1}
\expandafter\def\csname sga@target@IV.5.2.4\endcsname{ega4.c04.cor.5-2-4}
\expandafter\def\csname sga@target@IV.5.2.5\endcsname{ega4.c04.sec.5-2-5}
\expandafter\def\csname sga@target@IV.5.2.5.i\endcsname{ega4.c04.sec.5-2-5-i}
\expandafter\def\csname sga@target@IV.5.2.5.ii\endcsname{ega4.c04.sec.5-2-5-ii}
\expandafter\def\csname sga@target@IV.5.3.1\endcsname{ega4.c04.prop.5-3-1}
\expandafter\def\csname sga@target@IV.5.3.1.1\endcsname{ega4.c04.eq.5-3-1-1}
\expandafter\def\csname sga@target@IV.5.3.1.2\endcsname{ega4.c04.eq.5-3-1-2}
\expandafter\def\csname sga@target@IV.5.3.1.3\endcsname{ega4.c04.lem.5-3-1-3}
\expandafter\def\csname sga@target@IV.5.4.1\endcsname{ega4.c04.prop.5-4-1}
\expandafter\def\csname sga@target@IV.5.4.1.i\endcsname{ega4.c04.sec.5-4-1-i}
\expandafter\def\csname sga@target@IV.5.4.1.ii\endcsname{ega4.c04.sec.5-4-1-ii}
\expandafter\def\csname sga@target@IV.5.4.2\endcsname{ega4.c04.cor.5-4-2}
\expandafter\def\csname sga@target@IV.5.4.3\endcsname{ega4.c04.sec.5-4-3}
\expandafter\def\csname sga@target@IV.5.4.3.1\endcsname{ega4.c04.eq.5-4-3-1}
\expandafter\def\csname sga@target@IV.5.4.3.i\endcsname{ega4.c04.sec.5-4-3-i}
\expandafter\def\csname sga@target@IV.5.4.3.ii\endcsname{ega4.c04.sec.5-4-3-ii}
\expandafter\def\csname sga@target@IV.5.5\endcsname{ega4.c04.sec.5-5}
\expandafter\def\csname sga@target@IV.5.5.1\endcsname{ega4.c04.statement.5-5-1}
\expandafter\def\csname sga@target@IV.5.5.1.1\endcsname{ega4.c04.eq.5-5-1-1}
\expandafter\def\csname sga@target@IV.5.5.2\endcsname{ega4.c04.prop.5-5-2}
\expandafter\def\csname sga@target@IV.5.5.2.1\endcsname{ega4.c04.eq.5-5-2-1}
\expandafter\def\csname sga@target@IV.5.5.2.2\endcsname{ega4.c04.eq.5-5-2-2}
\expandafter\def\csname sga@target@IV.5.5.3\endcsname{ega4.c04.prop.5-5-3}
\expandafter\def\csname sga@target@IV.5.5.3.1\endcsname{ega4.c04.eq.5-5-3-1}
\expandafter\def\csname sga@target@IV.5.5.4\endcsname{ega4.c04.cor.5-5-4}
\expandafter\def\csname sga@target@IV.5.5.4.1\endcsname{ega4.c04.eq.5-5-4-1}
\expandafter\def\csname sga@target@IV.5.5.5\endcsname{ega4.c04.cor.5-5-5}
\expandafter\def\csname sga@target@IV.5.5.6\endcsname{ega4.c04.cor.5-5-6}
\expandafter\def\csname sga@target@IV.5.5.6.1\endcsname{ega4.c04.eq.5-5-6-1}
\expandafter\def\csname sga@target@IV.5.5.7\endcsname{ega4.c04.lem.5-5-7}
\expandafter\def\csname sga@target@IV.5.5.7.1\endcsname{ega4.c04.eq.5-5-7-1}
\expandafter\def\csname sga@target@IV.5.5.7.2\endcsname{ega4.c04.statement.5-5-7-2}
\expandafter\def\csname sga@target@IV.5.5.7.3\endcsname{ega4.c04.eq.5-5-7-3}
\expandafter\def\csname sga@target@IV.5.5.7.i\endcsname{ega4.c04.sec.5-5-7-i}
\expandafter\def\csname sga@target@IV.5.5.7.ii\endcsname{ega4.c04.sec.5-5-7-ii}
\expandafter\def\csname sga@target@IV.5.5.8\endcsname{ega4.c04.thm.5-5-8}
\expandafter\def\csname sga@target@IV.5.5.8.1\endcsname{ega4.c04.eq.5-5-8-1}
\expandafter\def\csname sga@target@IV.5.5.8.2\endcsname{ega4.c04.eq.5-5-8-2}
\expandafter\def\csname sga@target@IV.5.5.8.3\endcsname{ega4.c04.eq.5-5-8-3}
\expandafter\def\csname sga@target@IV.5.6\endcsname{ega4.c04.sec.5-6}
\expandafter\def\csname sga@target@IV.5.6.1\endcsname{ega4.c04.prop.5-6-1}
\expandafter\def\csname sga@target@IV.5.6.1.1\endcsname{ega4.c04.eq.5-6-1-1}
\expandafter\def\csname sga@target@IV.5.6.1.2\endcsname{ega4.c04.eq.5-6-1-2}
\expandafter\def\csname sga@target@IV.5.6.1.3\endcsname{ega4.c04.statement.5-6-1-3}
\expandafter\def\csname sga@target@IV.5.6.1.4\endcsname{ega4.c04.statement.5-6-1-4}
\expandafter\def\csname sga@target@IV.5.6.1.5\endcsname{ega4.c04.statement.5-6-1-5}
\expandafter\def\csname sga@target@IV.5.6.1.a\endcsname{ega4.c04.sec.5-6-1-a}
\expandafter\def\csname sga@target@IV.5.6.1.b\endcsname{ega4.c04.sec.5-6-1-b}
\expandafter\def\csname sga@target@IV.5.6.1.c\endcsname{ega4.c04.sec.5-6-1-c}
\expandafter\def\csname sga@target@IV.5.6.10\endcsname{ega4.c04.prop.5-6-10}
\expandafter\def\csname sga@target@IV.5.6.11\endcsname{ega4.c04.ex.5-6-11}
\expandafter\def\csname sga@target@IV.5.6.11.1\endcsname{ega4.c04.lem.5-6-11-1}
\expandafter\def\csname sga@target@IV.5.6.11.1.i\endcsname{ega4.c04.sec.5-6-11-1-i}
\expandafter\def\csname sga@target@IV.5.6.11.1.ii\endcsname{ega4.c04.sec.5-6-11-1-ii}
\expandafter\def\csname sga@target@IV.5.6.2\endcsname{ega4.c04.def.5-6-2}
\expandafter\def\csname sga@target@IV.5.6.3\endcsname{ega4.c04.sec.5-6-3}
\expandafter\def\csname sga@target@IV.5.6.3.i\endcsname{ega4.c04.sec.5-6-3-i}
\expandafter\def\csname sga@target@IV.5.6.3.ii\endcsname{ega4.c04.sec.5-6-3-ii}
\expandafter\def\csname sga@target@IV.5.6.3.iii\endcsname{ega4.c04.sec.5-6-3-iii}
\expandafter\def\csname sga@target@IV.5.6.3.iv\endcsname{ega4.c04.sec.5-6-3-iv}
\expandafter\def\csname sga@target@IV.5.6.4\endcsname{ega4.c04.prop.5-6-4}
\expandafter\def\csname sga@target@IV.5.6.5\endcsname{ega4.c04.prop.5-6-5}
\expandafter\def\csname sga@target@IV.5.6.5.1\endcsname{ega4.c04.eq.5-6-5-1}
\expandafter\def\csname sga@target@IV.5.6.5.2\endcsname{ega4.c04.eq.5-6-5-2}
\expandafter\def\csname sga@target@IV.5.6.5.3\endcsname{ega4.c04.eq.5-6-5-3}
\expandafter\def\csname sga@target@IV.5.6.6\endcsname{ega4.c04.cor.5-6-6}
\expandafter\def\csname sga@target@IV.5.6.6.1\endcsname{ega4.c04.eq.5-6-6-1}
\expandafter\def\csname sga@target@IV.5.6.6.2\endcsname{ega4.c04.eq.5-6-6-2}
\expandafter\def\csname sga@target@IV.5.6.7\endcsname{ega4.c04.cor.5-6-7}
\expandafter\def\csname sga@target@IV.5.6.7.1\endcsname{ega4.c04.eq.5-6-7-1}
\expandafter\def\csname sga@target@IV.5.6.8\endcsname{ega4.c04.cor.5-6-8}
\expandafter\def\csname sga@target@IV.5.6.8.1\endcsname{ega4.c04.eq.5-6-8-1}
\expandafter\def\csname sga@target@IV.5.6.8.2\endcsname{ega4.c04.eq.5-6-8-2}
\expandafter\def\csname sga@target@IV.5.6.9\endcsname{ega4.c04.sec.5-6-9}
\expandafter\def\csname sga@target@IV.5.6.9.i\endcsname{ega4.c04.sec.5-6-9-i}
\expandafter\def\csname sga@target@IV.5.6.9.ii\endcsname{ega4.c04.sec.5-6-9-ii}
\expandafter\def\csname sga@target@IV.5.7.1\endcsname{ega4.c04.def.5-7-1}
\expandafter\def\csname sga@target@IV.5.7.1.1\endcsname{ega4.c04.eq.5-7-1-1}
\expandafter\def\csname sga@target@IV.5.7.10\endcsname{ega4.c04.cor.5-7-10}
\expandafter\def\csname sga@target@IV.5.7.11\endcsname{ega4.c04.cor.5-7-11}
\expandafter\def\csname sga@target@IV.5.7.11.i\endcsname{ega4.c04.sec.5-7-11-i}
\expandafter\def\csname sga@target@IV.5.7.11.ii\endcsname{ega4.c04.sec.5-7-11-ii}
\expandafter\def\csname sga@target@IV.5.7.12\endcsname{ega4.c04.statement.5-7-12}
\expandafter\def\csname sga@target@IV.5.7.13\endcsname{ega4.c04.cor.5-7-13}
\expandafter\def\csname sga@target@IV.5.7.13.1\endcsname{ega4.c04.eq.5-7-13-1}
\expandafter\def\csname sga@target@IV.5.7.2\endcsname{ega4.c04.def.5-7-2}
\expandafter\def\csname sga@target@IV.5.7.2.1\endcsname{ega4.c04.eq.5-7-2-1}
\expandafter\def\csname sga@target@IV.5.7.2.2\endcsname{ega4.c04.eq.5-7-2-2}
\expandafter\def\csname sga@target@IV.5.7.3\endcsname{ega4.c04.sec.5-7-3}
\expandafter\def\csname sga@target@IV.5.7.3.i\endcsname{ega4.c04.sec.5-7-3-i}
\expandafter\def\csname sga@target@IV.5.7.3.ii\endcsname{ega4.c04.sec.5-7-3-ii}
\expandafter\def\csname sga@target@IV.5.7.3.iii\endcsname{ega4.c04.sec.5-7-3-iii}
\expandafter\def\csname sga@target@IV.5.7.3.iv\endcsname{ega4.c04.sec.5-7-3-iv}
\expandafter\def\csname sga@target@IV.5.7.3.v\endcsname{ega4.c04.sec.5-7-3-v}
\expandafter\def\csname sga@target@IV.5.7.3.vi\endcsname{ega4.c04.sec.5-7-3-vi}
\expandafter\def\csname sga@target@IV.5.7.4\endcsname{ega4.c04.prop.5-7-4}
\expandafter\def\csname sga@target@IV.5.7.4.1\endcsname{ega4.c04.eq.5-7-4-1}
\expandafter\def\csname sga@target@IV.5.7.4.2\endcsname{ega4.c04.eq.5-7-4-2}
\expandafter\def\csname sga@target@IV.5.7.4.i\endcsname{ega4.c04.sec.5-7-4-i}
\expandafter\def\csname sga@target@IV.5.7.4.ii\endcsname{ega4.c04.sec.5-7-4-ii}
\expandafter\def\csname sga@target@IV.5.7.5\endcsname{ega4.c04.prop.5-7-5}
\expandafter\def\csname sga@target@IV.5.7.6\endcsname{ega4.c04.cor.5-7-6}
\expandafter\def\csname sga@target@IV.5.7.7\endcsname{ega4.c04.cor.5-7-7}
\expandafter\def\csname sga@target@IV.5.7.8\endcsname{ega4.c04.sec.5-7-8}
\expandafter\def\csname sga@target@IV.5.7.9\endcsname{ega4.c04.prop.5-7-9}
\expandafter\def\csname sga@target@IV.5.7.9.1\endcsname{ega4.c04.eq.5-7-9-1}
\expandafter\def\csname sga@target@IV.5.7.9.2\endcsname{ega4.c04.eq.5-7-9-2}
\expandafter\def\csname sga@target@IV.5.8.1\endcsname{ega4.c04.statement.5-8-1}
\expandafter\def\csname sga@target@IV.5.8.2\endcsname{ega4.c04.def.5-8-2}
\expandafter\def\csname sga@target@IV.5.8.3\endcsname{ega4.c04.prop.5-8-3}
\expandafter\def\csname sga@target@IV.5.8.4\endcsname{ega4.c04.statement.5-8-4}
\expandafter\def\csname sga@target@IV.5.8.5\endcsname{ega4.c04.prop.5-8-5}
\expandafter\def\csname sga@target@IV.5.8.6\endcsname{ega4.c04.thm.5-8-6}
\expandafter\def\csname sga@target@IV.5.8.6.1\endcsname{ega4.c04.eq.5-8-6-1}
\expandafter\def\csname sga@target@IV.5.8.6.i\endcsname{ega4.c04.item.5-8-6-i}
\expandafter\def\csname sga@target@IV.5.8.6.ii\endcsname{ega4.c04.item.5-8-6-ii}
\expandafter\def\csname sga@target@IV.5.9.1\endcsname{ega4.c04.statement.5-9-1}
\expandafter\def\csname sga@target@IV.5.9.1.1\endcsname{ega4.c04.eq.5-9-1-1}
\expandafter\def\csname sga@target@IV.5.9.1.2\endcsname{ega4.c04.eq.5-9-1-2}
\expandafter\def\csname sga@target@IV.5.9.1.3\endcsname{ega4.c04.eq.5-9-1-3}
\expandafter\def\csname sga@target@IV.5.9.10\endcsname{ega4.c04.cor.5-9-10}
\expandafter\def\csname sga@target@IV.5.9.10.i\endcsname{ega4.c04.item.5-9-10-i}
\expandafter\def\csname sga@target@IV.5.9.10.ii\endcsname{ega4.c04.item.5-9-10-ii}
\expandafter\def\csname sga@target@IV.5.9.11\endcsname{ega4.c04.prop.5-9-11}
\expandafter\def\csname sga@target@IV.5.9.12\endcsname{ega4.c04.item.5-9-12}
\expandafter\def\csname sga@target@IV.5.9.12.i\endcsname{ega4.c04.item.5-9-12-i}
\expandafter\def\csname sga@target@IV.5.9.12.ii\endcsname{ega4.c04.item.5-9-12-ii}
\expandafter\def\csname sga@target@IV.5.9.12.iii\endcsname{ega4.c04.item.5-9-12-iii}
\expandafter\def\csname sga@target@IV.5.9.2\endcsname{ega4.c04.prop.5-9-2}
\expandafter\def\csname sga@target@IV.5.9.2.i\endcsname{ega4.c04.item.5-9-2-i}
\expandafter\def\csname sga@target@IV.5.9.2.ii\endcsname{ega4.c04.item.5-9-2-ii}
\expandafter\def\csname sga@target@IV.5.9.3\endcsname{ega4.c04.rem.5-9-3}
\expandafter\def\csname sga@target@IV.5.9.3.1\endcsname{ega4.c04.eq.5-9-3-1}
\expandafter\def\csname sga@target@IV.5.9.4\endcsname{ega4.c04.prop.5-9-4}
\expandafter\def\csname sga@target@IV.5.9.4.1\endcsname{ega4.c04.eq.5-9-4-1}
\expandafter\def\csname sga@target@IV.5.9.5\endcsname{ega4.c04.cor.5-9-5}
\expandafter\def\csname sga@target@IV.5.9.6\endcsname{ega4.c04.cor.5-9-6}
\expandafter\def\csname sga@target@IV.5.9.6.1\endcsname{ega4.c04.eq.5-9-6-1}
\expandafter\def\csname sga@target@IV.5.9.7\endcsname{ega4.c04.statement.5-9-7}
\expandafter\def\csname sga@target@IV.5.9.7.1\endcsname{ega4.c04.eq.5-9-7-1}
\expandafter\def\csname sga@target@IV.5.9.8\endcsname{ega4.c04.prop.5-9-8}
\expandafter\def\csname sga@target@IV.5.9.8.1\endcsname{ega4.c04.eq.5-9-8-1}
\expandafter\def\csname sga@target@IV.5.9.8.a\endcsname{ega4.c04.item.5-9-8-a}
\expandafter\def\csname sga@target@IV.5.9.8.a-prime\endcsname{ega4.c04.item.5-9-8-a-prime}
\expandafter\def\csname sga@target@IV.5.9.8.b\endcsname{ega4.c04.item.5-9-8-b}
\expandafter\def\csname sga@target@IV.5.9.8.b-prime\endcsname{ega4.c04.item.5-9-8-b-prime}
\expandafter\def\csname sga@target@IV.5.9.9\endcsname{ega4.c04.def.5-9-9}
\expandafter\def\csname sga@target@IV.6.1.1\endcsname{ega4.c04.prop.6-1-1}
\expandafter\def\csname sga@target@IV.6.1.1.1\endcsname{ega4.c04.eq.6-1-1-1}
\expandafter\def\csname sga@target@IV.6.1.2\endcsname{ega4.c04.cor.6-1-2}
\expandafter\def\csname sga@target@IV.6.1.2.1\endcsname{ega4.c04.eq.6-1-2-1}
\expandafter\def\csname sga@target@IV.6.1.3\endcsname{ega4.c04.cor.6-1-3}
\expandafter\def\csname sga@target@IV.6.1.4\endcsname{ega4.c04.cor.6-1-4}
\expandafter\def\csname sga@target@IV.6.1.4.1\endcsname{ega4.c04.eq.6-1-4-1}
\expandafter\def\csname sga@target@IV.6.1.5\endcsname{ega4.c04.prop.6-1-5}
\expandafter\def\csname sga@target@IV.6.10.1\endcsname{ega4.c04.statement.6-10-1}
\expandafter\def\csname sga@target@IV.6.10.2\endcsname{ega4.c04.prop.6-10-2}
\expandafter\def\csname sga@target@IV.6.10.3\endcsname{ega4.c04.prop.6-10-3}
\expandafter\def\csname sga@target@IV.6.10.3.1\endcsname{ega4.c04.eq.6-10-3-1}
\expandafter\def\csname sga@target@IV.6.10.3.2\endcsname{ega4.c04.eq.6-10-3-2}
\expandafter\def\csname sga@target@IV.6.10.3.3\endcsname{ega4.c04.eq.6-10-3-3}
\expandafter\def\csname sga@target@IV.6.10.4\endcsname{ega4.c04.cor.6-10-4}
\expandafter\def\csname sga@target@IV.6.10.5\endcsname{ega4.c04.prop.6-10-5}
\expandafter\def\csname sga@target@IV.6.10.5.1\endcsname{ega4.c04.eq.6-10-5-1}
\expandafter\def\csname sga@target@IV.6.10.5.2\endcsname{ega4.c04.eq.6-10-5-2}
\expandafter\def\csname sga@target@IV.6.10.5.3\endcsname{ega4.c04.eq.6-10-5-3}
\expandafter\def\csname sga@target@IV.6.10.5.4\endcsname{ega4.c04.eq.6-10-5-4}
\expandafter\def\csname sga@target@IV.6.10.6\endcsname{ega4.c04.prop.6-10-6}
\expandafter\def\csname sga@target@IV.6.10.6.1\endcsname{ega4.c04.eq.6-10-6-1}
\expandafter\def\csname sga@target@IV.6.10.6.2\endcsname{ega4.c04.eq.6-10-6-2}
\expandafter\def\csname sga@target@IV.6.10.6.3\endcsname{ega4.c04.eq.6-10-6-3}
\expandafter\def\csname sga@target@IV.6.11\endcsname{ega4.c04.subsec.6-11.key-iv-6-11}
\expandafter\def\csname sga@target@IV.6.11.1\endcsname{ega4.c04.lem.6-11-1}
\expandafter\def\csname sga@target@IV.6.11.2\endcsname{ega4.c04.prop.6-11-2}
\expandafter\def\csname sga@target@IV.6.11.2.1\endcsname{ega4.c04.eq.6-11-2-1}
\expandafter\def\csname sga@target@IV.6.11.2.2\endcsname{ega4.c04.eq.6-11-2-2}
\expandafter\def\csname sga@target@IV.6.11.3\endcsname{ega4.c04.cor.6-11-3}
\expandafter\def\csname sga@target@IV.6.11.4\endcsname{ega4.c04.statement.6-11-4}
\expandafter\def\csname sga@target@IV.6.11.5\endcsname{ega4.c04.prop.6-11-5}
\expandafter\def\csname sga@target@IV.6.11.5.1\endcsname{ega4.c04.eq.6-11-5-1}
\expandafter\def\csname sga@target@IV.6.11.6\endcsname{ega4.c04.prop.6-11-6}
\expandafter\def\csname sga@target@IV.6.11.6.1\endcsname{ega4.c04.lem.6-11-6-1}
\expandafter\def\csname sga@target@IV.6.11.7\endcsname{ega4.c04.cor.6-11-7}
\expandafter\def\csname sga@target@IV.6.11.8\endcsname{ega4.c04.prop.6-11-8}
\expandafter\def\csname sga@target@IV.6.11.8.1\endcsname{ega4.c04.eq.6-11-8-1}
\expandafter\def\csname sga@target@IV.6.11.9\endcsname{ega4.c04.statement.6-11-9}
\expandafter\def\csname sga@target@IV.6.12\endcsname{ega4.c04.subsec.6-12.key-iv-6-12}
\expandafter\def\csname sga@target@IV.6.12.1\endcsname{ega4.c04.statement.6-12-1}
\expandafter\def\csname sga@target@IV.6.12.10\endcsname{ega4.c04.prop.6-12-10}
\expandafter\def\csname sga@target@IV.6.12.2\endcsname{ega4.c04.prop.6-12-2}
\expandafter\def\csname sga@target@IV.6.12.3\endcsname{ega4.c04.cor.6-12-3}
\expandafter\def\csname sga@target@IV.6.12.4\endcsname{ega4.c04.thm.6-12-4}
\expandafter\def\csname sga@target@IV.6.12.4.1\endcsname{ega4.c04.lem.6-12-4-1}
\expandafter\def\csname sga@target@IV.6.12.5\endcsname{ega4.c04.cor.6-12-5}
\expandafter\def\csname sga@target@IV.6.12.6\endcsname{ega4.c04.cor.6-12-6}
\expandafter\def\csname sga@target@IV.6.12.7\endcsname{ega4.c04.thm.6-12-7}
\expandafter\def\csname sga@target@IV.6.12.8\endcsname{ega4.c04.cor.6-12-8}
\expandafter\def\csname sga@target@IV.6.12.9\endcsname{ega4.c04.prop.6-12-9}
\expandafter\def\csname sga@target@IV.6.13\endcsname{ega4.c04.subsec.6-13.key-iv-6-13}
\expandafter\def\csname sga@target@IV.6.13.1\endcsname{ega4.c04.statement.6-13-1}
\expandafter\def\csname sga@target@IV.6.13.2\endcsname{ega4.c04.prop.6-13-2}
\expandafter\def\csname sga@target@IV.6.13.3\endcsname{ega4.c04.cor.6-13-3}
\expandafter\def\csname sga@target@IV.6.13.4\endcsname{ega4.c04.prop.6-13-4}
\expandafter\def\csname sga@target@IV.6.13.5\endcsname{ega4.c04.cor.6-13-5}
\expandafter\def\csname sga@target@IV.6.13.6\endcsname{ega4.c04.prop.6-13-6}
\expandafter\def\csname sga@target@IV.6.13.7\endcsname{ega4.c04.prop.6-13-7}
\expandafter\def\csname sga@target@IV.6.13.7.1\endcsname{ega4.c04.lem.6-13-7-1}
\expandafter\def\csname sga@target@IV.6.14.1\endcsname{ega4.c04.prop.6-14-1}
\expandafter\def\csname sga@target@IV.6.14.1.1\endcsname{ega4.c04.lem.6-14-1-1}
\expandafter\def\csname sga@target@IV.6.14.1.2\endcsname{ega4.c04.lem.6-14-1-2}
\expandafter\def\csname sga@target@IV.6.14.1.2.env\endcsname{ega4.c04.statement.6-14-1-2-env}
\expandafter\def\csname sga@target@IV.6.14.1.3\endcsname{ega4.c04.eq.6-14-1-3}
\expandafter\def\csname sga@target@IV.6.14.1.4\endcsname{ega4.c04.lem.6-14-1-4}
\expandafter\def\csname sga@target@IV.6.14.2\endcsname{ega4.c04.cor.6-14-2}
\expandafter\def\csname sga@target@IV.6.14.3\endcsname{ega4.c04.cor.6-14-3}
\expandafter\def\csname sga@target@IV.6.14.4\endcsname{ega4.c04.prop.6-14-4}
\expandafter\def\csname sga@target@IV.6.14.4.1\endcsname{ega4.c04.lem.6-14-4-1}
\expandafter\def\csname sga@target@IV.6.14.5\endcsname{ega4.c04.cor.6-14-5}
\expandafter\def\csname sga@target@IV.6.15.1\endcsname{ega4.c04.statement.6-15-1}
\expandafter\def\csname sga@target@IV.6.15.10\endcsname{ega4.c04.prop.6-15-10}
\expandafter\def\csname sga@target@IV.6.15.11\endcsname{ega4.c04.item.6-15-11}
\expandafter\def\csname sga@target@IV.6.15.2\endcsname{ega4.c04.statement.6-15-2}
\expandafter\def\csname sga@target@IV.6.15.3\endcsname{ega4.c04.statement.6-15-3}
\expandafter\def\csname sga@target@IV.6.15.3.1\endcsname{ega4.c04.lem.6-15-3-1}
\expandafter\def\csname sga@target@IV.6.15.4\endcsname{ega4.c04.statement.6-15-4}
\expandafter\def\csname sga@target@IV.6.15.4.1\endcsname{ega4.c04.lem.6-15-4-1}
\expandafter\def\csname sga@target@IV.6.15.5\endcsname{ega4.c04.prop.6-15-5}
\expandafter\def\csname sga@target@IV.6.15.6\endcsname{ega4.c04.prop.6-15-6}
\expandafter\def\csname sga@target@IV.6.15.7\endcsname{ega4.c04.prop.6-15-7}
\expandafter\def\csname sga@target@IV.6.15.8\endcsname{ega4.c04.lem.6-15-8}
\expandafter\def\csname sga@target@IV.6.15.9\endcsname{ega4.c04.cor.6-15-9}
\expandafter\def\csname sga@target@IV.6.2.1\endcsname{ega4.c04.prop.6-2-1}
\expandafter\def\csname sga@target@IV.6.2.1.1\endcsname{ega4.c04.eq.6-2-1-1}
\expandafter\def\csname sga@target@IV.6.2.1.2\endcsname{ega4.c04.eq.6-2-1-2}
\expandafter\def\csname sga@target@IV.6.2.2\endcsname{ega4.c04.cor.6-2-2}
\expandafter\def\csname sga@target@IV.6.2.3\endcsname{ega4.c04.prop.6-2-3}
\expandafter\def\csname sga@target@IV.6.2.3.1\endcsname{ega4.c04.eq.6-2-3-1}
\expandafter\def\csname sga@target@IV.6.2.4\endcsname{ega4.c04.prop.6-2-4}
\expandafter\def\csname sga@target@IV.6.2.4.1\endcsname{ega4.c04.eq.6-2-4-1}
\expandafter\def\csname sga@target@IV.6.3.1\endcsname{ega4.c04.prop.6-3-1}
\expandafter\def\csname sga@target@IV.6.3.1.1\endcsname{ega4.c04.eq.6-3-1-1}
\expandafter\def\csname sga@target@IV.6.3.2\endcsname{ega4.c04.cor.6-3-2}
\expandafter\def\csname sga@target@IV.6.3.2.1\endcsname{ega4.c04.eq.6-3-2-1}
\expandafter\def\csname sga@target@IV.6.3.3\endcsname{ega4.c04.cor.6-3-3}
\expandafter\def\csname sga@target@IV.6.3.4\endcsname{ega4.c04.cor.6-3-4}
\expandafter\def\csname sga@target@IV.6.3.4.1\endcsname{ega4.c04.eq.6-3-4-1}
\expandafter\def\csname sga@target@IV.6.3.4.2\endcsname{ega4.c04.eq.6-3-4-2}
\expandafter\def\csname sga@target@IV.6.3.5\endcsname{ega4.c04.cor.6-3-5}
\expandafter\def\csname sga@target@IV.6.3.6\endcsname{ega4.c04.cor.6-3-6}
\expandafter\def\csname sga@target@IV.6.3.7\endcsname{ega4.c04.cor.6-3-7}
\expandafter\def\csname sga@target@IV.6.3.8\endcsname{ega4.c04.prop.6-3-8}
\expandafter\def\csname sga@target@IV.6.4.1\endcsname{ega4.c04.prop.6-4-1}
\expandafter\def\csname sga@target@IV.6.4.2\endcsname{ega4.c04.cor.6-4-2}
\expandafter\def\csname sga@target@IV.6.4.3\endcsname{ega4.c04.rem.6-4-3}
\expandafter\def\csname sga@target@IV.6.5.1\endcsname{ega4.c04.prop.6-5-1}
\expandafter\def\csname sga@target@IV.6.5.2\endcsname{ega4.c04.cor.6-5-2}
\expandafter\def\csname sga@target@IV.6.5.3\endcsname{ega4.c04.prop.6-5-3}
\expandafter\def\csname sga@target@IV.6.5.4\endcsname{ega4.c04.cor.6-5-4}
\expandafter\def\csname sga@target@IV.6.5.5\endcsname{ega4.c04.item.6-5-5}
\expandafter\def\csname sga@target@IV.6.6.1\endcsname{ega4.c04.prop.6-6-1}
\expandafter\def\csname sga@target@IV.6.6.2\endcsname{ega4.c04.rem.6-6-2}
\expandafter\def\csname sga@target@IV.6.7.1\endcsname{ega4.c04.prop.6-7-1}
\expandafter\def\csname sga@target@IV.6.7.1.1\endcsname{ega4.c04.lem.6-7-1-1}
\expandafter\def\csname sga@target@IV.6.7.2\endcsname{ega4.c04.cor.6-7-2}
\expandafter\def\csname sga@target@IV.6.7.3\endcsname{ega4.c04.cor.6-7-3}
\expandafter\def\csname sga@target@IV.6.7.4\endcsname{ega4.c04.prop.6-7-4}
\expandafter\def\csname sga@target@IV.6.7.4.1\endcsname{ega4.c04.lem.6-7-4-1}
\expandafter\def\csname sga@target@IV.6.7.5\endcsname{ega4.c04.item.6-7-5}
\expandafter\def\csname sga@target@IV.6.7.6\endcsname{ega4.c04.def.6-7-6}
\expandafter\def\csname sga@target@IV.6.7.7\endcsname{ega4.c04.prop.6-7-7}
\expandafter\def\csname sga@target@IV.6.7.8\endcsname{ega4.c04.cor.6-7-8}
\expandafter\def\csname sga@target@IV.6.8.1\endcsname{ega4.c04.def.6-8-1}
\expandafter\def\csname sga@target@IV.6.8.2\endcsname{ega4.c04.prop.6-8-2}
\expandafter\def\csname sga@target@IV.6.8.3\endcsname{ega4.c04.prop.6-8-3}
\expandafter\def\csname sga@target@IV.6.8.4\endcsname{ega4.c04.item.6-8-4}
\expandafter\def\csname sga@target@IV.6.8.5\endcsname{ega4.c04.prop.6-8-5}
\expandafter\def\csname sga@target@IV.6.8.6\endcsname{ega4.c04.thm.6-8-6}
\expandafter\def\csname sga@target@IV.6.8.7\endcsname{ega4.c04.cor.6-8-7}
\expandafter\def\csname sga@target@IV.6.8.8\endcsname{ega4.c04.rem.6-8-8}
\expandafter\def\csname sga@target@IV.6.9.1\endcsname{ega4.c04.thm.6-9-1}
\expandafter\def\csname sga@target@IV.6.9.2\endcsname{ega4.c04.lem.6-9-2}
\expandafter\def\csname sga@target@IV.6.9.3\endcsname{ega4.c04.cor.6-9-3}
\expandafter\def\csname sga@target@IV.7.1.1\endcsname{ega4.c04.def.7-1-1}
\expandafter\def\csname sga@target@IV.7.1.10\endcsname{ega4.c04.prop.7-1-10}
\expandafter\def\csname sga@target@IV.7.1.11\endcsname{ega4.c04.prop.7-1-11}
\expandafter\def\csname sga@target@IV.7.1.12\endcsname{ega4.c04.item.7-1-12}
\expandafter\def\csname sga@target@IV.7.1.13\endcsname{ega4.c04.cor.7-1-13}
\expandafter\def\csname sga@target@IV.7.1.2\endcsname{ega4.c04.prop.7-1-2}
\expandafter\def\csname sga@target@IV.7.1.3\endcsname{ega4.c04.prop.7-1-3}
\expandafter\def\csname sga@target@IV.7.1.3.1\endcsname{ega4.c04.eq.7-1-3-1}
\expandafter\def\csname sga@target@IV.7.1.3.2\endcsname{ega4.c04.eq.7-1-3-2}
\expandafter\def\csname sga@target@IV.7.1.3.3\endcsname{ega4.c04.eq.7-1-3-3}
\expandafter\def\csname sga@target@IV.7.1.3.4\endcsname{ega4.c04.eq.7-1-3-4}
\expandafter\def\csname sga@target@IV.7.1.3.5\endcsname{ega4.c04.eq.7-1-3-5}
\expandafter\def\csname sga@target@IV.7.1.3.6\endcsname{ega4.c04.eq.7-1-3-6}
\expandafter\def\csname sga@target@IV.7.1.3.7\endcsname{ega4.c04.eq.7-1-3-7}
\expandafter\def\csname sga@target@IV.7.1.4\endcsname{ega4.c04.cor.7-1-4}
\expandafter\def\csname sga@target@IV.7.1.5\endcsname{ega4.c04.cor.7-1-5}
\expandafter\def\csname sga@target@IV.7.1.6\endcsname{ega4.c04.rem.7-1-6}
\expandafter\def\csname sga@target@IV.7.1.7\endcsname{ega4.c04.lem.7-1-7}
\expandafter\def\csname sga@target@IV.7.1.8\endcsname{ega4.c04.thm.7-1-8}
\expandafter\def\csname sga@target@IV.7.1.9\endcsname{ega4.c04.def.7-1-9}
\expandafter\def\csname sga@target@IV.7.2.1\endcsname{ega4.c04.statement.7-2-1}
\expandafter\def\csname sga@target@IV.7.2.1.1\endcsname{ega4.c04.statement.7-2-1-1}
\expandafter\def\csname sga@target@IV.7.2.10\endcsname{ega4.c04.statement.7-2-10}
\expandafter\def\csname sga@target@IV.7.2.2\endcsname{ega4.c04.prop.7-2-2}
\expandafter\def\csname sga@target@IV.7.2.3\endcsname{ega4.c04.prop.7-2-3}
\expandafter\def\csname sga@target@IV.7.2.4\endcsname{ega4.c04.prop.7-2-4}
\expandafter\def\csname sga@target@IV.7.2.5\endcsname{ega4.c04.thm.7-2-5}
\expandafter\def\csname sga@target@IV.7.2.5.1\endcsname{ega4.c04.eq.7-2-5-1}
\expandafter\def\csname sga@target@IV.7.2.6\endcsname{ega4.c04.def.7-2-6}
\expandafter\def\csname sga@target@IV.7.2.7\endcsname{ega4.c04.cor.7-2-7}
\expandafter\def\csname sga@target@IV.7.2.8\endcsname{ega4.c04.cor.7-2-8}
\expandafter\def\csname sga@target@IV.7.2.9\endcsname{ega4.c04.cor.7-2-9}
\expandafter\def\csname sga@target@IV.7.3.1\endcsname{ega4.c04.statement.7-3-1}
\expandafter\def\csname sga@target@IV.7.3.10\endcsname{ega4.c04.statement.7-3-10}
\expandafter\def\csname sga@target@IV.7.3.11\endcsname{ega4.c04.rem.7-3-11}
\expandafter\def\csname sga@target@IV.7.3.12\endcsname{ega4.c04.prop.7-3-12}
\expandafter\def\csname sga@target@IV.7.3.12.1\endcsname{ega4.c04.eq.7-3-12-1}
\expandafter\def\csname sga@target@IV.7.3.13\endcsname{ega4.c04.statement.7-3-13}
\expandafter\def\csname sga@target@IV.7.3.14\endcsname{ega4.c04.prop.7-3-14}
\expandafter\def\csname sga@target@IV.7.3.15\endcsname{ega4.c04.prop.7-3-15}
\expandafter\def\csname sga@target@IV.7.3.16\endcsname{ega4.c04.prop.7-3-16}
\expandafter\def\csname sga@target@IV.7.3.16.1\endcsname{ega4.c04.eq.7-3-16-1}
\expandafter\def\csname sga@target@IV.7.3.16.2\endcsname{ega4.c04.lem.7-3-16-2}
\expandafter\def\csname sga@target@IV.7.3.17\endcsname{ega4.c04.prop.7-3-17}
\expandafter\def\csname sga@target@IV.7.3.18\endcsname{ega4.c04.prop.7-3-18}
\expandafter\def\csname sga@target@IV.7.3.19\endcsname{ega4.c04.statement.7-3-19}
\expandafter\def\csname sga@target@IV.7.3.2\endcsname{ega4.c04.lem.7-3-2}
\expandafter\def\csname sga@target@IV.7.3.3\endcsname{ega4.c04.cor.7-3-3}
\expandafter\def\csname sga@target@IV.7.3.4\endcsname{ega4.c04.statement.7-3-4}
\expandafter\def\csname sga@target@IV.7.3.5\endcsname{ega4.c04.item.7-3-5}
\expandafter\def\csname sga@target@IV.7.3.6\endcsname{ega4.c04.statement.7-3-6}
\expandafter\def\csname sga@target@IV.7.3.7\endcsname{ega4.c04.lem.7-3-7}
\expandafter\def\csname sga@target@IV.7.3.8\endcsname{ega4.c04.statement.7-3-8}
\expandafter\def\csname sga@target@IV.7.3.9\endcsname{ega4.c04.item.7-3-9}
\expandafter\def\csname sga@target@IV.7.4.1\endcsname{ega4.c04.statement.7-4-1}
\expandafter\def\csname sga@target@IV.7.4.2\endcsname{ega4.c04.lem.7-4-2}
\expandafter\def\csname sga@target@IV.7.4.3\endcsname{ega4.c04.cor.7-4-3}
\expandafter\def\csname sga@target@IV.7.4.4\endcsname{ega4.c04.thm.7-4-4}
\expandafter\def\csname sga@target@IV.7.4.4.1\endcsname{ega4.c04.lem.7-4-4-1}
\expandafter\def\csname sga@target@IV.7.4.4.2\endcsname{ega4.c04.lem.7-4-4-2}
\expandafter\def\csname sga@target@IV.7.4.5\endcsname{ega4.c04.cor.7-4-5}
\expandafter\def\csname sga@target@IV.7.4.6\endcsname{ega4.c04.prop.7-4-6}
\expandafter\def\csname sga@target@IV.7.4.7\endcsname{ega4.c04.cor.7-4-7}
\expandafter\def\csname sga@target@IV.7.4.8\endcsname{ega4.c04.rem.7-4-8}
\expandafter\def\csname sga@target@IV.7.5.0\endcsname{ega4.c04.statement.7-5-0}
\expandafter\def\csname sga@target@IV.7.5.1\endcsname{ega4.c04.thm.7-5-1}
\expandafter\def\csname sga@target@IV.7.5.1.1\endcsname{ega4.c04.lem.7-5-1-1}
\expandafter\def\csname sga@target@IV.7.5.2\endcsname{ega4.c04.cor.7-5-2}
\expandafter\def\csname sga@target@IV.7.5.3\endcsname{ega4.c04.item.7-5-3}
\expandafter\def\csname sga@target@IV.7.5.4\endcsname{ega4.c04.item.7-5-4}
\expandafter\def\csname sga@target@IV.7.5.5\endcsname{ega4.c04.lem.7-5-5}
\expandafter\def\csname sga@target@IV.7.5.6\endcsname{ega4.c04.prop.7-5-6}
\expandafter\def\csname sga@target@IV.7.5.7\endcsname{ega4.c04.cor.7-5-7}
\expandafter\def\csname sga@target@IV.7.5.7.1\endcsname{ega4.c04.lem.7-5-7-1}
\expandafter\def\csname sga@target@IV.7.5.8\endcsname{ega4.c04.rem.7-5-8}
\expandafter\def\csname sga@target@IV.7.6.1\endcsname{ega4.c04.prop.7-6-1}
\expandafter\def\csname sga@target@IV.7.6.2\endcsname{ega4.c04.cor.7-6-2}
\expandafter\def\csname sga@target@IV.7.6.3\endcsname{ega4.c04.cor.7-6-3}
\expandafter\def\csname sga@target@IV.7.6.4\endcsname{ega4.c04.thm.7-6-4}
\expandafter\def\csname sga@target@IV.7.6.4.1\endcsname{ega4.c04.lem.7-6-4-1}
\expandafter\def\csname sga@target@IV.7.6.5\endcsname{ega4.c04.cor.7-6-5}
\expandafter\def\csname sga@target@IV.7.6.6\endcsname{ega4.c04.cor.7-6-6}
\expandafter\def\csname sga@target@IV.7.6.7\endcsname{ega4.c04.item.7-6-7}
\expandafter\def\csname sga@target@IV.7.7.1\endcsname{ega4.c04.statement.7-7-1}
\expandafter\def\csname sga@target@IV.7.7.2\endcsname{ega4.c04.thm.7-7-2}
\expandafter\def\csname sga@target@IV.7.7.3\endcsname{ega4.c04.cor.7-7-3}
\expandafter\def\csname sga@target@IV.7.7.4\endcsname{ega4.c04.cor.7-7-4}
\expandafter\def\csname sga@target@IV.7.8.1\endcsname{ega4.c04.statement.7-8-1}
\expandafter\def\csname sga@target@IV.7.8.2\endcsname{ega4.c04.def.7-8-2}
\expandafter\def\csname sga@target@IV.7.8.3\endcsname{ega4.c04.statement.7-8-3}
\expandafter\def\csname sga@target@IV.7.8.3.1\endcsname{ega4.c04.statement.7-8-3-1}
\expandafter\def\csname sga@target@IV.7.8.4\endcsname{ega4.c04.item.7-8-4}
\expandafter\def\csname sga@target@IV.7.8.4.1\endcsname{ega4.c04.statement.7-8-4-1}
\expandafter\def\csname sga@target@IV.7.8.5\endcsname{ega4.c04.statement.7-8-5}
\expandafter\def\csname sga@target@IV.7.8.6\endcsname{ega4.c04.prop.7-8-6}
\expandafter\def\csname sga@target@IV.7.9.1\endcsname{ega4.c04.statement.7-9-1}
\expandafter\def\csname sga@target@IV.7.9.10\endcsname{ega4.c04.item.7-9-10}
\expandafter\def\csname sga@target@IV.7.9.2\endcsname{ega4.c04.statement.7-9-2}
\expandafter\def\csname sga@target@IV.7.9.3\endcsname{ega4.c04.prop.7-9-3}
\expandafter\def\csname sga@target@IV.7.9.3.1\endcsname{ega4.c04.lem.7-9-3-1}
\expandafter\def\csname sga@target@IV.7.9.4\endcsname{ega4.c04.cor.7-9-4}
\expandafter\def\csname sga@target@IV.7.9.5\endcsname{ega4.c04.prop.7-9-5}
\expandafter\def\csname sga@target@IV.7.9.6\endcsname{ega4.c04.rem.7-9-6}
\expandafter\def\csname sga@target@IV.7.9.7\endcsname{ega4.c04.statement.7-9-7}
\expandafter\def\csname sga@target@IV.7.9.8\endcsname{ega4.c04.prop.7-9-8}
\expandafter\def\csname sga@target@IV.7.9.8.1\endcsname{ega4.c04.statement.7-9-8-1}
\expandafter\def\csname sga@target@IV.7.9.9\endcsname{ega4.c04.cor.7-9-9}
\expandafter\def\csname sga@target@IV.8.1.1\endcsname{ega4.c04.statement.8-1-1}
\expandafter\def\csname sga@target@IV.8.1.2\endcsname{ega4.c04.statement.8-1-2}
\expandafter\def\csname sga@target@IV.8.1.3\endcsname{ega4.c04.statement.8-1-3}
\expandafter\def\csname sga@target@IV.8.10.1\endcsname{ega4.c04.prop.8-10-1}
\expandafter\def\csname sga@target@IV.8.10.2\endcsname{ega4.c04.cor.8-10-2}
\expandafter\def\csname sga@target@IV.8.10.3\endcsname{ega4.c04.prop.8-10-3}
\expandafter\def\csname sga@target@IV.8.10.3.1\endcsname{ega4.c04.rem.8-10-3-1}
\expandafter\def\csname sga@target@IV.8.10.4\endcsname{ega4.c04.prop.8-10-4}
\expandafter\def\csname sga@target@IV.8.10.5\endcsname{ega4.c04.thm.8-10-5}
\expandafter\def\csname sga@target@IV.8.10.5.1\endcsname{ega4.c04.prop.8-10-5-1}
\expandafter\def\csname sga@target@IV.8.10.5.2\endcsname{ega4.c04.lem.8-10-5-2}
\expandafter\def\csname sga@target@IV.8.11\endcsname{ega4.c04.subsec.8-11.key-iv-8-11}
\expandafter\def\csname sga@target@IV.8.11.1\endcsname{ega4.c04.thm.8-11-1}
\expandafter\def\csname sga@target@IV.8.11.2\endcsname{ega4.c04.thm.8-11-2}
\expandafter\def\csname sga@target@IV.8.11.3\endcsname{ega4.c04.rem.8-11-3}
\expandafter\def\csname sga@target@IV.8.11.4\endcsname{ega4.c04.statement.8-11-4}
\expandafter\def\csname sga@target@IV.8.11.5\endcsname{ega4.c04.prop.8-11-5}
\expandafter\def\csname sga@target@IV.8.11.5.1\endcsname{ega4.c04.rem.8-11-5-1}
\expandafter\def\csname sga@target@IV.8.11.6\endcsname{ega4.c04.prop.8-11-6}
\expandafter\def\csname sga@target@IV.8.12\endcsname{ega4.c04.subsec.8-12.key-iv-8-12}
\expandafter\def\csname sga@target@IV.8.12.1\endcsname{ega4.c04.lem.8-12-1}
\expandafter\def\csname sga@target@IV.8.12.10\endcsname{ega4.c04.cor.8-12-10}
\expandafter\def\csname sga@target@IV.8.12.10.1\endcsname{ega4.c04.lem.8-12-10-1}
\expandafter\def\csname sga@target@IV.8.12.11\endcsname{ega4.c04.cor.8-12-11}
\expandafter\def\csname sga@target@IV.8.12.12\endcsname{ega4.c04.cor.8-12-12}
\expandafter\def\csname sga@target@IV.8.12.2\endcsname{ega4.c04.lem.8-12-2}
\expandafter\def\csname sga@target@IV.8.12.3\endcsname{ega4.c04.prop.8-12-3}
\expandafter\def\csname sga@target@IV.8.12.4\endcsname{ega4.c04.cor.8-12-4}
\expandafter\def\csname sga@target@IV.8.12.5\endcsname{ega4.c04.statement.8-12-5}
\expandafter\def\csname sga@target@IV.8.12.6\endcsname{ega4.c04.thm.8-12-6}
\expandafter\def\csname sga@target@IV.8.12.6.1\endcsname{ega4.c04.eq.8-12-6-1}
\expandafter\def\csname sga@target@IV.8.12.7\endcsname{ega4.c04.rem.8-12-7}
\expandafter\def\csname sga@target@IV.8.12.8\endcsname{ega4.c04.cor.8-12-8}
\expandafter\def\csname sga@target@IV.8.12.9\endcsname{ega4.c04.cor.8-12-9}
\expandafter\def\csname sga@target@IV.8.13.1\endcsname{ega4.c04.prop.8-13-1}
\expandafter\def\csname sga@target@IV.8.13.1.1\endcsname{ega4.c04.eq.8-13-1-1}
\expandafter\def\csname sga@target@IV.8.13.2\endcsname{ega4.c04.cor.8-13-2}
\expandafter\def\csname sga@target@IV.8.13.2.1\endcsname{ega4.c04.eq.8-13-2-1}
\expandafter\def\csname sga@target@IV.8.13.3\endcsname{ega4.c04.statement.8-13-3}
\expandafter\def\csname sga@target@IV.8.13.4\endcsname{ega4.c04.statement.8-13-4}
\expandafter\def\csname sga@target@IV.8.13.4.1\endcsname{ega4.c04.eq.8-13-4-1}
\expandafter\def\csname sga@target@IV.8.13.4.2\endcsname{ega4.c04.eq.8-13-4-2}
\expandafter\def\csname sga@target@IV.8.13.5\endcsname{ega4.c04.prop.8-13-5}
\expandafter\def\csname sga@target@IV.8.13.6\endcsname{ega4.c04.rem.8-13-6}
\expandafter\def\csname sga@target@IV.8.14.1\endcsname{ega4.c04.statement.8-14-1}
\expandafter\def\csname sga@target@IV.8.14.2\endcsname{ega4.c04.prop.8-14-2}
\expandafter\def\csname sga@target@IV.8.14.2.1\endcsname{ega4.c04.eq.8-14-2-1}
\expandafter\def\csname sga@target@IV.8.14.2.2\endcsname{ega4.c04.cor.8-14-2-2}
\expandafter\def\csname sga@target@IV.8.14.2.3\endcsname{ega4.c04.eq.8-14-2-3}
\expandafter\def\csname sga@target@IV.8.14.2.4\endcsname{ega4.c04.eq.8-14-2-4}
\expandafter\def\csname sga@target@IV.8.14.3\endcsname{ega4.c04.rem.8-14-3}
\expandafter\def\csname sga@target@IV.8.2.1\endcsname{ega4.c04.statement.8-2-1}
\expandafter\def\csname sga@target@IV.8.2.1.1\endcsname{ega4.c04.eq.8-2-1-1}
\expandafter\def\csname sga@target@IV.8.2.10\endcsname{ega4.c04.cor.8-2-10}
\expandafter\def\csname sga@target@IV.8.2.11\endcsname{ega4.c04.cor.8-2-11}
\expandafter\def\csname sga@target@IV.8.2.12\endcsname{ega4.c04.cor.8-2-12}
\expandafter\def\csname sga@target@IV.8.2.12.1\endcsname{ega4.c04.eq.8-2-12-1}
\expandafter\def\csname sga@target@IV.8.2.13\endcsname{ega4.c04.prop.8-2-13}
\expandafter\def\csname sga@target@IV.8.2.14\endcsname{ega4.c04.rem.8-2-14}
\expandafter\def\csname sga@target@IV.8.2.2\endcsname{ega4.c04.statement.8-2-2}
\expandafter\def\csname sga@target@IV.8.2.3\endcsname{ega4.c04.prop.8-2-3}
\expandafter\def\csname sga@target@IV.8.2.3.1\endcsname{ega4.c04.eq.8-2-3-1}
\expandafter\def\csname sga@target@IV.8.2.4\endcsname{ega4.c04.lem.8-2-4}
\expandafter\def\csname sga@target@IV.8.2.5\endcsname{ega4.c04.prop.8-2-5}
\expandafter\def\csname sga@target@IV.8.2.6\endcsname{ega4.c04.lem.8-2-6}
\expandafter\def\csname sga@target@IV.8.2.7\endcsname{ega4.c04.rem.8-2-7}
\expandafter\def\csname sga@target@IV.8.2.8\endcsname{ega4.c04.statement.8-2-8}
\expandafter\def\csname sga@target@IV.8.2.9\endcsname{ega4.c04.prop.8-2-9}
\expandafter\def\csname sga@target@IV.8.3.1\endcsname{ega4.c04.statement.8-3-1}
\expandafter\def\csname sga@target@IV.8.3.10\endcsname{ega4.c04.statement.8-3-10}
\expandafter\def\csname sga@target@IV.8.3.10.1\endcsname{ega4.c04.eq.8-3-10-1}
\expandafter\def\csname sga@target@IV.8.3.11\endcsname{ega4.c04.thm.8-3-11}
\expandafter\def\csname sga@target@IV.8.3.12\endcsname{ega4.c04.cor.8-3-12}
\expandafter\def\csname sga@target@IV.8.3.13\endcsname{ega4.c04.prop.8-3-13}
\expandafter\def\csname sga@target@IV.8.3.2\endcsname{ega4.c04.thm.8-3-2}
\expandafter\def\csname sga@target@IV.8.3.2.1\endcsname{ega4.c04.eq.8-3-2-1}
\expandafter\def\csname sga@target@IV.8.3.3\endcsname{ega4.c04.cor.8-3-3}
\expandafter\def\csname sga@target@IV.8.3.4\endcsname{ega4.c04.cor.8-3-4}
\expandafter\def\csname sga@target@IV.8.3.5\endcsname{ega4.c04.cor.8-3-5}
\expandafter\def\csname sga@target@IV.8.3.6\endcsname{ega4.c04.cor.8-3-6}
\expandafter\def\csname sga@target@IV.8.3.7\endcsname{ega4.c04.cor.8-3-7}
\expandafter\def\csname sga@target@IV.8.3.7.1\endcsname{ega4.c04.eq.8-3-7-1}
\expandafter\def\csname sga@target@IV.8.3.8\endcsname{ega4.c04.prop.8-3-8}
\expandafter\def\csname sga@target@IV.8.3.9\endcsname{ega4.c04.statement.8-3-9}
\expandafter\def\csname sga@target@IV.8.3.9.1\endcsname{ega4.c04.eq.8-3-9-1}
\expandafter\def\csname sga@target@IV.8.3.9.2\endcsname{ega4.c04.eq.8-3-9-2}
\expandafter\def\csname sga@target@IV.8.3.9.3\endcsname{ega4.c04.eq.8-3-9-3}
\expandafter\def\csname sga@target@IV.8.3.9.4\endcsname{ega4.c04.eq.8-3-9-4}
\expandafter\def\csname sga@target@IV.8.3.9.5\endcsname{ega4.c04.eq.8-3-9-5}
\expandafter\def\csname sga@target@IV.8.4.1\endcsname{ega4.c04.prop.8-4-1}
\expandafter\def\csname sga@target@IV.8.4.2\endcsname{ega4.c04.prop.8-4-2}
\expandafter\def\csname sga@target@IV.8.4.2.1\endcsname{ega4.c04.lem.8-4-2-1}
\expandafter\def\csname sga@target@IV.8.4.3\endcsname{ega4.c04.cor.8-4-3}
\expandafter\def\csname sga@target@IV.8.4.4\endcsname{ega4.c04.prop.8-4-4}
\expandafter\def\csname sga@target@IV.8.4.5\endcsname{ega4.c04.cor.8-4-5}
\expandafter\def\csname sga@target@IV.8.4.6\endcsname{ega4.c04.rem.8-4-6}
\expandafter\def\csname sga@target@IV.8.5.1\endcsname{ega4.c04.statement.8-5-1}
\expandafter\def\csname sga@target@IV.8.5.1.1\endcsname{ega4.c04.eq.8-5-1-1}
\expandafter\def\csname sga@target@IV.8.5.1.2\endcsname{ega4.c04.eq.8-5-1-2}
\expandafter\def\csname sga@target@IV.8.5.1.3\endcsname{ega4.c04.eq.8-5-1-3}
\expandafter\def\csname sga@target@IV.8.5.1.4\endcsname{ega4.c04.eq.8-5-1-4}
\expandafter\def\csname sga@target@IV.8.5.10\endcsname{ega4.c04.statement.8-5-10}
\expandafter\def\csname sga@target@IV.8.5.10.1\endcsname{ega4.c04.eq.8-5-10-1}
\expandafter\def\csname sga@target@IV.8.5.10.2\endcsname{ega4.c04.eq.8-5-10-2}
\expandafter\def\csname sga@target@IV.8.5.11\endcsname{ega4.c04.prop.8-5-11}
\expandafter\def\csname sga@target@IV.8.5.2\endcsname{ega4.c04.thm.8-5-2}
\expandafter\def\csname sga@target@IV.8.5.2.1\endcsname{ega4.c04.lem.8-5-2-1}
\expandafter\def\csname sga@target@IV.8.5.2.2\endcsname{ega4.c04.eq.8-5-2-2}
\expandafter\def\csname sga@target@IV.8.5.2.3\endcsname{ega4.c04.eq.8-5-2-3}
\expandafter\def\csname sga@target@IV.8.5.2.4\endcsname{ega4.c04.cor.8-5-2-4}
\expandafter\def\csname sga@target@IV.8.5.2.5\endcsname{ega4.c04.cor.8-5-2-5}
\expandafter\def\csname sga@target@IV.8.5.3\endcsname{ega4.c04.statement.8-5-3}
\expandafter\def\csname sga@target@IV.8.5.4\endcsname{ega4.c04.cor.8-5-4}
\expandafter\def\csname sga@target@IV.8.5.5\endcsname{ega4.c04.prop.8-5-5}
\expandafter\def\csname sga@target@IV.8.5.6\endcsname{ega4.c04.prop.8-5-6}
\expandafter\def\csname sga@target@IV.8.5.7\endcsname{ega4.c04.cor.8-5-7}
\expandafter\def\csname sga@target@IV.8.5.8\endcsname{ega4.c04.cor.8-5-8}
\expandafter\def\csname sga@target@IV.8.5.9\endcsname{ega4.c04.lem.8-5-9}
\expandafter\def\csname sga@target@IV.8.6.1\endcsname{ega4.c04.statement.8-6-1}
\expandafter\def\csname sga@target@IV.8.6.2\endcsname{ega4.c04.statement.8-6-2}
\expandafter\def\csname sga@target@IV.8.6.2.1\endcsname{ega4.c04.eq.8-6-2-1}
\expandafter\def\csname sga@target@IV.8.6.2.2\endcsname{ega4.c04.eq.8-6-2-2}
\expandafter\def\csname sga@target@IV.8.6.2.3\endcsname{ega4.c04.eq.8-6-2-3}
\expandafter\def\csname sga@target@IV.8.6.3\endcsname{ega4.c04.prop.8-6-3}
\expandafter\def\csname sga@target@IV.8.6.3.1\endcsname{ega4.c04.cor.8-6-3-1}
\expandafter\def\csname sga@target@IV.8.6.4\endcsname{ega4.c04.cor.8-6-4}
\expandafter\def\csname sga@target@IV.8.7.1\endcsname{ega4.c04.prop.8-7-1}
\expandafter\def\csname sga@target@IV.8.7.2\endcsname{ega4.c04.prop.8-7-2}
\expandafter\def\csname sga@target@IV.8.7.3\endcsname{ega4.c04.cor.8-7-3}
\expandafter\def\csname sga@target@IV.8.7.4\endcsname{ega4.c04.rem.8-7-4}
\expandafter\def\csname sga@target@IV.8.8.1\endcsname{ega4.c04.statement.8-8-1}
\expandafter\def\csname sga@target@IV.8.8.1.1\endcsname{ega4.c04.eq.8-8-1-1}
\expandafter\def\csname sga@target@IV.8.8.2\endcsname{ega4.c04.thm.8-8-2}
\expandafter\def\csname sga@target@IV.8.8.2.1\endcsname{ega4.c04.lem.8-8-2-1}
\expandafter\def\csname sga@target@IV.8.8.2.2\endcsname{ega4.c04.lem.8-8-2-2}
\expandafter\def\csname sga@target@IV.8.8.2.3\endcsname{ega4.c04.lem.8-8-2-3}
\expandafter\def\csname sga@target@IV.8.8.2.4\endcsname{ega4.c04.cor.8-8-2-4}
\expandafter\def\csname sga@target@IV.8.8.2.5\endcsname{ega4.c04.cor.8-8-2-5}
\expandafter\def\csname sga@target@IV.8.8.3\endcsname{ega4.c04.statement.8-8-3}
\expandafter\def\csname sga@target@IV.8.9.1\endcsname{ega4.c04.prop.8-9-1}
\expandafter\def\csname sga@target@IV.8.9.2\endcsname{ega4.c04.statement.8-9-2}
\expandafter\def\csname sga@target@IV.8.9.3\endcsname{ega4.c04.prop.8-9-3}
\expandafter\def\csname sga@target@IV.8.9.4\endcsname{ega4.c04.prop.8-9-4}
\expandafter\def\csname sga@target@IV.8.9.4.1\endcsname{ega4.c04.lem.8-9-4-1}
\expandafter\def\csname sga@target@IV.8.9.5\endcsname{ega4.c04.cor.8-9-5}
\expandafter\def\csname sga@target@IV.9.1.1\endcsname{ega4.c04.prop.9-1-1}
\expandafter\def\csname sga@target@IV.9.1.2\endcsname{ega4.c04.cor.9-1-2}
\expandafter\def\csname sga@target@IV.9.1.3\endcsname{ega4.c04.rem.9-1-3}
\expandafter\def\csname sga@target@IV.9.1.4\endcsname{ega4.c04.prop.9-1-4}
\expandafter\def\csname sga@target@IV.9.1.5\endcsname{ega4.c04.cor.9-1-5}
\expandafter\def\csname sga@target@IV.9.1.6\endcsname{ega4.c04.prop.9-1-6}
\expandafter\def\csname sga@target@IV.9.2\endcsname{ega4.c04.subsec.9-2.key-iv-9-2}
\expandafter\def\csname sga@target@IV.9.2.1\endcsname{ega4.c04.def.9-2-1}
\expandafter\def\csname sga@target@IV.9.2.2\endcsname{ega4.c04.item.9-2-2}
\expandafter\def\csname sga@target@IV.9.2.3\endcsname{ega4.c04.prop.9-2-3}
\expandafter\def\csname sga@target@IV.9.2.4\endcsname{ega4.c04.cor.9-2-4}
\expandafter\def\csname sga@target@IV.9.2.5\endcsname{ega4.c04.rem.9-2-5}
\expandafter\def\csname sga@target@IV.9.2.6\endcsname{ega4.c04.prop.9-2-6}
\expandafter\def\csname sga@target@IV.9.2.6.1\endcsname{ega4.c04.prop.9-2-6-1}
\expandafter\def\csname sga@target@IV.9.2.6.2\endcsname{ega4.c04.cor.9-2-6-2}
\expandafter\def\csname sga@target@IV.9.3.1\endcsname{ega4.c04.prop.9-3-1}
\expandafter\def\csname sga@target@IV.9.3.2\endcsname{ega4.c04.prop.9-3-2}
\expandafter\def\csname sga@target@IV.9.3.3\endcsname{ega4.c04.prop.9-3-3}
\expandafter\def\csname sga@target@IV.9.3.4\endcsname{ega4.c04.rem.9-3-4}
\expandafter\def\csname sga@target@IV.9.3.5\endcsname{ega4.c04.prop.9-3-5}
\expandafter\def\csname sga@target@IV.9.3.6\endcsname{ega4.c04.statement.9-3-6}
\expandafter\def\csname sga@target@IV.9.4.1\endcsname{ega4.c04.statement.9-4-1}
\expandafter\def\csname sga@target@IV.9.4.2\endcsname{ega4.c04.prop.9-4-2}
\expandafter\def\csname sga@target@IV.9.4.3\endcsname{ega4.c04.cor.9-4-3}
\expandafter\def\csname sga@target@IV.9.4.3.1\endcsname{ega4.c04.eq.9-4-3-1}
\expandafter\def\csname sga@target@IV.9.4.4\endcsname{ega4.c04.prop.9-4-4}
\expandafter\def\csname sga@target@IV.9.4.5\endcsname{ega4.c04.cor.9-4-5}
\expandafter\def\csname sga@target@IV.9.4.6\endcsname{ega4.c04.cor.9-4-6}
\expandafter\def\csname sga@target@IV.9.4.7\endcsname{ega4.c04.prop.9-4-7}
\expandafter\def\csname sga@target@IV.9.4.7.1\endcsname{ega4.c04.lem.9-4-7-1}
\expandafter\def\csname sga@target@IV.9.4.7.2\endcsname{ega4.c04.rem.9-4-7-2}
\expandafter\def\csname sga@target@IV.9.4.8\endcsname{ega4.c04.prop.9-4-8}
\expandafter\def\csname sga@target@IV.9.4.9\endcsname{ega4.c04.rem.9-4-9}
\expandafter\def\csname sga@target@IV.9.5.1\endcsname{ega4.c04.prop.9-5-1}
\expandafter\def\csname sga@target@IV.9.5.2\endcsname{ega4.c04.cor.9-5-2}
\expandafter\def\csname sga@target@IV.9.5.3\endcsname{ega4.c04.prop.9-5-3}
\expandafter\def\csname sga@target@IV.9.5.3.case2\endcsname{ega4.c04.item.9-5-3-case2}
\expandafter\def\csname sga@target@IV.9.5.4\endcsname{ega4.c04.cor.9-5-4}
\expandafter\def\csname sga@target@IV.9.5.5\endcsname{ega4.c04.prop.9-5-5}
\expandafter\def\csname sga@target@IV.9.5.6\endcsname{ega4.c04.cor.9-5-6}
\expandafter\def\csname sga@target@IV.9.5.7\endcsname{ega4.c04.rem.9-5-7}
\expandafter\def\csname sga@target@IV.9.6.1\endcsname{ega4.c04.prop.9-6-1}
\expandafter\def\csname sga@target@IV.9.6.2\endcsname{ega4.c04.prop.9-6-2}
\expandafter\def\csname sga@target@IV.9.6.3\endcsname{ega4.c04.prop.9-6-3}
\expandafter\def\csname sga@target@IV.9.6.3.1\endcsname{ega4.c04.eq.9-6-3-1}
\expandafter\def\csname sga@target@IV.9.6.4\endcsname{ega4.c04.cor.9-6-4}
\expandafter\def\csname sga@target@IV.9.6.5\endcsname{ega4.c04.cor.9-6-5}
\expandafter\def\csname sga@target@IV.9.6.6\endcsname{ega4.c04.sec.9-6-6}
\expandafter\def\csname sga@target@IV.9.6.7\endcsname{ega4.c04.sec.9-6-7}
\expandafter\def\csname sga@target@IV.9.7\endcsname{ega4.c04.subsec.9-7.key-iv-9-7}
\expandafter\def\csname sga@target@IV.9.7.1\endcsname{ega4.c04.lem.9-7-1}
\expandafter\def\csname sga@target@IV.9.7.10\endcsname{ega4.c04.cor.9-7-10}
\expandafter\def\csname sga@target@IV.9.7.11\endcsname{ega4.c04.rem.9-7-11}
\expandafter\def\csname sga@target@IV.9.7.12\endcsname{ega4.c04.prop.9-7-12}
\expandafter\def\csname sga@target@IV.9.7.2\endcsname{ega4.c04.lem.9-7-2}
\expandafter\def\csname sga@target@IV.9.7.3\endcsname{ega4.c04.subsec.9-7-3}
\expandafter\def\csname sga@target@IV.9.7.4\endcsname{ega4.c04.lem.9-7-4}
\expandafter\def\csname sga@target@IV.9.7.5\endcsname{ega4.c04.lem.9-7-5}
\expandafter\def\csname sga@target@IV.9.7.5.1\endcsname{ega4.c04.lem.9-7-5-1}
\expandafter\def\csname sga@target@IV.9.7.6\endcsname{ega4.c04.prop.9-7-6}
\expandafter\def\csname sga@target@IV.9.7.7\endcsname{ega4.c04.thm.9-7-7}
\expandafter\def\csname sga@target@IV.9.7.8\endcsname{ega4.c04.prop.9-7-8}
\expandafter\def\csname sga@target@IV.9.7.9\endcsname{ega4.c04.cor.9-7-9}
\expandafter\def\csname sga@target@IV.9.8\endcsname{ega4.c04.subsec.9-8.key-iv-9-8}
\expandafter\def\csname sga@target@IV.9.8.1\endcsname{ega4.c04.prop.9-8-1}
\expandafter\def\csname sga@target@IV.9.8.1-proof\endcsname{ega4.c04.subsec.9-8-1-proof}
\expandafter\def\csname sga@target@IV.9.8.2\endcsname{ega4.c04.subsec.9-8-2}
\expandafter\def\csname sga@target@IV.9.8.3\endcsname{ega4.c04.prop.9-8-3}
\expandafter\def\csname sga@target@IV.9.8.3-proof\endcsname{ega4.c04.subsec.9-8-3-proof}
\expandafter\def\csname sga@target@IV.9.8.4\endcsname{ega4.c04.cor.9-8-4}
\expandafter\def\csname sga@target@IV.9.8.4-proof\endcsname{ega4.c04.item.9-8-4-proof}
\expandafter\def\csname sga@target@IV.9.8.5\endcsname{ega4.c04.prop.9-8-5}
\expandafter\def\csname sga@target@IV.9.8.5-proof\endcsname{ega4.c04.item.9-8-5-proof}
\expandafter\def\csname sga@target@IV.9.8.6\endcsname{ega4.c04.prop.9-8-6}
\expandafter\def\csname sga@target@IV.9.8.6-proof\endcsname{ega4.c04.item.9-8-6-proof}
\expandafter\def\csname sga@target@IV.9.8.7\endcsname{ega4.c04.cor.9-8-7}
\expandafter\def\csname sga@target@IV.9.8.7-proof\endcsname{ega4.c04.item.9-8-7-proof}
\expandafter\def\csname sga@target@IV.9.8.8\endcsname{ega4.c04.prop.9-8-8}
\expandafter\def\csname sga@target@IV.9.8.8-proof\endcsname{ega4.c04.item.9-8-8-proof}
\expandafter\def\csname sga@target@IV.9.8.9\endcsname{ega4.c04.rem.9-8-9}
\expandafter\def\csname sga@target@IV.9.8.9.1\endcsname{ega4.c04.item.9-8-9-1}
\expandafter\def\csname sga@target@IV.9.9\endcsname{ega4.c04.subsec.9-9.key-iv-9-9}
\expandafter\def\csname sga@target@IV.9.9.1\endcsname{ega4.c04.prop.9-9-1}
\expandafter\def\csname sga@target@IV.9.9.1-proof\endcsname{ega4.c04.item.9-9-1-proof}
\expandafter\def\csname sga@target@IV.9.9.1.1\endcsname{ega4.c04.eq.9-9-1-1}
\expandafter\def\csname sga@target@IV.9.9.2\endcsname{ega4.c04.prop.9-9-2}
\expandafter\def\csname sga@target@IV.9.9.2-proof\endcsname{ega4.c04.item.9-9-2-proof}
\expandafter\def\csname sga@target@IV.9.9.2.1\endcsname{ega4.c04.eq.9-9-2-1}
\expandafter\def\csname sga@target@IV.9.9.2.2\endcsname{ega4.c04.eq.9-9-2-2}
\expandafter\def\csname sga@target@IV.9.9.3\endcsname{ega4.c04.cor.9-9-3}
\expandafter\def\csname sga@target@IV.9.9.3-proof\endcsname{ega4.c04.statement.9-9-3-proof}
\expandafter\def\csname sga@target@IV.9.9.4\endcsname{ega4.c04.prop.9-9-4}
\expandafter\def\csname sga@target@IV.9.9.4-proof\endcsname{ega4.c04.item.9-9-4-proof}
\expandafter\def\csname sga@target@IV.9.9.5\endcsname{ega4.c04.cor.9-9-5}
\expandafter\def\csname sga@target@IV.9.9.5-proof\endcsname{ega4.c04.item.9-9-5-proof}
\expandafter\def\csname sga@target@IV.9.9.6\endcsname{ega4.c04.prop.9-9-6}
\expandafter\def\csname sga@target@IV.9.9.6-proof\endcsname{ega4.c04.item.9-9-6-proof}
\expandafter\def\csname sga@target@section:IV.1\endcsname{ega4.c04.sec.1}
\expandafter\def\csname sga@target@section:IV.10\endcsname{ega4.c04.sec.10}
\expandafter\def\csname sga@target@section:IV.11\endcsname{ega4.c04.sec.11}
\expandafter\def\csname sga@target@section:IV.12\endcsname{ega4.c04.sec.12}
\expandafter\def\csname sga@target@section:IV.13\endcsname{ega4.c04.sec.13}
\expandafter\def\csname sga@target@section:IV.14\endcsname{ega4.c04.sec.14}
\expandafter\def\csname sga@target@section:IV.15\endcsname{ega4.c04.sec.15}
\expandafter\def\csname sga@target@section:IV.16\endcsname{ega4.c04.sec.16}
\expandafter\def\csname sga@target@section:IV.17\endcsname{ega4.c04.sec.17}
\expandafter\def\csname sga@target@section:IV.18\endcsname{ega4.c04.sec.18}
\expandafter\def\csname sga@target@section:IV.19\endcsname{ega4.c04.sec.19}
\expandafter\def\csname sga@target@section:IV.2\endcsname{ega4.c04.sec.2}
\expandafter\def\csname sga@target@section:IV.21\endcsname{ega4.c04.sec.21}
\expandafter\def\csname sga@target@section:IV.3\endcsname{ega4.c04.sec.3}
\expandafter\def\csname sga@target@section:IV.4\endcsname{ega4.c04.sec.4}
\expandafter\def\csname sga@target@section:IV.5\endcsname{ega4.c04.sec.5}
\expandafter\def\csname sga@target@section:IV.6\endcsname{ega4.c04.sec.6}
\expandafter\def\csname sga@target@section:IV.7\endcsname{ega4.c04.sec.7}
\expandafter\def\csname sga@target@section:IV.8\endcsname{ega4.c04.sec.8}
\expandafter\def\csname sga@target@section:IV.9\endcsname{ega4.c04.sec.9}
\expandafter\def\csname sga@target@subsection:IV.1.1\endcsname{ega4.c04.subsec.1-1}
\expandafter\def\csname sga@target@subsection:IV.1.10\endcsname{ega4.c04.subsec.1-10}
\expandafter\def\csname sga@target@subsection:IV.1.2\endcsname{ega4.c04.subsec.1-2}
\expandafter\def\csname sga@target@subsection:IV.1.3\endcsname{ega4.c04.subsec.1-3}
\expandafter\def\csname sga@target@subsection:IV.1.4\endcsname{ega4.c04.subsec.1-4}
\expandafter\def\csname sga@target@subsection:IV.1.5\endcsname{ega4.c04.subsec.1-5}
\expandafter\def\csname sga@target@subsection:IV.1.6\endcsname{ega4.c04.subsec.1-6}
\expandafter\def\csname sga@target@subsection:IV.1.7\endcsname{ega4.c04.subsec.1-7}
\expandafter\def\csname sga@target@subsection:IV.1.8\endcsname{ega4.c04.subsec.1-8}
\expandafter\def\csname sga@target@subsection:IV.1.9\endcsname{ega4.c04.subsec.1-9}
\expandafter\def\csname sga@target@subsection:IV.10.1\endcsname{ega4.c04.subsec.10-1}
\expandafter\def\csname sga@target@subsection:IV.10.10\endcsname{ega4.c04.subsec.10-10}
\expandafter\def\csname sga@target@subsection:IV.10.2\endcsname{ega4.c04.subsec.10-2}
\expandafter\def\csname sga@target@subsection:IV.10.3\endcsname{ega4.c04.subsec.10-3}
\expandafter\def\csname sga@target@subsection:IV.10.4\endcsname{ega4.c04.subsec.10-4}
\expandafter\def\csname sga@target@subsection:IV.10.5\endcsname{ega4.c04.subsec.10-5}
\expandafter\def\csname sga@target@subsection:IV.10.6\endcsname{ega4.c04.subsec.10-6}
\expandafter\def\csname sga@target@subsection:IV.10.7\endcsname{ega4.c04.subsec.10-7}
\expandafter\def\csname sga@target@subsection:IV.10.8\endcsname{ega4.c04.subsec.10-8}
\expandafter\def\csname sga@target@subsection:IV.10.9\endcsname{ega4.c04.subsec.10-9}
\expandafter\def\csname sga@target@subsection:IV.11.1\endcsname{ega4.c04.subsec.11-1}
\expandafter\def\csname sga@target@subsection:IV.11.10\endcsname{ega4.c04.subsec.11-10}
\expandafter\def\csname sga@target@subsection:IV.11.2\endcsname{ega4.c04.subsec.11-2}
\expandafter\def\csname sga@target@subsection:IV.11.3\endcsname{ega4.c04.subsec.11-3}
\expandafter\def\csname sga@target@subsection:IV.11.4\endcsname{ega4.c04.subsec.11-4}
\expandafter\def\csname sga@target@subsection:IV.11.5\endcsname{ega4.c04.subsec.11-5}
\expandafter\def\csname sga@target@subsection:IV.11.6\endcsname{ega4.c04.subsec.11-6}
\expandafter\def\csname sga@target@subsection:IV.11.7\endcsname{ega4.c04.subsec.11-7}
\expandafter\def\csname sga@target@subsection:IV.11.8\endcsname{ega4.c04.subsec.11-8}
\expandafter\def\csname sga@target@subsection:IV.11.9\endcsname{ega4.c04.subsec.11-9}
\expandafter\def\csname sga@target@subsection:IV.12.0\endcsname{ega4.c04.subsec.12-0}
\expandafter\def\csname sga@target@subsection:IV.12.1\endcsname{ega4.c04.subsec.12-1}
\expandafter\def\csname sga@target@subsection:IV.12.2\endcsname{ega4.c04.subsec.12-2}
\expandafter\def\csname sga@target@subsection:IV.12.3\endcsname{ega4.c04.subsec.12-3}
\expandafter\def\csname sga@target@subsection:IV.13.1\endcsname{ega4.c04.subsec.13-1}
\expandafter\def\csname sga@target@subsection:IV.13.2\endcsname{ega4.c04.subsec.13-2}
\expandafter\def\csname sga@target@subsection:IV.13.3\endcsname{ega4.c04.subsec.13-3}
\expandafter\def\csname sga@target@subsection:IV.14.1\endcsname{ega4.c04.subsec.14-1}
\expandafter\def\csname sga@target@subsection:IV.14.2\endcsname{ega4.c04.subsec.14-2}
\expandafter\def\csname sga@target@subsection:IV.14.3\endcsname{ega4.c04.subsec.14-3}
\expandafter\def\csname sga@target@subsection:IV.14.4\endcsname{ega4.c04.subsec.14-4}
\expandafter\def\csname sga@target@subsection:IV.14.5\endcsname{ega4.c04.subsec.14-5}
\expandafter\def\csname sga@target@subsection:IV.15.1\endcsname{ega4.c04.subsec.15-1}
\expandafter\def\csname sga@target@subsection:IV.15.2\endcsname{ega4.c04.subsec.15-2}
\expandafter\def\csname sga@target@subsection:IV.15.3\endcsname{ega4.c04.subsec.15-3}
\expandafter\def\csname sga@target@subsection:IV.15.4\endcsname{ega4.c04.subsec.15-4}
\expandafter\def\csname sga@target@subsection:IV.15.5\endcsname{ega4.c04.subsec.15-5}
\expandafter\def\csname sga@target@subsection:IV.15.6\endcsname{ega4.c04.subsec.15-6}
\expandafter\def\csname sga@target@subsection:IV.15.7\endcsname{ega4.c04.subsec.15-7}
\expandafter\def\csname sga@target@subsection:IV.17.1\endcsname{ega4.c04.subsec.17-1}
\expandafter\def\csname sga@target@subsection:IV.17.10\endcsname{ega4.c04.subsec.17-10}
\expandafter\def\csname sga@target@subsection:IV.17.11\endcsname{ega4.c04.subsec.17-11}
\expandafter\def\csname sga@target@subsection:IV.17.12\endcsname{ega4.c04.subsec.17-12}
\expandafter\def\csname sga@target@subsection:IV.17.13\endcsname{ega4.c04.subsec.17-13}
\expandafter\def\csname sga@target@subsection:IV.17.14\endcsname{ega4.c04.subsec.17-14}
\expandafter\def\csname sga@target@subsection:IV.17.15\endcsname{ega4.c04.subsec.17-15}
\expandafter\def\csname sga@target@subsection:IV.17.16\endcsname{ega4.c04.subsec.17-16}
\expandafter\def\csname sga@target@subsection:IV.17.2\endcsname{ega4.c04.subsec.17-2}
\expandafter\def\csname sga@target@subsection:IV.17.3\endcsname{ega4.c04.subsec.17-3}
\expandafter\def\csname sga@target@subsection:IV.17.4\endcsname{ega4.c04.subsec.17-4}
\expandafter\def\csname sga@target@subsection:IV.17.5\endcsname{ega4.c04.subsec.17-5}
\expandafter\def\csname sga@target@subsection:IV.17.6\endcsname{ega4.c04.subsec.17-6}
\expandafter\def\csname sga@target@subsection:IV.17.7\endcsname{ega4.c04.subsec.17-7}
\expandafter\def\csname sga@target@subsection:IV.17.8\endcsname{ega4.c04.subsec.17-8}
\expandafter\def\csname sga@target@subsection:IV.17.9\endcsname{ega4.c04.subsec.17-9}
\expandafter\def\csname sga@target@subsection:IV.18.1\endcsname{ega4.c04.subsec.18-1}
\expandafter\def\csname sga@target@subsection:IV.18.10\endcsname{ega4.c04.subsec.18-10}
\expandafter\def\csname sga@target@subsection:IV.18.11\endcsname{ega4.c04.subsec.18-11}
\expandafter\def\csname sga@target@subsection:IV.18.12\endcsname{ega4.c04.subsec.18-12}
\expandafter\def\csname sga@target@subsection:IV.18.2\endcsname{ega4.c04.subsec.18-2}
\expandafter\def\csname sga@target@subsection:IV.18.3\endcsname{ega4.c04.subsec.18-3}
\expandafter\def\csname sga@target@subsection:IV.18.4\endcsname{ega4.c04.subsec.18-4}
\expandafter\def\csname sga@target@subsection:IV.18.5\endcsname{ega4.c04.subsec.18-5}
\expandafter\def\csname sga@target@subsection:IV.18.6\endcsname{ega4.c04.subsec.18-6}
\expandafter\def\csname sga@target@subsection:IV.18.7\endcsname{ega4.c04.subsec.18-7}
\expandafter\def\csname sga@target@subsection:IV.18.8\endcsname{ega4.c04.subsec.18-8}
\expandafter\def\csname sga@target@subsection:IV.18.9\endcsname{ega4.c04.subsec.18-9}
\expandafter\def\csname sga@target@subsection:IV.19.1\endcsname{ega4.c04.subsec.19-1}
\expandafter\def\csname sga@target@subsection:IV.19.2\endcsname{ega4.c04.subsec.19-2}
\expandafter\def\csname sga@target@subsection:IV.19.3\endcsname{ega4.c04.subsec.19-3}
\expandafter\def\csname sga@target@subsection:IV.19.4\endcsname{ega4.c04.subsec.19-4}
\expandafter\def\csname sga@target@subsection:IV.19.5\endcsname{ega4.c04.subsec.19-5}
\expandafter\def\csname sga@target@subsection:IV.19.6\endcsname{ega4.c04.subsec.19-6}
\expandafter\def\csname sga@target@subsection:IV.19.7\endcsname{ega4.c04.subsec.19-7}
\expandafter\def\csname sga@target@subsection:IV.19.8\endcsname{ega4.c04.subsec.19-8}
\expandafter\def\csname sga@target@subsection:IV.19.9\endcsname{ega4.c04.subsec.19-9}
\expandafter\def\csname sga@target@subsection:IV.2.1\endcsname{ega4.c04.subsec.2-1}
\expandafter\def\csname sga@target@subsection:IV.2.2\endcsname{ega4.c04.subsec.2-2}
\expandafter\def\csname sga@target@subsection:IV.2.3\endcsname{ega4.c04.subsec.2-3}
\expandafter\def\csname sga@target@subsection:IV.2.4\endcsname{ega4.c04.subsec.2-4}
\expandafter\def\csname sga@target@subsection:IV.2.5\endcsname{ega4.c04.subsec.2-5}
\expandafter\def\csname sga@target@subsection:IV.2.6\endcsname{ega4.c04.subsec.2-6}
\expandafter\def\csname sga@target@subsection:IV.2.7\endcsname{ega4.c04.subsec.2-7}
\expandafter\def\csname sga@target@subsection:IV.2.8\endcsname{ega4.c04.subsec.2-8}
\expandafter\def\csname sga@target@subsection:IV.21.1\endcsname{ega4.c04.subsec.21-1}
\expandafter\def\csname sga@target@subsection:IV.21.10\endcsname{ega4.c04.subsec.21-10}
\expandafter\def\csname sga@target@subsection:IV.21.11\endcsname{ega4.c04.subsec.21-11}
\expandafter\def\csname sga@target@subsection:IV.21.12\endcsname{ega4.c04.subsec.21-12}
\expandafter\def\csname sga@target@subsection:IV.21.13\endcsname{ega4.c04.subsec.21-13}
\expandafter\def\csname sga@target@subsection:IV.21.14\endcsname{ega4.c04.subsec.21-14}
\expandafter\def\csname sga@target@subsection:IV.21.15\endcsname{ega4.c04.subsec.21-15}
\expandafter\def\csname sga@target@subsection:IV.21.2\endcsname{ega4.c04.subsec.21-2}
\expandafter\def\csname sga@target@subsection:IV.21.3\endcsname{ega4.c04.subsec.21-3}
\expandafter\def\csname sga@target@subsection:IV.21.4\endcsname{ega4.c04.subsec.21-4}
\expandafter\def\csname sga@target@subsection:IV.21.5\endcsname{ega4.c04.subsec.21-5}
\expandafter\def\csname sga@target@subsection:IV.21.6\endcsname{ega4.c04.subsec.21-6}
\expandafter\def\csname sga@target@subsection:IV.21.7\endcsname{ega4.c04.subsec.21-7}
\expandafter\def\csname sga@target@subsection:IV.21.8\endcsname{ega4.c04.subsec.21-8}
\expandafter\def\csname sga@target@subsection:IV.21.9\endcsname{ega4.c04.subsec.21-9}
\expandafter\def\csname sga@target@subsection:IV.3.1\endcsname{ega4.c04.subsec.3-1}
\expandafter\def\csname sga@target@subsection:IV.3.2\endcsname{ega4.c04.subsec.3-2}
\expandafter\def\csname sga@target@subsection:IV.3.3\endcsname{ega4.c04.subsec.3-3}
\expandafter\def\csname sga@target@subsection:IV.3.4\endcsname{ega4.c04.subsec.3-4}
\expandafter\def\csname sga@target@subsection:IV.4.1\endcsname{ega4.c04.subsec.4-1}
\expandafter\def\csname sga@target@subsection:IV.4.2\endcsname{ega4.c04.subsec.4-2}
\expandafter\def\csname sga@target@subsection:IV.4.3\endcsname{ega4.c04.subsec.4-3}
\expandafter\def\csname sga@target@subsection:IV.4.4\endcsname{ega4.c04.subsec.4-4}
\expandafter\def\csname sga@target@subsection:IV.4.5\endcsname{ega4.c04.subsec.4-5}
\expandafter\def\csname sga@target@subsection:IV.4.6\endcsname{ega4.c04.subsec.4-6}
\expandafter\def\csname sga@target@subsection:IV.4.7\endcsname{ega4.c04.subsec.4-7}
\expandafter\def\csname sga@target@subsection:IV.4.8\endcsname{ega4.c04.subsec.4-8}
\expandafter\def\csname sga@target@subsection:IV.4.9\endcsname{ega4.c04.subsec.4-9}
\expandafter\def\csname sga@target@subsection:IV.5.1\endcsname{ega4.c04.subsec.5-1}
\expandafter\def\csname sga@target@subsection:IV.5.10\endcsname{ega4.c04.subsec.5-10}
\expandafter\def\csname sga@target@subsection:IV.5.11\endcsname{ega4.c04.subsec.5-11}
\expandafter\def\csname sga@target@subsection:IV.5.12\endcsname{ega4.c04.subsec.5-12}
\expandafter\def\csname sga@target@subsection:IV.5.13\endcsname{ega4.c04.subsec.5-13}
\expandafter\def\csname sga@target@subsection:IV.5.2\endcsname{ega4.c04.subsec.5-2}
\expandafter\def\csname sga@target@subsection:IV.5.3\endcsname{ega4.c04.subsec.5-3}
\expandafter\def\csname sga@target@subsection:IV.5.4\endcsname{ega4.c04.subsec.5-4}
\expandafter\def\csname sga@target@subsection:IV.5.5\endcsname{ega4.c04.subsec.5-5}
\expandafter\def\csname sga@target@subsection:IV.5.6\endcsname{ega4.c04.subsec.5-6}
\expandafter\def\csname sga@target@subsection:IV.5.7\endcsname{ega4.c04.subsec.5-7}
\expandafter\def\csname sga@target@subsection:IV.5.8\endcsname{ega4.c04.subsec.5-8}
\expandafter\def\csname sga@target@subsection:IV.5.9\endcsname{ega4.c04.subsec.5-9}
\expandafter\def\csname sga@target@subsection:IV.6.1\endcsname{ega4.c04.subsec.6-1}
\expandafter\def\csname sga@target@subsection:IV.6.10\endcsname{ega4.c04.subsec.6-10}
\expandafter\def\csname sga@target@subsection:IV.6.11\endcsname{ega4.c04.subsec.6-11}
\expandafter\def\csname sga@target@subsection:IV.6.12\endcsname{ega4.c04.subsec.6-12}
\expandafter\def\csname sga@target@subsection:IV.6.13\endcsname{ega4.c04.subsec.6-13}
\expandafter\def\csname sga@target@subsection:IV.6.14\endcsname{ega4.c04.subsec.6-14}
\expandafter\def\csname sga@target@subsection:IV.6.15\endcsname{ega4.c04.subsec.6-15}
\expandafter\def\csname sga@target@subsection:IV.6.2\endcsname{ega4.c04.subsec.6-2}
\expandafter\def\csname sga@target@subsection:IV.6.3\endcsname{ega4.c04.subsec.6-3}
\expandafter\def\csname sga@target@subsection:IV.6.4\endcsname{ega4.c04.subsec.6-4}
\expandafter\def\csname sga@target@subsection:IV.6.5\endcsname{ega4.c04.subsec.6-5}
\expandafter\def\csname sga@target@subsection:IV.6.6\endcsname{ega4.c04.subsec.6-6}
\expandafter\def\csname sga@target@subsection:IV.6.7\endcsname{ega4.c04.subsec.6-7}
\expandafter\def\csname sga@target@subsection:IV.6.8\endcsname{ega4.c04.subsec.6-8}
\expandafter\def\csname sga@target@subsection:IV.6.9\endcsname{ega4.c04.subsec.6-9}
\expandafter\def\csname sga@target@subsection:IV.7.1\endcsname{ega4.c04.subsec.7-1}
\expandafter\def\csname sga@target@subsection:IV.7.2\endcsname{ega4.c04.subsec.7-2}
\expandafter\def\csname sga@target@subsection:IV.7.3\endcsname{ega4.c04.subsec.7-3}
\expandafter\def\csname sga@target@subsection:IV.7.4\endcsname{ega4.c04.subsec.7-4}
\expandafter\def\csname sga@target@subsection:IV.7.5\endcsname{ega4.c04.subsec.7-5}
\expandafter\def\csname sga@target@subsection:IV.7.6\endcsname{ega4.c04.subsec.7-6}
\expandafter\def\csname sga@target@subsection:IV.7.7\endcsname{ega4.c04.subsec.7-7}
\expandafter\def\csname sga@target@subsection:IV.7.8\endcsname{ega4.c04.subsec.7-8}
\expandafter\def\csname sga@target@subsection:IV.7.9\endcsname{ega4.c04.subsec.7-9}
\expandafter\def\csname sga@target@subsection:IV.8.1\endcsname{ega4.c04.subsec.8-1}
\expandafter\def\csname sga@target@subsection:IV.8.10\endcsname{ega4.c04.subsec.8-10}
\expandafter\def\csname sga@target@subsection:IV.8.11\endcsname{ega4.c04.subsec.8-11}
\expandafter\def\csname sga@target@subsection:IV.8.12\endcsname{ega4.c04.subsec.8-12}
\expandafter\def\csname sga@target@subsection:IV.8.13\endcsname{ega4.c04.subsec.8-13}
\expandafter\def\csname sga@target@subsection:IV.8.14\endcsname{ega4.c04.subsec.8-14}
\expandafter\def\csname sga@target@subsection:IV.8.2\endcsname{ega4.c04.subsec.8-2}
\expandafter\def\csname sga@target@subsection:IV.8.3\endcsname{ega4.c04.subsec.8-3}
\expandafter\def\csname sga@target@subsection:IV.8.4\endcsname{ega4.c04.subsec.8-4}
\expandafter\def\csname sga@target@subsection:IV.8.5\endcsname{ega4.c04.subsec.8-5}
\expandafter\def\csname sga@target@subsection:IV.8.6\endcsname{ega4.c04.subsec.8-6}
\expandafter\def\csname sga@target@subsection:IV.8.7\endcsname{ega4.c04.subsec.8-7}
\expandafter\def\csname sga@target@subsection:IV.8.8\endcsname{ega4.c04.subsec.8-8}
\expandafter\def\csname sga@target@subsection:IV.8.9\endcsname{ega4.c04.subsec.8-9}
\expandafter\def\csname sga@target@subsection:IV.9.1\endcsname{ega4.c04.subsec.9-1}
\expandafter\def\csname sga@target@subsection:IV.9.2\endcsname{ega4.c04.subsec.9-2}
\expandafter\def\csname sga@target@subsection:IV.9.3\endcsname{ega4.c04.subsec.9-3}
\expandafter\def\csname sga@target@subsection:IV.9.4\endcsname{ega4.c04.subsec.9-4}
\expandafter\def\csname sga@target@subsection:IV.9.5\endcsname{ega4.c04.subsec.9-5}
\expandafter\def\csname sga@target@subsection:IV.9.6\endcsname{ega4.c04.subsec.9-6}
\expandafter\def\csname sga@target@subsection:IV.9.7\endcsname{ega4.c04.subsec.9-7}
\expandafter\def\csname sga@target@subsection:IV.9.8\endcsname{ega4.c04.subsec.9-8}
\expandafter\def\csname sga@target@subsection:IV.9.9\endcsname{ega4.c04.subsec.9-9}
\let\sgaOriginalRef\ref
\let\sgaOriginalPageref\pageref
\renewcommand{\ref}[1]{%
  \ifcsname sga@target@#1\endcsname
    \hyperlink{\csname sga@target@#1\endcsname}{\sgaOriginalRef*{#1}}%
  \else
    \sgaOriginalRef{#1}%
  \fi
}
\renewcommand{\pageref}[1]{%
  \ifcsname sga@target@#1\endcsname
    \hyperlink{\csname sga@target@#1\endcsname}{\sgaOriginalPageref*{#1}}%
  \else
    \sgaOriginalPageref{#1}%
  \fi
}
\renewcommand{\eqref}[1]{%
  \ifcsname sga@target@#1\endcsname
    \hyperlink{\csname sga@target@#1\endcsname}{\textup{\tagform@{\sgaOriginalRef*{#1}}}}%
  \else
    \textup{\tagform@{\sgaOriginalRef{#1}}}%
  \fi
}
% Stable semantic routing for the reader's prose reference helpers.
\newcommand{\sgaStableOrOriginalLink}[2]{%
  \ifcsname sga@target@#1\endcsname
    \hyperlink{\csname sga@target@#1\endcsname}{#2}%
  \else
    \@ifundefined{r@#1}{#2}{\hyperref[#1]{#2}}%
  \fi
}
\RenewDocumentCommand{\sref}{o m o}{%
  \IfNoValueTF{#1}%
    {\IfNoValueTF{#3}%
      {\sgaStableOrOriginalLink{#2}{\normalfont{(\StrBehind{#2}{.})}}}%
      {\sgaStableOrOriginalLink{#2}{\normalfont{(\StrBehind{#2}{.},~{#3})}}}}%
    {\IfNoValueTF{#3}%
      {\sgaStableOrOriginalLink{#2}{\normalfont{(\textbf{#1},~\StrBehind{#2}{.})}}}%
      {\sgaStableOrOriginalLink{#2}{\normalfont{(\textbf{#1},~\StrBehind{#2}{.},~{#3})}}}}%
}
\renewcommand{\eref}[1]{%
  \sgaStableOrOriginalLink{#1}{\normalfont{(\StrBehind{#1}{.})}}}
\makeatother
% END NATIVE INPUT EGA4_reference_v2_links
\setlength{\emergencystretch}{3em}

\hypersetup{
  pdftitle={Elements of Algebraic Geometry IV, Part Three -- English Edition},
  pdfauthor={Alexandre Grothendieck and Jean Dieudonne; English translation},
  pdfsubject={Standalone English reader of EGA IV, Part Three}
}

\begin{document}

\title{Elements of Algebraic Geometry IV\\Part Three}
\author{Alexandre Grothendieck and Jean Dieudonn\'e}
\date{}
\maketitle

% BEGIN SCHOLARLY TRANSLATION ATTRIBUTION

\begingroup
\small
\begin{quote}
\textbf{Translation sources and contributors.}
The original work is by Alexandre Grothendieck and Jean Dieudonn\'e.
The EGA English edition series draws on the human community translation
\texttt{github.com/ryankeleti/ega}, source version \texttt{3d7afe28}.
Its named contributors are Mateo Carmona, Tim Holzschuh, Tim Hosgood,
Ryan Keleti, Thiago Solovera \& Nery, and Wayne Yang.
This volume belongs to the wider EGA English edition series; the community translation's contribution to that series is described in the accompanying attribution notice.
The additional translation is AI-translated; subsequent corrections, editing,
typesetting and assembly are AI-produced. The publisher did not translate
the text. Recorded AI work includes OpenAI Codex,
GPT-6 Astra, Ultra.
The accompanying \texttt{ATTRIBUTION.md} gives source links, contribution
scope and rights information. Attribution does not assert endorsement
or a licence grant from the upstream contributors.

\medskip
\textbf{Rights and availability.}
Rights in the original works and in any protected contributions to these
derivative translations and transcriptions remain with their respective
rights holders. This project does not dedicate those works or contributions
to the public domain or grant a licence over them. Existing third-party
notices and any permissions already granted remain applicable.

These editions are available free of charge for non-commercial, purely
educational purposes, on the presumption that the respective rights holders
would not object to these translations and transcriptions being freely
available alongside the French originals on NUMDAM
(\texttt{www.numdam.org/journals/PMIHES/}). This presumption is not a statement
that permission has been obtained; free access to the originals does not
itself grant permission to translate, adapt or republish them.

Rights holders who object to the availability of their material may open
an issue at \texttt{github.com/KokunoYumeto/ega-en/issues}. The publisher
will take down the affected material at the rights holder's request.
\end{quote}
\endgroup
% END SCHOLARLY TRANSLATION ATTRIBUTION


\begin{center}
\large Standalone English reader
\end{center}

\tableofcontents
% BEGIN NATIVE INPUT EGA_IV3_standalone_reference_fallbacks_d18
% Standalone-only routing for references whose targets are outside EGA IV-3.
% Mathematical source files remain unchanged.  All fallback destinations land
% on this notice; the cumulative reader retains statement-level routing.
\clearpage
\phantomsection
\section*{Cross-volume references}

Some cross-volume citations open the numbered statement in the linked
cumulative English reader. Other cross-volume destinations refer to this
page. The cumulative reader provides the surrounding text.

\phantomsection
\label{0.2.1.8}
\label{0.2.2.2}
\label{0.6.3.4}
\label{0.7.3.5}
\label{0.9.2.5}

\makeatletter
\newcommand{\EGAIVThreeRegisterExternalTargets}[1]{%
  \@for\EGAIVThreeExternalTarget:=#1\do{%
    \edef\EGAIVThreeResolvedTarget{\EGAIVThreeExternalTarget}%
    \expandafter\sgaStableTarget\expandafter{\EGAIVThreeResolvedTarget}%
  }%
}
\makeatother

\EGAIVThreeRegisterExternalTargets{%
ega4.c04.cor.1-10-2,ega4.c04.cor.1-10-4,ega4.c04.cor.1-2-3,ega4.c04.cor.1-2-8,ega4.c04.cor.1-4-5,ega4.c04.cor.1-8-6,ega4.c04.cor.1-9-9,ega4.c04.cor.17-4-2,%
ega4.c04.cor.2-1-10,ega4.c04.cor.2-1-6,ega4.c04.cor.2-2-11,ega4.c04.cor.2-2-13,ega4.c04.cor.2-2-8,ega4.c04.cor.2-3-11,ega4.c04.cor.2-3-12,ega4.c04.cor.2-3-14,%
ega4.c04.cor.2-3-5,ega4.c04.cor.2-3-8,ega4.c04.cor.2-4-10,ega4.c04.cor.2-7-2,ega4.c04.cor.3-1-4,ega4.c04.cor.4-1-2-1,ega4.c04.cor.4-1-4,ega4.c04.cor.4-2-2,%
ega4.c04.cor.4-2-6,ega4.c04.cor.4-2-8,ega4.c04.cor.4-5-11,ega4.c04.cor.4-5-13-1,ega4.c04.cor.4-5-8,ega4.c04.cor.4-6-3,ega4.c04.cor.4-6-8,ega4.c04.cor.4-7-9,%
ega4.c04.cor.5-1-11,ega4.c04.cor.5-1-5,ega4.c04.cor.5-1-8,ega4.c04.cor.5-12-4,ega4.c04.cor.5-13-2,ega4.c04.cor.5-2-2,ega4.c04.cor.5-2-3,ega4.c04.cor.5-4-2,%
ega4.c04.cor.5-6-7,ega4.c04.cor.6-1-2,ega4.c04.cor.6-1-4,ega4.c04.cor.6-15-9,ega4.c04.cor.6-3-3,ega4.c04.cor.6-5-4,ega4.c04.cor.6-7-8,ega4.c04.cor.6-8-7,%
ega4.c04.cor.6-9-3,ega4.c04.cor.7-1-13,ega4.c04.cor.7-2-9,ega4.c04.def.1-6-1,ega4.c04.def.1-9-4,ega4.c04.def.17-3-7,ega4.c04.def.2-4-2,ega4.c04.def.3-2-2,%
ega4.c04.def.3-2-5,ega4.c04.def.4-1-1,ega4.c04.def.4-5-2,ega4.c04.def.4-6-17,ega4.c04.def.4-6-2,ega4.c04.def.4-6-22,ega4.c04.def.4-6-9,ega4.c04.def.4-7-12,%
ega4.c04.def.4-7-4,ega4.c04.def.4-7-5,ega4.c04.def.5-6-2,ega4.c04.def.5-7-1,ega4.c04.def.5-7-2,ega4.c04.def.6-7-6,ega4.c04.def.6-8-1,ega4.c04.eq.3-4-1-1,%
ega4.c04.eq.4-1-1-3,ega4.c04.eq.5-1-12-2,ega4.c04.eq.5-1-2-1,ega4.c04.eq.5-6-5-2,ega4.c04.eq.5-6-5-3,ega4.c04.eq.6-11-2-1,ega4.c04.item.6-15-11,ega4.c04.item.6-5-5,%
ega4.c04.lem.1-8-4-1,ega4.c04.lem.1-8-7-1,ega4.c04.lem.2-3-1,ega4.c04.lem.3-4-1-4,ega4.c04.lem.3-4-6-1,ega4.c04.lem.5-1-10-1,ega4.c04.lem.5-10-8-1,ega4.c04.lem.5-11-1-1,%
ega4.c04.lem.5-13-7-1,ega4.c04.lem.6-11-1,ega4.c04.lem.6-12-4-1,ega4.c04.lem.6-15-8,ega4.c04.lem.6-9-2,ega4.c04.prop.1-1-2,ega4.c04.prop.1-10-3,ega4.c04.prop.1-2-2,%
ega4.c04.prop.1-2-4,ega4.c04.prop.1-2-6,ega4.c04.prop.1-2-7,ega4.c04.prop.1-3-4,ega4.c04.prop.1-4-3,ega4.c04.prop.1-4-7,ega4.c04.prop.1-5-4,ega4.c04.prop.1-6-2,%
ega4.c04.prop.1-8-2,ega4.c04.prop.1-8-7,ega4.c04.prop.1-9-10,ega4.c04.prop.1-9-11,ega4.c04.prop.1-9-5,ega4.c04.prop.16-4-11,ega4.c04.prop.17-6-3,ega4.c04.prop.2-1-13,%
ega4.c04.prop.2-1-2,ega4.c04.prop.2-1-4,ega4.c04.prop.2-1-8,ega4.c04.prop.2-2-7,ega4.c04.prop.2-3-4,ega4.c04.prop.2-4-3,ega4.c04.prop.2-4-4,ega4.c04.prop.2-4-5,%
ega4.c04.prop.2-7-1,ega4.c04.prop.3-1-3,ega4.c04.prop.3-1-6,ega4.c04.prop.3-1-7,ega4.c04.prop.3-2-1,ega4.c04.prop.3-2-6,ega4.c04.prop.3-3-4,ega4.c04.prop.3-4-3,%
ega4.c04.prop.3-4-4,ega4.c04.prop.3-4-5,ega4.c04.prop.3-4-6,ega4.c04.prop.4-2-7,ega4.c04.prop.4-3-2,ega4.c04.prop.4-5-1,ega4.c04.prop.4-5-13,ega4.c04.prop.4-5-6,%
ega4.c04.prop.4-5-9,ega4.c04.prop.4-6-1,ega4.c04.prop.4-6-10,ega4.c04.prop.4-6-5,ega4.c04.prop.4-6-6,ega4.c04.prop.4-7-10,ega4.c04.prop.4-7-11,ega4.c04.prop.4-7-13,%
ega4.c04.prop.4-7-3,ega4.c04.prop.4-7-8,ega4.c04.prop.5-1-10,ega4.c04.prop.5-1-2,ega4.c04.prop.5-1-4,ega4.c04.prop.5-10-2,ega4.c04.prop.5-12-2,ega4.c04.prop.5-12-5,%
ega4.c04.prop.5-13-1,ega4.c04.prop.5-13-3,ega4.c04.prop.5-2-1,ega4.c04.prop.5-5-3,ega4.c04.prop.5-6-1,ega4.c04.prop.5-6-10,ega4.c04.prop.5-6-4,ega4.c04.prop.5-6-5,%
ega4.c04.prop.5-7-4,ega4.c04.prop.6-11-2,ega4.c04.prop.6-11-5,ega4.c04.prop.6-14-4,ega4.c04.prop.6-15-5,ega4.c04.prop.6-2-1,ega4.c04.prop.6-2-3,ega4.c04.prop.6-4-1,%
ega4.c04.prop.6-7-1,ega4.c04.rem.4-1-3,ega4.c04.rem.5-12-3,ega4.c04.rem.5-12-6,ega4.c04.sec.5-1-1-5,ega4.c04.sec.5-2-5,ega4.c04.sec.5-5,ega4.c04.sec.5-6,%
ega4.c04.sec.5-6-3,ega4.c04.sec.5-7-3,ega4.c04.statement.1-3-8,ega4.c04.statement.1-7-19,ega4.c04.statement.1-7-4,ega4.c04.statement.1-7-5,ega4.c04.statement.1-7-9,ega4.c04.statement.1-8-1,%
ega4.c04.statement.1-9-2-1,ega4.c04.statement.1-9-6,ega4.c04.statement.2-3-3,ega4.c04.statement.4-4-1,ega4.c04.statement.5-1-12,ega4.c04.statement.6-15-1,ega4.c04.statement.6-15-2,ega4.c04.statement.7-8-3,%
ega4.c04.statement.7-8-5,ega4.c04.subsec.3-4.key-iv-3-4,ega4.c04.subsec.4-5.key-iv-4-5,ega4.c04.subsec.4-6.key-iv-4-6,ega4.c04.subsec.4-7,ega4.c04.subsec.4-8,ega4.c04.subsec.5-12.key-iv-5-12,ega4.c04.thm.1-10-1,%
ega4.c04.thm.1-8-4,ega4.c04.thm.2-3-10,ega4.c04.thm.2-4-6,ega4.c04.thm.4-1-2,ega4.c04.thm.5-10-7,ega4.c04.thm.5-8-6,ega4.c04.thm.6-9-1%
}

% Selected cross-volume destinations in the cumulative English edition.
\makeatletter
\expandafter\def\csname EGAIVThreeExternal@0.2.1.8\endcsname{https://github.com/KokunoYumeto/ega-en/releases/download/ega-en-2026-10-23-r26/001_EGA_EN_CUMULATIVE.pdf\#page=24}
\expandafter\def\csname EGAIVThreeExternal@0.3.7.1\endcsname{https://github.com/KokunoYumeto/ega-en/releases/download/ega-en-2026-10-23-r26/001_EGA_EN_CUMULATIVE.pdf\#page=32}
\expandafter\def\csname EGAIVThreeExternal@0.4.3.1\endcsname{https://github.com/KokunoYumeto/ega-en/releases/download/ega-en-2026-10-23-r26/001_EGA_EN_CUMULATIVE.pdf\#page=37}
\expandafter\def\csname EGAIVThreeExternal@0.8.2.5\endcsname{https://github.com/KokunoYumeto/ega-en/releases/download/ega-en-2026-10-23-r26/001_EGA_EN_CUMULATIVE.pdf\#page=68}
\expandafter\def\csname EGAIVThreeExternal@0.9.1.1\endcsname{https://github.com/KokunoYumeto/ega-en/releases/download/ega-en-2026-10-23-r26/001_EGA_EN_CUMULATIVE.pdf\#page=69}
\expandafter\def\csname EGAIVThreeExternal@0.9.1.3\endcsname{https://github.com/KokunoYumeto/ega-en/releases/download/ega-en-2026-10-23-r26/001_EGA_EN_CUMULATIVE.pdf\#page=69}
\expandafter\def\csname EGAIVThreeExternal@IV.1.4.6\endcsname{https://github.com/KokunoYumeto/ega-en/releases/download/ega-en-2026-10-23-r26/001_EGA_EN_CUMULATIVE.pdf\#page=745}
\expandafter\def\csname EGAIVThreeExternal@IV.1.8.4\endcsname{https://github.com/KokunoYumeto/ega-en/releases/download/ega-en-2026-10-23-r26/001_EGA_EN_CUMULATIVE.pdf\#page=750}
\expandafter\def\csname EGAIVThreeExternal@IV.1.8.4.2\endcsname{https://github.com/KokunoYumeto/ega-en/releases/download/ega-en-2026-10-23-r26/001_EGA_EN_CUMULATIVE.pdf\#page=750}
\expandafter\def\csname EGAIVThreeExternal@IV.1.8.6\endcsname{https://github.com/KokunoYumeto/ega-en/releases/download/ega-en-2026-10-23-r26/001_EGA_EN_CUMULATIVE.pdf\#page=751}
\expandafter\def\csname EGAIVThreeExternal@IV.17.9.6\endcsname{https://github.com/KokunoYumeto/ega-en/releases/download/ega-en-2026-10-23-r26/001_EGA_EN_CUMULATIVE.pdf\#page=1202}
\let\EGAIVThreeOriginalStableOrOriginalLink\sgaStableOrOriginalLink
\renewcommand{\sgaStableOrOriginalLink}[2]{%
  \ifcsname EGAIVThreeExternal@#1\endcsname
    \href{\csname EGAIVThreeExternal@#1\endcsname}{#2}%
  \else
    \EGAIVThreeOriginalStableOrOriginalLink{#1}{#2}%
  \fi
}
\makeatother
% END NATIVE INPUT EGA_IV3_standalone_reference_fallbacks_d18

\setcounter{section}{7}
% BEGIN NATIVE INPUT source_aligned/ega4-8
\section{Projective limits of preschemes}
\label{section:IV.8}

\subsection{Introduction}
\label{subsection:IV.8.1}

\oldpage[IV-3]{5}

\begin{env}[8.1.1]
\label{IV.8.1.1}
In this paragraph, we will systematically study the following situation.
Let $\mathrm{I}$ be a filtered preordered increasing set, $(\mathrm{A}_\lambda,\varphi_{\mu\lambda})$ an inductive system of rings having $\mathrm{I}$ for index set, $\mathrm{A}=\varinjlim\mathrm{A}_\lambda$ its inductive limit.
For every $\alpha\in\mathrm{I}$ and every $\mathrm{A}_\alpha$-prescheme $\mathrm{X}_\alpha$, consider the $\mathrm{A}_\lambda$-preschemes $\mathrm{X}_\lambda=\mathrm{X}_\alpha\otimes_{\mathrm{A}_\alpha}\mathrm{A}_\lambda$ for $\lambda\geqslant\alpha$, and the $\mathrm{A}$-prescheme $\mathrm{X}=\mathrm{X}_\alpha\otimes_{\mathrm{A}_\alpha}\mathrm{A}$; it is clear that the preschemes $\mathrm{X}_\lambda$ (for $\lambda\geqslant\alpha$) form a \emph{projective system}, and we will see \sref{IV.8.2.5} that $\mathrm{X}$ is a \emph{projective limit} of this system in the category of preschemes.
We propose to find conditions on $\mathrm{X}_\alpha$ or on the $\mathrm{A}_\lambda$ allowing us to prove properties of the following type:
for $\mathrm{X}$ to have a property $\mathbf{P}$ (for example, the property of being proper over $\mathrm{S}=\operatorname{Spec}(\mathrm{A})$, or irreducible, or connected, etc.), it is necessary and sufficient that there exist an index $\mu\geqslant\alpha$ such that, for every $\lambda\geqslant\mu$, $\mathrm{X}_\lambda$ has (with respect to $\mathrm{S}_\lambda=\operatorname{Spec}(\mathrm{A}_\lambda)$, as the case may be) the same property $\mathbf{P}$.
We will obtain analogous statements for properties of $\sh{O}_\mathrm{X}$-Modules, of $\mathrm{A}$-morphisms of $\mathrm{A}$-preschemes, etc.
We will also show \sref{IV.8.9.1} that the datum of an $\mathrm{A}$-prescheme of finite presentation \sref{IV.1.6.1} is essentially equivalent to the datum of an $\mathrm{A}_\lambda$-prescheme of finite presentation $\mathrm{X}_\lambda$ for some sufficiently large $\lambda$, $\mathrm{X}$ being then \emph{isomorphic to} $\mathrm{X}_\lambda\otimes_{\mathrm{A}_\lambda}\mathrm{A}$.
There are also analogous statements for $\sh{O}_\mathrm{X}$-Modules of finite presentation, their homomorphisms, the $\mathrm{A}$-morphisms of $\mathrm{A}$-preschemes of finite presentation, etc.
\end{env}

\begin{env}[8.1.2]
\label{IV.8.1.2}
The utility of such results will appear, for example, in the following questions:
\begin{enumerate}[label=\emph{\alph*)}]
  \item
  Let $\mathrm{Y}$ be a prescheme, $y$ a point of $\mathrm{Y}$, $(\mathrm{U}_\lambda)$ the projective system of filtered decreasing open affine neighbourhoods of $y$ in $\mathrm{Y}$; if $\mathrm{A}_\lambda$ is the ring of $\mathrm{U}_\lambda$, the $\mathrm{A}_\lambda$ form a filtered increasing inductive system, whose inductive limit $\mathrm{A}$ is the local ring $\sh{O}_y$; moreover, if we denote by $\mathfrak{p}_\lambda$ the prime ideal $\mathfrak{j}_y$ in the ring $\mathrm{A}_\lambda$, the inductive system $(\mathrm{A}_\lambda)$ is cofinal in every inductive system $(\mathrm{A}_\alpha)_f$, where $f$ ranges over $\mathrm{A}_\alpha-\mathfrak{p}_\alpha$
\oldpage[IV-3]{6}
(for a fixed $\alpha$), since the $\mathrm{D}(f)$ form a fundamental system of neighbourhoods of $y$ in $\mathrm{U}_\alpha$, hence in $\mathrm{Y}$.
The results of the present paragraph will therefore imply that the algebraic geometry of $\sh{O}_y$-preschemes of finite presentation (and the theory of Modules of finite presentation on these preschemes) is essentially equivalent to the algebraic geometry of preschemes of finite presentation on ``sufficiently small'' open neighbourhoods of $y$.
Thus, the statement \sref{IV.8.10.5},~(xii) implies that if a morphism of finite presentation $f:\mathrm{X}\to\mathrm{Y}$ is such that $\mathrm{X}\times_\mathrm{Y}\operatorname{Spec}(\sh{O}_y)$ is proper over $\operatorname{Spec}(\sh{O}_y)$, it is necessary and sufficient that there exist an open neighbourhood $\mathrm{U}$ of $y$ in $\mathrm{Y}$ such that $f^{-1}(\mathrm{U})$ is proper over $\mathrm{U}$.

An important particular case, classical in a certain sense, is the one where $\mathrm{Y}$ is integral and where $y$ is its generic point, so that $\sh{O}_y$ is nothing other than \emph{the field} $\mathrm{R}(\mathrm{Y})=\mathrm{K}$ \emph{of rational functions on} $\mathrm{Y}$.
The results of the present paragraph then amount to interpreting the algebraic geometry over $\mathrm{K}$ in terms of algebraic geometry over ``sufficiently small'' non-empty open subsets of $\mathrm{Y}$, that is to say, intuitively, in terms of ``families'' of geometric objects indexed by the points of such an open set.
This point of view has moreover been commonly used for a long time, not only in Algebraic Geometry over algebraically closed fields, but also in the arithmetic study of varieties defined over a \emph{number field} $\mathrm{K}$ (a finite extension of $\mathbf{Q}$), considering the latter as the field of fractions of its ring of integers $\mathrm{A}$ (the ``theory of \emph{reduction modulo} $\mathfrak{p}$'' ($\mathfrak{p}$ a prime ideal of $\mathrm{A}$); cf.\ \sref[I]{I.3.7}).
The results of \textsection\textsection\,8 and~9 thus provide, among other things, some of the foundations of the language of ``reduction modulo $\mathfrak{p}$'' in arithmetic.

We note that in the example considered here, the morphisms $\mathrm{S}_\mu\to\mathrm{S}_\lambda$ (for $\lambda\leqslant\mu$) are the canonical \emph{open immersions} $\mathrm{U}_\mu\hookrightarrow\mathrm{U}_\lambda$, and \emph{a fortiori} are \emph{flat} morphisms (but not faithfully flat in general), which explains the interest of the statements that appeal to such a restriction.

  \item
  Suppose that the $\mathrm{A}_\lambda$ are \emph{fields}, so that $\mathrm{A}=\varinjlim\mathrm{A}_\lambda$ is also a field.
This case arises generally when one starts from geometric data over an arbitrary field $\mathrm{K}$, which one considers as an extension of a field $k$ (for example the prime subfield of $\mathrm{K}$).
There is then interest in considering $\mathrm{K}$ as an inductive limit of its sub-extensions that are \emph{of finite type over $k$}, which allows us to reduce many questions to the case where $\mathrm{K}$ is a finite type extension of $k$.
Using also the method sketched in \emph{a)}, one can then reduce generally to the case of a base ring $\mathrm{A}$ which is an \emph{integral algebra of finite type over $k$}.

We note that in this example, the morphisms $\mathrm{S}_\mu\to\mathrm{S}_\lambda$ are \emph{faithfully flat}.

  \item
  Suppose that we are interested in properties, local on $\mathrm{Y}$, of preschemes of finite presentation over an arbitrary prescheme $\mathrm{Y}$, which we can from now on suppose to be an affine scheme of ring $\mathrm{A}$.
There is then interest in considering $\mathrm{A}$ as an inductive limit of its subrings that are \emph{$\mathbf{Z}$-algebras of finite type}, which allows us to reduce many questions to the case where $\mathrm{Y}$ is the spectrum of such an algebra.
This is the explication of the ``Kroneckerian point of view'' according to which Algebraic Geometry reduces to the
\oldpage[IV-3]{7}
Algebraic Geometry of preschemes of finite type over $\mathbf{Z}$ (which one sometimes qualifies as ``Absolute Algebraic Geometry'').
This example shows us in particular that in most questions ``relative'' to a prescheme of base $\mathrm{Y}$, one can reduce to the case where $\mathrm{Y}$ is \emph{Noetherian}.

We note that in this example, in contrast to the previous ones, the morphisms $\mathrm{S}_\mu\to\mathrm{S}_\lambda$ have in general no particular regularity property.
\end{enumerate}

In what follows, when we apply the results that follow to any one of the three particular situations that have just been described, we will not go into detail about the procedure that allows us to make these applications, being content to refer to what precedes.
\end{env}

\begin{env}[8.1.3]
\label{IV.8.1.3}
In example \emph{a)} of \sref{IV.8.1.2}, we saw that if $\mathrm{Y}$ is an integral prescheme with generic point $y$, $f:\mathrm{X}\to\mathrm{Y}$ a morphism of finite presentation, and if the generic fibre $f^{-1}(y)$ is proper over $\boldsymbol{k}(y)$, then there exists an open neighbourhood $\mathrm{U}$ of $y$ in $\mathrm{Y}$ such that $f^{-1}(\mathrm{U})$ is proper over $\mathrm{U}$; it follows \emph{a fortiori} that for every $s\in\mathrm{U}$, $f^{-1}(s)$ is proper over $\boldsymbol{k}(s)$.
It turns out that one needs a converse, ensuring that if $f^{-1}(s)$ is proper over $\boldsymbol{k}(s)$ for ``sufficiently many'' points $s\in\mathrm{U}$, then $f^{-1}(y)$ is proper over $\boldsymbol{k}(y)$.
For example, suppose that $\mathrm{X}$ and $\mathrm{Y}$ are algebraic preschemes over an algebraically closed field $k$ (one can take for $k$ the field $\mathbf{C}$ of complex numbers, to fix ideas); one sometimes needs to know that if, for every $s\in\mathrm{Y}$, \emph{rational over $k$}, the fibre $f^{-1}(s)$ is proper over $\boldsymbol{k}(s)$, then $f^{-1}(y)$ is proper over $\boldsymbol{k}(y)$, and consequently $f^{-1}(\mathrm{U})$ is proper over $\mathrm{U}$ for a neighbourhood $\mathrm{U}$ of $y$\footnote{We note that such a statement is ultimately purely geometric in nature, in the sense that it only appeals to rational points over $k$, and not to generic points; for example, when $k=\mathbf{C}$, this statement has an obvious topological significance for the analyst, interpreting ``proper'' in the topological sense of the term, for the underlying analytic spaces formed by the points of $\mathrm{X}$ and $\mathrm{Y}$ rational over $\mathbf{C}$.}.
Now this statement will follow easily from the following: the set $\mathrm{E}$ of points $s\in\mathrm{Y}$ such that $f^{-1}(s)$ is proper over $\boldsymbol{k}(s)$ is \emph{constructible} (and therefore identical to all of $\mathrm{Y}$ if it contains the closed points of $\mathrm{Y}$, by the theorem of zeros of Hilbert \sref[I]{I.10.4.8}); this amounts again to saying that if $f^{-1}(y)$ \emph{is not proper} over $\boldsymbol{k}(y)$, then there exists an open neighbourhood $\mathrm{U}$ of $y$ such that $f^{-1}(s)$ is not proper over $\boldsymbol{k}(s)$ for every $s\in\mathrm{U}$ (cf.\ \sref{IV.9.6.1},~(iv))).
This example illustrates the interest of systematically developing \emph{criteria of constructibility}: this is what will be done in \textsection\,9.
\end{env}

\subsection{Projective limits of preschemes}
\label{subsection:IV.8.2}

\begin{env}[8.2.1]
\label{IV.8.2.1}
Let $\mathrm{S}_0$ be a ringed space, $\mathrm{L}$ a filtered preordered increasing set, $(\sh{A}_\lambda,\varphi_{\mu\lambda})$ an inductive system of $\sh{O}_{\mathrm{S}_0}$-Algebras (not necessarily commutative), having $\mathrm{L}$ for index set.
We know that, considered as an inductive system of $\sh{O}_{\mathrm{S}_0}$-Modules, $(\sh{A}_\lambda,\varphi_{\mu\lambda})$ admits an inductive limit $\sh{A}$; let us denote by $\varphi_\lambda:\sh{A}_\lambda\to\sh{A}$ the canonical homomorphism of $\sh{O}_{\mathrm{S}_0}$-Modules.
Let $m_\lambda:\sh{A}_\lambda\otimes\sh{A}_\lambda\to\sh{A}_\lambda$ be the homomorphism of $\sh{O}_{\mathrm{S}_0}$-Modules that defines multiplication in the $\sh{O}_{\mathrm{S}_0}$-Algebra $\sh{A}_\lambda$; the hypothesis on the $\varphi_{\mu\lambda}$ implies that the $m_\lambda$ form an inductive system of homomorphisms,
\oldpage[IV-3]{8}
and since the functor $\varinjlim$ commutes with tensor products, $m=\varinjlim m_\lambda$ is a homomorphism $\sh{A}\otimes\sh{A}\to\sh{A}$ of $\sh{O}_{\mathrm{S}_0}$-Modules; by passing to the limit on the commutative diagrams expressing the associativity of $m_\lambda$ and the existence of the identity section in $\sh{A}_\lambda$, one sees that $m$ defines on $\sh{A}$ a structure of $\sh{O}_{\mathrm{S}_0}$-Algebra and that $\varphi_\lambda$ is a homomorphism of $\sh{O}_{\mathrm{S}_0}$-Algebras for every $\lambda\in\mathrm{L}$.
Furthermore, $\sh{A}$ is the inductive limit of the system $(\sh{A}_\lambda,\varphi_{\mu\lambda})$ \emph{in the category of $\sh{O}_{\mathrm{S}_0}$-Algebras}, in other words, for every $\sh{O}_{\mathrm{S}_0}$-Algebra $\sh{B}$, the canonical map
\begin{equation}
  \operatorname{Hom}_{\sh{O}_{\mathrm{S}_0}\text{-}\mathrm{Alg.}}(\sh{A},\sh{B})\to\varprojlim\operatorname{Hom}_{\sh{O}_{\mathrm{S}_0}\text{-}\mathrm{Alg.}}(\sh{A}_\lambda,\sh{B})
  \tag{8.2.1.1}
\label{IV.8.2.1.1}
\end{equation}
which to every homomorphism $f:\sh{A}\to\sh{B}$ of $\sh{O}_{\mathrm{S}_0}$-Algebras associates the family $(f\circ\varphi_\lambda)$, is \emph{bijective}.
Indeed, we already know that it is injective and identifies $\operatorname{Hom}_{\sh{O}_{\mathrm{S}_0}\text{-}\mathrm{Alg.}}(\sh{A},\sh{B})$ with a subset of $\varprojlim\operatorname{Hom}_{\sh{O}_{\mathrm{S}_0}\text{-}\mathrm{Mod.}}(\sh{A}_\lambda,\sh{B})$; it all comes down to seeing that if $(f_\lambda)$ is an inductive system of $\sh{O}_{\mathrm{S}_0}$-Algebra homomorphisms, $f_\lambda:\sh{A}_\lambda\to\sh{B}$, its inductive limit $f:\sh{A}\to\sh{B}$, which is by definition a homomorphism of $\sh{O}_{\mathrm{S}_0}$-\emph{Modules}, is also a homomorphism of $\sh{O}_{\mathrm{S}_0}$-\emph{Algebras}; but this results from passage to the inductive limit in the commutative diagram of $\sh{O}_{\mathrm{S}_0}$-Modules
\[
\xymatrix{
  \sh{A}_\lambda\otimes\sh{A}_\lambda \ar[r]^{f_\lambda\otimes f_\lambda} \ar[d]_{m_\lambda} & \sh{B}\otimes\sh{B} \ar[d] \\
  \sh{A}_\lambda \ar[r]_{f_\lambda} & \sh{B}
}
\]
and from the fact that the functor $\varinjlim$ commutes with tensor products.
We note finally that if the $\sh{A}_\lambda$ are commutative $\sh{O}_{\mathrm{S}_0}$-Algebras, the same holds for $\sh{A}$.
\end{env}

\begin{env}[8.2.2]
\label{IV.8.2.2}
Suppose now that $\mathrm{S}_0$ is a \emph{prescheme}, and that the $\sh{A}_\lambda$ are (commutative) \emph{quasi-coherent} $\sh{O}_{\mathrm{S}_0}$-Algebras; we know then that $\sh{A}=\varinjlim\sh{A}_\lambda$ is a quasi-coherent $\sh{O}_{\mathrm{S}_0}$-Algebra \sref[I]{I.4.1.1}.
Let us denote by $\mathrm{S}_\lambda$ (resp.\ $\mathrm{S}$) the spectrum of the $\sh{O}_{\mathrm{S}_0}$-Algebra $\sh{A}_\lambda$ (resp.\ $\sh{A}$) \sref[II]{II.1.3.1}, and let $u_{\lambda\mu}:\mathrm{S}_\mu\to\mathrm{S}_\lambda$ (for $\lambda\leqslant\mu$) and $u_\lambda:\mathrm{S}\to\mathrm{S}_\lambda$ be the $\mathrm{S}_0$-morphisms corresponding to the homomorphisms $\varphi_{\mu\lambda}$ and $\varphi_\lambda$ respectively \sref[II]{II.1.2.7}; it is clear that $(\mathrm{S}_\lambda,u_{\lambda\mu})$ is a \emph{projective system} in the category of $\mathrm{S}_0$-preschemes.
We note that the $u_{\lambda\mu}$ and $u_\lambda$ are \emph{affine} morphisms \sref[II]{II.1.6.2}, hence \emph{quasi-compact} and \emph{separated}.
\end{env}

\begin{proposition}[8.2.3]
\label{IV.8.2.3}
With the notations of \sref{IV.8.2.2}, the morphisms $u_\lambda:\mathrm{S}\to\mathrm{S}_\lambda$ make $\mathrm{S}$ a projective limit of the projective system $(\mathrm{S}_\lambda,u_{\lambda\mu})$ in the category of preschemes.
Furthermore, if $h:\mathrm{S}_0\to\mathrm{T}$ is a morphism, making every $\mathrm{S}_0$-prescheme a $\mathrm{T}$-prescheme, $\mathrm{S}$ is also a projective limit of the system $(\mathrm{S}_\lambda,u_{\lambda\mu})$ in the category of $\mathrm{T}$-preschemes.
\end{proposition}

\begin{proof}
Let us first prove the second assertion of the statement in the case $\mathrm{T}=\mathrm{S}_0$.
\oldpage[IV-3]{9}
It all comes down to proving that if $\mathrm{X}$ is any $\mathrm{S}_0$-prescheme, the canonical map
\begin{equation}
  \operatorname{Hom}_{\mathrm{S}_0}(\mathrm{X},\mathrm{S})\to\varprojlim\operatorname{Hom}_{\mathrm{S}_0}(\mathrm{X},\mathrm{S}_\lambda)
  \tag{8.2.3.1}
\label{IV.8.2.3.1}
\end{equation}
which to every $\mathrm{S}_0$-morphism $v:\mathrm{X}\to\mathrm{S}$ associates the family $(u_\lambda\circ v)$, is \emph{bijective}.
Now, if $g:\mathrm{X}\to\mathrm{S}_0$ is the structural morphism and if we set $\sh{B}=g_*(\sh{O}_\mathrm{X})$, which is an $\sh{O}_{\mathrm{S}_0}$-Algebra, the map \sref{IV.8.2.3.1} identifies canonically with \sref{IV.8.2.1.1} \sref[II]{II.1.2.7}, and the conclusion follows from what we have seen in \sref{IV.8.2.1}.

The other assertions of \sref{IV.8.2.3} are consequences of the following general lemma:
\end{proof}

\begin{lemma}[8.2.4]
\label{IV.8.2.4}
Let $\mathbf{C}$ be a category, $\mathrm{T}$ an object of $\mathbf{C}$, $\mathbf{C}_\mathrm{T}$ the subcategory of objects of $\mathbf{C}$ over $\mathrm{T}$.
Let $(\mathrm{S}_\lambda,u_{\lambda\mu})$ be a projective system in $\mathbf{C}_\mathrm{T}$; then every projective limit of this system in $\mathbf{C}_\mathrm{T}$ is also a projective limit in $\mathbf{C}$, and conversely.
\end{lemma}

\begin{proof}
Let $f_\lambda:\mathrm{S}_\lambda\to\mathrm{T}$ be the structural morphism.
Suppose that $\mathrm{S}$ is a projective limit of $(\mathrm{S}_\lambda,u_{\lambda\mu})$ in $\mathbf{C}$, and let us denote by $u_\lambda:\mathrm{S}\to\mathrm{S}_\lambda$ the canonical corresponding morphisms.
Consider a projective system of $\mathrm{T}$-morphisms (in $\mathbf{C}$) $w_\lambda:\mathrm{Y}\to\mathrm{S}_\lambda$, where $\mathrm{Y}\in\mathbf{C}_\mathrm{T}$.
There exists by hypothesis a unique morphism (in $\mathbf{C}$) $w:\mathrm{Y}\to\mathrm{S}$ such that $w_\lambda=u_\lambda\circ w$ for every $\lambda$.
The hypothesis that the $u_{\lambda\mu}$ are $\mathrm{T}$-morphisms implies that the morphisms $f_\lambda\circ u_\lambda:\mathrm{S}\to\mathrm{T}$ are all the same, and this composite $f$ makes $\mathrm{S}$ a $\mathrm{T}$-object.
If $g:\mathrm{Y}\to\mathrm{T}$ is the structural morphism of $\mathrm{Y}$, one has $f\circ w=f_\lambda\circ u_\lambda\circ w=f_\lambda\circ w_\lambda=g$ for every $\lambda$, which proves that $w$ is a $\mathrm{T}$-morphism.
Conversely, suppose (with the same notations) that $\mathrm{S}$ is a projective limit of $(\mathrm{S}_\lambda,u_{\lambda\mu})$ in $\mathbf{C}_\mathrm{T}$, and consider now a projective system of morphisms (in $\mathbf{C}$) $w_\lambda:\mathrm{Y}\to\mathrm{S}_\lambda$.
The composites $f_\lambda\circ w_\lambda:\mathrm{Y}\to\mathrm{T}$ are then all the same: indeed, for any two indices $\lambda$, $\mu$, there is an index $\nu$ such that $\lambda\leqslant\nu$ and $\mu\leqslant\nu$, whence $f_\nu=f_\lambda\circ u_{\lambda\nu}=f_\mu\circ u_{\mu\nu}$; since $w_\lambda=u_{\lambda\nu}\circ w_\nu$ and $w_\mu=u_{\mu\nu}\circ w_\nu$, one has $f_\lambda\circ w_\lambda=f_\lambda\circ u_{\lambda\nu}\circ w_\nu=f_\nu\circ w_\nu$ and $f_\mu\circ w_\mu=f_\mu\circ u_{\mu\nu}\circ w_\nu=f_\nu\circ w_\nu$.
If $g:\mathrm{Y}\to\mathrm{T}$ is the unique morphism thus defined, $g$ makes $\mathrm{Y}$ a $\mathrm{T}$-object, and the $w_\lambda$ are $\mathrm{T}$-morphisms; they hence have a projective limit $w:\mathrm{Y}\to\mathrm{S}$ which is a $\mathrm{T}$-morphism, and \emph{a fortiori} a morphism of $\mathbf{C}$; moreover, the first part of the argument shows that any projective limit $w'$ (in $\mathbf{C}$) of the projective system $(w_\lambda)$ is also a $\mathrm{T}$-morphism, and therefore equal to $w$, which completes the proof of the lemma.
\end{proof}

\begin{proposition}[8.2.5]
\label{IV.8.2.5}
Under the conditions of \sref{IV.8.2.2}, let $\alpha$ be an element of $\mathrm{L}$, $\mathrm{X}_\alpha$ an $\mathrm{S}_\alpha$-prescheme.
For every $\lambda\geqslant\alpha$, set $\mathrm{X}_\lambda=\mathrm{X}_\alpha\times_{\mathrm{S}_\alpha}\mathrm{S}_\lambda$, and for $\alpha\leqslant\lambda\leqslant\mu$, set $v_{\lambda\mu}=1_{\mathrm{X}_\alpha}\times u_{\lambda\mu}$, so that $(\mathrm{X}_\lambda,v_{\lambda\mu})$ is a projective system of $\mathrm{X}_\alpha$-preschemes, whose index set is formed of the $\lambda\geqslant\alpha$ in $\mathrm{L}$.
Set likewise $\mathrm{X}=\mathrm{X}_\alpha\times_{\mathrm{S}_\alpha}\mathrm{S}$ and $v_\lambda=1_{\mathrm{X}_\alpha}\times u_\lambda$.
Then the $\mathrm{X}_\alpha$-morphisms $v_\lambda:\mathrm{X}\to\mathrm{X}_\lambda$ make $\mathrm{X}$ a projective limit of the projective system $(\mathrm{X}_\lambda,v_{\lambda\mu})$ in the category of $\mathrm{X}_\alpha$-preschemes, or in the category of all preschemes.
\end{proposition}

\begin{proof}
This will also follow from the following general lemma:
\end{proof}

\begin{lemma}[8.2.6]
\label{IV.8.2.6}
Let $\mathbf{C}$ be a category in which fibre products exist, $q:\mathrm{T}'\to\mathrm{T}$ a morphism of $\mathbf{C}$, $\mathbf{C}_\mathrm{T}$ (resp.\ $\mathbf{C}_{\mathrm{T}'}$) the category of objects of $\mathbf{C}$ over
\oldpage[IV-3]{10}
$\mathrm{T}$ (resp.\ $\mathrm{T}'$).
Let $(\mathrm{S}_\lambda,u_{\lambda\mu})$ be a projective system (not necessarily filtered) in $\mathbf{C}_\mathrm{T}$, and set $\mathrm{S}'_\lambda=\mathrm{S}_\lambda\times_\mathrm{T}\mathrm{T}'$, $u'_{\lambda\mu}=u_{\lambda\mu}\times 1_{\mathrm{T}'}$, so that $(\mathrm{S}'_\lambda,u'_{\lambda\mu})$ is a projective system in $\mathbf{C}_{\mathrm{T}'}$.
Then, if $(\mathrm{S},u_\lambda)$ is a projective limit of $(\mathrm{S}_\lambda,u_{\lambda\mu})$ in $\mathbf{C}_\mathrm{T}$, $(\mathrm{S}\times_\mathrm{T}\mathrm{T}',u_\lambda\times 1_{\mathrm{T}'})$ is a projective limit of $(\mathrm{S}'_\lambda,u'_{\lambda\mu})$ in $\mathbf{C}_{\mathrm{T}'}$.
\end{lemma}

\begin{proof}
By hypothesis, for every $\lambda$, one has a commutative diagram
\[
\xymatrix{
  \mathrm{S}' \ar[r]^{u'_\lambda} \ar[d]_p & \mathrm{S}'_\lambda \ar[r]^{h'_\lambda} \ar[d]_{p_\lambda} & \mathrm{T}' \ar[d]^q \\
  \mathrm{S} \ar[r]_{u_\lambda} & \mathrm{S}_\lambda \ar[r]_{h_\lambda} & \mathrm{T}
}
\]
where we have set $\mathrm{S}'=\mathrm{S}\times_\mathrm{T}\mathrm{T}'$, $u'_\lambda=u_\lambda\times 1_{\mathrm{T}'}$, $h'_\lambda=h_\lambda\times 1_{\mathrm{T}'}$.
Let $\mathrm{Y}$ be a $\mathrm{T}'$-object, $g':\mathrm{Y}\to\mathrm{T}'$, and consider a projective system of $\mathrm{T}'$-morphisms $w'_\lambda:\mathrm{Y}\to\mathrm{S}'_\lambda$.
Then $\mathrm{Y}$ is a $\mathrm{T}$-object for the morphism $g=q\circ g'$, and the $w_\lambda=p_\lambda\circ w'_\lambda$ are $\mathrm{T}$-morphisms since $h_\lambda\circ w_\lambda=h_\lambda\circ p_\lambda\circ w'_\lambda=q\circ h'_\lambda\circ w'_\lambda=q\circ g'$ by hypothesis.
Furthermore, one verifies immediately that $(w_\lambda)$ is a projective system.
There thus exists by hypothesis a unique $\mathrm{T}$-morphism $w:\mathrm{Y}\to\mathrm{S}$ such that $u_\lambda\circ w=w_\lambda$ for every $\lambda$.
By the definition of the fibre product, there is a unique $\mathrm{T}'$-morphism $w':\mathrm{Y}\to\mathrm{S}'$ such that $p\circ w'=w$.
Then $u_\lambda\circ p\circ w'=u_\lambda\circ w=w_\lambda=p_\lambda\circ w'_\lambda$, which can also be written $p_\lambda\circ u'_\lambda\circ w'=p_\lambda\circ w'_\lambda$; on the other hand, since $w'$ is a $\mathrm{T}'$-morphism, $h'_\lambda\circ u'_\lambda\circ w'=g'=h'_\lambda\circ w'_\lambda$. The definition of $\mathrm{S}'_\lambda$ as the fibre product $\mathrm{S}_\lambda\times_\mathrm{T}\mathrm{T}'$ therefore gives $u'_\lambda\circ w'=w'_\lambda$, and it is immediate that $w'$ is the unique $\mathrm{T}'$-morphism satisfying these relations, whence the lemma.
\end{proof}

\begin{remark}[8.2.7]
\label{IV.8.2.7}
Given a ringed space $\mathrm{S}$, the inductive limits with respect to a preordered set $\mathrm{L}$ (not necessarily filtered) exist in the category of commutative $\sh{O}_\mathrm{S}$-Algebras, since the filtered inductive limit exists by \sref{IV.8.2.1} and, on the other hand, for any two homomorphisms of $\sh{O}_\mathrm{S}$-Algebras $\sh{A}\to\sh{B}$, $\sh{A}\to\sh{C}$, the tensor product $\sh{B}\otimes_{\sh{A}}\sh{C}$ is the corresponding ``amalgamated sum'' in this category.

When $\mathrm{S}$ is a prescheme, we know that the tensor product $\sh{B}\otimes_{\sh{A}}\sh{C}$ is a quasi-coherent $\sh{O}_\mathrm{S}$-Algebra when this is so for $\sh{A}$, $\sh{B}$, and $\sh{C}$ \sref[I]{I.1.3.13}; from this one concludes that, in the category of quasi-coherent $\sh{O}_\mathrm{S}$-Algebras, inductive limits for any preordered index set always exist.
This allows us to generalise the definition of a projective limit of preschemes and propositions~\sref{IV.8.2.3} and \sref{IV.8.2.5} to the case where the preordered set $\mathrm{L}$ is not necessarily filtered.
\end{remark}

\begin{env}[8.2.8]
\label{IV.8.2.8}
The notations being those of \sref{IV.8.2.2}, set $u_{\lambda\mu}=(\psi_{\lambda\mu},\theta_{\lambda\mu})$ and $u_\lambda=(\psi_\lambda,\theta_\lambda)$; $(\mathrm{S}_\lambda,\psi_{\lambda\mu})$ is thus a projective system of topological spaces and $(\psi_\lambda)$ a projective system of continuous maps $\mathrm{S}\to\mathrm{S}_\lambda$ between the spaces underlying respectively the preschemes $\mathrm{S}$ and $\mathrm{S}_\lambda$.
\end{env}

\begin{proposition}[8.2.9]
\label{IV.8.2.9}
With the notations of \sref{IV.8.2.8}, the projective limit of the projective system $(\psi_\lambda)$ of continuous maps is a homeomorphism of the underlying space of $\mathrm{S}$ onto the projective limit of the projective system $(\mathrm{S}_\lambda,\psi_{\lambda\mu})$ of topological spaces.
\end{proposition}

\begin{proof}
\oldpage[IV-3]{11}
Let $\mathrm{T}$ be the topological limit space of the system $(\mathrm{S}_\lambda,\psi_{\lambda\mu})$ and set $\psi=\varprojlim\psi_\lambda:\mathrm{S}\to\mathrm{T}$.
We can restrict to the case where $\mathrm{S}_0=\mathrm{S}_\alpha$ for some $\alpha\in\mathrm{L}$, and $\lambda\geqslant\alpha$.

\medskip
\noindent
(i) Let us show first that the topology of $\mathrm{S}$ is the inverse image by $\psi$ of the topology of $\mathrm{T}$; in other words, if $\pi_\lambda:\mathrm{T}\to\mathrm{S}_\lambda$ is the canonical map, it amounts to showing that every open subset of $\mathrm{S}$ is a union of open subsets of the form $\psi^{-1}(\pi_\lambda^{-1}(\mathrm{U}_\lambda))$, where $\mathrm{U}_\lambda$ is open in $\mathrm{S}_\lambda$.
Now, every open subset of $\mathrm{S}$ is by definition a union of open subsets obtained in the following way: one considers an open affine $\mathrm{U}_0$ of $\mathrm{S}_0$, of ring $\mathrm{A}_0$, so that $\psi_\alpha^{-1}(\mathrm{U}_0)$ is the open affine subset of $\mathrm{S}$ of ring $\mathrm{A}=\Gamma(\mathrm{U}_0,\sh{A})$, then one takes an element $f\in\mathrm{A}$ and considers in $\operatorname{Spec}(\mathrm{A})$, identified with $\psi_\alpha^{-1}(\mathrm{U}_0)$, the open subset $\mathrm{D}(f)$.
These are the open subsets $\mathrm{D}(f)$ which form a base of the topology of $\mathrm{S}$ \sref[II]{II.1.3.1}.
Now, if one sets $\mathrm{A}_\lambda=\Gamma(\mathrm{U}_0,\sh{A}_\lambda)$, one has $\mathrm{A}=\varinjlim\mathrm{A}_\lambda$ \sref[I]{I.1.3.9}, hence there exists an index $\lambda$ such that $f$ is the canonical image of an element $f_\lambda\in\mathrm{A}_\lambda$; one has then $\mathrm{D}(f)=\psi_\lambda^{-1}(\mathrm{D}(f_\lambda))$ \sref[I]{I.1.2.2}, and since $\psi_\lambda=\pi_\lambda\circ\psi$, our assertion is proved.

\medskip
\noindent
(ii) Let us now show that $\psi$ is \emph{bijective}, which will complete the proof.
Since $\mathrm{S}$ is a Kolmogoroff space, it already follows from (i) that $\psi$ is injective, and it remains to show that $\psi$ is surjective.
We can evidently replace $\mathrm{T}$ by an open subset $\pi_\alpha^{-1}(\mathrm{U}_0)$, where $\mathrm{U}_0$ is an open affine in $\mathrm{S}_\alpha=\mathrm{S}_0$, hence we can restrict to the case where the $\mathrm{S}_\lambda$ and $\mathrm{S}$ are affine, in other words $\sh{A}_\lambda$ is the sheaf of algebras associated to an $\mathrm{A}_0$-algebra $\mathrm{A}_\lambda$, and $\sh{A}$ the sheaf of algebras associated to $\mathrm{A}=\varinjlim\mathrm{A}_\lambda$; we will still denote $\varphi_{\mu\lambda}:\mathrm{A}_\lambda\to\mathrm{A}_\mu$ and $\varphi_\lambda:\mathrm{A}_\lambda\to\mathrm{A}$ the canonical homomorphisms.
By definition, an element of $\mathrm{T}$ is a family $(\mathfrak{p}_\lambda)_{\lambda\in\mathrm{L}}$, where $\mathfrak{p}_\lambda$ is a prime ideal of $\mathrm{A}_\lambda$ and where $\mathfrak{p}_\lambda=\varphi_{\mu\lambda}^{-1}(\mathfrak{p}_\mu)$ for $\lambda\leqslant\mu$.
We know (\sref{IV.5.13.3} and \sref{IV.5.13.1}) that there exists a prime ideal $\mathfrak{p}$ of $\mathrm{A}$ such that $\mathfrak{p}_\lambda=\varphi_\lambda^{-1}(\mathfrak{p})$ for every $\lambda\in\mathrm{L}$, which completes the proof.
\end{proof}

In particular, we have thus proved the

\begin{corollary}[8.2.10]
\label{IV.8.2.10}
Let $(\mathrm{A}_\lambda)_{\lambda\in\mathrm{L}}$ be a filtered inductive system of rings, and let $\mathrm{A}=\varinjlim\mathrm{A}_\lambda$, $\varphi_\lambda:\mathrm{A}_\lambda\to\mathrm{A}$ be the canonical homomorphisms.
The canonical map $\mathfrak{p}\mapsto(\varphi_\lambda^{-1}(\mathfrak{p}))$ is a homeomorphism of $\operatorname{Spec}(\mathrm{A})$ onto the topological space $\varprojlim\operatorname{Spec}(\mathrm{A}_\lambda)$.
\end{corollary}

\begin{corollary}[8.2.11]
\label{IV.8.2.11}
With the notations of \sref{IV.8.2.8}, for every quasi-compact open subset $\mathrm{U}$ of $\mathrm{S}$, there exist an index $\lambda$ and a quasi-compact open subset $\mathrm{U}_\lambda$ of $\mathrm{S}_\lambda$ such that $\mathrm{U}=\psi_\lambda^{-1}(\mathrm{U}_\lambda)$.
\end{corollary}

\begin{proof}
This follows from the fact that, by definition of the projective limit topology, the $\psi_\lambda^{-1}(\mathrm{U}_\lambda)$ ($\mathrm{U}_\lambda$ a quasi-compact open subset of $\mathrm{S}_\lambda$) form a base of the topology of $\mathrm{S}$, and from the fact that the index set $\mathrm{L}$ is filtered.
\end{proof}

\begin{corollary}[8.2.12]
\label{IV.8.2.12}
With the notations of \sref{IV.8.2.8}, the inductive limit of the inductive system of homomorphisms $\theta_\lambda^\sharp:\psi_\lambda^*(\sh{O}_{\mathrm{S}_\lambda})\to\sh{O}_\mathrm{S}$ of sheaves of rings on $\mathrm{S}$ is an isomorphism
\begin{equation}
  \varinjlim_\lambda\psi_\lambda^*(\sh{O}_{\mathrm{S}_\lambda})\xrightarrow{\sim}\sh{O}_\mathrm{S}.
  \label{IV.8.2.12.1}
\tag{8.2.12.1}
\end{equation}
\end{corollary}

\begin{proof}
We can evidently suppose the $\mathrm{S}_\lambda$ affine; with the notations of the proof of \sref{IV.8.2.9}, it all comes down to seeing that the inductive limit of the inductive system of canonical maps $(\mathrm{A}_\lambda)_{\mathfrak{p}_\lambda}\to\mathrm{A}_\mathfrak{p}$ is an isomorphism, which is nothing other than \sref{IV.5.13.3},~(ii).
\end{proof}

\begin{proposition}[8.2.13]
\label{IV.8.2.13}
Suppose that the morphisms $u_{\lambda\mu}:\mathrm{S}_\mu\to\mathrm{S}_\lambda$ are open immersions, so that $\mathrm{S}_\mu$ identifies with the subscheme induced on an open subset of $\mathrm{S}_\lambda$ for $\lambda\leqslant\mu$.
Then, for every $\alpha\in\mathrm{L}$, the underlying space of the prescheme $\mathrm{S}$ identifies with the subspace of $\mathrm{S}_\alpha$ given by the intersection of the $\mathrm{S}_\lambda$ for $\lambda\geqslant\alpha$, and the structural sheaf $\sh{O}_\mathrm{S}$ identifies with the sheaf induced \textnormal{(}$\mathrm{G}$, II, $1.5$\textnormal{)} by $\sh{O}_{\mathrm{S}_\alpha}$ on this intersection; more generally, for every $\sh{O}_{\mathrm{S}_\alpha}$-Module $\sh{F}_\alpha$, $u_\alpha^*(\sh{F}_\alpha)$ identifies with the $\sh{O}_\mathrm{S}$-Module induced by $\sh{F}_\alpha$ on $\mathrm{S}$.
\end{proposition}

\begin{proof}
\oldpage[IV-3]{12}
The first assertion follows from \sref{IV.8.2.9}, in view of the definition of a projective limit of topological spaces; furthermore, all the $\psi_\lambda^*(\sh{O}_{\mathrm{S}_\lambda})$ are equal to the sheaf induced by $\sh{O}_{\mathrm{S}_\alpha}$ on $\mathrm{S}$ by definition \sref[$0_{\mathrm{I}}$]{0.3.7.1}, and with the notations of the proof of \sref{IV.8.2.9}, the homomorphisms $(\mathrm{A}_\lambda)_{\mathfrak{p}_\lambda}\to(\mathrm{A}_\mu)_{\mathfrak{p}_\mu}$ are bijective for a system $(\mathfrak{p}_\lambda)$ of prime ideals corresponding to a single point of $\mathrm{S}$; the assertion relative to $\sh{O}_\mathrm{S}$ is then deduced from \sref{IV.8.2.12}.
The last assertion then follows from the definition of $u_\alpha^*(\sh{F}_\alpha)$ \sref[$0_{\mathrm{I}}$]{0.4.3.1}.
\end{proof}

\begin{remark}[8.2.14]
\label{IV.8.2.14}
The results of \sref{IV.8.2.9} and \sref{IV.8.2.12} show that $\mathrm{S}$ is a projective limit of the projective system $(\mathrm{S}_\lambda,u_{\lambda\mu})$ in the category of \emph{all} ringed spaces (or of all locally ringed spaces).
Indeed, let $\mathrm{Y}$ be a ringed space, and consider a projective system of morphisms of ringed spaces $w_\lambda:\mathrm{Y}\to\mathrm{S}_\lambda$.
If one writes $w_\lambda=(\rho_\lambda,\omega_\lambda)$, the $\rho_\lambda$ form a projective system of continuous maps and, by \sref{IV.8.2.9}, their projective limit identifies with a continuous map $\rho:\mathrm{Y}\to\mathrm{S}$ such that $\rho_\lambda=\psi_\lambda\circ\rho$. On the other hand, the homomorphisms
\[
 \omega_\lambda^\sharp:\rho_\lambda^*(\sh{O}_{\mathrm{S}_\lambda})\longrightarrow\sh{O}_\mathrm{Y}
\]
form an inductive system of homomorphisms of sheaves of rings. Since
\[
 \rho_\lambda^*(\sh{O}_{\mathrm{S}_\lambda})
 =\rho^*\bigl(\psi_\lambda^*(\sh{O}_{\mathrm{S}_\lambda})\bigr)
\]
and the functor $\rho^*$ is exact, the inductive limit of the $\rho_\lambda^*(\sh{O}_{\mathrm{S}_\lambda})$ is $\rho^*(\sh{O}_\mathrm{S})$ by \sref{IV.8.2.12}. There is therefore a unique homomorphism
\[
 \omega^\sharp:\rho^*(\sh{O}_\mathrm{S})\longrightarrow\sh{O}_\mathrm{Y}
\]
such that $\omega_\lambda^\sharp=\omega^\sharp\circ\rho^*(\theta_\lambda^\sharp)$ for every $\lambda$, which proves the assertion.
\end{remark}

\subsection{Constructible subsets of a projective limit of preschemes}
\label{subsection:IV.8.3}

\begin{env}[8.3.1]
\label{IV.8.3.1}
In what follows throughout this paragraph, we suppose the conditions of \sref{IV.8.2.2} are satisfied, and we keep the notations.
\end{env}

\begin{theorem}[8.3.2]
\label{IV.8.3.2}
For every $\lambda$, let $\mathrm{E}_\lambda$, $\mathrm{F}_\lambda$ be two subsets of $\mathrm{S}_\lambda$.
Set
\begin{equation}
  \mathrm{E}=\bigcap_\lambda u_\lambda^{-1}(\mathrm{E}_\lambda),\qquad \mathrm{F}=\bigcup_\lambda u_\lambda^{-1}(\mathrm{F}_\lambda).
  \tag{8.3.2.1}
\label{IV.8.3.2.1}
\end{equation}
Suppose the following conditions are satisfied:
\begin{enumerate}[label=(\roman*)]
  \item For every $\lambda$, $\mathrm{E}_\lambda$ is pro-constructible and $\mathrm{F}_\lambda$ is ind-constructible \sref{IV.1.9.4}.
  \item For $\lambda\leqslant\mu$, we have
  \[
    \mathrm{E}_\mu\subset u_{\lambda\mu}^{-1}(\mathrm{E}_\lambda),\qquad \mathrm{F}_\mu\supset u_{\lambda\mu}^{-1}(\mathrm{F}_\lambda).
  \]
  \item There exists $\alpha\in\mathrm{L}$ such that $\mathrm{S}_\alpha$ is quasi-compact (which implies that $\mathrm{S}_\lambda$ is quasi-compact for $\lambda\geqslant\alpha$).
\end{enumerate}
Then the following properties are equivalent:
\begin{enumerate}[label=\emph{\alph*)}]
  \item $\mathrm{E}\subset\mathrm{F}$.
  \item There exists $\lambda\geqslant\alpha$ such that $u_\lambda^{-1}(\mathrm{E}_\lambda)\subset u_\lambda^{-1}(\mathrm{F}_\lambda)$ (and then we also have $u_\mu^{-1}(\mathrm{E}_\mu)\subset u_\mu^{-1}(\mathrm{F}_\mu)$ for
\oldpage[IV-3]{13}
$\mu\geqslant\lambda$).
  \item There exists $\lambda\geqslant\alpha$ such that $\mathrm{E}_\lambda\subset\mathrm{F}_\lambda$ (and then we also have $\mathrm{E}_\mu\subset\mathrm{F}_\mu$ for $\mu\geqslant\lambda$).
\end{enumerate}
The parenthetical remarks in \emph{b)} and \emph{c)} follow from (ii).
Set
\[
  \mathrm{G}_\lambda=\mathrm{E}_\lambda\cap(\mathrm{S}_\lambda-\mathrm{F}_\lambda),\qquad \mathrm{G}=\mathrm{E}\cap(\mathrm{S}-\mathrm{F}).
\]
Then $\mathrm{G}_\lambda$ is a pro-constructible subset of $\mathrm{S}_\lambda$ \sref{IV.1.9.5},~(ii), and by virtue of \sref{IV.8.3.2.1} and condition (ii) we have
\[
  \mathrm{G}_\mu\subset u_{\lambda\mu}^{-1}(\mathrm{G}_\lambda)\qquad\text{for }\lambda\leqslant\mu
\]
\[
  \mathrm{G}=\bigcap_\lambda u_\lambda^{-1}(\mathrm{G}_\lambda).
\]
We are thus reduced to proving the particular case of \sref{IV.8.3.2} corresponding to $\mathrm{F}_\lambda=\varnothing$ for every $\lambda$:
\end{theorem}

\begin{corollary}[8.3.3]
\label{IV.8.3.3}
For every $\lambda$, let $\mathrm{E}_\lambda$ be a pro-constructible subset of $\mathrm{S}_\lambda$ such that, for $\lambda\leqslant\mu$, $\mathrm{E}_\mu\subset u_{\lambda\mu}^{-1}(\mathrm{E}_\lambda)$.
Suppose that there exists $\alpha\in\mathrm{L}$ such that $\mathrm{S}_\alpha$ is quasi-compact.
Then the following conditions are equivalent:
\begin{enumerate}[label=\emph{\alph*)}]
  \item $\mathrm{E}=\bigcap_\lambda u_\lambda^{-1}(\mathrm{E}_\lambda)=\varnothing$.
  \item There exists $\lambda$ such that $u_\lambda^{-1}(\mathrm{E}_\lambda)=\varnothing$ (and then $u_\mu^{-1}(\mathrm{E}_\mu)=\varnothing$ for $\mu\geqslant\lambda$).
  \item There exists $\lambda$ such that $\mathrm{E}_\lambda=\varnothing$ (and then $\mathrm{E}_\mu=\varnothing$ for $\mu\geqslant\lambda$).
\end{enumerate}
\end{corollary}

\begin{proof}
It is clear that \emph{c)} implies \emph{a)}.
Let us prove that \emph{a)} implies \emph{b)}: since $\mathrm{S}_\lambda$ is quasi-compact, it is the same for $\mathrm{S}$ \sref{IV.8.2.2}; $\mathrm{E}_\lambda$ being pro-constructible, it is the same for $u_\lambda^{-1}(\mathrm{E}_\lambda)$ \sref{IV.1.9.5},~(vi); then the filtered decreasing family of pro-constructible subsets $u_\lambda^{-1}(\mathrm{E}_\lambda)$ has empty intersection, hence \sref{IV.1.9.9} one of them is empty.

Finally, let us show that \emph{b)} implies \emph{c)}.
Since $\mathrm{S}_\alpha$ is quasi-compact and $\mathrm{L}$ is filtered, we can replace $\mathrm{S}_\alpha$ by an open affine, hence suppose \sref{IV.8.2.2} that $\mathrm{S}$ and the $\mathrm{S}_\lambda$ for $\lambda\geqslant\alpha$ are affine; one has then \sref{IV.1.9.2.1}, for $\lambda\geqslant\alpha$,
\[
  u_\lambda(\mathrm{S})=\bigcap_{\mu\geqslant\lambda}u_{\lambda\mu}(\mathrm{S}_\mu)
\]
whence
\[
  \mathrm{E}_\lambda\cap u_\lambda(\mathrm{S})=\bigcap_{\mu\geqslant\lambda}(\mathrm{E}_\lambda\cap u_{\lambda\mu}(\mathrm{S}_\mu)).
\]
Now, since $u_\lambda$ and the $u_{\lambda\mu}$ are quasi-compact, $u_\lambda(\mathrm{S})$ and $u_{\lambda\mu}(\mathrm{S}_\mu)$ are pro-constructible in $\mathrm{S}_\lambda$ \sref{IV.1.9.5},~(viii), hence the sets $\mathrm{E}_\lambda\cap u_{\lambda\mu}(\mathrm{S}_\mu)$ for $\mu\geqslant\lambda$ form a filtered decreasing family of pro-constructible subsets of $\mathrm{S}_\lambda$.
Since $\mathrm{S}_\lambda$ is quasi-compact, the hypothesis \emph{b)} implies that the intersection of this family is empty, hence \sref{IV.1.9.9} one of the sets $\mathrm{E}_\lambda\cap u_{\lambda\mu}(\mathrm{S}_\mu)$ is empty, hence $\mathrm{E}_\mu\subset u_{\lambda\mu}^{-1}(\mathrm{E}_\lambda)$ is empty.
\end{proof}

\begin{corollary}[8.3.4]
\label{IV.8.3.4}
For every $\lambda$, let $\mathrm{F}_\lambda$ be an ind-constructible subset of $\mathrm{S}_\lambda$ such that for $\lambda\leqslant\mu$ one has $u_{\lambda\mu}^{-1}(\mathrm{F}_\lambda)\subset\mathrm{F}_\mu$.
Suppose that there exists $\alpha\in\mathrm{L}$ such that $\mathrm{S}_\alpha$ is quasi-compact.
Then the following conditions are equivalent:
\begin{enumerate}[label=\emph{\alph*)}]
  \item The set $\mathrm{F}=\bigcup_\lambda u_\lambda^{-1}(\mathrm{F}_\lambda)$ is equal to $\mathrm{S}$.
  \item There exists $\lambda$ such that $u_\lambda^{-1}(\mathrm{F}_\lambda)=\mathrm{S}$ (and then $u_\mu^{-1}(\mathrm{F}_\mu)=\mathrm{S}$ for $\mu\geqslant\lambda$).
\oldpage[IV-3]{14}
  \item There exists $\lambda$ such that $\mathrm{F}_\lambda=\mathrm{S}_\lambda$ (and then $\mathrm{F}_\mu=\mathrm{S}_\mu$ for $\mu\geqslant\lambda$).
\end{enumerate}
This follows immediately from \sref{IV.8.3.3} by passing to complements.
\end{corollary}

\begin{corollary}[8.3.5]
\label{IV.8.3.5}
For every $\lambda$, let $\mathrm{E}_\lambda$, $\mathrm{F}_\lambda$ be two constructible subsets of $\mathrm{S}_\lambda$ such that, for $\mu\geqslant\lambda$, one has $\mathrm{E}_\mu\subset u_{\lambda\mu}^{-1}(\mathrm{E}_\lambda)$ and $\mathrm{F}_\mu\supset u_{\lambda\mu}^{-1}(\mathrm{F}_\lambda)$.
Suppose that there exists $\alpha$ such that $\mathrm{S}_\alpha$ is quasi-compact.
Then, for $\mathrm{E}\subset\mathrm{F}$, it is necessary and sufficient that there exist $\lambda$ such that $\mathrm{E}_\lambda\subset\mathrm{F}_\lambda$, in which case one also has $\mathrm{E}_\mu\subset\mathrm{F}_\mu$ for $\mu\geqslant\lambda$.\footnote{The printed text also includes three parenthetical equality assertions; these do not follow from the stated one-sided transition hypotheses.}

This is a particular case of \sref{IV.8.3.2}.
\end{corollary}

In particular:

\begin{corollary}[8.3.6]
\label{IV.8.3.6}
Suppose that there exists $\alpha$ such that $\mathrm{S}_\alpha$ is quasi-compact.
For $\mathrm{S}=\varnothing$, it is necessary and sufficient that there exist $\lambda$ such that $\mathrm{S}_\lambda=\varnothing$.
\end{corollary}

\begin{corollary}[8.3.7]
\label{IV.8.3.7}
For every $\lambda$, we have
\begin{equation}
  u_\lambda(\mathrm{S})=\bigcap_{\mu\geqslant\lambda}u_{\lambda\mu}(\mathrm{S}_\mu).
  \tag{8.3.7.1}
\label{IV.8.3.7.1}
\end{equation}
\end{corollary}

\begin{proof}
It is clear that the first member of \sref{IV.8.3.7.1} is contained in the second.
Let $s$ be a point of $\mathrm{S}_\lambda$ and set $\mathrm{X}_\lambda=\operatorname{Spec}(\boldsymbol{k}(s))$; consider the projective system $(\mathrm{X}_\mu,v_{\mu\nu})$ where $\mathrm{X}_\mu=\mathrm{X}_\lambda\times_{\mathrm{S}_\lambda}\mathrm{S}_\mu$ and $v_{\mu\nu}=1\times u_{\mu\nu}$ for $\lambda\leqslant\mu\leqslant\nu$; its projective limit is $\mathrm{X}=\mathrm{X}_\lambda\times_{\mathrm{S}_\lambda}\mathrm{S}$, and $v_\lambda=1\times u_\lambda$ is the canonical map $\mathrm{X}\to\mathrm{X}_\lambda$ \sref{IV.8.2.5}.
If $s\in\bigcap_{\mu\geqslant\lambda}u_{\lambda\mu}(\mathrm{S}_\mu)$, this implies that $\mathrm{X}_\mu\neq\varnothing$ for every $\mu\geqslant\lambda$ \sref[I]{I.3.4.8}; it follows then from \sref{IV.8.3.6} that $\mathrm{X}\neq\varnothing$, hence $s\in u_\lambda(\mathrm{S})$ by \sref[I]{I.3.4.8}.
\end{proof}

\begin{proposition}[8.3.8]
\label{IV.8.3.8}
\begin{enumerate}[label=(\roman*)]
  \item For the morphism $u_\lambda:\mathrm{S}\to\mathrm{S}_\lambda$ to be dominant (resp.\ surjective), it is necessary and sufficient that for $\mu\geqslant\lambda$ the morphism $u_{\lambda\mu}:\mathrm{S}_\mu\to\mathrm{S}_\lambda$ be dominant (resp.\ surjective).
  \item If, for some index $\lambda$, the morphisms $u_{\lambda\mu}:\mathrm{S}_\mu\to\mathrm{S}_\lambda$ are flat (resp.\ faithfully flat) for every $\mu\geqslant\lambda$, then the morphism $u_\lambda:\mathrm{S}\to\mathrm{S}_\lambda$ is flat (resp.\ faithfully flat).
  \item Suppose that the morphisms $u_\mu:\mathrm{S}\to\mathrm{S}_\mu$ are surjective for $\mu\geqslant\lambda$.
For $u_\lambda$ to be an open morphism (resp.\ universally open), it is necessary and sufficient that, for every $\mu\geqslant\lambda$, $u_{\lambda\mu}$ be an open morphism (resp.\ universally open).
\end{enumerate}
\end{proposition}

\begin{proof}
\begin{enumerate}[label=(\roman*)]
  \item Since $u_\lambda(\mathrm{S})\subset u_{\lambda\mu}(\mathrm{S}_\mu)$ for every $\mu\geqslant\lambda$, the necessity of the conditions is trivial, and it follows immediately from \sref{IV.8.3.7.1} that if the $u_{\lambda\mu}$ are surjective, then so is $u_\lambda$.
Suppose now that the $u_{\lambda\mu}$ are dominant for $\mu\geqslant\lambda$, and consider in $\mathrm{S}_\lambda$ an open quasi-compact non-empty subset $\mathrm{U}_\lambda$; then the $\mathrm{U}_\mu=u_{\lambda\mu}^{-1}(\mathrm{U}_\lambda)$ for $\mu\geqslant\lambda$ form a projective system whose projective limit is $\mathrm{U}=u_\lambda^{-1}(\mathrm{U}_\lambda)$ \sref{IV.8.2.5}.
By hypothesis the $\mathrm{U}_\mu$ are all non-empty, hence it is the same for $\mathrm{U}$ by \sref{IV.8.3.6}, which proves that $u_\lambda$ is dominant.

  \item By virtue of (i), it suffices to consider the case where the $u_{\lambda\mu}$ are flat; one can then restrict to the case where $\mathrm{S}_\lambda$ is affine, hence also the $\mathrm{S}_\mu$ for $\mu\geqslant\lambda$ and $\mathrm{S}$, and the assertion follows from \sref{IV.2.1.2} and \sref[0]{0.6.2.3}.

  \item By virtue of \sref{IV.8.2.5} and of \sref[I]{I.3.5.2}, it suffices to treat the case of open morphisms.
Since $u_\lambda=u_{\lambda\mu}\circ u_\mu$ and $u_\mu$ is surjective, we know that if $u_\lambda$ is open then so is $u_{\lambda\mu}$ for $\mu\geqslant\lambda$ (Bourbaki, \emph{Top.\ g\'{e}n.}, chap.~I\textsuperscript{er}, 3\textsuperscript{e}~ed., \textsection\,5, n\textsuperscript{o}\,1, prop.~1).
\oldpage[IV-3]{15}
Conversely, to show that $u_\lambda$ is open when all the $u_{\lambda\mu}$ are open for $\mu\geqslant\lambda$, it suffices to see that for every quasi-compact open subset $\mathrm{V}$ of $\mathrm{S}$, $u_\lambda(\mathrm{V})$ is open in $\mathrm{S}_\lambda$; but there exists then $\mu\geqslant\lambda$ such that $\mathrm{V}=u_\mu^{-1}(\mathrm{V}_\mu)$, where $\mathrm{V}_\mu$ is open in $\mathrm{S}_\mu$ \sref{IV.8.2.11}; since $u_\mu$ is surjective, one has $\mathrm{V}_\mu=u_\mu(\mathrm{V})$ and $u_\lambda(\mathrm{V})=u_{\lambda\mu}(u_\mu(\mathrm{V}))$ is therefore open by hypothesis.
\end{enumerate}

One notes that it may happen that all the $u_{\lambda\mu}$ are open without $u_\lambda$ being open, when the $u_\mu$ are not surjective.
Indeed, one has an example by considering an integral ring $\mathrm{A}$ which is not a field, and its field of fractions $\mathrm{K}$, which is the inductive limit of the $\mathrm{A}_f$, where $f$ ranges over $\mathrm{A}-\{0\}$; setting $\mathrm{S}_1=\operatorname{Spec}(\mathrm{A})$, $\mathrm{S}_f=\operatorname{Spec}(\mathrm{A}_f)$, and $\mathrm{S}=\varprojlim\mathrm{S}_f=\operatorname{Spec}(\mathrm{K})$, the morphism $\mathrm{S}\to\mathrm{S}_1$ is not open, even though the morphisms $\mathrm{S}_f\to\mathrm{S}_1$ are open.
\end{proof}

\begin{env}[8.3.9]
\label{IV.8.3.9}
For every prescheme $\mathrm{X}$, we will denote as usual by $\mathfrak{P}(\mathrm{X})$ the set of all subsets of the underlying set of $\mathrm{X}$, by $\mathfrak{C}(\mathrm{X})$ (resp.\ $\mathfrak{D}_c(\mathrm{X})$, $\mathfrak{F}_c(\mathrm{X})$, $\mathfrak{LF}_c(\mathrm{X})$) the set of constructible (resp.\ constructible and open, constructible and closed, and constructible and locally closed) subsets of $\mathrm{X}$.
It is clear that $(\mathfrak{P}(\mathrm{S}_\lambda),u_{\lambda\mu}^{-1})$ is an \emph{inductive system of sets} and that the maps $u_\lambda^{-1}:\mathfrak{P}(\mathrm{S}_\lambda)\to\mathfrak{P}(\mathrm{S})$ form an inductive system of maps, whence, by passing to the inductive limit, a canonical map
\begin{equation}
  u_\mathfrak{P}:\varinjlim\mathfrak{P}(\mathrm{S}_\lambda)\to\mathfrak{P}(\mathrm{S}).
  \tag{8.3.9.1}
\label{IV.8.3.9.1}
\end{equation}

Furthermore, it follows from \sref{IV.1.8.2} that $u_{\lambda\mu}^{-1}$ sends $\mathfrak{C}(\mathrm{S}_\lambda)$ (resp.\ $\mathfrak{D}_c(\mathrm{S}_\lambda)$, $\mathfrak{F}_c(\mathrm{S}_\lambda)$, $\mathfrak{LF}_c(\mathrm{S}_\lambda)$) into $\mathfrak{C}(\mathrm{S}_\mu)$ (resp.\ $\mathfrak{D}_c(\mathrm{S}_\mu)$, $\mathfrak{F}_c(\mathrm{S}_\mu)$, $\mathfrak{LF}_c(\mathrm{S}_\mu)$) and that $u_\lambda^{-1}$ sends $\mathfrak{C}(\mathrm{S}_\lambda)$ (resp.\ $\mathfrak{D}_c(\mathrm{S}_\lambda)$, $\mathfrak{F}_c(\mathrm{S}_\lambda)$, $\mathfrak{LF}_c(\mathrm{S}_\lambda)$) into $\mathfrak{C}(\mathrm{S})$ (resp.\ $\mathfrak{D}_c(\mathrm{S})$, $\mathfrak{F}_c(\mathrm{S})$, $\mathfrak{LF}_c(\mathrm{S})$).
We thus have, by restriction of \sref{IV.8.3.9.1}, canonical maps
\begin{align}
  &\varinjlim\mathfrak{C}(\mathrm{S}_\lambda)\to\mathfrak{C}(\mathrm{S}) \tag{8.3.9.2}\label{IV.8.3.9.2}\\
  &\varinjlim\mathfrak{D}_c(\mathrm{S}_\lambda)\to\mathfrak{D}_c(\mathrm{S}) \tag{8.3.9.3}\label{IV.8.3.9.3}\\
  &\varinjlim\mathfrak{F}_c(\mathrm{S}_\lambda)\to\mathfrak{F}_c(\mathrm{S}) \tag{8.3.9.4}\label{IV.8.3.9.4}\\
  &\varinjlim\mathfrak{LF}_c(\mathrm{S}_\lambda)\to\mathfrak{LF}_c(\mathrm{S}). \tag{8.3.9.5}
\label{IV.8.3.9.5}
\end{align}
\end{env}

\begin{env}[8.3.10]
\label{IV.8.3.10}
Let $g_\alpha:\mathrm{X}_\alpha\to\mathrm{S}_\alpha$ be a morphism; with the notations of \sref{IV.8.2.5} one has, on the one hand, a canonical map $v_\mathfrak{P}:\varinjlim\mathfrak{P}(\mathrm{X}_\lambda)\to\mathfrak{P}(\mathrm{X})$; on the other hand, one has projection morphisms $g_\lambda:\mathrm{X}_\lambda\to\mathrm{S}_\lambda$ for every $\lambda\geqslant\alpha$ and a projection morphism $g:\mathrm{X}\to\mathrm{S}$.
It is clear that the $g_\lambda^{-1}:\mathfrak{P}(\mathrm{S}_\lambda)\to\mathfrak{P}(\mathrm{X}_\lambda)$ form an inductive system of maps, and that the diagrams
\[
\xymatrix{
  \mathfrak{P}(\mathrm{S}_\lambda) \ar[r]^{u_\lambda^{-1}} \ar[d]_{g_\lambda^{-1}} & \mathfrak{P}(\mathrm{S}) \ar[d]^{g^{-1}} \\
  \mathfrak{P}(\mathrm{X}_\lambda) \ar[r]_{v_\lambda^{-1}} & \mathfrak{P}(\mathrm{X})
}
\]
\oldpage[IV-3]{16}
are commutative; one deduces from this by passing to the inductive limit a commutative diagram
\begin{equation}
\xymatrix{
  \varinjlim\mathfrak{P}(\mathrm{S}_\lambda) \ar[r]^{u_\mathfrak{P}} \ar[d]_{\varinjlim g_\lambda^{-1}} & \mathfrak{P}(\mathrm{S}) \ar[d]^{g^{-1}} \\
  \varinjlim\mathfrak{P}(\mathrm{X}_\lambda) \ar[r]_{v_\mathfrak{P}} & \mathfrak{P}(\mathrm{X})
}
\tag{8.3.10.1}
\label{IV.8.3.10.1}
\end{equation}
and it follows from \sref{IV.1.8.2} that one has analogous diagrams replacing $\mathfrak{P}$ by $\mathfrak{C}$, $\mathfrak{D}_c$, $\mathfrak{F}_c$, or $\mathfrak{LF}_c$.
\end{env}

It follows from \sref{IV.8.3.5} that under the hypothesis that for some $\alpha\in\mathrm{L}$, $\mathrm{S}_\alpha$ is \emph{quasi-compact}, the canonical map \sref{IV.8.3.9.2} is \emph{injective}.
Furthermore:

\begin{theorem}[8.3.11]
\label{IV.8.3.11}
Suppose that there exists $\alpha\in\mathrm{L}$ such that $\mathrm{S}_\alpha$ is quasi-compact and quasi-separated.
Then the canonical maps \sref{IV.8.3.9.2}, \sref{IV.8.3.9.3}, \sref{IV.8.3.9.4}, and \sref{IV.8.3.9.5} are \emph{bijective}.
\end{theorem}

\begin{proof}
By virtue of the preceding remark, it remains to prove that these maps are \emph{surjective}; as every constructible subset of $\mathrm{S}$ is a finite union of sets of the form $\mathrm{U}\cap\complement\mathrm{V}$, where $\mathrm{U}$ and $\mathrm{V}$ are open and constructible, it will suffice to prove that \sref{IV.8.3.9.3} is surjective for it to be the same for \sref{IV.8.3.9.2} (and also for \sref{IV.8.3.9.4}, by passing to complements).
Now, since the morphisms $\mathrm{S}_\lambda\to\mathrm{S}_\alpha$ and $\mathrm{S}\to\mathrm{S}_\alpha$ are affine, hence separated, the $\mathrm{S}_\lambda$ for $\lambda\geqslant\alpha$ and $\mathrm{S}$ are quasi-compact and quasi-separated \sref{IV.1.2.2}, and one knows that the constructible open subsets of such a prescheme are none other than the quasi-compact open subsets \sref{IV.1.8.1}.
The conclusion then follows from \sref{IV.8.2.11} except for the map \sref{IV.8.3.9.5}.
To prove that this last map is surjective, consider a locally closed and constructible subset $\mathrm{Z}$ of $\mathrm{S}$; $\mathrm{Z}$ is therefore quasi-compact ($\mathbf{0}_{\mathrm{III}}$,~9.1.4).
Since every point $x\in\mathrm{Z}$ admits by hypothesis an open quasi-compact neighbourhood $\mathrm{V}_x$ in $\mathrm{S}$ such that $\mathrm{Z}\cap\mathrm{V}_x$ is closed in $\mathrm{V}_x$, we can cover $\mathrm{Z}$ by a finite number of the $\mathrm{V}_x$; in other words, there is a quasi-compact open subset $\mathrm{U}$ containing $\mathrm{Z}$ such that $\mathrm{Z}$ is closed in $\mathrm{U}$; since $\mathrm{Z}$ is constructible in $\mathrm{S}$, it is also constructible in $\mathrm{U}$ ($\mathbf{0}_{\mathrm{III}}$,~9.1.8).
We know \sref{IV.8.2.11} that there exist $\lambda$ and a quasi-compact open subset $\mathrm{U}_\lambda$ of $\mathrm{S}_\lambda$ such that $\mathrm{U}=u_\lambda^{-1}(\mathrm{U}_\lambda)$.
Applying to $\mathrm{U}$ (which is the projective limit of the $\mathrm{U}_\mu=u_{\lambda\mu}^{-1}(\mathrm{U}_\lambda)$ for $\mu\geqslant\lambda$) the fact that the map \sref{IV.8.3.9.4} is surjective, one sees that there exist $\mu\geqslant\lambda$ and a closed constructible subset $\mathrm{Z}_\mu$ in $\mathrm{U}_\mu$ such that $\mathrm{Z}=u_\mu^{-1}(\mathrm{Z}_\mu)$.
But since the canonical immersion $\mathrm{U}_\mu\to\mathrm{S}_\mu$ is quasi-compact by hypothesis \sref{IV.1.2.7}, it is of finite presentation \sref{IV.1.6.2},~(i), and $\mathrm{Z}_\mu$ is also a constructible subset \emph{of} $\mathrm{S}_\mu$ by virtue of \sref{IV.1.8.4} and \sref{IV.1.8.1}; since $\mathrm{Z}_\mu$ is evidently locally closed in $\mathrm{S}_\mu$, this completes the proof.
\end{proof}

\begin{corollary}[8.3.12]
\label{IV.8.3.12}
Suppose that there exists $\alpha$ such that $\mathrm{S}_\alpha$ is quasi-compact, and let, for every $\lambda$, $\mathrm{Z}_\lambda$ be a constructible subset of $\mathrm{S}_\lambda$ such that $\mathrm{Z}_\mu=u_{\lambda\mu}^{-1}(\mathrm{Z}_\lambda)$ for $\mu\geqslant\lambda$.
If $\mathrm{Z}=u_\lambda^{-1}(\mathrm{Z}_\lambda)$,
\oldpage[IV-3]{17}
then, for $\mathrm{Z}$ to be open (resp.\ closed, resp.\ locally closed) in $\mathrm{S}$, it is necessary and sufficient that there exist $\lambda\geqslant\alpha$ such that $\mathrm{Z}_\lambda$ be so in $\mathrm{S}_\lambda$.
\end{corollary}

\begin{proof}
Cover $\mathrm{S}_\alpha$ by a finite number of open affine subsets $\mathrm{U}_\alpha^{(i)}$; then the $\mathrm{U}^{(i)}=u_\alpha^{-1}(\mathrm{U}_\alpha^{(i)})$ form an open affine covering of $\mathrm{S}$, and for $\mathrm{Z}$ to be open (resp.\ closed, resp.\ locally closed) in $\mathrm{S}$, it is necessary and sufficient that each of the $\mathrm{Z}\cap\mathrm{U}^{(i)}$ be so in $\mathrm{U}^{(i)}$.
Since $\mathrm{L}$ is filtered and each $\mathrm{Z}\cap\mathrm{U}^{(i)}$ is constructible in $\mathrm{U}^{(i)}$ ($\mathbf{0}_{\mathrm{III}}$,~9.1.8), one can restrict to proving the corollary when $\mathrm{S}_\alpha$ is affine, hence quasi-compact and quasi-separated; but then it follows from \sref{IV.8.3.11}.
\end{proof}

\begin{proposition}[8.3.13]
\label{IV.8.3.13}
Suppose that the morphisms $u_{\lambda\mu}:\mathrm{S}_\mu\to\mathrm{S}_\lambda$ are flat for $\lambda\leqslant\mu$, and that there exists $\alpha$ such that $\mathrm{S}_\alpha$ is quasi-compact.
For every $\lambda$, let $\mathrm{Z}'_\lambda$, $\mathrm{Z}''_\lambda$ be two pro-constructible subsets of $\mathrm{S}_\lambda$, such that $\mathrm{Z}'_\mu=u_{\lambda\mu}^{-1}(\mathrm{Z}'_\lambda)$ and $\mathrm{Z}''_\mu=u_{\lambda\mu}^{-1}(\mathrm{Z}''_\lambda)$ for $\mu\geqslant\lambda$; suppose further that $\overline{\mathrm{Z}'_\alpha}$ is constructible in $\mathrm{S}_\alpha$.
Set $\mathrm{Z}'=u_\lambda^{-1}(\mathrm{Z}'_\lambda)$, $\mathrm{Z}''=u_\lambda^{-1}(\mathrm{Z}''_\lambda)$; for $\mathrm{Z}''\subset\overline{\mathrm{Z}'}$, it is necessary and sufficient that there exist $\lambda\geqslant\alpha$ such that $\mathrm{Z}''_\lambda\subset\overline{\mathrm{Z}'_\lambda}$.
\end{proposition}

\begin{proof}
Indeed, one knows that $u_\lambda:\mathrm{S}\to\mathrm{S}_\lambda$ is also a flat morphism for every $\lambda$ \sref{IV.8.3.8}; as $\mathrm{Z}'_\lambda$ is pro-constructible, the closure of $\mathrm{Z}'_\mu$ for $\mu\geqslant\lambda$ (resp.\ of $\mathrm{Z}'$) in $\mathrm{S}_\mu$ (resp.\ $\mathrm{S}$) is equal to $u_{\lambda\mu}^{-1}(\overline{\mathrm{Z}'_\lambda})$ (resp.\ $u_\lambda^{-1}(\overline{\mathrm{Z}'_\lambda})$) \sref{IV.2.3.10}.
Since the $u_{\lambda\mu}^{-1}(\overline{\mathrm{Z}'_\lambda})$ and $u_\lambda^{-1}(\overline{\mathrm{Z}'_\lambda})$ are constructible \sref{IV.1.8.2}, the conclusion follows from \sref{IV.8.3.2}.
\end{proof}

\subsection{Criteria of irreducibility and connectedness for projective limits of preschemes}
\label{subsection:IV.8.4}

\begin{proposition}[8.4.1]
\label{IV.8.4.1}
Suppose that there exists an index $\alpha$ such that $\mathrm{S}_\alpha$ is quasi-compact.
\begin{enumerate}[label=(\roman*)]
  \item If $\mathrm{S}$ is not irreducible and if in addition the underlying space of $\mathrm{S}$ is Noetherian and $\mathrm{S}_\alpha$ is quasi-separated, there exists $\lambda\geqslant\alpha$ such that, for $\mu\geqslant\lambda$, $\mathrm{S}_\mu$ is not irreducible.
  \item If $\mathrm{S}$ is not connected, there exists $\lambda\geqslant\alpha$ such that, for $\mu\geqslant\lambda$, $\mathrm{S}_\mu$ is not connected.
\end{enumerate}
\end{proposition}

\begin{proof}
Suppose that $\mathrm{S}$ is a union of two distinct closed subsets $\mathrm{S}'$, $\mathrm{S}''$ of $\mathrm{S}$ (resp.\ disjoint non-empty closed subsets).
In case (i), $\mathrm{S}'$ and $\mathrm{S}''$ are constructible since the space $\mathrm{S}$ is Noetherian.
By virtue of \sref{IV.8.3.11}, there exist $\lambda\geqslant\alpha$ and two constructible closed subsets $\mathrm{S}'_\lambda$, $\mathrm{S}''_\lambda$ of $\mathrm{S}_\lambda$ such that $\mathrm{S}'=u_\lambda^{-1}(\mathrm{S}'_\lambda)$, $\mathrm{S}''=u_\lambda^{-1}(\mathrm{S}''_\lambda)$; since $\mathrm{S}=\mathrm{S}'\cup\mathrm{S}''$, it follows from \sref{IV.8.3.11} that one can suppose $\mathrm{S}_\lambda=\mathrm{S}'_\lambda\cup\mathrm{S}''_\lambda$; since $\mathrm{S}'_\lambda$ and $\mathrm{S}''_\lambda$ are distinct from $\mathrm{S}_\lambda$, this proves that $\mathrm{S}_\lambda$ is not irreducible.

In case (ii), $\mathrm{S}'$ and $\mathrm{S}''$ are open and quasi-compact, hence, by virtue of \sref{IV.8.2.11}, there exist $\lambda\geqslant\alpha$ and two open quasi-compact subsets $\mathrm{S}'_\lambda$, $\mathrm{S}''_\lambda$ of $\mathrm{S}_\lambda$ such that $\mathrm{S}'=u_\lambda^{-1}(\mathrm{S}'_\lambda)$, $\mathrm{S}''=u_\lambda^{-1}(\mathrm{S}''_\lambda)$.
Moreover, since $\mathrm{S}'$ and $\mathrm{S}''$ are open and closed in $\mathrm{S}$, they are both pro-constructible and ind-constructible \sref{IV.1.9.6}, hence constructible \sref{IV.1.9.11}, and it follows from \sref{IV.8.3.5} that one can suppose $\lambda$ taken such that $\mathrm{S}_\lambda=\mathrm{S}'_\lambda\cup\mathrm{S}''_\lambda$ and $\mathrm{S}'_\lambda\cap\mathrm{S}''_\lambda=\varnothing$, which shows that $\mathrm{S}_\lambda$ is not connected.
\end{proof}

\begin{proposition}[8.4.2]
\label{IV.8.4.2}
Suppose that the underlying space of $\mathrm{S}$ is Noetherian and that one of the two following conditions is satisfied:
\begin{enumerate}[label=\emph{\alph*)}]
  \item For $\lambda\leqslant\mu$, $u_{\lambda\mu}:\mathrm{S}_\mu\to\mathrm{S}_\lambda$ is dominant, and there exists $\alpha$ such that $\mathrm{S}_\alpha$ is quasi-compact.
\oldpage[IV-3]{18}
  \item There exists $\alpha$ such that the underlying space of $\mathrm{S}_\alpha$ is Noetherian, and for $\mu\geqslant\lambda$, $u_{\lambda\mu}$ is a homeomorphism of $\mathrm{S}_\mu$ onto a subspace of $\mathrm{S}_\lambda$.
\end{enumerate}
Under these conditions, there exists $\lambda$ such that, for every $\mu\geqslant\lambda$:
\begin{enumerate}[label=(\roman*)]
  \item For every irreducible component $\mathrm{Y}_i$ of $\mathrm{S}$ ($1\leqslant i\leqslant m$), $\overline{u_\mu(\mathrm{Y}_i)}$ is an irreducible component of $\mathrm{S}_\mu$, and the map $\mathrm{Y}_i\mapsto\overline{u_\mu(\mathrm{Y}_i)}$ is a bijection of the set of irreducible components of $\mathrm{S}$ onto the set of irreducible components of $\mathrm{S}_\mu$.
  \item For every connected component $\mathrm{C}_j$ of $\mathrm{S}$ ($1\leqslant j\leqslant n$), $\overline{u_\mu(\mathrm{C}_j)}$ is a connected component of $\mathrm{S}_\mu$, and the map $\mathrm{C}_j\mapsto\overline{u_\mu(\mathrm{C}_j)}$ is a bijection of the set of connected components of $\mathrm{S}$ onto the set of connected components of $\mathrm{S}_\mu$.
\end{enumerate}
\end{proposition}

\begin{proof}
Let us first establish the

\begin{lemma}[8.4.2.1]
\label{IV.8.4.2.1}
Under condition \emph{a)} or \emph{b)} of \sref{IV.8.4.2}, there exists $\lambda$ such that, for $\mu\geqslant\lambda$, $u_\mu:\mathrm{S}\to\mathrm{S}_\mu$ is dominant.
\end{lemma}

In case \emph{a)}, this has already been proved without supposing the space $\mathrm{S}$ Noetherian \sref{IV.8.3.8},~(i).
In case \emph{b)}, set $\mathrm{Z}_\alpha=\overline{u_\alpha(\mathrm{S})}$; as a closed subset of the Noetherian space $\mathrm{S}_\alpha$, $\mathrm{Z}_\alpha$ is constructible, and since $u_\alpha^{-1}(\mathrm{Z}_\alpha)=\mathrm{S}$, it follows from \sref{IV.8.3.5} that there exists $\lambda\geqslant\alpha$ such that, for $\mu\geqslant\lambda$, $\mathrm{Z}_\mu=u_{\alpha\mu}^{-1}(\mathrm{Z}_\alpha)=\mathrm{S}_\mu$.
But since $u_{\alpha\mu}$ is a homeomorphism of $\mathrm{S}_\mu$ onto a subspace of $\mathrm{S}_\alpha$, and since the composite map $\mathrm{S}\to\mathrm{Z}_\mu\to\mathrm{Z}_\alpha$ is dominant, the same is true of $\mathrm{S}\to\mathrm{Z}_\mu=\mathrm{S}_\mu$.

This lemma being proved, one can suppose that for every $\lambda$, $u_\lambda$ is a dominant morphism.
\begin{enumerate}[label=(\roman*)]
  \item Each of the $\mathrm{S}_\lambda$ is a union of the $\overline{u_\lambda(\mathrm{Y}_i)}$, which are irreducible.
On the other hand, if $\mathrm{U}_i$ is the open subset of $\mathrm{S}$ complementary to the union of the $\mathrm{Y}_j$ of index $j\neq i$ ($1\leqslant i\leqslant m$), the $\mathrm{U}_i$ are pairwise disjoint and $\mathrm{Y}_i=\overline{\mathrm{U}_i}$ \sref[0]{0.2.1.6}.
Since the underlying space of $\mathrm{S}$ is Noetherian, the $\mathrm{U}_i$ are quasi-compact, hence there exist an index $\lambda$ and open subsets $\mathrm{U}_{i\lambda}$ of $\mathrm{S}_\lambda$ such that $\mathrm{U}_i=u_\lambda^{-1}(\mathrm{U}_{i\lambda})$ for $1\leqslant i\leqslant m$ \sref{IV.8.2.11}.
From this one concludes that if one sets $\mathrm{U}_{i\mu}=u_{\lambda\mu}^{-1}(\mathrm{U}_{i\lambda})$ for $\lambda\leqslant\mu$, the $\mathrm{U}_{i\mu}$ are pairwise disjoint, because the $\mathrm{U}_i=u_\mu^{-1}(\mathrm{U}_{i\mu})$ are so and $u_\mu$ is dominant.
Consequently, none of the closures $\overline{\mathrm{U}_{i\mu}}$ is contained in another, and $u_\mu(\mathrm{U}_i)$ is dense in $\mathrm{U}_{i\mu}$ since $u_\mu$ is dominant; hence $\overline{\mathrm{U}_{i\mu}}=\overline{u_\mu(\mathrm{Y}_i)}$, which proves that the $\overline{\mathrm{U}_{i\mu}}$ are the irreducible components of $\mathrm{S}_\mu$ \sref[0]{0.2.1.7} and completes the proof.

  \item Since the underlying space $\mathrm{S}$ is Noetherian, the $\mathrm{C}_j$ are open and closed in $\mathrm{S}$ and quasi-compact; the same reasoning as in (i) shows that there exist $\lambda$ and open subsets $\mathrm{V}_{j\lambda}$ of $\mathrm{S}_\lambda$ such that $\mathrm{C}_j=u_\lambda^{-1}(\mathrm{V}_{j\lambda})$ for $1\leqslant j\leqslant n$.
One sees also, as in (i), that if one sets $\mathrm{V}_{j\mu}=u_{\lambda\mu}^{-1}(\mathrm{V}_{j\lambda})$ for $\mu\geqslant\lambda$, the $\mathrm{V}_{j\mu}$ are pairwise disjoint, and $u_\mu(\mathrm{C}_j)$ is dense in $\mathrm{V}_{j\mu}$; this implies that $\mathrm{V}_{j\mu}$ is connected.
Furthermore, it follows from \sref{IV.8.3.4} that for $\mu$ sufficiently large, the union of the $\mathrm{V}_{j\mu}$ is $\mathrm{S}_\mu$, since every open subset of a prescheme is ind-constructible \sref{IV.1.9.6}.
The $\mathrm{V}_{j\mu}$ are thus the connected components of $\mathrm{S}_\mu$, and this completes the proof.
\end{enumerate}
\end{proof}

We note that if the morphisms $u_{\lambda\mu}$ are \emph{immersions}, they satisfy in particular condition \emph{b)} of \sref{IV.8.4.2}.

\oldpage[IV-3]{19}
\begin{corollary}[8.4.3]
\label{IV.8.4.3}
Suppose one of the conditions \emph{a)}, \emph{b)} of \sref{IV.8.4.2} is satisfied, the underlying space of $\mathrm{S}$ being Noetherian; then, for $\mathrm{S}$ to be irreducible, it is necessary and sufficient that there exist $\lambda$ such that $\mathrm{S}_\mu$ be so for every $\mu\geqslant\lambda$.
\end{corollary}

\begin{proposition}[8.4.4]
\label{IV.8.4.4}
Suppose that there exists $\alpha$ such that $\mathrm{S}_\alpha$ is quasi-compact and that, for $\lambda\leqslant\mu$, $u_{\lambda\mu}:\mathrm{S}_\mu\to\mathrm{S}_\lambda$ is dominant.
Then, for $\mathrm{S}$ to be connected, it is necessary and sufficient that there exist $\lambda$ such that $\mathrm{S}_\mu$ be so for every $\mu\geqslant\lambda$.
\end{proposition}

\begin{proof}
The condition is sufficient by virtue of \sref{IV.8.4.1}; on the other hand, we have seen \sref{IV.8.3.8},~(i) that $u_\lambda:\mathrm{S}\to\mathrm{S}_\lambda$ is dominant for $\lambda$ sufficiently large, hence, if $\mathrm{S}$ is connected, it is the same for $\mathrm{S}_\lambda$, since $u_\lambda(\mathrm{S})$ is dense in $\mathrm{S}_\lambda$ and is connected.
\end{proof}

\begin{corollary}[8.4.5]
\label{IV.8.4.5}
Let $k$ be a field, $\mathrm{X}$ a quasi-compact $k$-prescheme.
For $\mathrm{X}$ to be geometrically connected \sref{IV.4.5.2}, it is necessary and sufficient that, for every finite and separable extension $\mathrm{K}$ of $k$, $\mathrm{X}\otimes_k\mathrm{K}$ be connected.
\end{corollary}

\begin{proof}
The condition is trivially necessary.
To see that it is sufficient, we must prove that if $\Omega$ is an algebraic closure of $k$, $\mathrm{X}\otimes_k\Omega$ is connected \sref{IV.4.5.1}.
Now, $\Omega$ is a filtered inductive limit of the finite subextensions $k'$ of $k$, and for $k\subset k'\subset k''\subset\Omega$, the morphism $\mathrm{X}\otimes_k k''\to\mathrm{X}\otimes_k k'$ is surjective.
We are thus reduced, by virtue of \sref{IV.8.4.4}, to proving that $\mathrm{X}\otimes_k k'$ is connected for every finite extension $k'$ of $k$.
But if $\mathrm{K}$ is the largest separable extension contained in $k'$, the morphism $\mathrm{X}\otimes_k k'\to\mathrm{X}\otimes_k\mathrm{K}$ is finite, surjective, and radicial, hence \sref{IV.2.4.5} a homeomorphism, and since $\mathrm{X}\otimes_k\mathrm{K}$ is connected by hypothesis, so is $\mathrm{X}\otimes_k k'$.
\end{proof}

\begin{remark}[8.4.6]
\label{IV.8.4.6}
\begin{enumerate}[label=(\roman*)]
  \item The proof of \sref{IV.8.4.2} shows that the conclusion of this proposition is valid if one supposes that the underlying space of $\mathrm{S}$ is Noetherian, that there exists $\alpha$ such that $\mathrm{S}_\alpha$ is quasi-compact, and finally that the $u_\lambda:\mathrm{S}\to\mathrm{S}_\lambda$ are dominant.
  \item On the other hand, the conclusion of \sref{IV.8.4.2} can be at fault when the $u_\lambda$ are not dominant for $\lambda$ sufficiently large, even when the $\mathrm{S}_\lambda$ and $\mathrm{S}$ are Noetherian, as shown by the following example.
Take $\mathbf{N}$ for index set, all the $\mathrm{S}_n$ equal to $\operatorname{Spec}(\mathrm{A}\times\mathrm{K})=\operatorname{Spec}(\mathrm{A})\amalg\operatorname{Spec}(\mathrm{K})$, where $\mathrm{K}$ is a field, $\mathrm{A}$ any $\mathrm{K}$-algebra, and all the morphisms $u_{n,n+1}$ equal to the same morphism corresponding to the homomorphism $(x,y)\mapsto(j(y),y)$ of $\mathrm{A}\times\mathrm{K}$ into itself, where $j:\mathrm{K}\to\mathrm{A}$ is the canonical homomorphism.
One easily verifies that the inductive limit of this system of rings is $\mathrm{K}$, the canonical homomorphism $u_n$ corresponding to the second projection $\mathrm{A}\times\mathrm{K}\to\mathrm{K}$.
One thus sees that $\mathrm{S}=\operatorname{Spec}(\mathrm{K})$ is irreducible even though none of the $\mathrm{S}_n$ is connected.
\end{enumerate}
\end{remark}

\subsection{Modules of finite presentation on a projective limit of preschemes}
\label{subsection:IV.8.5}

\begin{env}[8.5.1]
\label{IV.8.5.1}
We always keep the notations of \sref{IV.8.2.2}; we will further restrict ourselves to the case where $\mathrm{S}_0$ is one of the $\mathrm{S}_\lambda$, to which one can always reduce.

\emph{Whenever, in this section, we consider a family $(\sh{F}_\lambda)$, where, for every $\lambda\in\mathrm{L}$, $\sh{F}_\lambda$ is an $\sh{O}_{\mathrm{S}_\lambda}$-Module, it will be understood that this family satisfies the condition}
\begin{equation}
  \sh{F}_\mu=u_{\lambda\mu}^*(\sh{F}_\lambda) \qquad\text{for }\lambda\leqslant\mu.
  \tag{8.5.1.1}
\label{IV.8.5.1.1}
\end{equation}
\oldpage[IV-3]{20}
We will then set
\begin{equation}
  \sh{F}=u_\lambda^*(\sh{F}_\lambda)
  \label{IV.8.5.1.2}
\tag{8.5.1.2}
\end{equation}
which is an $\sh{O}_\mathrm{S}$-Module not depending on the index $\lambda\in\mathrm{L}$, by virtue of the hypothesis \sref{IV.8.5.1.1}.

Let $(\sh{F}_\lambda)$, $(\sh{G}_\lambda)$ be two such families of $\sh{O}_{\mathrm{S}_\lambda}$-Modules.
It is clear that the maps $f_\lambda\mapsto u_{\lambda\mu}^*(f_\lambda)$ from $\operatorname{Hom}_{\sh{O}_{\mathrm{S}_\lambda}}(\sh{F}_\lambda,\sh{G}_\lambda)$ to $\operatorname{Hom}_{\sh{O}_{\mathrm{S}_\mu}}(\sh{F}_\mu,\sh{G}_\mu)$ define an \emph{inductive system} of abelian groups $(\operatorname{Hom}_{\sh{O}_{\mathrm{S}_\lambda}}(\sh{F}_\lambda,\sh{G}_\lambda),u_{\lambda\mu}^*)$, and that the maps $f_\lambda\mapsto u_\lambda^*(f_\lambda)$ form an inductive system of homomorphisms of abelian groups
\[
  u_\lambda^*:\operatorname{Hom}_{\sh{O}_{\mathrm{S}_\lambda}}(\sh{F}_\lambda,\sh{G}_\lambda)\to\operatorname{Hom}_\mathrm{S}(\sh{F},\sh{G});
\]
whence, by passing to the inductive limit, a canonical homomorphism of abelian groups
\begin{equation}
  u_{\sh{F},\sh{G}}:\varinjlim\operatorname{Hom}_{\sh{O}_{\mathrm{S}_\lambda}}(\sh{F}_\lambda,\sh{G}_\lambda)\to\operatorname{Hom}_\mathrm{S}(\sh{F},\sh{G}).
  \tag{8.5.1.3}
\label{IV.8.5.1.3}
\end{equation}

Note that when $\sh{F}_\lambda=\sh{O}_{\mathrm{S}_\lambda}$, the condition \sref{IV.8.5.1.1} is satisfied, and one has $\sh{F}=\sh{O}_\mathrm{S}$; the homomorphism \sref{IV.8.5.1.3} thus gives \sref[0]{0.5.1.1} a canonical homomorphism of abelian groups
\begin{equation}
  u_{\sh{G}}:\varinjlim\Gamma(\mathrm{S}_\lambda,\sh{G}_\lambda)\to\Gamma(\mathrm{S},\sh{G}).
  \tag{8.5.1.4}
\label{IV.8.5.1.4}
\end{equation}
\end{env}

\begin{theorem}[8.5.2]
\label{IV.8.5.2}
\begin{enumerate}[label=(\roman*)]
  \item Suppose $\mathrm{S}_0$ quasi-compact (resp.\ quasi-compact and quasi-separated) and that, for some $\lambda\in\mathrm{L}$, $\sh{F}_\lambda$ is quasi-coherent and of finite type (resp.\ of finite presentation) and $\sh{G}_\lambda$ quasi-coherent.
Then the homomorphism $u_{\sh{F},\sh{G}}$ is injective (resp.\ bijective).
  \item Suppose $\mathrm{S}_0$ quasi-compact and quasi-separated.
For every quasi-coherent $\sh{O}_\mathrm{S}$-Module $\sh{F}$ of finite presentation, there exist $\lambda\in\mathrm{L}$ and an $\sh{O}_{\mathrm{S}_\lambda}$-Module quasi-coherent of finite presentation $\sh{F}_\lambda$ such that $\sh{F}$ is isomorphic to $u_\lambda^*(\sh{F}_\lambda)$.
\end{enumerate}
\end{theorem}

\begin{proof}
\begin{enumerate}[label=(\roman*)]
  \item We can evidently restrict to the case where $\mathrm{S}_0=\mathrm{S}_\lambda$ since the morphisms $u_{0\lambda}:\mathrm{S}_\lambda\to\mathrm{S}_0$ are affine, hence quasi-compact and separated.
Consider first the case where $\mathrm{S}_0=\operatorname{Spec}(\mathrm{A}_0)$ is \emph{affine}.
Then assertion (i) is equivalent to the

\begin{lemma}[8.5.2.1]
\label{IV.8.5.2.1}
Let $\mathrm{A}_0$ be a ring, $(\mathrm{A}_\lambda)_{\lambda\in\mathrm{L}}$ an inductive system of $\mathrm{A}_0$-algebras, $\mathrm{A}=\varinjlim\mathrm{A}_\lambda$; let $\mathrm{M}_0$, $\mathrm{N}_0$ be two $\mathrm{A}_0$-modules, and set $\mathrm{M}_\lambda=\mathrm{M}_0\otimes_{\mathrm{A}_0}\mathrm{A}_\lambda$, $\mathrm{N}_\lambda=\mathrm{N}_0\otimes_{\mathrm{A}_0}\mathrm{A}_\lambda$, $\mathrm{M}=\mathrm{M}_0\otimes_{\mathrm{A}_0}\mathrm{A}=\varinjlim\mathrm{M}_\lambda$, $\mathrm{N}=\mathrm{N}_0\otimes_{\mathrm{A}_0}\mathrm{A}=\varinjlim\mathrm{N}_\lambda$.
If $\mathrm{M}_0$ is of finite type (resp.\ of finite presentation), the canonical homomorphism
\begin{equation}
  \varinjlim\operatorname{Hom}_{\mathrm{A}_\lambda}(\mathrm{M}_\lambda,\mathrm{N}_\lambda)\to\operatorname{Hom}_\mathrm{A}(\mathrm{M},\mathrm{N})
  \tag{8.5.2.2}
\label{IV.8.5.2.2}
\end{equation}
is injective (resp.\ bijective).
\end{lemma}

We know indeed (Bourbaki, \emph{Alg.}, chap.~II, 3\textsuperscript{e}~ed., \textsection\,5, n\textsuperscript{o}\,1) that one has canonical functorial isomorphisms
\[
  \operatorname{Hom}_{\mathrm{A}_\lambda}(\mathrm{M}_\lambda,\mathrm{N}_\lambda)\xrightarrow{\sim}\operatorname{Hom}_{\mathrm{A}_0}(\mathrm{M}_0,\mathrm{N}_\lambda),\qquad\operatorname{Hom}_\mathrm{A}(\mathrm{M},\mathrm{N})\xrightarrow{\sim}\operatorname{Hom}_{\mathrm{A}_0}(\mathrm{M}_0,\mathrm{N})
\]
so that the homomorphism \sref{IV.8.5.2.2} is none other, up to canonical isomorphisms, than the canonical homomorphism
\begin{equation}
  \varinjlim\operatorname{Hom}_{\mathrm{A}_0}(\mathrm{M}_0,\mathrm{N}_\lambda)\to\operatorname{Hom}_{\mathrm{A}_0}(\mathrm{M}_0,\varinjlim\mathrm{N}_\lambda),
  \tag{8.5.2.3}
\label{IV.8.5.2.3}
\end{equation}
\oldpage[IV-3]{21}
which, to every inductive system of homomorphisms of $\mathrm{A}_0$-modules $\theta_\lambda:\mathrm{M}_0\to\mathrm{N}_\lambda$, associates its inductive limit.

Now, if $\mathrm{M}_0$ is of finite type (resp.\ of finite presentation), one has an exact sequence $\mathrm{A}_0^m\to\mathrm{M}_0\to 0$ (resp.\ $\mathrm{A}_0^n\to\mathrm{A}_0^m\to\mathrm{M}_0\to 0$); since it is clear that \sref{IV.8.5.2.3} is bijective when $\mathrm{M}_0$ is of the form $\mathrm{A}_0^q$, it suffices to use the left-exactness of the functor $\mathrm{M}_0\mapsto\operatorname{Hom}_{\mathrm{A}_0}(\mathrm{M}_0,\mathrm{P})$ and the exactness of the functor $\varinjlim$ (in the category of abelian groups) to conclude.

\medskip
Let us pass to the case where $\mathrm{S}_0$ is quasi-compact, and let $(\mathrm{U}_i)$ be a finite covering of $\mathrm{S}_0$ by affine open subsets; for every $\lambda$, the $\mathrm{U}_{i\lambda}=u_{0\lambda}^{-1}(\mathrm{U}_i)$ form an open affine covering of $\mathrm{S}_\lambda$, and the $\mathrm{V}_i=u_0^{-1}(\mathrm{U}_i)$ an open affine covering of $\mathrm{S}$.
To show that $u_{\sh{F},\sh{G}}$ is injective, it suffices to show that if $f_\lambda:\sh{F}_\lambda\to\sh{G}_\lambda$ is such that $f=u_\lambda^*(f_\lambda)=0$, then there exists $\mu\geqslant\lambda$ such that $u_{\lambda\mu}^*(f_\lambda)=0$.
By virtue of the lemma \sref{IV.8.5.2.1}, for each $i$ there exists an index $\lambda_i$ such that $f_\mu|_{\mathrm{U}_{i\mu}}=0$ for $\mu\geqslant\lambda_i$.
It suffices therefore to take $\mu$ greater than all the $\lambda_i$.

Suppose further that $\mathrm{S}_0$ is quasi-separated and $\sh{F}_0$ is of finite presentation, and let $f:\sh{F}\to\sh{G}$ be a homomorphism of $\sh{O}_\mathrm{S}$-Modules.
By virtue of the lemma \sref{IV.8.5.2.1}, for every $i$ there exist an index $\lambda_i$ and a homomorphism $f_{\lambda_i}^{(i)}:\sh{F}_{\lambda_i}|\mathrm{U}_{i,\lambda_i}\to\sh{G}_{\lambda_i}|\mathrm{U}_{i,\lambda_i}$ such that $u_{\lambda_i}^*(f_{\lambda_i}^{(i)})=f|\mathrm{V}_i$.
Since $\mathrm{L}$ is filtered, we can further suppose that all the $\lambda_i$ are equal to a single $\lambda$.
Now $\mathrm{S}_\lambda$ is quasi-separated \sref{IV.1.2.3}, and $\sh{F}_\lambda$ is a quasi-coherent $\sh{O}_{\mathrm{S}_\lambda}$-Module of finite presentation \sref[0]{0.5.2.5}; since, for every pair of indices $i,j$, $\mathrm{U}_{ij\lambda}=\mathrm{U}_{i\lambda}\cap\mathrm{U}_{j\lambda}$ is quasi-compact and, by definition,
\[
  u_\lambda^*(f_\lambda^{(i)}|\mathrm{U}_{ij\lambda})=u_\lambda^*(f_\lambda^{(j)}|\mathrm{U}_{ij\lambda})=f|(\mathrm{V}_i\cap\mathrm{V}_j),
\]
it follows from what was proved above that there exists an index $\lambda_{ij}$ such that
\[
  u_{\lambda\mu}^*(f_\lambda^{(i)}|\mathrm{U}_{ij\lambda})=u_{\lambda\mu}^*(f_\lambda^{(j)}|\mathrm{U}_{ij\lambda})\quad\text{for }\mu\geqslant\lambda_{ij}.
\]
Taking $\mu$ greater than all the $\lambda_{ij}$, one sees that $u_{\lambda\mu}^*(f_\lambda^{(i)})$ and $u_{\lambda\mu}^*(f_\lambda^{(j)})$ coincide on $\mathrm{U}_{i\mu}\cap\mathrm{U}_{j\mu}$ for every pair $(i,j)$, and hence define a homomorphism $f_\mu:\sh{F}_\mu\to\sh{G}_\mu$ such that $f=u_\mu^*(f_\mu)$.

\end{enumerate}
\end{proof}

Before passing to the proof of (ii), let us note the following corollaries of (i):

\begin{corollary}[8.5.2.4]
\label{IV.8.5.2.4}
Suppose $\mathrm{S}_0$ quasi-compact, $\sh{F}_\lambda$ quasi-coherent of finite type, $\sh{G}_\lambda$ quasi-coherent of finite presentation.
Let $f_\lambda:\sh{F}_\lambda\to\sh{G}_\lambda$ be a homomorphism.
For $f=u_\lambda^*(f_\lambda):\sh{F}\to\sh{G}$ to be an isomorphism, it is necessary and sufficient that there exist $\mu\geqslant\lambda$ such that $f_\mu=u_{\lambda\mu}^*(f_\lambda):\sh{F}_\mu\to\sh{G}_\mu$ be so.
\end{corollary}

\begin{proof}
We can always suppose $\mathrm{S}_\lambda=\mathrm{S}_0$; the question being local on $\mathrm{S}_0$, we can further (since $\mathrm{S}_0$ is quasi-compact and $\mathrm{L}$ is filtered) reduce to the case where $\mathrm{S}_0$ is affine, hence quasi-separated.
The condition being trivially sufficient, it remains to prove that it is necessary: by hypothesis there is an $\sh{O}_\mathrm{S}$-homomorphism $g:\sh{G}\to\sh{F}$ such that $g\circ f=1_\sh{F}$ and $f\circ g=1_\sh{G}$.
Since $\sh{G}$ is of finite presentation, there exist an index $\nu\geqslant\lambda$ and a homomorphism $g_\nu:\sh{G}_\nu\to\sh{F}_\nu$ such that $g=u_\nu^*(g_\nu)$ by virtue of \sref{IV.8.5.2},~(i); one therefore has
\[
  u_\nu^*(g_\nu\circ f_\nu)=1_\sh{F}
\]
\oldpage[IV-3]{22}
and $u_\nu^*(f_\nu\circ g_\nu)=1_\sh{G}$; since $\sh{F}_\nu$ and $\sh{G}_\nu$ are of finite type, it follows from \sref{IV.8.5.2},~(i) that there exists $\mu\geqslant\nu$ such that $g_\mu\circ f_\mu=1_{\sh{F}_\mu}$ and $f_\mu\circ g_\mu=1_{\sh{G}_\mu}$, whence the corollary.
\end{proof}

\begin{corollary}[8.5.2.5]
\label{IV.8.5.2.5}
Suppose $\mathrm{S}_0$ quasi-compact and quasi-separated.
Suppose that $\sh{F}_\lambda$, $\sh{G}_\lambda$ are $\sh{O}_{\mathrm{S}_\lambda}$-Modules quasi-coherent of finite presentation.
For $\sh{F}$ and $\sh{G}$ to be isomorphic, it is necessary and sufficient that there exist $\mu\geqslant\lambda$ such that $\sh{F}_\mu$ and $\sh{G}_\mu$ are isomorphic.
Moreover, for every isomorphism $f:\sh{F}\xrightarrow{\sim}\sh{G}$, there exist $\nu\geqslant\mu$ and an isomorphism $f_\nu:\sh{F}_\nu\xrightarrow{\sim}\sh{G}_\nu$ such that $f=u_\nu^*(f_\nu)$.
\end{corollary}

\begin{proof}
This follows from \sref{IV.8.5.2.4} and from \sref{IV.8.5.2},~(i) since every homomorphism $f:\sh{F}\to\sh{G}$ is of the form $u_\mu^*(f_\mu)$ for some $\mu\geqslant\lambda$ and some homomorphism $f_\mu:\sh{F}_\mu\to\sh{G}_\mu$.
\end{proof}

\begin{proof}[Proof of \sref{IV.8.5.2},~(ii)]
Consider again first the case where $\mathrm{S}_0=\operatorname{Spec}(\mathrm{A}_0)$ is affine.
Then the assertion is equivalent to the lemma \sref{IV.5.13.7.1}.

In the general case, starting from a finite open affine covering $(\mathrm{U}_i)$ of $\mathrm{S}_0$, it follows from \sref{IV.5.13.7.1} that for every $i$ there exist an index $\lambda(i)$ and a quasi-coherent $\sh{O}_{\mathrm{U}_{i,\lambda(i)}}$-Module of finite presentation $\sh{F}_i$ such that $u_{\lambda(i)}^*(\sh{F}_i)=\sh{F}|_{\mathrm{V}_i}$ (with the notations of (i)).
Furthermore, since $\mathrm{L}$ is filtered, we can suppose that all the $\lambda_i$ are equal to a single $\lambda$.
Since $\mathrm{U}_{ij\lambda}=\mathrm{U}_{i\lambda}\cap\mathrm{U}_{j\lambda}$ is quasi-compact and quasi-separated \sref{IV.1.2.7}, it follows from \sref{IV.8.5.2.5} that for every pair $(i,j)$ there exist an index $\lambda_{ij}\geqslant\lambda$ and an isomorphism
\[
  \theta_{ij}:u_{\lambda,\lambda_{ij}}^*(\sh{F}_i|\mathrm{U}_{ij\lambda})\xrightarrow{\sim}u_{\lambda,\lambda_{ij}}^*(\sh{F}_j|\mathrm{U}_{ij\lambda})
\]
such that $u_{\lambda_{ij}}^*(\theta_{ij})$ is the identity automorphism of $\sh{F}|(\mathrm{V}_i\cap\mathrm{V}_j)$; one can further suppose all the $\lambda_{ij}$ equal to $\lambda$.
Changing the notations, one can therefore suppose that it holds for every pair $(i,j)$ an isomorphism $\theta_{ij}:\sh{F}_i|_{\mathrm{U}_{ij\lambda}}\xrightarrow{\sim}\sh{F}_j|_{\mathrm{U}_{ij\lambda}}$ such that $u_\lambda^*(\theta_{ij})$ is the identity automorphism of $\sh{F}|_{(\mathrm{V}_i\cap\mathrm{V}_j)}$.
Finally, for any three indices $i$, $j$, $k$, if one sets $\mathrm{U}_{ijk,\lambda}=\mathrm{U}_{i\lambda}\cap\mathrm{U}_{j\lambda}\cap\mathrm{U}_{k\lambda}$, $\mathrm{U}_{ijk,\lambda}$ is quasi-compact, and if $\theta'_{ij}$, $\theta'_{jk}$, and $\theta'_{ik}$ denote the restrictions of $\theta_{ij}$, $\theta_{jk}$, and $\theta_{ik}$ to $\mathrm{U}_{ijk,\lambda}$, one has $u_\lambda^*(\theta'_{jk}\circ\theta'_{ij})=u_\lambda^*(\theta'_{ik})$.
There is, hence, by virtue of (i), an index $\mu\geqslant\lambda$ such that $u_{\lambda\mu}^*(\theta'_{ik})=u_{\lambda\mu}^*(\theta'_{jk}\circ\theta'_{ij})$, hence the isomorphisms $u_{\lambda\mu}^*(\theta_{ij}):u_{\lambda\mu}^*(\sh{F}_i)|_{\mathrm{U}_{ij\mu}}\xrightarrow{\sim}u_{\lambda\mu}^*(\sh{F}_j)|_{\mathrm{U}_{ij\mu}}$ satisfy the gluing condition, and define on $\mathrm{S}_\mu$ a quasi-coherent $\sh{O}_{\mathrm{S}_\mu}$-Module $\sh{F}_\mu$ of finite presentation \sref[0]{0.3.3.1} such that $\sh{F}$ is isomorphic to $u_\mu^*(\sh{F}_\mu)$.
\end{proof}

\begin{env}[8.5.3]
\label{IV.8.5.3}
The result of \sref{IV.8.5.2} can also be expressed by saying that if $\mathrm{S}_0$ is quasi-compact and quasi-separated, the category of quasi-coherent $\sh{O}_\mathrm{S}$-Modules of finite presentation is determined up to equivalence by the data of the categories of quasi-coherent $\sh{O}_{\mathrm{S}_\lambda}$-Modules of finite presentation, the functors $u_{\lambda\mu}^*$ between these categories, and the transition isomorphisms $u_{\mu\nu}^*\circ u_{\lambda\mu}^*\xrightarrow{\sim}u_{\lambda\nu}^*$.
In a figurative way, one can say that the datum of a quasi-coherent $\sh{O}_\mathrm{S}$-Module of finite presentation $\sh{F}$ is ``functorially'' equivalent to the datum of a quasi-coherent $\sh{O}_{\mathrm{S}_\lambda}$-Module of finite presentation $\sh{F}'_\lambda$ for some sufficiently large $\lambda$ and that, if a quasi-coherent $\sh{O}_{\mathrm{S}_\mu}$-Module of finite presentation $\sh{F}''_\mu$ also has $\sh{F}$ as inverse image, then $\sh{F}'_\lambda$ and $\sh{F}''_\mu$ have the same inverse image in some suitable $\mathrm{S}_\nu$ ($\nu\geqslant\lambda$, $\nu\geqslant\mu$).

We will now interpret from this point of view various notions related to quasi-coherent $\sh{O}_\mathrm{S}$-Modules.
\end{env}

\oldpage[IV-3]{23}
\begin{corollary}[8.5.4]
\label{IV.8.5.4}
Suppose $\mathrm{S}_0$ quasi-compact and quasi-separated; then, for every quasi-coherent $\sh{O}_{\mathrm{S}_\lambda}$-Module $\sh{G}_\lambda$, the canonical homomorphism \sref{IV.8.5.1.4} is bijective.
\end{corollary}

\begin{proof}
Indeed, it suffices to apply \sref{IV.8.5.2},~(i) to the case where $\sh{F}_\lambda=\sh{O}_{\mathrm{S}_\lambda}$, which is of finite presentation.
\end{proof}

\begin{proposition}[8.5.5]
\label{IV.8.5.5}
Suppose $\mathrm{S}_0$ quasi-compact, and suppose that $\sh{F}_\lambda$ is an $\sh{O}_{\mathrm{S}_\lambda}$-Module quasi-coherent of finite presentation.
For $\sh{F}$ to be locally free (resp.\ locally free of rank $n$), it is necessary and sufficient that there exist $\mu\geqslant\lambda$ such that $\sh{F}_\mu$ be so.
\end{proposition}

\begin{proof}
The condition being trivially sufficient, let us prove that it is necessary.
If $\sh{F}$ is locally free (resp.\ locally free of rank $n$), there exists a finite open affine covering $(\mathrm{V}_i)$ of $\mathrm{S}$ such that $\sh{F}|_{\mathrm{V}_i}$ is isomorphic to $\sh{O}_\mathrm{S}^{n_i}|_{\mathrm{V}_i}$ (resp.\ $\sh{O}_\mathrm{S}^n|_{\mathrm{V}_i}$) for every $i$.
By virtue of \sref{IV.8.2.11}, there exist $\nu\geqslant\lambda$ and for each $i$ a quasi-compact open subset $\mathrm{U}_{i\nu}$ of $\mathrm{S}_\nu$ such that $\mathrm{V}_i=u_\nu^{-1}(\mathrm{U}_{i\nu})$.
Since $\mathrm{S}_\nu$ is quasi-compact, each $\mathrm{U}_{i\nu}$ is a finite union of open affines; we are thus reduced to the case where $\mathrm{S}_0$ is affine and $\sh{F}=\sh{O}_\mathrm{S}^n$; we then know that there exist $\mu\geqslant\lambda$ such that $\sh{F}_\mu$ is isomorphic to $\sh{O}_{\mathrm{S}_\mu}^n$ \sref{IV.8.5.2.5}.
\end{proof}

\begin{proposition}[8.5.6]
\label{IV.8.5.6}
Suppose $\mathrm{S}_0$ quasi-compact, and consider a sequence
\[
  \sh{F}_\lambda\to\sh{G}_\lambda\to\sh{H}_\lambda\to 0
\]
of homomorphisms of $\sh{O}_{\mathrm{S}_\lambda}$-Modules quasi-coherent, where $\sh{F}_\lambda$ and $\sh{G}_\lambda$ are of finite type and $\sh{H}_\lambda$ of finite presentation.
For the corresponding sequence $\sh{F}\to\sh{G}\to\sh{H}\to 0$ to be exact, it is necessary and sufficient that there exist $\mu\geqslant\lambda$ such that the sequence $\sh{F}_\mu\to\sh{G}_\mu\to\sh{H}_\mu\to 0$ be so (in which case it is the same for $\nu\geqslant\mu$).
\end{proposition}

\begin{proof}
The fact that the condition is sufficient and the last assertion result from the fact that the functor $u_\lambda^*$ (resp.\ $u_{\mu\nu}^*$) is right-exact.
To prove that the condition is necessary, note that it follows from the hypothesis and from \sref{IV.8.5.2},~(i) that there exists $\nu\geqslant\lambda$ such that the composite $\sh{F}_\nu\to\sh{G}_\nu\to\sh{H}_\nu$ is zero.
If one sets $\sh{H}'_\nu=\operatorname{Coker}(\sh{F}_\nu\to\sh{G}_\nu)$, one has therefore a homomorphism $f_\nu:\sh{H}'_\nu\to\sh{H}_\nu$; by hypothesis, $u_\nu^*(f_\nu)$ is an isomorphism, and it follows from \sref{IV.8.5.2.4} that there exists $\mu\geqslant\nu$ such that $u_{\nu\mu}^*(f_\nu)$ is an isomorphism, which completes the proof.
\end{proof}

\begin{corollary}[8.5.7]
\label{IV.8.5.7}
Suppose $\mathrm{S}_0$ is quasi-compact, $\sh{F}_\lambda$ quasi-coherent, $\sh{G}_\lambda$ quasi-coherent of finite type, and let $f_\lambda:\sh{F}_\lambda\to\sh{G}_\lambda$ be a homomorphism.
For $f=u_\lambda^*(f_\lambda):\sh{F}\to\sh{G}$ to be surjective, it is necessary and sufficient that there exist $\mu\geqslant\lambda$ such that $f_\mu=u_{\lambda\mu}^*(f_\lambda):\sh{F}_\mu\to\sh{G}_\mu$ be so.
\end{corollary}

\begin{proof}
This is the particular case of \sref{IV.8.5.6} applied to the sequence $0\to\sh{H}_\lambda\to 0\to 0$, where $\sh{H}_\lambda=\operatorname{Coker}(f_\lambda)$, which is quasi-coherent and of finite type (taking into account the fact that one has $\sh{H}=\operatorname{Coker}(f)$ and $\sh{H}_\mu=\operatorname{Coker}(f_\mu)$, by virtue of the right-exactness of the functors $u_\lambda^*$ and $u_{\lambda\mu}^*$).
\end{proof}

\begin{corollary}[8.5.8]
\label{IV.8.5.8}
Suppose that $\mathrm{S}_0$ is quasi-compact and that the morphisms $u_{\lambda\mu}:\mathrm{S}_\mu\to\mathrm{S}_\lambda$ are flat.
Then:
\begin{enumerate}[label=(\roman*)]
  \item Let $\sh{F}_\lambda\xrightarrow{f_\lambda}\sh{G}_\lambda\xrightarrow{g_\lambda}\sh{H}_\lambda$ be a sequence of homomorphisms of $\sh{O}_{\mathrm{S}_\lambda}$-Modules quasi-coherent, such that $\operatorname{Im} f_\lambda$ and $\operatorname{Ker} g_\lambda$ are of finite type.
For the corresponding sequence $\sh{F}\xrightarrow{f}\sh{G}\xrightarrow{g}\sh{H}$ to be
\oldpage[IV-3]{24}
exact, it is necessary and sufficient that there exist $\mu\geqslant\lambda$ such that the sequence $\sh{F}_\mu\xrightarrow{f_\mu}\sh{G}_\mu\xrightarrow{g_\mu}\sh{H}_\mu$ be so (in which case the same holds for $\nu\geqslant\mu$).
  \item Let $f_\lambda:\sh{F}_\lambda\to\sh{G}_\lambda$ be a homomorphism of quasi-coherent $\sh{O}_{\mathrm{S}_\lambda}$-Modules, and suppose that $\operatorname{Ker} f_\lambda$ is of finite type.
For $f=u_\lambda^*(f_\lambda):\sh{F}\to\sh{G}$ to be injective, it is necessary and sufficient that there exist $\mu\geqslant\lambda$ such that $f_\mu=u_{\lambda\mu}^*(f_\lambda):\sh{F}_\mu\to\sh{G}_\mu$ be so.
\end{enumerate}
\end{corollary}

\begin{proof}
(i) Taking account of \sref{IV.8.3.8},~(ii), observe that, by flatness, $\operatorname{Im} f$ and $\operatorname{Ker} g$ (resp. $\operatorname{Im} f_\mu$ and $\operatorname{Ker} g_\mu$ for $\mu\geqslant\lambda$) are the inverse images of $\operatorname{Im} f_\lambda$ and $\operatorname{Ker} g_\lambda$ \sref[0\textsubscript{I}]{0.6.7.2}. Suppose that the sequence $\sh{F}\xrightarrow{f}\sh{G}\xrightarrow{g}\sh{H}$ is exact. Since $\operatorname{Im} f_\lambda$ is of finite type, it follows from \sref{IV.8.5.2},~(i) that there exists $\mu\geqslant\lambda$ such that the composite $\sh{F}_\mu\xrightarrow{f_\mu}\sh{G}_\mu\xrightarrow{g_\mu}\sh{H}_\mu$ is zero. Changing notation, we may therefore suppose already that $g_\lambda\circ f_\lambda=0$. Since $\operatorname{Ker} g_\lambda$ is of finite type, it then follows from \sref{IV.8.5.7} that there exists $\mu\geqslant\lambda$ such that the homomorphism $\sh{F}_\mu\to\operatorname{Ker} g_\mu$ is surjective, which proves (i).

(ii) This is the special case of (i) applied to the sequence $0\to\sh{F}_\lambda\to\sh{G}_\lambda$.
\end{proof}

\begin{lemma}[8.5.9]
\label{IV.8.5.9}
Suppose that $\mathrm{S}_0$ is quasi-compact and that $\sh{F}_\lambda$ is quasi-coherent and of finite type. Let $\sh{G}'_\lambda$ and $\sh{G}''_\lambda$ be two quasi-coherent quotients of $\sh{F}_\lambda$, with $\sh{G}'_\lambda$ moreover supposed to be of finite presentation. For $\sh{G}''$ to be a quotient of $\sh{G}'$, it is necessary and sufficient that there exist $\mu\geqslant\lambda$ such that $\sh{G}''_\mu$ is a quotient of $\sh{G}'_\mu$.
\end{lemma}

\begin{proof}
By hypothesis, there are two surjective homomorphisms $p'_\lambda:\sh{F}_\lambda\to\sh{G}'_\lambda$ and $p''_\lambda:\sh{F}_\lambda\to\sh{G}''_\lambda$. By the right exactness of $u_\lambda^*$ and $u_{\lambda\mu}^*$, the homomorphisms $p'=u_\lambda^*(p'_\lambda)$, $p''=u_\lambda^*(p''_\lambda)$, $p'_\mu=u_{\lambda\mu}^*(p'_\lambda)$, and $p''_\mu=u_{\lambda\mu}^*(p''_\lambda)$ are also surjective. Moreover, if there exists a homomorphism $f:\sh{G}'\to\sh{G}''$ (resp. $f_\mu:\sh{G}'_\mu\to\sh{G}''_\mu$) such that $p''=f\circ p'$ (resp. $p''_\mu=f_\mu\circ p'_\mu$), that homomorphism is necessarily unique; this shows that the question is local on $\mathrm{S}_\lambda$, and consequently, since $\mathrm{S}_0$ is quasi-compact and $\mathbf{L}$ is filtered, that we may suppose $\mathrm{S}_\lambda$ affine and therefore quasi-separated. The condition in the statement is plainly sufficient. Conversely, since $\sh{G}'_\lambda$ is of finite presentation and $\mathrm{S}_\lambda$ is quasi-compact and quasi-separated, it follows from \sref{IV.8.5.2},~(i) that if there exists $f:\sh{G}'\to\sh{G}''$ such that $p''=f\circ p'$, then there exist $\mu\geqslant\lambda$ and $f_\mu:\sh{G}'_\mu\to\sh{G}''_\mu$ such that $f=u_\mu^*(f_\mu)$ and $p''_\mu=f_\mu\circ p'_\mu$, proving the lemma.
\end{proof}

\begin{env}[8.5.10]
\label{IV.8.5.10}
For the remainder of this section, for every quasi-coherent Module $\sh{F}$ on a prescheme, denote by $\mathfrak{Q}(\sh{F})$ the set of quotient Modules of $\sh{F}$ that are of finite presentation. If $\sh{F}_\lambda$ is quasi-coherent and $\sh{G}_\lambda\in\mathfrak{Q}(\sh{F}_\lambda)$, it follows from the right exactness of $u_{\lambda\mu}^*$ and $u_\lambda^*$ that $\sh{G}_\mu\in\mathfrak{Q}(\sh{F}_\mu)$ for $\mu\geqslant\lambda$ and $\sh{G}\in\mathfrak{Q}(\sh{F})$. It is clear that $(\mathfrak{Q}(\sh{F}_\lambda),u_{\lambda\mu}^*)$ is an inductive system of sets and that the $u_\lambda^*:\mathfrak{Q}(\sh{F}_\lambda)\to\mathfrak{Q}(\sh{F})$ form an inductive system of maps; hence, by passing to the inductive limit, there is a canonical map
\[
u_{\mathfrak{Q}}:\varinjlim\mathfrak{Q}(\sh{F}_\lambda)\longrightarrow\mathfrak{Q}(\sh{F}).
\tag{8.5.10.1}\label{IV.8.5.10.1}
\]
Furthermore, if $(\sh{F}'_\lambda)$ is a second family of quasi-coherent $\sh{O}_{\mathrm{S}_\lambda}$-Modules and, for every $\lambda$, $\sh{F}'_\lambda$ is a quotient of $\sh{F}_\lambda$, then $\sh{F}'$ is a quotient of $\sh{F}$, and one has a commutative diagram

\oldpage[IV-3]{25}

\[
\begin{tikzcd}
\varinjlim\mathfrak{Q}(\sh{F}_\lambda) \ar[r,"u_{\mathfrak{Q}}"] \ar[d] & \mathfrak{Q}(\sh{F}) \ar[d] \\
\varinjlim\mathfrak{Q}(\sh{F}'_\lambda) \ar[r,"u_{\mathfrak{Q}}"'] & \mathfrak{Q}(\sh{F}')
\end{tikzcd}
\tag{8.5.10.2}\label{IV.8.5.10.2}
\]
\end{env}

\begin{proposition}[8.5.11]
\label{IV.8.5.11}
Suppose that $\mathrm{S}_0$ is quasi-compact (resp. quasi-compact and quasi-separated). Suppose that $\sh{F}_\lambda$ is quasi-coherent and of finite type (resp. of finite presentation) for every $\lambda$; then the canonical map \eqref{IV.8.5.10.1} is injective (resp. bijective).
\end{proposition}

\begin{proof}
The first assertion follows more precisely from Lemma \sref{IV.8.5.9}. For the second, consider a quotient $\sh{O}_{\mathrm{S}}$-Module $\sh{G}$ of $\sh{F}$ that is of finite presentation. It follows from \sref{IV.8.5.2},~(ii) that there exist $\lambda'\in\mathbf{L}$ and a quasi-coherent $\sh{O}_{\mathrm{S}_{\lambda'}}$-Module $\sh{G}_{\lambda'}$ of finite presentation such that $\sh{G}=u_{\lambda'}^*(\sh{G}_{\lambda'})$; since $\mathbf{L}$ is filtered, we may replace $\lambda$ and $\lambda'$ by an index majorizing both and suppose that $\lambda'=\lambda$. Consider then the canonical homomorphism $p:\sh{F}\to\sh{G}$. It follows from \sref{IV.8.5.2},~(i) that there exist $\mu\geqslant\lambda$ and a homomorphism $p_\mu:\sh{F}_\mu\to\sh{G}_\mu$ such that $p=u_\mu^*(p_\mu)$. By \sref{IV.8.5.7}, we may moreover suppose that $\mu$ is large enough for $p_\mu$ to be surjective, which completes the proof.
\end{proof}

\subsection{Sub-preschemes of finite presentation of a projective limit of preschemes}
\label{subsection:IV.8.6}

\begin{env}[8.6.1]
\label{IV.8.6.1}
Given a prescheme $Y$, we denote in this section by $\mathfrak{Spr}(Y)$ the ordered set
\sref[I]{I.4.1.10} of sub-preschemes of $Y$ that are \emph{of finite presentation}
\sref{IV.1.6.1}, by $\mathfrak{Spro}(Y)$ (resp.\ $\mathfrak{Sprf}(Y)$) the subset of $\mathfrak{Spr}(Y)$ formed by the sub-preschemes induced on opens (resp.\ on closed sub-preschemes) of $Y$, of finite presentation over $Y$; it amounts to the same to say that a sub-prescheme $Z$ of $Y$ belongs to $\mathfrak{Spro}(Y)$ (resp.\ $\mathfrak{Sprf}(Y)$) or that it is induced on an open and that the underlying space is \emph{retrocompact} in $Y$ (resp.\ that it is closed and that the ideal $\sh{J}$ of $\sh{O}_Y$ that defines it is \emph{of finite type}, which also means that $j_*(\sh{O}_Z)\in\mathfrak{Q}(\sh{O}_Y)$ where $j:Z\to Y$ is the canonical injection) \sref{IV.1.6.1} and \sref{IV.1.4.5}.
\end{env}

\begin{env}[8.6.2]
\label{IV.8.6.2}
We always preserve the notations of \sref{IV.8.2.2} and suppose that $S_0$ is one of the $S_\lambda$.
Let $Y_\lambda$ be a sub-prescheme of $S_\lambda$; then $Y_\mu=u_{\lambda\mu}^{-1}(Y_\lambda)$ (resp.\ $Y=u_\lambda^{-1}(Y_\lambda)$) is a sub-prescheme of $S_\mu$ for $\mu\geqslant\lambda$ (resp.\ of $S$). It is induced on an open (resp.\ it is closed) if $Y_\lambda$ is so \sref[I]{I.4.3.2}, and is of finite presentation over $S_\mu$ (resp.\ $S$) if $Y_\lambda$ is of finite presentation over $S_\lambda$ \sref{IV.1.6.2},~(iii). It follows that $(\mathfrak{Spr}(S_\lambda),u_{\lambda\mu}^{-1})$ (resp.\ $(\mathfrak{Spro}(S_\lambda),u_{\lambda\mu}^{-1})$, $(\mathfrak{Sprf}(S_\lambda),u_{\lambda\mu}^{-1})$) is an \emph{inductive system} of sets and the maps $u_\lambda^{-1}:\mathfrak{Spr}(S_\lambda)\to\mathfrak{Spr}(S)$ (resp.\ $\mathfrak{Spro}(S_\lambda)\to\mathfrak{Spro}(S)$, $\mathfrak{Sprf}(S_\lambda)\to\mathfrak{Sprf}(S)$) form an inductive system of maps; whence, by passing to the inductive limit, canonical maps
\end{env}

\begin{align}
  \varinjlim \mathfrak{Spr}(S_\lambda) &\to \mathfrak{Spr}(S) \tag{8.6.2.1}\label{IV.8.6.2.1}\\
  \varinjlim \mathfrak{Spro}(S_\lambda) &\to \mathfrak{Spro}(S) \tag{8.6.2.2}\label{IV.8.6.2.2}\\
  \varinjlim \mathfrak{Sprf}(S_\lambda) &\to \mathfrak{Sprf}(S). \tag{8.6.2.3}\label{IV.8.6.2.3}
\end{align}

\oldpage[IV-3]{26}

Recall \sref[I]{I.4.1.10} that $\mathfrak{Spr}(X)$, for every prescheme $X$, is an ordered set by the relation ``$Z$ is a sub-prescheme of the sub-prescheme $Y$'' which we write $Z\leqslant Y$.
The maps $u_{\lambda\mu}^{-1}$ and $u_\lambda^{-1}$ are \emph{increasing} for the corresponding order relations in $\mathfrak{Spr}(S_\lambda)$, $\mathfrak{Spr}(S_\mu)$, $\mathfrak{Spr}(S)$.
Furthermore, by writing that $\eta\leqslant\eta'$ for two elements of this set whenever there exist $\lambda$ and two elements $Y_\lambda$, $Y'_\lambda$ of $\mathfrak{Spr}(S_\lambda)$, of which $\eta$ and $\eta'$ are the canonical images, and which are such that $Y_\lambda\leqslant Y'_\lambda$; one easily verifies that this does not depend on the choice of representatives, and that one thus has indeed a \emph{relation of order}.
The fact that the $u_\lambda^{-1}$ are increasing immediately implies that the canonical map \eqref{IV.8.6.2.1} is \emph{increasing}; the same is evidently true of \eqref{IV.8.6.2.2} and \eqref{IV.8.6.2.3}.

\begin{proposition}[8.6.3]
\label{IV.8.6.3}
Suppose $S_0$ quasi-compact (resp.\ quasi-compact and quasi-separated).
Then the maps \eqref{IV.8.6.2.1}, \eqref{IV.8.6.2.2}, \eqref{IV.8.6.2.3} are injective (resp.\ bijective).
\end{proposition}

\begin{proof}
Taking into account the remarks of \sref{IV.8.6.1}, the assertions relative to \eqref{IV.8.6.2.3} follow from \sref{IV.8.5.11} applied to $\sh{F}_\lambda=\sh{O}_{S_\lambda}$; similarly, the assertions relative to \eqref{IV.8.6.2.2} are special cases of \sref{IV.8.3.5} and \sref{IV.8.3.11}, taking into account that the $S_\lambda$ and $S$ are quasi-compact.
It remains to consider the map \eqref{IV.8.6.2.1}.
Let us first prove that it is surjective when $S_0$ is quasi-compact and quasi-separated.
Let $Z$ be a sub-prescheme of $S$, of finite presentation over $S$; as $S$ is quasi-compact, so is $Z$, and therefore there exists an open quasi-compact $U$ of $S$ such that $Z$ is a sub-prescheme of $U$, closed in $U$, of finite presentation over $U$.
There then exist an index $\lambda$ and an open quasi-compact $U_\lambda$ of $S_\lambda$ such that $U=u_\lambda^{-1}(U_\lambda)$ \sref{IV.8.2.11}; as $S_\lambda$ is quasi-separated, so is $U_\lambda$ \sref{IV.1.2.7}, and consequently one can restrict to the case where $U=S$; but in this case, one is reduced to the fact that \eqref{IV.8.6.2.3} is surjective.

Finally, to see that \eqref{IV.8.6.2.1} is injective when $S_0$ is quasi-compact, it will suffice to prove the following more precise result:
\end{proof}

\begin{corollary}[8.6.3.1]
\label{IV.8.6.3.1}
Suppose $S_0$ quasi-compact and let $Z'_\lambda$, $Z''_\lambda$ be two sub-preschemes of $S_\lambda$, of finite presentation.
For $Z'=u_\lambda^{-1}(Z'_\lambda)$ to be majorised by $Z''=u_\lambda^{-1}(Z''_\lambda)$ \sref[I]{I.4.1.10}, it is necessary and sufficient that there exist $\mu\geqslant\lambda$ such that $Z'_\mu=u_{\lambda\mu}^{-1}(Z'_\lambda)$ be majorised by $Z''_\mu=u_{\lambda\mu}^{-1}(Z''_\lambda)$.
\end{corollary}

\begin{proof}
It is trivial that the condition is sufficient.
To see that it is necessary, note first that the underlying sets $Z'_\lambda$ and $Z''_\lambda$ are locally constructible in $S_\lambda$ by hypothesis \sref{IV.1.8.4}, hence the hypothesis implies the existence of a $\nu\geqslant\lambda$ such that $Z'_\nu\subset Z''_\nu$ \sref{IV.8.3.5}; replacing $\lambda$ by $\nu$, we can therefore already suppose that $Z'_\lambda\subset Z''_\lambda$.
Furthermore, by hypothesis \sref{IV.1.6.1}, the sub-spaces $Z'_\lambda$ and $Z''_\lambda$ of $S_\lambda$ are quasi-compact; for every point $x\in Z'_\lambda$, there is an open quasi-compact neighbourhood $V(x)$ in $S_\lambda$ such that $V(x)\cap Z'_\lambda$ and $V(x)\cap Z''_\lambda$ are closed in $V(x)$.
By covering $Z'_\lambda$ by a finite number of such neighbourhoods, we see that there is an open quasi-compact $U_\lambda$ of $S_\lambda$ containing $Z'_\lambda$ and such that $Z'_\lambda$ and $Z''_\lambda\cap U_\lambda$ are closed in $U_\lambda$.
If we denote by $Y_\lambda$ the sub-prescheme of $S_\lambda$ induced by $Z''_\lambda$ on $U_\lambda\cap Z''_\lambda$, it is clear that with the usual notations, $Y_\mu$ (resp.\ $Y$) is induced by $Z''_\mu$ on $U_\mu\cap Z''_\mu$ for $\mu\geqslant\lambda$ (resp.\ by $Z''$ on $U\cap Z''$); moreover $Z'$ is majorised by $Y$ \sref[I]{I.4.4.1}, and as it suffices to prove that $Z'_\mu$ is majorised by $Y_\mu$ for $\mu$
\oldpage[IV-3]{27}
large enough, we see finally that one is reduced (replacing $S_\lambda$ by $U_\lambda$) to the case where $Z'_\lambda$ and $Z''_\lambda$ are closed sub-preschemes of $S_\lambda$.
But then, this has already been proved since \eqref{IV.8.6.2.3} is increasing and injective.
\end{proof}

\begin{corollary}[8.6.4]
\label{IV.8.6.4}
Suppose $S_0$ quasi-compact, and let $Z_\lambda$ be a sub-prescheme of $S_\lambda$, of finite presentation over $S_\lambda$.
For $Z=u_\lambda^{-1}(Z_\lambda)$ to be the subprescheme induced on an open subset of $S$ (resp.\ a closed subprescheme of $S$), it is necessary and sufficient that there exist $\mu\geqslant\lambda$ such that $Z_\mu=u_{\lambda\mu}^{-1}(Z_\lambda)$ is the subprescheme induced on an open subset of $S_\mu$ (resp.\ a closed subprescheme of $S_\mu$).
\end{corollary}

\begin{proof}
Let $(U_\lambda^{(i)})$ be a finite affine open covering of $S_\lambda$, and set $U^{(i)}_\mu=u_{\lambda\mu}^{-1}(U_\lambda^{(i)})$ for $\mu\geqslant\lambda$ and $U^{(i)}=u_\lambda^{-1}(U_\lambda^{(i)})$.
If $Z$ is open (resp.\ closed) in $S$, $Z\cap U^{(i)}$ is open (resp.\ closed) in $U^{(i)}$, and conversely if each of the $Z_\mu\cap U^{(i)}_\mu$ is open (resp.\ closed) in $U^{(i)}_\mu$, $Z_\mu$ is open (resp.\ closed) in $S_\mu$.
Since $L$ is filtered, it suffices to prove the corollary when $S_\lambda$ is affine, hence quasi-separated, and the result follows from the fact that the maps \eqref{IV.8.6.2.1}, \eqref{IV.8.6.2.2}, and \eqref{IV.8.6.2.3} are bijective.
\end{proof}

\subsection{Criteria for a projective limit of preschemes to be a reduced (resp.\ integral) prescheme}
\label{subsection:IV.8.7}

We preserve the hypotheses and notations of \sref{IV.8.2.2} and suppose always that $S_0$ is one of the $S_\lambda$.

\begin{proposition}[8.7.1]
\label{IV.8.7.1}
Suppose that $S$ is not reduced.
Then there exists $\lambda_0$ such that for $\lambda\geqslant\lambda_0$, $S_\lambda$ is not reduced.
\end{proposition}

\begin{proof}
The question being local on $S_0$, we can suppose $S_0=\operatorname{Spec}(A_0)$ affine, whence $S=\operatorname{Spec}(A)$, where $A=\varinjlim A_\lambda$ is the inductive limit of an inductive system of $A_0$-algebras $(A_\lambda)$.
We know then \sref{IV.5.13.2} that the nilradical of $A$ is the inductive limit of those of the $A_\lambda$; if it is not zero, one of the $A_\lambda$ therefore contains a nilpotent element $a_\lambda\neq 0$ whose image in $A$ is nilpotent and nonzero, and consequently the image of $a_\lambda$ in $A_\mu$ is nilpotent and nonzero for $\mu\geqslant\lambda$.
\end{proof}

\begin{proposition}[8.7.2]
\label{IV.8.7.2}
Suppose one of the following hypotheses satisfied:
\begin{enumerate}[label=\emph{\alph*)}]
\item $S_0$ is quasi-compact, the nilradical $\sh{N}_0$ of $\sh{O}_{S_0}$ is an ideal of finite type (which is satisfied for example when $S_0$ is Noetherian), and the morphisms $u_{\lambda\mu}:S_\mu\to S_\lambda$ are open immersions.
\item The morphisms $u_{\lambda\mu}:S_\mu\to S_\lambda$ are faithfully flat.
\end{enumerate}
Under these conditions, for $S$ to be reduced, it is necessary and sufficient that there exist $\lambda_0$ such that $S_\lambda$ be reduced for $\lambda\geqslant\lambda_0$.

Furthermore, in case b), if $S$ is reduced, all the $S_\lambda$ are reduced.
\end{proposition}

\begin{proof}
The last assertion follows from the fact that the morphism $S\to S_\lambda$ is then faithfully flat for every $\lambda$ \sref{IV.8.3.8} and from \sref{IV.2.1.13}.
On the other hand, \sref{IV.8.7.1} proves that the condition of the statement is sufficient (without hypothesis on $S_0$ or on the $u_{\lambda\mu}$).
It therefore remains to see that the condition is necessary in hypothesis $a)$; then by \sref{IV.8.2.13}, the underlying space of $S$ is identified with the intersection of the underlying spaces of the $S_\lambda$ (the $S_\lambda$ being identified with sub-preschemes induced on opens of $S_0$), and the structural sheaf $\sh{O}_S$ is identified with
\oldpage[IV-3]{28}
the sheaf induced on $S$ by all the $\sh{O}_{S_\lambda}$; in particular for every $s\in S$, the local ring $\sh{O}_s$ is the same for $S$ and for all the $S_\lambda$.
If $\sh{N}_\lambda$ is the nilradical of $\sh{O}_{S_\lambda}$, the nilradical $\sh{N}$ of $\sh{O}_S$ therefore has at each point of $S$ the same fibre $\sh{N}_s$ (the nilradical of $\sh{O}_s$) as $u_\lambda^*(\sh{N}_\lambda)$ (induced on $S$ by $\sh{N}_\lambda$).
The hypothesis that $S$ is reduced implies therefore $u_\lambda^*(\sh{N}_\lambda)=0$; as $\sh{N}_0$ is supposed of finite type, the same is true of the $\sh{N}_\lambda$, and the conclusion follows then from \sref{IV.8.5.7}.
\end{proof}

\begin{corollary}[8.7.3]
\label{IV.8.7.3}
Suppose one of the following hypotheses satisfied:
\begin{enumerate}[label=\emph{\alph*)}]
\item $S_0$ is a Noetherian prescheme and the morphisms $u_{\lambda\mu}:S_\mu\to S_\lambda$ are open immersions.
\item The morphisms $u_{\lambda\mu}:S_\mu\to S_\lambda$ are faithfully flat.
\end{enumerate}
Then, for $S$ to be integral, it is necessary and sufficient that there exist $\lambda_0$ such that $S_\lambda$ be integral for $\lambda\geqslant\lambda_0$.
\end{corollary}

\begin{proof}
To say that a prescheme is integral means that it is both reduced and irreducible. Reducedness follows from \sref{IV.8.7.2}. In case (a), irreducibility follows from \sref{IV.8.4.3}. In case (b), \sref{IV.8.3.8} shows that every morphism $S\to S_\lambda$ is faithfully flat, hence surjective; thus irreducibility of $S$ implies irreducibility of every $S_\lambda$. Conversely, eventual irreducibility of the $S_\lambda$ implies irreducibility of $S$ by the general projective-limit argument used in the proof of \sref{IV.8.4.1}.
\end{proof}

\begin{remark}[8.7.4]
\label{IV.8.7.4}
If one makes no hypothesis on the $u_{\lambda\mu}$, it can happen that $S$ is integral even though all the $S_\lambda$ are neither reduced nor connected, as the example \sref{IV.8.4.6},~(ii) shows, where one takes the ring $A$ non reduced.
\end{remark}

\subsection{Preschemes of finite presentation over a projective limit of preschemes}
\label{subsection:IV.8.8}

\begin{env}[8.8.1]
\label{IV.8.8.1}
Always preserving the notations and hypotheses of \sref{IV.8.2.2}, we shall suppose given in this section two $S_\alpha$-preschemes $X_\alpha$, $Y_\alpha$, which defines \sref{IV.8.2.5} two projective systems of preschemes $(X_\lambda,v_{\lambda\mu})$ and $(Y_\lambda,w_{\lambda\mu})$ by setting $X_\lambda=X_\alpha\times_{S_\alpha}S_\lambda$, $Y_\lambda=Y_\alpha\times_{S_\alpha}S_\lambda$, $v_{\lambda\mu}=1_{X_\alpha}\times u_{\lambda\mu}$, $w_{\lambda\mu}=1_{Y_\alpha}\times u_{\lambda\mu}$ (for $\alpha\leqslant\lambda\leqslant\mu$), whose projective limits are respectively $X=X_\alpha\times_{S_\alpha}S$, $Y=Y_\alpha\times_{S_\alpha}S$, with the canonical morphisms $v_\lambda=1_{X_\alpha}\times u_\lambda:X\to X_\lambda$ and $w_\lambda=1_{Y_\alpha}\times u_\lambda:Y\to Y_\lambda$.
For $\alpha\leqslant\lambda\leqslant\mu$, there is a canonical map $e_{\mu\lambda}:\operatorname{Hom}_{S_\lambda}(X_\lambda,Y_\lambda)\to\operatorname{Hom}_{S_\mu}(X_\mu,Y_\mu)$, which to every $S_\lambda$-morphism $f_\lambda:X_\lambda\to Y_\lambda$ associates $f_\mu=f_\lambda\times 1_{S_\mu}:X_\lambda\times_{S_\lambda}S_\mu\to Y_\lambda\times_{S_\lambda}S_\mu$, and it is clear that $(\operatorname{Hom}_{S_\lambda}(X_\lambda,Y_\lambda),e_{\mu\lambda})$ is an \emph{inductive system} of sets.
Similarly, there is a canonical map $e_\lambda:\operatorname{Hom}_{S_\lambda}(X_\lambda,Y_\lambda)\to\operatorname{Hom}_S(X,Y)$ which to $f_\lambda$ associates $f=f_\lambda\times 1_S:X_\lambda\times_{S_\lambda}S\to Y_\lambda\times_{S_\lambda}S$ and $(e_\lambda)$ is an inductive system of maps; whence, by passing to the inductive limit, a canonical map, functorial in $S_\alpha$, $X_\alpha$, and $Y_\alpha$:
\[
  e:\varinjlim\operatorname{Hom}_{S_\lambda}(X_\lambda,Y_\lambda)\to\operatorname{Hom}_S(X,Y).
  \tag{8.8.1.1}\label{IV.8.8.1.1}
\]
\end{env}

\begin{theorem}[8.8.2]
\label{IV.8.8.2}
\begin{enumerate}[label=(\roman*)]
\item Suppose $X_\alpha$ quasi-compact (resp.\ quasi-compact and quasi-separated), and $Y_\alpha$ locally of finite type (resp.\ locally of finite presentation) over $S_\alpha$.
Then the map \eqref{IV.8.8.1.1} is injective (resp.\ bijective).
\item Suppose $S_0$ quasi-compact and quasi-separated.
For every prescheme $X$ of finite presentation over $S$, there exist a $\lambda\in L$, a prescheme $X_\lambda$ of finite presentation over $S_\lambda$, and an $S$-isomorphism $X\xrightarrow{\sim}X_\lambda\times_{S_\lambda}S$.
\end{enumerate}
\end{theorem}

\begin{proof}
\oldpage[IV-3]{29}
(i) Consider first the case where $S_0=\operatorname{Spec}(A_0)$, $X_\alpha=\operatorname{Spec}(B_\alpha)$, and $Y_\alpha=\operatorname{Spec}(C_\alpha)$ are affine; then the $S_\lambda=\operatorname{Spec}(A_\lambda)$ and $S=\operatorname{Spec}(A)$ are also affine, with $A=\varinjlim A_\lambda$, and the assertions of (i) are equivalent to the

\begin{lemma}[8.8.2.1]
\label{IV.8.8.2.1}
Let $A_0$ be a ring, $(A_\lambda)_{\lambda\in L}$ an inductive system of $A_0$-algebras, $A=\varinjlim A_\lambda$; let $B_\alpha$ be an $A_\alpha$-algebra, $C_\alpha$ an $A_\alpha$-algebra of finite type (resp.\ of finite presentation).
Then the canonical homomorphism
\[
  \varinjlim_\lambda\operatorname{Hom}_{A_\lambda\text{-alg.}}(C_\alpha\otimes_{A_\alpha}A_\lambda,B_\alpha\otimes_{A_\alpha}A_\lambda)\to\operatorname{Hom}_{A\text{-alg.}}(C_\alpha\otimes_{A_\alpha}A,B_\alpha\otimes_{A_\alpha}A)
  \tag{8.8.2.2}\label{IV.8.8.2.2}
\]
is injective (resp.\ bijective).
\end{lemma}

\begin{proof}
We know that there are functorial canonical isomorphisms
\[
  \operatorname{Hom}_{A_\lambda\text{-alg.}}(C_\alpha\otimes_{A_\alpha}A_\lambda,B_\alpha\otimes_{A_\alpha}A_\lambda)\xrightarrow{\sim}\operatorname{Hom}_{A_\alpha\text{-alg.}}(C_\alpha,B_\alpha\otimes_{A_\alpha}A_\lambda)
\]
\[
  \operatorname{Hom}_{A\text{-alg.}}(C_\alpha\otimes_{A_\alpha}A,B_\alpha\otimes_{A_\alpha}A)\xrightarrow{\sim}\operatorname{Hom}_{A_\alpha\text{-alg.}}(C_\alpha,B_\alpha\otimes_{A_\alpha}A)
\]
by virtue of the universal property of the tensor product of two algebras.
It therefore suffices to prove the
\end{proof}

\begin{lemma}[8.8.2.3]
\label{IV.8.8.2.3}
Let $E$ be a ring, $G$ an $E$-algebra of finite type (resp.\ of finite presentation), $(F_\lambda)$ an inductive system of $E$-algebras.
Then the canonical homomorphism
\[
  \varinjlim\operatorname{Hom}_{E\text{-alg.}}(G,F_\lambda)\to\operatorname{Hom}_{E\text{-alg.}}(G,\varinjlim F_\lambda)
\]
which, to every inductive system of homomorphisms $\theta_\lambda:G\to F_\lambda$ of $E$-algebras, associates its inductive limit, is injective (resp.\ bijective).
\end{lemma}

\begin{proof}
Suppose first the $E$-algebra $G$ of finite type, and let $(t_i)_{1\leqslant i\leqslant n}$ be a system of generators of this $E$-algebra; let us show that if $(\theta'_\lambda)$, $(\theta''_\lambda)$ are two inductive systems of homomorphisms $G\to F_\lambda$, with $\varinjlim\theta'_\lambda=\varinjlim\theta''_\lambda$, there exists $\mu\geqslant\lambda$ such that $\theta'_\mu=\theta''_\mu$.
Indeed, if $\varphi_{\mu\lambda}:F_\lambda\to F_\mu$ and $\varphi_\lambda:F_\lambda\to F=\varinjlim F_\lambda$ are the homomorphisms of the inductive system $(F_\lambda)$, by hypothesis, for each index $i$, there exists an index $\lambda_i$ such that $\varphi_{\lambda_i}(\theta'_{\lambda_i}(t_i))=\varphi_{\lambda_i}(\theta''_{\lambda_i}(t_i))$, and one can suppose all the $\lambda_i$ equal to a single $\lambda$; it follows in the same way the existence of $\mu\geqslant\lambda$ such that $\varphi_{\mu\lambda}(\theta'_\lambda(t_i))=\varphi_{\mu\lambda}(\theta''_\lambda(t_i))$ for $1\leqslant i\leqslant n$, that is to say $\theta'_\mu(t_i)=\theta''_\mu(t_i)$ for $1\leqslant i\leqslant n$, whence $\theta'_\mu=\theta''_\mu$.

Suppose in the second place $G$ of finite presentation, so that one has
\[
  G=E[T_1,\ldots,T_n]/\mathfrak{J},
\]
where $\mathfrak{J}$ is an ideal of finite type, $t_i$ being the class of $T_i$ modulo $\mathfrak{J}$.
Let $(P_j)_{1\leqslant j\leqslant m}$ be a system of generators of $\mathfrak{J}$.
Suppose given a homomorphism of $E$-algebras $\theta:G\to F$; set $b_i=\theta(t_i)$; by definition, one has therefore $P_j(b_1,b_2,\ldots,b_n)=\theta(P_j(t_1,\ldots,t_n))=0$ for $1\leqslant j\leqslant m$.
Now, there exist $\lambda$ and elements $x_1,\ldots,x_n$ in $F_\lambda$ such that $b_i=\varphi_\lambda(x_i)$ for $1\leqslant i\leqslant n$; one has therefore $\varphi_\lambda(P_j(x_1,\ldots,x_n))=P_j(b_1,\ldots,b_n)=0$, and consequently there exists $\mu\geqslant\lambda$ such that $\varphi_{\mu\lambda}(P_j(x_1,\ldots,x_n))=P_j(\varphi_{\mu\lambda}(x_1),\ldots,\varphi_{\mu\lambda}(x_n))=0$ for $1\leqslant j\leqslant m$; one concludes that there exists a homomorphism of $E$-algebras $\theta_\mu:G\to F_\mu$ such that $\theta_\mu(t_i)=\varphi_{\mu\lambda}(x_i)$
\oldpage[IV-3]{30}
for $1\leqslant i\leqslant n$; one therefore deduces for every $\nu\geqslant\mu$ a homomorphism of $E$-algebras $\theta_\nu=\varphi_{\nu\mu}\circ\theta_\mu$ of $G$ into $F_\nu$, and it is clear that $\theta$ is the inductive limit of this inductive system of homomorphisms, which completes the proof of the lemma.
\end{proof}

Let us now pass to the case where $X_\alpha$ is quasi-compact and $Y_\alpha$ locally of finite type over $S_\alpha$.
Set $Z_\alpha=X_\alpha\times_{S_\alpha}Y_\alpha$ and introduce the projective system corresponding to $Z_\lambda=Z_\alpha\times_{S_\alpha}S_\lambda=X_\lambda\times_{S_\lambda}Y_\lambda$ and its limit $Z=Z_\alpha\times_{S_\alpha}S=X\times_S Y$; the canonical bijections \sref[I]{I.3.3.14} give commutative diagrams
\[
\xymatrix{
  \operatorname{Hom}_{S_\lambda}(X_\lambda,Y_\lambda) \ar[r]^-{e_{\mu\lambda}} \ar[d]^{\wr} & \operatorname{Hom}_{S_\mu}(X_\mu,Y_\mu) \ar[r]^-{e_\mu} \ar[d]^{\wr} & \operatorname{Hom}_S(X,Y) \ar[d]^{\wr}\\
  \operatorname{Hom}_{X_\lambda}(X_\lambda,Z_\lambda) \ar[r]_-{e_{\mu\lambda}} & \operatorname{Hom}_{X_\mu}(X_\mu,Z_\mu) \ar[r]_-{e_\mu} & \operatorname{Hom}_X(X,Z)
}
\]
and consequently we are reduced to proving that \eqref{IV.8.8.1.1} is injective in the particular case where $S_\alpha=X_\alpha$ (taking into account \sref[I]{I.1.3.4}).
Moreover, as $X_\alpha$ is quasi-compact, hence a finite union of open affines, we may suppose $X_\alpha$ affine (since $L$ is filtered). Suppose therefore given two $X_\alpha$-morphisms $f'_\alpha:X_\alpha\to Y_\alpha$, $f''_\alpha:X_\alpha\to Y_\alpha$ such that $f'$ and $f''$ are equal $X$-morphisms from $X$ to $Y$; we must prove that $f'_\mu=f''_\mu$ for $\mu\geqslant\alpha$ sufficiently large. As $X_\alpha$ is quasi-compact, so is $f'_\alpha(X_\alpha)\cup f''_\alpha(X_\alpha)$, and as $Y_\alpha$ is locally of finite type over $X_\alpha$, this union is contained in a finite union of affine opens $V_{i\alpha}$ of $Y_\alpha$, of finite type over $X_\alpha$. Set $U'_{i\alpha}=f_\alpha^{'-1}(V_{i\alpha})$, $U''_{i\alpha}=f_\alpha^{''-1}(V_{i\alpha})$, $U_{i\alpha}=U'_{i\alpha}\cap U''_{i\alpha}$, and $U_\alpha=\bigcup_iU_{i\alpha}$. The hypothesis $f'=f''$ implies $v_\alpha^{-1}(U'_{i\alpha})=v_\alpha^{-1}(U''_{i\alpha})$, these two sets being respectively equal to $f^{'-1}(w_\alpha^{-1}(V_{i\alpha}))$ and $f^{''-1}(w_\alpha^{-1}(V_{i\alpha}))$. Since $\bigcup_i V_{i\alpha}$ contains $f'_\alpha(X_\alpha)\cup f''_\alpha(X_\alpha)$, one has $v_\alpha^{-1}(U_\alpha)=X$. Since $X_\alpha$ is quasi-compact and every open subset of $X_\lambda$ is ind-constructible, it follows from \sref{IV.8.3.4} that there is an index $\mu\geqslant\alpha$ such that the $U_{i\mu}=v_{\alpha\mu}^{-1}(U_{i\alpha})$ form a covering of $X_\mu$. Replacing $\alpha$ by $\mu$, we may therefore suppose that the $U_{i\alpha}$ cover $X_\alpha$; this implies that for every $x\in X_\alpha$, there is an affine open neighbourhood $W(x)$ contained in one of the $U_{i\alpha}$, in other words such that the restrictions of $f'_\alpha$ and $f''_\alpha$ to $W(x)$ map $W(x)$ into the same affine open $V_{i\alpha}$. Since $X_\alpha$ is quasi-compact, it is covered by a finite number of affine opens $W(x_k)$; by virtue of Lemma~\sref{IV.8.8.2.1} and the fact that $L$ is filtered, there exists $\lambda\geqslant\alpha$ such that the restrictions of $f'_\lambda$ and $f''_\lambda$ to each of the opens $v_{\alpha\lambda}^{-1}(W(x_k))$ are equal, whence $f'_\lambda=f''_\lambda$.

Suppose now $X_\alpha$ quasi-compact and quasi-separated and $Y_\alpha$ locally of finite presentation over $S_\alpha$, and let us prove that \eqref{IV.8.8.1.1} is surjective.
Suppose therefore given an $X$-morphism $f:X\to Y$.
As $X$ is quasi-compact, so is
\oldpage[IV-3]{31}
$f(X)$, and consequently there is a quasi-compact open $Y'$ in $Y$ which contains $f(X)$; there therefore exist $\lambda\geqslant\alpha$ and a quasi-compact open $Y'_\lambda$ in $Y_\lambda$ such that $Y'=w_\lambda^{-1}(Y'_\lambda)$ \sref{IV.8.2.11}.
Replacing as needed $\alpha$ by $\lambda$ and $Y_\lambda$ by $Y'_\lambda$, we can therefore restrict to the case where $Y_\alpha$ is \emph{quasi-compact}, hence also $Y$ and the $Y_\lambda$.
As $Y$ is quasi-compact, it is a finite union of affine opens $V_i$, and consequently $X$ is a union of the opens $f^{-1}(V_i)$. Since every point of $X$ has a quasi-compact open neighbourhood contained in one of the $f^{-1}(V_i)$ and $X$ is quasi-compact, we may, taking account of \sref{IV.8.2.11} and replacing $\alpha$ if necessary by an index $\lambda\geqslant\alpha$, suppose that $Y$ is a finite union of opens $V_i=w_\alpha^{-1}(V_{i\alpha})$, where the $V_{i\alpha}$ are affine opens of $Y_\alpha$; consequently $X$ is a union of the opens $f^{-1}(V_i)$. Since every point of $X$ has a quasi-compact open neighbourhood contained in one of the $f^{-1}(V_i)$ and $X$ is quasi-compact, we can cover $X$ by a finite number of such neighbourhoods and, repeating some of the $V_i$ if necessary, suppose that there is a covering $(U_i)$ of $X$ by quasi-compact opens having the same index set as $(V_i)$ and such that $f(U_i)\subset V_i$ for every $i$. Moreover, using \sref{IV.8.2.11} (and replacing $\alpha$ if necessary by an index $\lambda\geqslant\alpha$), we may suppose that $U_i=v_\alpha^{-1}(U_{i\alpha})$, where the $U_{i\alpha}$ are quasi-compact opens in $X_\alpha$; furthermore, using \sref{IV.8.3.4} as above, we may suppose that $(U_{i\alpha})$ is a \emph{covering} of $X_\alpha$.

It will then suffice to show that there exist $\lambda\geqslant\alpha$ and for each $i$ a morphism $f_{i\lambda}:U_{i\lambda}\to V_{i\lambda}$ (where $U_{i\lambda}=v_{\alpha\lambda}^{-1}(U_{i\alpha})$ and $V_{i\lambda}=w_{\alpha\lambda}^{-1}(V_{i\alpha})$) such that the corresponding morphism $f_i=e_\lambda(f_{i\lambda}):U_i\to V_i$ is equal to the \emph{restriction} of $f$ to $U_i$.
Indeed, if this is so, as $X_\lambda$ is quasi-separated \sref[I]{I.1.2.3}, the $U_{i\lambda}\cap U_{j\lambda}$ are \emph{quasi-compact} and the uniqueness result proved above (which applies since $Y_\lambda$ is locally of finite type over $S_\lambda$) proves that there exists $\mu\geqslant\lambda$ such that $f_{i\mu}$ and $f_{j\mu}$ coincide on $U_{i\mu}\cap U_{j\mu}$ for every pair $(i,j)$, and consequently define an $S_\mu$-morphism $f_\mu:X_\mu\to Y_\mu$ such that $e_\mu(f_\mu)=f$.

We are thus reduced to the case where $Y_\alpha$ is \emph{affine}, and as in addition we may suppose that the images of the $V_{i\alpha}$ in $S_\alpha$ are contained in affine opens, we may also suppose that $S_\alpha$ is affine; thus let $S_\alpha=\operatorname{Spec}(A_\alpha)$, $Y_\alpha=\operatorname{Spec}(C_\alpha)$, with $C_\alpha$ an $A_\alpha$-algebra of finite presentation, and let $S=\operatorname{Spec}(A)$, $Y=\operatorname{Spec}(C)$, where $A=\varinjlim A_\lambda$ and $C=\varinjlim(C_\alpha\otimes_{A_\alpha}A_\lambda)=C_\alpha\otimes_{A_\alpha}A$.
One has then
\[
  \operatorname{Hom}_S(X,Y)=\operatorname{Hom}_{A\text{-alg.}}(C,\Gamma(X,\sh{O}_X))=\operatorname{Hom}_{A_\alpha\text{-alg.}}(C_\alpha,\Gamma(X,\sh{O}_X))
\]
\sref[I]{I.2.2.4} and similarly
\[
  \operatorname{Hom}_{S_\lambda}(X_\lambda,Y_\lambda)=\operatorname{Hom}_{A_\lambda\text{-alg.}}(C_\alpha\otimes_{A_\alpha}A_\lambda,\Gamma(X_\lambda,\sh{O}_{X_\lambda}))=\operatorname{Hom}_{A_\alpha\text{-alg.}}(C_\alpha,\Gamma(X_\lambda,\sh{O}_{X_\lambda})).
\]
But as $X_\alpha$ is quasi-compact and quasi-separated, one knows \sref{IV.8.5.4} that $\varinjlim\Gamma(X_\lambda,\sh{O}_{X_\lambda})=\Gamma(X,\sh{O}_X)$; since $C_\alpha$ is an $A_\alpha$-algebra of finite presentation, the fact that \eqref{IV.8.8.1.1} is bijective then follows from \sref{IV.8.8.2.3}.

Before passing to the proof of (ii), let us note the following corollaries of (i):
\end{proof}

\begin{corollary}[8.8.2.4]
\label{IV.8.8.2.4}
Suppose $S_0$ quasi-compact, $X_\alpha$ of finite presentation over $S_\alpha$, $Y_\alpha$ of finite type over $S_\alpha$ and quasi-separated over $S_\alpha$ (which is the case, for example, if $Y_\alpha$ is also
\oldpage[IV-3]{32}
of finite presentation over $S_\alpha$).
Let $f_\alpha:X_\alpha\to Y_\alpha$ be an $S_\alpha$-morphism.
For $f:X\to Y$ to be an isomorphism, it is necessary and sufficient that there exist $\lambda\geqslant\alpha$ such that $f_\lambda:X_\lambda\to Y_\lambda$ is an isomorphism.
\end{corollary}

\begin{proof}
The condition is evidently sufficient. To prove that it is necessary, first observe that, since the question is local on $S_0$ (because $L$ is filtered), we may always suppose $S_0$ affine and hence quasi-separated. By hypothesis there is an $S$-morphism $g:Y\to X$ such that $g\circ f=1_X$ and $f\circ g=1_Y$. Since $X_\alpha$ is of finite presentation over $S_\alpha$, and $Y_\alpha$ is quasi-compact and quasi-separated (because $S_\alpha$ is quasi-compact and quasi-separated), it follows from \sref{IV.8.8.2},~(i) that there exist $\mu\geqslant\alpha$ and an $S_\mu$-morphism $g_\mu:Y_\mu\to X_\mu$ such that $g=g_\mu\times 1_S$. It follows moreover from \sref{IV.8.8.2},~(i), together with the relations $g\circ f=1_X$ and $f\circ g=1_Y$, that there exists $\nu\geqslant\mu$ such that $g_\nu\circ f_\nu=1_{X_\nu}$ and $f_\nu\circ g_\nu=1_{Y_\nu}$, since $X_\alpha$ and $Y_\alpha$ are of finite type over $S_\alpha$ and quasi-compact. Thus $f_\nu:X_\nu\to Y_\nu$ is an isomorphism, proving the corollary.
\end{proof}

\begin{corollary}[8.8.2.5]
\label{IV.8.8.2.5}
Suppose $S_0$ quasi-compact and quasi-separated, $X_\alpha$ and $Y_\alpha$ of finite presentation over $S_\alpha$.
For $X$ and $Y$ to be $S$-isomorphic, it is necessary and sufficient that there exist $\lambda\geqslant\alpha$ such that $X_\lambda$ and $Y_\lambda$ are $S_\lambda$-isomorphic. Moreover, for every $S$-isomorphism $f:X\xrightarrow{\sim}Y$, there exist $\mu\geqslant\lambda$ and an $S_\mu$-isomorphism $f_\mu:X_\mu\xrightarrow{\sim}Y_\mu$ such that $f=f_\mu\times 1_S$.
\end{corollary}

\begin{proof}
The condition is evidently sufficient. Conversely, if there exists an $S$-isomorphism $f:X\to Y$, it follows from \sref{IV.8.8.2},~(i) that $f$ is of the form $f_\mu\times 1_S$ for some $\mu\geqslant\alpha$ and some $S_\mu$-morphism $f_\mu:X_\mu\to Y_\mu$. Since $f$ is an isomorphism, it follows from \sref{IV.8.8.2.4} that there exists $\nu\geqslant\mu$ such that $f_\nu:X_\nu\to Y_\nu$ is an isomorphism.
\end{proof}

\begin{proof}[Proof of \sref{IV.8.8.2}, (ii)]
Consider first the case where $S_0=\operatorname{Spec}(A_0)$ and $X=\operatorname{Spec}(B)$ are affine; then assertion~(ii) is equivalent to Lemma~\sref{IV.1.8.4.2}.

To prove (ii) in the general case, note that $S$ is quasi-compact and quasi-separated. Since $X$ is of finite presentation over $S$, and $S$ is affine over $S_0$, there is a finite open cover $(U_i)$ of $S_0$ and, if $(W_i)$ denotes the affine open cover of $S$ formed by the inverse images of the $U_i$ under $S\to S_0$, a finite cover $(X_r)$ of $X$ by affine opens such that the image of each $X_r$ under $X\to S$ is contained in some $W_{i(r)}$. The ring $A(X_r)$ is then an algebra of finite presentation over $A(W_{i(r)})$ \sref{IV.1.4.6}.

By Lemma~\sref{IV.1.8.4.2} and the fact that $L$ is filtered, there exist an index $\lambda\in L$ and, for every $r$, an affine scheme $Z_{r\lambda}$ of finite presentation over the inverse image $W_{i(r),\lambda}$ of $U_{i(r)}$ in $S_\lambda$, together with an $S$-isomorphism
\[
 g_r:Z_{r\lambda}\times_{S_\lambda}S\xrightarrow{\sim}X_r.
\]
Let $Z_{rs}$ be the inverse image under $g_r$ of the sub-prescheme induced on the open $X_r\cap X_s$ of $X_r$, which is quasi-compact because $X$ is quasi-separated, and let
\[
 g'_{rs}:Z_{rs}\xrightarrow{\sim}X_r\cap X_s
\]
be the restriction of $g_r$. By \sref{IV.8.2.11}, there exist $\mu\geqslant\lambda$ and, for every pair $(r,s)$, a quasi-compact open $Z_{rs\mu}$ of $Z_{r\mu}=v_{\lambda\mu}^{-1}(Z_{r\lambda})$ such that $Z_{rs}$ is the inverse image of $Z_{rs\mu}$. Moreover, since $S_\mu$ is quasi-separated and $W_{i(r),\mu}$ is a quasi-compact open of $S_\mu$, each $Z_{rs\mu}$ is of finite presentation over $S_\mu$ \sref[I]{I.1.6.2}.

For every pair $(r,s)$, consider the isomorphism
\[
 h_{sr}=(g'_{sr})^{-1}\circ g'_{rs}:Z_{rs}\xrightarrow{\sim}Z_{sr}.
\]
By \sref{IV.8.8.2.4}, there exist $\nu\geqslant\mu$ and, for every pair $(r,s)$, an isomorphism
\[
 h_{sr\nu}:Z_{rs\nu}\xrightarrow{\sim}Z_{sr\nu}
\]
such that $h_{sr}=h_{sr\nu}\times 1_S$. Finally, for every triple $(r,s,t)$, let $h'_{sr}$ denote the restriction of $h_{sr}$ to

\oldpage[IV-3]{33}

$Z_{rs}\cap Z_{rt}$. It follows immediately from the definitions that $h'_{sr}$ is an isomorphism from $Z_{rs}\cap Z_{rt}$ onto $Z_{sr}\cap Z_{st}$ and that
\[
 h'_{ts}\circ h'_{sr}=h'_{tr}.
\]
By \sref{IV.8.8.2},~(i), there therefore exists $\rho\geqslant\nu$ such that, for every triple $(r,s,t)$,
\[
 h'_{ts\rho}\circ h'_{sr\rho}=h'_{tr\rho}.
\]
Consequently the isomorphisms $h_{sr\rho}$ satisfy the gluing condition \sref[0\textsubscript{I}]{0.4.1.7} and hence define a prescheme $X_\rho$ such that $X$ is isomorphic to $X_\rho\times_{S_\rho}S$. Moreover, the $Z_{r\rho}$ are of finite presentation over $S_\rho$, and, identifying them with open subsets of $X_\rho$, their intersections $Z_{r\rho}\cap Z_{s\rho}$, which are isomorphic to the $Z_{rs\rho}$, are quasi-compact. Thus $X_\rho$ is of finite presentation over $S_\rho$ \sref[I]{I.1.6.3}.
\end{proof}

\begin{env}[8.8.3]
\label{IV.8.8.3}
The essential content of \sref{IV.8.8.2} is expressed again by saying that if $S_0$ is quasi-compact and quasi-separated, the category of $S$-preschemes of finite presentation is determined up to equivalence by the data of the categories of $S_\lambda$-preschemes of finite presentation, the functors $\rho_{\mu\lambda}:X_\lambda\mapsto X_\lambda\times_{S_\lambda}S_\mu$ (for $\lambda\leqslant\mu$) between these categories and the transitivity isomorphisms $\rho_{\nu\mu}\circ\rho_{\mu\lambda}\xrightarrow{\sim}\rho_{\nu\lambda}$.
Figuratively, one can say that the data of an $S$-prescheme of finite presentation $X$ is ``functorially'' equivalent to the data of an $S_\lambda$-prescheme of finite presentation $X'_\lambda$; if an $S_\mu$-prescheme of finite presentation $X''_\mu$ is such that $X$ is $S$-isomorphic to $X''_\mu\times_{S_\mu}S$, there exists $\nu$ such that $\lambda\leqslant\nu$, $\mu\leqslant\nu$, and $X'_\lambda\times_{S_\lambda}S_\nu$ and $X''_\mu\times_{S_\mu}S_\nu$ are $S_\nu$-isomorphic.
The fact that $L$ is filtered implies in addition that if one gives oneself a finite \emph{family} $(X^{(i)})_{i\in I}$ of $S$-preschemes of finite presentation, and a finite family $(f^{(j)})_{j\in J}$ of $S$-morphisms between these preschemes ($f^{(j)}$ being therefore of the form $X^{(\sigma(j))}\to X^{(\tau(j))}$ where $\sigma$ and $\tau$ are two maps from $J$ to $I$), then there is an index $\lambda\in L$, a family $(X^{(i)}_\lambda)_{i\in I}$ of $S_\lambda$-preschemes of finite presentation and a family $(f^{(j)}_\lambda)_{j\in J}$ of $S_\lambda$-morphisms such that $X^{(i)}$ is identified with $X^{(i)}_\lambda\times_{S_\lambda}S$ and $f^{(j)}$ with $f^{(j)}_\lambda\times 1_S$ for all $i$ and $j$; furthermore, if one has a relation $f^{(j)}=f^{(k)}$, one can suppose $\lambda$ chosen such that $f^{(j)}_\lambda=f^{(k)}_\lambda$.

Consider in particular three $S$-preschemes of finite presentation $X$, $Y$, $Z$ and two $S$-morphisms $f:X\to Z$, $g:Y\to Z$, so that the product $X\times_Z Y$ relative to these morphisms is also an $S$-prescheme of finite presentation \sref[I]{I.1.6.2}.
It then results from the preceding and from \sref[I]{I.3.3.11} that one can write $X\times_Z Y=(X_\lambda\times_{Z_\lambda}Y_\lambda)\times_{S_\lambda}S$ for a convenient $\lambda$, $X_\lambda$, $Y_\lambda$, $Z_\lambda$ being $S_\lambda$-preschemes of finite presentation; one can therefore say that the determination of $S$-preschemes of finite presentation by the data of $S_\lambda$-preschemes of finite presentation is ``\emph{compatible with fibre products}''.
We have seen on the other hand \sref{IV.8.6.3} that if $g$ is an immersion, we can suppose the same for $g_\lambda$; identifying $Y$ (resp.\ $Y_\lambda$) with a sub-prescheme of $Z$ (resp.\ $Z_\lambda$), we see that one can write, for a convenient $\lambda$, $f^{-1}(Y)=f_\lambda^{-1}(Y_\lambda)\times_{S_\lambda}S$ \sref[I]{I.4.4.1};
there is therefore also ``\emph{compatibility with the formation of inverse images of sub-preschemes}''.
More particularly, if $f_1$, $f_2$ are two $S$-morphisms from $X$ to $Y$, one calls the \emph{kernel} of this pair of morphisms the inverse image $N$ of the diagonal sub-prescheme of $Y\times_S Y$ by the $S$-morphism $(f_1,f_2)_S$; one deduces from what precedes that one will have for a convenient $\lambda$, $N=N_\lambda\times_{S_\lambda}S$, where $N_\lambda$ is the kernel of the couple of $S_\lambda$-morphisms $(f_{1\lambda},f_{2\lambda})$ of $X_\lambda$ into $Y_\lambda$.
These remarks extend to finite products and to ``kernels'' of arbitrary finite systems of $S$-morphisms from an $S$-prescheme of finite presentation into
\oldpage[IV-3]{34}
another; one concludes that in a general fashion the formation of $S$-preschemes of finite presentation by the data of $S_\lambda$-preschemes of finite presentation is ``\emph{compatible with finite projective limits}'', such a limit being by definition the kernel of a finite system of morphisms from one $S$-prescheme into a product of $S$-preschemes \sref[T]{T.1.8}.

We shall allow ourselves in what follows to freely make the translations implied by the preceding properties (as well as by \sref{IV.8.3.11}, \sref{IV.8.5.2}, and \sref{IV.8.6.3}) without always striving to justify them as above.
For example, the data of a ``group prescheme'' $G$ of finite presentation over $S$ is equivalent to the data of a group prescheme $G_\lambda$ of finite presentation over some $S_\lambda$ for a sufficiently large $\lambda$; indeed, writing the associativity condition for the $S$-morphism ``composition law'' $G\times_S G\to G$ of the group prescheme amounts to writing that the kernel of two $S$-morphisms of the form $G\times_S G\times_S G\to G$ is equal to $G\times_S G\times_S G$ \sref[II]{II.8.3.9}, and the other conditions entering into the definition of a group prescheme are interpreted similarly.
\end{env}

\subsection{First applications to the elimination of Noetherian hypotheses}
\label{subsection:IV.8.9}

\begin{proposition}[8.9.1]
\label{IV.8.9.1}
Let $A$ be a ring, $X$ an $A$-prescheme.
\begin{enumerate}[label=(\roman*)]
\item The following conditions are equivalent:
\begin{enumerate}[label=\emph{\alph*)}]
\item $X$ is of finite presentation over $A$.
\item There exist a Noetherian ring $A_0$, a prescheme $X_0$ of finite type over $A_0$, a ring homomorphism $A_0\to A$, and an $A$-isomorphism $X_0\otimes_{A_0}A\xrightarrow{\sim}X$.
\item There exist a subring $A_0$ of $A$, which is a $\mathbf{Z}$-algebra of finite type, a prescheme $X_0$ of finite type over $A_0$, and an $A$-isomorphism $X_0\otimes_{A_0}A\xrightarrow{\sim}X$.
\end{enumerate}
\item If $\sh{F}$ is a quasi-coherent $\sh{O}_X$-Module of finite presentation, one can suppose the subring $A_0$ such that there exists a coherent $\sh{O}_{X_0}$-Module $\sh{F}_0$ such that $\sh{F}$ is isomorphic to $\sh{F}_0\otimes_{A_0}A$; $\operatorname{Supp}(\sh{F})$ is constructible and closed in $X$, and there is a closed sub-prescheme $Z$ of $X$, having $\operatorname{Supp}(\sh{F})$ as underlying space, such that the canonical immersion $Z\to X$ be of finite presentation.
\item If $Y$ is a second $A$-prescheme of finite presentation, and $f:X\to Y$ an $A$-morphism, one can suppose the subring $A_0$ of $A$ such that there exists a prescheme $Y_0$ of finite type over $A_0$, an $A$-isomorphism $Y_0\otimes_{A_0}A\xrightarrow{\sim}Y$, and an $A_0$-morphism $f_0:X_0\to Y_0$ (necessarily of finite type) such that $f$ is identified with $f_0\otimes 1_A$.
\end{enumerate}
\end{proposition}

\begin{proof}
\begin{enumerate}[label=(\roman*)]
\item As $A$ is the inductive limit of its subrings which are of finite type over $\mathbf{Z}$, \emph{a)} implies \emph{c)} by virtue of \sref{IV.8.8.2},~(ii); \emph{c)} implies \emph{b)} since a $\mathbf{Z}$-algebra of finite type is a Noetherian ring; finally, if $A_0$ is Noetherian, an $A_0$-prescheme of finite type is of finite presentation \sref[I]{I.1.6.1}, hence \emph{b)} implies \emph{a)} by virtue of \sref[I]{I.1.6.2},~(iii).
\end{enumerate}
The assertion (ii) follows immediately from \sref{IV.8.5.2},~(ii), \sref{IV.8.3.11}, and \sref[I]{I.1.8.2}, and the assertion (iii) from \sref{IV.8.8.2},~(i).
\end{proof}

\begin{env}[8.9.2]
\label{IV.8.9.2}
The proposition \sref{IV.8.9.1} and the results of \sref{IV.8.5.2} and \sref{IV.8.8.2} allow one to extend in numerous cases to morphisms of finite presentation $X\to Y$ the results proved by Noetherian techniques by supposing $Y$ \emph{locally Noetherian}.
\oldpage[IV-3]{35}
We shall see numerous examples in what follows and we shall limit ourselves here to giving a few results of this type.
\end{env}

\begin{proposition}[8.9.3]
\label{IV.8.9.3}
Let $A$ be a ring, $M$ an $A$-module of finite presentation.
Then every surjective $A$-endomorphism of $M$ is bijective.
\end{proposition}

\begin{proof}
Indeed, consider $A$ as an inductive limit of its sub-$\mathbf{Z}$-algebras of finite type.
It follows from \sref{IV.8.5.2},~(ii) that there exists one of these sub-algebras $A_0$ and an $A_0$-module $M_0$ of finite presentation such that $M$ is $A$-isomorphic to $M_0\otimes_{A_0}A$; furthermore, if $f:M\to M$ is a surjective $A$-endomorphism, one can suppose \sref{IV.8.5.2},~(i) that there exists an $A_0$-endomorphism $f_0:M_0\to M_0$ such that $f=f_0\otimes 1_A$; finally, by \sref{IV.8.5.7}, one can suppose that $f_0$ is \emph{surjective}.
But as $A_0$ is Noetherian and $M_0$ an $A_0$-module of finite type, $M_0$ is a Noetherian module, hence (Bourbaki, \emph{Alg.}, chap.~VIII, \textsection\,2, n\textsuperscript{o}\,2, lemma~3) $f_0$ is bijective, and it is the same for $f$ by base change.
\end{proof}

\begin{proposition}[8.9.4]
\label{IV.8.9.4}
(``generic flatness theorem''). --- Let $Y$ be an integral prescheme, $u:X\to Y$ a morphism of finite type and locally of finite presentation, $\sh{F}$ an $\sh{O}_X$-Module quasi-coherent of finite presentation.
Then there exists an open non-empty $U$ of $Y$ such that $\sh{F}|u^{-1}(U)$ is flat over $U$.
\end{proposition}

\begin{proof}
The reasoning from the beginning of \sref{IV.6.9.1} leads to proving the

\begin{lemma}[8.9.4.1]
\label{IV.8.9.4.1}
Let $A$ be an integral ring, $B$ an $A$-algebra of finite presentation, $M$ a $B$-module of finite presentation.
Then there exists $f\neq 0$ in $A$ such that $M_f$ is a free $A_f$-module.
\end{lemma}

\begin{proof}
Indeed, by \sref{IV.8.9.1} there is a finitely generated $\mathbf{Z}$-subalgebra $A_0$ of $A$, a finite-type $A_0$-algebra $B_0$, and a finite $B_0$-module $M_0$ such that $B$ is $A$-isomorphic to $B_0\otimes_{A_0}A$ and $M$ is $B$-isomorphic to $M_0\otimes_{A_0}A$; by \sref{IV.6.9.2}, there exists $f_0\neq 0$ in $A_0$ such that $(M_0)_{f_0}$ is a free $(A_0)_{f_0}$-module.
As $M_{f_0}=(M_0)_{f_0}\otimes_{A_0}A$ and $A_{f_0}=(A_0)_{f_0}\otimes_{A_0}A$, $M_{f_0}$ is a free $A_{f_0}$-module.
\end{proof}
\end{proof}

\begin{corollary}[8.9.5]
\label{IV.8.9.5}
Let $S$ be a quasi-compact prescheme and quasi-separated, $u:X\to S$ a morphism of finite presentation, $\sh{F}$ an $\sh{O}_X$-Module quasi-coherent of finite presentation.
Then there exists a partition $(S_i)_{1\leqslant i\leqslant n}$ of $S$ into finitely many locally closed subsets of $S$ such that, for $1\leqslant i\leqslant n$, there is a subprescheme $S'_i$ of $S$, with underlying space $S_i$, which is of finite presentation over $S$, and such that, on putting $X_i=X\times_S S'_i$, the $\sh{O}_{X_i}$-Module $\sh{F}_i=\sh{F}\otimes_{\sh{O}_X}\sh{O}_{X_i}$ is flat over $S'_i$.
\end{corollary}

\begin{proof}
Consider a finite covering $(U_j)_{1\leqslant j\leqslant m}$ of $S$ by affine open subsets and define, inductively, $T_1=U_1$ and
\[
T_j=U_j\setminus\bigcup_{h<j}(U_j\cap U_h)\qquad(j\geqslant 2).
\]
Each $T_j$ is closed in the affine open $U_j$, and the $T_j$ are pairwise disjoint.
Moreover, the $U_j\cap U_h$ are quasi-compact since $S$ is quasi-separated and hence, since $S$ is also quasi-compact, are constructible in $S$ \sref{IV.1.8.1}; consequently each $T_j$ is constructible as well.
It will therefore suffice to prove the corollary for a suitable subprescheme $T'_j$ of $S$ having $T_j$ as underlying space and of finite presentation over $S$, and for the morphism and the Module deduced respectively from $u$ and $\sh{F}$ by the base change $T'_j\to S$.
For this, note that if $U_j$ is given the prescheme structure induced by that of $S$, the open immersion $U_j\to S$ is quasi-compact because $S$ is quasi-separated \sref{IV.1.2.7}, and hence is of
\oldpage[IV-3]{36}
finite presentation \sref{IV.1.6.2},~(i).
It therefore suffices that $T'_j$ be of finite presentation over $U_j$; in other words, we may already reduce to the case where $U_j=S$ and $T_j=T$ is closed and constructible in $S$.

Write $S=\operatorname{Spec}(A)$, and regard $A$ as the inductive limit of its finitely generated $\mathbf{Z}$-subalgebras, so that $S=\varprojlim S_\lambda$, where the $S_\lambda$ are affine and Noetherian.
By \sref{IV.8.3.11}, there exist $\lambda$ and a closed (necessarily constructible) subset $T_\lambda$ of $S_\lambda$ such that $T=u_\lambda^{-1}(T_\lambda)$, where $u_\lambda:S\to S_\lambda$ is the canonical morphism.
Give $T_\lambda$ a subprescheme structure and take $T'=T_\lambda\times_{S_\lambda}S$.
Since the canonical immersion $T_\lambda\to S_\lambda$ is of finite presentation \sref{IV.1.6.1}, so is the immersion $T'\to S$.
As $T'$ is affine, we see finally that it is enough to treat the case where $T'=S$ is affine.

Then, with the notation of \sref{IV.8.9.1}, there exist $\lambda$, a morphism of finite type $u_\lambda:X_\lambda\to S_\lambda$, and a coherent $\sh{O}_{X_\lambda}$-Module $\sh{F}_\lambda$ such that $X$ is isomorphic to $X_\lambda\times_{S_\lambda}S$ and $\sh{F}$ to $\sh{F}_\lambda\otimes_{\sh{O}_{X_\lambda}}\sh{O}_X$.
Applying \sref{IV.6.9.3} to $S_\lambda$, $X_\lambda$, and $\sh{F}_\lambda$, we obtain subpreschemes $S'_{\lambda i}$ of $S_\lambda$ whose underlying sets form a partition of $S_\lambda$ and which are such that, on putting $X_{\lambda i}=X_\lambda\times_{S_\lambda}S'_{\lambda i}$, the Module
\[
\sh{F}_{\lambda i}=\sh{F}_\lambda\otimes_{\sh{O}_{X_\lambda}}\sh{O}_{X_{\lambda i}}
\]
is flat over $S'_{\lambda i}$.
The subpreschemes $S'_i=S'_{\lambda i}\times_{S_\lambda}S$ of $S$ then have the required properties by \sref{IV.2.1.4}.
\end{proof}

\subsection{Properties of permanence of morphisms by passage to the projective limit}
\label{subsection:IV.8.10}

In this section, we preserve the general hypotheses and the notations of \sref{IV.8.8.1}.

\begin{proposition}[8.10.1]
\label{IV.8.10.1}
If there exists $\lambda$ such that, for $\mu\geqslant\lambda$, $f_\mu$ is an open morphism (resp.\ universally open), then $f$ is an open morphism (resp.\ universally open).
\end{proposition}

\begin{proof}
As $X=X_\lambda\times_{Y_\lambda}Y$, the assertion relative to universally open morphisms is a particular case of \sref{IV.2.4.3},~(iii).
Let us therefore suppose $f_\mu$ open for $\mu\geqslant\lambda$; it suffices to see that for every open \emph{quasi-compact} $U$ of $X$, $f(U)$ is open in $Y$.
Now, there exists $\mu\geqslant\lambda$ such that $U=v_\mu^{-1}(U_\mu)$, where $U_\mu$ is an open quasi-compact in $X_\mu$ \sref{IV.2.3.11}; one has then $f(v_\mu^{-1}(U_\mu))=w_\mu^{-1}(f_\mu(U_\mu))$ \sref[I]{I.3.4.8}, hence $f(U)$ is open in $Y$.
\end{proof}

\begin{corollary}[8.10.2]
\label{IV.8.10.2}
Let $f:X\to Y$ be a morphism.
For $f$ to be universally open, it suffices that, for every integer $n>0$, if one sets $Y_n=Y\otimes_\mathbf{Z}\mathbf{Z}[T_1,\ldots,T_n](=\mathbf{V}_Y^n)$, and $X_n=X\times_Y Y_n$, $f_n=f\times 1_{Y_n}:X_n\to Y_n$, the projection $f_n$ be an open morphism.
\end{corollary}

\begin{proof}
To prove that $f$ is universally open, it suffices to prove that it remains so for every affine $U$ of $Y$; by hypothesis, if $U_n=U\otimes_\mathbf{Z}\mathbf{Z}[T_1,\ldots,T_n]$ is the inverse image of $U$ in $Y_n$, the morphism $f_n^{-1}(U_n)\to U_n$, restriction of $f_n$, is open, and one sees that one can restrict to the case where $Y=\operatorname{Spec}(A)$ is affine.
Furthermore, it suffices evidently to show that for every morphism $Y'\to Y$, where $Y'=\operatorname{Spec}(A')$ is also \emph{affine}, $f'=f_{|Y'}$ is open.
Let us suppose first that $A'$ is an $A$-algebra of finite type, hence a \emph{quotient} of a polynomial algebra $A[T_1,\ldots,T_n]$; then $Y'$ is a closed sub-prescheme of $Y_n$ and $f'$ the restriction
\oldpage[IV-3]{37}
of $f_n$ to $f_n^{-1}(Y')$; but for every open $V$ of $X_n$, one has $f_n(V\cap f_n^{-1}(Y'))=f_n(V)\cap Y'$, and as by hypothesis $f_n(V)$ is open in $Y_n$, this shows that $f'$ is also an open morphism.
When $A'$ is arbitrary, it can be considered as an inductive limit of its sub-$A$-algebras of finite type $A'_\lambda$, and the fact that $f'$ is an open morphism follows from the preceding and from \sref{IV.8.10.1}.
\end{proof}

\begin{proposition}[8.10.3]
\label{IV.8.10.3}
Suppose that there exists $\alpha$ such that: $1^\circ$ $S_\alpha$ is quasi-compact; $2^\circ$ the morphisms $X_\alpha\to S_\alpha$, $Y_\alpha\to S_\alpha$ are quasi-compact and the morphism $Y_\alpha\to S_\alpha$ quasi-separated; $3^\circ$ for $\alpha\leqslant\lambda\leqslant\mu$, the morphisms $u_{\lambda\mu}:S_\mu\to S_\lambda$ are flat; $4^\circ$ $\overline{f_\alpha(X_\alpha)}$ is constructible in $Y_\alpha$.
Then, for $f$ to be dominant, it is necessary and sufficient that there exist $\lambda\geqslant\alpha$ such that $f_\lambda$ be dominant.
\end{proposition}

\begin{proof}
The hypotheses imply that $Y_\alpha$ is quasi-compact and that the morphism $f_\alpha$ is quasi-compact \sref{IV.1.2.4}; hence $\overline{f_\alpha(X_\alpha)}=Z_\alpha$ is pro-constructible \sref{IV.1.9.5},~(vii) in $Y_\alpha$.
If one sets $Z_\lambda=w_{\alpha\lambda}^{-1}(Z_\alpha)$ for $\lambda\geqslant\alpha$ and $Z=w_\alpha^{-1}(Z_\alpha)$, we have $Z_\lambda=\overline{f_\lambda(X_\lambda)}$ and $Z=\overline{f(X)}$ \sref[I]{I.3.4.8}, and $Z_\lambda$ is pro-constructible in $Y_\lambda$ \sref{IV.1.9.5},~(vii).
It suffices then to apply \sref{IV.8.3.13} by replacing $S_\lambda$, $Z'_\lambda$, and $Z''_\lambda$ by $Y_\lambda$, $Y_\lambda$, and $Z_\lambda$ respectively.
\end{proof}

\begin{proposition}[8.10.4]
\label{IV.8.10.4}
Suppose that there exists $\alpha$ such that $Y_\alpha$ is quasi-compact and $f_\alpha$ is of finite type and quasi-separated.
For the morphism $f$ to be separated, it is necessary and sufficient that there exist $\lambda\geqslant\alpha$ such that $f_\lambda$ be separated.
\end{proposition}

\begin{proof}
The question being local on $Y_\alpha$ (since $Y_\alpha$ is quasi-compact and $L$ filtered), one can restrict to the case where $Y_\alpha$ is affine, hence quasi-separated, and the hypothesis implies that $X_\alpha$ (hence the $X_\lambda$ and $X$) are quasi-compact and quasi-separated.
Set $X'_\lambda=X_\lambda\times_{Y_\lambda}X_\lambda$ for $\lambda\geqslant\alpha$ and $X'=X\times_Y X$; one has $X'_\lambda=X'_\alpha\times_{Y_\alpha}Y_\lambda$ and $X'=X'_\alpha\times_{Y_\alpha}Y$. The first projection $X'_\alpha\to X_\alpha$ is quasi-compact and quasi-separated by hypothesis \sref{IV.1.2.3},~(iii), hence $X'_\lambda$ is quasi-compact and quasi-separated.
Note now that if we denote by $\Delta_\lambda$ (resp.\ $\Delta$) the diagonal of $X_\lambda\times_{Y_\lambda}X_\lambda$ (resp.\ $X\times_Y X$), it results from \sref[I]{I.5.3.4} and \sref[I]{I.3.4.8} that $\Delta_\mu$ (resp.\ $\Delta$) is the inverse image of $\Delta_\lambda$ by the map $v'_{\lambda\mu}:X'_\mu\to X'_\lambda$ (resp.\ $v'_\lambda:X'\to X'_\lambda$).
On the other hand, $\Delta_\lambda$ is \emph{constructible} in $X'_\lambda$: indeed, since $f_\lambda$ is quasi-separated, the diagonal immersion $X_\lambda\to X'_\lambda$ is quasi-compact, and locally of finite presentation since $f_\lambda$ is of finite type \sref[I]{I.1.4.3} and \sref[I]{I.1.5.4},~(iii); it results then from \sref{IV.1.8.4.1} that $\Delta_\lambda$ is locally constructible, hence constructible since $X'_\lambda$ is quasi-compact.
One can now apply \sref{IV.8.3.12} by replacing $S_\lambda$ and $Z_\lambda$ by $X'_\lambda$ and $\Delta_\lambda$ respectively.
\end{proof}

\begin{theorem}[8.10.5]
\label{IV.8.10.5}
Suppose $S_0$ quasi-compact, $X_\alpha$ and $Y_\alpha$ of finite presentation over $S_\alpha$, and let $f_\alpha:X_\alpha\to Y_\alpha$ be an $S_\alpha$-morphism.
Consider, for a morphism, the property of being:
\begin{enumerate}[label=(\roman*)]
\item an isomorphism;
\item[\rm(i\,\emph{bis})] a monomorphism;
\item an immersion;
\item an open immersion;
\item a closed immersion;
\item separated;
\item surjective;
\oldpage[IV-3]{38}
\item radicial;
\item affine;
\item quasi-affine;
\item finite;
\item quasi-finite;
\item proper.
\end{enumerate}
Then, if $\mathbf{P}$ denotes one of the preceding properties, for $f$ to possess the property $\mathbf{P}$, it is necessary and sufficient that there exist $\lambda\geqslant\alpha$ such that $f_\lambda$ possess the property $\mathbf{P}$ (in which case $f_\mu$ also possesses it for $\mu\geqslant\lambda$).

If $S_0$ is furthermore supposed quasi-separated, the same conclusion holds when $\mathbf{P}$ is the property of being:
\begin{enumerate}[label=(\roman*)]
\setcounter{enumi}{12}
\item projective;
\item quasi-projective.
\end{enumerate}
\end{theorem}

\begin{proof}
The case where $\mathbf{P}$ is one of the properties (i) or (v) was included in the statement only as a reminder; in these cases, the theorem follows from what has been proved respectively in \sref{IV.8.8.2.4} and \sref{IV.8.10.4}.
Furthermore, taking into account \sref[I]{I.5.4.1} and \sref[I]{I.5.3.4}, (v) also results from (iv).
The case (i\,\emph{bis}) follows immediately from (i), by using \sref[I]{I.5.3.8} and by noting (as we have already used in \sref{IV.8.10.4}) that the diagonal $\Delta_f$ is deduced from $\Delta_{f_\lambda}$ by the base change $S\to S_\lambda$.

We will note on the other hand that (vi), (vii), and (xi) are in fact conditions on the \emph{fibres} $f^{-1}(y)$ of the morphisms considered, taking into account the transitivity of the fibres by base change \sref[I]{I.3.6.4}: condition (vi) means indeed that all the fibres must be non-empty, condition (vii) that they must be radicial \sref[I]{I.3.5.8}, and condition (xi) that they must be finite \sref[III]{III.6.2.2} and \sref[III]{III.6.2.3} and \sref[I]{I.6.4.4}, taking into account that $f$ and the $f_\lambda$ are morphisms of finite type by \sref{IV.1.5.4},~(v).
The theorem in these three cases will be yet another consequence of a general result on this type of properties, which will be established in \sref{IV.9.3.3}; we therefore postpone until then the proof of the theorem in the case (xi) (of course, the reader will be able to verify that, except for n\textsuperscript{os}~{\hyperlink{ega4.c04.subsec.8-11}{8.11}} and {\hyperlink{ega4.c04.subsec.8-12}{8.12}}, we will make no use of the theorem in these cases until \sref{IV.9.3.3} and that \sref{IV.8.11} and \sref{IV.8.12} will not be used before \sref{IV.9.3.3}).

For the cases that remain to be proved, one can restrict to proving that the condition of the statement is necessary, all the properties $\mathbf{P}$ considered being invariant by base change (see chap.~I\textsuperscript{er} and II for the references concerning each of these properties).
One can furthermore suppose that $S_\alpha=S_0$ and that $Y_\alpha=S_\alpha$, hence $Y_\lambda=S_\lambda$ for every $\lambda\geqslant\alpha$.
Finally, the properties (i) through (xii) are local on $S_0$, hence as $S_0$ is a finite union of open affines and $L$ filtered, one can restrict to proving them for $S_0=\operatorname{Spec}(A_0)$ \emph{affine} (hence quasi-separated).
We shall denote by $A_\lambda$ (resp.\ $A$) the ring of $S_\lambda$ (resp.\ $S$).

\emph{Cases} (ii), (iii), (iv): Suppose that $f$ is an immersion (resp.\ an open immersion, a closed immersion), and let $X'$ be the sub-prescheme (resp.\ induced on an open, resp.\ closed) of $S$ associated to $f$, which is therefore an $S$-prescheme of finite presentation.
\oldpage[IV-3]{39}
By virtue of \sref{IV.8.6.3}, there then exists $\lambda\geqslant\alpha$ and a sub-prescheme $X'_\lambda$ of $S_\lambda$, of finite presentation over $S_\lambda$, such that $X'$ is isomorphic to $X'_\lambda\times_{S_\lambda}S$.
For every $\mu\geqslant\lambda$, $X'_\mu=X'_\lambda\times_{S_\lambda}S_\mu$ is therefore a sub-prescheme (resp.\ induced on an open, resp.\ closed) of $S_\mu$, and it results therefore from \sref{IV.8.8.2.4} and \sref{IV.8.8.2.5} that there exists $\mu\geqslant\lambda$ and an isomorphism $g_\mu:X_\mu\to X'_\mu$ such that $g=g_\mu\times 1_S$ is the isomorphism $X\to X'$ associated to $f$; whence the conclusion.

\emph{Cases} (vi) \emph{and} (vii): We know \sref{IV.1.8.4.1} that $Z_\alpha=f_\alpha(X_\alpha)$ is constructible in $S_\alpha$; if we set $Z_\lambda=u_{\alpha\lambda}^{-1}(Z_\alpha)$ for $\lambda\geqslant\alpha$ and $Z=u_\alpha^{-1}(Z_\alpha)$, we have $Z_\lambda=f_\lambda(X_\lambda)$ and $Z=f(X)$ \sref[I]{I.3.4.8}.
As $Z_\alpha$ is pro-constructible in $S_\alpha$, and by virtue of \sref{IV.8.3.11} the canonical map $\varinjlim\mathfrak{C}(S_\lambda)\to\mathfrak{C}(S)$ is injective, the relation $f(X)=S$ implies the existence of $\lambda\geqslant\alpha$ such that $f_\lambda(X_\lambda)=S_\lambda$, which proves the theorem in the case (vi).
To prove it in case~(vii), note that the structural morphism $X_\alpha\times_{S_\alpha}X_\alpha\to S_\alpha$ is of finite presentation, since the first projection $X_\alpha\times_{S_\alpha}X_\alpha\to X_\alpha$ is so \sref{IV.1.6.2}. It is therefore enough, by \sref{IV.1.8.7.1}, to apply case~(vi) to the diagonal morphism
\[
  \Delta_{f_\alpha}:X_\alpha\longrightarrow X_\alpha\times_{S_\alpha}X_\alpha,
\]
noting that $\Delta_{f_\lambda}=\Delta_{f_\alpha}\times1_{S_\lambda}$ and $\Delta_f=\Delta_{f_\alpha}\times1_S$ \sref[I]{I.5.3.4} and \sref[I]{I.3.3.11}.

\emph{Cases} (viii) \emph{and} (ix): As $S=\operatorname{Spec}(A)$ is affine, to say that $f:X\to S$ is affine (resp.\ quasi-affine) means that there exists an integer $r$ and a closed immersion (resp.\ an immersion) $j:X\to\mathbf{V}^r_S=\operatorname{Spec}(A[T_1,\ldots,T_r])$ of $S$-preschemes, since $f$ is of finite type and $S$ quasi-compact \sref[II]{II.5.1.9}.
As $\mathbf{V}^r_S=\mathbf{V}^r_{S_0}\times_{S_0}S$, and $\mathbf{V}^r_{S_0}$ is an $S_0$-prescheme of finite presentation, it results from \sref{IV.8.8.2},~(i) applied to the $S$-morphism $j$, that there exists $\lambda$ and an $S_\lambda$-morphism $j_\lambda:X_\lambda\to\mathbf{V}^r_{S_\lambda}$ such that $j=j_\lambda\times 1_S$; applying (iv) (resp.\ (ii)) to $j$, one deduces that there exists $\mu\geqslant\lambda$ such that $j_\mu$ is a closed immersion (resp.\ an immersion); hence $f_\mu$ is affine (resp.\ quasi-affine).

\emph{Case} (x): By hypothesis, $X=\operatorname{Spec}(B)$, where $B$ is an $A$-algebra which is an $A$-module of finite presentation \sref{IV.1.4.7}; it results therefore from \sref{IV.8.5.2},~(ii) that there is $\lambda$ and an $A_\lambda$-module of finite presentation $B_\lambda$ such that the $A$-module $B$ is isomorphic to $B_\lambda\otimes_{A_\lambda}A$.
The structure of $A$-algebra of $B$ is defined by an $A$-homomorphism $m:B\otimes_A B\to B$; as $B\otimes_A B=(B_\lambda\otimes_{A_\lambda}B_\lambda)\otimes_{A_\lambda} A$, there exists by \sref{IV.8.5.2},~(i) $\mu\geqslant\lambda$ and an $A_\mu$-homomorphism $m_\mu:B_\mu\otimes_{A_\mu}B_\mu\to B_\mu$ such that $m=m_\mu\otimes 1$.
Considering the usual diagrams expressing the associativity and the commutativity of $m$, one sees by applying again \sref{IV.8.5.2},~(i) that one can suppose $\nu$ taken large enough so that $m_\nu$ defines on $B_\nu$ an associative and commutative multiplication.
In the same way one can suppose $\nu$ large enough so that $B_\nu$ thus defined admits a unit element.
If $X'_\nu=\operatorname{Spec}(B_\nu)$, it is clear then that $X$ is $S$-isomorphic to $X'_\nu\times_{S_\nu}S$, hence, by virtue of (i), there exists $\rho\geqslant\nu$ such that $X_\rho$ and $X'_\rho$ are $S_\rho$-isomorphic, which completes the proof in this case.

To prove the theorem in the case (xii), we shall first prove the following proposition:
\end{proof}

\begin{proposition}[8.10.5.1]
\label{IV.8.10.5.1}
(Chow's lemma for morphisms of finite presentation). --- Let $\mathrm{A}$ be a ring, let $X$ and $Y$ be two $\mathrm{A}$-preschemes of finite presentation, and let $f:X\to Y$ be a separated $\mathrm{A}$-morphism. Then there exist two $\mathrm{A}$-preschemes $X'$ and $P$ of finite presentation, and $\mathrm{A}$-morphisms $p:P\to Y$, $j:X'\to P$, and $g:X'\to X$, such that the diagram
\oldpage[IV-3]{40}
\[
\begin{tikzcd}[column sep=large,row sep=large]
X' \arrow[r,"j"] \arrow[d,"g"'] & P \arrow[d,"p"] \\
X \arrow[r,"f"'] & Y
\end{tikzcd}
\]
is commutative, and such that: $1^\circ$ $p$ is projective; $2^\circ$ $g$ is projective and surjective; $3^\circ$ $j$ is an open immersion.
\end{proposition}

\begin{proof}
Choose $\mathrm{A}_0\subset\mathrm{A}$, $X_0$, $Y_0$, and $f_0$ as in \sref{IV.8.9.1}, with $Y_0$ Noetherian and $f_0$ of finite type; by \sref{IV.8.10.4}, one may moreover suppose that $f_0$ is separated. Chow's lemma \sref[II]{II.5.6.1} then gives morphisms of finite type $p_0:P_0\to Y_0$, $g_0:X'_0\to X_0$, and $j_0:X'_0\to P_0$ such that
\[
\begin{tikzcd}[column sep=large,row sep=large]
X'_0 \arrow[r,"j_0"] \arrow[d,"g_0"'] & P_0 \arrow[d,"p_0"] \\
X_0 \arrow[r,"f_0"'] & Y_0
\end{tikzcd}
\]
is commutative, $p_0$ is projective, $g_0$ is projective and surjective, and $j_0$ is an open immersion. The assertions follow from the invariance of these properties under base change \sref[II]{II.5.5.5},~(iii), \sref[I]{I.3.5.2}, and \sref[I]{I.4.3.2}.
\end{proof}

\begin{proof}[Continuation of the proof of \sref{IV.8.10.5}]
\emph{Case (xii).} Apply \sref{IV.8.10.5.1} to $f_0:X_0\to S_0$. We obtain a commutative diagram
\[
\begin{tikzcd}[column sep=large,row sep=large]
X'_0 \arrow[r,"j_0"] \arrow[d,"g_0"'] & P_0 \arrow[d,"p_0"] \\
X_0 \arrow[r,"f_0"'] & S_0
\end{tikzcd}
\]
where $p_0$ is projective, $g_0$ is projective and surjective, and $j_0$ is an open immersion. For each $\lambda$ one obtains the analogous diagram in which $p_\lambda=p_0\times1_{S_\lambda}$, $g_\lambda=g_0\times1_{S_\lambda}$, and $j_\lambda=j_0\times1_{S_\lambda}$ have the same respective properties, and similarly for the inverse-limit morphisms $p=p_0\times1_S$, $g=g_0\times1_S$, and $j=j_0\times1_S$. Since $g$ is proper \sref[II]{II.5.5.3}, so is $p\circ j=f\circ g$ \sref[II]{II.5.4.2}; and since $p$ is separated, $j$ is proper, hence a closed immersion. Apply case~(iv) to $j$, noting that $X'_0$ and $P_0$ are $S_0$-preschemes of finite presentation by \sref{IV.8.10.5.1} and \sref{IV.1.6.2}. There exists $\lambda$ such that $j_\lambda$ is a closed immersion, hence proper \sref[II]{II.5.4.2}. Thus
\[
  f_\lambda\circ g_\lambda=p_\lambda\circ j_\lambda
\]
is proper \sref[II]{II.5.5.3} and \sref[II]{II.5.4.2}. Since $g_\lambda$ is surjective and, by the hypothesis on $f$ and case~(v), $f_\lambda$ may be supposed separated, \sref[II]{II.5.4.3} shows that $f_\lambda$ is proper.

\emph{Cases (xiii) and (xiv).} By case~(xii) and \sref[II]{II.5.5.3}, which applies because the $S_\lambda$ are quasi-compact and quasi-separated by \sref{IV.1.7.19}, it is enough to consider the case where $f$ is quasi-projective. Suppose there is an invertible $\sh{O}_X$-Module $\sh{L}$ which is ample for $f$. Since $S_0$ is quasi-compact and quasi-separated, so is $X_0$ \sref{IV.1.2.3}; hence there are $\lambda$ and a quasi-coherent $\sh{O}_{X_\lambda}$-Module $\sh{L}_\lambda$ of finite presentation such that $\sh{L}=v_\lambda^*(\sh{L}_\lambda)$ \sref{IV.8.5.2},~(ii). By \sref{IV.8.5.5}, one may moreover suppose $\sh{L}_\lambda$ invertible. The assertion now follows from the more precise lemma below.

\begin{lemma}[8.10.5.2]
\label{IV.8.10.5.2}
\oldpage[IV-3]{41}
Suppose that $S_0$ is quasi-compact, and let $\sh{L}_\lambda$ be an invertible $\sh{O}_{X_\lambda}$-Module. For $\sh{L}$ to be ample for $f$ (resp.\ very ample for $f$), it is necessary and sufficient that there exist $\mu\geqslant\lambda$ such that $\sh{L}_\mu$ is ample for $f_\mu$ (resp.\ very ample for $f_\mu$).
\end{lemma}

\begin{proof}
The condition is plainly sufficient \sref[II]{II.4.4.10} and \sref[II]{II.4.6.13}. Conversely, since the $S_\lambda$ are quasi-compact and the $f_\lambda$ are of finite type, after replacing $\sh{L}$ by a suitable tensor power, one may reduce to the case where $\sh{L}$ is very ample \sref[II]{II.4.6.11}. The question is local on $S_0$ by \sref[II]{II.4.4.5} and the fact that $L$ is filtered, so one may suppose that $S_0$, and hence $S$, is affine. By \sref[II]{II.4.4.1},~(ii) and \sref[II]{II.4.1.2}, there is an $S$-immersion
\[
  j:X\longrightarrow\mathbf{P}^r_S=P
\]
such that $\sh{L}\simeq j^*(\sh{O}_P(1))$. By \sref{IV.8.8.2},~(i), case~(ii), and \sref[II]{II.4.1.3}, there exist $\mu\geqslant\lambda$ and an immersion $j_\mu:X_\mu\to\mathbf{P}^r_{S_\mu}=P_\mu$ such that $j=j_\mu\times1_S$. Then \sref[II]{II.4.1.3.2} and \sref{IV.8.5.2.5} give $\nu\geqslant\mu$ such that
\[
  \sh{L}_\nu\simeq j_\nu^*(\sh{O}_{P_\nu}(1)),
\]
which shows that $\sh{L}_\nu$ is very ample for $f_\nu$ \sref[II]{II.4.4.2}.
\end{proof}
\end{proof}

\subsection{Application to quasi-finite morphisms}
\label{subsection:IV.8.11}
\label{IV.8.11}

We propose in this section to prove the two following theorems:

\begin{theorem}[8.11.1]
\label{IV.8.11.1}
Let $f:X\to Y$ be a proper morphism, locally of finite presentation, and quasi-finite.
Then the morphism $f$ is finite.
\end{theorem}

\begin{theorem}[8.11.2]
\label{IV.8.11.2}
Let $f:X\to Y$ be a morphism locally of finite presentation, quasi-finite and separated.
Then the morphism $f$ is quasi-affine, and a fortiori quasi-projective.
\end{theorem}

\begin{remark}[8.11.3]
\label{IV.8.11.3}
We will see below that one can reduce, for the proof of \sref{IV.8.11.1} and \sref{IV.8.11.2}, to the case where $Y$ is locally Noetherian; one will note that in this case one will thus obtain another proof of the theorem of Chevalley \sref[III]{III.4.4.2}.
\end{remark}

\begin{env}[8.11.4]
\label{IV.8.11.4}
The hypotheses and the conclusions of \sref{IV.8.11.1} and \sref{IV.8.11.2} are all local on $Y$.
Indeed, this follows from \sref{IV.1.6.1}, \sref{IV.1.2.6}, \sref[II]{II.5.1.1}, \sref[II]{II.5.4.1}, and \sref[II]{II.6.2.2}.
Therefore one can suppose $Y=\operatorname{Spec}(A)$ affine.
One knows that there then exists a subring $A_0$ of $A$, which is a $\bb{Z}$-algebra of finite type, and a morphism of finite type $f_0:X_0\to\operatorname{Spec}(A_0)$ such that $X$ identifies with $X_0\otimes_{A_0}A$ and $f$ with $f_0\times 1$ \sref{IV.8.9.1}.
Furthermore, $A$ can be considered as inductive limit of its subrings containing $A_0$ and which are $\bb{Z}$-algebras of finite type; using the method of \sref{IV.8.1.2},~$c$)) as well as \sref{IV.8.10.5},~(v),~(xi) and (xii)), one sees that it suffices to prove the theorems \sref{IV.8.11.1} and \sref{IV.8.11.2} for $f_0$.
Suppose henceforth $Y$ Noetherian; using now the method of \sref{IV.8.1.2},~$a$)) as well as \sref{IV.8.10.5},~(v),~(ix),~(x),~(xi) and (xii)) one can replace $Y$ by $\operatorname{Spec}(\sh{O}_y)$, where $y$ is an
\oldpage[IV-3]{42}
arbitrary point of $Y$, so one can finally suppose that $Y=\operatorname{Spec}(A)$, where $A$ is a \emph{local Noetherian ring}.
Let $\mathfrak{m}$ be the maximal ideal of $A$, and let $\hat{A}$ be the completion of $A$ for the $\mathfrak{m}$-adic topology. The ring $\hat{A}$ is Noetherian local and $\operatorname{Spec}(\hat{A})\to\operatorname{Spec}(A)$ is faithfully flat and quasi-compact \sref[0]{0.7.3.5}; applying \sref{IV.2.7.1},~(i), (vii), (xiv), (xv) and (xvi), we may therefore reduce to the case where $A$ is \emph{complete}.
It then follows from \sref[II]{II.6.2.6} that $X=X'\sqcup X''$, where $X'$ is a finite $Y$-scheme and $X''$ a quasi-finite $Y$-scheme such that $X''\cap f^{-1}(y)=\varnothing$.
\end{env}

Let us first place ourselves in the hypotheses of \sref{IV.8.11.1}; as $f$ is proper, $X''$, which is closed in $X$, is proper over $Y$ \sref[II]{II.5.4.10}, so $f(X'')$ is closed in $Y$; but $y$ is not contained in $f(X'')$, and moreover $y$ is adherent to every point of $Y$, so $f(X'')=\varnothing$, and consequently $X''$ is empty and $f$ is \emph{finite}.

Let us now place ourselves in the hypotheses of \sref{IV.8.11.2} and, restricting ourselves further (as one can do from what precedes) to the case where $Y=\operatorname{Spec}(A)$ is affine and Noetherian of finite dimension, let us moreover reason by induction on the dimension of $Y$.
Reducing as above to the case where $A$ is furthermore local and complete, one has $\dim(\sh{O}_y)=\dim(A)=\dim(Y)$ and for all $z\neq y$, $\dim(\sh{O}_z)<\dim(\sh{O}_y)$, so $\dim(Y-\{y\})<\dim(Y)$.
Now, by hypothesis one has $f(X'')\subset Y-\{y\}$ and the restriction of $f$ to $X''$ is evidently a quasi-finite and separated morphism; applying the hypothesis of induction to $Y-\{y\}$ and to $X''$, one sees that $X''$ is quasi-affine over $Y-\{y\}$; but the open subset $Y-\{y\}$ being quasi-affine over $Y$ since $Y$ is Noetherian, $X''$ is also quasi-affine over $Y$ \sref[II]{II.5.1.10},~(ii)); as moreover $X'$ is finite (and \emph{a fortiori} affine) over $Y$, $X$ is quasi-affine over $Y$ \sref[II]{II.4.6.17} and \sref[II]{II.5.1.2},~$c'$)).

\begin{proposition}[8.11.5]
\label{IV.8.11.5}
Let $f:X\to Y$ be a morphism of finite presentation.
The following properties are equivalent:
\begin{enumerate}[label=\emph{\alph*)}]
  \item $f$ is a closed immersion.
  \item $f$ is a proper monomorphism.
  \item $f$ is proper and for all $y\in Y$, $f^{-1}(y)$ is radicial and geometrically reduced over $k(y)$ (that is to say empty or $k(y)$-isomorphic to $\operatorname{Spec}(k(y))$).
\end{enumerate}
\end{proposition}

It is clear that \emph{a)} implies \emph{b)}.
To see that \emph{b)} implies \emph{c)}, let us note \sref[I]{I.3.3.12} that for all $y\in Y$, the morphism $f^{-1}(y)\to\operatorname{Spec}(k(y))$ deduced from $f$ by extension of the base is a monomorphism, hence injective, and consequently $f^{-1}(y)$ is empty or reduced to a point, and in every case affine.
Furthermore, if $A$ is the ring of $f^{-1}(y)$, the canonical homomorphism $A\otimes_k A\to A$ is bijective \sref[I]{I.5.3.8}.
This evidently implies that $A=k$, for otherwise there would exist an element $a\in A$ not in $k$ and the images of $a\otimes 1$ and of $1\otimes a$ in $A$ would both be equal to $a$, whereas $a\otimes 1\neq1\otimes a$ since $1$ and $a$ form a free system over $k$.

It remains to prove that \emph{c)} implies \emph{a)}.
It follows first from \sref{IV.8.11.1} that $f$ is a \emph{finite} morphism, so $X=\operatorname{Spec}(\sh{A})$, where $\sh{A}$ is a finite $\sh{O}_Y$-Algebra.
It suffices therefore to prove that the canonical homomorphism $\sh{O}_Y\to\sh{A}$ is \emph{surjective} \sref[II]{II.1.4.10}, or equivalently that for all $y\in Y$, the homomorphism $\sh{O}_y\to\sh{A}_y$ is surjective.
But by hypothesis
\oldpage[IV-3]{43}
$f^{-1}(y)=\operatorname{Spec}(\sh{A}_y/\mathfrak{m}_y\sh{A}_y)$ \sref[II]{II.1.5.5} is such that the corresponding homomorphism $k(y)=\sh{O}_y/\mathfrak{m}_y\to\sh{A}_y/\mathfrak{m}_y\sh{A}_y$ is surjective; as $\sh{A}_y$ is an $\sh{O}_y$-module of finite type, the lemma of Nakayama shows that $\sh{O}_y\to\sh{A}_y$ is surjective, which completes the proof.

\begin{remark}[8.11.5.1]
\label{IV.8.11.5.1}
One will note that the preceding reasoning proves that if $f$ is a \emph{monomorphism}, then, for all $y\in Y$, $f^{-1}(y)$ is empty or $k(y)$-isomorphic to $\operatorname{Spec}(k(y))$.
\end{remark}

\begin{proposition}[8.11.6]
\label{IV.8.11.6}
If a morphism $f:X\to Y$ of finite presentation is a universal homeomorphism, it is finite, surjective and radicial (the converse being true by \sref{IV.2.4.5},~(i))).
\end{proposition}

\begin{proof}
Indeed, $f$ being of finite type, universally closed, and separated by virtue of \sref{IV.2.4.4}, is proper by definition \sref[II]{II.5.4.1}.
As it is evidently quasi-finite \sref[II]{II.6.2.3}, it is finite by \sref{IV.8.11.1}.
One knows moreover that it is radicial \sref{IV.2.4.4}, and evidently surjective.
\end{proof}

\subsection{New proof and generalisation of the ``Main Theorem'' of Zariski}
\label{subsection:IV.8.12}
\label{IV.8.12}

\begin{lemma}[8.12.1]
\label{IV.8.12.1}
Let $f:X\to Y$ be a quasi-compact and quasi-separated morphism, $\sh{C}$ an $\sh{O}_Y$-Algebra quasi-coherent, $Z=\operatorname{Spec}(\sh{C})$, which is an affine $Y$-prescheme over $Y$.
Let $g:X\to Z$ be a $Y$-morphism, $\varphi=\sh{A}(g):\sh{C}\to f_*(\sh{O}_X)$ the corresponding $\sh{O}_Y$-homomorphism of $\sh{O}_Y$-Algebras \sref[II]{II.1.2.7}.
Suppose that $g$ is an immersion.
Then, for the closed image of $X$ by $g$ \sref[I]{I.9.5.3} to equal $Z$, it is necessary and sufficient that $\varphi$ be injective; $g$ is then an open immersion.
\end{lemma}

\begin{proof}
The hypothesis implies that $f_*(\sh{O}_X)$ is an $\sh{O}_Y$-Algebra quasi-coherent \sref{IV.1.7.5}; furthermore, as the canonical morphism $h:Z\to Y$ is affine, hence quasi-compact and separated, $g$ is a quasi-compact and quasi-separated morphism \sref{IV.1.2.2} and \sref{IV.1.1.2}, hence $g_*(\sh{O}_X)$ is an $\sh{O}_Z$-Algebra quasi-coherent \sref{IV.1.7.5}.
This being so, saying that the closed image of $X$ by $g$ equals $Z$ signifies \sref[I]{I.9.5.2} that the canonical homomorphism $\theta:\sh{O}_Z\to g_*(\sh{O}_X)$ is injective.
But one has $h_*(\sh{O}_Z)=\sh{C}$ by definition of $Z$ \sref[II]{II.1.3.1}, and $h_*(g_*(\sh{O}_X))=f_*(\sh{O}_X)$.
As $Z$ is affine over $Y$, it amounts to the same to say that the homomorphism $\theta:\sh{O}_Z\to g_*(\sh{O}_X)$ is injective or that the corresponding homomorphism $\varphi=h_*(\theta):\sh{C}\to f_*(\sh{O}_X)$ is injective \sref[I]{I.1.3.9}.
The fact that $g$ is then an open immersion follows from \sref[I]{I.9.5.10} and the hypothesis.
\end{proof}

\begin{lemma}[8.12.2]
\label{IV.8.12.2}
Let $Y$ be a prescheme quasi-compact and quasi-separated, $f:X\to Y$ a quasi-separated morphism of finite type, $\sh{C}$ an $\sh{O}_Y$-Algebra quasi-coherent, $Z=\operatorname{Spec}(\sh{C})$.
Let $g:X\to Z$ be a $Y$-morphism, $\varphi:\sh{C}\to f_*(\sh{O}_X)$ the corresponding $\sh{O}_Y$-homomorphism of $\sh{O}_Y$-Algebras.
Let $(\sh{C}_\lambda)$ be the filtered increasing family of sub-$\sh{O}_Y$-Algebras quasi-coherent of finite type of $\sh{C}$ (where $\sh{C}$ is the union \sref[I]{I.9.6.6} and \sref{IV.1.7.9}); set $Z_\lambda=\operatorname{Spec}(\sh{C}_\lambda)$ and let $g_\lambda$ be the morphism composed $X\xrightarrow{g}Z\to Z_\lambda$.
Then the following conditions are equivalent:
\begin{enumerate}[label=\emph{\alph*)}]
  \item $g$ is an immersion.
  \item There exists $\lambda$ such that $g_\lambda$ is an immersion.
\end{enumerate}
\oldpage[IV-3]{44}
Furthermore, when $g_\lambda$ is an immersion, the same is true of $g_\mu$ for $\mu\geqslant\lambda$.

It suffices to apply \sref[II]{II.3.8.4} by replacing $\sh{L}$ by $\sh{O}_X$ and $\sh{S}$ by $\sh{C}[\mathrm{T}]$, and taking into account \sref[II]{II.3.1.7}.
\end{lemma}

\begin{proposition}[8.12.3]
\label{IV.8.12.3}
Let $Y$ be a prescheme quasi-compact and quasi-separated, $f:X\to Y$ a separated morphism of finite type.
Let $\sh{B}=f_*(\sh{O}_X)$, which is an $\sh{O}_Y$-Algebra quasi-coherent \sref[I]{I.9.2.2}; let $\sh{C}$ be the integral closure of $\sh{O}_Y$ in $\sh{B}$, which is an $\sh{O}_Y$-Algebra quasi-coherent \sref[II]{II.6.3.4}; set $Z=\operatorname{Spec}(\sh{C})$, and let $g:X\to Z$ be the $Y$-morphism corresponding to the canonical injection $\varphi:\sh{C}\to\sh{B}=f_*(\sh{O}_X)$ \sref[II]{II.1.2.7}. Let $(\sh{C}_\lambda)$ be the filtered increasing family of quasi-coherent sub-$\sh{O}_Y$-Algebras of finite type of $\sh{C}$ (whose union is $\sh{C}$ \sref[I]{I.9.6.6} and \sref{IV.1.7.9}), and for all $\lambda$, let $g_\lambda:X\to Z_\lambda=\operatorname{Spec}(\sh{C}_\lambda)$ be the $Y$-morphism corresponding to the injection $\sh{C}_\lambda\to\sh{B}=f_*(\sh{O}_X)$.
Then the following conditions are equivalent:
\begin{enumerate}[label=\emph{\alph*)}]
  \item[\emph{a)}] There exists a factorisation of $f$ as
\[
X\xrightarrow{f'}Y'\xrightarrow{u}Y
\]
where $f'$ is an immersion and $u$ a finite morphism.
  \item[\emph{a$'$)}] There exists a factorisation $X\xrightarrow{f'}Y'\xrightarrow{u}Y$ of $f$, where $f'$ is an open immersion and $u$ a finite morphism.
  \item[\emph{b)}] The morphism $g:X\to Z$ is an immersion.
  \item[\emph{c)}] There exists $\lambda$ such that $g_\lambda:X\to Z_\lambda$ is an immersion.
\end{enumerate}
Furthermore, when $g_\lambda$ is an immersion, $g$ is an open immersion, $g(X)$ is dense in $Z$, and there exists $\lambda$ such that for $\mu\geqslant\lambda$, $g_\mu$ is an open immersion.
\end{proposition}

As the homomorphism $\varphi:\sh{C}\to f_*(\sh{O}_X)$ is injective, it follows from \sref{IV.8.12.1} that if $g$ is an immersion, it is an open immersion and $g(X)$ is dense in $Z$, and the same is true for $g_\lambda$.
The fact that \emph{a)} implies \emph{a$'$)} follows also from \sref{IV.8.12.1}, applied by replacing $Z$ by $Y'$ and $g$ by $f'$ (since $Y'$ is finite over $Y$, hence affine over $Y$): indeed if $Y''$ is the closed image of $X$ by $f'$, $Y''$ is finite over $Y$ and $f'$ factorises as $X\xrightarrow{f''}Y''\xrightarrow{j}Y'$, where $j$ is the canonical injection, and $f''$ is an immersion \sref[I]{I.4.1.10}; it follows therefore from \sref{IV.8.12.1} that $f''$ is an open immersion.

The equivalence of \emph{b)} and \emph{c)} follows from \sref{IV.8.12.2} as well as the fact that $g_\lambda$ is then an immersion for $\lambda$ large enough.
It is clear that \emph{c)} implies \emph{a)}, since $Z_\lambda$ is finite over $Y$ \sref[II]{II.6.3.4} and \sref[II]{II.6.1.2}.
Let us finally show that \emph{a)} implies \emph{c)}.
One has seen above that one can suppose that $Y'$ is the closed image of $X$ by $f'$, and it then follows from \sref{IV.8.12.1} that if one sets $\sh{B}'=u_*(\sh{O}_{Y'})$, so that $Y'$ identifies with $\operatorname{Spec}(\sh{B}')$, the homomorphism $\varphi':\sh{B}'\to\sh{B}=f_*(\sh{O}_X)$ is \emph{injective}.
But as by hypothesis $\sh{B}'$ is a finite $\sh{O}_Y$-algebra, it identifies by definition of $\sh{C}$ with one of the $\sh{O}_Y$-subalgebras $\sh{C}_\lambda$, which proves \emph{c)}.

We will say that a morphism $f:X\to Y$, where $Y$ is quasi-compact and quasi-separated, is \emph{pseudo-finite}, if it is of finite type and satisfies condition \emph{a)} of \sref{IV.8.12.3} (in which case it is necessarily separated).

\begin{corollary}[8.12.4]
\label{IV.8.12.4}
Let $Y$ be a prescheme quasi-compact and quasi-separated, $f:X\to Y$ a morphism.
\oldpage[IV-3]{45}
\begin{enumerate}[label=(\roman*)]
  \item Suppose $f$ pseudo-finite.
Then, for all morphisms $Y'\to Y$, where $Y'$ is quasi-compact and quasi-separated, $f_{(Y')}:X'=X_{(Y')}\to Y'$ is pseudo-finite.
  \item Let $(U_\lambda)$ be a covering of $Y$ by quasi-compact open subsets.
For $f$ to be pseudo-finite, it is necessary and sufficient that for all $\lambda$, the restriction $f_\lambda:f^{-1}(U_\lambda)\to U_\lambda$ of $f$ be a pseudo-finite morphism.
  \item Suppose furthermore $Y$ Noetherian, and $f$ of finite type.
Then, for $f$ to be pseudo-finite, it is necessary and sufficient that, for all $y\in Y$, the morphism $f_y=f\times 1:X_y=X\times_Y\operatorname{Spec}(\sh{O}_y)\to\operatorname{Spec}(\sh{O}_y)$ be pseudo-finite.
\end{enumerate}
\end{corollary}

\begin{proof}
(i) It is clear that $f_{(Y')}$ is of finite type \sref{IV.1.5.4}. Moreover, a factorisation $X\xrightarrow{g}Z\xrightarrow{u}Y$, with $g$ an immersion and $u$ finite, induces a factorisation $X'\xrightarrow{g'}Z'\xrightarrow{u'}Y'$ of $f_{(Y')}$, where $Z'=Z_{(Y')}$, $g'=g_{(Y')}$ and $u'=u_{(Y')}$. The morphism $g'$ is an immersion \sref[I]{I.4.3.2} and $u'$ is finite \sref[II]{II.6.1.5}; hence $f_{(Y')}$ is pseudo-finite.

(ii) The condition is necessary by virtue of (i), the $U_\lambda$ being quasi-separated since $Y$ is.
To see that it is sufficient, observe (with the notation of \sref{IV.8.12.3}) that if one sets $X_\lambda=f^{-1}(U_\lambda)$, then $\sh{B}|U_\lambda=(f_\lambda)_*(\sh{O}_{X_\lambda})$, and $\sh{C}|U_\lambda$ is the integral closure of $\sh{O}_{U_\lambda}$ in $\sh{B}|U_\lambda$. Consequently, if $h:Z\to Y$ is the canonical morphism, $Z_\lambda=\operatorname{Spec}(\sh{C}|U_\lambda)$ identifies with $h^{-1}(U_\lambda)$. For $g:X\to Z$ to be an immersion, it is necessary and sufficient that for every $\lambda$ its restriction $g_\lambda:f^{-1}(U_\lambda)\to h^{-1}(U_\lambda)$ be an immersion \sref[I]{I.4.2.4}. The conclusion now follows from \sref{IV.8.12.3}.

(iii) It suffices, in virtue of (ii), to prove that for all $y\in Y$ there exists a neighbourhood $U$ of $y$ such that the restriction $f^{-1}(U)\to U$ of $f$ be pseudo-finite.
Designate by $(U_\lambda)$ the decreasing projective filtering system of affine open neighbourhoods of $y$, and apply the method of \sref{IV.8.1.2},~$a$)).
Since $Y$ is Noetherian, the restrictions $f_\lambda:f^{-1}(U_\lambda)\to U_\lambda$ of $f$ are of finite presentation, and the same is true of $f_y$.
By hypothesis $f_y$ factorises as $X_y\xrightarrow{g_y}Z_y\xrightarrow{u_y}\operatorname{Spec}(\sh{O}_y)$, where $u_y$ is finite and $g_y$ is an immersion.
As $\sh{O}_y$ is Noetherian, so is $Z_y$; and since $Z_y$ is of finite presentation over $\operatorname{Spec}(\sh{O}_y)$, there exist $\lambda$ and a morphism of finite presentation $u_\lambda:Z_\lambda\to U_\lambda$ such that $Z_y$ identifies with $Z_\lambda\times_{U_\lambda}\operatorname{Spec}(\sh{O}_y)$ and $u_y$ with $u_\lambda\times 1$ \sref{IV.8.8.2},~(ii)). Moreover, there exists a morphism $g_\lambda:X_\lambda\to Z_\lambda$ such that $g_y=g_\lambda\times 1$ and $f_\lambda=u_\lambda\circ g_\lambda$ \sref{IV.8.8.2},~(i)). One may further take $\lambda$ so that $g_\lambda$ is an immersion and $u_\lambda$ a finite morphism \sref{IV.8.10.5},~(ii) and (x)), which proves that $f_\lambda$ is pseudo-finite.
\end{proof}

\begin{env}[8.12.5]
\label{IV.8.12.5}
We can now give the ``Main theorem'' of Zariski \sref[III]{III.4.4.3}), a proof not using the ``global'' cohomological results of chap.~III, but rather appealing to the finer properties of local Noetherian rings; we will furthermore generalise the statement of the theorem by removing the Noetherian hypotheses:
\end{env}

\begin{theorem}[8.12.6]
\label{IV.8.12.6}
(``Main theorem'' of Zariski). --- Let $Y$ be a prescheme quasi-compact and quasi-separated.
If a morphism $f:X\to Y$ is quasi-finite, separated and of finite presentation, there exists a factorisation of $f$
\begin{equation}
\label{IV.8.12.6.1}
X\xrightarrow{f'}Y'\xrightarrow{u}Y \tag{8.12.6.1}
\end{equation}
where $f'$ is an open immersion and $u$ a finite morphism.
\end{theorem}

\oldpage[IV-3]{46}

\begin{proof}
By virtue of \sref{IV.8.12.4},~(ii)) and the local character (on $Y$) of the notions of quasi-finite morphism, separated and of finite presentation, one can restrict to the case where $Y=\operatorname{Spec}(A)$ is affine.
Applying \sref{IV.8.9.1}), one can suppose that there is a subring $A_0$ of $A$, which is a $\bb{Z}$-algebra of finite type, and an $A$-isomorphism $X_0\otimes_{A_0}A\xrightarrow{\sim}X$, $f$ identifying by this isomorphism with $f_0\times 1$, where $f_0:X_0\to\operatorname{Spec}(A_0)$ is a morphism \emph{of finite type}; furthermore \sref{IV.8.10.5},~(v) and (xi)) one can suppose that $f_0$ is separated and quasi-finite; if one proves that $f_0$ is pseudo-finite, the same will be true of $f$ by \sref{IV.8.12.4},~(i)).
As $A_0$ is then Noetherian and the notions of morphism of finite type, separated and quasi-finite are preserved by base change, it follows from \sref{IV.8.12.4},~(iii)) that one can even suppose that $A$ is a local ring, \emph{essentially of finite type over} $\bb{Z}$ \sref{IV.1.3.8}).
Set $n=\dim(A)$, and let us proceed by induction on $n$; for $n=0$, the theorem is evident, $A$ being a field and the morphism $f$ being already finite \sref[II]{II.6.2.2}).
Set $B=\Gamma(X,\sh{O}_X)$; denote by $C$ the integral closure of $A$ in $B$, set $Z=\operatorname{Spec}(C)$ and let $g:X\to Z$ be the $Y$-morphism corresponding to the canonical injective $A$-homomorphism $C\to B$; by virtue of \sref{IV.8.12.3}), it suffices to show that $g$ is an open immersion.
Let $a$ be the closed point of $Y$, and $U=Y-\{a\}$; $U$ is Noetherian and all its local rings are essentially of finite type over $\bb{Z}$ and of dimension $<n$; taking into account the hypothesis of induction, one sees that the restriction $f^{-1}(U)\to U$ of $f$ is a pseudo-finite morphism.
One concludes from \sref{IV.8.12.3}) that, if $h:Z\to Y$ is the structural morphism, the restriction $f^{-1}(U)\to h^{-1}(U)$ of $g$ is an open immersion.
Set $A'=\hat{A}$, $Y'=\operatorname{Spec}(A')$, $X'=X_{(Y')}$, $f'=f_{(Y')}:X'\to Y'$.
As the canonical morphism $u:Y'\to Y$ is flat, it follows from \sref{IV.2.3.1}) that $B'=\Gamma(X',\sh{O}_{X'})$ identifies with the $A'$-algebra $B\otimes_A A'$.
On the other hand, as $A$ is a local ring excellent \sref{IV.7.8.3}), the morphism $u:Y'\to Y$ is regular, and consequently \sref{IV.6.14.4}) \emph{the integral closure of $A'$ in $B'$ is equal to} $C'=C\otimes_A A'$.
One sees therefore that $Z'=\operatorname{Spec}(C')$ is equal to $Z_{(Y')}$ and the morphism $g':X'\to Z'$ coming from the canonical injection $C'\to B'$ is equal to $g_{(Y')}$.
As $u:Y'\to Y$ is faithfully flat and quasi-compact, to prove that $g$ is an open immersion, it suffices to prove that $g'$ is an open immersion \sref{IV.2.7.1},~(x)).
Now $u^{-1}(a)$ consists of the closed point $a'$ of $Y'$, and consequently $U'=Y'-\{a'\}=u^{-1}(U)$.
If $h':Z'\to Y'$ is the canonical morphism, the fact that the restriction $f^{-1}(U)\to h^{-1}(U)$ of $g$ is an open immersion implies that the same is true of the restriction $f'^{-1}(U')\to h'^{-1}(U')$ of $g'$.
The morphism $f'$ is separated and quasi-finite \sref[II]{II.6.2.4}; as $A'$ is \emph{complete}, one deduces from \sref[II]{II.6.2.6} that $X'$ is $Y'$-isomorphic to a sum $X'_1\sqcup X'_2$, where the restriction $f'|X'_1=f'_1:X'_1\to Y'$ is a morphism \emph{finite}, and $X'_2\subset f'^{-1}(U')$.
It follows from this that $B'$ is the direct product of the two $A'$-algebras $\Gamma(X'_1,\sh{O}_{X'_1})=B'_1$ and $\Gamma(X'_2,\sh{O}_{X'_2})=B'_2$; one concludes immediately that the integral closure $C'$ of $A'$ in $B'$ is the direct product of the integral closures $C'_1$, $C'_2$ of $A'$ in $B'_1$, $B'_2$ respectively, whence $Z'=Z'_1\sqcup Z'_2$, where $Z'_i=\operatorname{Spec}(C'_i)$ ($i=1,2$); and the canonical morphism $g':X'\to Z'$ is such that $g'|X'_i$ is the canonical morphism $g'_i:X'_i\to Z'_i$ ($i=1,2$).
But as $B'_1$ is already a finite $A'$-algebra, one has $C'_1=B'_1$, and $g'_1$ is therefore an isomorphism.
On the other hand, as $X'_2\subset f'^{-1}(U')$ and is open in $f'^{-1}(U')$, one
\oldpage[IV-3]{47}
already knows that $g'_2$ is an open immersion.
One concludes that $g'$ is indeed an open immersion.
\end{proof}

\begin{remark}[8.12.7]
\label{IV.8.12.7}
When, in \sref{IV.8.12.6}), one supposes that $Y$ is an affine \emph{scheme}, the proof by reduction to the Noetherian case shows that, in the factorisation \sref{IV.8.12.6.1}), the morphisms $f'$ and $u$ are also morphisms \emph{of finite presentation} \sref{IV.1.6.2}).
\end{remark}

\begin{corollary}[8.12.8]
\label{IV.8.12.8}
Let $Y$ be a quasi-compact scheme such that there exists an $\sh{O}_Y$-Module ample \sref[II]{II.4.5.3}), $f:X\to Y$ a quasi-finite and quasi-projective morphism.
Then there exists a factorisation of $f$ as
\[
X\xrightarrow{f'}Y'\xrightarrow{u}Y
\]
where $f'$ is an open immersion and $u$ a finite morphism.
\end{corollary}

\begin{proof}
The hypothesis implies that $X$ identifies with a sub-$Y$-scheme quasi-compact of a $Y$-scheme of the form $Z=\mathbf{P}^r_Y$ \sref[II]{II.5.3.3}).
There is therefore a quasi-compact open neighbourhood $U$ of $X$ in $Z$ such that $X$ is closed in $U$; as $Z$ is a scheme, the canonical injection $U\to Z$ is a morphism of finite presentation (\sref{IV.1.2.7} and \sref{IV.1.6.2})), hence the composite morphism $g:U\to Z\to Y$ is also a morphism of finite presentation (the fact that $\mathbf{P}^r_Y$ is of finite presentation over $Y$ following immediately from the definition \sref[II]{II.4.1.1})).
Let $\sh{J}$ be the quasi-coherent Ideal of $\sh{O}_U$ defining the closed sub-prescheme $X$; as $U$ is a quasi-compact scheme, $\sh{J}$ is the filtered inductive limit of its sub-Ideals quasi-coherent of finite type $\sh{J}_\lambda$ \sref[I]{I.9.4.9}).
If $X_\lambda$ is the closed sub-prescheme of $U$ defined by $\sh{J}_\lambda$, one therefore has $X=\bigcap X_\lambda$.
For all $y\in Y$, one therefore has $f^{-1}(y)=\bigcap(X_\lambda\cap g^{-1}(y))$, and as the sets $X_\lambda\cap g^{-1}(y)$ are closed in the Noetherian space $g^{-1}(y)$, there exists for all $y$ an index $\lambda(y)$ such that $f^{-1}(y)=X_{\lambda(y)}\cap g^{-1}(y)$.
Denote by $E_\lambda$ the set of $y\in Y$ such that the fibre $X_\lambda\cap g^{-1}(y)$ of the restriction of $g$ to $X_\lambda$ is a $k(y)$-prescheme finite.
The hypothesis that $f$ is quasi-finite implies, by virtue of what precedes, that $Y=\bigcup_\lambda E_\lambda$.
Now, each of the $X_\lambda$ is, by definition, \emph{of finite presentation} over $Y$; it follows therefore from \sref{IV.9.2.3} and \sref{IV.9.2.6})~($^1$) that the $E_\lambda$ are \emph{constructible} sets in the scheme $Y$; as they form a filtered increasing family, there exists an index $\lambda$ such that $E_\lambda=Y$ \sref{IV.1.9.9}), and for this $\lambda$, the morphism $f_\lambda:X_\lambda\to Y$, restriction of $g$ to $X_\lambda$, is quasi-finite. Since it is of finite presentation and separated, one can apply \sref{IV.8.12.6}), and $f_\lambda$ factorises as
\[
X_\lambda\xrightarrow{j_\lambda}Y'_\lambda\xrightarrow{u_\lambda}Y
\]
where $j_\lambda$ is an open immersion and $u_\lambda$ a finite morphism.
As $X$ is a closed sub-prescheme of $X_\lambda$, one has thus proved that $f$ possesses the property \sref{IV.8.12.3},~$a$)), whence the corollary by virtue of the equivalence of \sref{IV.8.12.3},~$a$)) and \sref{IV.8.12.3},~$a'$)).
\end{proof}

\footnotetext{($^1$) The reader will verify that the corollaries \sref{IV.8.12.8} to \sref{IV.8.12.11} are not used in \textsection~9.}

\begin{corollary}[8.12.9]
\label{IV.8.12.9}
\oldpage[IV-3]{48}
Let $f:X\to Y$ be a locally quasi-finite morphism $(\mathbf{Err}_{\text{III}},~20)$.
For all $x\in X$ there exist an open neighbourhood $U$ of $x$ in $X$, an open neighbourhood $V$ of $y=f(x)$ in $Y$, such that $f(U)\subset V$ and a factorisation
\[
U\xrightarrow{f'}V'\xrightarrow{u}V
\]
of the restriction of $f$ to $U$, where $f'$ is an open immersion and $u$ a finite morphism.
\end{corollary}

\begin{proof}
It evidently suffices to take for $V$ an affine neighbourhood of $y$ in $Y$, for $U$ an affine neighbourhood of $x$ in $X$ contained in $f^{-1}(V)$ and such that $f|U$ is quasi-finite.
The morphism $U\to V$ restriction of $f$ being then affine (hence quasi-projective), one can apply \sref{IV.8.12.8}).
\end{proof}

\begin{corollary}[8.12.10]
\label{IV.8.12.10}
Let $Y$ be an integral normal prescheme, $X$ an integral prescheme, and $f:X\to Y$ a birational, locally quasi-finite morphism $(\mathbf{Err}_{\text{III}},~20)$.
Then $f$ is a local isomorphism; for $f$ to be an open immersion, it is necessary and sufficient that $f$ be separated in addition.
\end{corollary}

\begin{proof}
The second assertion follows immediately from the first and from \sref[I]{I.8.2.8}).
To prove the first assertion, one may suppose $X$ and $Y$ affine and $f$ quasi-finite; consider the factorisation $f=u\circ f'$ of \sref{IV.8.12.8}, which identifies $X$, via $f'$, with the subprescheme induced on an open subset of $Y'$.
As $X$ is integral, one can, by virtue of \sref[I]{I.5.2.3}), replace $Y'$ by the reduced sub-prescheme of $Y'$ having underlying space $\overline{X}$, hence one can also suppose that $Y'$ is \emph{integral}.
Furthermore, since $f$ is birational, $u$ is also.
The conclusion then follows from the following lemma:
\end{proof}

\begin{lemma}[8.12.10.1]
\label{IV.8.12.10.1}
Let $Y'$ be an integral prescheme, $Y$ an integral prescheme and normal; then a morphism $u:Y'\to Y$ finite and birational is an isomorphism.
\end{lemma}

\begin{proof}
Set $\sh{A}=u_*(\sh{O}_{Y'})$; then $\sh{A}$ is a finite $\sh{O}_Y$-algebra and $Y'$ identifies with $\operatorname{Spec}(\sh{A})$ \sref[II]{II.1.3.6}.
If $R(Y)$ is the field of rational functions of $Y$, one therefore has, for all $y\in Y$, $\sh{O}_{Y,y}\subset\sh{A}_y\subset R(Y)$; but as the ring $\sh{O}_{Y,y}$ is by hypothesis integrally closed and $R(Y)$ a field of fractions, one necessarily has $\sh{A}_y=\sh{O}_{Y,y}$, whence the lemma.
\end{proof}

\begin{corollary}[8.12.11]
\label{IV.8.12.11}
Let $Y$ be an integral prescheme and $X$ an integral normal prescheme, and let $f:X\to Y$ be a dominant, locally quasi-finite morphism.
Let $K$ and $L$ (an extension of $K$) be the fields of rational functions of $Y$ and $X$ respectively; and let $Y'$ be the integral closure of $Y$ relative to $L$ \sref[II]{II.6.3.4}); then $f$ factorises in a unique way as $f:X\xrightarrow{f'}Y'\xrightarrow{u}Y$, where $f'$ is birational and corresponds to the identity automorphism of $L$; $f'$ is then a local isomorphism, and for $f'$ to be an open immersion, it is necessary and sufficient that $f$ be separated.
\end{corollary}

\begin{proof}
The existence and uniqueness of the factorisation of $f$ follow from \sref[II]{II.6.3.9}).
It follows further from \sref[II]{II.6.2.4},~(v)), reducing to the affine case, that $f'$ is locally quasi-finite; furthermore, it follows from \sref[I]{I.5.5.1}) that, for $f'$ to be separated, it is necessary and sufficient that $f$ be, since $u$ is affine, hence separated; the last two assertions are therefore consequences of \sref{IV.8.12.10}) applied to $f'$.
\end{proof}

\oldpage[IV-3]{49}

\subsection{Translation in terms of pro-objects}
\label{subsection:IV.8.13}

The following proposition is essentially equivalent to \sref{IV.8.8.2},~(i):

\begin{proposition}[8.13.1]
\label{IV.8.13.1}
Let $S$ be a prescheme, and let $(X_\lambda,v_{\lambda\mu})$ be a filtered projective system of $S$-preschemes; suppose that there exists $\alpha$ such that $v_{\alpha\lambda}$ is an affine morphism for all $\lambda\geqslant\alpha$ (which implies by \sref[II]{II.1.6.2} that $v_{\lambda\mu}$ is affine for $\alpha\leqslant\lambda\leqslant\mu$), so that the projective limit $X=\varprojlim_\lambda X_\lambda$ exists in the category of $S$-preschemes \sref{IV.8.2.3}.
Let $Y$ be an $S$-prescheme, and, for all $\lambda\geqslant\alpha$, let $e_\lambda:\operatorname{Hom}_S(X_\lambda,Y)\to\operatorname{Hom}_S(X,Y)$ be the map which, to every $S$-morphism $f_\lambda:X_\lambda\to Y$, makes correspond $f=f_\lambda\circ v_\lambda$, where $v_\lambda:X\to X_\lambda$ is the canonical morphism.
The family $(e_\lambda)$ is an inductive system of maps, which therefore defines a canonical map
\begin{equation}
\label{IV.8.13.1.1}
\varinjlim_\lambda\operatorname{Hom}_S(X_\lambda,Y)\to\operatorname{Hom}_S(X,Y). \tag{8.13.1.1}
\end{equation}

Suppose that $X_\alpha$ is quasi-compact (resp. quasi-compact and quasi-separated), and that the structural morphism $Y\to S$ is locally of finite type (resp. locally of finite presentation).
Then the map \sref{IV.8.13.1.1} is injective (resp. bijective).
\end{proposition}

\begin{proof}
For $\lambda\geqslant\alpha$, set $Z_\lambda=Y\times_S X_\lambda$, so that $Z_\lambda=Z_\alpha\times_{X_\alpha}X_\lambda$.
Similarly, set $Z=Y\times_S X=Z_\alpha\times_{X_\alpha}X$; it then follows from \sref{IV.8.2.5} that, if one also sets $w_{\lambda\mu}=1\times v_{\lambda\mu}:Z_\mu\to Z_\lambda$ for $\alpha\leqslant\lambda\leqslant\mu$ and $w_\lambda=1\times v_\lambda:Z\to Z_\lambda$ for $\alpha\leqslant\lambda$, then $Z$ is the projective limit of the system $(Z_\lambda,w_{\lambda\mu})$, and the $w_\lambda$ are the corresponding canonical morphisms.
Furthermore, the morphism $Z_\alpha\to X_\alpha$ is locally of finite type (resp. locally of finite presentation) by \sref{IV.1.3.4} and \sref{IV.1.4.3}.
Finally, by \sref[I]{I.3.3.14}, one has
\[
\operatorname{Hom}_S(X_\lambda,Y)=\operatorname{Hom}_{X_\lambda}(X_\lambda,Z_\lambda) \quad\text{and}\quad \operatorname{Hom}_S(X,Y)=\operatorname{Hom}_X(X,Z)
\]
It now suffices to apply \sref{IV.8.8.2},~(i), taking $X_\lambda=S_\lambda$ and replacing $Y_\lambda$ by $Z_\lambda$.
\end{proof}

\begin{corollary}[8.13.2]
\label{IV.8.13.2}
With the notation of \sref{IV.8.13.1}, suppose that $X_\alpha$ is quasi-compact and quasi-separated, and that the $v_{\alpha\lambda}$ are affine for $\alpha\leqslant\lambda$; suppose furthermore that $Y=\varprojlim_\rho Y_\rho$, where $(Y_\rho,t_{\rho\sigma})$ is a filtered projective system of $S$-preschemes such that, for each $\rho$, the structural morphism $Y_\rho\to S$ is locally of finite presentation.
Then one has a canonical bijection
\begin{equation}
\label{IV.8.13.2.1}
\operatorname{Hom}_S(X,Y)\xrightarrow{\sim}\varprojlim_\rho\left(\varinjlim_\lambda\operatorname{Hom}_S(X_\lambda,Y_\rho)\right). \tag{8.13.2.1}
\end{equation}
\end{corollary}

\begin{proof}
Indeed, since $Y$ is the projective limit of the $Y_\rho$, the canonical map $\operatorname{Hom}_S(X,Y)\to\varprojlim_\rho\operatorname{Hom}_S(X,Y_\rho)$ is bijective; on the other hand, the hypotheses give, for each $\rho$, a canonical bijection $\operatorname{Hom}_S(X,Y_\rho)\xrightarrow{\sim}\varinjlim_\lambda\operatorname{Hom}_S(X_\lambda,Y_\rho)$ by \sref{IV.8.13.1}; whence the conclusion.
\end{proof}

\begin{env}[8.13.3]
\label{IV.8.13.3}
The preceding results allow one to interpret in the theory of preschemes the notions of ``pro-variety'' and ``pro-scheme'' which arise in certain applications (for example in local class field theory following the ideas of Serre~[39], or in N\'eron's theory of the reduction of
\oldpage[IV-3]{50}
abelian varieties~[32]).
Let us recall here rapidly the notion of \emph{pro-object} of a category, referring to chap.~V for more ample developments (we will moreover not use before chap.~V the interpretation that follows, and the reader can therefore omit it until the end of this number).
Given a category $\mathcal{C}$, the category $\mathbf{Pro}(\mathcal{C})$ of \emph{pro-objects} of $\mathcal{C}$ has as its objects the \emph{projective systems} (in the universe under consideration) $\mathbf{X}=(X_\mu)_{\mu\in M}$ of objects of $\mathcal{C}$ whose index sets (which depend on the projective system) are assumed to be filtered preordered sets.
Given two such pro-objects $\mathbf{X}=(X_\mu)_{\mu\in M}$ and $\mathbf{X}'=(X'_{\mu'})_{\mu'\in M'}$, the \emph{morphisms} from $\mathbf{X}$ to $\mathbf{X}'$ are by definition the elements of the set $\varprojlim_{\mu'}(\varinjlim_\mu\operatorname{Hom}(X_\mu,X'_{\mu'}))$; that these sets may serve as morphism sets is immediate. The composition of systems of morphisms $u_{\mu'}^\mu:X_\mu\to X'_{\mu'}$ and $u_{\mu''}^{\prime\mu'}:X'_{\mu'}\to X''_{\mu''}$, which are inductive in the upper index and projective in the lower index, is performed ``argument by argument'', that is, by considering the system $u_{\mu''}^{\prime\mu'}\circ u_{\mu'}^\mu$.
\end{env}

\begin{env}[8.13.4]
\label{IV.8.13.4}
Let $S$ be a quasi-compact and quasi-separated prescheme, and denote by $\mathcal{C}$ the full subcategory of the category $(\mathbf{Sch})_{/S}$ of $S$-preschemes, formed of the $S$-preschemes $X$ having the following property: the structural morphism $X\to S$ factorises as $X\xrightarrow{g}X_0\xrightarrow{f}S$, where $g:X\to X_0$ is affine and $f:X_0\to S$ is \emph{of finite presentation}; we will say for short that the preschemes of $\mathcal{C}$ are \emph{essentially affine over} $S$.
Consider on the other hand the full subcategory $\mathcal{C}'_0$ of $(\mathbf{Sch})_{/S}$ formed of the $S$-preschemes \emph{of finite presentation}, and the category $\mathbf{Pro}(\mathcal{C}'_0)$ of pro-objects of $\mathcal{C}'_0$.
We will say that an object $\mathbf{X}=(X_\mu)_{\mu\in M}$ of $\mathbf{Pro}(\mathcal{C}'_0)$ is \emph{essentially affine} if there exists a $\gamma\in M$ such that for $\mu\geqslant\gamma$, the transition morphism $v_{\gamma\mu}:X_\mu\to X_\gamma$ is affine (which implies that for $\gamma\leqslant\nu\leqslant\mu$, $v_{\nu\mu}:X_\mu\to X_\nu$ is affine).
One will note that an object of $\mathbf{Pro}(\mathcal{C}'_0)$ isomorphic to an essentially affine object \emph{is not} necessarily itself essentially affine.
We will denote by $\mathcal{C}'$ the full subcategory of $\mathbf{Pro}(\mathcal{C}'_0)$ formed of pro-objects essentially affine.

This being so, it follows from \sref{IV.8.2.2} and \sref{IV.8.2.3} that for every object $\mathbf{X}=(X_\mu)_{\mu\in M}$ of $\mathcal{C}'$, the $S$-prescheme $X=\varprojlim_\mu X_\mu$ exists; furthermore, since the morphism $X\to X_\mu$ is affine for all sufficiently large $\mu$ \sref{IV.8.2.2}, $X$ is \emph{essentially affine over} $S$ by definition.
Set $X=L(\mathbf{X})$; let us show that one has thus defined a canonical \emph{functor}
\begin{equation}
\label{IV.8.13.4.1}
L:\mathcal{C}'\to\mathcal{C} \tag{8.13.4.1}
\end{equation}

For two objects $\mathbf{X}=(X_\mu)$ and $\mathbf{X}'=(X'_{\mu'})$ of $\mathcal{C}'$, one has, for each $\mu'$, a canonical map
\[
\varinjlim_\mu\operatorname{Hom}_S(X_\mu,X'_{\mu'})\to\operatorname{Hom}_S(\varprojlim_\mu X_\mu,X'_{\mu'})
\]
defined in \sref{IV.8.13.1.1}; on the other hand, by definition of the projective limit, there is a canonical bijection
\[
\varprojlim_{\mu'}\operatorname{Hom}_S(\varprojlim_\mu X_\mu,X'_{\mu'})\xrightarrow{\sim}\operatorname{Hom}_S(\varprojlim_\mu X_\mu,\varprojlim_{\mu'}X'_{\mu'})
\]
\oldpage[IV-3]{51}

whence a canonical map
\begin{equation}
\label{IV.8.13.4.2}
\varprojlim_{\mu'}\varinjlim_\mu\operatorname{Hom}_S(X_\mu,X'_{\mu'})\to\operatorname{Hom}_S(\varprojlim_\mu X_\mu,\varprojlim_{\mu'}X'_{\mu'}) \tag{8.13.4.2}
\end{equation}
evidently functorial in $\mathbf{X}$ and $\mathbf{X}'$, and which completes the definition of the functor $L$.
\end{env}

\begin{proposition}[8.13.5]
\label{IV.8.13.5}
Under the hypotheses and notation of \sref{IV.8.13.4}, the functor $L$ is fully faithful.
If furthermore $S$ is a Noetherian prescheme (which already implies that $S$ is quasi-compact and quasi-separated \sref{IV.1.2.8}), $L$ is \emph{an equivalence of categories}.
\end{proposition}

\begin{proof}
To say that $L$ is fully faithful means that the map \sref{IV.8.13.4.2} is bijective, which is a special case of \sref{IV.8.13.2}.
Indeed, the structural morphisms $X_\mu\to S$ being of finite presentation, are in particular quasi-compact and quasi-separated, hence the $X_\mu$ are quasi-compact and quasi-separated.

To show that when $S$ is Noetherian $L$ is an equivalence of categories, it suffices, since one already knows that $L$ is fully faithful, to prove that every essentially affine $S$-prescheme $X$ is $S$-\emph{isomorphic} to an object of the form $L(\mathbf{X})$ where $\mathbf{X}\in\mathcal{C}'$ $(\mathbf{0}_{\text{III}},~8.1.5)$.
Now, by hypothesis there is a factorisation $X\xrightarrow{g}X_0\xrightarrow{f}S$ of the structural morphism, $f$ being of finite presentation and $g$ affine.
One can therefore write $X=\operatorname{Spec}(\sh{A})$, where $\sh{A}$ is a quasi-coherent $\sh{O}_{X_0}$-Algebra \sref[II]{II.1.3.1}.
Now, as $X_0$ is Noetherian (since it is of finite type over $S$ and $S$ is Noetherian), $\sh{A}$ is the filtered inductive limit of the family $(\sh{A}_\lambda)$ of its quasi-coherent sub-$\sh{O}_{X_0}$-Algebras of finite type \sref[I]{I.9.6.6}.
Set $X_\lambda=\operatorname{Spec}(\sh{A}_\lambda)$; the morphisms $X_\lambda\to X_0$ are of finite type, hence of finite presentation since $X_0$ is Noetherian, and consequently the same is true of the composite morphisms $X_\lambda\to X_0\to S$ \sref{IV.1.6.2}; in other words, the $X_\lambda$ belong to $\mathcal{C}'_0$, and since the morphisms $X_\lambda\to X_0$ are affine, $\mathbf{X}=(X_\lambda)$ is an object of $\mathcal{C}'$ whose projective limit exists and is $S$-isomorphic to $X$ by \sref{IV.8.2.2}.
This completes the proof.
\end{proof}

\begin{remark}[8.13.6]
\label{IV.8.13.6}
It follows from \sref{IV.1.6.2} and \sref[II]{II.1.6.2} that if $X$ and $Y$ are essentially affine over $S$, the same is true of $X\times_S Y$.
One concludes for example $(\mathbf{0}_{\text{III}},~8.2.5)$ that a $\mathcal{C}$-group is nothing other than a $(\mathbf{Sch})_{/S}$-group which is a prescheme essentially affine over $S$.
On the other hand, the finite products exist in the category $\mathcal{C}'$: indeed, if $\mathbf{X}=(X_\mu)_{\mu\in M}$, $\mathbf{Y}=(Y_\rho)_{\rho\in R}$ are two objects of $\mathcal{C}'$, the products $X_\mu\times_S Y_\rho$ are $S$-preschemes of finite presentation, and taking for transition morphisms $X_\nu\times_S Y_\sigma\to X_\mu\times_S Y_\rho$ the products of the transition morphisms $X_\nu\to X_\mu$ and $Y_\sigma\to Y_\rho$, one sees immediately that $(X_\mu\times_S Y_\rho)$ is a product of $\mathbf{X}$ and $\mathbf{Y}$ in $\mathbf{Pro}(\mathcal{C}'_0)$; furthermore, by \sref[II]{II.1.6.2}, the transition morphisms thus defined are affine for $\mu$ and $\rho$ large enough, hence the product $\mathbf{X}\times\mathbf{Y}$ also belongs to $\mathcal{C}'$.
One concludes then as above that a $\mathcal{C}'$-group is a $\mathbf{Pro}(\mathcal{C}'_0)$-group which is essentially affine.
One deduces from \sref{IV.8.13.5} that the categories of the $\mathcal{C}'$-groups and the $\mathcal{C}$-groups are \emph{equivalent} when $S$ is \emph{Noetherian}.
It seems plausible that when $S$ is the spectrum of a field $k$, the category of $\mathcal{C}$-groups is \emph{equivalent} to that of the $\boldsymbol{K}$-\emph{groups}, where $\boldsymbol{K}$ is the category of $S$-preschemes \emph{quasi-compact}; in other words, every prescheme in groups quasi-compact over $k$ would be essentially affine.
On the other hand, if one denotes
\oldpage[IV-3]{52}
by $\mathcal{C}'_0$-$\boldsymbol{gr}$ the category of the $\mathcal{C}'_0$-groups, it is very likely that the category of the $\mathcal{C}'$-groups is equivalent to the full subcategory of $\mathbf{Pro}(\mathcal{C}'_0\text{-}\boldsymbol{gr})$ formed of the ``proalgebraic groups essentially affine'', that is to say the projective systems $\mathbf{G}=(G_\mu)_{\mu\in M}$, where the $G_\mu$ are algebraic groups over $k$ and the transition morphisms $G_\nu\to G_\mu$ are affine for $\mu$ large enough (which is what one can also express by saying that $\mathbf{G}$ is an extension of an algebraic group by a pro-algebraic group affine).
The conjunction of these two conjectures is moreover equivalent to the following: every quasi-compact group prescheme over $k$ is an extension of an ``algebraic group'' (i.e.\ a group prescheme of finite type over $k$) by a group prescheme \emph{affine} over $k$.
The only pro-algebraic groups encountered in practice until now being in fact essentially affine, it would therefore be advantageous to substitute for the study of general pro-algebraic groups (introduced and studied by Serre~[40]) that of schemes in groups quasi-compact over $k$, whose definition is conceptually much simpler.
\end{remark}

\subsection{Characterisation of a prescheme locally of finite presentation over another, in terms of the functor it represents}
\label{subsection:IV.8.14}

\begin{env}[8.14.1]
\label{IV.8.14.1}
Given a prescheme $S$, we will say, as in \sref{IV.8.13.4}, that a filtered projective system $(X_\lambda,v_{\lambda\mu})$ of $S$-preschemes is \emph{essentially affine} if there exists $\alpha$ such that $v_{\alpha\lambda}$ is an affine morphism for $\lambda\geqslant\alpha$.

The following statement, which will be especially useful in chap.~V, makes \sref{IV.8.8.2},~(i) more precise by providing a converse:
\end{env}

\begin{proposition}[8.14.2]
\label{IV.8.14.2}
Let $S$ be a prescheme, $f:X\to S$ a morphism.
For every $S$-prescheme $T$, set
\[
h_X(T)=\operatorname{Hom}_S(T,X)
\]
so that $h_X$ is a contravariant functor from the category $(\mathbf{Sch})_{/S}$ of $S$-preschemes to the category $\mathbf{Ens}$ of sets $(\mathbf{0}_{\text{III}},~8.1.1)$, and $X$ is an object representing this functor $(\mathbf{0}_{\text{III}},~8.1.8)$.
The following conditions are equivalent:
\begin{enumerate}[label=\emph{\alph*)}]
  \item $f$ is locally of finite presentation.
  \item For every filtered projective system $(Z_\lambda)$ of $S$-preschemes, essentially affine \sref{IV.8.13.4} and formed of quasi-compact and quasi-separated preschemes, the canonical map \sref{IV.8.13.1.1}
\begin{equation}
\label{IV.8.14.2.1}
\varinjlim_\lambda h_X(Z_\lambda)\to h_X(\varprojlim_\lambda Z_\lambda) \tag{8.14.2.1}
\end{equation}
is bijective.
  \item For every filtered projective system $(Z_\lambda)$ of $S$-preschemes such that the $Z_\lambda$ are affine schemes, the map \sref{IV.8.14.2.1} is bijective.
  \item[\rm c$'$)] For every affine open subset $U$ of $S$ and every filtered projective system $(Z_\lambda)$ of $U$-preschemes such that the $Z_\lambda$ are affine schemes, the map \sref{IV.8.14.2.1} is bijective.
\end{enumerate}
\end{proposition}

\begin{proof}
The fact that \emph{a)} implies \emph{b)} is precisely \sref{IV.8.13.1}; it is trivial that \emph{b)} implies \emph{c)} and that \emph{c)} implies \emph{c$'$)}.
It remains to see that \emph{c$'$)} implies \emph{a)}, and as the property \emph{a)} is local on $S$, one can restrict to the case where $S$ is \emph{affine}.
\oldpage[IV-3]{53}

Suppose first that $X$ is also affine; the assertion to prove is then equivalent to the following:

\begin{corollary}[8.14.2.2]
\label{IV.8.14.2.2}
Let $A$ be a ring, $B$ an $A$-algebra.
In order that, for every filtered inductive system $(C_\lambda)$ of $A$-algebras, the canonical map
\begin{equation}
\label{IV.8.14.2.3}
\varinjlim_\lambda\operatorname{Hom}_{A\text{-}\mathrm{alg.}}(B,C_\lambda)\to\operatorname{Hom}_{A\text{-}\mathrm{alg.}}(B,\varinjlim_\lambda C_\lambda) \tag{8.14.2.3}
\end{equation}
be bijective, it is necessary and sufficient that $B$ be an $A$-algebra of finite presentation.
\end{corollary}

It remains only to show that the condition is necessary.
First, for $(C_\lambda)$ take the filtered inductive system of sub-$A$-algebras of finite type of $B$, so that $\varinjlim_\lambda C_\lambda=B$.
The fact that \sref{IV.8.14.2.3} is bijective implies in particular that the identity map $1_B:B\to B$ factorises as $B\to C_\lambda\to B$ for a suitable $\lambda$, which implies that $C_\lambda=B$, hence $B$ is an $A$-algebra \emph{of finite type}.
Write $B=C/\mathfrak{J}$, where $C=A[T_1,\ldots,T_n]$ and $\mathfrak{J}$ is an ideal of $C$. The ideal $\mathfrak{J}$ is the filtered inductive limit of its ideals of finite type $\mathfrak{J}_\lambda\subset\mathfrak{J}$; setting $C_\lambda=C/\mathfrak{J}_\lambda$ and using the exactness of filtered inductive limits, one sees that $B$ is again isomorphic to the inductive limit of the filtered inductive system $(C_\lambda)$.
There exists therefore a $\lambda$ and an $A$-homomorphism $u:B\to C_\lambda$ such that the composite $B\xrightarrow{u}C_\lambda\xrightarrow{p_\lambda}B$ (where $p_\lambda$ is the canonical homomorphism) is the identity.
Let $q_\lambda:C\to C_\lambda$ be the canonical homomorphism, and set
$t_i=p_\lambda(q_\lambda(T_i))$.
Then $p_\lambda(u(t_i))=p_\lambda(q_\lambda(T_i))$, in other words
$u(t_i)-q_\lambda(T_i)\in\mathfrak{J}/\mathfrak{J}_\lambda$.
There is therefore a $\mu\geqslant\lambda$ such that these $n$ elements belong to
$\mathfrak{J}_\mu/\mathfrak{J}_\lambda$ ($1\leqslant i\leqslant n$).
If $p_{\mu\lambda}:C_\lambda\to C_\mu$ is the canonical homomorphism, then
$p_{\mu\lambda}(u(t_i))=p_{\mu\lambda}(q_\lambda(T_i))=q_\mu(T_i)$.
Replacing $\lambda$ by $\mu$ and $u$ by $p_{\mu\lambda}\circ u$, we may therefore
suppose that $u(t_i)=q_\lambda(T_i)$ for every $i$; and if
$r=p_\lambda\circ q_\lambda$ is the canonical homomorphism
$C\to C/\mathfrak{J}=B$, we may write $u(r(T_i))=q_\lambda(T_i)$ for every $i$,
whence $q_\lambda=u\circ r$.
This necessarily implies that $\mathfrak{J}=\mathfrak{J}_\lambda$: if $z\in\mathfrak{J}$, then $r(z)=0$, hence $q_\lambda(z)=u(r(z))=0$ and $z\in\mathfrak{J}_\lambda$. Thus $B=C_\lambda$.

Let us now pass to the case where $S$ is affine and $X$ arbitrary; what amounts to proving is that an affine open subset $V$ of $X$ is of finite presentation over $S$, and by virtue of what has just been shown, it suffices to prove that for every filtered projective system $(Z_\lambda)$ of affine $S$-preschemes, the map
\begin{equation}
\label{IV.8.14.2.4}
\varinjlim_\lambda\operatorname{Hom}_S(Z_\lambda,V)\to\operatorname{Hom}_S(\varprojlim_\lambda Z_\lambda,V) \tag{8.14.2.4}
\end{equation}
is bijective.
It is immediate that this map is injective, for if $(v_\lambda)$, $(v'_\lambda)$ are two inductive systems of $S$-homomorphisms $v_\lambda:Z_\lambda\to V$, $v'_\lambda:Z_\lambda\to V$ such that the corresponding morphisms
\[
Z\xrightarrow{u_\lambda}Z_\lambda\xrightarrow{v_\lambda}V,\qquad
Z\xrightarrow{u_\lambda}Z_\lambda\xrightarrow{v'_\lambda}V
\]
are equal ($u_\lambda$ being the canonical morphism), then the morphisms
\[
Z\xrightarrow{u_\lambda}Z_\lambda\xrightarrow{v_\lambda}V\xrightarrow{j}X,\qquad Z\xrightarrow{u_\lambda}Z_\lambda\xrightarrow{v'_\lambda}V\xrightarrow{j}X
\]
(where $j$ is the canonical injection) are equal, which implies $j\circ v_\lambda=j\circ v'_\lambda$ by hypothesis for a suitable $\lambda$, hence $v_\lambda=v'_\lambda$.

\oldpage[IV-3]{54}

It remains to prove that \sref{IV.8.14.2.4} is surjective.
Let $v:Z\to V$ be an $S$-morphism.
By hypothesis there exist an index $\lambda$ and an $S$-morphism
$w_\lambda:Z_\lambda\to X$ such that $j\circ v$ factors as
\[
 Z\xrightarrow{u_\lambda}Z_\lambda\xrightarrow{w_\lambda}X.
\]
It suffices to prove that there is a $\mu\geqslant\lambda$ such that
\[
 Z_\mu\xrightarrow{u_{\lambda\mu}}Z_\lambda\xrightarrow{w_\lambda}X
\]
(where $u_{\lambda\mu}$ is the transition morphism) factors as
\[
 Z_\mu\xrightarrow{v_\mu}V\xrightarrow{j}X.
\]
For every $\mu\geqslant\lambda$, set
$U_\mu=u_{\lambda\mu}^{-1}(w_\lambda^{-1}(V))$.
Then
\[
 u_\mu^{-1}(U_\mu)=u_\lambda^{-1}(U_\lambda)
 =u_\lambda^{-1}(w_\lambda^{-1}(V))
 =v^{-1}(V)=Z.
\]
The $Z_\lambda$ are quasi-compact and the $U_\mu$, being open, are
ind-constructible \sref{IV.1.9.6}; hence \sref{IV.8.3.4} gives a
$\mu\geqslant\lambda$ for which $U_\mu=Z_\mu$.
\end{proof}

\begin{remark}[8.14.3]
\label{IV.8.14.3}
The fact that the map \sref{IV.8.14.2.1} is injective when $f$ is locally of finite
type \sref{IV.8.8.2},~(i) naturally leads one to ask whether this result has a converse.
It does not, even when $S$ and $X$ are affine, since there exist monomorphisms
$X\to S$ which are not of finite type \sref[I]{I.2.4.2}, and these disprove such a
conjecture.
\end{remark}
% END NATIVE INPUT source_aligned/ega4-8
% BEGIN NATIVE INPUT source_aligned/ega4-9
\section{Constructible properties}
\label{section:IV.9}

Let $S$ be a prescheme, $f:X\to S$ a morphism of \emph{finite presentation}~\sref{IV.1.6.1},
$\sh{F}$ a quasi-coherent $\sh{O}_X$-Module of \emph{finite presentation}.
We propose in this paragraph to give criteria ensuring, for example, that the set of $s\in S$ such that the $\mathbf{k}(s)$-prescheme $X_s=f^{-1}(s)=X\times_S\operatorname{Spec}(\mathbf{k}(s))$ has a certain property, or such that the $\sh{O}_{X_s}$-Module $\sh{F}_s=\sh{F}\otimes_{\sh{O}_X}\mathbf{k}(s)$ has a certain property, is \emph{locally constructible}, or at least \emph{ind-constructible}~\sref{IV.1.9.4}; we will see that this is the case for the majority of properties that arise in algebraic Geometry.
Assuming only $f$ \emph{locally of finite presentation}, we will also give~{\hyperlink{ega4.c04.subsec.9-9}{(9.9)}} criteria for the set of points $x\in X$ where the fibre $X_{f(x)}$ (or the $\sh{O}_{X_{f(x)}}$-Module $\sh{F}_{f(x)}$) has a certain property, to be \emph{locally constructible}.
We will see in \textsection12 that these results, together with the additional hypothesis that $f$ is \emph{flat} (resp.\ \emph{proper and flat}), allow us to prove that the sets considered in $X$ (resp.\ in $S$) are even \emph{open} in many cases.

\subsection{The finite extension principle}
\label{subsection:IV.9.1}

\begin{proposition}[9.1.1]
\label{IV.9.1.1}
\textnormal{(Finite extension principle).}
Let $k$ be a field, $\mathcal{C}$ a set of extensions of $k$.
Suppose that the following conditions are satisfied:
\begin{enumerate}[label=(\roman*)]
  \item If $K\in\mathcal{C}$ and if there exists a $k$-homomorphism $K\to K'$ (where $K'$ belongs to the universe where we have placed ourselves), then $K'\in\mathcal{C}$.
  \item If $K\in\mathcal{C}$, there exists a sub-extension $K'$ of $K$, of finite type over $k$, such that $K'\in\mathcal{C}$.
  \item If $K\in\mathcal{C}$ is the field of fractions of an integral $k$-algebra of finite type $A$, there exists $f\in A-\{0\}$ such that for every maximal ideal $\mathfrak{m}$ of $A$ we have $A_f/\mathfrak{m}\in\mathcal{C}$.
\end{enumerate}

\oldpage[IV-3]{55}

Let then $\Omega$ be an algebraically closed extension of $k$ (in the universe considered).
The following conditions are equivalent:
\begin{enumerate}[label=\emph{\alph*)}]
  \item $\mathcal{C}$ is non-empty.
  \item There exists $K'\in\mathcal{C}$ which is a finite extension of $k$.
  \item We have $\Omega\in\mathcal{C}$.
\end{enumerate}
\end{proposition}

The condition~(i) obviously implies that \emph{b)} entails \emph{c)}, and \emph{c)} trivially entails \emph{a)};
let us prove that \emph{a)} entails \emph{b)}.
By virtue of~(ii) and~(iii) there exists an extension $K'\in\mathcal{C}$ which is of the form $A/\mathfrak{m}$, where $A$ is a $k$-algebra of finite type over $k$ and $\mathfrak{m}$ is a maximal ideal of $A$.
We know, by Hilbert's Nullstellensatz (Bourbaki, \emph{Alg.\ comm.}, chap.~V, \textsection3, n\textsuperscript{o}~1, cor.~2 of th.~1), that $K'$ is a finite extension of $k$.

\begin{corollary}[9.1.2]
\label{IV.9.1.2}
Under the hypotheses~\emph{(i)}, \emph{(ii)}, \emph{(iii)} of~\sref{IV.9.1.1}, if $k$ is algebraically closed and if $\mathcal{C}$ is non-empty, $\mathcal{C}$ contains all the extensions of $k$ in the universe considered.
\end{corollary}

Indeed, as \emph{a)} entails \emph{c)}, we have $k\in\mathcal{C}$, and the conclusion follows from the hypothesis~(i).

\begin{remark}[9.1.3]
\label{IV.9.1.3}
In practice, one will verify the condition~(ii) of~\sref{IV.9.1.1} by noting that $K$ is the inductive limit of its sub-extensions of finite type $K_\alpha$, and by applying the results of \textsection8, taking into account eventually that for $K_\alpha\subset K_\beta$, $K_\beta$ is faithfully flat over $K_\alpha$.
Frequently the set $\mathcal{C}$ is formed of \emph{fields} belonging to a set $\mathcal{C}'$ of $k$-algebras satisfying the following condition:
\begin{enumerate}[label=(\roman*~bis)]
  \item If $A\in\mathcal{C}'$ and if there exists a homomorphism of $k$-algebras $A\to A'$ (where $A'$ belongs to the universe where we have placed ourselves), then $A'\in\mathcal{C}'$.
\end{enumerate}
\end{remark}

Since this is so, the condition~(i) is trivially satisfied, and one will verify generally the condition~(iii) of~\sref{IV.9.1.1} by noting that the field of fractions $K$ of $A$ is the inductive limit of the algebras $A_f$ (for the relation $D(f)\supset D(g)$), and by applying the results of \textsection8, taking into account eventually that the morphisms $D(g)\to D(f)$ are open immersions.

On the other hand, when (i~bis) is not satisfied, the proof of~(iii) is often more delicate, and is linked to constructibility criteria that will be developed further on.

Here are some typical examples of application of the finite extension principle:

\begin{proposition}[9.1.4]
\label{IV.9.1.4}
Let $k$ be a field, $\Omega$ an algebraically closed extension of $k$, $X$ and $Y$ two preschemes of finite type over $k$.
The following conditions are equivalent:
\begin{enumerate}[label=\emph{\alph*)}]
  \item There exists an $\Omega$-morphism $X_{(\Omega)}\to Y_{(\Omega)}$ \textnormal{(}resp.\ an $\Omega$-morphism possessing one (specified) of the properties~\emph{(i)} to~\emph{(xiv)} of~\sref{IV.8.10.5}\textnormal{)}.
  \item There exists a finite extension $k'$ of $k$ and a $k'$-morphism $X_{(k')}\to Y_{(k')}$ \textnormal{(}resp.\ a $k'$-morphism possessing the property considered\textnormal{)}.
  \item There exists an extension $K$ of $k$ and a $K$-morphism $X_{(K)}\to Y_{(K)}$ \textnormal{(}resp.\ a $K$-morphism possessing the property considered\textnormal{)}.
\end{enumerate}
\end{proposition}

We apply the remark~\sref{IV.9.1.3} taking for $\mathcal{C}'$ the set of all $k$-algebras $A$ (of the universe where we have placed ourselves) such that there exists an $A$-morphism $X\otimes_k A\to Y\otimes_k A$ (resp.\ a morphism having the particular property of~\sref{IV.8.10.5} under consideration).
The condition~(i~bis) of~\sref{IV.9.1.3} is then satisfied thanks to the fact that the
\oldpage[IV-3]{56}
properties considered in~\sref{IV.8.10.5} are all stable under base change.
The condition~(iii) of~\sref{IV.9.1.1} is thus satisfied by virtue of~\sref{IV.8.8.2}[i] (resp.\ \sref{IV.8.10.5}), since $\operatorname{Spec}(k)$ is quasi-compact and quasi-separated.
It remains to verify the condition~(ii) of~\sref{IV.9.1.1}, which also follows from~\sref{IV.8.8.2}[i] and~\sref{IV.8.10.5}, considering $K$ as an inductive limit of its sub-extensions of finite type.
We conclude therefore by~\sref{IV.9.1.1}.

In particular, if there exists an extension $K$ of $k$ and a $K$-isomorphism $X_{(K)}\xrightarrow{\sim} Y_{(K)}$, we say that $X$ and $Y$ are \emph{geometrically isomorphic}.

The following corollary generalises~(\textbf{II},~6.6.5):

\begin{corollary}[9.1.5]
\label{IV.9.1.5}
Let $k$ be a field, $X$ a $k$-prescheme.
If there exists an extension $K$ of $k$ such that $X_{(K)}$ is projective \textnormal{(}resp.\ quasi-projective\textnormal{)} over $K$, then $X$ is projective \textnormal{(}resp.\ quasi-projective\textnormal{)} over $k$.
\end{corollary}

The morphism $\operatorname{Spec}(K)\to\operatorname{Spec}(k)$ being faithfully flat and quasi-compact, it follows already from~\sref{IV.2.7.1}[v] that $X$ is of finite type over $k$.
The hypothesis means that there exists a closed immersion (resp.\ an immersion) $X_{(K)}\to\mathbf{P}_K^n=\mathbf{P}_k^n\otimes_k K$ (\textbf{II},~5.5.4,~(ii) and~5.5.3); applying~\sref{IV.9.1.4} for property~(iv) (resp.~(ii)) of~\sref{IV.8.10.5}, we deduce that for some finite extension $k'$ of $k$ there is a closed immersion (resp.\ an immersion) $X_{(k')}\to\mathbf{P}_{k'}^n$; in other words $X_{(k')}$ is projective (resp.\ quasi-projective) over $k'$.
We conclude then by~(\textbf{II},~6.6.5).

\begin{proposition}[9.1.6]
\label{IV.9.1.6}
Let $k$ be a field, $\Omega$ an algebraically closed extension of $k$, $X$ a prescheme of finite type over $k$, $\sh{F}$, $\sh{G}$ two coherent $\sh{O}_X$-Modules.
Suppose that there exists an isomorphism $\sh{F}\otimes_k\Omega\xrightarrow{\sim}\sh{G}\otimes_k\Omega$.
Then there exists a finite extension $k'$ of $k$ and an isomorphism $\sh{F}\otimes_k k'\xrightarrow{\sim}\sh{G}\otimes_k k'$.
\end{proposition}

The reasoning is the same as in~\sref{IV.9.1.4}, applying~\sref{IV.8.5.2}[i] (in the proof of property~(iii) of~\sref{IV.9.1.1}, we use here the fact that the morphisms $D(g)\to D(f)$ (with the notations of~\sref{IV.9.1.3}) are open immersions, and \emph{a fortiori} flat morphisms).

\subsection{Constructible and ind-constructible properties}
\label{subsection:IV.9.2}
\label{IV.9.2}

\begin{definition}[9.2.1]
\label{IV.9.2.1}
Let $\mathbf{P}(X,\sh{F},k)$ be a relation.
We say that $\mathbf{P}$ is a \emph{constructible} (resp.\ \emph{ind-constructible}) property of algebraic preschemes if the two following conditions are satisfied:
\begin{enumerate}[label=(\roman*)]
  \item If $k$ is a field, $X$ an algebraic prescheme over $k$, $\sh{F}$ a coherent $\sh{O}_X$-Module, $k'$ an extension of $k$, then for $\mathbf{P}(X,\sh{F},k)$ to be true, it is necessary and sufficient that $\mathbf{P}(X_{(k')},\sh{F}\otimes_k k',k')$ be true.
  \item Let $S$ be an integral Noetherian prescheme, of generic point $\eta$, $u:X\to S$ a morphism of finite type, $\sh{F}$ a coherent $\sh{O}_X$-Module.
For every $s\in S$, set $X_s=u^{-1}(s)=X\times_S\operatorname{Spec}(\mathbf{k}(s))$, $\sh{F}_s=\sh{F}\otimes_{\sh{O}_S}\mathbf{k}(s)$.
Let $E$ be the set of $s\in S$ such that $\mathbf{P}(X_s,\sh{F}_s,\mathbf{k}(s))$ is true.
Then one of the two sets $E$, $S-E$ (resp.\ the set $E$) contains a non-empty open set (and consequently is a neighbourhood of $\eta$) (resp.\ contains a non-empty open set if it contains $\eta$).
\end{enumerate}
\end{definition}

\oldpage[IV-3]{57}

\begin{remarks}[9.2.2]
\label{IV.9.2.2}
\begin{enumerate}[label=(\roman*)]
  \item This is of course a convention of meta-mathematical language and not a mathematical definition properly speaking.
We have analogous ``definitions'' for relations between $k$, one or more algebraic $k$-preschemes, $k$-morphisms between these preschemes, coherent Modules on these preschemes or homomorphisms between these Modules.
  \item We will also need to consider relations involving \emph{constructible parts} of preschemes.
For example, let $\mathbf{P}(X,Z,k)$ be a relation; we will say (by abuse of language) that $\mathbf{P}$ is a constructible (resp.\ ind-constructible) property of the constructible part $Z$ of $X$ if the following two conditions are satisfied:
\begin{enumerate}[label=$\arabic*^{\mathrm{o}}$]
  \item If $k$ is a field, $X$ an algebraic prescheme over $k$, $Z$ a \emph{constructible part} of $X$, $k'$ an extension of $k$, then for $\mathbf{P}(X,Z,k)$ to be true, it is necessary and sufficient that $\mathbf{P}(X_{(k')},p^{-1}(Z),k')$ be true ($p:X_{(k')}\to X$ being the canonical projection).
  \item Let $S$ be an integral Noetherian prescheme, of generic point $\eta$, $u:X\to S$ a morphism of finite type, $Z$ a \emph{constructible part of $X$}.
For every $s\in S$, set $X_s=u^{-1}(s)$, $Z_s=Z\cap X_s$.
Let $E$ be the set of $s\in S$ such that $\mathbf{P}(X_s,Z_s,\mathbf{k}(s))$ is true.
Then one of the two sets $E$, $S-E$ (resp.\ the set $E$) contains a non-empty open set (resp.\ contains a non-empty open set if it contains $\eta$).
\end{enumerate}

We will note that in condition~$2^\circ$ one must suppose that $Z$ is a constructible part of $X$, and not merely that $Z_s$ is a constructible part of $X_s$ for every $s$; the first of these two properties entails the second~\sref{IV.1.8.2}, but not conversely.

  \item If $\mathbf{P}$ is a constructible property, it is clear that the same is true of ``not $\mathbf{P}$''.
If $\mathbf{P}$, $\mathbf{Q}$ are two constructible properties (resp.\ ind-constructible), the same is true of the properties ``$\mathbf{P}$ or $\mathbf{Q}$'' and ``$\mathbf{P}$ and $\mathbf{Q}$''; indeed, under the hypotheses of~\sref{IV.9.2.1}[ii], if $E$, $E'$ are two parts of $S$ and $E$ contains a non-empty open set, the same is true of $E\cup E'$, and if $S-E$ and $S-E'$ each contain a non-empty open set, the same is true of $S-(E\cup E')=(S-E)\cap(S-E')$.

  \item Let $\mathbf{P}(X,\sh{F},k)$ be a relation satisfying the condition~\sref{IV.9.2.1}[i]; let $S$ be a prescheme, $u:X\to S$ a morphism of finite type; with the notations of~\sref{IV.9.2.1}[ii], let $E$ be the set of $s\in S$ such that $\mathbf{P}(X_s,\sh{F}_s,\mathbf{k}(s))$ is true.
Let on the other hand $g:S'\to S$ be any morphism, and set $X'=X_{(S')}$, $\sh{F}'=\sh{F}\otimes_{\sh{O}_S}\sh{O}_{S'}$; then it follows from the transitivity of fibres~(\textbf{I},~3.6.4) and from the condition~\sref{IV.9.2.1}[i] that \emph{the set of $s'\in S'$ such that $\mathbf{P}(X'_{s'},\sh{F}'_{s'},\mathbf{k}(s'))$ is true is equal to $g^{-1}(E)$}.
This extends immediately to the case where there appear in $\mathbf{P}$ several preschemes, Modules, morphisms of preschemes or homomorphisms of Modules, as well as to properties of the type considered in~(ii).

  \item As we will see in the remainder of this paragraph and in the rest of chap.~IV, the majority of properties $\mathbf{P}$ satisfying the condition~\sref{IV.9.2.1}[i] also satisfy the condition~\sref{IV.9.2.1}[ii].
As possible exceptions, let us note the property of being \emph{projective}, or \emph{quasi-projective}, or \emph{affine}, or \emph{quasi-affine} over the base field (for an algebraic prescheme); we will see~\sref{IV.9.6.2} that these properties are \emph{ind-constructible}, but we will give later an example where $S$ is a non-empty open part of $\operatorname{Spec}(\mathbf{Z})$ (or an open part of an elliptic curve over a finite field) and where all fibres $X_s$

\oldpage[IV-3]{58}
except the generic fibre $X_\eta$ are projective over $\mathbf{k}(s)$ (all preschemes $X_s$ being of dimension~2).
\end{enumerate}
\end{remarks}

\begin{proposition}[9.2.3]
\label{IV.9.2.3}
Let $\mathbf{P}$ be a constructible \textnormal{(}resp.\ ind-constructible\textnormal{)} property of algebraic preschemes, $S$ a prescheme, $X$ a prescheme of finite presentation over $S$, $\sh{F}$ a quasi-coherent $\sh{O}_X$-Module of finite presentation.
Then the set $E$ of $s\in S$ such that $\mathbf{P}(X_s,\sh{F}_s,\mathbf{k}(s))$ is true is locally constructible \textnormal{(}resp.\ ind-constructible\textnormal{)}.
Furthermore, if $S$ is irreducible of generic point $\eta$, one of the two sets $E$, $S-E$ is a neighbourhood of $\eta$ \textnormal{(}resp.\ $E$ is a neighbourhood of $\eta$ if it contains this point\textnormal{)}.
\end{proposition}

To prove these assertions, one can restrict to the case where $S=\operatorname{Spec}(A)$ is affine.
We know then that there exists a subring $A_0$ of $A$ which is a $\mathbf{Z}$-algebra of finite type, an $A_0$-prescheme of finite type $X_0$ and a coherent $\sh{O}_{X_0}$-Module $\sh{F}_0$ such that $X$ is isomorphic to $X_0\otimes_{A_0}A$ and $\sh{F}$ to $\sh{F}_0\otimes_{A_0}A$~\sref{IV.8.9.1}.
Let $p:S\to S_0=\operatorname{Spec}(A_0)$ be the morphism corresponding to the injection $A_0\to A$, and let $E_0$ be the set of $s_0\in S_0$ such that
\[
  \mathbf{P}((X_0)_{s_0},(\sh{F}_0)_{s_0},\mathbf{k}(s_0))
\]
is true; then, by virtue of~\sref{IV.9.2.2}[iv], we have $E=p^{-1}(E_0)$; by~\sref{IV.1.8.2}, we may therefore restrict to the case where $S$ is the spectrum of a $\mathbf{Z}$-algebra of finite type, hence a Noetherian scheme.
Use the constructibility criterion~($\mathbf{0}_{\mathrm{III}}$,~9.2.3) (resp.\ the ind-constructibility criterion~\sref{IV.1.9.10}); using~\sref{IV.9.2.2}[iv] as above and replacing $S$ by a closed integral subprescheme of $S$, we are reduced to the case where $S$ is Noetherian and integral and where it is necessary to prove that $E$ is nowhere dense in $S$ or contains a non-empty open subset of $S$ (resp.\ that $E$ contains a non-empty open subset of $S$ if it contains the generic point); this is guaranteed by condition~\sref{IV.9.2.1}[ii].

We will note that we have used~\sref{IV.9.2.1}[ii] only when $S$ is the spectrum of an integral ring of finite type over $\mathbf{Z}$.
It is clear, on the other hand, that the statement of~\sref{IV.9.2.3} also applies when there appear in $\mathbf{P}$ several preschemes, Modules on these preschemes, morphisms of preschemes or homomorphisms of Modules.
It applies also when there appear parts (in finite number) of the preschemes considered, provided that these parts are required to be \emph{locally constructible}.
Indeed, restriction to the case where $S$ is affine shows that one may restrict to the case where these parts are \emph{constructible}: one then applies~{\hyperlink{ega4.c04.thm.8-3-11}{(8.3.11)}}, which shows (with the preceding notation) that a constructible part of $X$ is the inverse image of a constructible part of $X_0$ for a suitable choice of $A_0$.

\begin{corollary}[9.2.4]
\label{IV.9.2.4}
Let $\mathbf{P}$ be a constructible \textnormal{(}resp.\ ind-constructible\textnormal{)} property of algebraic preschemes, $X$, $Y$ two $S$-preschemes of finite presentation, $f:X\to Y$ an $S$-morphism.
For every $s\in S$, set $X_s=X\times_S\operatorname{Spec}(\mathbf{k}(s))$, $Y_s=Y\times_S\operatorname{Spec}(\mathbf{k}(s))$, $f_s=f\times 1:X_s\to Y_s$.
Then the set $E$ of $s\in S$ such that for every $y\in Y_s$, the property $\mathbf{P}(f_s^{-1}(y),\mathbf{k}(y))$ is true, is locally constructible \textnormal{(}resp.\ ind-constructible\textnormal{)}.
\end{corollary}

Indeed, let $Z$ be the set of $y\in Y$ such that $\mathbf{P}(f^{-1}(y),\mathbf{k}(y))$ is true.
Since the fibres $f^{-1}(y)$ and $f_s^{-1}(y)$ are isomorphic, we see that $E$ is the set of $s\in S$ such that $Y_s\subset Z$; if $g:Y\to S$ is the structural morphism, we have therefore $E=S-g(Y-Z)$.

\oldpage[IV-3]{59}

Now, $f$ is of finite presentation~\sref{IV.1.6.2}[v], hence it follows from~\sref{IV.9.2.3} that $Z$ is locally constructible (resp.\ ind-constructible) in $Y$, hence $Y-Z$ is locally constructible (resp.\ pro-constructible) in $Y$.
As $g$ is of finite presentation, $g(Y-Z)$ is locally constructible (resp.\ pro-constructible) in $S$, by virtue of Chevalley's theorem~\sref{IV.1.8.4} (resp.\ of~\sref{IV.1.9.5}[vii]); hence $E$ is locally constructible (resp.\ ind-constructible) in $S$.

\begin{remark}[9.2.5]
\label{IV.9.2.5}
We will note that if $\mathbf{P}$ is a property of algebraic preschemes for which prop.~\sref{IV.9.2.3} is true, then $\mathbf{P}$ also satisfies the condition~\sref{IV.9.2.1}[ii]: this results indeed from the fact that in an irreducible Noetherian space, a constructible set is either nowhere dense or contains a non-empty open set~($\mathbf{0}_{\mathrm{III}}$,~9.2.3).
\end{remark}

\begin{proposition}[9.2.6]
\label{IV.9.2.6}
Let $\mathbf{P}$ denote one of the following properties of an algebraic $k$-prescheme $X$:
\begin{enumerate}[label=(\roman*)]
  \item $X$ is empty.
  \item $X$ is finite over $k$.
  \item $X$ is radicial over $k$.
  \item $\dim(X)$ belongs to a given part $\Phi$ of the set $\overline{\mathbf{Z}}=\mathbf{Z}\cup\{-\infty\}$.
\end{enumerate}
Then $\mathbf{P}$ is constructible.
\end{proposition}

It is clear that (i) and (ii) are particular cases of (iv), taking for $\Phi$ respectively the set $\{-\infty\}$ and the set $\{-\infty,0\}$.
We therefore need only prove (iii) and (iv).
In each of these two cases the condition~(i) of~\sref{IV.9.2.1} is satisfied by virtue of~\sref{IV.2.7.1}[xv] and~\sref{IV.4.1.4}.
On the other hand, in the case (iii), the property $\mathbf{P}$ satisfies the conclusion of~\sref{IV.9.2.3} by virtue of~\sref{IV.1.8.7}; it therefore remains to check that the same holds in the case (iv).
This will follow from the more precise proposition below:

\begin{proposition}[9.2.6.1]
\label{IV.9.2.6.1}
If $f:X\to S$ is a morphism of finite presentation, the function $s\mapsto\dim(f^{-1}(s))$ is locally constructible.
\end{proposition}

The question is local on $S$, so one can suppose that $S=\operatorname{Spec}(A)$ is affine and prove that for every $n$, the set $E$ of $s\in S$ such that $\dim(X_s)=n$ is constructible.
The same reasoning as in~\sref{IV.9.2.3} reduces to the case where $A$ is Noetherian and integral, and it suffices then to prove the following

\begin{corollary}[9.2.6.2]
\label{IV.9.2.6.2}
Let $S$ be an integral Noetherian prescheme of generic point $\eta$, $f:X\to S$ a morphism of finite type.
Then there exists a neighbourhood of $\eta$ in $S$ such that the function $s\mapsto\dim(X_s)$ is constant in this neighbourhood.
\end{corollary}

The images by $f$ of the irreducible components (in finite number) of $X$ that do not meet $X_\eta$ are contained in closed parts of $S$ not containing $\eta$ (since $S$ is integral~($\mathbf{0}_{\mathrm{I}}$,~2.1.5)); hence, after replacing $S$ by an open neighbourhood of $\eta$, we may restrict to the case where all the irreducible components $X_i$ of $X$ meet $X_\eta$.
Denote again by $X_i$ the reduced closed subprescheme of $X$ having $X_i$ as its underlying space; since $\dim(X_s)=\sup_i\dim((X_i)_s)$~\sref{IV.4.1.1}, we may restrict to the case where $X$ is irreducible.
There then exists a finite covering $(U_j)$ of $X$ by

\oldpage[IV-3]{60}
dense affine open sets, and the numbers $\dim((U_j)_\eta)$ are all equal to $n=\dim(X_\eta)$~\sref{IV.4.1.1.3}; one can therefore restrict to the case where $X$ is affine, hence also $X_\eta$.
There exists then, by virtue of~\sref{IV.4.1.2}, a non-empty open set $W$ of $X$ such that $W\supset X_\eta$, and a finite surjective $\mathbf{k}(\eta)$-morphism $h:W_\eta\to\mathbf{V}_{\mathbf{k}(\eta)}^n(=\operatorname{Spec}(\mathbf{k}(\eta)[T_1,\ldots,T_n]))$; applying~\sref{IV.8.8.2}[i] and the method of~\sref{IV.8.1.2}[a], we deduce (after replacing $S$ by a neighbourhood of $\eta$ if necessary) that $h=g_\eta$, where $g:W\to\mathbf{V}_A^n(=\operatorname{Spec}(A[T_1,\ldots,T_n]))$ is an $S$-morphism, and we may suppose this morphism finite and surjective by virtue of~\sref{IV.8.10.5}[vi] and~\sref{IV.8.10.5}[x].
We therefore conclude that for every $s\in S$, the morphism $g_s:W_s\to\mathbf{V}_{\mathbf{k}(s)}^n$ is finite and surjective, hence $\dim(W_s)=n$~\sref{IV.4.1.2}.

\subsection{Constructible properties of morphisms of algebraic preschemes}
\label{subsection:IV.9.3}

\begin{proposition}[9.3.1]
\label{IV.9.3.1}
Let $\mathbf{P}(X,k)$ be a constructible \textnormal{(}resp.\ ind-constructible\textnormal{)} property of algebraic preschemes.
Let $\mathbf{P}'(f,X,Y,k)$ denote the following relation: $f:X\to Y$ is a $k$-morphism of algebraic $k$-preschemes such that for every $y\in Y$, the property $\mathbf{P}(f^{-1}(y),\mathbf{k}(y))$ is true.
Then $\mathbf{P}'$ is a constructible \textnormal{(}resp.\ ind-constructible\textnormal{)} property.
\end{proposition}

Indeed, as $\mathbf{P}$ satisfies the condition~\sref{IV.9.2.1}[i], the same is true of $\mathbf{P}'$ by virtue of the transitivity of fibres~(\textbf{I},~3.6.4); furthermore the fact that $\mathbf{P}'$ satisfies the condition~\sref{IV.9.2.1}[ii] follows from~\sref{IV.9.2.4}, by reason of the remark~\sref{IV.9.2.5}.

\begin{proposition}[9.3.2]
\label{IV.9.3.2}
Let $\mathbf{P}$ denote one of the following properties of a $k$-morphism $f:X\to Y$ of algebraic preschemes:
\begin{enumerate}[label=(\roman*)]
  \item $f$ is surjective.
  \item $f$ is quasi-finite.
  \item $f$ is radicial.
  \item For every $y\in Y$, $\dim(f^{-1}(y))$ belongs to $\Phi$ \textnormal{(}notation of~\sref{IV.9.2.6}\textnormal{)}.
\end{enumerate}
Then $\mathbf{P}$ is constructible.
\end{proposition}

This follows immediately from~\sref{IV.9.3.1} and~\sref{IV.9.2.6} if one takes into account that $f$ is of finite type~\sref[I]{I.1.5.4}[v], the characterisation of radicial morphisms~\sref[I]{I.3.5.8}, and that of quasi-finite morphisms~\sref[II]{II.6.2.2}.

\begin{proposition}[9.3.3]
\label{IV.9.3.3}
Suppose that the hypotheses of~\sref{IV.8.8.1} are satisfied, and let us keep the same notations; suppose furthermore that $S_\alpha$ is quasi-compact, $X_\alpha$ and $Y_\alpha$ of finite presentation over $S_\alpha$, and let $f_\alpha:X_\alpha\to Y_\alpha$ be an $S_\alpha$-morphism.
Let $\mathbf{P}$ be an ind-constructible property of morphisms of algebraic preschemes.
For every $s\in S$ \textnormal{(}resp.\ $s_\lambda\in S_\lambda$\textnormal{)} set $X_s=X\times_S\operatorname{Spec}(\mathbf{k}(s))$, $Y_s=Y\times_S\operatorname{Spec}(\mathbf{k}(s))$, $f_s=f\times 1:X_s\to Y_s$ \textnormal{(}resp.\ $X_{\lambda,s_\lambda}=X_\lambda\times_{S_\lambda}\operatorname{Spec}(\mathbf{k}(s_\lambda))$, $Y_{\lambda,s_\lambda}=Y_\lambda\times_{S_\lambda}\operatorname{Spec}(\mathbf{k}(s_\lambda))$, $f_{\lambda,s_\lambda}=f_\lambda\times 1:X_{\lambda,s_\lambda}\to Y_{\lambda,s_\lambda}$\textnormal{)}.
Then, in order that for every $s\in S$ the property $\mathbf{P}(f_s)$ be true, it is necessary and sufficient that there exist $\lambda\geqslant\alpha$ such that for every $s_\lambda\in S_\lambda$ we have $\mathbf{P}(f_{\lambda,s_\lambda})$.
\end{proposition}

Indeed, let $E$ (resp.\ $E_\lambda$) be the set of $s\in S$ (resp.\ $s_\lambda\in S_\lambda$) such that the property $\mathbf{P}(f_s)$ (resp.\ $\mathbf{P}(f_{\lambda,s_\lambda})$) is true; it follows from~\sref{IV.9.2.2}[iv] that $E_\mu=u_{\lambda\mu}^{-1}(E_\lambda)$ for $\lambda\leqslant\mu$, and $E=u_\lambda^{-1}(E_\lambda)$; furthermore, by virtue of~\sref{IV.9.2.3}, $E$ (resp.\ $E_\lambda$) is ind-constructible in $S$ (resp.\ $S_\lambda$); the proposition then follows from~\sref{IV.8.3.4} applied to the ind-constructible part $E$ of $S$.

\oldpage[IV-3]{61}

This result extends without difficulty to properties $\mathbf{P}$ of the type considered in (9.2.3, (i) and (ii)).

\begin{remark}[9.3.4]
\label{IV.9.3.4}
The conjunction of~\sref{IV.9.3.3} and~\sref{IV.9.3.2}[ii] proves the assertion~\sref{IV.8.10.5}[xi].
\end{remark}

\begin{proposition}[9.3.5]
\label{IV.9.3.5}
Let $\mathbf{P}(X,Y,k)$ be the property: ``$X$ and $Y$ are two preschemes of finite type over the field $k$, and there exists an extension $k'$ of $k$ and a $k'$-morphism $g:X_{(k')}\to Y_{(k')}$ satisfying $\mathbf{Q}(g)$'', where $\mathbf{Q}$ is one of the properties~\emph{(i)} to~\emph{(xiv)} of~\sref{IV.8.10.5}.
Then $\mathbf{P}$ is an ind-constructible property.
\end{proposition}

The definition of $\mathbf{P}$ shows indeed that the condition~\sref{IV.9.2.1}[i] is satisfied, taking into account that the property $\mathbf{Q}(g)$ is stable under base change, and that two extensions of $k$ can always be considered as sub-extensions of a third extension of $k$.
To verify~\sref{IV.9.2.1}[ii], one can restrict to the case where $S=\operatorname{Spec}(A)$ is affine; if $K=\mathbf{k}(\eta)$ is the field of fractions of $A$, there exists by hypothesis and by virtue of~\sref{IV.9.1.4} a finite extension $K'$ of $K$ and a $K'$-morphism $g':(X_\eta)_{(K')}\to(Y_\eta)_{(K')}$ satisfying $\mathbf{Q}(g')$, and $K'$ is evidently the field of fractions of a finite integral $A$-algebra $A'$.
If we set $S'=\operatorname{Spec}(A')$, it follows then from~\sref{IV.8.10.5} that there is a neighbourhood $U'$ of the generic point $\eta'$ of $S'$ such that, if we set $X'=X\otimes_A A'$ and $Y'=Y\otimes_A A'$, then for every $s'\in U'$ there exists a morphism $X'_{s'}\to Y'_{s'}$ having property $\mathbf{Q}$.
But the morphism $h:S'\to S$ is finite, hence closed; since $h^{-1}(\eta)=\{\eta'\}$, $h(U')$ contains an open neighbourhood $U$ of $\eta$ in $S$.
For every $s\in U$, there therefore exists $s'\in U'$ such that $h(s')=s$, and as $X'_{s'}=X_s\otimes_{\mathbf{k}(s)}\mathbf{k}(s')$ and $Y'_{s'}=Y_s\otimes_{\mathbf{k}(s)}\mathbf{k}(s')$, the property $\mathbf{P}(X_s,Y_s,\mathbf{k}(s))$ is true for every $s\in U$.

\begin{env}[9.3.6]
\label{IV.9.3.6}
\textbf{Example.}
Take for example for $\mathbf{Q}$ the property of being an isomorphism.
Then, combining~\sref{IV.9.3.5} and~\sref{IV.9.3.3}, we obtain the following property: with the notation and hypotheses of~\sref{IV.8.8.1}, with $S_\alpha$ quasi-compact and $X_\alpha$ and $Y_\alpha$ of finite presentation over $S_\alpha$, in order that, for every $s\in S$, $X_s$ and $Y_s$ be \emph{geometrically isomorphic}~\sref{IV.9.1.4}, it is necessary and sufficient that there exist $\lambda\geqslant\alpha$ such that, for every $s_\lambda\in S_\lambda$, $X_{\lambda,s_\lambda}$ and $Y_{\lambda,s_\lambda}$ be geometrically isomorphic.
\end{env}

We have an analogous result when the preschemes considered are endowed with ``composition laws'' of a certain kind~($\mathbf{0}_{\mathrm{III}}$,~8.2.1), for example ``preschemes in groups'', ``preschemes in rings'', etc.~($\mathbf{0}_{\mathrm{III}}$,~8.2.3).
Then the preceding statement is still valid when by ``isomorphism'' one means isomorphisms of preschemes that are homomorphisms for the composition laws considered~($\mathbf{0}_{\mathrm{III}}$,~8.2.2); it suffices here to use not only~\sref{IV.8.10.5} but also~\sref{IV.8.8.2}[i], noting that the notion of homomorphism for a composition law is expressed by saying that diagrams of morphisms of preschemes are commutative (it must be understood that the transition morphisms $X_\mu\to X_\lambda$ and $Y_\mu\to Y_\lambda$ for $\lambda\leqslant\mu$ are homomorphisms for the composition laws considered).

One can also, instead of considering morphisms of preschemes as in~\sref{IV.9.3.5}, consider homomorphisms of Modules, using~\sref{IV.9.1.6} instead of~\sref{IV.9.1.5}.

\oldpage[IV-3]{62}

\subsection{Constructibility of certain properties of modules}
\label{subsection:IV.9.4}

\begin{env}[9.4.1]
\label{IV.9.4.1}
\textbf{Notations (9.4.1).}
In this n\textsuperscript{o} and the following ones until the end of \textsection9, we will systematically use the following notations: given a morphism $f:X\to S$, we will set, for every $s\in S$, $X_s=f^{-1}(s)=X\times_S\operatorname{Spec}(\mathbf{k}(s))$; for every $\sh{O}_X$-Module quasi-coherent $\sh{F}$, $\sh{F}_s$ will denote the $\sh{O}_{X_s}$-Module quasi-coherent $\sh{F}\otimes_{\sh{O}_S}\mathbf{k}(s)$, and for every homomorphism $u:\sh{F}\to\sh{G}$ of $\sh{F}$ into a quasi-coherent $\sh{O}_X$-Module $\sh{G}$, $u_s:\sh{F}_s\to\sh{G}_s$ will be the morphism $p^*(u)$, designating by $p$ the canonical projection $X_s\to X$.
For every section $g$ of $\sh{F}$ over $X$, we will denote by $g_s$ the image of $g$ by the canonical homomorphism $\Gamma(X,\sh{F})\to\Gamma(X_s,\sh{F}_s)$.
For every part $Z$ of $X$, we will denote $Z_s$ the inverse image $p^{-1}(Z)=Z\cap X_s$~(\textbf{I},~3.6.1).
Finally, if $Y$ is a second $S$-prescheme and $h:X\to Y$ an $S$-morphism, we will denote $h_s$ the morphism $h\times 1:X_s\to Y_s$.
\end{env}

\begin{proposition}[9.4.2]
\label{IV.9.4.2}
Let $S$ be an integral prescheme of generic point $\eta$, $f:X\to S$ a morphism of finite presentation, $\sh{F}$, $\sh{G}$, $\sh{H}$ three quasi-coherent $\sh{O}_X$-Modules of finite presentation.
Let $u:\sh{F}\to\sh{G}$, $v:\sh{G}\to\sh{H}$ be two homomorphisms of $\sh{O}_X$-Modules, and suppose that the sequence $\sh{F}_\eta\xrightarrow{u_\eta}\sh{G}_\eta\xrightarrow{v_\eta}\sh{H}_\eta$ is exact.
Then there exists a neighbourhood $U$ of $\eta$ in $S$ such that, for every $s\in U$, the sequence $\sh{F}_s\xrightarrow{u_s}\sh{G}_s\xrightarrow{v_s}\sh{H}_s$ is exact.
\end{proposition}

\begin{proof}
With the general notations of~\sref{IV.9.2.1}, it is a question here of the relation $\mathbf{P}(X,\sh{F},\sh{G},\sh{H},u,v,k)$: ``$X$ is an algebraic prescheme over the field $k$, $\sh{F}\xrightarrow{u}\sh{G}\xrightarrow{v}\sh{H}$ is an exact sequence of quasi-coherent $\sh{O}_X$-Modules''.
Since, for every extension $k'$ of $k$, the canonical projection $X_{(k')}\to X$ is a faithfully flat morphism~\sref{IV.2.2.13}, the condition~\sref{IV.9.2.1},~(i)) is satisfied~\sref{IV.2.2.7}.
By virtue of~\sref{IV.9.2.3}, one can restrict to the case where $S$ is integral and Noetherian, in which case $X$ is Noetherian, and $\sh{F}$, $\sh{G}$, $\sh{H}$ are coherent $\sh{O}_X$-Modules.
The hypothesis implies that there exists a neighbourhood $U$ of $\eta$ in $S$ such that $\sh{F}|f^{-1}(U)\to\sh{G}|f^{-1}(U)\to\sh{H}|f^{-1}(U)$ is \emph{exact}, by virtue of~\sref{IV.8.5.8},~(ii)) applied following the general method of~\sref{IV.8.1.2},~$a$)), and one can therefore already suppose that $\sh{F}\xrightarrow{u}\sh{G}\xrightarrow{v}\sh{H}$ is exact; it suffices evidently to prove that $\operatorname{Ker}(v_s)=(\operatorname{Ker} v)_s$ and $\operatorname{Im}(u_s)=(\operatorname{Im} u)_s$ for $s$ near $\eta$ in $S$.
Taking account of the coherence of $\sh{F}$, $\sh{G}$ and $\sh{H}$ and of
($\mathbf{0}_{\mathrm{I}}$,~5.3.4), we are reduced to the particular case where
\[
 0\longrightarrow\sh{F}\longrightarrow\sh{G}\longrightarrow\sh{H}
 \longrightarrow0
\]
is exact.
There is then an open neighbourhood $U$ of $\eta$ in $S$ such that
$\sh{H}|f^{-1}(U)$ is flat over $U$ \sref{IV.6.9.1}; hence \sref{IV.2.1.8}
shows that for every $s\in U$ the sequence
\[
 0\longrightarrow\sh{F}_s\longrightarrow\sh{G}_s\longrightarrow\sh{H}_s
 \longrightarrow0
\]
is exact.
\end{proof}

\begin{corollary}[9.4.3]
\label{IV.9.4.3}
Let $S$ be an integral prescheme, of generic point $\eta$, $f:X\to S$ a morphism of finite presentation.
Let $\sh{L}^\bullet=(\sh{L}^i)_{i\in\mathbf{Z}}$ be a complex whose terms are quasi-coherent $\sh{O}_X$-Modules of finite presentation.
For every $i$, there exists an open neighbourhood $U$ of $\eta$ in $S$ such that the canonical homomorphisms
\begin{equation}
\label{IV.9.4.3.1}
  (\sh{H}^i(\sh{L}^\bullet))_s\to\sh{H}^i(\sh{L}^\bullet_s) \tag{9.4.3.1}
\end{equation}
are bijective for every $s\in U$.
\end{corollary}

\oldpage[IV-3]{63}

\begin{proof}
It is enough to consider a three-term complex in degrees $-1$, $0$, $+1$,
\[
 \sh{M}\xrightarrow{u}\sh{N}\xrightarrow{v}\sh{P},
\]
and $i=0$.
The canonical homomorphism in question is then
\[
 (\operatorname{Ker}v/\operatorname{Im}u)_s
 \longrightarrow \operatorname{Ker}(v_s)/\operatorname{Im}(u_s).
\]
Using \sref{IV.8.9.1} and \sref{IV.8.5.2},~(i), we may reduce to the case
where $S$ (and hence $X$) is Noetherian; the Modules $\sh{M}$, $\sh{N}$ and
$\sh{P}$ are then coherent.
Thus $\operatorname{Im}u$ and $\operatorname{Ker}v$ are coherent
($\mathbf{0}_{\mathrm{I}}$,~5.3.4), and \sref{IV.9.4.2} gives a neighbourhood
$U$ of $\eta$ on which
\[
 \operatorname{Ker}(v_s)=(\operatorname{Ker}v)_s,
 \qquad
 \operatorname{Im}(u_s)=(\operatorname{Im}u)_s.
\]
The conclusion follows by applying \sref{IV.9.4.2} to the exact sequence
\[
 0\longrightarrow\operatorname{Im}u\longrightarrow\operatorname{Ker}v
 \longrightarrow\operatorname{Ker}v/\operatorname{Im}u\longrightarrow0,
\]
since $\sh{O}_\eta=\boldsymbol{k}(\eta)$ and consequently the corresponding
sequence over $\eta$ is exact.
\end{proof}

\begin{proposition}[9.4.4]
\label{IV.9.4.4}
Let $f:X\to S$ be a morphism of finite presentation, $\sh{F}$, $\sh{G}$, $\sh{H}$ three quasi-coherent $\sh{O}_X$-Modules of finite presentation.
Let $u:\sh{F}\to\sh{G}$, $v:\sh{G}\to\sh{H}$ be two homomorphisms of $\sh{O}_X$-Modules.
Then the set $E$ of $s\in S$ such that the sequence $\sh{F}_s\xrightarrow{u_s}\sh{G}_s\xrightarrow{v_s}\sh{H}_s$ is exact is locally constructible.
\end{proposition}

Taking into account~\sref{IV.9.2.3}, it suffices to establish that the property $\mathbf{P}$ considered in~\sref{IV.9.4.2} is constructible.
We have already remarked in the proof of~\sref{IV.9.4.2} that the condition~\sref{IV.9.2.1},~(i)) is satisfied for this property, and it remains to verify the condition~\sref{IV.9.2.1},~(ii)).
Let us therefore suppose $S$ integral Noetherian, of generic point $\eta$, and let us prove that $E$ or $S-E$ is a neighbourhood of $\eta$.
If $\eta\in E$, our assertion follows from~\sref{IV.9.4.2}, and one can therefore restrict to the case where $\eta\notin E$, in other words, the sequence $\sh{F}_\eta\xrightarrow{u_\eta}\sh{G}_\eta\xrightarrow{v_\eta}\sh{H}_\eta$ is not exact.
Let us then distinguish two cases:

$1^\circ$ Set $w=v\circ u$, and suppose first that $w_\eta=v_\eta\circ u_\eta\neq 0$.
As $\sh{F}$, $\sh{G}$, $\sh{H}$ are coherent, the same is true of $\sh{N}=\operatorname{Ker}(w)$~($\mathbf{0}_{\mathrm{I}}$,~5.3.4); it then follows from~\sref{IV.9.4.2} applied to the exact sequence $0\to\sh{N}\to\sh{F}\xrightarrow{w}\sh{H}$ that there exists a neighbourhood $U$ of $\eta$ in $S$ such that $\operatorname{Ker}(w_s)=\sh{N}_s$; restricting $S$, one can therefore suppose this relation satisfied for every $s\in S$.
Let $j$ be the canonical injection $\sh{N}\to\sh{F}$, and set $\sh{M}=\operatorname{Coker}(j)$; the right exactness of the functor $\sh{F}\mapsto\sh{F}_s$ entails that $\sh{M}_s=\operatorname{Coker}(j_s)$ for every $s\in S$.
The hypothesis $w_\eta\neq 0$ means that $\sh{M}_\eta\neq 0$; as $\sh{M}$ is coherent~($\mathbf{0}_{\mathrm{I}}$,~5.3.4), it follows from~\sref{IV.1.8.6} that there exists a neighbourhood $U$ of $\eta$ in $S$ such that $\sh{M}_s\neq 0$ for every $s\in U$, hence $w_s\neq 0$ for every $s\in U$, and \emph{a fortiori} $S-E$ is a neighbourhood of $\eta$.

$2^\circ$ Suppose that $w_\eta=0$; by virtue of~\sref{IV.8.5.2},~(i)), applied following the general method of~\sref{IV.8.1.2},~$a$)), there exists a neighbourhood $U$ of $\eta$ such that $w|f^{-1}(U)=0$; replacing $S$ by $U$, one can already suppose $w=0$ in $X$.
Then $\sh{F}\xrightarrow{u}\sh{G}\xrightarrow{v}\sh{H}$ is a complex of three terms, to which one can apply~\sref{IV.9.4.3}; by hypothesis we have $(\sh{H}^0(\sh{L}^\bullet))_\eta=\sh{H}^0(\sh{L}^\bullet_\eta)\neq 0$, and $\sh{H}^0(\sh{L}^\bullet)$ is coherent~($\mathbf{0}_{\mathrm{I}}$,~5.3.4 and~5.3.3), hence it follows from~\sref{IV.1.8.6} that there exists a neighbourhood $U$ of $\eta$ such that $(\sh{H}^0(\sh{L}^\bullet))_s\neq 0$ for every $s\in U$; but as one can suppose that $\sh{H}^0(\sh{L}^\bullet_s)=(\sh{H}^0(\sh{L}^\bullet))_s$ for every $s\in U$ by~\sref{IV.9.4.3}, we see again that $S-E$ is a neighbourhood of $\eta$ in $S$.

\oldpage[IV-3]{64}

\begin{corollary}[9.4.5]
\label{IV.9.4.5}
Let $f:X\to S$ be a morphism of finite presentation, $\sh{F}$, $\sh{G}$ two quasi-coherent $\sh{O}_X$-Modules of finite presentation, $u:\sh{F}\to\sh{G}$ a homomorphism of $\sh{O}_X$-Modules.
Then the set of $s\in S$ such that $u_s$ is injective \textnormal{(}resp.\ surjective, bijective, zero\textnormal{)} is locally constructible.
\end{corollary}

It suffices to apply~\sref{IV.9.4.4} to the sequences $0\to\sh{F}\xrightarrow{u}\sh{G}$, $\sh{F}\xrightarrow{u}\sh{G}\to 0$, $0\to\sh{F}\xrightarrow{u}\sh{G}\to 0$, $\sh{F}\xrightarrow{u}\sh{G}\xrightarrow{1_\sh{G}}\sh{G}$.

\begin{corollary}[9.4.6]
\label{IV.9.4.6}
Let $f:X\to S$ be a morphism of finite presentation, $\sh{F}$ a quasi-coherent $\sh{O}_X$-Module of finite presentation.
Let $h$ be a section of $\sh{F}$ over $X$; for every $s\in S$, let $h_s$ be the corresponding section of $\sh{F}_s$ over $X_s$ \textnormal{(}for the projection morphism $X_s\to X$\textnormal{)}.
Then the set of $s\in S$ such that $h_s=0$ is locally constructible.
\end{corollary}

It suffices to note that $h$ corresponds to a homomorphism $u:\sh{O}_X\to\sh{F}$~($\mathbf{0}_{\mathrm{I}}$,~5.1.1) and $h_s$ to the homomorphism $u_s$.

\begin{proposition}[9.4.7]
\label{IV.9.4.7}
Let $f:X\to S$ be a morphism of finite presentation, $\sh{F}$ a quasi-coherent $\sh{O}_X$-Module of finite presentation.
The set $E$ \textnormal{(}resp.\ $E'$\textnormal{)} of $s\in S$ such that $\sh{F}_s$ is a locally free $\sh{O}_{X_s}$-module \textnormal{(}resp.\ locally free of rank $n$\textnormal{)} is locally constructible.
\end{proposition}

If $X$ is an algebraic prescheme over a field $k$, $\sh{F}$ a coherent $\sh{O}_X$-Module, $k'$ an extension of $k$, then, for $\sh{F}$ to be locally free (resp.\ locally free of rank $n$), it is necessary and sufficient that it be so of $\sh{F}\otimes_k k'$, since the projection $X_{(k')}\to X$ is a faithfully flat morphism~\sref{IV.2.2.7}.
In other words, the condition~\sref{IV.9.2.1},~(i)) is satisfied for the properties we wish to prove constructible, and it remains to verify~\sref{IV.9.2.1},~(ii)); one can therefore still suppose that $S$ is affine, Noetherian and integral.
There are then four cases to consider:

$1^\circ$ $\eta\in E$.
By virtue of~\sref{IV.8.5.5}, applied following the general method of~\sref{IV.8.1.2},~$a$)), there exists a neighbourhood $U$ of $\eta$ in $S$ such that $\sh{F}|f^{-1}(U)$ is locally free; \emph{a fortiori} $\sh{F}_s$ is locally free for every $s\in U$.

$2^\circ$ $\eta\in E'$.
Same reasoning as in $1^\circ$.

$3^\circ$ $\eta\in S-E$.
As $\sh{F}_\eta$ is a coherent $\sh{O}_{X_\eta}$-module, to say that it is not locally free means that it is not flat over $\sh{O}_{X_\eta}$ (Bourbaki, \emph{Alg.\ comm.}, chap.~II, \textsection5, n\textsuperscript{o}~2, cor.~2 of th.~1).
The fact that $S-E$ is a neighbourhood of $\eta$ will then follow from the more general lemma below:

\begin{lemma}[9.4.7.1]
\label{IV.9.4.7.1}
Let $S$ be an integral Noetherian prescheme, of generic point $\eta$, $X$, $Y$ two $S$-preschemes of finite type over $S$, $g:X\to Y$ an $S$-morphism, $\sh{F}$ a coherent $\sh{O}_X$-Module.
If $\sh{F}_\eta$ is not $g_\eta$-flat, then there exists a neighbourhood $U$ of $\eta$ in $S$ such that for every $s\in U$, $\sh{F}_s$ is not $g_s$-flat.
\end{lemma}

Taking into account~\sref{IV.2.1.2} and Bourbaki, \emph{Alg.\ comm.}, chap.~I\textsuperscript{er}, \textsection2, n\textsuperscript{o}~3, Remarque~1, the hypothesis means that there exists a non-empty open $V$ of $Y_\eta$, and an injective homomorphism $v:\sh{M}\to\sh{N}$ of $\sh{O}_V$-Modules \emph{coherent}, such that the homomorphism $1\otimes v:\sh{F}_\eta\otimes_{\sh{O}_V}\sh{M}\to\sh{F}_\eta\otimes_{\sh{O}_V}\sh{N}$ is not injective.
We have $V=Y_\eta\cap W$, where $W$ is open in $Y$~(\textbf{I},~3.6.1), and it follows from~\sref{IV.8.5.2},~(i) and~(ii)), applied following the method of~\sref{IV.8.1.2},~$a$)), that there exists a neighbourhood $U_0$ of $\eta$ in $S$, two $\sh{O}_Z$-Modules
\oldpage[IV-3]{65}
coherent $\sh{G}$, $\sh{H}$ (where $Z=W\cap h^{-1}(U_0)$, $h:Y\to S$ being the structural morphism) and an $\sh{O}_Z$-homomorphism $u:\sh{G}\to\sh{H}$ such that $\sh{M}=\sh{G}_\eta$, $\sh{N}=\sh{H}_\eta$, and $v=u_\eta$; one may therefore take $U_0$ so that $u_s:\sh{G}_s\to\sh{H}_s$ is injective for every $s\in U_0$~\sref{IV.9.4.5}.
But for every $s\in U_0$, the homomorphism $1\otimes u_s:\sh{F}_s\otimes_{\sh{O}_{Z_s}}\sh{G}_s\to\sh{F}_s\otimes_{\sh{O}_{Z_s}}\sh{H}_s$ is nothing but $(1\otimes u)_s$; the hypothesis that $(1\otimes u)_\eta$ is not injective therefore entails~\sref{IV.9.4.5} the existence of a non-empty open $U\subset U_0$ such that for every $s\in U$, $(1\otimes u)_s$ is not injective, and hence $\sh{F}_s$ is not $g_s$-flat for every $s\in U$.

$4^\circ$ $\eta\in S-E'$.
It is clear that $S-E\subset S-E'$, and if $\eta\in S-E$, $S-E'$ is \emph{a fortiori} a neighbourhood of $\eta$ by $3^\circ$.
Let us therefore suppose $\eta\in E$, hence $\sh{F}_\eta$ is locally free; these hypotheses entail that $X_\eta$ is non-connected, and that the ranks of the $\sh{O}_{X_\eta}$-Module locally free $\sh{F}_\eta$ are not the same on the various connected components of $X_\eta$.
Now, it follows from \sref{IV.8.4.2}, applied by the method of \sref{IV.8.1.2},~$a)$, that one may suppose (after replacing $S$ by a neighbourhood of $\eta$) that $X$ and $X_\eta$ have the same number of connected components, the connected components of $X_\eta$ being the intersections of $X_\eta$ with the connected components of $X$.
The conclusion then follows from the reasoning given in $2^\circ$, applied to each of the connected components of $X$ (which are finite in number).

\begin{remark}[9.4.7.2]
\label{IV.9.4.7.2}
The lemma~\sref{IV.9.4.7.1} will be later generalised and freed of Noetherian hypotheses~\sref{IV.11.2.8}.
\end{remark}

\begin{proposition}[9.4.8]
\label{IV.9.4.8}
Let $S$ be a locally Noetherian prescheme, $f:X\to S$ a morphism of finite type, $\sh{F}$ a coherent $\sh{O}_X$-Module.
Suppose that for every $s\in S$, $X_s$ is a locally integral prescheme.
Then the set $E$ \textnormal{(}resp.\ $E'$\textnormal{)} of $s\in S$ such that $\sh{F}_s$ is a torsion $\sh{O}_{X_s}$-module \textnormal{(}resp.\ a torsion-free $\sh{O}_{X_s}$-module\textnormal{)} is locally constructible.
\end{proposition}

\begin{proof}
One can evidently suppose $S=\operatorname{Spec}(A)$ affine and Noetherian and prove that $E$ (resp.\ $E'$) is constructible by using the criterion~($\mathbf{0}_{\mathrm{III}}$,~9.2.3); replacing $S$ by the closed reduced sub-prescheme of $S$ having as underlying space a closed irreducible part $Y$ of $S$, and noting~(\textbf{I},~3.6.4) that for $s\in Y$, the fibre $(X_{(Y)})_s$ is canonically identified with $X_s$ and the sheaf $(\sh{F}\otimes_{\sh{O}_S}\sh{O}_Y)_s$ with $\sh{F}_s$, we see that one is reduced to the case where $S$ is integral of generic point $\eta$, and to proving that $E$ or $S-E$ (resp.\ $E'$ or $S-E'$) is a neighbourhood of $\eta$ in $S$.
Furthermore, $X$ is a finite union of affine open subsets $U_i$ of finite type over $S$, and each $(U_i)_s$ is an open subprescheme of $X_s$, hence locally integral. If $(U_i)_\eta$ is empty (resp.\ nonempty), then $(U_i)_s$ is also empty (resp.\ nonempty) in a neighbourhood of $\eta$ \sref{IV.9.2.6}. We may therefore suppose that all $(U_i)_\eta$ are nonempty and integral; then $\sh{F}_s$ is torsion (resp.\ torsion-free) if and only if each restriction $\sh{F}_s|_{(U_i)_s}$ is torsion (resp.\ torsion-free).

We are thus reduced to the case where $X=\operatorname{Spec}(B)$ is affine and $\sh{F}=\widetilde{M}$, with $M$ a finite $B$-module. Put $B_s=B\otimes_A\mathbf{k}(s)$ and $M_s=M\otimes_A\mathbf{k}(s)$, and suppose $B_\eta$ integral. There are four cases.

$1^\circ$ Suppose $\eta\in E$. Then $M_\eta$ is a finite torsion $B_\eta$-module, so there is an $h\neq 0$ in $B_\eta$ such that $hM_\eta=0$. By \sref{IV.8.5.2},~(i), applied as in \sref{IV.8.1.2},~$a)$,

\oldpage[IV-3]{66}

after shrinking $S$ we may write $h=g_\eta$ for some $g\in\Gamma(X,\sh{O}_X)$. Let $u:\sh{F}\to\sh{F}$ be multiplication by $g$. Since $u_\eta=0$, \sref{IV.9.4.5} shows that $u_s=0$ near $\eta$. On the other hand, multiplication by $g_\eta=h$ is injective on $\sh{O}_{X_\eta}$ because $B_\eta$ is integral and $h\neq 0$; again by \sref{IV.9.4.5}, multiplication by $g_s$ on $\sh{O}_{X_s}$ is injective near $\eta$. Thus $g_s$ is $\sh{O}_{X_s}$-regular and annihilates $\sh{F}_s$, so $\sh{F}_s$ is torsion near $\eta$.

$2^\circ$ Suppose $\eta\in S-E$. If $T$ is the torsion submodule of $M_\eta$, then $M_\eta/T\neq 0$ and, being a finite torsion-free $B_\eta$-module, it embeds into $B_\eta^n$ for some $n$. Hence there is a nonzero homomorphism $w:M_\eta\to B_\eta$. By \sref{IV.8.5.2},~(i), after shrinking $S$ it extends to a homomorphism $v:\sh{F}\to\sh{O}_X$. Since $v_\eta\neq 0$, \sref{IV.9.4.5} gives $v_s\neq 0$ near $\eta$; as $X_s$ is locally integral, $\sh{F}_s$ is not torsion for these $s$.

$3^\circ$ Suppose $\eta\in E'$. Since $M_\eta$ is a finite torsion-free $B_\eta$-module, there is an injective homomorphism $w:M_\eta\to B_\eta^n$. Using \sref{IV.8.5.2},~(i), and \sref{IV.9.4.5} as in $2^\circ$, after shrinking $S$ we obtain a homomorphism $v:\sh{F}\to\sh{O}_X^n$ such that $v_\eta=w$ and $v_s$ is injective near $\eta$. Thus $\sh{F}_s$ is torsion-free near $\eta$.

$4^\circ$ Suppose $\eta\in S-E'$. Let $T$ be the torsion submodule of $M_\eta$. Then $T\neq 0$ and is finite. By \sref{IV.8.5.2},~(i) and (ii), after shrinking $S$ there are a coherent $\sh{O}_X$-module $\sh{G}$ and an injective homomorphism $u:\sh{G}\to\sh{F}$ such that $\sh{G}_\eta=\widetilde{T}$ and $u_\eta$ is the canonical inclusion. By $1^\circ$ and \sref{IV.1.8.6}, near $\eta$ the module $\sh{G}_s$ is a nonzero torsion $\sh{O}_{X_s}$-module; by \sref{IV.9.4.5}, $u_s$ is injective near $\eta$. Therefore the torsion submodule of $\sh{F}_s$ is nonzero near $\eta$.
\end{proof}

\begin{remark}[9.4.9]
\label{IV.9.4.9}
The property ``$X$ is a locally integral algebraic $k$-prescheme'' does not satisfy condition \sref{IV.9.2.1},~(i), and it is therefore not certain that the statement of \sref{IV.9.4.8} remains valid when no hypothesis is imposed on $S$ and one assumes only that $f$ is a morphism of finite presentation and $\sh{F}$ is an $\sh{O}_X$-Module of finite presentation. Consider, however, the following special case. Let $S_0$ be a locally Noetherian prescheme, let $f_0:X_0\to S_0$ be a morphism of finite type whose fibres $(X_0)_{s_0}$ are locally integral for every $s_0\in S_0$, and let $\sh{F}_0$ be a coherent $\sh{O}_{X_0}$-Module. Let $g:S\to S_0$ be any morphism, set
$X=X_0\times_{S_0}S$, $\sh{F}=\sh{F}_0\otimes_{\sh{O}_{S_0}}\sh{O}_S$, and $f=f_0\times 1_S:X\to S$, and suppose that, for every $s\in S$, the fibre $X_s$ is again locally integral. Then the set $E$ (resp. $E'$) of $s\in S$ such that $\sh{F}_s$ is torsion (resp. torsion-free)

\oldpage[IV-3]{67}

is again locally constructible.

Indeed, let $s\in S$ and put $s_0=g(s)$. Taking account of \sref[I]{I.1.8.2}, it suffices to prove that $\sh{F}_s$ is torsion (resp. torsion-free) if and only if $(\sh{F}_0)_{s_0}$ is so. Now $\operatorname{Supp}(\sh{F}_s)$ is the inverse image of $\operatorname{Supp}((\sh{F}_0)_{s_0})$ under the projection $p:X_s\to(X_0)_{s_0}$ \sref[I]{I.9.1.13}. Since $p$ is faithfully flat and quasi-compact, saying that $\operatorname{Supp}(\sh{F}_s)$ contains a maximal point of $X_s$ is equivalent to saying that $\operatorname{Supp}((\sh{F}_0)_{s_0})$ contains a maximal point of $(X_0)_{s_0}$ \sref[I]{I.1.1.5} and \sref{IV.2.3.4}; this proves the assertion concerning $E$ \sref[I]{I.7.4.6}. If the torsion submodule $\sh{T}$ of $(\sh{F}_0)_{s_0}$ is nonzero, then $\sh{T}\otimes_{\boldsymbol{k}(s_0)}\boldsymbol{k}(s)$, which is torsion by what precedes, is nonzero and identifies with a submodule of $\sh{F}_s$ \sref{IV.2.2.7}; hence the torsion submodule of $\sh{F}_s$ is nonzero. Finally, if $(\sh{F}_0)_{s_0}$ is torsion-free, then after restricting to an affine open of $(X_0)_{s_0}$ we may suppose that it is isomorphic to a submodule of some $\sh{O}_{(X_0)_{s_0}}^n$; consequently $\sh{F}_s$ is isomorphic to a submodule of $\sh{O}_{X_s}^n$ \sref{IV.2.2.7}, which proves the assertion concerning $E'$.
\end{remark}

\subsection{Constructibility of topological properties}
\label{subsection:IV.9.5}

\begin{proposition}[9.5.1]
\label{IV.9.5.1}
Let $f:X\to S$ be a morphism of finite presentation, $Z$ a locally constructible part of $X$.
Then the set $E$ of $s\in S$ such that $Z_s\neq\varnothing$ is locally constructible.
\end{proposition}

\begin{proof}
Indeed, we have $E=f(Z)$, and it suffices to apply the theorem of Chevalley \sref{IV.1.8.4}.
\end{proof}

\begin{corollary}[9.5.2]
\label{IV.9.5.2}
If $Z'$, $Z''$ are two locally constructible parts of $X$, the set of $s\in S$ such that $Z'_s\subset Z''_s$ (resp.\ $Z'_s=Z''_s$) is locally constructible.
\end{corollary}

\begin{proof}
Indeed, the relation $Z'_s\subset Z''_s$ is equivalent to $(Z'\cap(\complement Z''))_s=\varnothing$, and $Z'\cap\complement Z''$ is locally constructible.
\end{proof}

\begin{proposition}[9.5.3]
\label{IV.9.5.3}
Let $f:X\to S$ be a morphism of finite presentation, $Z$, $Z'$ two locally constructible parts of $X$ such that $Z\subset Z'$.
Then the set $E$ of $s\in S$ such that $Z_s$ is dense in $Z'_s$ is locally constructible in $S$.
\end{proposition}

\begin{proof}
We must verify the two conditions of \sref{IV.9.2.2},~(iii).
As far as the first is concerned, consider a prescheme $X$ algebraic over a field $k$, two constructible parts $Z$, $Z'$ of $X$ such that $Z\subset Z'$, and an extension $k'$ of $k$.
Then the canonical projection $p:X_{(k')}\to X$ is a faithfully flat and quasi-compact morphism, and we have $p^{-1}(\overline{Z})=\overline{p^{-1}(Z)}$ and $p^{-1}(\overline{Z'})=\overline{p^{-1}(Z')}$; by virtue of \sref{IV.2.3.10}, as $p$ is surjective, the relation $\overline{Z}=\overline{Z'}$ is equivalent to $\overline{p^{-1}(Z)}=\overline{p^{-1}(Z')}$.

Let us verify the second condition, and suppose therefore $S$ affine, Noetherian, integral, of generic point $\eta$.
We distinguish two cases:
\begin{enumerate}[label=\upshape{$\arabic{enumi}^\circ$}]
  \item $\eta\in S-E$, in other words, $Z_\eta$ is not dense in $Z'_\eta$; then there exists in $X$ an open $V$ such that $V\cap Z_\eta=\varnothing$ and $V\cap Z'_\eta\neq\varnothing$.
  As $X$ is Noetherian, $V$ is locally constructible, and hence so is $V\cap Z$; by virtue of \sref{IV.9.5.1}, there is a neighbourhood $U$ of $\eta$ in $S$ such that for every $s\in U$, $(V\cap Z)_s=\varnothing$ and $(V\cap Z')_s\neq\varnothing$; this implies that $Z_s$ is not dense in $Z'_s$ for $s\in U$, in other words $U\subset S-E$.

\oldpage[IV-3]{68}

  \item\label{IV.9.5.3.case2} $\eta\in E$, so $Z_\eta$ is dense in $Z'_\eta$.
  Let us first show that we can suppose $Z'$ \emph{closed}.
  Indeed, $Z_\eta$ is dense in $\overline{Z'_\eta}$ (closure taken in $X_\eta$); setting $V_\eta=X_\eta-\overline{Z_\eta}$, which is open in $X_\eta$ and does not meet $Z'_\eta$, we can suppose $V_\eta$ of the form $V\cap X_\eta$, where $V$ is open (and hence constructible) in $X$, and the hypothesis $V_\eta\cap Z'_\eta=\varnothing$ implies then $V_s\cap Z'_s=\varnothing$ for every $s$ in a neighbourhood of $\eta$ in $S$ by virtue of \sref{IV.9.5.1}.
  Replacing $S$ by an open neighbourhood of $\eta$, we can therefore suppose $V\cap Z'=\varnothing$, whence $V\cap\overline{Z'}=\varnothing$ (closure taken in $X$), and consequently $(\overline{Z'})_\eta=\overline{Z'_\eta}$, whence our assertion.
  The decomposition of $Z'$ into its irreducible components being a finite union, and restricting further $S$ to a neighbourhood of $\eta$, we can suppose that all the irreducible components $Z'_i$ of $Z'$ meet $X_\eta$, whence it follows that $X_\eta$ contains the generic points of the $Z'_i$ \sref[0]{0.2.1.8}.
  To say that $Z_s$ is dense in $Z'_s$ is equivalent to saying that each of the $(Z\cap Z'_i)_s$ is dense in $(Z'_i)_s$, and one is thus reduced to the case where $Z'$ is \emph{irreducible}.
  Replacing then $X$ by the sub-prescheme having the reduced underlying space of $Z'$, we see that we can suppose $Z'=X$ and that $X$ is \emph{integral} and dominates $S$.
  Finally, covering $X$ by a finite number of affine opens $W_j$ and replacing $Z$ by $Z\cap W_j$, we can suppose that $X=\operatorname{Spec}(A)$, where $A$ is a Noetherian integral ring.
  As $X_\eta$ is Noetherian integral and $Z_\eta$ is constructible and dense in $X_\eta$, $Z_\eta$ contains a non-empty open of $X_\eta$~($\mathbf{0}_{\mathrm{III}}$,~9.2.2), which may be taken to be of the form $(D(t))_\eta$, where $t$ is a nonzero element of $A$.
  Replacing $S$ by a neighbourhood of $\eta$ in $S$ if necessary, by virtue of the relation $(D(t))_\eta\subset Z_\eta$, we can suppose in addition that $D(t)\subset Z$ \sref{IV.9.5.2}.
  Finally, since the homothety of ratio $t_\eta$ in $\sh{O}_{X_\eta}$ is injective, it follows from \sref{IV.9.4.5} that for $s$ close to $\eta$, $t_s$ is $\sh{O}_{X_s}$-regular, so $(X_s)_{t_s}$ is dense in $X_s$, and \emph{a fortiori} it is the same for $Z_s$ which contains $(X_s)_{t_s}$.
\end{enumerate}
\end{proof}

\begin{corollary}[9.5.4]
\label{IV.9.5.4}
Let $f:X\to S$ be a morphism of finite presentation, $Z$ a locally constructible part of $X$.
The set $E$ of $s\in S$ such that $Z_s$ is closed (resp.\ open, resp.\ locally closed) in $X_s$ is locally constructible in $S$.
\end{corollary}

\begin{proof}
To say that $Z_s$ is open in $X_s$ means that $(X-Z)_s=X_s-Z_s$ is closed in $X_s$, and as $X-Z$ is locally constructible, we can limit ourselves to considering the set of $s\in S$ such that $Z_s$ is closed and the set of $s\in S$ such that $Z_s$ is locally closed.

Let us verify in each case the two conditions of \sref{IV.9.2.2},~(ii).
The first results from the fact that $p:X_{(k')}\to X$ is faithfully flat and quasi-compact, and from \sref{IV.2.3.12} and \sref{IV.2.3.14}.
Let us therefore verify the second condition, $S$ being supposed affine, Noetherian, integral, of generic point $\eta$.
Set $Z'=\overline{Z}$; the same reasoning as in \sref{IV.9.5.3} shows that $Z'_\eta$ is equal to the closure of $Z_\eta$ in $X_\eta$; by virtue of \sref{IV.9.5.3}, there is thus a neighbourhood $U$ of $\eta$ such that for $s\in U$, $Z_s$ is dense in $Z'_s$, the latter being closed.
To say that $Z_s$ is closed in $X_s$ means that $Z''_s=\varnothing$, where $Z''=Z'-Z$; it follows from \sref{IV.9.5.1} that the set $E$ of $s\in S$ where $Z''_s=\varnothing$ is such that $E$ or $S-E$ contains a neighbourhood of $\eta$.
To say that $Z_s$ is locally closed in $X_s$ means that $Z''_s$ is closed in $X_s$; it suffices therefore to apply the preceding result replacing $Z$ by $Z''$, which is locally constructible in $X$.
\end{proof}

\oldpage[IV-3]{69}

\begin{proposition}[9.5.5]
\label{IV.9.5.5}
Let $f:X\to S$ be a morphism of finite presentation, $Z$ a locally constructible part of $X$ such that for every $s\in S$, $Z_s$ is locally closed in $X_s$.
For every $s\in S$, let $D(s)\subset\mathbf{Z}\cup\{-\infty\}$ be the set of dimensions of the irreducible components of $Z_s$.
Then the function $s\mapsto D(s)$ is locally constructible in $S$.
\end{proposition}

\begin{proof}
Let $\Phi$ be a finite part of $\mathbf{Z}\cup\{-\infty\}$; it is a question of showing that the set of $s\in S$ such that $D(s)=\Phi$ is locally constructible; taking into account \sref{IV.9.2.3}, we have once more to verify the two conditions of \sref{IV.9.2.2},~(iii).

As far as the first is concerned, it is a question of seeing that if $X$ is a prescheme algebraic over a field $k$, $Z$ a locally closed part of $X$, $k'$ an extension of $k$, $p:X_{(k')}\to X$ the canonical projection, then the set of dimensions of the irreducible components of $Z$ is the same as that of the dimensions of the irreducible components of $p^{-1}(Z)$; taking into account the existence of a sub-prescheme of $X$ having $Z$ as underlying space \sref[I]{I.5.2.1}, this follows from \sref{IV.4.2.8}.

For the second condition of \sref{IV.9.2.2},~(iii), we are in the case where $S$ is Noetherian, integral, of generic point $\eta$, and $f:X\to S$ is a morphism of finite type, so that $X$ is Noetherian.
The sub-space $Z_\eta$ of the Noetherian space $X_\eta$ is by hypothesis locally closed, therefore has a finite number of irreducible components $Z_{i\eta}$, which are locally closed in $X_\eta$.
It follows for each index $i$ that there exists a locally closed part $Z_i$ of $X$ such that $(Z_i)_\eta=Z_{i\eta}$, so if $Z'=\bigcup_i Z_i$, we have $Z'_\eta=Z_\eta$.
But as $Z$ and $Z'$ are locally constructible, we can, replacing $S$ by a neighbourhood of $\eta$, suppose that $Z'=Z$ \sref{IV.9.5.1}.
Furthermore, for $i\neq j$, $Z_{i\eta}\cap Z_{j\eta}$ is rare and \emph{closed} in $Z_{i\eta}$; therefore \sref{IV.9.5.3} and \sref{IV.9.5.4}, we can also suppose, restricting $S$, that for $j\neq i$, $(Z_i)_s\cap(Z_j)_s$ is rare and closed in $(Z_i)_s$.
As $Z_i$ is locally closed in $X$, there is a subprescheme of $X$ having $Z_i$ as underlying space (and still denoted $Z_i$), which is of finite type over $S$ \sref[I]{I.6.3.5}.
Set $U_i=Z_i-\bigcup_{j\neq i}(Z_i\cap Z_j)$, which is open in $Z_i$ and such that $(U_i)_s$ is open and everywhere dense in $(Z_i)_s$; furthermore, by construction, the $U_i$ have no two common points.
As $(Z_i)_s$ is a $\boldsymbol{k}(s)$-prescheme algebraic, the set of dimensions of the irreducible components of $Z_s=\bigcup_i(Z_i)_s$ is the same as the set of dimensions of the irreducible components of the reunion of the $(U_i)_s$ \sref{IV.4.1.1.3}, each of these components being already an irreducible component of one of the $(U_i)_s$.
We can therefore limit ourselves to the case where $Z=U_i$; furthermore, as $Z_\eta$ is then irreducible, there is only a single irreducible component of $Z$ meeting $X_\eta$~\sref[0]{0.2.1.8}, and one can evidently, restricting $Z$, suppose $Z$ irreducible.
We are thus finally reduced to proving the following particular case of \sref{IV.9.5.5}:
\end{proof}

\begin{corollary}[9.5.6]
\label{IV.9.5.6}
Let $S$ be a Noetherian prescheme, irreducible, of generic point $\eta$, $X$ an irreducible prescheme, $f:X\to S$ a dominant morphism of finite type.
Then there exists a neighbourhood $U$ of $\eta$ in $S$ such that, for every $s\in U$, all the irreducible components of $X_s$ are of dimension $n=\dim(X_\eta)$.
\end{corollary}

\begin{proof}
We can evidently reduce to the case where $S=\operatorname{Spec}(A)$ is affine, $A$ being therefore Noetherian; replacing $f$ by $f_{\mathrm{red}}$, which is of finite type \sref[I]{I.5.4},~(vii),

\oldpage[IV-3]{70}

we can suppose $A$ integral and $X$ reduced, therefore integral.
We know \sref{IV.4.1.2} that there exists an open $W$ dense in $X_\eta$ and a $\boldsymbol{k}(\eta)$-morphism \emph{finite} and \emph{surjective} $h:W\to\mathbf{V}((\boldsymbol{k}(\eta))^n)=\operatorname{Spec}(B')$, where $B'=\boldsymbol{k}(\eta)[\mathrm{T}_1,\ldots,\mathrm{T}_n]$ ($\mathrm{T}_i$ indeterminates).
If $V$ is an open of $X$ such that $V\cap X_\eta=W$, we know \sref{IV.9.5.3} that for $s$ close to $\eta$ in $S$, $V_s$ is an open dense in $X_s$, and one can therefore \sref{IV.4.1.1.3} reduce to the case where $V=X$, $W=X_\eta$.
Set $B=A[\mathrm{T}_1,\ldots,\mathrm{T}_n]$ and $Y=\operatorname{Spec}(B)=\mathbf{V}_A^n$, so that $\operatorname{Spec}(B')=Y_\eta$; it follows from \sref{IV.8.8.2},~(i) and \sref{IV.8.10.5},~(vi) and (x), applied following the method of \sref{IV.8.1.2},~\emph{a}), that replacing $S$ by a neighbourhood of $\eta$, we can suppose that $h=g_\eta$, where $g:X\to Y$ is a finite surjective morphism; in other words, we have $X=\operatorname{Spec}(C)$, where $C$ is a finite B-algebra and the homomorphism $B\to C$ corresponding to $g$ is \emph{injective}; as $C$ is an integral ring, $C$ is therefore a B-module of finite type \emph{without torsion}, and $C_\eta=C\otimes_A\boldsymbol{k}(\eta)$ is a module of finite type without torsion over $B_\eta=B'=B\otimes_A\boldsymbol{k}(\eta)=\boldsymbol{k}(\eta)[\mathrm{T}_1,\ldots,\mathrm{T}_n]$ (being a ring of fractions whose denominators lie in $A-\{0\}\subset B-\{0\}$).
It follows from \sref{IV.9.4.8} that there is a neighbourhood $U$ of $\eta$ in $S$ such that for $s\in S$, $C_s=C\otimes_A\boldsymbol{k}(s)$ is a module of finite type \emph{without torsion} over $B_s=B\otimes_A\boldsymbol{k}(s)=\boldsymbol{k}(s)[\mathrm{T}_1,\ldots,\mathrm{T}_n]$, and in particular the homomorphism $g_s:B_s\to C_s$ is injective.
As no element of $B_s$ is a divisor of zero in $C_s$, for every minimal prime ideal $\mathfrak{p}_i$ of $C_s$ (whose elements are divisors of zero in $C_s$), we necessarily have $\mathfrak{p}_i\cap B_s=0$, whence the canonical homomorphism $B_s\to C_s/\mathfrak{p}_i$ is injective.
We deduce that for each irreducible component $Z_i=\operatorname{Spec}(C_s/\mathfrak{p}_i)$ of $X_s$, the restriction to $Z_i$ of $g$ is a finite and dominant morphism $Z_i\to Y_s$, therefore \emph{surjective} \sref[II]{II.6.1.10}.
We conclude then from \sref{IV.4.1.2} that $\dim Z_i=n$, which completes the proof.
\end{proof}

\begin{remark}[9.5.7]
\label{IV.9.5.7}
One should note that under the hypotheses of \sref{IV.9.5.6}, it may happen that $X_s$ is irreducible for \emph{no} $s\neq\eta$ in a neighbourhood of $\eta$, in other words, the property ``$X$ is an irreducible $k$-prescheme algebraic'' is \emph{not constructible}.
Take for example $S=\operatorname{Spec}(k[\mathrm{T}])$, where $k$ is an algebraically closed field, $\mathrm{T}$ an indeterminate; we then have $\boldsymbol{k}(\eta)=K=k(\mathrm{T})$.
Let $L$ be a finite separable extension of $K$ of degree $>1$, and let $X$ be the \emph{integral closure} of $S$ in $L$ \sref[II]{II.6.3.4}; we then have $X=\operatorname{Spec}(B)$, where $B$ is the integral closure of $k[\mathrm{T}]$ in $L$.
We know that $B$ is a Dedekind ring, and that all but finitely many maximal ideals of $k[\mathrm{T}]$ are unramified in $B$. Moreover, the residue field of every maximal ideal of $B$ is necessarily $k$, since it is a finite extension of the algebraically closed field $k$. Thus, for almost all maximal ideals $\mathfrak{j}_s$ of $k[\mathrm{T}]$, the ring $B_s=B/\mathfrak{j}_s B$ is a direct product of $[L:K]$ copies of $k$; in particular $X_s$ is not irreducible, although $X_\eta=\operatorname{Spec}(L)$ is.
The same example shows that the property ``$X$ is an integral $k$-prescheme algebraic'' is not constructible.
Finally, the property ``$X$ is a reduced $k$-prescheme algebraic'' is also not constructible.
It suffices for this to take again $S=\operatorname{Spec}(k[\mathrm{T}])$, where this time $k$ is an algebraically closed field of characteristic $p>0$, and for $X$ the integral closure of $S$ in $L=K^{1/p}$ (where $K=k(\mathrm{T})$), so that $X=\operatorname{Spec}(B)$ with $B=k[\mathrm{T}^{1/p}]$ ($k$ being perfect); every maximal ideal of $k[\mathrm{T}]$ is of the form $(\mathrm{T}-\alpha)$ with $\alpha\in k$; the unique ideal of $B$ over the ideal $(\mathrm{T}-\alpha)$ is the principal ideal $(\mathrm{T}^{1/p}-\alpha^{1/p})$ and it is immediate
\oldpage[IV-3]{71}
that the quotient ring $B/(\mathrm{T}-\alpha)B$ has nilpotent elements; in other words, $X_s$ is not reduced for any $s\neq\eta$, whilst $X_\eta=\operatorname{Spec}(L)$ is integral.

We shall see a little further on \sref{IV.9.7} that one obtains on the other hand constructible properties when one considers the ``geometric'' notions corresponding to the notions of irreducible, reduced, or integral prescheme \sref{IV.4.5} and \sref{IV.4.6}.
\end{remark}

\subsection{Constructibility of certain properties of morphisms}
\label{subsection:IV.9.6}

\begin{proposition}[9.6.1]
\label{IV.9.6.1}
Let $X$, $Y$ be two S-preschemes of finite presentation, $f:X\to Y$ an S-morphism.
Let $E$ be the set of $s\in S$ for which $f_s$ has one of the following properties: being
\begin{enumerate}[label=(\roman*),widest=xii]
  \item surjective;
  \item dominant;
  \item separated;
  \item proper;
  \item radicial;
  \item finite;
  \item quasi-finite;
  \item an immersion;
  \item a closed immersion;
  \item an open immersion;
  \item an isomorphism;
  \item a monomorphism.
\end{enumerate}
Then $E$ is locally constructible in $S$.
\end{proposition}

\begin{proof}
The assertions concerning (i), (v), and (vii) are included only for reference, having already been established in \sref{IV.9.3.2}.

(ii): Let us note first that $f$ is of finite presentation \sref[I]{I.6.2},~(v), therefore, by virtue of the theorem of Chevalley \sref{IV.1.8.4}, $f(X)$ is locally constructible in $Y$.
Furthermore, we have $f_s(X_s)=(f(X))_s$ \sref[I]{I.3.6.1}; the set of $s$ such that $f_s$ is dominant is the set of $s$ such that $(f(X))_s$ is dense in $Y_s$; the conclusion follows therefore from \sref{IV.9.5.3}.

(iii): As $f:X\to Y$ is of finite presentation, the diagonal immersion $\Delta_f:X\to X\times_Y X$ is of finite presentation \sref[I]{I.6.2},~(iv) and (v); it follows from \sref{IV.1.8.4.1} that $\Delta_f(X)=\Delta_X$ is a locally constructible part of $X\times_Y X$.
To say that $f_s$ is separated means (taking into account \sref[I]{I.5.3.4}) that $(\Delta_X)_s$ is closed in $X_s\times_{Y_s} X_s=(X\times_Y X)_s$; the conclusion follows from \sref{IV.9.5.4}.

\oldpage[IV-3]{72}

(iv): Let us show that the property ``$X$ and $Y$ are preschemes algebraic over a field $k$, and $f:X\to Y$ a proper $k$-morphism'' is constructible.
The condition \sref{IV.9.2.1},~(i) is verified by virtue of \sref{IV.2.7.1},~(vii).
We can therefore reduce to the case where $S$ is affine, Noetherian, integral, of generic point $\eta$, and one must prove that $E$ or $S-E$ is a neighbourhood of $\eta$.
Suppose first that $\eta\in E$; it follows from \sref{IV.8.10.5},~(xii), applied following the method of \sref{IV.8.1.2},~\emph{a}), that replacing $S$ by a neighbourhood of $\eta$, we can
suppose that $f$ is itself a proper morphism; one then knows that it is the same for $f_s$ for every $s\in S$ \sref[II]{II.5.4.2},~(iii).
Suppose on the contrary that $\eta\in S-E$, and we distinguish two cases:
\begin{enumerate}[label=\upshape{$\arabic{enumi}^\circ$}]
  \item Suppose that $f_\eta$ is not separated; then it follows from (iii) that $f_s$ is not separated (and \emph{a fortiori} not proper) in a neighbourhood of $\eta$.

  \item Suppose $f_\eta$ separated; $Y$ is a finite union of affine opens $V_i$, and for $f_s$ to be proper, it suffices that each of its restrictions $f_s^{-1}((V_i)_s)\to(V_i)_s$ be so \sref[II]{II.5.4.1}; we can therefore reduce to the case where $Y$ is \emph{affine}, hence a scheme.
  To say that $f_\eta$ is not proper means \sref[II]{II.5.6.3} that there exists a morphism \emph{of finite type} $h:Z\to Y_\eta$ such that the morphism $(f_\eta)_{(Z)}:X_\eta\times_{Y_\eta} Z\to Z$ is not closed.
  As $Y$ is a scheme, we deduce from \sref{IV.8.8.2},~(i) and (ii) (restricting to $S$ if needed a neighbourhood of $\eta$) that there exists a morphism of finite type $g:Y'\to Y$ such that $Z$ is isomorphic to $Y'_\eta$ and that $g_\eta=h$;
  if we put $X'=X\times_Y Y'$ and $f'=f_{(Y')}:X'\to Y'$, then $(f_\eta)_{(Z)}=f'_\eta$, and by hypothesis there is therefore a closed part $M'$ of $X'_\eta$ such that $f'_\eta(M')$ is not closed in $Y'_\eta$.
  Now $M'$ is the trace on $X'_\eta$ of a closed part $N'$ of $X'$; as $X'$ is Noetherian and $f'$ of finite type, $f'(N')$ is constructible in $Y'$ \sref{IV.1.8.4} and by hypothesis $(f'(N'))_\eta=f'_\eta(N'_\eta)=f'_\eta(M')$ is not closed in $Y'_\eta$.
  We conclude from \sref{IV.9.5.4} that there exists a neighbourhood $U$ of $\eta$ in $S$ such that for $s\in U$, $(f'(N'))_s=f'_s(N'_s)$ is not closed in $Y'_s$; in other words, the morphism $f'_s$ is not closed, and \emph{a fortiori} the morphism $f_s$ is not proper.
\end{enumerate}

(vi): The property is a conjunction of properties (iv) and (vii) \sref{IV.8.11.1}, therefore the proposition follows from what has already been proved.

(ix): The verification of condition \sref{IV.9.2.1},~(i) follows from \sref{IV.2.7.1},~(xi).
We can therefore reduce to the case where $S$ is affine, Noetherian, integral, of generic point $\eta$, and it suffices to prove that $E$ or $S-E$ is a neighbourhood of $\eta$.
If $\eta\in E$, we can (replacing $S$ by a neighbourhood of $\eta$) suppose that $f$ is a closed immersion by virtue of \sref{IV.8.10.5},~(iv), and then it is clear that $f_s$ is a closed immersion for every $s\in S$ \sref[I]{I.4.3.2}.
Suppose therefore that $\eta\in S-E$, and we distinguish two cases:
\begin{enumerate}[label=\upshape{$\arabic{enumi}^\circ$}]
  \item $f_\eta$ is not a finite morphism.
  Then it follows from (vi) that in a neighbourhood of $\eta$, $f_s$ is not finite, nor \emph{a fortiori} a closed immersion.

  \item $f_\eta$ is finite; then, by virtue of \sref{IV.8.10.5},~(x), we can suppose (restricting $S$ if needed) that $f$ is itself a \emph{finite} morphism.
  In this case $\mathscr{A}=\mathscr{A}(X)=f_*(\sh{O}_X)$ is an $\sh{O}_Y$-Module coherent \sref[II]{II.6.1.3} and $f=\mathscr{A}(u)$, where $u:\sh{O}_Y\to\mathscr{A}$ is a homomorphism of $\sh{O}_Y$-Algebras \sref[II]{II.1.1.2}; as $f_\eta$ is not a closed immersion by hypothesis, $u_\eta$ is not surjective \sref[II]{II.1.4.10}, therefore \sref{IV.9.4.5} there exists a neighbourhood of $\eta$ in which $u_s$ is not surjective, and consequently \sref[I]{I.4.2.3} $f_s$ is not a closed immersion.
\end{enumerate}

(viii): The verification of condition \sref{IV.9.2.1},~(i) is done as in (ix), using this time \sref{IV.2.7.1},~(x) and the fact that every immersion of one Noetherian prescheme into another is quasi-compact.
We are therefore reduced to the case where $S$ is affine, Noetherian, integral, of generic point $\eta$, and it suffices to prove that $E$ or $S-E$ is a neighbourhood of $\eta$.
If $\eta\in E$, we conclude as in (ix) with the help of \sref{IV.8.10.5},~(ii).
If $\eta\in S-E$, we distinguish again two cases:
\begin{enumerate}[label=\upshape{$\arabic{enumi}^\circ$}]
  \item $f_\eta(X_\eta)$ is not a locally closed part of $Y_\eta$.
  As $f(X)$ is constructible in $Y$ \sref{IV.1.8.4} and $f_s(X_s)=(f(X))_s$, we deduce from \sref{IV.9.5.4} that for $s$ close to $\eta$, $f_s(X_s)$ is not locally closed in $Y_s$, and \emph{a fortiori} $f_s$ is not an immersion.

  \item\oldpage[IV-3]{73} $f_\eta(X_\eta)$ is locally closed in $Y_\eta$.
  As $f(X)$ is constructible in $Y$ \sref{IV.1.8.4} and since the same holds for $\overline{f(X)}$ since $Y$ is Noetherian, it follows from \sref{IV.8.3.11} that restricting to $S$ if necessary, we can suppose that $f(X)$ is locally closed in $Y$.
  There is then an open $V$ of $Y$ containing $f(X)$ and in which $f(X)$ is closed.
  Since $Y$ is Noetherian, $V$ is of finite type over $S$; replacing $Y$ by $V$, we may therefore reduce to the case where $f(X)$ is closed in $Y$.
  But then $f_s(X_s)=(f(X))_s$ is closed in $Y_s$ for every $s\in S$, and saying that $f_s$ is an immersion is equivalent to saying that it is a closed immersion.
  We are therefore reduced to case~(ix).
\end{enumerate}

(x): Using this time \sref{IV.2.7.1},~(ix) and \sref{IV.8.10.5},~(iii), we are reduced to the case where $S$ is affine, Noetherian, integral, of generic point $\eta$, and where $\eta\in S-E$.
We distinguish three cases:
\begin{enumerate}[label=\upshape{$\arabic{enumi}^\circ$}]
  \item $f_\eta(X_\eta)$ is not open in $Y_\eta$.
  As $f(X)$ is constructible in $Y$ \sref{IV.1.8.4}, we deduce from \sref{IV.9.5.4} that $f_s(X_s)$ is not open in $Y_s$ for $s$ close to $\eta$, and \emph{a fortiori} $f_s$ is not an open immersion.

  \item $f_\eta(X_\eta)$ is open in $Y_\eta$ but $f_\eta$ is not an immersion.
  It follows then from (viii) that for $s$ close to $\eta$, $f_s$ is not an immersion, nor \emph{a fortiori} an open immersion.

  \item $f_\eta(X_\eta)$ is open in $Y_\eta$ and $f_\eta$ is an immersion.
  As $f(X)$ is constructible in $Y$ \sref{IV.1.8.4}, it follows from \sref{IV.8.3.11} that restricting to $S$ if necessary, we can already suppose that $f(X)$ is open in $Y$.
  Since $Y$ is Noetherian, the sub-prescheme induced on $f(X)$ is of finite type over $S$, so we can reduce to the case where $f$ is \emph{surjective} by replacing $Y$ by $f(X)$.
  By hypothesis $f_\eta$ is a closed immersion, so, as in (ix), we can suppose that $f$ is a \emph{closed immersion}, in other words, that $X$ is a closed sub-prescheme of $Y$ defined by a coherent ideal $\mathscr{J}$ of $\sh{O}_Y$.
  By hypothesis $f_\eta$ is not an isomorphism, hence $\mathscr{J}_\eta\neq 0$; we conclude \sref{IV.9.4.5} that in a neighbourhood of $\eta$, $\mathscr{J}_s\neq 0$, and consequently $f_s$ is not an open immersion.
\end{enumerate}

(xi): The property is a conjunction of properties (i) and (x) and therefore follows from what has been proved.

(xii): By virtue of \sref[I]{I.5.3.8}, to say that $f_s$ is a monomorphism means that $\Delta_{f_s}=(\Delta_f)_s:X_s\to X_s\times_{Y_s} X_s=(X\times_Y X)_s$ is an isomorphism, and as $X\times_Y X$ is an S-prescheme of finite presentation \sref[I]{I.6.2},~(iv), the conclusion follows from (xi).
\end{proof}

\begin{proposition}[9.6.2]
\label{IV.9.6.2}
Let $X$, $Y$ be two S-preschemes of finite presentation, $f:X\to Y$ an S-morphism.
\begin{enumerate}[label=\upshape{I)}]
  \item Let $E$ be the set of $s\in S$ for which $f_s$ has one of the following properties: being
  \begin{enumerate}[label=(\roman*)]
    \item affine;
    \item quasi-affine;
    \item projective;
    \item quasi-projective.
  \end{enumerate}
  Then $E$ is ind-constructible in $S$.
\end{enumerate}

\oldpage[IV-3]{74}

\begin{enumerate}[label=\upshape{II)}]
  \item Let $\sh{L}$ be an $\sh{O}_X$-Module invertible.
  Then the set $E'$ of $s\in S$ such that $\sh{L}_s$ is an $\sh{O}_{X_s}$-Module ample (resp.\ very ample) relatively to $f_s$, is ind-constructible in $S$.
\end{enumerate}
\end{proposition}

\begin{proof}
I) Let us verify conditions (i) and (ii) of \sref{IV.9.2.1}.
As far as the condition \sref{IV.9.2.1},~(i) is concerned, it follows, for properties (i) and (ii), from \sref{IV.2.7.1},~(xiii) and (xiv); for properties (iii) and (iv), it follows from \sref{IV.9.1.5}.
Let us then verify condition \sref{IV.9.2.1},~(ii), supposing therefore $S$ Noetherian, integral, of generic point $\eta$, and that $f_\eta:X_\eta\to Y_\eta$ has one of the properties (i) to (iv) of the statement.
Applying \sref{IV.8.10.5},~(viii), (ix), (xiii), and (xiv) following the method of \sref{IV.8.1.2},~\emph{a}), we see first that there exists an open neighbourhood $U$ of $\eta$ such that, if $V$ and $W$ are the respective inverse images of $U$ in $X$ and $Y$ under the structural morphisms, the restriction $V\to W$ of $f$ has whichever of the properties (i) to (iv) is under consideration.
The conclusion then follows from the fact that all the properties considered are stable under base change.

II) We proceed in the same way.
For condition \sref{IV.9.2.1},~(i) it follows this time from \sref{IV.2.7.2}.
For condition \sref{IV.9.2.1},~(ii), with the same notations as in I), it follows from \sref{IV.8.10.5.2} that for a neighbourhood $U$ of $\eta$ in $S$, the restriction $\sh{L}|V$ is ample (resp.\ very ample) relatively to the restriction $V\to W$ of $f$.
The conclusion follows again from the stability of the properties considered under base change \sref[II]{II.4.6.13} and \sref[II]{II.4.4.10}.
\end{proof}

We can improve \sref{IV.9.6.2},~II) under certain conditions:

\begin{proposition}[9.6.3]
\label{IV.9.6.3}
Let $X$, $Y$ be two S-preschemes of finite presentation, $f:X\to Y$ a proper S-morphism, $\sh{L}$ an $\sh{O}_X$-Module invertible.
Then the set $E$ (resp.\ $E'$) of $s\in S$ such that $\sh{L}_s$ is an $\sh{O}_{X_s}$-Module ample (resp.\ very ample) relatively to $f_s:X_s\to Y_s$ is locally constructible in $S$.
\end{proposition}

\begin{proof}
Let $k$ be a field, $Z$, $T$ two preschemes algebraic over $k$, $\mathscr{M}$ an $\sh{O}_Z$-Module invertible, $g:Z\to T$ a $k$-morphism of finite type; then, if $\mathbf{P}(Z,T,\mathscr{M},g,k)$ denotes the relation ``$\mathscr{M}$ is ample (resp.\ very ample) relatively to $g$'', we have already noted in \sref{IV.9.6.2} that $\mathbf{P}(Z,T,\mathscr{M},g,k)$ satisfies condition \sref{IV.9.2.1},~(i) by virtue of \sref{IV.2.7.2}.
We know on the other hand that $E$ and $E'$ are ind-constructible.
It remains therefore to see that if $S$ is Noetherian, integral, of generic point $\eta$ and if $\eta\in S-E$ (resp.\ $\eta\in S-E'$), then $S-E$ (resp.\ $S-E'$) contains a neighbourhood of $\eta$.
We shall separately consider the cases $E$ and $E'$.

\oldpage[IV-3]{75}

\emph{Case of $E'$.} ---
Since $f$ is separated and $Y$ is quasi-compact, there is an integer $h$ such that
$R^qf_*(\sh{L})=0$ for $q>h$ \sref[III]{III.1.4.12}.
On the other hand, since $f$ is proper and $Y$ is Noetherian, all the
$R^qf_*(\sh{L})$ are coherent $\sh{O}_Y$-Modules \sref[III]{III.3.2.1}.
They vanish except for finitely many values of $q$, so the generic flatness theorem
\sref{IV.6.9.1} shows that, after restricting $S$ to a neighbourhood of $\eta$, we may
suppose that $\sh{L}$ and all the $R^qf_*(\sh{L})$ are flat over $S$.
It then follows from \sref[III]{III.6.9.9} that the canonical homomorphism
\begin{equation}
\label{IV.9.6.3.1}
 f_*(\sh{L})\otimes_{\sh{O}_S}\boldsymbol{k}(s)
 \longrightarrow (f_s)_*(\sh{L}_s)
 \tag{9.6.3.1}
\end{equation}
is an isomorphism.

It follows from \sref[II]{II.4.4.4} that to say that $\sh{L}_s$ is not very ample relatively to $f_s$ means that: \emph{either} the canonical homomorphism $(f_s)^*((f_s)_*(\sh{L}_s))\to\sh{L}_s$ is not surjective; \emph{or} the preceding homomorphism is surjective and the canonical morphism $r:X_s\to\mathbf{P}((f_s)_*(\sh{L}_s))$ is not an immersion.
Taking into account the isomorphism \sref{IV.9.6.3.1}, these conditions are written respectively in the form: 1$^\circ$ the canonical homomorphism $(f^*(f_*(\sh{L})))_s\to\sh{L}_s$ is not surjective; 2$^\circ$ the preceding homomorphism is surjective and the canonical morphism $X_s\to\mathbf{P}((f_*(\sh{L}))_s)$ is not an immersion.

Suppose first that the canonical homomorphism $(f^*(f_*(\sh{L})))_\eta\to\sh{L}_\eta$ is not surjective.
As $f_*(\sh{L})$ is coherent, so is $f^*(f_*(\sh{L}))$, and then it follows from \sref{IV.9.4.5} that for every $s$ in a neighbourhood of $\eta$, the homomorphism $(f^*(f_*(\sh{L})))_s\to\sh{L}_s$ is not surjective, which proves in this case that $S-E'$ is a neighbourhood of $\eta$.

Suppose next that the canonical homomorphism $(f^*(f_*(\sh{L})))_\eta\to\sh{L}_\eta$ is surjective but that the morphism $X_\eta\to\mathbf{P}((f_*(\sh{L}))_\eta)$ is not an immersion.
Then the same reasoning as above shows that for $s$ sufficiently close to $\eta$, the homomorphism $(f^*(f_*(\sh{L})))_s\to\sh{L}_s$ is surjective; on the other hand, by virtue of \sref{IV.9.6.1}, (viii), the morphism $X_s\to\mathbf{P}((f_*(\sh{L}))_s)$ is not an immersion for $s$ in a neighbourhood of $\eta$, which completes the proof in the case of $E'$.

\emph{Case of $E$.} --- Consider first the particular case where $Y=S$:

\begin{corollary}[9.6.4]
\label{IV.9.6.4}
Let $f:X\to S$ be a proper morphism of finite presentation, $\sh{L}$ an $\sh{O}_X$-Module invertible.
Then the set $E$ of $s\in S$ such that $\sh{L}_s$ is ample (relatively to $f_s$) is open in $S$, and $\sh{L}|f^{-1}(E)$ is ample relatively to the restriction $f^{-1}(E)\to E$ of $f$.
\end{corollary}

\begin{proof}
As condition \sref{IV.9.2.1},~(i) is verified by the property $\mathbf{P}$ defined above, the result of \sref{IV.9.2.2},~(iv) and the reasoning of \sref{IV.9.2.3} show that it suffices to reduce to the case where $S$ is \emph{Noetherian}; but then the result follows from \sref[III]{III.4.7.1} and the stability of amplitude under base change \sref[II]{II.4.6.13}.
\end{proof}

\begin{corollary}[9.6.5]
\label{IV.9.6.5}
Under the hypotheses of \sref{IV.9.6.4}, for $\sh{L}$ to be ample relatively to $f$, it is necessary and sufficient that, for every $s\in S$, $\sh{L}_s$ be ample relatively to $f_s$.
\end{corollary}

\phantomsection\label{IV.9.6.6}%
\textbf{(9.6.6) End of the proof of \sref{IV.9.6.3}.} ---
Let us return to the general case, $S$ being Noetherian, integral, of generic point $\eta\in S-E$.
As $f$ is proper, it follows from \sref{IV.9.6.4} that the set $V$ of $y\in Y$ such that $\sh{L}_y$ is an $\sh{O}_{X_y}$-Module ample, relatively to the morphism $f_y:X_y\to\operatorname{Spec}(\boldsymbol{k}(y))$, is open, and therefore $F=Y-V$ is closed in $Y$.
Since $f_s$ is proper and, for every $y\in Y$ above $s\in S$, one has
$\sh{L}_y=(\sh{L}_s)_y$, it follows from \sref{IV.9.6.5} that $s$ belongs to
$S-E$ if and only if $F_s\neq\varnothing$.
But $F$ is closed in $Y$ and $F_\eta\neq\varnothing$; hence \sref{IV.9.5.1} shows
that the set of $s\in S$ such that $F_s\neq\varnothing$ is a neighbourhood of
$\eta$ in $S$.
\end{proof}

\begin{remarks}[9.6.7]
\label{IV.9.6.7}
(i) For each of the properties $\mathbf{P}$ considered in \sref{IV.9.6.1}, proposition \sref{IV.9.3.3} is applicable, and these properties (for the morphisms $f_s$) are therefore ``stable'' under passage from the projective limit of an essentially affine projective system \sref{IV.8.13.4} of preschemes $X_\lambda$ to a suitable one of them.

(ii) Let $Z$ be a locally constructible part of $X$ such that, for every $s\in S$, $Z\cap X_s$ is \emph{open} in $X_s$, and denote by $Z_s$ the sub-prescheme of $X_s$ induced on the open $Z\cap X_s$.
\oldpage[IV-3]{76}
Then, in propositions \sref{IV.9.6.1} and \sref{IV.9.6.2},~(I), we can everywhere replace $f_s$ by its restriction $f_s|Z_s:Z_s\to Y_s$ without changing the conclusions.
Indeed, the verification of \sref{IV.9.2.1},~(i) proceeds as in \sref{IV.9.6.1} and \sref{IV.9.6.2}.
On the other hand, in the reduction to the Noetherian case, made in \sref{IV.9.2.3}, if $Z=q^{-1}(Z_0)$, where $q:X\to X_0$ is the canonical projection and $Z_0$ a constructible part of $X_0$ \sref{IV.8.3.11}, the fact that $(Z_0)_{s_0}$ is open in $(X_0)_{s_0}$, for $s_0=p(s)$, follows from \sref{IV.2.4.10} and from the fact that the projection $X_s\to(X_0)_{s_0}$ is surjective.
We are therefore reduced to verifying \sref{IV.9.2.1},~(ii) under the new hypotheses.
Now, as $Z_\eta$ is open in $X_\eta$, there is an open $Z'\subset X$ such that $Z_\eta=Z'\cap X_\eta$; as $X$ is then Noetherian, $Z'$ is constructible, and $Z$ is likewise constructible by hypothesis; we conclude from \sref{IV.9.5.2} and from ($\mathbf{0}_{\mathrm{III}}$,~9.2.2) that there is a neighbourhood $U$ of $\eta$ in $S$ such that $Z_s=Z'_s$ for $s\in U$.
Replacing $S$ by $U$, we can therefore reduce to the case where $Z$ is open in $X$, and then the result has been proved in \sref{IV.9.6.1} and \sref{IV.9.6.2}.
\end{remarks}

\subsection{Constructibility of properties of separability, geometric irreducibility, and geometric connectedness}
\label{subsection:IV.9.7}
\label{IV.9.7}

\begin{lemma}[9.7.1]
\label{IV.9.7.1}
Let $S$ be an irreducible prescheme of generic point $\eta$, $f:X\to S$ a morphism of finite presentation.
If $X_\eta$ has $n$ irreducible components (resp.\ $n'$ connected components), there exists an open neighbourhood $U$ of $\eta$ in $S$ such that for every $s\in U$, $X_s$ has at least $n$ irreducible components (resp.\ $n'$ connected components).
\end{lemma}

\begin{proof}
We can reduce to the case where $S$ is affine, so $X$ is quasi-compact and quasi-separated.
Let $V_i$ ($1\leqslant i\leqslant n$) be the interiors of the irreducible components of $X_\eta$, which are pairwise disjoint and quasi-compact, and whose union is dense in $X_\eta$; by virtue of \sref{IV.8.2.11}, applied following the method of \sref{IV.8.1.2},~\emph{a}), there exists for each $i$ a quasi-compact open $W_i$ in $X$ such that $W_i\cap X_\eta=V_i$; as $X$ is quasi-separated, the intersections $W_i\cap W_j$ are quasi-compact \sref{IV.1.2.7}, and therefore, replacing $S$ by a neighbourhood of $\eta$, we can suppose that $W_i\cap W_j=\varnothing$ for $i\neq j$ by virtue of \sref{IV.8.3.3}.
Furthermore, as the $W_i$ are constructible, there is a neighbourhood $U$ of $\eta$ in $S$ such that for every $s\in U$ the $(W_i)_s$ are non-empty \sref{IV.9.5.1} and \sref{IV.9.2.3} and the union of the $(W_i)_s$ is dense in $X_s$ \sref{IV.9.5.3} and \sref{IV.9.2.3}.
This being the case, as the irreducible components of each $(W_i)_s$ (being finite in number) are also irreducible components of the union of the $(W_i)_s$ \sref[0]{0.2.1.7}, the closures in $X_s$ of these components are the irreducible components of $X_s$ \sref[0]{0.2.1.6} and their number is evidently $\geqslant n$.

Let now $C_j$ ($1\leqslant j\leqslant n'$) be the connected components of $X_\eta$, which are quasi-compact opens of $X_\eta$, pairwise disjoint.
If we replace $V_i$ by $C_j$ in the preceding reasoning, we see (using twice \sref{IV.8.3.3}) that we can suppose $X$ a union of $n'$ quasi-compact opens $M_j$ ($1\leqslant j\leqslant n'$), pairwise disjoint, and such that, in a neighbourhood of $\eta$, the $(M_j)_s$ are non-empty.
As the union of the $(M_j)_s$ is $X_s$, the connected components of the $(M_j)_s$ are the connected components of $X_s$, and therefore their number is $\geqslant n'$.
\end{proof}

\oldpage[IV-3]{77}

\begin{lemma}[9.7.2]
\label{IV.9.7.2}
Let $S$ be an irreducible prescheme of generic point $\eta$, $f:X\to S$ a morphism of finite presentation.
If $X_\eta$ is not reduced, there exists an open neighbourhood $U$ of $\eta$ in $S$ such that, for every $s\in U$, $X_s$ is not reduced.
\end{lemma}

\begin{proof}
Indeed, let $\mathscr{N}$ be the Nilradical of $\sh{O}_X$; it follows from \sref{IV.8.2.13}, applied following the method of \sref{IV.8.1.2},~\emph{a}), that $\mathscr{N}_\eta$ is the nilradical of $\sh{O}_{X_\eta}$, and the hypothesis is that $\mathscr{N}_\eta\neq 0$.
We conclude from \sref{IV.9.4.5} and \sref{IV.9.2.3} that there is a neighbourhood $U$ of $\eta$ such that, for every $s\in U$, $\mathscr{N}_s$ identifies with an ideal of $\sh{O}_{X_s}$ and $\mathscr{N}_s\neq 0$; as $\mathscr{N}_s$ is evidently contained in the nilradical of $\sh{O}_{X_s}$, we see that $X_s$ is not reduced for $s\in U$.
\end{proof}

\textbf{(9.7.3)}
\label{IV.9.7.3}
Given a polynomial $F\in A[\mathrm{T}_1,\ldots,\mathrm{T}_n]$, where $A$ is a ring and the $\mathrm{T}_i$ are indeterminates, for a ring homomorphism $\rho:A\to B$, we denote by $F_{(\rho)}$ or $F_{(B)}$ the polynomial in $B[\mathrm{T}_1,\ldots,\mathrm{T}_n]$ obtained by replacing each coefficient of $F$ by its image under $\rho$.
If $k$ is a field, $F\in k[\mathrm{T}_1,\ldots,\mathrm{T}_n]$ a non-constant polynomial, and $X=\operatorname{Spec}(k[\mathrm{T}_1,\ldots,\mathrm{T}_n]/(F))$, to say that $X$ is integral (or that the ideal $(F)$ is prime) means that $F$ is \emph{irreducible} (i.e. that in every factorisation $F=F_1 F_2$ in polynomials of $k[\mathrm{T}_1,\ldots,\mathrm{T}_n]$, $F_1$ or $F_2$ is of degree $0$); this follows from the fact that the polynomial ring $k[\mathrm{T}_1,\ldots,\mathrm{T}_n]$ is factorial.
We deduce immediately from \sref{IV.4.6.2} and \sref{IV.4.5.2}:

\begin{lemma}[9.7.4]
\label{IV.9.7.4}
Let $k$ be a field, $\Omega$ an algebraically closed extension of $k$, $F$ a non-constant polynomial of $k[\mathrm{T}_1,\ldots,\mathrm{T}_n]$.
The following conditions are equivalent:
\begin{enumerate}[label=\emph{\alph*)}]
  \item $X=\operatorname{Spec}(k[\mathrm{T}_1,\ldots,\mathrm{T}_n]/(F))$ is geometrically integral.
  \item $F_{(K)}$ is irreducible for every extension $K$ of $k$.
  \item $F_{(\Omega)}$ is irreducible.
\end{enumerate}
We shall say in this case that $F$ is \emph{geometrically irreducible}.
\end{lemma}

\begin{lemma}[9.7.5]
\label{IV.9.7.5}
Let $A$ be an integral ring, $K$ the field of fractions of $A$, $F$ a non-constant polynomial of $A[\mathrm{T}_1,\ldots,\mathrm{T}_n]$ of degree $d$ such that $F_{(K)}$ is geometrically irreducible.
Then there exists $f\neq 0$ in $A$ such that for every $x\in D(f)$, $F_{(\boldsymbol{k}(x))}$ is geometrically irreducible.
\end{lemma}

\begin{proof}
Set $F=\sum_\alpha c_\alpha \mathrm{T}^\alpha$ as usual, with $c_\alpha=c_{\alpha_1\alpha_2\ldots\alpha_n}$, $\mathrm{T}^\alpha=\mathrm{T}_1^{\alpha_1}\mathrm{T}_2^{\alpha_2}\cdots\mathrm{T}_n^{\alpha_n}$, $|\alpha|=\alpha_1+\alpha_2+\cdots+\alpha_n\leqslant d$ (at least one of the $c_\alpha$ such that $|\alpha|=d$ being nonzero).
As the nonzero $c_\alpha$ are invertible in $K$, we can suppose, replacing $A$ by $A_g$ ($g\neq 0$ in $A$), that the nonzero $c_\alpha$ are invertible \emph{in} $A$.
It then follows that for every $x\in\operatorname{Spec}(A)$, $F_{(\boldsymbol{k}(x))}$ is of degree $d$.

We shall prove the lemma by means of a preliminary lemma.

\begin{lemma}[9.7.5.1]
\label{IV.9.7.5.1}
Let $A$ be an integral ring, $S=\operatorname{Spec}(A)$, $\eta$ the generic point of $S$, $B$ the polynomial ring $A[\mathrm{T}_1,\ldots,\mathrm{T}_m]$, $(P_i)_{i\in I}$ a finite family of elements of $B$, $\mathfrak{a}$ the ideal of $B$ generated by the $P_i$.
For every $s\in S$, let $\mathfrak{a}_s$ be the ideal of $\boldsymbol{k}(s)[\mathrm{T}_1,\ldots,\mathrm{T}_m]$ generated by the polynomials $(P_i)_{(\boldsymbol{k}(s))}$ ($i\in I$).
Then, if $V(\mathfrak{a}_\eta)=\varnothing$, there exists a neighbourhood $U$ of $\eta$ in $S$ such that $V(\mathfrak{a}_s)=\varnothing$ for every $s\in U$.
\end{lemma}

\begin{proof}
Indeed, let $Y=\operatorname{Spec}(B)$, and let $Z$ be the closed part $V(\mathfrak{a})$ of $Y$; as $\mathfrak{a}$ is an ideal of finite type in $B$, $Z$ is constructible in $Y$ \sref[0]{0.9.1.5} and, with the notations introduced in \sref{IV.9.4.1}, we have $V(\mathfrak{a}_s)=Z_s$ for every $s\in S$; as the structural morphism $Y\to S$ is of finite presentation, the conclusion of the lemma follows from \sref{IV.9.5.1} and \sref{IV.9.2.3}.
\end{proof}

\oldpage[IV-3]{78}
Let $(p,q)$ be a pair of integers $>0$ such that $p+q=d$; we introduce indeterminates $\mathrm{T}'_\beta$, $\mathrm{T}''_\gamma$ for all systems of integers $\beta$, $\gamma$ such that $|\beta|\leqslant p$ and $|\gamma|\leqslant q$; for every system of integers $\alpha$ such that $|\alpha|\leqslant d$, consider the polynomial of $B=A[\mathrm{T}'_\beta, \mathrm{T}''_\gamma \mid |\beta|\leqslant p, |\gamma|\leqslant q]: P_\alpha(\mathrm{T}'_\beta,\mathrm{T}''_\gamma)=\sum_{\beta+\gamma=\alpha}\mathrm{T}'_\beta \mathrm{T}''_\gamma-c_\alpha$.
Let $\Omega$ be an algebraic closure of $K$; to say that there exist two polynomials $F_1$, $F_2$ of $\Omega[\mathrm{T}_1,\ldots,\mathrm{T}_n]$, of respective degrees $p$ and $q$, and such that $F_1 F_2=F_{(\Omega)}$ means that the system of equations $P_\alpha(t'_\beta,t''_\gamma)=0$ ($|\alpha|\leqslant d$) admits a solution $(t'_\beta,t''_\gamma)$ ($|\beta|\leqslant p$, $|\gamma|\leqslant q$) formed of elements of $\Omega$.
Let $\mathfrak{a}$ be the ideal of $B$ generated by the $P_\alpha$; the preceding interpretation and the theorem of the zeros of Hilbert show that the hypothesis on $F_{(K)}$ implies that $V(\mathfrak{a}_\eta)=\varnothing$, designating by $\eta$ the generic point of $\operatorname{Spec}(A)$; lemma \sref{IV.9.7.5.1} then proves that in a neighbourhood of $\eta$, we have $V(\mathfrak{a}_s)=\varnothing$, and consequently for these values of $x$, $F_{(\boldsymbol{k}(x))}$ does not admit any factorisation $F_{(\boldsymbol{k}(x))}=G_1 G_2$, where $G_1$, $G_2$ are polynomials of respective degrees $p$ and $q$, whose coefficients lie in an algebraic closure of $\boldsymbol{k}(x)$.
It suffices to apply this result for all pairs of integers $(p,q)$ such that $p>0$, $q>0$ and $p+q=d$, to obtain the conclusion of lemma \sref{IV.9.7.5}.
\end{proof}

\begin{proposition}[9.7.6]
\label{IV.9.7.6}
Let $S$ be a Noetherian integral prescheme with generic point $\eta$, $f:X\to S$ a morphism of finite type, and $\sh{F}$ a coherent $\sh{O}_X$-Module.
If $\sh{F}_\eta$ has no embedded associated prime cycle, there exists a neighbourhood $U$ of $\eta$ in $S$ such that, for every $s\in U$, $\sh{F}_s$ has no embedded associated prime cycle.
\end{proposition}

\begin{proof}
As $X_\eta$ is Noetherian, there exists an injective homomorphism $v:\sh{F}_\eta\to\bigoplus_{i=1}^m\sh{G}_i$, where each $\sh{G}_i$ is irredundant and $\operatorname{Ass}(\sh{F}_\eta)$ is the set of the $x_i$, with $\{x_i\}=\operatorname{Ass}(\sh{G}_i)$ \sref{IV.3.2.6}; if $Z_i$ is the closure of $\{x_i\}$ in $X_\eta$, then $Z_i=\operatorname{Supp}(\sh{G}_i)$ \sref{IV.3.1.4}, and by hypothesis the $Z_i$ are the irreducible components of $Z=\operatorname{Supp}(\sh{F}_\eta)$.
It follows from \sref{IV.8.5.2},~(i) and (ii) (restricting to $S$ if necessary a neighbourhood of $\eta$) and \sref{IV.8.5.8} that there exist coherent $\sh{O}_X$-Modules $\sh{F}_i$ and an injective homomorphism $u:\sh{F}\to\bigoplus_{i=1}^m\sh{F}_i$ such that $(\sh{F}_i)_\eta=\sh{G}_i$ for every $i$ and $v=u_\eta$.
If $Y_i=\operatorname{Supp}(\sh{F}_i)$, we have $(Y_i)_s=\operatorname{Supp}((\sh{F}_i)_s)$ for every $s\in S$ \sref[I]{I.9.1.13} and in particular $(Y_i)_\eta=Z_i$ for every $i$.
The hypothesis implies that for $i\neq j$, $Z_i-(Z_i\cap Z_j)$ is open and dense in $Z_i$; we deduce therefore from \sref{IV.9.5.3} and \sref{IV.9.5.4} that in a neighbourhood of $\eta$, $(Y_i)_s-((Y_i)_s\cap(Y_j)_s)$ is open and dense in $(Y_i)_s$.
Suppose that one has proved that each of the $(\sh{F}_i)_s$ is without associated prime cycle for $s\in U$.
Then the elements of $\operatorname{Ass}((\sh{F}_i)_s)$ are the maximal points of $(Y_i)_s$; none of them can therefore belong to $(Y_j)_s$ for $i\neq j$, and the proposition will be proved. We may thus suppose that $\sh{F}_\eta$ is \emph{irredundant}; moreover, we may reduce to the case where $S=\operatorname{Spec}(A)$ is affine. Then $X$ is a union of finitely many affine opens $V_j$, and if $V_j\cap\operatorname{Supp}(\sh{F}_\eta)=\varnothing$, we also have $V_j\cap\operatorname{Supp}(\sh{F}_s)=\varnothing$ for $s$ near $\eta$ \sref{IV.9.5.1}. We may therefore reduce to the case where $X=\operatorname{Spec}(B)$ is also affine and $f:X\to S$ is dominant. Thus $A$ is a Noetherian integral ring, a subring of a finite-type $A$-algebra $B$, and $\sh{F}=\widetilde{M}$ with $M$ a finite $B$-module. By hypothesis, if $K$ is the field of fractions of $A$, the $B_{(K)}$-module $M_{(K)}$ is irredundant.
Let $\mathfrak{q}$ be the unique

\oldpage[IV-3]{79}

element of $\operatorname{Ass}(M_{(K)})$, and let $\mathfrak{p}$ be the prime ideal of $B$ which is the inverse image of $\mathfrak{q}$ by the canonical map $B\to B_{(K)}$.
We know \sref{IV.5.11.1.1} that there exists a finite filtration $(N_h)_{0\leqslant h\leqslant m}$ of $M_{(K)}$ such that $N_0=M_{(K)}$, $N_m=0$, and $N_h/N_{h+1}$ is isomorphic to a nonzero $B_{(K)}$-submodule of $B_{(K)}/\mathfrak{q}$.
Let $M_h$ be the inverse image of $N_h$ by the canonical map $M\to M_{(K)}$.
It follows from \sref{IV.9.4.4} that for $s$ sufficiently close to $\eta$, $(M_h)_{(\boldsymbol{k}(s))}$ identifies with a sub-$B_{(\boldsymbol{k}(s))}$-module of $M_{(\boldsymbol{k}(s))}$, and the quotient $(M_h)_{(\boldsymbol{k}(s))}/(M_{h+1})_{(\boldsymbol{k}(s))}$ with a nonzero $B_{(\boldsymbol{k}(s))}$-submodule of $B_{(\boldsymbol{k}(s))}/\mathfrak{p}_{(\boldsymbol{k}(s))}=(B/\mathfrak{p})_{(\boldsymbol{k}(s))}$.
Taking into account \sref{IV.3.1.7}, with the same notations, it remains to prove that if $B$ is \emph{integral}, then $B_{(\boldsymbol{k}(s))}$ has no embedded associated prime ideals for $s$ close to $\eta$.
Now, replacing $A$ by $A_g$ and $B$ by $B_g$ (where $g$ is a nonzero element of $A$), we can suppose that $B$ contains a polynomial sub-algebra $C=A[\mathrm{T}_1,\ldots,\mathrm{T}_n]$ such that $B$ is a $C$-module of finite type (Bourbaki, \emph{Alg.\ comm.}, chap.\ V, \S~3, n$^\circ$~1, cor.\ 1 of th.\ 1).
As $B_{(K)}$ is a torsion-free $C_{(K)}$-module, we may apply the preceding argument once more, replacing $B$, $M$, and $\mathfrak{q}$ by $C$, $B$, and $(0)$ respectively; it suffices to show that for $s$ close to $\eta$, $C_{(\boldsymbol{k}(s))}$ has no embedded associated prime ideals.
But this is obvious since $C_{(\boldsymbol{k}(s))}=\boldsymbol{k}(s)[\mathrm{T}_1,\ldots,\mathrm{T}_n]$ is an integral ring.
\end{proof}

\begin{theorem}[9.7.7]
\label{IV.9.7.7}
Let $f:X\to S$ be a morphism of finite presentation, and let $E$ be the set of $s\in S$ for which $X_s$ has one of the following properties: being
\begin{enumerate}[label=(\roman*)]
  \item geometrically irreducible;
  \item geometrically connected;
  \item geometrically reduced;
  \item geometrically integral.
\end{enumerate}
Then $E$ is locally constructible in $S$.
\end{theorem}

\begin{proof}
To see that the properties considered are constructible, let us note first that they trivially satisfy condition \sref{IV.9.2.1},~(i), by virtue of \sref{IV.4.5.6},~(i) and \sref{IV.4.6.5},~(i).
It remains therefore to verify \sref{IV.9.2.1},~(ii), so we can suppose $S=\operatorname{Spec}(A)$ affine, Noetherian, integral, of generic point $\eta$.
By virtue of \sref{IV.4.6.8}, there exists a finite extension $K'$ of $K=\boldsymbol{k}(\eta)$ such that the irreducible components (resp.\ connected components) of $(X_\eta)_{(K')}$ are geometrically irreducible (resp.\ geometrically connected) and $((X_\eta)_{(K')})_\mathrm{red}$ is geometrically reduced.
Furthermore, there exists a base of $K'$ over $K$ formed of elements integral over $A$, the integral ring $A'$ generated by these elements is finite over $A$ and has $K'$ as field of fractions.
Set $S'=\operatorname{Spec}(A')$; the morphism $g:S'\to S$ is finite and dominant, therefore surjective \sref[II]{II.6.1.10}. Set $X'=X_{(S')}$, so that, if $\eta'$ is the generic point of $S'$, then $X'_{\eta'}=(X_\eta)_{(K')}$. The set $E'$ of $s'\in S'$ such that $X'_{s'}$ has one of the properties (i), (ii), (iii), or (iv) is equal to $g^{-1}(E)$ \sref{IV.9.2.2},~(iv), where $E$ of course corresponds to the same property. Since $g$ is surjective, $E=g(E')$ and $S-E=g(S'-E')$. Moreover, $g$ is closed and $g^{-1}(\eta)=\{\eta'\}$ because $A'$ is integral and finite over $A$ (Bourbaki, \emph{Alg.\ comm.}, chap.\ V, \S~2, n$^\circ$~1, cor.\ 1 of prop.\ 1); hence the image under $g$ of every neighbourhood of $\eta'$ is a neighbourhood of $\eta$.
The theorem will therefore be proved if one proves that $E'$ or $S'-E'$ is a neighbourhood of $\eta'$.
In other words, we can henceforth suppose that the irreducible components (resp.\ connected components) of $X_\eta$ are geometrically

\oldpage[IV-3]{80}

irreducible (resp.\ geometrically connected) and that $(X_\eta)_\mathrm{red}$ is geometrically reduced.

Suppose first that $\eta\in S-E$ for one of the properties (i) to (iv).
If $X_\eta$ is not geometrically irreducible (resp.\ geometrically connected), it is not irreducible (resp.\ connected) by virtue of the preceding hypothesis, and therefore it is the same for $X_s$ for $s$ in a neighbourhood of $\eta$ \sref{IV.9.7.1}, and \emph{a fortiori} in this neighbourhood $X_s$ is not geometrically irreducible (resp.\ geometrically connected).
On the other hand, if $X_\eta$ is not geometrically reduced, it is not reduced (since otherwise it would be equal to $(X_\eta)_\mathrm{red}$ which is geometrically reduced by hypothesis); therefore, in a neighbourhood of $\eta$, $X_s$ is not reduced \sref{IV.9.7.2} and hence not geometrically reduced.
Finally, if $X_\eta$ is not geometrically integral, either it is not reduced, in which case we have just seen that $X_s$ is not reduced (nor \emph{a fortiori} integral) in a neighbourhood of $\eta$; or $X_\eta$ is reduced (therefore geometrically reduced by hypothesis), in which case it is not geometrically irreducible, and we have just seen above that the same then holds for $X_s$ for $s$ close to $\eta$; \emph{a fortiori} $X_s$ is not geometrically integral for these values of $s$.

We shall now suppose that $\eta\in E$ and examine separately each of the properties considered.

$1^\circ$ Suppose that $X_\eta$ is geometrically integral.
Let $L$ be the field of rational functions on $X_\eta$; the hypothesis on $X_\eta$ implies that $L$ is a separable extension of $K$ \sref{IV.4.6.3}, therefore a finite separable extension of a pure extension $K(\mathrm{T}_1,\ldots,\mathrm{T}_n)$ ($\mathrm{T}_i$ indeterminates); there is then an element $x\in L$, integral over the ring $K[\mathrm{T}_1,\ldots,\mathrm{T}_n]$, and such that $L=K(\mathrm{T}_1,\ldots,\mathrm{T}_n)(x)$; let $G\in K[\mathrm{T}_1,\ldots,\mathrm{T}_n,\mathrm{T}_{n+1}]$ be its minimal polynomial.
There exists $g\neq 0$ of $A$ such that all the coefficients $\neq 0$ of $G$ (which all belong to $K$) belong to $A_g$; replacing $A$ by $A_g$ (which amounts to replacing $S$ by a neighbourhood of $\eta$), we can therefore suppose that $G$ has its coefficients in $A$; designating by $F$ the polynomial $G$ considered as element of $A[\mathrm{T}_1,\ldots,\mathrm{T}_n,\mathrm{T}_{n+1}]$, we then have $G=F_{(\boldsymbol{k}(\eta))}$.
Set $Y=\operatorname{Spec}(A[\mathrm{T}_1,\ldots,\mathrm{T}_{n+1}]/(F))$, so that $Y_\eta=\operatorname{Spec}(K[\mathrm{T}_1,\ldots,\mathrm{T}_{n+1}]/(G))$; $Y_\eta$ is an integral scheme having $L$ as field of rational functions.
As $X_\eta$ and $Y_\eta$ are Noetherian, there exists a non-empty open $V\subset Y_\eta$ and an open immersion (necessarily dominant) $v:V\to X_\eta$ \sref[I]{I.6.5.1},~(ii) and \sref[I]{I.6.5.4},~(ii).
Let $W$ be an open of $Y$ such that $W\cap Y_\eta=V$; applying \sref{IV.8.8.2},~(i) and \sref{IV.8.10.5},~(iii) following the method of \sref{IV.8.1.2},~\emph{a}), replacing $S$ by a neighbourhood of $\eta$, we can suppose that $v=w_\eta$ where $w:W\to X$ is an open immersion.

We saw in \sref{IV.4.6.3} that the criterion for an integral algebraic prescheme to be geometrically integral depends only on its field of rational functions. Since $X_\eta$ is geometrically integral by hypothesis, the same is true of $Y_\eta$, and the definition of $G$ therefore implies that this polynomial is geometrically irreducible \sref{IV.9.7.4}.
Applying \sref{IV.9.7.5}, we see that there is a neighbourhood $U$ of $\eta$ in $S$ such that $F_{(\boldsymbol{k}(s))}$ is geometrically irreducible for every $s\in U$, therefore $Y_s$ is geometrically integral for $s\in U$ \sref{IV.9.7.4}; furthermore we can also suppose that for $s\in U$, $W_s$ is not empty \sref{IV.9.5.1}, and therefore is geometrically integral \sref{IV.4.6.3}; finally, we can also suppose that

\oldpage[IV-3]{81}

for $s\in U$, $w_s:W_s\to X_s$ is an open immersion \sref{IV.9.6.1},~(x), and that its image in $X_s$ is everywhere dense \sref{IV.9.5.3}.
In other words, for $s\in U$, there is in $X_s$ an open everywhere dense $W'_s$ which is geometrically integral; the criterion \sref{IV.4.5.9},~\emph{c}) therefore shows that $X_s$ is already geometrically irreducible for $s\in U$.
Finally, as $X_\eta$ is reduced and therefore has no embedded associated prime cycle \sref{IV.3.2.1}, we can also suppose that for $s\in U$, $X_s$ has no embedded associated prime cycle \sref{IV.9.7.6}; this being the case (for $s\in U$ fixed), let $\Omega$ be an algebraically closed extension of $\boldsymbol{k}(s)$, and let $p:(X_s)_{(\Omega)}\to X_s$ be the canonical projection; $p^{-1}(W'_s)$ is an open dense in $(X_s)_{(\Omega)}$ \sref{IV.2.3.10} and is integral by hypothesis; furthermore $(X_s)_{(\Omega)}$ has no embedded associated prime cycle \sref{IV.4.2.7}, therefore we conclude from \sref{IV.3.2.1} that $(X_s)_{(\Omega)}$ is reduced; this completes the proof that $X_s$ is geometrically integral \sref{IV.4.6.1}.

$2^\circ$ Suppose that $X_\eta$ is geometrically irreducible; as $X_{\mathrm{red}}$ is also of finite type over $S$ \sref{IV.1.5.4},~(vi), taking into account \sref[I]{I.5.1.8}, we can replace $X$ by $X_{\mathrm{red}}$; then $X_\eta$ is also integral, and as by hypothesis $X_\eta$ is geometrically reduced, it is geometrically integral. We are then in the conditions of $1^\circ$, and we conclude (returning to the initial hypotheses) that $X_s$ is geometrically irreducible for $s$ close to $\eta$.

$3^\circ$ Suppose $X_\eta$ geometrically connected, and let $Z_i$ ($1\leqslant i\leqslant n$) be the irreducible components of $X_\eta$; there exists (by virtue of \sref{IV.5.10.8.1} applied to $\mathfrak{T}=\{\varnothing\}$) a \emph{surjective} map $j\mapsto\nu(j)$ of an interval $[1,m]$ of $\mathbf{N}$ onto $[1,n]$ such that $Z_{\nu(j)}\cap Z_{\nu(j+1)}\neq\varnothing$ for $1\leqslant j\leqslant m$.
For every $i$, let $X_i$ be the closure of $Z_i$ in $X$, and let $Y$ be the union of the $X_i$; as $Y_\eta=X_\eta$ by definition, we can, by virtue of \sref{IV.9.5.1}, suppose that $Y_s=X_s$ for every $s\in S$, so that $X_s$ is the union of the $(X_i)_s$.
Now, considering the reduced sub-preschemes of $X$ having the $X_i$ as underlying spaces, we see by $2^\circ$ that there exists a neighbourhood $U$ of $\eta$ in $S$ such that for $s\in U$ the $(X_i)_s$ are \emph{geometrically irreducible} (since the $(X_i)_\eta$ can be supposed to be geometrically irreducible, as we saw at the beginning).
Furthermore, we can also suppose that for $s\in U$, $(X_{\nu(j)})_s\cap(X_{\nu(j+1)})_s\neq\varnothing$ \sref{IV.9.5.1} for $1\leqslant j\leqslant m$; we conclude immediately that $X_s$ is connected, therefore \sref{IV.4.5.13.1} geometrically connected for $s\in U$.

$4^\circ$ Suppose $X_\eta$ geometrically reduced; let $Z_i$ be the irreducible components of $X_\eta$, and let $W_i$ be the interior of $Z_i$ in $X_\eta$; there is for each $i$ an open $V_i$ of $X$ such that $W_i=V_i\cap X_\eta$ for every $i$; as the $W_i$ are opens, pairwise disjoint, and their union is dense in $X_\eta$, \sref{IV.9.5.1},~\sref{IV.9.5.3}, and \sref{IV.9.5.4} allow us to suppose that for $s$ close to $\eta$, the $(V_i)_s$ are opens, pairwise disjoint, in $X_s$ and their union is dense in $X_s$.
Furthermore, as the $W_i$ are geometrically reduced and have been supposed at the beginning geometrically irreducible, it follows from $1^\circ$ that for $s$ close to $\eta$, the $(V_i)_s$ are geometrically integral, and \emph{a fortiori} reduced.
On the other hand, we deduce from \sref{IV.9.7.6} that for $s$ close to $\eta$, $X_s$ has no embedded associated prime cycle, since the same is true of the reduced prescheme $X_\eta$ \sref{IV.3.2.1}; we therefore conclude from \sref{IV.3.2.1} that $X_s$ is reduced, and from \sref{IV.4.6.1} that $X_s$ is geometrically reduced.
\end{proof}

\oldpage[IV-3]{82}

The parts of the statement of \sref{IV.9.7.7} concerning properties (i) and (ii) generalize as follows.

\begin{proposition}[9.7.8]
\label{IV.9.7.8}
Let $S$ be an irreducible prescheme with generic point $\eta$, and let $f:X\to S$ be a morphism of finite presentation. Let $n$ (resp.\ $n'$) be the geometric number of irreducible components (resp.\ connected components) of $X_\eta$ \sref{IV.4.5.2}. Then there exists a neighborhood $U$ of $\eta$ in $S$ such that, for every $s\in U$, the geometric number of irreducible components (resp.\ connected components) of $X_s$ is equal to $n$ (resp.\ $n'$).
\end{proposition}

\begin{proof}
Taking into account that the geometric number of irreducible components (resp.\ connected components) of an algebraic prescheme is invariant under extension of the base field \sref{IV.4.5.6}, one sees by the method of \sref{IV.9.2.3} that one may reduce to the case where $S$ is affine, Noetherian, and integral. Moreover, arguing as at the beginning of the proof of \sref{IV.9.7.7}, one sees that one may suppose that the irreducible components (resp.\ connected components) of $X_\eta$ are geometrically irreducible (resp.\ geometrically connected). We already know from \sref{IV.9.7.1} that, for every $s$ near $\eta$, $X_s$ has at least $n$ irreducible components and $n'$ connected components.

On the other hand, if $Z_i$ (resp.\ $Z'_j$) are the irreducible components (resp.\ connected components) of $X_\eta$, let $X_i$ (resp.\ $X'_j$) be the reduced closed subprescheme of $X$ whose underlying space is the closure of $Z_i$ (resp.\ $Z'_j$) in $X$. It follows from \sref{IV.9.5.1} that, in a neighborhood of $\eta$, the $(X_i)_s$ (resp.\ $(X'_j)_s$) cover $X_s$, and that the $(X'_j)_s$ are pairwise disjoint. Since the $(X_i)_s$ (resp.\ $(X'_j)_s$) are geometrically irreducible (resp.\ geometrically connected) by \sref{IV.9.7.7} for $s$ near $\eta$, it follows that $X_s$ has at most $n$ irreducible components and at most $n'$ connected components, which proves the proposition.
\end{proof}

\begin{corollary}[9.7.9]
\label{IV.9.7.9}
Let $f:X\to S$ be a morphism of finite presentation. For every $s\in S$, let $n(s)$ (resp.\ $n'(s)$) be the geometric number of irreducible components (resp.\ connected components) of $f^{-1}(s)$ \sref{IV.4.5.2}. Then the function $s\mapsto n(s)$ (resp.\ $s\mapsto n'(s)$) is locally constructible on $S$.
\end{corollary}

\begin{proof}
It is enough to show that the property ``$X$ is an algebraic prescheme over a field $k$ and the geometric number of irreducible components (resp.\ connected components) of $X$ is equal to $n$ (resp.\ $n'$)'' is constructible. It follows from \sref{IV.4.5.6} that this property satisfies condition \sref{IV.9.2.1},~(i); hence one is reduced to the case where $S$ is affine, Noetherian, and integral with generic point $\eta$, and the conclusion then follows from \sref{IV.9.7.8}.
\end{proof}

\begin{corollary}[9.7.10]
\label{IV.9.7.10}
Let $X$ be a locally Noetherian prescheme such that, if $X'$ is the normalization of $X_{\mathrm{red}}$, the canonical morphism $f:X'\to X$ is finite. Then the set of points $x\in X$ such that $X$ is geometrically unibranch at $x$ is locally constructible in $X$.
\end{corollary}

\begin{proof}
Indeed, this set is, by definition, the set of points $x\in X$ such that the number of geometric points of $f^{-1}(x)$ is equal to $1$. Since $f$ is finite, this number is also the geometric number of irreducible components of the discrete space $f^{-1}(x)$ (taking account of the definition of the normalization \sref[II]{II.6.3.8} and of \sref{IV.4.5.11}); the conclusion therefore follows from \sref{IV.9.7.9}.
\end{proof}

\begin{remark}[9.7.11]
\label{IV.9.7.11}
Let $Z$ be a locally constructible subset of $X$ such that, for every $s\in S$, $Z\cap X_s$ is open in $X_s$, and write $Z_s$ for the subprescheme of $X_s$ induced

\oldpage[IV-3]{83}

on the open subset $Z\cap X_s$. Then, in \sref{IV.9.7.7}, one may replace $X_s$ everywhere by $Z_s$ without changing the conclusion; this follows by repeating the argument used in \sref{IV.9.6.7},~(ii).
\end{remark}

\begin{proposition}[9.7.12]
\label{IV.9.7.12}
Let $f:X\to S$ be a morphism of finite presentation, and let $g:S\to X$ be an $S$-section of $X$ \sref[I]{I.2.5.5}. For every $s\in S$, let $X_s^0$ be the connected component of $f^{-1}(s)$ containing $g(s)$. Then the union $X^0$ of the $X_s^0$, for $s\in S$, is a locally constructible subset of $X$.
\end{proposition}

\begin{proof}
We first show that one may reduce to the case where $S$ is affine and Noetherian. One may always suppose that $S=\operatorname{Spec}(A)$ is affine; with the notation of \sref{IV.9.2.3}, one has $f=(f_0)_{(S)}$, where $f_0:X_0\to S_0$ is a morphism of finite type, and one may moreover suppose that there exists an $S_0$-section $g_0:S_0\to X_0$ such that $g=(g_0)_{(S)}$ \sref{IV.8.9.1}. If $p:S\to S_0$ is the canonical morphism, then, for every $s_0\in S_0$, the connected component $(X_0)^0_{s_0}$ of $f_0^{-1}(s_0)$ containing $g_0(s_0)$ is geometrically connected \sref{IV.4.5.13}; consequently, if $s_0=p(s)$, one has $X_s^0=q^{-1}((X_0)^0_{s_0})$, where $q:X_s\to (X_0)_{s_0}$ is the canonical projection \sref{IV.4.5.8} and \sref{IV.4.4.1}. Our assertion therefore follows from \sref{IV.1.8.2}.

We now use the constructibility criterion \sref[0\textsubscript{III}]{0.9.2.3}. Let $x$ be a point of $X$, let $Z$ be the reduced subprescheme of $X$ whose underlying space is $\overline{\{x\}}$, and let $Y$ be the reduced subprescheme of $S$ whose underlying space is $\overline{f(Z)}$. Since the restriction of $f$ to $Z$ factors as $Z\to Y\to S$ \sref[I]{I.5.2.2}, one may replace $S$ by $Y$, that is, suppose that $f(x)=\eta$, the generic point of the integral prescheme $S$.

By hypothesis, $X_\eta$ is the sum of two subpreschemes $X_\eta^0$ and $X_\eta^1$ induced on complementary open subsets of $X_\eta$. By \sref{IV.9.5.4} and \sref{IV.9.5.1}, after replacing $S$, if necessary, by a neighborhood of $\eta$, one may suppose that $X$ is the union of two disjoint open subsets $X^0$ and $X^1$ such that $(X^0)_\eta=X_\eta^0$ and $(X^1)_\eta=X_\eta^1$. Since $g:S\to X$ is continuous and injective, $S$ is the union of the two disjoint open subsets $g^{-1}(X^0)$ and $g^{-1}(X^1)$; but $S$ is irreducible, and \emph{a fortiori} connected, so one of these two open subsets is empty. Since $g(\eta)\in X_\eta^0$ by definition, one has $g^{-1}(X^1)=\varnothing$.

In other words, $g$ is an $S$-section of $X^0$. On the other hand, since $(X^0)_\eta$ is geometrically connected \sref{IV.4.5.13}, the same is true of $(X^0)_s$ for every $s$ near $\eta$ by \sref{IV.9.7.7}; and since $g(s)\in (X^0)_s$, one indeed has $(X^0)_s=X_s^0$.
\end{proof}

\subsection{Primary decomposition in the neighbourhood of a generic fibre}
\label{subsection:IV.9.8}
\label{IV.9.8}

\begin{proposition}[9.8.1]
\label{IV.9.8.1}
Let $f:X\to S$ be a morphism of finite presentation, $\sh{F}$ a quasi-coherent $\sh{O}_X$-Module of finite presentation.
Then the set $E$ of $s\in S$ such that $\sh{F}_s$ is an $\sh{O}_{X_s}$-Module without associated immersed prime cycle is locally constructible.
\end{proposition}

\begin{proof}
\label{IV.9.8.1-proof}
We know that for quasi-coherent Modules on algebraic preschemes, the property of being without associated immersed prime cycle is invariant by change of base field \sref{IV.4.2.7}; we can thus restrict to the case where $S$ is affine, Noetherian and integral, with generic point $\eta$, and to proving that in this case $E$ or $S-E$ is a neighbourhood of $\eta$.
We saw in \sref{IV.9.7.6} that if $\eta\in E$, $E$ is a neighbourhood of $\eta$; it remains thus to consider the case where $\sh{F}_\eta$ has associated immersed prime cycles.
We can restrict to the case where $X$ is affine, and where there is a coherent sub-$\sh{O}_{X_\eta}$-Module $\sh{H}$ of $\sh{F}_\eta$ whose support $Z'$ is non-empty and \emph{rare} with respect to the support $Z$ of $\sh{F}$ \sref{IV.3.1.3}; then, by virtue of \sref{IV.8.5.2},~(i) and~(ii),
\oldpage[IV-3]{84}
and \sref{IV.8.5.8}, applied following the method of \sref{IV.8.1.2},~$a$), there exists a coherent sub-$\sh{O}_X$-Module $\sh{G}$ of $\sh{F}$ such that $\sh{G}_\eta=\sh{H}$; if $Y=\operatorname{Supp}(\sh{F})$ and $Y'=\operatorname{Supp}(\sh{G})$, then consequently $Y_s=\operatorname{Supp}(\sh{F}_s)$ and $Y'_s=\operatorname{Supp}(\sh{G}_s)$ for all $s\in S$ \sref[I]{I.9.1.13}, and in particular $Z=Y_\eta$ and $Z'=Y'_\eta$; as $Z-Z'$ is dense in $Z$ and $Z'\neq\varnothing$, there is a neighbourhood $U$ of $\eta$ such that for $s\in U$, $Y_s-Y'_s$ is dense in $Y_s$ and $Y'_s\neq\varnothing$ \sref{IV.9.5.1} and \sref{IV.9.5.3}; by considering a generic point of an irreducible component of $Y'_s$ and a sufficiently small neighbourhood of this point in $X_s$, we deduce immediately from \sref{IV.3.1.3} that $\sh{F}_s$ admits associated immersed prime cycles.
\end{proof}

\medskip
\noindent
\emph{\sref{IV.9.8.2}}
\label{IV.9.8.2}
Let $S$ be a Noetherian prescheme, \emph{integral}, of generic point $\eta$, $f:X\to S$ a morphism of finite type, $\sh{F}$ a coherent $\sh{O}_X$-Module.
Consider an irredundant reduced decomposition $(\sh{G}_i)_{i\in I}$ of $\sh{F}_\eta$ \sref{IV.3.2.6}; the $\sh{G}_i$ are thus quotients of $\sh{F}_\eta$ and there is an injective homomorphism $h:\sh{F}_\eta\to\bigoplus_i\sh{G}_i$; furthermore $\operatorname{Ass}(\sh{G}_i)$ is reduced to a point $x_i$.
Let $Z_i$ be the closure of $\{x_i\}$ in $X$, so that $(Z_i)_\eta=\operatorname{Supp}(\sh{G}_i)$.
Denote by $\sh{J}_i$ the coherent Ideal of $\sh{O}_X$ defining the closed sub-prescheme of $X$ of underlying space $Z_i$ (sub-prescheme still denoted $Z_i$).
By virtue of \sref{IV.8.5.2} and \sref{IV.8.5.8}, applied following the method of \sref{IV.8.1.2},~$a$), we can (by restricting $S$ to a neighbourhood of $\eta$ if needed) suppose that there exist quotients $\sh{F}_i$ of $\sh{F}$ ($i\in I$) such that $\sh{G}_i=(\sh{F}_i)_\eta$ for all $i$, and a homomorphism $g:\sh{F}\to\bigoplus_i\sh{F}_i$ such that $h=g_\eta$.
Furthermore, by \sref[I]{I.9.3.5}, there exists an integer $m$ such that $(\sh{J}_i^m\sh{F}_i)_\eta=(\sh{J}_i)_\eta^m\sh{G}_i=0$ and, by restricting $S$ again, we can suppose that $\sh{J}_i^m\sh{F}_i=0$ \sref{IV.8.5.2.5}, so that the support of $\sh{F}_i$ is contained in $Z_i$; but as it is closed and contains $x_i$, it is necessarily equal to $Z_i$.

\begin{proposition}[9.8.3]
\label{IV.9.8.3}
Under the conditions of \emph{\sref{IV.9.8.2}}, for all $s\in S$, and all $i\in I$, denote by $x_{is\alpha}$ ($\alpha\in J_{s,i}$) the maximal points of $(Z_i)_s$.
There exists a neighbourhood $U$ of $\eta$ in $S$ such that, for all $s\in U$, the $x_{is\alpha}$ (for $i\in I$ and, for each $i$, $\alpha\in J_{s,i}$) are pairwise distinct and that $\operatorname{Ass}(\sh{F}_s)$ is the set of $x_{is\alpha}$ (in other words, the prime cycles associated to $\sh{F}_s$ are the irreducible components of the $(Z_i)_s$).
Furthermore, one can take $U$ such that, for the closure of $\{x_{is\alpha}\}$ in $X_s$ to be a maximal associated prime cycle of $\sh{F}_s$, it is necessary and sufficient that $(Z_i)_\eta$ (closure of $\{x_i\}$ in $X_\eta$) be a maximal associated prime cycle of $\sh{F}_\eta$.
\end{proposition}

\begin{proof}
\label{IV.9.8.3-proof}
It follows from \sref{IV.3.1.3},~$c'$) that for each $i$, there exists an open $W_i$ in $X_\eta$ such that $W_i\cap(Z_i)_\eta$ is non-empty, and a coherent $\sh{O}_{X_\eta}$-Module $\sh{H}_i$, of support $W_i\cap(Z_i)_\eta$, such that there is an injective homomorphism $v_i:\sh{H}_i\to\sh{F}_\eta|W_i$.
Let $V_i$ be an open of $X$ such that $V_i\cap X_\eta=W_i$; applying, as in \sref{IV.9.8.2}, the results of \sref{IV.8.5.2} and \sref{IV.8.5.8}, one can (by restricting $S$) suppose that there exists a coherent Module $\sh{F}'_i$, of support $V_i\cap Z_i$, and a homomorphism $u_i:\sh{F}'_i\to\sh{F}|V_i$ such that $\sh{H}_i=(\sh{F}'_i)_\eta$ and $v_i=(u_i)_\eta$.
We will show that there is a neighbourhood $U$ of $\eta$ such that for $s\in U$, the following properties are satisfied:
\begin{enumerate}[label=(\roman*)]
  \item The $(\sh{F}_i)_s$ are without associated immersed prime cycle.
  \item The homomorphism $g_s:\sh{F}_s\to\bigoplus_i(\sh{F}_i)_s$ is injective.
  \item $(Z_i)_s\cap(V_i)_s$ is dense in $(Z_i)_s$, $(\sh{F}'_i)_s$ has support $(Z_i)_s\cap(V_i)_s$, and $(u_i)_s:(\sh{F}'_i)_s\to\sh{F}_s|(V_i)_s$ is injective.
\end{enumerate}
\oldpage[IV-3]{85}

\begin{enumerate}[label=(\roman*)]
  \item[(iv)] For $i\neq j$, every irreducible component of $(Z_i)_s$ is distinct from every irreducible component of $(Z_j)_s$.
\end{enumerate}

Now, (i) has already been seen \sref{IV.9.7.6}; (ii) is a particular case of \sref{IV.9.4.5}; (iii) follows from \sref{IV.9.5.3} and \sref{IV.9.4.5}.
Finally, if $i\neq j$, $(Z_i)_\eta\cap(Z_j)_\eta=(Z_i\cap Z_j)_\eta$ is rare in $(Z_i)_\eta$ \emph{or} in $(Z_j)_\eta$; suppose for example that $(Z_i\cap Z_j)_\eta$ is rare in $(Z_j)_\eta$; then it follows from \sref{IV.9.5.3} and \sref{IV.9.5.4} that for $s$ in a neighbourhood of $\eta$, $(Z_i\cap Z_j)_s=(Z_i)_s\cap(Z_j)_s$ is rare in $(Z_j)_s$, which shows that no irreducible component of $(Z_j)_s$ can be contained in an irreducible component of $(Z_i)_s$, nor \emph{a fortiori} be equal to it.

This being so, it follows from (ii) and \sref{IV.3.1.7} that for $s\in U$, we have
\[
  \operatorname{Ass}(\sh{F}_s)\subset\bigcup_i\operatorname{Ass}((\sh{F}_i)_s),
\]
from (i) that $\operatorname{Ass}((\sh{F}_i)_s)$ is the set of $x_{is\alpha}$ ($\alpha\in J_{s,i}$).
On the other hand, by virtue of (iii) and the criterion \sref{IV.3.1.3},~$c'$), each of the $x_{is\alpha}$ ($\alpha\in J_{s,i}$, $i\in I$) belongs to $\operatorname{Ass}(\sh{F}_s)$.
Finally (iv) signifies that the $x_{is\alpha}$ are pairwise distinct.

It remains to prove the last assertion of the statement.
Now, it follows from (i) that for $s$ and for given $i$, none of the sets $\{x_{is\alpha}\}$ can be contained in another for $\alpha\in J_{s,i}$.
On the other hand, if $(Z_j)_\eta\subset(Z_i)_\eta$, we can suppose $U$ taken such that $(Z_j)_s\subset(Z_i)_s$ for $s\in U$ \sref{IV.9.5.1}, so each $x_{js\beta}$ belongs to the closure of an $x_{is\alpha}$; on the contrary, if $(Z_j)_\eta\cap(Z_i)_\eta$ is rare in $(Z_j)_\eta$, we have seen by proving (iv) that $(Z_j)_s\cap(Z_i)_s$ is rare in $(Z_j)_s$, so none of the $x_{js\beta}$ is adherent to an $x_{is\alpha}$.
In particular, if $(Z_j)_\eta$ is maximal, this is equivalent to saying that $(Z_j)_\eta\cap(Z_i)_\eta$ is rare in $(Z_j)_\eta$ for \emph{all} $i\neq j$, and we conclude that $(Z_j)_s\cap(Z_i)_s$ is rare in $(Z_j)_s$ for all $i\neq j$, so that every $x_{js\alpha}$ ($\alpha\in J_{s,j}$) is \emph{maximal}.
\end{proof}

\begin{corollary}[9.8.4]
\label{IV.9.8.4}
The notations and hypotheses being those of \emph{\sref{IV.9.8.2}}, there exists a neighbourhood $U$ of $\eta$ in $S$ such that, for all $s\in U$, each of the $(\sh{F}_i)_s$ is without associated immersed prime cycle; furthermore, if $(\sh{F}_{is\alpha})_{\alpha\in J_{s,i}}$ is the unique irredundant reduced decomposition of $(\sh{F}_i)_s$, then, for all $s\in U$, the family $(\sh{F}_{is\alpha})_{i\in I,\,\alpha\in J_{s,i}}$ is an irredundant reduced decomposition of $\sh{F}_s$.
\end{corollary}

\begin{proof}
\label{IV.9.8.4-proof}
The first assertion follows from \sref{IV.9.8.1} and the definition of $\sh{F}_i$.
On the other hand, we have seen in \sref{IV.9.8.3} that the homomorphism $\sh{F}_s\to\bigoplus_i(\sh{F}_i)_s$ is injective and by definition it is the same for each of the homomorphisms $(\sh{F}_i)_s\to\bigoplus_\alpha\sh{F}_{is\alpha}$, so the homomorphism $\sh{F}_s\to\bigoplus_{i,\alpha}\sh{F}_{is\alpha}$ is injective.
As we can suppose \sref{IV.9.4.5} that for $s\in U$ each of the $(\sh{F}_i)_s$ is a quotient of $\sh{F}_s$, the $\sh{F}_{is\alpha}$ are quotients of $\sh{F}_s$, and it remains to verify \sref{IV.3.2.5} that the $x_{is\alpha}$ are pairwise distinct and belong to $\operatorname{Ass}(\sh{F}_s)$, which has been proved in \sref{IV.9.8.3}.
\end{proof}

\begin{proposition}[9.8.5]
\label{IV.9.8.5}
Let $f:X\to S$ be a morphism of finite presentation, $\sh{F}$ a quasi-coherent $\sh{O}_X$-Module of finite presentation.
For all $s\in S$, let $E(s)$ (resp.~$E'(s)$) be the set (finite subset of $\mathbf{Z}\cup\{-\infty\}$) of dimensions of the prime cycles associated to $\sh{F}_s$ (resp.~of the maximal associated prime cycles, i.e., the irreducible components of $\operatorname{Supp}(\sh{F}_s)$).
Then the functions $s\mapsto E(s)$ and $s\mapsto E'(s)$ are locally constructible in $S$.
\end{proposition}

\begin{proof}
\label{IV.9.8.5-proof}
\oldpage[IV-3]{86}
It follows from \sref{IV.4.2.7} and \sref{IV.4.2.8} that the condition \sref{IV.9.2.1},~(i) is satisfied for the properties whose constructibility we wish to show.
We are thus reduced to the case where $S$ is affine, Noetherian and integral, and to proving that $E$ and $E'$ are constant in a neighbourhood of the generic point $\eta$ of $S$; but this follows from \sref{IV.9.8.3} and \sref{IV.9.5.5}.
\end{proof}

\begin{proposition}[9.8.6]
\label{IV.9.8.6}
With the hypotheses and notations of \emph{\sref{IV.9.8.2}}, let $i\in I$ be such that $(Z_i)_\eta$ is maximal (in other words, an irreducible component of $\operatorname{Supp}(\sh{F}_\eta)$).
Then there exists a neighbourhood $U$ of $\eta$ in $S$ such that for all $s\in U$ and all $\alpha\in J_{s,i}$, the geometric length of $\sh{F}_s$ at $x_{is\alpha}$ (relative to $\mathbf{k}(s)$) \sref{IV.4.7.5} is equal to the geometric length of $\sh{F}_\eta$ at $x_i$ (relative to $\mathbf{k}(\eta)$).
\end{proposition}

\begin{proof}
\label{IV.9.8.6-proof}
We can evidently restrict to the case where $S=\operatorname{Spec}(A)$ is affine; let us first show that we can reduce to the case where the sub-prescheme $(Z_i)_\eta$ of $X_\eta$, which is reduced, is \emph{geometrically integral}.
Indeed there is a finite extension $K'$ of $K=\mathbf{k}(\eta)$ such that $(((Z_i)_\eta)_{(K')})_{\mathrm{red}}$ is geometrically reduced and that the irreducible components of $((Z_i)_\eta)_{(K')}$ are geometrically irreducible \sref{IV.4.6.8}.
We proceed as in the proof of \sref{IV.9.7.7} by considering a sub-$A$-algebra $A'$ of $K'$ having $K'$ as its field of fractions and finite over $A$.
We set $S'=\operatorname{Spec}(A')$ and consider the finite surjective morphism $g:S'\to S$; let $X'=X_{(S')}$ and $\sh{F}'=\sh{F}\otimes_{\sh{O}_S}\sh{O}_{S'}$, and let $\eta'$ be the generic point of $S'$.
For every $s'\in S'$, put $s=g(s')$; if $T$ is an irreducible component of $\operatorname{Supp}(\sh{F}_s)$, the irreducible components $T'_j$ of $T_{(\mathbf{k}(s'))}$ are irreducible components of $\operatorname{Supp}(\sh{F}'_{s'})$ and dominate $T$ \sref{IV.4.2.7}, and the radicial multiplicities of $T$ with respect to $\sh{F}_s$ and of each $T'_j$ with respect to $\sh{F}'_{s'}$ are the same \sref{IV.4.7.9}.
The reasoning from the first part of \sref{IV.9.7.7} then shows that we can restrict to proving the proposition for $X'$ and $\sh{F}'$; and by virtue of the choice of $K'$, the reduced sub-preschemes having for underlying spaces the irreducible components of $((Z_i)_\eta)_{(K')}$ are geometrically integral \sref{IV.4.6.1}.

Suppose now that $(Z_i)_\eta$ is geometrically integral; then, by \sref{IV.9.7.7}, the same is true of $(Z_i)_s$ for $s$ in a neighbourhood of $\eta$; the definition \sref{IV.4.7.5} shows that it will therefore suffice to prove that the length of the $\sh{O}_{X_\eta,x_i}$-module $(\sh{F}_\eta)_{x_i}$ is equal to that of the $\sh{O}_{X_s,x_{is}}$-module $(\sh{F}_s)_{x_{is}}$ (where we have dropped the index $\alpha$, now unnecessary by hypothesis).
The question being evidently local on $X$, we can suppose (by restricting to a neighbourhood of $x_i$) that $\sh{F}=\sh{F}_i$, so that $\sh{F}_\eta$ is \emph{irredundant} on $X=\operatorname{Spec}(B)$ affine, and one writes $Z$ instead of $Z_i$, and $x$ the generic point of $Z$ (and of $Z_\eta$).
The Noetherian ring $B$ contains $A$ as sub-ring, and $\sh{F}=\widetilde{M}$, where $M$ is a $B$-module of finite type; furthermore, if $K$ is the field of fractions of $A$, the $B_{(K)}$-module $M_{(K)}$ is $\neq 0$ and irredundant.
Let $q$ be the unique element of $\operatorname{Ass}(M_{(K)})$ and let $\mathfrak{p}$ be the prime ideal of $B$, inverse image of $q$.
Using \sref{IV.5.11.1.1} as in the proof of \sref{IV.9.7.6}, we reduce to the case where $B$ is \emph{integral} and $M$ a sub-module $\neq 0$ of $B$; then $\sh{F}_\eta$ is a sub-$\sh{O}_{X_\eta}$-Module non-zero of $\sh{O}_{Z_\eta}$, and by virtue of \sref{IV.9.4.5}, for $s$ in a neighbourhood of $\eta$, $\sh{F}_s$ is isomorphic to a sub-$\sh{O}_{X_s}$-Module non-zero of $\sh{O}_{Z_s}$; as $Z_s$ is geometrically integral, the lengths of $(\sh{F}_\eta)_x$ and $(\sh{F}_s)_{x_s}$ are both equal to $1$, which finishes the proof.
\end{proof}

\begin{corollary}[9.8.7]
\label{IV.9.8.7}
\oldpage[IV-3]{87}
Let $f:X\to S$ be a morphism of finite presentation, $\sh{F}$ a quasi-coherent $\sh{O}_X$-Module of finite presentation.
For all $s\in S$, let $M(s)$ be the set of pairs $(d,m)$ such that there exists an irreducible component of $\operatorname{Supp}(\sh{F}_s)$ of dimension $d$ and of radicial multiplicity $m$ for $\sh{F}_s$ \sref{IV.4.7.8}.
Then the function $s\mapsto M(s)$ is locally constructible in $S$.
\end{corollary}

\begin{proof}
\label{IV.9.8.7-proof}
It follows from \sref{IV.4.2.7} and \sref{IV.4.7.9} that the condition \sref{IV.9.2.1},~(i) is satisfied for the property whose constructibility we wish to show.
We are thus reduced to the case where $S$ is affine, Noetherian and integral, and to proving that $M$ is constant in a neighbourhood of the generic point of $S$; but this follows from the proposition, namely \sref{IV.9.8.3} and \sref{IV.9.8.6}.
\end{proof}

\begin{proposition}[9.8.8]
\label{IV.9.8.8}
Let $f:X\to S$ be a morphism of finite presentation, $\sh{F}$ a quasi-coherent $\sh{O}_X$-Module of finite presentation.
For all $s\in S$, let $l(s)$ be the sum of the total multiplicities for $\sh{F}_s$ (relative to $\mathbf{k}(s)$) of the generic points of the irreducible components of $\operatorname{Supp}(\sh{F}_s)$ \sref{IV.4.7.12}.
Then the function $s\mapsto l(s)$ is locally constructible in $S$.
\end{proposition}

\begin{proof}
\label{IV.9.8.8-proof}
Taking into account \sref{IV.4.7.12}, the condition \sref{IV.9.2.1},~(i) is satisfied for the property whose constructibility we wish to show, and we are thus reduced to the case where $S$ is affine, Noetherian and integral, and to showing that $l$ is constant in a neighbourhood of the generic point $\eta$.
Using the notations of \sref{IV.9.8.2}, this follows from the definition \sref{IV.4.7.12}, from the fact that the geometric number of irreducible components of each $(Z_i)_s$ is constant in a neighbourhood of $\eta$ \sref{IV.9.7.8}, that the geometric length of $\sh{F}_s$ at $x_{is\alpha}$ is also equal to that of $\sh{F}_\eta$ at $x_i$ for each $i$ such that $(Z_i)_\eta$ is maximal \sref{IV.9.8.6}, and finally from the fact that the closure of $\{x_{is\alpha}\}$ is a maximal associated prime cycle of $\sh{F}_s$ if and only if $(Z_i)_\eta$ is a maximal associated prime cycle of $\sh{F}_\eta$ \sref{IV.9.8.3}.
\end{proof}

\begin{remark}[9.8.9]
\label{IV.9.8.9}
One can make precise in various ways the preceding propositions; let us limit ourselves to a statement given as an example.
Let us say that a finite subset $P$ of an algebraic $k$-prescheme $X$ is \emph{saturated} if, for every pair of points $x$, $y$ of $P$, the generic points of the irreducible components of $\overline{\{x\}}\cap\overline{\{y\}}$ also belong to $P$; for every finite subset $Q$ of $X$, there exists a smallest finite subset $P$ of $X$ containing $Q$ and saturated; we say that $P$ is the \emph{saturate} of $Q$.
For every coherent $\sh{O}_X$-Module $\sh{F}$, we will call \emph{primary skeleton} of $\sh{F}$ the system $(P,Q,\omega,d,m)$ where $Q=\operatorname{Ass}(\sh{F})$, $P$ is the saturate of $Q$, $\omega$ the ordering relation $\overline{\{x\}}\subset\overline{\{y\}}$ on $P$, $d$ the function $x\mapsto\dim\overline{\{x\}}$ on $P$, and $m$ the function $x\mapsto\operatorname{length}_{\sh{O}_x}\sh{F}_x$ defined on the set of maximal elements of $Q$ for the relation $\omega$.
We call \emph{virtual skeleton} any system $(P,Q,\omega,d,m)$ where $P$ is a set, $Q$ a subset of $P$, $\omega$ an order relation on $P$, $d$ a map from $P$ to $\mathbf{N}$, and $m$ a map from the set of maximal elements of $Q$ to $\mathbf{N}$; one defines in an obvious way the notion of \emph{isomorphism} of two virtual skeletons.
Finally, with the preceding notations, we call \emph{primary type} of $\sh{F}$ the class (for the relation of isomorphism of virtual skeletons) of the primary skeleton of $\sh{F}$.
It follows from \sref{IV.4.2.6}, \sref{IV.4.2.7}, \sref{IV.4.2.8}, \sref{IV.4.5.1} and \sref{IV.4.7.9} that if, for an algebraically closed extension $k'$ of $k$, one sets $X'=X\otimes_k k'$ and $\sh{F}'=\sh{F}\otimes_k k'$, the \emph{primary type} of $\sh{F}'$ is independent of
\oldpage[IV-3]{88}
the algebraically closed extension $k'$ of $k$ considered; we say that this is the \emph{geometric primary type} of $\sh{F}$.
Given the definitions, the statement we have in view is the following:

\medskip
\noindent
\emph{\sref{IV.9.8.9.1}}
\label{IV.9.8.9.1}
Let $f:X\to S$ be a morphism of finite presentation, $\sh{F}$ a quasi-coherent $\sh{O}_X$-Module of finite presentation.
For all $s\in S$, let $T(s)$ be the geometric primary type of $\sh{F}_s$.
Then the function $s\mapsto T(s)$ is locally constructible in $S$.

\medskip
Taking into account the preceding remarks, we are reduced as usual to proving that if $S$ is affine, Noetherian and integral, with generic point $\eta$, then $T(s)$ is constant in a neighbourhood of $\eta$.
Reasoning as at the beginning of \sref{IV.9.7.7}, we can suppose that all the irreducible components of $X_\eta$ that intervene are geometrically irreducible, and then the proposition follows from \sref{IV.9.5.1}, \sref{IV.9.5.5}, \sref{IV.9.8.3} and \sref{IV.9.8.6}; we leave the details to the reader.
One could generalise by considering several coherent Modules and defining their ``simultaneous primary skeleton'', etc.
The general conclusion of this paragraph is that for all the properties of the type considered (and for irreducible $S$) \emph{the properties valid on the ``generic fibre'' are also valid on all neighbouring fibres}.
\end{remark}

\subsection{Constructibility of local properties of fibres}
\label{subsection:IV.9.9}
\label{IV.9.9}

\begin{proposition}[9.9.1]
\label{IV.9.9.1}
Let $f:X\to S$ be a morphism locally of finite presentation, $Z$ a locally constructible subset of $X$ such that for all $s\in S$, $Z_s$ is closed in $X_s$, $\Phi$ a finite subset of $\mathbf{Z}\cup\{\pm\infty\}$.
Then the following subsets of $X$ are locally constructible:
\begin{enumerate}[label=(\roman*)]
  \item The set of $x\in X$ such that $\dim_x(Z_{f(x)})$ belongs to $\Phi$.
  \item The set of $x\in X$ such that $\operatorname{codim}_x(Z_{f(x)},X_{f(x)})$ belongs to $\Phi$.
  \item The set of $x\in X$ such that the local ring $\sh{O}_{X_{f(x)},x}$ is equidimensional.
\end{enumerate}

We note that properties (i) and (ii) can also be expressed by saying that the functions $x\mapsto\dim_x(Z_{f(x)})$ and $x\mapsto\operatorname{codim}_x(Z_{f(x)},X_{f(x)})$ are \emph{locally constructible} in $X$ ($0_{\mathrm{III}}$, \sref{IV.9.3.1}).
\end{proposition}

\begin{proof}
\label{IV.9.9.1-proof}
The questions being local on $X$, we can restrict to the case where $S=\operatorname{Spec}(A)$ and $X=\operatorname{Spec}(B)$ are affine and where $f$ is a morphism of finite presentation; there exists then a sub-ring $A_0$ of $A$ that is a $\mathbf{Z}$-algebra of finite type, an $A_0$-prescheme of finite type $X_0$ and a constructible subset $Z_0$ of $X_0$ such that $X=X_0\otimes_{A_0}A$ and $Z=h^{-1}(Z_0)$, where $h:X\to X_0$ is the canonical projection (\sref{IV.8.9.1} and \sref{IV.8.3.11}).
Furthermore, for all $s\in S$, if $s_0$ is the projection of $s$ in $S_0=\operatorname{Spec}(A_0)$, we have
\[
  X_s=(X_0)_{s_0}\otimes_{\mathbf{k}(s_0)}\mathbf{k}(s),
\]
and, if $h_s$ is the projection $X_s\to(X_0)_{s_0}$, we have $Z_s=h_s^{-1}((Z_0)_{s_0})$.
As the morphism $h_s$ is faithfully flat and quasi-compact, the hypothesis that $Z_s$ is closed in $X_s$ implies that $(Z_0)_{s_0}$ is closed in $(X_0)_{s_0}$ \sref{IV.2.3.12}.

This being so, the transitivity of fibres \sref[I]{I.3.6.4} and the prop.~\sref{IV.4.2.7} imply that the set of dimensions of the irreducible components of $Z_s$ containing $x$ is the same as the set of dimensions of the irreducible components of $(Z_0)_{s_0}$ containing $x_0=h_s(x)$.
In particular, we have $\dim_x(Z_s)=\dim_{x_0}((Z_0)_{s_0})$.
On the other hand, if $Z_s^{(\beta)}$ are the irreducible components of $Z_s$ containing $x$ and $X_s^{(\alpha)}$ the irreducible components of $X_s$
\oldpage[IV-3]{89}
containing $x$, we have
\[
  \operatorname{codim}_x(Z_s,X_s)
  =\inf_{\beta}\sup_{\alpha}\operatorname{codim}(Z_s^{(\beta)},X_s^{(\alpha)}),
\]
where $(\alpha,\beta)$ ranges over the pairs such that $x\in Z_s^{(\beta)}\subset X_s^{(\alpha)}$ ($0$, \sref{IV.14.2.6}).
As irreducible algebraic preschemes are biequidimensional \sref{IV.5.2.1}, we can write, by virtue of ($0$, \sref{IV.14.3.3.1})
\begin{equation}
\tag{9.9.1.1}
\label{IV.9.9.1.1}
  \operatorname{codim}_x(Z_s,X_s)=\inf_{\beta}\sup_{\alpha}\bigl(\dim(X_s^{(\alpha)})-\dim(Z_s^{(\beta)})\bigr)
\end{equation}
with the same choice of pairs $(\alpha,\beta)$.
As $h_s$ is faithfully flat and quasi-compact, for every pair formed of an irreducible component $(X_0)_{s_0}^{(\lambda)}$ of $(X_0)_{s_0}$ containing $x_0$ and an irreducible component $(Z_0)_{s_0}^{(\mu)}$ of $(Z_0)_{s_0}$ containing $x_0$ and contained in $(X_0)_{s_0}^{(\lambda)}$, there exists a pair $(X_s^{(\alpha)},Z_s^{(\beta)})$ of the type described above such that $Z_s^{(\beta)}$ dominates $(Z_0)_{s_0}^{(\mu)}$ and $X_s^{(\alpha)}$ dominates $(X_0)_{s_0}^{(\lambda)}$ \sref{IV.2.3.5}.
The formula \sref{IV.9.9.1.1} (and the analogous formula applied to $(X_0)_{s_0}$) shows then, by virtue of \sref{IV.4.2.7}, that we have
\[
  \operatorname{codim}_x(Z_s,X_s)=\operatorname{codim}_{x_0}((Z_0)_{s_0},(X_0)_{s_0}).
\]

We thus see that if $E$ (resp.~$E_0$) is the set of $x\in X$ (resp.~$x_0\in X_0$) satisfying one of the conditions (i), (ii), (iii) of the statement (resp.~the same condition), then $E=h^{-1}(E_0)$; by \sref{IV.1.8.2}, we may therefore restrict to the case where $A$, and hence also $B$, is Noetherian.
Taking account of ($0_{\mathrm{III}}$, \sref{IV.9.2.3}) and of \sref{IV.9.9.1.1}, it remains to show that for every $x\in X$ there is a neighbourhood $V$ of $x$ in $\overline{\{x\}}$ such that, for every $x'\in V$, the sets of dimensions of the irreducible components of $X_{f(x')}$ and of $Z_{f(x')}$ containing $x'$ are constant, and likewise the set of pairs
\[
  \bigl(\dim Z_{f(x')}^{(\beta)},\dim X_{f(x')}^{(\alpha)}\bigr),
\]
where $X_{f(x')}^{(\alpha)}$ is an irreducible component of $X_{f(x')}$ and $Z_{f(x')}^{(\beta)}$ is an irreducible component of $Z_{f(x')}$ contained in it and containing $x'$.
For this purpose we may replace $S$ by the reduced closed subprescheme $S'$ whose underlying space is $\overline{\{f(x)\}}$, and $X$ by $X'=f^{-1}(S')$, since the fibres of $X$ and $X'$ at points of $S'$ are the same.
Thus we may suppose that $S$ is integral and that $\eta=f(x)$ is its generic point.

By hypothesis $Z_\eta$ is closed in $X_\eta$; as $Z$ is constructible, it follows from \sref{IV.8.3.11}, applied by the method of \sref{IV.8.1.2},~(a), that after replacing $S$ by an open neighbourhood of $\eta$ we may suppose that $Z$ is \emph{closed} in $X$.
Let $X_i$ (resp.~$Z_j$) be the irreducible components of $X$ (resp.~$Z$) containing $x$.
By ($0_I$, \sref{IV.2.1.8}), the $X_i\cap X_\eta$ (resp.~$Z_j\cap X_\eta$) are the irreducible components of $X_\eta$ (resp.~$Z_\eta$) containing $x$; by \sref{IV.9.5.1}, after shrinking $S$ about $\eta$, we may further suppose that the $X_i$ (resp.~$Z_j$) are exactly the irreducible components of $X$ (resp.~$Z$) meeting $\overline{\{x\}}$, and that $(X_i)_s\cap\overline{\{x\}}$ and $(Z_j)_s\cap\overline{\{x\}}$ are nonempty for every $s\in S$.
This being so, it follows still from \sref{IV.9.7.1} that we can suppose, by restricting $S$, that the couples $(i,j)$ such that $(Z_j)_s\subset(X_i)_s$ are the same for all $s\in S$.
The conclusion follows then from \sref{IV.9.5.6}: for all $s$ sufficiently close to $\eta$, all the irreducible components of $(X_i)_s$ (resp.~$(Z_j)_s$) have the same dimension equal to that of $(X_i)_\eta$ (resp.~$(Z_j)_\eta$).
Furthermore, if $(i,j)$ is a pair such that $Z_j\not\subset X_i$, $(X_i)_\eta$ does not contain the
\oldpage[IV-3]{90}
generic point of $(Z_j)_\eta$, hence $\dim((X_i\cap Z_j)_\eta)<\dim((Z_j)_\eta)$; consequently, the common dimension of the irreducible components of $(X_i)_s\cap(Z_j)_s$ is, for all $s\in S$, strictly less than the common dimension of the irreducible components of $(Z_j)_s$, which proves that none of the irreducible components of $(Z_j)_s$ is contained in an irreducible component of $(X_i)_s$ for $s\in S$.
\end{proof}

\begin{proposition}[9.9.2]
\label{IV.9.9.2}
Let $f:X\to S$ be a morphism locally of finite presentation, $\sh{F}$ a quasi-coherent $\sh{O}_X$-Module of finite presentation, $\Phi$ a finite subset of $\mathbf{Z}\cup\{\pm\infty\}$.
The following subsets of $X$ are locally constructible:
\begin{enumerate}[label=(\roman*)]
  \item The set of $x\in X$ such that $\sh{F}_{f(x)}$ is equidimensional at the point $x$ \sref{IV.5.1.12}.
  \item The set of $x\in X$ such that the geometric lengths of $\sh{F}_{f(x)}$ relative to $\mathbf{k}(f(x))$, at the generic points of the irreducible components of $\operatorname{Supp}(\sh{F}_{f(x)})$ containing $x$ \sref{IV.4.7.5}, belong to $\Phi$.
  \item The set of $x\in X$ such that the dimensions of the prime cycles associated to $\sh{F}_{f(x)}$ and containing $x$ belong to $\Phi$.
  \item The set of $x\in X$ such that $\sh{F}_{f(x)}$ is geometrically reduced at the point $x$ \sref{IV.4.6.17}.
  \item The set of $x\in X$ such that $\sh{F}_{f(x)}$ is geometrically integral at the point $x$ \sref{IV.4.6.22}.
  \item The set of $x\in X$ such that $\dim.\operatorname{proj}((\sh{F}_{f(x)})_x)\in\Phi$.
  \item The set of $x\in X$ such that $\operatorname{coprof}((\sh{F}_{f(x)})_x)\in\Phi$.
  \item The set of $x\in X$ such that $\sh{F}_{f(x)}$ possesses property ($S_k$) at the point $x$ \sref{IV.5.7.2}.
  \item The set of $x\in X$ such that $\sh{F}_{f(x)}$ is a Cohen-Macaulay $\sh{O}_{X_{f(x)}}$-Module at the point $x$ \sref{IV.5.7.1}.
\end{enumerate}
\end{proposition}

\begin{proof}
\label{IV.9.9.2-proof}
(i) The support $Z$ of $\sh{F}$ is locally constructible and closed in $X$ \sref{IV.8.9.1}, and by considering a sub-prescheme of $X$ having $Z$ for underlying space and which is of finite presentation over $S$ \sref{IV.8.9.1}, we see that the assertion (i) is a particular case of \sref{IV.9.9.1},~(iii).

(ii) All the properties considered are local on $X$, and we may therefore suppose that $X=\operatorname{Spec}(B)$ and $S=\operatorname{Spec}(A)$ are affine and that $f$ is of finite presentation.
We retain the notation from the beginning of the proof of \sref{IV.9.9.1}, and choose $A_0$ so that there is a coherent $\sh{O}_{X_0}$-Module $\sh{F}_0$ with $\sh{F}\simeq\sh{F}_0\otimes_{A_0}A$.
Then \sref{IV.4.2.7} and \sref{IV.4.7.9} show that the set of geometric lengths of $(\sh{F}_0)_{f(x_0)}$ at the generic points of the irreducible components of its support containing $x_0$ is the same as the corresponding set for $\sh{F}_{f(x)}$ at $x$.
Thus, if $E$ (resp.~$E_0$) denotes the set of $x\in X$ (resp.~$x_0\in X_0$) satisfying condition (ii), then $E=h^{-1}(E_0)$, and by \sref{IV.1.8.2} we may restrict to the case where $A$ is Noetherian.
As in the proof of \sref{IV.9.9.1}, it remains to show that for every $x\in X$ there is a neighbourhood $V$ of $x$ in $\overline{\{x\}}$ on which the set of geometric lengths of $\sh{F}_{f(x')}$ at the generic points of the irreducible components of its support containing $x'$ is constant.
If $S'$ is the reduced closed subprescheme of $S$ with underlying space $\overline{\{f(x)\}}$, and if $X'=f^{-1}(S')$, the fibres of $X$ and $X'$ at the points of $S'$ are the same,
\oldpage[IV-3]{91}
so we may replace $S$ by $S'$ and $X$ by $X'$, in other words suppose that $S$ is integral and that $\eta=f(x)$ is its generic point.
If $Z=\operatorname{Supp}(\sh{F})$ and $Z_i$ are the irreducible components of $Z$ containing $x$, the argument of \sref{IV.9.9.1} lets us suppose that these are exactly the irreducible components of $Z$ meeting $\overline{\{x\}}$ and that $(Z_i)_s\cap\overline{\{x\}}\neq\varnothing$ for every $s\in S$.
The conclusion follows from \sref{IV.9.8.3} and \sref{IV.9.8.6}.

(iii) Using \sref{IV.4.2.7}, we reduce as in (ii) to the case where $S$ is Noetherian and integral and $\eta=f(x)$ is its generic point.
As in the proof of \sref{IV.9.9.1}, it remains to find a neighbourhood $V$ of $x$ in $\overline{\{x\}}$ on which the set of dimensions of the prime cycles associated to $\sh{F}_{f(x')}$ and containing $x'$ is constant.
Let $Z'_i$ be the closures in $X$ of the prime cycles associated to $\sh{F}_\eta$ which contain $x$ (cf.~\sref{IV.9.8.2}).
By \sref{IV.9.8.3} and \sref{IV.9.5.1}, for every $s$ sufficiently near $\eta$, the prime cycles associated to $\sh{F}_s$ which meet $\overline{\{x\}}$ are irreducible components of the $(Z'_i)_s$; by \sref{IV.9.5.6}, all these components have dimension $\dim((Z'_i)_\eta)$, which proves the assertion.

(iv) We reason as in (iii), now using \sref{IV.4.7.11}.
With the same notation, the prime cycles associated to $\sh{F}_s$ which are irreducible components of $(Z'_i)_s$ are embedded if and only if $(Z'_i)_\eta$ is an embedded associated prime cycle of $\sh{F}_\eta$.
It follows that, according as $x$ belongs to an embedded associated prime cycle of $\sh{F}_\eta$ or to none, the points $x'\in\overline{\{x\}}$ having the analogous property for $\sh{F}_{f(x')}$ form a neighbourhood of $x$ in $\overline{\{x\}}$.
The conclusion follows from this observation, the characterization of geometrically reduced Modules in \sref{IV.4.7.10}, and (ii).

(v) Taking \sref{IV.4.7.11} into account, we reduce again to the case where $S$ is Noetherian and integral and where $\eta=f(x)$ is its generic point.
Let $Y$ be a closed subprescheme of $X$ whose underlying space is $\operatorname{Supp}(\sh{F})$.
As at the beginning of the proof of \sref{IV.9.7.7}, there is a finite integral $A$-algebra $A'$ (so that, if $S'=\operatorname{Spec}(A')$, the morphism $S'\to S$ is finite and surjective) such that, if $Y'=Y_{(S')}$ and $\eta'$ is the generic point of $S'$, the irreducible components of $Y'_{\eta'}$ are geometrically irreducible.
If $X'=X_{(S')}$, the projection $h:X'\to X$ is finite and surjective, hence closed; consequently, for a point $x'$ above $x$, one has $h(\overline{\{x'\}})=\overline{\{x\}}$, and the image under $h$ of a neighbourhood of $x'$ in $\overline{\{x'\}}$ is a neighbourhood of $x$ in $\overline{\{x\}}$.
Writing $\sh{F}'=\sh{F}_{(S')}$ and using \sref{IV.4.7.11}, we are reduced to proving that, according as $\sh{F}'_{\eta'}$ is or is not geometrically integral at $x'$, the points $y'\in\overline{\{x'\}}$ where $\sh{F}'_{f'(y')}$ (where $f'=f_{(S')}:X'\to S'$) is or is not geometrically integral at $y'$ form a neighbourhood of $x'$ in $\overline{\{x'\}}$.
We may therefore suppose that $X'=X$ and that the irreducible components of $Y_\eta$ are geometrically irreducible.
Choose closed subsets $Z_i$ of $X$ such that the $Z_i\cap X_\eta$ are the
  \oldpage[IV-3]{92}
  irreducible components of $Y_\eta$.
By \sref{IV.9.7.7}, \sref{IV.9.7.8}, and \sref{IV.9.5.3}, for every $s$ near $\eta$ the $(Z_i)_s$ are the irreducible components of $Y_s$ and are geometrically irreducible.
At a point $y\in X_s$, the Module $\sh{F}_s$ is geometrically integral exactly when it is geometrically reduced there and $y$ belongs to only one of the $(Z_i)_s$ \sref{IV.4.6.22}.
The conclusion follows from (iv) and from \sref{IV.9.5.1} applied to $\overline{\{x\}}\cap Z_i$ for each $i$.

(vi) Keeping the notations of (ii), \sref{IV.6.2.1} gives
\[
  \dim.\operatorname{proj}((\sh{F}_s)_x)
  =\dim.\operatorname{proj}(((\sh{F}_0)_{s_0})_{x_0}),
\]
so we may again restrict to the case where $A$ is Noetherian. It remains to show that for every $x\in X$ there is a neighbourhood $V$ of $x$ in $\overline{\{x\}}$ such that
\[
  \dim.\operatorname{proj}((\sh{F}_{f(x')})_{x'})
  =\dim.\operatorname{proj}((\sh{F}_{f(x)})_x)
\]
for every $x'\in V$. As above, we may suppose that $S$ is integral and that $\eta=f(x)$ is its generic point, so $(\sh{F}_\eta)_x=\sh{F}_x$. By generic flatness \sref{IV.6.9.1}, after replacing $S$ by an open neighbourhood of $\eta$, we may suppose that $f$ is flat and that $\sh{F}$ is $f$-flat. Then \sref{IV.6.2.3} gives
\[
  \dim.\operatorname{proj}((\sh{F}_{f(x')})_{x'})
  =\dim.\operatorname{proj}(\sh{F}_{x'})
\]
for every $x'\in X$. By \sref{IV.6.11.1}, after replacing $X$ by an open neighbourhood of $x$, we may suppose that
\[
  \dim.\operatorname{proj}(\sh{F}_{x'})
  \leqslant\dim.\operatorname{proj}(\sh{F}_x)
\]
for every $x'\in X$.

If $\dim.\operatorname{proj}(\sh{F}_x)=n$, there is a finite $\sh{O}_x$-module $M$ such that $\operatorname{Ext}^n_{\sh{O}_x}(\sh{F}_x,M)\neq 0$ \sref[0]{0.17.2.4}.
After replacing $X$ by an open neighbourhood of $x$, there exists a coherent $\sh{O}_X$-module $\sh{G}$ such that $M=\sh{G}_x$ \sref[0]{0.5.3.8};
by virtue of (T, \sref{IV.4.2.2}), we have $(\shExt^n_{\sh{O}_X}(\sh{F},\sh{G}))_x\neq 0$.
But $\shExt^n_{\sh{O}_X}(\sh{F},\sh{G})$ is a coherent $\sh{O}_X$-Module ($0_{\mathrm{III}}$, \sref{IV.12.3.3}), hence its support $Y$ is closed ($0_I$, \sref{IV.5.2.2}); as it contains $x$, it contains also $\overline{\{x\}}$, whence we conclude (by applying again (T, \sref{IV.4.2.2})) that $\dim.\operatorname{proj}(\sh{F}_{x'})\geqslant n$ for all $x'\in\overline{\{x\}}$, which achieves the proof of the assertion in case (vi).

(vii) As $B$ is an $A$-algebra of finite type, $X$ is $S$-isomorphic to a \emph{closed sub-scheme} of an $S$-scheme of the form $Y=\operatorname{Spec}(A[T_1,\ldots,T_r])$; if $j:X\to Y$ is the canonical injection and $\sh{G}=j_*(\sh{F})$, then $\sh{G}_s=(j_s)_*(\sh{F}_s)$ for all $s\in S$, and, by virtue of ($0$, \sref{IV.16.4.11}), we can restrict to proving the assertion for $Y$ and $\sh{G}$.
In other words, we can suppose that $B=A[T_1,\ldots,T_r]$, so that the schemes $X_s=\operatorname{Spec}(\mathbf{k}(s)[T_1,\ldots,T_r])$ are \emph{regular} ($0$, \sref{IV.17.3.7}).
Let $W=\operatorname{Supp}(\sh{F})$ and $W_s=\operatorname{Supp}(\sh{F}_s)$ \sref[I]{I.9.1.13}; we have, by \sref{IV.6.11.2.1}
\begin{equation}
\tag{9.9.2.1}
\label{IV.9.9.2.1}
  \operatorname{coprof}((\sh{F}_{f(x)})_x)=\dim.\operatorname{proj}((\sh{F}_{f(x)})_x)-\operatorname{codim}_x(W_{f(x)},X_{f(x)}).
\end{equation}

But as $W$ is constructible \sref{IV.8.9.1} and each $W_s$ is closed, it follows from (vi) and \sref{IV.9.9.1},~(ii), that the two functions on the right-hand side of \sref{IV.9.9.2.1} are constructible; hence so is their difference, which proves the assertion in case (vii).

(viii) Let $U_n$ be the set of $x\in X$ such that $\operatorname{coprof}((\sh{F}_{f(x)})_x)\leqslant n$, and set $Z_n=X-U_n$; it follows from (vii) that the $Z_n$ are \emph{constructible}.
\oldpage[IV-3]{93}
Moreover, since the function $x\mapsto\dim_x(W_{f(x)})$ is constructible by \sref{IV.9.9.1},~(i), it takes only finitely many values, and the numbers $\dim(W_{f(x)})$ therefore have a finite upper bound $m$ as $x$ ranges over $X$.
Since $\operatorname{coprof}((\sh{F}_{f(x)})_x)\leqslant\dim((\sh{F}_{f(x)})_x)\leqslant\dim(W_{f(x)})$, we have $Z_n=\varnothing$ for $n\geqslant m$.
Finally, \sref{IV.6.11.2},~(i), shows that $(Z_n)_s$ is closed in $X_s$ for every $n$ and every $s\in S$.
By \sref{IV.5.7.4}, the set of $x\in X$ where $(\sh{F}_{f(x)})_x$ possesses property ($S_k$) is the set of $x\in X$ satisfying all the relations
\begin{equation}
\tag{9.9.2.2}
\label{IV.9.9.2.2}
  \operatorname{codim}_x((Z_n)_{f(x)},W_{f(x)})>n+k
\end{equation}
for all $n\geqslant 0$; as this relation is automatically satisfied for $n\geqslant m$, we need only consider the relations \sref{IV.9.9.2.2} for $0\leqslant n<m$.
But by virtue of \sref{IV.9.9.1},~(ii), the set $V_{n,k}$ of $x$ satisfying \sref{IV.9.9.2.2} is constructible, and so is the intersection of these sets for $0\leqslant n<m$.

(ix) The assertion follows immediately from (vii), the set of $x\in X$ where $(\sh{F}_{f(x)})_x$ is a Cohen-Macaulay module being defined by the relation $\operatorname{coprof}((\sh{F}_{f(x)})_x)=0$.
\end{proof}

\begin{corollary}[9.9.3]
\label{IV.9.9.3}
Let $f:X\to S$ be a morphism of finite presentation, $\sh{F}$ a quasi-coherent $\sh{O}_X$-Module of finite presentation, $\mathbf{P}$ one of the properties (i) to (ix) of \sref{IV.9.9.2}.
Then the set of $s\in S$ such that the property $\mathbf{P}$ is true at all the points $x\in X_s$, is locally constructible in $S$.
\end{corollary}

\begin{proof}
\label{IV.9.9.3-proof}
Indeed, its complement is the image by $f$ of the complement of the set $E$ of points where $\mathbf{P}$ is true.
As $E$ is locally constructible in $X$, so is $X-E$, and hence $f(X-E)$ is locally constructible in $S$ by virtue of the theorem of Chevalley \sref{IV.1.8.4}.
\end{proof}

\begin{proposition}[9.9.4]
\label{IV.9.9.4}
Let $f:X\to S$ be a morphism locally of finite presentation.
The set of $x\in X$ such that $X_{f(x)}$ has at the point $x$ one of the (determined) following properties:
\begin{enumerate}[label=(\roman*)]
  \item being geometrically regular;
  \item possessing property ($R_k$) geometric;
  \item being geometrically normal;
  \item being geometrically reduced (i.e.\ separable);
  \item being geometrically pointwise integral;
\end{enumerate}
is a locally constructible subset of $X$.
\end{proposition}

\begin{proof}
\label{IV.9.9.4-proof}
For properties (iv) and (v), this follows from \sref{IV.9.9.2},~(iv) and~(v) applied to $\sh{F}=\sh{O}_X$.
For the other properties, taking into account \sref{IV.6.7.8}, we reduce, as at the beginning of \sref{IV.9.9.2},~(ii), to the case where $S$ is Noetherian and integral and of generic point $\eta=f(x)$.
Furthermore, by virtue of the generic flatness theorem \sref{IV.6.9.1}, we can, by replacing $S$ by an open neighbourhood of $\eta$, suppose that $f$ is a \emph{flat} morphism.
To say that $X_{f(y)}$ is geometrically regular at the point $y$ means that $f$ is \emph{regular} at the point $y$ \sref{IV.6.8.1}, and we know that the set $L$ of these points is \emph{open} in $X$ \sref{IV.6.8.7}, which proves the proposition in case (i).

To prove case (ii), set $F=X-L$, which is \emph{closed} in $X$.
To say that $X_s$
\oldpage[IV-3]{94}
satisfies the geometric property ($R_k$) at the point $y\in f^{-1}(s)$ means either that $y\in L$, or that the generic points $z_i$ of the irreducible components of the closed set $F_s$ which contain $y$ satisfy the relation $\dim(\sh{O}_{X_s,z_i})\geqslant k+1$ (taking into account \sref{IV.4.2.7} and \sref{IV.5.2.3}); in other words, the points $y\in F_s$ where $X_s$ satisfies property ($R_k$) geometric are those such that $\operatorname{codim}_y(F_s,X_s)\geqslant k+1$ \sref{IV.5.1.2}.
The conclusion follows thus from (i) and \sref{IV.9.9.1},~(ii).

Finally, in case (iii), the conclusion follows from (ii), from \sref{IV.9.9.2},~(viii), from the fact that property ($S_k$) is stable by extension of the base field \sref{IV.6.7.1} and finally from the criterion of Serre \sref{IV.5.8.6}.
\end{proof}

\begin{corollary}[9.9.5]
\label{IV.9.9.5}
Let $f:X\to S$ be a morphism of finite presentation, and denote by $\mathbf{P}$ one of the properties (i) to (v) of \sref{IV.9.9.4}.
Then the set of $s\in S$ such that the property $\mathbf{P}$ is true at all the points $x\in X_s$, is locally constructible in $S$.
\end{corollary}

\begin{proof}
\label{IV.9.9.5-proof}
The proof from \sref{IV.9.9.4} is the same as that of \sref{IV.9.9.3} from \sref{IV.9.9.2}.
\end{proof}

\begin{proposition}[9.9.6]
\label{IV.9.9.6}
Let $f:X\to S$ be a morphism locally of finite presentation, $\sh{L}_\bullet$ a complex formed of $\sh{O}_X$-Modules quasi-coherent of finite presentation; for all $s\in S$, let $(\sh{L}_\bullet)_s$ be the complex $\sh{L}_\bullet\otimes_{\sh{O}_S}\mathbf{k}(s)$ of $\sh{O}_{X_s}$-Modules coherent of finite type.
Then, for a given integer $n$, the set of $x\in X$ such that $(\sh{H}_n((\sh{L}_\bullet)_{f(x)}))_x=0$ is locally constructible in $X$.
\end{proposition}

\begin{proof}
\label{IV.9.9.6-proof}
We can restrict to the case where $\sh{L}_i=0$ except for $i=0$, $1$ or $2$, and where $n=1$.
Furthermore, the question being local on $X$, we can restrict to the case where $S=\operatorname{Spec}(A)$ and $X=\operatorname{Spec}(B)$ are affine, $B$ being a finitely presented $A$-algebra.
There exists then a Noetherian sub-ring $A_0$ of $A$, a finite type $A_0$-prescheme $X_0$ and a complex $\sh{L}_\bullet^{(0)}$ of $\sh{O}_{X_0}$-Modules coherent, zero except in dimensions $0$, $1$ and $2$ and such that $X=X_0\otimes_{A_0}A$ and $\sh{L}_\bullet=\sh{L}_\bullet^{(0)}\otimes_{A_0}A$.
For every $s\in S$, if $s_0$ is the projection of $s$ in $S_0=\operatorname{Spec}(A_0)$, then
\[
  X_s=(X_0)_{s_0}\otimes_{\mathbf{k}(s_0)}\mathbf{k}(s),
\]
and the projection $X_s\to(X_0)_{s_0}$ is faithfully flat.
Consequently,
\[
  \sh{H}_n((\sh{L}_\bullet)_s)
  =\sh{H}_n((\sh{L}_\bullet^{(0)})_{s_0})\otimes_{\mathbf{k}(s_0)}\mathbf{k}(s).
\]
If $E$ (resp.~$E_0$) is the set of $x\in X$ (resp.~$x_0\in X_0$) for which
\[
  \begin{aligned}
  (\sh{H}_n((\sh{L}_\bullet)_{f(x)}))_x&=0\\
  \bigl(\text{resp. }(\sh{H}_n((\sh{L}_\bullet^{(0)})_{f_0(x_0)}))_{x_0}&=0\bigr),
  \end{aligned}
\]
then $E=h^{-1}(E_0)$, where $h:X\to X_0$ is the canonical projection.
By \sref{IV.1.8.2}, we may therefore restrict to the case where $A$ is Noetherian.
It remains to show ($0_{\mathrm{III}}$, \sref{IV.9.2.3}) that if $x\in X$ is such that $(\sh{H}_n((\sh{L}_\bullet)_{f(x)}))_x=0$ (resp.~$\neq 0$), there is an open neighbourhood $V$ of $x$ in $\overline{\{x\}}$ such that for every $x'\in V$ one has $(\sh{H}_n((\sh{L}_\bullet)_{f(x')}))_{x'}=0$ (resp.~$\neq 0$).
Replacing $S$ by the reduced closed subprescheme having $\overline{\{f(x)\}}$ as its underlying space, we may suppose that $S$ is integral and that $f(x)=\eta$ is its generic point.
Then, by restricting $S$ to an open neighbourhood of $\eta$, we may suppose that for every $s\in S$,
$(\sh{H}_n(\sh{L}_\bullet))_s=\sh{H}_n((\sh{L}_\bullet)_s)$ \sref{IV.9.4.3}; consequently, if $Z$ is the support of $\sh{H}_n(\sh{L}_\bullet)$, then the support of $\sh{H}_n((\sh{L}_\bullet)_s)$ is $Z_s=Z\cap X_s$ \sref[I]{I.9.1.13.1}.
The hypothesis is that $Z_\eta\cap\overline{\{x\}}=\varnothing$ (resp.~$\neq\varnothing$).
As $Z\cap\overline{\{x\}}$ is closed in the Noetherian space $X$, we conclude from \sref{IV.9.5.1} that there is a neighbourhood $U$ of $\eta$ in $S$ such that, for all $s\in U$, we have $Z_s\cap\overline{\{x\}}=\varnothing$ (resp.~$\neq\varnothing$).
The set $V=f^{-1}(U)\cap\overline{\{x\}}$ thus answers the question.
\end{proof}
% END NATIVE INPUT source_aligned/ega4-9
% BEGIN NATIVE INPUT source_aligned/ega4-10
\oldpage[IV-3]{95}
\section{Jacobson preschemes}
\label{section:IV.10}

We have already had occasion to observe \sref{IV.5.2.5} that even
\emph{excellent} preschemes \sref{IV.7.8.5} do not always behave like the
``varieties'' of classical Algebraic Geometry, especially with respect to
questions of dimension; thus, if $X$ is the spectrum of a complete discrete
valuation ring, the set of closed points of $X$ (reduced to a single element)
is not everywhere dense in $X$, while its complement (reduced to a single
element) is an everywhere dense open subset of $X$ whose dimension is zero,
hence $<\dim(X)$. In this section we examine general conditions under which
such phenomena do not occur; as a result, the relations between dimension,
codimension, depth, and codepth in these preschemes behave more satisfactorily
in certain respects (10.6 and 10.8). Moreover, in the preschemes considered,
the fact that the set of closed points is everywhere dense (and even ``very
dense'' \sref{IV.10.1.3}) makes it possible to consider only these points in
many proofs; we thus rejoin the classical point of view of ``algebraic
varieties'', which, from our point of view, are the sets of closed points of
\emph{algebraic preschemes over a field}, and establish the link between the
language of schemes and that of Serre's ``varieties'' or ``algebraic spaces''
(10.9 and 10.10).

\subsection{Very dense subsets of a topological space}
\label{subsection:IV.10.1}

\begin{env}[10.1.1]
\label{IV.10.1.1}
We say that a subset of a topological space $X$ is \emph{quasi-constructible} if it is a finite union of locally closed subsets of $X$.
We say that a subset $T$ of $X$ is \emph{locally quasi-constructible} if for every $x\in X$, there exists an open neighbourhood $V$ of $x$ such that $T\cap V$ is quasi-constructible in $V$.
It is clear that every quasi-constructible subset of $X$ is locally quasi-constructible; the converse is true if $X$ is quasi-compact.
The argument of \sref[$0_{\mathrm{III}}$]{0.9.1.3}, where one removes the word ``retrocompact'', shows that the set of quasi-constructible (resp.\ locally quasi-constructible) subsets of $X$, which we will denote by $\mathfrak{Qc}(X)$ (resp.\ $\mathfrak{Lqc}(X)$), is stable under finite union, finite intersection, and passage to complements.
If $f:X\to Y$ is a continuous map, it follows immediately from these definitions that, for every quasi-constructible (resp.\ locally quasi-constructible) subset $Z$ of $Y$, $f^{-1}(Z)$ is quasi-constructible (resp.\ locally quasi-constructible) in $X$.

The constructible (resp.\ locally constructible) subsets of $X$ are evidently quasi-constructible (resp.\ locally quasi-constructible); the converse is true when $X$ is Noetherian (resp.\ locally Noetherian).

In what follows, we will denote respectively by $\mathfrak{O}(X)$, $\mathfrak{F}(X)$, $\mathfrak{Lf}(X)$, $\mathfrak{C}(X)$, $\mathfrak{Lc}(X)$ the set of subsets of $X$ that are respectively open, closed, locally closed, constructible, locally constructible.
\end{env}

\oldpage[IV-3]{96}

\begin{proposition}[10.1.2]
\label{IV.10.1.2}
Let $X$ be a topological space, $X_0$ a subset of $X$.
The following conditions are equivalent:
\begin{enumerate}[label=\emph{\alph*)}]
  \item For every locally closed subset $Z\neq\varnothing$ of $X$, we have $Z\cap X_0\neq\varnothing$.
  \item[\rm a$'$)] For every closed subset $Z$ of $X$, we have $Z=\overline{Z\cap X_0}$.
  \item For every locally quasi-constructible subset $Z\neq\varnothing$ of $X$, we have $Z\cap X_0\neq\varnothing$.
  \item[\rm b$'$)] For every locally quasi-constructible subset $Z$ of $X$, we have $Z\subset\overline{Z\cap X_0}$ (in other words, $Z\cap X_0$ is dense in $Z$).
  \item The map $U\mapsto U\cap X_0$ of $\mathfrak{O}(X)$ into $\mathfrak{O}(X_0)$ is injective (hence bijective).
  \item[\rm c$'$)] The map $F\mapsto F\cap X_0$ of $\mathfrak{F}(X)$ into $\mathfrak{F}(X_0)$ is injective (hence bijective).
  \item[\rm c$''$)] The map $Z\mapsto Z\cap X_0$ of $\mathfrak{Qc}(X)$ into $\mathfrak{Qc}(X_0)$ is injective (hence bijective).
  \item[\rm c$'''$)] The map $Z\mapsto Z\cap X_0$ of $\mathfrak{Lqc}(X)$ into $\mathfrak{Lqc}(X_0)$ is injective.
\end{enumerate}
Furthermore, when these conditions are satisfied, the map $Z\mapsto Z\cap X_0$ of $\mathfrak{Lqc}(X)$ into $\mathfrak{Lqc}(X_0)$ is bijective.
\end{proposition}

\begin{proof}
Let us note that the assertions of surjectivity of $c)$ and $c')$ are trivial; they mean that every locally closed subset of $X_0$ is the trace on $X_0$ of a locally closed subset of $X$, and that the map $\mathfrak{Qc}(X)\to\mathfrak{Qc}(X_0)$ defined in $c'')$ is also surjective.

We will prove the implications
\[
  c''')\Rightarrow c'')\Rightarrow c')\Rightarrow c)\Rightarrow b)\Rightarrow b')\Rightarrow a')\Rightarrow a)\Rightarrow c''').
\]

The first three are trivial.
To see that $c)$ implies $b)$, let us first note that $c)$ means that $X_0$ is \emph{dense} in $X$.
Replacing $X$ and $X_0$ by $U$ and $U\cap X_0$ respectively, one can therefore restrict to the case where $Z$ is locally closed in $X$, hence $Z=V\cap\complement W$, where $V$ and $W$ are open in $X$; the hypothesis $Z\neq\varnothing$ means that $V\not\subset W$, or equivalently $V\cup W\neq W$.
By virtue of $c)$, we have $(V\cup W)\cap X_0\neq W\cap X_0$, hence $V\cap X_0\not\subset W$, and consequently $(V\cap\complement W)\cap X_0\neq\varnothing$.

To see that $b)$ implies $b')$, it suffices to apply $b)$ to $Z\cap U$, where $U$ is an open neighbourhood of an arbitrary point of $Z$.
As $\overline{Z\cap X_0}\subset\overline{Z}$, it is trivial that $b')$ implies $a')$.
To show that $a')$ implies $a)$, let us note that if $Z$ is locally closed in $X$, we can write $Z=F-F'$ where $F$ and $F'$ are closed in $X$ and $F'\subset F$; whence $Z\cap X_0=(F\cap X_0)-(F'\cap X_0)$.
If $Z\cap X_0=\varnothing$, then, by virtue of $a')$, we would have $F\cap X_0=F'\cap X_0$, hence $F=F'$, i.e.\ $Z=\varnothing$.

To see that $a)$ implies $c''')$, it suffices to show that if $Z\neq\varnothing$ is locally quasi-constructible, then $Z\cap X_0\neq\varnothing$: indeed, the relation $Z'\cap X_0=Z''\cap X_0$ is equivalent to $((Z'\cup Z'')-(Z'\cap Z''))\cap X_0=\varnothing$.
In other words, it suffices to prove that $a)$ implies $b)$; indeed, replacing $X$ and $X_0$ by $U$ and $U\cap X_0$ respectively, one is reduced to the case where $Z$ is locally closed in $X$, whence the conclusion.

It remains to show that the map $\mathfrak{Lqc}(X)\to\mathfrak{Lqc}(X_0)$ is surjective.
So let $Z_0$ be a locally quasi-constructible subset of $X_0$: there is a covering $(V_\alpha)$ of $X_0$ by open subsets of $X_0$ such that $Z_0\cap V_\alpha$ is quasi-constructible in $V_\alpha$ (and
\oldpage[IV-3]{97}
hence also in $X_0$).
By virtue of $c)$, for every $\alpha$ there is a unique open subset $U_\alpha$ of $X$ such that $X_0\cap U_\alpha=V_\alpha$, and by virtue of $c'')$ a unique quasi-constructible subset $Z_\alpha$ of $X$ such that $Z_0\cap V_\alpha=X_0\cap Z_\alpha$.
If $\alpha$ and $\beta$ are any two indices, we have $Z_\beta\cap U_\alpha\cap X_0=U_\alpha\cap Z_0\cap V_\beta=V_\alpha\cap Z_0\cap V_\beta=Z_\alpha\cap X_0\cap V_\beta\subset Z_\alpha\cap X_0$; as $Z_\beta\cap U_\alpha$ and $Z_\alpha$ are quasi-constructible in $X$, it follows from $c'')$ that $Z_\beta\cap U_\alpha\subset Z_\alpha$.
If we set $Z=\bigcup_\alpha Z_\alpha$, then $Z\cap U_\alpha=Z_\alpha$ for every $\alpha$; moreover, since the $V_\alpha$ cover $X_0$, it follows from $c)$ that the $U_\alpha$ cover $X$, and we see therefore that $Z$ is locally quasi-constructible in $X$ and $Z_0=Z\cap X_0$.
\end{proof}

\begin{definition}[10.1.3]
\label{IV.10.1.3}
When a subset $X_0$ of a topological space satisfies the equivalent conditions of \sref{IV.10.1.2}, we say that $X_0$ is \emph{very dense} in $X$.
\end{definition}

We have already seen in the course of the proof of \sref{IV.10.1.2} that $X_0$ is then dense in $X$.

\begin{corollary}[10.1.4]
\label{IV.10.1.4}
If $X_0$ is very dense in $X$ and $U$ an open subset of $X$, $U\cap X_0$ is very dense in $U$.
Conversely, if $(U_\alpha)$ is an open covering of $X$ such that $U_\alpha\cap X_0$ is very dense in $U_\alpha$ for every $\alpha$, then $X_0$ is very dense in $X$.
\end{corollary}

As every locally closed subset of $U$ is locally closed in $X$, the first assertion follows from criterion $a)$ of \sref{IV.10.1.2}; and the same is true of the second, since if $Z\neq\varnothing$ is locally closed in $X$, $Z\cap U_\alpha$ is locally closed in $U_\alpha$ for every $\alpha$, and $Z\cap U_\alpha\neq\varnothing$ for at least one $\alpha$.

\subsection{Quasi-homeomorphisms}
\label{subsection:IV.10.2}

\begin{proposition}[10.2.1]
\label{IV.10.2.1}
Let $X_0$, $X$ be two topological spaces, $f:X_0\to X$ a continuous map.
The following conditions are equivalent:
\begin{enumerate}[label=\emph{\alph*)}]
  \item The map $U\mapsto f^{-1}(U)$ of $\mathfrak{O}(X)$ into $\mathfrak{O}(X_0)$ is bijective.
  \item[\rm a$'$)] The map $F\mapsto f^{-1}(F)$ of $\mathfrak{F}(X)$ into $\mathfrak{F}(X_0)$ is bijective.
  \item The topology of $X_0$ is the inverse image by $f$ of that of $X$, and $f(X_0)$ is very dense \sref{IV.10.1.3} in $X$.
  \item The functor $\sh{F}\mapsto f^*(\sh{F})$ from the category of sheaves of sets (resp.\ sheaves of abelian groups) on $X$ to the category of sheaves of sets (resp.\ sheaves of abelian groups) on $X_0$ is an equivalence of categories.
\end{enumerate}
\end{proposition}

\begin{proof}
It is clear that $a)$ and $a')$ are equivalent, and $a)$ implies that the topology of $X_0$ is the inverse image by $f$ of that of $X$.
On the other hand, if $f(X_0)$ is not very dense in $X$, there are two distinct open subsets $U_1$, $U_2$ of $X$ such that $U_1\cap f(X_0)=U_2\cap f(X_0)$, and hence $f^{-1}(U_1)=f^{-1}(U_2)$, which completes showing that $a)$ implies $b)$.
Conversely, the condition $b)$ implies that the maps $U\mapsto U\cap f(X_0)$ of $\mathfrak{O}(X)$ into $\mathfrak{O}(f(X_0))$ and $V\mapsto f^{-1}(V)$ of $\mathfrak{O}(f(X_0))$ into $\mathfrak{O}(X_0)$ are bijective, hence so is the composite $U\mapsto f^{-1}(U)$.

To see that $a)$ implies $c)$, it suffices to apply the definition of $f^*(\sh{F})$ \sref[0]{0.3.7.1} and the axioms for sheaves: $a)$ means that for every open $U$ of $X$, the canonical map $\Gamma(U,\sh{F})\to\Gamma(f^{-1}(U),f^*(\sh{F}))$ is a functorial bijection in $\sh{F}$, whence $c)$.
It remains to show that $c)$ implies $b)$.

\oldpage[IV-3]{98}
Suppose first that $f(X_0)$ is not very dense in $X$; there then exist two distinct closed subsets $Y\supset Y'$ of $X$ such that $Y\cap f(X_0)=Y'\cap f(X_0)$.
Let $\sh{F}$ (resp.\ $\sh{F}'$) be the sheaf of abelian groups on $X$ that is the direct image by the canonical injection $Y\to X$ (resp.\ $Y'\to X$) of the simple sheaf associated to the constant presheaf $\mathbf{Z}$ on $Y$ (resp.\ $Y'$); the definition of $f^*$ \sref[0]{0.3.7.1} shows that $f^*(\sh{F})$ is isomorphic to $f^*(\sh{F}')$, but $\sh{F}$ is not isomorphic to $\sh{F}'$, hence condition $c)$ is not satisfied.
(One notes that the functor $f^*$ is not then even \emph{faithful}, for if $u$ and $v$ are the identity automorphism and the zero endomorphism of $\sh{F}/\sh{F}'$, $f^*(u)$ and $f^*(v)$ are equal.)

Let us now show that $c)$ implies that the topology of $X_0$ is the inverse image of that of $X$ by $f$.
Let us note that if condition $c)$ is satisfied for the category of sheaves of sets, it is also satisfied for the category of sheaves of abelian groups, for it follows immediately from the definitions \sref[$0_{\mathrm{III}}$]{0.8.2.5} that the category of \emph{objects in abelian groups} in the category of sheaves of sets is none other than the category of sheaves of abelian groups.
It will therefore suffice to prove our assertion supposing $c)$ satisfied for the categories of sheaves of abelian groups on $X$ and $X_0$.
Now, if $\mathbf{Z}_X$ denotes the simple sheaf on $X$ associated to the constant presheaf $\mathbf{Z}$, there is a canonical isomorphism $\Gamma(X,\sh{F})\xrightarrow{\sim}\operatorname{Hom}(\mathbf{Z}_X,\sh{F})$ functorial in $\sh{F}$ (by the same argument as in \sref[0]{0.5.1.1}); as it is clear by definition \sref[0]{0.3.7.1} that $f^*(\mathbf{Z}_X)=\mathbf{Z}_{X_0}$, the hypothesis that $f^*$ is an equivalence means that the canonical homomorphism $\Gamma(X,\sh{F})\to\Gamma(X_0,f^*(\sh{F}))$ is \emph{bijective} for every sheaf of abelian groups $\sh{F}$ on $X$.
Now, let $Y_0$ be a closed subset of $X_0$; let $\sh{F}_0$ be the sheaf on $X_0$ that is the direct image by the canonical injection $Y_0\to X_0$ of the simple sheaf associated to the constant presheaf $\mathbf{Z}$ on $Y_0$.
As $f^*$ is an equivalence, $\sh{F}_0$ is isomorphic to a sheaf of the form $f^*(\sh{F})$, where $\sh{F}$ is a sheaf of abelian groups on $X$.
Consider then a section $s_0$ of $\sh{F}_0$ over $X_0$ such that $(s_0)_{x_0}=1_{x_0}$ if $x_0\in Y_0$, $(s_0)_{x_0}=0_{x_0}$ if $x_0\notin Y_0$.
The preceding remarks show that there exists a section $s$ of $\sh{F}$ over $X$ such that $(s_0)_{x_0}=s_{f(x_0)}$ for every $x_0\in X_0$ (the fibres $(\sh{F}_0)_{x_0}$ and $\sh{F}_{f(x_0)}$ being canonically identified \sref[0]{0.3.7.1}); this means that $Y_0$ is the inverse image by $f$ of the set $Y$ of $x\in X$ such that $s_x\neq 0_x$, and $Y$ is closed in $X$.
\end{proof}

\begin{definition}[10.2.2]
\label{IV.10.2.2}
When a continuous map $f:X_0\to X$ satisfies the equivalent conditions of \sref{IV.10.2.1}, we say that $f$ is a \emph{quasi-homeomorphism} of $X_0$ into $X$.
\end{definition}

By virtue of \sref{IV.10.2.1},~$b)$, to say that a subset $X_0$ of a topological space $X$ is very dense in $X$ means that the canonical injection $X_0\to X$ is a quasi-homeomorphism.

\begin{corollary}[10.2.3]
\label{IV.10.2.3}
The composite of two quasi-homeomorphisms is a quasi-homeomorphism.
\end{corollary}

This follows immediately from \sref{IV.10.2.1},~$a)$.

\begin{corollary}[10.2.4]
\label{IV.10.2.4}
Let $f:X\to Y$ be a quasi-homeomorphism, $Y'$ a locally quasi-constructible subset of $Y$, $X'=f^{-1}(Y')$; then the restriction $f'=f|X':X'\to Y'$ is a quasi-homeomorphism.
\end{corollary}

\begin{proof}
It is clear by virtue of \sref{IV.10.2.1},~$b)$ that the topology induced on $X'$ by that of $X$ is the inverse image by $f'$ of the topology induced on $Y'$ by that of $Y$.
On the other hand,
\oldpage[IV-3]{99}
let $Z\neq\varnothing$ be a closed subset \emph{in} $Y'$, and $z\in Z$; there is an open neighbourhood $U$ of $z$ in $Y$ such that $U\cap Y'$ is the union of a finite number of closed subsets $Y'_j$ of $U$; if $j$ is an index such that $z\in Y'_j$, $Z\cap Y'_j$ is therefore closed in $U$.
Since $f(X)\cap U$ is very dense in $U$ \sref{IV.10.1.4}, $Z\cap Y'_j\cap f(X)$ is not empty \sref{IV.10.1.2},~$a)$, nor, \emph{a fortiori}, is $Z\cap f(X)$; but as $Z\subset Y'$, we have $Z\cap f(X)=Z\cap f'(X')$; the criterion \sref{IV.10.1.2},~$a)$ shows that $f'(X')$ is very dense in $Y'$, and we conclude by \sref{IV.10.2.1},~$b)$.
\end{proof}

\begin{corollary}[10.2.5]
\label{IV.10.2.5}
Let $f:X\to Y$ be a continuous map, $(V_\alpha)$ an open covering of $Y$.
If, for every $\alpha$, the restriction $f^{-1}(V_\alpha)\to V_\alpha$ of $f$ is a quasi-homeomorphism, then $f$ is a quasi-homeomorphism.
\end{corollary}

This follows immediately from the criterion \sref{IV.10.2.1},~$b)$ and from \sref{IV.10.1.4}.

\begin{corollary}[10.2.6]
\label{IV.10.2.6}
Let $f:X\to Y$ be a quasi-homeomorphism, $Y'$ a locally quasi-constructible subset of $Y$, $X'=f^{-1}(Y')$.
For $Y'$ to be quasi-compact (resp.\ Noetherian, resp.\ retrocompact in $Y$), it is necessary and sufficient that $X'$ be quasi-compact (resp.\ Noetherian, resp.\ retrocompact in $X$).
\end{corollary}

\begin{proof}
Let us prove the first two assertions; by virtue of \sref{IV.10.2.4}, one can restrict to the case where $Y'=Y$.
To say that $X$ is quasi-compact (resp.\ Noetherian) means that for every increasing filtered family $(U_\alpha)$ in $\mathfrak{O}(X)$ having $X$ as greatest element (resp.\ for every increasing filtered family $(U_\alpha)$ in $\mathfrak{O}(X)$), there exists $\gamma$ such that $U_\alpha=U_\gamma$ for $\alpha\geq\gamma$.
As $U\mapsto f^{-1}(U)$ is a bijection of $\mathfrak{O}(Y)$ onto $\mathfrak{O}(X)$, our assertion follows immediately from the preceding remark.

The quasi-compact open subsets of $X$ are therefore the sets of the form $f^{-1}(U)$ where $U$ is an open quasi-compact subset of $Y$, by virtue of \sref{IV.10.2.1},~$a)$ and of what precedes.
For $X'$ to be retrocompact in $X$, it is necessary and sufficient that for every quasi-compact open $U$ in $Y$, $f^{-1}(U)\cap X'=f^{-1}(U\cap Y')$ be quasi-compact \sref[$0_{\mathrm{III}}$]{0.9.1.1}; the first part of the proof shows that this amounts to saying that $U\cap Y'$ is quasi-compact for every open quasi-compact $U$, i.e.\ that $Y'$ is retrocompact in $Y$.
\end{proof}

\begin{proposition}[10.2.7]
\label{IV.10.2.7}
Let $f:X\to Y$ be a quasi-homeomorphism.
Then the map $Z\mapsto f^{-1}(Z)$ of $\mathfrak{P}(Y)$ into $\mathfrak{P}(X)$ defines by restriction the bijections (cf.\ \sref{IV.10.1.1} for the notation):
\[
  \mathfrak{O}(Y)\xrightarrow{\sim}\mathfrak{O}(X)
\]
\[
  \mathfrak{F}(Y)\xrightarrow{\sim}\mathfrak{F}(X)
\]
\[
  \mathfrak{Lf}(Y)\xrightarrow{\sim}\mathfrak{Lf}(X)
\]
\[
  \mathfrak{Qc}(Y)\xrightarrow{\sim}\mathfrak{Qc}(X)
\]
\[
  \mathfrak{Lqc}(Y)\xrightarrow{\sim}\mathfrak{Lqc}(X)
\]
\[
  \mathfrak{C}(Y)\xrightarrow{\sim}\mathfrak{C}(X)
\]
\[
  \mathfrak{Lc}(Y)\xrightarrow{\sim}\mathfrak{Lc}(X).
\]
\end{proposition}

\begin{proof}
\oldpage[IV-3]{100}
For the first two, this is none other than the definition \sref{IV.10.2.2}; as the topology of $X$ is the inverse image by $f$ of that of $f(X)$, the five remaining maps, where one replaces $Y$ by $f(X)$, are bijective.
It therefore suffices (by \sref{IV.10.2.1},~$b)$) to restrict to the case where $X$ is a very dense subspace of $Y$, and that the maps $\mathfrak{Lf}(Y)\to\mathfrak{Lf}(X)$, $\mathfrak{Qc}(Y)\to\mathfrak{Qc}(X)$, $\mathfrak{Lqc}(Y)\to\mathfrak{Lqc}(X)$ are bijective has already been proved \sref{IV.10.1.2}.
Every locally constructible subset being locally quasi-constructible, the maps $\mathfrak{C}(Y)\to\mathfrak{C}(X)$ and $\mathfrak{Lc}(Y)\to\mathfrak{Lc}(X)$ are injective; furthermore, for every open $U\subset Y$, the restriction $f^{-1}(U)\to U$ of $f$ is a quasi-homeomorphism \sref{IV.10.2.4}, hence if we show that $\mathfrak{C}(Y)\to\mathfrak{C}(X)$ is surjective, the same will be true of $\mathfrak{Lc}(Y)\to\mathfrak{Lc}(X)$.
But by virtue of \sref{IV.10.2.6}, every retrocompact open subset $Z$ in $X$ is of the form $f^{-1}(T)$, where $T$ is retrocompact open in $Y$; this evidently proves the surjectivity of $\mathfrak{C}(Y)\to\mathfrak{C}(X)$.
\end{proof}

\begin{remarks}[10.2.8]
\label{IV.10.2.8}
(i) We saw in the proof of \sref{IV.10.2.1} that if $f$ is a quasi-homeomorphism, the canonical map
\begin{equation}
\label{IV.10.2.8.1}
\tag{10.2.8.1}
\Gamma(Y,\sh{F})\to\Gamma(X,f^*(\sh{F}))
\end{equation}
is a functorial isomorphism in $\sh{F}$ in the category of sheaves of abelian groups on $Y$.
As $f^*$ is exact in this category \sref[0]{0.3.7.2}, this implies that the canonical homomorphisms $(T,3.2.2)$
\[
  \mathrm{H}^i(Y,\sh{F})\xrightarrow{\sim}\mathrm{H}^i(X,f^*(\sh{F}))
\]
are \emph{bijective} for every $i$.

(ii) If $f:X\to Y$ is a continuous map and $\sh{F}$, $\sh{G}$ two sheaves of abelian groups on $Y$, we have $f^*(\sh{F}\otimes_{\mathbf{Z}_Y}\sh{G})=f^*(\sh{F})\otimes_{\mathbf{Z}_X}f^*(\sh{G})$, up to a canonical isomorphism \sref[0]{0.4.3.3}.
We conclude that if $f$ is a quasi-homeomorphism, $f^*$ is also an \emph{equivalence} of the category of sheaves of rings on $Y$, and of the category of sheaves of rings on $X$; the data of a structure of ringed space on $Y$ is therefore equivalent to that of a structure of ringed space on $X$.

Given two ringed spaces $(X,\sh{O}_X)$, $(Y,\sh{O}_Y)$, we will say that a morphism of ringed spaces $f=(\psi,\theta):(X,\sh{O}_X)\to(Y,\sh{O}_Y)$ is a \emph{quasi-isomorphism} if $\psi$ is a quasi-homeomorphism of $X$ into $Y$ and $\theta^\sharp:\psi^*(\sh{O}_Y)\to\sh{O}_X$ an isomorphism of sheaves of rings; when this is so, the ringed space $(X,\sh{O}_X)$ is entirely determined, up to isomorphism, by $(Y,\sh{O}_Y)$, the space $X$, and the quasi-homeomorphism $\psi$ (which can be chosen arbitrarily).
When $f$ is a quasi-isomorphism of ringed spaces, the functor
\[
  \sh{F}\mapsto f^*(\sh{F})
\]
is an \emph{equivalence} of the category of $\sh{O}_Y$-Modules and of that of $\sh{O}_X$-Modules, since $f^*(\sh{F})$ is canonically identified here with $\psi^*(\sh{F})$.
We conclude for example isomorphisms of bi-$\partial$-functors
\[
  \operatorname{Ext}^i_{\sh{O}_Y}(\sh{F},\sh{G})\xrightarrow{\sim}\operatorname{Ext}^i_{\sh{O}_X}(f^*(\sh{F}),f^*(\sh{G})).
\]

In general, one can say that the usual constructions of the theory of sheaves and of homological algebra, carried out on the ringed space $Y$ or the ringed space $X$, are equivalent.
\end{remarks}

\subsection{Jacobson spaces}
\label{subsection:IV.10.3}

\oldpage[IV-3]{101}

\begin{definition}[10.3.1]
\label{IV.10.3.1}
We say that a topological space $X$ is a \emph{Jacobson space} if the set $X_0$ of closed points of $X$ is very dense in $X$ (in other words, if the canonical injection $X_0\to X$ is a quasi-homeomorphism).
\end{definition}

This means therefore \sref{IV.10.1.2} that every locally closed (or only locally quasi-constructible) subset $Z\neq\varnothing$ of $X$ contains a closed point of $X$, or equivalently that \emph{every closed subset of $X$ is the closure of the set of its closed points}.

\begin{proposition}[10.3.2]
\label{IV.10.3.2}
Let $X$ be a Jacobson space, $Z$ a locally quasi-constructible subset of $X$; then the subspace $Z$ is a Jacobson space, and for a point of $Z$ to be closed in $Z$, it is necessary and sufficient that it be closed in $X$.
\end{proposition}

\begin{proof}
If $X_0$ is the set of closed points of $X$, $Z\cap X_0$ is very dense in $Z$ by virtue of \sref{IV.10.2.4} applied to the injection $i:X_0\to X$; as the set $Z_0$ of closed points of $Z$ evidently contains $Z\cap X_0$, it is \emph{a fortiori} very dense in $Z$, hence $Z$ is a Jacobson space.
Let us now show that we have $Z_0=Z\cap X_0$; let $x\in Z$ be a closed point in $Z$; let $\overline{\{x\}}$ be its closure in $X$ and $\overline{\{x\}}\cap Z=\{x\}$, which is therefore a locally quasi-constructible subset of $X$, and as its intersection with $X_0$ is not empty \sref{IV.10.1.2}, we have $x\in X_0$.
\end{proof}

\begin{proposition}[10.3.3]
\label{IV.10.3.3}
Let $X$ be a topological space, $(U_\alpha)$ an open covering of $X$.
For $X$ to be a Jacobson space, it is necessary and sufficient that each of the subspaces $U_\alpha$ be a Jacobson space.
\end{proposition}

\begin{proof}
The condition is necessary by virtue of \sref{IV.10.3.2}.
Conversely, let us show first that the hypothesis that the $U_\alpha$ are Jacobson spaces implies that for a point $x\in U_\alpha$ to be closed in $X$, it suffices that it be closed in $U_\alpha$.
It suffices indeed to see that this condition implies that $x$ is also closed in each of the $U_\beta$ that contain it; but $U_\alpha\cap U_\beta$ is open in $U_\alpha$, hence $x$ is closed in $U_\alpha\cap U_\beta$, and by virtue of \sref{IV.10.3.2}, $x$ is also closed in $U_\beta$, which completes the proof.
\end{proof}

\subsection{Jacobson preschemes and Jacobson rings}
\label{subsection:IV.10.4}

\begin{definition}[10.4.1]
\label{IV.10.4.1}
We say that a prescheme $X$ is a \emph{Jacobson prescheme} if the underlying topological space of $X$ is a Jacobson space.
We say that a ring $A$ is a \emph{Jacobson ring} if $\operatorname{Spec}(A)$ is a Jacobson prescheme.
\end{definition}

Every closed subset of $X=\operatorname{Spec}(A)$ is of the form $Z=V(\mathfrak{a})$, where $\mathfrak{a}$ is an ideal equal to its radical, and the set $X_0$ of closed points of $X$ is the set of maximal ideals of $A$; to say that $Z\cap X_0$ is dense in $Z$ means therefore that $\mathfrak{a}$ is an \emph{intersection of maximal ideals} \sref[I]{I.1.1.4}; as $\mathfrak{a}$ is an intersection of prime ideals, this amounts to saying that every prime ideal of $A$ is an intersection of maximal ideals; by virtue of \sref{IV.10.3.1} and \sref{IV.10.1.2}, the usual definition of Jacobson rings (Bourbaki, \emph{Alg.\ comm.}, chap.~V, \S\,3, n$^\circ$\,4, def.~1) thus coincides with the definition \sref{IV.10.4.1}.

\begin{proposition}[10.4.2]
\label{IV.10.4.2}
\oldpage[IV-3]{102}
Let $X$ be a prescheme, $(U_\alpha)$ a covering of $X$ formed by affine open subsets.
For $X$ to be a Jacobson prescheme, it is necessary and sufficient that the rings $\Gamma(U_\alpha,\sh{O}_X)$ of the $U_\alpha$ be Jacobson rings.
\end{proposition}

This follows from \sref{IV.10.3.3} and the definition \sref{IV.10.4.1}.

\begin{env}[10.4.3]
\label{IV.10.4.3}
A discrete prescheme is a Jacobson prescheme; an Artinian ring is therefore a Jacobson ring.
Every principal ideal ring having an infinity of maximal ideals (for example $\mathbf{Z}$) is a Jacobson ring; a Noetherian local ring $A$ is a Jacobson ring if and only if its maximal ideal is its only prime ideal, i.e.\ if $A$ is Artinian.
Every sub-prescheme of a Jacobson prescheme is a Jacobson prescheme by virtue of \sref{IV.10.3.2}.
\end{env}

\begin{proposition}[10.4.4]
\label{IV.10.4.4}
Let $B$ be an integral ring.
The following conditions are equivalent:
\begin{enumerate}[label=\emph{\alph*)}]
  \item There exists $f\neq 0$ in $B$ such that $B_f$ is a field.
  \item The field of fractions of $B$ is a $B$-algebra of finite type.
  \item There exists a field $K$ containing $B$ and which is a $B$-algebra of finite type.
  \item The generic point of $\operatorname{Spec}(B)$ is isolated.
\end{enumerate}
\end{proposition}

\begin{proof}
It is clear that $d)$ is equivalent to $a)$, since $d)$ means that there exists $f\neq 0$ in $B$ such that $D(f)$ is reduced to the generic point of $\operatorname{Spec}(B)$.
It is trivial that $a)$ implies $b)$ and that $b)$ implies $c)$.
Finally, $c)$ implies $a)$, by virtue of Bourbaki, \emph{Alg.\ comm.}, chap.~V, \S\,3, n$^\circ$\,1, cor.~2 of th.~1.
\end{proof}

\begin{proposition}[10.4.5]
\label{IV.10.4.5}
Let $A$ be a ring.
The following conditions are equivalent:
\begin{enumerate}[label=\emph{\alph*)}]
  \item $A$ is a Jacobson ring.
  \item For every non-maximal prime ideal $\mathfrak{p}$ of $A$ and every $f\neq 0$ in $B=A/\mathfrak{p}$, $B_f$ is not a field.
  \item[\rm b$'$)] Every $A$-algebra of finite type $K$ which is a field is a finite $A$-algebra.
\end{enumerate}
\end{proposition}

\begin{proof}
We know that $a)$ implies $b')$ (Bourbaki, \emph{Alg.\ comm.}, chap.~V, \S\,3, n$^\circ$\,4, cor.~3 of th.~3).
Furthermore, the kernel of the homomorphism $A\to K$ is then a \emph{maximal} ideal $\mathfrak{m}$ of $A$, and $K$ is a finite extension of $A/\mathfrak{m}$ (\emph{loc.\ cit.}).
It is trivial that $b')$ implies $b)$, since $A/\mathfrak{p}=B$ is not a field and $B_f$ is a $B$-algebra of finite type.
It remains to show that $b)$ implies $a)$.
We will use the following lemma:
\end{proof}

\begin{lemma}[10.4.5.1]
\label{IV.10.4.5.1}
Let $X$ be a topological space having the following property: for every locally closed subset $Z\neq\varnothing$ of $X$, there exists a subset $Z'$ of $Z$, locally closed in $X$ (or in $Z$, which amounts to the same), and a point $x\in Z'$, closed in $Z'$.
Then, for $X$ to be a Jacobson space, it is necessary and sufficient that no non-closed point $x$ of $X$ be isolated in $\overline{\{x\}}$.
\end{lemma}

\begin{proof}
If $X$ is a Jacobson space, a non-closed point $x$ of $X$ cannot be isolated in $\overline{\{x\}}$, for this would mean that there exists an open $U$ of $X$ such that $U\cap\overline{\{x\}}=\{x\}$; but $U\cap\overline{\{x\}}$ is locally closed in $X$ and would contain no closed point in $X$, which contradicts the hypothesis that $X$ is a Jacobson space.
Conversely, suppose the conditions of the statement are satisfied; then the set $X_0$ of closed points of $X$ is identical to the set of $x\in X$ that are isolated in $\overline{\{x\}}$; but it follows from \sref{IV.5.1.10.1} that
\oldpage[IV-3]{103}
this set is \emph{very dense} in $X$, hence $X$ is a Jacobson space by definition.
\end{proof}

Let us recall \sref{IV.5.1.10} that the hypothesis made on $X$ in \sref{IV.10.4.5.1} is always satisfied when the underlying space of $X$ is a \emph{prescheme}.

Returning to the proof of \sref{IV.10.4.5}, condition $b)$ means that for every non-closed point $x$ of $\operatorname{Spec}(A)$, the generic point $x$ of $\operatorname{Spec}(A/\mathfrak{j}_x)=V(\mathfrak{j}_x)=\overline{\{x\}}$ is not isolated in $\overline{\{x\}}$, hence $b)$ implies $a)$ by the lemma \sref{IV.10.4.5.1} and \sref{IV.10.4.4}.

\begin{corollary}[10.4.6]
\label{IV.10.4.6}
Every algebra of finite type $B$ over a Jacobson ring $A$ is a Jacobson ring, and the inverse image in $A$ of every maximal ideal of $B$ is a maximal ideal of $A$.
In particular, every algebra of finite type over a field or over $\mathbf{Z}$ is a Jacobson ring.
\end{corollary}

\begin{proof}
A $B$-algebra of finite type $K$ which is a field is also an $A$-algebra of finite type, hence an $A$-module of finite type and \emph{a fortiori} a $B$-module of finite type, whence the first assertion; the second was proved in the course of the proof of \sref{IV.10.4.5}, applied to $K=B/\mathfrak{m}$.
\end{proof}

\begin{corollary}[10.4.7]
\label{IV.10.4.7}
If $X$ is a Jacobson prescheme, $f:Y\to X$ a morphism locally of finite type, then $Y$ is a Jacobson prescheme and the image by $f$ of every closed point in $Y$ is a closed point in $X$.
\end{corollary}

\begin{proof}
The question being local on $X$ and on $Y$, one reduces to the case where $X$ and $Y$ are affine, and the corollary then follows from \sref{IV.10.4.6}.
\end{proof}

\begin{corollary}[10.4.8]
\label{IV.10.4.8}
If $k$ is an algebraically closed field, $X$ a $k$-prescheme locally of finite type, the set of rational points of $X$ over $k$ is very dense in $X$.
\end{corollary}

\begin{proof}
Indeed, $X$ is a Jacobson prescheme \sref{IV.10.4.7} and the closed points are exactly the rational points over $k$ \sref[I]{I.6.4.2}.
\end{proof}

\begin{env}[10.4.9]
\label{IV.10.4.9}
The fact that preschemes locally of finite type over a field or over $\mathbf{Z}$ are Jacobson preschemes is particularly important, because of the possibility of reducing to this case in many questions of algebraic geometry \sref{IV.8.1.2},~$c)$.
We will give two examples:
\end{env}

\begin{env}[10.4.10]
\label{IV.10.4.10}
\textbf{Applications (10.4.10) :~I~: Proof of \sref{IV.6.15.9}.}
Let $k$ be a separably closed field, $X$ a $k$-prescheme locally of finite type over $k$ and \emph{unibranch}.
We know that the integral closure of a $k$-algebra of finite type $A$ that is an integral domain in a finite extension of its field of fractions is a finite $A$-algebra (Bourbaki, \emph{Alg.\ comm.}, chap.~V, \S\,3, n$^\circ$\,2, th.~2), hence every $k$-algebra of finite type is a universally Japanese ring; we conclude that the set of points $x\in X$ where $X$ is geometrically unibranch is locally constructible \sref{IV.9.7.10}.
But the hypothesis and the lemma \sref{IV.6.15.8} imply that this set contains all the \emph{closed} points of $X$.
The conclusion therefore follows from \sref{IV.10.4.6}, \sref{IV.10.3.1} and the bijectivity of the canonical map $\mathfrak{Lc}(X)\to\mathfrak{Lc}(X_0)$ (where $X_0$ is the set of closed points of $X$) \sref{IV.10.2.7}.
\end{env}

\begin{env}[10.4.11]
\label{IV.10.4.11}
\oldpage[IV-3]{104}
\textbf{Applications :~II~: Proposition (10.4.11).}
Let $S$ be a prescheme, $X$ an $S$-prescheme of finite type.
Every $S$-endomorphism of $X$ which is radicial is surjective (hence bijective).
\end{env}

\begin{proof}
Let $f:X\to S$ be the structural morphism, $g:X\to X$ the $S$-endomorphism considered, and, for every $s\in S$, let $g_s$ be the morphism deduced from $g$ by the base change $\operatorname{Spec}(k(s))\to S$, which is a $k(s)$-endomorphism of the fibre $f^{-1}(s)=X\times_S\operatorname{Spec}(k(s))$.

To prove that $g$ is surjective, it suffices to prove that $g_s$ is surjective for every $s\in S$, hence one can (by virtue of \sref[I]{I.3.5.7}) reduce to the case where $S=\operatorname{Spec}(k)$ is the spectrum of a field $k$, in which case $f$ is a morphism of finite presentation, since $S$ is Noetherian.
Applying \sref{IV.8.9.1} and \sref{IV.8.10.5},~(vii), one reduces to the case where $S=\operatorname{Spec}(A)$, where $A$ is a sub-$\mathbf{Z}$-algebra of finite type of $k$.
But then $X$ is a Jacobson prescheme \sref{IV.10.4.7}, and $g(X)$ is constructible in $X$ \sref[I]{I.1.8.5}, hence, to prove that $g(X)=X$, it suffices to show that $g(X)$ contains all the \emph{closed} points of $X$ \sref{IV.10.3.1}.

\begin{lemma}[10.4.11.1]
\label{IV.10.4.11.1}
Let $Y$ be a $\mathbf{Z}$-prescheme of finite type.
\begin{enumerate}[label=(\roman*)]
  \item For a point $y\in Y$ to be closed in $Y$, it is necessary and sufficient that $k(y)$ be a finite field.
  \item For every prime number $p$ and every integer $d\geq 1$, the set of points $y\in Y$ such that $k(y)$ is an extension of $\mathbf{F}_p$ of degree dividing $d$ is finite.
\end{enumerate}
\end{lemma}

The assertion (i) follows from the fact that the image of a closed point $y\in Y$ in $\operatorname{Spec}(\mathbf{Z})$ is a closed point \sref{IV.10.4.7}, in other words a prime number, and from \sref[I]{I.6.4.2}.
On the other hand, as $Y$ is a finite union of affine open subsets of finite type over $\mathbf{Z}$, one can restrict to proving (ii) in the case where $Y=\operatorname{Spec}(C)$, where $C$ is a $\mathbf{Z}$-algebra of finite type.
Now, the points $y\in Y$ such that the degree of $k(y)$ over $\mathbf{F}_p$ divides $d$ are precisely the kernels of the homomorphisms $C\to\mathbf{F}_{p^d}$; but if $(t_i)$ is a system of generators of the $\mathbf{Z}$-algebra $C$, every homomorphism of $C$ is determined by its values on the elements $t_i$, and there are therefore only a finite number of homomorphisms of $C$ into a finite field.

This being so, since $X$ is a $\mathbf{Z}$-prescheme of finite type, let $T_{p,d}$ be the set of closed points $z\in X$ such that $k(z)$ is an extension of $\mathbf{F}_p$ of degree dividing $d$; it follows from \sref{IV.10.4.11.1} that the set $T_{p,d}$ is finite and that the set of closed points of $X$ is the union of the $T_{p,d}$.
Furthermore, if $z\in T_{p,d}$ and if $h$ is any endomorphism of $X$, $k(h(z))$ is isomorphic to a subfield of $k(z)$, hence $h(z)\in T_{p,d}$; in other words $T_{p,d}$ is \emph{stable} under every endomorphism of $X$.
As $g$ is by hypothesis injective, its restriction to $T_{p,d}$ is \emph{bijective} since $T_{p,d}$ is finite, which completes proving the proposition.

We will see further on \sref{IV.17.9.6} that when one further supposes that $X$ is an $S$-prescheme of finite presentation, and on the other hand that $g$ is a monomorphism, then one can assert that $g$ is an \emph{automorphism} of $X$.
\end{proof}

\subsection{Noetherian Jacobson preschemes}
\label{subsection:IV.10.5}

\begin{proposition}[10.5.1]
\label{IV.10.5.1}
Let $B$ be a Noetherian integral ring.
The conditions equivalent to $a)$ through $d)$ of \sref{IV.10.4.4} are then also equivalent to the following:
\begin{enumerate}[label=\emph{\alph*)}]
  \setcounter{enumi}{4}
  \item $\operatorname{Spec}(B)$ is finite.
  \item $B$ is a semi-local ring of dimension $\leq 1$.
\end{enumerate}
\end{proposition}

\begin{proof}
\oldpage[IV-3]{105}
It follows indeed from the Artin--Tate theorem \sref[0]{0.16.3.3} that conditions $a)$ and $f)$ are equivalent.
Condition $f)$ implies that $\operatorname{Spec}(B)$ is the union of the set of its closed points and of its generic point, hence $f)$ implies $e)$ without assuming $B$ to be Noetherian.
Finally, $e)$ implies $d)$ without assuming $B$ to be Noetherian, for the generic point $x$ of $X$ is the complement of the union of the closures $\overline{\{y\}}$, where $y$ ranges over the set of points $y\neq x$, and as these points are finite in number, $\{x\}$ is the complement of a closed set in $X$.
\end{proof}

\begin{corollary}[10.5.2]
\label{IV.10.5.2}
For a Noetherian ring $A$ to be a Jacobson ring, it is necessary and sufficient that there exist no prime ideal $\mathfrak{p}$ of $A$ such that $A/\mathfrak{p}$ is a semi-local ring of dimension $1$.
\end{corollary}

\begin{proof}
This follows immediately from \sref{IV.10.5.1} and from condition $b)$ of \sref{IV.10.4.5}, the prime ideals $\mathfrak{p}$ of $A$ such that $A/\mathfrak{p}$ is semi-local of dimension~$0$ being the maximal ideals of $A$.
\end{proof}

\begin{corollary}[10.5.3]
\label{IV.10.5.3}
Let $X$ be an irreducible Noetherian prescheme.
The following conditions are equivalent:
\begin{enumerate}[label=\emph{\alph*)}]
  \item The generic point of $X$ is isolated.
  \item $X$ is finite.
  \item $X$ is of dimension $\leq 1$ and the set of its closed points is finite.
\end{enumerate}
\end{corollary}

\begin{proof}
There exists by hypothesis a finite number of irreducible affine open subsets $U_i$ ($1\leq i\leq n$) covering $X$, and each of which contains the generic point of $X$; it suffices to prove the equivalence of $a)$, $b)$, and $c)$ for each of the $(U_i)_{\mathrm{red}}$ (taking into account \sref[0]{0.14.1.7}).
But this equivalence follows from \sref{IV.10.5.1}.
\end{proof}

\begin{remark}[10.5.4]
\label{IV.10.5.4}
A Noetherian prescheme $X$ satisfying the equivalent conditions of \sref{IV.10.5.3} is not necessarily an affine scheme; in fact it can even be \emph{non-separated}, as one sees by replacing, in the example \sref[I]{I.5.5.11} of the ``affine line with doubled point'', $X_1$ and $X_2$ by the spectrum of the discrete valuation ring $(K[s])_{(s)}$, $U_{12}$ and $U_{21}$ by the open subset reduced to the generic point in $X_1$ and $X_2$ respectively; the non-separated prescheme $X$ that one obtains has exactly $3$ points.
\end{remark}

\begin{proposition}[10.5.5]
\label{IV.10.5.5}
Let $X$ be a Noetherian prescheme.
\begin{enumerate}[label=(\roman*)]
  \item The set $X_0$ of points $x\in X$ such that $\overline{\{x\}}$ is finite is very dense in $X$.
  \item For $X$ to be a Jacobson prescheme, it is necessary and sufficient that there exist no sub-prescheme of $X$ isomorphic to the spectrum of a semi-local integral ring of dimension $1$.
\end{enumerate}
\end{proposition}

\begin{proof}
The condition that $\overline{\{x\}}$ is finite is equivalent here, by \sref{IV.10.5.3}, to the fact that $x$ is isolated in $\overline{\{x\}}$, it being the underlying space of a sub-prescheme (Noetherian) of $X$; the assertion (i) therefore follows from \sref{IV.5.1.10.1}.
Similarly, taking into account \sref{IV.10.4.5.1}, to show the assertion (ii), let us note that for a non-closed point $x$ of $X$ to belong to $X_0$, it is necessary and sufficient that the closed integral sub-prescheme $Y$ of $X$ having $\overline{\{x\}}$ as underlying space be of dimension $1$ and finite, hence a finite union of sub-preschemes which are open (in $Y$) and affine $U_i$ which are spectra of semi-local integral rings of dimension $1$.
Conversely, if there is a sub-prescheme $Z$ of $X$ which is the spectrum of a semi-local integral ring of dimension $1$, $Z$ is not a Jacobson prescheme \sref{IV.10.5.2}, hence neither is $X$ \sref{IV.10.3.2}.
\end{proof}

\oldpage[IV-3]{106}

\begin{remark}[10.5.6]
\label{IV.10.5.6}
The assertion (ii) of \sref{IV.10.5.5} is still valid when $X$ is \emph{locally Noetherian}: indeed, if $(V_\alpha)$ is a covering of $X$ by affine open subsets (Noetherian), every sub-prescheme of one of the $V_\alpha$ is a sub-prescheme of $X$; conversely, if a sub-prescheme $Z$ of $X$ is isomorphic to the spectrum of a semi-local integral ring of dimension $1$, there is an $\alpha$ such that $V_\alpha\cap Z$ contains an affine open $U$ of $Z$ which is non-reduced to the generic point of $Z$, and which is therefore also the spectrum of a semi-local integral ring of dimension $1$.
We conclude by \sref{IV.10.3.3}.
\end{remark}

\begin{proposition}[10.5.7]
\label{IV.10.5.7}
Let $X$ be a locally Noetherian prescheme, $Y$ a closed subset of $X$ such that every non-empty closed subset of $X$ meets $Y$.
Then the subprescheme induced on the open subset $X-Y$ is a Jacobson prescheme.
\end{proposition}

\begin{proof}
Apply the criterion \sref{IV.10.5.5},~(ii), and suppose that there is in $X-Y$ a sub-prescheme $Z$ which is the spectrum of a semi-local integral ring of dimension $1$. The generic point $z$ of $Z$ is then isolated in $Z$ (or in the closure $\overline{Z}$ of $Z$ in $X$), and $Z$ is distinct from $\{z\}$.
Let $y\neq z$ be a point of $Z$; as it does not belong to $Y$, it is not a closed point in $X$, and its closure $\overline{\{y\}}$ in $X$ meets $Y$ at a point $x\neq y$, which is therefore a specialisation of $y$.
The existence of the chain $\{x\}\subset\overline{\{y\}}\subset\overline{\{z\}}$ then shows that the dimension of $\overline{\{z\}}$ would be $\geq 2$, and the same would be true of the dimension of $U\cap\overline{\{z\}}$, where $U$ is an affine (hence Noetherian) open neighbourhood of $x$ in $X$.
But this contradicts the fact that the generic point $z$ of $U\cap\overline{\{z\}}$ is isolated in $U\cap\overline{\{z\}}$ \sref{IV.10.5.3}.
\end{proof}

\begin{corollary}[10.5.8]
\label{IV.10.5.8}
Let $A$ be a Noetherian ring; for every element $f$ of the radical $\mathfrak{R}$ of $A$, the ring $A_f$ is a Jacobson ring, and the open subset $\operatorname{Spec}(A)-V(\mathfrak{R})$ is a Jacobson scheme.
\end{corollary}

\begin{proof}
If $X=\operatorname{Spec}(A)$, $Y=V(\mathfrak{J})$, where $\mathfrak{J}$ is an ideal of $A$, to say that every non-empty closed subset of $X$ meets $Y$ is equivalent \sref[0]{0.2.1.3} to saying that $Y$ contains all the \emph{closed points} of $X$, or equivalently that $\mathfrak{J}$ is \emph{contained in the radical} $\mathfrak{R}$ of $A$.
If $f\in\mathfrak{R}$, the open $D(f)$ does not meet $V(\mathfrak{R})$, and is a Jacobson space by virtue of \sref{IV.10.5.7}.
\end{proof}

\begin{corollary}[10.5.9]
\label{IV.10.5.9}
Let $A$ be a Noetherian local ring, $\mathfrak{m}$ its maximal ideal, $X=\operatorname{Spec}(A)$, $Y=X-\{\mathfrak{m}\}$; then $Y$ is a Jacobson scheme, whose closed points are the prime ideals $\mathfrak{p}\in\operatorname{Spec}(A)$ such that $\dim(A/\mathfrak{p})=1$.
\end{corollary}

\begin{proof}
The first assertion is a particular case of \sref{IV.10.5.8}; on the other hand the closed points of $Y$ are the prime ideals $\mathfrak{p}$ of $A$ which are maximal elements in the set of prime ideals $\neq\mathfrak{m}$, which, by definition of the dimension, means that $\dim(A/\mathfrak{p})=1$.
\end{proof}

\begin{proposition}[10.5.10]
\label{IV.10.5.10}
Let $A$ be a Noetherian local ring, reduced, complete, and not a field.
Then the finite intersections of the kernels of the local homomorphisms of $A$ into discrete valuation rings $V$, which make $V$ a finite $A$-algebra, form a filter base tending to $0$ for the adic topology of $A$.
\end{proposition}

\begin{proof}
It suffices (Bourbaki, \emph{Alg.\ comm.}, chap.~III, \S\,2, n$^\circ$\,7, prop.~8) to prove that the intersection of the kernels considered in the statement is reduced to $0$.
Suppose first that $A$ is integral and of dimension $1$; by virtue of the Nagata theorem \sref[0]{0.23.1.5} and \sref[0]{0.23.1.6}, the integral closure $A'$ of $A$ is a complete local integral ring and
\oldpage[IV-3]{107}
integrally closed and of dimension $1$ and a finite $A$-algebra; it is therefore a discrete valuation ring, and the proposition follows immediately in this case.

Let us pass to the general case; let $x$ be the closed point of $X=\operatorname{Spec}(A)$, and set $Y=X-\{x\}$; we know \sref{IV.10.5.9} that $Y$ is a Jacobson prescheme.
Let us show that the intersection of the prime ideals $\mathfrak{p}$ of $A$ such that $\dim(A/\mathfrak{p})=1$ is reduced to $0$; the proposition will then follow since, for each of these ideals $\mathfrak{p}$, the intersection of the kernels of the local homomorphisms of $A/\mathfrak{p}$ into discrete valuation rings $V$ which make $V$ a finite $(A/\mathfrak{p})$-algebra is reduced to $0$.
But to say that the intersection of these prime ideals is reduced to $0$ means that the set of these ideals is dense in $X$, or equivalently in $Y$ (since $Y$ is dense in $X$), and this follows immediately from \sref{IV.10.5.9}.
\end{proof}

\subsection{Dimension in Jacobson preschemes}
\label{subsection:IV.10.6}

The results of this section make more precise, in certain cases, and generalise results from \S\,5.

\begin{proposition}[10.6.1]
\label{IV.10.6.1}
Let $S$ be a locally Noetherian prescheme, satisfying in addition the conditions: $1^\circ$~$S$ is a Jacobson prescheme; $2^\circ$~for every $s\in S$, $\sh{O}_{S,s}$ is universally catenary \sref{IV.5.6.2}; $3^\circ$~every irreducible component $S'$ of $S$ is equicodimensional (in other words, for every closed point $s$ of $S'$ and every sub-prescheme of $S$ having $S'$ as underlying space, we have $\dim(S')=\dim(\sh{O}_{S',s})$).
We then have the following properties:
\begin{enumerate}[label=(\roman*)]
  \item For every morphism locally of finite type $g:X\to S$, $X$ satisfies conditions $1^\circ$, $2^\circ$, and $3^\circ$ above.
In particular, if $X$ is equidimensional (for example if $X$ is irreducible), $X$ is biequidimensional (in other words $X$ is catenary and for every closed point $x$ of $X$, we have $\dim(\sh{O}_{X,x})=\dim(X)$ \sref[0]{0.14.3.3}).
  \item Let $X$, $Y$ be two $S$-preschemes locally of finite type over $S$, $f:X\to Y$ an $S$-morphism; suppose $X$ irreducible and $f$ dominant.
If $\xi$ (resp.\ $\eta$) is the generic point of $X$ (resp.\ $Y$) and $e=\dim(f^{-1}(\eta))=\deg.\mathrm{tr}_{k(\eta)}k(\xi)$, we have
\begin{equation}
\label{IV.10.6.1.1}
\tag{10.6.1.1}
\dim(X)=\dim(Y)+e.
\end{equation}
  \item Let $X$, $Y$ be two $S$-preschemes locally of finite type over $S$, $f:X\to Y$ an $S$-morphism, $n$ a non-negative integer such that $\dim(f^{-1}(y))\geq n$ (resp.\ $\dim(f^{-1}(y))\leq n$) for every $y\in Y$.
Then we have
\begin{equation}
\label{IV.10.6.1.2}
\tag{10.6.1.2}
\dim(X)\geq\dim(Y)+n
\end{equation}
(resp.
\begin{equation}
\label{IV.10.6.1.3}
\tag{10.6.1.3}
\dim(X)\leq\dim(Y)+n).
\end{equation}
\end{enumerate}
\end{proposition}

\begin{proof}
(i) The property $1^\circ$ for $X$ follows from \sref{IV.10.4.7}.
For every $x\in X$, $\sh{O}_{X,x}$ is the local ring of an $\sh{O}_{S,g(x)}$-algebra of finite type at a prime ideal, and the homomorphism $\sh{O}_{S,g(x)}\to\sh{O}_{X,x}$
\oldpage[IV-3]{108}
is local; hence \sref{IV.5.6.3},~(iv) $\sh{O}_{X,x}$ is universally catenary.
To show that $X$ satisfies condition $3^\circ$, consider several cases:

\emph{a)} $X$ is a closed irreducible sub-prescheme of $S$; let $S'$ be an irreducible component of $S$ containing $X$, $\xi$ the generic point of $X$, $x$ a closed point of $X$; for every $s\in S'$, $\sh{O}_{S',s}$, a quotient of $\sh{O}_{S,s}$, is catenary \sref{IV.5.6.1}, hence conditions $2^\circ$ and $3^\circ$ imply that $S'$ is biequidimensional \sref[0]{0.14.3.3}.
By virtue of \sref{IV.5.1.2} and of \sref[0]{0.14.3.3.2}, we therefore have
\[
  \dim(\sh{O}_{X,x})=\dim(\sh{O}_{S',x})-\dim(\sh{O}_{S',\xi})=\dim(S')-\dim(\sh{O}_{S',\xi})
\]
which proves that $\dim(\sh{O}_{X,x})$ does not depend on the closed point $x$ considered, whence the assertion in this case \sref{IV.5.1.4}.

\emph{b)} $X$ is irreducible and $g$ dominant.
Then, for every closed point $x$ of $X$, $g(x)$ is closed in $S$ by \sref{IV.10.4.7}; as $\sh{O}_{S,g(x)}$ is universally catenary, it follows from \sref{IV.5.6.5.3} that we have
\[
  \dim(\sh{O}_{X,x})=\dim(\sh{O}_{S,g(x)})+e
\]
where $e=\dim(g^{-1}(\zeta))$, $\zeta$ being the generic point of $S$.
As $\dim(S)=\dim(\sh{O}_{S,g(x)})$ by virtue of condition $3^\circ$ for $S$, we have $\dim(\sh{O}_{X,x})=\dim(S)+e$ for every closed point $x\in X$; this proves condition $3^\circ$ for $X$ \sref{IV.5.1.4}, and at the same time the formula
\begin{equation}
\label{IV.10.6.1.4}
\tag{10.6.1.4}
\dim(X)=\dim(S)+e.
\end{equation}

\emph{c) General case.}
Considering a reduced sub-prescheme $X'$ of $X$ having for underlying space an irreducible component of $X$, one reduces to the case where $X$ is integral; using \sref[I]{I.5.2.2} and the case $a)$ proved above, one can then replace $S$ by the reduced sub-prescheme having $\overline{g(X)}$ as underlying space; one is then brought to case $b)$, and this completes proving (i).

(ii) The morphism $f$ being locally of finite type \sref[I]{I.1.3.4}, one can apply the results of (i) replacing $S$ by $Y$; furthermore, as $X$ is irreducible and $f$ dominant, one can also replace $S$ by $Y$ in \sref{IV.10.6.1.4}, which gives \sref{IV.10.6.1.1}.

(iii) The relative assertion where $\dim(f^{-1}(y))\leq n$ for every $y\in Y$ has already been proved under more general hypotheses in \sref{IV.5.6.7}.
Suppose that $\dim(f^{-1}(y))\geq n$ for every $y\in Y$, and consider a generic point $\eta$ of an irreducible component $Y'$ of $Y$; there exists at least one irreducible component $Z$ of $f^{-1}(\eta)$ of dimension $\geq n$; if $\xi$ is the generic point of $Z$, $\xi$ is also a generic point of an irreducible component $X'$ of $X$ such that $f(X')$ is dense in $Y'$ \sref[0]{0.2.1.8}.
Consider the reduced sub-preschemes of $X$, $Y$ having respectively for underlying spaces $X'$, $Y'$, and the restriction $X'\to Y'$ of $f$ \sref[I]{I.5.2.2}; it then follows from (ii) that $\dim(X')\geq\dim(Y')+n$, and \emph{a fortiori} $\dim(X)\geq\dim(Y')+n$; this being true for every irreducible component $Y'$ of $Y$, we conclude the inequality \sref{IV.10.6.1.2}.
\end{proof}

\begin{corollary}[10.6.2]
\label{IV.10.6.2}
Suppose that $X$ satisfies conditions $1^\circ$, $2^\circ$, and $3^\circ$ of \sref{IV.10.6.1}.
Then, for every dense open $U$ in $X$, we have $\dim(U)=\dim(X)$.
\end{corollary}

\begin{proof}
\oldpage[IV-3]{109}
We know indeed that $\dim(U)\leq\dim(X)$ \sref[0]{0.14.1.4}; furthermore, as $X$ is a Jacobson prescheme, $U$ contains a closed point of every irreducible component of $X$, hence $\dim(U)=\dim(X)$ by virtue of \sref{IV.10.6.1} and \sref[0]{0.14.1.2.1}.
\end{proof}

\begin{proposition}[10.6.3]
\label{IV.10.6.3}
Suppose that $X$ satisfies conditions $1^\circ$, $2^\circ$, and $3^\circ$ of \sref{IV.10.6.1}, and let $Y$ be a closed subset of $X$.
For every $x\in Y$, and every open neighbourhood $U$ of $x$ in $X$ not meeting the irreducible components of $Y$ that do not contain $x$, we have
\begin{equation}
\label{IV.10.6.3.1}
\tag{10.6.3.1}
\dim_x(Y)=\dim(U\cap Y)=\dim\overline{\{x\}}+\operatorname{codim}(\overline{\{x\}},Y)=\sup_i\dim(Y_i)
\end{equation}
where $Y_i$ ($1\leq i\leq m$) are the irreducible components of $Y$ containing $x$.
\end{proposition}

\begin{proof}
Considering a closed sub-prescheme of $X$ having $Y$ as underlying space, one can restrict, by virtue of \sref{IV.10.6.1},~(i), to the case where $Y=X$.
By the choice of $U$, $U$ is the union of the $U\cap X_i$, hence $\dim(U)=\sup_i\dim(U\cap X_i)$; but by \sref{IV.10.6.2} applied to a closed sub-prescheme of $X$ having $X_i$ as underlying space, we have (taking into account \sref{IV.10.6.1},~(i)), $\dim(U\cap X_i)=\dim(X_i)$; this proves that the second and the fourth terms of \sref{IV.10.6.3.1} are equal.
As $U\cap X_i$ is biequidimensional \sref{IV.10.6.1},~(i) (since the immersion $U\cap X_i\to X_i$ is of finite type \sref[I]{I.6.3.5}), we have

\begin{equation}
\label{IV.10.6.3.2}
\tag{10.6.3.2}
\dim(X_i)=\dim(U\cap X_i)=\dim(U\cap\overline{\{x\}})+\operatorname{codim}(U\cap\overline{\{x\}},U\cap X_i)
\end{equation}
\sref[0]{0.14.3.5}.
By \sref{IV.10.6.2} applied to a closed sub-prescheme of $X$ having $\overline{\{x\}}$ as underlying space, $\dim(U\cap\overline{\{x\}})=\dim(\overline{\{x\}})$; as we also have, for the same reasons,

\begin{equation}
\label{IV.10.6.3.3}
\tag{10.6.3.3}
\dim(X_i)=\dim(\overline{\{x\}})+\operatorname{codim}(\overline{\{x\}},X_i)
\end{equation}
we obtain
\[
  \dim(U)=\dim(\overline{\{x\}})+\sup_i\operatorname{codim}(\overline{\{x\}},X_i)=\dim(\overline{\{x\}})+\operatorname{codim}(\overline{\{x\}},X)
\]
by definition of the codimension \sref[0]{0.14.2.1}.
This shows that $\dim(U)$ is \emph{independent} of the open neighbourhood $U$ of $x$ satisfying the conditions of the statement, hence is equal to $\dim_x(X)$, by \sref[0]{0.14.1.4.1}.
\end{proof}

\begin{corollary}[10.6.4]
\label{IV.10.6.4}
Under the hypotheses of \sref{IV.10.6.3}, let $\sh{F}$ be a coherent $\sh{O}_X$-module, $Y$ its support.
For every $x\in Y$, we have
\begin{equation}
\label{IV.10.6.4.1}
\tag{10.6.4.1}
\dim(\overline{\{x\}})+\dim(\sh{F}_x)=\dim_x Y.
\end{equation}
\end{corollary}

\begin{proof}
Indeed, this follows from \sref{IV.10.6.3.1} and from the formula $\dim(\sh{F}_x)=\operatorname{codim}(\overline{\{x\}},Y)$ \sref{IV.5.1.12.2}.
\end{proof}

\subsection{Examples and counter-examples}
\label{subsection:IV.10.7}

\begin{env}[10.7.1]
\label{IV.10.7.1}
Let $S$ be a locally Noetherian prescheme of dimension $\leq 1$ and suppose that $S$ is a \emph{Jacobson prescheme}; when $S$ is Noetherian, this amounts to saying that the irreducible components of $S$ of dimension $1$ are infinite, for every $x\in S$ which is not a closed point is the generic point of such a component \sref{IV.10.4.5} and \sref{IV.10.5.4}.
Then $S$ also satisfies conditions $2^\circ$ and $3^\circ$ of \sref{IV.10.6.1}: indeed, every local ring $\sh{O}_{S,s}$ is of dimension $0$ or $1$, and hence universally catenary \sref{IV.7.2.9}; on the other hand, an irreducible component $S'$ of $S$ is either reduced
\oldpage[IV-3]{110}
to a point or of dimension $1$, and for every closed point $s\in S'$, $\sh{O}_{S',s}$ is necessarily of dimension $1$.

From these remarks and from \sref{IV.10.6.1} it follows that every prescheme locally of finite type over $S$ also satisfies the properties $1^\circ$, $2^\circ$, and $3^\circ$; it is in particular the case of preschemes \emph{locally of finite type over a field or over $\mathbf{Z}$}.
\end{env}

\begin{env}[10.7.2]
\label{IV.10.7.2}
Let $A$ be a Noetherian local ring \emph{universally catenary} and let $S$ be the complementary of \emph{the closed point} $a$ in $X=\operatorname{Spec}(A)$.
Then $S$ satisfies conditions $1^\circ$, $2^\circ$, and $3^\circ$ of \sref{IV.10.6.1}: indeed, one has already seen that $S$ is a Jacobson prescheme \sref{IV.10.5.9}; as $A$ is universally catenary, so are all the local rings $A_\mathfrak{p}$ at the prime ideals of $A$ \sref{IV.5.6.3}.
On the other hand, an irreducible component $S'$ of $S$ is the complement of $a$ in an irreducible component $X'$ of $X$. For every closed point $x$ of $S$, the closure of $x$ in $X$ is $\{x,a\}$, so $\dim(\overline{\{x\}})=1$. Since $X'$ is biequidimensional by hypothesis \sref[0]{0.14.3.3}, it follows that $\dim(\sh{O}_{X',x})=\dim(X')-1$ \sref[0]{0.14.3.2}, proving that $S'$ is equicodimensional \sref{IV.5.1.4}.
\end{env}

\begin{env}[10.7.3]
\label{IV.10.7.3}
Let $A$ be a discrete valuation ring; let us show that in $B=A[T_1,\ldots,T_n]$, with $n$ indeterminates, there exist two maximal ideals $\mathfrak{m}$, $\mathfrak{n}$ of \emph{respective heights} $n$ and $n+1$.
We have seen \sref{IV.5.2.5},~(i) for $n=1$; let us prove it by induction on $n$.
As $A[T_1,\ldots,T_{n+1}]$ is a free, hence faithfully flat, $A[T_1,\ldots,T_n]$-module, there are maximal ideals $\mathfrak{m}'$ and $\mathfrak{n}'$ of $A[T_1,\ldots,T_{n+1}]$ lying respectively over $\mathfrak{m}$ and $\mathfrak{n}$ \sref[0]{0.6.5.1}; by \sref{IV.5.5.3}, these ideals have respective heights $n+1$ and $n+2$, proving the assertion.
Suppose $n\geq 2$ in what follows.
Let $\mathfrak{J}=\mathfrak{m}\cap\mathfrak{n}=\mathfrak{m}\mathfrak{n}$ and let $R=1+\mathfrak{J}$, which is a multiplicative subset of $B$.
If $B'=R^{-1}B$, the ideal $\mathfrak{J}'=R^{-1}\mathfrak{J}$ is contained in the radical of $B'$ (Bourbaki, \emph{Alg\`ebre commutative}, Chap.~III, \S3, no.~5, Prop.~12).
As a topological space, $X'=\operatorname{Spec}(B')$ is identified with a subspace of $X=\operatorname{Spec}(B)$, and at the points $x$ of $X'$ the local rings $\sh{O}_{X,x}$ and $\sh{O}_{X',x}$ are the same \sref[I]{I.1.6.2}.
Consider in $X'$ the closed subset $Y'=V(\mathfrak{J}')$, and put $S=X'\setminus Y'$.
Then $S$ is a Jacobson prescheme \sref{IV.10.5.7}, and is evidently irreducible and Noetherian; moreover, for every $s\in S$, the local ring $\sh{O}_{S,s}=\sh{O}_{X,s}$ is universally catenary by \sref{IV.5.6.3}, since $A$ is universally catenary \sref{IV.5.6.4}.
Nevertheless, there are two closed points $a,b$ of $S$ such that $\sh{O}_{S,a}$ and $\sh{O}_{S,b}$ \emph{do not have the same dimension}; in other words, $S$ does not satisfy condition $3^\circ$ of \sref{IV.10.6.1}.
To see this, consider the two maximal ideals $\mathfrak{m}'=R^{-1}\mathfrak{m}$ and $\mathfrak{n}'=R^{-1}\mathfrak{n}$ of $B'$, which have heights $n$ and $n+1$, respectively.
We have $\mathfrak{J}'=\mathfrak{m}'\cap\mathfrak{n}'$, and hence $\mathfrak{J}'$ is contained in no prime ideal of $B'$ distinct from $\mathfrak{m}'$ and $\mathfrak{n}'$; these are consequently the only maximal ideals of $B'$.
Let $a',b'$ be the only closed points of $X'$, corresponding to $\mathfrak{m}'$ and $\mathfrak{n}'$.
There exists in $B'$ a nonmaximal prime ideal $\mathfrak{p}'\subset\mathfrak{m}'$ which is not contained in $\mathfrak{n}'$: it suffices to show that in $B$ there is a nonmaximal prime ideal contained in $\mathfrak{m}$ and not in $\mathfrak{n}$; for this one may, for example, consider the fibre of $\mathfrak{m}$ for the morphism corresponding to the injection $A[T_1]\longrightarrow B=A[T_1,\ldots,T_n]$ and apply \sref{IV.6.1.2}.
By considering a maximal chain of prime ideals between $\mathfrak{p}'$ and $\mathfrak{m}'$, and replacing $\mathfrak{p}'$ by the penultimate ideal in this chain, we may therefore suppose that the point $a$ of $X'$ corresponding to $\mathfrak{p}'$ has closure $\{a,a'\}$ in $X'$.
Since $B'_{\mathfrak{m}'}$ is biequidimensional, $\mathfrak{p}'$ then has height $n-1$.
In the same way one constructs a nonmaximal prime ideal $\mathfrak{q}'$ of $B'$ of height $n$ such that, if $b$ is the corresponding point in $X'$, the closure of $b$ in $X'$ is $\{b,b'\}$.
Thus $a$ and $b$ belong to $S$ and are closed in $S$, and consequently answer the question.
\end{env}

\subsection{Rectified depth}
\label{subsection:IV.10.8}

\begin{definition}[10.8.1]
\label{IV.10.8.1}
Let $X$ be a locally Noetherian prescheme, $\sh{F}$ a coherent $\sh{O}_X$-Module.
For every $x\in X$, we call \emph{rectified depth} of $\sh{F}$ at the point $x$ and denote $\operatorname{prof}^*_x(\sh{F})$ the number (integer $\geq 0$ or $+\infty$) equal to
\begin{equation}
\label{IV.10.8.1.1}
\tag{10.8.1.1}
\operatorname{prof}^*_x(\sh{F})=\operatorname{prof}(\sh{F}_x)+\dim(\overline{\{x\}})
\end{equation}
where $\overline{\{x\}}$ is the closure of the point $x$ in $X$.
For every subset $Z$ of $X$, we call \emph{rectified depth of $\sh{F}$ along $Z$} and denote $\operatorname{prof}^*_Z(\sh{F})$ the number
\begin{equation}
\label{IV.10.8.1.2}
\tag{10.8.1.2}
\operatorname{prof}^*_Z(\sh{F})=\inf_{x\in Z}\operatorname{prof}^*_x(\sh{F}).
\end{equation}
\end{definition}

In other terms, for every integer $n$, the relation $\operatorname{prof}^*_Z(\sh{F})\geq n$ is equivalent to $\operatorname{prof}^*_x(\sh{F})\geq n$ for every $x\in Z$.
If $Z=X$, one writes $\operatorname{prof}^*(\sh{F})$ instead of $\operatorname{prof}^*_X(\sh{F})$.

\oldpage[IV-3]{111}

\begin{remarks}[10.8.2]
\label{IV.10.8.2}
(i) At every closed point $x\in X$, the rectified depth equals the depth.

(ii) Let $Y$ be a closed sub-prescheme of $X$, $j:Y\to X$ the canonical injection, and $\sh{G}$ a coherent $\sh{O}_Y$-module.
We know \sref{IV.5.7.3},~(vi) that $\operatorname{prof}(\sh{G}_x)=\operatorname{prof}(j_*(\sh{G})_x)$ for every $x\in Y$; we therefore deduce that $\operatorname{prof}^*_x(\sh{G})=\operatorname{prof}^*_x(j_*(\sh{G}))$ for every $x\in Y$.

(iii) The notion of rectified depth is of interest only when it is of a \emph{local} character, i.e.\ when it does not change when one replaces $X$ by an arbitrary open neighbourhood $U$ of $x$.
This obviously requires that $x$ not be isolated in $\overline{\{x\}}$ when $x$ is not a closed point of $X$, and consequently that $X$ be a \emph{Jacobson prescheme} \sref{IV.10.4.5.1}; most often, it will also be necessary to know that $\dim(U)=\dim(X)$ for every dense open $U$ in $X$, and one must therefore also suppose that $X$ satisfies conditions $2^\circ$ and $3^\circ$ of \sref{IV.10.6.1}.
\end{remarks}

\begin{lemma}[10.8.3]
\label{IV.10.8.3}
Let $X$ be a regular and biequidimensional prescheme, $\sh{F}$ a coherent $\sh{O}_X$-Module.
Then, for every $x\in X$, we have
\begin{equation}
\label{IV.10.8.3.1}
\tag{10.8.3.1}
\operatorname{prof}^*_x(\sh{F})=\dim(X)-\dim.\operatorname{proj}(\sh{F}_x).
\end{equation}
\end{lemma}

\begin{proof}
Indeed, as $X$ is biequidimensional, we have \sref[0]{0.14.3.5.1}
\[
  \dim(\overline{\{x\}})=\dim(X)-\operatorname{codim}(\overline{\{x\}},X)=\dim(X)-\dim(\sh{O}_{X,x})
\]
by virtue of \sref{IV.5.1.2}.
On the other hand, as $X$ is regular, we have by \sref[0]{0.17.3.4}
\[
  \operatorname{prof}(\sh{F}_x)=\dim(\sh{O}_{X,x})-\dim.\operatorname{proj}(\sh{F}_x)
\]
whence the lemma.
\end{proof}

\begin{corollary}[10.8.4]
\label{IV.10.8.4}
Under the hypotheses of \sref{IV.10.8.3}, the function $x\mapsto\operatorname{prof}^*_x(\sh{F})$ is lower semi-continuous.
\end{corollary}

\begin{proof}
This follows from \sref{IV.10.8.3.1}, since $x\mapsto\dim.\operatorname{proj}(\sh{F}_x)$ is upper semi-continuous \sref{IV.6.11.1}.
\end{proof}

\begin{proposition}[10.8.5]
\label{IV.10.8.5}
Let $S$ be a locally Noetherian prescheme, $X$ a prescheme locally of finite type over $S$, $\sh{F}$ a coherent $\sh{O}_X$-module.
Suppose that $S$ satisfies the conditions: $1^\circ$~$S$ is a Jacobson prescheme; $2^\circ$~$S$ is regular; $3^\circ$~the irreducible components of $S$ are equicodimensional.
Then the function $x\mapsto\operatorname{prof}^*_x(\sh{F})$ is lower semicontinuous; in other terms, for every integer $n$, the set $U_n$ of $x\in X$ such that $\operatorname{prof}^*_x(\sh{F})\geq n$ is open.
\end{proposition}

\begin{proof}
As the local rings $\sh{O}_{S,s}$ of $S$ are regular, they are universally catenary \sref{IV.5.6.4}; in other words, $S$ satisfies conditions $1^\circ$, $2^\circ$, and $3^\circ$ of \sref{IV.10.6.1}, hence the same is true of $X$ \sref{IV.10.6.1},~(i).
The notion of rectified depth being then of local character \sref{IV.10.8.2},~(iii), one can restrict to the case where $S=\operatorname{Spec}(A)$ and $X=\operatorname{Spec}(B)$ are affine, $A$ being a regular ring and $B$ an $A$-algebra of finite type, hence a quotient of a polynomial ring $C=A[T_1,\ldots,T_n]$, and this last ring is regular \sref[0]{0.17.3.7}.
One can therefore suppose that $X$ is a closed sub-prescheme of a regular prescheme $Y$ also satisfying conditions $1^\circ$, $2^\circ$, and $3^\circ$ of \sref{IV.10.6.1}; taking into account the remark \sref{IV.10.8.2},~(ii), one is thus brought to the case where $X$ is in addition regular and Noetherian.
But as the local rings of $X$ are then integral, the irreducible components of $X$ are
\oldpage[IV-3]{112}
open \sref[I]{I.6.1.10}, and one can therefore also suppose $X$ irreducible.
Then, as the local rings of $X$ are catenary \sref[0]{0.16.5.12}, the hypothesis $3^\circ$ of \sref{IV.10.6.1} implies that $X$ is biequidimensional (\sref{IV.5.1.5} and \sref[0]{0.14.3.3}); it therefore suffices to apply \sref{IV.10.8.4}.
\end{proof}

One notes that if $S$ is the spectrum of a field or of $\mathbf{Z}$, it satisfies the conditions of \sref{IV.10.8.5}.

\begin{corollary}[10.8.6]
\label{IV.10.8.6}
The hypotheses being those of \sref{IV.10.8.5}, for every $x\in X$, the number $\operatorname{prof}^*_x(\sh{F})$ is the unique integer $n$ having the following property: there exists an open neighbourhood $U$ of $x$ such that for every point $x'\in U\cap\overline{\{x\}}$, closed in $U$, we have $\operatorname{prof}(\sh{F}_{x'})=n$.
In particular, for $\operatorname{prof}^*_x(\sh{F})\geq m$ (resp.\ $\operatorname{prof}^*_x(\sh{F})\leq m$), it is necessary and sufficient that there exist an open neighbourhood $V$ of $x$ in $X$ such that, for every point $x'\in V\cap\overline{\{x\}}$, closed in $V$, we have $\operatorname{prof}(\sh{F}_{x'})\geq m$ (resp.\ $\operatorname{prof}(\sh{F}_{x'})\leq m$).
\end{corollary}

\begin{proof}
Indeed, one can restrict to the case where $x$ is not a closed point; if $\operatorname{prof}^*_x(\sh{F})=n$, the set of $y\in X$ such that $\operatorname{prof}^*_y(\sh{F})\leq n$ is closed by virtue of \sref{IV.10.8.5}, hence contains $\overline{\{x\}}$, and by virtue of the lower semi-continuity of $\operatorname{prof}^*_y(\sh{F})$, there exists an open neighbourhood $U$ of $x$ such that $\operatorname{prof}^*_y(\sh{F})=n$ for every $y\in U\cap\overline{\{x\}}$, hence $\operatorname{prof}(\sh{F}_{x'})=n$ if $x'\in U\cap\overline{\{x\}}$ is closed (since the notion of rectified depth is local).
\end{proof}

For preschemes satisfying the hypotheses of \sref{IV.10.8.5}, the notion of rectified depth can therefore be defined by means of the values of the depth at the \emph{closed points} of $X$ (the latter forming a set which is very dense in every closed subset of $X$).

\begin{proposition}[10.8.7]
\label{IV.10.8.7}
Let $S$ be a prescheme locally Noetherian satisfying conditions $1^\circ$, $2^\circ$, and $3^\circ$ of \sref{IV.10.6.1}.
Let $X$ be a prescheme locally of finite type over $S$, $\sh{F}$ a coherent $\sh{O}_X$-Module, $Y=\operatorname{Supp}(\sh{F})$.
Then, for every $x\in Y$, we have
\begin{equation}
\label{IV.10.8.7.1}
\tag{10.8.7.1}
\operatorname{prof}^*_x(\sh{F})=\dim_x(Y)-\operatorname{coprof}(\sh{F}_x).
\end{equation}
\end{proposition}

\begin{proof}
Indeed, by definition, $\operatorname{coprof}(\sh{F}_x)=\dim(\sh{F}_x)-\operatorname{prof}(\sh{F}_x)$, and it follows from \sref{IV.10.6.4} that $\dim_x(Y)=\dim(\overline{\{x\}})+\dim(\sh{F}_x)$; whence \sref{IV.10.8.7.1} by definition of $\operatorname{prof}^*_x(\sh{F})$.
\end{proof}

\begin{corollary}[10.8.8]
\label{IV.10.8.8}
Suppose the hypotheses on $f$ and $\sh{F}$ are those of \sref{IV.9.9.1}. Then the function $x\mapsto\operatorname{prof}^*_x(\sh{F}_{f(x)})$ is constructible.
\end{corollary}

\begin{proof}
For every $s\in S$, the fibre $X_s$, being locally of finite type over $k(s)$, is a Jacobson prescheme \sref{IV.10.4.7}. Since $\operatorname{Spec}(k(s))$ satisfies the conditions of \sref{IV.10.6.1}, one has
\[
  \operatorname{prof}^*_x(\sh{F}_{f(x)})
  =\dim_x(\operatorname{Supp}(\sh{F}_{f(x)}))
   -\operatorname{coprof}((\sh{F}_{f(x)})_x)
\]
by \sref{IV.10.8.7}.
Now, if $Z=\operatorname{Supp}(\sh{F})$, we have $Z_{f(x)}=\operatorname{Supp}(\sh{F}_{f(x)})$ \sref[I]{I.9.1.13} and $Z$ is locally constructible \sref{IV.8.9.1}; therefore the functions $x\mapsto\dim_x(\operatorname{Supp}(\sh{F}_{f(x)}))$ and $x\mapsto\operatorname{coprof}((\sh{F}_{f(x)})_x)$ are locally constructible (\sref{IV.9.9.1} and \sref{IV.9.9.3}), which proves the proposition.
\end{proof}

\subsection{Maximal spectra and ultra-preschemes}
\label{subsection:IV.10.9}

The results of this section will not be used in what follows.

\begin{env}[10.9.1]
\label{IV.10.9.1}
Let $X$ be a Jacobson prescheme, and let $S(X)$ be the \emph{ringed space} whose underlying space is the subspace of closed points of $X$, and the sheaf of rings the sheaf induced on this subspace by $\sh{O}_X$, in other words the sheaf $\psi^*(\sh{O}_X)$, where $\psi:S(X)\to X$ denotes the canonical injection.
As $\psi$ is a quasi-
\oldpage[IV-3]{113}
homeomorphism, we have seen \sref{IV.10.2.8},~(ii) that if $\theta:\sh{O}_X\to\psi_*(\sh{O}_{S(X)})$ is the homomorphism of sheaves of rings such that $\theta^\sharp:\psi^*(\sh{O}_X)\to\sh{O}_{S(X)}$ is the identity, then $j_X=(\psi,\theta)$ is a \emph{quasi-isomorphism} of ringed spaces, and $\sh{F}\mapsto j_X^*(\sh{F})$ is an equivalence of the category of $\sh{O}_X$-Modules and that of $\sh{O}_{S(X)}$-Modules.
It is clear that under this equivalence, locally free (resp.\ coherent) $\sh{O}_X$-modules correspond to locally free (resp.\ coherent) $\sh{O}_{S(X)}$-modules. Furthermore, if $\sh{L}$ is a locally free $\sh{O}_X$-module and $U$ is an open subset of $X$ such that $\sh{L}|_U\simeq(\sh{O}_X|_U)^n$, then
\[
  j_X^*(\sh{L})|_{U\cap S(X)}
  \simeq \bigl(\sh{O}_{S(X)}|_{U\cap S(X)}\bigr)^n.
\]
\end{env}

\begin{env}[10.9.2]
\label{IV.10.9.2}
Let $X$, $Y$ be two Jacobson preschemes and $f=(\rho,\lambda):X\to Y$ a morphism \emph{locally of finite type}.
One has seen \sref{IV.10.4.7} that $\rho(S(X))\subset S(Y)$ and by restriction of $\rho$ to $S(X)$, one defines therefore a continuous map $S(\rho):S(X)\to S(Y)$.
For every open $V$ of $Y$, one defines by composition a homomorphism of rings
\[
  \Gamma(V\cap S(Y),\sh{O}_{S(Y)})\to\Gamma(V,\sh{O}_Y)\xrightarrow{\lambda_V}\Gamma(f^{-1}(V),\sh{O}_X)\to\Gamma(f^{-1}(V)\cap S(X),\sh{O}_{S(X)})
\]
where the two extreme isomorphisms have been defined in \sref{IV.10.9.1}; it is clear that this defines a homomorphism of sheaves of rings $S(\lambda):\sh{O}_{S(Y)}\to\rho_*(\sh{O}_{S(X)})$ (recalling that the open subsets of $S(X)$ (resp.\ $S(Y)$) correspond bijectively to those of $X$ (resp.\ $Y$) \sref{IV.10.2.1}); one obtains thus a morphism of ringed spaces $S(f)=(S(\rho),S(\lambda)):(S(X),\sh{O}_{S(X)})\to(S(Y),\sh{O}_{S(Y)})$ such that the diagram
\[
  \xymatrix{
    S(X)\ar[r]^{S(f)}\ar[d]_{j_X} & S(Y)\ar[d]^{j_Y}\\
    X\ar[r]^{f} & Y
  }
\]
is commutative; furthermore, if $Z$ is a third Jacobson prescheme and $g:Y\to Z$ a morphism locally of finite type, $S(g\circ f)=S(g)\circ S(f)$.
One has thus defined a \emph{covariant functor} $S:\mathbf{C}\to\mathbf{C}'$, where $\mathbf{C}'$ is the category of \emph{ringed spaces in local rings}, and $\mathbf{C}$ the category whose objects are the \emph{Jacobson preschemes} and the morphisms are the morphisms locally of finite type between Jacobson preschemes.
\end{env}

\begin{env}[10.9.3]
\label{IV.10.9.3}
Let us now determine the subcategory $\mathbf{C}''$ of $\mathbf{C}'$ formed by the ringed spaces isomorphic to $S(X)$ and whose morphisms come from the $S(f)$.
Suppose first that $X=\operatorname{Spec}(A)$, where $A$ is a Jacobson ring; then $S(X)$ is the set of \emph{maximal} ideals of $A$, equipped: $1^\circ$~with the topology induced by that of $X$, so that a base for this topology is formed by the $D^m(h)=D(h)\cap S(X)$, i.e.\ the set of maximal ideals $\mathfrak{m}$ of $A$ such that $h\notin\mathfrak{m}$; $2^\circ$~with the sheaf of rings $\sh{O}_{S(X)}$ such that $\Gamma(D^m(h),\sh{O}_{S(X)})=A_h$.
We will say that this ringed space is the \emph{maximal spectrum} of the Jacobson ring $A$ and denote it $\operatorname{Spm}(A)$.

One notes that $j:D(h)\to X$ is the canonical injection, the ringed space induced on $D^m(h)$ by $S(X)$ and $S(D(h))$ are canonically isomorphic, and that the canonical injection $D^m(h)\to S(X)$ of ringed spaces is equal to $S(j)$.

Let $B$ be a second Jacobson ring, $Y=\operatorname{Spec}(B)$, $\varphi:B\to A$ a ring homomorphism making $A$ a $B$-algebra of \emph{finite type}, $f=({}^a\varphi,\widetilde{\varphi}):X\to Y$ the corresponding morphism of preschemes, and
\[
  S(f):\operatorname{Spm}(A)\to\operatorname{Spm}(B)
\]
the morphism of ringed spaces corresponding to $f$.
It is clear that $S(f)=(\psi,\theta)$ is a morphism of ringed spaces in local rings, i.e.\ $(\mathrm{Err}_{\mathrm{II}},1.8.2)$ for every $x\in\operatorname{Spm}(A)$, $\theta^\sharp_x$ is a \emph{local} homomorphism.
Conversely:
\end{env}

\begin{proposition}[10.9.4]
\label{IV.10.9.4}
Let $A$, $B$ be two Jacobson rings.
If $u=(\psi,\theta):\operatorname{Spm}(A)\to\operatorname{Spm}(B)$ is a morphism of ringed spaces in local rings such that $\Gamma(\theta):B\to A$ makes $A$ a $B$-algebra of finite type, there exists a morphism of preschemes $f:\operatorname{Spec}(A)\to\operatorname{Spec}(B)$ and a unique one such that $u=S(f)$.
\end{proposition}

\begin{proof}
The uniqueness of $f$ is evident, since if $f=({}^a\varphi,\widetilde{\varphi})$, one must have $\varphi=\Gamma(\theta)$; it remains to show that $S(f)$ is defined and that $u=S(f)$.
Now, the first assertion follows from the fact that $\varphi$ is supposed to make $A$ a $B$-algebra of finite type, and hence $f(\operatorname{Spm}(A))\subset\operatorname{Spm}(B)$; the fact that $\theta^\sharp_x$ is a local homomorphism for every $x\in\operatorname{Spm}(A)$ allows one to repeat the argument of \sref[I]{I.1.7.3}, restricting to points $x$ of $\operatorname{Spm}(A)$, to show thus successively that $\psi(x)={}^a\varphi(x)$ for every $x\in\operatorname{Spm}(A)$, then that $\theta^\sharp_x=\varphi_x:B_{\psi(x)}\to A_x$, which completes proving that $u=S(f)$.
\end{proof}

\oldpage[IV-3]{114}

\begin{env}[10.9.5]
\label{IV.10.9.5}
Consider now a ringed space $(X,\sh{O}_X)$; we will say that an open subset $U$ of $X$ is \emph{ultra-affine} if the ringed space induced $(U,\sh{O}_X|U)$ is isomorphic to a maximal spectrum $\operatorname{Spm}(A)$, where $A$ is a Jacobson ring.
We will say that $X$ is an \emph{ultra-prescheme} if every point of $X$ admits an ultra-affine open neighbourhood.
One shows, as in \sref[I]{I.2.1.3} and \sref[I]{I.2.1.4}, that the ultra-affine open subsets form a \emph{base} for the topology of $X$ and that $X$ is a Kolmogoroff space.
If $Y$ is a second ultra-prescheme, we will say that a morphism of ringed spaces $f:X\to Y$ is a \emph{morphism of ultra-preschemes} if it satisfies the following conditions: $1^\circ$~$f$ is a morphism of ringed spaces in local rings; $2^\circ$~for every $x\in X$, there is an ultra-affine open neighbourhood $V$ of $f(x)$ in $Y$ and an ultra-affine open $U$ of $x$ in $X$ such that $f(U)\subset V$ and the homomorphism $\Gamma(V,\sh{O}_Y)\to\Gamma(U,\sh{O}_X)$ corresponding to $f$ makes $\Gamma(U,\sh{O}_X)$ a $\Gamma(V,\sh{O}_Y)$-\emph{algebra of finite type}.

It is immediate that one defines well the composite of two such morphisms, being a third such by the final remark of \sref{IV.10.9.3}.
It is clear that the category $C''_0$ thus defined is a subcategory of $C'$ that contains $C''$; we now propose to show that $C''=C''_0$, in other words:
\end{env}

\begin{proposition}[10.9.6]
\label{IV.10.9.6}
The functor $X\mapsto S(X)$ of $\mathbf{C}$ into $\mathbf{C}''_0$ is an equivalence of categories.
\end{proposition}

\begin{proof}
$1^\circ$ Let us first show that the functor $X\mapsto S(X)$ is \emph{fully faithful}, in other words that for $X$, $Y$ in $\mathbf{C}$, the canonical map
\[
  \operatorname{Hom}_{\mathbf{C}}(X,Y)\to\operatorname{Hom}_{\mathbf{C}''_0}(S(X),S(Y))
\]
is \emph{bijective}.
In the first place it is \emph{injective}: let $f$, $g$ be two morphisms locally of finite type of $X$ into $Y$, and suppose that $S(f)=S(g)$; we have $f^{-1}(V)\cap S(X)=g^{-1}(V)\cap S(X)$ for every open $V$ of $Y$, hence \sref{IV.10.3.1} and \sref{IV.10.2.7} $f^{-1}(V)=g^{-1}(V)$; it suffices therefore to prove that for every affine open $V$ of $Y$, $f$ and $g$ coincide in $f^{-1}(V)=g^{-1}(V)$; in other words, one reduces to the case where $Y=\operatorname{Spec}(B)$ is the spectrum of a Jacobson ring $B$.
It suffices \sref[I]{I.2.2.4} to show that the homomorphisms of rings $B\to\Gamma(X,\sh{O}_X)$ corresponding to $f$ and $g$ are the same.
Now, for every $s\in B$, the images of $s$ under these homomorphisms are two sections of $\sh{O}_X$ over $X$ which, by hypothesis, induce the \emph{same} section over $S(X)$; we know \sref{IV.10.2.8} that these sections are identical, whence our assertion.
In the second place let us prove that every morphism $h:S(X)\to S(Y)$ (for the category $\mathbf{C}''_0$) is of the form $S(f)$, where $f:X\to Y$ is a morphism locally of finite type.
By hypothesis, there exists an ultra-affine open covering $(U'_\lambda)$ (resp.\ $(V'_\mu)$) of $S(X)$ (resp.\ $S(Y)$) such that for every $\lambda$, $h(U'_\lambda)$ is contained in one of the $V'_\mu$ and the homomorphism of rings corresponding to the restriction $h_\lambda:U'_\lambda\to V'_\mu$ makes the second ring an algebra of finite type over the first.
We may suppose that $U'_\lambda=U_\lambda\cap S(X)$ and $V'_\mu=V_\mu\cap S(Y)$, where $U_\lambda$ and $V_\mu$ are uniquely determined affine open subsets of $X$ and $Y$, respectively. If, for every $\lambda$, $h_\lambda=S(f_\lambda)$ for a morphism of finite type $f_\lambda:U_\lambda\to V_\mu$, then the first part of the argument, applied to the restrictions of $f_\alpha$ and $f_\beta$ to $U_\alpha\cap U_\beta$, shows that the $f_\lambda$ are restrictions of one morphism $f:X\to Y$, and clearly $h=S(f)$. We are thus reduced to the case in which $X$ and $Y$ are affine, and the conclusion follows from \sref{IV.10.9.4}.

$2^\circ$ It remains to prove that every ultra-prescheme $X'$ is of the form $S(X)$ for a Jacobson prescheme $X$ (which will necessarily be unique to isomorphism near, by virtue of $1^\circ$).
There are ultra-affine open subsets $(U'_\alpha)$ of $X'$, each of which is of the form $S(U_\alpha)$, $U_\alpha$ being the spectrum of a Jacobson ring.
For every couple of indices $\alpha$, $\beta$, consider the unique open $U_{\alpha\beta}$ of $U_\alpha$ whose trace on $U'_\alpha$ is $U'_\alpha\cap U'_\beta$; by virtue of $1^\circ$, the identical self-morphism of $U'_\alpha\cap U'_\beta$ is of the form $S(\theta_{\alpha\beta})$, where $\theta_{\alpha\beta}:U_{\alpha\beta}\to U_{\beta\alpha}$ is an isomorphism of preschemes.
One easily verifies (by virtue of $1^\circ$) that the family $(\theta_{\alpha\beta})$ satisfies the glueing condition \sref[0]{0.4.1.7}, and that this family defines a prescheme $X$; it is clear then that we have $X'=S(X)$, which completes the proof.
\end{proof}

\subsection{Algebraic spaces of Serre}
\label{subsection:IV.10.10}

\begin{env}[10.10.1]
\label{IV.10.10.1}
The language introduced by Serre in (FAC) is sometimes convenient, notably in questions where the main interest is attached to the rational points over the base field $k$ (algebraically closed by hypothesis) of the ``algebraic varieties'' over $k$ that one considers.
We will sketch here this language by linking it to the considerations that precede, to allow the reader to translate the statements of Serre in the language of schemes.
It is moreover possible to develop the language of Serre also for preschemes over a non-algebraically closed field (and even over an Artinian ring); but this introduces considerable technical complications, and moreover, on a field of arbitrary base, the advantages (above all psychological ones) of the point of view of Serre disappear; therefore we will confine ourselves to the framework fixed by Serre.
The present section, just like the preceding one, will not be used in the sequel of this Treatise, and we will therefore content ourselves with brief indications.
\end{env}

\oldpage[IV-3]{115}

\begin{env}[10.10.2]
\label{IV.10.10.2}
Given an ultra-prescheme $R$, one can naturally (as in the whole category) define the notion of $R$-\emph{ultra-prescheme}.
Consider in particular an algebraically closed field $k$; $\operatorname{Spm}(k)$ is then identical to $\operatorname{Spec}(k)$; we will say that a $\operatorname{Spm}(k)$-ultra-prescheme $X'$ is a \emph{prealgebraic space over $k$}: it is a $k$-ringed space in local rings whose every point has an open neighbourhood isomorphic to the maximal spectrum of a $k$-algebra of finite type; it follows that it is the same as saying (by \sref{IV.10.9.6}) that $X'=S(X)$, where $X$ is a $k$-prescheme locally of finite type.
If $X$, $Y$, $Z$ are three $k$-preschemes locally of finite type, the same is true of $X\times_ZY$ \sref[I]{subsection:I.3.4}; hence, by \sref{IV.10.9.6}, products exist in the category of prealgebraic spaces over $k$, both the product ``over $k$'' $X'\times_kY'$ and the ``fibre product'' $X'\times_{Z'}Y'$, where $X'$, $Y'$, $Z'$ are three prealgebraic spaces over $k$.
One can therefore define the diagonal morphism $\Delta_{X'}:X'\to X'\times_k X'$ (which is none other than $S(\Delta_X)$); one says that $X'$ is an \emph{algebraic space over $k$} if $X$ is a scheme, and it amounts to the same as saying that the image of $\Delta_{X'}$ is a \emph{closed} subset of $X'\times_k X'$.
\end{env}

\begin{env}[10.10.3]
\label{IV.10.10.3}
The simplifications that come from the hypothesis that $k$ is algebraically closed amount first of all to this: for every $k$-prescheme $X$ locally of finite type, there is a bijective correspondence between \emph{closed points of $X$}, \emph{points of $X$ with values in $k$} \sref[I]{I.3.4.4}, and \emph{rational points of $X$ over $k$} \sref[I]{I.3.4.5}, by virtue of \sref[I]{I.6.4.2}.
This shows in particular \sref[I]{I.3.4.3.1} that for two $k$-prealgebraic spaces $X'$, $Y'$, the underlying set of the product $X'\times_kY'$ is identical to the \emph{product set} $X'\times Y'$ of the underlying sets (but of course the topology of the space underlying $X'\times_kY'$ is not the \emph{product topology} of the topologies of the spaces underlying $X'$ and $Y'$; in general, it is strictly finer than that topology).

On the other hand, the local rings $\sh{O}_x$ at the points of a prealgebraic space $X'$ over $k$ are $k$-algebras whose residue field is \emph{isomorphic to $k$}; if $A$ and $B$ are two such $k$-algebras, a $k$-homomorphism $\varphi:A\to B$ is necessarily \emph{local}: indeed, if an element $x$ of the maximal ideal of $A$ were such that $\varphi(x)$ is invertible, there would exist $\lambda\in k$ nonzero such that $\varphi(x-\lambda.1)$ belongs to the maximal ideal of $B$, which is absurd since $x-\lambda.1$ is invertible.
We conclude immediately from this that if $X'=S(X)$, $Y'=S(Y)$ are two prealgebraic spaces over $k$, every morphism $X'\to Y'$ of \emph{$k$-ringed spaces} is also a morphism of $k$-ringed spaces \emph{in local rings} (\sref[I]{I.1.8.2}; cf.\ $\mathrm{Err}_{\mathrm{II}}$). Moreover, with the preceding notation, if $A$ and $B$ are $k$-algebras of finite type, the corresponding homomorphism makes $B$ an $A$-algebra of finite type; hence every morphism $X'\to Y'$ of $k$-ringed spaces is, by \sref{IV.10.9.6}, of the form $S(f)$, where $f:X\to Y$ is a morphism of $k$-preschemes.

Finally, for every open $U$ of $X'$, every section $s\in\Gamma(U,\sh{O}_{X'})$, and every $x\in U$, $s(x)$ \sref[0]{0.5.5.1} is identified with an element of $k$; thus $s$ determines a map $\overline{s}:x\mapsto s(x)$ from $U$ to $k$, in other words a section over $U$ of the sheaf $\sh{A}(X')$ of germs of maps from $X'$ to $k$. Since $h_U:s\mapsto\overline{s}$ is evidently a homomorphism of rings $\Gamma(U,\sh{O}_{X'})\to\Gamma(U,\sh{A}(X'))$ and commutes with restriction to an open $V\subset U$, the $h_U$ define a homomorphism of sheaves of rings $h:\sh{O}_{X'}\to\sh{A}(X')$.
If one takes $U$ to be ultra-affine, $\operatorname{Spm}(A)$, where $A$ is a Jacobson ring, to say that $\overline{s}=0$ means that for every maximal ideal $\mathfrak{m}$ of $A$, $s$ belongs to $\mathfrak{m}$, or equivalently that $s$ is in the \emph{radical} of $A$; but as $A$ is a Jacobson ring, its radical is its nilradical. Consequently $h_U$ is injective if and only if $A$ is reduced.
\end{env}

\begin{env}[10.10.4]
\label{IV.10.10.4}
One says that the $k$-prealgebraic space $X'=S(X)$ is \emph{reduced} if $X$ is reduced. Since the set of points $x\in X$ at which $X$ is reduced is open \sref[0]{0.5.2.2}, its complement, if non-empty, contains at least one closed point \sref{IV.5.1.11}; hence it amounts to the same to say that $X'$ is reduced or that each of its local rings $\sh{O}_x$ ($x\in X'$) is reduced. We have just seen in \sref{IV.10.10.3} that for the homomorphism $h:\sh{O}_{X'}\to\sh{A}(X')$ to be \emph{injective}, it is necessary and sufficient that $X'$ be \emph{reduced}.
In (FAC), Serre restricts himself in fact to prealgebraic \emph{reduced} spaces, which allows him to define $\sh{O}_{X'}$ as a \emph{sub-sheaf} of $\sh{A}(X')$.
One notes that if $X'$ and $Y'$ are two reduced $k$-prealgebraic spaces, the same is true of $X'\times_k Y'$: indeed, this amounts to saying that if $A$ and $B$ are two reduced $k$-algebras of finite type, $A\otimes_k B$ is also reduced; but we have just seen that the radicals of $A$ and $B$ are $0$, and as $k$ is algebraically closed, $A$ and $B$ are ``separable'' algebras over $k$ in the sense of Bourbaki (\emph{Alg.}, chap.~VIII, \S\,7, n$^\circ$\,5, prop.~5); hence $A\otimes_k B$ is without radical (\emph{loc.\ cit.}, n$^\circ$\,6, cor.~3 of th.~3), and since it is a Jacobson ring, it is reduced.
However, if $Z'$ is a third prealgebraic space reduced over $k$, the ``fibre product'' $X'\times_{Z'}Y'$ of two prealgebraic reduced spaces over $Z'$ is not in general reduced, which implies that the category of these spaces is insufficient in many questions (notably in the theory of algebraic groups).
But as we have seen above, one can keep the language of Serre without restricting oneself, as this latter (who in addition considers only quasi-compact \emph{prealgebraic} spaces), to the case of reduced prealgebraic spaces.
\end{env}

\oldpage[IV-3]{116}

\begin{env}[10.10.5]
\label{IV.10.10.5}
Finally, one can also consider ultra-preschemes over an arbitrary field $k$ while
retaining a language close to Serre's, and introducing, as in Weil, a fixed
algebraically closed extension $K$ of $k$ (chosen sufficiently large---for example,
of infinite transcendence degree over $k$---so as to have enough ``generic points''
in Weil's sense).
To every prescheme $X$ locally of finite type over $k$ one then associates the set
\[
 S_K(X)=S(X\otimes_k K)
\]
of points of $X$ with values in $K$.
There is a canonical map $j:S_K(X)\to X$, which is a quasi-homeomorphism when
$S_K(X)$ is given the inverse-image topology from $X$ under $j$; endowing $S_K(X)$
with the sheaf of $k$-algebras $j^*(\sh{O}_X)$ gives a subcategory of the category
of $k$-ringed spaces in local rings, which one might call the category of
\emph{$(k,K)$-prealgebraic spaces}.
Products can still be defined there, and the results of \sref{IV.10.10.3} and
\sref{IV.10.10.4} can be generalized (with $\sh{A}(X')$ replaced here by the sheaf
$\sh{A}_K(X')$ of germs of maps from $X'$ into $K$).
This point of view, however, assigns an artificial role to an arbitrarily chosen
overfield of $k$, and we mention it only to reject it.
\end{env}
% END NATIVE INPUT source_aligned/ega4-10
% BEGIN NATIVE INPUT source_aligned/ega4-11
\setcounter{section}{10}
\section{Topological properties of finitely presented flat morphisms. Flatness criteria}
\label{section:IV.11}

\oldpage[IV-3]{116}

Whereas in \S~2 we considered statements concerning flatness that do not depend
on any finiteness hypothesis, and \S~6 studies the notion of flatness in the
framework of locally Noetherian preschemes (but without a finiteness hypothesis
on the morphisms), the present section is devoted to the notion of
$f$-flatness when the morphism $f:X\to Y$ is \emph{locally of finite
presentation}. The interest of the notion of a flat morphism of finite
presentation lies in the fact that it seems to express technically, in the
most adequate way, the intuitive notion of a ``family of algebraic preschemes
parametrised by a scheme $Y$'', whose study is one of the principal objects of
Algebraic Geometry. Moreover, even if one were initially interested only in
the case of a Noetherian base scheme, it is indispensable, for certain
technical reasons (for example for certain applications of the theory of
``descent'', which leads to the introduction of schemes that are not
necessarily Noetherian), not to restrict oneself to that case when dealing
with essentially relative problems connected with morphisms locally of finite
presentation. We will systematically follow this principle, already supported
by the results of \S\S~8 and 9, throughout the rest of this Chapter and even
throughout the rest of this Treatise, though this may occasionally sacrifice
the simplicity of certain proofs that Noetherian hypotheses sometimes permit
one to lighten.\footnote{This principle is also inspired by the necessity of
granting a place, as ``parameter spaces'' for families of algebraic schemes,
to arbitrary ringed spaces (and even arbitrary ringed ``topoi''), for which
Noetherian hypotheses can no longer in general be contemplated. It seems
fairly clear that this new extension of Algebraic Geometry cannot be avoided
much longer, and it is appropriate from now on to develop notions and
techniques of a ``relative'' nature in the theory of schemes so that they may
be adapted practically unchanged to this more general setting.} In the
present section, this leads us to take up again, in the context of ``finite
presentation'' (notably in n$^\circ$~3), certain statements about flatness
already obtained in the Noetherian context. The essential technical tool for
the reduction to the Noetherian case is the theorem on compatibility of
flatness with projective limits of preschemes \sref{IV.11.2.6}, completing the
general results of \S~8. We also prove in passing \sref{IV.11.3.1} a result
often used

\oldpage[IV-3]{117}
later, implying that the set of flatness points of a morphism locally of finite
presentation is open.

In n$^\circ$s~4 through 8, we study the question of the ``descent'' of
flatness, consisting in finding useful conditions on a base-change morphism
$Y'\to Y$ (not flat in general) under which one may conclude, if
$X\times_YY'$ is flat over $Y'$, that $X$ is flat over $Y$. These results,
more technical than those of n$^\circ$s~1 through 3, are used less frequently
in the sequel; they nevertheless play an important role in the non-projective
construction techniques of the following chapter. The only result of
n$^\circ$s~4 through 8 used later in Chapter IV is the valuative criterion for
flatness \sref{IV.11.8}, which will be applied in \sref{IV.15.2}.

Finally, n$^\circ$s~9 and 10 are devoted to the study of a notion that refines,
in scheme theory, density in the topological sense, namely the notion of a
family of subpreschemes \emph{schematically dense} in a given prescheme, and
in particular the study of the behaviour of this notion under base change
(flat or arbitrary). For the time being, this notion is used above all in the
study of group schemes.


\subsection{Flatness loci (Noetherian case)}
\label{subsection:IV.11.1}

\begin{theorem}[11.1.1]
\label{IV.11.1.1}
Let $\mathrm{Y}$ be a locally Noetherian prescheme, $f:\mathrm{X}\to\mathrm{Y}$ a morphism locally of finite type, $\sh{F}$ an $\sh{O}_\mathrm{X}$-Module coherent. Then the set $\mathrm{U}$ of $x\in\mathrm{X}$ such that $\sh{F}$ is $f$-flat at the point $x$ is open in $\mathrm{X}$.
\end{theorem}

\begin{proof}
The question being evidently local on $\mathrm{X}$ and $\mathrm{Y}$, we can suppose $\mathrm{Y}=\Spec(\mathrm{A})$, $\mathrm{X}=\Spec(\mathrm{B})$, $\mathrm{A}$ being Noetherian and $\mathrm{B}$ an A-algebra of finite type. We then have $\sh{F}=\widetilde{\mathrm{M}}$, where $\mathrm{M}$ is a B-module of finite type. Applying the criterion $(\mathbf{0}_{\mathrm{III}},~9.2.6)$: it suffices therefore to prove the following assertion:

\medskip
\phantomsection\label{IV.11.1.1.1}
\noindent\textbf{(11.1.1.1)}
\emph{Let $\mathrm{A}$ be a Noetherian ring, $\mathrm{B}$ an $\mathrm{A}$-algebra of finite type, $\mathrm{M}$ a $\mathrm{B}$-module of finite type, $\mathfrak{q}$ a prime ideal of $\mathrm{B}$, $\mathfrak{p}$ its inverse image in $\mathrm{A}$. Suppose that $\mathrm{M}_\mathfrak{q}$ is a flat $\mathrm{A}_\mathfrak{p}$-module. Then there exists $g\in\mathrm{B}-\mathfrak{q}$ such that, for every prime ideal $\mathfrak{q}'\supset\mathfrak{q}$ of $\mathrm{B}$ such that $g\notin\mathfrak{q}'$, $\mathrm{M}_{\mathfrak{q}'}$ is a flat $\mathrm{A}_{\mathfrak{p}'}$-module, where $\mathfrak{p}'$ denotes the inverse image of $\mathfrak{q}'$ in $\mathrm{A}$ (it amounts to the same $(\mathbf{0}_\mathrm{I},~6.3.1)$ to say that $\mathrm{M}_{\mathfrak{q}'}$ is a flat $\mathrm{A}$-module).}
\medskip

To this end, regard $\mathrm{B}_{\mathfrak{q}'}$ as an $\mathrm{A}$-algebra; clearly
$\mathfrak{p}\mathrm{B}_{\mathfrak{q}'}\subset
\mathfrak{q}'\mathrm{B}_{\mathfrak{q}'}$. By
$(\mathbf{0}_{\mathrm{III}},~10.2.2)$, the $\mathrm{A}$-module
$\mathrm{M}_{\mathfrak{q}'}$ is flat if and only if
$\mathrm{M}_{\mathfrak{q}'}/\mathfrak{p}\mathrm{M}_{\mathfrak{q}'}$ is a
flat $(\mathrm{A}/\mathfrak{p})$-module and
\[
  \operatorname{Tor}_1^\mathrm{A}
  (\mathrm{M}_{\mathfrak{q}'},\mathrm{A}/\mathfrak{p})=0.
\]
Now $\mathrm{M}_{\mathfrak{q}'}=
\mathrm{M}\otimes_{\mathrm{B}}\mathrm{B}_{\mathfrak{q}'}$, and
$\mathrm{B}_{\mathfrak{q}'}$ is flat over $\mathrm{B}$; hence
\[
  \mathrm{M}_{\mathfrak{q}'}/\mathfrak{p}\mathrm{M}_{\mathfrak{q}'}
    =(\mathrm{M}/\mathfrak{p}\mathrm{M})_{\mathfrak{q}'},
  \qquad
  \operatorname{Tor}_1^\mathrm{A}
  (\mathrm{M}_{\mathfrak{q}'},\mathrm{A}/\mathfrak{p})
    =\bigl(\operatorname{Tor}_1^\mathrm{A}
      (\mathrm{M},\mathrm{A}/\mathfrak{p})\bigr)_{\mathfrak{q}'}.
\]
Here the second equality may be obtained by defining Tor from a projective
resolution of $\mathrm{A}/\mathfrak{p}$. For the same reason, since we require
$g\notin\mathfrak{q}'$, the ring $\mathrm{B}_{\mathfrak{q}'}$ is flat over
$\mathrm{B}_g$, and therefore
\[
  (\mathrm{M}/\mathfrak{p}\mathrm{M})_{\mathfrak{q}'}
    =\bigl((\mathrm{M}/\mathfrak{p}\mathrm{M})_g
      \bigr)_{\mathfrak{q}'\mathrm{B}_g},
  \qquad
  \bigl(\operatorname{Tor}_1^\mathrm{A}
      (\mathrm{M},\mathrm{A}/\mathfrak{p})\bigr)_{\mathfrak{q}'}
    =\bigl(\bigl(\operatorname{Tor}_1^\mathrm{A}
      (\mathrm{M},\mathrm{A}/\mathfrak{p})\bigr)_g
      \bigr)_{\mathfrak{q}'\mathrm{B}_g},
\]
where $\mathrm{M}/\mathfrak{p}\mathrm{M}$ and
$\operatorname{Tor}_1^\mathrm{A}
(\mathrm{M},\mathrm{A}/\mathfrak{p})$ are regarded as
$\mathrm{B}$-modules. In view of $(\mathbf{0}_\mathrm{I},~6.3.2)$, we are
thus reduced to proving the following lemma.

\begin{lemma}[11.1.1.2]
\label{IV.11.1.1.2}
Under the conditions of \emph{{\hyperlink{ega4.c04.sec.11-1-1-1}{(11.1.1.1)}}}, there exists $g\in\mathrm{B}-\mathfrak{q}$ such that:
\begin{enumerate}[label=\emph{\roman*)}]
  \item $(\mathrm{M}/\mathfrak{p}\mathrm{M})_g$ is a flat $(\mathrm{A}/\mathfrak{p})$-module;
  \item $(\operatorname{Tor}_1^\mathrm{A}(\mathrm{M},\mathrm{A}/\mathfrak{p}))_g=0$.
\end{enumerate}
\end{lemma}

\oldpage[IV-3]{118}

By the generic flatness theorem $(6.9.1)$, applied to the integral ring
$\mathrm{A}/\mathfrak{p}$, the finite-type
$(\mathrm{A}/\mathfrak{p})$-algebra
$\mathrm{B}/\mathfrak{p}\mathrm{B}$, and the finite-type
$(\mathrm{B}/\mathfrak{p}\mathrm{B})$-module
$\mathrm{M}/\mathfrak{p}\mathrm{M}$, there exists
$h\in\mathrm{A}-\mathfrak{p}$ such that, if $\bar h$ denotes its canonical
image in $\mathrm{A}/\mathfrak{p}$,
$(\mathrm{M}/\mathfrak{p}\mathrm{M})_{\bar h}$ is a flat
$(\mathrm{A}/\mathfrak{p})$-module. On the other hand,
$\mathrm{M}_{\mathfrak q}$ is flat over $\mathrm{A}_{\mathfrak p}$ and
hence, by $(\mathbf{0}_\mathrm{I},~6.3.1)$, flat over $\mathrm{A}$.
Consequently,
\[
  \operatorname{Tor}_1^\mathrm{A}
  (\mathrm{M}_{\mathfrak q},\mathrm{A}/\mathfrak p)=0,
  \qquad\text{or equivalently}\qquad
  \bigl(\operatorname{Tor}_1^\mathrm{A}
  (\mathrm{M},\mathrm{A}/\mathfrak p)\bigr)_{\mathfrak q}=0.
\]
Since $\mathrm{A}$ and $\mathrm{B}$ are Noetherian,
$\operatorname{Tor}_1^\mathrm{A}
(\mathrm{M},\mathrm{A}/\mathfrak p)$ is a finite-type
$\mathrm{B}$-module. Thus $(\mathbf{0}_\mathrm{I},~5.2.2)$ gives an element
$g\in\mathrm{B}-\mathfrak q$ such that
\[
  \bigl(\operatorname{Tor}_1^\mathrm{A}
  (\mathrm{M},\mathrm{A}/\mathfrak p)\bigr)_g=0.
\]
Moreover,
$(\mathrm{M}/\mathfrak p\mathrm{M})_{\bar h}
=(\mathrm{M}/\mathfrak p\mathrm{M})_h$, where on the right the module is
regarded as an $\mathrm{A}$-module. Since
$(\mathrm{M}/\mathfrak p\mathrm{M})_h$ is also a $\mathrm{B}$-module,
$(\mathrm{M}/\mathfrak p\mathrm{M})_{hg}$ remains flat over
$\mathrm{A}/\mathfrak p$: it is
$(\mathrm{M}/\mathfrak p\mathrm{M})_{\overline{hg}}$, where
$\overline{hg}$ is the canonical image of $hg$ in
$\mathrm{B}/\mathfrak p\mathrm{B}$, and one applies
$(\mathbf{0}_\mathrm{I},~6.3.2)$. Finally,
\[
  \bigl(\operatorname{Tor}_1^\mathrm{A}
  (\mathrm{M},\mathrm{A}/\mathfrak p)\bigr)_{hg}=0,
\]
and $hg\in\mathrm{B}-\mathfrak q$ because
$h\in\mathrm{A}-\mathfrak p$. This proves the lemma and the theorem.
\end{proof}

\begin{corollary}[11.1.2]
\label{IV.11.1.2}
Let $\mathrm{Y}$ be a locally Noetherian prescheme, $f:\mathrm{X}\to\mathrm{Y}$ a morphism locally of finite type, $\sh{F}$, $\sh{F}'$ two $\sh{O}_\mathrm{X}$-Modules coherent, $u:\sh{F}'\to\sh{F}$ a homomorphism of $\sh{O}_\mathrm{X}$-Modules. Suppose that $\sh{F}$ is $f$-flat. Then the set $\mathrm{U}$ of $x\in\mathrm{X}$ such that, setting $y=f(x)$, the homomorphism $u_x\otimes 1:\sh{F}'_x\otimes_{\sh{O}_y}\boldsymbol{k}(y)\to\sh{F}_x\otimes_{\sh{O}_y}\boldsymbol{k}(y)$ is injective, is open in $\mathrm{X}$.
\end{corollary}

\begin{proof}
Let $\sh{N}$ (resp.\ $\sh{P}$) be the kernel (resp.\ cokernel) of $u$.
Apply $(\mathbf{0}_{\mathrm{III}},~10.2.4)$\footnote{Here is a proof of
$(\mathbf{0}_{\mathrm{III}},~10.2.4)$ which was not published in
N.~Bourbaki's \emph{Alg\`ebre commutative}. In view of
$(\mathbf{0}_\mathrm{I},~6.1.2)$, it is enough to show that b) implies a).
Set $\mathrm{P}=\operatorname{Im}(u)$,
$\mathrm{Q}=\operatorname{Coker}(u)$, and
$\mathrm{R}=\operatorname{Ker}(u)$. The composite
$\mathrm{M}\otimes k\to\mathrm{P}\otimes k\to\mathrm{N}\otimes k$ is
injective, while $\mathrm{M}\otimes k\to\mathrm{P}\otimes k$ is
surjective. Hence $\mathrm{P}\otimes k\to\mathrm{N}\otimes k$ is
injective and $\mathrm{M}\otimes k\to\mathrm{P}\otimes k$ is bijective.
The exact sequence $0\to\mathrm{P}\to\mathrm{N}\to\mathrm{Q}\to0$ gives
the exact sequence
\[
0=\operatorname{Tor}_1^\mathrm{A}(\mathrm{N},k)\to\operatorname{Tor}_1^\mathrm{A}(\mathrm{Q},k)\to\mathrm{P}\otimes k\to\mathrm{N}\otimes k\to\mathrm{Q}\otimes k\to 0
\]
and, since $\mathrm{P}\otimes k\to\mathrm{N}\otimes k$ is injective,
$\operatorname{Tor}_1^\mathrm{A}(\mathrm{Q},k)=0$. As $\mathrm{Q}$ is a
finite-type $\mathrm{B}$-module,
$(\mathbf{0}_{\mathrm{III}},~10.2.2)$ shows that it is flat over
$\mathrm{A}$; then $\mathrm{P}$ is also flat over $\mathrm{A}$ by
$(\mathbf{0}_\mathrm{I},~6.1.2)$. From the exact sequence
$0\to\mathrm{R}\to\mathrm{M}\to\mathrm{P}\to0$, the same result gives an
exact sequence
$0\to\mathrm{R}\otimes k\to\mathrm{M}\otimes k\to
\mathrm{P}\otimes k\to0$. Since the last map is bijective,
$\mathrm{R}\otimes k=0$. But $\mathrm{B}$ is Noetherian, so
$\mathrm{R}$ is a finite-type $\mathrm{B}$-module; Nakayama's lemma
therefore gives $\mathrm{R}=0$.}
to the local rings $\sh{O}_y$ and $\sh{O}_x$, and to the
$\sh{O}_x$-modules $\sh{F}'_x$ and $\sh{F}_x$. The injectivity of
$u_x\otimes1$ is equivalent to the two conditions
$\sh{N}_x=0$ and that $\sh{P}$ be $f$-flat at $x$. Since
$\sh{N}$ and $\sh{P}$ are coherent $(\mathbf{0}_\mathrm{I},~5.3.4)$,
the locus where $\sh{N}_x=0$ is open
$(\mathbf{0}_\mathrm{I},~5.2.2)$, and the locus where $\sh{P}$ is
$f$-flat is open by \sref{IV.11.1.1}. The result follows.

In particular:
\end{proof}

\begin{corollary}[11.1.3]
\label{IV.11.1.3}
Let $\mathrm{Y}$ be a locally Noetherian prescheme,
$f:\mathrm{X}\to\mathrm{Y}$ a flat morphism locally of finite type, and
$g$ a section of $\sh{O}_\mathrm{X}$ over $\mathrm{X}$. The set of
$x\in\mathrm{X}$ such that $g_x\otimes 1$ is a non-zero-divisor in
$\sh{O}_x\otimes_{\sh{O}_{f(x)}}\boldsymbol{k}(f(x))$ is open in
$\mathrm{X}$.
\end{corollary}

\begin{proof}
It suffices to apply \sref{IV.11.1.2} to the endomorphism of $\sh{O}_\mathrm{X}$ defined by $g$.
\end{proof}

\begin{corollary}[11.1.4]
\label{IV.11.1.4}
Let $\mathrm{Y}$ be a locally Noetherian prescheme, $f:\mathrm{X}\to\mathrm{Y}$ a morphism locally of finite type, $\sh{F}$ an $\sh{O}_\mathrm{X}$-Module coherent and $f$-flat. Let $(g_i)_{1\leqslant i\leqslant n}$ be a sequence of sections of $\sh{O}_\mathrm{X}$ above $\mathrm{X}$. The set $\mathrm{U}$ of $x\in\mathrm{X}$ such that the sequence $((g_i)_x\otimes 1)$ is $(\sh{F}_x\otimes_{\sh{O}_{f(x)}}\boldsymbol{k}(f(x)))$-regular is open in $\mathrm{X}$.
\end{corollary}

\begin{proof}
\oldpage[IV-3]{119}
Since $\sh{F}$ is $f$-flat, it follows from
$(\mathbf{0},~15.1.16)$ that $\mathrm{U}$ is also the set of points
$x\in\mathrm{X}$ for which the sequence $((g_i)_x)$ is
$\sh{F}_x$-regular and the $\sh{O}_x$-module
\[
  \sh{G}_x,\qquad
  \sh{G}=\sh{F}\big/\Bigl(\sum_{i=1}^n g_i\sh{F}\Bigr),
\]
is $f$-flat. The modules $\sh{F}$ and $\sh{G}$ are coherent, so the
corollary follows from \sref{IV.11.1.1} and
$(\mathbf{0},~15.2.4)$.
\end{proof}

\begin{corollary}[11.1.5]
\label{IV.11.1.5}
Let $\mathrm{X}$, $\mathrm{Y}$, $\mathrm{Z}$ be three locally Noetherian preschemes, $f:\mathrm{X}\to\mathrm{Y}$, $g:\mathrm{Y}\to\mathrm{Z}$ two morphisms of finite type, $\sh{F}$ an $\sh{O}_\mathrm{X}$-Module coherent. Then the set $\mathrm{U}$ of $y\in\mathrm{Y}$ such that, for every generization $y'$ of $y$, $\sh{F}$ is $(g\circ f)$-flat at all the points of $f^{-1}(y')$ (i.e.\ such that $\sh{F}\otimes_\mathrm{Y}\Spec(\sh{O}_{\mathrm{Y},y})$ is flat relative to the morphism $\mathrm{X}\times_\mathrm{Y}\Spec(\sh{O}_{\mathrm{Y},y})\to\mathrm{Z}$) is open in $\mathrm{Y}$.
\end{corollary}

\begin{proof}
If $\mathrm{V}$ is the set of points $x\in\mathrm{X}$ where $\sh{F}$ is
$(g\circ f)$-flat, then $\mathrm{U}$ is the set of $y\in\mathrm{Y}$ such
that $f^{-1}(y')\subset\mathrm{V}$ for every generization $y'$ of $y$.
Now $\mathrm{V}$ is open by \sref{IV.11.1.1}, hence locally
constructible in $\mathrm{X}$, and the set $\mathrm{W}$ of
$y\in\mathrm{Y}$ such that $f^{-1}(y)\subset\mathrm{V}$ is
\[
  \mathrm{W}=\mathrm{Y}-f(\mathrm{X}-\mathrm{V}).
\]
It is therefore locally constructible in $\mathrm{Y}$ by Chevalley's
theorem $(1.8.4)$. It follows from
$(\mathbf{0}_{\mathrm{III}},~9.2.5)$ that the points of $\mathrm{U}$ are
the interior points of $\mathrm{W}$, and hence that $\mathrm{U}$ is open.
\end{proof}

\begin{corollary}[11.1.6]
\label{IV.11.1.6}
Let $\mathrm{A}$ be a Noetherian ring, $\mathrm{B}$ an $\mathrm{A}$-algebra of finite type, $\mathrm{C}$ a $\mathrm{B}$-algebra of finite type, $\mathrm{M}$ a $\mathrm{C}$-module of finite type. Then the set of $\mathfrak{q}\in\Spec(\mathrm{B})$ such that $\mathrm{M}_\mathfrak{q}$ is a flat $\mathrm{A}$-module is open in $\Spec(\mathrm{B})$.
\end{corollary}

\begin{proof}
Taking into account $(2.1.2)$, this is a consequence of \sref{IV.11.1.5} applied to $\mathrm{Z}=\Spec(\mathrm{A})$, $\mathrm{Y}=\Spec(\mathrm{B})$, $\mathrm{X}=\Spec(\mathrm{C})$, $\sh{F}=\widetilde{\mathrm{M}}$.

The results of this number will be freed from the Noetherian hypotheses in \sref{IV.11.3}.
\end{proof}

\subsection{Flatness of a projective limit of preschemes}
\label{subsection:IV.11.2}

\oldpage[IV-3]{119}

\begin{env}[11.2.1]
\label{IV.11.2.1}
Let $\mathrm{A}$ be a ring, $\mathrm{M}$, $\mathrm{N}$ two A-modules, $\mathrm{A}'$ an A-algebra; set $\mathrm{M}'=\mathrm{M}\otimes_\mathrm{A}\mathrm{A}'$, $\mathrm{N}'=\mathrm{N}\otimes_\mathrm{A}\mathrm{A}'$. Recall $(\textbf{III},~6.3.8)$ that one defines, for every $i$, a canonical homomorphism of A-modules
\begin{equation}
  \psi_i:\operatorname{Tor}_i^\mathrm{A}(\mathrm{M},\mathrm{N})\to\operatorname{Tor}_i^{\mathrm{A}'}(\mathrm{M}',\mathrm{N}')
  \label{IV.11.2.1.1}
\tag{11.2.1.1}
\end{equation}
in the following way: one considers a left resolution of $\mathrm{M}$ by free A-modules
\begin{equation}
  \ldots\longrightarrow\mathrm{L}_{i+1}\xrightarrow{f_i}\mathrm{L}_i\longrightarrow\ldots\longrightarrow\mathrm{L}_0\xrightarrow{\varepsilon}\mathrm{M}\longrightarrow 0
  \label{IV.11.2.1.2}
\tag{11.2.1.2}
\end{equation}
whence one deduces by tensoring with $\mathrm{A}'$ a complex of $\mathrm{A}'$-modules
\begin{equation}
  \ldots\longrightarrow\mathrm{L}'_{i+1}\xrightarrow{f'_i}\mathrm{L}'_i\longrightarrow\ldots\longrightarrow\mathrm{L}'_0\xrightarrow{\varepsilon'}\mathrm{M}'\longrightarrow 0
  \label{IV.11.2.1.3}
\tag{11.2.1.3}
\end{equation}
where we have set $\mathrm{L}'_i=\mathrm{L}_i\otimes_\mathrm{A}\mathrm{A}'$, $f'_i=f_i\otimes 1_{\mathrm{A}'}$, $\varepsilon'=\varepsilon\otimes 1_{\mathrm{A}'}$. Consider on the other hand a left resolution of $\mathrm{M}'$ by free $\mathrm{A}'$-modules
\begin{equation}
  \ldots\longrightarrow\mathrm{L}''_{i+1}\xrightarrow{f''_i}\mathrm{L}''_i\longrightarrow\ldots\longrightarrow\mathrm{L}''_0\xrightarrow{\varepsilon''}\mathrm{M}'\longrightarrow 0
  \label{IV.11.2.1.4}
\tag{11.2.1.4}
\end{equation}
\end{env}

\oldpage[IV-3]{120}

As the $\mathrm{L}'_i$ are free $\mathrm{A}'$-modules, we know $(\mathrm{M},~\mathrm{V},~1.1)$ that there are $\mathrm{A}'$-homomorphisms $u_i$ forming a commutative diagram
\[
\xymatrix{
  \ldots \ar[r] & \mathrm{L}'_{i+1} \ar[r]^{f'_i} \ar[d]_{u_{i+1}} & \mathrm{L}'_i \ar[r] \ar[d]_{u_i} & \ldots \ar[r] & \mathrm{L}'_0 \ar[r]^{\varepsilon'} \ar[d]_{u_0} & \mathrm{M}' \ar[r] \ar[d]_{1_{\mathrm{M}'}} & 0 \\
  \ldots \ar[r] & \mathrm{L}''_{i+1} \ar[r]_{f''_i} & \mathrm{L}''_i \ar[r] & \ldots \ar[r] & \mathrm{L}''_0 \ar[r]_{\varepsilon''} & \mathrm{M}' \ar[r] & 0
}
\]

Composing the morphism of complexes
$u_\bullet:\mathrm{L}'_\bullet\to\mathrm{L}''_\bullet$ with the canonical
morphism $\mathrm{L}_\bullet\to\mathrm{L}'_\bullet$ gives a morphism of
complexes of $\mathrm{A}$-modules
$\mathrm{L}_\bullet\to\mathrm{L}''_\bullet$. Since
\[
  (\mathrm{L}_i\otimes_\mathrm{A}\mathrm{N})
    \otimes_\mathrm{A}\mathrm{A}'
  =\mathrm{L}'_i\otimes_{\mathrm{A}'}\mathrm{N}',
\]
we obtain a morphism of complexes
\[
  \mathrm{L}_\bullet\otimes_\mathrm{A}\mathrm{N}
  \longrightarrow
  \mathrm{L}''_\bullet\otimes_{\mathrm{A}'}\mathrm{N}'.
\]
Passing to homology yields the canonical homomorphisms
\eref{IV.11.2.1.1}. Since the $u_i$ are determined up to
homotopy $(\mathrm{M},~\mathrm{V},~1.1)$, these homomorphisms depend
neither on the choice of the $u_i$ nor on the choice of the free
resolutions $\mathrm{L}_\bullet$ and $\mathrm{L}''_\bullet$.

As the $\operatorname{Tor}_i^{\mathrm{A}'}(\mathrm{M}',\mathrm{N}')$ are $\mathrm{A}'$-modules, one deduces canonically from \eref{IV.11.2.1.1} $\mathrm{A}'$-homomorphisms
\begin{equation}
  \varphi_i:\operatorname{Tor}_i^\mathrm{A}(\mathrm{M},\mathrm{N})\otimes_\mathrm{A}\mathrm{A}'\to\operatorname{Tor}_i^{\mathrm{A}'}(\mathrm{M}',\mathrm{N}').
  \label{IV.11.2.1.5}
\tag{11.2.1.5}
\end{equation}

Consider now two ring homomorphisms
\[
  \rho:\mathrm{A}\to\mathrm{A}^{(1)},\quad \sigma:\mathrm{A}^{(1)}\to\mathrm{A}^{(2)},
\]
and their composite $\sigma\circ\rho:\mathrm{A}\to\mathrm{A}^{(2)}$; set $\mathrm{M}^{(j)}=\mathrm{M}\otimes_\mathrm{A}\mathrm{A}^{(j)}$, $\mathrm{N}^{(j)}=\mathrm{N}\otimes_\mathrm{A}\mathrm{A}^{(j)}$ for $j=1,2$. Then the canonical composite homomorphism
\[
  \operatorname{Tor}_i^\mathrm{A}(\mathrm{M},\mathrm{N})
  \longrightarrow
  \operatorname{Tor}_i^{\mathrm{A}^{(1)}}
    (\mathrm{M}^{(1)},\mathrm{N}^{(1)})
  \longrightarrow
  \operatorname{Tor}_i^{\mathrm{A}^{(2)}}
    (\mathrm{M}^{(2)},\mathrm{N}^{(2)})
\]
is the same as the canonical homomorphism deduced from $\sigma\circ\rho$; this follows from the fact that, if $\mathrm{L}_\bullet^{(j)}$ is a free resolution of $\mathrm{M}^{(j)}$, the diagram
\[
\xymatrix{
  \mathrm{L}_\bullet \ar[r] & \mathrm{L}_\bullet\otimes_\mathrm{A}\mathrm{A}^{(1)} \ar[r] \ar[d] & \mathrm{L}_\bullet^{(1)} \ar[d] \\
  & \mathrm{L}_\bullet\otimes_\mathrm{A}\mathrm{A}^{(2)} \ar[r] &
  \mathrm{L}_\bullet^{(1)}\otimes_{\mathrm{A}^{(1)}}\mathrm{A}^{(2)} \ar[r] & \mathrm{L}_\bullet^{(2)}
}
\]
is commutative.

\begin{env}[11.2.2]
\label{IV.11.2.2}
With the notation of \sref{IV.11.2.1}, let
$(\mathrm{A}_\alpha,g_{\beta\alpha})$ be a filtered inductive system of
$\mathrm{A}$-algebras and set
\[
  \mathrm{M}_\alpha=\mathrm{M}\otimes_\mathrm{A}\mathrm{A}_\alpha,
  \qquad
  \mathrm{N}_\alpha=\mathrm{N}\otimes_\mathrm{A}\mathrm{A}_\alpha.
\]
For each $i$, the modules
$\operatorname{Tor}_i^{\mathrm{A}_\alpha}
(\mathrm{M}_\alpha,\mathrm{N}_\alpha)$, with the canonical maps furnished
by \sref{IV.11.2.1}, form an inductive system. Put
\[
  \mathrm{A}'=\varinjlim_\alpha\mathrm{A}_\alpha,\qquad
  \mathrm{M}'=\varinjlim_\alpha\mathrm{M}_\alpha
    =\mathrm{M}\otimes_\mathrm{A}\mathrm{A}',\qquad
  \mathrm{N}'=\varinjlim_\alpha\mathrm{N}_\alpha
    =\mathrm{N}\otimes_\mathrm{A}\mathrm{A}'.
\]
If $g_\alpha:\mathrm{A}_\alpha\to\mathrm{A}'$ denotes the canonical
homomorphism, we obtain canonical maps
\[
  \psi_{i\alpha}:
  \operatorname{Tor}_i^{\mathrm{A}_\alpha}
    (\mathrm{M}_\alpha,\mathrm{N}_\alpha)
  \longrightarrow
  \operatorname{Tor}_i^{\mathrm{A}'}(\mathrm{M}',\mathrm{N}')
\]
which again form an inductive system. We complete the result of
$(\mathrm{M},~\mathrm{V},~9.4^*)$ by showing that
\begin{equation}
  \psi_i=\varinjlim\psi_{i\alpha}:\varinjlim_\alpha\operatorname{Tor}_i^{\mathrm{A}_\alpha}(\mathrm{M}_\alpha,\mathrm{N}_\alpha)\to\operatorname{Tor}_i^{\mathrm{A}'}(\mathrm{M}',\mathrm{N}')
  \label{IV.11.2.2.1}
\tag{11.2.2.1}
\end{equation}
are \emph{isomorphisms of $\mathrm{A}'$-modules}. To this end, we proceed as in $(\mathrm{M},~\mathrm{V},~9.5^*)$, associating to each $\mathrm{M}_\alpha$ its \emph{canonical free resolution}. It amounts (taking into account the exactness of the functor $\varinjlim$) to proving the
\end{env}

\oldpage[IV-3]{121}

\begin{lemma}[11.2.2.2]
\label{IV.11.2.2.2}
Let $(\mathrm{A}_\alpha,g_{\beta\alpha})$ be a filtered inductive system
of rings and $(\mathrm{M}_\alpha,h_{\beta\alpha})$ a filtered inductive
system of sets. Put
$\mathrm{A}'=\varinjlim\mathrm{A}_\alpha$ and
$\mathrm{M}'=\varinjlim\mathrm{M}_\alpha$, and let
$g_\alpha:\mathrm{A}_\alpha\to\mathrm{A}'$ and
$h_\alpha:\mathrm{M}_\alpha\to\mathrm{M}'$ be the canonical maps. For
each $\alpha$, let $\mathrm{F}(\mathrm{M}_\alpha)$ be the
$\mathrm{A}_\alpha$-module of formal linear combinations of elements of
$\mathrm{M}_\alpha$; similarly, let $\mathrm{F}(\mathrm{M}')$ be the
$\mathrm{A}'$-module of formal linear combinations of elements of
$\mathrm{M}'$. If
$h_{\beta\alpha}^0:\mathrm{F}(\mathrm{M}_\alpha)\to
\mathrm{F}(\mathrm{M}_\beta)$ (for $\alpha\leqslant\beta$) and
$h'_\alpha:\mathrm{F}(\mathrm{M}_\alpha)\to\mathrm{F}(\mathrm{M}')$
are the $\mathrm{A}_\alpha$-homomorphisms induced by
$h_{\beta\alpha}$ and $h_\alpha$, respectively, then
$(\mathrm{F}(\mathrm{M}_\alpha),h^0_{\beta\alpha})$ is an inductive
system of $\mathrm{A}_\alpha$-modules and $(h'_\alpha)$ is an inductive
system of homomorphisms. Moreover,
\[
  h''=\varinjlim h'_\alpha:\varinjlim\mathrm{F}(\mathrm{M}_\alpha)\to\mathrm{F}(\mathrm{M}')=\mathrm{F}(\varinjlim\mathrm{M}_\alpha)
\]
\emph{is an isomorphism.}
\end{lemma}

\begin{proof}
For the proof, see Bourbaki, \emph{Alg.}, chap.~II, 3\textsuperscript{e}~\'{e}d., \textsection\,6, n\textsuperscript{o}\,6, cor.\ de la prop.~10.
\end{proof}

\begin{env}[11.2.3]
\label{IV.11.2.3}
Resume the notation of \sref{IV.11.2.1} and consider the case $i=1$.
Set
$\mathrm{T}=\operatorname{Tor}_1^\mathrm{A}(\mathrm{M},\mathrm{N})$,
$\mathrm{T}'=\mathrm{T}\otimes_\mathrm{A}\mathrm{A}'$, and
$\mathrm{T}''=\operatorname{Tor}_1^{\mathrm{A}'}
(\mathrm{M}',\mathrm{N}')$. Then $\psi_1$ is the homomorphism
\[
  \frac{\operatorname{Ker}(f_0\otimes1_\mathrm{N})}
       {\operatorname{Im}(f_1\otimes1_\mathrm{N})}
  \longrightarrow
  \frac{\operatorname{Ker}(f''_0\otimes1_{\mathrm{N}'})}
       {\operatorname{Im}(f''_1\otimes1_{\mathrm{N}'})}
\]
induced on the quotients by the restriction
\[
  \operatorname{Ker}(f_0\otimes1_\mathrm{N})
  \longrightarrow
  \operatorname{Ker}(f''_0\otimes1_{\mathrm{N}'})
\]
of
$u_1\otimes1:\mathrm{L}_1\otimes_\mathrm{A}\mathrm{N}\to
\mathrm{L}''_1\otimes_{\mathrm{A}'}\mathrm{N}'$.

Set $\mathrm{R}=\operatorname{Ker}(\varepsilon)$, $\mathrm{R}''=\operatorname{Ker}(\varepsilon'')$, so that one has the exact sequences
\begin{equation}
  0\longrightarrow\mathrm{R}\xrightarrow{j}\mathrm{L}_0\xrightarrow{\varepsilon}\mathrm{M}\longrightarrow 0
  \quad\text{and}\quad
  0\longrightarrow\mathrm{R}''\xrightarrow{j''}\mathrm{L}''_0\xrightarrow{\varepsilon''}\mathrm{M}'\longrightarrow 0,
\end{equation}
whence one deduces the exact homology sequences
\begin{equation}
  0=\operatorname{Tor}_1^\mathrm{A}(\mathrm{L}_0,\mathrm{N})\to\mathrm{T}\xrightarrow{\partial}\mathrm{R}\otimes_\mathrm{A}\mathrm{N}\xrightarrow{j\otimes 1}\mathrm{L}_0\otimes_\mathrm{A}\mathrm{N}\xrightarrow{\varepsilon\otimes 1}\mathrm{M}\otimes_\mathrm{A}\mathrm{N}\to 0
  \label{IV.11.2.3.1}
\tag{11.2.3.1}
\end{equation}
\begin{equation}
  0=\operatorname{Tor}_1^{\mathrm{A}'}(\mathrm{L}''_0,\mathrm{N}')\to\mathrm{T}''\xrightarrow{\partial''}\mathrm{R}''\otimes_{\mathrm{A}'}\mathrm{N}'\xrightarrow{j''\otimes 1}\mathrm{L}''_0\otimes_{\mathrm{A}'}\mathrm{N}'\xrightarrow{\varepsilon''\otimes 1}\mathrm{M}'\otimes_{\mathrm{A}'}\mathrm{N}'\to 0
  \label{IV.11.2.3.2}
\tag{11.2.3.2}
\end{equation}

On the other hand there is a homomorphism of $\mathrm{A}'$-modules
\[
  v:\mathrm{R}'=\operatorname{Ker}(\varepsilon)\otimes_\mathrm{A}\mathrm{A}'\longrightarrow\operatorname{Ker}(\varepsilon')\xrightarrow{u_0}\operatorname{Ker}(\varepsilon'')=\mathrm{R}''.
\]
\end{env}

\oldpage[IV-3]{122}

Let us show that the diagram
\begin{equation}
\xymatrix{
  & \mathrm{T}' \ar[r]^-{\partial\otimes 1} \ar[d]_{\varphi_1} & \mathrm{R}'\otimes_{\mathrm{A}'}\mathrm{N}' \ar[r] \ar[d]_{v\otimes 1} & \mathrm{L}'_0\otimes_{\mathrm{A}'}\mathrm{N}' \ar[r] \ar[d]_{u_0\otimes 1} & \mathrm{M}'\otimes_{\mathrm{A}'}\mathrm{N}' \ar[r] \ar@{=}[d] & 0 \\
  0 \ar[r] & \mathrm{T}'' \ar[r]_{\partial''} & \mathrm{R}''\otimes_{\mathrm{A}'}\mathrm{N}' \ar[r] & \mathrm{L}''_0\otimes_{\mathrm{A}'}\mathrm{N}' \ar[r] & \mathrm{M}'\otimes_{\mathrm{A}'}\mathrm{N}' \ar[r] & 0
}
  \label{IV.11.2.3.3}
\tag{11.2.3.3}
\end{equation}
is \emph{commutative}. For this, one verifies immediately $(\mathrm{M},~\mathrm{IV},~1)$ that the homomorphism $\partial:\mathrm{T}\to\mathrm{R}\otimes_\mathrm{A}\mathrm{N}$ comes (in the present case) by passage to the quotient from the homomorphism $w:\operatorname{Ker}(f_0\otimes 1_\mathrm{N})\to\mathrm{R}\otimes_\mathrm{A}\mathrm{N}$, restriction of the homomorphism $g_0\otimes 1_\mathrm{N}:\mathrm{L}_1\otimes_\mathrm{A}\mathrm{N}\to\mathrm{R}\otimes_\mathrm{A}\mathrm{N}$, where $g_0:\mathrm{L}_1\to\mathrm{R}$ is such that $f_0=j\circ g_0$; similarly for $\partial''$. It therefore suffices to see that the diagram
\[
\xymatrix{
  \operatorname{Ker}(f_0\otimes 1_\mathrm{N}) \ar[r]^{w} \ar[d] & \mathrm{R}\otimes_\mathrm{A}\mathrm{N} \ar[r] & \mathrm{R}'\otimes_{\mathrm{A}'}\mathrm{N}'=(\mathrm{R}\otimes_\mathrm{A}\mathrm{N})\otimes_\mathrm{A}\mathrm{A}' \ar[d]_{v\otimes 1} \\
  \operatorname{Ker}(f''_0\otimes 1_{\mathrm{N}'}) \ar[rr] & & \mathrm{R}''\otimes_{\mathrm{A}'}\mathrm{N}'
}
\]
is commutative, which follows from the commutativity of the diagram
\[
\xymatrix{
  \mathrm{L}_1 \ar[r] \ar[d] & \mathrm{R} \ar[d] \\
  \mathrm{L}''_1 \ar[r] & \mathrm{R}''
}
\]

\begin{lemma}[11.2.4]
\label{IV.11.2.4}
Let $\mathrm{A}$ be a ring, $\mathrm{C}$ an $\mathrm{A}$-algebra, $\mathrm{M}$ an $\mathrm{A}$-module, $\mathrm{N}$ a $\mathrm{C}$-module, $\mathrm{A}'$ an $\mathrm{A}$-algebra. Set $\mathrm{C}'=\mathrm{C}\otimes_\mathrm{A}\mathrm{A}'$, $\mathrm{M}'=\mathrm{M}\otimes_\mathrm{A}\mathrm{A}'$, $\mathrm{N}'=\mathrm{N}\otimes_\mathrm{A}\mathrm{A}'=\mathrm{N}\otimes_\mathrm{C}\mathrm{C}'$. Suppose that $\mathrm{M}\otimes_\mathrm{A}\mathrm{N}$ is a flat $\mathrm{C}$-module. Then the canonical homomorphism
\[
  \varphi_1:\operatorname{Tor}_1^\mathrm{A}(\mathrm{M},\mathrm{N})\otimes_\mathrm{A}\mathrm{A}'\to\operatorname{Tor}_1^{\mathrm{A}'}(\mathrm{M}',\mathrm{N}')
\]
(cf.\ \eref{IV.11.2.1.5}) \emph{is surjective}.
\end{lemma}

\begin{proof}
Keep the notation of \sref{IV.11.2.3}. Right exactness of tensor
products shows that
\[
  \mathrm{L}'_1\xrightarrow{f'_0}\mathrm{L}'_0
  \xrightarrow{\varepsilon'}\mathrm{M}'\longrightarrow0
\]
is exact. Since $\mathrm{L}'_0$ and $\mathrm{L}'_1$ are free
$\mathrm{A}'$-modules, we may suppose that
$\mathrm{L}''_0=\mathrm{L}'_0$ and
$\mathrm{L}''_1=\mathrm{L}'_1$, that $u_0$ and $u_1$ are the identity
maps, and that $f''_0=f'_0$. Since
$\operatorname{Ker}(\varepsilon')=\operatorname{Im}(f'_0)$ and
$\mathrm{R}''=\operatorname{Im}(f''_0)=\operatorname{Im}(f'_0)$, the
homomorphism $v$ is surjective, and hence so is $v\otimes1$. The proof
will therefore be complete if the first row of
\eref{IV.11.2.3.3} is exact, because $v\otimes1$ is
surjective and $u_0\otimes1$ is injective (Bourbaki,
\emph{Alg.\ comm.}, chap.~I, \S\,1, no.~4, Cor.~2 to Prop.~2). For this
we use

\oldpage[IV-3]{123}

the hypothesis that $\mathrm{M}\otimes_\mathrm{A}\mathrm{N}$ is a flat C-module. Setting $\mathrm{Q}=\operatorname{Ker}(\varepsilon\otimes 1)=\operatorname{Im}(j\otimes 1)$, by the exact sequence \eref{IV.11.2.3.1}, one has the two exact sequences $0\to\mathrm{Q}\to\mathrm{L}_0\otimes_\mathrm{A}\mathrm{N}\to\mathrm{M}\otimes_\mathrm{A}\mathrm{N}\to 0$ and $\mathrm{T}\to\mathrm{R}\otimes_\mathrm{A}\mathrm{N}\to\mathrm{Q}\to 0$, where the homomorphisms are homomorphisms of \emph{C-modules}; using the hypothesis of flatness, one deduces $(\mathbf{0}_\mathrm{I},~6.1.2)$ the exact sequence
\[
  0\to\mathrm{Q}\otimes_\mathrm{C}\mathrm{C}'\to(\mathrm{L}_0\otimes_\mathrm{A}\mathrm{N})\otimes_\mathrm{C}\mathrm{C}'\to(\mathrm{M}\otimes_\mathrm{A}\mathrm{N})\otimes_\mathrm{C}\mathrm{C}'\to 0
\]
and on the other hand, the tensor product being right exact, one has the exact sequence
\[
  \mathrm{T}\otimes_\mathrm{C}\mathrm{C}'\to(\mathrm{R}\otimes_\mathrm{A}\mathrm{N})\otimes_\mathrm{C}\mathrm{C}'\to\mathrm{Q}\otimes_\mathrm{C}\mathrm{C}'\to 0
\]
whence finally the exact sequence
\[
  \mathrm{T}\otimes_\mathrm{C}\mathrm{C}'\to(\mathrm{R}\otimes_\mathrm{A}\mathrm{N})\otimes_\mathrm{C}\mathrm{C}'\to(\mathrm{L}_0\otimes_\mathrm{A}\mathrm{N})\otimes_\mathrm{C}\mathrm{C}'\to(\mathrm{M}\otimes_\mathrm{A}\mathrm{N})\otimes_\mathrm{C}\mathrm{C}'\to 0
\]
But by definition, for every C-module $\mathrm{P}$, $\mathrm{P}\otimes_\mathrm{C}\mathrm{C}'=\mathrm{P}\otimes_\mathrm{A}\mathrm{A}'$, whence the conclusion.
\end{proof}

\begin{lemma}[11.2.5]
\label{IV.11.2.5}
Let $\mathrm{A}$ be a ring, $\mathfrak{J}$ an ideal of $\mathrm{A}$,
$\mathrm{B}$ an $\mathrm{A}$-algebra, $\mathrm{M}$ a
$\mathrm{B}$-module, and $\mathrm{A}\to\mathrm{A}'$ a ring
homomorphism. Set
$\mathfrak{J}'=\mathfrak{J}\mathrm{A}'$,
$\mathrm{B}'=\mathrm{B}\otimes_\mathrm{A}\mathrm{A}'$, and
$\mathrm{M}'=\mathrm{M}\otimes_\mathrm{A}\mathrm{A}'
=\mathrm{M}\otimes_\mathrm{B}\mathrm{B}'$. Let
$\mathfrak{p}'$ be a prime ideal of $\mathrm{B}'$ containing
$\mathfrak{J}'\mathrm{B}'$. Suppose that one of the following
hypotheses holds:
\begin{enumerate}[label=\emph{\alph*)}]
  \item $\mathfrak{J}$ is nilpotent.
  \item $\mathrm{A}'$ is Noetherian, $\mathrm{B}'$ is an $\mathrm{A}'$-algebra of finite type, $\mathrm{M}'$ a $\mathrm{B}'$-module of finite type.
\end{enumerate}
Under these conditions, suppose the two following properties are verified:
\begin{enumerate}[label=\emph{(\roman*)}]
  \item $\mathrm{M}\otimes_\mathrm{A}(\mathrm{A}/\mathfrak{J})$ is a flat $(\mathrm{A}/\mathfrak{J})$-module.
  \item The canonical composite homomorphism
\[
  \operatorname{Tor}_1^\mathrm{A}(\mathrm{M},\mathrm{A}/\mathfrak{J})\xrightarrow{\psi_1}\operatorname{Tor}_1^{\mathrm{A}'}(\mathrm{M}',\mathrm{A}'/\mathfrak{J}')\xrightarrow{\theta}(\operatorname{Tor}_1^{\mathrm{A}'}(\mathrm{M}',\mathrm{A}'/\mathfrak{J}'))_{\mathfrak{p}'}
\]
(where $\psi_1$ is the homomorphism \eref{IV.11.2.1.1} and $\theta$ the canonical homomorphism of a $\mathrm{B}'$-module localised at $\mathfrak{p}'$) is zero.
\end{enumerate}

Then $\mathrm{M}'_{\mathfrak{p}'}$ is a flat $\mathrm{A}'$-module.
\end{lemma}

\begin{proof}
Under hypothesis b), $\mathrm{M}'_{\mathfrak{p}'}$ is a finite-type
$\mathrm{B}'_{\mathfrak{p}'}$-module,
$\mathrm{B}'_{\mathfrak{p}'}$ is a Noetherian
$\mathrm{A}'$-algebra, and
$\mathfrak{J}'\mathrm{B}'_{\mathfrak{p}'}$ is contained in the radical
$\mathfrak{p}'\mathrm{B}'_{\mathfrak{p}'}$ of
$\mathrm{B}'_{\mathfrak{p}'}$. Under hypothesis a),
$\mathfrak{J}'$ is nilpotent. We may therefore apply the flatness
criterion $(\mathbf{0}_{\mathrm{III}},~10.2.1)$ or
$(\mathbf{0}_{\mathrm{III}},~10.2.2)$, according as a) or b) holds.
First,
\[
\begin{split}
  \mathrm{M}'_{\mathfrak{p}'}\otimes_{\mathrm{A}'}
    (\mathrm{A}'/\mathfrak{J}')
  &=(\mathrm{M}'\otimes_{\mathrm{A}'}
    (\mathrm{A}'/\mathfrak{J}'))_{\mathfrak{p}'}\\
  &=\bigl((\mathrm{M}\otimes_\mathrm{A}
    (\mathrm{A}/\mathfrak{J}))
    \otimes_{\mathrm{A}/\mathfrak{J}}
    (\mathrm{A}'/\mathfrak{J}')\bigr)_{\mathfrak{p}'}.
\end{split}
\]
Hypothesis (i) therefore implies, by
$(\mathbf{0}_\mathrm{I},~6.2.1\text{ and }6.3.2)$, that this is a flat
$(\mathrm{A}'/\mathfrak{J}')$-module. It remains to prove that
\[
  \operatorname{Tor}_1^{\mathrm{A}'}
  (\mathrm{M}'_{\mathfrak{p}'},\mathrm{A}'/\mathfrak{J}')=0.
\]
By flatness of $\mathrm{B}'_{\mathfrak{p}'}$ over $\mathrm{B}'$, this
module is
\[
  \bigl(\operatorname{Tor}_1^{\mathrm{A}'}
  (\mathrm{M}',\mathrm{A}'/\mathfrak{J}')\bigr)_{\mathfrak{p}'}.
\]
Hypothesis (ii) says that $\theta\circ\varphi_1$ is zero. On the other
hand, \sref{IV.11.2.4}, applied with
$\mathrm{C}=\mathrm{A}/\mathfrak{J}$ and $\mathrm{N}=\mathrm{C}$,
shows in view of (i) that $\varphi_1$ is surjective
(since $\mathrm{A}'/\mathfrak{J}'=\mathrm{C}\otimes_\mathrm{A}
\mathrm{A}'$). Hence $\theta$ is zero. Since the image under $\theta$
of $\operatorname{Tor}_1^{\mathrm{A}'}
(\mathrm{M}',\mathrm{A}'/\mathfrak{J}')$ generates its localization as
a $\mathrm{B}'_{\mathfrak{p}'}$-module, that localization is zero.
\end{proof}

\begin{theorem}[11.2.6]
\label{IV.11.2.6}
The notations being those of $(8.5.1)$ and $(8.8.1)$, suppose $\mathrm{S}_\alpha$ quasi-compact, $\mathrm{X}_\alpha$ and $\mathrm{Y}_\alpha$ of finite presentation over $\mathrm{S}_\alpha$; let $f_\alpha:\mathrm{X}_\alpha\to\mathrm{Y}_\alpha$ be an $\mathrm{S}_\alpha$-morphism, $\sh{F}_\alpha$ an $\sh{O}_{\mathrm{X}_\alpha}$-Module quasi-coherent of finite presentation.

\oldpage[IV-3]{124}

\begin{enumerate}[label=\emph{(\roman*)}]
  \item \emph{Let $x$ be a point of $\mathrm{X}$, $x_\lambda$ its canonical projection in $\mathrm{X}_\lambda$. For $\sh{F}$ to be $f$-flat at the point $x$, it is necessary and sufficient that there exist $\lambda\geqslant\alpha$ such that $\sh{F}_\lambda$ is $f_\lambda$-flat at the point $x_\lambda$.}
  \item \emph{For $\sh{F}$ to be $f$-flat, it is necessary and sufficient that there exist $\lambda\geqslant\alpha$ such that $\sh{F}_\lambda$ is $f_\lambda$-flat.}
\end{enumerate}
\end{theorem}

\begin{proof}
One may suppose that $\mathrm{S}_\alpha=\mathrm{S}_0$; as $\mathrm{Y}_0$ is of finite presentation over $\mathrm{S}_0$, $\mathrm{Y}_0$ is quasi-compact, and $f_0:\mathrm{X}_0\to\mathrm{Y}_0$ is a morphism of finite presentation $(1.6.2,~(\mathrm{v}))$, hence one can also restrict to the case where $\mathrm{S}_0=\mathrm{Y}_0$. Moreover, by the quasi-compactness of $\mathrm{S}_0$ and of the fact that the set of indices $\mathrm{L}$ is filtered, one can restrict by the same reasoning to the case where $\mathrm{S}_0=\Spec(\mathrm{A}_0)$ is affine. Furthermore $\mathrm{X}_0$ is also quasi-compact, hence by the same reasoning one can also suppose $\mathrm{X}_0=\Spec(\mathrm{B}_0)$ affine; one then has $\sh{F}_0=\widetilde{\mathrm{M}}_0$, where $\mathrm{M}_0$ is a $\mathrm{B}_0$-module of finite presentation, and the statement \sref{IV.11.2.6} is in this case equivalent to the following (by $(\mathbf{0}_\mathrm{I},~6.3.1)$):

\begin{corollary}[11.2.6.1]
\label{IV.11.2.6.1}
Let $\mathrm{A}_0$ be a ring,
$(\mathrm{A}_\lambda)_{\lambda\in\mathrm{L}}$ a filtered inductive
system of $\mathrm{A}_0$-algebras, $\mathrm{B}_0$ an
$\mathrm{A}_0$-algebra of finite presentation, and $\mathrm{M}_0$ a
$\mathrm{B}_0$-module of finite presentation. Put
\[
\begin{gathered}
  \mathrm{B}_\lambda
    =\mathrm{B}_0\otimes_{\mathrm{A}_0}\mathrm{A}_\lambda,\qquad
  \mathrm{M}_\lambda
    =\mathrm{M}_0\otimes_{\mathrm{A}_0}\mathrm{A}_\lambda
    =\mathrm{M}_0\otimes_{\mathrm{B}_0}\mathrm{B}_\lambda,\\
  \mathrm{A}=\varinjlim_\lambda\mathrm{A}_\lambda,\qquad
  \mathrm{B}=\varinjlim_\lambda\mathrm{B}_\lambda
    =\mathrm{B}_0\otimes_{\mathrm{A}_0}\mathrm{A},\\
  \mathrm{M}=\varinjlim_\lambda\mathrm{M}_\lambda
    =\mathrm{M}_0\otimes_{\mathrm{A}_0}\mathrm{A}
    =\mathrm{M}_0\otimes_{\mathrm{B}_0}\mathrm{B}.
\end{gathered}
\]
\begin{enumerate}[label=\emph{(\roman*)}]
  \item \emph{Let $\mathfrak{p}$ be a prime ideal of $\mathrm{B}$, and for every $\lambda$, let $\mathfrak{p}_\lambda$ be its inverse image in $\mathrm{B}_\lambda$. For $\mathrm{M}_\mathfrak{p}$ to be a flat $\mathrm{A}$-module, it is necessary and sufficient that there exist $\lambda$ such that $(\mathrm{M}_\lambda)_{\mathfrak{p}_\lambda}$ be a flat $\mathrm{A}_\lambda$-module.}
  \item \emph{For $\mathrm{M}$ to be a flat $\mathrm{A}$-module, it is necessary and sufficient that there exist $\lambda$ such that $\mathrm{M}_\lambda$ be a flat $\mathrm{A}_\lambda$-module.}
\end{enumerate}
\end{corollary}

One needs only prove that the conditions are necessary $(2.1.4)$. We will proceed in several steps.

\medskip
\noindent\textbf{I)} \emph{Reduction to the case where the
$\mathrm{A}_\lambda$ are Noetherian.} --- By $(8.9.1)$ there exist a
subring $\mathrm{A}'_0\subset\mathrm{A}_0$ which is a finite-type
$\mathbf{Z}$-algebra, a finite-type $\mathrm{A}'_0$-algebra
$\mathrm{B}'_0$, and a finite-type $\mathrm{B}'_0$-module
$\mathrm{M}'_0$, such that
\[
  \mathrm{B}_0=\mathrm{B}'_0\otimes_{\mathrm{A}'_0}\mathrm{A}_0,
  \qquad
  \mathrm{M}_0=\mathrm{M}'_0\otimes_{\mathrm{A}'_0}\mathrm{A}_0.
\]
Since
$\mathrm{B}_\lambda=\mathrm{B}'_0\otimes_{\mathrm{A}'_0}
\mathrm{A}_\lambda$ and
$\mathrm{M}_\lambda=\mathrm{M}'_0\otimes_{\mathrm{A}'_0}
\mathrm{A}_\lambda$, we may replace
$\mathrm{A}_0,\mathrm{B}_0,\mathrm{M}_0$ by
$\mathrm{A}'_0,\mathrm{B}'_0,\mathrm{M}'_0$, regarding the
$\mathrm{A}_\lambda$ as $\mathrm{A}'_0$-algebras. We may therefore
already suppose that $\mathrm{A}_0$ is Noetherian.

Let $\mathrm{H}$ be the set of pairs $(\lambda,\mathrm{C}_\lambda)$,
where $\mathrm{C}_\lambda$ is a finite-type
$\mathrm{A}_0$-subalgebra of $\mathrm{A}_\lambda$. Order $\mathrm{H}$ by
declaring $(\lambda,\mathrm{C}_\lambda)\leqslant
(\mu,\mathrm{D}_\mu)$ if $\lambda\leqslant\mu$ and
$\varphi_{\mu\lambda}(\mathrm{C}_\lambda)\subset\mathrm{D}_\mu$.
This makes $\mathrm{H}$ filtered: two pairs are dominated by choosing
$\nu\geqslant\lambda,\mu$ and taking for $\mathrm{E}_\nu$ the
$\mathrm{A}_0$-subalgebra generated by
$\varphi_{\nu\lambda}(\mathrm{C}_\lambda)$ and
$\varphi_{\nu\mu}(\mathrm{D}_\mu)$. For
$\xi=(\lambda,\mathrm{C}_\lambda)$ put
$\mathrm{A}_\xi=\mathrm{C}_\lambda$; the evident restricted transition
maps form a filtered inductive system of $\mathrm{A}_0$-algebras.
Put
$\mathrm{B}_\xi=\mathrm{B}_0\otimes_{\mathrm{A}_0}\mathrm{A}_\xi$
and
$\mathrm{M}_\xi=\mathrm{M}_0\otimes_{\mathrm{A}_0}\mathrm{A}_\xi$.
The rings $\mathrm{A}_\xi$ are Noetherian, and the double-inductive-limit
formula (Bourbaki, \emph{Alg.}, chap.~II, 3rd ed., \S\,6, no.~4,
Prop.~7) gives
\[
  \varinjlim_{\mathrm H}\mathrm{A}_\xi=\mathrm{A},\qquad
  \varinjlim_{\mathrm H}\mathrm{B}_\xi=\mathrm{B},\qquad
  \varinjlim_{\mathrm H}\mathrm{M}_\xi=\mathrm{M}.
\]
Suppose \sref{IV.11.2.6.1} proved for the system
$(\mathrm{A}_\xi)_{\xi\in\mathrm H}$, and let $\mathfrak p$ be a prime
ideal of $\mathrm B$ such that $\mathrm M_{\mathfrak p}$ is flat over
$\mathrm A$. Then there exists $\xi\in\mathrm H$ such that
\oldpage[IV-3]{125}
$(\mathrm{M}_\xi)_{\mathfrak{p}_\xi}$ is flat over
$\mathrm{A}_\xi$, where $\mathfrak p_\xi$ is the inverse image of
$\mathfrak p$. Write $\xi=(\lambda,\mathrm{C}_\lambda)$. Since
$\mathfrak{p}_\lambda$ is the inverse image of $\mathfrak p$ in
$\mathrm{B}_\lambda$, $\mathfrak p_\xi$ is the inverse image of
$\mathfrak p_\lambda$ in $\mathrm B_\xi$, and
\[
  (\mathrm{M}_\lambda)_{\mathfrak p_\lambda}
  =(\mathrm{M}_\xi\otimes_{\mathrm{C}_\lambda}
    \mathrm{A}_\lambda)_{\mathfrak p_\lambda}
  =\bigl((\mathrm{M}_\xi)_{\mathfrak p_\xi}
    \otimes_{\mathrm{C}_\lambda}\mathrm{A}_\lambda
    \bigr)_{\mathfrak p_\lambda}.
\]
Thus $(\mathrm{M}_\lambda)_{\mathfrak p_\lambda}$ is flat over
$\mathrm{A}_\lambda$ by
$(\mathbf{0}_\mathrm{I},~6.2.1\text{ and }6.3.2)$. Statement (ii) is
handled in the same way. Henceforth we may therefore suppose that all
$\mathrm{A}_\lambda$ are Noetherian, although $\mathrm A$ itself need
not be.

\medskip
\noindent\textbf{II)} \emph{Reduction of the global statement (ii) to the pointwise statement (i).} --- Suppose $\sh{F}$ is $f$-flat. For every $\lambda$, let $\mathrm{U}_\lambda$ be the set of $x_\lambda\in\mathrm{X}_\lambda$ such that $\sh{F}_\lambda$ is $f_\lambda$-flat at the point $x_\lambda$; we know that $\mathrm{U}_\lambda$ is open in $\mathrm{X}_\lambda$, since $\mathrm{S}_\lambda$ is Noetherian and $f_\lambda$ is of finite type \sref{IV.11.1.1}; let $\mathrm{V}_\lambda=v_\lambda^{-1}(\mathrm{U}_\lambda)$ its inverse image in $\mathrm{X}$. By hypothesis, for every $x\in\mathrm{X}$, the projection of $x$ in $\mathrm{X}_\lambda$ belongs to $\mathrm{U}_\lambda$ for some $\lambda$, which means that $\mathrm{X}=\bigcup\mathrm{V}_\lambda$. Moreover, by $(2.1.4)$, for $\lambda\leqslant\mu$, $\mathrm{V}_\lambda\subset\mathrm{V}_\mu$, hence, since $\mathrm{X}$ is quasi-compact, there exists an index $\mu$ such that $\mathrm{X}=\mathrm{V}_\mu$. As the $\mathrm{X}_\lambda$ are quasi-compact, it follows from $(8.3.4)$ that there exists a $\nu\geqslant\mu$ such that $v_{\mu\nu}^{-1}(\mathrm{U}_\mu)=\mathrm{X}_\nu$; but by $(2.1.4)$, this implies that $\sh{F}_\nu$ is $f_\nu$-flat.

\medskip
\noindent\textbf{III)} \emph{End of the proof.} --- It remains to prove
(i), with $\mathrm S_0$ affine and Noetherian. Let $y_0$ be the
projection of $x$ in $\mathrm S_0$. By $(2.1.4)$ and
$(\mathbf I,~3.6.5)$ we may replace $\mathrm S_0$ by
$\Spec(\sh O_{\mathrm S_0,y_0})$; thus we may assume that
$\mathrm A_0$ is a Noetherian local ring with maximal ideal
$\mathfrak m_0=\mathfrak j_{y_0}$. By definition,
$\mathfrak m_0$ is the inverse image of the prime ideal
$\mathfrak p=\mathfrak j_x$ of $\mathrm B$, and
$\mathrm M_{\mathfrak p}$ is flat over $\mathrm A$. Hence
\[
  \operatorname{Tor}_1^\mathrm A
  (\mathrm A/\mathfrak m_0\mathrm A,\mathrm M_{\mathfrak p})=0.
\]
Since the Tor modules in question are $\mathrm B$-modules and
$\mathrm B_{\mathfrak p}$ is flat over $\mathrm B$, this is equivalent
to
\[
  \bigl(\operatorname{Tor}_1^\mathrm A
  (\mathrm A/\mathfrak m_0\mathrm A,\mathrm M)\bigr)_{\mathfrak p}=0.
\]

The $\mathrm B_0$-module
$\operatorname{Tor}_1^{\mathrm A_0}
(\mathrm A_0/\mathfrak m_0,\mathrm M_0)$ is of finite type. Indeed, it
may be computed from a resolution of $\mathrm A_0/\mathfrak m_0$ by
finite free $\mathrm A_0$-modules; tensoring with $\mathrm M_0$ gives a
complex of finite-type $\mathrm B_0$-modules, whose homology is again
of finite type because $\mathrm B_0$ is Noetherian. Choose generators
$t_i^0$ $(1\leqslant i\leqslant m)$ of this module, and let $t_i$ be
their canonical images in
$\operatorname{Tor}_1^\mathrm A
(\mathrm A/\mathfrak m_0\mathrm A,\mathrm M)$. The preceding
localization equality gives an element
$h\in\mathrm B-\mathfrak p$ such that $ht_i=0$ for all $i$.

By \eref{IV.11.2.2.1},
\[
  \operatorname{Tor}_1^\mathrm A
  (\mathrm A/\mathfrak m_0\mathrm A,\mathrm M)
  =\varinjlim_\lambda
  \operatorname{Tor}_1^{\mathrm A_\lambda}
  (\mathrm A_\lambda/\mathfrak m_0\mathrm A_\lambda,
   \mathrm M_\lambda).
\]
Consequently, for some $\lambda$ there is an element
$h_\lambda\in\mathrm B_\lambda$ mapping to $h$ such that
$h_\lambda t_i^\lambda=0$ for every $i$, where $t_i^\lambda$ is the
image of $t_i^0$. If $\mathfrak p_\lambda$ is the inverse image of
$\mathfrak p$ in $\mathrm B_\lambda$, then
$h_\lambda\notin\mathfrak p_\lambda$. Thus the images of all the
$t_i^0$ in
\[
  \bigl(\operatorname{Tor}_1^{\mathrm A_\lambda}
  (\mathrm A_\lambda/\mathfrak m_0\mathrm A_\lambda,
   \mathrm M_\lambda)\bigr)_{\mathfrak p_\lambda}
\]
vanish, and the canonical map from
$\operatorname{Tor}_1^{\mathrm A_0}
(\mathrm A_0/\mathfrak m_0,\mathrm M_0)$ to this localization is zero.
The hypotheses of \sref{IV.11.2.5} are therefore satisfied; condition
(i) there is automatic because $\mathrm A_0/\mathfrak m_0$ is a field.
It follows that $(\mathrm M_\lambda)_{\mathfrak p_\lambda}$ is flat
over $\mathrm A_\lambda$, completing the proof.

\begin{corollary}[11.2.7]
\label{IV.11.2.7}
Let $\mathrm{S}=\Spec(\mathrm{A})$ be an affine scheme, $f:\mathrm{X}\to\mathrm{S}$ a morphism, $\sh{F}$ an $\sh{O}_\mathrm{X}$-Module quasi-coherent, $x$ a point of $\mathrm{X}$. The following conditions are equivalent:
\begin{enumerate}[label=\emph{\alph*)}]
  \item $f$ is a morphism of finite presentation, $\sh{F}$ is an $\sh{O}_\mathrm{X}$-Module of finite presentation and $\sh{F}$ is $f$-flat at the point $x$ (resp.\ $f$-flat).
\oldpage[IV-3]{126}
  \item There exists an affine Noetherian scheme $\mathrm{S}_0=\Spec(\mathrm{A}_0)$, a morphism of finite type $f_0:\mathrm{X}_0\to\mathrm{S}_0$, an $\sh{O}_{\mathrm{X}_0}$-Module coherent $\sh{F}_0$, a morphism $\mathrm{S}\to\mathrm{S}_0$ such that the S-prescheme $\mathrm{X}_0\otimes_{\mathrm{S}_0}\mathrm{S}$ is S-isomorphic to $\mathrm{X}$ and that $\sh{F}$ is isomorphic to $\sh{F}_0\otimes_{\sh{O}_{\mathrm{S}_0}}\sh{O}_\mathrm{S}$, and $\sh{F}_0$ is $f_0$-flat at the point $x_0$ the projection of $x$ in $\mathrm{X}_0$ (resp.\ $f_0$-flat).
  \item The conditions of $b)$ are satisfied and in addition $\mathrm{A}_0$ is a sub-$\mathbf{Z}$-algebra of finite type of $\mathrm{A}$, the morphism $\mathrm{S}\to\mathrm{S}_0$ corresponding to the canonical injection $\mathrm{A}_0\to\mathrm{A}$.
\end{enumerate}
\end{corollary}

\begin{proof}
It is clear that $c)$ implies $b)$ and that $b)$ implies $a)$ by
$(2.1.4)$. On the other hand, we may regard $\mathrm{A}$ as the filtered
direct limit of its finite-type $\mathbf{Z}$-subalgebras, and by $(8.9.1)$
there exist such a subalgebra $\mathrm{A}_0$ and a finite-type morphism
$f_0:\mathrm{X}_0\to\mathrm{S}_0$ such that $\mathrm{X}$ is
$\mathrm{S}$-isomorphic to
$\mathrm{X}_0\times_{\mathrm{S}_0}\mathrm{S}$ and $\sh{F}$ is isomorphic
to $\sh{F}_0\otimes_{\sh{O}_{\mathrm{S}_0}}\sh{O}_{\mathrm{S}}$. We may
therefore write
\[
  \mathrm{X}=\varprojlim_\lambda \mathrm{X}_\lambda,\qquad
  \mathrm{X}_\lambda=\mathrm{X}_0\otimes_{\mathrm{A}_0}\mathrm{A}_\lambda,
\]
where $\mathrm{A}_\lambda$ runs through the finite-type
$\mathbf{Z}$-subalgebras of $\mathrm{A}$ containing $\mathrm{A}_0$, and
$\sh{F}=v_\lambda^*(\sh{F}_\lambda)$, where
$v_\lambda:\mathrm{X}\to\mathrm{X}_\lambda$ is the canonical projection
and
$\sh{F}_\lambda=\sh{F}_0\otimes_{\mathrm{A}_0}\mathrm{A}_\lambda$. It
follows from \sref{IV.11.2.6} that there exists $\lambda$ such that
$\sh{F}_\lambda$ is $f_\lambda$-flat at
$x_\lambda=v_\lambda(x)$ (resp.\ $f_\lambda$-flat); the subring
$\mathrm{A}_\lambda$ of $\mathrm{A}$ then satisfies condition $c)$.
\end{proof}

\begin{proposition}[11.2.8]
\label{IV.11.2.8}
Let $g:\mathrm{X}\to\mathrm{S}$, $h:\mathrm{Y}\to\mathrm{S}$ be two morphisms of finite presentation, $f:\mathrm{X}\to\mathrm{Y}$ an S-morphism, $\sh{F}$ an $\sh{O}_\mathrm{X}$-Module quasi-coherent of finite presentation; for every $s\in\mathrm{S}$, let $\mathrm{X}_s=g^{-1}(s)=\mathrm{X}\times_\mathrm{S}\Spec(\boldsymbol{k}(s))$, $\mathrm{Y}_s=h^{-1}(s)=\mathrm{Y}\times_\mathrm{S}\Spec(\boldsymbol{k}(s))$, $f_s$ the morphism $f\times 1:\mathrm{X}_s\to\mathrm{Y}_s$, $\sh{F}_s=\sh{F}\otimes_{\sh{O}_\mathrm{S}}\boldsymbol{k}(s)$. Then the set $\mathrm{E}$ of $s\in\mathrm{S}$ such that $\sh{F}_s$ is $f_s$-flat is locally constructible.
\end{proposition}

\begin{proof}
The property whose constructibility we wish to prove satisfies condition
$(9.2.1,~(\mathrm{i}))$, by $(2.2.13)$ and $(2.5.1)$. In view of
$(9.2.3)$, we may therefore restrict to the case where $\mathrm{S}$ is
affine, Noetherian, and integral, with generic point $\eta$, and it remains
to prove that $\mathrm{E}$ or $\mathrm{S}-\mathrm{E}$ is a neighbourhood
of $\eta$ in $\mathrm{S}$. If
$\eta\in\mathrm{S}-\mathrm{E}$, this follows immediately from Lemma
$(9.4.7.1)$. If, on the contrary, $\eta\in\mathrm{E}$, it follows first
from \sref{IV.11.2.6}, applied by the method of $(8.1.2,~a))$, that there
is an open neighbourhood $\mathrm{U}$ of $\eta$ in $\mathrm{S}$ such that
$\sh{F}|_{g^{-1}(\mathrm{U})}$ is flat over $h^{-1}(\mathrm{U})$; \emph{a
fortiori}, by $(2.1.4)$, $\sh{F}_s$ is $f_s$-flat for every
$s\in\mathrm{U}$, which completes the proof.
\end{proof}

The following theorem generalizes $(11.2.6.1,~(\mathrm{ii}))$.

\begin{theorem}[11.2.9]
\label{IV.11.2.9}
(Raynaud). --- Let $\mathrm{A}_0$ be a ring, $(\mathrm{A}_\lambda)_{\lambda\in\mathrm{L}}$ a filtered inductive system of $\mathrm{A}_0$-algebras, $\mathrm{B}_0$ an $\mathrm{A}_0$-algebra of finite presentation, $\mathfrak{J}_0$ a finite-type ideal of $\mathrm{B}_0$, and $\mathrm{M}_0$ a finitely presented $\mathrm{B}_0$-module. Set $\mathrm{B}_\lambda=\mathrm{B}_0\otimes_{\mathrm{A}_0}\mathrm{A}_\lambda$, $\mathfrak{J}_\lambda=\mathfrak{J}_0\mathrm{B}_\lambda$, $\mathrm{M}_\lambda=\mathrm{M}_0\otimes_{\mathrm{A}_0}\mathrm{A}_\lambda=\mathrm{M}_0\otimes_{\mathrm{B}_0}\mathrm{B}_\lambda$, $\mathrm{A}=\varinjlim\mathrm{A}_\lambda$, $\mathrm{B}=\varinjlim\mathrm{B}_\lambda=\mathrm{B}_0\otimes_{\mathrm{A}_0}\mathrm{A}$, $\mathfrak{J}=\mathfrak{J}_0\mathrm{B}$, and $\mathrm{M}=\varinjlim\mathrm{M}_\lambda=\mathrm{M}_0\otimes_{\mathrm{A}_0}\mathrm{A}=\mathrm{M}_0\otimes_{\mathrm{B}_0}\mathrm{B}$. For $\operatorname{gr}^*_\mathfrak{J}(\mathrm{M})$ to be a flat $\mathrm{A}$-module, it is necessary and sufficient that there exist $\lambda$ such that $\operatorname{gr}^*_{\mathfrak{J}_\lambda}(\mathrm{M}_\lambda)$ is a flat $\mathrm{A}_\lambda$-module; the canonical homomorphisms
\begin{equation}
  \begin{cases}
    \operatorname{gr}^*_{\mathfrak{J}_\lambda}(\mathrm{M}_\lambda)\otimes_{\mathrm{A}_\lambda}\mathrm{A} \to \operatorname{gr}^*_\mathfrak{J}(\mathrm{M}) \\
    \operatorname{gr}^*_{\mathfrak{J}_\lambda}(\mathrm{M}_\lambda)\otimes_{\mathrm{A}_\lambda}\mathrm{A}_\mu \to \operatorname{gr}^*_{\mathfrak{J}_\mu}(\mathrm{M}_\mu)
  \end{cases}
  \qquad\text{for }\lambda\leqslant\mu
  \label{IV.11.2.9.1}
\tag{11.2.9.1}
\end{equation}
are then bijective and $\operatorname{gr}^*_{\mathfrak{J}_\mu}(\mathrm{M}_\mu)$ is a flat $\mathrm{A}_\mu$-module for $\lambda\leqslant\mu$.
\end{theorem}

\medskip
\noindent\emph{Proof.}
Let us first show that the conditions are sufficient, which amounts to proving
the bijectivity of $(11.2.9.1)$. This follows from the following lemma:

\begin{lemma}[11.2.9.2]
\label{IV.11.2.9.2}
Let $\mathrm{A}$ be a ring, $\mathrm{B}$ an $\mathrm{A}$-algebra, $\mathrm{M}$ a $\mathrm{B}$-module, $\mathfrak{J}$ an ideal of $\mathrm{B}$. Let $\mathrm{A}'$ be an $\mathrm{A}$-algebra; set $\mathrm{B}'=\mathrm{B}\otimes_\mathrm{A}\mathrm{A}'$, $\mathrm{M}'=\mathrm{M}\otimes_\mathrm{A}\mathrm{A}'=\mathrm{M}\otimes_\mathrm{B}\mathrm{B}'$, $\mathfrak{J}'=\mathfrak{J}\mathrm{B}'$. If $\operatorname{gr}^*_\mathfrak{J}(\mathrm{M})$ is a flat $\mathrm{A}$-module, the canonical homomorphism
\[
  (\operatorname{gr}^*_\mathfrak{J}(\mathrm{M}))\otimes_\mathrm{A}\mathrm{A}'\to\operatorname{gr}^*_{\mathfrak{J}'}(\mathrm{M}')
\]
\emph{is bijective.}
\end{lemma}

\oldpage[IV-3]{127}

Indeed, by induction on $k$, the hypothesis that the $\mathfrak{J}^k\mathrm{M}/\mathfrak{J}^{k+1}\mathrm{M}$ are flat A-modules implies first, by $(\mathbf{0}_\mathrm{I},~6.1.2)$, that $\mathrm{M}/\mathfrak{J}^{k+1}\mathrm{M}$ is a flat A-module; moreover, by $(\mathbf{0}_\mathrm{I},~6.1.2)$, the sequence
\[
  0\to(\mathfrak{J}^{k+1}\mathrm{M})\otimes_\mathrm{A}\mathrm{A}'\to\mathrm{M}\otimes_\mathrm{A}\mathrm{A}'\to(\mathrm{M}/\mathfrak{J}^{k+1}\mathrm{M})\otimes_\mathrm{A}\mathrm{A}'\to 0
\]
is then exact, in other words $(\mathfrak{J}^{k+1}\mathrm{M})\otimes_\mathrm{A}\mathrm{A}'$ identifies with its canonical image $\mathfrak{J}'^{k+1}\mathrm{M}'$ in $\mathrm{M}'=\mathrm{M}\otimes_\mathrm{A}\mathrm{A}'$. On the other hand, always by $(\mathbf{0}_\mathrm{I},~6.1.2)$, the sequence
\[
  0\to(\mathfrak{J}^{k+1}\mathrm{M})\otimes_\mathrm{A}\mathrm{A}'\to(\mathfrak{J}^k\mathrm{M})\otimes_\mathrm{A}\mathrm{A}'\to(\mathfrak{J}^k\mathrm{M}/\mathfrak{J}^{k+1}\mathrm{M})\otimes_\mathrm{A}\mathrm{A}'\to 0
\]
is exact, which proves the lemma.

To prove the necessity of the conditions of \sref{IV.11.2.9}, we proceed in several steps.

\medskip
\noindent\textbf{I)} \emph{Reduction to the case where the $\mathrm{A}_\lambda$
are Noetherian.} --- One proceeds as in reduction I) of
\sref{IV.11.2.6.1}, retaining the notation; it is enough to begin by replacing
$\mathrm{A}'_0$ by a finite-type sub-$\mathrm{A}'_0$-algebra
$\mathrm{A}''_0$ of $\mathrm{A}_0$ such that, if one sets
$\mathrm{B}''_0=\mathrm{B}'_0\otimes_{\mathrm{A}'_0}\mathrm{A}''_0$, there is
in $\mathrm{B}''_0$ a finitely generated ideal $\mathfrak{J}''_0$ such that
$\mathfrak{J}''_0\mathrm{B}_0=\mathfrak{J}_0$. For this, consider the
finite-type sub-$\mathrm{A}'_0$-algebras $\mathrm{A}'_\alpha$ of
$\mathrm{A}_0$, which form a filtered family; one has
$\mathrm{B}_0=\varinjlim\mathrm{B}'_\alpha$, where
$\mathrm{B}'_\alpha=\mathrm{B}'_0\otimes_{\mathrm{A}'_0}\mathrm{A}'_\alpha$.
Hence there is an index $\beta$ such that a finite set of generators of
$\mathfrak{J}_0$ consists of the images in $\mathrm{B}_0$ of elements of
$\mathrm{B}'_\beta$. Take $\mathrm{A}''_0=\mathrm{A}'_\beta$,
$\mathrm{B}''_0=\mathrm{B}'_\beta$, and take for $\mathfrak{J}''_0$ the ideal
generated by these elements. We may therefore suppose that $\mathrm{A}_0$
(and hence also $\mathrm{B}_0$) is Noetherian. One then defines, as in
\emph{loc.\ cit.}, the filtered set $\mathrm{H}$ and the $\mathrm{A}_\xi$,
$\mathrm{B}_\xi$, $\mathrm{M}_\xi$ for $\xi\in\mathrm{H}$; set also
$\mathfrak{J}_\xi=\mathfrak{J}''_0\mathrm{B}_\xi$ for all
$\xi\in\mathrm{H}$. Suppose that one has shown that there exists
$\xi=(\lambda,\mathrm{C}_\lambda)$ such that
$\operatorname{gr}^*_{\mathfrak{J}_\xi}(\mathrm{M}_\xi)$ is a flat
$\mathrm{C}_\lambda$-module. Since
$\mathrm{M}_\lambda=\mathrm{M}_\xi\otimes_{\mathrm{C}_\lambda}
\mathrm{A}_\lambda$, it follows from \sref{IV.11.2.9.2} that
$\operatorname{gr}^*_{\mathfrak{J}_\lambda}(\mathrm{M}_\lambda)$ is a flat
$\mathrm{A}_\lambda$-module.

\medskip
\noindent\textbf{II)} \emph{Preliminary lemmas.}

\begin{lemma}[11.2.9.3]
\label{IV.11.2.9.3}
Let $\mathrm{A}$ be a ring, let
$\mathrm{B}=\bigoplus_{i\geqslant 0}\mathrm{B}_i$ be an
$\mathrm{A}$-algebra graded in nonnegative degrees, and let
$\mathrm{M}=\bigoplus_{i\in\mathbf{Z}}\mathrm{M}_i$ be a graded
$\mathrm{B}$-module. Suppose that $\mathrm{B}_0$ is a local ring and
that each $\mathrm{M}_i$ is a finite-type $\mathrm{B}_0$-module. For
$\mathrm{M}$ to be a finite-type $\mathrm{B}$-module, it is necessary
and sufficient that, if $\mathfrak{q}$ is the inverse image in
$\mathrm{A}$ of the maximal ideal $\mathfrak{m}$ of $\mathrm{B}_0$,
$\mathrm{M}\otimes_\mathrm{A}\boldsymbol{k}(\mathfrak{q})$ be a
finite-type
$(\mathrm{B}\otimes_\mathrm{A}\boldsymbol{k}(\mathfrak{q}))$-module.
\end{lemma}

\oldpage[IV-3]{128}

\begin{proof}
The condition is evidently necessary; let us prove that it is
sufficient. Since $\mathrm{B}_0$ (and hence $\mathrm{B}$) is an
$\mathrm{A}_{\mathfrak{q}}$-algebra, we may replace $\mathrm{A}$ by
$\mathrm{A}_{\mathfrak{q}}$; in other words, we may suppose that
$\mathrm{A}$ is also a local ring with maximal ideal
$\mathfrak{q}$. By hypothesis, there is an integer $\mathrm{N}$ such
that the canonical homomorphism
\[
 \left(\bigoplus_{-\mathrm{N}\leqslant i\leqslant\mathrm{N}}
   \mathrm{M}_i\otimes_{\mathrm{B}_0}\mathrm{B}\right)
   \otimes_\mathrm{A}\boldsymbol{k}(\mathfrak{q})
 \longrightarrow
 \mathrm{M}\otimes_\mathrm{A}\boldsymbol{k}(\mathfrak{q})
\]
is surjective. This also means that, for every integer $n$, the
canonical homomorphism of $\boldsymbol{k}(\mathfrak{q})$-modules
\[
 \left(\bigoplus_{-\mathrm{N}\leqslant i\leqslant\mathrm{N}}
   \mathrm{M}_i\otimes_{\mathrm{B}_0}\mathrm{B}_{n-i}\right)
   \otimes_\mathrm{A}\boldsymbol{k}(\mathfrak{q})
 \longrightarrow
 \mathrm{M}_n\otimes_\mathrm{A}\boldsymbol{k}(\mathfrak{q})
\]
is surjective. Now $\mathrm{M}_n$ is a finite-type
$\mathrm{B}_0$-module, $\mathrm{B}_0$ is local, and
$\mathfrak{q}\mathrm{B}_0\subset\mathfrak{m}$; Nakayama's lemma
therefore shows that each canonical homomorphism
\[
 \bigoplus_{-\mathrm{N}\leqslant i\leqslant\mathrm{N}}
   \mathrm{M}_i\otimes_{\mathrm{B}_0}\mathrm{B}_{n-i}
 \longrightarrow \mathrm{M}_n
\]
is surjective, which proves the assertion.
\end{proof}

\begin{corollary}[11.2.9.4]
\label{IV.11.2.9.4}
Under the hypotheses of \sref{IV.11.2.9.3}, suppose moreover that each of the $\mathrm{B}_i$ and each $\mathrm{M}_i$ is a finitely presented $\mathrm{B}_0$-module, and that $\mathrm{M}$ is a flat $\mathrm{A}$-module. For $\mathrm{M}$ to be a finitely presented $\mathrm{B}$-module, it is necessary and sufficient that $\mathrm{M}\otimes_\mathrm{A}\boldsymbol{k}(\mathfrak{q})$ be a finitely presented $(\mathrm{B}\otimes_\mathrm{A}\boldsymbol{k}(\mathfrak{q}))$-module.
\end{corollary}

\begin{proof}
By \sref{IV.11.2.9.3}, if the condition in the statement holds,
$\mathrm{M}$ is a finite-type $\mathrm{B}$-module. Hence there is a
finite-type free graded $\mathrm{B}$-module $\mathrm{L}$ (thus having
a finite basis of homogeneous elements) and a surjective graded
homomorphism of degree $0$,
\[
 u:\mathrm{L}\longrightarrow\mathrm{M}.
\]
Let $\mathrm{R}=\operatorname{Ker}(u)$, which is a graded
$\mathrm{B}$-module. There are finitely many integers $m_j$
($1\leqslant j\leqslant r$) such that, for every integer $i$,
$\mathrm{R}_i$ is the kernel of a surjective homomorphism
\[
 \bigoplus_{1\leqslant j\leqslant r}\mathrm{B}_{i+m_j}
 \longrightarrow \mathrm{M}_i.
\]
The hypotheses on the $\mathrm{M}_i$ and the $\mathrm{B}_i$ then
imply that $\mathrm{R}_i$ is a finite-type
$\mathrm{B}_0$-module (Bourbaki, \emph{Alg.\ comm.}, chap.~I,
\textsection\,2, n\textsuperscript{o}\,8, lemma~9). To prove that
$\mathrm{R}$ is a finite-type $\mathrm{B}$-module, observe that,
by the flatness hypothesis and $(\mathbf{0}_\mathrm{I},~6.1.2)$, the
sequence
\[
 0\longrightarrow
 \mathrm{R}\otimes_\mathrm{A}\boldsymbol{k}(\mathfrak{q})
 \longrightarrow
 \mathrm{L}\otimes_\mathrm{A}\boldsymbol{k}(\mathfrak{q})
 \longrightarrow
 \mathrm{M}\otimes_\mathrm{A}\boldsymbol{k}(\mathfrak{q})
 \longrightarrow 0
\]
is exact. The hypothesis therefore implies (Bourbaki,
\emph{loc.\ cit.}) that
$\mathrm{R}\otimes_\mathrm{A}\boldsymbol{k}(\mathfrak{q})$ is a
finite-type
$(\mathrm{B}\otimes_\mathrm{A}\boldsymbol{k}(\mathfrak{q}))$-module.
It is therefore enough to apply \sref{IV.11.2.9.3} to $\mathrm{R}$.
\end{proof}

\medskip
\noindent\textbf{III)} \emph{Reduction to the case where the homomorphisms of transitions $\varphi_{\mu\lambda}:\mathrm{A}_\lambda\to\mathrm{A}_\mu$ (for $\lambda\leqslant\mu$) are injective.} --- Let $\mathrm{A}'_\lambda$ be the image of $\mathrm{A}_\lambda$ by the canonical homomorphism $\varphi_\lambda:\mathrm{A}_\lambda\to\mathrm{A}$; it is clear that the $\mathrm{A}'_\lambda$ form an inductive system of sub-rings \emph{Noetherian} of $\mathrm{A}$, and the homomorphisms of transition $\mathrm{A}'_\lambda\to\mathrm{A}'_\mu$ (for $\lambda\leqslant\mu$) are injective. Set $\mathrm{B}'_\lambda=\mathrm{B}_0\otimes_{\mathrm{A}_0}\mathrm{A}'_\lambda$, $\mathfrak{J}'_\lambda=\mathfrak{J}_0\mathrm{B}'_\lambda$, $\mathrm{M}'_\lambda=\mathrm{M}_0\otimes_{\mathrm{A}_0}\mathrm{A}'_\lambda=\mathrm{M}_\lambda\otimes_{\mathrm{A}_\lambda}\mathrm{A}'_\lambda$; one has still $\mathrm{B}=\varinjlim\mathrm{B}'_\lambda$, $\mathrm{M}=\varinjlim\mathrm{M}'_\lambda$. Suppose one has proved that there exists $\lambda$ such that, for $\mu\geqslant\lambda$, $\operatorname{gr}^*_{\mathfrak{J}'_\mu}(\mathrm{M}'_\mu)$ is a flat $\mathrm{A}'_\mu$-module. Let $a_\lambda$ be the kernel of the homomorphism $\varphi_\lambda$, which is an ideal of the Noetherian ring $\mathrm{A}_\lambda$. It follows from the definition of the inductive limit that there exists an index $\mu\geqslant\lambda$ such that $\varphi_{\mu\lambda}(a_\lambda)=0$, in other words $\varphi_{\mu\lambda}$ factorises as $\varphi_{\mu\lambda}:\mathrm{A}_\lambda\to\mathrm{A}'_\lambda\to\mathrm{A}_\mu$, and one can write $\mathrm{B}_\mu=\mathrm{B}'_\lambda\otimes_{\mathrm{A}'_\lambda}\mathrm{A}_\mu$, $\mathfrak{J}_\mu=\mathfrak{J}'_\lambda\mathrm{B}_\mu$ and $\mathrm{M}_\mu=\mathrm{M}'_\lambda\otimes_{\mathrm{A}'_\lambda}\mathrm{A}_\mu$; one deduces then from \sref{IV.11.2.9.2} that $\operatorname{gr}^*_{\mathfrak{J}_\mu}(\mathrm{M}_\mu)$ is a flat $\mathrm{A}_\mu$-module.

\medskip
\noindent\textbf{IV)} \emph{Reduction to the case where $\mathrm{B}_0=\mathrm{A}_0[\mathrm{T}_1,\ldots,\mathrm{T}_n]$ and $\operatorname{gr}^*_{\mathfrak{J}_0}(\mathrm{B}_0)=\mathrm{B}_0$.} --- By virtue of the hypothesis, there is a system $(t_i)_{1\leqslant i\leqslant n}$ of generators of the finite-type $\mathrm{A}_0$-algebra $\mathrm{B}_0$, such that the $t_i$ for $i\leqslant m$ generate the ideal $\mathfrak{J}_0$ of $\mathrm{B}_0$. Set $\mathrm{B}'_0=\mathrm{A}_0[\mathrm{T}_1,\ldots,\mathrm{T}_n]$, a polynomial algebra (hence Noetherian), and let $\mathfrak{J}'_0$ be the ideal of $\mathrm{B}'_0$ generated by the $\mathrm{T}_i$ with $i\leqslant m$; there is then a surjective $\mathrm{A}_0$-homomorphism $u_0:\mathrm{B}'_0\to\mathrm{B}_0$ taking each $\mathrm{T}_i$ to $t_i$ ($1\leqslant i\leqslant n$), and hence satisfying $u_0(\mathfrak{J}'_0)=\mathfrak{J}_0$. This homomorphism allows us to regard $\mathrm{M}_0$ as a finite-type $\mathrm{B}'_0$-module. Set
\[
  \mathrm{B}'_\lambda=\mathrm{B}'_0\otimes_{\mathrm{A}_0}\mathrm{A}_\lambda
  =\mathrm{A}_\lambda[\mathrm{T}_1,\ldots,\mathrm{T}_n],
  \qquad
  \mathfrak{J}'_\lambda=\mathfrak{J}'_0\mathrm{B}'_\lambda,
\]
\oldpage[IV-3]{129}
and also
\[
  \mathrm{B}'=\mathrm{B}'_0\otimes_{\mathrm{A}_0}\mathrm{A}
  =\mathrm{A}[\mathrm{T}_1,\ldots,\mathrm{T}_n],
  \qquad
  \mathfrak{J}'=\mathfrak{J}'_0\mathrm{B}'.
\]
Thus $\mathfrak{J}'$ is the ideal of $\mathrm{B}'$ generated by the $\mathrm{T}_i$ with $i\leqslant m$. Moreover, for every integer $k\geqslant 0$ one clearly has
\[
  \mathfrak{J}'^k\mathrm{M}=\mathfrak{J}^k\mathrm{M},
  \qquad
  \mathfrak{J}'_\lambda{}^k\mathrm{M}_\lambda
  =\mathfrak{J}_\lambda^k\mathrm{M}_\lambda
\]
for every $\lambda$. Hence
\[
  \operatorname{gr}^*_{\mathfrak{J}'}(\mathrm{M})
  =\operatorname{gr}^*_{\mathfrak{J}}(\mathrm{M})
\]
as $\mathrm{A}$-modules, and
\[
  \operatorname{gr}^*_{\mathfrak{J}'_\lambda}(\mathrm{M}_\lambda)
  =\operatorname{gr}^*_{\mathfrak{J}_\lambda}(\mathrm{M}_\lambda)
\]
as $\mathrm{A}_\lambda$-modules. We may therefore substitute $\mathrm{B}'$, $\mathrm{B}'_\lambda$, $\mathfrak{J}'$, and $\mathfrak{J}'_\lambda$ for $\mathrm{B}$, $\mathrm{B}_\lambda$, $\mathfrak{J}$, and $\mathfrak{J}_\lambda$, respectively, in the proof. Note also that, by construction, $\operatorname{gr}^*_{\mathfrak{J}'}(\mathrm{B}')$ identifies with $\mathrm{B}'$, and $\operatorname{gr}^*_{\mathfrak{J}'_\lambda}(\mathrm{B}'_\lambda)$ with $\mathrm{B}'_\lambda$.

\medskip
\noindent\textbf{V)} \emph{Proof that $\operatorname{gr}^*_\mathfrak{J}(\mathrm{M})$ is a $\operatorname{gr}^*_\mathfrak{J}(\mathrm{B})$-module of finite presentation.} --- We thus suppose from now on that $\mathrm{A}_0$ and the $\mathrm{A}_\lambda$ are Noetherian, that the transition homomorphisms $\mathrm{A}_\lambda\to\mathrm{A}_\mu$ are injective, that $\mathrm{B}_0=\mathrm{A}_0[\mathrm{T}_1,\ldots,\mathrm{T}_n]$, and that $\mathfrak{J}_0$ is generated by the $\mathrm{T}_i$ with $i\leqslant m$. The ring $\mathrm{C}_0=\operatorname{gr}^*_{\mathfrak{J}_0}(\mathrm{B}_0)$ therefore identifies with $\mathrm{B}_0$, and $\mathrm{C}=\operatorname{gr}^*_\mathfrak{J}(\mathrm{B})$ identifies with $\mathrm{B}$; at first, in what follows, we shall use only the fact that $\mathrm{C}\cong\mathrm{C}_0\otimes_{\mathrm{A}_0}\mathrm{A}$.

We will need the following variant of $(6.9.3)$:

\begin{lemma}[11.2.9.5]
\label{IV.11.2.9.5}
Let $\mathrm{A}_0$ be a Noetherian ring, $\mathrm{B}_0$ an
$\mathrm{A}_0$-algebra of finite type, $\mathfrak{J}_0$ an ideal of
$\mathrm{B}_0$, $\mathrm{M}_0$ a $\mathrm{B}_0$-module of finite type. There
exists a sequence $(\mathrm{S}_{0i})_{1\leqslant i\leqslant m}$ of
subschemes of $\mathrm{S}_0=\Spec(\mathrm{A}_0)$ having the following
properties:
\begin{enumerate}[label=\arabic{enumi}\textsuperscript{o}]
  \item The underlying spaces of the $\mathrm{S}_{0i}$ are pairwise disjoint and form a covering of $\mathrm{S}_0$.
  \item For each $i$, the set $\mathrm{S}_{0i}$ is open in $\bigcup_{j\geqslant i}\mathrm{S}_{0j}$.
  \item Each scheme $\mathrm{S}_{0i}$ is affine.
  \item If $\mathrm{A}_{0i}$ is the ring of $\mathrm{S}_{0i}$ and one sets $\mathrm{B}_{0i}=\mathrm{B}_0\otimes_{\mathrm{A}_0}\mathrm{A}_{0i}$, $\mathfrak{J}_{0i}=\mathfrak{J}_0\mathrm{B}_{0i}$, then $\operatorname{gr}^*_{\mathfrak{J}_{0i}}(\mathrm{M}_0\otimes_{\mathrm{A}_0}\mathrm{A}_{0i})$ is a flat $\mathrm{A}_{0i}$-module.
\end{enumerate}
\end{lemma}

\begin{proof}
We proceed as in $(6.9.3)$ by Noetherian induction, assuming the lemma known after replacing $\mathrm{A}_0$ by
\[
  \mathrm{A}'_0=\mathrm{A}_0/\mathfrak{a}_0,
\]
where $\mathfrak{a}_0$ is an ideal such that $\mathrm{V}(\mathfrak{a}_0)\neq\mathrm{S}_0$, and replacing $\mathrm{B}_0$, $\mathfrak{J}_0$, and $\mathrm{M}_0$ by
\[
  \mathrm{B}'_0=\mathrm{B}_0\otimes_{\mathrm{A}_0}\mathrm{A}'_0,
  \qquad
  \mathfrak{J}'_0=\mathfrak{J}_0\mathrm{B}'_0,
  \qquad
  \mathrm{M}_0\otimes_{\mathrm{A}_0}\mathrm{A}'_0,
\]
respectively. Let $\mathfrak{N}_0$ be the nilradical of $\mathrm{A}_0$. It
is plainly enough to prove the lemma after replacing $\mathrm{A}_0$ by
$\mathrm{A}_0/\mathfrak{N}_0$ and replacing $\mathrm{B}_0$,
$\mathfrak{J}_0$, and $\mathrm{M}_0$ by the corresponding objects just
described. In other words, we may suppose that $\mathrm{A}_0$ is
\emph{reduced}. Now
$\operatorname{gr}^*_{\mathfrak{J}_0}(\mathrm{B}_0)$ is a finite-type
$\mathrm{A}_0$-algebra and
$\operatorname{gr}^*_{\mathfrak{J}_0}(\mathrm{M}_0)$ is a finite-type
$\operatorname{gr}^*_{\mathfrak{J}_0}(\mathrm{B}_0)$-module. By the generic
flatness theorem $(6.9.1)$, there exists an open subset
$\mathrm{D}(t_0)\subset\mathrm{S}_0$ such that, on setting
$\mathrm{A}_{01}=(\mathrm{A}_0)_{t_0}$,
\[
  \operatorname{gr}^*_{\mathfrak{J}_0}(\mathrm{M}_0)
  \otimes_{\mathrm{A}_0}\mathrm{A}_{01}
\]
is a flat $\mathrm{A}_{01}$-module. Since $\mathrm{A}_{01}$ is a flat $\mathrm{A}_0$-module, this module identifies with
\[
  \operatorname{gr}^*_{\mathfrak{J}_{01}}
  (\mathrm{M}_0\otimes_{\mathrm{A}_0}\mathrm{A}_{01}),
\]
where
\[
  \mathrm{B}_{01}=\mathrm{B}_0\otimes_{\mathrm{A}_0}\mathrm{A}_{01},
  \qquad
  \mathfrak{J}_{01}=\mathfrak{J}_0\mathrm{B}_{01}.
\]
The complement of $\mathrm{D}(t_0)$ in $\mathrm{S}_0$ is then of the form $\mathrm{V}(\mathfrak{a}_0)$, and the conclusion follows by applying the Noetherian-induction hypothesis.
\end{proof}

Apply this lemma to the present situation, keeping the same notations; set $\mathrm{U}_{0i}=\bigcup_{j\leqslant i}\mathrm{S}_{0j}$, which is an open quasi-compact of $\mathrm{S}_0$. There is then a family $(f_{ik})_{1\leqslant i\leqslant m,\,1\leqslant k\leqslant n}$ of elements of $\mathrm{A}_0$ such that for each $i$, $\mathrm{U}_{0i}$ is the union of the $\mathrm{D}(f_{ik})$ ($1\leqslant k\leqslant n$).

\oldpage[IV-3]{130}

The $\mathrm{C}_0$-module $\operatorname{gr}^*_{\mathfrak{J}_0}(\mathrm{M}_0)$ is generated by finitely many homogeneous elements of degree $0$, so there is an epimorphism of graded $\mathrm{C}_0$-modules
\[
  u_0:\mathrm{C}_0^r\longrightarrow\operatorname{gr}^*_{\mathfrak{J}_0}(\mathrm{M}_0).
\]
As $\mathrm{C}=\mathrm{C}_0\otimes_{\mathrm{A}_0}\mathrm{A}$ and the canonical homomorphism $\operatorname{gr}^*_{\mathfrak{J}_0}(\mathrm{M}_0)\otimes_{\mathrm{A}_0}\mathrm{A}\to\operatorname{gr}^*_\mathfrak{J}(\mathrm{M})$ is surjective, one therefore deduces from $u_0$ an epimorphism of graded $\mathrm{C}$-modules
\[
  u:\mathrm{C}^r\xrightarrow{u_0\otimes 1_{\mathrm{A}}}\operatorname{gr}^*_{\mathfrak{J}_0}(\mathrm{M}_0)\otimes_{\mathrm{A}_0}\mathrm{A}\longrightarrow\operatorname{gr}^*_\mathfrak{J}(\mathrm{M}),
\]
and it remains to show that the graded $\mathrm{C}$-module
$\mathrm{Q}=\operatorname{Ker}(u)$ is \emph{of finite type}.
\\[0.4\baselineskip]
\begin{lemma}[11.2.9.6]
\label{IV.11.2.9.6}
Under the preceding hypotheses, let $\mathfrak{R}_0$ be an ideal of $\mathrm{A}_0$; set
\[
  \mathrm{A}'_0=\mathrm{A}_0/\mathfrak{R}_0,\qquad
  \mathrm{B}'_0=\mathrm{B}_0\otimes_{\mathrm{A}_0}\mathrm{A}'_0
  =\mathrm{B}_0/\mathfrak{R}_0\mathrm{B}_0,\qquad
  \mathfrak{J}'_0=\mathfrak{J}_0\mathrm{B}'_0,
\]
and $\mathfrak{R}=\mathfrak{R}_0\mathrm{A}$, $\mathrm{A}'=\mathrm{A}/\mathfrak{R}$, and suppose moreover that
\[
  \operatorname{gr}^*_{\mathfrak{J}'_0}
  (\mathrm{M}_0\otimes_{\mathrm{A}_0}\mathrm{A}'_0)
\]
is a flat $\mathrm{A}'_0$-module. Then there exist finitely many elements $q_h$ ($1\leqslant h\leqslant s$) of $\mathrm{Q}$ such that, for every prime ideal $\mathfrak{p}\supset\mathfrak{R}\mathrm{B}$ of $\mathrm{B}$, the canonical images of the $q_h$ in $\mathrm{Q}_\mathfrak{p}$ generate the $\mathrm{C}_\mathfrak{p}$-module $\mathrm{Q}_\mathfrak{p}$.
\end{lemma}

Suppose the lemma proved, and note that it can still be applied after replacing
$\mathrm{A}_0$ by $(\mathrm{A}_0)_{t_0}$, and $\mathrm{B}_0$,
$\mathfrak{J}_0$, $\mathrm{M}_0$, $\mathrm{A}_\lambda$, and $\mathrm{A}$ by
the corresponding localized objects, for any $t_0\in\mathrm{A}_0$, since
$(\mathrm{A}_0)_{t_0}$ is flat over $\mathrm{A}_0$. Apply the lemma
successively after replacing $\mathrm{A}_0$ by each of the rings
$(\mathrm{A}_0)_{f_{ik}}$, and $\mathfrak{R}_0$ by the ideal
$\mathfrak{R}_{ik}$ defining the closed subscheme of $\mathrm{D}(f_{ik})$
induced by $\mathrm{S}_{0i}$ on the open subset
$\mathrm{S}_{0i}\cap\mathrm{D}(f_{ik})$ of $\mathrm{S}_{0i}$. Since the
flatness hypothesis of \sref{IV.11.2.9.6} is satisfied in each of these cases
by \sref{IV.11.2.9.5}, for every pair $(i,k)$ one obtains a finite family
$(a_{ikh})_{1\leqslant h\leqslant l}$ of elements of
$\mathrm{Q}_{f_{ik}}$ whose images in $\mathrm{Q}_{\mathfrak{p}_{ik}}$
generate this $\mathrm{C}_{\mathfrak{p}_{ik}}$-module for every prime ideal
$\mathfrak{p}_{ik}$ of $\mathrm{B}_{f_{ik}}$ containing
$\mathfrak{R}_{ik}\mathrm{B}_{f_{ik}}$. For a suitable integer $d>0$, one
may write
\[
  a_{ikh}=\frac{b_{ikh}}{f_{ik}^{d}}
\]
for all $i,k,h$, with $b_{ikh}\in\mathrm{Q}$. Since the $\mathrm{S}_{0i}$ cover $\mathrm{S}_0$, every prime ideal $\mathfrak{p}$ of $\mathrm{B}$ has image in some $\mathrm{S}_{0i}\cap\mathrm{D}(f_{ik})$, and hence its image $\mathfrak{p}_{ik}$ in $\mathrm{B}_{f_{ik}}$ contains $\mathfrak{R}_{ik}\mathrm{B}_{f_{ik}}$. The images of the $b_{ikh}$ ($1\leqslant h\leqslant l$) in $\mathrm{Q}_{\mathfrak{p}_{ik}}=\mathrm{Q}_{\mathfrak{p}}$ therefore generate this $\mathrm{C}_{\mathfrak{p}}$-module. It follows that the $b_{ikh}$ ($1\leqslant i\leqslant m$, $1\leqslant k\leqslant n$, $1\leqslant h\leqslant l$) generate the $\mathrm{C}$-module $\mathrm{Q}$ (Bourbaki, \emph{Alg.\ comm.}, chap.~II, \S3, no.~3, th.~1).

It remains to prove Lemma \sref{IV.11.2.9.6}. Set
\[
  \mathrm{C}'=\mathrm{C}\otimes_\mathrm{A}\mathrm{A}'=\mathrm{C}/\mathfrak{R}\mathrm{C},\qquad
  \mathrm{B}'=\mathrm{B}\otimes_\mathrm{A}\mathrm{A}',\qquad
  \mathfrak{J}'=\mathfrak{J}\mathrm{B}',\qquad
  \mathrm{M}'=\mathrm{M}\otimes_\mathrm{A}\mathrm{A}'.
\]
By hypothesis $\operatorname{gr}^*_\mathfrak{J}(\mathrm{M})$ is a flat $\mathrm{A}$-module, hence $\operatorname{gr}^*_{\mathfrak{J}'}(\mathrm{M}')$ identifies with $\operatorname{gr}^*_\mathfrak{J}(\mathrm{M})\otimes_\mathrm{A}\mathrm{A}'$ \sref{IV.11.2.9.2}, and one has the exact sequence $(\mathbf{0}_\mathrm{I},~6.1.2)$
\begin{equation}
  0\longrightarrow\mathrm{Q}\otimes_\mathrm{A}\mathrm{A}'
  \longrightarrow(\mathrm{C}')^r
  \xrightarrow{u\otimes 1_{\mathrm{A}'}}
  \operatorname{gr}^*_{\mathfrak{J}'}(\mathrm{M}')
  \longrightarrow 0.
  \label{IV.11.2.9.7}
\tag{11.2.9.7}
\end{equation}

Set
\[
  \mathrm{C}'_0=\mathrm{C}_0\otimes_{\mathrm{A}_0}\mathrm{A}'_0
  =\mathrm{C}_0/\mathfrak{R}_0\mathrm{C}_0
\]
and consider the epimorphism of graded $\mathrm{C}'_0$-modules
\[
  u'_0:(\mathrm{C}'_0)^r
  \xrightarrow{u_0\otimes 1_{\mathrm{A}'_0}}
  \operatorname{gr}^*_{\mathfrak{J}_0}(\mathrm{M}_0)\otimes_{\mathrm{A}_0}\mathrm{A}'_0
  \longrightarrow
  \operatorname{gr}^*_{\mathfrak{J}'_0}
  (\mathrm{M}_0\otimes_{\mathrm{A}_0}\mathrm{A}'_0).
\]

Let $\mathrm{Q}_0=\operatorname{Ker}(u'_0)$. Using the fact that
\[
  \operatorname{gr}^*_{\mathfrak{J}'_0}
  (\mathrm{M}_0\otimes_{\mathrm{A}_0}\mathrm{A}'_0)
\]
is a flat $\mathrm{A}'_0$-module, one sees that its tensor product over $\mathrm{A}'_0$ with $\mathrm{A}'$ identifies with $\operatorname{gr}^*_{\mathfrak{J}'}(\mathrm{M}')$ \sref{IV.11.2.9.2}, and one has the exact sequence $(\mathbf{0}_\mathrm{I},~6.1.2)$
\begin{equation}
  0\longrightarrow\mathrm{Q}_0\otimes_{\mathrm{A}'_0}\mathrm{A}'
  \longrightarrow(\mathrm{C}')^r
  \xrightarrow{u\otimes 1_{\mathrm{A}'}}
  \operatorname{gr}^*_{\mathfrak{J}'}(\mathrm{M}')
  \longrightarrow 0.
  \label{IV.11.2.9.8}
\tag{11.2.9.8}
\end{equation}

\oldpage[IV-3]{131}

whence, by comparison with {\hyperlink{ega4.c04.thm.11-2-9}{(11.2.9.7)}}, an isomorphism of graded $\mathrm{C}'$-modules
\[
  \mathrm{Q}\otimes_\mathrm{A}\mathrm{A}'
  \xrightarrow{\sim}
  \mathrm{Q}_0\otimes_{\mathrm{A}'_0}\mathrm{A}'.
\]

Since $\mathrm{C}'_0$ is Noetherian, $\mathrm{Q}_0$ is a finite-type $\mathrm{C}'_0$-module; hence $\mathrm{Q}_0\otimes_{\mathrm{A}'_0}\mathrm{A}'$ is a finite-type graded $\mathrm{C}'$-module, and therefore so is $\mathrm{Q}\otimes_\mathrm{A}\mathrm{A}'$. Let $q_h$ ($1\leqslant h\leqslant s$) be homogeneous elements of $\mathrm{Q}$ whose images in $\mathrm{Q}\otimes_\mathrm{A}\mathrm{A}'$ generate this $\mathrm{C}'$-module.

Now consider an arbitrary prime ideal $\mathfrak{p}\supset\mathfrak{R}\mathrm{B}$ of $\mathrm{B}$. By flatness,
\[
  \mathrm{C}_\mathfrak{p}
  =\operatorname{gr}^*_{\mathfrak{J}_\mathfrak{p}}(\mathrm{B}_\mathfrak{p}),
\]
and $\operatorname{gr}^0_{\mathfrak{J}_\mathfrak{p}}(\mathrm{B}_\mathfrak{p})=\mathrm{B}_\mathfrak{p}/\mathfrak{J}_\mathfrak{p}$ is either the zero ring or a local ring; to prove the lemma, it plainly suffices to consider the second case. Each $(\mathrm{B}_\mathfrak{p}/\mathfrak{J}_\mathfrak{p})$-module $\operatorname{gr}^i_{\mathfrak{J}_\mathfrak{p}}(\mathrm{B}_\mathfrak{p})$ is then of finite presentation. On the other hand, for every $i$,
\[
  \mathrm{M}/\mathfrak{J}^i\mathrm{M}
  \simeq
  (\mathrm{M}_0/\mathfrak{J}_0^i\mathrm{M}_0)
  \otimes_{\mathrm{A}_0}\mathrm{A},
\]
and is therefore a finitely presented $\mathrm{B}$-module. By induction on $i$, the hypothesis that $\operatorname{gr}^*_\mathfrak{J}(\mathrm{M})$ is flat over $\mathrm{A}$ implies, by $(\mathbf{0}_\mathrm{I},~6.1.2)$, that $\mathrm{M}/\mathfrak{J}^i\mathrm{M}$ is flat over $\mathrm{A}$. Applying \sref{IV.11.2.6.1}, with $\mathrm{M}_0$ replaced by $\mathrm{M}_0/\mathfrak{J}_0^k\mathrm{M}_0$ for $k\leqslant i$, shows that there is an index $\lambda_i$ such that, for $\mu\geqslant\lambda_i$, every $\mathrm{M}_\mu/\mathfrak{J}_\mu^{k+1}\mathrm{M}_\mu$ is flat over $\mathrm{A}_\mu$; consequently, again by $(\mathbf{0}_\mathrm{I},~6.1.2)$, every $\operatorname{gr}^k_{\mathfrak{J}_\mu}(\mathrm{M}_\mu)$ is flat over $\mathrm{A}_\mu$ for $k\leqslant i$. It follows from \sref{IV.11.2.9.2} that every $(\mathrm{B}/\mathfrak{J})$-module $\operatorname{gr}^i_\mathfrak{J}(\mathrm{M})$ is of finite presentation, and hence that $\operatorname{gr}^i_{\mathfrak{J}_\mathfrak{p}}(\mathrm{M}_\mathfrak{p})$ is a finitely presented $(\mathrm{B}_\mathfrak{p}/\mathfrak{J}_\mathfrak{p})$-module.

Moreover, the images of the $q_h$ in
\[
  \mathrm{Q}_\mathfrak{p}\otimes_{\mathrm{B}_\mathfrak{p}}\boldsymbol{k}(\mathfrak{p})
  =
  (\mathrm{Q}\otimes_\mathrm{A}\mathrm{A}')_\mathfrak{p}
  \otimes_{\mathrm{A}'}\boldsymbol{k}(\mathfrak{p})
\]
generate this module over $\mathrm{C}_\mathfrak{p}\otimes_{\mathrm{B}_\mathfrak{p}}\boldsymbol{k}(\mathfrak{p})$, since
\[
  \mathrm{C}_\mathfrak{p}\otimes_{\mathrm{B}_\mathfrak{p}}\boldsymbol{k}(\mathfrak{p})
  =
  \mathrm{C}'_\mathfrak{p}\otimes_{\mathrm{A}'}\boldsymbol{k}(\mathfrak{p}).
\]
Since there is an exact sequence
\[
  \mathrm{Q}_\mathfrak{p}\otimes_{\mathrm{B}_\mathfrak{p}}\boldsymbol{k}(\mathfrak{p})
  \longrightarrow
  \mathrm{C}_\mathfrak{p}^{\,r}\otimes_{\mathrm{B}_\mathfrak{p}}\boldsymbol{k}(\mathfrak{p})
  \longrightarrow
  \operatorname{gr}^*_{\mathfrak{J}_\mathfrak{p}}(\mathrm{M}_\mathfrak{p})
  \otimes_{\mathrm{B}_\mathfrak{p}}\boldsymbol{k}(\mathfrak{p})
  \longrightarrow 0,
\]
it follows that $\operatorname{gr}^*_{\mathfrak{J}_\mathfrak{p}}(\mathrm{M}_\mathfrak{p})\otimes_{\mathrm{B}_\mathfrak{p}}\boldsymbol{k}(\mathfrak{p})$ is a finitely presented module over $\mathrm{C}_\mathfrak{p}\otimes_{\mathrm{B}_\mathfrak{p}}\boldsymbol{k}(\mathfrak{p})$. Applying \sref{IV.11.2.9.4}, with $\mathrm{A}$, $\mathrm{B}$, and $\mathrm{M}$ replaced respectively by $\mathrm{B}_\mathfrak{p}$, $\mathrm{C}_\mathfrak{p}$, and $\operatorname{gr}^*_{\mathfrak{J}_\mathfrak{p}}(\mathrm{M}_\mathfrak{p})$, we conclude that $\mathrm{Q}_\mathfrak{p}$ is a finite-type $\mathrm{C}_\mathfrak{p}$-module. Nakayama's lemma (and the fact that $\mathrm{C}=\mathrm{B}$) now shows that the images of the $q_h$ in $\mathrm{Q}_\mathfrak{p}$ generate this $\mathrm{C}_\mathfrak{p}$-module. This completes the proof of \sref{IV.11.2.9.6}, and also proves that $\operatorname{gr}^*_\mathfrak{J}(\mathrm{M})$ is a finitely presented $\mathrm{C}$-module.

\medskip
\noindent\textbf{VI)} \emph{End of the proof.} --- For brevity, set $\mathrm{C}_\lambda=\mathrm{C}_0\otimes_{\mathrm{A}_0}\mathrm{A}_\lambda$ (in fact equal to $\mathrm{B}_\lambda$), $\mathrm{N}_\lambda=\operatorname{gr}^*_{\mathfrak{J}_\lambda}(\mathrm{M}_\lambda)$, and $\mathrm{N}=\operatorname{gr}^*_\mathfrak{J}(\mathrm{M})$. First note that, for every integer $k$, $\mathfrak{J}^k\mathrm{M}$ identifies with $\varinjlim_\lambda\mathfrak{J}_\lambda^k\mathrm{M}_\lambda$ by exactness of the functor $\varinjlim_\lambda$: the image of $\mathfrak{J}_\lambda^k\mathrm{M}_\lambda$ in $\mathrm{M}_\mu$ for $\lambda\leqslant\mu$ (respectively in $\mathrm{M}$) generates $\mathfrak{J}_\mu^k\mathrm{M}_\mu$ (respectively $\mathfrak{J}^k\mathrm{M}$). Using exactness of $\varinjlim$ once more, we conclude that $\mathrm{N}$ identifies canonically with $\varinjlim_\lambda\mathrm{N}_\lambda$. With this identification, we shall first prove:
\begin{equation}
\label{IV.11.2.9.9}
\tag{11.2.9.9}
\text{\emph{For $\lambda$ sufficiently large, the canonical homomorphism $v_\lambda:\mathrm{N}_\lambda\otimes_{\mathrm{A}_\lambda}\mathrm{A}\to\mathrm{N}$ is bijective.}}
\end{equation}
Since the $\mathrm{C}_\lambda$ are Noetherian and $\mathrm{N}_\lambda$ is a finite-type $\mathrm{C}_\lambda$-module, the $\mathrm{C}$-modules $\mathrm{N}_\lambda\otimes_{\mathrm{A}_\lambda}\mathrm{A}$ are of finite presentation and form a filtered inductive system whose inductive limit identifies canonically with $\mathrm{N}$, since filtered colimits commute with tensor products. Moreover, the transition homomorphisms $v_{\mu\lambda}:\mathrm{N}_\lambda\otimes_{\mathrm{A}_\lambda}\mathrm{A}\to\mathrm{N}_\mu\otimes_{\mathrm{A}_\mu}\mathrm{A}$

\oldpage[IV-3]{132}

for $\lambda\leqslant\mu$, as well as the homomorphisms $v_\lambda$, are
\emph{surjective}. Fix $\lambda$, and regard
\[
  \operatorname{Ker}(v_{\mu\lambda})=\mathrm{N}'_\mu
\]
as a $\mathrm{C}$-submodule of
$\mathrm{N}_\lambda\otimes_{\mathrm{A}_\lambda}\mathrm{A}$. By the
definition of the inductive limit, the $\mathrm{N}'_\mu$ form an
increasing filtered family whose union is $\operatorname{Ker}(v_\lambda)$.
We proved in V) that $\mathrm{N}$ is a finitely presented
$\mathrm{C}$-module; consequently
$\operatorname{Ker}(v_\lambda)$ is a finite-type $\mathrm{C}$-module
(Bourbaki, \emph{Alg.\ comm.}, chap.~I, \S~2, no.~8, lemma~9). Hence
there is an index $\mu\geqslant\lambda$ such that
$\mathrm{N}'_\mu=\operatorname{Ker}(v_\lambda)$. Since $v_\mu$ is
surjective, this equality says that $v_\mu$ is also injective, and
therefore proves {\hyperlink{ega4.c04.thm.11-2-9}{(11.2.9.9)}}.

After replacing $\mathrm{A}_0$ and $\mathrm{M}_0$ by
$\mathrm{A}_\lambda$ and $\mathrm{M}_\lambda$ for a sufficiently large
$\lambda$, we may therefore assume that the canonical homomorphism
$v_\lambda$ is bijective for \emph{every} $\lambda$. Put
\[
  \mathrm{P}_\lambda
  =\mathrm{N}_0\otimes_{\mathrm{C}_0}\mathrm{C}_\lambda
  =\mathrm{N}_0\otimes_{\mathrm{A}_0}\mathrm{A}_\lambda,
\]
so that $\mathrm{N}=\varinjlim\mathrm{P}_\lambda$. Since
$\mathrm{C}_0$ is a finitely presented $\mathrm{A}_0$-algebra and
$\mathrm{P}_0$ is a finitely presented $\mathrm{C}_0$-module,
\sref{IV.11.2.6.1}, applied with
$\mathrm{B}_0=\mathrm{C}_0$ and $\mathrm{M}_0=\mathrm{P}_0$, shows that
there is an index $\lambda$ such that $\mathrm{P}_\mu$ is a flat
$\mathrm{A}_\mu$-module for every $\mu\geqslant\lambda$. For such a
$\mu$ we have a commutative diagram
\[
\xymatrix{
  \mathrm{P}_\mu \ar[r] \ar[d]^{w_\mu} & \mathrm{N} \ar@{=}[d] \\
  \mathrm{N}_\mu \ar[r] & \mathrm{N}
}
\]
and the definitions show that $w_\mu$ is \emph{surjective}. By III), the
homomorphism $\mathrm{A}_\mu\to\mathrm{A}$ is injective; since
$\mathrm{P}_\mu$ is flat over $\mathrm{A}_\mu$, the canonical
homomorphism
\[
  \mathrm{P}_\mu\longrightarrow
  \mathrm{N}
  =\mathrm{P}_\mu\otimes_{\mathrm{A}_\mu}\mathrm{A}
\]
is therefore \emph{injective}. The preceding diagram now shows that
$w_\mu$ is injective as well, hence bijective. Thus
$\mathrm{N}_\mu$ is a flat $\mathrm{A}_\mu$-module for
$\mu\geqslant\lambda$, which completes the proof of
\sref{IV.11.2.9}.
\end{proof}

\begin{remark}[11.2.10]
\label{IV.11.2.10}
We do not know whether the analogue of Raynaud's theorem for the
generalization of \sref{IV.11.2.6},~(i), is valid.
\end{remark}

\subsection{Application to the elimination of Noetherian hypotheses}
\label{subsection:IV.11.3}
\label{IV.11.3}

\oldpage[IV-3]{132}

\begin{theorem}[11.3.1]
\label{IV.11.3.1}
Let $f:\mathrm{X}\to\mathrm{Y}$ be a morphism locally of finite
presentation, and let $\sh{F}$ be a quasi-coherent
$\sh{O}_\mathrm{X}$-Module of finite presentation. Then the set
$\mathrm{U}$ of points $x\in\mathrm{X}$ at which $\sh{F}$ is
$f$-flat is open in $\mathrm{X}$. If, in addition,
$\operatorname{Supp}(\sh{F})=\mathrm{X}$, the restriction
$f|_\mathrm{U}$ is an open morphism.
\end{theorem}

\begin{proof}
The second assertion has already been proved in $(2.4.6)$ and is
included here only for reference.

Since the question is local on $\mathrm{X}$ and $\mathrm{Y}$, we may
assume that both are affine and that $f$ is of finite presentation. Let
$x\in\mathrm{X}$ be a point at which $\sh{F}$ is $f$-flat. Applying
\sref{IV.11.2.7}, we may assume that
\[
  \mathrm{X}=\mathrm{X}_0\times_{\mathrm{Y}_0}\mathrm{Y},
  \qquad
  f=f_0\times 1,
  \qquad
  \sh{F}
  =\sh{F}_0\otimes_{\sh{O}_{\mathrm{Y}_0}}\sh{O}_\mathrm{Y},
\]
where $\mathrm{Y}_0$ is Noetherian,
$f_0:\mathrm{X}_0\to\mathrm{Y}_0$ is of finite type, and
$\sh{F}_0$ is a coherent $\sh{O}_{\mathrm{X}_0}$-Module. Moreover, if
$x_0$ is the canonical projection of $x$ in $\mathrm{X}_0$, we may
assume that $\sh{F}_0$ is $f_0$-flat at $x_0$. By
\sref{IV.11.1.1}, the set $\mathrm{U}_0$ of points of
$\mathrm{X}_0$ at which $\sh{F}_0$ is $f_0$-flat is therefore a
neighbourhood of $x_0$. It follows from $(2.1.4)$ that $\sh{F}$ is
$f$-flat at every point of the inverse image of $\mathrm{U}_0$ in
$\mathrm{X}$; thus $\mathrm{U}$ is a neighbourhood of $x$.
\end{proof}

\oldpage[IV-3]{133}

\begin{corollary}[11.3.2]
\label{IV.11.3.2}
Let $f:\mathrm{X}\to\mathrm{Y}$, $g:\mathrm{Y}\to\mathrm{Z}$ be two morphisms of finite presentation, $\sh{F}$ an $\sh{O}_\mathrm{X}$-Module quasi-coherent of finite presentation. Then the set $\mathrm{U}$ of $y\in\mathrm{Y}$ such that, for every generization $y'$ of $y$, $\sh{F}$ is $(g\circ f)$-flat at all the points of $f^{-1}(y')$ (i.e.\ such that $\sh{F}\otimes_\mathrm{Y}\Spec(\sh{O}_{\mathrm{Y},y})$ is flat relative to the morphism $\mathrm{X}\times_\mathrm{Y}\Spec(\sh{O}_{\mathrm{Y},y})\to\mathrm{Z}$) is open in $\mathrm{Y}$.
\end{corollary}

\begin{proof}
The same reasoning as in \sref{IV.11.1.5} shows that if $\mathrm{V}$ is the set of $x\in\mathrm{X}$ where $\sh{F}$ is $(g\circ f)$-flat, $\mathrm{U}$ is the set of $y\in\mathrm{Y}$ whose every generization belongs to $\mathrm{E}=\mathrm{Y}-f(\mathrm{X}-\mathrm{V})$. As $g\circ f$ is of finite presentation, $\mathrm{V}$ is open in $\mathrm{X}$ by virtue of \sref{IV.11.3.1}, and hence $\mathrm{V}$ is ind-constructible $(1.9.6)$, and hence $\mathrm{X}-\mathrm{V}$ is pro-constructible; but as $f$ is quasi-compact, $f(\mathrm{X}-\mathrm{V})$ is pro-constructible in $\mathrm{Y}$ $(1.9.5,~(\mathrm{vii}))$, and hence $\mathrm{E}$ is ind-constructible in $\mathrm{Y}$. It then follows from $(1.10.1)$ that $\mathrm{U}$ is the interior of $\mathrm{E}$, and is therefore open in $\mathrm{Y}$.
\end{proof}

\begin{corollary}[11.3.3]
\label{IV.11.3.3}
Let $\mathrm{A}$ be a ring, $\mathrm{B}$ an $\mathrm{A}$-algebra of finite presentation, $\mathrm{C}$ a $\mathrm{B}$-algebra of finite presentation, $\mathrm{M}$ a $\mathrm{C}$-module of finite presentation. Then the set of $\mathfrak{q}\in\Spec(\mathrm{B})$ such that $\mathrm{M}_\mathfrak{q}$ is a flat $\mathrm{A}$-module is open in $\Spec(\mathrm{B})$.
\end{corollary}

\begin{proposition}[11.3.4]
\label{IV.11.3.4}
Let $f:\mathrm{X}\to\mathrm{S}$ be a morphism locally of finite presentation, $\mathscr{J}$ a quasi-coherent Ideal of finite type of $\sh{O}_\mathrm{X}$, $\mathrm{Y}$ the closed sub-prescheme of $\mathrm{X}$ defined by $\mathscr{J}$, $\sh{F}$ an $\sh{O}_\mathrm{X}$-Module of finite presentation. Suppose that $\operatorname{gr}^*_\mathscr{J}(\sh{F})$ is $f$-flat. Under these conditions:
\begin{enumerate}[label=\emph{(\roman*)}]
  \item $\sh{F}$ is $f$-flat at the points of $\mathrm{Y}$.
  \item If $\operatorname{gr}^*_\mathscr{J}(\sh{O}_\mathrm{X})$ is $f$-flat, $\operatorname{gr}^*_\mathscr{J}(\sh{O}_\mathrm{X})$ is an $(\sh{O}_\mathrm{X}/\mathscr{J})$-Algebra of finite presentation and $\operatorname{gr}^*_\mathscr{J}(\sh{F})$ is an $(\operatorname{gr}^*_\mathscr{J}(\sh{O}_\mathrm{X}))$-Module of finite presentation.
  \item Suppose $\operatorname{gr}^*_\mathscr{J}(\sh{O}_\mathrm{X})$ is $f$-flat (which implies that $\sh{O}_\mathrm{X}/\mathscr{J}$ is $f$-flat). Then the set $\mathrm{W}$ of $y\in\mathrm{Y}$ such that $(\operatorname{gr}^*_\mathscr{J}(\sh{F}))_y$ is a flat $(\sh{O}_\mathrm{X}/\mathscr{J})_y$-module is open in $\mathrm{Y}$.
  \item Suppose $\operatorname{gr}^*_\mathscr{J}(\sh{O}_\mathrm{X})$ is $f$-flat. Let $\mathrm{S}'\to\mathrm{S}$ be any morphism, let $\mathrm{X}'=\mathrm{X}\times_\mathrm{S}\mathrm{S}'$, and let $p:\mathrm{X}'\to\mathrm{X}$ be the canonical projection. Set $\mathrm{Y}'=p^{-1}(\mathrm{Y})=\mathrm{Y}\times_\mathrm{S}\mathrm{S}'$, the closed sub-prescheme of $\mathrm{X}'$ defined by $\mathscr{J}'=p^*(\mathscr{J})\sh{O}_{\mathrm{X}'}$, and set $\sh{F}'=p^*(\sh{F})=\sh{F}\otimes_{\sh{O}_\mathrm{S}}\sh{O}_{\mathrm{S}'}$. Then, if $\mathrm{W}'=p^{-1}(\mathrm{W})$, one has $\operatorname{gr}^*_{\mathscr{J}'}(\sh{F}')|_{\mathrm{W}'}\cong p^*(\operatorname{gr}^*_\mathscr{J}(\sh{F})|_\mathrm{W})$, and $\operatorname{gr}^*_{\mathscr{J}'}(\sh{F}')$ is a flat $(\sh{O}_{\mathrm{X}'}/\mathscr{J}')$-Module at the points of $\mathrm{W}'$.
\end{enumerate}
\end{proposition}

\begin{proof}
The questions being local on $\mathrm{X}$ and $\mathrm{S}$, we may suppose that
$\mathrm{S}=\Spec(\mathrm{A})$ and $\mathrm{X}=\Spec(\mathrm{B})$ are affine,
where $\mathrm{B}$ is an $\mathrm{A}$-algebra of finite presentation, that
$\sh{F}=\widetilde{\mathrm{M}}$, where $\mathrm{M}$ is a
$\mathrm{B}$-module of finite presentation, and that
$\mathscr{J}=\widetilde{\mathfrak{J}}$, where $\mathfrak{J}$ is an ideal
of finite type of $\mathrm{B}$. By virtue of $(8.9.1)$ and $(8.5.11)$,
there exist a Noetherian subring $\mathrm{A}_0$ of $\mathrm{A}$, an
$\mathrm{A}_0$-algebra of finite type $\mathrm{B}_0$ such that
$\mathrm{B}=\mathrm{B}_0\otimes_{\mathrm{A}_0}\mathrm{A}$, an ideal
$\mathfrak{J}_0$ of $\mathrm{B}_0$ such that
$\mathfrak{J}=\mathfrak{J}_0\mathrm{B}$, and a finite-type
$\mathrm{B}_0$-module $\mathrm{M}_0$ such that
$\mathrm{M}=\mathrm{M}_0\otimes_{\mathrm{A}_0}\mathrm{A}
=\mathrm{M}_0\otimes_{\mathrm{B}_0}\mathrm{B}$. Moreover,
$\mathrm{A}$ is the inductive limit of its finite-type
$\mathrm{A}_0$-subalgebras $\mathrm{A}_\lambda$. Set
$\mathrm{B}_\lambda=\mathrm{B}_0\otimes_{\mathrm{A}_0}\mathrm{A}_\lambda$,
$\mathrm{M}_\lambda=\mathrm{M}_0\otimes_{\mathrm{A}_0}\mathrm{A}_\lambda
=\mathrm{M}_0\otimes_{\mathrm{B}_0}\mathrm{B}_\lambda$, and
$\mathfrak{J}_\lambda=\mathfrak{J}_0\mathrm{B}_\lambda$, so that
$\mathrm{B}=\varinjlim\mathrm{B}_\lambda$ and
$\mathrm{M}=\varinjlim\mathrm{M}_\lambda$. By hypothesis,
$\operatorname{gr}^*_\mathfrak{J}(\mathrm{M})$ is a flat
$\mathrm{A}$-module; hence it follows from \sref{IV.11.2.9} that there
exists $\lambda$ such that
$\operatorname{gr}^*_{\mathfrak{J}_\lambda}(\mathrm{M}_\lambda)$ is a
flat $\mathrm{A}_\lambda$-module. To prove that $\sh{F}$ is $f$-flat at
the points of $\mathrm{Y}$, we may therefore restrict to the case where
$\mathrm{S}$ is Noetherian; but then $(\mathbf{0}_\mathrm{I},~6.1.2)$,
applied by induction on $n$, proves that the
$\sh{F}/\mathscr{J}^{n+1}\sh{F}$ are $f$-flat, and one concludes by
$(\mathbf{0}_{\mathrm{III}},~10.2.6)$.
\oldpage[IV-3]{134}

The proof of (ii) reduces similarly to the case where $\mathrm{S}$ is Noetherian, using \sref{IV.11.2.9}; the conclusion is then evident, since in this case $\operatorname{gr}^*_\mathfrak{J}(\mathrm{B})$ is a $(\mathrm{B}/\mathfrak{J})$-algebra of finite type and $\operatorname{gr}^*_\mathfrak{J}(\mathrm{M})$ is an $\operatorname{gr}^*_\mathfrak{J}(\mathrm{B})$-module of finite type.

To prove (iii), one reduces as in (i) to the case where $\mathrm{S}$ and $\mathrm{X}$ are affine. With the same notation, one may suppose, by virtue of \sref{IV.11.2.9}, that $\operatorname{gr}^*_{\mathfrak{J}_\lambda}(\mathrm{B}_\lambda)$ and $\operatorname{gr}^*_{\mathfrak{J}_\lambda}(\mathrm{M}_\lambda)$ are flat $\mathrm{A}_\lambda$-modules and that
\[
\operatorname{gr}^*_\mathfrak{J}(\mathrm{B})
=\operatorname{gr}^*_{\mathfrak{J}_\lambda}(\mathrm{B}_\lambda)\otimes_{\mathrm{A}_\lambda}\mathrm{A},
\qquad
\operatorname{gr}^*_\mathfrak{J}(\mathrm{M})
=\operatorname{gr}^*_{\mathfrak{J}_\lambda}(\mathrm{M}_\lambda)\otimes_{\mathrm{A}_\lambda}\mathrm{A}.
\]
Consequently, $\operatorname{gr}^*_\mathfrak{J}(\mathrm{B})$ is a $(\mathrm{B}/\mathfrak{J})$-algebra of finite presentation and $\operatorname{gr}^*_\mathfrak{J}(\mathrm{M})$ is an $\operatorname{gr}^*_\mathfrak{J}(\mathrm{B})$-module of finite presentation. If $\mathrm{W}^m$ is the set of $\mathfrak{p}\in\Spec(\mathrm{B}/\mathfrak{J})$ such that $\operatorname{gr}^m_\mathfrak{J}(\mathrm{M})_{\mathfrak{p}}$ is a flat $(\mathrm{B}/\mathfrak{J})_{\mathfrak{p}}$-module, then $\mathrm{W}^m$ is open in $\Spec(\mathrm{B}/\mathfrak{J})$ by \sref{IV.11.3.1}, and $\mathrm{W}=\bigcap_m\mathrm{W}^m$; thus $\mathrm{W}$ is stable under generization. Assertion (iii) now follows from \sref{IV.11.3.2}, applied with $\mathrm{X}=\Spec(\operatorname{gr}^*_\mathfrak{J}(\mathrm{B}))$, $\mathrm{Y}=\mathrm{Z}=\Spec(\mathrm{B}/\mathfrak{J})$, and $\sh{F}=(\operatorname{gr}^*_\mathfrak{J}(\mathrm{M}))^\sim$.

Finally, (iv) follows immediately from \sref{IV.11.2.9.2}.

Generalising the definition of $(6.10.1)$, let $\mathrm{Y}$ be a closed sub-prescheme of a prescheme $\mathrm{X}$, defined by a quasi-coherent ideal $\mathscr{J}$ of $\sh{O}_\mathrm{X}$, and let $\sh{F}$ be a quasi-coherent $\sh{O}_\mathrm{X}$-Module. One says that $\sh{F}$ is \emph{normally flat along $\mathrm{Y}$ at a point $y$} if $(\operatorname{gr}^*_\mathscr{J}(\sh{F}))_y$ is a flat $\sh{O}_{\mathrm{Y},y}$-module. One says that $\sh{F}$ is \emph{normally flat along $\mathrm{Y}$} if this is so at every point of $\mathrm{Y}$.
\end{proof}

\begin{corollary}[11.3.5]
\label{IV.11.3.5}
Under the general hypotheses of \sref{IV.11.3.4}, suppose that $\operatorname{gr}^*_\mathscr{J}(\sh{O}_\mathrm{X})$ and $\operatorname{gr}^*_\mathscr{J}(\sh{F})$ are $f$-flat. Then the set $\mathrm{W}$ of $y\in\mathrm{Y}$ such that $\sh{F}$ is normally flat along $\mathrm{Y}$ at the point $y$ is open in $\mathrm{Y}$, and (with the notations of \sref{IV.11.3.4},~(iv))) $\sh{F}'$ is normally flat along $\mathrm{Y}'$ at all points of $\mathrm{W}'=p^{-1}(\mathrm{W})$.
\end{corollary}

\begin{proposition}[11.3.6]
\label{IV.11.3.6}
The notations being those of $(8.5.1)$ and $(8.8.1)$, suppose $\mathrm{S}_\alpha$ quasi-compact, $\mathrm{X}_\alpha$ of finite presentation over $\mathrm{S}_\alpha$, $\mathrm{Y}_\alpha$ a closed sub-prescheme of $\mathrm{X}_\alpha$ defined by a quasi-coherent Ideal $\mathscr{J}_\alpha$ of finite type of $\sh{O}_{\mathrm{X}_\alpha}$; suppose that $\operatorname{gr}^*_{\mathscr{J}_\alpha}(\sh{O}_{\mathrm{X}_\alpha})$ is flat on $\mathrm{S}_\alpha$; finally suppose that $\sh{F}_\alpha$ is an $\sh{O}_{\mathrm{X}_\alpha}$-Module of finite presentation. For $\sh{F}$ to be normally flat along $\mathrm{Y}$, it is necessary and sufficient that there exist $\lambda$ such that $\sh{F}_\lambda$ be normally flat along $\mathrm{Y}_\lambda$.
\end{proposition}

\begin{proof}
Note that $\mathrm{Y}$ (resp.\ $\mathrm{Y}_\lambda$) is the closed sub-prescheme of $\mathrm{X}$ (resp.\ $\mathrm{X}_\lambda$) defined by $\mathscr{J}=\mathscr{J}_\alpha\sh{O}_\mathrm{X}$ (resp.\ $\mathscr{J}_\lambda=\mathscr{J}_\alpha\sh{O}_{\mathrm{X}_\lambda}$). By the hypothesis and \sref{IV.11.2.9.2}, one has
\[
\operatorname{gr}^*_\mathscr{J}(\sh{O}_\mathrm{X})
=\operatorname{gr}^*_{\mathscr{J}_\alpha}(\sh{O}_{\mathrm{X}_\alpha})
 \otimes_{\sh{O}_{\mathrm{S}_\alpha}}\sh{O}_\mathrm{S},
\qquad
\operatorname{gr}^*_{\mathscr{J}_\lambda}(\sh{O}_{\mathrm{X}_\lambda})
=\operatorname{gr}^*_{\mathscr{J}_\alpha}(\sh{O}_{\mathrm{X}_\alpha})
 \otimes_{\sh{O}_{\mathrm{S}_\alpha}}\sh{O}_{\mathrm{S}_\lambda}
\quad(\lambda\geqslant\alpha),
\]
and $\operatorname{gr}^*_\mathscr{J}(\sh{O}_\mathrm{X})$ (resp.\ $\operatorname{gr}^*_{\mathscr{J}_\lambda}(\sh{O}_{\mathrm{X}_\lambda})$) is flat over $\mathrm{S}$ (resp.\ $\mathrm{S}_\lambda$). It follows that $\mathrm{Y}$ is flat over $\mathrm{S}$ (resp.\ $\mathrm{Y}_\lambda$ is flat over $\mathrm{S}_\lambda$). If $\sh{F}_\lambda$ is normally flat along $\mathrm{Y}_\lambda$, then $\left.\operatorname{gr}^*_{\mathscr{J}_\lambda}(\sh{F}_\lambda)\right|_{\mathrm{Y}_\lambda}$ is flat over $\mathrm{Y}_\lambda$, hence also over $\mathrm{S}_\lambda$; and since
\[
\left.\operatorname{gr}^n_{\mathscr{J}_\lambda}(\sh{F}_\lambda)\right|_{\mathrm{X}_\lambda\setminus\mathrm{Y}_\lambda}=0
\quad\text{for }n\geqslant0,
\]
$\operatorname{gr}^*_{\mathscr{J}_\lambda}(\sh{F}_\lambda)$ is flat over $\mathrm{S}_\lambda$. By \sref{IV.11.2.9.2},
\[
\operatorname{gr}^*_\mathscr{J}(\sh{F})
=\operatorname{gr}^*_{\mathscr{J}_\lambda}(\sh{F}_\lambda)
 \otimes_{\sh{O}_{\mathrm{S}_\lambda}}\sh{O}_\mathrm{S},
\]
so $\sh{F}$ is normally flat along $\mathrm{Y}$.

Conversely, one may suppose that $\mathrm{S}_\alpha$ and $\mathrm{X}_\alpha$ are affine and adopt the notation of \sref{IV.11.2.9}. Since $\operatorname{gr}^*_\mathfrak{J}(\mathrm{B})$ is a flat $\mathrm{A}$-module, the same is true of $\operatorname{gr}^*_\mathfrak{J}(\mathrm{M})$; hence, by \sref{IV.11.2.9}, there exists $\lambda\geqslant\alpha$ such that $\operatorname{gr}^*_{\mathfrak{J}_\lambda}(\mathrm{M}_\lambda)$ is a flat $\mathrm{A}_\lambda$-module and
\[
\operatorname{gr}^*_\mathfrak{J}(\mathrm{M})
=\operatorname{gr}^*_{\mathfrak{J}_\lambda}(\mathrm{M}_\lambda)
 \otimes_{\mathrm{A}_\lambda}\mathrm{A}.
\]
Moreover, by \sref{IV.11.3.4},~(ii), $\operatorname{gr}^*_{\mathfrak{J}_\lambda}(\mathrm{M}_\lambda)$ is an $\operatorname{gr}^*_{\mathfrak{J}_\lambda}(\mathrm{B}_\lambda)$-module of finite presentation and $\operatorname{gr}^*_{\mathfrak{J}_\lambda}(\mathrm{B}_\lambda)$ is a $(\mathrm{B}_\lambda/\mathfrak{J}_\lambda)$-algebra of finite presentation. The conclusion now follows from \sref{IV.11.2.6}.
\end{proof}

\oldpage[IV-3]{135}
\begin{proposition}[11.3.7]
\label{IV.11.3.7}
Let $f:\mathrm{X}\to\mathrm{Y}$ be a morphism locally of finite presentation, $\sh{F}$, $\sh{F}'$ two $\sh{O}_\mathrm{X}$-Modules quasi-coherent of finite presentation, $u:\sh{F}'\to\sh{F}$ an $\sh{O}_\mathrm{X}$-homomorphism, $x$ a point of $\mathrm{X}$ and $y=f(x)$; suppose that $\sh{F}$ is $f$-flat at the point $x$. The following conditions are equivalent:
\begin{enumerate}[label=\emph{\alph*)}]
  \item $(\operatorname{Ker} u)_x=0$ and $\operatorname{Coker} u$ is an $\sh{O}_\mathrm{X}$-Module $f$-flat at the point $x$.
  \item The homomorphism $(u\otimes 1)_x:(\sh{F}'\otimes_{\sh{O}_\mathrm{Y}}\boldsymbol{k}(y))_x\to(\sh{F}\otimes_{\sh{O}_\mathrm{Y}}\boldsymbol{k}(y))_x$ is injective.
\end{enumerate}
If in addition $\sh{F}$ is $f$-flat, the set of points $x$ satisfying the preceding equivalent conditions is open in $\mathrm{X}$.
\end{proposition}

\begin{proof}
Condition $a)$ implies $b)$ by virtue of $(\mathbf{0}_\mathrm{I},~6.1.2)$, without any hypothesis on $\sh{F}$. To prove the converse, we may restrict ourselves to the case where $\mathrm{X}$ and $\mathrm{Y}$ are affine and then, by virtue of $(8.9.1)$ and \sref{IV.11.2.7}, suppose that
\[
  \mathrm{X}=\mathrm{X}_0\times_{\mathrm{Y}_0}\mathrm{Y},\qquad
  f=f_0\times 1_\mathrm{Y},\qquad
  \sh{F}=\sh{F}_0\otimes_{\sh{O}_{\mathrm{Y}_0}}\sh{O}_\mathrm{Y},\qquad
  \sh{F}'=\sh{F}'_0\otimes_{\sh{O}_{\mathrm{Y}_0}}\sh{O}_\mathrm{Y},\qquad
  u=u_0\otimes 1_\mathrm{Y},
\]
where $\mathrm{Y}_0$ is Noetherian, $f_0:\mathrm{X}_0\to\mathrm{Y}_0$ is a morphism of finite type, $\sh{F}_0$ and $\sh{F}'_0$ are two coherent $\sh{O}_{\mathrm{X}_0}$-Modules, and $u_0:\sh{F}'_0\to\sh{F}_0$ is a homomorphism. Moreover, if $x_0$ is the projection of $x$ in $\mathrm{X}_0$, we may suppose that $\sh{F}_0$ is $f_0$-flat at $x_0$. Set $y_0=f_0(x_0)$. By the transitivity of fibres $(\mathbf{I},~3.6.4)$, the projection $f^{-1}(y)\to f_0^{-1}(y_0)$ is faithfully flat $(2.2.13)$; and since
\[
  (u\otimes 1)_x=((u_0\otimes 1)_{x_0})\otimes 1_{\boldsymbol{k}(y)},
\]
the hypothesis that $(u\otimes 1)_x$ is injective implies that $(u_0\otimes 1)_{x_0}$ is injective $(2.2.7)$. Applying $(\mathbf{0}_{\mathrm{III}},~10.2.4)$ to the Noetherian local rings $\sh{O}_{\mathrm{Y}_0,y_0}$ and $\sh{O}_{\mathrm{X}_0,x_0}$ and to the $\sh{O}_{\mathrm{X}_0,x_0}$-modules $(\sh{F}'_0)_{x_0}$ and $(\sh{F}_0)_{x_0}$, the second of which is flat over $\sh{O}_{\mathrm{Y}_0,y_0}$, we obtain
\[
  (\operatorname{Ker}u_0)_{x_0}=0
  \quad\text{and}\quad
  \operatorname{Coker}u_0\ \text{is $f_0$-flat at }x_0.
\]
It follows first that $\operatorname{Coker}u$ is $f$-flat at $x$ $(2.1.4)$. Moreover, by right exactness of the tensor product,
\[
  \operatorname{Coker}u=(\operatorname{Coker}u_0)
  \otimes_{\sh{O}_{\mathrm{Y}_0}}\sh{O}_\mathrm{Y}.
\]
Applying $(\mathbf{0}_\mathrm{I},~6.1.2)$ to the exact sequence
\[
  0\longrightarrow(\sh{F}'_0)_{x_0}
  \longrightarrow(\sh{F}_0)_{x_0}
  \longrightarrow(\operatorname{Coker}u_0)_{x_0}
  \longrightarrow0,
\]
we conclude that $u_x$ is injective, which proves condition $a)$.

Finally, it follows from \sref{IV.11.1.2} that the set $\mathrm{U}_0$ of points $z_0\in\mathrm{X}_0$ such that $(u_0\otimes 1)_{z_0}$ is injective is open. By flatness, for every $z\in\mathrm{X}$ above a point $z_0\in\mathrm{U}_0$, the homomorphism $(u\otimes 1)_z$ is injective. The set of such points therefore contains the inverse image $\mathrm{U}$ of $\mathrm{U}_0$ in $\mathrm{X}$ and is consequently a neighbourhood of $x$, which completes the proof.
\end{proof}

\begin{theorem}[11.3.8]
\label{IV.11.3.8}
Let $f:\mathrm{X}\to\mathrm{Y}$ be a morphism locally of finite presentation, $\sh{F}$ an $\sh{O}_\mathrm{X}$-Module quasi-coherent of finite presentation, $(g_i)_{1\leqslant i\leqslant n}$ a sequence of sections of $\sh{O}_\mathrm{X}$ above $\mathrm{X}$; set $\sh{F}_i=\sh{F}/\sum_{j=1}^{i}g_j\sh{F}$ for $0\leqslant i\leqslant n$ (with $\sh{F}_0=\sh{F}$). Let $x$ be a point of $\mathrm{X}$, $y=f(x)$, and set $\mathrm{X}_y=f^{-1}(y)=\mathrm{X}\times_\mathrm{Y}\Spec(\boldsymbol{k}(y))$, $\sh{F}_y=\sh{F}\otimes_{\sh{O}_y}\boldsymbol{k}(y)$, which is an $\sh{O}_{\mathrm{X}_y}$-Module. Suppose that the $(g_i)_x$ belong to the maximal ideal $\mathfrak{m}_x$. The following conditions are equivalent:
\begin{enumerate}[label=\emph{\alph*)}]
  \item The sequence $((g_i)_x)$ is $\sh{F}_x$-regular and the $\sh{F}_i$ ($0\leqslant i\leqslant n$) are $f$-flat at the point $x$.
  \item The sequence $((g_i)_x)$ is $\sh{F}_x$-regular and $\sh{F}_n$ is $f$-flat at the point $x$.
  \item[\emph{b$'$)}] There exists a neighbourhood $\mathrm{U}$ of $x$ such that the sequence $(g_i|\mathrm{U})$ is $(\sh{F}|\mathrm{U})$-quasi-regular, and $\sh{F}_n$ is $f$-flat at the point $x$.
\oldpage[IV-3]{136}
  \item $\sh{F}$ is $f$-flat at the point $x$, and the sequence $((g_i\otimes 1)_x)$ of elements of $\sh{O}_{\mathrm{X}_y,x}$ images of the $(g_i)_x$ is $(\sh{F}_y)_x$-regular.
  \item $\sh{F}$ is $f$-flat at the point $x$, and for every morphism $\mathrm{Y}'\to\mathrm{Y}$ and every point $x'\in\mathrm{X}'=\mathrm{X}\times_\mathrm{Y}\mathrm{Y}'$ above $x$, if one sets $\sh{F}'=\sh{F}\otimes_\mathrm{Y}\mathrm{Y}'$, the sequence $(g_i\otimes 1)_{x'}$ of elements of $\sh{O}_{\mathrm{X}',x'}$ is $\sh{F}'_{x'}$-regular.
\end{enumerate}
Moreover, the set of $x\in\mathrm{X}$ satisfying these conditions is open in the set $\mathrm{Z}$ of $x$ such that $g_i(x)=0$ for every $i$.
\end{theorem}

\begin{proof}
The fact that $a)$ implies $d)$ follows immediately from $(\mathbf{0},~15.1.15)$, and $c)$ is a particular case of $d)$; moreover, $a)$ implies $b)$ trivially. Let us show that $b)$ or $c)$ implies $a)$. Since the $\sh{F}_i$ are of finite presentation, the fact that $c)$ implies $a)$ follows immediately from \sref{IV.11.3.7} by induction on $n$, and this also shows that the set of $x\in\mathrm{X}$ satisfying $c)$ is open in $\mathrm{Z}$.

To show that $b)$ implies $a)$, induction on $n$ reduces us immediately to the case $n=1$; we write $g$ instead of $g_1$. Since the question is local on $\mathrm{X}$ and $\mathrm{Y}$, we may suppose that $\mathrm{Y}=\Spec(\mathrm{A})$ and $\mathrm{X}=\Spec(\mathrm{B})$ are affine, that $\mathrm{B}$ is an $\mathrm{A}$-algebra of finite presentation, and that $\sh{F}=\widetilde{\mathrm{M}}$, where $\mathrm{M}$ is a $\mathrm{B}$-module of finite presentation. By $(8.9.1)$, we may then write
\[
  \mathrm{B}=\mathrm{B}_0\otimes_{\mathrm{A}_0}\mathrm{A},
  \qquad
  \mathrm{M}=\mathrm{M}_0\otimes_{\mathrm{A}_0}\mathrm{A},
\]
where $\mathrm{A}_0$ is a Noetherian subring of $\mathrm{A}$, $\mathrm{B}_0$ is an $\mathrm{A}_0$-algebra of finite type, and $\mathrm{M}_0$ is a $\mathrm{B}_0$-module of finite type. Moreover, $\mathrm{A}$ is the filtered inductive limit of its subrings $\mathrm{A}_\lambda$ that are $\mathrm{A}_0$-algebras of finite type (and hence Noetherian). If we set
\[
  \mathrm{B}_\lambda=\mathrm{B}_0\otimes_{\mathrm{A}_0}\mathrm{A}_\lambda,
  \qquad
  \mathrm{M}_\lambda=\mathrm{M}_0\otimes_{\mathrm{A}_0}\mathrm{A}_\lambda,
\]
then $\mathrm{B}_\lambda$ is an $\mathrm{A}_\lambda$-algebra of finite type, $\mathrm{M}_\lambda$ is a $\mathrm{B}_\lambda$-module of finite type, and
\[
  \mathrm{B}=\varinjlim_\lambda\mathrm{B}_\lambda,
  \qquad
  \mathrm{M}=\varinjlim_\lambda\mathrm{M}_\lambda.
\]
There therefore exist an index $\lambda$ and an element $g_\lambda\in\mathrm{B}_\lambda$ such that $g=g_\lambda\otimes 1$. Returning to geometric notation and setting
\[
  \mathrm{Y}_\lambda=\Spec(\mathrm{A}_\lambda),
  \qquad
  \mathrm{X}_\lambda=\Spec(\mathrm{B}_\lambda),
  \qquad
  \sh{F}_\lambda=\widetilde{\mathrm{M}_\lambda},
\]
it will suffice to prove that there exists $\mu\geqslant\lambda$ such that, at the point $x_\mu\in\mathrm{X}_\mu$ that is the projection of $x$, $(g_\mu)_{x_\mu}$ is $(\sh{F}_\mu)_{x_\mu}$-regular and $\sh{F}_\mu/g_\mu\sh{F}_\mu$ is flat over $\mathrm{Y}_\mu$ at $x_\mu$. Indeed, it will then follow from $(\mathbf{0},~15.1.16)$ that $\sh{F}_\mu$ is flat over $\mathrm{Y}_\mu$ at $x_\mu$, and hence that $\sh{F}$ is flat over $\mathrm{Y}$ at $x$.

The fact that $b)$ implies $a)$ will therefore follow from the following proposition.

\begin{proposition}[11.3.9]
\label{IV.11.3.9}
With the notation and general hypotheses of $(8.5.1)$ and $(8.8.1)$, suppose that $f_\alpha:\mathrm{X}_\alpha\to\mathrm{Y}_\alpha$ is a morphism locally of finite presentation. Let $\sh{F}_\alpha$ and $\sh{G}_\alpha$ be two quasi-coherent $\sh{O}_{\mathrm{X}_\alpha}$-Modules of finite presentation, let $h_\alpha:\sh{F}_\alpha\to\sh{G}_\alpha$ be an $\sh{O}_{\mathrm{X}_\alpha}$-homomorphism, and let $\sh{H}_\alpha$ be its cokernel. Finally, let $x$ be a point of $\mathrm{X}$, and let $x_\lambda$ be its projection in $\mathrm{X}_\lambda$ for $\lambda\geqslant\alpha$. In order that $h_x$ be injective and that $\sh{H}=\operatorname{Coker}h$ be $f$-flat at $x$, it is necessary and sufficient that there exist $\lambda\geqslant\alpha$ such that $(h_\lambda)_{x_\lambda}$ is injective and $\sh{H}_\lambda=\operatorname{Coker}h_\lambda$ is $f_\lambda$-flat at $x_\lambda$. Moreover, the set of $x\in\mathrm{X}$ having the preceding properties is open in $\mathrm{X}$.
\end{proposition}

\begin{proof}
Recall that, by the right exactness of the tensor product,
\[
  \sh{H}_\mu=v^*_{\lambda\mu}(\sh{H}_\lambda)
  \quad\text{for }\lambda\leqslant\mu,
  \qquad
  \sh{H}=v^*_\lambda(\sh{H}_\lambda),
\]
which justifies the notation.

The sufficiency of the condition follows from the fact that, if the sequence
\[
  0\longrightarrow(\sh{F}_\lambda)_{x_\lambda}
  \longrightarrow(\sh{G}_\lambda)_{x_\lambda}
  \longrightarrow(\sh{H}_\lambda)_{x_\lambda}
  \longrightarrow0
\]
is exact and $(\sh{H}_\lambda)_{x_\lambda}$ is a flat $\sh{O}_{f_\lambda(x_\lambda)}$-module, then $\sh{H}_x$ is a flat $\sh{O}_{f(x)}$-module by base change $(2.1.4)$ and the sequence
\[
  0\longrightarrow\sh{F}_x
  \longrightarrow\sh{G}_x
  \longrightarrow\sh{H}_x
  \longrightarrow0
\]
is exact $(2.1.8)$. To prove that

\oldpage[IV-3]{137}

the condition is necessary, note that $\sh{H}$ is of finite presentation. Since the question is local on $\mathrm{X}$ and $\mathrm{Y}$, we may suppose that $\mathrm{X}$ and $\mathrm{Y}$ are affine and, by \sref{IV.11.3.1}, that $\sh{H}$ is $f$-flat. We now record the following lemma.

\begin{lemma}[11.3.9.1]
\label{IV.11.3.9.1}
Let $f:\mathrm{X}\to\mathrm{Y}$ be a morphism locally of finite presentation, $\sh{G}$, $\sh{H}$ two $\sh{O}_\mathrm{X}$-Modules quasi-coherent of finite presentation such that $\sh{H}$ is $f$-flat, $p:\sh{G}\to\sh{H}$ a surjective homomorphism. Then $\operatorname{Ker}(p)$ is an $\sh{O}_\mathrm{X}$-Module of finite presentation.
\end{lemma}

\begin{proof}
We may indeed suppose that $\mathrm{X}$ and $\mathrm{Y}$ are affine and that there exist a morphism $\mathrm{Y}\to\mathrm{Y}_0$, where $\mathrm{Y}_0$ is Noetherian, a morphism $f_0:\mathrm{X}_0\to\mathrm{Y}_0$ of finite type such that $\mathrm{X}=\mathrm{X}_0\times_{\mathrm{Y}_0}\mathrm{Y}$ and $f=(f_0)_{(\mathrm{Y})}$, two coherent $\sh{O}_{\mathrm{X}_0}$-Modules $\sh{G}_0$, $\sh{H}_0$, and a homomorphism $p_0:\sh{G}_0\to\sh{H}_0$ such that $\sh{G}$, $\sh{H}$, and $p$ are obtained from $\sh{G}_0$, $\sh{H}_0$, and $p_0$ by base change $(8.9.1\text{ and }8.5.2)$. We may moreover suppose that $p_0$ is surjective $(8.5.7)$ and that $\sh{H}_0$ is $f_0$-flat \sref{IV.11.2.7}. Then, if $\sh{K}_0=\operatorname{Ker}(p_0)$, $\sh{K}_0$ is a coherent $\sh{O}_{\mathrm{X}_0}$-Module $(\mathbf{0}_\mathrm{I},~5.3.4)$, and by $(\mathbf{0}_\mathrm{I},~6.1.2)$, $\sh{K}=\operatorname{Ker}(p)$ is obtained from $\sh{K}_0$ by base change and is therefore of finite presentation.
\end{proof}

This lemma being proved, set $\sh{K}=\operatorname{Ker}(\sh{G}\to\sh{H})$, which is therefore of finite presentation. There is a canonical homomorphism $q:\sh{F}\to\sh{K}$, and by hypothesis $q_x:\sh{F}_x\to\sh{K}_x$ is an isomorphism. Hence, by $(\mathbf{0}_\mathrm{I},~5.2.7)$, there exists a neighbourhood $\mathrm{U}$ of $x$ such that $q|_{\mathrm{U}}$ is an isomorphism; after restricting $\mathrm{X}$, we may therefore suppose that the sequence
\[
  0\to\sh{F}\xrightarrow{h}\sh{G}\to\sh{H}\to 0
\]
is \emph{exact}. This being so, it follows from \sref{IV.11.2.6} that, for all sufficiently large $\lambda$, $\sh{H}_\lambda$ is $f_\lambda$-flat. Set $\sh{K}_\lambda=\operatorname{Ker}(\sh{G}_\lambda\to\sh{H}_\lambda)$ and $\sh{K}_\mu=v^*_{\lambda\mu}(\sh{K}_\lambda)$ for $\mu\geqslant\lambda$. It follows from $(\mathbf{0}_\mathrm{I},~6.1.2)$ that
\[
  \sh{K}_\mu=\operatorname{Ker}(\sh{G}_\mu\to\sh{H}_\mu)
  \quad\text{and}\quad
  \sh{K}=v^*_\lambda(\sh{K}_\lambda)
  =\operatorname{Ker}(\sh{G}\to\sh{H})=\sh{F}.
\]
On the other hand, for every $\mu\geqslant\lambda$ there is a canonical homomorphism $q_\mu:\sh{F}_\mu\to\sh{K}_\mu$ such that $q_\mu=v^*_{\lambda\mu}(q_\lambda)$ and $q=v^*_\mu(q_\mu)$. Since $q$ is an isomorphism, it follows from $(8.5.2.4)$ and \sref{IV.11.3.9.1} that there exists $\mu\geqslant\lambda$ for which $q_\mu$ is an isomorphism, and hence $h_\mu:\sh{F}_\mu\to\sh{G}_\mu$ is injective.
\end{proof}

Let us return to the proof of \sref{IV.11.3.8}.

Since the set of $x\in\mathrm{X}$ satisfying $b)$ is open in $\mathrm{X}$, it is clear that $b)$ implies $b')$. Let us finally show that $b')$ implies $c)$. First, $\sh{F}_n$ is $f$-flat in a neighbourhood of $x$ \sref{IV.11.3.1}; we may therefore reduce to the case where $\mathrm{U}=\mathrm{X}$ and $\sh{F}_n$ is $f$-flat. Since, by definition, $\operatorname{gr}^*_\mathscr{J}(\sh{F})$ is isomorphic to $\sh{F}_n\otimes_{\sh{O}_\mathrm{Y}}\sh{O}_\mathrm{Y}[\mathrm{T}_1,\ldots,\mathrm{T}_n]$ $(\mathbf{0},~15.2.2)$, it is $f$-flat, and it follows from \sref{IV.11.3.4},~(i), that $\sh{F}$ itself is $f$-flat in a neighbourhood of $x$. On the other hand, if $(\mathscr{J}_y)_x$ is the ideal of $\sh{O}_{\mathrm{X}_y,x}$ generated by the $(g_i\otimes 1)_x$, it follows from \sref{IV.11.2.9.2} that, in the diagram
\[
\xymatrix{
  ((\operatorname{gr}^0_{\mathscr{J}_x}(\sh{F}_x))[\mathrm{T}_1,\ldots,\mathrm{T}_n])\otimes_{\sh{O}_y}\boldsymbol{k}(y) \ar[r] \ar[d] & \operatorname{gr}^*_{\mathscr{J}_x}(\sh{F}_x)\otimes_{\sh{O}_y}\boldsymbol{k}(y) \ar[d] \\
  (\operatorname{gr}^0_{(\mathscr{J}_y)_x}((\sh{F}_y)_x))[\mathrm{T}_1,\ldots,\mathrm{T}_n] \ar[r] & \operatorname{gr}^*_{(\mathscr{J}_y)_x}((\sh{F}_y)_x)
}
\]
\oldpage[IV-3]{138}
the vertical arrows are isomorphisms. Since the first horizontal arrow is an isomorphism, so is the second; hence the sequence $((g_i\otimes 1)_x)$ is $(\sh{F}_y)_x$-quasi-regular, and therefore regular, since $\mathrm{X}_y$ is locally of finite type over $\boldsymbol{k}(y)$.
\end{proof}

\begin{theorem}[11.3.10]
\label{IV.11.3.10}
(Flatness criterion by fibres). --- Let $\mathrm{S}$ be a prescheme, $g:\mathrm{X}\to\mathrm{S}$, $h:\mathrm{Y}\to\mathrm{S}$ two morphisms, $f:\mathrm{X}\to\mathrm{Y}$ an $\mathrm{S}$-morphism, $\sh{F}$ a quasi-coherent $\sh{O}_\mathrm{X}$-Module, $x$ a point of $\mathrm{X}$, $y=f(x)$, $s=h(y)=g(x)$. Suppose that one of the following two hypotheses holds:
\begin{enumerate}[label=\arabic{enumi}\textsuperscript{o}]
  \item $\mathrm{S}$, $\mathrm{X}$ and $\mathrm{Y}$ are locally Noetherian and $\sh{F}$ is coherent.
  \item $g$ and $h$ are locally of finite presentation and $\sh{F}$ is of finite presentation.
\end{enumerate}
Then, with the notation of $(9.4.1)$, if $\sh{F}_x\neq 0$, the following conditions are equivalent:
\begin{enumerate}[label=\emph{\alph*)}]
  \item $\sh{F}$ is $g$-flat at the point $x$ and $\sh{F}_s$ is $f_s$-flat at the point $x$.
  \item $h$ is flat at the point $y$ and $\sh{F}$ is $f$-flat at the point $x$.
\end{enumerate}
Moreover, under hypothesis 2\textsuperscript{o}, the set of $x\in\mathrm{X}$ satisfying condition $b)$ is open in $\mathrm{X}$.
\end{theorem}

\begin{proof}
The last assertion results from \sref{IV.11.3.1} applied to $\sh{O}_\mathrm{Y}$ and to the morphism $h$ on the one hand, and to $\sh{F}$ and to the morphism $f$ (which is locally of finite presentation) on the other.

As $g=h\circ f$, it is clear that $b)$ implies $a)$ without assuming 1\textsuperscript{o} or 2\textsuperscript{o} $(2.1.6$ and $2.1.4)$. To prove that $a)$ implies $b)$, we may reduce to the case where $\mathrm{S}$, $\mathrm{X}$, and $\mathrm{Y}$ are affine; under hypothesis 2\textsuperscript{o}, by applying \sref{IV.11.2.7}, we reduce further to the case where $\mathrm{S}$ is \emph{Noetherian}, that is, to the case where hypothesis 1\textsuperscript{o} holds. The assertion is then equivalent to the following lemma, which improves $(\mathbf{0}_{\mathrm{III}},~10.2.5)$:

\begin{lemma}[11.3.10.1]
\label{IV.11.3.10.1}
Let $\rho:\mathrm{A}\to\mathrm{B}$, $\sigma:\mathrm{B}\to\mathrm{C}$ be homomorphisms of local Noetherian rings, $k$ the residue field of $\mathrm{A}$, $\mathrm{M}$ a $\mathrm{C}$-module $\neq 0$ of finite type. Then the following conditions are equivalent:
\begin{enumerate}[label=\emph{\alph*)}]
  \item $\mathrm{M}$ is a flat $\mathrm{A}$-module and $\mathrm{M}\otimes_\mathrm{A} k$ is a flat $(\mathrm{B}\otimes_\mathrm{A} k)$-module.
  \item $\mathrm{B}$ is a flat $\mathrm{A}$-module and $\mathrm{M}$ is a flat $\mathrm{B}$-module.
\end{enumerate}
\end{lemma}

We will first establish the more general following lemma:

\begin{lemma}[11.3.10.2]
\label{IV.11.3.10.2}
Let $\mathrm{A}$ be a commutative ring, $\mathrm{B}$ a commutative $\mathrm{A}$-algebra, $\mathfrak{J}$ an ideal of $\mathrm{A}$, and $\mathrm{M}$ a $\mathrm{B}$-module. Consider first the following conditions:
\begin{enumerate}[label=\emph{(\roman*)}]
  \item $\mathfrak{J}$ is nilpotent.
  \item $\mathrm{B}$ is Noetherian and $\mathrm{M}$ is ideally separated for the $\mathfrak{J}\mathrm{B}$-preadic topology $(\mathbf{0}_{\mathrm{III}},~10.2.1)$.
  \item $\mathfrak{J}\mathrm{B}$ is contained in the radical of $\mathrm{B}$.
\end{enumerate}
Consider, on the other hand, the following four properties:
\begin{enumerate}[label=\emph{\alph*)}]
  \item $\mathrm{M}$ is a flat $\mathrm{B}$-module.
  \item $\mathrm{B}$ is a flat $\mathrm{A}$-module.
  \item $\mathrm{M}$ is a flat $\mathrm{A}$-module and $\mathrm{M}/\mathfrak{J}\mathrm{M}$ is a flat $(\mathrm{B}/\mathfrak{J}\mathrm{B})$-module.
  \item $\mathrm{B}/\mathfrak{J}\mathrm{B}$ is a flat $(\mathrm{A}/\mathfrak{J})$-module and for every maximal ideal $\mathfrak{m}\supset\mathfrak{J}\mathrm{B}$ of $\mathrm{B}$ one has $\mathfrak{m}\mathrm{M}\neq\mathrm{M}$.
\end{enumerate}
Then:
\begin{enumerate}[label=\arabic{enumi}\textsuperscript{o}]
  \item If one of the conditions (i) or (ii) holds, the conjunction of $a)$ and $b)$ implies $c)$, and $c)$ implies $a)$.
  \item \oldpage[IV-3]{139} If the condition (i) or the conjunction of (ii) and (iii) is verified, the conjunction of $c)$ and $d)$ implies the conjunction of $a)$ and $b)$.
\end{enumerate}
\end{lemma}

\begin{proof}
1\textsuperscript{o} The first assertion is immediate $(\mathbf{0}_\mathrm{I},~6.2.1)$. Suppose therefore $c)$ verified, and let us prove $a)$. Consider the graded rings $\operatorname{gr}^*_\mathfrak{J}(\mathrm{A})$, $\operatorname{gr}^*_\mathfrak{J}(\mathrm{B})$ and the graded module $\operatorname{gr}^*_\mathfrak{J}(\mathrm{M})$ (over both $\operatorname{gr}^*_\mathfrak{J}(\mathrm{A})$ and $\operatorname{gr}^*_\mathfrak{J}(\mathrm{B})$) relative to the $\mathfrak{J}$-preadic filtrations, as well as the surjective canonical maps $(\mathbf{0}_{\mathrm{III}},~10.1.1.2)$
\begin{align*}
  u &: \operatorname{gr}^0(\mathrm{B})\otimes_{\operatorname{gr}^0(\mathrm{A})}\operatorname{gr}^*(\mathrm{A})\to\operatorname{gr}^*(\mathrm{B}) \\
  v &: \operatorname{gr}^0(\mathrm{M})\otimes_{\operatorname{gr}^0(\mathrm{A})}\operatorname{gr}^*(\mathrm{A})\to\operatorname{gr}^*(\mathrm{M}) \\
  w &: \operatorname{gr}^0(\mathrm{M})\otimes_{\operatorname{gr}^0(\mathrm{B})}\operatorname{gr}^*(\mathrm{B})\to\operatorname{gr}^*(\mathrm{M})
\end{align*}
It is clear there is a commutative diagram
\begin{equation}
\xymatrix{
  \operatorname{gr}^0(\mathrm{M})\otimes_{\operatorname{gr}^0(\mathrm{A})}\operatorname{gr}^*(\mathrm{A}) \ar[rr]^{v} \ar[dr]_{1\otimes u} & & \operatorname{gr}^*(\mathrm{M}) \\
  & \operatorname{gr}^0(\mathrm{M})\otimes_{\operatorname{gr}^0(\mathrm{B})}\operatorname{gr}^*(\mathrm{B}) \ar[ur]_{w} &
}
  \label{IV.11.3.10.3}
\tag{11.3.10.3}
\end{equation}

The hypothesis $c)$ implies that $v$ is bijective $(\mathbf{0}_{\mathrm{III}},~10.2.1)$; since the other two maps in the diagram are surjective, they too are bijective. Since, by hypothesis $c)$, $\mathrm{M}/\mathfrak{J}\mathrm{M}$ is a flat $(\mathrm{B}/\mathfrak{J}\mathrm{B})$-module, it follows from $(\mathbf{0}_{\mathrm{III}},~10.2.1)$ that $\mathrm{M}$ is a flat $\mathrm{B}$-module.

2\textsuperscript{o} Either of the conditions (i), (iii) implies that every maximal ideal of $\mathrm{B}$ contains $\mathfrak{J}\mathrm{B}$. It therefore follows from 1\textsuperscript{o} and the conjunction of $c)$ and $d)$ that $\mathrm{M}$ is a \emph{faithfully flat} $\mathrm{B}$-module, and consequently that $\operatorname{gr}^0(\mathrm{M})$ is a \emph{faithfully flat} $\operatorname{gr}^0(\mathrm{B})$-module $(\mathbf{0}_{\mathrm{I}},~6.2.1)$. We saw in 1\textsuperscript{o} that hypothesis $c)$ implies that the three maps in diagram \textup{(11.3.10.3)} are bijective; the faithful flatness of $\operatorname{gr}^0(\mathrm{M})$ over $\operatorname{gr}^0(\mathrm{B})$ therefore implies that $u$ is bijective $(\mathbf{0}_{\mathrm{I}},~6.4.1)$. On the other hand, conditions (ii) and (iii) imply that $\mathrm{B}$, as an $\mathrm{A}$-module, is ideally separated for the $\mathfrak{J}$-preadic filtration (Bourbaki, \emph{Alg.\ comm.}, chap.~III, \textsection\,5, n\textsuperscript{o}\,4, prop.~2). It follows again from $(\mathbf{0}_{\mathrm{III}},~10.2.1)$ that, if condition (i), or the conjunction of (ii) and (iii), holds, $\mathrm{B}$ is a flat $\mathrm{A}$-module.
\end{proof}

\phantomsection\label{IV.11.3.10.4}
\textbf{(11.3.10.4)} Lemma \sref{IV.11.3.10.2} being established, one deduces \sref{IV.11.3.10.1} by taking for $\mathfrak{J}$ the maximal ideal of $\mathrm{A}$ and observing that conditions (ii) and (iii) of \sref{IV.11.3.10.2} are then satisfied (Bourbaki, \emph{Alg.\ comm.}, chap.~III, \textsection\,5, n\textsuperscript{o}\,4, prop.~2). This therefore also completes the proof of \sref{IV.11.3.10}.
\end{proof}

\begin{corollary}[11.3.11]
\label{IV.11.3.11}
Let $g:\mathrm{X}\to\mathrm{S}$, $h:\mathrm{Y}\to\mathrm{S}$ be two morphisms locally of finite presentation, $f:\mathrm{X}\to\mathrm{Y}$ an $\mathrm{S}$-morphism. The following conditions are equivalent:
\begin{enumerate}[label=\emph{\alph*)}]
  \item $g$ is flat and $f_s:\mathrm{X}_s\to\mathrm{Y}_s$ is flat for every $s\in\mathrm{S}$.
  \item $h$ is flat at all the points of $f(\mathrm{X})$ and $f$ is flat.
\end{enumerate}
\end{corollary}

\begin{proof}
It suffices to apply \sref{IV.11.3.10} for $\sh{F}=\sh{O}_\mathrm{X}$.
\end{proof}

\begin{remark}[11.3.12]
\label{IV.11.3.12}
It would be interesting to be able to give proofs of \sref{IV.11.3.2} and \sref{IV.11.3.10} not using passage to the inductive limit; it would suffice for this to prove the following criterion:

Let $\mathrm{A}$ be a ring, $\mathrm{B}$ an $\mathrm{A}$-algebra of finite presentation, $\mathrm{M}$ a $\mathrm{B}$-module of finite presentation, $\mathfrak{J}$ an \oldpage[IV-3]{140} ideal of $\mathrm{A}$, and $\mathfrak{p}$ a prime ideal of $\mathrm{B}$ containing $\mathfrak{J}\mathrm{B}$. Then the following two conditions are equivalent:
\begin{enumerate}[label=\emph{\alph*)}]
  \item $\mathrm{M}_\mathfrak{p}$ is a flat $\mathrm{A}$-module.
  \item $\mathrm{M}_\mathfrak{p}/\mathfrak{J}\mathrm{M}_\mathfrak{p}$ is a flat $(\mathrm{A}/\mathfrak{J})$-module and $\operatorname{Tor}_1^\mathrm{A}(\mathrm{A}/\mathfrak{J},\mathrm{M}_\mathfrak{p})=0$.
\end{enumerate}
When $\mathrm{A}$ is \emph{Noetherian}, this follows from $(\mathbf{0}_{\mathrm{III}},~10.2.2)$ applied to the Noetherian $\mathrm{A}$-algebra $\mathrm{B}_\mathfrak{p}$; can the general statement be deduced from this by passage to the inductive limit?

It should be noted in this connection that such a generalization is certainly not possible if condition $b)$ above is replaced by either of conditions $c)$, $d)$ of $(\mathbf{0}_{\mathrm{III}},~10.2.1)$, with $\mathrm{M}$ replaced by $\mathrm{M}_\mathfrak{p}$. For example, let $\mathrm{A}$ be a local ring whose maximal ideal $\mathfrak{J}$ is principal and such that
\[
  \mathfrak{R}=\bigcap_{n\geqslant 0}\mathfrak{J}^{n+1}
\]
is nonzero (for instance, the ring of germs of indefinitely differentiable real-valued functions in a neighbourhood of $0$ in $\mathbf{R}$). Take $\mathrm{B}=\mathrm{A}$, $\mathfrak{p}=\mathfrak{J}$, and $\mathrm{M}=\mathrm{A}/\mathfrak{N}$, where $\mathfrak{N}$ is a cyclic nonzero submodule of $\mathfrak{R}$; $\mathrm{M}$ is therefore of finite presentation. It is clear that $\mathrm{M}$ is not a flat $\mathrm{A}$-module, for, being of finite presentation, it would then be free (Bourbaki, \emph{Alg.\ comm.}, chap.~II, \textsection\,3, n\textsuperscript{o}\,2, cor.~2 to prop.~5), which is absurd since $\mathfrak{N}\neq 0$. Nevertheless,
\[
  \mathfrak{J}^k\mathrm{M}/\mathfrak{J}^{k+n}\mathrm{M}\simeq \mathrm{A}/\mathfrak{J}^n
\]
for all positive integers $k$ and $n$, so $\mathrm{M}$ does satisfy conditions $c)$ and $d)$ of $(\mathbf{0}_{\mathrm{III}},~10.2.1)$, since $\mathrm{A}/\mathfrak{J}$ is a field.
\end{remark}

\begin{proposition}[11.3.13]
\label{IV.11.3.13}
Let $f:\mathrm{X}\to\mathrm{Y}$ be a morphism of preschemes and let $x$ be a point of $\mathrm{X}$ such that $f$ is flat at $x$; set $y=f(x)$.
\begin{enumerate}[label=\emph{(\roman*)}]
  \item If $\mathrm{X}$ is reduced (resp.\ normal) at the point $x$, then $\mathrm{Y}$ is reduced (resp.\ normal) at the point $y$.
  \item Suppose in addition that $f$ is of finite presentation at the point $x$. Then, if $\mathrm{Y}$ is reduced (resp.\ normal) at the point $y$ and if the morphism $f$ is reduced (resp.\ normal) at the point $x$ $(6.8.1)$, $\mathrm{X}$ is reduced (resp.\ normal) at the point $x$.
\end{enumerate}
\end{proposition}

\begin{proof}
The first assertion is recalled here only for convenience, having already been proved in $(2.1.13)$. To prove (ii), one may restrict to the case where $\mathrm{X}=\Spec(\mathrm{B})$ and $\mathrm{Y}=\Spec(\mathrm{A})$, where $\mathrm{A}$ is a reduced local ring (resp.\ an integral and integrally closed local ring) and $\mathrm{B}$ is an $\mathrm{A}$-algebra of finite presentation. One may then write $(8.9.1)$
\[
  \mathrm{B}=\mathrm{B}_0\otimes_{\mathrm{A}_0}\mathrm{A},
\]
where $\mathrm{A}_0$ is a finite type $\mathbf{Z}$-subalgebra of $\mathrm{A}$ and $\mathrm{B}_0$ is a finite type $\mathrm{A}_0$-algebra. Let $(\mathrm{C}_\alpha)$ be the filtered increasing family of finite type $\mathrm{A}_0$-subalgebras of $\mathrm{A}$, which are therefore finite type $\mathbf{Z}$-algebras; one has $\mathrm{A}=\varinjlim\mathrm{C}_\alpha$. We distinguish the two cases:

1\textsuperscript{o} Suppose that $\mathrm{A}$ is reduced and that $f$ is reduced at $x$. If $\mathfrak{m}$ is the maximal ideal of $\mathrm{A}$, let $\mathfrak{p}_\alpha=\mathfrak{m}\cap\mathrm{C}_\alpha$ and set $\mathrm{A}'_\alpha=(\mathrm{C}_\alpha)_{\mathfrak{p}_\alpha}$, so that one also has $\mathrm{A}=\varinjlim\mathrm{A}'_\alpha$ $(5.13.3)$. Set $\mathrm{Y}_\alpha=\Spec(\mathrm{A}'_\alpha)$ and $\mathrm{X}_\alpha=\Spec(\mathrm{B}_0\otimes_{\mathrm{A}_0}\mathrm{A}'_\alpha)$, and let $f_\alpha:\mathrm{X}_\alpha\to\mathrm{Y}_\alpha$, $x_\alpha$ (resp.\ $y_\alpha$) be the projection of $x$ (resp.\ $y$) in $\mathrm{X}_\alpha$ (resp.\ $\mathrm{Y}_\alpha$). Since $f$ is reduced at $x$, there exists $\alpha_0$ such that $f_\alpha$ is reduced at $x_\alpha$ for $\alpha\geqslant\alpha_0$ $((6.7.8)$ and \sref{IV.11.2.6}); since $\mathrm{Y}_\alpha$ and $\mathrm{X}_\alpha$ are Noetherian and $\mathrm{A}'_\alpha$ is reduced (because $\mathrm{C}_\alpha$ is), it follows from $(3.3.5)$ that $\mathrm{X}_\alpha$ is reduced at $x_\alpha$. But
\[
  \mathrm{B}=\varinjlim(\mathrm{B}_0\otimes_{\mathrm{A}_0}\mathrm{A}'_\alpha)
\]
and consequently $\sh{O}_{\mathrm{X},x}=\varinjlim\sh{O}_{\mathrm{X}_\alpha,x_\alpha}$ $(5.13.3)$; hence $\sh{O}_{\mathrm{X},x}$ is reduced $(5.13.2)$.

2\textsuperscript{o} Suppose that $\mathrm{A}$ is integrally closed and that $f$ is normal at $x$. Since $\mathrm{C}_\alpha$ is universally Japanese $(7.7.4)$, its integral closure $\mathrm{C}'_\alpha$ is a finite $\mathrm{C}_\alpha$-algebra, evidently contained in $\mathrm{A}$. Let $\mathfrak{p}'_\alpha=\mathfrak{m}\cap\mathrm{C}'_\alpha$ and set $\mathrm{A}''_\alpha=(\mathrm{C}'_\alpha)_{\mathfrak{p}'_\alpha}$, so that $\mathrm{A}''_\alpha$ is a local Noetherian integral and integrally closed ring, and $\mathrm{A}=\varinjlim\mathrm{A}''_\alpha$ $(5.13.3)$. Set $\mathrm{Y}'_\alpha=\Spec(\mathrm{A}''_\alpha)$ and $\mathrm{X}'_\alpha=\Spec(\mathrm{B}_0\otimes_{\mathrm{A}_0}\mathrm{A}''_\alpha)$, and let $f'_\alpha:\mathrm{X}'_\alpha\to\mathrm{Y}'_\alpha$, $x'_\alpha$ (resp.\ $y'_\alpha$) be the projection of $x$ (resp.\ $y$) in $\mathrm{X}'_\alpha$ (resp.\ $\mathrm{Y}'_\alpha$). Since $f$ is normal at $x$, there exists $\alpha_0$ such that, for $\alpha\geqslant\alpha_0$, $f'_\alpha$ is normal at $x'_\alpha$ $((6.7.8)$ and \sref{IV.11.2.6}); \oldpage[IV-3]{141} since $\mathrm{X}'_\alpha$ and $\mathrm{Y}'_\alpha$ are Noetherian and $\mathrm{A}''_\alpha$ is integrally closed, it follows from $(6.5.4)$ that $\mathrm{X}'_\alpha$ is normal at $x'_\alpha$. But the morphisms
\[
  \Spec(\mathrm{A}''_\beta)\longrightarrow\Spec(\mathrm{A}''_\alpha)
  \qquad(\alpha\leqslant\beta)
\]
are dominant; hence, since by virtue of \sref{IV.11.3.1} we may suppose the $f'_\alpha$ flat for $\alpha\geqslant\alpha_0$, it follows from $(2.3.7)$ that every irreducible component of $\mathrm{X}'_\beta$ dominates an irreducible component of $\mathrm{X}'_\alpha$. We then conclude from $(5.13.4)$, applied to the schemes
\[
  \Spec\!\left(\sh{O}_{\mathrm{X}_0,x_0}\otimes_{\mathrm{A}_0}\mathrm{A}''_\alpha\right),
\]
that $\mathrm{X}$ is normal at $x$.
\end{proof}

\begin{corollary}[11.3.14]
\label{IV.11.3.14}
Let $f:\mathrm{X}\to\mathrm{Y}$ be a morphism locally of finite presentation, $x$ a point of $\mathrm{X}$, and $y=f(x)$. If $\mathrm{Y}$ is geometrically unibranch at $y$ $(6.15.1)$ and if $f$ is normal at $x$ $(6.8.1)$, then $\mathrm{X}$ is geometrically unibranch at $x$.
\end{corollary}

\begin{proof}
We may evidently reduce to the case where $\mathrm{Y}=\Spec(\mathrm{A})$, with $\mathrm{A}$ a geometrically unibranch local integral domain and $y$ the closed point of $\mathrm{Y}$. Let $\mathrm{A}'$ be the integral closure of $\mathrm{A}$; set $\mathrm{Y}'=\Spec(\mathrm{A}')$ and let $y'$ be the closed point of $\mathrm{Y}'$. Thus the morphism $g:\mathrm{Y}'\to\mathrm{Y}$ is radicial at $y$ $(6.15.3)$, integral, and birational. Set $\mathrm{X}'=\mathrm{X}\times_\mathrm{Y}\mathrm{Y}'$, and let $f':\mathrm{X}'\to\mathrm{Y}'$ and $g':\mathrm{X}'\to\mathrm{X}$ be the projections. The morphism $g'$ is integral and radicial at $x$ $(6.15.3.1)$. Let $x'$ therefore be the unique point of $g'^{-1}(x)$ lying over $y'$. The morphism $f'$ is of finite presentation and normal at $x'$ $(6.7.8)$; consequently the local ring $\sh{O}_{\mathrm{X}',x'}$ is an integrally closed integral domain \sref{IV.11.3.13}. On the other hand, we may suppose $f$ flat \sref{IV.11.3.1}, and hence $g'$ is birational $(6.15.4.1)$. It follows that $\Spec(\sh{O}_{\mathrm{X},x})$ is irreducible; since $\sh{O}_{\mathrm{X},x}$ is reduced \sref{IV.11.3.13}, it is an integral domain and is geometrically unibranch by $(6.15.5)$.
\end{proof}

\begin{proposition}[11.3.15]
\label{IV.11.3.15}
Let $\mathrm{A}$ be a ring, $\mathrm{B}$ an $\mathrm{A}$-algebra of finite presentation, and $\mathrm{M}$ a $\mathrm{B}$-module of finite presentation which is flat as an $\mathrm{A}$-module. Then there exists a finite sequence $(f_i)_{1\leqslant i\leqslant n}$ of elements of $\mathrm{A}$ such that the ideal generated by the $f_i$ is $\mathrm{A}$ and, for $0\leqslant i<n$,
\[
  \frac{\mathrm{M}_{f_{i+1}}}
       {\sum_{j=1}^{i}f_j\mathrm{M}_{f_{i+1}}}
\]
is a free module over
\[
  \frac{\mathrm{A}_{f_{i+1}}}
       {\sum_{j=1}^{i}f_j\mathrm{A}_{f_{i+1}}}.
\]
\end{proposition}

Equivalently, the $\mathrm{D}(f_i)$ form an open covering of $\mathrm{X}=\Spec(\mathrm{A})$ and, if one sets $\mathrm{Z}_1=\mathrm{D}(f_1)$ and then, inductively,
\[
  \mathrm{Z}_{i+1}
  =\mathrm{D}(f_{i+1})
   \cap\mathrm{V}(f_1\mathrm{A}+\cdots+f_i\mathrm{A}),
\]
the $\mathrm{Z}_i$ form a partition of $\mathrm{X}$ into locally closed subsets, while
\[
  \frac{\mathrm{A}_{f_{i+1}}}
       {\sum_{j=1}^{i}f_j\mathrm{A}_{f_{i+1}}}
\]
is the ring of an affine subscheme of $\mathrm{X}$ whose underlying space is $\mathrm{Z}_{i+1}$.

\begin{proof}
By virtue of \sref{IV.11.2.7}, we may first suppose that there exist a Noetherian subring $\mathrm{A}_0$ of $\mathrm{A}$, a finite type $\mathrm{A}_0$-algebra $\mathrm{B}_0$, and a finite type $\mathrm{B}_0$-module $\mathrm{M}_0$, flat over $\mathrm{A}_0$, such that
\[
  \mathrm{B}=\mathrm{B}_0\otimes_{\mathrm{A}_0}\mathrm{A},
  \qquad
  \mathrm{M}=\mathrm{M}_0\otimes_{\mathrm{A}_0}\mathrm{A}.
\]
It clearly suffices to prove the proposition for $\mathrm{A}_0$, $\mathrm{B}_0$, and $\mathrm{M}_0$: if elements $f_i\in\mathrm{A}_0$ satisfy the conditions in the statement in that case, then they also satisfy them for $\mathrm{A}$, $\mathrm{B}$, and $\mathrm{M}$, since
\[
  \frac{\mathrm{M}_{f_{i+1}}}
       {\sum_{j=1}^{i}f_j\mathrm{M}_{f_{i+1}}}
  =
  \left(
    \frac{(\mathrm{M}_0)_{f_{i+1}}}
         {\sum_{j=1}^{i}f_j(\mathrm{M}_0)_{f_{i+1}}}
  \right)\otimes_{\mathrm{A}_0}\mathrm{A}
\]
(Bourbaki, \emph{Alg.\ comm.}, chap.~II, \textsection\,2, n\textsuperscript{o}\,7, prop.~18). We may therefore henceforth restrict to the case where $\mathrm{A}$ is Noetherian.

Let us now note that if $\mathrm{C}$ is a Noetherian ring, $\mathfrak{N}$ its nilradical, and $\mathrm{P}$ a flat $\mathrm{C}$-module, then it follows from $(\mathbf{0}_{\mathrm{III}},10.1.2)$ that, for $\mathrm{P}$ to be a free $\mathrm{C}$-module, \oldpage[IV-3]{142} it is necessary and sufficient that $\mathrm{P}\otimes_\mathrm{C}(\mathrm{C}/\mathfrak{N})$ be a free $(\mathrm{C}/\mathfrak{N})$-module. Note also that if $\mathrm{C}$ is a reduced Noetherian ring, there exists $g\in\mathrm{C}$ such that $\mathrm{C}_g$ is an integral domain.

We now apply Lemma $(6.9.2)$. A sequence $(f_i)_{i\geqslant1}$ of elements of $\mathrm{A}$ may be defined inductively as follows:
\begin{enumerate}
\item[$1^\circ$] $f_1$ is chosen so that
\[
  \mathrm{A}_1=(\mathrm{A}_{\mathrm{red}})_{f_1}
\]
is an integral domain and $\mathrm{M}\otimes_\mathrm{A}\mathrm{A}_1$ is a free $\mathrm{A}_1$-module;
\item[$2^\circ$] if the ideal $\mathfrak{J}_i$ generated by $f_1,\ldots,f_i$ is equal to $\mathrm{A}$, set $f_{i+1}=f_i$;
\item[$3^\circ$] if, on the contrary, $\mathfrak{J}_i\neq\mathrm{A}$, choose $f_{i+1}\notin\mathfrak{J}_i$ so that
\[
  \mathrm{A}_{i+1}
  =\bigl((\mathrm{A}/\mathfrak{J}_i)_{\mathrm{red}}\bigr)_{f_{i+1}}
\]
is an integral domain and $\mathrm{M}\otimes_\mathrm{A}\mathrm{A}_{i+1}$ is a free $\mathrm{A}_{i+1}$-module.
\end{enumerate}
Since $\mathrm{A}$ is Noetherian, the increasing sequence $(\mathfrak{J}_i)$ is stationary. Hence there exists $n$ such that $f_1,\ldots,f_n$ generate the unit ideal of $\mathrm{A}$, and these elements satisfy the required conditions because
\[
  \bigl((\mathrm{A}/\mathfrak{J}_i)_{\mathrm{red}}\bigr)_{f_{i+1}}
  =
  \bigl((\mathrm{A}/\mathfrak{J}_i)_{f_{i+1}}\bigr)_{\mathrm{red}}.
\]
\end{proof}

\begin{proposition}[11.3.16]
\label{IV.11.3.16}
Let $f:\mathrm{X}\to\mathrm{Y}$ be a faithfully flat morphism of finite presentation, and let $g:\mathrm{Y}\to\mathrm{Z}$ be a morphism such that $g\circ f:\mathrm{X}\to\mathrm{Z}$ is of finite type (resp.\ of finite presentation). Then $g$ is of finite type (resp.\ of finite presentation).
\end{proposition}

\begin{proof}
Since $f$ is surjective and $g\circ f$ is quasi-compact, $g$ is quasi-compact $(1.1.3)$. Let us first show that if $g\circ f$ is of finite presentation, then $g$ is quasi-separated. Indeed, let $\mathrm{W}\subset\mathrm{Z}$ be an affine open subset and let $(\mathrm{V}_i)$ be a finite affine open covering of $g^{-1}(\mathrm{W})$; it is enough to show that the $\mathrm{V}_i\cap\mathrm{V}_j$ are quasi-compact $(1.2.6,\ 1.2.7)$. For each $i$, let $(\mathrm{U}_{ih})$ be a finite affine open covering of $f^{-1}(\mathrm{V}_i)$. Since $g\circ f$ is of finite presentation, all the intersections $\mathrm{U}_{ih}\cap\mathrm{U}_{jk}$ are quasi-compact. Since $f$ is surjective, $\mathrm{V}_i\cap\mathrm{V}_j$ is the union, as $h$ and $k$ vary, of the images $f(\mathrm{U}_{ih}\cap\mathrm{U}_{jk})$; it is therefore quasi-compact because $f$ is continuous.

The question is thus local on $\mathrm{Y}$, and we may suppose that $\mathrm{Y}=\Spec(\mathrm{B})$ and $\mathrm{Z}=\Spec(\mathrm{A})$ are affine, while $\mathrm{X}$ is the union of finitely many affine open subsets $\mathrm{X}_i=\Spec(\mathrm{C}_i)$, where the $\mathrm{C}_i$ are $\mathrm{B}$-algebras of finite type (resp.\ of finite presentation). Let $\mathrm{X}'$ be the disjoint union of the $\mathrm{X}_i$, and let $p:\mathrm{X}'\to\mathrm{X}$ be the morphism which on each $\mathrm{X}_i$ is the canonical open immersion. The morphism $p$ is faithfully flat and of finite presentation $(1.6.5)$; consequently $g\circ f\circ p$ is of finite type (resp.\ of finite presentation), and $f\circ p$ is faithfully flat and of finite presentation. We may therefore suppose in addition that $\mathrm{X}=\Spec(\mathrm{C})$ is affine, and the proposition is reduced to the following corollary.
\end{proof}

\begin{corollary}[11.3.17]
\label{IV.11.3.17}
Let $\mathrm{A}$ be a ring, $\mathrm{B}$ an $\mathrm{A}$-algebra, and $\mathrm{C}$ a $\mathrm{B}$-algebra of finite presentation which is faithfully flat as a $\mathrm{B}$-module. If $\mathrm{C}$ is an $\mathrm{A}$-algebra of finite type (resp.\ of finite presentation), then $\mathrm{B}$ is an $\mathrm{A}$-algebra of finite type (resp.\ of finite presentation).
\end{corollary}

\begin{proof}
Suppose first that $\mathrm{C}$ is of finite type over $\mathrm{A}$. Let $(\mathrm{B}_\lambda)$ be the increasing filtered family of finite type $\mathrm{A}$-subalgebras of $\mathrm{B}$. By $(8.8.2)$, there exists an index $\alpha$ such that
\[
  \mathrm{C}=\mathrm{C}_\alpha\otimes_{\mathrm{B}_\alpha}\mathrm{B},
\]
where $\mathrm{C}_\alpha$ is a finitely presented $\mathrm{B}_\alpha$-algebra. Moreover, by $(8.10.5)$ and \sref{IV.11.2.6}, we may suppose that $\mathrm{C}_\alpha$ is faithfully flat as a $\mathrm{B}_\alpha$-module. For $\lambda\geqslant\alpha$, put
\[
  \mathrm{C}_\lambda
  =\mathrm{C}_\alpha\otimes_{\mathrm{B}_\alpha}\mathrm{B}_\lambda.
\]
Then $\mathrm{C}=\mathrm{C}_\lambda\otimes_{\mathrm{B}_\lambda}\mathrm{B}$ and $\mathrm{C}_\lambda$ is faithfully flat as a $\mathrm{B}_\lambda$-module. Since $\mathrm{B}_\lambda\to\mathrm{B}$ is injective, so is $\mathrm{C}_\lambda\to\mathrm{C}$. Furthermore, $\mathrm{C}=\varinjlim\mathrm{C}_\lambda$, so every element of $\mathrm{C}$ belongs to some $\mathrm{C}_\lambda$; since $\mathrm{C}$ is of finite type over $\mathrm{A}$, there exists $\lambda$ for which $\mathrm{C}_\lambda\to\mathrm{C}$ \oldpage[IV-3]{143} is bijective. Faithful flatness then implies that $\mathrm{B}_\lambda\to\mathrm{B}$ is bijective. Thus $\mathrm{B}=\mathrm{B}_\lambda$ is of finite type over $\mathrm{A}$.

Suppose now that $\mathrm{C}$ is of finite presentation over $\mathrm{A}$. By what has just been proved, $\mathrm{B}$ is of finite type over $\mathrm{A}$; write
\[
  \mathrm{B}
  =\mathrm{A}[\mathrm{T}_1,\ldots,\mathrm{T}_n]/\widetilde{\mathfrak{J}}
\]
for an ideal $\widetilde{\mathfrak{J}}$. Let $(\widetilde{\mathfrak{J}}_\lambda)$ be the filtered family of finite type ideals of $\mathrm{A}[\mathrm{T}_1,\ldots,\mathrm{T}_n]$ contained in $\widetilde{\mathfrak{J}}$. Thus
\[
  \widetilde{\mathfrak{J}}
  =\varinjlim\widetilde{\mathfrak{J}}_\lambda,
  \qquad
  \mathrm{B}
  =\varinjlim\mathrm{B}_\lambda,
  \qquad
  \mathrm{B}_\lambda
  =\mathrm{A}[\mathrm{T}_1,\ldots,\mathrm{T}_n]/
   \widetilde{\mathfrak{J}}_\lambda.
\]
Applying $(8.8.2)$, $(8.10.5)$, and \sref{IV.11.2.6} as above, we have
\[
  \mathrm{C}
  =\mathrm{C}_\alpha\otimes_{\mathrm{B}_\alpha}\mathrm{B},
\]
where $\mathrm{C}_\alpha$ is a finitely presented $\mathrm{B}_\alpha$-algebra and a faithfully flat $\mathrm{B}_\alpha$-module. For $\lambda\geqslant\alpha$, put
\[
  \mathrm{C}_\lambda
  =\mathrm{C}_\alpha\otimes_{\mathrm{B}_\alpha}\mathrm{B}_\lambda.
\]
By flatness,
\[
  \mathrm{C}_\lambda
  =
  \mathrm{C}_\alpha\Big/
  \Bigl(
    \mathrm{C}_\alpha\otimes_{\mathrm{B}_\alpha}
    (\widetilde{\mathfrak{J}}_\lambda/
     \widetilde{\mathfrak{J}}_\alpha)
  \Bigr),
\]
and similarly
\[
  \mathrm{C}
  =
  \mathrm{C}_\alpha\Big/
  \Bigl(
    \mathrm{C}_\alpha\otimes_{\mathrm{B}_\alpha}
    (\widetilde{\mathfrak{J}}/
     \widetilde{\mathfrak{J}}_\alpha)
  \Bigr).
\]
Since both $\mathrm{C}$ and $\mathrm{C}_\alpha$ are finitely presented $\mathrm{A}$-algebras, the ideal
\[
  \mathrm{C}_\alpha\otimes_{\mathrm{B}_\alpha}
  (\widetilde{\mathfrak{J}}/
   \widetilde{\mathfrak{J}}_\alpha)
\]
is of finite type in $\mathrm{C}_\alpha$ $(1.4.4)$. It is the union of the finite type ideals
\[
  \mathrm{C}_\alpha\otimes_{\mathrm{B}_\alpha}
  (\widetilde{\mathfrak{J}}_\lambda/
   \widetilde{\mathfrak{J}}_\alpha).
\]
Consequently there exists $\lambda\geqslant\alpha$ such that
\[
  \mathrm{C}_\alpha\otimes_{\mathrm{B}_\alpha}
  (\widetilde{\mathfrak{J}}/
   \widetilde{\mathfrak{J}}_\alpha)
  =
  \mathrm{C}_\alpha\otimes_{\mathrm{B}_\alpha}
  (\widetilde{\mathfrak{J}}_\lambda/
   \widetilde{\mathfrak{J}}_\alpha).
\]
Faithful flatness then gives
$\widetilde{\mathfrak{J}}=\widetilde{\mathfrak{J}}_\lambda$, proving that $\mathrm{B}$ is of finite presentation over $\mathrm{A}$.
\end{proof}

\subsection{Descent of flatness by arbitrary morphisms: the case of an Artinian base}
\label{subsection:IV.11.4}

\begin{theorem}[11.4.1]
\label{IV.11.4.1}
Let $\mathrm{A}$ be a ring, $\mathfrak{J}$ a nilpotent ideal of $\mathrm{A}$, $u_\lambda:\mathrm{A}\to\mathrm{B}_\lambda$ $(\lambda\in\mathrm{L})$
a family of homomorphisms of rings such that the intersection of the kernels of the $u_\lambda$ is reduced to $0$.
Let $\mathrm{M}$ be an $\mathrm{A}$-module such that for all $\lambda\in\mathrm{L}$, $\mathrm{M}\otimes_\mathrm{A}\mathrm{B}_\lambda$ is a $\mathrm{B}_\lambda$-module free and that $\mathrm{M}\otimes_\mathrm{A}(\mathrm{A}/\mathfrak{J})$
is an $(\mathrm{A}/\mathfrak{J})$-module free. Then $\mathrm{M}$ is an $\mathrm{A}$-module free. If moreover the index set $\mathrm{L}$ is
finite, one can replace everywhere ``free'' by ``flat'' in the preceding statement.
\end{theorem}

\begin{proof}
In both cases it will suffice to prove that $\mathrm{M}$ is an $\mathrm{A}$-module \emph{flat}: indeed,
when $\mathrm{M}\otimes_\mathrm{A}(\mathrm{A}/\mathfrak{J})$ is an $(\mathrm{A}/\mathfrak{J})$-module free, it will follow that $\mathrm{M}$ is an $\mathrm{A}$-module
free by virtue of $(\mathbf{0}_{\mathrm{III}},~10.1.2)$.

We will use the following lemma, which generalises $(\mathbf{0}_{\mathrm{III}},~10.2.1)$:

\begin{lemma}[11.4.1.1]
\label{IV.11.4.1.1}
Let $\mathrm{A}$ be a ring equipped with a finite filtration $(\mathfrak{J}_n)_{0\leq n\leq\mathrm{N}+1}$ with
$\mathfrak{J}_0=\mathrm{A}$, $\mathfrak{J}_{\mathrm{N}+1}=0$. Let $\mathrm{M}$ be an $\mathrm{A}$-module equipped with the filtration $(\mathfrak{J}_n\mathrm{M})_{0\leq n\leq\mathrm{N}+1}$, and let us denote
by $\operatorname{gr}(\mathrm{A})$ and $\operatorname{gr}(\mathrm{M})$ the ring and the graded module respectively. Suppose that $\mathrm{M}\otimes_\mathrm{A}(\mathrm{A}/\mathfrak{J}_1)$
is an $(\mathrm{A}/\mathfrak{J}_1)$-module flat and that the canonical homomorphism
\[
\label{IV.11.4.1.2}
\operatorname{gr}_0(\mathrm{M})\otimes_{\operatorname{gr}_0(\mathrm{A})}\operatorname{gr}(\mathrm{A}) \to \operatorname{gr}(\mathrm{M})
\tag{11.4.1.2}
\]
is injective. Then $\mathrm{M}$ is an $\mathrm{A}$-module flat.
\end{lemma}

The canonical homomorphism {\hyperlink{ega4.c04.lem.11-4-1-2}{(11.4.1.2)}} is defined in the same way as that
of $(\mathbf{0}_{\mathrm{III}},~10.1.1)$ as being in degree $n$ the homomorphism
\[
(\mathrm{M}/\mathfrak{J}_1\mathrm{M})\otimes(\mathfrak{J}_n/\mathfrak{J}_{n+1})=(\mathrm{M}\otimes\mathfrak{J}_n)/(\operatorname{Im}(\mathrm{M}\otimes\mathfrak{J}_{n+1})+\operatorname{Im}(\mathfrak{J}_1\mathrm{M}\otimes\mathfrak{J}_n))\to\mathfrak{J}_n\mathrm{M}/\mathfrak{J}_{n+1}\mathrm{M}
\]
coming from the canonical homomorphism $\mathrm{M}\otimes\mathfrak{J}_n\to\mathfrak{J}_n\mathrm{M}$ by passage to quotients.
The lemma is proved by induction on $\mathrm{N}$, since there is nothing to prove for $\mathrm{N}=0$.
\oldpage[IV-3]{144}
The conditions on $\mathrm{M}$ imply, by the induction hypothesis, that $\mathrm{M}\otimes_\mathrm{A}(\mathrm{A}/\mathfrak{J}_\mathrm{N})$
is an $(\mathrm{A}/\mathfrak{J}_\mathrm{N})$-module flat. Let us now note that $\mathfrak{J}_\mathrm{N}^2\subset\mathfrak{J}_{\mathrm{N}+1}=0$ if $\mathrm{N}\geq 1$, and
$(\mathrm{M}/\mathfrak{J}_1\mathrm{M})\otimes(\mathfrak{J}_\mathrm{N}/\mathfrak{J}_{\mathrm{N}+1})=(\mathrm{M}/\mathfrak{J}_\mathrm{N}\mathrm{M})\otimes(\mathfrak{J}_\mathrm{N}/\mathfrak{J}_{\mathrm{N}+1})=(\mathrm{M}/\mathfrak{J}_\mathrm{N}\mathrm{M})\otimes\mathfrak{J}_\mathrm{N}$; hence the canonical homomorphism
\[
(\mathrm{M}/\mathfrak{J}_\mathrm{N}\mathrm{M})\otimes\mathfrak{J}_\mathrm{N} \to \mathfrak{J}_\mathrm{N}\mathrm{M}=\mathfrak{J}_\mathrm{N}\mathrm{M}/\mathfrak{J}_\mathrm{N}^2\mathrm{M}
\]
is injective. Applying $(\mathbf{0}_{\mathrm{III}},~10.2.1)$ to the $\mathfrak{J}_\mathrm{N}$-preadic filtration, we conclude that
$\mathrm{M}$ is an $\mathrm{A}$-module flat.

To apply this lemma to \sref{IV.11.4.1}, we will denote by $\mathfrak{J}_n$ the ideal of $\mathrm{A}$ which is the intersection
of the inverse images $u_\lambda^{-1}(\mathfrak{J}^n\mathrm{B}_\lambda)$: it is immediate that $\mathfrak{J}_0=\mathrm{A}$, $\mathfrak{J}_m\mathfrak{J}_n\subset\mathfrak{J}_{m+n}$
for $m\geq 0$, $n\geq 0$; moreover, if $\mathfrak{J}^{\mathrm{N}+1}=0$, we also have $\mathfrak{J}_{\mathrm{N}+1}=0$ since the intersection of the
kernels of the $u_\lambda$ is reduced to $0$.

Let us equip $\mathrm{A}$ with the filtration $(\mathfrak{J}_n)$, $\mathrm{M}$ with the filtration $(\mathfrak{J}_n\mathrm{M})$, and, for each $\lambda$,
let us equip $\mathrm{B}_\lambda$ and $\mathrm{N}_\lambda=\mathrm{M}\otimes_\mathrm{A}\mathrm{B}_\lambda$ with the $\mathfrak{J}$-preadic filtrations; let us consider for each $\lambda$ the
commutative diagram
\[
\xymatrix{
\operatorname{gr}^0(\mathrm{M})\otimes_{\operatorname{gr}^0(\mathrm{A})}\operatorname{gr}(\mathrm{A}) \ar[r]^-{\varphi} \ar[d]_{f_\lambda} & \operatorname{gr}(\mathrm{M}) \ar[d]^{g_\lambda} \\
\operatorname{gr}^0_\mathfrak{J}(\mathrm{N}_\lambda)\otimes_{\operatorname{gr}^0_\mathfrak{J}(\mathrm{B}_\lambda)}\operatorname{gr}_\mathfrak{J}(\mathrm{B}_\lambda) \ar[r]_-{\varphi_\lambda} & \operatorname{gr}_\mathfrak{J}(\mathrm{N}_\lambda)
}
\]
where the horizontal arrows are the canonical homomorphisms {\hyperlink{ega4.c04.lem.11-4-1-2}{(11.4.1.2)}} and the vertical
arrows are deduced from the canonical homomorphisms $\mathrm{A}\to\mathrm{B}_\lambda$ and $\mathrm{M}\to\mathrm{N}_\lambda$. The hypothesis
that $\mathrm{N}_\lambda$ is a $\mathrm{B}_\lambda$-module flat implies that $\varphi_\lambda$ is injective $(\mathbf{0}_{\mathrm{III}},~10.2.1)$, hence
$\operatorname{Ker}(\varphi)\subset\operatorname{Ker}(f_\lambda)$. Setting $\mathrm{A}_0=\mathrm{A}/\mathfrak{J}_1$, $\mathrm{M}_0=\mathrm{M}\otimes_\mathrm{A}\mathrm{A}_0$, let us note that
\[
\operatorname{gr}^0_\mathfrak{J}(\mathrm{N}_\lambda)=\mathrm{N}_\lambda/\mathfrak{J}\mathrm{N}_\lambda=(\mathrm{M}/\mathfrak{J}\mathrm{M})\otimes_\mathrm{A}\mathrm{B}_\lambda=(\mathrm{M}/\mathfrak{J}_1\mathrm{M})\otimes_{\mathrm{A}}(\mathrm{B}_\lambda/\mathfrak{J}\mathrm{B}_\lambda)
\]
and therefore
\[
\operatorname{gr}^0_\mathfrak{J}(\mathrm{N}_\lambda)\otimes_{\operatorname{gr}^0_\mathfrak{J}(\mathrm{B}_\lambda)}\operatorname{gr}_\mathfrak{J}(\mathrm{B}_\lambda)=\mathrm{M}_0\otimes_{\mathrm{A}_0}\operatorname{gr}_\mathfrak{J}(\mathrm{B}_\lambda).
\]
\oldpage[IV-3]{145}
The homomorphism $f_\lambda$ can therefore be written
\[
1\otimes\operatorname{gr}(u_\lambda) : \mathrm{M}_0\otimes_{\mathrm{A}_0}\operatorname{gr}(\mathrm{A}) \to \mathrm{M}_0\otimes_{\mathrm{A}_0}\operatorname{gr}_\mathfrak{J}(\mathrm{B}_\lambda),
\]
and since by hypothesis $\mathrm{M}_0$ is an $\mathrm{A}_0$-module \emph{flat}, the kernel of $f_\lambda$ is equal to $\mathrm{M}_0\otimes_{\mathrm{A}_0}\mathrm{R}_\lambda$,
where $\mathrm{R}_\lambda$ is the kernel of $\operatorname{gr}(u_\lambda):\operatorname{gr}(\mathrm{A})\to\operatorname{gr}_\mathfrak{J}(\mathrm{B}_\lambda)$. It therefore suffices to prove that
$\bigcap_{\lambda\in\mathrm{L}}(\mathrm{M}_0\otimes_{\mathrm{A}_0}\mathrm{R}_\lambda)=0$. Now, by definition of the $\mathfrak{J}_n$, the intersection of the kernels of the homomorphisms
$\mathfrak{J}_n\mathrm{A}/\mathfrak{J}_{n+1}\mathrm{A}\to\mathfrak{J}^n\mathrm{B}_\lambda/\mathfrak{J}^{n+1}\mathrm{B}_\lambda$, as $\lambda$ ranges over $\mathrm{L}$, is reduced to $0$; in other
words, $\bigcap_{\lambda\in\mathrm{L}}\mathrm{R}_\lambda=0$. This being the case, suppose first that $\mathrm{M}_0$ is an $\mathrm{A}_0$-module free; taking
a basis of $\mathrm{M}_0$, we see immediately that $\mathrm{M}_0\otimes_{\mathrm{A}_0}(\bigcap_{\lambda\in\mathrm{L}}\mathrm{R}_\lambda)=\bigcap_{\lambda\in\mathrm{L}}(\mathrm{M}_0\otimes_{\mathrm{A}_0}\mathrm{R}_\lambda)$, whence the
proposition in this case. When $\mathrm{L}$ is finite, the preceding formula is still true under
the sole hypothesis that $\mathrm{M}_0$ is an $\mathrm{A}_0$-module flat $(\mathbf{0}_\mathrm{I},~6.1.3)$, which finishes the
proof.
\end{proof}

\begin{remark}[11.4.2]
\label{IV.11.4.2}
The conclusion of \sref{IV.11.4.1} may be inexact if $\mathrm{L}$ is infinite
and one only supposes that $\mathrm{M}\otimes_\mathrm{A}(\mathrm{A}/\mathfrak{J})$ is an $(\mathrm{A}/\mathfrak{J})$-module flat. For example,
let $\mathrm{V}$ be a discrete valuation ring, $\mathrm{K}$ its field of fractions, and let $\mathrm{A}=\mathrm{V}[\mathrm{T}]/(\mathrm{T}^2)$
($\mathrm{T}$ an indeterminate); let us take for $\mathfrak{J}$ the image of $(\mathrm{T})$ in $\mathrm{A}$, so that $\mathrm{A}/\mathfrak{J}=\mathrm{V}$, and
let us take $\mathrm{M}=\mathrm{K}$, which is an $(\mathrm{A}/\mathfrak{J})$-module, hence equal to $\mathrm{M}\otimes_\mathrm{A}(\mathrm{A}/\mathfrak{J})$; moreover $\mathrm{M}$ is
an $(\mathrm{A}/\mathfrak{J})$-module flat, but not an $\mathrm{A}$-module flat: tensoring the inclusion $\mathfrak{J}\hookrightarrow\mathrm{A}$ with $\mathrm{K}$
gives the zero map $\mathfrak{J}\otimes_{\mathrm{A}}\mathrm{K}\simeq\mathrm{K}\longrightarrow\mathrm{K}$,
which is not injective since $\mathrm{K}\neq0$. On the other hand, consider the maximal ideal $\mathfrak{m}$ of the local Noetherian
ring $\mathrm{A}$; we have $\mathfrak{m}\mathrm{K}=\mathrm{K}$, hence $\mathrm{M}\otimes_\mathrm{A}(\mathrm{A}/\mathfrak{m}^n)=0$ for every integer $n$; the $(\mathrm{A}/\mathfrak{m}^n)$-modules
$\mathrm{M}\otimes_\mathrm{A}(\mathrm{A}/\mathfrak{m}^n)$ are therefore flat for all $n$, and the intersection of $\mathfrak{m}^n$ is reduced to $0$.
\end{remark}

\begin{corollary}[11.4.3]
\label{IV.11.4.3}
Let $\mathrm{A}$ be a semi-local ring whose radical $\mathfrak{J}$ is nilpotent (for
example an Artinian ring), $u_\lambda:\mathrm{A}\to\mathrm{B}_\lambda$ $(\lambda\in\mathrm{L})$ a family of homomorphisms of rings such
that the intersection of the kernels of the $u_\lambda$ is reduced to $0$. For an $\mathrm{A}$-module $\mathrm{M}$ to be flat, it is necessary and
sufficient that for all $\lambda\in\mathrm{L}$, $\mathrm{M}\otimes_\mathrm{A}\mathrm{B}_\lambda$ be a $\mathrm{B}_\lambda$-module flat.
\end{corollary}

\begin{proof}
Since $\mathrm{A}/\mathfrak{J}$ is a direct product of a finite number of fields (Bourbaki, \emph{Alg. comm.},
chap.~II, \S~3, n\textsuperscript{o}~5, prop.~16) and that $\mathfrak{J}$ is nilpotent, $\mathrm{A}$ is a direct product of a
finite number of \emph{local} rings $\mathrm{A}_i$ whose radical is nilpotent (\emph{loc.\ cit.}, \S~4, n\textsuperscript{o}~3, cor.~1
of prop.~15) and $\mathrm{M}$ is therefore a direct sum of $\mathrm{A}_i$-modules $\mathrm{M}_i$, each $\mathrm{M}_i$ being
annihilated by the $\mathrm{A}_j$ of index $j\neq i$; for $\mathrm{M}$ to be a flat $\mathrm{A}$-module, it is necessary and it
suffices that each of the $\mathrm{M}_i$ be a flat $\mathrm{A}_i$-module; moreover the intersection of the kernels
of the homomorphisms $\mathrm{A}_i\longrightarrow\mathrm{A}\xrightarrow{u_\lambda}\mathrm{B}_\lambda$ is reduced to $0$, and $\mathrm{M}_i\otimes_{\mathrm{A}_i}\mathrm{B}_\lambda$ is a direct factor
of $\mathrm{M}\otimes_\mathrm{A}\mathrm{B}_\lambda$. One can therefore restrict to the case where $\mathrm{A}$ is also \emph{local}. Then $\mathrm{A}/\mathfrak{J}$ is a
field, hence $\mathrm{M}\otimes_\mathrm{A}(\mathrm{A}/\mathfrak{J})$ is an $(\mathrm{A}/\mathfrak{J})$-module free, and it suffices to see that for
all $\lambda$, $\mathrm{M}\otimes_\mathrm{A}\mathrm{B}_\lambda$ is a $\mathrm{B}_\lambda$-module \emph{free}, by virtue of \sref{IV.11.4.1}. But if one sets $\mathfrak{J}_\lambda=\mathfrak{J}\mathrm{B}_\lambda$,
$(\mathrm{M}\otimes_\mathrm{A}\mathrm{B}_\lambda)\otimes_{\mathrm{B}_\lambda}(\mathrm{B}_\lambda/\mathfrak{J}_\lambda)=(\mathrm{M}\otimes_\mathrm{A}(\mathrm{A}/\mathfrak{J}))\otimes_{\mathrm{A}/\mathfrak{J}}(\mathrm{B}_\lambda/\mathfrak{J}_\lambda)$ is a $(\mathrm{B}_\lambda/\mathfrak{J}_\lambda)$-module free, and $\mathfrak{J}_\lambda$ is
nilpotent. The conclusion therefore follows from the hypothesis that $\mathrm{M}\otimes_\mathrm{A}\mathrm{B}_\lambda$ is a $\mathrm{B}_\lambda$-module
flat and from $(\mathbf{0}_{\mathrm{III}},~10.1.2)$.
\end{proof}

\begin{corollary}[11.4.4]
\label{IV.11.4.4}
Let $\mathrm{A}$ be a ring, $\mathrm{M}$ an $\mathrm{A}$-module; suppose that there exists
a nilpotent ideal $\mathfrak{J}$ of $\mathrm{A}$ such that $\mathrm{M}\otimes_\mathrm{A}(\mathrm{A}/\mathfrak{J})$ is an $(\mathrm{A}/\mathfrak{J})$-module free. Then the set of
ideals $\mathfrak{K}$ of $\mathrm{A}$ such that $\mathrm{M}\otimes_\mathrm{A}(\mathrm{A}/\mathfrak{K})$ is an $(\mathrm{A}/\mathfrak{K})$-module free admits a smallest element $\mathfrak{K}_0$
\oldpage[IV-3]{146}
(which is also the smallest of the ideals $\mathfrak{K}$ such that $\mathrm{M}\otimes_\mathrm{A}(\mathrm{A}/\mathfrak{K})$ is an $(\mathrm{A}/\mathfrak{K})$-module \emph{flat}).
\end{corollary}

\begin{proof}
For any homomorphism $u:\mathrm{A}\to\mathrm{A}'$ such that $\mathrm{M}\otimes_\mathrm{A}\mathrm{A}'$ is an $\mathrm{A}'$-module free (resp.\ an
$\mathrm{A}'$-module flat), it is necessary and it suffices that $u$ factorise through $\mathrm{A}\to\mathrm{A}/\mathfrak{K}_0\to\mathrm{A}'$ (or equivalently that $\mathfrak{K}_0\mathrm{A}'=0$).

The fact that the intersection $\mathfrak{K}_0$ of the ideals $\mathfrak{K}$ for which $\mathrm{M}\otimes_\mathrm{A}(\mathrm{A}/\mathfrak{K})$ is an
$(\mathrm{A}/\mathfrak{K})$-module free is itself the smallest of these ideals follows from \sref{IV.11.4.1} applied to the
ring $\mathrm{A}/\mathfrak{K}_0$, to a nilpotent ideal $\mathfrak{J}/\mathfrak{K}_0$ and to the homomorphisms $\mathrm{A}/\mathfrak{K}_0\to\mathrm{A}/\mathfrak{K}$,
whose kernels have $0$ for intersection. If $\mathrm{A}'$ is an $(\mathrm{A}/\mathfrak{K}_0)$-algebra, then
$\mathrm{M}\otimes_\mathrm{A}\mathrm{A}'=(\mathrm{M}\otimes_\mathrm{A}(\mathrm{A}/\mathfrak{K}_0))\otimes_{\mathrm{A}/\mathfrak{K}_0}\mathrm{A}'$, and hence $\mathrm{M}\otimes_\mathrm{A}\mathrm{A}'$ is an $\mathrm{A}'$-module free. Conversely, if $\mathrm{A}'$ is an $\mathrm{A}$-algebra such that $\mathrm{M}\otimes_\mathrm{A}\mathrm{A}'$ is an $\mathrm{A}'$-module free and if $\mathfrak{K}$ is the kernel
of the homomorphism $\mathrm{A}\to\mathrm{A}'$, it follows from \sref{IV.11.4.1} applied to the ring $\mathrm{A}/\mathfrak{K}$, to the
$(\mathrm{A}/\mathfrak{K})$-module $\mathrm{M}\otimes_\mathrm{A}(\mathrm{A}/\mathfrak{K})$, to the nilpotent ideal $(\mathfrak{J}+\mathfrak{K})/\mathfrak{K}$ and to the \emph{injective} homomorphism
$\mathrm{A}/\mathfrak{K}\to\mathrm{A}'$, that $\mathrm{M}\otimes_\mathrm{A}(\mathrm{A}/\mathfrak{K})$ is an $(\mathrm{A}/\mathfrak{K})$-module free, hence $\mathfrak{K}\supset\mathfrak{K}_0$. The fact that
one can replace ``free'' by ``flat'' in what precedes (while naturally preserving
the hypothesis that $\mathrm{M}\otimes_\mathrm{A}(\mathrm{A}/\mathfrak{J})$ is an $(\mathrm{A}/\mathfrak{J})$-module \emph{free}) follows from $(\mathbf{0}_{\mathrm{III}},~10.1.2)$,
as we saw at the beginning of the proof of \sref{IV.11.4.1}.
\end{proof}

\begin{proposition}[11.4.5]
\label{IV.11.4.5}
Let $\mathrm{Y}$ be an irreducible prescheme, $f:\mathrm{X}\to\mathrm{Y}$ a morphism of
finite presentation, $\sh{F}$ an $\sh{O}_\mathrm{X}$-module of finite presentation. There exists an open non-empty $\mathrm{U}$
in $\mathrm{Y}$ and a closed sub-prescheme $\mathrm{Z}$ of $\mathrm{U}$, of finite presentation over $\mathrm{U}$, having the following property:
for every base change $\mathrm{U}'\to\mathrm{U}$, if we set $\mathrm{X}'=\mathrm{X}\times_\mathrm{Y}\mathrm{U}'$, $f'=f_{(\mathrm{U}')}$ and $\sh{F}'=\sh{F}\otimes_{\sh{O}_\mathrm{Y}}\sh{O}_{\mathrm{U}'}$,
then, for $\sh{F}'$ to be $f'$-flat, it is necessary and sufficient that the morphism $\mathrm{U}'\to\mathrm{U}$ factorise
through $\mathrm{U}'\to\mathrm{Z}\to\mathrm{U}$. Such a scheme $\mathrm{Z}$ has the same underlying space as $\mathrm{U}$. Suppose moreover that $\mathrm{Y}$ is affine,
and let $(\mathrm{W}_i)$ be a finite affine open covering of $\mathrm{X}$; then, one can suppose $\mathrm{U}$ chosen
such that, if $\mathrm{U}'=\operatorname{Spec}(\mathrm{A}')$ is an affine scheme and $\mathrm{U}'\to\mathrm{U}$ a morphism that factorises
through $\mathrm{U}'\to\mathrm{Z}\to\mathrm{U}$, each of the $\Gamma(\mathrm{W}_i\times_\mathrm{Y}\mathrm{U}',\sh{F}')$ is a free $\mathrm{A}'$-module.
\end{proposition}

\begin{proof}
One can evidently restrict to the case where $\mathrm{Y}=\operatorname{Spec}(\mathrm{A})$ is affine. Using
\sref{IV.8.9.1}, there is a Noetherian sub-ring $\mathrm{A}_1$ of $\mathrm{A}$, a morphism of finite type
$f_1:\mathrm{X}_1\to\mathrm{Y}_1=\operatorname{Spec}(\mathrm{A}_1)$ and a coherent $\sh{O}_{\mathrm{X}_1}$-module $\sh{F}_1$ such that $f=f_{(\mathrm{Y})}$ and $\sh{F}=\sh{F}_1\otimes_{\sh{O}_{\mathrm{Y}_1}}\sh{O}_\mathrm{Y}$;
one can moreover suppose that the $\mathrm{W}_i$ are inverse images of affine opens of $\mathrm{X}_1$.
Let us note moreover that $\mathrm{Y}_1$ is irreducible, the morphism $\mathrm{Y}\to\mathrm{Y}_1$ being dominant (\textbf{I},~$1.2.7$).
Suppose the proposition proved for $\mathrm{Y}_1$, $\mathrm{X}_1$ and $\sh{F}_1$, and let $\mathrm{U}_1$ and $\mathrm{Z}_1$ be the open set of $\mathrm{Y}_1$
and the closed sub-prescheme of $\mathrm{U}_1$ having the required properties. Set $\mathrm{U}=\mathrm{U}_1\times_{\mathrm{Y}_1}\mathrm{Y}$ and $\mathrm{Z}=\mathrm{Z}_1\times_{\mathrm{Y}_1}\mathrm{Y}$, their inverse images. Then, if $\mathrm{U}'\to\mathrm{U}$ is a base change
such that $\sh{F}'=\sh{F}\otimes_{\sh{O}_\mathrm{Y}}\sh{O}_{\mathrm{U}'}$ is $f'$-flat, the morphism $\mathrm{U}'\to\mathrm{U}_1$ factorises
through $\mathrm{U}'\to\mathrm{Z}_1\to\mathrm{U}_1$; but since $\mathrm{U}'\to\mathrm{U}\to\mathrm{U}_1$ also factorises through $\mathrm{U}'\to\mathrm{U}\to\mathrm{U}_1$, the definition
of the product of preschemes shows that $\mathrm{U}'\to\mathrm{U}$ factorises through $\mathrm{U}'\to\mathrm{Z}\to\mathrm{U}$.

One can therefore restrict to the case where $\mathrm{A}$ is \emph{Noetherian}; let $\mathfrak{N}$ be its nilradical, which is
nilpotent, and set $\mathrm{A}_0=\mathrm{A}_{\mathrm{red}}=\mathrm{A}/\mathfrak{N}$, $\mathrm{Y}_0=\operatorname{Spec}(\mathrm{A}_0)=\mathrm{Y}_{\mathrm{red}}$, $\mathrm{X}_0=\mathrm{X}\times_\mathrm{Y}\mathrm{Y}_0$, $\sh{F}_0=\sh{F}\otimes_{\sh{O}_\mathrm{Y}}\sh{O}_{\mathrm{Y}_0}$,
$\mathrm{W}_{i0}=\mathrm{W}_i\times_\mathrm{Y}\mathrm{Y}_0$; if $\mathrm{M}_i=\Gamma(\mathrm{W}_i,\sh{F})$, $\mathrm{M}_{i0}=\Gamma(\mathrm{W}_{i0},\sh{F}_0)$ is therefore equal to $\mathrm{M}_i\otimes_\mathrm{A}\mathrm{A}_0$.
Since $\mathrm{A}_0$ is \emph{integral}, one can, by virtue of the generic flatness theorem $(6.9.2)$,
replacing $\mathrm{Y}$ by an open non-empty set of $\mathrm{Y}$ as needed, suppose that the $\mathrm{M}_{i0}$ are
free $\mathrm{A}_0$-modules. By virtue of \sref{IV.11.4.4}, there is therefore for each $i$ an ideal $\mathfrak{J}_i\subset\mathfrak{N}$ such
\oldpage[IV-3]{147}
that the $\mathrm{A}$-algebras $\mathrm{A}'$ for which the $\mathrm{M}_i\otimes_\mathrm{A}\mathrm{A}'$ are $\mathrm{A}'$-modules free (or flat)
are exactly those for which $\mathfrak{J}_i\mathrm{A}'=0$. Let $\mathfrak{J}=\sum_i\mathfrak{J}_i$, which is an ideal contained
in $\mathfrak{N}$; to say that $\sh{F}\otimes_{\mathrm{A}}\mathrm{A}'$ is $\mathrm{A}'$-flat amounts to saying that the $\mathrm{M}_i\otimes_\mathrm{A}\mathrm{A}'$ are all
$\mathrm{A}'$-modules flat, hence that $\mathfrak{J}\mathrm{A}'=0$. It follows that if one takes $\mathrm{Z}=\mathrm{V}(\mathfrak{J})$, one
answers the question, since for any morphism $\mathrm{U}'\to\mathrm{U}$ having the property of the statement,
it is necessary and it suffices that for every affine open $\mathrm{V}'\subset\mathrm{U}'$, the morphism $\mathrm{V}'\to\mathrm{U}$
have the same property.
\end{proof}

\begin{corollary}[11.4.6]
\label{IV.11.4.6}
Let $\mathrm{A}$ be a ring such that $\operatorname{Spec}(\mathrm{A})$ is irreducible, $\mathfrak{p}$ its unique
minimal prime ideal, $\mathrm{B}$ an $\mathrm{A}$-algebra of finite presentation, $\mathrm{M}$ a $\mathrm{B}$-module of finite presentation.
Suppose that $\mathrm{M}_\mathfrak{p}$ is a flat $\mathrm{A}_\mathfrak{p}$-module; then there exists $t\in\mathrm{A}-\mathfrak{p}$ such that $\mathrm{M}_t$ is a free $\mathrm{A}_t$-module.
\end{corollary}

\begin{proof}
Applying \sref{IV.11.4.5} with $\mathrm{Y}=\operatorname{Spec}(\mathrm{A})$, $\mathrm{X}=\operatorname{Spec}(\mathrm{B})$, $\sh{F}=\widetilde{\mathrm{M}}$, one can (replacing
$\mathrm{A}$ by $\mathrm{A}_h$ as needed, where $h\in\mathrm{A}-\mathfrak{p}$) suppose that there exists an ideal of finite type $\mathfrak{J}$ in $\mathrm{A}$
such that the $\mathrm{A}$-algebras $\mathrm{A}'$ for which $\mathrm{M}\otimes_\mathrm{A}\mathrm{A}'$ is an $\mathrm{A}'$-module free (or flat) are
exactly those such that $\mathfrak{J}\mathrm{A}'=0$. In particular, the hypothesis implies that $\mathrm{M}_\mathfrak{p}$ is a flat
$\mathrm{A}_\mathfrak{p}$-module, hence $\mathfrak{J}\mathrm{A}_\mathfrak{p}=0$, or $\mathfrak{J}_\mathfrak{p}=0$. But since $\mathfrak{J}$ is of finite type, there
exists $t\in\mathrm{A}-\mathfrak{p}$ such that $\mathfrak{J}_t=0$, or equivalently $\mathfrak{J}\mathrm{A}_t=0$, and therefore $\mathrm{M}_t$ is a free $\mathrm{A}_t$-module.
\end{proof}

\begin{proposition}[11.4.7]
\label{IV.11.4.7}
Let $\mathrm{A}$ be a Noetherian ring, $\mathfrak{J}$ an ideal of $\mathrm{A}$, $\mathrm{M}$ an $\mathrm{A}$-module
ideally separated $(\mathbf{0}_{\mathrm{III}},~10.2.1)$ for the $\mathfrak{J}$-preadic topology. Let $(\mathfrak{p}_i)_{1\leq i\leq r}$ be a finite family
of prime ideals of $\mathrm{A}$ containing $\mathfrak{J}$; for every integer $n\geq 0$, let $\mathfrak{p}_i^{(n)}$ be the $n$-th symbolic
power (i.e.\ the kernel of the homomorphism $\mathrm{A}\to\mathrm{A}_{\mathfrak{p}_i}/\mathfrak{p}_i^n\mathrm{A}_{\mathfrak{p}_i}$); set $\mathfrak{J}_n=\bigcap_{i=1}^{r}\mathfrak{p}_i^{(n)}$, so that $\mathfrak{J}_n\supset\mathfrak{J}^n$,
and suppose that the topology defined by the filtration $(\mathfrak{J}_n)$ is identical to the $\mathfrak{J}$-preadic topology
(in other words, for every $n$, there exists $m$ such that $\mathfrak{J}^n\supset\mathfrak{J}_m$). For $\mathrm{M}$ to be an $\mathrm{A}$-module flat,
it is necessary and sufficient that $\mathrm{M}\otimes_\mathrm{A}(\mathrm{A}/\mathfrak{J})$ be an $(\mathrm{A}/\mathfrak{J})$-module flat and that, for all $i$, $\mathrm{M}\otimes_\mathrm{A}\mathrm{A}_{\mathfrak{p}_i}=\mathrm{M}_{\mathfrak{p}_i}$
be a flat $\mathrm{A}_{\mathfrak{p}_i}$-module.
\end{proposition}

\begin{proof}
Since $\mathrm{M}$ is ideally separated, it suffices, by virtue of $(\mathbf{0}_{\mathrm{III}},~10.2.1)$, to show
that, for every $n\geq 0$, $\mathrm{M}\otimes_\mathrm{A}(\mathrm{A}/\mathfrak{J}^n)$ is an $(\mathrm{A}/\mathfrak{J}^n)$-module flat; since every $\mathfrak{J}^n$ contains
a $\mathfrak{J}_m$, it amounts to the same to prove that for every $n$, $\mathrm{M}\otimes_\mathrm{A}(\mathrm{A}/\mathfrak{J}_n)$ is an $(\mathrm{A}/\mathfrak{J}_n)$-module
flat. Now, since $\mathfrak{J}_1\supset\mathfrak{J}$, $\mathrm{M}\otimes_\mathrm{A}(\mathrm{A}/\mathfrak{J}_1)$ is an $(\mathrm{A}/\mathfrak{J}_1)$-module flat; in the ring $\mathrm{A}/\mathfrak{J}_n$,
the ideal $\mathfrak{J}_1/\mathfrak{J}_n$ is nilpotent and finally the intersection of the kernels of the homomorphisms
$\mathrm{A}/\mathfrak{J}_n\to\mathrm{A}_{\mathfrak{p}_i}/\mathfrak{p}_i^n\mathrm{A}_{\mathfrak{p}_i}$ is zero in $\mathrm{A}/\mathfrak{J}_n$, by definition of $\mathfrak{J}_n$. It therefore suffices, by \sref{IV.11.4.1},
to verify that $\mathrm{M}\otimes_\mathrm{A}(\mathrm{A}_{\mathfrak{p}_i}/\mathfrak{p}_i^n\mathrm{A}_{\mathfrak{p}_i})$ is an $(\mathrm{A}_{\mathfrak{p}_i}/\mathfrak{p}_i^n\mathrm{A}_{\mathfrak{p}_i})$-module flat, which follows from
the hypothesis that $\mathrm{M}\otimes_\mathrm{A}\mathrm{A}_{\mathfrak{p}_i}$ is a flat $\mathrm{A}_{\mathfrak{p}_i}$-module.
\end{proof}

\begin{remark}[11.4.8]
\label{IV.11.4.8}
The hypothesis made in \sref{IV.11.4.7} on the topology defined
by the $\mathfrak{J}_n$ is satisfied if, for every $n$ sufficiently large, $\operatorname{Ass}(\mathrm{A}/\mathfrak{J}^n)$ \emph{is contained in the set
of the $\mathfrak{p}_i$}. Indeed, $\mathfrak{J}^n$ is then an intersection of primary ideals for the $\mathfrak{p}_i$, each of which
contains a symbolic power of $\mathfrak{p}_i$, whence the conclusion. In particular:
\end{remark}

\begin{corollary}[11.4.9]
\label{IV.11.4.9}
Let $\mathrm{A}$ be a Noetherian ring, $\mathfrak{J}$ a nilpotent ideal of $\mathrm{A}$, $\mathrm{M}$ an
$\mathrm{A}$-module. For $\mathrm{M}$ to be an $\mathrm{A}$-module flat (resp.\ free), it is necessary and sufficient that $\mathrm{M}\otimes_\mathrm{A}(\mathrm{A}/\mathfrak{J})$
be an $(\mathrm{A}/\mathfrak{J})$-module flat (resp.\ free) and that for every prime ideal $\mathfrak{p}\in\operatorname{Ass}(\mathrm{A})$, $\mathrm{M}_\mathfrak{p}$ be a flat
$\mathrm{A}_\mathfrak{p}$-module.
\end{corollary}

The assertion concerning the case where $\mathrm{M}\otimes_\mathrm{A}(\mathrm{A}/\mathfrak{J})$ is free follows again from the assertion concerning the case where $\mathrm{M}\otimes_\mathrm{A}(\mathrm{A}/\mathfrak{J})$ is flat by $(\mathbf{0}_{\mathrm{III}},10.1.2)$.

\begin{corollary}[11.4.10]
\label{IV.11.4.10}
Let $\mathrm{A}$ be a Noetherian ring, $\mathfrak{J}$ an ideal of $\mathrm{A}$, $\mathrm{M}$ an $\mathrm{A}$-module.
Suppose that $\mathrm{M}$ is ideally separated for the $\mathfrak{J}$-preadic topology and that $\operatorname{gr}_\mathfrak{J}(\mathrm{A})$ is an
$(\mathrm{A}/\mathfrak{J})$-module flat. For $\mathrm{M}$ to be an $\mathrm{A}$-module flat, it is necessary and sufficient that $\mathrm{M}\otimes_\mathrm{A}(\mathrm{A}/\mathfrak{J})$ be
an $(\mathrm{A}/\mathfrak{J})$-module flat and that for every $\mathfrak{p}\in\operatorname{Ass}(\mathrm{A}/\mathfrak{J})$, $\mathrm{M}_\mathfrak{p}$ be a flat $\mathrm{A}_\mathfrak{p}$-module.
\end{corollary}

\begin{proof}
\oldpage[IV-3]{148}
Taking into account \sref{IV.11.4.8}, it suffices to show that $\operatorname{Ass}(\mathrm{A}/\mathfrak{J}^n)$ is contained in
$\operatorname{Ass}(\mathrm{A}/\mathfrak{J})$ for every $n$. Now, if $a\in\mathrm{A}$ does not belong to any $\mathfrak{p}\in\operatorname{Ass}(\mathrm{A}/\mathfrak{J})$, the homothety
of ratio $a$ is injective in $\mathrm{A}/\mathfrak{J}$; since each of the $\mathfrak{J}^k/\mathfrak{J}^{k+1}$ is an $(\mathrm{A}/\mathfrak{J})$-module
flat, $a$ is also a $(\mathfrak{J}^k/\mathfrak{J}^{k+1})$-regular element, hence $a$ is a $(\mathrm{A}/\mathfrak{J}^n)$-regular element for every $n$
(Bourbaki, \emph{Alg.\ comm.}, chap.~III, \S~2, n\textsuperscript{o}~8, cor.~1 of th.~1), and therefore it does not belong
to any prime ideal associated to $\mathrm{A}/\mathfrak{J}^n$, whence the corollary (Bourbaki, \emph{Alg.\ comm.}, chap.~II,
\S~1, n\textsuperscript{o}~1, prop.~2).
\end{proof}

\begin{proposition}[11.4.11]
\label{IV.11.4.11}
Let $\mathrm{A}$ be an Artinian local ring, $\mathfrak{m}$ its maximal ideal,
$\mathrm{Y}=\operatorname{Spec}(\mathrm{A})$, $y$ the unique point of $\mathrm{Y}$, $f:\mathrm{X}\to\mathrm{Y}$ a morphism locally of finite type, $\sh{F}$ a
coherent $\sh{O}_\mathrm{X}$-module. Let $(\mathrm{A}'_\alpha)$ be a family of local rings, and for each $\alpha$, $u_\alpha:\mathrm{A}\to\mathrm{A}'_\alpha$
a homomorphism of rings (necessarily local). Set $\mathrm{Y}'_\alpha=\operatorname{Spec}(\mathrm{A}'_\alpha)$, $\mathrm{X}'_\alpha=\mathrm{X}\times_\mathrm{Y}\mathrm{Y}'_\alpha$,
$\sh{F}'_\alpha=\sh{F}\otimes_\mathrm{A}\mathrm{A}'_\alpha$. Let $x$ be a point
of $\mathrm{X}$, and suppose the following conditions are satisfied:
\begin{enumerate}[label=(\roman*)]
  \item The intersection of the kernels of the $u_\alpha$ is reduced to $0$.
  \item The extension $k(x)$ of the residue field $k$ of $\mathrm{A}$ is primary $(4.3.1)$ (a condition automatically
  satisfied if $k$ is separably closed).
\end{enumerate}
\begin{enumerate}[label=(\roman*)]
  \setcounter{enumi}{2}
  \item For all $\alpha$, there exists a point $x'_\alpha\in\mathrm{X}'_\alpha$ whose projections in $\mathrm{X}$ and $\mathrm{Y}'_\alpha$
  are $x$ and the closed point $y'_\alpha$ of $\mathrm{Y}'_\alpha$, and such that $\sh{F}'_\alpha$ is $\mathrm{Y}'_\alpha$-flat at the point $x'_\alpha$.
\end{enumerate}

Under these conditions, $\sh{F}$ is $f$-flat at the point $x$.
\end{proposition}

\begin{proof}
The question being local on $\mathrm{X}$, one can evidently restrict to the case where $f$ is
a morphism of finite type and suppose $\sh{F}_x\neq 0$. We will proceed in several steps.

\medskip
\noindent\textbf{I)} \emph{Reduction to the case where $\mathrm{A}'_\alpha$ is a local ring of a $\mathrm{Y}$-prescheme of finite type and where the
residue field $k'_\alpha$ of $\mathrm{A}'_\alpha$ is a finite extension of $k$.}

\medskip
Since the reduction is done separately on each $\mathrm{A}'_\alpha$, one can suppress the index $\alpha$ in
this part. Let $\mathfrak{m}'$ be the maximal ideal of $\mathrm{A}'$. Let us consider $\mathrm{A}'$ as an
inductive limit of its sub-$\mathrm{A}$-algebras of finite type $\mathrm{B}_\lambda$, and set $\mathfrak{n}_\lambda=\mathfrak{m}'\cap\mathrm{B}_\lambda$; $\mathrm{A}'$ is also an
inductive limit of its local sub-rings $\mathrm{B}'_\lambda=(\mathrm{B}_\lambda)_{\mathfrak{n}_\lambda}$ $(5.13.3)$, and one has evidently
$\mathfrak{m}'\cap\mathrm{B}'_\lambda=\mathfrak{n}'_\lambda$, maximal ideal of $\mathrm{B}'_\lambda$. It follows therefore \sref{IV.11.2.6} that in
setting $\mathrm{Z}'_\lambda=\operatorname{Spec}(\mathrm{B}'_\lambda)$, $\sh{F}\otimes_\mathrm{A}\mathrm{B}'_\lambda$ is $\mathrm{Z}'_\lambda$-flat at the point $x'_\lambda$, projection of $x'$, the projection
of $x'_\lambda$ on $\mathrm{Z}'_\lambda$ being the closed point $z'_\lambda$ of $\mathrm{Z}'_\lambda$.

One can therefore suppose that there exists an affine scheme $\mathrm{Z}'$ of finite type over $\mathrm{Y}$ and a
point $z'$ of $\mathrm{Z}'$ such that $\mathrm{A}'=\sh{O}_{\mathrm{Z}',z'}$, and that, if one sets $\mathrm{T}'=\mathrm{X}\times_\mathrm{Y}\mathrm{Z}'$, there exists a point
$t'\in\mathrm{T}'$ whose projections in $\mathrm{Z}'$ and $\mathrm{X}$ are $z'$ and $x$, and such that $\sh{F}\otimes_\mathrm{Y}\mathrm{Z}'$ is $\mathrm{Z}'$-flat at the
point $t'$. Let $\mathrm{Z}'_1$ (resp.\ $\mathrm{X}_1$) be a closed sub-prescheme of $\mathrm{Z}'$ (resp.\ $\mathrm{X}$) having for underlying space
the closure of $\{z'\}$ (resp.\ $\{x\}$), and set $\mathrm{T}'_1=\mathrm{X}_1\times_\mathrm{Y}\mathrm{Z}'_1$. The set $\mathrm{U}$ of the
points of $\mathrm{T}'$ where $\sh{F}\otimes_\mathrm{Y}\mathrm{Z}'$ is $\mathrm{Z}'$-flat is open in $\mathrm{T}'$ \sref{IV.11.1.1} and contains $t'$, hence
\oldpage[IV-3]{149}
$\mathrm{V}=\mathrm{U}\cap\mathrm{T}'_1$ is a nonempty open subset of $\mathrm{T}'_1$. The ring $\mathrm{A}$, being Artinian, is a Jacobson
ring; hence, by $(10.4.6)$, $\mathrm{Z}'_1$ and $\mathrm{T}'_1$ are Jacobson preschemes. There is therefore
a point $t'_0$ in $\mathrm{V}$ that is closed in $\mathrm{V}$, and its image $z'_0$ in $\mathrm{Z}'_1$ is a \emph{closed} point of $\mathrm{Z}'_1$ $(10.4.7)$.
Let $f_1$ be the restriction of $f$ to $\mathrm{X}_1$, and let $f'_1:\mathrm{T}'_1\to\mathrm{Z}'_1$ be the canonical projection. Then
$\mathrm{V}\cap f_1'^{-1}(z'_0)$ is a nonempty open subset of
$f_1^{-1}(y)\otimes_{k(y)}k(z'_0)$; since this latter prescheme is flat over $f_1^{-1}(y)$, a maximal point $t'_1$ of
$\mathrm{V}\cap f_1'^{-1}(z'_0)$ necessarily lies above the unique maximal point $x$ of $\mathrm{X}_1$ $(2.3.4)$.
Finally, $k(z'_0)$ is a \emph{finite} extension of $k$ (\textbf{I},~$6.4.2$), and the homomorphism
$\mathrm{A}=\sh{O}_{\mathrm{Y},y}\to\mathrm{A}'=\sh{O}_{\mathrm{Z}',z'}$ factors as
$\sh{O}_{\mathrm{Y},y}\to\sh{O}_{\mathrm{Z}',z'_0}\to\sh{O}_{\mathrm{Z}',z'}$; its kernel is therefore contained in that of
$\sh{O}_{\mathrm{Y},y}\to\sh{O}_{\mathrm{Z}',z'_0}$. This completes the announced reduction.

\medskip
\noindent\textbf{II)} \emph{Reduction to the case where there are only finitely many $\mathrm{A}'_\alpha$ and they are finite $\mathrm{A}$-algebras.} --- Let $\mathfrak{m}'_\alpha$ be
the maximal ideal of $\mathrm{A}'_\alpha$. Since $\mathrm{A}'_\alpha$ is a Noetherian local ring, the intersection of the $\mathfrak{m}'^n_\alpha$
is $0$ $(\mathbf{0}_\mathrm{I},~7.3.5)$. The intersection of the $u_\alpha^{-1}(\mathfrak{m}'^n_\alpha)$, over all indices $n$ and $\alpha$, is therefore
equal to the intersection of the kernels of the $u_\alpha$, and hence is reduced to $0$ by hypothesis (i). Since $\mathrm{A}$
is Artinian, finitely many of these ideals already have intersection $0$; denote them by
$u_i^{-1}(\mathfrak{m}'^{n_i}_i)$ $(1\leq i\leq r)$. Since the $\mathrm{A}'_i$ are Noetherian, the $\mathfrak{m}'^h_i/\mathfrak{m}'^{h+1}_i$ are
$(\mathrm{A}'_i/\mathfrak{m}'_i)$-modules of finite length, and since $\mathrm{A}'_i/\mathfrak{m}'_i$ is a finite-dimensional $k$-vector space,
$\mathrm{A}''_i=\mathrm{A}'_i/\mathfrak{m}'^{n_i}_i$ is a finite $\mathrm{A}$-algebra and an Artinian local ring. The announced reduction
therefore follows from $(2.1.4)$.

\medskip
\noindent\textbf{III)} \emph{End of the proof.} --- Suppose now that there are only finitely many $\mathrm{A}'_i$ $(1\leq i\leq r)$ and that
they are finite $\mathrm{A}$-algebras. For every $i$, the residue field $k_i$ of $\mathrm{A}'_i$ is a finite extension of $k$; using
hypothesis (ii), one concludes that the inverse image of $x$ in $\mathrm{X}'_i$ consists of the \emph{single} point $x'_i$ $(4.3.2)$.
Let $\mathrm{Y}'$ be the disjoint union of the $\mathrm{Y}'_i$, let $\mathrm{X}'=\mathrm{X}\times_\mathrm{Y}\mathrm{Y}'$ be the disjoint union of the
$\mathrm{X}'_i$, and put $\sh{F}'=\sh{F}\otimes_\mathrm{Y}\mathrm{Y}'$. The hypothesis implies that $\sh{F}'$ is $\mathrm{Y}'$-flat at every
point of the inverse image of $x$ under the projection $p:\mathrm{X}'\to\mathrm{X}$. Since $p$ is of finite type, there is therefore an
open subset $\mathrm{U}'\supset p^{-1}(x)$ such that $\sh{F}'$ is $\mathrm{Y}'$-flat at every point of $\mathrm{U}'$ \sref{IV.11.1.1}.
Moreover, the morphism $\mathrm{Y}'\to\mathrm{Y}$ is finite because the $\mathrm{A}'_i$ are finite $\mathrm{A}$-algebras. Thus $p$ is a finite
morphism (\textbf{II},~$6.1.5$), and hence is closed (\textbf{II},~$6.1.10$); consequently there is an affine neighborhood $\mathrm{U}$ of
$x$ in $\mathrm{X}$ such that $p^{-1}(\mathrm{U})\subset\mathrm{U}'$. Let $\mathrm{B}$ and $\mathrm{A}'$ be the rings of the schemes
$\mathrm{U}$ and $\mathrm{Y}'$, respectively (so $\mathrm{A}'$ is the direct product of the $\mathrm{A}'_i$). Replacing $\mathrm{X}$ by
$\mathrm{U}$, one has $\sh{F}=\widetilde{\mathrm{M}}$, where $\mathrm{M}$ is a $\mathrm{B}$-module, and by hypothesis
$\mathrm{M}'=\mathrm{M}\otimes_\mathrm{A}\mathrm{A}'$ is a flat $\mathrm{A}'$-module $(2.1.2)$. Since the homomorphism
$\mathrm{A}\to\mathrm{A}'$ is injective by construction, one can apply \sref{IV.11.4.3}, which proves that $\mathrm{M}$ is a flat
$\mathrm{A}$-module.
\end{proof}

\begin{proposition}[11.4.12]
\label{IV.11.4.12}
Let $\mathrm{A}$ be an Artinian local ring with residue field $k$,
$\mathrm{Y}=\operatorname{Spec}(\mathrm{A})$, $f:\mathrm{X}\to\mathrm{Y}$ a morphism locally of finite type, and $\sh{F}$ a coherent $\sh{O}_\mathrm{X}$-module.
Let $(\mathrm{A}'_\alpha)$ be a family of local rings, and for each $\alpha$, $u_\alpha:\mathrm{A}\to\mathrm{A}'_\alpha$ a homomorphism
of rings. Set $\mathrm{Y}'_\alpha=\operatorname{Spec}(\mathrm{A}'_\alpha)$, $\mathrm{X}'_\alpha=\mathrm{X}\times_\mathrm{Y}\mathrm{Y}'_\alpha$,
$f'_\alpha=f_{(\mathrm{Y}'_\alpha)}$, and $\sh{F}'_\alpha=\sh{F}\otimes_\mathrm{A}\mathrm{A}'_\alpha$. Let $x$ be a point
of $\mathrm{X}$, and suppose that the following conditions are satisfied:
\begin{enumerate}[label=(\roman*)]
  \item The intersection of the kernels of the $u_\alpha$ is reduced to $0$.
  \item For all $\alpha$, $\sh{F}'_\alpha$ is $f'_\alpha$-flat at every point $x'_\alpha\in\mathrm{X}'_\alpha$ whose projections in $\mathrm{X}$ and $\mathrm{Y}'_\alpha$
  are $x$ and the closed point $y'_\alpha$ of $\mathrm{Y}'_\alpha$.
\end{enumerate}

Then $\sh{F}$ is $f$-flat at the point $x$.
\end{proposition}

\begin{proof}
\oldpage[IV-3]{150}
By hypothesis, the intersection of the kernels of the $u_\alpha$ is reduced to $0$; since $\mathrm{A}$ is Artinian,
there are already a \emph{finite} number of these kernels whose intersection is $0$, hence one can
restrict to the case where the family $(\mathrm{A}'_\alpha)$ is \emph{finite}. Let $k'$ be an algebraic closure of $k$; one knows
$(\mathbf{0}_{\mathrm{III}},~10.3.1)$ that there exists a local homomorphism $\mathrm{A}\to\mathrm{B}$ making $\mathrm{B}$ a flat $\mathrm{A}$-module,
such that $\mathrm{B}$ is an Artinian local ring, integral over $\mathrm{A}$, and $\mathrm{B}\otimes_\mathrm{A} k$ is isomorphic to $k'$.
By flatness, the kernels of the homomorphisms $\mathrm{B}\to\mathrm{B}'_\alpha=\mathrm{B}\otimes_\mathrm{A}\mathrm{A}'_\alpha$ are deduced from those of
$u_\alpha$ by tensorisation with $\mathrm{B}$, and since they are in finite number, their intersection is
reduced to $0$ $(\mathbf{0}_\mathrm{I},~6.1.3)$. Let us consider the rings $\mathrm{B}'_{\alpha\beta}$, localised of $\mathrm{B}'_\alpha$ at its \emph{maximal} ideals;
one knows (Bourbaki, \emph{Alg.\ comm.}, chap.~II, \S~3, n\textsuperscript{o}~3, cor.~2 of th.~1) that the intersection
of the kernels of the homomorphisms $\mathrm{B}'_\alpha\to\mathrm{B}'_{\alpha\beta}$ (for a given $\alpha$) is reduced to $0$; one therefore concludes
that the intersection of the kernels of the composite homomorphisms $v_{\alpha\beta}:\mathrm{B}\to\mathrm{B}'_\alpha\to\mathrm{B}'_{\alpha\beta}$ ($\alpha$ and
$\beta$ variable) is reduced to $0$. On the other hand, since $\mathrm{B}$ is integral over $\mathrm{A}$, $\mathrm{B}'_\alpha$ is integral over $\mathrm{A}'_\alpha$,
hence its maximal ideals are above the maximal ideal of $\mathrm{A}'_\alpha$. If one sets
$\mathrm{Z}'_{\alpha\beta}=\operatorname{Spec}(\mathrm{B}'_{\alpha\beta})$,
$\mathrm{T}'_{\alpha\beta}=\mathrm{X}'_\alpha\times_{\mathrm{Y}'_\alpha}\mathrm{Z}'_{\alpha\beta}$,
$f'_{\alpha\beta}=f_{(\mathrm{Z}'_{\alpha\beta})}$,
$\sh{G}=\sh{F}\otimes_\mathrm{A}\mathrm{B}$, and
$\sh{G}'_{\alpha\beta}=\sh{G}\otimes_\mathrm{B}\mathrm{B}'_{\alpha\beta}$, one sees
that hypotheses (i) and (ii) are still satisfied when one replaces
$\mathrm{A}$, $\mathrm{A}'_\alpha$, $u_\alpha$, $\sh{F}$ and $x$ by $\mathrm{B}$, $\mathrm{B}'_{\alpha\beta}$, $v_{\alpha\beta}$, $\sh{G}$ and a point $t$ of $\mathrm{T}=\mathrm{X}\otimes_\mathrm{A}\mathrm{B}$ above $x$.
Since the residue field of $\mathrm{B}$ is separably closed, one deduces from \sref{IV.11.4.11} that
$\sh{G}$ is flat over $\mathrm{B}$ at the point $t$. But since $\mathrm{B}$ is a faithfully flat $\mathrm{A}$-module, it follows from $(2.1.4)$
that $\sh{F}$ is flat over $\mathrm{A}$ at the point $x$, which proves \sref{IV.11.4.12}.
\end{proof}

\begin{corollary}[11.4.13]
\label{IV.11.4.13}
Let $\mathrm{A}$ be an Artinian local ring, $\mathrm{Y}=\operatorname{Spec}(\mathrm{A})$, $f:\mathrm{X}\to\mathrm{Y}$ a
morphism locally of finite type, $\sh{F}$ a coherent $\sh{O}_\mathrm{X}$-module. Let $(\mathrm{Y}'_\alpha)$ be a family of
$\mathrm{Y}$-preschemes, and for all $\alpha$, set $\mathrm{X}'_\alpha=\mathrm{X}\times_\mathrm{Y}\mathrm{Y}'_\alpha$, $f'_\alpha=f_{(\mathrm{Y}'_\alpha)}$, and $\sh{F}'_\alpha=\sh{F}\otimes_\mathrm{Y}\mathrm{Y}'_\alpha$. Let $x$ be a point
of $\mathrm{X}$ and suppose the following hypotheses satisfied:
\begin{enumerate}[label=(\roman*)]
  \item The intersection of the kernels of the homomorphisms $\mathrm{A}\to\Gamma(\mathrm{Y}'_\alpha,\sh{O}_{\mathrm{Y}'_\alpha})$ corresponding to the
  structural morphisms is reduced to $0$.
  \item For all $\alpha$, $\sh{F}'_\alpha$ is $f'_\alpha$-flat at every point $x'_\alpha\in\mathrm{X}'_\alpha$ whose projection in $\mathrm{X}$ is $x$.
\end{enumerate}

Then $\sh{F}$ is $f$-flat at the point $x$.
\end{corollary}

\begin{proof}
Indeed, for every $y'_\alpha\in\mathrm{Y}'_\alpha$, consider the local scheme
$\mathrm{Y}''_{\alpha,y'_\alpha}=\operatorname{Spec}(\sh{O}_{\mathrm{Y}'_\alpha,y'_\alpha})$; by
virtue of $(2.1.4)$, $\sh{F}'_\alpha\otimes_{\mathrm{Y}'_\alpha}\mathrm{Y}''_{\alpha,y'_\alpha}$ is
$\mathrm{Y}''_{\alpha,y'_\alpha}$-flat at the points whose projections on $\mathrm{X}$ and
$\mathrm{Y}''_{\alpha,y'_\alpha}$ are $x$ and the closed point of $\mathrm{Y}''_{\alpha,y'_\alpha}$.
Moreover, the kernel of the homomorphism
$\mathrm{A}\to\Gamma(\mathrm{Y}'_\alpha,\sh{O}_{\mathrm{Y}'_\alpha})$ is the intersection of the kernels of the homomorphisms
$\mathrm{A}\to\Gamma(\mathrm{Y}''_{\alpha,y'_\alpha},\sh{O}_{\mathrm{Y}''_{\alpha,y'_\alpha}})$:
indeed, one is immediately reduced to the case where $\mathrm{Y}'_\alpha$ is affine, and it then suffices to apply Bourbaki, \emph{Alg.\ comm.},
chap.~II, \S~3, n\textsuperscript{o}~3, cor.~2 of th.~1. Replacing the family $(\mathrm{Y}'_\alpha)$ by the family of the $\mathrm{Y}''_{\alpha,y'_\alpha}$,
one is therefore brought back to \sref{IV.11.4.12}.
\end{proof}

\subsection{Descent of flatness by arbitrary morphisms: the general case}
\label{subsection:IV.11.5}

\begin{theorem}[11.5.1]
\label{IV.11.5.1}
Let $\mathrm{Y}$ be a locally Noetherian prescheme, $f:\mathrm{X}\to\mathrm{Y}$ a
morphism locally of finite type, $\sh{F}$ a coherent $\sh{O}_\mathrm{X}$-module, $x$ a point of $\mathrm{X}$, $y=f(x)$. Let $(\mathrm{Y}'_\alpha)$
be a family of local $\mathrm{Y}$-preschemes $\mathrm{Y}'_\alpha=\operatorname{Spec}(\mathrm{A}'_\alpha)$ such that the images of the closed points $y'_\alpha$
of $\mathrm{Y}'_\alpha$ are all equal to $y$. For all $\alpha$, let $\mathfrak{m}'_\alpha$ be the maximal ideal of $\mathrm{A}'_\alpha$, and $u_\alpha:\sh{O}_y\to\mathrm{A}'_\alpha$
\oldpage[IV-3]{151}
the canonical homomorphism (\textbf{I},~$2.4.4$); suppose that the finite intersections of the ideals
$u_\alpha^{-1}(\mathfrak{m}'_\alpha{}^{n_\alpha})$
form a fundamental system of neighbourhoods of $0$ in $\sh{O}_y$. Set $\mathrm{X}'_\alpha=\mathrm{X}\times_\mathrm{Y}\mathrm{Y}'_\alpha$, $f'_\alpha=f_{(\mathrm{Y}'_\alpha)}$,
$\sh{F}'_\alpha=\sh{F}\otimes_\mathrm{Y}\mathrm{Y}'_\alpha$, and suppose that one of the following hypotheses holds:
\begin{enumerate}[label=(\roman*)]
  \item For all $\alpha$, $\sh{F}'_\alpha$ is $f'_\alpha$-flat at every point whose projection in $\mathrm{X}$ is equal
  to $x$ and whose projection in $\mathrm{Y}'_\alpha$ is equal to $y'_\alpha$.
  \item For all $\alpha$, there exists a point $x'_\alpha\in\mathrm{X}'_\alpha$ whose projection on $\mathrm{X}$ is $x$ and whose projection in $\mathrm{Y}'_\alpha$
  is equal to $y'_\alpha$, such that $\sh{F}'_\alpha$ is $f'_\alpha$-flat at the point $x'_\alpha$, and $k(x)$ is a primary extension of $k(y)$.
\end{enumerate}

Under these conditions, $\sh{F}$ is $f$-flat at the point $x$.
\end{theorem}

\begin{proof}
Let $\mathfrak{m}_y$ be the maximal ideal of $\sh{O}_y$; since $\sh{O}_y$ and $\sh{O}_x$ are Noetherian and $\sh{O}_y\to\sh{O}_x$ a
homomorphism of local rings, it suffices, by virtue of $(\mathbf{0}_{\mathrm{III}},~10.2.2)$, to prove that for
every $n>0$, $\sh{F}_x/\mathfrak{m}_y^n\sh{F}_x$ is a flat $(\sh{O}_y/\mathfrak{m}_y^n)$-module. Let us denote by $(\mathfrak{J}_\lambda)$ the family of
finite intersections of the $u_\alpha^{-1}(\mathfrak{m}'_\alpha{}^{n_\alpha})$; by hypothesis, there exists $\lambda$ such that
$\mathfrak{J}_\lambda\subset\mathfrak{m}_y^n$, and since
$\sh{F}_x\otimes_{\sh{O}_y}(\sh{O}_y/\mathfrak{m}_y^n)=(\sh{F}_x\otimes_{\sh{O}_y}(\sh{O}_y/\mathfrak{J}_\lambda))\otimes_{\sh{O}_y/\mathfrak{J}_\lambda}(\sh{O}_y/\mathfrak{m}_y^n)$, it will suffice to prove that $\sh{F}_x/\mathfrak{J}_\lambda\sh{F}_x$ is
a flat $(\sh{O}_y/\mathfrak{J}_\lambda)$-module. Now, for each $\alpha$ such that
$\mathfrak{J}_\lambda\subset u_\alpha^{-1}(\mathfrak{m}'_\alpha{}^{n_\alpha})$, passage to quotients gives
a homomorphism $u'_\alpha:\sh{O}_y/\mathfrak{J}_\lambda\to\mathrm{A}'_\alpha/\mathfrak{m}'_\alpha{}^{n_\alpha}$, and the intersection of the kernels
of the $u'_\alpha$ is reduced to $0$. Taking into account (\textbf{I},~$3.6.1$), one sees that one is brought back to \sref{IV.11.4.11}
in case (ii), and to \sref{IV.11.4.12} in case (i).
\end{proof}

\begin{corollary}[11.5.2]
\label{IV.11.5.2}
Let $\mathrm{Y}$ be a locally Noetherian prescheme, $f:\mathrm{X}\to\mathrm{Y}$ a
morphism locally of finite type, $\sh{F}$ a coherent $\sh{O}_\mathrm{X}$-module, $x$ a point of $\mathrm{X}$, $y=f(x)$; set
$\mathrm{A}=\sh{O}_y$. Let $u:\mathrm{A}\to\mathrm{A}'$ be a homomorphism from $\mathrm{A}$ into a Zariski ring $\mathrm{A}'$, such that
the inverse image $u^{-1}(\mathfrak{r}')$ of the radical $\mathfrak{r}'$ of $\mathrm{A}'$ is the maximal ideal $\mathfrak{m}$ of $\mathrm{A}$; suppose moreover
that the homomorphism $\hat{u}:\hat{\mathrm{A}}\to\hat{\mathrm{A}}'$ is injective. Set $\mathrm{Y}'=\operatorname{Spec}(\mathrm{A}')$, $\mathrm{X}'=\mathrm{X}\times_\mathrm{Y}\mathrm{Y}'$, $f'=f_{(\mathrm{Y}')}$,
$\sh{F}'=\sh{F}\otimes_\mathrm{Y}\mathrm{Y}'$. For $\sh{F}$ to be $f$-flat at the point $x$, it is necessary and sufficient that $\sh{F}'$ be $f'$-flat at every
point whose projection in $\mathrm{X}$ is equal to $x$ and whose projection in $\mathrm{Y}'$ is equal to a closed point
of $\mathrm{Y}'$.

If moreover $\mathrm{A}'$ is a finite $\mathrm{A}$-algebra, one can in what precedes replace the hypothesis that $\hat{u}$
is injective by the hypothesis that $u$ is injective.
\end{corollary}

\begin{proof}
Since $\mathrm{A}$ (resp.\ $\mathrm{A}'$) is identified with a subring of $\hat{\mathrm{A}}$ (resp.\ $\hat{\mathrm{A}}'$) (Bourbaki,
\emph{Alg.\ comm.}, chap.~III, \S~3, n\textsuperscript{o}~3, prop.~6), one sees first that $u$ is itself injective
and that $\hat{u}$ is its extension by continuity to $\hat{\mathrm{A}}$.

Let $(\mathfrak{m}'_\alpha)$ be the family of maximal ideals of $\mathrm{A}'$; since one has
\[
\mathfrak{m}'_\alpha{}^n\hat{\mathrm{A}}'=\hat{\mathfrak{m}}'_\alpha{}^n, \qquad \text{and} \qquad \hat{\mathfrak{m}}'_\alpha{}^n\cap\mathrm{A}'=\mathfrak{m}'_\alpha{}^n,
\]
and since the $\mathfrak{m}'_\alpha{}^n$ are open in $\mathrm{A}'$, one has
$u^{-1}(\mathfrak{m}'_\alpha{}^n)=\mathrm{A}\cap\hat{u}^{-1}(\hat{\mathfrak{m}}'_\alpha{}^n)$, and it will suffice to prove that in
$\hat{\mathrm{A}}$ the finite intersections of the $\hat{u}^{-1}(\hat{\mathfrak{m}}'_\alpha{}^{n_\alpha})$ form a fundamental system of neighbourhoods
of $0$. This will permit the application of \sref{IV.11.5.1} to the composite homomorphisms
$v_\alpha\circ u:\mathrm{A}\to\mathrm{A}'_{\mathfrak{m}'_\alpha}$, where
$v_\alpha:\mathrm{A}'\to\mathrm{A}'_{\mathfrak{m}'_\alpha}$ is the canonical homomorphism. Since $\hat{\mathrm{A}}$ is complete, it will suffice to
prove that the intersection of the $\hat{u}^{-1}(\hat{\mathfrak{m}}'_\alpha{}^{n_\alpha})$ is reduced to $0$ (Bourbaki,
\emph{Alg.\ comm.}, chap.~III, \S~2, n\textsuperscript{o}~7, prop.~8, where in the proof the decreasing sequence may be replaced by an arbitrary
filtered set). Now, for every $\alpha$
\oldpage[IV-3]{152}
fixed, the intersection of the ideals $\hat{\mathfrak{m}}'_\alpha{}^n\hat{\mathrm{A}}'_{\hat{\mathfrak{m}}'_\alpha}$ for $n>0$ is reduced to $0$ in the
Noetherian local ring $\hat{\mathrm{A}}'_{\hat{\mathfrak{m}}'_\alpha}$.
On the other hand the $\hat{\mathfrak{m}}'_\alpha$ are the maximal ideals of $\hat{\mathrm{A}}'$, hence the canonical homomorphism
$\hat{\mathrm{A}}'\to\prod_\alpha\hat{\mathrm{A}}'_{\hat{\mathfrak{m}}'_\alpha}$ is injective (Bourbaki, \emph{Alg.\ comm.}, chap.~II, \S~3, n\textsuperscript{o}~3, cor.~2 of th.~1), and
since by hypothesis $\hat{u}:\hat{\mathrm{A}}\to\hat{\mathrm{A}}'$ is also injective, this achieves the proof in
the general case. The last assertion follows from the fact that $\hat{\mathrm{A}}$ is an $\mathrm{A}$-module faithfully
flat $(\mathbf{0}_\mathrm{I},~7.3.5)$ and $\hat{\mathrm{A}}'=\mathrm{A}'\otimes_\mathrm{A}\hat{\mathrm{A}}$ since $\mathrm{A}'$ is an $\mathrm{A}$-module of finite type
(Bourbaki, \emph{Alg.\ comm.}, chap.~III, \S~3, n\textsuperscript{o}~4, th.~3 and chap.~IV, \S~2, n\textsuperscript{o}~5, cor.~3 of prop.~9).
\end{proof}

\begin{proposition}[11.5.3]
\label{IV.11.5.3}
Let $f:\mathrm{X}\to\mathrm{Y}$ be a morphism locally of finite presentation,
$\sh{F}$ a quasi-coherent $\sh{O}_\mathrm{X}$-module of finite presentation, $x$ a point of $\mathrm{X}$, $y=f(x)$, $g=(\psi,\theta):\mathrm{Y}'\to\mathrm{Y}$
a proper morphism of finite presentation. Suppose that:
\begin{enumerate}[label=(\roman*)]
  \item The homomorphism $\theta_y:\sh{O}_y\to(g_*(\sh{O}_{\mathrm{Y}'}))_y$ is injective.
  \item For every $x'\in\mathrm{X}'=\mathrm{X}\times_\mathrm{Y}\mathrm{Y}'$ whose projection in $\mathrm{X}$ is $x$, $\sh{F}'=\sh{F}\otimes_\mathrm{Y}\mathrm{Y}'$ is
  $\mathrm{Y}'$-flat at the point $x'$.
\end{enumerate}

Then $\sh{F}$ is $f$-flat at the point $x$.
\end{proposition}

\begin{proof}
The question being local on $\mathrm{X}$, one can suppose $f$ of finite presentation. Let
$p:\mathrm{X}'\to\mathrm{X}$ be the canonical projection. Since
$f'=f_{|\mathrm{Y}'}:\mathrm{X}'\to\mathrm{Y}'$ is of finite presentation
$(1.6.2)$ and $\sh{F}'$ is an $\sh{O}_{\mathrm{X}'}$-module of finite presentation $(\mathbf{0}_\mathrm{I},~5.2.5)$, it follows
from \sref{IV.11.3.1} that the set $\mathrm{U}'$ of points of $\mathrm{X}'$ where $\sh{F}'$ is $f'$-flat is open. Since $\mathrm{U}'$
contains $p^{-1}(x)$ by hypothesis, and $p$ is proper, hence \emph{closed}, $\mathrm{U}'$ contains a set of the form
$p^{-1}(\mathrm{U})$, where $\mathrm{U}$ is a neighbourhood of $x$. Replacing $\mathrm{X}$ by $\mathrm{U}$, one can therefore
suppose that $\sh{F}'$ is $f'$-flat. On the other hand, taking into account (\textbf{I},~$3.6.5$), (\textbf{II},~$5.4.2$)
and $(2.1.4)$, one can also replace $\mathrm{Y}$ by $\operatorname{Spec}(\sh{O}_y)$, i.e.\ suppose that $\mathrm{Y}=\operatorname{Spec}(\mathrm{A})$,
where $\mathrm{A}$ is a local ring. Under these conditions, we shall prove that $\sh{F}$ is $f$-flat. By virtue
of $(5.13.3)$, $\mathrm{A}$ is the inductive limit of local Noetherian sub-rings $\mathrm{A}_\lambda$ such that
the canonical injection $\mathrm{A}_\lambda\to\mathrm{A}$ is a \emph{local} homomorphism. By virtue of $(8.9.1)$, one
can suppose that $\mathrm{X}=\mathrm{X}_\lambda\otimes_{\mathrm{A}_\lambda}\mathrm{A}$, $f=f_\lambda\otimes 1_\mathrm{A}$, $\sh{F}=\sh{F}_\lambda\otimes_{\mathrm{A}_\lambda}\mathrm{A}$,
for a suitable $\lambda$, where $f_\lambda:\mathrm{X}_\lambda\to\mathrm{Y}_\lambda=\operatorname{Spec}(\mathrm{A}_\lambda)$ is a morphism of finite type and
$\sh{F}_\lambda$ is a coherent $\sh{O}_{\mathrm{X}_\lambda}$-module.
Similarly, one can suppose that $\mathrm{Y}'=\mathrm{Y}'_\lambda\otimes_{\mathrm{A}_\lambda}\mathrm{A}$, $g=g_\lambda\otimes 1_\mathrm{A}$, where $g_\lambda:\mathrm{Y}'_\lambda\to\mathrm{Y}_\lambda$ is a
morphism of finite presentation; moreover $(8.10.5,~(\text{xii}))$, one can suppose $g_\lambda$
\emph{proper}. Since by hypothesis the homomorphism $\mathrm{A}\to\Gamma(\mathrm{Y}',\sh{O}_{\mathrm{Y}'})$ is injective and that $\mathrm{A}_\lambda$
is a sub-ring of $\mathrm{A}$, the homomorphism $\mathrm{A}_\lambda\to\Gamma(\mathrm{Y}',\sh{O}_{\mathrm{Y}'})$ is also injective. Finally,
by virtue of \sref{IV.11.2.6}, one can suppose $\lambda$ taken such that $\sh{F}'_\lambda=\sh{F}_\lambda\otimes_{\mathrm{Y}_\lambda}\mathrm{Y}'_\lambda$ is $f'_\lambda$-flat,
since $\sh{F}'=\sh{F}'_\lambda\otimes_{\mathrm{Y}'_\lambda}\mathrm{Y}'$. These remarks show that one can henceforth suppose that
the ring $\mathrm{A}$ is \emph{Noetherian}, the other hypotheses of \sref{IV.11.5.3} being verified. Let us then set
$\mathrm{B}=\hat{\mathrm{A}}$, $\mathrm{Z}=\operatorname{Spec}(\mathrm{B})$; since $\mathrm{B}$ is an $\mathrm{A}$-module faithfully flat $(\mathbf{0}_\mathrm{I},~7.3.5)$, it amounts
to the same to say that $\sh{F}$ is $f$-flat or that $\sh{F}\otimes_\mathrm{Y}\mathrm{Z}$ is $\mathrm{Z}$-flat $(2.1.4)$; similarly, if one
sets $\mathrm{Z}'=\mathrm{Y}'\times_\mathrm{Y}\mathrm{Z}$, the morphism $\mathrm{Z}'\to\mathrm{Y}'$ is faithfully flat $(2.2.13)$, hence
it amounts to the same to say that $\sh{F}'$ is $f'$-flat or that $\sh{F}'\otimes_{\mathrm{Y}'}\mathrm{Z}'$ is $\mathrm{Z}'$-flat; finally $h=g_{(\mathfrak{m})}:\mathrm{Z}'\to\mathrm{Z}$
is proper (\textbf{II},~$5.4.2$) and of finite type $(1.5.2)$, $\hat{\mathrm{A}}$ is Noetherian, and if $z$ is its
\oldpage[IV-3]{153}
closed point, the homomorphism $\sh{O}_z\to(h_*(\sh{O}_{\mathrm{Z}'}))_z$ is \emph{injective}, since it follows from $(2.3.1)$ that
$h_*(\sh{O}_{\mathrm{Z}'})=g_*(\sh{O}_{\mathrm{Y}'})\otimes_\mathrm{Y}\mathrm{Z}$, and our assertion follows from hypothesis (i) and from the definition of
flat modules $(\mathbf{0}_\mathrm{I},~6.1.1)$.

One can therefore now suppose the local Noetherian ring $\mathrm{A}$ \emph{complete}; the proof
will be finished if one proves that the intersection of the kernels of the homomorphisms
$\mathrm{A}=\sh{O}_y\to\sh{O}_{y'}$ (where $y'$ ranges over $g^{-1}(y)$) is reduced to $0$. Indeed, the $\sh{O}_{y'}$ are local
Noetherian rings, hence for each $y'$ the intersection of the $\mathfrak{m}_{y'}^n$ $(n\geq 0)$ is reduced to $0$;
if $\mathfrak{a}_{n,y'}$ is the inverse image of $\mathfrak{m}_{y'}^n$ in $\mathrm{A}$, the intersections of finitely many of the $\mathfrak{a}_{n,y'}$ are
neighbourhoods of $0$ in $\mathrm{A}$ and the intersection of all the $\mathfrak{a}_{n,y'}$ is reduced to $0$; the intersections
of finitely many of the $\mathfrak{a}_{n,y'}$ will therefore form a fundamental system of neighbourhoods of $0$ in $\mathrm{A}$
(Bourbaki, \emph{Alg.\ comm.}, chap.~III, \S~2, n\textsuperscript{o}~7, prop.~8, where one can in the proof
replace the decreasing sequence by an arbitrary filtered set); one will then be able to
apply \sref{IV.11.5.1}. Now, let $s\in\mathrm{A}$ be an element belonging to the kernel of each of
the homomorphisms $\mathrm{A}\to\sh{O}_{y'}$; the image $s'$ of $s$ in $\Gamma(\mathrm{Y}',\sh{O}_{\mathrm{Y}'})$ is therefore a section
of $\sh{O}_{\mathrm{Y}'}$ above $\mathrm{Y}'$ such that $s'_{y'}=0$ for every $y'\in g^{-1}(y)$; then there is a neighbourhood $\mathrm{V}'$
of $g^{-1}(y)$ in $\mathrm{Y}'$ such that $s'|\mathrm{V}'=0$. But since $g$ is \emph{closed}, $\mathrm{V}'$ contains a set
of the form $g^{-1}(\mathrm{V})$, where $\mathrm{V}$ is an open neighbourhood of $y$ in $\mathrm{Y}$; but $\mathrm{Y}$ is a local scheme
and $y$ is its closed point, so $\mathrm{V}=\mathrm{Y}$ and consequently $\mathrm{V}'=\mathrm{Y}'$,
$s'=0$, and since $\mathrm{A}\to\Gamma(\mathrm{Y}',\sh{O}_{\mathrm{Y}'})$ is injective by hypothesis, $s=0$.
\end{proof}

The following particular case of \sref{IV.11.5.3} will be useful in chap.~V:

\begin{corollary}[11.5.4]
\label{IV.11.5.4}
Let $f=(\psi,\theta):\mathrm{X}\to\mathrm{Y}$ be a proper morphism of finite presentation,
and let $p:\mathrm{X}\times_\mathrm{Y}\mathrm{X}\to\mathrm{X}$ be the first projection. Suppose that $\theta:\sh{O}_\mathrm{Y}\to f_*(\sh{O}_\mathrm{X})$ is injective.
Then, for $f$ to be flat, it is necessary and sufficient that $p$ be flat.
\end{corollary}

\begin{proposition}[11.5.5]
\label{IV.11.5.5}
Let $\mathrm{A}$ be a ring, $\mathrm{Y}=\operatorname{Spec}(\mathrm{A})$, $f:\mathrm{X}\to\mathrm{Y}$ a morphism
locally of finite presentation, $\sh{F}$ a quasi-coherent $\sh{O}_\mathrm{X}$-module of finite presentation, $x$ a point
of $\mathrm{X}$. Let $\mathrm{A}\to\mathrm{A}'$ be an injective homomorphism making $\mathrm{A}'$ an integral algebra over $\mathrm{A}$. Set
$\mathrm{Y}'=\operatorname{Spec}(\mathrm{A}')$, $\mathrm{X}'=\mathrm{X}\times_\mathrm{Y}\mathrm{Y}'$, $f'=f\times 1$, $\sh{F}'=\sh{F}\otimes_\mathrm{Y}\mathrm{Y}'$. Then, if $\sh{F}'$ is $f'$-flat at every point
of $\mathrm{X}'$ whose projection in $\mathrm{X}$ is $x$, $\sh{F}$ is $f$-flat at the point $x$.
\end{proposition}

\begin{proof}
Suppose first that $\mathrm{A}'$ is a \emph{finite} $\mathrm{A}$-algebra of \emph{finite presentation}; then the
morphism $\mathrm{Y}'\to\mathrm{Y}$ is proper (\textbf{II},~$6.1.11$) and of finite presentation, hence the hypotheses
of \sref{IV.11.5.3} are satisfied, whence the conclusion. In the general case, the proposition results
from this particular case, from the fact that $\mathrm{A}'$ is the inductive limit of its \emph{finite} sub-$\mathrm{A}$-algebras
$\mathrm{A}'_\lambda$ and of the two following lemmas:
\end{proof}

\begin{lemma}[11.5.5.1]
\label{IV.11.5.5.1}
Every finite $\mathrm{A}$-algebra $\mathrm{A}'$ is the inductive limit of $\mathrm{A}$-algebras $\mathrm{A}'_\lambda$
which are finite and of finite presentation.
\end{lemma}

\begin{proof}
One reasons as in $(1.9.3.1)$. Indeed $\mathrm{A}'=\mathrm{B}/\mathfrak{J}$, where $\mathrm{B}$ is a finite $\mathrm{A}$-algebra
that is free as an $\mathrm{A}$-module, and $\mathfrak{J}$ is an ideal of $\mathrm{B}$ $(1.4.7.1)$. Now, $\mathfrak{J}$ is the inductive
limit of ideals $\mathfrak{J}_\lambda\subset\mathfrak{J}$ of $\mathrm{B}$ which are of finite type (and hence, \emph{a fortiori}, are $\mathrm{A}$-modules \emph{of finite type}),
hence,
by the exactness of the functor $\varinjlim$, $\mathrm{A}'$ is the inductive limit of the $\mathrm{A}$-algebras $\mathrm{B}/\mathfrak{J}_\lambda$; now, $\mathrm{B}/\mathfrak{J}_\lambda$
is by definition an $\mathrm{A}$-module of finite presentation, hence also $(1.4.7)$ a finite $\mathrm{A}$-algebra
of finite presentation.
\end{proof}

\oldpage[IV-3]{154}

To apply this lemma to the situation of \sref{IV.11.5.5}, we note in addition
that if the homomorphism $\mathrm{A}\to\mathrm{A}'$ is injective, then
\emph{a fortiori} so is $\mathrm{A}\to\mathrm{A}'_\lambda$, for every
$\lambda$.

\begin{lemma}[11.5.5.2]
\label{IV.11.5.5.2}
Let $\mathrm{A}$ be a ring, $\mathrm{A}'$ an $\mathrm{A}$-algebra, $(\mathrm{A}'_\lambda)$ an inductive system
of $\mathrm{A}$-algebras such that $\mathrm{A}'=\varinjlim\mathrm{A}'_\lambda$; set $\mathrm{Y}=\operatorname{Spec}(\mathrm{A})$, $\mathrm{Y}'=\operatorname{Spec}(\mathrm{A}')$, $\mathrm{Y}'_\lambda=\operatorname{Spec}(\mathrm{A}'_\lambda)$.
Let $f:\mathrm{X}\to\mathrm{Y}$ be a morphism of finite presentation, $\sh{F}$ a quasi-coherent $\sh{O}_\mathrm{X}$-module of finite
presentation; set $\mathrm{X}'=\mathrm{X}\times_\mathrm{Y}\mathrm{Y}'$, $f'=f_{|\mathrm{Y}'}$, $\sh{F}'=\sh{F}\otimes_\mathrm{Y}\mathrm{Y}'$, $\mathrm{X}'_\lambda=\mathrm{X}\times_\mathrm{Y}\mathrm{Y}'_\lambda$, $f'_\lambda=f_{|\mathrm{Y}'_\lambda}$,
$\sh{F}'_\lambda=\sh{F}\otimes_\mathrm{Y}\mathrm{Y}'_\lambda$. Let $x$ be a point of $\mathrm{X}$, such that $\sh{F}'$ is $f'$-flat at every point $x'\in\mathrm{X}'$ above $x$;
then there exists $\lambda$ such that $\sh{F}'_\lambda$ is $f'_\lambda$-flat at every point $x'\in\mathrm{X}'_\lambda$ above $x$.
\end{lemma}

\begin{proof}
Let $\mathrm{U}'$ be the set of $x'\in\mathrm{X}'$ such that $\sh{F}'$ is $f'$-flat at the point $x'$; one knows \sref{IV.11.3.1}
that $\mathrm{U}'$ is open in $\mathrm{X}'$ since $f'$ is of finite presentation $(1.6.2)$; similarly the
set $\mathrm{U}'_\lambda$ of points of $\mathrm{X}'_\lambda$ where $\sh{F}'_\lambda$ is $f'_\lambda$-flat is open in $\mathrm{X}'_\lambda$, and one knows moreover
\sref{IV.11.2.6} that $\mathrm{U}'$ is the union of the $v_\lambda^{-1}(\mathrm{U}'_\lambda)$, where $v_\lambda:\mathrm{X}'\to\mathrm{X}'_\lambda$ is the canonical projection.
Let us consider the scheme $\mathrm{T}=\operatorname{Spec}(k(x))$; set $\mathrm{T}'=\mathrm{T}\times_\mathrm{Y}\mathrm{Y}'$, $\mathrm{T}'_\lambda=\mathrm{T}\times_\mathrm{Y}\mathrm{Y}'_\lambda$, and let
$w_\lambda:\mathrm{T}'\to\mathrm{T}'_\lambda$, $g':\mathrm{T}'\to\mathrm{X}'$, $g'_\lambda:\mathrm{T}'_\lambda\to\mathrm{X}'_\lambda$ be the canonical projections. Let us set $\mathrm{V}'_\lambda=g_\lambda'^{-1}(\mathrm{U}'_\lambda)$,
$\mathrm{V}'=g'^{-1}(\mathrm{U}')=\bigcup w_\lambda^{-1}(\mathrm{V}'_\lambda)$. By hypothesis (and taking into account (\textbf{I},~$3.6.1$)), $\mathrm{V}'=\mathrm{T}'$;
as $\mathrm{T}'$ is quasi-compact, there exists $\lambda$ such that $w_\lambda^{-1}(\mathrm{V}'_\lambda)=\mathrm{T}'$. One deduces then from $(8.3.3)$
applied to the closed quasi-compact subsets $\mathrm{T}'_\lambda-\mathrm{V}'_\lambda$ of $\mathrm{T}'_\lambda$, that there exists $\mu\geq\lambda$ such
that $\mathrm{T}'_\mu=\mathrm{V}'_\mu$; this means that $\sh{F}'_\mu$ is $f'_\mu$-flat at every point of $\mathrm{X}'_\mu$ whose projection
to $\mathrm{X}$ is $x$.
\end{proof}

\subsection{Descent of flatness by arbitrary morphisms: the case of a unibranch base prescheme}
\label{subsection:IV.11.6}

\begin{theorem}[11.6.1]
\label{IV.11.6.1}
Let $\mathrm{A}$ be an integral geometrically unibranch local ring
$(\mathbf{0},~23.2.1)$, $\mathrm{Y}=\operatorname{Spec}(\mathrm{A})$, $f:\mathrm{X}\to\mathrm{Y}$ a morphism locally of finite presentation, $\sh{F}$ a
quasi-coherent $\sh{O}_\mathrm{X}$-module of finite presentation. Let $\mathrm{A}'$ be a local ring, $u:\mathrm{A}\to\mathrm{A}'$ a local
injective homomorphism; set $\mathrm{Y}'=\operatorname{Spec}(\mathrm{A}')$, $\mathrm{X}'=\mathrm{X}\times_\mathrm{Y}\mathrm{Y}'$, $f'=f_{|\mathrm{Y}'}$, $\sh{F}'=\sh{F}\otimes_\mathrm{Y}\mathrm{Y}'$. Let $x$ be
a point of $\mathrm{X}$ whose projection $f(x)=y$ is the closed point of $\mathrm{Y}$, $x'$ a point of $\mathrm{X}'$ whose projections
in $\mathrm{X}$ and $\mathrm{Y}'$ are respectively $x$ and the closed point $y'$ of $\mathrm{Y}'$. Then, if $\sh{F}'$ is $f'$-flat at the point $x'$,
$\sh{F}$ is $f$-flat at the point $x$.
\end{theorem}

\begin{proof}
One can restrict to the case where $f$ is of finite presentation, the question being local
on $\mathrm{X}$. We will proceed in several steps.

\medskip
\noindent\textbf{I)} \emph{Reduction to the case where $\mathrm{A}$ and $\mathrm{A}'$ are integrally closed local rings.}

\medskip
Since $u$ is injective and $\mathrm{A}$ is integral, there exists a prime ideal $\mathfrak{p}'$ of $\mathrm{A}'$ such that
$u^{-1}(\mathfrak{p}')=(0)$ $(\mathbf{0}_\mathrm{I},~1.5.8)$; the composite homomorphism $\mathrm{A}\to\mathrm{A}'\to\mathrm{A}''=\mathrm{A}'/\mathfrak{p}'$ is therefore injective
and local, and if $\mathrm{Y}''=\operatorname{Spec}(\mathrm{A}'')$, $\mathrm{X}''=\mathrm{X}'\otimes_{\mathrm{A}'}\mathrm{A}''=\mathrm{X}\times_\mathrm{Y}\mathrm{Y}''$, $\sh{F}''=\sh{F}'\otimes_{\mathrm{A}'}\mathrm{A}''$, $\sh{F}''$ is $\mathrm{Y}''$-flat
at the points of $\mathrm{X}''$ above $x'$ $(2.1.4)$; replacing $\mathrm{A}'$ by $\mathrm{A}''$ as needed and taking
into account (\textbf{I},~$3.4.7$), one can therefore first suppose $\mathrm{A}'$ integral. If $\mathrm{K}'$ is the field of
fractions of $\mathrm{A}'$, there is a valuation ring $\mathrm{B}'$ in $\mathrm{K}'$ that dominates $\mathrm{A}'$; the composite
homomorphism $\mathrm{A}\xrightarrow{u}\mathrm{A}'\to\mathrm{B}'$ being injective and local, the same reasoning as
\oldpage[IV-3]{155}
the preceding argument allows one to replace $\mathrm{A}'$ by $\mathrm{B}'$; one may therefore suppose that the local ring $\mathrm{A}'$ is integrally closed,
$\mathrm{A}$ being a local subring of $\mathrm{A}'$ dominated by $\mathrm{A}'$. Let $\mathrm{A}_1$ be the integral closure of $\mathrm{A}$; it is clear that
$\mathrm{A}\subset\mathrm{A}_1\subset\mathrm{A}'$, and by hypothesis $\mathrm{A}_1$ is a \emph{local} ring;
if $\mathfrak{m}$, $\mathfrak{m}_1$, $\mathfrak{m}'$ are the maximal ideals of $\mathrm{A}$, $\mathrm{A}_1$, $\mathrm{A}'$, one has $\mathfrak{m}\subset\mathfrak{m}_1\subset\mathfrak{m}'$; indeed, $\mathfrak{m}_1$ is
the \emph{sole} prime ideal of $\mathrm{A}_1$ above $\mathfrak{m}$, since $\mathrm{A}_1$ is a local ring (Bourbaki,
\emph{Alg.\ comm.}, chap.~V, \S~2, n\textsuperscript{o}~1, prop.~1); since $\mathfrak{m}'\cap\mathrm{A}=\mathfrak{m}$, one has $\mathfrak{m}'\cap\mathrm{A}_1\cap\mathrm{A}=\mathfrak{m}$,
hence $\mathfrak{m}'\cap\mathrm{A}_1=\mathfrak{m}_1$. Set $\mathrm{Y}_1=\operatorname{Spec}(\mathrm{A}_1)$, $\mathrm{X}_1=\mathrm{X}\times_\mathrm{Y}\mathrm{Y}_1$, $f_1=f_{(\mathrm{Y}_1)}$, $\sh{F}_1=\sh{F}\otimes_\mathrm{Y}\mathrm{Y}_1$, and
let $x_1$ be the projection of $x'$ in $\mathrm{X}_1$; let us denote furthermore by $y_1$ the unique closed point
of $\mathrm{Y}_1$, so that $y_1=f_1(x_1)$. By hypothesis the morphism $\operatorname{Spec}(k(y_1))\to\operatorname{Spec}(k(y))$
is \emph{radicial}, whence one concludes, by the transitivity of fibres (\textbf{I},~$3.6.4$) and (\textbf{I},~$3.5.7$), that
the morphism $f_1^{-1}(y_1)\to f^{-1}(y)$ is radicial, and in particular $x_1$ is the \emph{sole} point of $\mathrm{X}_1$
whose projections in $\mathrm{X}$ and $\mathrm{Y}_1$ are $x$ and $y_1$ respectively; moreover, as one has seen, $y_1$
is the \emph{sole} point of $\mathrm{Y}_1$ whose projection in $\mathrm{Y}$ is $y$, hence $x_1$ is the \emph{sole} point of $\mathrm{X}_1$
whose projection in $\mathrm{X}$ is $x$. If one proves that $\sh{F}_1$ is $f_1$-flat at the point $x_1$, one will be able to
apply \sref{IV.11.5.5}, whence the conclusion. One is therefore brought back to the case where $\mathrm{A}$
is also \emph{integrally closed}.

\medskip
\noindent\textbf{II)} \emph{Reduction to the case where $\mathrm{A}$ and $\mathrm{A}'$ are local rings of $\mathbf{Z}$-algebras of finite type integrally
closed.}

\medskip
One may regard $\mathrm{A}'$ as the filtered inductive limit of its finite-type sub-$\mathbf{Z}$-algebras
$\mathrm{B}'_\lambda$ that are integral domains; moreover, since $\mathrm{A}'$ is integrally closed and the integral
closure of a $\mathbf{Z}$-algebra of finite type is also a $\mathbf{Z}$-algebra of finite type $(7.8.3)$, one
sees that $\mathrm{A}'$ is the inductive limit of its sub-$\mathbf{Z}$-algebras of finite type $\mathrm{B}'_\lambda$ \emph{integrally closed};
if $\mathfrak{m}'_\lambda=\mathfrak{m}'\cap\mathrm{B}'_\lambda$, $\mathrm{A}'$ is also the inductive limit of the local subrings $(\mathrm{B}'_\lambda)_{\mathfrak{m}'_\lambda}=\mathrm{A}'_\lambda$
dominated by $\mathrm{A}'$ $(5.13.3)$. For every $\lambda$, $\mathrm{B}'_\lambda\cap\mathrm{A}$ is likewise the inductive limit of its finite-type
sub-$\mathbf{Z}$-algebras $\mathrm{B}_{\alpha\lambda}$, hence $\mathrm{A}=\varinjlim_{\alpha,\lambda}\mathrm{B}_{\alpha\lambda}$; as above one may replace $\mathrm{B}_{\alpha\lambda}$
in this formula by its integral closure (which by hypothesis is contained in $\mathrm{B}'_\lambda\cap\mathrm{A}$), and then by the local ring
$\mathrm{A}_{\alpha\lambda}=(\mathrm{B}_{\alpha\lambda})_{\mathfrak{m}_{\alpha\lambda}}$, where
$\mathfrak{m}_{\alpha\lambda}=\mathfrak{m}\cap\mathrm{B}_{\alpha\lambda}=\mathfrak{m}'_\lambda\cap\mathrm{B}_{\alpha\lambda}$, so that $\mathrm{A}_{\alpha\lambda}$ is dominated by $\mathrm{A}'_\lambda$ and by $\mathrm{A}$.
Set $\mathrm{Y}_{\alpha\lambda}=\operatorname{Spec}(\mathrm{A}_{\alpha\lambda})$; it follows from $(8.9.1)$ that there exist a sufficiently large pair
$(\alpha,\lambda)$, a morphism $f_{\alpha\lambda}:\mathrm{X}_{\alpha\lambda}\to\mathrm{Y}_{\alpha\lambda}$ of finite type, and a
coherent $\sh{O}_{\mathrm{X}_{\alpha\lambda}}$-module $\sh{F}_{\alpha\lambda}$ such that
$\mathrm{X}=\mathrm{X}_{\alpha\lambda}\times_{\mathrm{Y}_{\alpha\lambda}}\mathrm{Y}$,
$f=f_{\alpha\lambda}\times 1_\mathrm{Y}$, and $\sh{F}=\sh{F}_{\alpha\lambda}\otimes_{\mathrm{Y}_{\alpha\lambda}}\mathrm{Y}$; if $x_{\alpha\lambda}$
is the projection of $x$ in $\mathrm{X}_{\alpha\lambda}$, it will suffice to show that $\sh{F}_{\alpha\lambda}$ is $f_{\alpha\lambda}$-flat at the point $x_{\alpha\lambda}$.
Since $\mathrm{A}'$ is the inductive limit of the $\mathrm{A}'_\mu$ for $\mu\geq\lambda$, $\mathrm{X}'$ is the projective limit of the
$\mathrm{X}_{\alpha\lambda}\times_{\mathrm{Y}_{\alpha\lambda}}\mathrm{Y}'_\mu=\mathrm{X}'_\mu$, and one also has
$\sh{F}'=\sh{F}'_{\alpha\mu}\otimes_{\mathrm{Y}'_\mu}\mathrm{Y}'$, where
$\sh{F}'_{\alpha\mu}=\sh{F}_{\alpha\lambda}\otimes_{\mathrm{Y}_{\alpha\lambda}}\mathrm{Y}'_\mu$. Applying
\sref{IV.11.2.6}, one sees that one can take $\mu$ large enough for $\sh{F}'_{\alpha\mu}$ to be $\mathrm{Y}'_\mu$-flat at the
point $x'_\mu$, projection of $x'$ in $\mathrm{X}'_\mu$, and moreover, by construction of the $\mathrm{A}'_\mu$, the projection
of $x'_\mu$ in $\mathrm{Y}'_\mu$ is the closed point of $\mathrm{Y}'_\mu$.

\medskip
\noindent\textbf{III)} \emph{Reduction to the case where the residue field $k'$ of $\mathrm{A}'$ is a finite extension and radicial of the
residue field $k$ of $\mathrm{A}$.}

\medskip
One can first of all repeat the reasoning of part I) of the
proof of \sref{IV.11.4.11}, taking into account the fact that $\mathbf{Z}$ is a
\oldpage[IV-3]{156}
Jacobson ring; one is thus reduced to the case where $k'$ is a finite extension of $k$, which we suppose from now on.
Let $k''$ be the largest separable extension of $k$ contained in $k'$, and let $k_1$ be a finite Galois extension of $k$ containing $k''$, so that $k''\otimes_k k_1$ is a direct product of fields isomorphic to $k_1$.
Since $k'$ is a radicial extension of $k''$, $k'\otimes_k k_1$ is therefore a direct product of radicial extensions of $k_1$.

There exists a local ring $\mathrm{A}_1$ which is an $\mathrm{A}$-algebra and a free $\mathrm{A}$-module of finite type, such that $\mathrm{A}_1/\mathfrak{m}\mathrm{A}_1$ is $k$-isomorphic to $k_1$ $(\mathbf{0}_{\mathrm{III}},~10.3.1.2)$.
More precisely, write $k_1=k[\mathrm{T}]/(r)$, where $r$ is a monic irreducible separable polynomial of degree $n$, and let $\mathrm{R}\in\mathrm{A}[\mathrm{T}]$ be a monic polynomial whose canonical image is $r$; one may take $\mathrm{A}_1=\mathrm{A}[\mathrm{T}]/(\mathrm{R})$.
If $\mathrm{K}$ is the field of fractions of $\mathrm{A}$, it is clear that $\mathrm{R}$ is an irreducible separable polynomial of $\mathrm{K}[\mathrm{T}]$; it follows first that $\mathrm{A}_1$ is an integral ring whose field of fractions $\mathrm{K}_1=\mathrm{K}[\mathrm{T}]/(\mathrm{R})$ is a separable extension of $\mathrm{K}$.
Moreover, if $t$ is the canonical image of $\mathrm{T}$ in $\mathrm{A}_1$, the $t^j$ $(0\leq j<n)$ form a basis of the $\mathrm{A}$-module $\mathrm{A}_1$, and their images in $k_1$ form a basis over $k$; hence $d=\det\bigl(\operatorname{Tr}_{\mathrm{A}_1/\mathrm{A}}(t^{i+j})\bigr)$ is an element of $\mathrm{A}$ whose class in $k$ is nonzero, and which is consequently invertible.
The same reasoning as in \sref{IV.6.12.4.1},~I) then proves that the morphism $\operatorname{Spec}(\mathrm{A}_1)\to\operatorname{Spec}(\mathrm{A})$ is flat and has regular fibres; one concludes from \sref{IV.6.5.4},~(ii), that $\mathrm{A}_1$ is \emph{integrally closed}.

Consider now the ring $\mathrm{A}'_1=\mathrm{A}'\otimes_\mathrm{A}\mathrm{A}_1$.
It is a free $\mathrm{A}'$-module of finite type, hence a semilocal ring (Bourbaki, \emph{Alg.\ comm.}, chap.~IV, \S~2, n\textsuperscript{o}~5, cor.~3 of prop.~9); moreover, the maximal ideals of this finite $\mathrm{A}'$-algebra all lie above the maximal ideal $\mathfrak{m}'$ of $\mathrm{A}'$, and \emph{a fortiori} contain $\mathfrak{m}\mathrm{A}'_1$.
But
\[
\mathrm{A}'_1/\mathfrak{m}\mathrm{A}'_1=(\mathrm{A}'/\mathfrak{m}\mathrm{A}')\otimes_k k_1,
\]
and, since $k_1$ is a finite separable extension of $k$, the radical of $\mathrm{A}'_1/\mathfrak{m}\mathrm{A}'_1$ is equal to $(\mathfrak{m}'/\mathfrak{m}\mathrm{A}')\otimes_k k_1$ (Bourbaki, \emph{Alg.}, chap.~VIII, \S~7, n\textsuperscript{o}~2, cor.~2 of prop.~3).
If $\mathfrak{n}_i$ $(1\leq i\leq r)$ are the maximal ideals of $\mathrm{A}'_1$, the fields $\mathrm{A}'_1/\mathfrak{n}_i$ are therefore the fields composing the algebra $k'\otimes_k k_1$; in other words, they are \emph{finite radicial extensions of $k_1$}.

Since $\mathrm{A}\to\mathrm{A}'$ is injective and $\mathrm{A}_1$ is a flat $\mathrm{A}$-module, the homomorphism $\mathrm{A}_1\to\mathrm{A}'_1$ is injective.
The canonical homomorphism
\[
\mathrm{A}'_1\longrightarrow\prod_{i=1}^{r}(\mathrm{A}'_1)_{\mathfrak{n}_i}
\]
is also injective (Bourbaki, \emph{Alg.\ comm.}, chap.~II, \S~3, n\textsuperscript{o}~3, th.~1), and consequently so is the composite from $\mathrm{A}_1$ to this product.
Now $\mathrm{A}_1$ is integral and there are only finitely many kernels of the homomorphisms $\mathrm{A}_1\to(\mathrm{A}'_1)_{\mathfrak{n}_i}$; since their intersection is zero, one of them is already zero.
In other words, for some $i$ there is a ring $\mathrm{B}_1=(\mathrm{A}'_1)_{\mathfrak{n}_i}$ such that the homomorphism $\mathrm{A}_1\to\mathrm{B}_1$ is \emph{injective} and \emph{local}.

Set $\mathrm{Y}'_1=\operatorname{Spec}(\mathrm{B}_1)$, $\mathrm{X}'_1=\mathrm{X}'\times_{\mathrm{Y}'}\mathrm{Y}'_1$, and $\sh{F}'_1=\sh{F}'\otimes_{\mathrm{Y}'}\mathrm{Y}'_1$.
The module $\sh{F}'_1$ is $\mathrm{Y}'_1$-flat at all points of $\mathrm{X}'_1$ above $x'$; moreover, the maximal ideal of $\mathrm{B}_1$ is the only one above $\mathfrak{m}'$, and hence all these points project to the closed point $y'_1$ of $\mathrm{Y}'_1$.
Let $x'_1$ be one of these points.
On the other hand, set $\mathrm{Y}_1=\operatorname{Spec}(\mathrm{A}_1)$, $\mathrm{X}_1=\mathrm{X}\times_\mathrm{Y}\mathrm{Y}_1$, and $\sh{F}_1=\sh{F}\otimes_\mathrm{Y}\mathrm{Y}_1$; if $x_1$ is the projection of $x'_1$ in $\mathrm{X}_1$, its projections in $\mathrm{X}$ and $\mathrm{Y}_1$ are respectively $x$ and the closed point $y_1$.
If one proves that $(\sh{F}_1)_{x_1}$ is a flat $\sh{O}_{y_1}$-module, it follows that $\sh{F}_x$ is a flat $\sh{O}_y$-module: indeed, $\sh{O}_{y_1}$ is a flat $\sh{O}_y$-module, so $(\sh{F}_1)_{x_1}$ is a flat $\sh{O}_y$-module $(\mathbf{0}_\mathrm{I},~6.2.1)$.

\oldpage[IV-3]{157}
But $(\sh{F}_1)_{x_1}=\sh{F}_x\otimes_{\sh{O}_x}\sh{O}_{x_1}$, and $\sh{O}_{x_1}$ is a faithfully flat $\sh{O}_x$-module; hence \sref{IV.2.2.11},~(iii), $\sh{F}_x$ is a flat $\sh{O}_y$-module.
Since $\mathrm{X}'_1=\mathrm{X}_1\times_{\mathrm{Y}_1}\mathrm{Y}'_1$ and $\sh{F}'_1=\sh{F}_1\otimes_{\mathrm{Y}_1}\mathrm{Y}'_1$, we have indeed reduced to the situation of \sref{IV.11.6.1}, with $\mathrm{A}$ and $\mathrm{A}'$ replaced respectively by $\mathrm{A}_1$ and $\mathrm{B}_1$.

\medskip
\noindent\textbf{IV)} \emph{End of the proof.} --- We are finally reduced to proving \sref{IV.11.6.1} under the following supplementary hypotheses:
\begin{enumerate}[label=(\roman*)]
  \item $\mathrm{A}$ and $\mathrm{A}'$ are local rings of $\mathbf{Z}$-algebras of finite type (hence \emph{excellent} rings \sref{IV.7.8.3});
  \item $\mathrm{A}$ is integrally closed;
  \item the residue field $k'$ of $\mathrm{A}'$ is a finite radicial extension of the residue field $k$ of $\mathrm{A}$.
\end{enumerate}
It is then known (\sref{IV.7.8.3} and \sref{IV.2.3.8}) that, under conditions~(i) and~(ii), the $\mathfrak{m}$-adic topology on $\mathrm{A}$ is induced by the $\mathfrak{m}'$-adic topology on $\mathrm{A}'$.
The homomorphism on completions $\widehat{u}:\widehat{\mathrm{A}}\to\widehat{\mathrm{A}'}$ is therefore injective.
On the other hand, since the morphism $\operatorname{Spec}(k')\to\operatorname{Spec}(k)$ is radicial, the same holds for $f'^{-1}(y')\to f^{-1}(y)$ \sref[I]{I.3.5.7}; there is therefore only one point $x'$ of $\mathrm{X}'$ whose projections in $\mathrm{X}$ and $\mathrm{Y}'$ are respectively $x$ and $y'$.
One can therefore apply \sref{IV.11.5.2}.
\end{proof}

\begin{corollary}[11.6.2]
\label{IV.11.6.2}
Let $\mathrm{A}$ be a unibranch integral local ring, $\mathrm{Y}=\operatorname{Spec}(\mathrm{A})$, $f:\mathrm{X}\to\mathrm{Y}$ a morphism locally of finite presentation, and $\sh{F}$ a quasi-coherent $\sh{O}_\mathrm{X}$-module of finite presentation.
Let $\mathrm{A}'$ be a local ring and $\mathrm{A}\to\mathrm{A}'$ an injective local homomorphism; set $\mathrm{Y}'=\operatorname{Spec}(\mathrm{A}')$, $\mathrm{X}'=\mathrm{X}\times_\mathrm{Y}\mathrm{Y}'$, $f'=f_{(\mathrm{Y}')}$, and $\sh{F}'=\sh{F}\otimes_\mathrm{Y}\mathrm{Y}'$.
Let $x$ be a point of $\mathrm{X}$ whose projection $f(x)=y$ is the closed point of $\mathrm{Y}$, and suppose that $\sh{F}'$ is $f'$-flat at all points $x'$ of $\mathrm{X}'$ whose projections in $\mathrm{X}$ and $\mathrm{Y}'$ are respectively $x$ and the closed point $y'$ of $\mathrm{Y}'$.
Then $\sh{F}$ is $f$-flat at $x$.
\end{corollary}

\begin{proof}
One may repeat part~I) of the proof of \sref{IV.11.6.1}, which shows (with the same notation) that if $\sh{F}_1$ is $f_1$-flat at \emph{all} points $x_1$ of $\mathrm{X}_1$ whose projections in $\mathrm{X}$ and $\mathrm{Y}_1$ are respectively $x$ and $y_1$, then $\sh{F}$ is $f$-flat at $x$.
One is thus reduced to the case where $\mathrm{A}$ is integrally closed, hence geometrically unibranch, and the conclusion follows from \sref{IV.11.6.1}.
\end{proof}

\subsection{Counter-examples}
\label{subsection:IV.11.7}

\begin{env}[11.7.1]
\label{IV.11.7.1}
Let us first consider the case where $\mathrm{A}$ is an \emph{Artinian} local ring and where the hypotheses of \sref{IV.11.4.11} are satisfied \emph{except} for condition~(ii) concerning the residue field $k$ of $\mathrm{A}$. We shall see that the conclusion of \sref{IV.11.4.11} can then fail.
Let $k$ be a field admitting a Galois extension $k'$ of degree $[k':k]>1$, and let $\Gamma$ denote its Galois group.
Let $\mathrm{A}$ be a $k$-algebra having a basis of three elements $1,a,b$, with multiplication table $a^2=ab=ba=b^2=0$, so that $\mathrm{A}$ is an Artinian local ring whose maximal ideal $\mathfrak{m}=ka+kb$ has square zero.
Let $\mathrm{A}'=\mathrm{A}\otimes_k k'$, which is a $k'$-algebra with basis $1,a,b$ and an Artinian local ring with square-zero maximal ideal $\mathfrak{m}'=k'a+k'b=\mathfrak{m}\mathrm{A}'$; the ring $\mathrm{A}$ identifies canonically with a subring of $\mathrm{A}'$.
Let $\mathfrak{J}$ be the sub-$k'$-vector space of $\mathfrak{m}'$ generated by $a+\gamma b$, where $\gamma\in k'$ does not belong to $k$; it is clear that $\mathfrak{J}$ is an ideal of $\mathrm{A}'$.
Set $\mathrm{B}=\mathrm{A}'/\mathfrak{J}$; this is an Artinian ring which is a non-flat $\mathrm{A}$-module. Otherwise (Bourbaki, \emph{Alg.\ comm.}, chap.~II, \S~3, n\textsuperscript{o}~2, cor.~2 of prop.~5), $\mathrm{B}$ would be a free $\mathrm{A}$-module. Since $\mathrm{A}'$ is likewise a free $\mathrm{A}$-module and the canonical homomorphism $\mathrm{A}'/\mathfrak{m}\mathrm{A}'\to\mathrm{B}/\mathfrak{m}\mathrm{B}$ is bijective, the canonical homomorphism $\mathrm{A}'\to\mathrm{B}=\mathrm{A}'/\mathfrak{J}$ would also
\oldpage[IV-3]{158}
be bijective (loc.~cit., n\textsuperscript{o}~2, cor.~of prop.~6), which is absurd. In other words, if $\mathrm{Y}=\Spec(\mathrm{A})$ and $\mathrm{X}=\Spec(\mathrm{B})$, then $\sh{O}_\mathrm{X}$ is not $\mathrm{Y}$-flat at the unique point $x$ of $\mathrm{X}$.
Now put $\mathrm{A}_1=\mathrm{B}$ and $\mathrm{Y}_1=\Spec(\mathrm{A}_1)$, and consider the canonical homomorphism $\mathrm{A}\to\mathrm{A}_1$, which is local and injective since $\mathfrak{J}\cap\mathrm{A}=0$. If $\mathrm{B}_1=\mathrm{B}\otimes_\mathrm{A}\mathrm{A}_1=\mathrm{B}\otimes_\mathrm{A}\mathrm{B}$, we shall see that there is a point of $\mathrm{X}_1=\mathrm{X}\times_\mathrm{Y}\mathrm{Y}_1=\Spec(\mathrm{B}_1)$ at which $\sh{O}_{\mathrm{X}_1}$ is $\mathrm{Y}_1$-flat.
For this, note that
\[
\mathrm{B}\otimes_\mathrm{A}\mathrm{B}
=\frac{\mathrm{A}'\otimes_\mathrm{A}\mathrm{A}'}
{\operatorname{Im}(\mathfrak{J}\otimes_\mathrm{A}\mathrm{A}')+\operatorname{Im}(\mathrm{A}'\otimes_\mathrm{A}\mathfrak{J})}.
\]
The structure of $\mathrm{C}'=\mathrm{A}'\otimes_\mathrm{A}\mathrm{A}'$ is obtained easily. Consider the product $\mathrm{A}$-algebra $\mathrm{C}=\prod_{\sigma\in\Gamma}\mathrm{A}'_\sigma$, where all the $\mathrm{A}'_\sigma$ are equal to $\mathrm{A}'$, and the canonical map $\varphi:\mathrm{A}'\otimes_\mathrm{A}\mathrm{A}'\to\mathrm{C}$ defined by $\varphi(x\otimes y)=(\sigma(x)y)_{\sigma\in\Gamma}$, where $\Gamma$ acts canonically on $\mathrm{A}'$. Passing to the quotients gives a homomorphism $\mathrm{C}'/\mathfrak{m}\mathrm{C}'\to\mathrm{C}/\mathfrak{m}\mathrm{C}$ which is none other than the canonical homomorphism $k'\otimes_k k'\to\prod_\sigma k'_\sigma$, with $k'_\sigma=k'$ for all $\sigma\in\Gamma$. This last homomorphism is bijective (Bourbaki, \emph{Alg.}, chap.~VIII, \S~8, prop.~4), and hence so is $\varphi$, since $\mathrm{C}'$ and $\mathrm{C}$ are free $\mathrm{A}$-modules (Bourbaki, \emph{Alg.\ comm.}, chap.~II, \S~3, n\textsuperscript{o}~2, cor.~of prop.~6).
It follows that $\mathrm{B}\otimes_\mathrm{A}\mathrm{B}$ is a semilocal $\mathrm{A}$-algebra which is the direct product of the local $\mathrm{A}$-algebras $\mathrm{A}'/(\sigma(\mathfrak{J})+\mathfrak{J})$. The factor corresponding to the identity element of $\Gamma$ is isomorphic to $\mathrm{A}'/\mathfrak{J}$, and is therefore a flat $\mathrm{A}_1$-module. Since $\gamma\notin k$, however, there is at least one $\sigma\in\Gamma$ such that $\sigma(\mathfrak{J})\ne\mathfrak{J}$; then $\sigma(\mathfrak{J})+\mathfrak{J}=\mathfrak{m}'$, and $\mathrm{A}'/\mathfrak{m}'$ is not a flat $(\mathrm{A}'/\mathfrak{J})$-module, since it is the quotient of $\mathrm{A}'/\mathfrak{J}$ by a nonzero ideal.
\end{env}

\begin{env}[11.7.2]
\label{IV.11.7.2}
We shall now show that the result of \sref{IV.11.5.4} loses its validity when $f$ is no longer assumed to be proper (and \emph{a fortiori} \sref{IV.11.5.3} ceases to be exact when $g$ is no longer assumed proper).
Let $k$ be a field, let $\mathrm{A}_0=k[\mathrm{S},\mathrm{T}]$, and let $\mathrm{A}=\mathrm{A}_0/\mathrm{A}_0\mathrm{S}\mathrm{T}$. Thus $\mathrm{Y}=\Spec(\mathrm{A})$ is the reducible curve formed by the two ``coordinate axes'' in the affine plane $\mathrm{V}_k^2=\Spec(\mathrm{A}_0)$.
The ring $\mathrm{A}$ has two minimal prime ideals $\mathfrak{p}_1=\mathrm{A}_0\mathrm{S}/\mathrm{A}_0\mathrm{S}\mathrm{T}$ and $\mathfrak{p}_2=\mathrm{A}_0\mathrm{T}/\mathrm{A}_0\mathrm{S}\mathrm{T}$. Since $\mathrm{A}$ is reduced, it embeds canonically into $\mathrm{B}=\mathrm{B}_1\oplus\mathrm{B}_2$, where $\mathrm{B}_1=\mathrm{A}/\mathfrak{p}_1$ and $\mathrm{B}_2=\mathrm{A}/\mathfrak{p}_2$. Moreover, $\mathrm{B}_1$ identifies with $k[\mathrm{T}]$ and $\mathrm{B}_2$ with $k[\mathrm{S}]$; these are integrally closed integral rings, and consequently $\mathrm{Z}=\Spec(\mathrm{B})$ is precisely the normalization of $\mathrm{Y}$ relative to $\mathrm{R}(\mathrm{Y})$ \sref[II]{II.6.3.8}, the sum of the two schemes $\mathrm{Z}_1=\Spec(\mathrm{B}_1)$ and $\mathrm{Z}_2=\Spec(\mathrm{B}_2)$.
Denote by $y$ the ``double point'' of $\mathrm{Y}$ corresponding to the maximal ideal $\mathfrak{p}_1+\mathfrak{p}_2=\mathfrak{m}$ of $\mathrm{A}$, and by $z_1,z_2$ the points of $\mathrm{Z}$ above $y$, corresponding respectively to the maximal ideals $\mathfrak{m}_1=(\mathrm{T})$ and $\mathfrak{m}_2=(\mathrm{S})$ of $\mathrm{B}_1$ and $\mathrm{B}_2$.
Let $\mathrm{X}$ be the subprescheme of $\mathrm{Z}$ induced on the complement of $z_2$. Thus $\mathrm{X}=\Spec(\mathrm{B}_1\oplus(\mathrm{B}_2)_\mathrm{S})$. It is immediate that the homomorphism $\mathrm{A}\to\mathrm{A}'=\mathrm{B}_1\oplus(\mathrm{B}_2)_\mathrm{S}$ is injective, but the corresponding morphism $f:\mathrm{X}\to\mathrm{Y}$ is not closed, since $f(\mathrm{Z}_2-\{z_2\})$ is not closed in $\mathrm{Y}$ although $\mathrm{Z}_2-\{z_2\}$ is closed in $\mathrm{X}$; \emph{a fortiori}, $f$ is not proper.
We shall now see that $f$ is not flat at $z_1$. It is enough to show \sref[0]{0.6.6.2} that $\sh{O}_{z_1}$ is not a faithfully flat $\sh{O}_y$-module, and for this it suffices to observe that the canonical homomorphism $\sh{O}_y\to\sh{O}_{z_1}$ is not injective. This is immediate, since $\sh{O}_{z_1}$ is an integral ring whereas $\sh{O}_y$ has zero-divisors.
Nevertheless, the first projection $p:\mathrm{X}\times_\mathrm{Y}\mathrm{X}\to\mathrm{X}$ is an isomorphism. Indeed,
\[
\mathrm{B}_1\otimes_\mathrm{A}\mathrm{B}_1
=(\mathrm{A}/\mathfrak{p}_1)\otimes_\mathrm{A}(\mathrm{A}/\mathfrak{p}_1)
=\mathrm{A}/\mathfrak{p}_1,
\]
while $(\mathrm{B}_2)_\mathrm{S}\otimes_\mathrm{A}(\mathrm{B}_2)_\mathrm{S}=(\mathrm{B}_2\otimes_\mathrm{A}\mathrm{B}_2)_\mathrm{S}=(\mathrm{B}_2)_\mathrm{S}$ for the same reason; finally, $\mathrm{B}_1\otimes_\mathrm{A}(\mathrm{B}_2)_\mathrm{S}=0$, since the canonical image of $\mathrm{S}$ in $\mathrm{B}_1$ is zero.
\end{env}

\begin{env}[11.7.3]
\label{IV.11.7.3}
The preceding example can be generalized. Over a field $k$, consider a reduced algebraic curve $\mathrm{Y}$ admitting a single ``ordinary double point'' $y$ (a notion that will be defined later in general), together with its normalization $\mathrm{Z}$, such that $g:\mathrm{Z}\to\mathrm{Y}$ is finite, the restriction of $g$ to $\mathrm{Z}-g^{-1}(y)$ is an isomorphism onto $\mathrm{Y}-\{y\}$, and $g^{-1}(y)$ consists of two ``simple'' points $z_1,z_2$. Moreover, the prescheme $g^{-1}(y)$ is the sum of the two preschemes $\Spec(k(z_1))$ and $\Spec(k(z_2))$, each canonically isomorphic to $\Spec(k(y))$.
Let $\mathrm{X}$ be the subprescheme of $\mathrm{Z}$ induced on the open set $\mathrm{Z}-\{z_2\}$. The morphism $f:\mathrm{X}\to\mathrm{Y}$ obtained by restricting $g$ is not proper; otherwise \sref[II]{II.5.4.3} would imply the same for the canonical injection $j:\mathrm{X}\to\mathrm{Z}$, which is not closed.
The morphism $f$ is radicial: for every $y'\in\mathrm{Y}$, the fibre $f^{-1}(y')$ consists of a single point $x'$, and $k(y')\to k(x')$ is an isomorphism. It follows first that the diagonal $\Delta_f:\mathrm{X}\to\mathrm{X}\times_\mathrm{Y}\mathrm{X}$ is bijective \sref{IV.1.7.7.1}. On the other hand, since $f$ is unramified \sref{IV.17.4.2}[(d')], $\Delta_f$ is an open immersion \sref{IV.17.4.2}[(b')]. Consequently $\Delta_f$ is an isomorphism, and the first projection $p:\mathrm{X}\times_\mathrm{Y}\mathrm{X}\to\mathrm{X}$ is its inverse.
Nevertheless, $f$ is not flat at $z_1$. Otherwise $\sh{O}_{z_1}$ would be a faithfully flat $\sh{O}_y$-module \sref[0]{0.6.6.2}; since $\sh{O}_y$ contains two distinct minimal prime ideals $\mathfrak{p}_1,\mathfrak{p}_2$, corresponding to the two branches of $\mathrm{Y}$ at $y$, there would then be two distinct prime ideals of $\sh{O}_{z_1}$ whose inverse images under $u:\sh{O}_y\to\sh{O}_{z_1}$ were $\mathfrak{p}_1$ and $\mathfrak{p}_2$ \sref[0]{0.6.5.1}. This is absurd, since $\sh{O}_{z_1}$ has only two distinct prime ideals, $0$ and its maximal ideal $\mathfrak{m}_1$, while $u^{-1}(\mathfrak{m}_1)$ is the maximal ideal $\mathfrak{m}$ of $\sh{O}_y$.
\end{env}

\begin{env}[11.7.4]
\label{IV.11.7.4}
One will note that in the preceding example the homomorphism $u$ is \emph{injective} when $\mathrm{Y}$ is \emph{irreducible} (one may for example take $\mathrm{Y}=\operatorname{Spec}(k[\mathrm{S},\mathrm{T}]/(\mathrm{S}(\mathrm{S}^2+\mathrm{T}^2)-(\mathrm{S}^2-\mathrm{T}^2)))$, a ``cubic with a double point''). Indeed, after replacing $\mathrm{Y}$ by an affine neighbourhood of $y$, one may suppose that $\mathrm{Y}=\operatorname{Spec}(\mathrm{A})$, where $\mathrm{A}$ is integral, and hence that $\mathrm{Z}=\operatorname{Spec}(\mathrm{B})$, where $\mathrm{B}$ is the integral closure of $\mathrm{A}$. Since $\mathrm{B}$, $\sh{O}_y$, and $\sh{O}_{z_1}$ then identify with subrings of
\oldpage[IV-3]{159}
the field of fractions of $\mathrm{A}$, $u$ is obviously injective.
On the other hand, the homomorphism $\widehat{u}:\widehat{\sh{O}}_y\to\widehat{\sh{O}}_{z_1}$ is not injective: $\widehat{\sh{O}}_{z_1}$ is an integral local ring, since $z_1$ is a simple point, whereas $\widehat{\sh{O}}_y$ has two distinct minimal prime ideals corresponding to the two branches of $\mathrm{Y}$ and therefore has zero-divisors.
This gives an example showing that in the statement of \sref{IV.11.5.2}, the hypothesis that $\widehat{u}$ is injective cannot be replaced by the hypothesis that $u$ itself is injective, even when $\mathrm{A}'$ is a local ring.
Indeed, with the preceding notation, it suffices to take $\mathrm{A}=\sh{O}_y$, $\mathrm{A}'=\sh{O}_{z_1}$, $\mathrm{Y}=\operatorname{Spec}(\mathrm{A})$, and $\mathrm{X}=\operatorname{Spec}(\mathrm{A}')$; the reasoning of \sref{IV.11.7.3} still proves that the first projection $p:\mathrm{X}\times_\mathrm{Y}\mathrm{X}\to\mathrm{X}$ is an isomorphism, although $f$ is not flat at $z_1$.
\end{env}

\begin{env}[11.7.5]
\label{IV.11.7.5}
The examples of \sref{IV.11.7.2} and \sref{IV.11.7.3} explain the restriction to \emph{unibranch} local rings in \sref{IV.11.6.1} and \sref{IV.11.6.2}.
We shall now see that in \sref{IV.11.6.1}, one cannot weaken the hypothesis on $\mathrm{A}$ by supposing only that $\mathrm{A}$ is unibranch.
Consider the complete local integral ring $\mathrm{A}=\mathbf{R}[[\mathrm{U},\mathrm{V}]]/(\mathrm{U}^2+\mathrm{V}^2)$, which is unibranch but not geometrically unibranch \sref{IV.6.15.11}.
One knows (\emph{loc.\ cit.}) that if $u$ and $v$ are the images of $\mathrm{U}$ and $\mathrm{V}$ in $\mathrm{A}$, the integral closure of $\mathrm{A}$ is $\bar{\mathrm{A}}=\mathrm{A}[t]$, where $t=v/u$ and $t^2=-1$, so that $\bar{\mathrm{A}}$ is isomorphic to $\mathbf{C}[[\mathrm{U}]]$.
Set $\mathrm{Y}=\operatorname{Spec}(\mathrm{A})$ and $\mathrm{X}=\operatorname{Spec}(\bar{\mathrm{A}})$, the normalization of $\mathrm{Y}$ \sref[II]{II.6.3.8}, and let $y$ and $x$ be the closed points of $\mathrm{Y}$ and $\mathrm{X}$ respectively.
We shall show that for a suitable local $\mathrm{A}$-algebra $\mathrm{A}'$, if one sets $\mathrm{Y}'=\operatorname{Spec}(\mathrm{A}')$ and $\mathrm{X}'=\mathrm{X}\times_\mathrm{Y}\mathrm{Y}'$, and if $y'$ denotes the closed point of $\mathrm{Y}'$, then $\sh{O}_{\mathrm{X}'}$ is $\mathrm{Y}'$-flat at one point of $\mathrm{X}'$ whose projections in $\mathrm{X}$ and $\mathrm{Y}'$ are respectively $x$ and $y'$, but is not $\mathrm{Y}'$-flat at every point having these projections.
It will then follow from \sref{IV.2.1.4} that $\sh{O}_\mathrm{X}$ is not $\mathrm{Y}$-flat at $x$ (which is moreover trivial \emph{a priori}, since $\bar{\mathrm{A}}$ is not a free $\mathrm{A}$-module).

Let $\mathrm{B}=\mathrm{A}\otimes_\mathbf{R}\mathbf{C}$, which is isomorphic to $\mathbf{C}[[\mathrm{U},\mathrm{V}]]/((\mathrm{U}+i\mathrm{V})(\mathrm{U}-i\mathrm{V}))$.
The ring $\mathrm{B}$ has two minimal prime ideals $\mathfrak{p}'$ and $\mathfrak{p}''$, generated respectively by $u+iv$ and $u-iv$, and $\mathfrak{n}=\mathfrak{p}'+\mathfrak{p}''$ is the maximal ideal of the complete local ring $\mathrm{B}$.
Let $\bar{\mathrm{B}}=\bar{\mathrm{A}}\otimes_\mathbf{R}\mathbf{C}$; it is the direct product of two algebras isomorphic to $\mathbf{C}[[\mathrm{U}]]$, generated by the idempotents $e'=(1\otimes 1+t\otimes i)/2$ and $e''=(1\otimes 1-t\otimes i)/2$.
Since the homomorphism $\mathrm{A}\to\bar{\mathrm{A}}$ is injective, so is $\mathrm{B}\to\bar{\mathrm{B}}$; the images of $u+iv$ and $u-iv$ under this injection are respectively $2ue'$ and $2ue''$.
It follows immediately that $\bar{\mathrm{B}}$ identifies canonically with $(\mathrm{B}/\mathfrak{p}')\oplus(\mathrm{B}/\mathfrak{p}'')$.
Now take for $\mathrm{A}'$ the local ring $\mathrm{B}/\mathfrak{p}'$; then $\bar{\mathrm{A}}\otimes_\mathrm{A}\mathrm{A}'$ identifies with $\bar{\mathrm{B}}\otimes_\mathrm{B}\mathrm{A}'$.
But $(\mathrm{B}/\mathfrak{p}')\otimes_\mathrm{B}(\mathrm{B}/\mathfrak{p}')=\mathrm{B}/\mathfrak{p}'$ and $(\mathrm{B}/\mathfrak{p}')\otimes_\mathrm{B}(\mathrm{B}/\mathfrak{p}'')=\mathrm{B}/\mathfrak{n}$, so $\bar{\mathrm{A}}\otimes_\mathrm{A}\mathrm{A}'$ is isomorphic to $\mathrm{A}'\oplus(\mathrm{B}/\mathfrak{n})$.
This establishes our assertion, since $\mathrm{B}/\mathfrak{n}=\mathrm{A}'/(\mathfrak{n}/\mathfrak{p}')$ is not a flat $\mathrm{A}'$-module; otherwise it would be a free $\mathrm{A}'$-module (Bourbaki, \emph{Alg.\ comm.}, chap.~II, \S~3, no.~2, cor.~2 of prop.~5), which is absurd.
\end{env}

\subsection{A valuative criterion for flatness}
\label{subsection:IV.11.8}

\begin{theorem}[11.8.1]
\label{IV.11.8.1}
Let $f:\mathrm{X}\to\mathrm{Y}$ be a locally of finite presentation morphism, $\mathscr{F}$ a quasi-coherent $\sh{O}_\mathrm{X}$-Module of finite presentation, $x$ a point of $\mathrm{X}$, $y=f(x)$.
Suppose that the local ring $\sh{O}_y$ is integral (resp.\ reduced and Noetherian).
For $\mathscr{F}$ to be $f$-flat at the point $x$, it is necessary and sufficient that, for every valuation ring (resp.\ every discrete valuation ring) $\mathrm{A}'$ and every local homomorphism $\sh{O}_y\to\mathrm{A}'$, the following condition be satisfied: setting $\mathrm{Y}'=\operatorname{Spec}(\mathrm{A}')$, $\mathrm{X}'=\mathrm{X}\times_\mathrm{Y}\mathrm{Y}'$, $f'=f_{(\mathrm{Y}')}$, the $\sh{O}_{\mathrm{X}'}$-Module $\mathscr{F}'=\mathscr{F}\otimes_\mathrm{Y}\mathrm{Y}'$ is $f'$-flat in all the points $x'$ of $\mathrm{X}'$ whose projections in $\mathrm{X}$ and $\mathrm{Y}'$ are respectively $x$ and the closed point $y'$ of $\mathrm{Y}'$.
\end{theorem}

\begin{proof}
The condition being evidently necessary \sref{IV.2.1.4}, it remains to prove that it is sufficient.
One can evidently (\sref[I]{I.3.6.5} and \sref[I]{I.2.4.4}) restrict to the case where $\mathrm{Y}=\operatorname{Spec}(\mathrm{A})$ is the spectrum of a \emph{local} ring $\mathrm{A}$ and $y$ is the closed point of $\mathrm{Y}$.
\medskip

(i) \emph{Case where $\mathrm{A}$ is integral.} --- Let $\mathrm{K}$ be the field of fractions of $\mathrm{A}$, $\mathrm{A}_1$ the integral closure of $\mathrm{A}$; if one sets $\mathrm{Y}_1=\operatorname{Spec}(\mathrm{A}_1)$, $\mathrm{X}_1=\mathrm{X}\times_\mathrm{Y}\mathrm{Y}_1$, $f_1=f_{(\mathrm{Y}_1)}$, it suffices, by virtue of \sref{IV.11.5.5}, to show that $\mathscr{F}_1=\mathscr{F}\otimes_\mathrm{Y}\mathrm{Y}_1$ is $f_1$-flat in \emph{all} the points $x_1$ of $\mathrm{X}_1$ whose projection $x$ is in $\mathrm{X}$.
Now, if $f_1(x_1)=y_1$, then $\mathfrak{i}_{y_1}=\mathfrak{m}_1$, where $\mathfrak{m}_1$ is a prime ideal of $\mathrm{A}_1$ whose trace on $\mathrm{A}$ is $\mathfrak{m}=\mathfrak{i}_y$ (and $\mathfrak{m}_1$ is moreover necessarily a maximal ideal). Let $\mathrm{A}'$ be a valuation ring for $\mathrm{K}$ that dominates $\sh{O}_{y_1}=(\mathrm{A}_1)_{\mathfrak{m}_1}$; since the homomorphism $\mathrm{A}\to\sh{O}_{y_1}$ is local, so is $\mathrm{A}\to\mathrm{A}'$.
There exists then at least one
\oldpage[IV-3]{160}
point $x'\in\mathrm{X}'$ whose projections in $\mathrm{X}_1$ and $\mathrm{Y}'$ are respectively $x_1$ and $y'$ \sref[I]{I.3.4.7}; as the hypothesis states that $\mathscr{F}'$ is $f'$-flat at the point $x'$, and $\sh{O}_{y_1}$ is integrally closed, one can apply \sref{IV.11.6.1}, and one deduces that $\mathscr{F}_1$ is $f_1$-flat at the point $x_1$, whence the theorem in this case.
\medskip

(ii) \emph{Case where $\mathrm{A}$ is reduced and Noetherian.} --- Let $\mathfrak{p}_i$ ($1\leq i\leq m$) be the minimal prime ideals of $\mathrm{A}$; since $\mathrm{A}$ is reduced, it identifies canonically with a subring of the product of the $\mathrm{A}_i=\mathrm{A}/\mathfrak{p}_i$, which are local Noetherian rings; setting $\mathrm{Y}_i=\operatorname{Spec}(\mathrm{A}_i)$, $\mathrm{X}_i=\mathrm{X}\times_\mathrm{Y}\mathrm{Y}_i$, $f_i=f_{(\mathrm{Y}_i)}$, it follows from \sref{IV.11.5.2} that it suffices to show that for each $i$, $\mathscr{F}_i=\mathscr{F}\otimes_\mathrm{Y}\mathrm{Y}_i$ is $f_i$-flat in every point $x_i$ of $\mathrm{X}_i$ whose projections in $\mathrm{X}$ and $\mathrm{Y}_i$ are $x$ and the closed point $y_i$ of $\mathrm{Y}_i$ respectively.
Now, since $\mathrm{A}_i$ is a local Noetherian integral ring, there exists in its field of fractions $\mathrm{K}_i$ a sub-ring $\mathrm{A}''_i$ containing $\mathrm{A}_i$, which is a finite $\mathrm{A}_i$-algebra (hence a semilocal Noetherian ring) and whose local rings are \emph{geometrically unibranch} (\sref{IV.6.15.5} and \sref[0]{0.23.2.5}).
Since the maximal ideals of $\mathrm{A}''_i$ are necessarily above the maximal ideal of $\mathrm{A}_i$, one deduces further from \sref[I]{I.3.4.7} and \sref{IV.11.5.2} that it suffices, setting $\mathrm{Y}''_i=\operatorname{Spec}(\mathrm{A}''_i)$, $\mathrm{X}''_i=\mathrm{X}\times_\mathrm{Y}\mathrm{Y}''_i$, $f''_i=f_{(\mathrm{Y}''_i)}$, to prove that $\mathscr{F}''_i=\mathscr{F}\otimes_\mathrm{Y}\mathrm{Y}''_i$ is $f''_i$-flat in every point $x''_i$ of $\mathrm{X}''_i$ whose projections in $\mathrm{X}$ and $\mathrm{Y}''_i$ are $x$ and the closed point $y''_i$ of $\mathrm{Y}''_i$ respectively.
Now, let $\mathrm{A}'$ be a discrete valuation ring for $\mathrm{K}_i$ that dominates $\mathrm{A}''_i$, and let $x'$ be a point of $\mathrm{X}'$ whose projections in $\mathrm{X}''_i$ and $\mathrm{Y}'$ are $x''_i$ and $y'$ respectively \sref[I]{I.3.4.7}; as $\sh{O}_{y''_i}$ is geometrically unibranch, one can again apply \sref{IV.11.6.1} and one deduces that $\mathscr{F}''_i$ is indeed $f''_i$-flat at the point $x''_i$.
\end{proof}

\begin{remarks}[11.8.2]
\label{IV.11.8.2}
\medskip\noindent
\begin{enumerate}[label=(\roman*)]
  \item In the statement of \sref{IV.11.8.1}, one may restrict to assuming that the condition on $\mathscr{F}'$ is satisfied for the \emph{complete} valuation rings $\mathrm{A}'$ whose residue field is \emph{algebraically closed}.
  Indeed, it is known that every valuation ring $\mathrm{A}'$ is dominated by such a ring $\mathrm{A}''$ \sref[II]{II.7.1.2}, and that if $\mathrm{A}'$ is a discrete valuation ring, it is the same for $\mathrm{A}''$ ($\mathbf{0}_{\mathrm{III}}$,~10.3.1).
  \item The proof of \sref{IV.11.8.1} simplifies when one supposes not only that $\mathrm{A}$ is integral and Noetherian, but that its completion $\hat{\mathrm{A}}$ is also \emph{integral}.
  Replacing $\mathrm{X}$ by $\mathrm{X}\otimes_\mathrm{A}\hat{\mathrm{A}}$ and reasoning as in the proof of \sref{IV.11.5.3}, one can in fact reduce to proving \sref{IV.11.8.1} when $\mathrm{A}=\sh{O}_y$ is integral, Noetherian, and \emph{complete}.
  Now, one knows \sref[II]{II.7.1.7} that such a ring $\mathrm{A}$ is dominated by a discrete complete valuation ring; the conclusion follows therefore directly from \sref{IV.11.5.2}.
\end{enumerate}
\end{remarks}

\subsection{Separating and universally separating families of homomorphisms of sheaves of modules}
\label{subsection:IV.11.9}

\begin{env}[11.9.1]
\label{IV.11.9.1}
Let $\mathrm{X}$ be a prescheme, $(f_\lambda)_{\lambda\in\mathrm{L}}$ a family of morphisms $f_\lambda:\mathrm{Z}_\lambda\to\mathrm{X}$, $\mathscr{F}$ a quasi-coherent $\sh{O}_\mathrm{X}$-Module; for every $\lambda\in\mathrm{L}$, suppose given a quasi-coherent $\sh{O}_{\mathrm{Z}_\lambda}$-Module $\sh{G}_\lambda$ and a homomorphism
\[
  u_\lambda:\mathscr{F}\to(f_\lambda)_*(\sh{G}_\lambda).
\label{IV.11.9.1.1}
\]
\end{env}

\oldpage[IV-3]{161}
We say that the family $(u_\lambda)$ (or the corresponding family of $u_\lambda^{\#}:(f_\lambda)^*(\mathscr{F})\to\sh{G}_\lambda$) is \emph{separating} if the intersection of the kernels of the $u_\lambda$ is \emph{zero}.
In other words, this signifies that for every open $\mathrm{U}$ of $\mathrm{X}$, and every section $t$ of $\mathscr{F}$ above $\mathrm{U}$, such that, for \emph{every} $\lambda$, the section $u_\lambda(t)$ (which, by definition, is a section of $\sh{G}_\lambda$ above $f_\lambda^{-1}(\mathrm{U})$) is zero, then $t$ is itself zero.

\begin{env}[11.9.2]
\label{IV.11.9.2}
With the notations of \sref{IV.11.9.1}, let $\mathrm{M}$ be a second set of indices; for every $\lambda\in\mathrm{L}$, let $(g_{\lambda\mu})_{\mu\in\mathrm{M}}$ be a family of morphisms $g_{\lambda\mu}:\mathrm{Z}'_{\lambda\mu}\to\mathrm{Z}_\lambda$; for every pair $(\lambda,\mu)$, suppose given a quasi-coherent $\sh{O}_{\mathrm{Z}'_{\lambda\mu}}$-Module $\sh{H}_{\lambda\mu}$ and a homomorphism $v_{\lambda\mu}:\sh{G}_\lambda\to(g_{\lambda\mu})_*(\sh{H}_{\lambda\mu})$; set $h_{\lambda\mu}=f_\lambda\circ g_{\lambda\mu}:\mathrm{Z}'_{\lambda\mu}\to\mathrm{X}$ and consider the composite homomorphism
\[
  w_{\lambda\mu}:\mathscr{F}\xrightarrow{u_\lambda}(f_\lambda)_*(\sh{G}_\lambda)\xrightarrow{(f_\lambda)_*(v_{\lambda\mu})}(h_{\lambda\mu})_*(\sh{H}_{\lambda\mu}).
\]
Suppose that, for every $\lambda\in\mathrm{L}$, the family $(v_{\lambda\mu})_{\mu\in\mathrm{M}}$ is separating: then, it is the same for the family $(f_\lambda)_*(v_{\lambda\mu})$ ($\mu\in\mathrm{M}$), as one sees at once.
One concludes that, for the family $(u_\lambda)$ to be separating, it is necessary and sufficient that the family $(w_{\lambda\mu})_{(\lambda,\mu)\in\mathrm{L}\times\mathrm{M}}$ be so.
\end{env}

\begin{env}[11.9.3]
\label{IV.11.9.3}
To verify that the family $(u_\lambda)$ is separating (with the notations of \sref{IV.11.9.1}), one can evidently reduce first to the case where $\mathrm{X}$ is \emph{affine}, the property being local on $\mathrm{X}$.
One can moreover suppose that $\mathrm{Z}_\lambda=\mathrm{X}$ \emph{for every} $\lambda\in\mathrm{L}$.
Indeed, let $\mathrm{M}$ be an index set, sum of a family $(\mathrm{M}_\lambda)_{\lambda\in\mathrm{L}}$, and for every $\lambda\in\mathrm{L}$, let $(\mathrm{Y}_{\lambda\mu})_{\mu\in\mathrm{M}_\lambda}$ be an open affine covering of $\mathrm{Z}_\lambda$; let $j_{\lambda\mu}:\mathrm{Y}_{\lambda\mu}\to\mathrm{Z}_\lambda$ be the canonical injection and set $\sh{H}_{\lambda\mu}=j_{\lambda\mu}^*(\sh{G}_\lambda)=\sh{G}_\lambda|\mathrm{Y}_{\lambda\mu}$.
If one considers the canonical homomorphism $v_{\lambda\mu}:\sh{G}_\lambda\to(j_{\lambda\mu})_*(j_{\lambda\mu}^*(\sh{G}_\lambda))=(j_{\lambda\mu})_*(\sh{H}_{\lambda\mu})$ relative to $j_{\lambda\mu}$ \sref[0]{0.4.4.3.2}, it is immediate that for each $\lambda\in\mathrm{L}$, the family $(v_{\lambda\mu})_{\mu\in\mathrm{M}_\lambda}$ is separating.
By virtue of \sref{IV.11.9.2}, one is therefore reduced to proving that the family of composite homomorphisms $((f_\lambda)_*(v_{\lambda\mu}))\circ u_\lambda$, where the $f_\lambda$ are \emph{affine}, is separating.
But then the $(f_\lambda)_*(\sh{G}_\lambda)$ are quasi-coherent $\sh{O}_\mathrm{X}$-Modules \sref[I]{I.1.6.2} and, by virtue of the definition, one can replace the $\mathrm{Z}_\lambda$ by $\mathrm{X}$ and the $\sh{G}_\lambda$ by $(f_\lambda)_*(\sh{G}_\lambda)$, whence our assertion.

One will note further that if $\mathrm{L}$ is \emph{finite} and the $f_\lambda$ are \emph{quasi-compact}, one can, in the preceding reduction, suppose that the $\mathrm{M}_\lambda$ are also \emph{finite}, whence in this case one is brought back to verifying that a \emph{finite} family of homomorphisms of $\mathscr{F}$ in quasi-coherent $\sh{O}_\mathrm{X}$-Modules is separating.
\end{env}

\begin{env}[11.9.4]
\label{IV.11.9.4}
Consider therefore the case where $\mathrm{Z}_\lambda=\mathrm{X}$ for every $\lambda$, and where $\mathrm{X}=\operatorname{Spec}(\mathrm{A})$ is affine; then $\mathscr{F}=\widetilde{\mathrm{M}}$ and $\sh{G}_\lambda=\widetilde{\mathrm{N}}_\lambda$, where $\mathrm{M}$ and $\mathrm{N}_\lambda$ are $\mathrm{A}$-modules, and $u_\lambda=\widetilde{\varphi}_\lambda$, where the $\varphi_\lambda:\mathrm{M}\to\mathrm{N}_\lambda$ are $\mathrm{A}$-module homomorphisms.
To say that the family $(u_\lambda)$ is separating signifies then that, for every $s\in\mathrm{A}$, the intersection of the kernels of $(\varphi_\lambda)_s:\mathrm{M}_s\to(\mathrm{N}_\lambda)_s$ is zero.
One then also says that the family $(\varphi_\lambda)$ is \emph{separating}.
One notes that if $\mathrm{L}$ is \emph{finite}, this is equivalent to saying that the intersection of the kernels of the $\varphi_\lambda$ is zero, since then
\[
  \bigcap_{\lambda\in\mathrm{L}}\operatorname{Ker}((\varphi_\lambda)_s)
  =\left(\bigcap_{\lambda\in\mathrm{L}}\operatorname{Ker}(\varphi_\lambda)\right)_s
  \quad \sref[0]{0.1.3.2}.
\]
But this relation is no longer exact in general
\oldpage[IV-3]{162}
when $\mathrm{L}$ is \emph{infinite}, and the fact that the intersection of the kernels of the $\varphi_\lambda$ is zero \emph{does not therefore entail},
in general, that the family $(\varphi_\lambda)$ is separating.
For example, suppose $\mathrm{A}$ is a discrete valuation ring with maximal ideal $\mathfrak{m}$, and consider the family of homomorphisms $\varphi_k:\mathrm{A}\to\mathrm{A}/\mathfrak{m}^k$; their kernels are the ideals $\mathfrak{m}^k$, whose intersection is zero. This family is nevertheless not separating, for the fibres of all the $\mathfrak{m}^k$ at the generic point $x$ of $\mathrm{X}=\operatorname{Spec}(\mathrm{A})$ (which is open in $\mathrm{X}$) are equal to $\sh{O}_x=k(x)$, the field of fractions of $\mathrm{A}$, and hence their intersection is not zero.
\end{env}

\begin{env}[11.9.5]
\label{IV.11.9.5}
We will mainly be concerned in what follows with the problem of \emph{change of base} for separating families.
The notations being those of \sref{IV.11.9.1}, consider a morphism $g:\mathrm{X}'\to\mathrm{X}$ and set $\mathscr{F}'=\mathscr{F}\otimes_{\sh{O}_\mathrm{X}}\sh{O}_{\mathrm{X}'}=g^*(\mathscr{F})$, and, for every $\lambda$, $\mathrm{Z}'_\lambda=\mathrm{Z}_\lambda\times_\mathrm{X}\mathrm{X}'$, $f'_\lambda=f_\lambda\times 1$, $\sh{G}'_\lambda=\sh{G}_\lambda\otimes_{\sh{O}_{\mathrm{Z}_\lambda}}\sh{O}_{\mathrm{Z}'_\lambda}=(g'_\lambda)^*(\sh{G}_\lambda)$, where $g'_\lambda:\mathrm{Z}'_\lambda\to\mathrm{Z}_\lambda$ is the projection.
For every $\lambda$, we denote by $u'_\lambda:\mathscr{F}'\to(f'_\lambda)_*(\sh{G}'_\lambda)$ the homomorphism obtained as follows: let
\[
  \psi:(f_\lambda)_*(\sh{G}_\lambda)\to(f_\lambda)_*((g'_\lambda)_*((g'_\lambda)^*(\sh{G}_\lambda)))=g_*((f'_\lambda)_*(\sh{G}'_\lambda))
\]
be the homomorphism $(f_\lambda)_*(\rho_{\sh{G}_\lambda})$, where $\rho_{\sh{G}_\lambda}$ is the canonical homomorphism \sref[0]{0.4.4.3.2} corresponding to $g'_\lambda$.
Then $u'_\lambda$ is defined as the composite
\[
  g^*(\mathscr{F})\xrightarrow{g^*(u_\lambda)}g^*((f_\lambda)_*(\sh{G}_\lambda))\xrightarrow{g^*(\psi)}g^*g_*((f'_\lambda)_*(\sh{G}'_\lambda))\xrightarrow{\sigma}(f'_\lambda)_*(\sh{G}'_\lambda)
\]
where $\sigma=\sigma_{(f'_\lambda)_*(\sh{G}'_\lambda)}$ is the canonical homomorphism \sref[0]{0.4.4.3.3} corresponding to $g$.
We will say for short that $u'_\lambda$ is \emph{deduced from $u_\lambda$ by the change of base} $g$.
When $f_\lambda:\mathrm{Z}_\lambda\to\mathrm{X}$ is an \emph{affine} morphism, $(f'_\lambda)_*(\sh{G}'_\lambda)=g^*((f_\lambda)_*(\sh{G}_\lambda))$, and in this case one has therefore simply $u'_\lambda=g^*(u_\lambda)$.

One can also interpret $u'_\lambda$ in the following manner: it suffices to know the value of $u'_\lambda(t')$, where $t'$ is a section of $\mathscr{F}'$ above an open $\mathrm{U}'$ of $\mathrm{X}'$, of the following particular type: $t'$ is the restriction to $\mathrm{U}'$ of the canonical image by $\Gamma(\mathrm{U},\mathscr{F})\to\Gamma(g^{-1}(\mathrm{U}),\mathscr{F}')$ of a section $t$ of $\mathscr{F}$ above an open $\mathrm{U}$ of $\mathrm{X}$ containing $g(\mathrm{U}')$ (these sections generate the $\sh{O}_{\mathrm{X}'}$-Module $\mathscr{F}'$ \sref[0]{0.3.7.1}).
Consider the section $u_\lambda(t)$ of $\sh{G}_\lambda$ above $f_\lambda^{-1}(\mathrm{U})$, and its canonical image $t''$ by $\Gamma(f_\lambda^{-1}(\mathrm{U}),\sh{G}_\lambda)\to\Gamma(g_\lambda'^{-1}(f_\lambda^{-1}(\mathrm{U})),\sh{G}'_\lambda)$; then $u'_\lambda(t')$ is the restriction of $t''$ to $f_\lambda'^{-1}(\mathrm{U}')$.
Consider in particular the case where $\mathrm{Z}_\lambda$ is the sub-prescheme induced on an open subset of $\mathrm{X}$, where $\sh{G}_\lambda=j_\lambda^*(\mathscr{F})$, with $j_\lambda:\mathrm{Z}_\lambda\to\mathrm{X}$ the canonical injection, and where $u_\lambda$ is the canonical homomorphism $\rho_\mathscr{F}:\mathscr{F}\to(j_\lambda)_*(j_\lambda^*(\mathscr{F}))=(j_\lambda)_*(\sh{G}_\lambda)$ \sref[0]{0.4.4.3.2} corresponding to $j_\lambda$.
Then $\mathrm{Z}'_\lambda$ is induced on an open subset of $\mathrm{X}'$, and the preceding interpretation shows that $u'_\lambda$ is none other than the canonical homomorphism $\rho_{\mathscr{F}'}:\mathscr{F}'\to(j'_\lambda)_*((j'_\lambda)^*(\mathscr{F}'))=(j'_\lambda)_*(\sh{G}'_\lambda)$ corresponding to the canonical injection $j'_\lambda:\mathrm{Z}'_\lambda\to\mathrm{X}'$.
\end{env}

\begin{env}[11.9.6]
\label{IV.11.9.6}
Under the conditions of \sref{IV.11.9.5}, suppose that $\mathrm{X}$ and $\mathrm{X}'$ are both affine, and that one wishes to prove that every section $t'$ of $\mathscr{F}'$ above $\mathrm{X}'$ whose images by all the $u'_\lambda$ are zero is itself zero.
Then, one can also restrict to the case where $\mathrm{Z}_\lambda=\mathrm{X}$ for every $\lambda\in\mathrm{L}$.
Indeed, by \sref{IV.11.9.3} and \sref{IV.11.9.5}, if one sets $\mathrm{Y}'_{\lambda\mu}=g_\lambda'^{-1}(\mathrm{Y}_{\lambda\mu})$, the homomorphism $v'_{\lambda\mu}$ deduced from $v_{\lambda\mu}$ by the change of base $g$ is none other than the canonical homomorphism
\oldpage[IV-3]{163}
$\rho_{\sh{G}'_\lambda}:\sh{G}'_\lambda\to(j'_{\lambda\mu})_*(j_{\lambda\mu}'^*(\sh{G}'_\lambda))$ corresponding to the canonical injection $j'_{\lambda\mu}:\mathrm{Y}'_{\lambda\mu}\to\mathrm{Z}'_\lambda$, as we saw in \sref{IV.11.9.5}.
The assertion follows then from the reasonings of \sref{IV.11.9.2} and \sref{IV.11.9.3}, $\mathrm{Y}_{\lambda\mu}$ and $v_{\lambda\mu}$ being replaced by $\mathrm{Y}'_{\lambda\mu}$ and $v'_{\lambda\mu}$.

One has an analogous reduction when one wishes to prove that the family $(u'_\lambda)$ is \emph{separating} ($\mathrm{X}$ and $\mathrm{X}'$ being affine): this follows again from \sref{IV.11.9.2} and \sref{IV.11.9.3}.
\end{env}

\begin{proposition}[11.9.7]
\label{IV.11.9.7}
With the notations of \sref{IV.11.9.1} and \sref{IV.11.9.5}, suppose $\mathrm{X}=\operatorname{Spec}(\mathrm{A})$ and $\mathrm{X}'=\operatorname{Spec}(\mathrm{A}')$ affine, and moreover that $\mathrm{A}'$ is a projective $\mathrm{A}$-module.
Then, if the family $(u_\lambda)$ is separating, every section $t'$ of $\mathscr{F}'$ above $\mathrm{X}'$ whose images by all the $u'_\lambda$ are zero, is itself zero.
\end{proposition}

\begin{proof}
One has seen in \sref{IV.11.9.6} that one can restrict to the case where all the $\mathrm{Z}_\lambda$ are \emph{equal to} $\mathrm{X}$.
The proposition is then a consequence of the following lemma:
\end{proof}

\begin{lemma}[11.9.7.1]
\label{IV.11.9.7.1}
Let $\mathrm{A}$ be a ring, $(\mathrm{M}_\lambda)_{\lambda\in\mathrm{L}}$ a family of $\mathrm{A}$-modules, $\mathrm{M}$ an $\mathrm{A}$-module, and for each $\lambda$, $u_\lambda:\mathrm{M}\to\mathrm{M}_\lambda$ a homomorphism.
Suppose that the intersection of the kernels of the $u_\lambda$ is reduced to zero.
Then, for every projective $\mathrm{A}$-module $\mathrm{N}$, the intersection of the kernels of the homomorphisms $u_\lambda\otimes 1:\mathrm{M}\otimes_\mathrm{A}\mathrm{N}\to\mathrm{M}_\lambda\otimes_\mathrm{A}\mathrm{N}$ is reduced to zero.
\end{lemma}

\begin{proof}
Indeed, $\mathrm{N}$ is a direct factor of a free $\mathrm{A}$-module $\mathrm{L}$, and it suffices evidently to prove that the intersection of the kernels of the homomorphisms $v_\lambda:\mathrm{M}\otimes_\mathrm{A}\mathrm{L}\to\mathrm{M}_\lambda\otimes_\mathrm{A}\mathrm{L}$ is reduced to zero, since $u_\lambda\otimes 1$ is the restriction of $v_\lambda$.
But the assertion follows then trivially from the hypothesis.
\end{proof}

\begin{remarks}[11.9.8]
\label{IV.11.9.8}
We do not know whether, under the hypotheses of proposition \sref{IV.11.9.7}, the family $(u'_\lambda)$ is \emph{separating}: indeed, by \sref{IV.11.9.4} one would have to prove that a section $t'$ of $\mathscr{F}'$ above an open subset $\mathrm{D}(h')\subset\mathrm{X}'$ (where $h'\in\mathrm{A}'$) such that all the $u'_\lambda(t')$ are zero is itself zero.
Now, one cannot apply proposition \sref{IV.11.9.7} to $\mathrm{D}(h')=\operatorname{Spec}(\mathrm{A}'_{h'})$, for the fact that $\mathrm{A}'$ is a projective $\mathrm{A}$-module (or even a free one), does not imply that $\mathrm{A}'_{h'}$ is a projective $\mathrm{A}$-module.
For example, one can take for $\mathrm{A}$ a discrete valuation ring, for $\mathrm{A}'$ a discrete valuation ring that is an $\mathrm{A}$-module free of rank $2$, and for $\mathrm{A}'_{h'}$ the fraction field of $\mathrm{A}'$.
\end{remarks}

One has however the following result:

\begin{corollary}[11.9.9]
\label{IV.11.9.9}
Let $\mathrm{X}$ be an Artinian prescheme, $g:\mathrm{X}'\to\mathrm{X}$ a flat morphism (note that these two conditions are satisfied if $\mathrm{X}$ is the spectrum of a \emph{field} and $g$ is an arbitrary morphism).
Then, with the notations of \sref{IV.11.9.1} and \sref{IV.11.9.5}, if the family $(u_\lambda)$ is separating, it is the same for the family $(u'_\lambda)$.
\end{corollary}

\begin{proof}
One can evidently restrict to the case where $\mathrm{X}=\operatorname{Spec}(\mathrm{A})$ is the spectrum of an Artinian local ring $\mathrm{A}$ \sref[I]{I.6.2.2}; one notes then that for every open affine $\mathrm{U}'=\operatorname{Spec}(\mathrm{A}')$ of $\mathrm{X}'$, $\mathrm{A}'$ is a flat $\mathrm{A}$-module, hence \emph{projective} ($\mathbf{0}_{\mathrm{III}}$,~10.1.3).
It suffices therefore to apply \sref{IV.11.9.7} to every open affine of $\mathrm{X}'$ to obtain the corollary.
\end{proof}

\begin{theorem}[11.9.10]
\label{IV.11.9.10}
Let $\mathrm{X}$ be a prescheme, $(u_\lambda)$ a family of homomorphisms \sref{IV.11.9.1.1}, $g:\mathrm{X}'\to\mathrm{X}$ a morphism, $(u'_\lambda)$ the family of homomorphisms deduced from $(u_\lambda)$ by the change of base $g$ \sref{IV.11.9.5}.
\begin{enumerate}[label=(\roman*)]
  \item If $g$ is a faithfully flat morphism and if the family $(u'_\lambda)$ is separating, then the family $(u_\lambda)$ is separating.
  \item Suppose that $g$ is a flat morphism and moreover that one of the two following conditions is satisfied:
  \begin{enumerate}[label=\emph{\alph*)}]
    \item $\mathrm{L}$ is finite and the $f_\lambda$ are quasi-compact.
    \item The morphism $g$ is locally of finite presentation.
  \end{enumerate}
  Then, if the family $(u_\lambda)$ is separating, it is the same for the family $(u'_\lambda)$.
\end{enumerate}
\end{theorem}

\begin{proof}
\oldpage[IV-3]{164}
\begin{enumerate}[label=(\roman*)]
  \item By virtue of \sref{IV.2.2.8}, it suffices to show that if a section $t$ of $\mathscr{F}$ above an open $\mathrm{U}$ of $\mathrm{X}$ belongs to the kernel of each of the $u_\lambda$, its image $t'$ by the canonical homomorphism $\Gamma(\rho):\Gamma(\mathrm{U},\mathscr{F})\to\Gamma(g^{-1}(\mathrm{U}),g^*(\mathscr{F}))=\Gamma(\mathrm{U},g_*(g^*(\mathscr{F})))$ is zero.
  Now the images of $t'$ by the $g^*(u_\lambda)$ are the images of $u_\lambda(t)$ by the homomorphism $\Gamma(\mathrm{U},(f_\lambda)_*(\sh{G}_\lambda))\to\Gamma(g^{-1}(\mathrm{U}),g^*((f_\lambda)_*(\sh{G}_\lambda)))$, whence they are zero, and \emph{a fortiori} $u'_\lambda(t')=0$ for every $\lambda$, whence $t'=0$ by hypothesis, which proves~(i).

  \item The question being local on $\mathrm{X}$ and $\mathrm{X}'$, one can restrict to the case where $\mathrm{X}=\operatorname{Spec}(\mathrm{A})$ and $\mathrm{X}'=\operatorname{Spec}(\mathrm{A}')$ are affine, $\mathrm{A}'$ being a flat $\mathrm{A}$-module, and it suffices to prove that every section $z'$ of $\mathscr{F}'$ above $\mathrm{X}'$ whose images by all the $u'_\lambda$ are zero is itself zero.
  We have seen in \sref{IV.11.9.6} that one can also restrict to the case where $\mathrm{Z}_\lambda=\mathrm{X}$ for every $\lambda\in\mathrm{L}$.
  Let us distinguish now the two cases.
  \medskip

  \emph{a)} If $\mathrm{L}$ is finite and the $f_\lambda$ are quasi-compact, we have seen in \sref{IV.11.9.3} that one can reduce further to the case where $\mathrm{Z}_\lambda=\mathrm{X}$ for every $\lambda\in\mathrm{L}$, and where moreover $\mathrm{L}$ is finite.
  It is then equivalent to say that the intersection of the kernels of the $u_\lambda:\mathscr{F}\to\sh{G}_\lambda$ is zero, or that the homomorphism $u=(u_\lambda):\mathscr{F}\to\sh{G}=\bigoplus_\lambda\sh{G}_\lambda$ is injective.
  Since $u'=g^*(u)$ is injective because $g$ is flat, the proposition is proved in this case.

  \emph{b)} Let $\mathscr{F}=\widetilde{\mathrm{M}}$, $\sh{G}_\lambda=\widetilde{\mathrm{N}}_\lambda$, and set $\mathrm{M}'=\mathrm{M}\otimes_\mathrm{A}\mathrm{A}'$, $\mathrm{N}'_\lambda=\mathrm{N}_\lambda\otimes_\mathrm{A}\mathrm{A}'$, so that $\mathscr{F}'=\widetilde{\mathrm{M}'}$, $\sh{G}'_\lambda=\widetilde{\mathrm{N}'_\lambda}$; by abuse of language, let us still denote by $u_\lambda$ the homomorphism $\mathrm{M}\to\mathrm{N}_\lambda$, and by $u'_\lambda$ the homomorphism $u_\lambda\otimes 1:\mathrm{M}'\to\mathrm{N}'_\lambda$.
  Let $z'\in\mathrm{M}'$ be an element such that $u'_\lambda(z')=0$ for every $\lambda$; it remains to prove that $z'=0$.
  Now, the hypothesis that $g$ is flat and of finite presentation implies, by \sref{IV.11.3.15}, that there exists a finite sequence $(s_i)_{1\leq i\leq n}$ of elements of $\mathrm{A}$, such that, if one sets $\mathfrak{J}_k=s_1\mathrm{A}+\ldots+s_k\mathrm{A}$, and $\mathrm{A}_i=\mathrm{A}_{s_i}/\mathfrak{J}_{i-1}\mathrm{A}_{s_i}$, the ring $\mathrm{A}'_i=\mathrm{A}'\otimes_\mathrm{A}\mathrm{A}_i$ is a \emph{free} $\mathrm{A}_i$-module for $1\leq i\leq n$, and $\mathfrak{J}_n=\mathrm{A}$.
  The proposition will be established if we prove the following assertion for $1\leq i\leq n$:
  \begin{enumerate}[label=$(*_i)$]
    \item \emph{There exists an integer $m_i>0$ such that $s_j^{m_i}z'=0$ for $j\leq i$.}
  \end{enumerate}
  Indeed, setting $k=m_n$ and noting that the $s_i^k$ ($1\leq i\leq n$) generate the unit ideal of $\mathrm{A}$, the assertion $(*_n)$ will show that $z'$, as a linear combination of the $s_i^k z'$, is zero.

  Let us prove $(*_i)$ by induction on $i$, the assertion being trivially true for $i=0$.
  Suppose $i>0$ and let $m$ be a common multiple of the $m_j$ for $j\leq i-1$.
  We note (for $1\leq h\leq n$) that if $\mathfrak{J}'_h$ is the ideal generated by the $s_j^m$ ($1\leq j\leq h$), then $\mathfrak{J}_h/\mathfrak{J}'_h$ is nilpotent.
  Replacing the $s_j$ by $s_j^m$ for $1\leq j\leq n$ therefore amounts to replacing, for $1\leq h\leq n$, $\mathrm{A}_h$ by
  $\mathrm{B}_h=\mathrm{A}_{s_h}/\mathfrak{J}'_{h-1}\mathrm{A}_{s_h}$, so that
  $\mathrm{A}_h=\mathrm{B}_h/((\mathfrak{J}_{h-1}/\mathfrak{J}'_{h-1})\mathrm{B}_h)$; since $\mathrm{A}'$ is a flat $\mathrm{A}$-module, it follows from ($\mathbf{0}_{\mathrm{III}}$,~10.1.2) that $\mathrm{B}'_h=\mathrm{A}'\otimes_\mathrm{A}\mathrm{B}_h$ is again a free $\mathrm{B}_h$-module.
  One can therefore
  \oldpage[IV-3]{165}
  replace all the $s_j$ ($1\leq j\leq n$) by $s_j^m$ without changing the properties of $\mathfrak{J}_h$ and $\mathrm{A}'_h$, and suppose $m=1$.
  Then $z'$, being annihilated by $\mathfrak{J}_{i-1}$, identifies with an element of $\operatorname{Hom}_{\mathrm{A}'}(\mathrm{A}'/\mathfrak{J}_{i-1}\mathrm{A}',\mathrm{M}')$; and since $\mathfrak{J}_{i-1}\mathrm{A}'$ is an ideal of finite type of $\mathrm{A}'$, this module of homomorphisms identifies in turn with $\operatorname{Hom}_\mathrm{A}(\mathrm{A}/\mathfrak{J}_{i-1},\mathrm{M})\otimes_\mathrm{A}\mathrm{A}'$ ($\mathbf{0}_{\mathrm{I}}$,~6.2.2).
  Let $v_\lambda=\operatorname{Hom}(1,u_\lambda):\operatorname{Hom}_\mathrm{A}(\mathrm{A}/\mathfrak{J}_{i-1},\mathrm{M})\to\operatorname{Hom}_\mathrm{A}(\mathrm{A}/\mathfrak{J}_{i-1},\mathrm{N}_\lambda)$, the homomorphism deduced from $u_\lambda$; the family $(v_\lambda)$ is also \emph{separating}.
  Indeed, for every $t\in\mathrm{A}$,
  $(\operatorname{Hom}_\mathrm{A}(\mathrm{A}/\mathfrak{J}_{i-1},\mathrm{M}))_t
  =\operatorname{Hom}_{\mathrm{A}_t}(\mathrm{A}_t/(\mathfrak{J}_{i-1})_t,\mathrm{M}_t)$,
  and likewise with $\mathrm{M}$ replaced by $\mathrm{N}_\lambda$, since the ideal $\mathfrak{J}_{i-1}$ is of finite type ($\mathbf{0}_{\mathrm{I}}$,~1.3.5); as by hypothesis the intersection of the kernels of the $(u_\lambda)_t$ is zero, the same holds for the intersection of the kernels of the $(v_\lambda)_t=\operatorname{Hom}(1,(u_\lambda)_t)$, whence our assertion by \sref{IV.11.9.4}.
  Replacing $\mathrm{A}$ by $\mathrm{A}/\mathfrak{J}_{i-1}$, $\mathrm{M}$ by $\operatorname{Hom}_\mathrm{A}(\mathrm{A}/\mathfrak{J}_{i-1},\mathrm{M})$, $\mathrm{N}_\lambda$ by $\operatorname{Hom}_\mathrm{A}(\mathrm{A}/\mathfrak{J}_{i-1},\mathrm{N}_\lambda)$ (which are $(\mathrm{A}/\mathfrak{J}_{i-1})$-modules), $u_\lambda$ by $v_\lambda$, and finally $\mathrm{A}'$ by $\mathrm{A}'/\mathfrak{J}_{i-1}\mathrm{A}'$, one sees that one can reduce to the case where, in the initial situation, the element $s=s_i\in\mathrm{A}$ is such that $\mathrm{A}'_s$ is a free $\mathrm{A}_s$-module.
  Now the family $(u_\lambda)_s:\mathrm{M}_s\to(\mathrm{N}_\lambda)_s$ is separating by hypothesis; since $(u'_\lambda)_s(z'/1)=0$, it follows from \sref{IV.11.9.7} that $z'/1=0$ in $\mathrm{M}'_s$, which signifies that there exists an integer $r$ such that $s^r z'=0$ in $\mathrm{M}'$, and this finishes the proof of $(*_i)$ by induction.
\end{enumerate}
\end{proof}

\begin{remark}[11.9.11]
\label{IV.11.9.11}
Let us restrict, for simplicity, to the case where $\mathrm{Z}_\lambda=\mathrm{X}$ for every $\lambda$.
It should be noted that if the family of homomorphisms $u_\lambda:\mathscr{F}\to\sh{G}_\lambda$ is separating, it does \emph{not necessarily} follow that, for every $x\in\mathrm{X}$, the intersection of the kernels of the homomorphisms $(u_\lambda)_x:\mathscr{F}_x\to(\sh{G}_\lambda)_x$ is reduced to $0$.
For example, let $\mathrm{X}$ be a locally Noetherian prescheme of Jacobson, of dimension $\geq 1$, and $\mathscr{F}$ a coherent $\sh{O}_\mathrm{X}$-Module; for every closed point $x$ of $\mathrm{X}$, and every integer $n\geq 0$, $\mathscr{F}/\mathfrak{m}_x^n\mathscr{F}$ is a coherent $\sh{O}_\mathrm{X}$-Module with support contained in $\{x\}$.
The family of canonical homomorphisms $\mathscr{F}\to\mathscr{F}/\mathfrak{m}_x^n\mathscr{F}$ (where $x$ runs through the set $\mathrm{X}_0$ of closed points of $\mathrm{X}$ and $n$ through the set of integers $\geq 0$) is \emph{separating}: indeed, if $t$ is a section of $\mathscr{F}$ above an open $\mathrm{U}$ of $\mathrm{X}$ whose images in all the $\Gamma(\mathrm{U},\mathscr{F}/\mathfrak{m}_x^n\mathscr{F})$ are zero, it follows at once that $t_x=0$ for every closed point $x\in\mathrm{U}$; as the set of closed points contained in $\mathrm{U}$ is very dense in $\mathrm{U}$, this implies $t=0$ \sref{IV.10.2.1}.
However, if one takes $\mathscr{F}=\sh{O}_\mathrm{X}$ and if $y\in\mathrm{X}$ is a \emph{non-closed} point, one has $(\sh{O}_\mathrm{X}/\mathfrak{m}_x^n\sh{O}_\mathrm{X})_y=0$ for every closed point $x$ of $\mathrm{X}$, but $\sh{O}_{\mathrm{X},y}\neq 0$, and the intersection of the kernels of the homomorphisms $\sh{O}_{\mathrm{X},y}\to(\sh{O}_\mathrm{X}/\mathfrak{m}_x^n\sh{O}_\mathrm{X})_y$ is equal to $\sh{O}_{\mathrm{X},y}$.
\end{remark}

\begin{lemma}[11.9.12]
\label{IV.11.9.12}
Let the notations be those of \sref{IV.11.9.1} and \sref{IV.11.9.5}, and suppose the family $(u_\lambda)$ to be separating; suppose in addition that $\mathrm{X}$ is an $\mathrm{S}$-prescheme, where $\mathrm{S}=\operatorname{Spec}(\mathrm{A})$ is an affine scheme, and that $\mathrm{X}'=\mathrm{X}\times_\mathrm{S}\mathrm{S}'$, where $\mathrm{S}'=\operatorname{Spec}(\mathrm{A}')$, $\mathrm{A}'$ being an $\mathrm{A}$-algebra; suppose finally that the $\sh{G}_\lambda$ are $\mathrm{S}$-flat.
Let $(\mathrm{A}'_\alpha)_{\alpha\in\mathrm{I}}$ be the filtered family of sub-$\mathrm{A}$-algebras of $\mathrm{A}'$ of finite type, and, for every $\alpha\in\mathrm{I}$, let $\mathrm{S}''_\alpha=\operatorname{Spec}(\mathrm{A}'_\alpha)$, $\mathrm{X}''_\alpha=\mathrm{X}\times_\mathrm{S}\mathrm{S}''_\alpha$, $\mathrm{Z}''_{\alpha\lambda}=\mathrm{Z}_\lambda\times_\mathrm{X}\mathrm{X}''_\alpha$, and let $f''_{\alpha\lambda}:\mathrm{Z}''_{\alpha\lambda}\to\mathrm{X}''_\alpha$, $\mathscr{F}''_\alpha$, $\sh{G}''_{\alpha\lambda}$, and $u''_{\alpha\lambda}$ be the morphisms, Modules, and homomorphisms of Modules deduced from $f_\lambda$, $\mathscr{F}$, $\sh{G}_\lambda$, and $u_\lambda$ by means of the change of base $\mathrm{X}''_\alpha\to\mathrm{X}$.
Then, if for every $\alpha\in\mathrm{I}$, the family $(u''_{\alpha\lambda})_{\lambda\in\mathrm{L}}$ is separating, it is the same for $(u'_\lambda)_{\lambda\in\mathrm{L}}$.
\end{lemma}

\begin{proof}
It suffices to prove that if $t'$ is a section of $\mathscr{F}'$ above an open affine $\mathrm{U}'$ of $\mathrm{X}'$ such that $u'_\lambda(t')=0$ for every $\lambda\in\mathrm{L}$, then $t'=0$.
If $h_\alpha:\mathrm{X}'\to\mathrm{X}''_\alpha$ is the canonical projection,
\oldpage[IV-3]{166}
it follows from \sref{IV.8.2.11} that there exists an index $\alpha$ and an open quasi-compact $\mathrm{U}''_\alpha\subset\mathrm{X}''_\alpha$ such that $\mathrm{U}'=h_\alpha^{-1}(\mathrm{U}''_\alpha)$; moreover, by virtue of \sref{IV.8.5.2},~(i), one can suppose that $t'$ is the canonical image of a section $t''_\alpha$ of $\mathscr{F}''_\alpha$ above $\mathrm{U}''_\alpha$.
Replacing, if necessary, $\mathrm{S}$, $\mathrm{X}$, $\mathrm{Z}_\lambda$, $f_\lambda$, $\mathscr{F}$, $\sh{G}_\lambda$, and $u_\lambda$ by $\mathrm{S}''_\alpha$, $\mathrm{U}''_\alpha$, $(f''_{\alpha\lambda})^{-1}(\mathrm{U}''_\alpha)$, and the corresponding restrictions of $\mathscr{F}''_\alpha$, $\sh{G}''_{\alpha\lambda}$, and $u''_{\alpha\lambda}$, we may therefore suppose that $\mathrm{U}'=\mathrm{X}'$, that $t'$ is the canonical image of a section $t$ of $\mathscr{F}$ above $\mathrm{X}$, and that the homomorphism $\mathrm{A}\to\mathrm{A}'$ is injective; equivalently, if $p:\mathrm{S}'\to\mathrm{S}$ is the corresponding morphism, the homomorphism $\sh{O}_\mathrm{S}\to p_*(\sh{O}_{\mathrm{S}'})$ occurring in the definition of $p$ is \emph{injective}.
It follows at once by virtue of the flatness of $\sh{G}_\lambda$ on $\mathrm{S}$ (and by reducing for example to the case where $\mathrm{Z}_\lambda$ is affine on $\mathrm{S}$ \sref[I]{I.1.6.3} and \sref[I]{I.1.6.5}) that the canonical homomorphism $\rho:\sh{G}_\lambda\to(g'_\lambda)_*(\sh{G}'_\lambda)$ is also \emph{injective}.
But the composite homomorphism
\[
  \Gamma(\mathrm{X},\mathscr{F})\xrightarrow{\Gamma(u_\lambda)}\Gamma(\mathrm{Z}_\lambda,\sh{G}_\lambda)\xrightarrow{\Gamma(\rho)}\Gamma(\mathrm{Z}'_\lambda,\sh{G}'_\lambda)
\]
is equal by definition to the composite homomorphism
\[
  \Gamma(\mathrm{X},\mathscr{F})\xrightarrow{\Gamma(\rho)}\Gamma(\mathrm{X}',\mathscr{F}')\xrightarrow{\Gamma(u'_\lambda)}\Gamma(\mathrm{Z}'_\lambda,\sh{G}'_\lambda)
\]
whence the image of $t$ by these composite homomorphisms is $u'_\lambda(t')=0$; by virtue of the injectivity of the homomorphism $\Gamma(\mathrm{Z}_\lambda,\sh{G}_\lambda)\xrightarrow{\Gamma(\rho)}\Gamma(\mathrm{Z}'_\lambda,\sh{G}'_\lambda)$, one concludes that $u_\lambda(t)=0$ for every $\lambda\in\mathrm{L}$, whence $t=0$ by hypothesis, and finally $t'=0$.
\end{proof}

\begin{proposition}[11.9.13]
\label{IV.11.9.13}
Let the notations be those of \sref{IV.11.9.1} and \sref{IV.11.9.5}, and suppose that $\mathrm{X}$ is a prescheme over a field $k$, and that, setting $\mathrm{S}=\operatorname{Spec}(k)$, one has $\mathrm{X}'=\mathrm{X}\times_\mathrm{S}\mathrm{S}'$, where $\mathrm{S}'$ is an arbitrary $k$-prescheme.
Then, if the family $(u_\lambda)_{\lambda\in\mathrm{L}}$ is separating, it is the same for $(u'_\lambda)$.
\end{proposition}

\begin{proof}
One can restrict to the case where $\mathrm{S}'=\operatorname{Spec}(\mathrm{A}')$ is affine.
If $\mathrm{A}'$ is a $k$-algebra of finite type, the morphism $g:\mathrm{X}'\to\mathrm{X}$ is flat and of finite presentation, and one is therefore under the hypotheses of \sref{IV.11.9.10},~(ii),~$b$).
In the general case, one regards $\mathrm{A}'$ as the inductive limit of its sub-$k$-algebras $\mathrm{A}'_\alpha$ of finite type and applies to each $\mathrm{A}'_\alpha$ the result of \sref{IV.11.9.10},~(ii),~$b$; one then concludes by means of Lemma \sref{IV.11.9.12}, since the $\sh{G}_\lambda$ are $\mathrm{S}$-flat.
\end{proof}

\begin{env}[11.9.14]
\label{IV.11.9.14}
Let us keep the notations of \sref{IV.11.9.1} and \sref{IV.11.9.5}, and suppose that $\mathrm{X}$ is an $\mathrm{S}$-prescheme.
If for every change of base $g:\mathrm{X}\times_\mathrm{S}\mathrm{S}'\to\mathrm{X}$, where $\mathrm{S}'\to\mathrm{S}$ is an arbitrary morphism, the corresponding family $(u'_\lambda)$ is separating, we say that the family $(u_\lambda)$ is \emph{universally separating relatively to} $\mathrm{S}$.
When the family $(u_\lambda)$ is reduced to a single element $u$, we say further that $u$ is \emph{universally injective}, relatively to $\mathrm{S}$.
It is clear then that for every morphism $h:\mathrm{S}'\to\mathrm{S}$, the family $(u'_\lambda)$ corresponding is universally separating relative to $\mathrm{S}'$; conversely, if $h$ is \emph{faithfully flat} and if $(u'_\lambda)$ is universally separating relative to $\mathrm{S}'$, then $(u_\lambda)$ is universally separating relative to $\mathrm{S}$, as follows immediately from \sref{IV.11.9.10},~(i) and from the fact that for every morphism $\mathrm{S}''\to\mathrm{S}$, the corresponding morphism $\mathrm{S}''\times_\mathrm{S}\mathrm{S}'\to\mathrm{S}''$ is faithfully flat.
\end{env}

\begin{proposition}[11.9.15]
\label{IV.11.9.15}
Let the notations be those of \sref{IV.11.9.1}, and suppose that $\mathrm{X}$ is an $\mathrm{S}$-prescheme, the $\sh{G}_\lambda$ being $\mathrm{S}$-flat.
Let $\mathrm{S}_0$ be a closed sub-prescheme of $\mathrm{S}$ defined by a nilpotent quasi-coherent Ideal $\sh{J}$ of $\sh{O}_\mathrm{S}$, such that the $(\sh{O}_\mathrm{S}/\sh{J})$-Modules $\sh{J}^k/\sh{J}^{k+1}$ are all locally free.
\oldpage[IV-3]{167}
Let $(u_{\lambda 0})$ be the family of homomorphisms obtained from $(u_\lambda)$ by the change of base $\mathrm{X}\leftarrow\mathrm{X}_0=\mathrm{X}\times_\mathrm{S}\mathrm{S}_0$.
Then, if the family $(u_{\lambda 0})$ is separating (resp.\ universally separating relative to $\mathrm{S}_0$), the family $(u_\lambda)$ is separating (resp.\ universally separating relative to $\mathrm{S}$).
\end{proposition}

\begin{proof}
Let us note that if $\mathrm{S}'\to\mathrm{S}$ is an arbitrary change of base and $\mathrm{S}'_0=\mathrm{S}'\times_\mathrm{S}\mathrm{S}_0$, then $\mathrm{S}'_0$ is a closed sub-prescheme of $\mathrm{S}'$ defined by a quasi-coherent nilpotent Ideal $\sh{J}'$ of $\sh{O}_{\mathrm{S}'}$, such that $\sh{J}'^k/\sh{J}'^{k+1}$ is a locally free $(\sh{O}_{\mathrm{S}'}/\sh{J}')$-Module for every $k$ \sref{IV.2.1.8},~(i); as moreover the $\sh{G}'_\lambda=\sh{G}_\lambda\otimes_{\sh{O}_\mathrm{S}}\sh{O}_{\mathrm{S}'}$ are $\mathrm{S}'$-flat, one sees that the assertion relative to universally separating families is a consequence of that relative to separating families.
To prove the latter, one can reduce by \sref{IV.11.9.3} to the case where $\mathrm{S}=\operatorname{Spec}(\mathrm{A})$ and $\mathrm{X}=\operatorname{Spec}(\mathrm{B})$ are affine, $\mathrm{Z}_\lambda=\mathrm{X}$ for every $\lambda$, $\mathscr{F}=\widetilde{\mathrm{M}}$, and $\sh{G}_\lambda=\widetilde{\mathrm{N}}_\lambda$, where $\mathrm{M}$ and the $\mathrm{N}_\lambda$ are $\mathrm{B}$-modules, the $\mathrm{N}_\lambda$ being flat $\mathrm{A}$-modules.
Moreover, the question being local on $\mathrm{X}$ and $\mathrm{S}$, it suffices to prove that if $t\in\mathrm{M}$ is such that $u_\lambda(t)=0$ for every $\lambda$, then $t=0$.
One has $\sh{J}=\widetilde{\mathfrak{J}}$, where $\mathfrak{J}$ is a nilpotent ideal of $\mathrm{A}$ such that the $\mathfrak{J}^k/\mathfrak{J}^{k+1}$ are free $(\mathrm{A}/\mathfrak{J})$-modules, and by hypothesis the homomorphisms $u_{\lambda 0}:\mathrm{M}/\mathfrak{J}\mathrm{M}\to\mathrm{N}_\lambda/\mathfrak{J}\mathrm{N}_\lambda$ form a separating family.
Suppose $\mathfrak{J}^{n+1}=0$ ($n$ an integer $\geq 0$) and let us reason by induction on $n$, the assertion being trivially true for $n=0$.
If $\bar{t}$ is the class of $t$ in $\mathrm{M}/\mathfrak{J}^n\mathrm{M}$, the class of $u_\lambda(t)$ in $\mathrm{N}_\lambda/\mathfrak{J}^n\mathrm{N}_\lambda$ is zero for every $\lambda$, whence, by the hypothesis of induction, $\bar{t}=0$, otherwise said $t\in\mathfrak{J}^n\mathrm{M}$.
Now $\mathfrak{J}^n=\mathfrak{J}^n/\mathfrak{J}^{n+1}$ is a free $(\mathrm{A}/\mathfrak{J})$-module; if $(e_\alpha)$ is a basis of this module, one can therefore write $t=\sum_\alpha e_\alpha t_\alpha^0$, with $t_\alpha^0\in\mathrm{M}/\mathfrak{J}\mathrm{M}$, zero for all but a finite number of indices.
On the other hand, since $\mathrm{N}_\lambda$ is a flat $\mathrm{A}$-module, $\mathfrak{J}^n\otimes_{(\mathrm{A}/\mathfrak{J})}(\mathrm{N}_\lambda/\mathfrak{J}\mathrm{N}_\lambda)$ identifies with $\mathfrak{J}^n\mathrm{N}_\lambda$ and one can therefore write $u_\lambda(t)=\sum_\alpha e_\alpha u_{\lambda 0}(t_\alpha^0)=\sum_\alpha e_\alpha\otimes u_{\lambda 0}(t_\alpha^0)$.
Since by hypothesis $u_\lambda(t)=0$, one deduces that $u_{\lambda 0}(t_\alpha^0)=0$ for every $\alpha$ and every $\lambda$; whence $t_\alpha^0=0$ for every $\alpha$, and hence $t=0$.
\end{proof}

\begin{theorem}[11.9.16]
\label{IV.11.9.16}
Let the notations be those of \sref{IV.11.9.1}, and suppose that $\mathrm{X}$ is a locally Noetherian $\mathrm{S}$-prescheme, $\mathscr{F}$ a coherent $\sh{O}_\mathrm{X}$-Module, and that the $\sh{G}_\lambda$ are $\mathrm{S}$-flat.
For every $s\in\mathrm{S}$, let $((u_\lambda)_s)_{\lambda\in\mathrm{L}}$ be the family obtained from $(u_\lambda)$ by the change of base $\mathrm{X}_s=\mathrm{X}\times_\mathrm{S}\operatorname{Spec}(k(s))\to\mathrm{X}$.
For the family $(u_\lambda)$ to be universally separating relative to $\mathrm{S}$, it is necessary and sufficient that for every $s\in\mathrm{S}$, the family $((u_\lambda)_s)$ be separating.
\end{theorem}

\begin{proof}
The necessity of the condition follows trivially from the definitions.
Conversely, suppose the condition satisfied, and let us prove first that the family $(u_\lambda)$ is separating.
One can by \sref{IV.11.9.3} reduce to the case where $\mathrm{S}=\operatorname{Spec}(\mathrm{A})$, $\mathrm{X}=\operatorname{Spec}(\mathrm{B})$ are affine, $\mathrm{Z}_\lambda=\mathrm{X}$ for every $\lambda$, $\mathscr{F}=\widetilde{\mathrm{M}}$, $\sh{G}_\lambda=\widetilde{\mathrm{N}}_\lambda$, where $\mathrm{B}$ is Noetherian, $\mathrm{M}$ is a $\mathrm{B}$-module of finite type, and the $\mathrm{N}_\lambda$ are flat $\mathrm{A}$-modules, and restrict to proving that, if $t\in\mathrm{M}$ is such that $u_\lambda(t)=0$ for every $\lambda$, then $t=0$.
To show that $t=0$, it suffices to prove that the image $t_\mathfrak{p}$ of $t$ in $\mathrm{M}_\mathfrak{p}$ is zero for every maximal ideal $\mathfrak{p}$ of $\mathrm{B}$ (Bourbaki, \emph{Alg.\ comm.}, chap.~II, \S~3, no.~3, cor.~1 of th.~1).
One may therefore restrict to proving that the intersection
\oldpage[IV-3]{168}
of the kernels of the $(u_\lambda)_\mathfrak{p}:\mathrm{M}_\mathfrak{p}\to(\mathrm{N}_\lambda)_\mathfrak{p}$ deduced from $(u_\lambda)$ by the change of base $\operatorname{Spec}(\mathrm{B}_\mathfrak{p})\to\operatorname{Spec}(\mathrm{B})$ is reduced to $0$.
Otherwise stated, one is brought back to the case where $\mathrm{B}$ is a local Noetherian ring; and, on considering the prime ideal of $\mathrm{A}$ which is the inverse image of the maximal ideal of $\mathrm{B}$, one may also suppose that $\mathrm{A}$ is a local ring with maximal ideal $\mathfrak{m}$.
Then, since $\mathfrak{m}\mathrm{B}$ is contained in the maximal ideal of $\mathrm{B}$, and $\mathrm{M}$ is a $\mathrm{B}$-module of finite type, the intersection of the $\mathfrak{m}^{n}\mathrm{M}$ is reduced to $0$ \sref[0]{0.7.3.5}, whence it suffices to prove that for every $n$, the image of $t$ in $\mathrm{M}/\mathfrak{m}^{n+1}\mathrm{M}$ is zero.
It suffices therefore to prove that the family deduced from $(u_\lambda)$ by the change of base $\operatorname{Spec}(\mathrm{B}/\mathfrak{m}^{n+1}\mathrm{B})\to\operatorname{Spec}(\mathrm{B})$ is separating, which signifies that one can restrict to the case where $\mathrm{A}$ is a local ring whose maximal ideal $\mathfrak{m}$ is nilpotent.
But then the $\mathfrak{m}^k/\mathfrak{m}^{k+1}$ are $(\mathrm{A}/\mathfrak{m})$-modules free, and, by virtue of the hypothesis on the $(u_\lambda)_s$, one is precisely in the conditions of application of \sref{IV.11.9.15}, whence the conclusion announced.

Let now $h:\mathrm{S}'\to\mathrm{S}$ be a change-of-base morphism, and let $(u'_\lambda)$ be the family obtained from $(u_\lambda)$ by the change of base $h':\mathrm{X}\times_\mathrm{S}\mathrm{S}'\to\mathrm{X}$; let us prove that $(u'_\lambda)$ is also separating.
Suppose first that $h$ is \emph{locally of finite type}; then the same is true of $h'$, and hence $\mathrm{X}'=\mathrm{X}\times_\mathrm{S}\mathrm{S}'$ is locally Noetherian. Moreover, if $s'\in\mathrm{S}'$ lies above $s\in\mathrm{S}$, it follows from \sref{IV.11.9.13}, applied to $\mathrm{X}_s$ and to $\mathrm{X}_{s'}=\mathrm{X}_s\otimes_{k(s)}k(s')$, that the family $((u'_\lambda)_{s'})$ is separating for every $s'\in\mathrm{S}'$; one may therefore conclude from the first part of the proof that $(u'_\lambda)$ is separating in this case.

Finally, if $h$ is arbitrary, one can evidently restrict to the case where $\mathrm{S}=\operatorname{Spec}(\mathrm{A})$ and $\mathrm{S}'=\operatorname{Spec}(\mathrm{A}')$ are affine, and consider $\mathrm{A}'$ as an inductive limit of its sub-$\mathrm{A}$-algebras of finite type.
As the $\sh{G}_\lambda$ are $\mathrm{S}$-flat, it suffices to apply what precedes and the lemma \sref{IV.11.9.12} to complete the proof.
\end{proof}

\begin{proposition}[11.9.17]
\label{IV.11.9.17}
Let $f:\mathrm{X}\to\mathrm{S}$ be a morphism locally of finite presentation, $\mathscr{F}$ a quasi-coherent $\sh{O}_\mathrm{X}$-Module of finite presentation and $f$-flat, $\mathrm{U}$ an open of $\mathrm{X}$, $j:\mathrm{U}\to\mathrm{X}$ the canonical injection, $u:\mathscr{F}\to j_*(j^*(\mathscr{F}))$ the canonical homomorphism \sref[0]{0.4.4.3.2}.
For every $s\in\mathrm{S}$, let $\mathrm{X}_s$ be the fibre $\mathrm{X}\times_\mathrm{S}\operatorname{Spec}(k(s))$, $\mathrm{U}_s$ the open $\mathrm{U}\cap\mathrm{X}_s$ of $\mathrm{X}_s$, $\mathscr{F}_s=\mathscr{F}\otimes_{\sh{O}_\mathrm{S}}k(s)$, $j_s:\mathrm{U}_s\to\mathrm{X}_s$ the canonical injection, $u_s:\mathscr{F}_s\to(j_s)_*((j_s)^*(\mathscr{F}_s))$ the canonical homomorphism.
For $u$ to be universally injective relative to $\mathrm{S}$, it is necessary and sufficient that $u_s$ be injective for every $s\in\mathrm{S}$.
\end{proposition}

\begin{proof}
There is nothing more than to prove the sufficiency of the condition.
When $\mathrm{X}$ is locally Noetherian, the proposition is an immediate corollary of \sref{IV.11.9.16}.
We shall now reduce to this case in two stages, restricting, as we clearly may, to the case where $\mathrm{S}=\operatorname{Spec}(\mathrm{A})$ and $\mathrm{X}$ are affine.
\medskip

\emph{A) Case where $\mathrm{U}$ is quasi-compact.} --- We will use the following lemma:

\begin{lemma}[11.9.17.1]
\label{IV.11.9.17.1}
Under the general hypotheses of \sref{IV.11.9.17}, and supposing in addition $\mathrm{S}$ and $\mathrm{X}$ affine and $\mathrm{U}$ quasi-compact, the set $\mathrm{E}$ of the $s\in\mathrm{S}$ such that $u_s$ is injective is constructible.
\end{lemma}

\begin{proof}
Indeed, the fibres $\mathrm{X}_s$ are locally Noetherian preschemes, whence $\mathrm{E}$ can also, by virtue of \sref{IV.5.10.2}, be defined as the set of $s\in\mathrm{S}$ such that $\operatorname{Ass}(\mathscr{F}_s)\subset\mathrm{U}_s$.
Note moreover that $\mathrm{U}$, being quasi-compact in an affine scheme, is constructible. Thus verification of condition \sref{IV.9.2.1},~(i) follows immediately from \sref{IV.4.2.7}; on the other hand,
\oldpage[IV-3]{169}
verification of \sref{IV.9.2.1},~(ii) is carried out easily by using the study of associated prime cycles in a neighbourhood of the generic fibre \sref{IV.9.8.3}, as well as \sref{IV.9.5.2} and \sref{IV.9.5.3}.
\end{proof}

This lemma being established, one can, by virtue of \sref{IV.8.9.1} and \sref{IV.8.2.11}, suppose that there exists a Noetherian sub-ring $\mathrm{A}_0$ of $\mathrm{A}$, a prescheme $\mathrm{X}_0$ of finite type over $\mathrm{S}_0=\operatorname{Spec}(\mathrm{A}_0)$, an open $\mathrm{U}_0$ of $\mathrm{X}_0$, and a coherent $\sh{O}_{\mathrm{X}_0}$-Module $\mathscr{F}_0$, such that $\mathrm{X}=\mathrm{X}_0\times_{\mathrm{S}_0}\mathrm{S}$ and, if $p_0:\mathrm{X}\to\mathrm{X}_0$ is the canonical projection, $\mathrm{U}=p_0^{-1}(\mathrm{U}_0)$ and $\mathscr{F}=p_0^*(\mathscr{F}_0)$.
Let $(\mathrm{A}_\alpha)$ be the filtered family of sub-rings of $\mathrm{A}$ that are $\mathrm{A}_0$-algebras of finite type, and set $\mathrm{S}_\alpha=\operatorname{Spec}(\mathrm{A}_\alpha)$, $\mathrm{X}_\alpha=\mathrm{X}_0\times_{\mathrm{S}_0}\mathrm{S}_\alpha$, $\mathrm{U}_\alpha=\mathrm{U}_0\times_{\mathrm{S}_0}\mathrm{S}_\alpha$, and $\mathscr{F}_\alpha=\mathscr{F}_0\otimes_{\mathrm{S}_0}\mathrm{S}_\alpha$; let $u_\alpha$ be the canonical homomorphism relative to $\mathscr{F}_\alpha$ and $\mathrm{U}_\alpha$, defined as in \sref{IV.11.9.17}.
For every $s\in\mathrm{S}$, the hypothesis that $u_s$ is injective implies that $(u_\alpha)_{s_\alpha}$ is injective \sref{IV.11.9.10},~(i), where $s_\alpha=q_\alpha(s)$ is the image of $s$ by the morphism $q_\alpha:\mathrm{S}\to\mathrm{S}_\alpha$.
If $\mathrm{E}_\alpha$ is the set of $s_\alpha\in\mathrm{S}_\alpha$ such that $(u_\alpha)_{s_\alpha}$ is injective, one therefore has $\mathrm{S}=q_\alpha^{-1}(\mathrm{E}_\alpha)$, and the $\mathrm{E}_\alpha$ form a projective system of sets.
But the lemma \sref{IV.11.9.17.1}, applied to $u_\alpha$, shows that $\mathrm{E}_\alpha$ is constructible in $\mathrm{S}_\alpha$; one therefore deduces from \sref{IV.8.3.4} that there exists an index $\alpha$ such that $\mathrm{E}_\alpha=\mathrm{S}_\alpha$.
But then, since $\mathrm{X}_\alpha$ is Noetherian, $u_\alpha$ is universally injective by virtue of \sref{IV.11.9.16}, whence $u$ is also universally injective.
\medskip

\emph{B) General case.} --- The open $\mathrm{U}$ is a filtered increasing union of quasi-compact opens $\mathrm{U}_\lambda$; if $j_\lambda:\mathrm{U}_\lambda\to\mathrm{X}$ is the canonical injection and $u_\lambda:\mathscr{F}\to(j_\lambda)_*(j_\lambda^*(\mathscr{F}))$ the canonical homomorphism corresponding, it follows from the lemma \sref{IV.11.9.17.1} applied to $u_\lambda$, that the set $\mathrm{E}_\lambda$ of the $s\in\mathrm{S}$ such that $(u_\lambda)_s$ is injective is constructible.
On the other hand, for every $s\in\mathrm{S}$, to say that $u_s$ (resp.\ $(u_\lambda)_s$) is injective signifies that $\operatorname{Ass}(\mathscr{F}_s)\subset\mathrm{U}\cap\mathrm{X}_s$ (resp.\ $\operatorname{Ass}(\mathscr{F}_s)\subset\mathrm{U}_\lambda\cap\mathrm{X}_s$) \sref{IV.5.10.2}.
As $\operatorname{Ass}(\mathscr{F}_s)$ is finite \sref{IV.3.1.6}, to say that $u_s$ is injective signifies therefore that there exists $\lambda$ such that $s\in\mathrm{E}_\lambda$.
The hypothesis therefore signifies that $\mathrm{S}=\bigcup_\lambda\mathrm{E}_\lambda$.
By virtue of \sref{IV.1.9.9}, there exists an index $\lambda$ such that $\mathrm{S}=\mathrm{E}_\lambda$, whence one concludes by the first part of the reasoning that $u_\lambda$ is universally injective.
It follows then from the factorisation of $u_\lambda$:
\[
  \mathscr{F}\xrightarrow{u}\,j_*(\mathscr{F}|\mathrm{U})\to(j_\lambda)_*(\mathscr{F}|\mathrm{U}_\lambda)
\]
that $u$ is also universally injective.
\end{proof}

\begin{remarks}[11.9.18]
\label{IV.11.9.18}
Let $\mathrm{X}$ be a prescheme, $\mathscr{F}$, $\sh{G}$ two quasi-coherent $\sh{O}_\mathrm{X}$-Modules of finite presentation, $\sh{G}$ being moreover supposed \emph{locally free}, $u:\mathscr{F}\to\sh{G}$ a homomorphism.
The following conditions are equivalent:
\begin{enumerate}[label=\emph{\alph*)}]
  \item for every morphism $g:\mathrm{X}'\to\mathrm{X}$, the homomorphism $g^*(u):g^*(\mathscr{F})\to g^*(\sh{G})$ is injective;
  \item for every $x\in\mathrm{X}$, the homomorphism $u\otimes 1_{k(x)}:\mathscr{F}\otimes_{\sh{O}_{\mathrm{X}}}k(x)\to\sh{G}\otimes_{\sh{O}_\mathrm{X}}k(x)$ is injective;
  \item for every $x\in\mathrm{X}$, there exists a neighbourhood $\mathrm{U}$ of $x$ such that $u|\mathrm{U}:\mathscr{F}|\mathrm{U}\to\sh{G}|\mathrm{U}$ is an isomorphism of $\mathscr{F}|\mathrm{U}$ onto a direct factor of $\sh{G}|\mathrm{U}$.
\end{enumerate}

It is clear that $c)$ implies $a)$, and $a)$ implies $b)$.
The fact that $b)$ implies $c)$ follows from \sref[0]{0.19.1.12}, \sref[0]{0.5.2.5} and \sref[0]{0.5.5.5}.

When the preceding equivalent conditions are satisfied, one says that $u$ is \emph{universally injective}.
\end{remarks}

\oldpage[IV-3]{170}
\subsection{Schematically dominant families of morphisms and families of schematically dense sub\mbox{-}preschemes}
\label{subsection:IV.11.10}

\begin{proposition}[11.10.1]
\label{IV.11.10.1}
Let $\mathrm{X}$ be a prescheme, $(Z_\lambda)_{\lambda\in\mathrm{L}}$ a family of preschemes, and for every $\lambda\in\mathrm{L}$, let $f_\lambda=(\psi_\lambda,\theta_\lambda):\mathrm{Z}_\lambda\to\mathrm{X}$ be a morphism.
The following conditions are equivalent:
\begin{enumerate}[label=\emph{\alph*)}]
  \item The family of homomorphisms $\theta_\lambda:\sh{O}_\mathrm{X}\to(f_\lambda)_*(\sh{O}_{\mathrm{Z}_\lambda})$ is separating (in other words \sref{IV.11.9.1}, the intersection of the kernels of the $\theta_\lambda$ is \emph{zero}).
  \item For every open $\mathrm{U}$ of $\mathrm{X}$, every section $t$ of $\sh{O}_\mathrm{X}$ above $\mathrm{U}$ whose images by all the canonical homomorphisms
\end{enumerate}
\[
  (\theta_\lambda)_\mathrm{U}:\Gamma(\mathrm{U},\sh{O}_\mathrm{X})\to\Gamma(f_\lambda^{-1}(\mathrm{U}),\sh{O}_{\mathrm{Z}_\lambda})
\label{IV.11.10.1.1}
\]
are zero, is itself zero.
\begin{enumerate}[label=\emph{\alph*)}]
\setcounter{enumi}{2}
  \item For every open $\mathrm{U}$ of $\mathrm{X}$ and every closed sub-prescheme $\mathrm{Y}$ of $\mathrm{U}$ such that for every $\lambda\in\mathrm{L}$, there exists a factorisation
\end{enumerate}
\[
  f_\lambda^{-1}(\mathrm{U})\xrightarrow{g_\lambda}\mathrm{Y}\xrightarrow{j}\mathrm{U}
\label{IV.11.10.1.2}
\]
of the restriction $f_\lambda^{-1}(\mathrm{U})\to\mathrm{U}$ of $f_\lambda$ (where $j$ is the canonical injection), one has $\mathrm{Y}=\mathrm{U}$.

\emph{If moreover $\mathrm{X}$ is an $\mathrm{S}$-prescheme, these conditions are also equivalent to the following:}
\begin{enumerate}[label=\emph{\alph*)}]
\setcounter{enumi}{3}
  \item For every separated morphism $g:\mathrm{X}'\to\mathrm{S}$ and every pair $u$, $v$ of $\mathrm{S}$-morphisms from an open $\mathrm{U}$ of $\mathrm{X}$ to $\mathrm{X}'$, such that for every $\lambda$, the composites of $u$ and $v$ with the restriction of $f_\lambda$ to $f_\lambda^{-1}(\mathrm{U})$ are equal, one has $u=v$.
\end{enumerate}
\end{proposition}

\begin{proof}
The equivalence of $a)$ and $b)$ follows from the definitions.
To see that $b)$ implies $c)$, it suffices to consider the quasi-coherent Ideal $\sh{J}$ of $\sh{O}_\mathrm{U}$ defining $\mathrm{Y}$, and to note that, for every open $\mathrm{V}\subset\mathrm{U}$, the hypothesis implies that the morphism $(\theta_\lambda)_\mathrm{V}$ factorises as
\[
  \Gamma(\mathrm{V},\sh{O}_\mathrm{U})\to\Gamma(\mathrm{V},\sh{O}_\mathrm{U}/\sh{J})\to\Gamma(f_\lambda^{-1}(\mathrm{V}),\sh{O}_{\mathrm{Z}_\lambda}).
\]
One concludes that every section $t$ of $\sh{J}$ above $\mathrm{V}$ has image $0$ in all the $\Gamma(f_\lambda^{-1}(\mathrm{V}),\sh{O}_{\mathrm{Z}_\lambda})$, whence, by virtue of $b)$, $\sh{J}=0$ and $\mathrm{Y}=\mathrm{U}$.
Conversely, if $c)$ is satisfied, to prove $b)$, it suffices to apply $c)$ to the closed sub-prescheme $\mathrm{Y}$ of $\mathrm{U}$ defined by the Ideal $t\sh{O}_\mathrm{U}$: the hypothesis that the images of $t$ by the $(\theta_\lambda)_\mathrm{U}$ are all zero implies that one has the factorisations \sref{IV.11.10.1.2} for every $\lambda$ \sref[I]{I.4.1.9}.
To prove that $c)$ implies $d)$, it suffices to apply $c)$ to the closed sub-prescheme $\mathrm{Y}$ of $\mathrm{U}$, the inverse image under $(u,v)_\mathrm{S}$ of the diagonal $\mathrm{X}'$ in $\mathrm{X}'\times_\mathrm{S}\mathrm{X}'$, and to use \sref[I]{I.4.4.1}.
Conversely, one deduces $b)$ from $d)$ by considering the $\mathrm{S}$-scheme $\mathrm{X}'=\mathrm{S}[\mathrm{T}]=\mathrm{S}\otimes_\mathbf{Z}\mathbf{Z}[\mathrm{T}]$ ($\mathrm{T}$ indeterminate) and noting that the sections of $\sh{O}_\mathrm{U}$ above $\mathrm{U}$ correspond bijectively to the $\mathrm{S}$-morphisms $\mathrm{U}\to\mathrm{S}[\mathrm{T}]$ \sref[I]{I.3.3.15}; to say that two sections of $\sh{O}_\mathrm{U}$ above $\mathrm{U}$ have the same images by all the $(\theta_\lambda)_\mathrm{U}$ is equivalent to saying that the composites of the two corresponding morphisms with all the morphisms $f_\lambda^{-1}(\mathrm{U})\to\mathrm{U}$ are equal.
\end{proof}

\oldpage[IV-3]{171}
\begin{definition}[11.10.2]
\label{IV.11.10.2}
When the equivalent conditions of \sref{IV.11.10.1} are satisfied, one says that the family $(f_\lambda)$ is \emph{schematically dominant}.
When the $f_\lambda$ are the canonical immersions in $\mathrm{X}$ of a family $(\mathrm{Z}_\lambda)$ of sub-preschemes of $\mathrm{X}$, one says also that the family $(\mathrm{Z}_\lambda)$ is \emph{schematically dense}.
\end{definition}

\begin{remarks}[11.10.3]
\label{IV.11.10.3}
\medskip\noindent
\begin{enumerate}[label=(\roman*)]
  \item The notion of a schematically dominant family on $\mathrm{X}$, as follows for example from a condition of the form $b)$ of \sref{IV.11.10.1}, is \emph{local} on $\mathrm{X}$: if $(\mathrm{W}_\alpha)$ is an open covering of $\mathrm{X}$, the family $(f_\lambda)$ is schematically dominant if and only if each of the families formed of the morphisms $f_\lambda^{-1}(\mathrm{W}_\alpha)\to\mathrm{W}_\alpha$, restrictions of the $f_\lambda$, is schematically dominant.
  \item If $\mathrm{Z}$ is the prescheme sum of the $\mathrm{Z}_\lambda$, $f:\mathrm{Z}\to\mathrm{X}$ the morphism coinciding with $f_\lambda$ on each $\mathrm{Z}_\lambda$, it amounts to the same to say that the family $(f_\lambda)$ is schematically dominant, or that the family reduced to the single element $f$ is.
  \item Let $\mathrm{M}$ be a second set of indices, and, for every $\lambda\in\mathrm{L}$, let $(g_{\lambda\mu})_{\mu\in\mathrm{M}}$ be a family of morphisms $g_{\lambda\mu}:\mathrm{Z}'_{\lambda\mu}\to\mathrm{Z}_\lambda$; if, for every $\lambda\in\mathrm{L}$, the family $(g_{\lambda\mu})_{\mu\in\mathrm{M}}$ is schematically dominant, then, it amounts to the same to say that the family $(f_\lambda)_{\lambda\in\mathrm{L}}$ is schematically dominant, or that the family $(f_\lambda\circ g_{\lambda\mu})_{(\lambda,\mu)\in\mathrm{L}\times\mathrm{M}}$ is schematically dominant \sref{IV.11.9.2}.
  \item Let $f:\mathrm{Z}\to\mathrm{X}$ be a morphism such that $f_*(\sh{O}_\mathrm{Z})$ is a quasi-coherent $\sh{O}_\mathrm{X}$-Module (for example a quasi-compact and quasi-separated morphism \sref{IV.1.7.4}).
  Then, to say that $f$ is schematically dominant signifies that the \emph{closed image} of $\mathrm{Z}$ \emph{by} $f$ \sref[I]{I.9.5.3} \emph{is identical to} $\mathrm{X}$.
\end{enumerate}
\end{remarks}

\begin{proposition}[11.10.4]
\label{IV.11.10.4}
If the family of morphisms $f_\lambda:\mathrm{Z}_\lambda\to\mathrm{X}$ is schematically dominant, the union of the $f_\lambda(\mathrm{Z}_\lambda)$ is dense in $\mathrm{X}$.
Conversely, if this union is dense in $\mathrm{X}$ and if moreover $\mathrm{X}$ is \emph{reduced}, the family $(f_\lambda)$ is schematically dominant.
\end{proposition}

\begin{proof}
The first assertion follows immediately from \sref{IV.11.10.1},~$b$)).
On the other hand, if $\mathrm{X}$ is reduced, and if the union of the $f_\lambda(\mathrm{Z}_\lambda)$ is dense in $\mathrm{X}$, then, if one has the factorisations \sref{IV.11.10.1.2} for every $\lambda\in\mathrm{L}$, one also has factorisations $(f_\lambda^{-1}(\mathrm{U}))_\mathrm{red}\xrightarrow{(g_\lambda)_\mathrm{red}}\mathrm{Y}_\mathrm{red}\xrightarrow{j_\mathrm{red}}\mathrm{U}$, and the hypothesis implies that the underlying space of $\mathrm{Y}$ is identical to $\mathrm{U}$, whence $\mathrm{Y}=\mathrm{Y}_\mathrm{red}=\mathrm{U}$ since $\mathrm{U}$ is reduced.
\end{proof}

The results of {\hyperlink{ega4.c04.subsec.11-9}{\textsection~11.9}} on separating families translate into results on schematically dominant families:

\begin{theorem}[11.10.5]
\label{IV.11.10.5}
Let $(f_\lambda)$ be a family of morphisms $f_\lambda:\mathrm{Z}_\lambda\to\mathrm{X}$, $g:\mathrm{X}'\to\mathrm{X}$ a morphism, and set $\mathrm{Z}'_\lambda=\mathrm{Z}_\lambda\times_\mathrm{X}\mathrm{X}'$, $f'_\lambda=(f_\lambda)_{(\mathrm{X}')}:\mathrm{Z}'_\lambda\to\mathrm{X}'$.
\begin{enumerate}[label=(\roman*)]
  \item If $g$ is faithfully flat and if the family $(f'_\lambda)$ is schematically dominant, then the family $(f_\lambda)$ is schematically dominant.
  \item Suppose that $g$ is a flat morphism, and moreover that one of the two following conditions is satisfied:
  \begin{enumerate}[label=\emph{\alph*)}]
    \item $\mathrm{L}$ is finite and the $f_\lambda$ are quasi-compact.
    \item The morphism $g$ is locally of finite presentation.
  \end{enumerate}
  Then, if the family $(f_\lambda)$ is schematically dominant, it is the same for the family $(f'_\lambda)$.
\end{enumerate}
\end{theorem}

\oldpage[IV-3]{172}
\begin{proposition}[11.10.6]
\label{IV.11.10.6}
Let the notations be those of \sref{IV.11.10.5}, and suppose that $\mathrm{X}$ is a prescheme over a field $k$, and that, setting $\mathrm{S}=\operatorname{Spec}(k)$, one has $\mathrm{X}'=\mathrm{X}\times_\mathrm{S}\mathrm{S}'$, where $\mathrm{S}'$ is an arbitrary $k$-prescheme.
Then, if the family $(f_\lambda)$ is schematically dominant, it is the same for $(f'_\lambda)$.
\end{proposition}

\begin{corollary}[11.10.7]
\label{IV.11.10.7}
Let $\mathrm{X}$ be a prescheme over a field $k$, $(Z_\lambda)_{\lambda\in\mathrm{L}}$ a family of $k$-preschemes, geometrically reduced over $k$, and for every $\lambda\in\mathrm{L}$, let $f_\lambda:\mathrm{Z}_\lambda\to\mathrm{X}$ be a $k$-morphism.
Let $\mathrm{Y}$ denote the reduced sub-prescheme of $\mathrm{X}$ having for underlying space the closure of $\bigcup_{\lambda\in\mathrm{L}}f_\lambda(\mathrm{Z}_\lambda)$.
Let $k'$ be an extension of $k$, $\mathrm{X}'$, $\mathrm{Z}'_\lambda$, $f'_\lambda:\mathrm{Z}'_\lambda\to\mathrm{X}'$ the preschemes deduced from $\mathrm{X}$, $\mathrm{Z}_\lambda$ and $f_\lambda$ by extension of the base to $k'$.
Then, if $\mathrm{Y}'$ is the reduced sub-prescheme of $\mathrm{X}'$ having for underlying space the closure of $\bigcup_{\lambda\in\mathrm{L}}f'_\lambda(\mathrm{Z}'_\lambda)$, one has $\mathrm{Y}'=\mathrm{Y}\otimes_k k'$.
In particular, $\mathrm{Y}$ is geometrically reduced over $k$.
\end{corollary}

\begin{proof}
Since the $\mathrm{Z}_\lambda$ are reduced, the morphisms $f_\lambda$ factorise as $\mathrm{Z}_\lambda\xrightarrow{g_\lambda}\mathrm{Y}\xrightarrow{j}\mathrm{X}$, where $j$ is the canonical injection \sref[I]{I.5.2.2}.
It follows then from \sref{IV.11.10.4} that $(g_\lambda)$ is a schematically dominant family.
Let $\mathrm{Y}'_1=\mathrm{Y}\otimes_k k'$, a closed sub-prescheme of $\mathrm{X}'$, and let $g'_\lambda$ be the morphism deduced from $g_\lambda$ by extension of the base to $k'$, so that $f'_\lambda$ factorises as $\mathrm{Z}'_\lambda\xrightarrow{g'_\lambda}\mathrm{Y}'_1\xrightarrow{j'}\mathrm{X}'$, where $j'$ is the canonical injection.
It follows from \sref{IV.11.10.6} that the family $(g'_\lambda)$ is schematically dominant.
But by hypothesis the $\mathrm{Z}'_\lambda$ are \emph{reduced}, whence \sref[I]{I.5.2.2} the $g'_\lambda$ factorise as $\mathrm{Z}'_\lambda\xrightarrow{h'_\lambda}\mathrm{Y}'\to\mathrm{Y}'_1$; one concludes therefore from \sref{IV.11.10.2} that $\mathrm{Y}'=\mathrm{Y}'_1$ and $h'_\lambda=g'_\lambda$, which establishes the corollary.
\end{proof}

\begin{definition}[11.10.8]
\label{IV.11.10.8}
Let the notations be those of \sref{IV.11.10.5}, and suppose that $\mathrm{X}$ is an $\mathrm{S}$-prescheme.
One says that the family $(f_\lambda)$ is \emph{universally schematically dominant} (relative to $\mathrm{S}$) if, for every change of base $\mathrm{S}'\to\mathrm{S}$, the family $(f'_\lambda)$ corresponding to the change of base $\mathrm{X}'=\mathrm{X}\times_\mathrm{S}\mathrm{S}'\to\mathrm{X}$ is schematically dominant.
\end{definition}

When $\mathrm{S}$ is the spectrum of a field, \sref{IV.11.10.6} then shows that a schematically dominant family is universally schematically dominant (relative to $\mathrm{S}$).

When the $f_\lambda$ are canonical immersions $\mathrm{Z}_\lambda\to\mathrm{X}$, one says also that the family of sub-preschemes $\mathrm{Z}_\lambda$ is \emph{universally schematically dense in} $\mathrm{X}$ (relative to $\mathrm{S}$) instead of saying that the family $(f_\lambda)$ is universally schematically dominant (relative to $\mathrm{S}$).

\begin{theorem}[11.10.9]
\label{IV.11.10.9}
Let the notations be those of \sref{IV.11.10.5}, and suppose that $\mathrm{X}$ is a locally Noetherian $\mathrm{S}$-prescheme and that the $\mathrm{Z}_\lambda$ are all $\mathrm{S}$-flat.
For every $s\in\mathrm{S}$, let $\mathrm{X}_s=\mathrm{X}\times_\mathrm{S}\operatorname{Spec}(k(s))$, $(\mathrm{Z}_\lambda)_s=\mathrm{Z}_\lambda\times_\mathrm{S}\operatorname{Spec}(k(s))$, $(f_\lambda)_s=f_\lambda\times 1:(\mathrm{Z}_\lambda)_s\to\mathrm{X}_s$.
For the family $(f_\lambda)$ to be universally schematically dominant relative to $\mathrm{S}$, it is necessary and sufficient that for every $s\in\mathrm{S}$, the family $((f_\lambda)_s)$ be schematically dominant.
\end{theorem}

\begin{proposition}[11.10.10]
\label{IV.11.10.10}
Let $f:\mathrm{X}\to\mathrm{S}$ be a flat locally of finite presentation morphism, $\mathrm{U}$ an open of $\mathrm{X}$.
For $\mathrm{U}$ to be universally schematically dense in $\mathrm{X}$ relative to $\mathrm{S}$, it is necessary and sufficient that, for every $s\in\mathrm{S}$, $\mathrm{U}_s=\mathrm{U}\cap\mathrm{X}_s$ be schematically dense in $\mathrm{X}_s$ (or, what amounts to the same, that one has $\operatorname{Ass}(\sh{O}_{\mathrm{X}_s})\subset\mathrm{U}_s$).
\end{proposition}
% END NATIVE INPUT source_aligned/ega4-11
% BEGIN NATIVE INPUT source_aligned/ega4-12
\oldpage[IV-3]{173}

\section{Study of fibres of flat morphisms of finite presentation}
\label{section:IV.12}

\setcounter{subsection}{-1}
\subsection{Introduction}
\label{subsection:IV.12.0}

\emph{We will use throughout this paragraph the general notation of} \sref{IV.9.4.1}.

\begin{env}[12.0.1]
\label{IV.12.0.1}
Given a morphism \emph{locally of finite presentation} $f:X\to Y$, we have seen \sref{IV.9.9} that for certain \emph{local} properties $\mathbf{P}$ of preschemes over fields or of Modules over these preschemes, the set $\mathrm{E}$ of $x\in\mathrm{X}$ such that the property $\mathbf{P}$ holds, for the \emph{fibre} $\mathrm{X}_{f(x)}$, at the point $x$ of this fibre, is \emph{locally constructible in} $\mathrm{X}$.
We now propose to show that, for the majority of these properties, if one further supposes that the morphism $f$ is \emph{flat}, then the set $\mathrm{E}$ is in fact \emph{open in} $\mathrm{X}$.
Similarly (\sref{IV.9.2} to \sref{IV.9.8}) we have shown that if $f$ is \emph{of finite presentation}, and if this time $\mathbf{P}$ denotes certain \emph{global} properties of preschemes over fields or of Modules over these preschemes, the set $\mathrm{F}$ of $y\in\mathrm{Y}$ such that the property $\mathbf{P}$ holds for the fibre $\mathrm{X}_y$ is \emph{locally constructible in} $\mathrm{Y}$.
We will show that if one further supposes that the morphism $f$ is \emph{proper} and \emph{flat}, $\mathrm{F}$ is in fact \emph{open in} $\mathrm{Y}$.
\end{env}

\begin{env}[12.0.2]
\label{IV.12.0.2}
The general method of proof of the properties in question involves three stages.
One first reduces to the case where $\mathrm{Y}$ is \emph{affine} and $\mathrm{X}$ of finite presentation; then:

\medskip
\begin{enumerate}[label=\emph{\Alph*)}]
  \item With the help of \sref{IV.8.9.1} (and possibly other results of \textsection\,8) and of \sref{IV.11.2.6}, one reduces to the case where $\mathrm{X}$ and $\mathrm{Y}$ are \emph{Noetherian}.

  \item One applies the results of \textsection\,9 recalled in \sref{IV.12.0.1} to show that $\mathrm{E}$ (resp.\ $\mathrm{F}$) is \emph{constructible}.

  \item To see that $\mathrm{E}$ is open, it suffices, by virtue of \sref[0]{0.9.2.5}, to show that if $x\in\mathrm{E}$, then every \emph{generization} $x'$ of $x$ also belongs to $\mathrm{E}$.
Using \sref[II]{II.7.1.7}, one sees that, since $\mathrm{Y}$ is \emph{Noetherian}, there exists a spectrum of a \emph{discrete valuation ring} $\mathrm{Y}_1$ and a morphism $h:\mathrm{Y}_1\to\mathrm{X}$ such that, if $y_1$ (resp.\ $y_1'$) is the closed point (resp.\ the generic point) of $\mathrm{Y}_1$, one has $h(y_1)=x$, $h(y_1')=x'$.
One then performs the base change $g=f\circ h:\mathrm{Y}_1\to\mathrm{Y}$; taking into account the results of \textsection\textsection\,4 and 6 on locally Noetherian preschemes over fields and the changes of base field, one is reduced to proving the assertion in question for a point $x_1$ of $\mathrm{X}_1=\mathrm{X}\times_{\mathrm{Y}}\mathrm{Y}_1$ above $y_1$ and for a generization $x_1'$ of $x_1$ above $y_1'$.
Since $\mathrm{Y}_1=\Spec(\mathrm{A})$, where $\mathrm{A}$ is a discrete valuation ring, with uniformiser $t$, one must in the end, for an $\mathrm{A}$-module $\mathrm{M}$, prove a property of $\mathrm{M}$, knowing that $\mathrm{M}/t\mathrm{M}$ has the same property and that $t$ is $\mathrm{M}$-\emph{regular} (which follows from the flatness hypothesis); one uses for this the results of \sref{IV.3.4} and of \sref{IV.5.12}.
One proceeds similarly for the set $\mathrm{F}\subset\mathrm{Y}$.
\end{enumerate}
\end{env}

\oldpage[IV-3]{174}

\subsection{Local properties of fibres of a flat morphism locally of finite presentation}
\label{subsection:IV.12.1}

\begin{theorem}[12.1.1]
\label{IV.12.1.1}
Let $f:X\to Y$ be a morphism locally of finite presentation, $\sh{F}$ an $\sh{O}_X$-Module $f$-flat and of finite presentation, $\Phi$ a finite subset of $\mathbf{Z}\cup\{\pm\infty\}$, $k$ an integer.
The following subsets of $X$ are open:
\begin{enumerate}[label=(\roman*)]
  \item The set of $x\in X$ such that the dimensions of the prime cycles associated to $\sh{F}_{f(x)}$ and containing $x$ are elements of $\Phi$.
  \item The set of $x\in X$ such that all the prime cycles associated to $\sh{F}_{f(x)}$ containing $x$ have dimension $k$.
  \item The set of $x\in X$ such that $x$ does not belong to any associated immersed prime cycle of $\sh{F}_{f(x)}$.
  \item The set of $x\in X$ such that $\sh{F}_{f(x)}$ is equidimensional at the point $x$ and possesses the property $(\mathrm{S}_k)$ at the point $x$ \sref{IV.5.7.2}.
  \item The set of $x\in X$ such that $\operatorname{coprof}((\sh{F}_{f(x)})_x)\leqslant k$ \sref[0]{0.16.4.9}.
  \item The set of $x\in X$ such that $\sh{F}_{f(x)}$ is a Cohen--Macaulay $\sh{O}_{X_{f(x)}}$-Module at the point $x$ \sref{IV.5.7.1}.
  \item The set of $x\in X$ such that $\sh{F}_{f(x)}$ is geometrically reduced at the point $x$ \sref{IV.4.6.22}.
  \item The set of $x\in X$ such that $\sh{F}_{f(x)}$ is geometrically punctually integral at the point $x$ \sref{IV.4.6.22}.
\end{enumerate}

The questions being local on $X$, one first reduces to the case where $X=\Spec(B)$ and $Y=\Spec(A)$ are affine, $f$ a morphism of finite presentation.
There is then a Noetherian subring $A_0$ of $A$, an $A_0$-prescheme of finite type $X_0$, and a coherent $\sh{O}_{X_0}$-Module $\sh{F}_0$ such that $\sh{F}_0\otimes_{A_0}A$ is isomorphic to $\sh{F}$ \sref{IV.8.9.1}; in addition, by virtue of \sref{IV.11.2.6}, one can suppose that $\sh{F}_0$ is $Y_0$-flat (with $Y_0=\Spec(A_0)$).
If $h:X\to X_0$ is the canonical projection, the set $E$ of points $x\in X$ where one of the properties (i) to (viii) is satisfied is equal to $h^{-1}(E_0)$, where $E_0$ is the set of $x_0\in X_0$ where the corresponding property for $\sh{F}_0$ is satisfied: this follows, for properties (i) to (iii), from \sref{IV.4.2.7}; for properties (iv) to (vi), from \sref{IV.6.7.1}; for properties (vii) and (viii), from \sref{IV.4.7.11}.
\end{theorem}

One is thus reduced to the case where $A$ and $B$ are Noetherian, $f$ of finite type, and $\sh{F}=\widetilde{M}$, where $M$ is a $B$-module of finite type.
In virtue of \sref{IV.9.9.2}, one knows that the set $E$ is \emph{constructible} for all the properties considered, and it remains in each case the step~C) of \sref{IV.12.0.2}, where one must prove that $E$ is \emph{stable under generization}.

\begin{env}[12.1.1.1]
\label{IV.12.1.1.1}
Let us begin with case~(iii), which is the simplest.
Let $x'\neq x$ be a generization of $x$ in $X$.
Set $y=f(x)$, $y'=f(x')$; there exists a spectrum of a discrete valuation ring $Y_1$ and a morphism $u:Y_1\to X$ such that, if $y_1$ and $y_1'$ are the closed point and the generic point of $Y_1$ respectively, one has $u(y_1)=x$ and $u(y_1')=x'$ \sref[II]{II.7.1.7}; set $g=f\circ u:Y_1\to Y$; if $X_1=X\times_Y Y_1$, there exists a $Y_1$-section $u_1:Y_1\to X_1$ such that $u=g_1\circ u_1$, where $g_1:X_1\to X$ is the canonical projection.
If $x_1=u_1(y_1)$, $x_1'=u_1(y_1')$, one has $g_1(x_1)=x$ and $g_1(x_1')=x'$, and $x_1'$ is a generization of $x_1$.
Applying again \sref{IV.4.2.7}, one sees
\oldpage[IV-3]{175}
that one is reduced to the case where $Y=\Spec(A)$ is the spectrum of a discrete valuation ring, $y$ being the closed point and $y'$ the generic point of $Y$.
If $t$ is a uniformiser of $A$, the hypothesis that $\sh{F}$ is $f$-flat implies that $t$ is $\sh{F}$-regular \sref[0]{0.6.3.4}.
By hypothesis, none of the associated immersed prime cycles of $\sh{F}/t\sh{F}$ contains $x$.
Then, it follows from \sref{IV.3.4.4} that $x$ does not belong to any associated immersed prime cycle of $\sh{F}$, and the same is true of every generization of $x$.
In particular, $x'$ does not belong to any associated immersed prime cycle of $\sh{F}$, nor \emph{a fortiori} to any of those associated to $\sh{F}_{y'}$ ($X_{y'}$ being a sub-prescheme inducing an open subset of $X$).
\end{env}

\begin{env}[12.1.1.2]
\label{IV.12.1.1.2}
Let us now consider cases~(iv) and~(v) ((vi) is only a particular case of~(v)).
One proceeds as above (using \sref{IV.6.7.1} in place of \sref{IV.4.2.7}) and one is reduced to the case where $Y$ is the spectrum of a discrete valuation ring~$A$.

The ring $C=\sh{O}_{X,x}$ is then the localisation of an $A$-algebra of finite type, hence is catenary \sref{IV.5.6.4}, and by hypothesis the $C$-module $M_x/tM_x$ satisfies $(\mathrm{S}_k)$ and is equidimensional (resp.\ one has $\operatorname{coprof}(M_x/tM_x)\leqslant k$); one deduces from \sref{IV.5.12.2} (resp.\ from \sref[0]{0.16.4.10}) that $M_x$ satisfies $(\mathrm{S}_k)$ and is equidimensional (resp.\ that $\operatorname{coprof}(M_x)\leqslant k$) since $t$ is $M_x$-regular and belongs to the maximal ideal of $C$.
Whence the conclusions since $x'$ is a generization of $x$ and $(\sh{F}_{y'})_{x'}=\sh{F}_{x'}$.
\end{env}

\begin{env}[12.1.1.3]
\label{IV.12.1.1.3}
Let us pass to cases~(vii) and~(viii).
One can evidently replace $Y$ by the integral sub-prescheme having $\overline{\{y'\}}$ as underlying space, and $X$ by $f^{-1}(\overline{\{y'\}})$, without changing the fibres in $y$ and $y'$, so one can suppose $A$ integral, with field of fractions $K=k(y')$.
One knows (\sref{IV.4.5.11} and \sref{IV.4.7.8}) that there exists a finite extension $K'$ of $K$ such that, for the $(\sh{O}_X\otimes_A K')$-Module $\sh{F}'=\sh{F}\otimes_A K'$, the associated prime cycles are geometrically irreducible and the geometric lengths of $\sh{F}'$ at the maximal points $z_i'$ of its support are respectively equal to the lengths of $\sh{F}'_{z_i'}$ at these points; it amounts therefore to the same to say that $\sh{F}$ is geometrically reduced (resp.\ geometrically punctually integral) at a point $x'$ of $X_{y'}=X\otimes_A K$, or to say that $\sh{F}'$ is \emph{reduced} (resp.\ \emph{integral}) at the points of $X\otimes_A K'$ lying above $x'$, taking into account \sref{IV.4.2.7}.
Let $A'$ be a finite $A$-algebra whose field of fractions is $K'$; it follows from \sref{IV.4.7.11} that at every point of $X\otimes_A A'$ above $x$, $\sh{F}\otimes_A A'$ has property~(vii) (resp.~(viii)).
Replacing $A$ by $A'$ and $X$ by $X\otimes_A A'$, one sees that one can reduce to proving that $\sh{F}$ is \emph{reduced} (resp.\ \emph{integral}) at every generization $x'$ of $x$.
One proceeds then as in \sref{IV.12.1.1.1}, and using this time \sref{IV.4.7.11}, in the case where $A$ is a discrete valuation ring, $y$ being the closed point, $y'$ the generic point of $Y$, and the uniformiser $t$ of $A$ being an $M$-\emph{regular} element.
Since by hypothesis $\sh{F}/t\sh{F}$ is reduced (resp.\ integral) at the point $x$, it follows from \sref{IV.3.4.6} (resp.\ \sref{IV.3.4.5}) that $\sh{F}$ is reduced (resp.\ integral) at the point $x$, hence also in a neighbourhood of $x$, and in particular at the point $x'$, which finishes the proof in cases~(vii) and~(viii).
\end{env}

\begin{env}[12.1.1.4]
\label{IV.12.1.1.4}
It remains to examine cases~(i) and~(ii).
Replacing $Y$ by the integral sub-prescheme having $\overline{\{y'\}}$ as underlying space, one can reduce to the case where $Y$ is irreducible and $y'$ is its generic point. Let $z'$ be a generic point of an associated prime cycle of $\sh{F}_{y'}$
\oldpage[IV-3]{176}
containing $x'$, and let $Z'$ be the closure of $z'$ in $X$, so that one has $x\in Z'$.
To treat case~(i), it will suffice to prove the

\begin{proposition}[12.1.1.5]
\label{IV.12.1.1.5}
Let $Y$ be a locally Noetherian prescheme, $f:X\to Y$ a morphism locally of finite type, $\sh{F}$ a coherent $\sh{O}_X$-Module and $f$-flat; for every $y\in Y$, set $\sh{F}_y=\sh{F}\otimes_{\sh{O}_Y}k(y)$.
Let $z'$ be a point of $X$, $y'=f(z')$, and suppose that $z'\in\operatorname{Ass}(\sh{F}_{y'})$.
Let $Z'$ \textnormal{(}resp.\ $Y'$\textnormal{)} be the reduced sub-prescheme of $X$ \textnormal{(}resp.\ $Y$\textnormal{)} having as underlying space the closure of $\{z'\}$ in $X$ \textnormal{(}resp.\ of $\{y'\}$ in $Y$\textnormal{)}.
Then, for every $y\in f(Z')$, the dimensions of all the irreducible components of $Z'_y$ are equal to $\dim(Z'_{y'})$ \textnormal{(}dimension of the closure of $z'$ in $X_{y'}$\textnormal{)} \textnormal{(}which we will express more precisely later \sref{IV.13.2.2} by saying that the restriction $Z'\to Y'$ of $f$ is an equidimensional morphism\textnormal{)}; furthermore, at every maximal point $z$ of $Z'_y$, one has $z\in\operatorname{Ass}(\sh{F}_y)$.
\end{proposition}

For this, we will now reduce to the case where $Y$ is the spectrum of a discrete valuation ring.
Let us take a spectrum of a discrete valuation ring $Y_1$ and a morphism $h:Y_1\to X$ such that $h(y_1)=z$, $h(y_1')=z'$, $y_1$ and $y_1'$ being the closed point and the generic point of $Y_1$ respectively \sref[II]{II.7.1.7}.
Set $g=f\circ h:Y_1\to Y$, $X_1=X\times_Y Y_1$, and let $f_1:X_1\to Y_1$, $g_1:X_1\to X$ be the canonical projections; there is a $Y_1$-section $h_1:Y_1\to X_1$ such that $h=g_1\circ h_1$ \sref[I]{I.3.3.14}.
If $Z_1'=Z'_{y'}\otimes_{k(y')}k(y_1')$, one knows \sref{IV.4.2.6} that the irreducible components of $Z_1'$ dominate $Z'_{y'}$; let $z_1'$ be the generic point of one of these components which contains $h_1(y_1')$, so that $f_1(z_1')=y_1'$ and $g_1(z_1')=z'$.
Since $h_1(y_1)$ is a specialisation of $h_1(y_1')$, it is \emph{a fortiori} a specialisation of $z_1'$; let $z_1$ be a generic point of one of the irreducible components of $\overline{\{z_1'\}}\cap(X_1)_{y_1}$ containing $h_1(y_1)$.
Then $g_1(z_1)$ is a specialisation of $g_1(z_1')=z'$ in $X$, and $z=h(y_1)=g_1(h_1(y_1))$ is a specialisation of $g_1(z_1)$; but since $f(g_1(z_1))=y$ and $z$ is a generic point of $Z'_y$, one necessarily has $z=g_1(z_1)$.

Suppose now that one has proved that, if one sets $\sh{F}_1=\sh{F}\otimes_{\sh{O}_Y}\sh{O}_{Y_1}$, then $z_1\in\operatorname{Ass}((\sh{F}_1)_{y_1})$; it will follow from \sref{IV.4.2.7} that $z\in\operatorname{Ass}(\sh{F}_y)$; furthermore, the dimensions of $\overline{\{z_1\}}\cap(X_1)_{y_1}$ and of $\overline{\{z\}}\cap X_y$ are equal, and likewise the dimensions of $\overline{\{z_1'\}}\cap(X_1)_{y_1'}$ and of $\overline{\{z'\}}\cap X_{y'}$ are equal \sref{IV.4.2.7}; one has therefore indeed reduced, as announced, to the case where $Y$ is the spectrum of a discrete valuation ring $A$, $y$ and $y'$ being its closed point and its generic point respectively.

Since $z'\in\operatorname{Ass}(\sh{F}_{y'})$, one has \emph{a fortiori} $z'\in\operatorname{Ass}(\sh{F})$.
If $t$ is a uniformiser of $A$, $t$ is $\sh{F}$-regular by flatness, so by \sref{IV.3.4.3} one has $z\in\operatorname{Ass}(\sh{F}/t\sh{F})=\operatorname{Ass}(\sh{F}_y)$.
As for the assertion relative to dimensions, it follows from \sref{IV.7.1.13}, applied to a closed sub-prescheme of $X$ having $Z'$ as underlying space.

Let us finally tackle case~(ii) of \sref{IV.12.1.1}.
With the same notation, it is necessary to prove that if all the prime cycles associated to $\sh{F}_y$ and containing $x$ have dimension $k$, the same is true of all the prime cycles associated to $\sh{F}_{y'}$ and containing $x'$.
Now, one has just seen that \emph{every} prime cycle associated to $\sh{F}_{y'}$ and containing $x'$ has the same dimension as \emph{one} of the prime cycles associated to $\sh{F}_y$ and containing $x$; this shows~(ii).
\end{env}

\begin{remarks}[12.1.2]
\label{IV.12.1.2}
\begin{enumerate}[label=(\roman*)]
  \item Under the conditions of \sref{IV.12.1.1.5} with $\sh{F}=\sh{O}_X$ (so that $f$ is \emph{flat}), one cannot in general assert that the restriction $Z'\to Y'$ of $f$ is a flat morphism nor even an \emph{open} morphism.
This is what the example \sref{IV.6.5.5},
\oldpage[IV-3]{177}
(ii) shows, where one takes for $Z'$ one of the two irreducible components of $X=\Spec(B)$.
It is immediate that this restriction morphism is not open at the points of $Z'$ above the ``double point'' of $Y'=Y$.

  \item Under the hypotheses of \sref{IV.12.1.1.5}, with $\sh{F}=\sh{O}_X$, one cannot weaken the condition ``$f$ is flat'' to ``$f$ is universally open'' (cf.\ \sref{IV.2.4.6}), as we will see further on with an example \sref{IV.14.4.10},~(i)).
\end{enumerate}
\end{remarks}

\begin{corollary}[12.1.3]
\label{IV.12.1.3}
Under the hypotheses of \sref{IV.12.1.1}, the function $x\mapsto\operatorname{coprof}((\sh{F}_{f(x)})_x)$ is upper semi-continuous in $X$ and the function $x\mapsto\operatorname{prof}_x^*(\sh{F}_{f(x)})$ \sref{IV.10.8.1} is lower semi-continuous in $X$.
\end{corollary}

The first assertion is nothing other than an equivalent formulation of \sref{IV.12.1.1},~(v).
To prove the second, one can first, by taking into account \sref{IV.10.8.7} and \sref{IV.10.8.8}, reduce as in \sref{IV.12.1.1} to the case where $Y$ is the spectrum of a discrete valuation ring $A$ of closed point $y$ and of generic point $y'$, $X=\Spec(B)$ being affine, $\sh{F}=\widetilde{M}$, where $M$ is a $B$-module of finite type.
Since $\sh{F}$ is $f$-flat, every irreducible component of $\operatorname{Supp}(\sh{F})$ which contains $x\in X_y$ dominates $Y$ \sref{IV.2.3.4}; in other words, its generic point $z'$ belongs to $X_{y'}$. The irreducible components $Z$ of $\operatorname{Supp}(\sh{F})$ which contain a generization $x'$ of $x$ belonging to $X_{y'}$ are therefore exactly those which contain $x$; furthermore, by \sref{IV.12.1.1.5}, one has $\dim(Z_{y'})=\dim(Z_y)$, whence, if $T=\operatorname{Supp}(\sh{F})$, $\dim_x(T_y)=\dim_{x'}(T_{y'})$.
Furthermore, for a uniformiser $t$ of $A$, one has seen in \sref{IV.12.1.1},~(v) that $\operatorname{coprof}(M_x)=\operatorname{coprof}(M_x/tM_x)$, and since on the other hand $\operatorname{coprof}(M_{x'})\leqslant\operatorname{coprof}(M_x)$ \sref{IV.6.11.5}, one sees that $\operatorname{coprof}((\sh{F}_{y'})_{x'})\leqslant\operatorname{coprof}((\sh{F}_y)_x)$; the relation $\operatorname{prof}_x^*(\sh{F}_y)\leqslant\operatorname{prof}_{x'}^*(\sh{F}_{y'})$ then follows from \sref{IV.10.8.7}.

\begin{remark}[12.1.4]
\label{IV.12.1.4}
One also deduces from \sref{IV.12.1.1},~(i) that the function $x\mapsto\dim_x(X_{f(x)})$ is upper semi-continuous in $X$, but we will see further on \sref{IV.13.1.3} that this property is true even without supposing $f$ flat.
\end{remark}

\begin{corollary}[12.1.5]
\label{IV.12.1.5}
Let $Y$ be a locally Noetherian prescheme, $f:X\to Y$ a morphism locally of finite type, $\sh{F}$ a coherent $\sh{O}_X$-Module which is $f$-flat, $\sh{G}$ a coherent $\sh{O}_Y$-Module, $x$ a point of $X$, $y=f(x)$.
Suppose that $\sh{G}$ has the property $(\mathrm{S}_k)$ at the point $y$, and that $\sh{F}_y$ has the property $(\mathrm{S}_k)$ at the point $x$ and is equidimensional at the point $x$.
Then:
\begin{enumerate}[label=(\roman*)]
  \item $\sh{F}\otimes_{\sh{O}_Y}\sh{G}$ has the property $(\mathrm{S}_k)$ at the point $x$.
  \item If furthermore $Y$ is locally immersible in a regular scheme, or is an excellent prescheme \sref{IV.7.8.5}, there exists a neighbourhood $V$ of $x$ in $X$ such that $\sh{F}\otimes_{\sh{O}_Y}\sh{G}$ has the property $(\mathrm{S}_k)$ in this neighbourhood.
\end{enumerate}
\end{corollary}

\begin{proof}
Indeed, by \sref{IV.12.1.1},~(iv), there exists an open neighbourhood $V$ of $x$ such that for every $x'\in V$, $\sh{F}_{f(x')}$ satisfies $(\mathrm{S}_k)$ at the point $x'$.
On the other hand, $\sh{G}$ satisfies $(\mathrm{S}_k)$ at all points of $\Spec(\sh{O}_y)$ \sref{IV.5.7.2}.
Replacing $Y$ by $Y'=\Spec(\sh{O}_y)$ and $X$ by $V\times_Y Y'$, one is reduced (taking into account \sref[I]{I.3.6.5}) to the case where $\sh{G}$ satisfies $(\mathrm{S}_k)$ throughout $Y$ and where, for every $y\in f(X)$, $\sh{F}_y$ has the property $(\mathrm{S}_k)$ at all points of $f^{-1}(y)$; since $\sh{F}$ is $f$-flat, one knows then \sref{IV.6.4.1} that $\sh{F}\otimes_{\sh{O}_Y}\sh{G}$ has the property $(\mathrm{S}_k)$ in $X$, which proves~(i).
To prove~(ii), it suffices to observe that, under the hypotheses made, $\sh{G}$ has the
\oldpage[IV-3]{178}
property $(\mathrm{S}_k)$ in a neighbourhood $U$ of $y$ in $Y$ \sref{IV.6.11.2} and \sref{IV.7.8.3},~(iv)); one concludes then in the same manner by replacing $Y$ this time by $U$ and $X$ by $V\times_Y U$.
\end{proof}

\begin{theorem}[12.1.6]
\label{IV.12.1.6}
Let $f:X\to Y$ be a flat morphism locally of finite presentation, and let $k$ be an integer.
The following subsets of $X$ are open:
\begin{enumerate}[label=(\roman*)]
  \item The set of points $x\in X$ such that $X_{f(x)}$ has the property $(\mathrm{S}_k)$ at the point $x$.
  \item The set of $x\in X$ such that $X_{f(x)}$ satisfies the property $(\mathrm{R}_k)$ geometrically at the point $x$, is equidimensional at the point $x$, and such that, furthermore, $x$ does not belong to any associated immersed prime cycle of $X_{f(x)}$.
  \item The set of $x\in X$ such that $X_{f(x)}$ is geometrically regular (i.e.\ smooth) at the point $x$.
  \item The set of $x\in X$ such that $X_{f(x)}$ is geometrically normal at the point $x$.
\end{enumerate}

Steps~A) and~B) of \sref{IV.12.0.2} are done here as in \sref{IV.12.1.1}; for step~A), one uses \sref{IV.6.7.8}, and also \sref{IV.4.2.7} for~(ii); for step~B), one uses \sref{IV.9.9.2} and \sref{IV.9.9.4}.
It remains to examine step~C) in each case.
\end{theorem}

\begin{proof}
\begin{enumerate}[label=(\roman*)]
  \item[(i)] As in \sref{IV.12.1.1.2}, one reduces (using \sref{IV.6.7.8}) to the case where $Y$ is the spectrum of a discrete valuation ring $A$, $X=\Spec(B)$, where $B$ is an $A$-algebra of finite type, and $x,x'$ are two points of $X$ such that $f(x)=y$ is the closed point and $y'=f(x')$ is the generic point of $Y$, with $x'$ in addition a generization of $x$.
As one wishes to prove that $X$ satisfies $(\mathrm{S}_k)$ at the point $x'$, one may replace $X$ by $\Spec(\sh{O}_{X,x})$; that is, one may suppose that the ring $B$ is local, that the homomorphism $A\to B$ is local, and that $B$ is an $A$-algebra essentially of finite type \sref{IV.1.3.8}.
Then, by hypothesis, if $t$ is a uniformiser of $A$, $t$ is a $B$-regular element by flatness, and $B/tB$ satisfies $(\mathrm{S}_k)$.
But since $A$ is a universally catenary ring \sref{IV.5.6.4}, $B$ is catenary, so, by \sref{IV.5.12.4}, it follows that $B$ satisfies $(\mathrm{S}_k)$, which finishes the proof in case~(i).

  \item[(iii)] Here step~C) is unnecessary; since $f$ is flat, one knows indeed \sref{IV.6.8.7} that when $Y$ is locally Noetherian and $f$ locally of finite type, the set of $x\in X$ such that $X_{f(x)}$ is geometrically regular at the point $x$ is open in $X$.

  \item[(ii)] By reasoning as in \sref{IV.12.1.1.3}, to prove that the property considered is stable under generization, one can first, by considering an arbitrary finite extension $K'$ of $k(y')$, and taking into account the definition \sref{IV.6.7.6} of the geometric property $(\mathrm{R}_k)$, as well as the invariance under change of base field of the two other properties appearing in~(ii), replace in~(ii) the geometric property $(\mathrm{R}_k)$ by the property $(\mathrm{R}_k)$.
Proceeding then as in~(i), one reduces to the case where $Y$ is the spectrum of a discrete valuation ring $A$, and since $A$ is catenary, it suffices to apply \sref{IV.5.12.5} to conclude.

  \item[(iv)] The set of points $x\in X$ such that $X_{f(x)}$ is geometrically normal at the point $x$ is contained in that of $x\in X$ such that $X_{f(x)}$ is geometrically punctually integral and satisfies $(\mathrm{S}_2)$ at the point $x$, and this latter set is open in $X$ by virtue of \sref{IV.12.1.1},~(viii).
One can therefore already suppose that, for every $x\in X$, $X_{f(x)}$ is geometrically punctually integral and satisfies $(\mathrm{S}_2)$ at the point $x$, and \emph{a fortiori} it is equidimensional.
On the other hand, by the Serre criterion \sref{IV.5.8.6}, to say that $X_{f(x)}$ is geometrically normal at the point $x$ means that $X_{f(x)}$ satisfies $(\mathrm{S}_2)$ and the property $(\mathrm{R}_1)$ geometrically at $x$;
\oldpage[IV-3]{179}
but by virtue of~(ii) and the preceding remarks, this set is the intersection of two open subsets of $X$, hence is open in $X$.
\end{enumerate}
\end{proof}

\begin{corollary}[12.1.7]
\label{IV.12.1.7}
Let $f:X\to Y$ be a morphism locally of finite presentation.
Then the set of points $x\in X$ at which $f$ has one of the following properties \sref{IV.6.8.1}:
\begin{enumerate}[label=(\roman*)]
  \item satisfies the property $(\mathrm{S}_k)$;
  \item has coprofundity $\leqslant n$;
  \item is Cohen--Macaulay;
  \item is regular (or smooth, which amounts to the same);
  \item is normal;
  \item is reduced;
\end{enumerate}
is open in $X$.
\end{corollary}

\begin{proof}
Indeed, it follows from \sref{IV.11.3.1} that the set of $x\in X$ where $f$ is \emph{flat} is \emph{open}.
One can therefore restrict to the case where $f$ is flat, and then the corollary follows from \sref{IV.12.1.1},~(iv),~(vi), and~(vii), and from \sref{IV.12.1.6},~(i),~(ii), and~(iv).
\end{proof}

\begin{remarks}[12.1.8]
\label{IV.12.1.8}
\begin{enumerate}[label=(\roman*)]
  \item In \sref{IV.12.1.6},~(ii), one cannot suppress the hypothesis that $x$ does not belong to any associated immersed prime cycle of $X_{f(x)}$.
This is shown by the example \sref{IV.5.12.3}: one takes $X=\Spec(A)$, with $Y$ the spectrum of the local ring $A_0$ of $k[T]$ corresponding to the prime ideal $(T)$, and the morphism $X\to Y$ corresponding to the injective homomorphism $A_0\to A$ obtained, by localization and passage to the quotient, from the injection $k[T]\to k[T,U,V]$.
This morphism is flat since $A$ is a torsion-free $A_0$-module \sref[0]{0.6.3.4}.
Here the fibre $X_y$ at the closed point $y$ of $Y$, equal to $\Spec(A/tA)$, is irreducible, of dimension~1 and satisfies the condition $(\mathrm{R}_0)$ geometrically since the local ring at its generic point is a field.
In contrast, at the generic point $y'$ of $Y$, the fibre $X_{y'}$ has two irreducible components of dimensions~0 and~1, and that of dimension~0 is not reduced, so $X_{y'}$ does not satisfy the condition $(\mathrm{R}_0)$.

  \item In \sref{IV.12.1.6},~(ii), one cannot suppress the hypothesis that $X_{f(x)}$ is equidimensional at the point $x$ either.
This is seen from the example \sref{IV.5.12.6}, with $X=\Spec(A)$ and $Y=\Spec(A_0)$, where $A_0$ is defined as in~(i); the morphism $X\to Y$ is again obtained from the injection $k[T]\to k[T,U,V,W]$ by localization and passage to the quotient, and is flat since $A$ is a torsion-free $A_0$-module \sref[0]{0.6.3.4}.
The fibre $X_y$ at the closed point $y$ of $Y$ is reduced (hence satisfies $(\mathrm{S}_1)$) and satisfies $(\mathrm{R}_1)$, but has two irreducible components of dimensions~2 and~1, while the fibre $X_{y'}$ at the generic point $y'$ of $Y$ does not satisfy $(\mathrm{R}_1)$.
\end{enumerate}
\end{remarks}


\subsection{Local and global properties of fibres of a proper, flat morphism of finite presentation}
\label{subsection:IV.12.2}

\oldpage[IV-3]{179}

\begin{theorem}[12.2.1]
\label{IV.12.2.1}
Let $f:X\to Y$ be a \emph{proper} morphism of finite presentation, $\sh{F}$ an $f$-flat $\sh{O}_X$-module of finite presentation, $\Phi$ a finite subset of $\mathbf{Z}\cup\{\pm\infty\}$, and $k$ an integer.
The following subsets of $Y$ are open:
\begin{enumerate}[label=(\roman*)]
  \item The set of $y\in Y$ such that the set of dimensions of the prime cycles associated to $\sh{F}_y$ is contained in $\Phi$.
  \item The set of $y\in Y$ such that the set of dimensions of the irreducible components of $\operatorname{Supp}(\sh{F}_y)$ contains $\Phi$.
  \item The set of $y\in Y$ such that all the prime cycles associated to $\sh{F}_y$ have the same dimension, equal to $k$.
  \item The set of $y\in Y$ such that $\sh{F}_y$ has no associated immersed prime cycle and the set of dimensions of the irreducible components of $\operatorname{Supp}(\sh{F}_y)$ is equal to $\Phi$.
\oldpage[IV-3]{180}
  \item The set of $y\in Y$ such that $\sh{F}_y$ has the property $(\mathrm{S}_k)$ and is equidimensional at every point of $X_y$.
  \item The set of $y\in Y$ such that $\operatorname{coprof}(\sh{F}_y)\leqslant k$.
  \item The set of $y\in Y$ such that $\sh{F}_y$ is a Cohen--Macaulay $\sh{O}_{X_y}$-Module.
  \item The set of $y\in Y$ such that $\sh{F}_y$ is geometrically reduced.
  \item The set of $y\in Y$ such that $\sh{F}_y$ is geometrically punctually integral at every point of $X_y$.
  \item The set of $y\in Y$ such that $\sh{F}_y$ is geometrically integral.
  \item The set of $y\in Y$ such that $\sh{F}_y$ has no associated immersed prime cycle and that the sum $l(y)$ of the total multiplicities \sref{IV.4.7.12} of $\sh{F}_y$ at the maximal points of $\operatorname{Supp}(\sh{F}_y)$ is $\leqslant k$.
\end{enumerate}
\end{theorem}

\begin{proof}
With the exception of~(ii), (iii), (iv), (x) and~(xi), the properties $\mathbf{P}(y)$ considered are of the form: ``for every $x\in X_y$, $\sh{F}_y$ has the property $\mathbf{Q}(x)$'', where $\mathbf{Q}(x)$ is one of the properties~(i) to~(viii) of \sref{IV.12.1.1}.
It follows from \sref{IV.12.1.1} that the set $U$ of $x\in X$ such that $\mathbf{Q}(x)$ is true is open, and that the set $V$ of $y\in Y$ such that $\mathbf{P}(y)$ is true is precisely $Y-f(X-U)$; in all these cases, the theorem is therefore already proved under the sole assumption that the morphism $f$ is \emph{closed}.
For~(iii), one applies~(i) with $\Phi$ reduced to a single element.
It remains to examine cases~(ii), (iv), and~(xi) separately ((x) follows immediately from~(xi) and \sref{IV.4.7.13}),\footnote{The printed source cites (4.7.14), but {\hyperlink{ega4.c04.subsec.4-7}{Section 4.7}} ends with {\hyperlink{ega4.c04.prop.4-7-13}{Proposition 4.7.13}} before {\hyperlink{ega4.c04.subsec.4-8}{Section 4.8}} begins; {\hyperlink{ega4.c04.prop.4-7-13}{Proposition 4.7.13}} is the relevant geometric-integrality criterion.} always applying the method described in \sref{IV.12.0.2}.
\end{proof}

\begin{env}[12.2.1.1]
\label{IV.12.2.1.1}
Case~(ii): Steps~A) and~B) of the proof proceed as at the start of \sref{IV.12.1.1}; for step~A), one uses \sref{IV.8.9.1}, \sref{IV.11.2.6} as well as \sref{IV.4.2.7}; for step~B), one uses \sref{IV.9.5.5} applied to $\operatorname{Supp}(\sh{F})$.
It remains to show that if (supposing $Y$ Noetherian) the set of dimensions of the irreducible components of $\operatorname{Supp}(\sh{F}_y)$ contains $\Phi$, the same is true of the set of dimensions of the irreducible components of $\operatorname{Supp}(\sh{F}_{y'})$, for every \emph{generization} $y'$ of $y$ in $Y$.
Let $Y_1$ be a spectrum of a discrete valuation ring such that, if $y_1$ and $y_1'$ are the closed point and the generic point of $Y_1$, there exists a morphism $h:Y_1\to Y$ such that $h(y_1)=y$, $h(y_1')=y'$ \sref[II]{II.7.1.7}.
Applying \sref{IV.4.2.7}, one sees that one can replace $Y$ by $Y_1$ and $X$ by $X_1=X\times_Y Y_1$, in other words, one can restrict to the case where $Y$ is the spectrum of a discrete valuation ring, $y$ its closed point and $y'$ its generic point.

Using (\textbf{Err}\textsubscript{\textbf{III}},~30), one can restrict to the case where $\operatorname{Supp}(\sh{F})=X$, so that $f$ is quasi-flat \sref{IV.2.3.3}; the irreducible components $Z_i$ of $X$ then dominate $Y$ \sref{IV.2.3.4}, in other words their generic points belong to $X_{y'}$, and the $Z_i\cap X_{y'}$ are the irreducible components of $X_{y'}$.
But every irreducible component of $X_y$ is contained in one of the $Z_i$, hence is an irreducible component of $(Z_i)_y$; now, one knows \sref{IV.7.1.13} that the dimensions of all the irreducible components of $(Z_i)_y$ are equal to $\dim((Z_i)_{y'})$, which finishes the proof of~(ii).
\end{env}

\begin{env}[12.2.1.2]
\label{IV.12.2.1.2}
Case~(iv): The same reasoning as at the start of \sref{IV.12.2.1.1}, using \sref{IV.12.1.1},~(iii), shows that one can already suppose (replacing $Y$ by an open subset of $Y$) that for every $y\in Y$, $\sh{F}_y$ has no associated immersed prime cycle.
The assertion of case~(iv) is then an immediate consequence of the assertions of cases~(i) and~(ii).
\end{env}

\begin{env}[12.2.1.3]
\label{IV.12.2.1.3}
\oldpage[IV-3]{181}
Case~(xi): For step~A) of the reasoning, one uses \sref{IV.8.9.1}, \sref{IV.11.2.6}, {\hyperlink{ega4.c04.thm.8-10-5}{(8.10.5,~(xii))}} (to keep the hypothesis that $f$ is proper), as well as \sref{IV.4.2.7}, \sref{IV.4.5.6} and \sref{IV.4.7.9}.
For step~B), one sees, as in case~(iv), that one can suppose that for every $y\in Y$, $\sh{F}_y$ has no associated immersed prime cycle, and one applies \sref{IV.9.8.8} which proves that the function $z\mapsto l(z)$ is constructible.
It remains therefore to see (supposing $Y$ Noetherian and $f$ proper) that for every generization $y'$ of a point $y\in Y$, $l(y')\leqslant l(y)$.
By reasoning as in \sref{IV.12.1.1.3}, one sees (with the help of \sref{IV.4.7.8} and \sref{IV.4.5.11}) that one can suppose the irreducible components of $\operatorname{Supp}(\sh{F}_{y'})$ are geometrically irreducible and that $l(y')$ is the \emph{sum of the lengths} of $(\sh{F}_{y'})_{z_j'}$ at the maximal points $z_j'$ of $\operatorname{Supp}(\sh{F}_{y'})$.
Applying again \sref{IV.4.2.7}, \sref{IV.4.5.6} and \sref{IV.4.7.9}, one reduces, as in \sref{IV.12.2.1.1}, to the case where $Y$ is a spectrum of a discrete valuation ring, $y$ its closed point and $y'$ its generic point.
The fact that $l(y')\leqslant l(y)$ will then be a consequence of the

\begin{lemma}[12.2.1.4]
\label{IV.12.2.1.4}
Let $Y$ be a spectrum of a discrete valuation ring, $y$ its closed point, $y'$ its generic point, $f:X\to Y$ a proper morphism, $\sh{F}$ a coherent $\sh{O}_X$-Module and $f$-flat, $z_i$ \textnormal{(}resp.\ $z_j'$\textnormal{)} the maximal points of $\operatorname{Supp}(\sh{F}_y)$ \textnormal{(}resp.\ $\operatorname{Supp}(\sh{F}_{y'})$\textnormal{)}.
Suppose that $\sh{F}_y$ has no associated immersed prime cycle.
Then one has
\[
  \sum_j \operatorname{long}((\sh{F}_{y'})_{z_j'}) \leqslant \sum_i \operatorname{long}((\sh{F}_y)_{z_i}).
\]
\end{lemma}

One has $Y=\Spec(A)$, where $A$ is a discrete valuation ring, whose uniformiser we will denote by $t$, so that $\sh{F}_y=\sh{F}/t\sh{F}$.
Since $\sh{F}$ is $f$-flat, the $z_j'$ are also the maximal points of $\operatorname{Supp}(\sh{F})$ \sref{IV.2.3.4}; for every $z_i$, denote by $z_{ik}'$ those of the $z_j'$ which are generizations of $z_i$; it follows from \sref{IV.3.4.1.1} that one has
\[
  \operatorname{long}((\sh{F}_y)_{z_i}) \geqslant \sum_k \operatorname{long}((\sh{F}_{y'})_{z_{ik}'})
\]
whence by summing
\[
  \sum_i \operatorname{long}((\sh{F}_y)_{z_i}) \geqslant \sum_i \sum_k \operatorname{long}((\sh{F}_{y'})_{z_{ik}'}).
\]

The lemma will therefore be proved if we establish that for every $z_j'$, there is at least one index $i$ such that $z_j'$ is one of the $z_{ik}'$.
Now, since $f$ is proper (hence closed) and $f(z_j')=y'$, there exists $x\in X$ which is a specialisation of $z_j'$ and is such that $f(x)=y$, in other words $\overline{\{z_j'\}}\cap X_y$ is non-empty.
Since $t$ is $\sh{F}$-regular by flatness, one deduces from \sref{IV.3.4.3} that there is at least one point of $\operatorname{Ass}(\sh{F}_y)$ of which $z_j'$ is a generization; but since the points of $\operatorname{Ass}(\sh{F}_y)$ are, by hypothesis, the $z_i$, this finishes the proof of the lemma.
\end{env}

\begin{corollary}[12.2.2]
\label{IV.12.2.2}
Under the hypotheses of \sref{IV.12.2.1}:
\begin{enumerate}[label=(\roman*)]
  \item The function $y\mapsto\dim(\operatorname{Supp}(\sh{F}_y))$ is continuous (hence locally constant) in $Y$.
  \item The function $y\mapsto\operatorname{coprof}(\sh{F}_y)$ \sref{IV.5.7.1} is upper semi-continuous in $Y$.
  \item The function $y\mapsto l(y)$ \sref{IV.12.2.1},~(xi) is upper semi-continuous in $Y$ when the $\sh{F}_y$ have no associated immersed prime cycle.
\oldpage[IV-3]{182}
  \item The function $y\mapsto\operatorname{prof}^*(\sh{F}_y)$ \sref{IV.10.8.1} is lower semi-continuous in $Y$.
\end{enumerate}
\end{corollary}

\begin{proof}
\begin{enumerate}[label=(\roman*)]
  \item It follows from \sref{IV.12.2.1},~(i) applied to $\Phi$ equal to the smallest interval of $\mathbf{N}$ containing the dimensions of the prime cycles associated to $\sh{F}_y$, that $z\mapsto\dim(\operatorname{Supp}(\sh{F}_z))$ is upper semi-continuous at the point $y$; it follows on the other hand from \sref{IV.12.2.1},~(ii) applied to $\Phi$ equal to the set of dimensions of the irreducible components of $\operatorname{Supp}(\sh{F}_y)$, that this same function is lower semi-continuous at the point $y$, whence the conclusion.
The assertions~(ii) and~(iii) are nothing but other formulations of \sref{IV.12.2.1},~(vi) and~(xi).
Finally, by \sref{IV.12.1.3}, the set of $x\in X$ such that $\operatorname{prof}_{f(x)}^*(\sh{F}_{f(x)})\leqslant k$ is open, and the conclusion of~(iv) follows by the same reasoning as at the beginning of \sref{IV.12.2.1}.
\end{enumerate}
\end{proof}

\begin{remarks}[12.2.3]
\label{IV.12.2.3}
\begin{enumerate}[label=(\roman*)]
  \item One will observe that for \sref{IV.12.2.1},~(ii), one can dispense with the hypothesis that $f$ is \emph{proper}.

  \item In \sref{IV.12.2.1},~(ii), one cannot replace the condition that the set $E(y)$ of dimensions of the irreducible components of $\operatorname{Supp}(\sh{F}_y)$ contains $\Phi$, by the condition that $E(y)$ be \emph{equal to} $\Phi$.
Indeed, let $k$ be a field, and consider the projective space $P=\mathbf{P}_k^3=\operatorname{Proj}(C)$ of dimension~$3$, where $C=k[t_0,t_1,t_2,t_3]$ (the $t_i$ being indeterminates). In $P$, let $X_0$ be the closed sub-prescheme which is the ``union of the line $X_1$ defined by $t_1=t_3=0$ and the plane $X_2$ defined by $t_2-t_3=0$''. In geometric terms, this says that $X_0=\operatorname{Proj}(C/\mathfrak{a})$, where the graded ideal $\mathfrak{a}=\mathfrak{b}\mathfrak{c}$, with $\mathfrak{b}=Ct_1+Ct_3$ and $\mathfrak{c}=C(t_2-t_3)$. Consider the morphism $f_0:X_0\to X_1$ which, geometrically, is the ``projection of $X_0$ onto $X_1$ from the line at infinity $D$ defined by $t_0=t_2=0$''. Algebraically, $f_0$ corresponds to the homomorphism of graded algebras
\[
  k[t_0,t_2]\longrightarrow k[t_0,t_1,t_2,t_3]/\mathfrak{a}
\]
obtained, by passage to the quotient, from the canonical injection $k[t_0,t_2]\to k[t_0,t_1,t_2,t_3]$.
It is clear that $f_0$ is a projective morphism, hence \emph{proper}; the same is therefore true of its restriction $f:X\to Y$, where $Y=D_+(t_0)$ and $X=f_0^{-1}(Y)$. To see that $f$ is \emph{flat}, it suffices to note that $k[t_2]$ is a principal ideal domain and that the ring
\[
  C'=k[t_1,t_2,t_3]/(t_2-t_3)((t_1)+(t_3))
\]
is a torsion-free $k[t_2]$-module \sref[0]{0.6.3.4}.
For every $y\in Y$ distinct from the point $y_0$ defined by $t_2=0$, $f^{-1}(y)$ then has two irreducible components, of dimensions~0 and~1, whereas $f^{-1}(y_0)$ has only one irreducible component, of dimension~1.

  \item In \sref{IV.12.2.1},~(i), one cannot also replace the condition that the set $E'(y)\supset E(y)$ of dimensions of the prime cycles associated to $\sh{F}_y$ be contained in $\Phi$ by the condition that it be \emph{equal to} $\Phi$.
Indeed, let $k$ be a field, let $A$ be the polynomial ring $k[t]$, and let $B$ be the subring $k[t^4,t^3u,tu^3,u^4]$ of the polynomial ring $k[t,u]$ (where $t,u$ are indeterminates). The ring $B$ is not integrally closed: the element $t^2u^2$ belongs to its integral closure but not to $B$. If $\mathfrak{m}$ is the maximal ideal of $B$ generated by $t^4,t^3u,tu^3,u^4$, then $\mathfrak{m}$ is an embedded associated prime ideal of $B/Bt^4$.\footnote{The printed source has $A/At^4$; the construction and the following fibre computation require $B/Bt^4$.}
Set $Y=\Spec(A)$, $X=\Spec(B)$, and let $f:X\to Y$ be the morphism corresponding to the homomorphism $A\to B$ which sends $t$ to $t^4$.
Since this homomorphism makes $B$ into an $A$-module without torsion, $f$ is flat \sref[0]{0.6.3.4}.
If $y$ is the point of $Y$ corresponding to the maximal ideal $tA$, the prescheme $f^{-1}(y)$ is irreducible and of dimension~1, but has an associated immersed prime cycle of dimension~0. In contrast, if $y'$ is the generic point of $Y$, then $f^{-1}(y')$ has no associated immersed prime cycle, since it is the spectrum of an integral $k(y')$-algebra of dimension~1. In this example $f$ is an affine, nonproper morphism; one obtains at once an analogous example in which $f$ is proper and flat by considering $X$ as an everywhere dense open subset of a projective scheme over $Y$ \sref[II]{II.5.3.4} and \sref[II]{II.5.3.2}, or by proceeding directly as in Remark~(ii).

Remarks~(ii) and~(iii) show the necessity of including in \sref{IV.12.2.1},~(iv) the condition that $\sh{F}_y$ have no associated immersed prime cycle.

  \item In \sref{IV.12.2.1},~(xi), one cannot suppress the hypothesis that $f$ be proper.
Indeed, let $k$ be a field, let $Y$ be the ``affine line'', the spectrum of the polynomial ring $k[t]$ ($t$ an indeterminate), and let $X$ be the disjoint union of $Y$ and the open complement of the closed point $y$ of $Y$ defined by $t=0$. It is clear that the morphism $f:X\to Y$ which is equal to the canonical injections on the two components of $X$ is flat and that, on taking $\sh{F}=\sh{O}_X$, the Module $\sh{F}_z$ has no associated immersed prime cycle for any $z\in Y$. Nevertheless, one has $l(y)=1$ and $l(z)=2$ for every $z\neq y$.

  \item In \sref{IV.12.2.1},~(xi), one cannot also suppress the hypothesis that $\sh{F}_y$ has no associated immersed prime cycle.
This is shown by the example of Remark~(ii): one sees immediately that, with $\sh{F}=\sh{O}_X$, $l(y_0)=1$ and $l(y)=2$ for $y\neq y_0$ in $Y$.

  \item We do not know if, in \sref{IV.12.2.1},~(xi), one can replace the hypothesis that $\sh{F}$ is $f$-flat by the hypothesis that $\operatorname{Supp}(\sh{F})=X$ and that $f$ is universally open.
By taking up the proof of \sref{IV.12.2.1.4}, one sees (using also \sref{IV.14.3.6}) that one would have to solve the following question: let $A$ be a local Noetherian ring, $M$ an $A$-module of finite type,
\oldpage[IV-3]{183}
$t$ an element of the maximal ideal $\mathfrak{m}$ of $A$, which is not contained in any minimal prime ideal of $A$; if there exists one of these minimal prime ideals $\mathfrak{p}$ such that $\dim(A/\mathfrak{p})=1$, is it true that $\mathfrak{m}$ is an associated prime ideal of $M/tM$ ($t$ not being supposed $M$-regular)?

One will note, however, that in \sref{IV.12.2.1},~(xi) one cannot simply suppress the hypothesis that $\sh{F}$ is $f$-flat. Take here the $P$ of Remark~(ii), and in $P$ let $X_0$ be the closed sub-prescheme which is the union of the three lines $X_1,X_2,X_3$ defined respectively by $t_1=t_3=0$, $t_1-t_0=t_3=0$, and $t_2=t_3=0$. Define the projection $f_0$ of $X_0$ onto $X_1$ as above, and its restriction $f:X\to Y$, where $Y=D_+(t_0)$ is the affine line and $X=f_0^{-1}(Y)$. The morphism $f$ is proper but not flat, and, on taking $\sh{F}=\sh{O}_X$, the Module $\sh{F}_y$ has no associated immersed prime cycle for any $y\in Y$. But one has $l(y_0)=1$ and $l(y)=2$ for $y\neq y_0$ in $Y$.
\end{enumerate}
\end{remarks}

\begin{theorem}[12.2.4]
\label{IV.12.2.4}
Let $f:X\to Y$ be a proper, flat morphism of finite presentation, and let $k$ be an integer $\geqslant 1$.
The following subsets of $Y$ are open:
\begin{enumerate}[label=(\roman*)]
  \item The set of $y\in Y$ such that $X_y$ possesses the property $(\mathrm{S}_k)$.
  \item The set of $y\in Y$ such that $X_y$ satisfies the property $(\mathrm{R}_k)$ geometrically, is equidimensional at every point, and has no associated immersed prime cycle.
  \item The set of $y\in Y$ such that $X_y$ is geometrically regular (i.e.\ smooth over $k(y)$).
  \item The set of $y\in Y$ such that $X_y$ is geometrically normal.
  \item The set of $y\in Y$ such that $X_y$ is geometrically reduced.
  \item The set of $y\in Y$ such that $X_y$ is geometrically reduced and that the number of geometric connected components of $X_y$ is equal to $k$.
  \item The set of $y\in Y$ such that $X_y$ is geometrically punctually integral \sref{IV.4.6.9}.
  \item The set of $y\in Y$ such that $X_y$ is geometrically integral.
  \item The set of $y\in Y$ such that $X_y$ has no associated immersed prime cycle and that the total multiplicity \sref{IV.4.7.4} of $X_y$ over $k(y)$ is $\leqslant k$.
\end{enumerate}
\end{theorem}

\begin{proof}
Except for~(vi), these assertions are special cases of assertions of \sref{IV.12.2.1}, or are deduced from assertions of \sref{IV.12.1.6} as at the start of the proof of \sref{IV.12.2.1}.
For~(vi), one reduces as always (taking into account the invariance of the geometric number of connected components by base change \sref{IV.4.5.6}) to the case where $Y$ is Noetherian.
The set $U$ of $y\in Y$ such that $X_y$ is geometrically reduced is open in $Y$ by virtue of~(v); it follows then from \sref[III]{III.7.8.7} and~\sref[III]{III.7.8.6} that for every $y\in U$, there is a neighbourhood $V\subset U$ of $y$ and an integer $m$ such that, for \emph{every} $z\in V$, $\Gamma(X_z,\sh{O}_{X_z})$ is isomorphic to $(k(z))^m$; but by virtue of \sref[III]{III.4.3.4}, $m$ is then the geometric number of connected components of $X_z$, whence the conclusion.
We will give another proof of~(vi) in \sref{IV.15.5.9}.
\end{proof}


\subsection{Local cohomological properties of fibres of a flat morphism, locally of finite presentation}
\label{subsection:IV.12.3}

\begin{lemma}[12.3.1]
\label{IV.12.3.1}
Let $A$ be a ring, $B$ an $A$-algebra of finite presentation which is a flat $A$-module, $M$ a $B$-module of finite presentation, which is a flat $A$-module.
Then the following conditions are equivalent:
\begin{enumerate}[label=\emph{\alph*)}]
  \item $M$ is a projective $B$-module.
  \item $M$ is a flat $B$-module.
  \item For every $y\in\Spec(A)$, $M\otimes_A k(y)$ is a projective $(B\otimes_A k(y))$-module.
\end{enumerate}
\end{lemma}

\begin{proof}
\oldpage[IV-3]{184}

The equivalence of \emph{a)} and \emph{b)} follows from Bourbaki, \emph{Alg.\ comm.}, chap.~II, \textsection\,5, n\textsuperscript{o}\,2, cor.~2 of th.~1.
Since \emph{a)} implies \emph{c)} trivially, it remains to prove that \emph{c)} implies \emph{b)}, which follows from the criterion of flatness by fibres \sref{IV.11.3.10}, applied with $g=h$, $f=1_X$.
\end{proof}

\begin{proposition}[12.3.2]
\label{IV.12.3.2}
Let $A$ be a ring, $B$ an $A$-algebra of finite presentation which is a flat $A$-module, $M$ a $B$-module of finite presentation which is a flat $A$-module.
Then:
\begin{enumerate}[label=(\roman*)]
  \item There exists a left resolution of $M$ by $B$-modules free of finite type.
  \item One has
  \begin{equation}
  \tag{12.3.2.1}
  \label{IV.12.3.2.1}
    \operatorname{dim.proj}_B(M)
    = \sup_{y\in\Spec(A)} \operatorname{dim.proj}_{B\otimes_A k(y)}(M\otimes_A k(y))
    = \operatorname{Tor.dim}_B(M).
  \end{equation}
  Here $\operatorname{Tor.dim}_B(M)$ is the smallest integer $i$ such that $\operatorname{Tor}_j^B(M,N)=0$ for every $j>i$ and every $B$-module $N$ (and is $+\infty$ if no such integer $i$ exists).
\end{enumerate}
\end{proposition}

\begin{proof}
\begin{enumerate}[label=(\roman*)]
  \item By \sref{IV.8.9.1}, there exist a Noetherian subring $A_0$ of $A$, an $A_0$-algebra of finite type $B_0$, and a $B_0$-module of finite type $M_0$, such that $B\simeq B_0\otimes_{A_0}A$ and $M\simeq M_0\otimes_{A_0}A$.
  By \sref{IV.11.2.7}, we may moreover suppose that $B_0$ and $M_0$ are flat $A_0$-modules.
  Let $(L_0)_\bullet$ be a left resolution of $M_0$ by free $B_0$-modules of finite type.
  Since $M_0$ and the $(L_0)_i$ are flat $A_0$-modules, $(L_0)_\bullet\otimes_{A_0}A$ is a left resolution of $M$ by free $B$-modules of finite type \sref{IV.2.1.10}.

  \item By \sref[0]{0.17.2.2},~(ii),
  $\operatorname{Tor.dim}_B(M)\leqslant\operatorname{dim.proj}_B(M)$; and the definition of projective dimension immediately gives
  \[
    \operatorname{dim.proj}_{B\otimes_A k(y)}(M\otimes_A k(y))
    \leqslant \operatorname{dim.proj}_B(M)
  \]
  for every $y\in\Spec(A)$.
  To prove the reverse inequalities, consider a left resolution $(L_i)$ of $M$ by $B$-modules free of finite type, and suppose that $\operatorname{Tor.dim}_B(M)=n$ (resp.\ $\operatorname{dim.proj}_{B\otimes_A k(y)}(M\otimes_A k(y))\leqslant n$ for every $y\in\Spec(A)$).
  Then
  $R=\operatorname{Im}(L_n\to L_{n-1})=\operatorname{Ker}(L_{n-1}\to L_{n-2})$
  is a $B$-module of finite type and a flat $A$-module, by the hypotheses on $M$ and $B$ and by \sref{IV.2.1.10}.
  Moreover,
  \[
    \operatorname{Tor}_{n+1}^B(M,N)=\operatorname{Tor}_1^B(R,N)
  \]
  for every $B$-module $N$ (M, V, 7).
  The hypothesis $\operatorname{Tor.dim}_B(M)=n$ therefore implies $\operatorname{Tor}_1^B(R,N)=0$ for every $B$-module $N$; that is, $R$ is a flat $B$-module, hence projective by \sref{IV.12.3.1}.
  Thus $\operatorname{dim.proj}_B(M)\leqslant n$.
  On the other hand, the hypothesis $\operatorname{dim.proj}_{B\otimes_A k(y)}(M\otimes_A k(y))\leqslant n$ for every $y\in\Spec(A)$ implies, upon tensoring with $k(y)$, that in each of the following sequences (which are exact by the flatness over $A$ of $M$, the $L_i$, and $R$ \sref{IV.2.1.10})
  \[
    0\to R\otimes_A k(y)\to L_{n-1}\otimes_A k(y)\to\cdots\to L_0\otimes_A k(y)\to M\otimes_A k(y)\to 0
  \]
  the $(B\otimes_A k(y))$-module $R\otimes_A k(y)$ is projective, for every $y\in\Spec(A)$.
  It follows again from \sref{IV.12.3.1} that $R$ is a projective $B$-module, and hence that $\operatorname{dim.proj}_B(M)\leqslant n$.
\end{enumerate}
\end{proof}

\begin{proposition}[12.3.3]
\label{IV.12.3.3}
Let $f:X\to S$ be a morphism locally of finite presentation, and let $\sh{L}_\bullet$ be a complex of quasi-coherent $\sh{O}_X$-modules of finite presentation; for every $s\in S$, let $(\sh{L}_\bullet)_s$ denote the complex $\sh{L}_\bullet\otimes_{\sh{O}_S}k(s)$ of $\sh{O}_{X_s}$-modules of finite type.
Suppose that $\sh{L}_{n-1}$ is $f$-flat.
Then the set $U$ of $x\in X$ such that $(\sh{H}_n((\sh{L}_\bullet)_{f(x)}))_x=0$ is open in $X$.
If, in addition, $\sh{L}_n$ is $f$-flat, then $\sh{H}_n(\sh{L}_\bullet)|_U=0$.
\end{proposition}

\begin{proof}
It plainly suffices to consider the case $n=1$ in which the complex $\sh{L}_\bullet$ reduces to
$0\to\sh{L}_2\xrightarrow{u}\sh{L}_1\xrightarrow{v}\sh{L}_0\to 0$.
We may first reduce to the case where $S=\Spec(A)$ and $X$
\oldpage[IV-3]{185}
are affine, and then to the case where $S$ is \emph{Noetherian}. Indeed, by \sref{IV.8.9.1} and \sref{IV.8.5.2}, there exist a Noetherian prescheme $S'=\Spec(A')$, where $A'$ is a subring of $A$, a morphism of finite type $f':X'\to S'$, and three coherent $\sh{O}_{X'}$-modules $\sh{L}'_i$ ($0\leqslant i\leqslant 2$), such that
$X=X'\times_{S'}S$, $f=f'_{(S)}$, and $\sh{L}_i=\sh{L}'_i\otimes_{S'}S$, together with two homomorphisms $u'$, $v'$ such that $u=u'\otimes 1$, $v=v'\otimes 1$, and $v'\circ u'=0$.
We may moreover suppose that $\sh{L}'_0$ is $f'$-flat \sref{IV.11.2.7}, and that $\sh{L}'_1$ is $f'$-flat whenever $\sh{L}_1$ is supposed $f$-flat.
By faithful flatness, the hypothesis
$(\sh{H}_n((\sh{L}_\bullet)_{f(x)}))_x=0$
is equivalent to
$(\sh{H}_n((\sh{L}'_\bullet)_{f'(x')}))_{x'}=0$,
where $x'$ is the projection of $x$ in $X'$.
Thus, if $U'$ is the set of $x'\in X'$ such that
$(\sh{H}_n((\sh{L}'_\bullet)_{f'(x')}))_{x'}=0$,
then $U$ is the inverse image of $U'$ under the projection $X\to X'$; this reduces the first assertion to the case where $S$ is Noetherian.

For the second assertion, note also that $A$ is the filtered colimit of its subrings $A_\lambda$ which are $A'$-algebras of finite type.
Set $S_\lambda=\Spec(A_\lambda)$,
$X_\lambda=X'\times_{S'}S_\lambda$, and
$\sh{L}^{(\lambda)}_i=\sh{L}'_i\otimes_{S'}S_\lambda$; let
$u_\lambda=u'\otimes 1:\sh{L}_2^{(\lambda)}\to\sh{L}_1^{(\lambda)}$ and
$v_\lambda=v'\otimes 1:\sh{L}_1^{(\lambda)}\to\sh{L}_0^{(\lambda)}$, so that
$\sh{H}_1(\sh{L}_\bullet^{(\lambda)})=\operatorname{Ker}(v_\lambda)/\operatorname{Im}(u_\lambda)$.
Since filtered colimits are exact in the category of commutative groups,
\[
  \operatorname{Ker}(v)=\varinjlim_\lambda\operatorname{Ker}(v_\lambda),\qquad
  \operatorname{Im}(u)=\varinjlim_\lambda\operatorname{Im}(u_\lambda),\qquad
  \sh{H}_1(\sh{L}_\bullet)=\varinjlim_\lambda\sh{H}_1(\sh{L}_\bullet^{(\lambda)}).
\]
If $\sh{L}_1$ has been supposed $f$-flat and, after reducing to the case $X=U$, one has proved that
$\sh{H}_1(\sh{L}_\bullet^{(\lambda)})=0$ for every $\lambda$, it follows that
$\sh{H}_1(\sh{L}_\bullet)=0$.

\medskip
I) Suppose henceforth that $S$ is Noetherian.
We know (without any flatness hypothesis on $\sh{L}_0$) that the set $U$ is \emph{constructible} in $X$ \sref{IV.9.9.6}.
By \hyperref[0.9.2.5]{($\mathbf{0}_{\mathrm{III}}$,~9.2.5)}, it remains to show that every \emph{generization} $x'$ in $X$ of a point $x\in U$ also belongs to $U$.
The method set out in \sref{IV.12.0.2} applies without change (taking into account that, for a prescheme $Z$ over a field $k$ and an extension $k'$ of $k$, the projection $Z\otimes_k k'\to Z$ is faithfully flat).
We may therefore suppose that $S$ is the spectrum of a discrete valuation ring $A$, with uniformizer $t$, that $x$ lies over the closed point of $S$, and that $x'$ lies over the generic point of $S$.
The hypothesis that $\sh{L}_0$ is $f$-flat implies that $t$ is $\sh{L}_0$-regular \hyperref[0.6.3.4]{($\mathbf{0}_{\mathrm{I}}$,~6.3.4)}.
We are thus reduced to proving the following lemma, in which $B$, $M$, $N$, and $P$ will be replaced respectively by $\sh{O}_x$, $(\sh{L}_2)_x$, $(\sh{L}_1)_x$, and $(\sh{L}_0)_x$.

\begin{lemma}[12.3.3.1]
\label{IV.12.3.3.1}
Let $B$ be a local Noetherian ring, let $M$, $N$, and $P$ be three $B$-modules, with $N$ of finite type, and let
$M\xrightarrow{u}N\xrightarrow{v}P$ be two homomorphisms such that $v\circ u=0$.
Let $t$ be an element of the maximal ideal of $B$ which is $P$-regular.
If the sequence
\begin{equation}
\tag{12.3.3.2}
\label{IV.12.3.3.2}
M/tM \xrightarrow{u\otimes 1_{B/tB}} N/tN \xrightarrow{v\otimes 1_{B/tB}} P/tP
\end{equation}
is exact, then the sequence $M\xrightarrow{u}N\xrightarrow{v}P$ is exact.
\end{lemma}

First note that if $M$ is replaced by its image $Q$ in $N$ and $u$ by the injection $j:Q\to N$, the image of $j\otimes 1$ is the same as that of $u\otimes 1$; hence, without changing either the hypothesis or the conclusion, we may suppose that $u$ is injective.
On the other hand, if $R$ is the image of $v$ and $\rho:N\to R$ the canonical surjection, then $t$ is plainly $R$-regular, $\rho\circ u=0$, and the kernel of $\rho\otimes 1$ is contained in that of $v\otimes 1$.
It is therefore equal to it if \sref{IV.12.3.3.2} is exact; since moreover $\operatorname{Ker}(\rho)=\operatorname{Ker}(v)$,
\oldpage[IV-3]{186}
we see that, in proving the lemma, we may also suppose that $v$ is surjective.
The lemma is then a consequence of the following.

\begin{lemma}[12.3.3.3]
\label{IV.12.3.3.3}
Let $B$ be a ring, let $M$, $N$, and $P$ be three $B$-modules, and let
$u:M\to N$, $v:N\to P$ be two homomorphisms such that $u$ is injective, $v$ is surjective, and $v\circ u=0$.
Let $t$ be an element of $B$ which is $P$-regular.
Then one has, up to canonical isomorphism,
\begin{equation}
\tag{12.3.3.4}
\label{IV.12.3.3.4}
\operatorname{Ker}(v\otimes 1)/\operatorname{Im}(u\otimes 1) = (\operatorname{Ker}(v)/\operatorname{Im}(u))\otimes_B(B/tB)
\end{equation}
\end{lemma}

Indeed, the hypothesis that the sequence \sref{IV.12.3.3.2} is exact implies then that $(\operatorname{Ker}\,v/\operatorname{Im}\,u)\otimes_B(B/tB)=0$, and since $t$ belongs to the maximal ideal of $B$ and $\operatorname{Ker}(v)$ is a $B$-module of finite type (since $B$ is supposed Noetherian in \sref{IV.12.3.3.1}), the Nakayama lemma proves that $\operatorname{Ker}(v)=\operatorname{Im}(u)$.

It remains therefore to prove \sref{IV.12.3.3.3}.
Set $u'=u\otimes 1$, $v'=v\otimes 1$, $Z=\operatorname{Ker}(v)$, $H=Z/\operatorname{Im}(u)$, so that the sequences
\begin{gather*}
  0\to M\to Z\to H\to 0\\
  0\to Z\to N\xrightarrow{v} P\to 0
\end{gather*}
whence, by tensorising by $B/tB$ and using the lemma \sref{IV.3.4.1.4} and the fact that $t$ is $P$-regular, the exact sequences
\begin{gather*}
  M/tM \xrightarrow{w'} Z/tZ \to H/tH \to 0\\
  0\to Z/tZ \to N/tN \xrightarrow{v'} P/tP \to 0
\end{gather*}
whence $\operatorname{Ker}(v')=Z/tZ$, and since $\operatorname{Im}(w')=\operatorname{Im}(u')$, one obtains \sref{IV.12.3.3.4}.

\medskip
II) Suppose now furthermore that $\sh{L}_1$ is $f$-flat; replacing in addition $X$ by $U$, one can suppose that $X=U$.
It is a question of seeing that for every $x\in X$, $(\sh{H}_1(\sh{L}_\bullet))_x=0$.
Set $s=f(x)$, and let $\mathfrak{n}$ be the ideal $\mathfrak{m}_s\sh{O}_{X,x}$, which is contained in the maximal ideal of $\sh{O}_{X,x}$.
Since $(\sh{H}_1(\sh{L}_\bullet))_x$ is an $\sh{O}_{X,x}$-module of finite type ($X$ being locally Noetherian), it is separated for the $\mathfrak{n}$-adic topology \hyperref[0.7.3.5]{($\mathbf{0}_{\mathrm{I}}$,~7.3.5)}, and it therefore suffices to show that its separated completion for this topology is $0$.
Now, by virtue of \sref[III]{III.7.4.7.2}, this separated completion is equal to
$\varprojlim_k\sh{H}_1(\sh{L}_\bullet\otimes_{\sh{O}_S}(\sh{O}_{S,s}/\mathfrak{m}_s^{k+1}))$, and it will therefore suffice to prove that each term of this projective system is zero.
This is true by hypothesis for $k=0$; we therefore reason by induction on $k$.
The conclusion will follow from the following more general lemma, which we apply with $A=\sh{O}_{S,s}/\mathfrak{m}_s^{k+1}$ and $\mathfrak{J}$ equal to the maximal ideal $\mathfrak{m}_s/\mathfrak{m}_s^{k+1}$ of $A$.

\begin{lemma}[12.3.3.5]
\label{IV.12.3.3.5}
Let $A$ be a ring, let $\mathfrak{J}$ be a nilpotent ideal of $A$, and let
$L_\bullet:P\xrightarrow{u}Q\xrightarrow{v}R$ be a complex of $A$-modules.
Set $A_0=A/\mathfrak{J}$ and
$L_\bullet^{(0)}=L_\bullet\otimes_A A_0:P_0\xrightarrow{u_0}Q_0\xrightarrow{v_0}R_0$.
Suppose that $\operatorname{gr}_{\mathfrak{J}}^{\bullet}(A)$ is a flat $A_0$-module, and that $Q$ and $R$ are flat $A$-modules.
Then the relation $\mathrm{H}_1(L_\bullet^{(0)})=0$ implies $\mathrm{H}_1(L_\bullet)=0$.
\end{lemma}

Set
$L_\bullet^{(k)}=L_\bullet\otimes_A(A/\mathfrak{J}^{k+1}):P_k\xrightarrow{u_k}Q_k\xrightarrow{v_k}R_k$.
Since there exists an integer $k\geqslant 1$ such that $L_\bullet^{(k)}=L_\bullet$, it will suffice to prove, by induction on $k$, that the sequence
$P_k\xrightarrow{u_k}Q_k\xrightarrow{v_k}R_k$ is
\oldpage[IV-3]{187}
exact (this assertion following from the hypothesis for $k=0$).
Let then $x\in Q_k$ be an element such that $v_k(x)=0$.
Since, by the induction hypothesis, the sequence
$P_{k-1}\xrightarrow{u_{k-1}}Q_{k-1}\xrightarrow{v_{k-1}}R_{k-1}$ is exact, the canonical image of $x$ in
$Q_{k-1}=Q_k/(\mathfrak{J}^kQ/\mathfrak{J}^{k+1}Q)$ belongs to $\operatorname{Im}(u_{k-1})$.
There consequently exists $z\in P_k$ such that
$x=u_k(z)+x'$ with $x'\in\mathfrak{J}^kQ/\mathfrak{J}^{k+1}Q$.
Since $v_k\circ u_k=0$, the relation $v_k(x)=0$ implies $v_k(x')=0$; moreover one plainly has
$v_k(x')\in\mathfrak{J}^kR/\mathfrak{J}^{k+1}R$.
Everything is therefore reduced to proving that the sequence
\[
  \mathfrak{J}^kP/\mathfrak{J}^{k+1}P
  \xrightarrow{u'}
  \mathfrak{J}^kQ/\mathfrak{J}^{k+1}Q
  \xrightarrow{v'}
  \mathfrak{J}^kR/\mathfrak{J}^{k+1}R
\]
where $u'$ and $v'$ are induced by $u$ and $v$ by restriction and passage to the quotients, is exact.
Now, by hypothesis, $\mathfrak{J}^k/\mathfrak{J}^{k+1}$ is a flat $(A/\mathfrak{J})$-module, so the sequence
\[
  (P/\mathfrak{J}P)\otimes_{A/\mathfrak{J}}(\mathfrak{J}^k/\mathfrak{J}^{k+1})
  \xrightarrow{u''}
  (Q/\mathfrak{J}Q)\otimes_{A/\mathfrak{J}}(\mathfrak{J}^k/\mathfrak{J}^{k+1})
  \xrightarrow{v''}
  (R/\mathfrak{J}R)\otimes_{A/\mathfrak{J}}(\mathfrak{J}^k/\mathfrak{J}^{k+1})
\]
is exact.
But
$(Q/\mathfrak{J}Q)\otimes_{A/\mathfrak{J}}(\mathfrak{J}^k/\mathfrak{J}^{k+1})
=Q\otimes_A(\mathfrak{J}^k/\mathfrak{J}^{k+1})$
identifies with $\mathfrak{J}^kQ/\mathfrak{J}^{k+1}Q$ by the flatness of $Q$ over $A$; likewise
$(R/\mathfrak{J}R)\otimes_{A/\mathfrak{J}}(\mathfrak{J}^k/\mathfrak{J}^{k+1})$
identifies with $\mathfrak{J}^kR/\mathfrak{J}^{k+1}R$.
Finally, the image of $u'$ identifies with that of $u''$, which completes the proof.
\end{proof}

\begin{corollary}[12.3.4]
\label{IV.12.3.4}
Let $f:X\to S$ be a morphism flat and locally of finite presentation, and let $\sh{F}$ and $\sh{G}$ be two $\sh{O}_X$-modules of finite presentation which are $f$-flat.
Let $n$ be an integer $\geq 0$, $U$ the set of $x\in X$ such that $(\mathscr{E}\!xt^n_{\sh{O}_{X_{f(x)}}}(\sh{F}_{f(x)},\sh{G}_{f(x)}))_x=0$ (resp.\ $(\mathscr{T}\!or_n^{\sh{O}_{X_{f(x)}}}(\sh{F}_{f(x)},\sh{G}_{f(x)}))_x=0$).
Then $U$ is open, and one has
$\mathscr{E}\!xt^n_{\sh{O}_X}(\sh{F},\sh{G})|_U=0$
(resp.\ $\mathscr{T}\!or_n^{\sh{O}_X}(\sh{F},\sh{G})|_U=0$).
\end{corollary}

\begin{proof}
We may plainly restrict to the case where $S$ and $X$ are affine.
Then \sref{IV.12.3.2} gives a left resolution $\sh{L}_\bullet=(\sh{L}_i)$ of $\sh{F}$ formed by $\sh{O}_X$-modules \emph{free of finite type}.
If we set $\sh{M}^\bullet=\mathscr{H}\!om(\sh{L}_\bullet,\sh{G})$ (resp.\ $\sh{N}_\bullet=\sh{L}_\bullet\otimes_{\sh{O}_X}\sh{G}$), it follows that each $\sh{M}^i$ (resp.\ $\sh{N}_i$) is isomorphic to an $\sh{O}_X$-module of the form $\sh{G}^m$; hence the $\sh{M}^i$ and the $\sh{N}_i$ are $\sh{O}_X$-modules of finite presentation and are $f$-flat.
Moreover, for every base change $S'\to S$, if we set
$X'=X\times_S S'$, $\sh{F}'=\sh{F}\otimes_S S'$, $\sh{G}'=\sh{G}\otimes_S S'$, and
$\sh{L}'_\bullet=\sh{L}_\bullet\otimes_S S'$, then $\sh{L}'_\bullet$ is still a left resolution of $\sh{F}'$ by $\sh{O}_{X'}$-modules free of finite type \sref{IV.2.1.10}; furthermore,
$\sh{M}'^\bullet=\sh{M}^\bullet\otimes_S S'$ is equal to $\mathscr{H}\!om(\sh{L}'_\bullet,\sh{G}')$
(resp.\ $\sh{N}'_\bullet=\sh{N}_\bullet\otimes_S S'$ is equal to $\sh{L}'_\bullet\otimes_{\sh{O}_{X'}}\sh{G}'$), by what precedes.
In particular, for every $s\in S$, one has
$\mathscr{E}\!xt^n_{\sh{O}_{X_s}}(\sh{F}_s,\sh{G}_s)=\sh{H}^n((\sh{M}^\bullet)_s)$
(resp.\ $\mathscr{T}\!or_n^{\sh{O}_{X_s}}(\sh{F}_s,\sh{G}_s)=\sh{H}_n((\sh{N}_\bullet)_s)$).
Applying \sref{IV.12.3.3} to the $f$-flat complexes $\sh{M}^\bullet$ and $\sh{N}_\bullet$, one deduces immediately the corollary.
\end{proof}
% END NATIVE INPUT source_aligned/ega4-12
% BEGIN NATIVE INPUT source_aligned/ega4-13
\section{Equidimensional morphisms}
\label{section:IV.13}

\oldpage[IV-3]{187}

This section is devoted to the study of the variation of the \emph{dimension} of the fibres of a morphism locally of finite type $f:X\to Y$ (which has already intervened with respect to the ``dimension formula'' in~\sref{IV.5.5} and~\sref{IV.5.6}).
We first prove \sref{IV.13.1.3} that the function $x\mapsto\dim_x f^{-1}(f(x))$ is always \emph{upper semi-continuous} in $X$ (Chevalley's semi-continuity theorem).
We then study more particularly the morphisms, called ``equidimensional'', for which this function is \emph{locally constant}.
\oldpage[IV-3]{188}
Unfortunately, the notion of equidimensional morphism is not stable under base change; this is why in many questions it is more convenient to work with the notion of universally open morphism, the study of which is the subject of \textsection\textsection~14 and~15.

\subsection{Chevalley's semi-continuity theorem}
\label{subsection:IV.13.1}

\begin{lemma}[13.1.1]
\label{IV.13.1.1}
Let $Y$ be an irreducible locally Noetherian prescheme, $X$ an irreducible prescheme, $f:X\to Y$ a dominant morphism of finite type.
Let $\xi$ (resp.\ $\eta$) be the generic point of $X$ (resp.\ $Y$) and let $e=\dim(f^{-1}(\eta))$.
Then, for every $x\in X$, all the irreducible components of $f^{-1}(f(x))$ are of dimension $\geq e$.
\end{lemma}

\begin{proof}
The assertion is immediate when, for every $y\in f(X)$, $\sh{O}_y$ is a \emph{universally catenary} ring: indeed, if $x$ is the generic point of an irreducible component $Z$ of $f^{-1}(y)$, it follows from \sref{IV.5.6.5}, combined with \sref{IV.5.2.1}, that $e+\dim(\sh{O}_y)=\dim(Z)+\dim(\sh{O}_x)$; but by \sref[0]{0.16.3.9}, one has $\dim(\sh{O}_x)\leq\dim(\sh{O}_y)$, whence the conclusion in this case.

We shall now reduce the general case to this particular case.
The question is evidently local on $Y$ and, taking into account \sref{IV.4.1.1.3}, it is also local on $X$; we can therefore restrict to the case where $Y=\Spec(A)$ and $X=\Spec(B)$ are affine and irreducible, $B$ being an $A$-algebra of finite type.
Moreover, by \sref{IV.1.5.4}, we may suppose $X$ and $Y$ reduced, so that $A$ and $B$ are integral domains.
As $f$ is dominant, $A$ is then a \emph{subring} of $B$.
Consider $A$ as an inductive limit of its sub-$\bb{Z}$-algebras of finite type; it then follows from \sref{IV.8.9.1} that there exists such a subalgebra $A_0$ and an $A_0$-algebra of finite type $B_0$ such that $B=B_0\otimes_{A_0}A$.
Set $Y_0=\Spec(A_0)$, $X_0=\Spec(B_0)$, and let $f_0:X_0\to Y_0$ be the morphism corresponding to the homomorphism $A_0\to B_0$, so that $f=f_0\times 1_Y$.
It is not \emph{a priori} evident that the prescheme $X_0$ is integral, but we shall see that we can reduce to this case.
Let $\eta_0$ be the generic point of $Y_0$, so that, if $g$ denotes the morphism $Y\to Y_0$, one has $g(\eta)=\eta_0$; by transitivity of fibres \sref[I]{I.3.6.4}, one has $f^{-1}(\eta)=f_0^{-1}(\eta_0)\otimes_{k(\eta_0)}k(\eta)$, and, since $f^{-1}(\eta)$ is irreducible by hypothesis, so is $f_0^{-1}(\eta_0)$ \sref{IV.4.4.1}.
Our assertion will then follow from the following lemma:

\begin{lemma}[13.1.2]
\label{IV.13.1.2}
Let $Y_0$, $Y$ be two integral preschemes with generic points $\eta_0$, $\eta$, let $g:Y\to Y_0$ be a dominant morphism, and let $f_0:X_0\to Y_0$ be a dominant morphism such that $f_0^{-1}(\eta_0)$ is irreducible.
Let $X_0'$ be the unique irreducible component of $X_0$ meeting $f_0^{-1}(\eta_0)$ \hyperref[0.2.1.8]{($\mathbf{0}_{\mathrm{I}}$,~2.1.8)}, and also denote by $X_0'$ the closed reduced sub-prescheme of $X_0$ having $X_0'$ as underlying space.
Suppose that the prescheme $X=X_0\times_{Y_0}Y$ is integral; then $X$ is isomorphic to $X_0'\times_{Y_0}Y$.
\end{lemma}

Indeed, if $j_0:X_0'\hookrightarrow X_0$ is the canonical injection, which is a closed immersion, $j=j_0\times 1_Y:X_0'\times_{Y_0}Y\to X_0\times_{Y_0}Y=X$ is a closed immersion.
On the other hand, $X_0'$ contains the fibre $f_0^{-1}(\eta_0)$, so $X_0'\times_{Y_0}Y$ contains $f^{-1}(\eta)$; note also that $f^{-1}(\eta)$ is not empty \sref[I]{I.3.4.7}, so it contains the generic point of $X$; consequently, the image of $j$ is necessarily all of $X$.
But as $X$ is integral, the only closed sub-prescheme of $X$ having $X$ as underlying space is $X$ itself, so $j$ is an isomorphism.

This lemma being established, we can therefore suppose that $X_0$ is integral; for every $y_0\in Y_0$, $\sh{O}_{y_0}$ is a $\bb{Z}$-algebra essentially of finite type, hence a universally
\oldpage[IV-3]{189}
catenary ring \sref{IV.5.6.4}; consequently, for every $y_0\in f_0(X_0)$, all the irreducible components of $f_0^{-1}(y_0)$\footnote{The printed source has $f_0^{-1}(y)$ here; the quantified point and the argument require $f_0^{-1}(y_0)$.} have dimension $\geq e$, since we know that $\dim(f_0^{-1}(\eta_0))=e$ by transitivity of fibres \sref{IV.4.1.4}.
For every $y\in f(X)$, we then have $y_0=g(y)\in f_0(X_0)$ and, by transitivity of fibres and \sref{IV.4.2.7}, this completes the proof.
\end{proof}

\begin{theorem}[13.1.3]
\label{IV.13.1.3}
\textnormal{(Chevalley).} ---
Let $f:X\to Y$ be a morphism locally of finite type.
For every integer $n$, the set $\mathrm{F}_n(X)$ of $x\in X$ such that $\dim_x(f^{-1}(f(x)))\geq n$ is closed; in other words, the function $x\mapsto\dim_x(f^{-1}(f(x)))$ is upper semi-continuous in $X$.
\end{theorem}

\begin{proof}
I) Suppose first that $f$ is locally of finite presentation.
The question is evidently local on $X$ and on $Y$, and we can therefore suppose $Y=\Spec(A)$, $X=\Spec(B)$ affine, $B$ being an $A$-algebra of finite presentation.
We know then \sref{IV.8.9.1} that there exist a Noetherian subring $A_0$ of $A$ and an $A_0$-algebra of finite type $B_0$ such that, on setting $Y_0=\Spec(A_0)$ and $X_0=\Spec(B_0)$, we have
\[
X=X_0\times_{Y_0}Y,\qquad f=f_0\times 1_Y,
\]
where $f_0:X_0\to Y_0$ is the morphism corresponding to the homomorphism $A_0\to B_0$.
Let $g:Y\to Y_0$ be the morphism corresponding to the canonical injection $A_0\to A$, let $y$ be a point of $Y$, and set $y_0=g(y)$; we know that $f^{-1}(y)=f_0^{-1}(y_0)\otimes_{k(y_0)}k(y)$, and it follows from \sref{IV.4.2.7} that if $p$ is the canonical projection $f^{-1}(y)\to f_0^{-1}(y_0)$, the irreducible components of $f^{-1}(y)$ are the irreducible components of the spaces $p^{-1}(Z_0)$, where $Z_0$ runs over the set of irreducible components of $f_0^{-1}(y_0)$, and each irreducible component of $p^{-1}(Z_0)$ dominates $Z_0$ and has dimension $\dim(Z_0)$.
Taking into account \sref[0]{0.14.1.5}, we see therefore that if $g':X\to X_0$ is the canonical projection, $x$ a point of $X$ and $x_0=g'(x)$, we have $\dim_x(f^{-1}(f(x)))=\dim_{x_0}(f_0^{-1}(f_0(x_0)))$, whence $\mathrm{F}_n(X)={g'}^{-1}(\mathrm{F}_n(X_0))$, and we are reduced to proving the theorem when $Y$ is \emph{Noetherian}, which we will henceforth suppose.

We can evidently suppose $X$ and $Y$ reduced \sref{IV.1.5.4}.
By considering the set of closed subsets $Y'$ of $Y$ such that the theorem is true for the closed sub-prescheme of $Y$ having $Y'$ as underlying space and for $X'=f^{-1}(Y')$, we can reason by Noetherian recurrence \hyperref[0.2.2.2]{($\mathbf{0}_{\mathrm{I}}$,~2.2.2)}, and suppose that the theorem is true for every closed subset $Y'\neq Y$ of $Y$.
If $X_i$ ($1\leq i\leq m$) are the closed reduced sub-preschemes of $X$ whose underlying spaces are the irreducible components of $X$, we have $\mathrm{F}_n(X)=\bigcup_i\mathrm{F}_n(X_i)$ by virtue of \sref[0]{0.14.1.5}, and we can therefore restrict to proving the theorem for each $X_i$, in other words, we can suppose $X$ \emph{irreducible}.
If $Z$ is the closed sub-prescheme of $Y$ having $\overline{f(X)}$ as underlying space, $f$ factorises as $X\xrightarrow{g}Z\xrightarrow{j}Y$, where $j$ is the canonical injection \sref[I]{I.5.2.2} and $g$ is of finite type \sref{IV.1.5.4}; it therefore suffices to prove the theorem for $Z$ and $g$. By the recurrence hypothesis, we are thus reduced to the case $Z=Y$, in other words to the case in which $Y$ is irreducible and $f$ is dominant.
Let then $\eta$ be the generic point of $Y$ and set $e=\dim(f^{-1}(\eta))$; it follows from \sref{IV.13.1.1} that for $n\leq e$ we have $\mathrm{F}_n(X)=X$, and consequently we can restrict to the case where $n>e$.
But then \sref{IV.9.5.6}, there is an open neighbourhood $U$ of $\eta$ in $Y$ such that $\mathrm{F}_n(X)\subset f^{-1}(Y-U)$; as $Y-U\neq Y$, the hypothesis of recurrence implies that $\mathrm{F}_n(X)$ is closed.

II) We now pass to the general case, still writing $S=\Spec(A)$ for the target and supposing $X=\Spec(B)$ affine, with $B$ an $A$-algebra of finite type, hence of the form
\oldpage[IV-3]{190}
$B=A[T_1,\ldots,T_n]/\mathfrak{J}$.
Let $(\mathfrak{J}_\lambda)$ be the family of finitely generated ideals of $A[T_1,\ldots,T_n]$ contained in $\mathfrak{J}$, so that $\mathfrak{J}$ is the filtered union of the $\mathfrak{J}_\lambda$; if $X$ and the $X_\lambda=\Spec(A[T_1,\ldots,T_n]/\mathfrak{J}_\lambda)$ are considered as closed sub-preschemes of $Z=\Spec(A[T_1,\ldots,T_n])$, we therefore have, for the underlying spaces, $X=\bigcap_\lambda X_\lambda$.
If $f_\lambda:X_\lambda\to S$ is the structural morphism, we deduce that $f^{-1}(s)=\bigcap_\lambda f_\lambda^{-1}(s)$ for every $s\in S$; and since the sets $f_\lambda^{-1}(s)$ are closed in the Noetherian space $Z_s$, there exists an index $\lambda$ (depending on $s$) such that $f^{-1}(s)=f_\lambda^{-1}(s)$.
For every $x\in X$, set $d(x)=\dim_x(f^{-1}(f(x)))$ and $d_\lambda(x)=\dim_x(f_\lambda^{-1}(f_\lambda(x)))$. What precedes shows that $d(x)=\inf_\lambda d_\lambda(x)$; since the functions $d_\lambda$ are upper semi-continuous by the first part of the proof, so is $d$.
\end{proof}

\begin{corollary}[13.1.4]
\label{IV.13.1.4}
Under the hypotheses of \sref{IV.13.1.3}, the set of $x\in X$ such that $x$ is isolated in $f^{-1}(f(x))$ is open in $X$.
\end{corollary}

\begin{proof}
Indeed, it is the complement of $\mathrm{F}_1(X)$ \sref[0]{0.14.1.10}.
\end{proof}

We note that we thus recover, under more general hypotheses, the consequence \sref[III]{III.4.4.10} of the ``Main theorem'' of Zariski.

\begin{corollary}[13.1.5]
\label{IV.13.1.5}
Under the hypotheses of \sref{IV.13.1.3}, suppose in addition that $f$ is a closed morphism.
Then, for every integer $n$, the set of $y\in Y$ such that $\dim(f^{-1}(y))\geq n$ is closed; in other words, the map $y\mapsto\dim(f^{-1}(y))$ is upper semi-continuous; in particular, if $y$ is a specialisation of $y'$, we have $\dim(f^{-1}(y))\geq\dim(f^{-1}(y'))$.
\end{corollary}

\begin{proof}
Indeed, to say that $\dim(f^{-1}(y))\geq n$ means that $y\in f(\mathrm{F}_n(X))$ \sref[0]{0.14.1.6}.
\end{proof}

\begin{corollary}[13.1.6]
\label{IV.13.1.6}
Let $X$, $Y$ be two irreducible preschemes, $f:X\to Y$ a dominant morphism locally of finite type.
Let $\xi$ (resp.\ $\eta$) be the generic point of $X$ (resp.\ $Y$) and let $e=\dim(f^{-1}(\eta))$.
Then, for every $x\in X$, we have $\dim_x(f^{-1}(f(x)))\geq e$.
\end{corollary}

\begin{proof}
Indeed, the set of $x$ such that $\dim_x(f^{-1}(f(x)))<e$ is open by \sref{IV.13.1.3}, and as it cannot contain $\xi$, it is empty.
\end{proof}

Let us note finally the following easy result:

\begin{proposition}[13.1.7]
\label{IV.13.1.7}
Let $Y$ be a quasi-compact prescheme, $f:X\to Y$ a morphism of finite type.
There exists an integer $n$ such that, for every $y\in Y$, we have $\dim(f^{-1}(y))\leq n$.
\end{proposition}

\begin{proof}
As there is a finite affine open covering $(V_i)$ of $Y$ such that each $f^{-1}(V_i)$ is a finite union of affine open sets, we are immediately reduced to the case where $Y=\Spec(A)$ and $X=\Spec(B)$ are affine, $B$ being an $A$-algebra of finite type.
If $B$ admits a system of $n$ generators, then, for every $y\in Y$, $B\otimes_A k(y)$ is a $k(y)$-algebra admitting $n$ generators, hence $\dim(f^{-1}(y))\leq n$ by virtue of \sref{IV.4.1.1}.
\end{proof}

\subsection{Equidimensional morphisms: case of dominant morphisms of irreducible preschemes}
\label{subsection:IV.13.2}

\begin{env}[13.2.1]
\label{IV.13.2.1}
Let $Y$ be an irreducible prescheme, $X$ an irreducible prescheme, $f:X\to Y$ a \emph{dominant} morphism locally of finite type; let $\eta$ be the generic point of $Y$.
We know \sref{IV.13.1.6} that for every $x\in X$, we have $\dim_x(f^{-1}(f(x)))\geq\dim(f^{-1}(\eta))$.
\end{env}

\oldpage[IV-3]{191}
\begin{definition}[13.2.2]
\label{IV.13.2.2}
Under the hypotheses of \sref{IV.13.2.1}, we say that $f$ is \emph{equidimensional at the point} $x$ (or that $X$ is \emph{equidimensional over} $Y$ \emph{at the point} $x$) if
\[
  \dim_x(f^{-1}(f(x)))=\dim(f^{-1}(\eta)).
\]
We say that $f$ is \emph{equidimensional} (or that $X$ is \emph{equidimensional over} $Y$) if $f$ is equidimensional at every point $x\in X$.
\end{definition}

It follows from Chevalley's theorem \sref{IV.13.1.3} that the set of $x\in X$ where $f$ is equidimensional is \emph{open} and non-empty.
Furthermore, if $f$ is equidimensional at the point $x$, all the irreducible components of $f^{-1}(f(x))$ containing $x$ have the \emph{same dimension}, since each of them has dimension $\geq\dim(f^{-1}(\eta))$ \sref{IV.13.1.6} and $\leq\dim_x(f^{-1}(f(x)))$ by virtue of \sref[0]{0.14.1.5}.

\begin{proposition}[13.2.3]
\label{IV.13.2.3}
Let $Y$ be an irreducible locally Noetherian prescheme, $X$ an irreducible prescheme, $f:X\to Y$ a dominant morphism locally of finite type; let $\eta$ be the generic point of $Y$, $x$ a point of $X$, $y=f(x)$, and suppose that we have
\[
\label{IV.13.2.3.1}
\tag{13.2.3.1}
  \dim(\sh{O}_x)=\dim(\sh{O}_y)+\dim(\sh{O}_x\otimes_{\sh{O}_y}k(y)).
\]
Then $f$ is equidimensional at the point $x$.
The converse holds if equality holds in \sref{IV.5.6.5.2}, in particular if $\sh{O}_y$ is universally catenary.
\end{proposition}

\begin{proof}
This follows immediately from \sref{IV.5.6.5.2} and from the inequality $\delta(x)\geq 0$ \sref{IV.13.1.1}.
\end{proof}

\begin{env}[13.2.4]
\label{IV.13.2.4}
Now, more generally, let $Y$ be a locally Noetherian prescheme, $f:X\to Y$ a morphism of finite type, $x$ a point of $X$, and $X'$ a closed irreducible subset of $X$ containing $x$; set $y=f(x)$ and $Y'=\overline{f(X')}$, and let $\eta'$ be the generic point of $Y'$.
Let $X'$ and $Y'$ also denote the closed reduced sub-preschemes of $X$ and $Y$ respectively having $X'$ and $Y'$ as underlying spaces; the restriction $X'\to Y$ of $f$ factorises as $X'\xrightarrow{f'}Y'\xrightarrow{j}Y$, where $j$ is the canonical injection \sref[I]{I.5.2.2}, and $f'$ is of finite type \sref{IV.1.5.4}.
Set, for brevity,
\[
\label{IV.13.2.4.1}
\tag{13.2.4.1}
  A=\sh{O}_{Y,y},\quad B=\sh{O}_{X,x},\quad A'=\sh{O}_{Y',y},\quad B'=\sh{O}_{X',x}.
\]
The formula \sref{IV.5.6.5.2} applied to the dominant morphism $f'$ and to the irreducible preschemes $X'$, $Y'$, gives
\[
\label{IV.13.2.4.2}
\tag{13.2.4.2}
  \dim(B')\leq\dim(A')+\dim(B'\otimes_{A'}k(y))-(\dim_x(f'^{-1}(y))-\dim(f'^{-1}(\eta')))
\]
On the other hand, the local ring $A'$ (resp.\ $B'$, $B'\otimes_{A'}k(y)$) is a quotient of $A$ (resp.\ $B$, $B\otimes_A k(y)$), hence by \sref[0]{0.16.1.2.1} we have
\[
\tag{13.2.4.3}
\label{IV.13.2.4.3}
  \dim(A')\leq\dim(A),\quad \dim(B')\leq\dim(B),\quad \dim(B'\otimes_{A'}k(y))\leq\dim(B\otimes_A k(y)).
\]
We deduce therefore first from \sref{IV.13.2.4.3} and \sref[0]{0.16.3.9}
\[
\label{IV.13.2.4.4}
\tag{13.2.4.4}
  \dim(B')\leq\dim(A')+\dim(B'\otimes_{A'}k(y))\leq\dim(A)+\dim(B\otimes_A k(y)).
\]
Furthermore, by virtue of \sref[0]{0.16.3.9}, we also have the inequalities
\[
\label{IV.13.2.4.5}
\tag{13.2.4.5}
  \dim(B')\leq\dim(B)\leq\dim(A)+\dim(B\otimes_A k(y)).
\]
\end{env}

\oldpage[IV-3]{192}

Comparing these inequalities therefore gives:

\begin{lemma}[13.2.5]
\label{IV.13.2.5}
With the notations of \sref{IV.13.2.4}, the following conditions are equivalent:
\begin{enumerate}[label=\emph{\alph*)}]
  \item $\dim(B')=\dim(A)+\dim(B\otimes_A k(y))$.
  \item The following relations hold simultaneously:
  \begin{enumerate}[label=(\roman*)]
    \item $\dim(A')=\dim(A)$.
    \item $\dim(B'\otimes_{A'}k(y))=\dim(B\otimes_A k(y))$, in other words
    \[
      \dim_x(X'\cap f^{-1}(y))=\dim_x(f^{-1}(y)).
    \]
    \item $\dim_x(f'^{-1}(y))=\dim(f'^{-1}(\eta'))$, in other words $f':X'\to Y'$ is equidimensional \sref{IV.13.2.2} at the point $x$.
    \item We have the equality
    \[
      \dim(B')=\dim(A')+\dim(B'\otimes_{A'}k(y))-(\dim_x(f'^{-1}(y))-\dim(f'^{-1}(\eta')))
    \]
  \end{enumerate}
  (a relation which always holds when $A=\sh{O}_{Y,y}$ is a universally catenary ring, by virtue of \sref{IV.5.6.5}).
  \item The following relations hold simultaneously:
  \begin{enumerate}[label=(\roman*)]
    \item $\dim(B')=\dim(B)$.
    \item $\dim(B)=\dim(A)+\dim(B\otimes_A k(y))$.
  \end{enumerate}
\end{enumerate}
\end{lemma}

\begin{env}[13.2.6]
\label{IV.13.2.6}
Recall now that the irreducible components $X_i$ of $X$ containing $x$ are finite in number and that, by \sref{IV.5.1.2.1} and \sref[0]{0.14.2.1.1},
\[
\label{IV.13.2.6.1}
\tag{13.2.6.1}
  \dim(\sh{O}_{X,x})=\sup_i\dim(\sh{O}_{X_i,x}).
\]
The equivalence of conditions $b)$ and $c)$ in \sref{IV.13.2.5} therefore implies, taking into account \sref[0]{0.16.3.9} and \sref[0]{0.14.2.1}:
\end{env}

\begin{proposition}[13.2.7]
\label{IV.13.2.7}
Let $Y$ be a locally Noetherian prescheme, $f:X\to Y$ a morphism of finite type, $x$ a point of $X$, $f(x)=y$; we have
\[
\tag{13.2.7.1}
\label{IV.13.2.7.1}
  \dim(\sh{O}_{X,x})\leq\dim(\sh{O}_{Y,y})+\dim(\sh{O}_{X,x}\otimes_{\sh{O}_{Y,y}}k(y)).
\]
For the two sides of \sref{IV.13.2.7.1} to be equal, it is necessary and sufficient that there exist an irreducible closed subset $X'$ of $X$ containing $x$ and satisfying simultaneously the following conditions:
\begin{enumerate}[label=(\roman*)]
  \item If $Y'=\overline{f(X')}$, we have $\dim(\sh{O}_{Y',y})=\dim(\sh{O}_{Y,y})$.
  \item $\dim_x(X'\cap f^{-1}(y))=\dim_x(f^{-1}(y))$ (in other words, $X'$ contains one of the irreducible components of $f^{-1}(y)$ of maximal dimension among those containing $x$). This amounts to saying that $\dim(\sh{O}_{X',x}\otimes_{\sh{O}_{Y',y}}k(y))=\dim(\sh{O}_{X,x}\otimes_{\sh{O}_{Y,y}}k(y))$.
  \item $X'$ is equidimensional over $Y'$ at the point $x$, in other words, we have
  \[
    \dim_x(X'\cap f^{-1}(y))=\dim(X'\cap f^{-1}(\eta')),
  \]
  where $\eta'$ is the generic point of $Y'$ (and consequently, all the irreducible components of $X'\cap f^{-1}(y)$ containing $x$ have the same dimension \sref{IV.13.2.2}).
\oldpage[IV-3]{193}
  \item We have the equality
  \[
    \dim(\sh{O}_{X',x})=\dim(\sh{O}_{Y',y})+\dim(\sh{O}_{X',x}\otimes_{\sh{O}_{Y',y}}k(y))
  \]
  (a condition always implied by (iii) when $\sh{O}_{Y,y}$ is a universally catenary ring).
\end{enumerate}
Moreover, $X'$ is then an irreducible component of $X$, and $Y'$ an irreducible component of $Y$.
\end{proposition}

In the statement, the local rings $\sh{O}_{X',x}$ and $\sh{O}_{Y',y}$ are relative to the closed reduced sub-preschemes of $X$, $Y$ respectively having $X'$ and $Y'$ as underlying spaces.

Moreover, this also proves:

\begin{corollary}[13.2.8]
\label{IV.13.2.8}
If the two sides of \sref{IV.13.2.7.1} are equal, the irreducible closed subsets $X'$ of $X$ containing $x$ and satisfying conditions (i) to (iv) of \sref{IV.13.2.7} are exactly those for which $\dim(\sh{O}_{X',x})=\dim(\sh{O}_{X,x})$ (which necessarily implies that $X'$ is an irreducible component of $X$).
\end{corollary}

Proposition \sref{IV.13.2.7} implies:

\begin{corollary}[13.2.9]
\label{IV.13.2.9}
Let $Y$ be a locally Noetherian prescheme, $f:X\to Y$ a morphism of finite type, $x$ a point of $X$, $f(x)=y$.
Suppose that $\sh{O}_{Y,y}$ is a universally catenary ring.
Then the following conditions are equivalent:
\begin{enumerate}[label=\emph{\alph*)}]
  \item The two members of \sref{IV.13.2.7.1} are equal.
  \item There exists an irreducible component $Z$ of $f^{-1}(y)$ containing $x$, of dimension $\dim_x(f^{-1}(y))$, such that, for every $x'$ in a neighbourhood of $x$ in $Z$, we have
  \[
  \label{IV.13.2.9.1}
\tag{13.2.9.1}
    \dim(\sh{O}_{X,x'})=\dim(\sh{O}_{Y,y})+\dim(\sh{O}_{X,x'}\otimes_{\sh{O}_{Y,y}}k(y)).
  \]
  \item There exists an irreducible component $Z$ of $f^{-1}(y)$ containing $x$, of dimension $\dim_x(f^{-1}(y))$, such that for the generic point $z$ of $Z$, we have
  \[
  \label{IV.13.2.9.2}
\tag{13.2.9.2}
    \dim(\sh{O}_{X,z})=\dim(\sh{O}_{Y,y}).
  \]
\end{enumerate}
\end{corollary}

Let us show that $a)$ implies $b)$.
Let $\dim_x(f^{-1}(y))=n$; by virtue of \sref{IV.13.2.7}, there exists an irreducible component $X'$ of $X$ satisfying conditions (i) to (iv) of \sref{IV.13.2.7}; let $Z$ be an irreducible component of dimension $n$ of $f^{-1}(y)$, containing $x$ and contained in $X'$.
As $f^{-1}(y)$ is locally Noetherian, there exists an open neighbourhood $U$ of $x$ in $Z$ such that $U$ meets no irreducible component of $f^{-1}(y)$ other than those containing $x$, hence \sref{IV.4.1.1.3} $\dim_{x'}(f^{-1}(y))=n$ for every $x'\in U$; it is then clear that the conditions (i) to (iii) of \sref{IV.13.2.7} are satisfied when we replace $x$ by any point $x'\in U$, and it follows similarly for condition (iv) since $\sh{O}_{Y,y}$ is universally catenary; whence the conclusion by \sref{IV.13.2.7}.
The condition $b)$ implies trivially $c)$ by \sref{IV.5.1.2}.
Finally, if $c)$ is satisfied and if $X''$ is an irreducible component of $X$ containing $Z$ and such that the conditions (i) to (iv) of \sref{IV.13.2.7} are satisfied when we replace $X'$ by $X''$ and $x$ by $z$, it is clear that these conditions are also satisfied when we replace $X'$ by $X''$ and $x$ by $x$ since $\sh{O}_{Y,y}$ is universally catenary, hence $c)$ implies $a)$.

\begin{proposition}[13.2.10]
\label{IV.13.2.10}
Let $Y$ be an irreducible locally Noetherian prescheme, $\eta$ its generic point, $f:X\to Y$ a morphism of finite type, $y$ a point of $f(X)$.
Let $Z_i$ be the
\oldpage[IV-3]{194}
irreducible components of $f^{-1}(y)$, $z_i$ the generic point of $Z_i$, and consider the following conditions:
\begin{enumerate}[label=\emph{\alph*)}]
  \item For every $x\in f^{-1}(y)$, we have the relation
  \[
  \tag{13.2.10.1}
\label{IV.13.2.10.1}
    \dim(\sh{O}_{X,x})=\dim(\sh{O}_{Y,y})+\dim(\sh{O}_{X,x}\otimes_{\sh{O}_{Y,y}}k(y)).
  \]
  \item For every $i$, we have
  \[
  \label{IV.13.2.10.2}
\tag{13.2.10.2}
    \dim(\sh{O}_{X,z_i})=\dim(\sh{O}_{Y,y}).
  \]
  \item For every $i$, there exists an irreducible component $X_i$ of $X$ containing $Z_i$ and such that $\dim(X_i\cap f^{-1}(\eta))=\dim(Z_i)$ (in other words, the \emph{closed reduced sub-prescheme} $X_i$ of $X$ is equidimensional over $Y$ at the point $z_i$).
\end{enumerate}
Then $a)$ implies $b)$ and $b)$ implies $c)$; furthermore, if $\sh{O}_{Y,y}$ is universally catenary, the three conditions $a)$, $b)$, $c)$ are equivalent.
\end{proposition}

The ring $\sh{O}_{X,z_i}\otimes_{\sh{O}_{Y,y}}k(y)$ being of dimension $0$ \sref{IV.5.1.2}, $a)$ implies $b)$ evidently; $b)$ implies $c)$ by virtue of \sref{IV.13.2.7} applied at the point $z_i$.
Conversely, suppose that $\sh{O}_{Y,y}$ is universally catenary; as the condition $c)$ implies that $f(X_i)$ is dense in $Y$, it follows from $c)$ that the conditions (i) to (iv) of \sref{IV.13.2.7} are satisfied when we replace $X'$ by $X_i$ and $x$ by $z_i$, hence $c)$ implies $b)$; finally $b)$ implies $a)$ by \sref{IV.13.2.9}.

\begin{corollary}[13.2.11]
\label{IV.13.2.11}
With the notation of \sref{IV.13.2.10}, suppose that $\sh{O}_{Y,y}$ is universally catenary.
For every $y'\in Y$, let $E(y')$ be the set of dimensions of the irreducible components of $f^{-1}(y')$ and set $d(y')=\dim(f^{-1}(y'))=\sup(E(y'))$.
Then, if the equivalent conditions $a)$, $b)$, $c)$ of \sref{IV.13.2.10} are satisfied, we have $E(y)\subset E(\eta)$, whence $d(y)\leq d(\eta)$.
\end{corollary}

Indeed, with the notations of \sref{IV.13.2.10}, $X_i\cap f^{-1}(\eta)$ is non-empty and is consequently an irreducible component of $f^{-1}(\eta)$ \sref[0]{0.2.1.8}.

\begin{remarks}[13.2.12]
\label{IV.13.2.12}
\begin{enumerate}[label=(\roman*)]
  \item Recall from \sref{IV.6.1.2} that the relation \sref{IV.13.2.10.1} is always satisfied when the morphism $f$ is \emph{flat} at the point $x$.
  \item With the notations of \sref{IV.13.2.10}, suppose that $\sh{O}_{Y,y}$ is universally catenary and in addition that the morphism $f$ is \emph{proper}; then, if the equivalent conditions $a)$, $b)$, $c)$ of \sref{IV.13.2.10} are satisfied, we have $d(y)=d(\eta)$, since it follows from \sref{IV.13.1.5} that $d(y)\geq d(\eta)$.
  \item The morphism of \sref{IV.12.2.3},~(ii) is \emph{proper} and \emph{flat} and all the local rings of $Y$ are universally catenary; furthermore, the two irreducible components $X_1$, $X_2$ of $X$ are \emph{equidimensional} over $Y$ at every point; but $E(y)$ has \emph{two} elements for every $y\neq y_0$, whereas $E(y_0)$ is reduced to a \emph{single} element, hence $E(y)$ is \emph{not constant} in $Y$.
\end{enumerate}
\end{remarks}

\subsection{Equidimensional morphisms: general case}
\label{subsection:IV.13.3}

\begin{proposition}[13.3.1]
\label{IV.13.3.1}
Let $Y$ be a prescheme, $f:X\to Y$ a morphism locally of finite type, $x$ a point of $X$, $y=f(x)$.
Let $Y_\alpha$ denote the irreducible components of $Y$ containing $y$.
Then the following conditions are equivalent:
\begin{enumerate}[label=\emph{\alph*)}]
\oldpage[IV-3]{195}
  \item There exist an integer $e$ and an open neighbourhood $U$ of $x$ such that the image by $f$ of every irreducible component of $U$ is dense in some $Y_\alpha$, and such that, for every $x'\in U$, the space $U\cap f^{-1}(f(x'))$ is equidimensional and of dimension $e$.
  \item[\emph{a$'$)}] There exist an integer $e$ and an open neighbourhood $U$ of $x$ such that the image by $f$ of every irreducible component of $U$ is dense in some $Y_\alpha$, and such that, if we denote by $y_\alpha$ the generic point of $Y_\alpha$, every irreducible component of the spaces $U\cap f^{-1}(y)$, $U\cap f^{-1}(y_\alpha)$ is of dimension $e$.
  \item[\emph{a$''$)}] There exist an integer $e$ and an open neighbourhood $U$ of $x$ such that, for each of the irreducible components $U_\lambda$ of $U$, $f(U_\lambda)$ is dense in some $Y_\alpha$ and such that, for every $x'\in U_\lambda$, the irreducible components of $U_\lambda\cap f^{-1}(f(x'))$ are all of dimension $e$.
  \item There exist an integer $e$, an open neighbourhood $U$ of $x$ and a $Y$-morphism quasi-finite $g:U\to Y\otimes_{\bb{Z}}\bb{Z}[T_1,\ldots,T_e]$ (a prescheme we will also denote $Y[T_1,\ldots,T_e]$, to abbreviate) such that the image by $g$ of every irreducible component of $U$ is dense in an irreducible component of $Y[T_1,\ldots,T_e]$.
\end{enumerate}
\end{proposition}

\begin{proof}
It is immediate that $a'')$ implies $a)$, since for every $x\in U$, the irreducible components of $U\cap f^{-1}(f(x))$ are each an irreducible component of one of the $U_\lambda\cap f^{-1}(f(x))$, whence the conclusion by \sref[0]{0.14.1.4}.
The condition $a)$ trivially implies $a')$.
Let us show that $a')$ implies $a'')$; we can restrict to the case where $f$ is of finite type and $X$ and $Y$ are reduced \sref{IV.1.5.4}.
Let $U_\lambda$ be an irreducible component of $U$, and suppose that $f(U_\lambda)$ is dense in $Y_\alpha$; then the restriction of $f$ to $U_\lambda$ factorises as $U_\lambda\xrightarrow{f_\alpha}Y_\alpha\longrightarrow Y$, where $f_\alpha$ is of finite type and dominant \sref[I]{I.5.2.2}.
Let $x_\lambda$ be the generic point of $U_\lambda$; by virtue of \sref[0]{0.2.1.8}, $U_\lambda\cap f_\alpha^{-1}(y_\alpha)$ is the unique irreducible component of $U\cap f^{-1}(y_\alpha)$ containing $x_\lambda$ and is by hypothesis of dimension $e$, equal to the dimensions of all the irreducible components of $U_\lambda\cap f^{-1}(y)$, by virtue of the hypothesis and of \sref{IV.13.1.1}.
But by Chevalley's theorem \sref{IV.13.1.3} and \sref{IV.13.1.1}, the set $V_\lambda$ of $x'\in U_\lambda$ such that $\dim_{x'}(f_\lambda^{-1}(f_\lambda(x')))=e$ is open and contains $x$, and it suffices to take the union of the $V_\lambda$ as an open set satisfying the conditions of $a'')$.

Let us now show that $a)$ implies $b)$; we can restrict to the case where $Y=\Spec(A)$ and $X=\Spec(B)$ are affine and $U=X$.
Let us first prove the following lemma:

\begin{lemma}[13.3.1.1]
\label{IV.13.3.1.1}
Let $Y=\Spec(A)$, $X=\Spec(B)$ be two affine schemes, $f:X\to Y$ a morphism of finite type, $y$ a point of $f(X)$, and set $e=\dim(f^{-1}(y))$.
Then there exists a $Y$-morphism $g:X\to Y[T_1,\ldots,T_e]$ such that (if we set $X_y=X\times_Y\Spec(k(y))$), the morphism $g_y:X_y\to\Spec(k(y)[T_1,\ldots,T_e])$ is finite.
Furthermore, for any such morphism $g$, $g_y$ is necessarily surjective, and there exists an open neighbourhood $U$ of $f^{-1}(y)$ in $X$ such that $g|U$ is quasi-finite.
\end{lemma}

\begin{proof}
Let $\mathfrak{p}=\mathfrak{j}_y$; the ring $B\otimes_A k(\mathfrak{p})$ is a $k(\mathfrak{p})$-algebra of finite type, hence by the normalisation lemma (Bourbaki, \emph{Alg.\ comm.}, chap.~V, \textsection~3, n$^{\text{o}}$~1, th.~1) there exist in $B\otimes_A k(\mathfrak{p})$ a finite sequence $(t_i)_{1\leq i\leq r}$ of elements algebraically independent over $k(\mathfrak{p})$ and such that if we set $C'=k(\mathfrak{p})[t_1,\ldots,t_r]$, $B\otimes_A k(\mathfrak{p})$ is a $C'$-algebra \emph{finite}; we therefore have $\dim(B\otimes_A k(\mathfrak{p}))=\dim(C')$ \sref[0]{0.16.1.5} and as $\dim(C')=r$ \sref{IV.5.2.1}, we have $r=e$.
As $B\otimes_A k(\mathfrak{p})=(B/\mathfrak{p}B)\otimes_{A/\mathfrak{p}}k(\mathfrak{p})$, we can, by multiplying the $t_i$ by a suitable element $\neq0$
\oldpage[IV-3]{196}
of $A/\mathfrak{p}$, suppose that each $t_i$ is the canonical image in $B\otimes_A k(\mathfrak{p})$ of an element $s_i\in B$.
Let then $u:A[T_1,\ldots,T_e]\to B$ be the homomorphism such that $u(T_i)=s_i$ for every $i$, and let $g:X\to Y[T_1,\ldots,T_e]$ be the corresponding morphism.
It is clear that, by the choice of the $s_i$, $g_y$ is a finite morphism.
For every morphism $g:X\to Y[T_1,\ldots,T_e]$, such that $g_y$ is finite, it follows from \sref{IV.5.4.2} and \sref{IV.4.1.2.1} that $g_y$ is necessarily surjective.
On the other hand, by \sref{IV.13.1.4}, the set $U$ of $x\in X$ that are isolated in their fibre $g^{-1}(g(x))$ is open in $X$ and contains $f^{-1}(y)$, and by definition the restriction $g|U$ is a quasi-finite morphism.
\end{proof}

This lemma being established, to prove that $a)$ implies $b)$, it remains to show that if $x_\lambda$ is a maximal point of $U$, then $g(x_\lambda)=z_\lambda$ is a maximal point of $Z=Y[T_1,\ldots,T_e]$.
By hypothesis $a)$, we may (after shrinking $U$ if necessary) suppose that $f(x_\lambda)$ is one of the generic points $y_\alpha$ of the irreducible components of $Y$ containing $y$. If $p:Z\to Y$ is the structural morphism, then $p(z_\lambda)=y_\alpha$, and $g$ consequently induces a quasi-finite $k(y_\alpha)$-morphism
\[
h:U\cap f^{-1}(y_\alpha)\longrightarrow p^{-1}(y_\alpha).
\]
Now $p^{-1}(y_\alpha)=\Spec(k(y_\alpha)[T_1,\ldots,T_e])$ is integral and of dimension $e$. If $z_\lambda$ were not a maximal point of $Z$, it would not be a maximal point of $p^{-1}(y_\alpha)$ \sref[0]{0.2.1.8}, and its closure $Z'_\lambda$ in $p^{-1}(y_\alpha)$ would therefore have dimension $<e$ \sref{IV.4.1.2.1}. But since $h$ is quasi-finite, it follows from \sref{IV.4.1.2} and hypothesis $a)$ that $\dim(Z'_\lambda)\geq e$ (the restriction of $h$ to $\overline{\{x_\lambda\}}$ factorising as
\[
\overline{\{x_\lambda\}}\longrightarrow Z'_\lambda\longrightarrow p^{-1}(y_\alpha)
\]
by \sref[I]{I.5.2.2}). This is a contradiction, and proves that $a)$ implies $b)$.

Let us finally prove that $b)$ implies $a)$.
Note that the structural morphism $p:Z=Y[T_1,\ldots,T_e]\to Y$ is faithfully flat; hence \sref{IV.2.3.4} the images under $p$ of the maximal points of $Z$ are the maximal points of $Y$. This already proves that the $f(x_\lambda)$ are generic points of the $Y_\alpha$.
Furthermore, if $f(x_\lambda)=y_\alpha$, the $k(y_\alpha)$-morphism $h:U\cap f^{-1}(y_\alpha)\to p^{-1}(y_\alpha)$ deduced from $g$ is dominant and quasi-finite by hypothesis; hence \sref{IV.4.1.2} we have $\dim(U_\lambda\cap f^{-1}(y_\alpha))=e$ ($U_\lambda$ being the irreducible component of $U$ with generic point $x_\lambda$).
Similarly, for every $x'\in U_\lambda$, the morphism $U_\lambda\cap f^{-1}(f(x'))\to p^{-1}(f(x'))$ deduced from $g$ is a $k(f(x'))$-morphism quasi-finite, hence \sref{IV.4.1.2} we have $\dim(f^{-1}(f(x'))\cap U_\lambda)\leq e$; but by \sref{IV.13.1.6} all the irreducible components of $U_\lambda\cap f^{-1}(f(x'))$ are of dimension $\geq e$; we see therefore that these components are of exactly dimension $e$, hence $b)$ implies $a)$.
\end{proof}

\begin{definition}[13.3.2]
\label{IV.13.3.2}
Let $Y$ be a prescheme, $f:X\to Y$ a morphism locally of finite type, $x$ a point of $X$.
We say that $f$ is \emph{equidimensional at the point} $x$ (or that $X$ is \emph{equidimensional over} $Y$ \emph{at the point} $x$) if the equivalent conditions of \sref{IV.13.3.1} are satisfied.
We say that $f$ is \emph{equidimensional} (or that $X$ is \emph{equidimensional over} $Y$) if $f$ is equidimensional at every point $x\in X$.
\end{definition}

It is clear that, for $f$ to be equidimensional at a point $x$, it is necessary and sufficient that $f_{\mathrm{red}}$ be so.
Furthermore, the conditions of \sref{IV.13.3.1} show that the set of points where $f$ is equidimensional is \emph{open} in $X$.

We note that when $X$ and $Y$ are \emph{irreducible}, to say that $f$ is equidimen-
\oldpage[IV-3]{197}
sional at the point $x$ means, by virtue of \sref{IV.13.3.1},~$a'')$, that $f$ is dominant and that $\dim_x(f^{-1}(f(x)))=\dim(f^{-1}(\eta))$ (where $\eta$ is the generic point of $Y$); the definition \sref{IV.13.3.2} thus coincides in this case with the definition \sref{IV.13.2.2}.

\begin{proposition}[13.3.3]
\label{IV.13.3.3}
Let $Y$ be a locally Noetherian prescheme, $f:X\to Y$ a morphism locally of finite type, $x$ a point of $X$, $X_j$ the irreducible components of $X$ (in finite number) containing $x$.
For $f$ to be equidimensional at the point $x$, it is necessary and sufficient that, for every $j$, $Y_j=\overline{f(X_j)}$ is an irreducible component of $Y$ and, denoting by $X_j$ and $Y_j$ the reduced closed sub-preschemes of $X$ and $Y$ with underlying spaces $X_j$ and $Y_j$, by $f_j:X_j\to Y_j$ the morphism deduced from $f$ \sref[I]{I.5.2.2}, by $y_j$ the generic point of $Y_j$, that $f_j$ is equidimensional at the point $x$ and that all the numbers $\dim(f_j^{-1}(y_j))$ are equal.
\end{proposition}

\begin{proof}
This follows immediately from \sref{IV.13.3.1},~$a')$).
\end{proof}

\begin{corollary}[13.3.4]
\label{IV.13.3.4}
With the notations of \sref{IV.13.3.3}, set $y=f(x)$.
If $\sh{O}_{Y,y}$ is a universally catenary ring and if $f$ is equidimensional at the point $x$, we have
\[
\tag{13.3.4.1}
\label{IV.13.3.4.1}
  \dim(\sh{O}_{X_j,x})=\dim(\sh{O}_{Y_j,y})+e-\deg.\mathrm{tr}_{k(y)}k(x)
\]
where $e$ is the common value of the numbers $\dim(f_j^{-1}(y_j))$.
\end{corollary}

\begin{proof}
As each of the $\sh{O}_{Y_j,y}$, a quotient of $\sh{O}_{Y,y}$, is a universally catenary ring \sref{IV.5.6.1}, the equality is a consequence of \sref{IV.5.6.5}.
\end{proof}

\begin{corollary}[13.3.5]
\label{IV.13.3.5}
With the notations of \sref{IV.13.3.4}, suppose that $f$ is equidimensional at the point $x$ and that $\sh{O}_{Y,y}$ is a universally catenary ring.
If the ring $\sh{O}_{Y,y}$ is equidimensional, then so is $\sh{O}_{X,x}$; the converse is true if the image under $f$ of the union of the $X_j$ is dense in a neighbourhood of $y$.
\end{corollary}

\begin{proof}
Indeed, to say that $\sh{O}_{Y,y}$ is equidimensional means that the numbers $\dim(\sh{O}_{Y'_i,y})$ are equal for all the irreducible components $Y'_i$ of $Y$ containing $y$, as follows from \sref{IV.5.1.1.5} applied to the local scheme $\Spec(\sh{O}_{Y,y})$; since we have the relation \sref{IV.13.3.4.1}, it follows by the same reasoning that $\sh{O}_{X,x}$ is then equidimensional.
Conversely, if $\sh{O}_{X,x}$ is equidimensional, all the numbers $\dim(\sh{O}_{Y_j,y})$ are equal by \sref{IV.13.3.4.1}; we deduce that $\sh{O}_{Y,y}$ is equidimensional if the $Y_j$ are all the irreducible components of $Y$ containing $y$, which follows from the supplementary hypothesis.
\end{proof}

\begin{proposition}[13.3.6]
\label{IV.13.3.6}
Let $Y$ be a locally Noetherian prescheme, $f:X\to Y$ a morphism locally of finite type, $x$ a point of $X$, $y=f(x)$.
Suppose that the ring $\sh{O}_y$ is equidimensional. Then, if $\sh{O}_x$ is equidimensional and if we have the equality
\[
\tag{13.3.6.1}
\label{IV.13.3.6.1}
  \dim(\sh{O}_x)=\dim(\sh{O}_y)+\dim(\sh{O}_x\otimes_{\sh{O}_y}k(y))
\]
(cf.\ \sref{IV.13.2.7.1}) $f$ is equidimensional at the point $x$, and the converse is true if $\sh{O}_y$ is a universally catenary ring.
\end{proposition}

\begin{proof}
Keeping the notations of \sref{IV.13.3.3}, it follows from \sref{IV.13.2.8} and from the hypothesis that $\sh{O}_x$ is equidimensional that each $Y_j$ is an irreducible component of $Y$ and that each $f_j$ is equidimensional at the point $x$; moreover, by \sref{IV.13.2.8}, we have (taking into account \sref{IV.5.6.5})
\[
  \dim(\sh{O}_{X_j,x})=\dim(\sh{O}_{Y_j,y})+\dim_x(X_j\cap f^{-1}(y))-\deg.\mathrm{tr}_{k(y)}k(x).
\]

\oldpage[IV-3]{198}

Now, since $\sh{O}_y$ is assumed equidimensional, this equality can be written
\[
\tag{13.3.6.2}
\label{IV.13.3.6.2}
  \dim_x(X_j\cap f^{-1}(y))=\dim(\sh{O}_{X,x})-\dim(\sh{O}_{Y,y})+\deg.\mathrm{tr}_{k(y)}k(x).
\]
The first member of \sref{IV.13.3.6.2} is therefore independent of $j$; but as $f_j$ is equidimensional at the point $x$, we have $\dim_x(X_j\cap f^{-1}(y))=\dim(f_j^{-1}(y_j))$, hence the criterion of \sref{IV.13.3.3} shows that $f$ is equidimensional at the point $x$.

Conversely, suppose that $\sh{O}_{Y,y}$ is a universally catenary ring and that $f$ is equidimensional at the point $x$; then \sref{IV.13.3.5} $\sh{O}_{X,x}$ is equidimensional, and it follows from \sref{IV.13.3.4} that we have the relation
\[
  \dim(\sh{O}_{X,x})=\dim(\sh{O}_{Y,y})+e-\deg.\mathrm{tr}_{k(y)}k(x)
\]
where $e$ is the common value of the numbers $\dim(f_j^{-1}(y_j))=\dim_x(X_j\cap f^{-1}(y))$; but by definition $e$ is also equal to $\dim_x(f^{-1}(y))$, whence the relation \sref{IV.13.3.6.1}, taking into account \sref{IV.5.6.5.2}.
\end{proof}

\begin{proposition}[13.3.7]
\label{IV.13.3.7}
Let $Y$ be a prescheme, $f:X\to Y$ a morphism locally of finite type and equidimensional, $Z$ a closed subset of $X$.
Then the function
\[
\label{IV.13.3.7.1}
\tag{13.3.7.1}
  x\mapsto\operatorname{codim}_x(Z\cap f^{-1}(f(x)),f^{-1}(f(x)))
\]
is lower semi-continuous in $X$.
\end{proposition}

\begin{proof}
As all the irreducible components of $f^{-1}(f(x))$ containing $x$ have the same dimension by hypothesis \sref{IV.13.3.1}, we have, by virtue of \sref{IV.5.2.1},
\[
  \operatorname{codim}_x(Z\cap f^{-1}(f(x)),f^{-1}(f(x)))=\dim_x(f^{-1}(f(x)))-\dim_x(Z\cap f^{-1}(f(x))).
\]
But by hypothesis the first term of the right-hand side is a continuous function of $x$ \sref{IV.13.3.1} and the second is an upper semi-continuous function of $x$ by \sref{IV.13.1.3}\footnote{The printed source cites (13.3.3) here; the upper semi-continuity invoked is Chevalley's theorem (13.1.3).}; whence the conclusion.
\end{proof}

\begin{proposition}[13.3.8]
\label{IV.13.3.8}
Let $Y$ be a prescheme, $f:X\to Y$ a morphism locally of finite type, $x$ a point of $X$.
Let $Y'$ be a prescheme, $g:Y'\to Y$ a flat morphism, $X'=X\times_Y Y'$, $f'=f_{(Y')}:X'\to Y'$; if $f$ is equidimensional at the point $x$, then $f'$ is equidimensional at every point $x'\in X'$ above $x$.
\end{proposition}

\begin{proof}
The question being local on $X$ and $Y$, we can restrict to the case where every irreducible component of $X$ (resp.\ $Y$) contains $x$ (resp.\ $y$); as the image by $f$ of every irreducible component of $X$ is dense in an irreducible component of $Y$ \sref{IV.13.3.1}, we know \sref{IV.2.3.5} that the image by $f'$ of every irreducible component of $X'$ is dense in an irreducible component of $Y'$.
On the other hand, by transitivity of fibres \sref[I]{I.3.6.4}, it follows from \sref{IV.4.2.8} and \sref{IV.13.3.1} that, for every $z'\in X'$, the set of dimensions of the irreducible components of ${f'}^{-1}(f'(z'))$ is the same as the set of dimensions of the irreducible components of $f^{-1}(f(z))$, where $z$ is the projection of $z'$ onto $X$; whence the conclusion by virtue of \sref{IV.13.3.1},~$a)$).
\end{proof}

\begin{remarks}[13.3.9]
\label{IV.13.3.9}
The property of equidimensionality of a morphism $f:X\to Y$ is \emph{not} stable by arbitrary base change $g:Y'\to Y$, even when $g$ is the canonical injection of an irreducible component $Y'$ of $Y$ into $Y$ ($Y'$ being
\oldpage[IV-3]{199}
considered as a reduced closed sub-prescheme of $Y$).
For example, let $k$ be a field, let $A_0=k[S,T]$ be the polynomial ring in two indeterminates, and set $\mathfrak{a}=\mathfrak{p}_1\mathfrak{p}_2$, where $\mathfrak{p}_1$ and $\mathfrak{p}_2$ are the prime ideals $A_0S$ and $A_0T$ of $A_0$; let $A=A_0/\mathfrak{a}$ and $Y=\Spec(A)$, which has two irreducible components $Y_1=\Spec(A/\mathfrak{p}_1)$ and $Y_2=\Spec(A/\mathfrak{p}_2)$; take $X=Y_1$, with $f:X\to Y$ the canonical injection, which is evidently an equidimensional morphism.
On the other hand, take $Y'=Y_2$, with $g:Y'\to Y$ the canonical injection; then $X'=X\times_Y Y'=\Spec\bigl((A/\mathfrak{p}_1)\otimes_A(A/\mathfrak{p}_2)\bigr)=\Spec\bigl(A/(\mathfrak{p}_1+\mathfrak{p}_2)\bigr)$, and $\mathfrak{p}_1+\mathfrak{p}_2$ is a maximal ideal of $A$, so $X'$ is reduced to a point; the morphism $f':X'\to Y'$ is \emph{not dominant}, the image under $f'$ of the unique point of $X'$ being a closed point of $Y'$; consequently, $f'$ is not equidimensional.

One can also give a counterexample in which $X$ and $Y$ are \emph{integral}, $f$ is \emph{finite} and \emph{birational} (and \emph{a fortiori} equidimensional by \sref{IV.13.3.1},~$b)$), and $g:Y'\to Y$ is \emph{finite} and \emph{dominant}.
Let $A$ and $\overline{A}$ be the local rings defined in \sref{IV.11.7.5}, and take $Y=\Spec(A)$ and $X=\Spec(\overline{A})$; on the other hand, with the notation of \sref{IV.11.7.5}, take $Y'=\Spec(\overline{B})$; then $X'=\Spec(\overline{A}\otimes_A\overline{B})=\Spec(\overline{B}\otimes_B\overline{B})$; but one verifies at once that $\overline{B}\otimes_B\overline{B}$ is the direct product of the rings $B/\mathfrak{p}'$, $B/\mathfrak{p}''$, and two rings isomorphic to $B/\mathfrak{n}$, whose spectra are therefore reduced to a point; as the projections of these points are closed points of $Y'$, here again one sees that $f'$ does not carry an irreducible component of $X'$ onto a dense subset of an irreducible component of $Y'$, and hence $f'$ is not equidimensional.

In this example, the ring $A$ is not geometrically unibranch; we shall see in the following paragraph {\hyperlink{ega4.c04.cor.14-4-6}{(14.4.6)}} that such phenomena cannot occur when the points of $Y$ are \emph{geometrically unibranch}.
The lack of stability of the notion of equidimensional morphism greatly restricts its interest, in favour of the notion of universally open morphism, which will be studied in detail in the following paragraph.
\end{remarks}
% END NATIVE INPUT source_aligned/ega4-13
% BEGIN NATIVE INPUT source_aligned/ega4-14
\section{Universally open morphisms}
\label{section:IV.14}

Sections~\textsection\textsection~14 and~15 are devoted to the study of the notion of \emph{universally open} morphism \sref{IV.2.4.2}.
We have already seen \sref{IV.2.4.6} that a \emph{flat} morphism locally of finite presentation is universally open, the converse being false.
In~\textsection~14, we first examine the properties of the dimensions of the fibres of a universally open morphism $f:X\to Y$; when $X$ and $Y$ are locally Noetherian, $f$ behaves in this regard \sref{IV.14.2.1} like a flat morphism (cf.~\sref{IV.6.1.2}), and is in particular \emph{equidimensional} when it is locally of finite type and dominant and $X$ is irreducible \sref{IV.14.2.2}.
Conversely, an equidimensional morphism locally of finite type $f:X\to Y$ is universally open when one further assumes that $Y$ is \emph{geometrically unibranch}, and in particular when $Y$ is \emph{normal} (Chevalley's criterion, \sref{IV.14.4.4}).
We also show that universally open morphisms locally of finite type $f:X\to Y$ (when $X$ and $Y$ are
\oldpage[IV-3]{200}
locally Noetherian and $Y$ irreducible) admit ``sufficiently many'' \emph{quasi-sections}, i.e.\ in the neighbourhood of a closed point $x$ of a fibre $f^{-1}(y)$, there is a closed sub-prescheme $X'$ of $X$ containing $x$ and such that the restriction $X'\to Y$ of $f$ is a \emph{quasi-finite} morphism (hence with \emph{discrete} fibres) and dominant \sref{IV.14.5.3}.

In~\textsection~15 we study various properties of the fibres of universally open morphisms $f:X\to Y$, locally of finite type, notably when $X$ and $Y$ are locally Noetherian.
We thus obtain in particular a criterion for a point $x\in X$ to belong to a \emph{single} irreducible component of $X$, in terms of properties of the \emph{fibre} $f^{-1}(f(x))$ of $x$: it suffices that $Y$ be geometrically unibranch (for example normal) at the point $f(x)$ and that $f^{-1}(f(x))$ be geometrically punctually integral at the point $x$ \sref{IV.15.3.3}; if moreover $Y$ is locally integral at the point $f(x)$, then $f$ is flat at the point $x$ and $X$ is locally integral at the point $x$.
We also study the variation of the \emph{number of geometric connected components} of a fibre $f^{-1}(y)$; for example, if $f$ is universally open and \emph{proper}, and the fibres of $f$ geometrically reduced, this number is \emph{locally constant} \sref{IV.15.5.7}.
Finally, when $f$ admits a \emph{section} $g$ (which will be the case when $X$ is a \emph{group scheme over $Y$}) and, for every $y\in Y$, one denotes by $X_y^0$ the connected component of the fibre $X_y=f^{-1}(y)$ at the point $g(y)$ (the ``neutral component'' in the case of groups), one studies the union $X^0$ of the $X_y^0$ for $y\in Y$, and one shows \sref{IV.15.6.4} that if $f$ is universally open and the fibres $X_y$ are geometrically reduced, then $X^0$ is an \emph{open} set in $X$.

\subsection{Open morphisms}
\label{subsection:IV.14.1}

\begin{env}[14.1.1]
\label{IV.14.1.1}
Recall \sref{IV.1.10.2} that a continuous map $\psi:X\to Y$ is said to be \emph{open at a point} $x\in X$ if the image by $\psi$ of \emph{every} neighbourhood of $x$ in $X$ is a neighbourhood of $\psi(x)$ in $Y$.

We note that this does \emph{not} imply that there exists a fundamental system of neighbourhoods of $x$ whose images are \emph{open} in $Y$.
\end{env}

\begin{proposition}[14.1.2]
\label{IV.14.1.2}
Let $X$, $Y$ be two topological spaces, $\psi:X\to Y$ a continuous map, $x$ a point of $X$, $y=\psi(x)$.
\begin{enumerate}[label=(\roman*)]
  \item If $\psi$ is open at $x$, then for every subset $Y'$ of $Y$ containing $\psi(x)$ the restriction $\psi^{-1}(Y')\to Y'$ of $\psi$ is open at the point $x$.
  \item Suppose that $Y$ is the union of a locally finite family of closed subsets $(Y_i)$ and that for every $i$ such that $\psi(x)\in Y_i$, the restriction $\psi^{-1}(Y_i)\to Y_i$ of $\psi$ is open at the point $x$; then $\psi$ is open at the point $x$.
  \item Let $\gamma:X'\to X$ be a continuous map, $x'$ a point of $X'$; if the composite map $\psi\circ\gamma:X'\to Y$ is open at the point $x'$, $\psi$ is open at the point $\gamma(x')$.
\end{enumerate}
\end{proposition}

\begin{proof}
If $U$ is a neighbourhood of $x$, we have $\psi(U\cap\psi^{-1}(Y'))=\psi(U)\cap Y'$; whence (i) follows immediately.
To prove (ii), let us note that there is a neighbourhood $W$ of $\psi(x)$ in $Y$ meeting only finitely many of the closed subsets $Y_i$ that \emph{contain} $\psi(x)$; hence $W$ is the union of the $W\cap Y_i$ for these indices.
Now, if $U$ is a neighbourhood of $x$ such that $\psi(U)\cap Y_i$ is a neighbourhood of $\psi(x)$ in $Y_i$, there exists a neighbourhood $V\subset W$ of $\psi(x)$ in $Y$ such that, for every $i$ for which $\psi(x)\in Y_i$,
\oldpage[IV-3]{201}
we have $V\cap Y_i\subset\psi(U)\cap Y_i$ and since the union of the $V\cap Y_i$ for these indices is $V$, we have $V\subset\psi(U)$, hence $\psi(U)$ is a neighbourhood of $\psi(x)$ in $Y$.
Assertion (iii) is trivial.
\end{proof}

\begin{remarks}[14.1.3]
\label{IV.14.1.3}
\begin{enumerate}[label=(\roman*)]
  \item The set $Z$ of points $x\in X$ where a morphism is open is \emph{not necessarily open}.
  For example, let $K$ be a field, $A$ the polynomial ring $K[S,T]$, $V$ the affine plane $\Spec(A)$, $X$ the closed sub-prescheme of $V$ which is the ``union'' of the line $X_1$ defined by $T=0$ and of the line $X_2$ defined by $S=0$, i.e.\ $\Spec(A/\mathfrak{a})$, where $\mathfrak{a}=AST$; let $Y=X_2=\Spec(A/AS)$ and let $f$ be the projection corresponding to the canonical injection $K[T]\to A/\mathfrak{a}$; then we have $Z=X_2$, which is not open in $X$.

  \item Let $X$, $Y$ be two Noetherian irreducible preschemes of generic points $\xi$, $\eta$ respectively, $f:X\to Y$ a morphism \emph{dominant locally of finite type}; then $f$ \emph{is open at the point $\xi$}.
  Indeed, we can evidently restrict to the case where $X$ and $Y$ are reduced (hence integral) \sref{IV.1.5.4} and affine; by virtue of the generic flatness theorem \sref{IV.6.9.1}, there exists a non-empty open set $V$ of $Y$ such that the restriction $f^{-1}(V)\to V$ of $f$ is a \emph{flat} morphism, and we conclude from \sref{IV.2.4.6} that this restriction is an open morphism.
  However, there are dominant morphisms of finite type $f:X\to Y$ (where $X$ and $Y$ are \emph{irreducible}) that are not open (see for example \sref[II]{II.8.1}), the set of points where such a morphism is open is not necessarily \emph{closed} in $X$.

  \item We do not know whether or not, when $X$ and $Y$ are locally Noetherian and $f:X\to Y$ is a morphism locally of finite type, the set of points of $X$ where $f$ is open is \emph{locally constructible}.
\end{enumerate}
\end{remarks}

\begin{proposition}[14.1.4]
\label{IV.14.1.4}
Let $X$, $Y$ be two topological spaces, $\psi:X\to Y$ a continuous map.
For every $y\in Y$, the set of $x\in\psi^{-1}(y)$ where $\psi$ is open is a closed subset of $\psi^{-1}(y)$.
\end{proposition}

\begin{proof}
Indeed, suppose that $\psi$ is not open at a point $x\in\psi^{-1}(y)$; there exists then a neighbourhood $V$ of $x$ in $X$ such that $\psi(V)$ is not a neighbourhood of $y$; it follows that for every $x'\in V\cap\psi^{-1}(y)$, $\psi$ is not open at the point $x'$.
\end{proof}

\begin{remark}[14.1.5]
\label{IV.14.1.5}
Even if $X$ and $Y$ are locally Noetherian and $f$ a \emph{finite} morphism, it can happen that $f$ is open at every point of a fibre $f^{-1}(y)$, without there existing a \emph{neighbourhood} of $f^{-1}(y)$ at every point of which $f$ is open.
Take for example $K$ a field, $A$ the polynomial ring $K[\mathrm{T}_1,\mathrm{T}_2,\mathrm{T}_3]$, $V=\Spec(A)$ the affine $3$-space over $K$, $X=\Spec(A/\mathfrak{a})$, where $\mathfrak{a}=\mathfrak{b}\mathfrak{c}$, with $\mathfrak{b}=(\mathrm{T}_3)$ and $\mathfrak{c}=(\mathrm{T}_1)+(\mathrm{T}_2-\mathrm{T}_3)$ in $A$, so that $X$ is the union of the plane $X_1=\Spec(A/\mathfrak{b})$ (``plane with equation $\mathrm{T}_3=0$'') and of the line $X_2=\Spec(A/\mathfrak{c})$ (``line with equations $\mathrm{T}_1=0$, $\mathrm{T}_2=\mathrm{T}_3$~'') which are its irreducible components.
Take $Y=X_1$ and let $f:X\to Y$ be the projection corresponding to the canonical injection $K[\mathrm{T}_1,\mathrm{T}_2]\to A/\mathfrak{a}$; if $y$ is the point common to $X_1$ and $X_2$, $f^{-1}(y)$ is reduced to $y$ and $f$ is open at this point but is not open at any point of $X_2$ in a neighbourhood of $y$ and distinct from $y$.
\end{remark}

\begin{env}[14.1.6]
\label{IV.14.1.6}
In what follows, the essential role will be played by the criterion \sref{IV.1.10.3} characterising the morphisms \emph{locally of finite presentation} $f:X\to Y$ that are open at a point $x$ by the ``lifting of generizations'': for every generization $y'$ of $y=f(x)$, there exists $x'\in X$, a generization of $x$, such that $y'=f(x')$.
\end{env}

\oldpage[IV-3]{202}
\subsection{Open morphisms and the dimension formula}
\label{subsection:IV.14.2}

\begin{theorem}[14.2.1]
\label{IV.14.2.1}
Let $X$, $Y$ be two locally Noetherian preschemes, $f:X\to Y$ a morphism, $x$ a point of $X$, $y=f(x)$.
Suppose that $f$ is open at the generic points of the irreducible components of $f^{-1}(y)$ containing $x$.
Then we have the relation
\begin{equation}
\tag{14.2.1.1}
\label{IV.14.2.1.1}
  \dim(\sh{O}_x)=\dim(\sh{O}_y)+\dim(\sh{O}_x\otimes_{\sh{O}_y}k(y)).
\end{equation}
\end{theorem}

\begin{proof}
Let $y'$ be any generization of $y$ distinct from $y$, and consider the closed reduced sub-prescheme $Y'$ of $Y$ having $\overline{\{y'\}}$ as underlying space; then no irreducible component $X'$ of $f^{-1}(Y')$ containing $x$ can be contained in $f^{-1}(y)$: indeed, if $x'$ is the generic point of $X'$, we would have $f(x')=y\neq y'$, and as $x'$ is its only generization in $f^{-1}(Y')$, this would contradict the hypothesis that $f$ is open at $x'$, by virtue of the fact that $a)$ implies $c)$ in \sref{IV.1.10.3}.
We can therefore apply \sref{IV.6.1.2} to the local homomorphism of Noetherian local rings $\sh{O}_y\to\sh{O}_x$, whence the conclusion.
\end{proof}

\begin{corollary}[14.2.2]
\label{IV.14.2.2}
Let $Y$ be a locally Noetherian prescheme, $f:X\to Y$ a morphism locally of finite type, $x$ a point of $X$, $y=f(x)$.
Suppose $f$ is open at the generic points of the irreducible components of $f^{-1}(y)$ containing $x$.
Then $f$ is equidimensional at the point $x$ in each of the two following cases:
\begin{enumerate}[label=(\roman*)]
  \item $X$ is irreducible and $f$ is dominant.
  \item The rings $\sh{O}_x$ and $\sh{O}_y$ are equidimensional.
\end{enumerate}
\end{corollary}

\begin{proof}
For (i), this follows from \sref{IV.14.2.1} and from \sref{IV.13.2.3}.
For (ii), this follows from \sref{IV.14.2.1} and from \sref{IV.13.3.6}.
\end{proof}

\begin{proposition}[14.2.3]
\label{IV.14.2.3}
Let $Y$ be an irreducible locally Noetherian prescheme, $\eta$ its generic point, $f:X\to Y$ a morphism locally of finite type, $y$ a point of $f(X)$, $Z$ an irreducible component of $f^{-1}(y)$ such that $f$ is open at the generic point of $Z$.
Then $Z$ is contained in an irreducible component $X'$ of $X$ dominant over $Y$ and such that $\dim(Z)=\dim(X'\cap f^{-1}(y))=\dim(X'\cap f^{-1}(\eta))$.
\end{proposition}

\begin{proof}
This follows indeed from \sref{IV.14.2.1} applied at the generic point of $Z$, and from \sref{IV.13.2.7}.
\end{proof}

\begin{corollary}[14.2.4]
\label{IV.14.2.4}
With the notations of \sref{IV.14.2.3}, suppose that $f$ is open at all the points of $f^{-1}(y)$.
For every $z\in Y$, let $\mathrm{E}(z)$ be the set of dimensions of the irreducible components of $f^{-1}(z)$ and $d(z)=\sup\mathrm{E}(z)=\dim(f^{-1}(z))$.
Then $\mathrm{E}(y)\subset\mathrm{E}(\eta)$, whence $d(y)\leq d(\eta)$.
\end{corollary}

\begin{proof}
This follows immediately from \sref{IV.14.2.3}.
\end{proof}

\begin{corollary}[14.2.5]
\label{IV.14.2.5}
Let $Y$ be a locally Noetherian prescheme, $f:X\to Y$ a proper morphism, $y\in Y$ a point such that $f$ is open at all the points of $f^{-1}(y)$.
Then the function $z\mapsto\dim(f^{-1}(z))$ is constant in a neighbourhood of $y$.
\end{corollary}

\begin{proof}
Let $Y_i$ ($1\leq i\leq n$) be the closed reduced sub-preschemes of $Y$ having as underlying spaces the irreducible components of $Y$ containing $y$; if $f_i:f^{-1}(Y_i)\to Y_i$ is the restriction of $f$, we know that each of the $f_i$ is proper \sref[II]{II.5.4.5} and open at all the points of $f^{-1}(y)$ \sref{IV.14.1.2}.
As the union of the $Y_i$ is a neighbourhood of $y$ in $Y$,
\oldpage[IV-3]{203}
we see that one can restrict to proving the corollary when $Y$ is irreducible; let $\eta$ be its generic point.
It follows from \sref{IV.14.2.4} that $d(y)\leq d(\eta)$; on the other hand, as $f$ is proper, we deduce from \sref{IV.13.1.5} that $z\mapsto\dim(f^{-1}(z))$ is upper semi-continuous; in particular $d(y)\geq d(\eta)$, hence $d(y)=d(\eta)$.
Moreover, there is a neighbourhood $V$ of $y$ such that $d(z)\leq d(y)$ for every $z\in V$; but as every $z$ is a specialisation of $\eta$, we have on the other hand $d(z)\geq d(\eta)$, whence finally $d(z)=d(y)$ for $z\in V$.
\end{proof}

\begin{corollary}[14.2.6]
\label{IV.14.2.6}
Let $Y$ be a locally Noetherian prescheme, $f:X\to Y$ an open morphism of finite type, $y$ a point of $Y$.
Let $y_i$ ($1\leq i\leq n$) be the generic points of the irreducible components of $Y$ containing $y$.
Then, with the notations of \sref{IV.14.2.4}:
\begin{enumerate}[label=(\roman*)]
  \item If for $1\leq i\leq n$, we have $\mathrm{E}(y)=\mathrm{E}(y_i)$ (resp.\ $d(y)=d(y_i)$), the function $\mathrm{E}$ (resp.\ $d$) is constant in a neighbourhood of $y$.
  \item There exists an open neighbourhood $U$ of $f^{-1}(y)$ such that the function $z\mapsto\dim(U\cap f^{-1}(z))$ is constant in a neighbourhood of $y$.
\end{enumerate}
\end{corollary}

\begin{proof}
\begin{enumerate}[label=(\roman*)]
  \item The same reasoning as in \sref{IV.14.2.5} shows that one can restrict to considering the case where $Y$ is irreducible, of generic point $\eta$.
  If $y'$ is any generization of $y$, we have then (applying \sref{IV.14.2.4} to the restriction $f^{-1}(Y')\to Y'$ of $f$, where $Y'=\overline{\{y'\}}$, and using \sref{IV.14.1.2}) $\mathrm{E}(y)\subset\mathrm{E}(y')\subset\mathrm{E}(\eta)$ and $d(y)\leq d(y')\leq d(\eta)$; this shows that the relation $\mathrm{E}(y)=\mathrm{E}(\eta)$ (resp.\ $d(y)=d(\eta)$) implies $\mathrm{E}(y)=\mathrm{E}(y')$ (resp.\ $d(y)=d(y')$) for every generization $y'$ of $y$.
  Now, by virtue of \sref{IV.9.5.5}, the functions $\mathrm{E}$ and $d$ are locally constructible, hence it follows from $(\mathbf{0}_{\mathrm{III}},~9.2.5)$ applied to the set of $z$ such that $\mathrm{E}(z)=\mathrm{E}(y)$ (resp.\ $d(z)=d(y)$) that this set is a neighbourhood of $y$.

  \item For each of the $y_i$, $f^{-1}(y_i)$ is a Noetherian space since $f$ is of finite type; let $x_{ij}$ ($1\leq j\leq m_i$) be the generic points of its irreducible components, and set $X_{ij}'=\overline{\{x_{ij}\}}$; let us show that the complement $U$ in $X$ of the union of the $X_{ij}'$ that do not meet $f^{-1}(y)$ (which is evidently an open neighbourhood of $f^{-1}(y)$) answers the question.
  Indeed, the restriction $f|U:U\to Y$ of $f$ is an open morphism of finite type \sref[I]{I.6.3.5}; for every pair $(i,j)$ such that $x_{ij}\in U$, $x_{ij}$ is a maximal point of $U\cap f^{-1}(y_i)$, and it follows from \sref{IV.13.1.1} applied to the restriction $X_{ij}'\to Y_i=\overline{\{y_i\}}$ of $f$ (taking into account \sref[I]{I.5.2.2} and \sref{IV.1.5.4}) that all the irreducible components of $X_{ij}'\cap f^{-1}(y)$ have dimension $\geq\dim(X_{ij}'\cap f^{-1}(y_i))$.
  Taking into account \sref{IV.14.2.3} applied to the restriction $f^{-1}(Y_i)\to Y_i$ of $f$, which is an open morphism \sref{IV.14.1.2}, $f^{-1}(y)$ is, for each $i$, the union of the irreducible components of $X_{ij}'\cap f^{-1}(y)$ ($1\leq j\leq m_i$), hence $\dim(f^{-1}(y))\geq\dim(U\cap f^{-1}(y_i))$; but as $f|U$ is open, we also have $\dim(f^{-1}(y))\leq\dim(f^{-1}(y_i)\cap U)$ in virtue of \sref{IV.14.2.4}; we see therefore that one can apply (i) to $f|U$, whence the conclusion.
\end{enumerate}
\end{proof}

\begin{remark}[14.2.7]
\label{IV.14.2.7}
The example \sref{IV.13.2.12},~(iii)) shows that under the hypotheses of \sref{IV.14.2.5} or \sref{IV.14.2.6},~(ii)), one cannot, in the conclusion, replace the function $z\mapsto\dim(f^{-1}(z))$ by the function $z\mapsto\mathrm{E}(z)$; indeed, in this example, $f$ is proper and flat (hence universally open \sref{IV.2.4.6}) and every neighbourhood of the unique point of $f^{-1}(y_0)$ contains the generic points of the irreducible components of $f^{-1}(\eta)$.
\end{remark}

\oldpage[IV-3]{204}

\subsection{Universally open morphisms}
\label{subsection:IV.14.3}
\label{IV.14.3}

\begin{env}[14.3.1]
\label{IV.14.3.1}
Recall \sref{IV.2.4.2} that to say that a morphism $f:X\to Y$ is universally open means that for every morphism $g:Y'\to Y$, $f_{(Y')}:X_{(Y')}\to Y'$ is open; it suffices moreover that this be so when $Y'=Y[\mathrm{T}_1,\ldots,\mathrm{T}_e]$, for every integer $e$ \sref{IV.8.10.2} (and if $Y$ is locally Noetherian, it suffices therefore that it be so for every $Y'$ locally Noetherian).
\end{env}

\begin{proposition}[14.3.2]
\label{IV.14.3.2}
Let $f:X\to Y$ be a morphism of preschemes.
If $f$ is universally open, then, for every morphism $g:Y'\to Y$, where $Y'$ is irreducible, the image by $f'=f_{(Y')}:X_{(Y')}\to Y'$ of every irreducible component of $X_{(Y')}$ is dense in $Y'$.
Conversely, if this condition is satisfied for every $Y'$ irreducible and every morphism of finite type $g:Y'\to Y$, and if moreover $f$ is locally of finite presentation, then $f$ is universally open.
\end{proposition}

\begin{proof}
Indeed, it follows from \sref{IV.1.10.4} that if $f$ is universally open, it satisfies the condition of the statement.
Conversely, suppose that $f$ is locally of finite presentation, and let us show that for every integer $e$, if one sets $Y''=Y[\mathrm{T}_1,\ldots,\mathrm{T}_e]$, $f_{(Y'')}$ is an open morphism.
Indeed, let $Y'$ be a closed sub-prescheme of $Y''$ having as underlying space a closed irreducible subset of $Y''$; the composite morphism $Y'\to Y''\to Y$ is of finite type, hence every irreducible component of $X_{(Y')}$ dominates $Y'$ by hypothesis; we deduce therefore from \sref{IV.1.10.4} that $f_{(Y'')}$ is an open morphism.
\end{proof}

This proposition shows that the definition \sref[III]{III.4.3.9} coincides in the case considered with the general definition of universally open morphisms given in \sref{IV.2.4.2}.

The notion of an open morphism at a point \sref{IV.14.1.1} gives rise in the same way to the following:

\begin{definition}[14.3.3]
\label{IV.14.3.3}
Let $f:X\to Y$ be a morphism of preschemes, $x$ a point of $X$.
We say that $f$ is \emph{universally open at the point $x$} if, for every morphism $g:Y'\to Y$, setting $X'=X\times_Y Y'$, the morphism $f'=f_{(Y')}:X'\to Y'$ is open at every point $x'$ of $X'$ whose projection in $X$ is $x$.
\end{definition}

\begin{remarks}[14.3.3.1]
\label{IV.14.3.3.1}
\begin{enumerate}[label=(\roman*)]
  \item The reasoning of \sref{IV.8.10.1} shows (with the same notations) that if $x$ is a point of $X$, $x_\lambda$ its projection in $X_\lambda$, then, if $f_\mu$ is open at the point $x_\mu$ for every $\mu\geq\lambda$, $f$ is open at the point $x$; it suffices to restrict to opens $U$ of $X$ containing $x$ and to remark that the hypothesis implies that $f_\mu(U_\mu)$ is a neighbourhood of $f_\mu(x_\mu)$, hence $f(v_\mu^{-1}(U_\mu))$ is a neighbourhood of $f(x)$.
  We deduce that the statement \sref{IV.8.10.2} is still exact when one replaces ``universally open'' by ``universally open at the point $x$'', and ``open morphism'' by ``morphism open at every point $x_n$ of $X_n$ whose projection in $X$ is $x$'': it suffices in the proof to limit to opens $V$ containing $x_n$.

  \item The result of \sref{IV.14.1.4} is again valid for a morphism $f:X\to Y$, replacing ``open'' by ``universally open''.
  Indeed, suppose that $f$ is not universally open at a point $x\in f^{-1}(y)$; there is then a morphism $Y'\to Y$ and a point $x'\in X'$ projecting to $x$, such that $f'$ is not open at the point $x'$.
\end{enumerate}
\end{remarks}

\oldpage[IV-3]{205}

Now, if $y'=f'(x')$, the projection $f'^{-1}(y')\to f^{-1}(y)$ is an open morphism \sref{IV.2.4.10}, and there is, by virtue of \sref{IV.14.1.4}, a neighbourhood $V'$ of $x'$ in $f'^{-1}(y')$ where $f'$ is not open; hence $f$ is not universally open at any point of the image of $V'$ in $f^{-1}(y)$, which is a neighbourhood of $x$ in $f^{-1}(y)$.

\begin{proposition}[14.3.4]
\label{IV.14.3.4}
\begin{enumerate}[label=(\roman*)]
  \item Let $f:X\to Y$, $g:Y\to Z$ be two morphisms, $x$ a point of $X$, $y=f(x)$.
  If $f$ is universally open at the point $x$ and $g$ universally open at the point $y$, then $g\circ f$ is universally open at the point $x$.
  Conversely, if $g\circ f$ is universally open at the point $x$, $g$ is universally open at the point $y$.

  \item If $f:X\to Y$ is an $S$-morphism universally open at the point $x\in X$, then, for every base change $S'\to S$, $f_{(S')}:X_{(S')}\to Y_{(S')}$ is universally open at every point of $X_{(S')}$ above $x$.

  \item For $f:X\to Y$ to be universally open at the point $x\in X$, it is necessary and sufficient that $f_{\mathrm{red}}$ be so.

  \item Let $f:X\to Y$ be a morphism locally of finite presentation, $x$ a point of $X$, $y=f(x)$; set $Y_1=\Spec(\sh{O}_y)$, $X_1=X\times_Y Y_1$, $f_1=f_{(Y_1)}$.
  For $f$ to be universally open at the point $x$, it is necessary and sufficient that $f_1$ be so (recall \sref[I]{I.3.6.5} that $X_1$ is canonically identified with a subspace of $X$).
\end{enumerate}
\end{proposition}

\begin{proof}
Indeed (ii) is an evident consequence of the definition \sref{IV.14.3.3}; it also follows from the definition that in order to prove assertion (i), it suffices to do so dropping everywhere the word ``universally'', and this follows then from \sref{IV.14.1.2}.
Assertion (iii) follows from (ii) and from the fact that the canonical morphism $X_{\mathrm{red}}\to X$ is surjective.
Finally, condition (iv) is trivially necessary.
On the other hand, if it is satisfied, and if $g:Y'\to Y$ is any morphism, $X'=X_{(Y')}$, $f'=f_{(Y')}$, $x'$ a point of $X'$ above $x$, then $f'$ is locally of finite presentation, and to see that it is open at the point $x'$, it suffices to apply the criterion {\hyperlink{ega4.c04.prop.1-10-3}{(1.10.3, $b)$)}}.
Let $y'=f'(x')$, $Y_1'=\Spec(\sh{O}_{y'})$, $X_1'=X'\times_{Y'}Y_1'$; since $g(y')=y$, the composite morphism $Y_1'\to Y'\xrightarrow{g}Y$ factorises as $Y_1'\to Y_1\to Y$ \sref[I]{I.2.4.4}, hence $X_1'=X_1\times_{Y_1}Y_1'$ and $f_1'=(f_1)_{(Y_1')}$; the conclusion follows then immediately from the hypothesis that $f_1'$ is open at the point $x'$ and from \sref{IV.1.10.3}.
\end{proof}

\begin{proposition}[14.3.5]
\label{IV.14.3.5}
Let $X$, $Y$ be two preschemes, $f:X\to Y$ a morphism, $x$ a point of $X$.
Let $(Y_i)_{1\leq i\leq n}$ be a locally finite family of closed sub-preschemes of $Y$ such that the space $Y$ is the union of the $Y_i$, and suppose that for every $i$ such that $f(x)\in Y_i$, the restriction $f^{-1}(Y_i)\to Y_i$ of $f$ is universally open at the point $x$; then $f$ is universally open at the point $x$.
\end{proposition}

\begin{proof}
Taking into account the definition, this follows from the analogous proposition {\hyperlink{ega4.c04.prop.14-1-2}{(14.1.2, (ii))}} for open morphisms at a point.
\end{proof}

\begin{proposition}[14.3.6]
\label{IV.14.3.6}
Let $Y$ be a locally Noetherian prescheme, $f:X\to Y$ a morphism locally of finite type.
For $f$ to be universally open at a point $x\in X$, it is necessary and sufficient that the following condition be satisfied: for every morphism $g:Y'\to Y$, where $Y'=\Spec(A)$ is the spectrum of a discrete valuation ring, such that the image $g(y')$ of the closed point $y'$ of $Y'$ is equal to $y=f(x)$, and for every point $x'\in X'=X\times_Y Y'$ whose projections on $X$ and $Y'$ are $x$ and $y'$ respectively, there exists a generization $z'$ of $x'$ in $X'$ whose projection in $Y'$ is the generic point of $Y'$ (in other words, \emph{there is an irreducible component of $X'$ containing $x'$ and dominating $Y'$}).

\oldpage[IV-3]{206}

Moreover, in the preceding condition, one may restrict to the case where $A$ is complete, its residue field is algebraically closed, and $x'$ is rational over $k(y')$.
\end{proposition}

\begin{proof}
If $Y'$ is as in the statement, the necessity of the condition follows from the fact that $f'$ must be open at the point $x'$, and the criterion \sref{IV.1.10.3}.
To see that the condition is sufficient, consider a morphism of finite type $g:Y''\to Y$\footnote{The printed source has $g:Y''\to U$ here, although $U$ is undefined; the subsequent fibre product is over $Y$.}, and let $X''=X\times_Y Y''$, $f''=f_{(Y'')}:X''\to Y''$, and $x''$ a point of $X''$ above $x$.
Set $y''=f''(x'')$, and let $t$ be a generization of $y''$ in $Y''$, distinct from $y''$.
Since $Y''$ is locally Noetherian, it follows from \sref[II]{II.7.1.9} and $(\mathbf{0}_{\mathrm{III}},~10.3.1)$ that there exists a scheme $Y'=\Spec(A)$, where $A$ is a complete discrete valuation ring whose residue field is an algebraic closure of $k(x'')$, and a morphism $g:Y'\to Y''$ such that, if $s$ and $y'$ are the generic point and the closed point respectively, we have $g(s)=t$ and $g(y')=y''$.
There is then a point $x'$ of $X'=X\times_Y Y'=X''\times_{Y''}Y'$ whose projections in $X''$ and $Y'$ are $x''$ and $y'$ respectively, and which is rational over $k(y')$ \sref[I]{I.3.4.9}.
The hypothesis implies that there is a generization $z'$ of $x'$ in $X'$ whose projection in $Y'$ is $s$; if $z''$ is the projection of $z'$ in $X''$, $z''$ is a generization of $x''$ and its projection in $Y''$ is $t$; we conclude therefore from \sref{IV.1.10.3} that $f''$ is open at the point $x''$, hence $f$ is universally open at the point $x$ {\hyperlink{ega4.c04.subsec.14-3-3-1}{(14.3.3.1, (i))}}.
\end{proof}

\begin{corollary}[14.3.7]
\label{IV.14.3.7}
The notations being those of \sref{IV.14.3.6}:
\begin{enumerate}[label=(\roman*)]
  \item Given a point $y\in Y$, for $f$ to be universally open at all the points of $f^{-1}(y)$, it is necessary and sufficient that for every morphism $g:Y'\to Y$, where $Y'$ is the spectrum of a discrete valuation ring and where the image $g(y')$ of the closed point $y'$ of $Y'$ is $y$, every irreducible component of $X'=X\times_Y Y'$ dominates $Y'$.
  \item For $f$ to be universally open, it is necessary and sufficient that for every morphism $g:Y'\to Y$, where $Y'$ is the spectrum of a discrete valuation ring, every irreducible component of $X'=X\times_Y Y'$ dominates $Y'$.
\end{enumerate}
\end{corollary}

\begin{proof}
It is clear that it suffices to prove (i); the necessity of (i) follows from \sref{IV.14.3.2} and its sufficiency from \sref{IV.14.3.6}.
\end{proof}

\begin{proposition}[14.3.8]
\label{IV.14.3.8}
Let $Y$ be a locally Noetherian prescheme, irreducible, regular, and of dimension $1$ (for example the spectrum of a Dedekind ring), $f:X\to Y$ a morphism locally of finite type, $y$ a point of $Y$.
The following conditions are equivalent:
\begin{enumerate}[label=\emph{\alph*)}]
  \item $f_{\mathrm{red}}$ is flat at every point of $f^{-1}(y)$.
  \item $f$ is universally open in a neighbourhood of $f^{-1}(y)$.
  \item $f$ is open in a neighbourhood of $f^{-1}(y)$.
  \item Every irreducible component of $X$ meeting $f^{-1}(y)$ dominates $Y$.
\end{enumerate}
\end{proposition}

\begin{proof}
As $f_{\mathrm{red}}$ is locally of finite type \sref{IV.1.3.4}, $a)$ implies that $f_{\mathrm{red}}$ is flat in a neighbourhood of $f^{-1}(y)$ \sref{IV.11.1.1}, and it suffices to apply \sref{IV.2.4.6} in such a neighbourhood to see that $a)$ implies $b)$.
The implication $b)\Rightarrow c)$ is trivial, and the implication $c)\Rightarrow d)$ follows from \sref{IV.1.10.4} applied to a neighbourhood of $f^{-1}(y)$.
It remains to see that $d)$ implies $a)$.
We can evidently, by virtue of \sref{IV.1.3.4}, restrict to the case where $X$ is reduced.
The question being moreover local on $X$ and on $Y$, we can suppose $Y=\Spec(A)$ and $X=\Spec(B)$
\oldpage[IV-3]{207}
affine; if $X_i$ ($1\leq i\leq n$) are the closed integral sub-preschemes of $X$ whose underlying spaces are the irreducible components of $X$, then, for every $x\in X$, $\sh{O}_{X,x}$ is a subring of the direct product of the $\sh{O}_{X_i,x}$; if $y=f(x)$, it suffices to show that each $\sh{O}_{X_i,x}$ is an $\sh{O}_y$-module \emph{without torsion}, for the same will then be true of $\sh{O}_{X,x}$; since by hypothesis $\sh{O}_y$ is a regular local ring of dimension $1$, i.e.\ a discrete valuation ring \sref[II]{II.7.1.6}, it then follows from $(\mathbf{0}_1,~6.3.4)$ that $\sh{O}_{X,x}$ is a \emph{flat} $\sh{O}_y$-module.
But if $X_i=\Spec(B_i)$, where $B_i$ is an integral ring, the hypothesis implies that the homomorphism $A\to B_i$ is \emph{injective} \sref[I]{I.1.2.7}; hence $B_i$ is an $A$-module without torsion, and \emph{a fortiori} $\sh{O}_{X_i,x}$ is an $\sh{O}_y$-module without torsion.
\end{proof}

\begin{remarks}[14.3.9]
\label{IV.14.3.9}
\begin{enumerate}[label=(\roman*)]
  \item In the statement of \sref{IV.14.3.8}, one cannot dispense with the hypothesis that $Y$ is \emph{regular}.
  With the notations of \sref{IV.11.7.5}, take indeed $Y=\Spec(A)$, $X=\Spec(\overline{A})$, so that $f:X\to Y$ is a \emph{finite surjective} morphism; since $A$ is an integral local ring of dimension $1$, as is $\overline{A}$, it follows immediately from \sref{IV.1.10.4} that the morphism $f$ is \emph{open}.
  However $f$ is not universally open (nor \emph{a fortiori} flat) \sref{IV.11.7.5}.
  One could give an analogous example by taking for $Y$ the local scheme at an ``ordinary double point'' of an algebraic curve, and for $X$ the normalization of $Y$.

  \item The example of {\hyperlink{ega4.c04.item.14-1-3}{(14.1.3, (i))}} shows that the set of points of $X$ at which a morphism is universally open is not necessarily open, since the morphism $f$ in that example is universally open at every point at which it is open.
  The example $f:X\to Y$ considered above in (i) likewise shows that the set of points at which a morphism is universally open is not necessarily closed, since it is immediate that $f$ is universally open at all but one point of $X$ (indeed, it is even a local isomorphism there).
  It would be interesting to know if the set of points where a morphism is universally open is locally constructible.
\end{enumerate}
\end{remarks}

The following two propositions were pointed out to us by M.\ Artin:

\begin{lemma}[14.3.10]
\label{IV.14.3.10}
Let $A$ be a valuation ring (not necessarily discrete), $k$ its residue field, $K$ its field of fractions, and set $S=\Spec(A)$.
Let $X$ be an irreducible prescheme, $f:X\to S$ a dominant morphism of finite type; let $X_0=X\otimes_A k$ (resp.\ $X_1=X\otimes_A K$) be the fibre of $f$ at the closed point (resp.\ at the generic point) of $S$.
Then, if $X_0\neq\varnothing$, we have $\dim(X_0)=\dim(X_1)$.
\end{lemma}

\begin{proof}
We can restrict to the case where $X$ is affine, replacing if necessary $X$ by an open affine containing a generic point of an irreducible component of $X_0$ of maximal dimension.
Let $n=\dim(X_0)\geq 0$; using \sref{IV.4.1.1.3}, it follows from \sref{IV.13.3.1.1} that there exists a neighbourhood $U$ of $X_0$ in $X$ and a quasi-finite $S$-morphism $g:U\to S[\mathrm{T}_1,\ldots,\mathrm{T}_n]=Z$ such that the restriction $g_0:U_0=U\cap X_0\to Z_0=\Spec(k[\mathrm{T}_1,\ldots,\mathrm{T}_n])$ is finite and \emph{surjective}.
By base change $\Spec(K)\to S$ and restriction to the open subset $U_1=U\cap X_1$ of $X_1$, we deduce from $g$ a quasi-finite morphism $g_1:U_1\to Z_1=\Spec(K[\mathrm{T}_1,\ldots,\mathrm{T}_n])$.
As $U_1$ is dense in $X_1$ \sref{IV.4.1.1.3}, we have $\dim(U_1)=\dim(X_1)$; if we can prove that the morphism $g_1$ is \emph{dominant}, the proposition will be established, since then $\dim(U_1)\leq\dim(Z_1)=n$ and on the other hand $\dim(X_1)\geq n$ by \sref{IV.4.1.2}.
Suppose the contrary; there would then be a polynomial $\mathrm{F}_1\neq 0$ in $K[\mathrm{T}_1,\ldots,\mathrm{T}_n]$ such that $g_1(U_1)\subset\mathrm{V}(\mathrm{F}_1)$.
If $\omega$ is a valuation on $K$ associated to $A$, and if $(a_\alpha)$ is the family of coefficients of $\mathrm{F}_1$, after multiplication of $\mathrm{F}_1$ by a nonzero element of $K$, we can suppose that $\inf_\alpha(\omega(a_\alpha))=0$; in other words, $\mathrm{F}_1$ comes from a polynomial $\mathrm{F}\in A[\mathrm{T}_1,\ldots,\mathrm{T}_n]$ (with which it is identified) such that the image $\mathrm{F}_0$ of $\mathrm{F}$ in $k[\mathrm{T}_1,\ldots,\mathrm{T}_n]$ is \emph{nonzero}.
Consider then the closed subset $\mathrm{V}(\mathrm{F})$ of $Z$; we have $\mathrm{V}(\mathrm{F})\cap Z_1=\mathrm{V}(\mathrm{F}_1)$, and as $U_1$
\oldpage[IV-3]{208}
is dense in $U$ (since it contains the generic point of $X$), $g(U)\subset\mathrm{V}(\mathrm{F})$; in particular, $g_0(U_0)\subset\mathrm{V}(\mathrm{F})\cap Z_0=\mathrm{V}(\mathrm{F}_0)$; but since $\mathrm{F}_0\neq 0$, $\mathrm{V}(\mathrm{F}_0)$ is a closed subset of $Z_0$ distinct from $Z_0$, and one arrives at a contradiction.
\end{proof}

\begin{proposition}[14.3.11]
\label{IV.14.3.11}
Let $f:X\to S$ be a morphism of finite type, $(g_\lambda)_{\lambda\in L}$ a family of universally open morphisms $g_\lambda:Y_\lambda\to S$, and for every $\lambda\in L$, let $u_\lambda:Y_\lambda\to X$ be an $S$-morphism.
For every $s\in S$, set $X_s=X\times_S\Spec(k(s))$, $(Y_\lambda)_s=Y_\lambda\times_S\Spec(k(s))$, and let $(u_\lambda)_s:(Y_\lambda)_s\to X_s$ be the morphism $u_\lambda\times 1$ induced from $u_\lambda$ by base change.
Let $Z(s)$ be the closure in $X_s$ of the union of the sets $(u_\lambda)_s((Y_\lambda)_s)$, and set $d(s)=\dim(Z(s))$.
Then, for every generization $s'$ of a point $s\in S$, we have $d(s)\leq d(s')$.
\end{proposition}

\begin{proof}
We know from \sref[II]{II.7.1.4} that there exists a valuation ring $A'$ and a morphism $h:S'=\Spec(A')\to S$ such that, if $a$ (resp.\ $b$) is the closed point (resp.\ the generic point) of $S'$, we have $h(a)=s$, $h(b)=s'$.
Furthermore, the projection $p:X_s\otimes_{k(s)}k(a)\to X_s$ is surjective and open \sref{IV.2.4.10}, and thus makes $X_s$ a quotient space of $X_s\otimes_{k(s)}k(a)$ by an open equivalence relation; consequently, for every subset $M$ of $X_s$, $p^{-1}(\overline{M})=\overline{p^{-1}(M)}$ (Bourbaki, \emph{General Topology}, Chap.~I, 4th ed., \S5, no.~3, Prop.~7).
One reasons in the same way for $X_{s'}$; taking account of \sref[I]{I.3.4.8}, \sref{IV.4.2.7}, and the fact that the $g_\lambda$ are universally open, we see that it suffices to prove the proposition after the base change $S'\to S$.
Suppose therefore $S'=S$, $s$ being the closed point and $s'$ the generic point of $S$.
The hypothesis that $g_\lambda$ is open implies that every irreducible component of $Y_\lambda$ dominates $S$ \sref{IV.1.10.4}, hence its generic point is a maximal point of $(Y_\lambda)_{s'}$; if $Z$ denotes the closure in $X$ of $Z(s')$, we have therefore $u_\lambda(Y_\lambda)\subset Z$, and consequently $(u_\lambda)_s((Y_\lambda)_s)\subset Z_s=Z\cap X_s$; in other words, $Z(s)\subset Z_s$, whence $\dim(Z(s))\leq\dim(Z_s)$.
But applying \sref{IV.14.3.10} to a reduced sub-prescheme of $X$ having as underlying space an irreducible component of $Z$, we obtain $\dim(Z_s)\leq\dim(Z_{s'})$ (the inequality coming from the fact that there can be irreducible components of $Z$ not meeting $X_s$).
On the other hand, since $Z(s')$ is closed in $X_{s'}$ by definition, we have $Z_{s'}=Z(s')$, which completes the proof.
\end{proof}

\begin{remark}[14.3.12]
\label{IV.14.3.12}
The case considered by M.\ Artin was the case where $Y_\lambda=S$ for every $\lambda$, in other words the case where $(u_\lambda)$ is a family of $S$-\emph{sections} of $X$.
Another useful case is that where the family $(u_\lambda)$ is reduced to a single element; one can then always reduce to this case by considering the prescheme $Y$ which is the sum of the $Y_\lambda$ and the morphisms $g:Y\to S$ and $u:Y\to X$ whose restrictions to each $Y_\lambda$ are respectively $g_\lambda$ and $u_\lambda$.
\end{remark}

\begin{proposition}[14.3.13]
\label{IV.14.3.13}
Let $f:X\to Y$ be a morphism locally of finite type, $y$ a point of $Y$, $x$ a maximal point of the fibre $X_y=X\times_Y\Spec(k(y))$.
Consider the following conditions:
\begin{enumerate}[label=\emph{\alph*)}]
  \item $f$ is universally open at the point $x$ (or even at every point of the irreducible component of $X_y$ with generic point $x$ {\hyperlink{ega4.c04.subsec.14-3-3-1}{(14.3.3.1, (ii))}}).
  \item For every irreducible component $Y_0$ of $Y$ containing $y$, there exists an irreducible component $Z$ of $X$ containing $x$, dominant over $Y_0$ and such that $\dim_x(X_y)=\dim_x(Z\cap X_y)=\dim(Z\cap X_\eta)$, where $\eta$ is the generic point of $Y_0$ (which implies that $Z$ is \emph{equidimensional} over $Y_0$ at the point $x$ \sref{IV.13.2.2}).
\end{enumerate}

\oldpage[IV-3]{209}

\begin{enumerate}[label=\emph{\alph*$'$)}]
  \item[b$'$)] For every open neighbourhood $U$ of $x$ in $X$ and every generization $y'$ of $y$, we have $\dim(U\cap X_{y'})\geq\dim_x(U\cap X_y)$.
\end{enumerate}

Then we have the implications $a)\Rightarrow b)\Leftrightarrow b')$.
\end{proposition}

\begin{proof}
To show that $b)$ implies $b'$), it suffices to remark that $y'$ belongs to an irreducible component $Y_0$ of $Y$ containing $y$, of generic point $\eta$; taking $Z$ as in $b)$ and noting that the generic point of $Z$ (which is also that of $Z\cap X_\eta$ $(\mathbf{0}_1,~2.1.8)$) is contained in $U$, we have $\dim(U\cap X_{y'})\geq\dim(U\cap Z\cap X_{y'})$, and, by virtue of \sref{IV.13.1.6}, $\dim(U\cap Z\cap X_{y'})\geq\dim(Z\cap X_\eta)$; whence the assertion, since $\dim_x(U\cap X_y)=\dim_x(X_y)$ \sref{IV.4.1.1.3}.

To prove that $b')$ implies $b)$, one can first replace $X$ by $f^{-1}(Y_0)$, hence suppose $Y_0=Y$; we can restrict to the case where $X$ and $Y$ are affine, since $\dim_x(U\cap X_y)=\dim_x(X_y)$ \sref{IV.4.1.1.3}.
The irreducible components $W_i$ of the Noetherian prescheme $X_\eta$ are then finite in number, and the complement of the union of the $\overline{W}_i$ (closures \emph{in} $X$) that do not contain $x$ is an open neighbourhood $V$ of $x$.
Replacing $X$ by $V$, we can therefore suppose that $x\in\overline{W}_i$ for every $i$ (the hypothesis $b'$) implies $\dim(V\cap X_\eta)\geq 0$, hence $x$ belongs to at least one of the $\overline{W}_i$ for at least one $i$).
Furthermore, the $\overline{W}_i$ are exactly the irreducible components of $X$ that dominate $Y$ $(\mathbf{0}_1,~2.1.8)$; if one had, for each of these components $Z$, $\dim(Z\cap X_\eta)<\dim_x(X_y)$, one would conclude $\dim(X_\eta)<\dim_x(X_y)$, contrary to the hypothesis $b'$).
The relation $\dim_x(Z\cap X_y)=\dim(Z\cap X_\eta)$ follows then from \sref{IV.13.1.6}.

It remains to show that $a)$ implies $b'$).
Taking into account \sref[II]{II.7.1.4} and the invariance of the hypotheses and of the conclusion by base change (by virtue of \sref{IV.4.2.7}), we can restrict to the case where $Y$ is the spectrum of a valuation ring, with closed point $y$ and generic point $y'$, and where $U=X$.
The hypothesis that $f$ is open at the point $x$ implies that there exists an irreducible component $Z$ of $X$ containing $x$ and dominant over $Y$ \sref{IV.1.10.3}.
Applying \sref{IV.14.3.10} to a neighbourhood of $x$ in $Z$, we conclude that $\dim(Z\cap X_{y'})=\dim_x(Z\cap X_y)$; but as $x$ is maximal in $X_y$, $Z\cap X_y$ contains the irreducible component of $X_y$ of generic point $x$, hence $\dim_x(Z\cap X_y)=\dim_x(X_y)$; on the other hand, we have $\dim(X_{y'})\geq\dim(Z\cap X_{y'})$, which completes the proof of $b'$).
\end{proof}

\begin{remark}[14.3.14]
\label{IV.14.3.14}
We do not know whether, in \sref{IV.14.3.13}, the conclusion remains valid when one replaces hypothesis $a)$ by the weaker hypothesis that $f$ is \emph{open} at the point $x$.
One could however easily show that it would suffice to treat the case where $Y$ is the spectrum of an integral local ring, whose generic point is isolated, and where $X$ is a closed sub-prescheme of the vector bundle $Y[\mathrm{T}]$.
\end{remark}
\subsection{Chevalley's criterion for universally open morphisms}
\label{subsection:IV.14.4}

\begin{theorem}[14.4.1]
\label{IV.14.4.1}
Let $f:X\to Y$ be a morphism locally of finite type, $y$ a point of $Y$, $x$ a maximal point of the fibre $X_y=X\times_Y\Spec(k(y))$.
Suppose $y$ geometrically unibranch.
Then the following conditions are equivalent:
\begin{enumerate}[label=\emph{\alph*)}]
  \item $f$ is universally open at $x$ (or equivalently at every point of the irreducible component of $X_y$ with generic point $x$ {\hyperlink{ega4.c04.subsec.14-3-3-1}{(14.3.3.1, (iii))}}).
\oldpage[IV-3]{210}
  \item If $Y_0$ is the unique irreducible component of $Y$ containing $y$, $\eta$ its generic point, there exists an irreducible component $Z$ of $X$ containing $x$, dominating $Y_0$ and such that $\dim_x(Z\cap X_y)=\dim(Z\cap X_\eta)$ (which means that $Z$ is equidimensional over $Y_0$ at the point $x$ \sref{IV.13.2.2}).
  \item[\emph{b$'$)}] For every open neighbourhood $U$ of $x$ in $X$ and every generization $y'$ of $y$, we have $\dim(U\cap X_{y'})\geq\dim_x(U\cap X_y)$.
\end{enumerate}
If moreover $Y$ is locally Noetherian, these conditions are also equivalent to the following:
\begin{enumerate}[label=\emph{\alph*)}]
  \item[\emph{c)}] $f$ is open at $x$.
\end{enumerate}
\end{theorem}

Let us first note that since $(\sh{O}_y)_{\mathrm{red}}$ is integral, $y$ belongs to only one irreducible component $Y_0$ of $Y$.
The fact that \emph{b)} and \emph{b$'$)} are equivalent and that \emph{a)} implies \emph{b$'$)} results from \sref{IV.14.3.13}; on the other hand, if $Y$ is locally Noetherian, we have seen in \sref{IV.14.2.3} that \emph{c)} implies \emph{b)}.
It remains to show that when $y$ is geometrically unibranch, \emph{b)} implies \emph{a)}.

\begin{lemma}[14.4.1.1]
\label{IV.14.4.1.1}
Let $Y$ be a prescheme, $Y'=Y[\mathrm{T}_1,\ldots,\mathrm{T}_n]$, $y'$ a point of $Y'$, $y$ its image in $Y$.
If $Y$ is geometrically unibranch at the point $y$, then $Y'$ is geometrically unibranch at the point $y'$.
\end{lemma}

Indeed, for every $y\in Y$, $\Spec(\kappa(y)[\mathrm{T}_1,\ldots,\mathrm{T}_n])$ is geometrically regular over $\kappa(y)$ \sref[0]{0.17.3.7}, the structural morphism $Y'\to Y$ is smooth \sref{IV.6.8.1}, and it suffices to apply \sref{IV.11.3.14}.

\begin{lemma}[14.4.1.2]
\label{IV.14.4.1.2}
Let $A$ be a local integral ring, unibranch, $B$ an integral ring containing and integral over $A$, and $\mathfrak{n}$ a prime ideal of $B$ above the maximal ideal $\mathfrak{m}$ of $A$.
Then the morphism $\Spec(B_\mathfrak{n})\to\Spec(A)$ is surjective, in other words, for every prime ideal $\mathfrak{p}$ of $A$, there exists a prime ideal $\mathfrak{q}$ of $B$ such that $\mathfrak{q}\subset\mathfrak{n}$ and $\mathfrak{q}\cap A=\mathfrak{p}$.
\end{lemma}

Let $K$ (resp.\ $L$) be the field of fractions of $A$ (resp.\ $B$), $A'$ the integral closure of $A$, $B'$ the subring of $L$ generated by $A'$ and $B$, so that we have a commutative diagram of canonical injections
\[
  \xymatrix{
    B \ar[r] & B' \\
    A \ar[u] \ar[r] & A' \ar[u]
  }
\]
As $B'$ is integral over $B$, there exists a prime ideal $\mathfrak{n}'$ of $B'$ such that $\mathfrak{n}'\cap B=\mathfrak{n}$ (Bourbaki, \emph{Alg.\ comm.}, chap.~V, \S~2, n\textsuperscript{o}~1, th.~1), and (for the same reason) $\Spec(A')\to\Spec(A)$ is surjective.
On the other hand, as $A$ is unibranch, $A'$ is a \emph{local} ring, so $\mathfrak{n}'\cap A'$, which is above the maximal ideal $\mathfrak{m}$ of $A$, is necessarily equal to the unique maximal ideal $\mathfrak{m}'$ of $A'$.
By virtue of the second theorem of Cohen--Seidenberg (\emph{loc.\ cit.}, \S~2, n\textsuperscript{o}~4, th.~3) the morphism $\Spec(B'_{\mathfrak{n}'})\to\Spec(A')$ is surjective, so it is the same for the composite $\Spec(B'_{\mathfrak{n}'})\to\Spec(A')\to\Spec(A)$; but this morphism is also the composite $\Spec(B'_{\mathfrak{n}'})\to\Spec(B_\mathfrak{n})\to\Spec(A)$, so the morphism $\Spec(B_\mathfrak{n})\to\Spec(A)$ is surjective.

These lemmas being established, let us return to the proof of the implication $b)\Rightarrow a)$ in \sref{IV.14.4.1}.
By virtue of {\hyperlink{ega4.c04.subsec.14-3-3-1}{(14.3.3.1, (i))}}, it suffices to prove that, for every integer $n\geq 0$ and every point $x'$ of $X'=X[\mathrm{T}_1,\ldots,\mathrm{T}_n]$ above $x$, the morphism
\oldpage[IV-3]{211}
$f':X'\to Y'=Y[\mathrm{T}_1,\ldots,\mathrm{T}_n]$, deduced from $f$ by base change, is open at the point $x'$.
Taking into account the lemma \sref{IV.14.4.1.1}, \sref{IV.2.3.4} and \sref{IV.4.2.7}, we are thus reduced to proving that $f$ is \emph{open} at the point $x$: furthermore, it plainly suffices, by {\hyperlink{ega4.c04.prop.14-1-2}{(14.1.2, (iii))}}, to show that the restriction of $f$ to a closed sub-prescheme of $X$ having $Z$ as underlying space is open at the point $x$, so that we can restrict to the case where $X=Z$ is irreducible.
Replacing $X$ by an open neighbourhood $V$ of $x$ such that $V\cap X_y$ is irreducible, we may suppose, by virtue of \sref{IV.13.3.1}, that the morphism $f$ factorises as $X\xrightarrow{g}Y''=Y[\mathrm{T}_1,\ldots,\mathrm{T}_m]\to Y$, where $g$ is quasi-finite, dominant and locally of finite type.
As the structural morphism $Y''\to Y$ is open \sref{IV.2.4.6}, we are reduced to proving that $g:X\to Y''$ is open at the point $x$.
By \sref{IV.14.4.1.1}, $g(x)$ is a geometrically unibranch point of $Y''$.
We are thus reduced to proving the following lemma:

\begin{lemma}[14.4.1.3]
\label{IV.14.4.1.3}
Let $X$, $Y$ be two irreducible preschemes, $f:X\to Y$ a morphism locally quasi-finite and dominant.
If $x\in X$ is such that $y=f(x)$ is a unibranch point of $Y$, then $f$ is open at the point $x$.
\end{lemma}

It suffices to prove that $f(\Spec(\sh{O}_{X,x}))=\Spec(\sh{O}_{Y,y})$ \sref{IV.1.10.3}.
One can thus restrict, by the base change $\Spec(\sh{O}_{Y,y})\to Y$, to the case where $Y=\Spec(A)$, where $A$ is a local ring and $y$ is the closed point of $Y$ (taking into account \sref[I]{I.3.6.5} and \sref[$0_{\mathrm{I}}$]{0.2.1.8}), which proves that $X\times_Y\Spec(\sh{O}_{Y,y})$ is irreducible; replacing $f$ by $f_{\mathrm{red}}$, we may suppose $X$ and $Y$ reduced, hence integral.
Replacing $X$ and $Y$ eventually by affine neighbourhoods of $x$ and $y$ respectively, we may suppose \sref{IV.8.12.9} that the morphism $f$ factorises as $X\xrightarrow{j}X_1\xrightarrow{g}Y$, where $j$ is an \emph{open immersion} and $g$ a \emph{finite} morphism (evidently dominant); as $X$ and $X_1$ are affine, $j$ is affine, hence separated and quasi-compact, and consequently $j$ factorises as $X\xrightarrow{u}X_2\xrightarrow{h}X_1$, where $X_2$ is the closed image of $X$ by $j$, $h$ the canonical injection and $u$ an open immersion \sref[I]{I.9.5.3}.
In other terms, one can suppose that $X_1$ is integral, or even that $X_1=\Spec(B)$, where $B$ is an integral and finite $A$-algebra containing $A$, since $g$ is dominant. If $\mathfrak{n}$ is the prime ideal of $B$ corresponding to the point $x$, the hypothesis that $A$ is unibranch then implies \sref{IV.14.4.1.2} that the morphism $\Spec(B_\mathfrak{n})\to\Spec(A)$ is surjective, that is, $\Spec(\sh{O}_{X,x})\to\Spec(\sh{O}_{Y,y})$ is surjective.
Q.E.D.

\begin{corollary}[14.4.2]
\label{IV.14.4.2}
Let $f:X\to Y$ be a morphism locally of finite type, $y$ a geometrically unibranch point of $Y$, $\eta$ the generic point of the unique irreducible component $Y_0$ of $Y$ containing $y$.
The following conditions are equivalent:
\begin{enumerate}[label=\emph{\alph*)}]
  \item $f$ is universally open at all points of $X_y$ (or, what amounts to the same, at the maximal points of $X_y$ {\hyperlink{ega4.c04.subsec.14-3-3-1}{(14.3.3.1, (ii))}}).
  \item For every $x\in X_y$, there exists an irreducible component $Z$ of $X$ containing $x$ and equidimensional over $Y$ at the point $x$ \sref{IV.13.2.2}.
  \item[\emph{b$'$)}] For every $x\in X_y$ and every open neighbourhood $U$ of $x$ in $X$, we have
    \[
      \dim(U\cap X_{\eta})\geq\dim_x(U\cap X_y).
    \]
  \item[\emph{b$''$)}] For every open set $U$ of $X$, we have $\dim(U\cap X_{\eta})\geq\dim(U\cap X_y)$.
\end{enumerate}
\end{corollary}
\oldpage[IV-3]{212}
When moreover $Y$ is locally Noetherian, these conditions are also equivalent to the following:
\begin{enumerate}[label=\emph{\alph*)}]
  \item[\emph{c)}] $f$ is open at all points of $X_y$ (or, what amounts to the same, at the maximal points of $X_y$ {\hyperlink{ega4.c04.subsec.14-3-3-1}{(14.3.3.1, (ii))}}).
\end{enumerate}
The equivalence of \emph{a)} and \emph{c)} when $Y$ is locally Noetherian results from \sref{IV.14.4.1}; the conditions \emph{b)} or \emph{b$'$)}, applied to the \emph{maximal} points of $X_y$, imply \emph{a)} by virtue also of {\hyperlink{ega4.c04.subsec.14-3-3-1}{(14.3.3.1, (ii))}}; finally, \emph{b$'$)} and \emph{b$''$)} are equivalent, since
\[
  \dim(U\cap X_y) = \sup_{x\in X_y}(\dim_x(U\cap X_y)).
\]

It remains to see that condition \emph{a)} implies \emph{b)} and \emph{b$'$)} at \emph{every} point $x\in X_y$.
Let $d=\dim_x(X_y)$, and let $x'$ be the generic point of an irreducible component of $X_y$ containing $x$ and of dimension $d$.
By virtue of \emph{a)} and of \sref{IV.14.4.1}, there is an irreducible component $Z$ of $X$ containing $x'$ and equidimensional over $Y$ at the point $x'$, hence such that $\dim_{x'}(Z\cap X_y)=\dim(Z\cap X_\eta)$.
But by construction $\dim_{x'}(Z\cap X_y)=\dim_x(X_y)=d$, and $\dim_x(Z\cap X_y)\leq\dim(Z\cap X_y)=d$; taking into account \sref{IV.13.1.6}, this proves that $Z$ is equidimensional over $Y$ at the point $x$; hence \emph{a)} implies \emph{b)}.
Moreover, we have $\dim(X_\eta)\geq\dim(Z\cap X_\eta)=d=\dim_x(X_y)$.
Replacing $X$ by an open neighbourhood $U$ of $x$, we thus see that \emph{a)} implies \emph{b$'$)}.
Q.E.D.

\begin{corollary}[14.4.3]
\label{IV.14.4.3}
Let $Y$ be a geometrically unibranch prescheme, $f:X\to Y$ a morphism locally of finite type.
The following conditions are equivalent:
\begin{enumerate}[label=\emph{\alph*)}]
  \item $f$ is universally open.
  \item For every open set $U$ of $X$, every $y\in Y$ and every generization $y'$ of $y$, we have
    \[
      \dim(U\cap X_{y'})\geq\dim(U\cap X_y).
    \]
\end{enumerate}
If moreover $Y$ is locally Noetherian, these conditions are also equivalent to the following:
\begin{enumerate}[label=\emph{\alph*)}]
  \item[\emph{c)}] $f$ is open.
\end{enumerate}
\end{corollary}

\begin{corollary}[14.4.4]
\label{IV.14.4.4}
(Chevalley's criterion).
Let $f:X\to Y$ be a morphism locally of finite type.
\begin{enumerate}[label=(\roman*)]
  \item If $f$ is equidimensional at a point $x\in X$ \sref{IV.13.3.2} and if $y=f(x)$ is a geometrically unibranch point of $Y$, $f$ is universally open at the point $x$.
  \item If $Y$ is geometrically unibranch, $f$ is universally open at all the points of $X$ where $f$ is equidimensional, and the set of these points is open in $X$.
In particular, if $f$ is equidimensional, it is universally open.
\end{enumerate}
\end{corollary}

Assertion (ii) is a trivial consequence of (i), since we already know that the set of points where $f$ is equidimensional is open \sref{IV.13.3.2}.
As to assertion (i), it follows from the fact that the hypothesis implies that condition \emph{b)} of \sref{IV.14.4.1} is satisfied at the generic point of an irreducible component of $X_y$ containing $x$, taking into account \sref{IV.13.3.1}; it therefore suffices to apply \sref{IV.14.4.1}.

\begin{remark}[14.4.5]
\label{IV.14.4.5}
One can prove that if $Y$ is locally Noetherian, and if all the generizations of $y=f(x)$ are geometrically unibranch points of $Y$ (cf.\ \sref{IV.6.15.2}), then, if $f$ is equidimensional at the point $x$, it is universally open \emph{in a neighbourhood of $x$}.
\end{remark}

\begin{corollary}[14.4.6]
\label{IV.14.4.6}
\oldpage[IV-3]{213}
Let $Y$ be a locally Noetherian prescheme, $f:X\to Y$ a morphism locally of finite type.
Let $y$ be a geometrically unibranch point of $Y$, and suppose furthermore that for every $x\in f^{-1}(y)$, the ring $\sh{O}_{X,x}$ is equidimensional.
Then the following conditions are equivalent:
\begin{enumerate}[label=\emph{\alph*)}]
  \item $f$ is equidimensional \sref{IV.13.3.2} at all points of $f^{-1}(y)$.
  \item $f$ is open at all points of $f^{-1}(y)$.
  \item $f$ is universally open at all points of $f^{-1}(y)$.
\end{enumerate}
\end{corollary}

Indeed, \emph{c)} implies trivially \emph{b)}, and \emph{a)} implies \emph{c)} by virtue of \sref{IV.14.4.4}; finally, because of the hypotheses on $\sh{O}_y$ and $\sh{O}_x$, \emph{b)} implies \emph{a)} by \sref{IV.14.2.2}.
More generally:

\begin{proposition}[14.4.7]
\label{IV.14.4.7}
Let $Y$ be a locally Noetherian prescheme, $f:X\to Y$ a morphism locally of finite type, $x$ a point of $X$, such that $y=f(x)$ is geometrically unibranch.
The following conditions are equivalent:
\begin{enumerate}[label=\emph{\alph*)}]
  \item $f$ is equidimensional \sref{IV.13.3.2} at the point $x$.
  \item The ring $\sh{O}_x$ is equidimensional and $f$ is open at the generic points of the irreducible components of $f^{-1}(y)$ containing $x$ (and hence also at every point of such a component).
  \item The ring $\sh{O}_x$ is equidimensional, and $f$ is universally open at the generic points of the irreducible components of $f^{-1}(y)$ containing $x$ (and hence also at every point of such a component).
\end{enumerate}
Furthermore, when these conditions are satisfied, for every reduced closed sub-prescheme $X_i$ of $X$ having as underlying space an irreducible component of $X$ containing $x$, the restriction $f_i:X_i\to Y$ of $f$ is an equidimensional morphism at the point $x$, and universally open at all the points of the irreducible components of $f^{-1}(y)\cap X_i$ which contain $x$.
\end{proposition}

Condition \emph{a)} implies that $f$ is equidimensional at all the generic points of the irreducible components of $f^{-1}(y)$ containing $x$ \sref{IV.13.3.1}, and consequently \sref{IV.14.4.4} universally open at these points; the same argument applied to each $f_i$ (taking into account \sref{IV.13.3.3}) proves the last assertion of the proposition, in view of {\hyperlink{ega4.c04.subsec.14-3-3-1}{(14.3.3.1, (iii))}}.
Furthermore, by \sref{IV.14.2.1}, we have the relations
\begin{align}
  \dim(\sh{O}_{X_i,x}) &= \dim(\sh{O}_y) + \dim(\sh{O}_{X_i,x}\otimes_{\sh{O}_y}\kappa(y))
  \tag{14.4.7.1}\label{IV.14.4.7.1}\\
  \dim(\sh{O}_{X,x}) &= \dim(\sh{O}_y) + \dim(\sh{O}_{X,x}\otimes_{\sh{O}_y}\kappa(y))
  \tag{14.4.7.2}\label{IV.14.4.7.2}
\end{align}
and since $f$ is equidimensional at the point $x$, it follows from \sref{IV.13.3.1} that $\dim_x(X_i\cap f^{-1}(y))=\dim_x(f^{-1}(y))$, hence \sref{IV.5.2.3} $\dim(\sh{O}_{X_i,x}\otimes_{\sh{O}_y}\kappa(y))=\dim(\sh{O}_{X,x}\otimes_{\sh{O}_y}\kappa(y))$.
We conclude that $\dim(\sh{O}_{X_i,x})=\dim(\sh{O}_{X,x})$ for every $i$, in other words $\sh{O}_{X,x}$ is equidimensional, and this finishes showing that \emph{a)} implies \emph{c)}.
It is clear that \emph{c)} implies \emph{b)}; finally, \emph{b)} implies the relation \eqref{IV.14.4.7.2} by virtue of \sref{IV.14.2.1}; it then follows from \sref{IV.13.3.6} that \emph{b)} implies \emph{a)}.

\begin{proposition}[14.4.8]
\label{IV.14.4.8}
Let $Y$ be a Noetherian prescheme, $f:X\to Y$ a morphism locally of finite type, $y$ a point of $Y$, $x$ a maximal point of $X_y=f^{-1}(y)$.
The following conditions are equivalent:
\begin{enumerate}[label=\emph{\alph*)}]
  \item The morphism $f$ is universally open at $x$, in other words, for every change of base $g:Y'\to Y$, the property $\mathbf{P}(Y')$ holds: for every point $x'$ of $X'=X\times_Y Y'$ above $x$, the morphism $f'=f_{(Y')}:X'\to Y'$ is open at $x'$.
\oldpage[IV-3]{214}
  \item[\emph{a$'$)}] The property $\mathbf{P}(Y')$ holds for every finite morphism $g:Y'\to Y$.
  \item[\emph{a$''$)}] The property $\mathbf{P}(Y'')$ holds for the normalisation $Y''$ of $Y_{\mathrm{red}}$ \sref[II]{II.6.3.8}.
  \item For every point $x''$ of $X''=X\times_Y Y''$ above $x$, there exists an irreducible component $Z''$ of $X''$ containing $x''$ and equidimensional over $Y''$ at the point $x''$.
\end{enumerate}
\end{proposition}

It is trivial that \emph{a)} implies \emph{a$'$)}.
To show that \emph{a$'$)} implies \emph{a$''$)}, note that we can write $Y''=\Spec(\sh{B})$, where $\sh{B}$ is a quasi-coherent $\sh{O}_Y$-Algebra integral over $\sh{O}_Y$; as $Y$ is Noetherian, $\sh{B}$ is the inductive limit of its sub-$\sh{O}_Y$-Algebras $\sh{B}_\lambda$, which are quasi-coherent and of finite type \sref[I]{I.9.6.6}; but then the $\sh{B}_\lambda$ are \emph{finite} $\sh{O}_Y$-Algebras \sref[II]{II.6.1.2}; we can therefore write $Y''=\varprojlim Y'_\lambda$, where $Y'_\lambda=\Spec(\sh{B}_\lambda)$, and hence $X''=X\times_Y Y''=\varprojlim X'_\lambda$, with $X'_\lambda=X\times_Y Y'_\lambda$.
By virtue of \emph{a$'$)}, the morphisms $f'_\lambda:X'_\lambda\to Y'_\lambda$ are open at all the points of $X'_\lambda$ above $x$; we conclude that $f''$ is open at all the points of $X''$ above $x$, by \sref{IV.8.10.1} and {\hyperlink{ega4.c04.subsec.14-3-3-1}{(14.3.3.1, (i))}}.

As the prescheme $Y''$ is \emph{normal} by definition, the fact that \emph{b)} implies \emph{a$''$)} follows from \sref{IV.14.4.4}, applied to the equidimensional irreducible component in the statement and to the restriction of $f''$ to that component.
It therefore remains to show that \emph{a$''$)} implies \emph{a)} and \emph{b)}.
Taking into account \sref{IV.1.10.3}, one may restrict to the case where $Y=\Spec(\sh{O}_y)$, noting that the canonical morphism $\Spec(\sh{O}_y)\to Y$ is universally bicontinuous \sref[I]{I.3.6.5}, and on the other hand that $Y''\times_Y\Spec(\sh{O}_y)$ is the normalisation of $(\Spec(\sh{O}_y))_{\mathrm{red}}$ as follows from the commutativity of the operations of integral closure and localisation (Bourbaki, \emph{Alg.\ comm.}, chap.~V, \S~1, n\textsuperscript{o}~5, prop.~16).
Suppose therefore that $Y=\Spec(A)$, where $A$ is a local Noetherian ring, and $y$ the closed point of $Y$; we know \sref[0]{0.23.2.5} that there exists a factorisation
\[
  Y'' \xrightarrow{u} Y_1 \xrightarrow{v} Y
\]
of the structural morphism, such that $v$ is a \emph{finite} and surjective morphism, while $u$ is integral, \emph{radicial} and dominant, hence surjective \sref[II]{II.6.1.10} and consequently a universal \emph{homeomorphism} \sref{IV.2.4.5}.
If we set $X_1=X\times_Y Y_1$, the projection $X''\to X_1$ is therefore a homeomorphism, and the hypothesis \emph{a$''$)} implies that $f_1=f_{(Y_1)}:X_1\to Y_1$ is open at all points of $X_1$ above $x$.
Furthermore, this shows that in order to prove property \emph{b)}, it suffices to prove the same property where we replace $Y''$, $X''$ and $x''$ by $Y_1$, $X_1$ and a point $x_1$ of $X_1$ above $x$.
But $Y_1$ is Noetherian and furthermore geometrically unibranch since $Y''$ is normal and $u$ is radicial \sref{IV.6.15.1}; the property \emph{b)} therefore results from \sref{IV.14.4.1}.
It remains to show that $f$ is universally open at the point $x$, which will result from the following lemma:

\begin{lemma}[14.4.8.1]
\label{IV.14.4.8.1}
Let $v:Y_1\to Y$ be a closed (resp.\ universally closed) and surjective morphism.
For a morphism $f:X\to Y$ to be open (resp.\ universally open) at a point $x\in X$, it suffices that $f_1=f_{(Y_1)}:X_1=X\times_Y Y_1\to Y_1$ be open (resp.\ universally open) at all the points of $X_1$ above $x$.
\end{lemma}

The second assertion results trivially from the first and from the fact that for every base change $Y'\to Y$, the morphism $v_{(Y')}:Y_1\times_Y Y'\to Y'$ is still surjective and universally closed.
To prove the first assertion, consider
\oldpage[IV-3]{215}
an open neighbourhood $U$ of $x$ in $X$; as $v$ is closed and surjective, for $f(U)$ to be a neighbourhood of $y=f(x)$, it is necessary and sufficient that $v^{-1}(f(U))$ be a neighbourhood of $v^{-1}(y)$.
But if $p:X_1\to X$ is the canonical projection, we have $v^{-1}(f(U))=f_1(p^{-1}(U))$ \sref[I]{I.3.4.8}, and the hypothesis implies that $f_1(p^{-1}(U))$ is a neighbourhood of $f_1(p^{-1}(x))=v^{-1}(y)$ \sref[I]{I.3.4.8}.

\begin{corollary}[14.4.9]
\label{IV.14.4.9}
Let $Y$ be a Noetherian prescheme, $f:X\to Y$ a morphism locally of finite type.
The following conditions are equivalent:
\begin{enumerate}[label=\emph{\alph*)}]
  \item $f$ is universally open, in other words, for every base change $Y'\to Y$, the morphism $f_{(Y')}:X\times_Y Y'\to Y'$ is open.
  \item[\emph{a$'$)}] For every finite morphism $Y_1\to Y$, $f_{(Y_1)}$ is open.
  \item[\emph{a$''$)}] If $Y''$ is the normalisation of $Y_{\mathrm{red}}$, $f_{(Y'')}$ is open.
  \item For every point $x''$ of $X''=X\times_Y Y''$, there exists an irreducible component $Z''$ of $X''$ containing $x''$ and equidimensional over $Y''$ at the point $x''$ (cf.\ {\hyperlink{ega4.c04.item.14-4-10}{(14.4.10, (ii))}}).
\end{enumerate}
\end{corollary}

This results immediately from \sref{IV.14.4.8} and \sref{IV.14.1.4}.

\begin{remarks}[14.4.10]
\label{IV.14.4.10}
(i) The equivalence of conditions \emph{a)} and \emph{b)} in \sref{IV.14.4.8} (resp.\ \sref{IV.14.4.9}) is also valid for any prescheme $Y$ and any morphism $f$ locally of finite type.
Indeed, \emph{a)} implies \emph{b)} by virtue of \sref{IV.14.4.1}; conversely, \emph{b)} implies that $f''$ is universally open at the points of $X''$ above $x$ (resp.\ at every point of $X''$) by virtue of \sref{IV.14.4.1}, and we conclude property \emph{a)} by applying the lemma \sref{IV.14.4.8.1} to the integral and surjective morphism $Y''\to Y$.

It may be that, in \sref{IV.14.4.1}, for the equivalence of \emph{a)} and \emph{c)}, the supplementary hypothesis that $Y$ is Noetherian is superfluous (cf.\ \sref{IV.14.3.14}).
If so, the Noetherian hypotheses are also superfluous in \sref{IV.14.4.2}, \sref{IV.14.4.3}, \sref{IV.14.4.8} and \sref{IV.14.4.9}.

(ii) One can give examples of morphisms $f:X\to Y$ with the following properties: $Y$ is Noetherian, regular and of dimension~2, $f$ is \emph{universally open} and of finite type, $X$ has two irreducible components $X_1$, $X_2$, but the restriction $X_1\to Y$ of $f$ to one of them is \emph{not an open morphism}.
The principle of the construction relies on the general method of ``recollement'' which will be presented in chap.~V, and which can therefore only be sketched here.
Start with a closed point $y$ of $Y$, and consider the $Y$-scheme $Y_1$ obtained by blowing up $y$ \sref[II]{II.8.1.3}. If $f_1:Y_1\to Y$ is the structural morphism, we know that the restriction of $f_1$ to $f_1^{-1}(Y-\{y\})$ is an \emph{isomorphism} over $Y-\{y\}$ (\emph{loc.\ cit.}), while the fibre $f_1^{-1}(y)$ is isomorphic to $\operatorname{Proj}(S)$, where $S=\bigoplus_{k=0}^{\infty}\mathfrak{m}_y^k/\mathfrak{m}_y^{k+1}$ \sref[II]{II.3.5.3}, i.e.\ here $\mathbf{P}^1_{\kappa(y)}$: it follows from \sref{IV.14.4.1} that $f_1$ is not open at the generic point of $f_1^{-1}(y)$.
On the other hand, let us set $Y_2=\mathbf{P}^1_Y$, and let $f_2:Y_2\to Y$ be the structural morphism. It follows from \sref[II]{II.8.4.4} that $f_2$ is flat, hence universally open \sref{IV.2.4.6}; moreover, by \sref[II]{II.3.5.3}, $f_2^{-1}(y)$ is isomorphic to $\mathbf{P}^1_{\kappa(y)}$. It then suffices to ``glue'' $Y_1$ and $Y_2$ along the isomorphic fibres $f_1^{-1}(y)$ and $f_2^{-1}(y)$; this gives a morphism $f:X\to Y$ whose irreducible components $X_1$, $X_2$ identify canonically with $Y_1$, $Y_2$, respectively, and whose restrictions to these components are $f_1$, $f_2$.

Let us recall however \sref{IV.12.1.1.5} that if $Y$ is locally Noetherian, $f:X\to Y$ is of finite type and \emph{flat}, then every irreducible component of $X$ is equidimensional over $Y$ (and consequently the restriction of $f$ to such a component is \emph{universally open} if all the points of $Y$ are \emph{geometrically unibranch}).

(iii) Recall {\hyperlink{ega4.c04.item.12-1-2}{(12.1.2, (i))}} that there are morphisms $f:X\to Y$ with the following properties: $Y$ is Noetherian (not geometrically unibranch), $f$ is finite and flat (and even \'etale \sref{IV.17.6.3}), but the restriction of $f$ to an irreducible component of $X$ is \emph{not an open morphism} (although $f$ itself is so by \sref{IV.2.4.6}).

(iv) The Chevalley criterion \sref{IV.14.4.4} explains the importance of the notion of universally open morphism.
This notion allows in effect, in numerous results, more or less classical, to replace a hypothesis of \emph{normality} by the hypothesis that a certain morphism is universally open; the more general statement
\oldpage[IV-3]{216}
thus obtained will in particular apply to \emph{flat} morphisms \sref{IV.2.4.6}, whose importance in algebraic geometry is growing.
One can consider that the statements involving the hypothesis that a morphism is universally open are common generalisations of statements involving a hypothesis of normality and of statements involving a hypothesis of flatness.
\end{remarks}


\subsection{Universally open morphisms and quasi-sections}
\label{subsection:IV.14.5}
\label{IV.14.5}

\begin{lemma}[14.5.1]
\label{IV.14.5.1}
Let $Y$ be an irreducible locally Noetherian prescheme, $f:X\to Y$ a morphism locally of finite type, $x$ a point of $X$, such that $f$ is equidimensional \sref{IV.13.3.2} at the point $x$; set $y=f(x)$, $e=\dim_x(f^{-1}(y))$.
Let $X'$ be a closed irreducible subspace of $X$ containing $x$, $n$ an integer such that $\codim(X',X)\leq n$ and $\dim_x(X'\cap f^{-1}(y))\leq e-n$.
Then necessarily $\codim(X',X)=n$, $\dim_x(X'\cap f^{-1}(y))=e-n$, and the restriction $X'\to Y$ of $f$ is an equidimensional morphism at $x$ (and a fortiori dominant).
\end{lemma}

The question being local on $X$, we can suppose that $f$ is of finite type and \emph{equidimensional}, so that for \emph{every} $z\in Y$, \emph{all} the irreducible components of $f^{-1}(z)$ are of dimension $e$ {\hyperlink{ega4.c04.prop.13-3-1}{(13.3.1, $a''$))}}.
Let $x'$ be the generic point of $X'$, $y'=f(x')$, $Y'=\overline{\{y'\}}=\overline{f(X')}$ and set $Z=f^{-1}(Y')$.
By virtue of \sref[0]{0.14.2.2}, we have
\begin{equation}
  \codim(X',Z)\leq n
  \tag{14.5.1.1}\label{IV.14.5.1.1}
\end{equation}
and if the two sides are equal, we necessarily have $\codim(X',X)=n$ and $Z$ contains an irreducible component of $X$, which implies \sref{IV.13.3.1} that $f(X')$ is dense in $Y$ and hence $Z=X$; thus the equality $\codim(X',Z)=n$ is equivalent to the conjunction of the equality $\codim(X',X)=n$ and the relation $Z=X$.

On the other hand, reasoning in the reduced preschemes of $Y$ and $X$ having $Y'$ and $Z$ respectively as underlying spaces, we deduce from \sref{IV.5.1.2} and from \sref[I]{I.3.6.5} that
\begin{equation}
  \codim(X'\cap f^{-1}(y'),f^{-1}(y'))=\codim(X',Z).
  \tag{14.5.1.2}\label{IV.14.5.1.2}
\end{equation}
By hypothesis, $f^{-1}(y')$ is biequidimensional \sref{IV.5.2.1} and of dimension $e$, hence \sref[0]{0.14.3.5}, by virtue of \eqref{IV.14.5.1.2} and \eqref{IV.14.5.1.1}
\begin{equation}
  e'=\dim(X'\cap f^{-1}(y'))=e-\codim(X',Z)\geq e-n
  \tag{14.5.1.3}\label{IV.14.5.1.3}
\end{equation}
the equality holding if and only if $\codim(X',X)=n$ and if $f(X')$ is dense in $Y$.
Finally, \sref{IV.13.1.1}, we have $\dim_x(X'\cap f^{-1}(y))\geq e'$, whence, by \eqref{IV.14.5.1.3},
\begin{equation}
  \dim_x(X'\cap f^{-1}(y))\geq e'\geq e-n.
  \tag{14.5.1.4}\label{IV.14.5.1.4}
\end{equation}
But, by hypothesis, we also have $\dim_x(X'\cap f^{-1}(y))\leq e-n$, whence the conclusions of the proposition.

\begin{corollary}[14.5.2]
\label{IV.14.5.2}
The hypotheses on $f$, $X$, $Y$, $x$ being those of \sref{IV.14.5.1}, suppose moreover that $x$ is not a maximal point of $f^{-1}(y)$.
Then there exists an affine open neighbourhood $U$ of $x$ in $X$ and a section $g\in\Gamma(U,\sh{O}_X)$ such that the set $X'$ of $x'\in U$ such that $g(x')=0$ contains $x$ and does not contain any maximal point of $f^{-1}(y)$.
For every $g$ having these properties, $X'$ is equidimensional over $Y$ at the point $x$, and we have
\begin{equation}
  \dim_x(X'\cap f^{-1}(y))=e-1 \qquad\text{and}\qquad \codim(X',X)=1.
  \tag{14.5.2.1}\label{IV.14.5.2.1}
\end{equation}
\end{corollary}

\oldpage[IV-3]{217}
One can restrict to the case where $X=U$ is an open affine neighbourhood of $x$ such that all the irreducible components of $f^{-1}(y)$ contain $x$.
These components correspond to the minimal prime ideals of $\sh{O}_x/\mathfrak{m}_y\sh{O}_x$, and by hypothesis these ideals are distinct from $\mathfrak{m}_x/\mathfrak{m}_y\sh{O}_x$ (Bourbaki, \emph{Alg.\ comm.}, chap.~II, \S~1, n\textsuperscript{o}~1, prop.~2); to obtain a $g\in\Gamma(U,\sh{O}_X)$ satisfying the conditions of the statement, it suffices to take $g\in\mathfrak{j}_x$ such that the image of $g$ in $\mathfrak{m}_x$ does not belong to any of the preceding prime ideals.
Moreover, as $\codim(X',X)\leq 1$ \sref{IV.5.1.8}, and as $X'$ does not contain any of the irreducible components of $f^{-1}(y)$ of dimension $e$, and since these are all the irreducible components of $f^{-1}(y)$, we have \sref[0]{0.14.2.2.2} $\dim_x(X'\cap f^{-1}(y))\leq e-1$.
It then suffices to apply \sref{IV.14.5.1}.

\begin{proposition}[14.5.3]
\label{IV.14.5.3}
Let $Y$ be an irreducible locally Noetherian prescheme, $f:X\to Y$ a morphism locally of finite type, $x$ a point of $X$.
Suppose that $f$ is equidimensional at the point $x$ and that $x$ is closed in $f^{-1}(f(x))$.
Then there exists an irreducible subset $X'$ of $X$, locally closed in $X$, containing $x$, such that the restriction $X'\to Y$ of $f$ (where $X'$ is endowed with the reduced subscheme structure induced from $X$) is a \emph{quasi-finite and dominant} morphism.
\end{proposition}

Indeed, with the notation of \sref{IV.14.5.2}, the hypothesis that $x$ is closed in $f^{-1}(y)$ implies that $x$ is not a maximal point of $X'\cap f^{-1}(y)$ as long as $e-1\geq 1$.
It therefore suffices to apply \sref{IV.14.5.2} by descending induction on $e=\dim_x(f^{-1}(f(x)))$ until one arrives at $e=1$; the application of \sref{IV.14.5.2} in this last case gives a subspace $X'$ such that $X'\cap f^{-1}(f(x))$ is Noetherian and of dimension $0$, hence finite and discrete; as $X'$ is equidimensional over $Y$, $X'\cap f^{-1}(f(x'))$ is of dimension $0$ for every $x'\in X'$, which implies that the restriction $X'\to Y$ of $f$ is a quasi-finite morphism \sref[II]{II.6.2.2}.

\begin{corollary}[14.5.4]
\label{IV.14.5.4}
Under the hypotheses of \sref{IV.14.5.3}, suppose moreover that $Y=\Spec(A)$, where $A$ is a local Noetherian ring, integral and complete, and that $y=f(x)$ is the unique closed point of $Y$.
Then there exists a local integral ring $A'$, containing $A$, which is a finite $A$-algebra, and such that, if we set $Y'=\Spec(A')$ and $X'=X\times_Y Y'$, there exists a $Y'$-section $h:Y'\to X'$ such that the composite morphism $Y'\xrightarrow{h}X'\to X$ is an immersion whose image contains $x$.
\end{corollary}

Replacing if necessary $X$ by a reduced closed irreducible sub-prescheme of $X$, we may, by virtue of \sref{IV.14.5.3}, restrict to the case where the morphism $f$ is already quasi-finite and dominant.
Using \sref[II]{II.6.2.5}, we deduce that $\sh{O}_x=A'$ is an integral ring which is a finite $A$-algebra, and that $X$ is a disjoint sum of the closed sub-prescheme $Y'=\Spec(A')$ and of a sub-prescheme $Y''$; the scheme $Y'$ satisfies the question, the composite morphism $Y'\to X'\to X$ being nothing other than the canonical morphism $\Spec(\sh{O}_x)\to X$.

\begin{remark}[14.5.5]
\label{IV.14.5.5}
If, in the statement of \sref{IV.14.5.4}, one does not require that the morphism $Y'\to X$ be an immersion, one can suppose that $A'$ is moreover \emph{integrally closed}: it suffices in effect to replace $A'$ by its \emph{integral closure} $A_1$, since we know \sref[0]{0.23.1.5} that $A_1$ is an $A$-module of finite type.
\end{remark}

\begin{proposition}[14.5.6]
\label{IV.14.5.6}
Let $Y$ be a locally Noetherian prescheme, irreducible, regular and of dimension $1$, $f:X\to Y$ a morphism locally of finite type, $y$ a point of $Y$.
The following conditions are equivalent:
\begin{enumerate}[label=\emph{\alph*)}]
  \item $f_{\mathrm{red}}$ is flat at every point of $f^{-1}(y)$.
\oldpage[IV-3]{218}
  \item $f$ is universally open in a neighbourhood of $f^{-1}(y)$.
  \item $f$ is open in a neighbourhood of $f^{-1}(y)$.
  \item Every irreducible component of $X$ meeting $f^{-1}(y)$ dominates $Y$.
  \item For every point $x\in f^{-1}(y)$, closed in $f^{-1}(y)$, and every irreducible component $X_i$ of $X$ containing $x$, there exists an irreducible subset $X'$ of $X$, locally closed in $X$, containing $x$, contained in $X_i$, and such that the restriction $X'\to Y$ of $f$ is a quasi-finite and dominant morphism.
\end{enumerate}
\end{proposition}

The equivalence of \emph{a)}, \emph{b)}, \emph{c)} and \emph{d)} has already been shown \sref{IV.14.3.8}.
It is clear that \emph{e)} implies \emph{d)}, since for every irreducible component $X_i$ of $X$ meeting $f^{-1}(y)$, there exists in $X_i\cap f^{-1}(y)$ a closed point \sref{IV.5.1.11}.
Finally, to prove that \emph{c)} implies \emph{e)}, we may restrict to the case where $f$ is a morphism of finite type; consider a closed point $x$ of $f^{-1}(y)$ and let $X_i$ be an irreducible component of $X$ containing $x$.
As the restriction $f_i:X_i\to Y$ of $f$ to $X_i$ is a dominant morphism, it follows from the equivalence of \emph{c)} and \emph{d)} for $f_i$ that this morphism is open at the generic points of $X_i\cap f^{-1}(y)$. It follows from \sref{IV.14.2.2} that $f_i$ is equidimensional at the point $x$, and we conclude with the aid of \sref{IV.14.5.3}.

\begin{remark}[14.5.7]
\label{IV.14.5.7}
If, in the statement of \sref{IV.14.5.5}, we suppose that $Y=\Spec(A)$, where $A$ is a complete discrete valuation ring, and that $y$ is the closed point of $Y$, we can moreover suppose that $X'=\Spec(A')$, where $A'$ is a discrete valuation ring which is a finite $A$-algebra, as follows from the proof of \sref{IV.14.5.4} and from the fact that a local integral ring, regular and of dimension $1$, is a discrete valuation ring \sref[II]{II.7.1.6}.
\end{remark}

\begin{proposition}[14.5.8]
\label{IV.14.5.8}
Let $Y$ be a locally Noetherian prescheme, $f:X\to Y$ a morphism locally of finite type, $y$ a point of $Y$.
For every $Y$-prescheme $Y'$, set $X(Y')=\Hom_Y(Y',X)$.
Let $x$ be a point of $f^{-1}(y)$; in order that $f$ be universally open at $x$, it is necessary and sufficient that the following condition be satisfied:

For every complete discrete valuation ring $A$ with algebraically closed residue field, every morphism $g:Y'=\Spec(A)\to Y$ such that the image under $g$ of the closed point $y'$ of $Y'$ is equal to $y$, and every element $u_0\in X(\Spec(\kappa(y')))$ such that $u_0(y')=x$, there exist a discrete valuation ring $B$, a local homomorphism $A\to B$ making $B$ into a finite $A$-algebra, and, if we set $Z'=\Spec(B)$, an element $u\in X(Z')$ such that, if $z'$ is the closed point of $Z'$, the diagram
\[
  \xymatrix{
    \Spec(\kappa(z')) \ar[r] \ar[d]^{\ell} & Z' \ar[d]^{u} \\
    \Spec(\kappa(y')) \ar[r]_-{u_0} & X
  }
\]
is commutative.
\end{proposition}

\oldpage[IV-3]{219}
Let us note that if $A$ and $B$ satisfy the conditions of the statement, $B$ is a complete discrete valuation ring (Bourbaki, \emph{Alg.\ comm.}, chap.~III, \S~3, n\textsuperscript{o}~3, prop.~7 and chap.~IV, \S~2, n\textsuperscript{o}~2, cor.~3 of prop.~9) whose residue field is isomorphic to that of $A$, and hence algebraically closed. Moreover, since $u(z')=x$, there is in $X''=X\times_Y Z'$ a point $x''$ whose projections in $X$ and $Z'$ are $x$ and $z'$, respectively, and which is rational over $\kappa(z')$; furthermore, since there exists a $Z'$-section $v$ of $X''$ such that $v(z')=x''$, the image under $v$ of the generic point $s$ of $Z'$ is a generization $t$ of $x''$ whose projection in $Z'$ is $s$. Applying \sref{IV.14.3.6}, we see that the condition of the statement is \emph{sufficient}.
Let us now prove that it is \emph{necessary}.
Set $X'=X\times_Y Y'$, $f'=f_{(Y')}:X'\to Y'$.
By hypothesis there is a point $x'\in X'$ above $x$ and $y'$ and rational over $\kappa(y')$ \sref[I]{I.3.3.14}, hence closed in $f'^{-1}(y')$.
By virtue of \sref{IV.14.3.6}, there is an irreducible component $T'$ of $X'$ containing $x'$ and dominant over $Y'$; hence, by \sref{IV.14.3.8} and \sref{IV.14.3.13}, the restriction $T'\to Y'$ of $f'$ is equidimensional at $x'$.
It follows from \sref{IV.14.5.5} that there is a finite $A$-algebra $B$ which is an integrally closed local domain dominating $A$ (and hence is a discrete valuation ring) and, setting $Z'=\Spec(B)$ and $X''=X\times_Y Z'=X'\times_{Y'}Z'$, a $Z'$-section $v:Z'\to X''$ such that the image of the closed point $z'$ of $Z'$ under the composite $Z'\xrightarrow{v}X''\to X'$ is $x'$; the composite morphism $u:Z'\xrightarrow{v}X''\to X$ therefore answers the question.

\begin{proposition}[14.5.9]
\label{IV.14.5.9}
Let $Y$ be an irreducible locally Noetherian prescheme, $f:X\to Y$ a morphism locally of finite type, and $y$ a geometrically unibranch point of $Y$.
Then the following conditions are equivalent:
\begin{enumerate}[label=\emph{\alph*)}]
  \item $f$ is universally open at every point of $f^{-1}(y)$.
  \item $f$ is open at every point of $f^{-1}(y)$.
  \item For every irreducible component $Z$ of $f^{-1}(y)$, with generic point $z$, there exists an irreducible component $X_i$ of $X$ containing $z$ and equidimensional over $Y$ at $z$.
  \item For every closed point $x$ of $f^{-1}(y)$, there exists an irreducible subset $X'$ of $X$, locally closed in $X$, containing $x$, and such that the restriction $X'\to Y$ of $f$ is a quasi-finite dominant morphism.
\end{enumerate}
\end{proposition}

The equivalence of \emph{a)}, \emph{b)} and \emph{c)} has already been shown \sref{IV.14.4.2}.
To prove that \emph{a)} implies \emph{d)}, note that by \sref{IV.14.4.2}, \emph{a)} implies that there exists an irreducible component $Z$ of $X$ containing $x$ and equidimensional over $Y$ at the point $x$; the existence of $X'$ then comes from \sref{IV.14.5.3} applied to the restriction $Z\to Y$ of $f$.
Conversely, suppose \emph{d)} is satisfied; by the Chevalley criterion \sref{IV.14.4.4}, the restriction $X'\to Y$ of $f$ is universally open at $x$, and \emph{a fortiori} $f$ is open at $x$. Thus $f$ is open at every closed point of $X_y=f^{-1}(y)$.
But $X_y$ is a $\kappa(y)$-prescheme locally of finite type, hence a Jacobson prescheme; the set of closed points of $X_y$ is therefore \emph{dense} in $X_y$ \sref{IV.10.3.1}, and it follows from \sref{IV.14.1.4} that $f$ is open at \emph{all} points of $X_y$, which finishes proving that \emph{d)} implies \emph{b)}.

The following result was brought to light by D.\ Mumford:

\begin{proposition}[14.5.10]
\label{IV.14.5.10}
\oldpage[IV-3]{220}
Let $Y$ be a Noetherian prescheme, $f:X\to Y$ a morphism universally open, surjective and locally of finite type.
Then there exists a finite surjective morphism $g:Y'\to Y$ such that, if we set $X'=X\times_Y Y'$ and $f'=f_{(Y')}:X'\to Y'$, every point $y'\in Y'$ admits an open neighbourhood $U'$ such that there exists a $U'$-section of $f'^{-1}(U')$.
\end{proposition}

We will prove the proposition in several steps.

\emph{I) Reduction to the case where $Y$ is integral.}
If we have proved the proposition for each of the reduced sub-preschemes $Y_i$ having as underlying space an irreducible component of $Y$, and for the inverse images $f^{-1}(Y_i)$, it is clear that the prescheme $Y'$ which is the \emph{sum} of the corresponding $Y'_i$ will answer the question.
We can therefore suppose $Y$ \emph{integral} and what follows will be limited to this case.
Furthermore, in the conclusion, we will also be able to take $Y'$ integral (up to replacing it by a suitable irreducible component).

\emph{II) Local character on $Y$.}
We will show that if $Y$ can be covered by finitely many open sets $U_j$ such that, for every $j$, the conclusion of the proposition holds for the morphism $f^{-1}(U_j)\to U_j$, the restriction of $f$, then the conclusion also holds for $f$.
Indeed, we can evidently suppose the $U_j$ affine, so that $U_j=\Spec(A_j)$, where $A_j$ is an integral Noetherian ring whose field of fractions $K=R(Y)$ is the field of rational functions on $Y$.
For every $j$, by hypothesis there is a finite integral $A_j$-algebra $A'_j$ such that the homomorphism $A_j\to A'_j$ is injective \sref[I]{I.1.2.7} and that the corresponding morphism $g_j:U'_j=\Spec(A'_j)\to\Spec(A_j)=U_j$ satisfies the conditions of the statement (for $U_j$ and $f^{-1}(U_j)$).
Let $K'$ be a finite extension of $K$ containing the fields of fractions of all the $A'_j$ (which are finite extensions of $K$).
Consider the normalisation $Y''$ of $Y$ in $K'$ \sref[II]{II.6.3.8}; it is of the form $\Spec(\sh{B})$, where $\sh{B}$ is an integral quasi-coherent $\sh{O}_Y$-Algebra, the integral closure of $\sh{O}_Y$ in $K'$ \sref[II]{II.6.3.4}.
These definitions show that, for every $j$,
\[
 A'_j=\Gamma(U'_j,\sh{O}_{U'_j})
     =\Gamma\bigl(U_j,(g_j)_*(\sh{O}_{U'_j})\bigr)
\]
identifies with a finite sub-$A_j$-algebra of $\Gamma(U_j,\sh{B})$; in other words,
\[
 (g_j)_*(\sh{O}_{U'_j})=\sh{A}'_j
\]
is a coherent $\sh{O}_{U_j}$-Algebra, a sub-Algebra of $\sh{B}|U_j$.
There consequently exists a coherent sub-$\sh{O}_Y$-Module $\sh{C}'_j$ of $\sh{B}$ such that $\sh{C}'_j|U_j=\sh{A}'_j$ \sref[I]{I.9.4.7}.
If we set $\sh{C}'=\sum_j\sh{C}'_j$, the sub-$\sh{O}_Y$-Algebra $\sh{B}'$ of $\sh{B}$ generated by $\sh{C}'$ is coherent, since $\sh{B}$ is an integral $\sh{O}_Y$-Algebra.
We now show that $Y'=\Spec(\sh{B}')$ answers the question.
Indeed, it is clear that the morphism $g:Y'\to Y$ is finite and surjective and that $Y'$ is integral; moreover, for every $j$, $\Gamma(U_j,\sh{B}')$ is a finite $A'_j$-algebra.
In other words, the restriction $g^{-1}(U_j)\to U_j$ of $g$ factors as
\[
 g^{-1}(U_j)\longrightarrow U'_j\longrightarrow U_j,
\]
and since the local existence of sections is stable under base change, this proves our assertion, every $y\in Y$ belonging to some $U_j$.

\emph{III) Reduction to the case where $Y$ is integral, local and geometrically unibranch.}
Suppose first that the proposition has been proved when $Y$ is integral and \emph{local} (with $Y'$ integral), and show that it is valid when $Y$ is integral (Noetherian) affine.
Indeed, by virtue of the reduction II), it suffices to prove that for every point $y\in Y$, the proposition holds for an open affine neighbourhood $V$ of $y$ in $Y$.
Let $Y=\Spec(A)$, $\mathfrak{p}=\mathfrak{j}_y$, let $Y_1=\Spec(A_\mathfrak{p})$, and set $X_1=X\times_Y Y_1$.
By hypothesis, there exists a finite surjective morphism $g_1:Y'_1\to Y_1$, where $Y'_1=\Spec(B_1)$ and $B_1$ is a finite integral $A_\mathfrak{p}$-algebra, hence a semilocal ring, such that $g_1$ satisfies the conditions of the statement for $Y_1$ and $X_1$.
If $y'_j$ ($1\leq j\leq r$) are the closed points of $Y'_1$, there is therefore a covering of $Y'_1$ by open sets $U'_j$ such that $y'_j\in U'_j$ and there exists a $U'_j$-section $h'_j$ of $X\times_Y U'_j$ ($1\leq j\leq r$).
The $A_\mathfrak{p}$-module $B_1$ admits a finite system of generators of the form $z_i/s$ (with $s\in A-\mathfrak{p}$,
\oldpage[IV-3]{221}
$z_i$ integral over $A$), which one may suppose (multiplying $s$, if necessary, by an element of $A$) to be elements of the field of fractions of $B_1$ integral over $A_s$, so that if $V$ is the open affine set $D(s)=\Spec(A_s)\subset Y$, $Y'_1$ identifies with $Y_1\times_V V'$, where $V'$ is the spectrum of the finite $A_s$-algebra generated by the $z_i/s$; $g:V'\to V$ is therefore a finite surjective morphism and $g_1=g_{(Y_1)}$.
Moreover, applying the method of \sref{IV.8.1.2},~$a$), one may suppose that each $U'_j$ is the inverse image under $p:Y'_1\to V'$ of an open set $W'_j$ of $V'$, that the $W'_j$ cover $V'$ \sref{IV.8.3.11}, and that each section $h'_j$ is of the form $(v'_j)_{(Y_1)}$, where $v'_j$ is a $W'_j$-section of $X\times_Y W'_j$ \sref{IV.8.8.2},~(i).
We are thus indeed reduced to proving the proposition when $Y=\Spec(A)$, where $A$ is a local integral Noetherian ring.
We then know that there exists a finite integral $A$-algebra $B$, having the same field of fractions as $A$, such that $A\subset B$ and $\Spec(B)$ is geometrically unibranch \sref[0]{0.23.2.5} and \sref{IV.6.15.5}; as the morphism $\Spec(B)\to\Spec(A)$ is surjective, one may evidently replace $Y$ by $\Spec(B)$ and $X$ by $X\otimes_A B$ to prove the proposition.
Arguing as at the beginning of the reduction III), we may therefore suppose that $A$ is local and integral and that $Y=\Spec(A)$ is geometrically unibranch.

\emph{IV) Reduction to the case where $X$ is integral, affine, and $f$ quasi-finite, surjective, birational and universally open.}
Suppose therefore $Y=\Spec(A)$ integral, local and geometrically unibranch.
There exists then an irreducible sub-prescheme $X_0$ of $X$ such that the restriction $f_0:X_0\to Y$ of $f$ is a quasi-finite and \emph{dominant} morphism and $f_0(X_0)$ contains the closed point $y$ of $Y$ \sref{IV.14.5.9}; as $Y$ is geometrically unibranch, it follows from \sref{IV.14.4.1} that $f_0$ is still \emph{universally open}.
Furthermore, we may suppose $X_0$ reduced, hence integral, and we see that, by replacing $X$ by $X_0$, we may suppose $X$ integral, $f$ \emph{quasi-finite and dominant}, and $f(X)$ containing $y$; since $f$ is \emph{open} and every open subset of $Y$ containing $y$ is equal to $Y$, $f$ is surjective.

Let $\xi$, $\eta$ be the generic points of $X$ and $Y$ respectively; $\kappa(\xi)$ is therefore a finite extension of $K=\kappa(\eta)$.
By \sref{IV.4.6.8}, there is a \emph{finite} extension $K'$ of $K$ such that $(\Spec(\kappa(\xi)\otimes_K K'))_{\mathrm{red}}$ is geometrically reduced over $K'$ and its irreducible components are geometrically irreducible; it follows from \sref{IV.4.5.9} and \sref{IV.4.6.1} that the residue fields of $\Spec(\kappa(\xi)\otimes_K K')$ are finite extensions of $K'$ which are primary and separable, and are therefore equal to $K'$.
Applying once more \sref[0]{0.23.2.5} and \sref{IV.6.15.5}, there is a finite integral $A$-algebra $B$, having $K'$ as field of fractions, containing $A$ and such that $\Spec(B)$ is geometrically unibranch; as the morphism $\Spec(B)\to\Spec(A)$ is surjective, one may therefore again replace $Y$ by $\Spec(B)$ and $X$ by $(X\otimes_A B)_{\mathrm{red}}$ in order to prove the proposition; the reduction III) allows then to suppose once more $A$ local, integral and $Y=\Spec(A)$ geometrically unibranch, with closed point $y$.
Furthermore, $f$ is still a quasi-finite, \emph{surjective} and \emph{universally open} morphism; each of the irreducible components $X_i$ of $X$ ($1\leq i\leq m$) dominates $Y$ \sref{IV.1.10.4} and is \emph{birational} over $Y$.
So let $x$ be a point of $f^{-1}(y)$ and let $X_i$ be an irreducible component of $X$ containing $x$; denote also by $X_i$ the reduced sub-prescheme (hence integral) of $X$ having $X_i$ as underlying space, and let $f_i:X_i\to Y$ be the restriction of $f$.
Applying once more \sref{IV.14.4.1} to the quasi-finite and dominant morphism $f_i$, we see that $f_i$ is universally open; if $U$ is an affine open set of $X_i$ containing $x$, $f_i(U)$ is therefore
\oldpage[IV-3]{222}
open in $Y$ and contains $y$, hence is equal to $Y$.
One may thus replace $X$ by $U$ in order to prove the proposition.

\emph{V) End of the proof.}
We suppose therefore $Y$ local, integral and geometrically unibranch, $X$ integral and affine, $f$ quasi-finite, surjective, birational and universally open; it follows that $f$ is automatically \emph{separated}.
By virtue of the Main theorem \sref{IV.8.12.6}, there exists a factorisation $X\xrightarrow{j}Z\xrightarrow{u}Y$, where $j$ is an open immersion and $u$ a \emph{finite} morphism.
Replacing moreover $Z$ by the closed image of $j$ \sref[I]{I.9.5.10}, we may suppose that $Z$ is \emph{integral}; as $u$ is finite and birational and $Y$ geometrically unibranch, it follows from \sref[III]{III.4.3.5} and \sref[III]{III.4.3.4} that $u$ (and consequently $f$) is a \emph{radicial} morphism; as $f$ is universally open and surjective, it is therefore a universal \emph{homeomorphism}.
Consequently, by \sref{IV.2.4.5},~(ii), $f$ is a \emph{finite} morphism.
But then we satisfy the conditions of the statement with $Y'$ \emph{integral} by simply taking $Y'=X$ and $g=f$, the $Y'$-section being the diagonal morphism $\Delta_f$.

This result can be used to develop criteria of ``descent'' for the properties of universally open and surjective morphisms.
Let us point out in particular the following criterion, due to D.\ Mumford:

\begin{corollary}[14.5.11]
\label{IV.14.5.11}
Let $Y$ be a locally Noetherian prescheme, $f:X\to Y$ a morphism locally of finite type, $g:Y'\to Y$ a morphism locally of finite type, universally open and surjective; set $X'=X\times_Y Y'$, $f'=f_{(Y')}:X'\to Y'$.
Then, for $f$ to be affine, it is necessary and sufficient that $f'$ be so.
\end{corollary}

One need only prove that the condition is sufficient.
The question being local on $Y$, we may suppose $Y$ affine, hence Noetherian.
By virtue of \sref{IV.14.5.10}, there exists a finite surjective morphism $h:Y_1\to Y$ such that, if we set $Y'_1=Y'\times_Y Y_1$, and $g'=g_{(Y_1)}:Y'_1\to Y_1$, every point of $Y_1$ admits an open neighbourhood $U_1$ such that there exists a $U_1$-section of $g'^{-1}(U_1)$.
If we set $X_1=X\times_Y Y_1$, the canonical projection $p:X_1\to X$ is a finite surjective morphism, hence, if we prove that $X_1$ is an affine scheme, it will follow first of all that $X$ is quasi-compact, hence Noetherian, then that $X$ is affine by virtue of the theorem of Chevalley \sref[II]{II.6.7.1}.
It suffices therefore to prove that the morphism $f_1=f_{(Y_1)}:X_1\to Y_1$ is affine.
Or, if we set
\[
  X'_1=X_1\times_{Y_1}Y'_1=X'\times_{Y'}Y'_1 \qquad\text{and}\qquad f'_1=(f_1)_{(Y'_1)}=f'_{(Y'_1)}:X'_1\to Y'_1,
\]
$f'_1$ is affine by virtue of the hypothesis.
We are therefore reduced to proving the corollary by replacing $Y$, $X$, $f$ and $Y'$ by $Y_1$, $X_1$, $f_1$ and $Y'_1$; in other words, it suffices to prove that $f$ is affine when, in the statement of \sref{IV.14.5.11}, one assumes in addition that every point of $Y$ admits an open neighbourhood $U$ such that there exists a $U$-section of $g^{-1}(U)$.
The question being local on $Y$, one may even suppose that there exists a $Y$-section $s$ of $Y'$.
Now we have the following elementary lemma (valid in any category admitting fibre products):

\begin{lemma}[14.5.11.1]
\label{IV.14.5.11.1}
Let $f:X\to S$, $g:Y\to S$ be two morphisms, $p_1:X\times_S Y\to X$, $p_2:X\times_S Y\to Y$ the canonical projections.
If $s:S\to Y$ is a section of $g$, then $s'=(1,s\circ f)_S$ is a section of $p_1$, and $X$, equipped with the morphisms $f:X\to S$ and $s':X\to X\times_S Y$, identifies with the product of the $Y$-preschemes $S$ and $X\times_S Y$ for the morphisms $s:S\to Y$ and $p_2:X\times_S Y\to Y$.
\end{lemma}

\begin{proof}
This is a particular case of \sref[I]{I.3.3.11}, where one replaces the diagram by
\[
\xymatrix{
  X \ar[r]^-{s'} \ar[d]_f & X\times_S Y \ar[r]^-{p_1} \ar[d]^{p_2} & X \ar[d]^f \\
  S \ar[r]_-{s} & Y \ar[r]_-{g} & S
}
\]
\end{proof}

Applying this lemma with $S$, $Y$ replaced by $Y$, $Y'$, one sees that one can write $f=(f')_{(Y)}$ for the base change $s:Y\to Y'$, hence $f$ is affine since $f'$ is.
Q.E.D.

A variant of this criterion is the following:

\begin{corollary}[14.5.12]
\label{IV.14.5.12}
\oldpage[IV-3]{223}
With the notations and general hypotheses of \sref{IV.14.5.11}, let $\sh{L}$ be an invertible $\sh{O}_X$-Module.
Suppose that there exists a closed subset $Z$ of $X$, proper over $Y$ \sref[II]{II.5.4.10}, and such that the prescheme induced by $X$ on the open subset $X-Z$ is normal.
Then, for $\sh{L}$ to be ample relative to $f$, it is necessary and sufficient that $\sh{L}'=\sh{L}\otimes_{\sh{O}_Y}\sh{O}_{Y'}$ be ample relative to $f'$.
\end{corollary}

\begin{proof}
Keeping the notations of the proof of \sref{IV.14.5.11}, set $\sh{L}_1=\sh{L}\otimes_{\sh{O}_Y}\sh{O}_{Y_1}$.
By \sref[III]{III.2.6.2}, it will suffice to prove that $\sh{L}_1$ is ample relative to $f\circ p$; but $f\circ p=h\circ f_1$, and, in view of \sref[II]{II.4.6.13},~(v), it will suffice to prove that $\sh{L}_1$ is ample relative to $f_1$.
The question being local on $Y_1$, one may again suppose that $g'$ admits a section $s$.
If $\sh{L}'_1=\sh{L}'\otimes_{\sh{O}_{Y'}}\sh{O}_{Y'_1}$, one can then write, by virtue of Lemma \sref{IV.14.5.11.1}, $\sh{L}_1=\sh{L}'_1\otimes_{\sh{O}_{Y'_1}}\sh{O}_{Y_1}$ for the base change $s:Y_1\to Y'_1$; the conclusion follows therefore from two applications of \sref[II]{II.4.6.13},~(iii).
\end{proof}
% END NATIVE INPUT source_aligned/ega4-14
% BEGIN NATIVE INPUT source_aligned/ega4-15
\section{Study of fibres of a universally open morphism}
\label{section:IV.15}

\subsection{Multiplicities of fibres of a universally open morphism}
\label{subsection:IV.15.1}

\begin{proposition}[15.1.1]
\label{IV.15.1.1}
\oldpage[IV-3]{223}
Let $\mathrm{Y}$ be a locally Noetherian irreducible prescheme, with generic point $\eta$, $f:\mathrm{X}\to\mathrm{Y}$ a morphism locally of finite type, $y$ a point of $\mathrm{Y}$, $\sh{F}$ a coherent $\sh{O}_\mathrm{X}$-Module;
we set $\sh{F}_y=\sh{F}\otimes_{\sh{O}_\mathrm{Y}}\mathbf{k}(y)$ (sheaf of modules on the fibre $f^{-1}(y)$).
Let $z$ be the generic point of an irreducible component of $\operatorname{Supp}(\sh{F}_y)$, and let $\mathrm{X}_i$ ($1\leqslant i\leqslant n$) be those of the reduced closed sub-preschemes of $\mathrm{X}$ whose underlying space is an irreducible component of $\operatorname{Supp}(\sh{F})$ containing $z$, and which are such that the restriction $\mathrm{X}_i\to\mathrm{Y}$ of $f$ is a universally open morphism at the point $z$; let $z_i$ be the generic point of $\mathrm{X}_i$ ($1\leqslant i\leqslant n$).
If $\lambda_x(\sh{F}_y)$ (resp.~$\lambda_t(\sh{F}_\eta)$) denotes the geometric length of $\sh{F}_y$ at a point $x\in f^{-1}(y)$ (resp.~that of $\sh{F}_\eta$ at a point $t\in f^{-1}(\eta)$) \sref{IV.4.7.5}, then we have
\[
\tag{15.1.1.1}
\label{IV.15.1.1.1}
  \lambda_z(\sh{F}_y)\geqslant\sum_{i=1}^{n}\lambda_{z_i}(\sh{F}_\eta).
\]
If furthermore we suppose the ring $\sh{O}_y$ regular, then we have
\[
\tag{15.1.1.2}
\label{IV.15.1.1.2}
  \operatorname{length}((\sh{F}_y)_z)\geqslant\sum_{i=1}^{n}\operatorname{length}((\sh{F}_\eta)_{z_i}).
\]
\end{proposition}

\begin{proof}
\label{IV.15.1.1-proof}
\oldpage[IV-3]{224}
The fibres $f^{-1}(y')$ of $f$ (for all $y'\in\mathrm{Y}$) do not change when we replace $f$ by $f_\mathrm{red}$, so we can suppose that $\mathrm{X}$ and $\mathrm{Y}$ are reduced \sref{IV.2.4.3},~(vi), and hence $\mathrm{Y}$ integral; we will set $\mathrm{K}=\mathbf{k}(\eta)$.
We can furthermore replace $f$ by the restriction $\operatorname{Supp}(\sh{F})\to\mathrm{Y}$ of $f$ to the reduced closed sub-prescheme of $\mathrm{X}$ having $\operatorname{Supp}(\sh{F})$ as underlying space, and hence restrict to the case where $\mathrm{X}=\operatorname{Supp}(\sh{F})$.
Finally, if $y=\eta$, there is only a single irreducible component $\mathrm{X}_i$ of $\mathrm{X}$ containing $z$ \sref[0]{0.2.1.8}, and we have $z_i=z$, which renders \sref{IV.15.1.1.1} and \sref{IV.15.1.1.2} trivial (without hypotheses on $f$).
We can thus restrict to the case $y\neq\eta$; it follows then from the hypothesis and from \sref{IV.1.10.3} that the $z_i$ belong to $f^{-1}(\eta)$, so the second members of \sref{IV.15.1.1.1} and \sref{IV.15.1.1.2} are well-defined.
Let $\mathrm{Z}_i$ be the reduced closed sub-prescheme of $f^{-1}(\eta)$ having $\mathrm{X}_i\cap f^{-1}(\eta)$ as underlying space; we know \sref{IV.4.6.6} that there exists a finite radicial extension $\mathrm{K}'$ of $\mathrm{K}$ such that, for all $i$, the $\mathrm{K}'$-prescheme $(\mathrm{Z}_i\otimes_\mathrm{K}\mathrm{K}')_\mathrm{red}$ is geometrically reduced over $\mathrm{K}'$.
Furthermore, the projection morphism $f^{-1}(\eta)\otimes_\mathrm{K}\mathrm{K}'\to f^{-1}(\eta)$ is finite, dominant, and radicial \sref[I]{I.3.5.7}, hence is a homeomorphism \sref{IV.2.4.5}.
If $z_i'$ denotes the unique point of $f^{-1}(\eta)\otimes_\mathrm{K}\mathrm{K}'$ above $z_i$, it follows from \sref{IV.4.7.9} and \sref{IV.4.7.5} that we have
\[
\tag{15.1.1.3}
\label{IV.15.1.1.3}
  \lambda_{z_i}(\sh{F}_\eta)=\lambda_{z_i'}(\sh{F}_\eta\otimes_\mathrm{K}\mathrm{K}')=\operatorname{length}((\sh{F}_\eta\otimes_\mathrm{K}\mathrm{K}')_{z_i'}).
\]

Let $\mathrm{V}$ be a discrete valuation ring, with field of fractions $\mathrm{K}'$, dominating $\sh{O}_y$ \sref[II]{II.7.1.7}; set $\mathrm{Y}'=\operatorname{Spec}(\mathrm{V})$, $\mathrm{X}'=\mathrm{X}\times_\mathrm{Y}\mathrm{Y}'$, $\sh{F}'=\sh{F}\otimes_\mathrm{Y}\mathrm{Y}'$, $f'=f_{(\mathrm{Y}')}:\mathrm{X}'\to\mathrm{Y}'$; if $g:\mathrm{Y}'\to\mathrm{Y}$ is the structural morphism, $\eta'$ the generic point of $\mathrm{Y}'$ and $y'$ its closed point, then we have $g(y')=y$, $g(\eta')=\eta$; the morphism $\operatorname{Spec}(\mathbf{k}(y'))\to\operatorname{Spec}(\mathbf{k}(y))$ being faithfully flat, it is the same for the projection $g':f'^{-1}(y')\to f^{-1}(y)$ \sref{IV.2.2.13}, and hence \sref{IV.2.3.4}, there is a generic point $z'$ of an irreducible component of $f'^{-1}(y')$ such that $g'(z')=z$.
By construction, we have $\mathbf{k}(\eta')=\mathrm{K}'$, hence $f'^{-1}(\eta')=f^{-1}(\eta)\otimes_\mathrm{K}\mathrm{K}'$; on the other hand, if we set $\mathrm{X}_i'=\mathrm{X}_i\times_\mathrm{Y}\mathrm{Y}'$, the restriction $\mathrm{X}_i'\to\mathrm{Y}'$ of $f'$ is a universally open morphism at the point $z'$, and hence \sref{IV.1.10.3} there exists a point $z_i''\in f'^{-1}(\eta')$ which is a generization of $z'$, and one can evidently suppose that $z_i''$ is the generic point of an irreducible component of $\mathrm{X}_i'\cap f'^{-1}(\eta')$; as the projection of $\mathrm{X}_i'\cap f'^{-1}(\eta')$ in $f^{-1}(\eta)$ is $\mathrm{Z}_i$, we have $z_i''=z_i'$.
By virtue of \sref{IV.4.7.9} and \sref{IV.4.7.5}, we have
\[
\tag{15.1.1.4}
\label{IV.15.1.1.4}
  \lambda_z(\sh{F}_y)=\lambda_{z'}(\sh{F}'_{y'})\geqslant\operatorname{length}((\sh{F}'_{y'})_{z'})
\]
and it follows from this inequality and from \sref{IV.15.1.1.3} that it suffices, in order to establish \sref{IV.15.1.1.1}, to prove the inequality
\[
\tag{15.1.1.5}
\label{IV.15.1.1.5}
  \operatorname{length}((\sh{F}_{y'}')_{z'})\geqslant\sum_{i=1}^{n}\operatorname{length}((\sh{F}_{\eta'}')_{z_i'}).
\]

Let us now consider the second inequality \sref{IV.15.1.1.2}; we will use the following lemma:

\begin{lemma}[15.1.1.6]
\label{IV.15.1.1.6}
Let $\mathrm{A}$ be a regular local ring that is not a field, $\mathrm{K}$ its field of fractions, $k$ its residue field.
There exists a discrete valuation ring $\mathrm{V}$ dominating $\mathrm{A}$, having $\mathrm{K}$ as field of fractions, and whose residue field is a pure transcendental extension of $k$.
\end{lemma}

\oldpage[IV-3]{225}
Set $\mathrm{S}=\operatorname{Spec}(\mathrm{A})$, and consider the $\mathrm{S}$-scheme $\mathrm{P}$ obtained by blowing up the closed point $s$ of $\mathrm{S}$ \sref[II]{II.8.1.3}; if $\mathfrak{m}$ is the maximal ideal of $\mathrm{A}$, we have by definition $\mathrm{P}=\operatorname{Proj}(\mathrm{B})$, where $\mathrm{B}$ is the $\mathrm{A}$-algebra graded by $\bigoplus_{n\geqslant 0}\mathfrak{m}^n$.
If $h:\mathrm{P}\to\mathrm{S}$ is the structural morphism, the fibre $h^{-1}(s)$ is isomorphic to $\operatorname{Proj}(\mathrm{B}\otimes_\mathrm{A} k)$ \sref[II]{II.2.8.10}, and by definition $\mathrm{B}\otimes_\mathrm{A} k$ is the $k$-algebra graded by $\bigoplus_{n\geqslant 0}\mathfrak{m}^n/\mathfrak{m}^{n+1}$; but as $\mathrm{A}$ is regular, this algebra is isomorphic to a polynomial algebra $k[t_1,\ldots,t_e]$ ($t_i$ indeterminates) \sref[0]{0.17.1.1} and it follows that $h^{-1}(s)$ is isomorphic to $\mathbf{P}_k^{e-1}$.
Let $\zeta$ be the generic point of $h^{-1}(s)$; if $t_i'$ ($1\leqslant i\leqslant e$) are the elements of $\mathfrak{m}$ whose classes $\bmod\mathfrak{m}^2$ are the $t_i$, $\mathrm{B}_{(t_e')}$ is the ring of an affine open neighbourhood of $\zeta$ in $\mathrm{P}$, and we have $\mathbf{k}(\zeta)=k(t_1,\ldots,t_{e-1})$; on the other hand, as $\mathrm{A}$ is integrally closed, one easily verifies that $\mathrm{B}$ is as well (Bourbaki, \emph{Alg.\ comm.}, chap.~V, \S\,1, no~8, cor.~1 of prop.~21), so $\mathrm{B}_{(t_e')}$ is also integrally closed, and hence so is the local ring $\sh{O}_{\mathrm{P},\zeta}$.
Finally, as $\mathrm{A}=\sh{O}_s$ is regular, hence a universally catenary ring \sref{IV.5.6.4}, we have, by virtue of \sref{IV.5.6.5}, $\dim(\sh{O}_{\mathrm{P},\zeta})=\dim(\sh{O}_s)-(e-1)$, since $h$ is a birational morphism, $\dim(h^{-1}(s))=e-1$, and the local ring of the fibre $h^{-1}(s)$ at its generic point $\zeta$ is of dimension $0$.
But $\dim(\sh{O}_s)=e$, hence $\dim(\sh{O}_{\mathrm{P},\zeta})=1$, and $\sh{O}_{\mathrm{P},\zeta}=\mathrm{V}$ is a discrete valuation ring, Noetherian, of dimension $1$ and integrally closed \sref[II]{II.7.1.6}; it clearly answers the question.

This lemma being established, one can resume the reasoning for the inequality \sref{IV.15.1.1.1} with $\mathrm{Y}'=\operatorname{Spec}(\mathrm{V})$; this time $\mathbf{k}(y')$ is a \emph{separable} extension of $\mathbf{k}(y)$, and by \sref{IV.4.7.9}, we have $\operatorname{length}((\sh{F}_y)_z)=\operatorname{length}((\sh{F}_{y'}')_{z'})$; as furthermore $f'^{-1}(\eta')=f^{-1}(\eta)$ and $\sh{F}_{\eta'}'=\sh{F}_\eta$, one is again reduced to proving the inequality \sref{IV.15.1.1.5}.
Now, if $t$ is a uniformiser of $\sh{O}_{y'}$, $f'^{-1}(y')$ is the closed sub-prescheme of $\mathrm{X}'$ defined by the ideal $t\sh{O}_{\mathrm{X}'}$; on the other hand, we have $\sh{F}_{y'}'=\sh{F}'/t\sh{F}'$; one also verifies immediately \sref[I]{I.9.1.12} that $(\sh{F}_\eta')_{z_i'}=\sh{F}_{z_i'}'$; the inequality to prove \sref{IV.15.1.1.5} is thus nothing but a particular case of \sref{IV.3.4.1.1}.
\end{proof}

\begin{corollary}[15.1.2]
\label{IV.15.1.2}
The hypotheses being those of \sref{IV.15.1.1}, suppose furthermore that $\lambda_z(\sh{F}_y)=1$ (resp.~that $\sh{O}_y$ is regular and $\operatorname{length}((\sh{F}_y)_z)=1$).
Then there exists at most one irreducible component $\mathrm{X}_i$ of $\operatorname{Supp}(\sh{F})$ containing $z$ and such that the restriction $\mathrm{X}_i\to\mathrm{Y}$ of $f$ is universally open at the point $z$, and moreover, if $z_i$ is the generic point of this component, then $\lambda_{z_i}(\sh{F}_\eta)=1$ (resp.~$\operatorname{length}((\sh{F}_\eta)_{z_i})=1$).
\end{corollary}

This follows immediately from \sref{IV.15.1.1}, noting that one has necessarily, by definition of $\mathrm{X}_i$, $(\sh{F}_\eta)_{z_i}\neq 0$ \sref[I]{I.9.1.13}.

\begin{corollary}[15.1.3]
\label{IV.15.1.3}
Let $\mathrm{Y}$ be an irreducible locally Noetherian prescheme of generic point $\eta$, $f:\mathrm{X}\to\mathrm{Y}$ a morphism locally of finite type, $\sh{F}$ a coherent $\sh{O}_\mathrm{X}$-module of support $\mathrm{X}$, $x$ a point of $\mathrm{X}$, $y=f(x)$.
Suppose that: $1^\circ$~$x$ belongs to a single irreducible component $\mathrm{Z}$ of $f^{-1}(y)$; $2^\circ$~$\sh{F}_y=\sh{F}\otimes_{\sh{O}_\mathrm{Y}}\mathbf{k}(y)$ is geometrically reduced over $\mathbf{k}(y)$, in other words, if $z$ is the generic point of $\mathrm{Z}$, $\lambda_z(\sh{F}_y)=1$ (resp.~$\sh{O}_y$ is regular and $\operatorname{length}((\sh{F}_y)_z)=1$).
Then there exists at most one irreducible component $\mathrm{X}'$ of $\mathrm{X}$ containing $x$, such that $\dim_x(\mathrm{X}'\cap f^{-1}(y))=\dim_x(f^{-1}(y))$ and that the restriction $\mathrm{X}'\to\mathrm{Y}$ of $f$ is universally open at the generic points of the irreducible components of $\mathrm{X}'\cap f^{-1}(y)$ containing $x$.
Moreover, if such a component $\mathrm{X}'$ exists and if $z'$ is its generic point, then $\lambda_{z'}(\sh{F}_\eta)=1$ (resp.~$\operatorname{length}((\sh{F}_\eta)_{z'})=1$).
\end{corollary}

\oldpage[IV-3]{226}
Indeed, an irreducible component $\mathrm{X}'$ of $\mathrm{X}$ containing $x$ and such that $\dim_x(\mathrm{X}'\cap f^{-1}(y))=\dim_x(f^{-1}(y))$ necessarily contains $\mathrm{Z}$, since an irreducible component of $\mathrm{X}'\cap f^{-1}(y)$ containing $x$ must be of the same dimension as the sole irreducible component $\mathrm{Z}$ of $f^{-1}(y)$ containing $x$.
One can then apply \sref{IV.15.1.2} to $z$.

\begin{remarks}[15.1.4]
\label{IV.15.1.4}
\begin{enumerate}[label=(\roman*)]
  \item In the statement of \sref{IV.15.1.1}, one cannot replace ``universally open'' by ``equidimensional'', as the example \sref{IV.14.4.10},~(ii) shows, where we set $\sh{F}=\sh{O}_\mathrm{X}$; the fibres of $f$ are then reduced Artinian schemes, and hence (with the notations introduced \emph{loc.\ cit.}) we have $\lambda_{x'}(\sh{F}_y)=1$, but there are two irreducible components $\mathrm{X}_1$, $\mathrm{X}_2$ of $\mathrm{X}$ containing $x'$, and the right-hand side of \sref{IV.15.1.1.1} is thus equal to $2$, even though the restriction of $f$ to each of the components $\mathrm{X}_1$, $\mathrm{X}_2$ is an equidimensional morphism at every point.
This same example also shows that in \sref{IV.15.1.2} and \sref{IV.15.1.3}, one cannot replace the hypothesis ``universally open'' by ``equidimensional''.
  \item It is plausible that for the validity of the inequality \sref{IV.15.1.1.2}, one cannot suppress the hypothesis that $\sh{O}_y$ is \emph{regular}.
\end{enumerate}
\end{remarks}

\subsection{Flatness of universally open morphisms with geometrically reduced fibres}
\label{subsection:IV.15.2}

\begin{env}[15.2.1]
\label{IV.15.2.1}
We saw in \sref{IV.2.4.6} that a morphism which is \emph{flat} and locally of finite presentation is \emph{universally open}.
We have a partial converse of this result under additional hypotheses:
\end{env}

\begin{theorem}[15.2.2]
\label{IV.15.2.2}
Let $\mathrm{Y}$ be a locally Noetherian prescheme, $f:\mathrm{X}\to\mathrm{Y}$ a morphism locally of finite type, $\sh{F}$ a coherent $\sh{O}_\mathrm{X}$-module, $x$ a point of $\mathrm{X}$, $y=f(x)$.
Suppose the following conditions are satisfied:
\begin{enumerate}[label=(\roman*)]
  \item $\operatorname{Supp}(\sh{F})=\mathrm{X}$.
  \item $f$ is universally open at the generic points of the irreducible components of $f^{-1}(y)$ containing $x$.
  \item $\sh{F}_y=\sh{F}\otimes_{\sh{O}_\mathrm{Y}}\mathbf{k}(y)$ is geometrically reduced at the point $x$ \sref{IV.4.6.22}.
  \item The ring $\sh{O}_y$ is reduced.
\end{enumerate}
Then $\sh{F}$ is $f$-flat at the point $x$.
\end{theorem}

\begin{proof}
\label{IV.15.2.2-proof}
As $\sh{O}_y$ is reduced and Noetherian, we will apply the valuative criterion of flatness \sref{IV.11.8.1}.
Consider thus a local homomorphism $\sh{O}_y\to\mathrm{A}'$, where $\mathrm{A}'$ is a discrete valuation ring; set $\mathrm{Y}'=\operatorname{Spec}(\mathrm{A}')$, $\mathrm{X}'=\mathrm{X}\times_\mathrm{Y}\mathrm{Y}'$, $f'=f_{(\mathrm{Y}')}:\mathrm{X}'\to\mathrm{Y}'$, and let $x'$ be a point of $\mathrm{X}'$ whose projections in $\mathrm{X}$ and $\mathrm{Y}'$ are respectively $x$ and the closed point $y'$ of $\mathrm{Y}'$; if we set $\sh{F}'=\sh{F}\otimes_\mathrm{Y}\mathrm{Y}'$, then $\operatorname{Supp}(\sh{F}')=\mathrm{X}'$ by \sref[I]{I.9.1.13}; every generic point $z'$ of an irreducible component of $f'^{-1}(y')$ containing $x'$ projects in $f^{-1}(y)$ to a generic point $z$ of an irreducible component of $f^{-1}(y)$ containing $x$, by virtue of \sref{IV.4.2.6}; hence $f'$ is universally open at the point $z'$.
Finally, it follows from \sref{IV.4.7.11} that $\sh{F}_{y'}'$ is geometrically reduced at the point $x'$.
We see thus that we are reduced to proving the following

\oldpage[IV-3]{227}
\begin{lemma}[15.2.2.1]
\label{IV.15.2.2.1}
Let $\mathrm{Y}=\operatorname{Spec}(\mathrm{A})$ be the spectrum of a discrete valuation ring, $y$ its closed point, $f:\mathrm{X}\to\mathrm{Y}$ a morphism locally of finite type, $x$ a point of $f^{-1}(y)$, $\sh{F}$ a coherent $\sh{O}_\mathrm{X}$-module.
Suppose the following conditions are satisfied:
\begin{enumerate}[label=(\roman*)]
  \item $\operatorname{Supp}(\sh{F})=\mathrm{X}$.
  \item $f$ is open at the generic points of the irreducible components of $f^{-1}(y)$ containing $x$.
  \item $\sh{F}_y=\sh{F}\otimes_{\sh{O}_y}\mathbf{k}(y)$ is reduced at the point $x$ \sref{IV.3.2.2}.
\end{enumerate}
Then $\sh{F}$ is $f$-flat at the point $x$.
\end{lemma}

Let us denote by $t$ a uniformiser of $\mathrm{A}$, set $\mathrm{B}=\sh{O}_x$ and $\sh{F}_x=\mathrm{M}$; condition~(i) implies that $\operatorname{Supp}(\mathrm{M})=\operatorname{Spec}(\mathrm{B})$ \sref[I]{I.9.1.13}.
Condition~(ii) implies, by virtue of \sref{IV.14.3.2}, that every irreducible component of $\mathrm{X}$ containing $x$ dominates $\mathrm{Y}$; in other words, the inverse image in $\mathrm{A}$ of every minimal prime ideal $\mathfrak{p}_i$ of $\mathrm{B}$ is $0$, which also means that $t\notin\mathfrak{p}_i$ for every $i$. Finally, (iii) means that $\mathrm{M}/t\mathrm{M}$ is a $(\mathrm{B}/t\mathrm{B})$-module \emph{reduced} \sref{IV.3.2.2}.
It remains to show that under these conditions $\mathrm{M}$ is a \emph{torsion-free} $\mathrm{A}$-module \sref[0]{0.6.3.4}, and for this it suffices to show that $t$ is an $\mathrm{M}$-\emph{regular} element; but this follows from \sref{IV.3.4.6.1}.
\end{proof}

\begin{corollary}[15.2.3]
\label{IV.15.2.3}
Let $\mathrm{Y}$ be a locally Noetherian prescheme, $f:\mathrm{X}\to\mathrm{Y}$ a morphism locally of finite type, $x$ a point of $\mathrm{X}$, $y=f(x)$.
Suppose the following conditions are satisfied:
\begin{enumerate}[label=(\roman*)]
  \item $f$ is universally open at the generic points of the irreducible components of $f^{-1}(y)$ containing $x$.
  \item $f^{-1}(y)$ is geometrically reduced over $\mathbf{k}(y)$ at the point $x$ \sref{IV.4.6.9}.
  \item The ring $\sh{O}_y$ is reduced.
\end{enumerate}
Under these conditions $f$ is flat at the point $x$.
It suffices to apply \sref{IV.15.2.2} with $\sh{F}=\sh{O}_\mathrm{X}$ \sref{IV.4.6.22}.
\end{corollary}

\begin{remarks}[15.2.4]
\label{IV.15.2.4}
\begin{enumerate}[label=(\roman*)]
  \item The first three conditions of \sref{IV.15.2.2} do not change when we replace $f$ by $f_\mathrm{red}$; but if $\mathfrak{N}$ is the nilradical of a Noetherian local ring $\mathrm{A}$, $\mathrm{A}/\mathfrak{N}$ is flat over $\mathrm{A}/\mathfrak{N}$ but not over $\mathrm{A}$ itself when $\mathfrak{N}\neq 0$ (Bourbaki, \emph{Alg.\ comm.}, chap.~II, \S\,3, no~2, cor.~2 of prop.~5).
One thus sees that one cannot, in \sref{IV.15.2.2}, suppress condition~(iv).
  \item It follows from \sref{IV.2.4.6} that if the conclusion of \sref{IV.15.2.2} is true, and also hypothesis~(i), $f$ is universally open in a neighbourhood of $x$, and in particular at the generic points of the irreducible components of $f^{-1}(y)$ containing $x$.
Even if conditions~(i), (iii), and (iv) are satisfied, one cannot replace (ii) by the weaker hypothesis that $f$ is universally open \emph{at the point $x$ only}.
This is what the example \sref{IV.14.1.3},~(i) shows, where $f$ is evidently universally open at every point of $\mathrm{X}_2$, hence in particular at the point $x$, intersection of $\mathrm{X}_1$ and $\mathrm{X}_2$; conditions~(ii) and (iii) of \sref{IV.15.2.3} are also satisfied by this example.
However, $\sh{O}_x$ is not a torsion-free $\sh{O}_y$-module: $\sh{O}_y$ identifies with a subring of $\sh{O}_x$ that contains zero divisors in $\sh{O}_x$; hence $f$ is not flat at the point $x$.
  \item In \sref{IV.15.2.3}, the conclusion is not necessarily valid when we replace hypothesis~(ii) by the weaker hypothesis that $f^{-1}(y)$, considered as a $\mathbf{k}(y)$-prescheme, \emph{has no immersed associated prime cycles}.
An example is provided by taking for $\mathrm{Y}$ a curve having a cusp, for example $\mathrm{Y}=\operatorname{Spec}(\mathrm{K}[\mathrm{S},\mathrm{T}]/(\mathrm{S}^3-\mathrm{T}^2))$, $\mathrm{K}$ being an algebraically closed field, and for $\mathrm{X}$ its \emph{normalisation} \sref[II]{II.6.3.8}.
If $s$ and $t$ are the classes of $\mathrm{S}$ and $\mathrm{T}$ in $\mathrm{A}=\mathrm{K}[\mathrm{S},\mathrm{T}]/(\mathrm{S}^3-\mathrm{T}^2)$, one verifies immediately that $\mathrm{X}=\operatorname{Spec}(\mathrm{B})$, where $\mathrm{B}=\mathrm{K}[s,t,u]$, $u$ being the element $t/s$ of the field of fractions of $\mathrm{A}$, and since $s=u^2$, $t=u^3$, we also have $\mathrm{B}=\mathrm{K}[u]$, isomorphic to a polynomial ring in one indeterminate over $\mathrm{K}$.
The only point $x$ of $\mathrm{X}$ above the point $y$ corresponding to the maximal ideal $(s)+(t)$ corresponds to the maximal ideal $(u)$; but since the class $\bar{u}$ of $u$ in $\sh{O}_x/\mathfrak{m}_y\sh{O}_x$ satisfies $\bar{u}^2=0$, $f^{-1}(y)$ is not a reduced $\mathbf{k}(y)$-prescheme.
Here $f$ is a finite, surjective and radicial morphism, hence a universal homeomorphism \sref{IV.2.4.5}, but it is not flat at the point $x$, for if $\sh{O}_x$ were a flat $\sh{O}_y$-module, it would be an $\sh{O}_y$-module
\oldpage[IV-3]{228}
free of rank $1$ generated by the element $1$ of $\sh{O}_x$ (Bourbaki, \emph{Alg.\ comm.}, chap.~II, \S\,3, no~2, prop.~5), which is absurd.
  \item Let us finally show that in \sref{IV.15.2.2} or \sref{IV.15.2.3}, one cannot replace ``geometrically reduced'' by ``reduced''.
We will in fact define two rings $\mathrm{A}$, $\mathrm{A}'$ having the following properties:
\begin{enumerate}[label=\arabic*)]
  \item $\mathrm{A}$ is a Noetherian local ring of maximal ideal $\mathfrak{m}$, integral, complete, of dimension $1$ and geometrically unibranch.
  \item $\mathrm{A}'$ is the integral closure of $\mathrm{A}$, a finite $\mathrm{A}$-algebra whose maximal ideal is $\mathfrak{m}\mathrm{A}'$, and its residue field $k'$ is a non-trivial finite radicial extension of the residue field $k=\mathrm{A}/\mathfrak{m}$ of $\mathrm{A}$.
\end{enumerate}
One can then take $\mathrm{Y}=\operatorname{Spec}(\mathrm{A})$, $\mathrm{X}=\operatorname{Spec}(\mathrm{A}')$, $\sh{F}=\sh{O}_\mathrm{X}$, $y$ and $x$ being the closed points of $\mathrm{Y}$ and $\mathrm{X}$ respectively; the hypotheses~(i) and (iv) of \sref{IV.15.2.2} are trivially satisfied; as $\mathrm{A}$ is of dimension $1$, $f:\mathrm{X}\to\mathrm{Y}$ is radicial, finite, and surjective, hence a universal homeomorphism \sref{IV.2.4.5}, which proves hypothesis~(i) of \sref{IV.15.2.3}.
Finally, it is clear that $f^{-1}(y)$ is reduced.
However, $f$ is not flat at the point $x$, for if $\mathrm{A}'$ were a flat $\mathrm{A}$-module, it would be a free $\mathrm{A}$-module of rank $1$ (since it has the same field of fractions as $\mathrm{A}$) generated by the element $1$ of $\mathrm{A}$ (Bourbaki, \emph{Alg.\ comm.}, chap.~II, \S\,3, no~2, prop.~5), which is absurd.

To construct the rings $\mathrm{A}$ and $\mathrm{A}'$, let us start from an imperfect field $k$ of characteristic $p>0$; let $\mathrm{K}_0=k((\mathrm{T}))$ be the field of formal Laurent series over $k$, and let $\mathrm{A}_0$ be its valuation ring for the discrete valuation given by the order of a formal Laurent series; the residue field of $\mathrm{A}_0$ is thus $k$.
Let $\alpha$ be an element of $k$ that is not a $p$-th power, and let us take $k'=k(\alpha^{1/p})$; $\mathrm{K}=k'((\mathrm{T}))$ is then a finite radicial extension of $\mathrm{K}_0$, and $\mathrm{A}'=\mathrm{A}_0\otimes_k k'$ is the discrete valuation ring for $\mathrm{K}$.
If $\mathfrak{m}'$ is the maximal ideal of $\mathrm{A}'$, conditions~$1)$ and $2)$ are satisfied by taking $\mathrm{A}=\mathrm{A}_0+\mathfrak{m}'$: indeed, as $\mathrm{A}_0$ is Noetherian and $\mathrm{A}'$ is an $\mathrm{A}_0$-module of finite type, $\mathrm{A}$ is a finite $\mathrm{A}_0$-algebra, hence a Noetherian ring; since $\mathrm{A}'$ is finite over $\mathrm{A}$, $\mathfrak{m}'$ is the only maximal ideal of $\mathrm{A}$, which is thus a complete local ring, evidently of dimension $1$ (being finite over $\mathrm{A}_0$) and geometrically unibranch, since $\mathrm{A}'$ is integrally closed and has the same field of fractions $\mathrm{K}$ as $\mathrm{A}$.
  \item The proof of \sref{IV.15.2.2} shows however that one can weaken condition~(iii) by replacing ``geometrically reduced'' by ``reduced'' when one can apply the valuative criterion of flatness \sref{IV.11.8.1} using only discrete valuation rings $\mathrm{A}'$ whose residue field $k'$ is \emph{separable} over $\mathbf{k}(y)$; indeed, if $x'$ and $z'$ have the same meanings as in \sref{IV.15.2.2}, the radicial multiplicities $\mu_i(\mathbf{k}(z')|k')$ and $\mu_i(\mathbf{k}(z)|\mathbf{k}(y))$ are the same, by virtue of \sref{IV.4.7.3}, hence it follows from \sref{IV.4.7.9} that $\operatorname{length}((\sh{F}_y)_z)$ and $\operatorname{length}((\sh{F}_{y'}')_{z'})$ are equal, and we are again reduced to the case where $\mathrm{Y}$ is the spectrum of a discrete valuation ring and $y$ its closed point.
For example, one obtains this stronger result when $\sh{O}_y$ is \emph{unibranch} and dominated by a discrete valuation ring whose residue field is a separable extension of $\mathbf{k}(y)$, by virtue of \sref{IV.11.6.2}.\footnote{The printed text cites (11.6.4), which does not occur in {\hyperlink{ega4.c04.subsec.11-6}{\S11.6}}; the applicable unibranch descent result is Corollary~{\hyperlink{ega4.c04.cor.11-6-2}{(11.6.2)}}.}
This is the case, for instance, when $\sh{O}_y$ is a \emph{regular} ring, by \sref{IV.15.1.1.6}.
But, as Hironaka has pointed out, there are Noetherian local rings that are \emph{integrally closed}, of dimension $2$ (coming from algebraic local schemes over imperfect fields) that do not satisfy the preceding condition.
\end{enumerate}
\end{remarks}

\subsection{Applications: criteria for reduction and irreducibility}
\label{subsection:IV.15.3}

\begin{proposition}[15.3.1]
\label{IV.15.3.1}
Let $\mathrm{Y}$ be a locally Noetherian prescheme, $f:\mathrm{X}\to\mathrm{Y}$ a morphism locally of finite type, $\sh{F}$ a coherent $\sh{O}_\mathrm{X}$-module, $x$ a point of $\mathrm{X}$, $y=f(x)$.
Suppose conditions~\emph{(i)}, \emph{(ii)}, and \emph{(iii)} of \sref{IV.15.2.2} are satisfied.
Then:
\begin{enumerate}[label=(\roman*)]
  \item There exists a neighbourhood $\mathrm{U}$ of $x$ in $\mathrm{X}$ such that $f|\mathrm{U}$ is universally open and that, for all $x'\in\mathrm{U}$, $\sh{F}_{f(x')}$ is geometrically reduced over $\mathbf{k}(f(x'))$ at the point $x'$.
  \item If furthermore $\mathrm{Y}$ is reduced at the point $y$, there exists a neighbourhood $\mathrm{U}_1\subset\mathrm{U}$ of $x$ such that $\sh{F}|\mathrm{U}_1$ is $f$-flat and reduced.
\end{enumerate}
\end{proposition}

\begin{proof}
\label{IV.15.3.1-proof}
Taking into account \sref[I]{I.5.1.8}, \sref{IV.4.2.7}, and \sref{IV.4.7.11}, neither the hypotheses nor conclusion~(i) of the statement change when we replace $\mathrm{Y}$ by $\mathrm{Y}_\mathrm{red}$ and $\mathrm{X}$ by $\mathrm{X}\times_\mathrm{Y}\mathrm{Y}_\mathrm{red}$, so we may suppose $\mathrm{Y}$ \emph{reduced}.
It then follows from \sref{IV.15.2.2} that $\sh{F}$ is $f$-flat at the point $x$, hence, by \sref{IV.11.1.1}, in a neighbourhood of $x$; \emph{a~fortiori}, $f$ is universally open in this neighbourhood \sref{IV.2.4.6}.
The fact that $\sh{F}_{f(x')}$ is then geometrically reduced at the point $x'$
\oldpage[IV-3]{229}
for all the points of a neighbourhood of $x$ follows from \sref{IV.12.1.1},~(vii).
Assertion~(ii) follows from this and from \sref{IV.3.3.4}.
\end{proof}

\begin{proposition}[15.3.2]
\label{IV.15.3.2}
The notations being those of \sref{IV.15.3.1}, suppose that conditions~\emph{(i)}, \emph{(ii)}, \emph{(iii)} of \sref{IV.15.2.2} are satisfied, and furthermore that $f^{-1}(y)$ is equidimensional at the point $x$.
Then $f$ is equidimensional at the point $x$.
\end{proposition}

\begin{proof}
\label{IV.15.3.2-proof}
We have, for all $y'\in\mathrm{Y}$, $\operatorname{Supp}(\sh{F}_{y'})=f^{-1}(y')$.
By hypothesis $\sh{F}_y$ is reduced (and satisfies $(S_1)$) and is equidimensional at the point $x$.
By virtue of \sref{IV.12.1.1},~(ii), one can thus suppose that $\sh{F}_{f(x')}$ is equidimensional and of constant dimension $e=\dim_x(f^{-1}(y))$ at the point $x'$ for all $x'\in\mathrm{U}$, which means that $f^{-1}(f(x'))$ is equidimensional and of dimension $e$ at the point $x'$.
Since moreover $\sh{F}|\mathrm{U}$ is $f$-flat, it follows from \sref{IV.2.3.4} and \sref{IV.13.3.1},~$a)$ that $f$ is equidimensional at the point $x$.
\end{proof}

\begin{corollary}[15.3.3]
\label{IV.15.3.3}
The notations being those of \sref{IV.15.3.1}, suppose the following conditions are satisfied:
\begin{enumerate}[label=(\roman*)]
  \item $\operatorname{Supp}(\sh{F})=\mathrm{X}$.
  \item $f$ is universally open at the generic points of the irreducible components of $f^{-1}(y)$ containing $x$.
  \item $\sh{F}_y$ is geometrically pointwise integral over $\mathbf{k}(y)$ at the point $x$ \sref{IV.4.6.22}.
  \item $\mathrm{Y}$ is geometrically unibranch at the point $y$.
\end{enumerate}
Then $x$ belongs to a single irreducible component of $\mathrm{X}$.
If furthermore $\mathrm{Y}$ is integral at the point $y$ and $\sh{F}=\sh{O}_\mathrm{X}$, then $\mathrm{X}$ is integral at the point $x$ and $f$ is flat at the point $x$.
\end{corollary}

\begin{proof}
\label{IV.15.3.3-proof}
Note first that (i) and (iii) imply that $x$ belongs to a single irreducible component $\mathrm{Z}$ of $f^{-1}(y)$.
It follows thus from \sref{IV.14.4.7} and from \sref{IV.15.3.2} that for every irreducible component $\mathrm{X}_i$ of $\mathrm{X}$ containing $x$ (and hence containing $\mathrm{Z}$), the restriction $f_i$ of $f$ to $\mathrm{X}_i$ is equidimensional at the point $x$ and universally open at the generic point of $\mathrm{Z}$.
The first assertion of the statement then follows from \sref{IV.15.1.3} and the second from \sref{IV.15.3.1}.
\end{proof}

\begin{remarks}[15.3.4]
\label{IV.15.3.4}
\begin{enumerate}[label=(\roman*)]
  \item The example \sref{IV.14.4.10},~(ii) shows that, in \sref{IV.15.3.3}, one cannot suppress the hypothesis that $\mathrm{Y}$ is geometrically unibranch at the point $y$.
  \item The remark \sref{IV.15.2.4},~(v) shows that the conclusion of \sref{IV.15.3.3} is still true under the following hypotheses: $\sh{O}_y$ is regular, $\operatorname{Supp}(\sh{F})=\mathrm{X}$, and $\sh{F}_y$ is integral at the point $x$.
It is unknown whether, in this statement, the hypothesis that $\sh{O}_y$ is regular can be replaced by the hypothesis that it is geometrically unibranch, or even integral and integrally closed.
\end{enumerate}
\end{remarks}

\subsection{Complements on Cohen-Macaulay morphisms}
\label{subsection:IV.15.4}

\begin{proposition}[15.4.1]
\label{IV.15.4.1}
Let $\mathrm{Y}$ be a locally Noetherian prescheme, $f:\mathrm{X}\to\mathrm{Y}$ a morphism locally of finite type, $\sh{F}$ a coherent $\sh{O}_\mathrm{X}$-module such that $\operatorname{Supp}(\sh{F})=\mathrm{X}$, $x$ a point of $\mathrm{X}$, $y=f(x)$.
We set, for all $z\in\mathrm{X}$, $\mathrm{X}_{f(z)}=f^{-1}(f(z))$, $\sh{F}_{f(z)}=\sh{F}\otimes_{\sh{O}_{f(z)}}\mathbf{k}(f(z))$.
Suppose that $\sh{F}$ is $f$-flat at the point $x$ and that $\sh{F}_y$ is a Cohen-Macaulay $\sh{O}_{\mathrm{X}_y}$-module at the point $x$.
Then $f$ is
\oldpage[IV-3]{230}
equidimensional at the point $x$.
Indeed, by \sref{IV.11.1.1} and \sref{IV.12.1.1},~(vi), there exists a neighbourhood $\mathrm{U}$ of $x$ such that $\sh{F}|\mathrm{U}$ is $f$-flat and, for every $x'\in\mathrm{U}$, $\sh{F}_{f(x')}$ is a Cohen-Macaulay $\sh{O}_{\mathrm{X}_{f(x')}}$-module at the point $x'$; this latter property implies that $\sh{F}_{f(x')}$ is equidimensional at the point $x'$ and that $x'$ belongs to no embedded prime cycle associated with $\sh{F}_{f(x')}$ \sref[0]{0.16.5.4}.
Furthermore, by virtue of \sref{IV.12.1.1},~(ii), one can suppose that the dimensions of the irreducible components of $\operatorname{Supp}(\sh{F}_{f(x')})|(\mathrm{U}\cap\mathrm{X}_{f(x')})$ have a (common) value independent of $x'$.
As $\operatorname{Supp}(\sh{F}_{f(x')})=\mathrm{X}_{f(x')}$ by hypothesis, it follows from \sref{IV.2.3.4} and \sref{IV.13.3.1},~$a)$ that $f$ is equidimensional at the point $x$.
\end{proposition}

\begin{proposition}[15.4.2]
\label{IV.15.4.2}
Let $\mathrm{Y}$ be a locally Noetherian prescheme, $f:\mathrm{X}\to\mathrm{Y}$ a morphism locally of finite type, $\sh{F}$ a coherent $\sh{O}_\mathrm{X}$-module such that $\operatorname{Supp}(\sh{F})=\mathrm{X}$, $x$ a point of $\mathrm{X}$, $y=f(x)$.
Suppose that $\sh{O}_y$ is regular and that $\sh{F}$ is a Cohen-Macaulay $\sh{O}_\mathrm{X}$-module at the point $x$.
Then (with the notations of \sref{IV.15.4.1}) the following conditions are equivalent:
\begin{enumerate}[label=\emph{\alph*})]
  \item $\sh{F}$ is $f$-flat at the point $x$ and $\sh{F}_y$ is a Cohen-Macaulay $\sh{O}_{\mathrm{X}_y}$-module at the point $x$.
  \item $\sh{F}$ is $f$-flat at the point $x$.
  \item $f$ is universally open in a neighbourhood of $x$.
  \item $f$ is open at the generic points of the irreducible components of $\mathrm{X}_y$ containing $x$.
  \item $\dim(\sh{F}_x)=\dim(\sh{O}_y)+\dim((\sh{F}_y)_x)$.
  \item[\emph{e}$'$)] $\dim(\sh{O}_{\mathrm{X},x})=\dim(\sh{O}_y)+\dim(\sh{O}_{\mathrm{X},x}\otimes_{\sh{O}_y}\mathbf{k}(y))$.
\end{enumerate}
As $\operatorname{Supp}(\sh{F})=\mathrm{X}$, conditions $e)$ and $e')$ are equivalent by definition \sref{IV.5.1.12}.
Condition $a)$ trivially implies $b)$; $b)$ implies that $\sh{F}$ is $f$-flat at the points of a neighbourhood of $x$ \sref{IV.11.1.1}, hence $b)$ implies $c)$ \sref{IV.2.4.6}; $c)$ trivially implies $d)$, and $d)$ implies $e')$ by \sref{IV.14.2.1}.
Finally, since $\sh{O}_y$ is regular and $\sh{F}_x$ is a Cohen-Macaulay $\sh{O}_{\mathrm{X},x}$-module, it follows from \sref{IV.6.1.4} that $e)$ implies $a)$.
\end{proposition}

\begin{proposition}[15.4.3]
\label{IV.15.4.3}
Let $f:\mathrm{X}\to\mathrm{Y}$ be a morphism locally of finite presentation, $\sh{F}$ an $\sh{O}_\mathrm{X}$-module of finite presentation, such that $\operatorname{Supp}(\sh{F})=\mathrm{X}$, $x$ a point of $\mathrm{X}$, $y=f(x)$.
Then (with the notations of \sref{IV.15.4.1}) the following conditions are equivalent:
\begin{enumerate}[label=\emph{\alph*})]
  \item $\sh{F}$ is $f$-flat at the point $x$ and $\sh{F}_y$ is a Cohen-Macaulay $\sh{O}_{\mathrm{X}_y}$-module at the point $x$.
  \item There exists an open neighbourhood $\mathrm{U}$ of $x$ in $\mathrm{X}$ and a quasi-finite $\mathrm{Y}$-morphism $g:\mathrm{U}\to\mathrm{Y}'=\mathrm{Y}[\mathrm{T}_1,\ldots,\mathrm{T}_n]$, where the $\mathrm{T}_i$ are indeterminates, such that $\sh{F}|\mathrm{U}$ is $g$-flat at the point $x$.
\end{enumerate}
\end{proposition}

\begin{proof}
\label{IV.15.4.3-proof}
Let us first show that $b)$ implies $a)$.
Set $h:\mathrm{Y}'\to\mathrm{Y}$; the $\mathbf{k}(y)$-prescheme $\mathrm{Y}_y'=h^{-1}(y)$ is regular \sref[0]{0.17.3.7}; let $\mathrm{A}$ be the local ring of this prescheme at the point $y'=g(x)$, $k$ its residue field, and $\mathrm{B}$ the local ring of $\mathrm{X}_y$ at the point $x$; set further $(\sh{F}_y)_x=\mathrm{M}$, which is a $\mathrm{B}$-module of finite type.
The hypothesis that the morphism $g$ is quasi-finite implies the same for $g_y:\mathrm{U}_y\to\mathrm{Y}_y'$ \sref[II]{II.6.2.4}, hence $\mathrm{M}\otimes_\mathrm{A} k$ is a $(\mathrm{B}\otimes_\mathrm{A} k)$-module of finite length \sref[II]{II.6.2.2}; since $\mathrm{M}$ is a flat $\mathrm{A}$-module by hypothesis and $\mathrm{A}$ is a Cohen-Macaulay ring, $\mathrm{M}$ is a Cohen-Macaulay $\mathrm{B}$-module \sref{IV.6.3.3}.
The fact that $\sh{F}$ is $f$-flat at the point $x$ follows from its being $g$-flat and from the fact that the morphism $h$ is flat \sref{IV.2.1.6}.

Let us now show that $a)$ implies $b)$.
The question being local on $\mathrm{X}$ and $\mathrm{Y}$, we may suppose that $\mathrm{X}$ and $\mathrm{Y}$ are affine; using \sref{IV.8.9.1} and \sref{IV.11.2.7}, we reduce to the case where $\mathrm{X}$ and $\mathrm{Y}$ are Noetherian, taking into account \sref{IV.6.7.1}.
\oldpage[IV-3]{231}
The hypothesis implies that $f$ is equidimensional at the point $x$ \sref{IV.15.4.1}, whence there exists a quasi-finite morphism $g:\mathrm{U}\to\mathrm{Y}'$ such that every irreducible component of $\mathrm{U}$ dominates $\mathrm{Y}'$ \sref{IV.13.3.1},~$b)$ and such that $g_y:\mathrm{U}_y\to\mathrm{Y}'_y$ is finite and surjective \sref{IV.13.3.1.1}.
On the other hand, to prove that $\sh{F}|\mathrm{U}$ is $g$-flat at the point $x$, it suffices, since $h:\mathrm{Y}'\to\mathrm{Y}$ is flat, to show by the fibrewise criterion for flatness \sref{IV.11.3.10} that $\sh{F}_y$ is $g_y$-flat at the point $x$.
But by hypothesis $\sh{F}_y$ is a Cohen-Macaulay $\sh{O}_{\mathrm{X}_y}$-module at the point $x$, and $\mathrm{Y}'_y$ is a regular prescheme; moreover, since $g_y$ is finite and surjective, it satisfies condition $e')$ of \sref{IV.15.4.2} by \sref{IV.5.6.10}.
It therefore suffices, in order to conclude, to apply \sref{IV.15.4.2} to $g_y$ and $\sh{F}_y$.
\end{proof}
\subsection{Separable rank of fibres of a quasi-finite and universally open morphism. Application to geometric connected components of fibres of a proper morphism}
\label{subsection:IV.15.5}

\begin{proposition}[15.5.1]
\label{IV.15.5.1}
Let $Y$ be a locally Noetherian prescheme, $f:X\to Y$ a separated and quasi-finite morphism.
For every $z\in Y$, let $n(z)$ be the geometric number of points of $f^{-1}(z)$ (or \emph{separable rank} of $f^{-1}(z)$ over $\kappa(z)$; cf.\ \sref[I]{I.6.4.8}).
We have the following properties:
\begin{enumerate}[label=(\roman*)]
  \item The function $z\mapsto n(z)$ is lower semi-continuous at every point $y\in Y$ such that $f$ is universally open at the points of $f^{-1}(y)$.
  \item Let $y$ be a point of $Y$ such that $f$ is universally open at the points of $f^{-1}(y)$; if the function $z\mapsto n(z)$ is constant in a neighbourhood of $y$, then $f$ is proper at the point $y$ \sref{IV.15.7.1}.
  \item Suppose that $f$ is universally open and that every point of $Y$ is geometrically unibranch; then, if the function $z\mapsto n(z)$ is locally constant in $Y$, the irreducible components of $X$ are pairwise disjoint.
\end{enumerate}
\end{proposition}

\begin{proof}
(i)
Note that since $f$ is quasi-finite, the number $n(z)$ is the geometric number of connected components of $f^{-1}(z)$; we know that the function $z\mapsto n(z)$ is locally constructible \sref{IV.9.7.9}; taking into account $(\mathbf{0}_{\mathrm{III}},~9.3.4)$, we are reduced to proving the

\begin{corollary}[15.5.2]
\label{IV.15.5.2}
Under the general hypotheses of \sref{IV.15.5.1}, let $y$ be a point of $Y$ such that $f$ is universally open at the points of $f^{-1}(y)$.
Then, if $y'$ is a generization of $y$, we have
\end{corollary}

\begin{equation}
\tag{15.5.2.1}
\label{IV.15.5.2.1}
  n(y) \leq n(y').
\end{equation}

For every base change $g:Z\to Y$, $f_{(\mathrm{Z})}:X_{(\mathrm{Z})}\to Z$ is separated, quasi-finite, and universally open at the points of $X_{(\mathrm{Z})}$ projecting into $f^{-1}(y)$; moreover, if $z$, $z'$ are two points of $Z$ such that $g(z)=y$, $g(z')=y'$ and $z'$ is a generization of $z$, we have $n(z)=n(y)$ and $n(z')=n(y')$ (\sref[I]{I.6.4.12}); it therefore suffices to prove \sref{IV.15.5.2.1} for a suitable $g$.
We can evidently suppose $y\neq y'$.
We will use the following lemma:

\begin{lemma}[15.5.2.2]
\label{IV.15.5.2.2}
Under the general hypotheses of \sref{IV.15.5.1}, if $y$ is a point of $Y$ and $y'$ is a generization of $y$ distinct from $y$, there exists a scheme $Z$ that is the spectrum of a discrete valuation ring, with closed point $z$ and generic point $z'$, and a morphism $g:Z\to Y$ such that:
$1^\circ$ $g(z)=y$, $g(z')=y'$;
$2^\circ$ if one sets $f'=f_{(Z)}$, $n(y)$ is the number of points of $f'^{-1}(z)$ and $n(y')$ is the number of points of $f'^{-1}(z')$.
\end{lemma}

Indeed, there exists a discrete valuation ring $A_1$ and a morphism $g_1$ from $Z_1=\Spec(A_1)$ to $Y$ such that if $z_1$ and $z_1'$ are the closed point and the generic point of $Z_1$,
\oldpage[IV-3]{232}
we have $g_1(z_1)=y$ and $g_1(z_1')=y'$ (\sref[II]{II.7.1.9}); we can therefore already suppose that $Y$ is the spectrum of a discrete valuation ring $A$, $y$ its closed point and $y'$ its generic point.
For every $x\in f^{-1}(y)$, $\kappa(x)$ is a finite extension of $\kappa(y)$ by hypothesis (\sref[II]{II.6.2.2}); we deduce from \sref{IV.4.5.11} that there exists a finite algebraic extension $k$ of $\kappa(y)$ such that the \emph{number of points} of $f^{-1}(y)\otimes_{\kappa(y)}k$ is equal to $n(y)$.
Since there exists a discrete valuation ring $A_2$ dominating $A$ whose residue field is finite over $\kappa(y)$ and contains $k$ (\sref[II]{II.7.1.2}), we can in the second place suppose that $n(y)$ is the number of points of $f^{-1}(y)$.
Finally, the same reasoning shows that there is a \emph{finite} extension $K$ of $\kappa(y')$ such that the number of points of $f^{-1}(y')\otimes_{\kappa(y')}K$ is equal to $n(y')$; if $w$ is a valuation of $\kappa(y')$ associated to $A$, there is a discrete valuation of $K$ extending $w$, and the ring $A_3$ of this valuation dominates $A$ and has $K$ as field of fractions, and satisfies the conditions of the lemma.

We can therefore henceforth suppose that $Y$ is the spectrum of a discrete valuation ring, $y$ its closed point, $y'$ its generic point, and that $n(y)$ and $n(y')$ are the numbers of points of $f^{-1}(y)$ and $f^{-1}(y')$, so that for every $x\in f^{-1}(y)$ (resp.\ every $x'\in f^{-1}(y')$) we have $\kappa(x)=\kappa(y)$ (resp.\ $\kappa(x')=\kappa(y')$).

Denote then by $X_i$ the reduced sub-preschemes of $X$ whose underlying spaces are the irreducible components of $X$ ($1\leq i\leq m$).
By virtue of \sref{IV.1.10.4}, the restriction $X_i\to Y$ of $f$ to each $X_i$ is a \emph{dominant} morphism, and since $f^{-1}(y')$ is discrete, each intersection $X_i\cap f^{-1}(y')$ consists only of the generic point of $X_i$, whence (denoting by $n_i(z)$ the geometric number of points of $X_i\cap f^{-1}(z)$ for every $z\in Y$), $n_i(y')=1$ for every $i$, and $n(y')=\sum_{i=1}^{m} n_i(y')=m$.
Furthermore, $f^{-1}(y)$ is the union of the $X_i\cap f^{-1}(y)$, whence
\begin{equation}
\tag{15.5.2.3}
\label{IV.15.5.2.3}
  n(y) \leq \sum_{i=1}^{m} n_i(y),
\end{equation}
and to prove \sref{IV.15.5.2.1}, it suffices to establish that $n_i(y)\leq 1$ for every $i$.
In other words, we can suppose $X$ \emph{integral}, and the morphism $f$ birational.
But then, for every $x\in f^{-1}(y)$, we have $\sh{O}_y\subset\sh{O}_x\subset K$, where $K=\kappa(y')$ is the field of fractions of $\sh{O}_y$.
Since $\sh{O}_y$ is a valuation ring and $\sh{O}_x$ dominates $\sh{O}_y$, we necessarily have $\sh{O}_y=\sh{O}_x$ (\sref[II]{II.7.1.1}).
There exists then a $Y$-section $h:Y\to X$ such that $h(y)=x$ (\sref[I]{I.2.4.4}); since $Y$ is reduced and $X$ is separated over $Y$, such a section is \emph{unique} (\sref[I]{I.7.2.2}), so the set $f^{-1}(y)$ contains only a single point, which completes the proof of (i).

(ii)
As at the beginning of (i), we are reduced to showing that if the two members of \sref{IV.15.5.2.1} are equal for \emph{every} generization $y'$ of $y$, $f$ is proper at the point $y$.
Thanks to the criterion of local properness \sref{IV.15.7.5} and to the remarks at the beginning we can further suppose that $Y=\Spec(A)$ is the spectrum of a discrete valuation ring $A$, $y$ its closed point and $y'$ its generic point, and it remains to show that $f$ is proper.

Let us use \sref{IV.15.5.2.2} and note that if $A'$ is a discrete valuation ring dominating $A$, $A'$ is a torsion-free $A$-module, so the morphism $\Spec(A')\to\Spec(A)$ is faithfully flat and quasi-compact, and it therefore amounts to the same to say that $f$ is proper or that the morphism deduced from $f$ by the base change $\Spec(A')\to\Spec(A)$ is
\oldpage[IV-3]{233}
proper \sref{IV.2.7.1},~(vii); we can therefore suppose that $n(y)$ and $n(y')$ are the numbers of points of the fibres $f^{-1}(y)$ and $f'^{-1}(y')$ respectively.
With the notations of (i), we have then by \sref{IV.15.5.2.3} and \sref{IV.15.5.2.1}
\begin{equation}
\tag{15.5.2.4}
\label{IV.15.5.2.4}
  n(y) \leq \sum_i n_i(y) \leq \sum_i n_i(y') = n(y')
\end{equation}
and for the extreme members to be equal, we need $n_i(y)=n_i(y')=1$ for every $i$.
As we saw in (i), if $n_i(y)=1$ for every $i$, $X_i\cap f^{-1}(y)$ is nonempty and the restriction $X_i\to Y$ of $f$ is an isomorphism; since the $X_i$ are closed sub-preschemes of $X$, this implies that $f$ is proper (\sref[II]{II.5.4.5}).

(iii)
We can restrict to the case where $Y$ is affine and reduced, and since $Y$ is by hypothesis locally integral (\sref[I]{I.5.1.4}), we can suppose $Y$ integral, and the function $y\mapsto n(y)$ constant in $Y$.
Let again $X_i$ ($1\leq i\leq m$) be the irreducible components of $X$.
It follows from \sref{IV.1.10.4} that if $y'$ is the generic point of $Y$, each of the $X_i\cap f^{-1}(y')$ consists of a single point; the restriction $f_i:X_i\to Y$ of $f$ being a quasi-finite and dominant morphism and every point of $Y$ being geometrically unibranch, it follows from the criterion of Chevalley \sref{IV.14.4.4} that $f_i$ is \emph{universally open}.
So then for any point $y$ of $Y$, we have $n_i(y)\leq n_i(y')$; on the other hand, since the $X_i\cap f^{-1}(y')$\footnote{The printed source has $X_i\cap f^{-1}(y)$ here; the following equality $\sum_i n_i(y')=n(y')$ and the generic-fibre argument require $y'$.} are pairwise disjoint, $\sum_i n_i(y')=n(y')$; we therefore still have the relations \sref{IV.15.5.2.4}.
But by hypothesis the extreme members are equal, so $n(y)=\sum_i n_i(y)$, which implies that the $X_i\cap f^{-1}(y)$ are pairwise disjoint, and this completes the proof.
\end{proof}

Combining this proposition with the connectivity theorem of Zariski and its consequences (\sref[III]{III.4.3}), we obtain the following results that complement (\sref[III]{III.4.3.7} and \sref[III]{III.4.3.10}):

\begin{proposition}[15.5.3]
\label{IV.15.5.3}
Let $Y$ be an irreducible locally Noetherian prescheme, $f:X\to Y$ a proper morphism; for every $y\in Y$, let $n(y)$ be the geometric number of connected components \sref{IV.4.5.2} of $f^{-1}(y)$.
Suppose that the restriction of $f$ to every irreducible component of $X$ is dominant onto $Y$.
Then, if $y_0$ is a geometrically unibranch point of $Y$, the function $y\mapsto n(y)$ is lower semi-continuous at the point $y_0$.
\end{proposition}

\begin{proof}
Consider the Stein factorisation $f=g\circ f'$ (\sref[III]{III.4.3.3}), where $g:Y'\to Y$ is a finite morphism and $f':X\to Y'$\footnote{The printed source has $f':X\to Y$ here; the stated Stein factorisation and the remainder of the proof require target $Y'$.} is a surjective morphism whose fibres are geometrically connected (\sref[III]{III.4.3.4}).
Since $f'$ is surjective, the inverse image under $f'$ of every irreducible component of $Y'$ contains at least one irreducible component of $X$; therefore each irreducible component of $Y'$ dominates $Y$.
We can therefore apply the criterion of Chevalley \sref{IV.14.4.4} to the restrictions of $g$ to each of these irreducible components, and we conclude that $g$ is \emph{universally open at every point} of $g^{-1}(y_0)$.
The conclusion therefore follows from \sref{IV.15.5.2} and from the fact that the number of geometrically connected components of $f^{-1}(y)$ is equal to the geometric number of points of $g^{-1}(y)$ (\sref[III]{III.4.3.4}).
\end{proof}

\begin{corollary}[15.5.4]
\label{IV.15.5.4}
Let $Y$ be a locally Noetherian prescheme, $f:X\to Y$ a proper morphism; let $y$ be a point of $Y$ such that $f$ is universally open at the points of $f^{-1}(y)$;
\oldpage[IV-3]{234}
then, with the notations of \sref{IV.15.5.3}, the function $z\mapsto n(z)$ is lower semi-continuous at the point $y$.
\end{corollary}

\begin{proof}
With the notations of \sref{IV.15.5.3}, note that in the Stein factorisation $f=g\circ f'$, $f'$ is surjective; since $f$ is universally open at the points of $f^{-1}(y)$, $g$ is universally open at the points of $g^{-1}(y)$ (\sref{IV.14.3.4},~(i)); it then suffices to apply to $g$ the proposition \sref{IV.15.5.1},~(i) in view of (\sref[III]{III.4.3.4}).
\end{proof}

\begin{remarks}[15.5.5]
\label{IV.15.5.5}
(i)
Even if $Y$ is integral and normal, $X$ integral, $f$ finite and surjective (hence universally open by virtue of the criterion of Chevalley \sref{IV.14.4.4}), the function $y\mapsto n(y)$ is not necessarily locally constant.
We have an example by taking $Y=\Spec(k[\mathrm{T}])$ where $k$ is algebraically closed, and $X=\Spec(A)$, where $A=k[\mathrm{S},\mathrm{T}]/(\mathrm{S}^2+\mathrm{T}^2-1)$; at the points $y'$, $y''$ of $Y$ corresponding to the maximal ideals $(\mathrm{T}-1)$ and $(\mathrm{T}+1)$, we have $n(y')=n(y'')=1$, but $n(y)=2$ at all other points of $Y$.
We shall give below an additional condition ensuring that $y\mapsto n(y)$ is locally constant \sref{IV.15.5.7}.

(ii)
The example \sref{IV.14.4.10} (ii) shows that in \sref{IV.15.5.1},~(iii), one cannot suppress the hypothesis that the points of $Y$ are geometrically unibranch.

(iii)
The example \sref{IV.11.7.5} shows that the conclusion of \sref{IV.15.5.1},~(i) is no longer valid if one only supposes that the morphism $f$ is \emph{open} (even if, as is the case in the cited example, $f$ is finite and surjective).

(iv)
Finally, in \sref{IV.15.5.1}, one cannot dispense with the hypothesis that the morphism $f$ is \emph{separated}, as shown by the example where $X$ is the affine line in which a point $x$ has been ``doubled'' (\sref[I]{I.5.5.11}), and $Y$ is the affine line: here $f$ is a local isomorphism (hence flat) and is quasi-finite, but $z\mapsto n(z)$ is not lower semi-continuous.
\end{remarks}

\begin{lemma}[15.5.6]
\label{IV.15.5.6}
Let $Y$ be the spectrum of a discrete valuation ring, $a$ its closed point, $b$ its generic point.
Let $f:X\to Y$ be an open morphism locally of finite type.
Suppose that $X$ is connected and that the fibre $f^{-1}(a)$ is a reduced prescheme.
Then $f^{-1}(b)$ is connected.
\end{lemma}

\begin{proof}
We will use the following purely topological lemma (a special case of a more general result from Ch.\ III, 3rd Part):

\begin{lemma}[15.5.6.1]
\label{IV.15.5.6.1}
Let $X$ be a locally Noetherian topological space, every closed irreducible subset of which admits a generic point.
Let $Z$ be a nowhere dense closed subset of $X$ having the following property: for every $x\in Z$, if one denotes by $\mathrm{T}_x$ the set of generizations of $x$ in $X$, then $\mathrm{T}_x-\{x\}$ is connected.
Under these conditions, if $X$ is connected, $X-Z$ is connected.
\end{lemma}

\begin{proof}
Let us give an independent proof.
Arguing by contradiction, suppose that the open set $U=X-Z$, which is dense in $X$ since $Z$ is nowhere dense, is the union of two disjoint nonempty open sets $U'$ and $U''$, and denote by $X'$ and $X''$ the closures in $X$ of $U'$ and $U''$.
Since $X=\overline{U}=X'\cup X''$ and $X$ is connected, we have $X'\cap X''\neq\varnothing$, and it is clear that $X'\cap X''\subset Z$.
Let $x$ be a generic point of an irreducible component of $X'\cap X''$.
Since $X$ is locally Noetherian, there is an open neighbourhood $V$ of $x$ in $X$ such that $V\cap U'$ and $V\cap U''$ have only a finite number of irreducible components; since $x$ is adherent to $U'$ and to $U''$, it is necessarily in the closure of an irreducible component $Z'$ of $V\cap U'$ and in the closure of an irreducible component $Z''$ of $V\cap U''$.
But $\overline{Z'}$ (resp.\ $\overline{Z''}$) is closed and irreducible in $X$, hence admits a generic point $z'$ (resp.\ $z''$), and we necessarily have $z'\in U'$ and $z''\in U''$.
This proves that the intersections of $\mathrm{T}_x-\{x\}$ with $U'$ and $U''$ are nonempty.
But by definition $(X'\cap X'')\cap(\mathrm{T}_x-\{x\})=\varnothing$; therefore the intersections of $X'$ and of $X''$ with $\mathrm{T}_x-\{x\}$ are two nonempty disjoint closed sets in $\mathrm{T}_x-\{x\}$, whose union is $\mathrm{T}_x-\{x\}$; this contradicts the hypothesis that $\mathrm{T}_x-\{x\}$ is connected.
\end{proof}
\end{proof}

\oldpage[IV-3]{235}
\begin{proof}[Proof of {\rm {\hyperlink{ega4.c04.lem.15-5-6}{(15.5.6)}}}]
To prove \sref{IV.15.5.6}, we apply the lemma \sref{IV.15.5.6.1} to $X$ and to its nowhere dense closed subset $Z=f^{-1}(a)$, which is nowhere dense because $f$ is open.
Note that by hypothesis, taking into account \sref{IV.15.2.2.1}, this implies that $f$ is \emph{flat} at the points of $f^{-1}(a)$.
For any such point $x$, $\sh{O}_x$ is therefore an $\sh{O}_a$-module without torsion ($\mathbf{0}_{\mathrm{I}},~6.3.4$), and in particular, if $t$ is a uniformiser of $\sh{O}_a$, $t$ is an $\sh{O}_x$-regular element; moreover $\sh{O}_x/t\sh{O}_x$ is reduced by hypothesis. If it is of depth $0$, it is therefore a field ($\mathbf{0},~16.4.7$), and as $t\sh{O}_x$ is then the maximal ideal of $\sh{O}_x$, $\sh{O}_x$ is a discrete valuation ring ($\mathbf{0},~17.1.4$); hence $\Spec(\sh{O}_x)-\{x\}$ consists of a single point, and \emph{a fortiori} is connected.
If on the contrary $\operatorname{depth}(\sh{O}_x/t\sh{O}_x)\geq 1$, we have $\operatorname{depth}(\sh{O}_x)\geq 2$ since $t$ is $\sh{O}_x$-regular ($\mathbf{0},~16.4.6$); it then follows from the theorem of Hartshorne \sref{IV.5.10.7} that $\Spec(\sh{O}_x)-\{x\}$ is connected, which completes the proof.
\end{proof}

\begin{proposition}[15.5.7]
\label{IV.15.5.7}
Let $Y$ be a locally Noetherian prescheme, $f:X\to Y$ a proper morphism.
For every $z\in Y$, let $n(z)$ be the geometric number of connected components of $f^{-1}(z)$.
Let $y$ be a point of $Y$ such that $f$ is universally open at the points of $f^{-1}(y)$ and that $f^{-1}(y)$ is geometrically reduced over $\kappa(y)$ \sref{IV.4.6.2}.
Then the function $z\mapsto n(z)$ is constant in a neighbourhood of $y$.
\end{proposition}

\begin{proof}
We already know from \sref{IV.15.5.4} that under the conditions of the statement, $z\mapsto n(z)$ is lower semi-continuous at the point $y$; it remains to see that it is also upper semi-continuous, and by the same reasoning as at the beginning of the proof of \sref{IV.15.5.1},~(i), it suffices to show that if $y'$ is a generization of $y$, we have
\begin{equation}
\tag{15.5.7.1}
\label{IV.15.5.7.1}
  n(y') \leq n(y).
\end{equation}

Using the argument at the beginning of the proof of \sref{IV.15.5.2} and Lemma \sref{IV.15.5.2.2}, applied to the finite morphism $g$ in the Stein factorisation $f=g\circ f'$ (recalling that the geometric number of points of $g^{-1}(y)$ is equal to the number of connected components of $f^{-1}(y)$ (\sref[III]{III.4.3.4})), we reduce to the case where $n(y)$ and $n(y')$ are respectively the numbers of connected components of $f^{-1}(y)$ and $f^{-1}(y')$, and where $Y$ is the spectrum of a discrete valuation ring, $y$ is its closed point, and $y'$ is its generic point.
We can moreover replace $X$ by any one of its connected components, these being open and closed in $X$, and proper over $Y$.
Suppose therefore that $X$ is connected. The hypothesis that $f$ is open at the points of $f^{-1}(y)$ implies that if $f^{-1}(y)\neq\varnothing$, then $f^{-1}(y')\neq\varnothing$ (1.10.3); the hypothesis that $f$ is closed implies that if $f^{-1}(y')\neq\varnothing$, then $f^{-1}(y)\neq\varnothing$ (\sref[II]{II.7.2.1}).
So, if $X\neq\varnothing$ (which we can evidently suppose, the proposition being trivially true in the contrary case), $f^{-1}(y)$ and $f^{-1}(y')$ are both nonempty.
Moreover, the hypothesis implies that the prescheme $f^{-1}(y)$ is \emph{reduced} \sref{IV.4.6.1}; since $f$ is open at the points of $f^{-1}(y)$, the lemma \sref{IV.15.5.6} implies that $f^{-1}(y')$ is connected, in other words $n(y')=1$.
Since $f^{-1}(y)$ is not empty, this proves \sref{IV.15.5.7.1}.
\end{proof}

\begin{remarks}[15.5.8]
\label{IV.15.5.8}
The relation \sref{IV.15.5.7.1} is not necessarily valid if $f$ is not supposed proper, even if it satisfies the other hypotheses of \sref{IV.15.5.7}.
With the notations of \sref{IV.15.5.5},~(i), it suffices to consider the restriction of $f$ to $X-\{x'\}$, where $x'$ is one of the two points of $X$ above a point $y$ distinct from $y'$ and $y''$; we then have $n(y)=1$, whereas if $\eta$ is the generic point of $Y$, $n(\eta)=2$.
\end{remarks}

\oldpage[IV-3]{236}

\begin{proposition}[15.5.9]
\label{IV.15.5.9}
Let $f:X\to Y$ be a flat and locally of finite presentation morphism; for every $y\in Y$, let $n(y)$ be the geometric number of connected components of $f^{-1}(y)$.
\begin{enumerate}[label=(\roman*)]
  \item If $f$ is separated and quasi-finite, the function $y\mapsto n(y)$ (then equal to the geometric number of points of $f^{-1}(y)$) is lower semi-continuous on $Y$; if it is constant in a neighbourhood of $y_0$, $f$ is proper at $y_0$.
  \item If $f$ is proper, the function $y\mapsto n(y)$ is lower semi-continuous on $Y$; if in addition $f^{-1}(y_0)$ is geometrically reduced over $\kappa(y_0)$, $n(y)$ is constant in a neighbourhood of $y_0$.
\end{enumerate}
\end{proposition}

\begin{proof}
In all cases, the questions are local on $Y$, so we can suppose $Y$ affine and $f$ of finite presentation; using \sref{IV.8.9.1}, \sref{IV.8.10.5} and \sref{IV.11.2.7}, we reduce to the case where $Y$ is Noetherian (and using moreover \sref{IV.8.2.11} for the second assertion of (i)).
Since $f$ is then moreover universally open \sref{IV.2.4.6}, it suffices to apply \sref{IV.15.5.1}, \sref{IV.15.5.4} and \sref{IV.15.5.7} to conclude.
\end{proof}

\begin{remarks}[15.5.10]
\label{IV.15.5.10}
We note that we have thus obtained another proof of \sref{IV.12.2.4},~(vi).
Conversely, one can deduce \sref{IV.15.5.7} from \sref{IV.12.2.4},~(vi): indeed, one can restrict to the case where $Y$ is reduced (replacing $f$ by $f_\mathrm{red}$), and it then follows from the hypotheses of \sref{IV.15.5.7} and from \sref{IV.15.2.3} that $f$ is flat in an open neighbourhood of $f^{-1}(y)$; since moreover $f$ is proper, we can suppose that this neighbourhood is of the form $f^{-1}(U)$, where $U$ is an open neighbourhood of $y$ in $Y$.
We conclude then by \sref{IV.12.2.4},~(vi).
The proof of \sref{IV.15.5.7} given above has the advantage of exhibiting the result \sref{IV.15.5.6}, which is of independent interest.
\end{remarks}


\subsection{Connected components of fibres along a section}
\label{subsection:IV.15.6}

\begin{proposition}[15.6.1]
\label{IV.15.6.1}
Let $Y$ be a locally Noetherian prescheme, $f:X\to Y$ a locally of finite type morphism, $g:Y\to X$ a $Y$-section of $X$ (\sref[I]{I.2.5.5}).
For every $y\in Y$, denote by $X_y^0$ the connected component of $f^{-1}(y)$ containing the point $g(y)$.
Let $y$ be a point of $Y$, and let $y'$ be a generization of $y$ in $Y$.
Then:
\begin{enumerate}[label=(\roman*)]
  \item If there exists an irreducible component $Z$ of $X_{y'}^0$, of dimension equal to $\dim(X_{y'}^0)$ and containing $g(y')$ (as will be the case if $X_{y'}^0$ is irreducible), then
\begin{equation}
\tag{15.6.1.1}
\label{IV.15.6.1.1}
  \dim(X_y^0) \geq \dim(X_{y'}^0).
\end{equation}
  \item If $X_y^0$ is moreover irreducible and if the two members of \sref{IV.15.6.1.1} are equal, then $X_y^0\subset\overline{X_{y'}^0}$ (closure in $X$).
\end{enumerate}
\end{proposition}

\begin{proof}
(i)
Let $\overline{Z}$ be the closure of $Z$ in $X$; since $y$ is a specialization of $y'$ and $g$ is continuous, $g(y)\in\overline{Z}$.
Since $Z=\overline{Z}\cap f^{-1}(y')$, it follows from the semi-continuity theorem of Chevalley (\sref{IV.13.1.3}) applied to the restriction of $f$ to the closed reduced sub-prescheme of $X$ having $\overline{Z}$ as underlying space, that $\dim_{g(y)}(\overline{Z}\cap f^{-1}(y))\geq\dim(Z)$; but the irreducible components of $\overline{Z}\cap f^{-1}(y)$ containing $g(y)$ are evidently contained in $X_y^0$, whence \emph{a fortiori} the inequality \sref{IV.15.6.1.1}.

(ii)
The reasoning of (i) shows moreover that if the two members of \sref{IV.15.6.1.1} are equal, we necessarily have $\dim_{g(y)}(\overline{Z}\cap f^{-1}(y))=\dim_{g(y)}(X_y^0)$; if $X_y^0$ is irreducible, $\overline{Z}\cap f^{-1}(y)=X_y^0$, so $X_y^0$ is contained in $\overline{Z}$, and \emph{a fortiori} in $\overline{X_{y'}^0}$.
\end{proof}

\oldpage[IV-3]{237}
\begin{corollary}[15.6.2]
\label{IV.15.6.2}
With the notations of \sref{IV.15.6.1}, suppose moreover that for every $y\in Y$, $X_y^0$ is irreducible.
Then the function $z\mapsto\dim(X_z^0)$ is upper semi-continuous on $Y$.
If moreover this function is constant in a neighbourhood of a point $y\in Y$, then $X_y^0\subset\overline{X_{y'}^0}$ for every generization $y'$ of $y$.
\end{corollary}

\begin{proof}
Let $y$ be a point of $Y$, and let $U$ be an open affine neighbourhood of $g(y)$ in $X$; then there is an open neighbourhood $V$ of $y$ in $Y$ such that $g(V)\subset U$, and for every $z\in V$, $U\cap X_z^0$ is dense in $X_z^0$, hence $\dim(X_z^0)=\dim(U\cap X_z^0)$ (\sref{IV.4.1.1.3}).
Replacing $X$ by $U$, we can suppose that $f$ is a morphism of finite type.
We then know from \sref{IV.9.7.10} that the function $z\mapsto\dim(X_z^0)$ is locally constructible, and the assertions of the corollary follow from \sref{IV.15.6.1} and from $(\mathbf{0}_{\mathrm{III}},~9.3.4)$.
\end{proof}

\begin{proposition}[15.6.3]
\label{IV.15.6.3}
With the notations of \sref{IV.15.6.1}, suppose that for every $y\in Y$, $X_y^0$ is geometrically irreducible \sref{IV.4.5.2}.
Then, for every $y\in Y$ such that the function $z\mapsto\dim(X_z^0)$ is constant in a neighbourhood of $y$, $f$ is universally open at the points of $X_y^0$.
\end{proposition}

\begin{proof}
Apply the criterion \sref{IV.14.3.7}.
Let therefore $h:Y'\to Y$ be a morphism, with $Y'$ the spectrum of a discrete valuation ring whose closed and generic points we denote by $t$ and $t'$ respectively, and suppose that $h(t)=y$, so that $y'=h(t')$ is a generization of $y$.
Set $X'=X\times_Y Y'$, $f'=f_{(Y')}:X'\to Y'$, $g'=g_{(Y')}:Y'\to X'$.
With the same notations as in \sref{IV.15.6.1}, the hypothesis on $X_y^0$ implies that $X_t^{\prime 0}=X_y^0\otimes_{\kappa(y)}\kappa(t)$ and $X_{t'}^{\prime 0}=X_{y'}^0\otimes_{\kappa(y')}\kappa(t')$, and that $X_t^{\prime 0}$ and $X_{t'}^{\prime 0}$ are irreducible \sref{IV.4.4.1}; since $\dim(X_y^0)=\dim(X_t^{\prime 0})$ and $\dim(X_{y'}^0)=\dim(X_{t'}^{\prime 0})$ \sref{IV.4.1.4}, we see that it suffices to prove the proposition for $f'$ and $X_t^{\prime 0}$. In other words, we can restrict to the case where $Y$ is the spectrum of a discrete valuation ring, $y$ its closed point and $y'$ its generic point; if $x'$ is the generic point of $X_{y'}^0$, it follows from \sref{IV.15.6.2} that every point of $X_y^0$ is a specialization of $x'$, whence the conclusion.
\end{proof}

\begin{proposition}[15.6.4]
\label{IV.15.6.4}
Under the general hypotheses of \sref{IV.15.6.1}, let $X^0$ be the union of the $X_y^0$ as $y$ runs through $Y$.
If $y\in Y$ is such that $X_y^0$ is geometrically reduced over $\kappa(y)$ \sref{IV.4.6.2} and if $f$ is universally open at every point of $X_y^0$, then $X^0$ is a neighbourhood of $X_y^0$ in $X$.
In particular, if $f$ is universally open and if, for every $y\in Y$, $X_y^0$ is geometrically reduced over $\kappa(y)$, $X^0$ is open in $X$.
\end{proposition}

\begin{proof}
Let us first show that we can reduce to the case where $f$ is a morphism of \emph{finite type}.
Let $x$ be a point of $X_y^0$; it follows from \sref{IV.5.10.8.1} (with $\mathfrak{F}=\{\varnothing\}$) that there is a finite sequence $(Z_i)_{0\leq i\leq n}$ of irreducible components of $f^{-1}(y)$ such that $g(y)\in Z_0$, $x\in Z_n$ and $Z_i\cap Z_{i+1}\neq\varnothing$ for $0\leq i\leq n-1$; the $Z_i$ being irreducible, there is a finite sequence $(U_i)$ ($1\leq i\leq n$) of open affine subsets of $f^{-1}(y)$ such that $g(y)\in U_1$, $x\in U_n$, and $U_i$ is an open affine neighbourhood of a point $x_i$ of $Z_i\cap Z_{i-1}$ for $1\leq i\leq n$.
There exists for each $i$ a quasi-compact open set $W_i$ in $X$ such that $U_i=f^{-1}(y)\cap W_i$; let $W$ be the quasi-compact open set of $X$ which is the union of the $W_i$.
Since $g$ is continuous, there is an open affine neighbourhood $V$ of $y$ in $Y$ such that $g(V)\subset W$; if $X'=W\cap f^{-1}(V)$, the restriction of $f$ to $X'$ is of finite type. Moreover, for every $z\in V$, $X_z^{\prime 0}$ is the connected component of $X'\cap f^{-1}(z)$ containing $g(z)$, and we have $X_z^{\prime 0}\subset X_z^0$, while $x$ belongs to $X_y^{\prime 0}$.
Indeed, each $U_i$ contains the two irreducible sets $U_i\cap Z_{i-1}$ and $U_i\cap Z_i$, which meet, and the irreducible sets $U_i\cap Z_{i-1}$ and $U_{i-1}\cap Z_{i-1}$
\oldpage[IV-3]{238}
meet for $2\leq i\leq n$; hence the union of the $U_i\cap Z_{i-1}$ and the $U_i\cap Z_i$, for $1\leq i\leq n$, is connected and contains $x$ and $g(y)$.
Since $f|_{X'}$ is universally open at the points of $X_y^{\prime 0}$, this completes the proof of our assertion: if the proposition has been proved when $f$ is of finite type, the union $X^{\prime 0}$ of the $X_z^{\prime 0}$, for $z\in V$, is a neighbourhood of $x$, and therefore so is $X^0$.

Suppose therefore that $f$ is of finite type.
Then $X^0$ is locally constructible \sref{IV.9.7.12}; it suffices therefore to show that $X^0$ contains every generization $x'$ of $x$ ($\mathbf{0}_{\mathrm{III}},~9.2.5$).
Set $y'=f(x')$; there exists a discrete valuation ring $A$ and a morphism $u$ from $Y'=\Spec(A)$ to $X$ such that, if $z$ is the closed point and $z'$ the generic point of $Y'$, we have $u(z)=x$ and $u(z')=x'$ (\sref[II]{II.7.1.9}); so $h=f\circ u$, and therefore $h(z)=y$ and $h(z')=y'$.
Set $X'=X\times_Y Y'$, $f'=f_{(Y')}:X'\to Y'$, which is a morphism of finite type universally open at the points lying over $X_y^0$,\footnote{The printed source says that $f'$ is universally open at the points of $f'^{-1}(z)$; the hypothesis only gives universal openness over the distinguished connected component $X_y^0$.} and $g'=g_{(Y')}:Y'\to X'$, which is a $Y'$-section of $X'$.
If $p_y:f'^{-1}(z)\to f^{-1}(y)$ (resp.\ $p_{y'}:f'^{-1}(z')\to f^{-1}(y')$) is the canonical projection, $p_y^{-1}(X_y^0)$ (resp.\ $p_{y'}^{-1}(X_{y'}^0)$) is the connected component of $f'^{-1}(z)$ (resp.\ $f'^{-1}(z')$) containing $g'(z)$ (resp.\ $g'(z')$), since the connected components $X_t^0$ are geometrically connected \sref{IV.4.5.13} and it suffices to apply \sref{IV.4.4.1}.
Moreover, the morphism $u$ corresponds to a $Y'$-section $v$ of $X'$ such that $p_y(v(z))=x$ and $p_{y'}(v(z'))=x'$, so that $v(z')$ is a generization of $v(z)$; it will therefore suffice to prove that $v(z')$ belongs to $p_{y'}^{-1}(X_{y'}^0)$.\footnote{The printed source writes $v(y)$ and $v(y')$ in this sentence; the closed and generic points of $Y'$ were denoted by $z$ and $z'$.}
Finally, it is clear that $p_y^{-1}(X_y^0)$ is geometrically reduced over $\kappa(z)$.

In other words, we are reduced to the case where $Y$ is the spectrum of a discrete valuation ring, $y$ its closed point, $y'$ its generic point.
Since $f^{-1}(y)-X_y^0$ is then closed in $X$, replacing $X$ by the open set $X-(f^{-1}(y)-X_y^0)$, we can suppose that $X_y^0=f^{-1}(y)$; similarly, one may replace $X$ by the open subset that is the connected component of $X$ containing $g(Y)$, this connected component evidently containing $X^0$; in other words, we can suppose that $X$ is \emph{connected}.
The proposition will then be established if we show that $X^0=X$, or equivalently that $X_{y'}$ is connected; but $f^{-1}(y')$ is geometrically reduced over $\kappa(y')$, and \emph{a fortiori} \emph{reduced}; moreover, $f$ is open at the points of $f^{-1}(y)$, so we can apply the lemma \sref{IV.15.5.6}, whence the conclusion.
\end{proof}

\begin{corollary}[15.6.5]
\label{IV.15.6.5}
Let $f:X\to Y$ be a flat morphism of finite presentation, $g$ a $Y$-section of $X$; for every $y\in Y$, let $X_y^0$ be the connected component of $f^{-1}(y)$ containing $g(y)$, and let $X^0$ be the union of the $X_y^0$ as $y$ runs through $Y$.
Then, for every $y\in Y$ such that $X_y^0$ is geometrically reduced over $\kappa(y)$, $X^0$ is a neighbourhood of $X_y^0$ in $X$.
In particular, if $f$ is reduced \sref{IV.6.8.1}, $X^0$ is open in $X$.
\end{corollary}

\begin{proof}
The question being local on $Y$, we can suppose $Y=\Spec(A)$ affine.
There exist therefore a Noetherian subring $A_0$ and a flat morphism of finite type $f_0:X_0\to Y_0=\Spec(A_0)$ such that $X=X_0\times_{Y_0}Y$ and $f=(f_0)_{(Y)}$ (\sref{IV.8.9.1} and \sref{IV.11.2.7}); moreover, we can suppose that there exists a $Y_0$-section $g_0$ of $X_0$ such that $g=(g_0)_{(Y)}$ \sref{IV.8.8.2}.
For every $y\in Y$, let $y_0$ be its projection in $Y_0$; it follows from \sref{IV.4.5.13} that $(X_0)_{y_0}^0$ is geometrically connected, so, if $p:f^{-1}(y)\to f_0^{-1}(y_0)$ is the canonical projection, we have $X_y^0=p^{-1}((X_0)_{y_0}^0)$ \sref{IV.4.4.1}; moreover, since $f$ is flat, the hypothesis that $X_y^0$ is geometrically reduced over $\kappa(y)$ implies that $(X_0)_{y_0}^0$ is geometrically reduced over $\kappa(y_0)$ \sref{IV.4.6.10}.
We are thus reduced to the case where
\oldpage[IV-3]{239}
$Y$ is moreover \emph{Noetherian}; since $f$ is universally open (\sref{IV.11.1.1}), it suffices then to apply \sref{IV.15.6.4}.
\end{proof}

\begin{proposition}[15.6.6]
\label{IV.15.6.6}
Let $f:X\to Y$ be a morphism such that $Y$ is locally Noetherian and $f$ is locally of finite type, or such that $f$ is of finite presentation.
Let $g$ be a $Y$-section of $X$; for every $y\in Y$, let $X_y^0$ be the connected component of $f^{-1}(y)$ containing $g(y)$; finally, let $X^0$ be the union of the $X_y^0$ as $y$ runs through $Y$.
Suppose that, for every $y\in Y$, $X_y^0$ is geometrically irreducible.
Then:
\begin{enumerate}[label=(\roman*)]
  \item The function $y\mapsto\dim(X_y^0)$ is upper semi-continuous in $Y$.
  \item Let $x_0\in X^0$, $y_0=f(x_0)$.
If the function $y\mapsto\dim(X_y^0)$ is constant in a neighbourhood of $y_0$, there exists an open neighbourhood $U$ of $y_0$ such that $f$ is universally open at all points of $X^0\cap f^{-1}(U)$.
If moreover $Y$ is reduced and $f^{-1}(y_0)$ geometrically reduced over $\kappa(y_0)$ at the point $x_0$, $f$ is flat at the point $x_0$.
  \item Conversely, suppose that $f$ is universally open at the generic point of $X_{y_0}^0$ and in addition that one of the following conditions is satisfied:
  \begin{enumerate}
    \item[$\alpha$)] The fibre $f^{-1}(y_0)$ is geometrically reduced over $\kappa(y_0)$ at the point $g(y_0)$ and the ring $\sh{O}_{Y,y_0}$ is Noetherian.
    \item[$\beta$)] For every generic point $y'$ of an irreducible component of $Y$ containing $y_0$, we have $\dim(f^{-1}(y'))=\dim(X_{y'}^0)$.
  \end{enumerate}
Then the function $y\mapsto\dim(X_y^0)$ is constant in a neighbourhood of $y_0$.
  \item The set $W$ of points $x\in X^0$ such that the function $y\mapsto\dim(X_y^0)$ is constant in a neighbourhood of $f(x)$ and $X_{f(x)}^0$ is geometrically reduced over $\kappa(f(x))$, is open in $X$.
\end{enumerate}
\end{proposition}

\begin{proof}
The questions being local on $Y$, we can suppose that $Y=\Spec(A)$ is affine.

When $f$ is of finite presentation, there exists a Noetherian subring $A_1$ of $A$ and a morphism of finite type $f_1:X_1\to Y_1=\Spec(A_1)$ such that $X=X_1\times_{Y_1}Y$ and $f=(f_1)_{(Y)}$ \sref{IV.8.9.1}; moreover, we can suppose that there exists a $Y_1$-section $g_1$ of $X_1$ such that $g=(g_1)_{(Y)}$ \sref{IV.8.8.2}.
For every $y\in Y$, let $y_1$ be its projection in $Y_1$; it follows from \sref{IV.4.5.13} that $(X_1)_{y_1}^0$ is geometrically connected, so, if $p:f^{-1}(y)\to f_1^{-1}(y_1)$ is the canonical projection, $X_y^0=p^{-1}((X_1)_{y_1}^0)$ \sref{IV.4.4.1}.
Note moreover that one can, by virtue of \sref{IV.9.7.11} and \sref{IV.9.7.12}, apply \sref{IV.9.3.3} to the property of being geometrically irreducible; we can therefore suppose $A_1$ chosen so that $(X_1)_{y_1}^0$ is geometrically irreducible for every $y_1\in Y_1$.
We have $\dim(X_y^0)=\dim((X_1)_{y_1}^0)$ \sref{IV.4.1.4}, and, applying \sref{IV.9.3.3} once more to the property of having a prescribed dimension (and using \sref{IV.9.5.5} and \sref{IV.9.7.12}), we can suppose (restricting as needed $Y$ to a neighbourhood of $y_0$) that if $\dim(X_y^0)$ is constant in $Y$, then $\dim((X_1)_{y_1}^0)$ is constant in $Y_1$.
If $Y$ is reduced, so is $Y_1$ since $A_1\subset A$; moreover, if $f^{-1}(y_0)$ is geometrically reduced over $\kappa(y_0)$ at the point $x_0$, then $f_1^{-1}(y_{01})$ is geometrically reduced at the point $x_{01}$ ($x_{01}$ and $y_{01}$ being the respective projections of $x_0$ and $y_0$ in $X_1$ and $Y_1$) \sref{IV.4.6.10}.

These remarks show that to prove (i) and (ii), we can restrict to the case where $Y$ is Noetherian and $f$ locally of finite type.
But in this case, assertion (i) follows from \sref{IV.15.6.2} and the first assertion of (ii) in \sref{IV.15.6.3}.
Moreover, if $Y$ is reduced and $f^{-1}(y_0)$ is geometrically reduced over $\kappa(y_0)$ at the point $x_0$ (with $Y$ still assumed
\oldpage[IV-3]{240}
Noetherian), then, since $X_{y_0}^0$ is the unique irreducible component of $f^{-1}(y_0)$ containing $x_0$,
\footnote{The printed source has $X_y^0$, $f^{-1}(y)$, and $x$ here; assertion~(ii) concerns the fixed points $x_0$ and $y_0=f(x_0)$.}
we deduce from \sref{IV.15.6.3} and \sref{IV.15.2.3} that if $\dim(X_y^0)$ is constant in a neighbourhood of $y_0$, $f$ is flat at the point $x_0$.

To prove (iii), let us first remark that the set $X^0$ is locally constructible in $X$ \sref{IV.9.7.12}, so, by \sref{IV.9.5.5}, the set $E$ of $y\in Y$ such that $\dim(X_y^0)=\dim(X_{y_0}^0)$ is locally constructible in $Y$.
Apply then \sref{IV.1.10.1} to $E$: it suffices to prove, taking (i) into account, that $\dim(X_{y'}^0)=\dim(X_{y_0}^0)$ for the generic point $y'$ of an irreducible component of $Y$ containing $y_0$; this already allows us to suppose that $Y=\Spec(\sh{O}_{Y,y_0})$.

If we are in case $\alpha$), then, by virtue of \sref[II]{II.7.1.9}, and using \sref{IV.4.5.13} and \sref{IV.4.4.1} as in \sref{IV.15.6.4}, we reduce immediately to the case where $Y$ is the spectrum of a discrete valuation ring, $y_0$ its closed point and $y'$ its generic point.
But then the hypotheses imply, by virtue of \sref{IV.15.2.2.1}, that $f$ is \emph{flat} at the point $g(y_0)$, and also in a neighbourhood $V$ of this point in $X$ (\sref{IV.11.1.1}).
To prove our assertion, we can replace $X$ by $V$, for $V\cap X_{y'}^0$ is nonempty by virtue of \sref{IV.2.3.4}, hence is an everywhere dense open set in $X_{y'}^0$ and therefore has the same dimension \sref{IV.4.1.1.3}; moreover, since $X_{y'}^0$ is also the connected component of $f^{-1}(y')$ containing $g(y')$, we can suppose $V$ chosen such that it meets no other irreducible component of $f^{-1}(y_0)$ or of $f^{-1}(y')$, in other words, we can suppose that $X^0=X$; but then the conclusion follows from \sref{IV.12.1.1},~(i), since by hypothesis $X_{y_0}^0$ is integral.

Suppose now that we are in the case $\beta$).
Apply this time \sref[II]{II.7.1.4} in the same way as \sref[II]{II.7.1.9} in the case $\alpha$): we are then reduced to the case where $Y$ is the spectrum of a valuation ring (not necessarily discrete), $y_0$ its closed point and $y'$ its generic point.
By virtue of \sref{IV.14.3.13}, there exists an irreducible component $Z$ of $X$ containing $X_{y_0}^0$ and dominant $Y$, and such that $\dim(Z\cap f^{-1}(y'))=\dim(X_{y_0}^0)$; but the hypothesis $\beta$) and the fact that $\dim(Z\cap f^{-1}(y'))\leq\dim(f^{-1}(y'))$ show that $\dim(X_{y_0}^0)\leq\dim(X_{y'}^0)$ whence $\dim(X_{y_0}^0)=\dim(X_{y'}^0)$ by virtue of (i).

It remains to prove (iv).
Note first that the sets envisaged do not change when we replace $f$ by $f_{(Y_\mathrm{red})}:X\times_Y Y_\mathrm{red}\to Y_\mathrm{red}$, the projection $X\times_Y Y_\mathrm{red}\to X$ being a homeomorphism.
In other words, we can suppose $Y$ reduced, and then it follows from (ii) that $f$ is \emph{flat} at the points of $W$.
Consider a point $x_0$ of $W$ and let us prove that $W$ is a neighbourhood of $x_0$; proceeding as at the beginning of the proof, and using \sref{IV.11.2.7}, we can also suppose $Y$ \emph{Noetherian}; it then follows from (ii) and from \sref{IV.15.6.4} that $X^0$ is a neighbourhood of $x_0$ in $X$, and from \sref{IV.12.1.1},~(vii) that $W$ is also a neighbourhood of $x_0$ in $X$.
\end{proof}

\begin{corollary}[15.6.7]
\label{IV.15.6.7}
Assume that the preliminary conditions of \sref{IV.15.6.6} are satisfied by $f$, and suppose moreover that for every $y\in Y$, $X_y^0$ is geometrically integral over $\kappa(y)$.
Then the following conditions are equivalent:
\begin{enumerate}[label=\emph{\alph*)}]
  \item The function $y\mapsto\dim(X_y^0)$ is locally constant in $Y$.
  \item The morphism $f_{(Y_\mathrm{red})}:X\times_Y Y_\mathrm{red}\to Y_\mathrm{red}$ deduced from $f$ by base change, is flat at the points of $X^0$.
\end{enumerate}
\oldpage[IV-3]{241}
Moreover, these conditions imply that $X^0$ is open in $X$.
Finally, when $Y$ is locally Noetherian, $a)$ and $b)$ are also equivalent to
\begin{enumerate}
  \item[\rm b$'$)] $f$ is universally open at the points of $X^0$.
\end{enumerate}
\end{corollary}

\begin{proof}
The fact that $a)$ implies $b)$ follows from \sref{IV.15.6.6},~(ii), together with the fact that $X^0$ is then open in $X$.
If $b)$ is satisfied, one can restrict to the case where $Y$ is reduced and $f$ flat; then, we reduce, as at the beginning of \sref{IV.15.6.6}, and using moreover \sref{IV.11.2.7}, to the case where $Y$ is Noetherian, case in which the conclusion follows from \sref{IV.15.6.6},~(iii),~case $\alpha$)).
The equivalence of $b)$ and $b')$ has already been proved when $Y$ is locally Noetherian \sref{IV.15.2.2.1}, taking into account the fact that the morphism $X\times_Y Y_\mathrm{red}\to X$ is a universal homeomorphism.
\end{proof}

\begin{proposition}[15.6.8]
\label{IV.15.6.8}
Let $f:X\to Y$ be a separated morphism of finite presentation, $g$ a $Y$-section of $X$; for every $y\in Y$, let $X_y^0$ be the connected component of $f^{-1}(y)$ containing $g(y)$, and let $X^0$ be the union of the $X_y^0$ as $y$ runs through $Y$.
Then, if $y\in Y$ is such that $X_y^0$ is proper over $\kappa(y)$, $X^0$ is proper over $Y$ at the point $y$ (i.e., \sref{IV.15.7.1}, there exists an open neighbourhood $U$ of $y$ in $Y$ such that $X^0\cap f^{-1}(U)$ is closed in $f^{-1}(U)$ and proper over $U$).
\end{proposition}

\begin{proof}
Proceeding as at the beginning of \sref{IV.15.6.6} (where one uses \sref{IV.8.10.5},~(v), and where one replaces \sref{IV.9.7.11} by \sref{IV.9.6.7}), we reduce to the case where $Y$ is Noetherian and $f$ of finite type. We then know from \sref{IV.9.7.12} that $X^0$ is locally constructible in $X$; moreover, the $X_z^0$ are geometrically connected (for every $z\in Y$) by virtue of \sref{IV.4.5.13}; finally, as $g(Y)\subset X^0$, $X^0$ is universally submersive over $Y$ \sref{IV.15.7.8}.
It suffices then to apply to $X^0$ the criterion \sref{IV.15.7.9}.
\end{proof}

\begin{corollary}[15.6.9]
\label{IV.15.6.9}
Let $f:X\to Y$ be a separated morphism such that $Y$ is locally Noetherian and $f$ locally of finite type, or that $f$ is of finite presentation.
Let $g$ be a $Y$-section of $X$, and for every $y\in Y$, let $X_y^0$ be the connected component of $f^{-1}(y)$ containing $g(y)$.
Let $y$ be a point of $Y$ such that $X_y^0$ is proper over $\kappa(y)$.
Then, for every generization $y'$ of $y$, $\dim(X_y^0)\geq\dim(X_{y'}^0)$.
\end{corollary}

\begin{proof}
Suppose first that $f$ is of finite presentation; then, replacing $Y$ by a neighbourhood of $y$ if necessary, we may suppose that $f|_{X^0}$ is \emph{proper} \sref{IV.15.6.8}, and it suffices to apply \sref{IV.13.1.5} to this restriction.
If $Y$ is locally Noetherian and $f$ locally of finite type, $X$ is locally Noetherian; moreover, the hypothesis that $X_y^0$ is proper over $\kappa(y)$ implies that $X_y^0$ is Noetherian. There therefore exists a Noetherian open subset $V\subset X$ containing $X_y^0$; hence there is an open neighbourhood $W$ of $y$ in $Y$ such that $g(W)\subset V$.
Moreover, if $z$ is a maximal point of $X_{y'}^0$, we may suppose that $z\in V$. The restriction of $f$ to $V\cap f^{-1}(W)$ is then a morphism of finite type, and $g|_W$ is a $W$-section; applying to this morphism the result already obtained, we see that if $Z$ is the irreducible component of $X_{y'}^0$ with generic point $z$, then $\dim(Z)\leq\dim(X_y^0)$. Since $z$ is an arbitrary maximal point of $X_{y'}^0$, we indeed have $\dim(X_{y'}^0)\leq\dim(X_y^0)$.
\end{proof}

\begin{remarks}[15.6.10]
\label{IV.15.6.10}
(i)
The supplementary hypothesis on the existence of an irreducible component of $X_{y'}^0$ containing $g(y')$ and of dimension equal to $\dim(X_{y'}^0)$, made in \sref{IV.15.6.1}, is not superfluous, as shown by the following example.
Let $A$ be a discrete valuation ring, $\pi$ a uniformiser of $A$, $K$ the field of fractions of $A$, and $k$ its residue field. Set $Y=\Spec(A)$, and denote by $y$ and $y'$ the closed point and the generic point of $Y$.
Consider the two $Y$-schemes: $X_1=\Spec(A[t])$, $X_2=\Spec(K[u,v])$, where $t$, $u$, $v$ are indeterminates (we note that the projection of $X_2$ to $Y$ is therefore $\{y'\}$).
In the ring $A[t]$, the principal ideal $(\pi t-1)$ is \emph{maximal}, and thus corresponds to a closed point $x_1$ of $X_1$, which projects to the \emph{generic point} $y'$ of $Y$ and is such that $\kappa(x_1)=K$.
\oldpage[IV-3]{242}
Consider also the closed point $x_2$ of $X_2$ corresponding to the maximal ideal $\left(u-\frac{1}{\pi}\right)+(v)$ of $K[u,v]$; we have also $\kappa(x_2)=K$.
As one will see in Ch.\ V, one can therefore \emph{reglue} $X_1$ and $X_2$ along two closed sub-preschemes $Z_1$ and $Z_2$ of $X_1$ and $X_2$ having respectively $\{x_1\}$ and $\{x_2\}$ as underlying spaces, by means of the unique $Y$-isomorphism from $Z_1$ onto $Z_2$.
Denote by $X$ the $Y$-prescheme thus obtained, and by $x$ the point of $X$ which is the image of $x_1$ and $x_2$.
The ``null section'' $g$ of $X_1$ (\sref[II]{II.1.7.9}) is then also a $Y$-section of $X$; $f^{-1}(y)$ is equal to $(X_1)_y$, whereas $f^{-1}(y')$ has two irreducible components, respectively isomorphic to $(X_1)_{y'}^0$ and $(X_2)_{y'}^0$, having the unique point $x$ in common and having respective dimensions~1 and~2. See, however, Proposition~\sref{IV.15.6.9}.

(ii)
Consider the morphism $f$ defined in \sref{IV.12.2.3},~(ii), which is proper and flat; moreover, the restriction of $f$ to $X_1\cap f_0^{-1}(Y)$ is an isomorphism, so the inverse morphism $g:Y\to X$ is a $Y$-section of $X$.
We then have $\dim(X_{y_0})=1$, whereas $\dim(X_y^0)=0$ for $y\neq y_0$, although $X_y^0$ is geometrically irreducible for every $y\in Y$ (but $X_{y_0}^0$ is not reduced); one therefore sees that in \sref{IV.15.6.6},~(iii), the hypotheses $\alpha$) and $\beta$) cannot be omitted.
Moreover, $X^0$ is not a neighbourhood of $X_{y_0}^0$ in $X$, which shows that in \sref{IV.15.6.4} one cannot dispense with the hypothesis that $X_{y_0}^0$ is reduced.

(iii)
In Ch.\ VI, we will apply the preceding results to $Y$-\emph{group preschemes} locally of finite type over a locally Noetherian prescheme $Y$.
If $G$ is such a prescheme, there exists the canonical $Y$-section $g$, the ``unit section'', and one shows that $G_y^0$ is always geometrically irreducible and $\dim(G_y^0)=\dim(G_y)$, setting $G_y=f^{-1}(y)$.
It then follows from \sref{IV.15.6.2} and \sref{IV.15.6.3} that the function $z\mapsto\dim(G_z)$ is upper semi-continuous, and that if it is constant in a neighbourhood of a point $y$, $f$ is universally open at the points of $G_y^0$.
If moreover $G_y$ is geometrically reduced over $\kappa(y)$ at one of its points, one shows that $G_y$ is \emph{smooth} \sref{IV.6.8.1} over $\kappa(y)$ at all its points; it will follow then from \sref{IV.15.6.6} that if moreover $\sh{O}_y$ is reduced, then also $f$ is smooth at all the points of $G_y^0$ (though not necessarily at all points of $G_y$).
These results will apply for example in the theory of Picard schemes, where we shall use a general theorem ensuring that, under certain conditions, the function $z\mapsto\dim(G_z)$ is locally constant.
Let us remark that it is with applications of this nature in view that the statements such as \sref{IV.15.6.1} are given for morphisms \emph{locally} of finite type, and not only for morphisms of finite type.

We note also that in the case of a $Y$-group prescheme $G$, the hypothesis $\beta$) of \sref{IV.15.6.6} is always satisfied.
\end{remarks}


\subsection{Appendix: Valuative criteria for local properness}
\label{subsection:IV.15.7}

This section gives complements to the valuative criterion of properness proved in (\sref[II]{II.7.3.10}); it is independent of the rest of \S~15.

\begin{definition}[15.7.1]
\label{IV.15.7.1}
Let $f:X\to Y$ be a morphism of preschemes, $y$ a point of $Y$.
We say that $f$ is \emph{proper at the point $y$} if there exists an open neighbourhood $U$ of $y$ in $Y$ such that the restriction $f^{-1}(U)\to U$ of $f$ is a proper morphism.
We say that a subset $Z$ of $X$ is \emph{proper over $Y$ at the point $y$} if there exists an open neighbourhood $U$ of $y$ such that $Z\cap f^{-1}(U)$ is closed in $f^{-1}(U)$ and proper over $U$ (\sref[II]{II.5.4.10}) (which amounts to saying that for every closed sub-prescheme of $X$ having $Z$ as underlying space, the restriction of $f$ to $Z$ is proper at the point $y$).
\end{definition}

Let $g:Y'\to Y$ be any morphism; set $X'=X\times_Y Y'$, $f'=f_{(Y')}$ and, if $p:X'\to X$ is the canonical projection, $Z'=p^{-1}(Z)$.
Then, if $Z$ is proper over $Y$ at the point $y$, $Z'$ is proper over $Y'$ at every point $y'$ above $y$ (\sref[II]{II.5.4.2}).

\begin{proposition}[15.7.2]
\label{IV.15.7.2}
Let $Y$ be an integral locally Noetherian prescheme, $X$ an integral prescheme, $f:X\to Y$ a separated, dominant morphism of finite type, and $y$ a point of $Y$.
The following conditions are equivalent:
\begin{enumerate}[label=\emph{\alph*)}]
  \item $f$ is proper at the point $y$.
  \item If $Y_1=\Spec(\sh{O}_y)$, $X_1=X\times_Y Y_1$, the morphism $f_1=f_{(Y_1)}:X_1\to Y_1$ is proper.
  \item Every discrete valuation ring having $R(X)$ as field of fractions and dominating $\sh{O}_y$\footnote{The printed source writes only $\mathcal O$ here; the local ring required by the criterion is $\mathcal O_y$.} dominates a local ring $\sh{O}_x$ of $X$ (in which case necessarily $f(x)=y$).
\end{enumerate}
\end{proposition}

\oldpage[IV-3]{243}
\begin{proof}
The fact that $a)$ implies $b)$ follows from (\sref[II]{II.5.4.2}), and the implication $b)\Rightarrow c)$ follows from (\sref[II]{II.7.3.10}).
It remains therefore to show that $c)$ implies $a)$.
The question is local on $Y$, so we may suppose that $Y$ is affine, hence Noetherian. By Chow's lemma (\sref[II]{II.5.6.1}), there exist an integral prescheme $X'$, a projective morphism $p:P\to Y$, a dominant open immersion $j:X'\to P$, and a projective birational (hence surjective) morphism $g:X'\to X$ such that the diagram
\[
  \xymatrix{
    P \ar[d]_-{p} & X' \ar[l]_-{j} \ar[d]^-{g} \\
    Y & X \ar[l]^-{f}
  }
\]
is commutative.
Since $X'$ is integral and $j$ is dominant, $P$ is irreducible, and, replacing $P$ by $P_\mathrm{red}$, we may suppose that $P$ is integral (\sref[I]{I.5.2.2} and \sref[II]{II.5.5.5},~(vi)).
It suffices to prove that there exists an open neighbourhood $U$ of $y$ such that $p^{-1}(U)\subset j(X')$, for then the restriction of $j$ to
\[
  j^{-1}\bigl(p^{-1}(U)\bigr)=g^{-1}\bigl(f^{-1}(U)\bigr)=V
\]
will be an isomorphism onto $p^{-1}(U)$. Hence the restriction $V\to U$ of $f\circ g$ will be proper, and, since the restriction $V\to f^{-1}(U)$ of $g$ is surjective, the restriction $f^{-1}(U)\to U$ of $f$ will be proper (\sref[II]{II.5.4.3}).
Since $T=P-j(X')$ is closed in $P$, $p(T)$ is closed in $Y$, and it suffices to show that $y\notin p(T)$, or even that every $z'\in P$ such that $p(z')=y$ \emph{belongs} to $j(X')$.
Now, the field of rational functions $R(P)$ is equal to $R(X')$, hence to $R(X)$ by construction; since one can restrict to the case where $z'$ is not the generic point of $P$, there is a discrete valuation ring $A$ having $R(X)$ as field of fractions and dominating $\sh{O}_{z'}$ (\sref[II]{II.7.1.7}); since $p(z')=y$, $A$ dominates $\sh{O}_y$, hence by hypothesis there exists $x\in X$ such that $f(x)=y$ and $A$ dominates $\sh{O}_x$.
Since $g$ is proper, there exists $x'\in X'$ such that $g(x')=x$ and $A$ dominates $\sh{O}_{x'}$ (\sref[II]{II.7.3.10}); therefore $z'$ and $j(x')$ are two points of $P$ whose local rings $\sh{O}_{z'}$ and $\sh{O}_{j(x')}=\sh{O}_{x'}$ are related by (\sref[I]{I.8.1.4}); since $P$ is a scheme, this implies $z'=j(x')$ (\sref[I]{I.8.2.2}), which completes the proof.
\end{proof}

\begin{remark}[15.7.3]
\label{IV.15.7.3}
In \sref{IV.15.7.2}, one may omit the hypothesis that $Y$ is locally Noetherian by replacing, in condition $c)$, the discrete valuation rings by arbitrary valuation rings; the proof is then unchanged, taking (\sref[II]{II.7.3.10}) into account.
\end{remark}

\begin{corollary}[15.7.4]
\label{IV.15.7.4}
Let $Y$ be a Noetherian prescheme, $f:X\to Y$ a separated morphism of finite type, $Z$ a subset of $X$ such that the maximal points of its closure $\overline{Z}$ belong to $Z$ (which is the case when $Z$ is finite, or when $Z$ is constructible ($\mathbf{0}_{\mathrm{III}},~9.2.2$)).
Let $y$ be a point of $Y$.
The following conditions are equivalent:
\begin{enumerate}[label=\emph{\alph*)}]
  \item The set $\overline{Z}$ is proper over $Y$ at the point $y$.
  \item If $Y_1=\Spec(\sh{O}_y)$, $X_1=X\times_Y Y_1$, and if $g_1:X_1\to X$ is the canonical projection and $Z_1=g_1^{-1}(Z)$, the closure $\overline{Z_1}$ of $Z_1$ in $X_1$ is proper over $Y_1$.
  \item For every scheme $Y'=\Spec(A')$, where $A'$ is a discrete valuation ring, and every morphism $g:Y'\to Y$, such that the image $g(y')$ of the closed point $y'$ of $Y'$ is $y$, we have the following property: setting $X'=X\times_Y Y'$, $f'=f_{(Y')}:X'\to Y'$, $g'=g_{(\mathrm{X})}:X'\to X$, $Z'=g'^{-1}(Z)$, and denoting by $\eta'$
\oldpage[IV-3]{244}
the generic point of $Y'$, for every point $x'\in Z'\cap f'^{-1}(\eta')$ rational over $\kappa(\eta')$, there exists a $Y'$-section of $X'$ containing $x'$.
\end{enumerate}
\end{corollary}

\begin{proof}
The fact that $a)$ implies $b)$ follows from the fact that $\overline{Z_1}\subset g_1^{-1}(\overline{Z})$ and from the fact that $g_1^{-1}(\overline{Z})$ is proper over $Y_1$ (\sref{IV.15.7.1}).
The fact that $b)$ implies $c)$ follows from (\sref[II]{II.7.3.3}).
It remains therefore to see that $c)$ implies $a)$.
It clearly suffices to prove that every irreducible component of $\overline{Z}$ is proper over $Y$ at the point $y$, and the hypothesis allows us therefore to restrict to the case where $Z$ is reduced to a single point $z$.
We may also, taking into account (\sref[I]{I.5.2.2}) and (\sref[II]{II.5.4.6}), suppose that $X$ and $Y$ are reduced; then, by the definition of a subset proper over $Y$ (\sref[II]{II.5.4.10}), suppose that $\overline{Z}=X=\overline{\{z\}}$ and, applying (\sref[I]{I.5.2.2}) once more, that $f$ is dominant. Thus $X$ and $Y$ are integral and Noetherian, and it remains to prove that $f$ is proper at the point $y$ of $Y$. For this, let us show that $f$ satisfies condition $c)$ of (\sref{IV.15.7.2}).
Let $A'$ be a discrete valuation ring having $R(X)=\kappa(z)$ as field of fractions and dominating $\sh{O}_y$; with the notation of $c)$, we have $g(y')=y$, and $g(\eta')$ is the generic point $\eta=f(z)$ of $Y$.
Since by definition $\kappa(\eta')=\kappa(z)$, there exists a $\kappa(\eta)$-morphism $\Spec(\kappa(\eta'))\to\Spec(\kappa(z))$ transforming $\eta'$ into $z$, hence a $\kappa(\eta')$-section of $f'^{-1}(\eta')$ (\sref[I]{I.3.3.14}); if $z'$ is the image of $\eta'$ by this section, $z'$ is therefore rational over $\kappa(\eta')$ and $g'(z')=z$.
We conclude from hypothesis $c)$ that there exists a $Y'$-section $h'$ of $X'$ such that $h'(\eta')=z'$; let $h=g'\circ h':Y'\to X$ be the corresponding $Y$-morphism, and let $x=h(y')$. It is clear that $f(x)=y$ and that $A'=\sh{O}_{y'}$ dominates $\sh{O}_x$, which completes the proof.
\end{proof}

\begin{corollary}[15.7.5]
\label{IV.15.7.5}
Let $Y$ be a locally Noetherian prescheme, $f:X\to Y$ a separated morphism of finite type, $y$ a point of $Y$.
The following conditions are equivalent:
\begin{enumerate}[label=\emph{\alph*)}]
  \item $f$ is proper at the point $y$.
  \item If $Y_1=\Spec(\sh{O}_y)$ and $X_1=X\times_Y Y_1$, the morphism $f_1=f_{(Y_1)}:X_1\to Y_1$ is proper.
  \item For every scheme $Y'=\Spec(A')$, where $A'$ is a discrete valuation ring, and every morphism $g:Y'\to Y$ such that the image $g(y')$ of the closed point $y'$ of $Y'$ is $y$, the following property holds: setting $X'=X\times_Y Y'$, $f'=f_{(Y')}:X'\to Y'$, and denoting by $\eta'$ the generic point of $Y'$, for every $x'\in f'^{-1}(\eta')$ rational over $\kappa(\eta')$, there exists a $Y'$-section of $X'$ containing $x'$.
\end{enumerate}
\end{corollary}

\begin{corollary}[15.7.6]
\label{IV.15.7.6}
Let $Y$ be a locally Noetherian prescheme, $f:X\to Y$ a morphism of finite type.
The following conditions are equivalent:
\begin{enumerate}[label=\emph{\alph*)}]
  \item There exists a proper morphism $g:X\to Z$ and an immersion $h:Z\to Y$ such that $f=h\circ g$.
  \item The subspace $f(X)$ is locally closed in $Y$; if $Z'$ is the closed-image sub-prescheme of $Y$ determined by $X$ under $f$ (\sref[I]{I.9.5.3} and \sref[I]{I.9.5.1}), and $Z''$ is the prescheme induced by $Z'$ on the open subset $f(X)$ of $Z'$, so that $f$ factors uniquely as $f=j\circ g$ (where $j:Z''\to Y$ is the canonical injection), then $g:X\to Z''$ is proper.
  \item For every $y\in f(X)$, $f$ is proper at the point $y$.
  \item For every scheme $Y'=\Spec(A')$, where $A'$ is a discrete valuation ring, and every morphism $g:Y'\to Y$ such that $g(Y')\subset f(X)$, every rational $Y'$-section (\sref[I]{I.7.1.2}) of $X'=X\times_Y Y'$ extends uniquely to a $Y'$-section of $X'$.
\end{enumerate}
\end{corollary}

\begin{proof}
It is clear that $b)$ implies $a)$. To prove that $a)$ implies $c)$, first note that $a)$ implies that $f(X)$ is locally closed in $Y$. For every $y\in f(X)$, there is therefore an open neighbourhood $U$
\oldpage[IV-3]{245}
of $y$ in $Y$ such that $f(X)\cap U$ is closed in $U$; the restriction $h^{-1}(U)\to U$ of $h$ is then a closed immersion, and the restriction
\[
  f^{-1}(U)=g^{-1}\bigl(h^{-1}(U)\bigr)\longrightarrow h^{-1}(U)
\]
of $g$ is proper. Consequently the restriction $f^{-1}(U)\to U$ of $f$ is proper (\sref[II]{II.5.4.2}).

To prove that $c)$ implies $b)$, first note that, for every $y\in f(X)$, there exists by $c)$ an open neighbourhood $U$ of $y$ in $Y$ such that $f(X)\cap U$ is closed in $U$. Thus $f(X)$ is locally closed, and hence open in its closure. Since this closure is the underlying space of $Z'$, and since $f$ factors as $f=j_0\circ g_0$, where $j_0:Z'\to Y$ is the canonical injection and $g_0:X\to Z'$, the equality $g_0(X)=Z''$ (which is open in $Z'$) gives the factorisation in the statement of $b)$. The fact that $g$ is proper then follows from the local character of this property (on $Z''$) and from \sref[II]{II.5.4.3}.

It is clear that $c)$ implies $d)$ by \sref{IV.15.7.5}; the uniqueness of the extension of a rational $Y'$-section follows from the hypothesis that $f$ is separated (\sref[I]{I.7.2.2}).
Finally, let us prove that $d)$ implies $c)$. First, it follows from $d)$, taking into account \sref[II]{II.7.2.3} and \sref[II]{II.7.2.4}, that $f$ is separated; we may then apply the criterion \sref{IV.15.7.5},~$c$), which shows that $f$ is proper at every point of $f(X)$.
\end{proof}

\begin{remarks}[15.7.7]
\label{IV.15.7.7}
(i)
In \sref{IV.15.7.4},~$c$), \sref{IV.15.7.5},~$c$), and \sref{IV.15.7.6},~$d$), one may restrict to the case where the discrete valuation ring $A'$ is \emph{complete} and has an algebraically closed residue field.
Indeed, in the proof of \sref{IV.15.7.4}, consider a complete discrete valuation ring $A''$ dominating $A'$ and having an algebraically closed residue field ($\mathbf{0}_{\mathrm{III}},~10.3.1$). Condition $c)$ is then satisfied for $Y''=\Spec(A'')$, $X''=X\times_Y Y''$, and $f''=f_{(Y'')}$; but $X''=X'\times_{Y'}Y''$, and it follows from the remark \sref[II]{II.7.3.9},~(i), that $f'$ is proper. Consequently $Y'$, $X'$, and $f'$ also satisfy condition $c)$ (\sref[II]{II.7.3.8}).

(ii)
Taking into account (\sref[II]{II.7.3.8}), condition $c)$ of \sref{IV.15.7.5} can be replaced by the condition that $f'$ is \emph{proper}.
\end{remarks}

\begin{env}[15.7.8]
\label{IV.15.7.8}
We say that a morphism of preschemes $f:X\to Y$ is \emph{submersive} if it is surjective and if the topology of $Y$ is equal to the quotient of the topology of $X$ by the equivalence relation defined by $f$.
We say that $f$ is \emph{universally submersive} if for every base change $Y'\to Y$, the morphism $f'=f_{(Y')}:X\times_Y Y'\to Y'$ is submersive.
Every surjective morphism that is \emph{universally open}, or \emph{universally closed}, or \emph{faithfully flat and quasi-compact} is universally submersive \sref{IV.2.3.12}.
Given a morphism $f:X\to Y$, we say that a subset $Z$ of $X$ is \emph{submersive over $f(Z)$} if the topology of the sub-space $f(Z)$ of $Y$ is equal to the quotient of the topology of the sub-space $Z$ of $X$ by the equivalence relation defined by $f|Z$.
We say that $Z$ is \emph{universally submersive over $f(Z)$} if, for every base change $g:Y'\to Y$, the inverse image $Z'=p^{-1}(Z)$ of $Z$ by the projection $p:X\times_Y Y'\to X$ is submersive over $f'(Z')=g^{-1}(f(Z))$.
We note that if the morphism $f$ admits a $Y$-section $h:Y\to X$, every set $Z$ containing $h(Y)$ is universally submersive over $Y$.
\end{env}

\begin{proposition}[15.7.9]
\label{IV.15.7.9}
Let $Y$ be a locally Noetherian prescheme, $y$ a point of $Y$, $f:X\to Y$ a separated morphism of finite type.
Let $Z$ be a locally constructible subset of $X$ having the following properties:
\begin{enumerate}[label=(\roman*)]
  \item For every generization $y'$ of $y$, $Z_{y'}=Z\cap f^{-1}(y')$ is a connected component of $f^{-1}(y')$ and is geometrically connected over $\kappa(y')$ \sref{IV.4.5.2}.
\oldpage[IV-3]{246}
  \item $Z_y$ is proper over $\Spec(\kappa(y))$.
  \item $Z$ is universally submersive over $f(Z)$ \sref{IV.15.7.8}.
\end{enumerate}
Then $Z$ is proper over $Y$ at the point $y$ (in other words \sref{IV.15.7.1} there exists an open neighbourhood $U$ of $y$ such that $Z\cap f^{-1}(U)$ is closed in $f^{-1}(U)$ and proper over $U$).
\end{proposition}

\begin{proof}
Using the method of (\sref{IV.8.1.2},~$a$)) and (\sref{IV.8.10.5},~(xii)), we are reduced to proving that when $Y=\Spec(A)$ is the spectrum of a \emph{local} Noetherian ring, the hypotheses (i), (ii) and (iii) imply that $Z$ is \emph{proper} over $Y$.

Let us first note that if $A'$ is a local Noetherian ring, $\varphi:A\to A'$ a local homomorphism, $Y'=\Spec(A')$, and $g:Y'\to Y$ the morphism corresponding to $\varphi$, the inverse image $Z'$ of $Z$ under the canonical projection $p:X\times_Y Y'\to X$ has, with respect to $Y'$, properties (i), (ii), and (iii) of the statement, taking into account \sref{IV.1.8.2}. Using \sref{IV.2.7.1},~(vii), it is therefore equivalent to prove that $Z'$ is proper over $Y'$, provided that $A'$ is a faithfully flat $A$-module.
We shall make use of this remark by taking $A'=\hat{A}$ ($\mathbf{0}_{\mathrm{I}},~7.3.5$), in other words we shall now limit ourselves to the case where $A$ is \emph{complete}.
Since $Z_y$ is a connected component of $f^{-1}(y)$, proper over $\kappa(y)$, it follows from \sref[III]{III.5.5.2} that $X$ is the \emph{sum} of two sub-preschemes $X_0$, $X_1$ induced on open and closed subsets of $X$, such that $X_0$ is proper over $Y$ and $X_0\cap f^{-1}(y)=Z_y$.
Set $Z_0=X_0\cap Z$, $Z_1=X_1\cap Z$, so that these sets form a partition of $Z$ into two open and closed subsets.
Since the $Z_{y'}$ are connected, we necessarily have $Z_{y'}\subset Z_0$ or $Z_{y'}\subset Z_1$, hence $f(Z_0)\cap f(Z_1)=\varnothing$.
But since $Z_0$ and $Z_1$ are saturated for the equivalence relation defined by $f|Z$, it follows from (iii) that $f(Z_0)$ and $f(Z_1)$ are both open in $f(Z)$.
But $f(Z_0)$ contains $y$, and every neighbourhood of $y$ in $Y$ is equal to $Y$, so $f(Z_0)=f(Z)$, and hence $f(Z_1)=\varnothing$, so $Z_1=\varnothing$ and $Z\subset X_0$.

We can therefore henceforth suppose additionally that $f$ is \emph{proper}, and it remains to prove that $Z$ is closed in $X$.
In other words, it will suffice to prove that a \emph{constructible} subset $Z$ of a prescheme $X$ \emph{proper} over $Y$, satisfying only conditions (i) and (iii), is \emph{closed} in $X$.
By virtue of ($\mathbf{0}_{\mathrm{III}},~9.2.5$), it suffices therefore to prove that for every $x'\in Z$ and every specialisation $x$ of $x'$ in $X$, we have $x\in Z$.
Let $A_1$ be a complete discrete valuation ring, $u$ a morphism of $Y_1=\Spec(A_1)$ into $X$ such that, if $y_1$ and $y_1'$ are the closed point and the generic point of $Y_1$, then $u(y_1)=x$ and $u(y_1')=x'$ (\sref[II]{II.7.1.7}); set $g=f\circ u:Y_1\to Y$ and $X_1=X\times_Y Y_1$. There exists a $Y_1$-section $u_1:Y_1\to X_1$ of $f_1=f_{(Y_1)}$ such that $u=p\circ u_1$, where $p:X_1\to X$ is the canonical projection.
Moreover, $x_1=u_1(y_1)$, $x_1'=u_1(y_1')$, and then $p(x_1)=x$ and $p(x_1')=x'$, and $x_1$ is a specialisation of $x_1'$.
Moreover, $x_1'\in Z_1=p^{-1}(Z)$. Since we have remarked that conditions (i) and (iii) are stable under \emph{every} base change, as is the property of being constructible, we are reduced to the following situation: $Y$ is the spectrum of a complete discrete valuation ring, $X$ is proper over $Y$, $f(x)=y$ is the closed point of $Y$, $y'=f(x')$ is its generic point, there exists a $Y$-section $h$ of $X$ such that $h(y)=x$ and $h(y')=x'$, and finally $x'\in Z$; we must prove that $x\in Z$.
Now, $Z_y$ is a connected component of $f^{-1}(y)$, \emph{proper} over $\Spec(\kappa(y))$ since $X$ is proper over $Y$.
Applying once more \sref[III]{III.5.5.2}, one sees, as above, that there is an open and closed subset $X'$ of $X$ such that $X'\cap f^{-1}(y)=Z_y$; since $Z_y$ is connected,
\oldpage[IV-3]{247}
we may suppose that $X'$ is connected (replacing $X'$, if necessary, by the connected component that contains $Z_y$).
But then the same reasoning as at the beginning proves that necessarily $Z\subset X'$; since $h(Y)$ is connected and contains $x'$, it is necessarily contained in $X'$, hence $x=h(y)\in X'\cap f^{-1}(y)=Z_y$, which completes the proof.
\end{proof}

\begin{corollary}[15.7.10]
\label{IV.15.7.10}
Let $Y$ be a locally Noetherian prescheme, $f:X\to Y$ a separated morphism of finite type, universally submersive \sref{IV.15.7.8} and such that all the fibres $X_y=f^{-1}(y)$ are geometrically connected.
Then, if $y\in Y$ is such that $X_y$ is proper over $\kappa(y)$, $f$ is proper at the point $y$.
\end{corollary}

\begin{proof}
Taking into account \sref{IV.15.7.5}, we are reduced to the case where $Y=\Spec(\sh{O}_y)$ since the hypotheses are stable under base change.
But in this case it suffices to apply \sref{IV.15.7.9} to $Z=X$.
\end{proof}

\begin{corollary}[15.7.11]
\label{IV.15.7.11}
Let $f:X\to Y$ be a separated morphism of finite presentation; suppose that there exists a surjective morphism of finite presentation $g:Y'\to Y$, which is moreover proper or flat, and such that, if one sets $X'=X\times_Y Y'$, there exists a $Y'$-section of $X'$ (or even a $Y$-morphism $Y'\to X$ (\sref[I]{I.3.3.14})).
Suppose finally that the fibres $X_y=f^{-1}(y)$ are geometrically connected.
Then the set $U$ of $y\in Y$ such that $X_y$ is proper over $\kappa(y)$ is open in $Y$, and the restriction $f^{-1}(U)\to U$ of $f$ is proper.
\end{corollary}

\begin{proof}
The question being local on $Y$, we can suppose $Y=\Spec(A)$ affine.
Applying \sref{IV.8.9.1} to $f$ and to $g$, one sees that there exists a Noetherian subring $A_0$ of $A$ and two morphisms of finite type $f_0:X_0\to Y_0=\Spec(A_0)$, $g_0:Y_0'\to Y_0$ such that $X=X_0\times_{Y_0}Y$, $Y'=Y_0'\times_{Y_0}Y$, $f=(f_0)_{(Y)}$ and $g=(g_0)_{(Y)}$; moreover, using \sref{IV.8.10.5},~(v),~(vi) and (xii) and \sref{IV.11.2.7}, we can suppose $f_0$ separated, and $g_0$ surjective and proper (resp.\ flat) if $g$ is proper (resp.\ flat); using \sref{IV.8.8.2}, we can further suppose that there exists a $Y_0'$-section of $X_0$.

On the other hand, using \sref{IV.9.7.7}, we may apply \sref{IV.9.3.3} to the property of being geometrically connected, and may therefore suppose $A_0$ chosen so that the $f_0^{-1}(y_0)$ are geometrically connected for every $y_0\in Y_0$.
Moreover, using \sref{IV.2.7.1},~(vii), it amounts to the same to say that $f^{-1}(y)$ is proper over $\kappa(y)$ or that $f_0^{-1}(y_0)$ is proper over $\kappa(y_0)$, where $y_0$ is the projection of $y$ in $Y_0$.
These remarks show that we are reduced to proving the corollary when $Y$ is \emph{Noetherian}.
Let us now prove the

\begin{lemma}[15.7.11.1]
\label{IV.15.7.11.1}
Let $f:X\to Y$ be a morphism, and $g:Y'\to Y$ a surjective morphism which is moreover proper or flat and quasi-compact.
Set $X'=X\times_Y Y'$, $f'=f_{(Y')}:X'\to Y'$.
Then, if $f'$ is submersive (resp.\ universally submersive), $f$ is also.
\end{lemma}

\begin{proof}
We can restrict to the case where $f'$ is submersive.
Indeed, suppose it proved in this case, and let us show that if $f'$ is universally submersive, then $f$ is also.
Let $Y_1\to Y$ be any morphism, and set $Y_1'=Y'\times_Y Y_1$; if $X_1'=X'\times_{Y'} Y_1'$, we have $f_1'=(f')_{(Y_1')}:X_1'\to Y_1'$; $f_1'$ is submersive by hypothesis; moreover, the morphism $g_1:Y_1'\to Y_1$ is surjective, and proper (resp.\ flat and quasi-compact) if $g$ is.
Hence, if one sets $X_1=X\times_Y Y_1$, $f_1=f_{(Y_1)}$, it follows from the hypothesis and from the fact that $X_1'=X_1\times_{Y_1} Y_1'$, that $f_1$ is submersive, hence $f$ universally submersive.
Suppose therefore only that $f'$ is submersive.
It is clear first of all that $f$ is surjective.
Let $E$ be a subset of $Y$ such that $f^{-1}(E)$ is closed in $X$; then, if $p:X'\to X$ is the canonical projection,
\oldpage[IV-3]{248}
$p^{-1}(f^{-1}(E))=f'^{-1}(g^{-1}(E))$ is closed in $X'$, hence, since $f'$ is submersive, $g^{-1}(E)$ is closed in $Y'$; but as $g$ is also universally submersive \sref{IV.15.7.8}, $E$ is closed in $Y$, which proves the lemma.
\end{proof}

Since there exists by hypothesis a $Y'$-section of $X'$, $f'$ is universally submersive \sref{IV.15.7.8}, hence $f$ is also by the lemma \sref{IV.15.7.11.1}.
It suffices then to apply \sref{IV.15.7.10}, noting that the set of $y\in Y$ such that $f$ is proper at $y$ is open in $Y$ by definition.
\end{proof}

% BEGIN NATIVE INPUT source_aligned/ega4-15-backmatter.tex
% EGA IV-3 standalone-English terminal matter, printed pp. 249--255.
% Source-grounded independent English edition.

\begin{center}
\emph{(To be continued.)}
\end{center}

\clearpage
\oldpage[IV-3]{249}
\phantomsection
\addcontentsline{toc}{section}{Bibliography (continued)}
\section*{BIBLIOGRAPHY (\emph{continued})}

\begingroup
\small
\setlength{\parindent}{0pt}
\setlength{\parskip}{3pt}
[39] J.-P.~\textsc{Serre}, \emph{Local Fields}, Paris (Hermann), 1962.

[40] J.-P.~\textsc{Serre}, \emph{Proalgebraic Groups}, \emph{Publ. Inst. Hautes Études Sci.}, no.~7 (1960).

\endgroup

\clearpage
\oldpage[IV-3]{250}
\phantomsection
\addcontentsline{toc}{section}{Index of notation}
\section*{INDEX OF NOTATION}

\begingroup
\small
\setlength{\parindent}{0pt}
\setlength{\parskip}{2pt}

$\mathfrak C(X)$, $\mathfrak D_c(X)$, $\mathfrak F_c(X)$,
$\mathfrak{LF}_c(X)$ ($X$ a prescheme) : 8.3.9.

$\mathfrak Q(\mathcal F)$ ($\mathcal F$ an $\mathcal O_X$-module) : 8.5.10.

$\mathfrak{Spr}(Y)$, $\mathfrak{Spro}(Y)$, $\mathfrak{Sprf}(Y)$
($Y$ a prescheme) : 8.6.1.

$\mathbf{Pro}(\mathcal C)$ ($\mathcal C$ a category) : 8.13.3.

$X_s$, $\mathcal F_s$, $u_s$, $g_s$, $Z_s$, $h_s$
($X$ an $S$-prescheme, $s\in S$, $\mathcal F$ an $\mathcal O_X$-module,
$u$ a homomorphism of $\mathcal O_X$-modules, $g$ a section of an
$\mathcal O_X$-module, $Z$ a subset of $X$, $h:X\to Y$ an $S$-morphism) : 9.4.1.

$\mathfrak{Qc}(X)$, $\mathfrak{Lqc}(X)$, $\mathfrak O(X)$,
$\mathfrak F(X)$, $\mathfrak{Lf}(X)$, $\mathfrak{Lc}(X)$
($X$ a topological space) : 10.1.1.

$\operatorname{prof}^{*}_{x}(\mathcal F)$,
$\operatorname{prof}^{*}_{Z}(\mathcal F)$,
$\operatorname{prof}^{*}(\mathcal F)$
($\mathcal F$ an $\mathcal O_X$-module) : 10.8.1.

$S(X)$ ($X$ a Jacobson prescheme) : 10.9.1.

$S(f)$ ($f$ a morphism of Jacobson preschemes) : 10.9.2.

$\operatorname{Spm}(A)$, $D^m(h)$ ($A$ a Jacobson ring, $h\in A$) : 10.9.3.

$\operatorname{Tor.dim}_B(M)$ ($M$ a $B$-module) : 12.3.2.

\endgroup

\clearpage
\oldpage[IV-3]{251}
\phantomsection
\addcontentsline{toc}{section}{Terminological index}
\section*{TERMINOLOGICAL INDEX}

\begingroup
\footnotesize
\setlength{\parindent}{0pt}
\setlength{\parskip}{1.1pt}
\newcommand{\EGAIVthreeTerm}[2]{#1 : #2.\par}

\EGAIVthreeTerm{Flatness criterion by fibres}{11.3.10}
\EGAIVthreeTerm{Algebraic space over $k$, $k$-algebraic space, prealgebraic space over $k$, $k$-prealgebraic space}{10.10.2}
\EGAIVthreeTerm{Jacobson space}{10.3.1}
\EGAIVthreeTerm{Reduced $k$-prealgebraic space}{10.10.4}
\EGAIVthreeTerm{Schematically dominant family of morphisms}{11.10.2}
\EGAIVthreeTerm{Universally schematically dominant family of morphisms}{11.10.7}
\EGAIVthreeTerm{Schematically dense family of subpreschemes}{11.10.2}
\EGAIVthreeTerm{Separating family of homomorphisms of sheaves of modules}{11.9.1}
\EGAIVthreeTerm{Separating family of module homomorphisms}{11.9.4}
\EGAIVthreeTerm{Universally separating family of homomorphisms of sheaves of modules}{11.9.14}
\EGAIVthreeTerm{Geometrically isomorphic ($k$-preschemes)}{9.1.4}
\EGAIVthreeTerm{Universally injective homomorphism of modules}{11.9.14 and 11.9.18}
\EGAIVthreeTerm{Chow's lemma for morphisms of finite presentation}{8.10.5.1}
\EGAIVthreeTerm{Zariski's ``Main Theorem'' for morphisms of finite presentation}{8.12.6}
\EGAIVthreeTerm{Equidimensional morphism at a point, equidimensional morphism}{13.2.2 and 13.3.2}
\EGAIVthreeTerm{Morphism proper at a point}{15.7.1}
\EGAIVthreeTerm{Pseudo-finite morphism}{8.12.3}
\EGAIVthreeTerm{Submersive morphism, universally submersive morphism}{15.7.8}
\EGAIVthreeTerm{Universally open morphism at a point}{14.3.3}
\EGAIVthreeTerm{Ultra-affine open subset}{10.9.5}
\EGAIVthreeTerm{Saturated finite subset of an algebraic $k$-prescheme}{9.8.9}
\EGAIVthreeTerm{Subset proper at a point (of an $S$-prescheme)}{15.7.1}
\EGAIVthreeTerm{Quasi-constructible subset, locally quasi-constructible subset of a topological space}{10.1.1}
\EGAIVthreeTerm{Very dense subset of a topological space}{10.1.3}
\EGAIVthreeTerm{Geometrically irreducible polynomial}{9.7.4}
\EGAIVthreeTerm{Jacobson prescheme}{10.4.1}
\EGAIVthreeTerm{Prescheme equidimensional over another at a point, equidimensional over another}{13.2.2 and 13.3.2}
\EGAIVthreeTerm{Prescheme essentially affine over another}{8.13.4}
\EGAIVthreeTerm{Finite extension principle}{9.1.1}
\EGAIVthreeTerm{Rectified depth}{10.8.1}
\EGAIVthreeTerm{Pro-object, morphism of pro-objects}{8.13.3}
\EGAIVthreeTerm{Essentially affine pro-object}{8.13.4, 8.14.1}
\EGAIVthreeTerm{Constructible property, ind-constructible property}{9.2.1 and 9.2.2, (i) and (ii)}
\EGAIVthreeTerm{Quasi-homeomorphism of topological spaces}{10.2.2}
\EGAIVthreeTerm{Quasi-isomorphism of ringed spaces}{10.2.8}
\EGAIVthreeTerm{Saturate of a finite subset of an algebraic $k$-prescheme}{9.8.9}
\EGAIVthreeTerm{Subset of an $S$-prescheme submersive, universally submersive over its image}{15.7.8}
\EGAIVthreeTerm{Primary skeleton of a module, virtual skeleton}{9.8.9}
\EGAIVthreeTerm{Maximal spectrum of a Jacobson ring}{10.9.3}
\EGAIVthreeTerm{Primary type of a module, geometric primary type of a module}{9.8.9}
\EGAIVthreeTerm{Ultra-prescheme, morphism of ultra-preschemes}{10.9.5}
\EGAIVthreeTerm{$R$-ultra-prescheme}{10.10.2}

\let\EGAIVthreeTerm\relax
\endgroup

\clearpage
% Printed page 252 is blank.
\thispagestyle{empty}
\null

\clearpage
\oldpage[IV-3]{253}
\phantomsection
\addcontentsline{toc}{section}{Table of contents}
\section*{TABLE OF CONTENTS}

\begingroup
\small
\setlength{\parindent}{0pt}
\setlength{\parskip}{1.5pt}
\newcommand{\EGAIVthreeContentsLine}[3]{\textbf{#1}\quad #2\unskip\nobreak\hspace{0.4em}\nobreak\dotfill\nobreak\mbox{#3}\par}

\EGAIVthreeContentsLine{CHAPTER IV.}{\textbf{Local study of schemes and morphisms of schemes} \emph{(continued)}}{5}
\EGAIVthreeContentsLine{\S~8.}{Projective limits of preschemes}{5}
\EGAIVthreeContentsLine{8.1.}{Introduction}{5}
\EGAIVthreeContentsLine{8.2.}{Projective limits of preschemes}{7}
\EGAIVthreeContentsLine{8.3.}{Constructible subsets in a projective limit of preschemes}{12}
\EGAIVthreeContentsLine{8.4.}{Irreducibility and connectedness criteria for projective limits of preschemes}{17}
\EGAIVthreeContentsLine{8.5.}{Modules of finite presentation on a projective limit of preschemes}{19}
\EGAIVthreeContentsLine{8.6.}{Subpreschemes of finite presentation of a projective limit of preschemes}{25}
\EGAIVthreeContentsLine{8.7.}{Criteria for a projective limit of preschemes to be a reduced (resp. integral) prescheme}{27}
\EGAIVthreeContentsLine{8.8.}{Preschemes of finite presentation on a projective limit of preschemes}{28}
\EGAIVthreeContentsLine{8.9.}{First applications to the elimination of Noetherian hypotheses}{34}
\EGAIVthreeContentsLine{8.10.}{Stability properties of morphisms under passage to the projective limit}{36}
\EGAIVthreeContentsLine{8.11.}{Application to quasi-finite morphisms}{41}
\EGAIVthreeContentsLine{8.12.}{New proof and generalization of Zariski's ``Main Theorem''}{43}
\EGAIVthreeContentsLine{8.13.}{Translation in terms of pro-objects}{49}
\EGAIVthreeContentsLine{8.14.}{Characterization of a prescheme locally of finite presentation over another, in terms of the functor it represents}{52}

\EGAIVthreeContentsLine{\S~9.}{Constructible properties}{54}
\EGAIVthreeContentsLine{9.1.}{The finite extension principle}{54}
\EGAIVthreeContentsLine{9.2.}{Constructible and ind-constructible properties}{56}
\EGAIVthreeContentsLine{9.3.}{Constructible properties of morphisms of algebraic preschemes}{60}

\clearpage
\oldpage[IV-3]{254}
\EGAIVthreeContentsLine{9.4.}{Constructibility of certain properties of modules}{62}
\EGAIVthreeContentsLine{9.5.}{Constructibility of topological properties}{67}
\EGAIVthreeContentsLine{9.6.}{Constructibility of certain properties of morphisms}{71}
\EGAIVthreeContentsLine{9.7.}{Constructibility of separability, geometric irreducibility, and geometric connectedness}{76}
\EGAIVthreeContentsLine{9.8.}{Primary decomposition near a generic fibre}{83}
\EGAIVthreeContentsLine{9.9.}{Constructibility of local properties of fibres}{88}

\EGAIVthreeContentsLine{\S~10.}{Jacobson preschemes}{95}
\EGAIVthreeContentsLine{10.1.}{Very dense subsets of a topological space}{95}
\EGAIVthreeContentsLine{10.2.}{Quasi-homeomorphisms}{97}
\EGAIVthreeContentsLine{10.3.}{Jacobson spaces}{101}
\EGAIVthreeContentsLine{10.4.}{Jacobson preschemes and Jacobson rings}{101}
\EGAIVthreeContentsLine{10.5.}{Noetherian Jacobson preschemes}{104}
\EGAIVthreeContentsLine{10.6.}{Dimension in Jacobson preschemes}{107}
\EGAIVthreeContentsLine{10.7.}{Examples and counterexamples}{109}
\EGAIVthreeContentsLine{10.8.}{Rectified depth}{110}
\EGAIVthreeContentsLine{10.9.}{Maximal spectra and ultrapreschemes}{112}
\EGAIVthreeContentsLine{10.10.}{Serre algebraic spaces}{114}

\EGAIVthreeContentsLine{\S~11.}{Topological properties of flat morphisms of finite presentation. Flatness criteria}{116}
\EGAIVthreeContentsLine{11.1.}{Flatness loci (Noetherian case)}{117}
\EGAIVthreeContentsLine{11.2.}{Flatness of a projective limit of preschemes}{119}
\EGAIVthreeContentsLine{11.3.}{Application to the elimination of Noetherian hypotheses}{132}
\EGAIVthreeContentsLine{11.4.}{Descent of flatness by arbitrary morphisms: case of an Artinian base prescheme}{143}
\EGAIVthreeContentsLine{11.5.}{Descent of flatness by arbitrary morphisms: general case}{150}
\EGAIVthreeContentsLine{11.6.}{Descent of flatness by arbitrary morphisms: case of a unibranch base prescheme}{154}
\EGAIVthreeContentsLine{11.7.}{Counterexamples}{157}
\EGAIVthreeContentsLine{11.8.}{A valuative criterion for flatness}{159}
\EGAIVthreeContentsLine{11.9.}{Separating and universally separating families of homomorphisms of sheaves of modules}{160}
\EGAIVthreeContentsLine{11.10.}{Schematically dominant families of morphisms and schematically dense families of subpreschemes}{170}

\EGAIVthreeContentsLine{\S~12.}{Study of the fibres of flat morphisms of finite presentation}{173}
\EGAIVthreeContentsLine{12.0.}{Introduction}{173}
\EGAIVthreeContentsLine{12.1.}{Local properties of fibres of a flat morphism locally of finite presentation}{174}

\clearpage
\oldpage[IV-3]{255}
\EGAIVthreeContentsLine{12.2.}{Local and global properties of fibres of a proper, flat morphism of finite presentation}{179}
\EGAIVthreeContentsLine{12.3.}{Local cohomological properties of fibres of a flat morphism locally of finite presentation}{183}

\EGAIVthreeContentsLine{\S~13.}{Equidimensional morphisms}{187}
\EGAIVthreeContentsLine{13.1.}{Chevalley's semicontinuity theorem}{188}
\EGAIVthreeContentsLine{13.2.}{Equidimensional morphisms: case of dominant morphisms of irreducible preschemes}{190}
\EGAIVthreeContentsLine{13.3.}{Equidimensional morphisms: general case}{194}

\EGAIVthreeContentsLine{\S~14.}{Universally open morphisms}{199}
\EGAIVthreeContentsLine{14.1.}{Open morphisms}{200}
\EGAIVthreeContentsLine{14.2.}{Open morphisms and the dimension formula}{202}
\EGAIVthreeContentsLine{14.3.}{Universally open morphisms}{204}
\EGAIVthreeContentsLine{14.4.}{Chevalley's criterion for universally open morphisms}{209}
\EGAIVthreeContentsLine{14.5.}{Universally open morphisms and quasi-sections}{216}

\EGAIVthreeContentsLine{\S~15.}{Study of the fibres of a universally open morphism}{223}
\EGAIVthreeContentsLine{15.1.}{Multiplicities of the fibres of a universally open morphism}{223}
\EGAIVthreeContentsLine{15.2.}{Flatness of universally open morphisms with geometrically reduced fibres}{226}
\EGAIVthreeContentsLine{15.3.}{Application: criteria of reducedness and irreducibility}{228}
\EGAIVthreeContentsLine{15.4.}{Complements on Cohen-Macaulay morphisms}{229}
\EGAIVthreeContentsLine{15.5.}{Separable rank of the fibres of a quasi-finite and universally open morphism. Application to the geometric connected components of the fibres of a proper morphism}{231}
\EGAIVthreeContentsLine{15.6.}{Connected components of fibres along a section}{236}
\EGAIVthreeContentsLine{15.7.}{Appendix: Valuative criteria of local properness}{242}
\EGAIVthreeContentsLine{}{BIBLIOGRAPHY \emph{(continued)}}{249}
\EGAIVthreeContentsLine{}{INDEX OF NOTATION}{250}
\EGAIVthreeContentsLine{}{TERMINOLOGICAL INDEX}{251}

\vspace{2em}
\begin{center}
\emph{Received November 15, 1964.}
\end{center}

\let\EGAIVthreeContentsLine\relax
\endgroup
% END NATIVE INPUT source_aligned/ega4-15-backmatter.tex
% END NATIVE INPUT source_aligned/ega4-15

\clearpage
\phantomsection
\addcontentsline{toc}{section}{Consolidated bibliography}
\section*{Consolidated bibliography}
\nocite{*}
\bibliography{the}
\bibliographystyle{amsalpha}

\end{document}
